% Complete English EGA IV2 editable text.
%
\PassOptionsToPackage{hyperfootnotes=false}{hyperref}
% BEGIN NATIVE INPUT preamble
\documentclass[oneside]{amsart}

\usepackage[all]{xy}
\usepackage{tikz-cd}
\usepackage[T1]{fontenc}
\usepackage{xstring}
\usepackage{xparse}
\usepackage{xcolor}
\definecolor{brightmaroon}{rgb}{0.76, 0.13, 0.28}
\usepackage[linktocpage=true,colorlinks=true,hyperindex,hypertexnames=false,citecolor=blue,linkcolor=brightmaroon]{hyperref}
\usepackage[left=1.25in,right=1.25in,top=0.75in,bottom=0.75in]{geometry}
%\usepackage[charter,greekfamily=didot]{mathdesign}
%\usepackage{Baskervaldx}
\usepackage{amssymb}
\usepackage{mathrsfs}
\usepackage{mathpazo}
\linespread{1.05}

\usepackage[nobottomtitles]{titlesec}
\usepackage{marginnote}
\usepackage[shortlabels]{enumitem}
\usepackage{longtable}
\usepackage{multicol}
\usepackage{aurical}
\usepackage{microtype}

\newtheoremstyle{ega-env-style}%
  {}{}{\rmfamily}{}{\bfseries}{.}{ }{\thmnote{(#3)}}%

\newtheoremstyle{ega-thm-env-style}%
  {}{}{\itshape}{}{\bfseries}{. --- }{ }{\thmname{#1}\thmnote{ (#3)}}%

\newtheoremstyle{ega-defn-env-style}%
  {}{}{\rmfamily}{}{\bfseries}{. --- }{ }{\thmname{#1}\thmnote{ (#3)}}%

\theoremstyle{ega-env-style}
\newtheorem*{env}{---}

\theoremstyle{ega-thm-env-style}
\newtheorem*{theorem}{Theorem}
\newtheorem*{proposition}{Proposition}
\newtheorem*{lemma}{Lemma}
\newtheorem*{corollary}{Corollary}

\theoremstyle{ega-defn-env-style}
\newtheorem*{definition}{Definition}
\newtheorem*{example}{Example}
\newtheorem*{examples}{Examples}
\newtheorem*{remark}{Remark}
\newtheorem*{remarks}{Remarks}
\newtheorem*{scholium}{Scholium}
\newtheorem*{notation}{Notation}

% indent subsections, see https://tex.stackexchange.com/questions/177290/.
% also make section titles bigger.
% also add § to \thesection, https://tex.stackexchange.com/questions/119667/ and https://tex.stackexchange.com/questions/308737/.
\makeatletter
\def\l@subsection{\@tocline{2}{0pt}{2.5pc}{2.2pc}{}}
\def\section{\@startsection{section}{1}%
  \z@{.7\linespacing\@plus\linespacing}{.5\linespacing}%
  {\normalfont\bfseries\Large\scshape\centering}}
\renewcommand{\@seccntformat}[1]{%
  \ifnum\pdfstrcmp{#1}{section}=0\textsection\fi%
  \csname the#1\endcsname.~}
\makeatother

%\allowdisplaybreaks[1]
%\binoppenalty=9999
%\relpenalty=9999

% for Chapter 0, Chapter I, etc.
% credit for ZeroRoman https://tex.stackexchange.com/questions/211414/
% added into scripts/make_book.py
%\newcommand{\ZeroRoman}[1]{\ifcase\value{#1}\relax 0\else\Roman{#1}\fi}
%\renewcommand{\thechapter}{\ZeroRoman{chapter}}

\def\mathcal{\mathscr}
% BEGIN NATIVE INPUT preamble-base
\def\sh{\mathcal}                   % sheaf font
\def\bb{\mathbf}                    % bold font
\def\cat{\mathtt}                   % category font
\def\leq{\leqslant}                 % <=
\def\geq{\geqslant}                 % >=
\def\setmin{-}                      % set minus
\def\rad{\mathfrak{r}}              % radical
\def\nilrad{\mathfrak{N}}           % nilradical
\def\emp{\varnothing}               % empty set
\def\vphi{\phi}                     % for switching \phi and \varphi, change if needed
\def\HH{\mathrm{H}}                 % cohomology H
\def\CHH{\check{\HH}}               % Čech cohomology H
\def\RR{\mathrm{R}}                 % right derived R
\def\LL{\mathrm{L}}                 % left derived L
\def\dual#1{{#1}^\vee}              % dual
\def\kres{k}                        % residue field k
\def\C{\cat{C}}                     % category C
\def\op{^\cat{op}}                  % opposite category
\def\Set{\cat{Set}}                 % category of sets
\def\CHom{\cat{Hom}}                % functor category
\def\supertilde{{\,\widetilde{\,}\,}}   % use \supertilde instead of ^\sim
\def\GL{\bb{GL}}
\def\bbGamma{\boldsymbol\Gamma}
\def\red{\mathrm{red}}
\def\rg{\operatorname{rg}}
\def\gr{\operatorname{gr}}
\def\Gr{\operatorname{Gr}}
\def\Hom{\operatorname{Hom}}
\def\Proj{\operatorname{Proj}}
\def\Tor{\operatorname{Tor}}
\def\Ext{\operatorname{Ext}}
\def\Supp{\operatorname{Supp}}
\def\Ker{\operatorname{Ker}}
\def\Im{\operatorname{Im}}
\def\Coker{\operatorname{Coker}}
\def\Spec{\operatorname{Spec}}
\def\Spm{\operatorname{Spm}}
\def\Spf{\operatorname{Spf}}
\def\grad{\operatorname{grad}}
\def\dimc{\operatorname{dimc}}
\def\codim{\operatorname{codim}}
\def\id{\operatorname{id}}
\def\Der{\operatorname{Der}}
\def\Diff{\operatorname{Diff}}
\def\Hyp{\operatorname{Hyp}}
\def\Tr{\operatorname{Tr}}

\renewcommand{\to}{\mathchoice{\longrightarrow}{\rightarrow}{\rightarrow}{\rightarrow}}
\newcommand{\from}{\mathchoice{\longleftarrow}{\leftarrow}{\leftarrow}{\leftarrow}}
\let\mapstoo\mapsto
\renewcommand{\mapsto}{\mathchoice{\longmapsto}{\mapstoo}{\mapstoo}{\mapstoo}}
\def\isoto{\simeq} % isomorphism
\def\simto{\xrightarrow{\sim}}

\renewcommand{\hat}[1]{\widehat{#1}}
\renewcommand{\tilde}[1]{\widetilde{#1}}
% END NATIVE INPUT preamble-base

\def\shHom{\sh{H}\textup{\kern-2.2pt{\Fontauri\slshape om}}}   % sheaf Hom
\def\shProj{\sh{P}\textup{\kern-2.2pt{\Fontauri\slshape roj}}} % sheaf Proj
\def\shExt{\sh{E}\textup{\kern-2.2pt{\Fontauri\slshape xt}}}   % sheaf Ext
\def\shTor{\sh{T}\textup{\kern-2.2pt{\Fontauri\slshape or}}}   % sheaf Tor
\def\shGr{\sh{G}\textup{\kern-2.2pt{\Fontauri\slshape r}}}     % sheaf Gr
\def\shDer{\sh{D}\,\textup{\kern-2.2pt{\Fontauri\slshape er}}} % sheaf Der
\def\shDiff{\sh{D}\,\textup{\kern-2.2pt{\Fontauri\slshape if{}f}}\,} % sheaf Diff
\def\shHomcont{\sh{H}\textup{\kern-2.2pt{\Fontauri\slshape om.\,cont}}}   % sheaf Hom.cont
\def\shAut{\sh{A}\textup{\kern-2.2pt{\Fontauri\slshape ut}}}   % sheaf Aut

% if unsure of a translation
%\newcommand{\unsure}[2][]{\hl{#2}\marginpar{#1}}
%\newcommand{\completelyunsure}{\unsure{[\ldots]}}
\def\unsure#1{#1 {\color{red}(?)}}
\def\completelyunsure{{\color{red}(???)}}

% use to mark where original page starts
\newcommand{\oldpage}[2][]{{\marginnote{\normalfont{\textbf{#1}~|~#2}}}\ignorespaces}
\def\sectionbreak{\begin{center}***\end{center}}

% Give every source label its own zero-size PDF destination.  The inherited
% unnumbered theorem environments otherwise leave many labels pointing to the
% preceding subsection or display anchor rather than to their exact statement.
\let\EGAOriginalLabel\label
\renewcommand{\label}[1]{\phantomsection\EGAOriginalLabel{#1}}

% for referencing environments.
% use as \sref{chapter-number.x.y.z}, with optional args
% for volume and indices, e.g. \sref[volume]{chapter-number.x.y.z}[i].
\makeatletter
\newcommand{\EGAUnavailableHyperref}[2]{%
  \begingroup
  \let\@latex@warning\@gobble
  \let\@latex@warning@no@line\@gobble
  \let\G@refundefinedtrue\relax
  \hyperref[#1]{#2}%
  \endgroup
}
\NewDocumentCommand{\sref}{o m o}{%
  \@ifundefined{r@#2}%
    {\IfNoValueTF{#1}%
      {\IfNoValueTF{#3}%
        {\EGAUnavailableHyperref{#2}{\normalfont{(\StrBehind{#2}{.})}}}%
        {\EGAUnavailableHyperref{#2}{\normalfont{(\StrBehind{#2}{.},~{#3})}}}}%
      {\IfNoValueTF{#3}%
        {\EGAUnavailableHyperref{#2}{\normalfont{(\textbf{#1},~\StrBehind{#2}{.})}}}%
        {\EGAUnavailableHyperref{#2}{\normalfont{(\textbf{#1},~\StrBehind{#2}{.},~{#3})}}}}}%
    {\IfNoValueTF{#1}%
      {\IfNoValueTF{#3}%
        {\hyperref[#2]{\normalfont{(\StrBehind{#2}{.})}}}%
        {\hyperref[#2]{\normalfont{(\StrBehind{#2}{.},~{#3})}}}}%
      {\IfNoValueTF{#3}%
        {\hyperref[#2]{\normalfont{(\textbf{#1},~\StrBehind{#2}{.})}}}%
        {\hyperref[#2]{\normalfont{(\textbf{#1},~\StrBehind{#2}{.},~{#3})}}}}}%
}
\makeatother

% for marking changes following errata
% use as \erratum[volume]{correction} to say where the erratum is given
\newcommand{\erratum}[2][]{{#2}\marginpar{\footnotesize\textbf{Err}\textsubscript{#1}}}

% for referencing equations, use as \eref{eq:chapter-number.x.y.z}.
\newcommand{\eref}[1]{\hyperref[#1]{\normalfont{(\StrBehind{#1}{.})}}}
% END NATIVE INPUT preamble
% BEGIN NATIVE INPUT EGA4_reference_v2_registry
% Generated stable semantic destinations for the EGA IV reader.
% Visible text and numbering are unchanged.
\makeatletter
\newcommand{\sgaStableTarget}[1]{%
  \Hy@raisedlink{\hyper@anchorstart{#1}\hyper@anchorend}}
\expandafter\def\csname sga@target@IV.1.1.1\endcsname{ega4.c04.def.1-1-1}
\expandafter\def\csname sga@target@IV.1.1.2\endcsname{ega4.c04.prop.1-1-2}
\expandafter\def\csname sga@target@IV.1.1.3\endcsname{ega4.c04.prop.1-1-3}
\expandafter\def\csname sga@target@IV.1.1.4\endcsname{ega4.c04.cor.1-1-4}
\expandafter\def\csname sga@target@IV.1.1.5\endcsname{ega4.c04.prop.1-1-5}
\expandafter\def\csname sga@target@IV.1.1.6\endcsname{ega4.c04.prop.1-1-6}
\expandafter\def\csname sga@target@IV.1.10.1\endcsname{ega4.c04.thm.1-10-1}
\expandafter\def\csname sga@target@IV.1.10.2\endcsname{ega4.c04.cor.1-10-2}
\expandafter\def\csname sga@target@IV.1.10.3\endcsname{ega4.c04.prop.1-10-3}
\expandafter\def\csname sga@target@IV.1.10.4\endcsname{ega4.c04.cor.1-10-4}
\expandafter\def\csname sga@target@IV.1.2.1\endcsname{ega4.c04.def.1-2-1}
\expandafter\def\csname sga@target@IV.1.2.2\endcsname{ega4.c04.prop.1-2-2}
\expandafter\def\csname sga@target@IV.1.2.3\endcsname{ega4.c04.cor.1-2-3}
\expandafter\def\csname sga@target@IV.1.2.4\endcsname{ega4.c04.prop.1-2-4}
\expandafter\def\csname sga@target@IV.1.2.4.1\endcsname{ega4.c04.eq.1-2-4-1}
\expandafter\def\csname sga@target@IV.1.2.5\endcsname{ega4.c04.prop.1-2-5}
\expandafter\def\csname sga@target@IV.1.2.6\endcsname{ega4.c04.prop.1-2-6}
\expandafter\def\csname sga@target@IV.1.2.7\endcsname{ega4.c04.prop.1-2-7}
\expandafter\def\csname sga@target@IV.1.2.8\endcsname{ega4.c04.cor.1-2-8}
\expandafter\def\csname sga@target@IV.1.2.9\endcsname{ega4.c04.prop.1-2-9}
\expandafter\def\csname sga@target@IV.1.3.1\endcsname{ega4.c04.statement.1-3-1}
\expandafter\def\csname sga@target@IV.1.3.10\endcsname{ega4.c04.statement.1-3-10}
\expandafter\def\csname sga@target@IV.1.3.11\endcsname{ega4.c04.prop.1-3-11}
\expandafter\def\csname sga@target@IV.1.3.2\endcsname{ega4.c04.def.1-3-2}
\expandafter\def\csname sga@target@IV.1.3.3\endcsname{ega4.c04.prop.1-3-3}
\expandafter\def\csname sga@target@IV.1.3.4\endcsname{ega4.c04.prop.1-3-4}
\expandafter\def\csname sga@target@IV.1.3.5\endcsname{ega4.c04.cor.1-3-5}
\expandafter\def\csname sga@target@IV.1.3.6\endcsname{ega4.c04.prop.1-3-6}
\expandafter\def\csname sga@target@IV.1.3.7\endcsname{ega4.c04.prop.1-3-7}
\expandafter\def\csname sga@target@IV.1.3.8\endcsname{ega4.c04.statement.1-3-8}
\expandafter\def\csname sga@target@IV.1.3.9\endcsname{ega4.c04.prop.1-3-9}
\expandafter\def\csname sga@target@IV.1.4.1\endcsname{ega4.c04.statement.1-4-1}
\expandafter\def\csname sga@target@IV.1.4.2\endcsname{ega4.c04.def.1-4-2}
\expandafter\def\csname sga@target@IV.1.4.3\endcsname{ega4.c04.prop.1-4-3}
\expandafter\def\csname sga@target@IV.1.4.3.1\endcsname{ega4.c04.cor.1-4-3-1}
\expandafter\def\csname sga@target@IV.1.4.4\endcsname{ega4.c04.prop.1-4-4}
\expandafter\def\csname sga@target@IV.1.4.4.1\endcsname{ega4.c04.lem.1-4-4-1}
\expandafter\def\csname sga@target@IV.1.4.5\endcsname{ega4.c04.cor.1-4-5}
\expandafter\def\csname sga@target@IV.1.4.6\endcsname{ega4.c04.prop.1-4-6}
\expandafter\def\csname sga@target@IV.1.4.7\endcsname{ega4.c04.prop.1-4-7}
\expandafter\def\csname sga@target@IV.1.4.7.1\endcsname{ega4.c04.lem.1-4-7-1}
\expandafter\def\csname sga@target@IV.1.5.1\endcsname{ega4.c04.prop.1-5-1}
\expandafter\def\csname sga@target@IV.1.5.2\endcsname{ega4.c04.def.1-5-2}
\expandafter\def\csname sga@target@IV.1.5.3\endcsname{ega4.c04.prop.1-5-3}
\expandafter\def\csname sga@target@IV.1.5.4\endcsname{ega4.c04.prop.1-5-4}
\expandafter\def\csname sga@target@IV.1.5.5\endcsname{ega4.c04.cor.1-5-5}
\expandafter\def\csname sga@target@IV.1.5.6\endcsname{ega4.c04.cor.1-5-6}
\expandafter\def\csname sga@target@IV.1.5.7\endcsname{ega4.c04.prop.1-5-7}
\expandafter\def\csname sga@target@IV.1.6.1\endcsname{ega4.c04.def.1-6-1}
\expandafter\def\csname sga@target@IV.1.6.2\endcsname{ega4.c04.prop.1-6-2}
\expandafter\def\csname sga@target@IV.1.6.3\endcsname{ega4.c04.cor.1-6-3}
\expandafter\def\csname sga@target@IV.1.6.4\endcsname{ega4.c04.rem.1-6-4}
\expandafter\def\csname sga@target@IV.1.6.5\endcsname{ega4.c04.prop.1-6-5}
\expandafter\def\csname sga@target@IV.1.7.1\endcsname{ega4.c04.statement.1-7-1}
\expandafter\def\csname sga@target@IV.1.7.10\endcsname{ega4.c04.statement.1-7-10}
\expandafter\def\csname sga@target@IV.1.7.11\endcsname{ega4.c04.statement.1-7-11}
\expandafter\def\csname sga@target@IV.1.7.12\endcsname{ega4.c04.statement.1-7-12}
\expandafter\def\csname sga@target@IV.1.7.13\endcsname{ega4.c04.statement.1-7-13}
\expandafter\def\csname sga@target@IV.1.7.14\endcsname{ega4.c04.statement.1-7-14}
\expandafter\def\csname sga@target@IV.1.7.15\endcsname{ega4.c04.statement.1-7-15}
\expandafter\def\csname sga@target@IV.1.7.16\endcsname{ega4.c04.statement.1-7-16}
\expandafter\def\csname sga@target@IV.1.7.17\endcsname{ega4.c04.statement.1-7-17}
\expandafter\def\csname sga@target@IV.1.7.18\endcsname{ega4.c04.statement.1-7-18}
\expandafter\def\csname sga@target@IV.1.7.19\endcsname{ega4.c04.statement.1-7-19}
\expandafter\def\csname sga@target@IV.1.7.2\endcsname{ega4.c04.statement.1-7-2}
\expandafter\def\csname sga@target@IV.1.7.20\endcsname{ega4.c04.statement.1-7-20}
\expandafter\def\csname sga@target@IV.1.7.21\endcsname{ega4.c04.statement.1-7-21}
\expandafter\def\csname sga@target@IV.1.7.22\endcsname{ega4.c04.statement.1-7-22}
\expandafter\def\csname sga@target@IV.1.7.3\endcsname{ega4.c04.statement.1-7-3}
\expandafter\def\csname sga@target@IV.1.7.4\endcsname{ega4.c04.statement.1-7-4}
\expandafter\def\csname sga@target@IV.1.7.5\endcsname{ega4.c04.statement.1-7-5}
\expandafter\def\csname sga@target@IV.1.7.6\endcsname{ega4.c04.statement.1-7-6}
\expandafter\def\csname sga@target@IV.1.7.7\endcsname{ega4.c04.statement.1-7-7}
\expandafter\def\csname sga@target@IV.1.7.8\endcsname{ega4.c04.statement.1-7-8}
\expandafter\def\csname sga@target@IV.1.7.9\endcsname{ega4.c04.statement.1-7-9}
\expandafter\def\csname sga@target@IV.1.8.1\endcsname{ega4.c04.statement.1-8-1}
\expandafter\def\csname sga@target@IV.1.8.2\endcsname{ega4.c04.prop.1-8-2}
\expandafter\def\csname sga@target@IV.1.8.3\endcsname{ega4.c04.lem.1-8-3}
\expandafter\def\csname sga@target@IV.1.8.4\endcsname{ega4.c04.thm.1-8-4}
\expandafter\def\csname sga@target@IV.1.8.4.1\endcsname{ega4.c04.lem.1-8-4-1}
\expandafter\def\csname sga@target@IV.1.8.4.2\endcsname{ega4.c04.lem.1-8-4-2}
\expandafter\def\csname sga@target@IV.1.8.5\endcsname{ega4.c04.cor.1-8-5}
\expandafter\def\csname sga@target@IV.1.8.5.1\endcsname{ega4.c04.lem.1-8-5-1}
\expandafter\def\csname sga@target@IV.1.8.6\endcsname{ega4.c04.cor.1-8-6}
\expandafter\def\csname sga@target@IV.1.8.7\endcsname{ega4.c04.prop.1-8-7}
\expandafter\def\csname sga@target@IV.1.8.7.1\endcsname{ega4.c04.lem.1-8-7-1}
\expandafter\def\csname sga@target@IV.1.9.1\endcsname{ega4.c04.lem.1-9-1}
\expandafter\def\csname sga@target@IV.1.9.10\endcsname{ega4.c04.prop.1-9-10}
\expandafter\def\csname sga@target@IV.1.9.11\endcsname{ega4.c04.prop.1-9-11}
\expandafter\def\csname sga@target@IV.1.9.12\endcsname{ega4.c04.cor.1-9-12}
\expandafter\def\csname sga@target@IV.1.9.13\endcsname{ega4.c04.statement.1-9-13}
\expandafter\def\csname sga@target@IV.1.9.14\endcsname{ega4.c04.prop.1-9-14}
\expandafter\def\csname sga@target@IV.1.9.15\endcsname{ega4.c04.prop.1-9-15}
\expandafter\def\csname sga@target@IV.1.9.16\endcsname{ega4.c04.statement.1-9-16}
\expandafter\def\csname sga@target@IV.1.9.2\endcsname{ega4.c04.prop.1-9-2}
\expandafter\def\csname sga@target@IV.1.9.2.1\endcsname{ega4.c04.statement.1-9-2-1}
\expandafter\def\csname sga@target@IV.1.9.2.i\endcsname{ega4.c04.statement.1-9-2-i}
\expandafter\def\csname sga@target@IV.1.9.2.ii\endcsname{ega4.c04.statement.1-9-2-ii}
\expandafter\def\csname sga@target@IV.1.9.3\endcsname{ega4.c04.prop.1-9-3}
\expandafter\def\csname sga@target@IV.1.9.3.1\endcsname{ega4.c04.lem.1-9-3-1}
\expandafter\def\csname sga@target@IV.1.9.3.2\endcsname{ega4.c04.lem.1-9-3-2}
\expandafter\def\csname sga@target@IV.1.9.4\endcsname{ega4.c04.def.1-9-4}
\expandafter\def\csname sga@target@IV.1.9.5\endcsname{ega4.c04.prop.1-9-5}
\expandafter\def\csname sga@target@IV.1.9.6\endcsname{ega4.c04.statement.1-9-6}
\expandafter\def\csname sga@target@IV.1.9.7\endcsname{ega4.c04.thm.1-9-7}
\expandafter\def\csname sga@target@IV.1.9.8\endcsname{ega4.c04.cor.1-9-8}
\expandafter\def\csname sga@target@IV.1.9.9\endcsname{ega4.c04.cor.1-9-9}
\expandafter\def\csname sga@target@IV.10.1.1\endcsname{ega4.c04.statement.10-1-1}
\expandafter\def\csname sga@target@IV.10.1.2\endcsname{ega4.c04.prop.10-1-2}
\expandafter\def\csname sga@target@IV.10.1.3\endcsname{ega4.c04.def.10-1-3}
\expandafter\def\csname sga@target@IV.10.1.4\endcsname{ega4.c04.cor.10-1-4}
\expandafter\def\csname sga@target@IV.10.10.1\endcsname{ega4.c04.statement.10-10-1}
\expandafter\def\csname sga@target@IV.10.10.2\endcsname{ega4.c04.statement.10-10-2}
\expandafter\def\csname sga@target@IV.10.10.3\endcsname{ega4.c04.statement.10-10-3}
\expandafter\def\csname sga@target@IV.10.10.4\endcsname{ega4.c04.statement.10-10-4}
\expandafter\def\csname sga@target@IV.10.10.5\endcsname{ega4.c04.statement.10-10-5}
\expandafter\def\csname sga@target@IV.10.2.1\endcsname{ega4.c04.prop.10-2-1}
\expandafter\def\csname sga@target@IV.10.2.2\endcsname{ega4.c04.def.10-2-2}
\expandafter\def\csname sga@target@IV.10.2.3\endcsname{ega4.c04.cor.10-2-3}
\expandafter\def\csname sga@target@IV.10.2.4\endcsname{ega4.c04.cor.10-2-4}
\expandafter\def\csname sga@target@IV.10.2.5\endcsname{ega4.c04.cor.10-2-5}
\expandafter\def\csname sga@target@IV.10.2.6\endcsname{ega4.c04.cor.10-2-6}
\expandafter\def\csname sga@target@IV.10.2.7\endcsname{ega4.c04.prop.10-2-7}
\expandafter\def\csname sga@target@IV.10.2.8\endcsname{ega4.c04.item.10-2-8}
\expandafter\def\csname sga@target@IV.10.2.8.1\endcsname{ega4.c04.eq.10-2-8-1}
\expandafter\def\csname sga@target@IV.10.3.1\endcsname{ega4.c04.def.10-3-1}
\expandafter\def\csname sga@target@IV.10.3.2\endcsname{ega4.c04.prop.10-3-2}
\expandafter\def\csname sga@target@IV.10.3.3\endcsname{ega4.c04.prop.10-3-3}
\expandafter\def\csname sga@target@IV.10.4.1\endcsname{ega4.c04.def.10-4-1}
\expandafter\def\csname sga@target@IV.10.4.10\endcsname{ega4.c04.statement.10-4-10}
\expandafter\def\csname sga@target@IV.10.4.11\endcsname{ega4.c04.statement.10-4-11}
\expandafter\def\csname sga@target@IV.10.4.11.1\endcsname{ega4.c04.lem.10-4-11-1}
\expandafter\def\csname sga@target@IV.10.4.12\endcsname{ega4.c04.cor.10-4-12}
\expandafter\def\csname sga@target@IV.10.4.12.1\endcsname{ega4.c04.sec.10-4-12-1}
\expandafter\def\csname sga@target@IV.10.4.12.1.errata\endcsname{ega4.c04.eq.10-4-12-1-errata}
\expandafter\def\csname sga@target@IV.10.4.13\endcsname{ega4.c04.cor.10-4-13}
\expandafter\def\csname sga@target@IV.10.4.14\endcsname{ega4.c04.rem.10-4-14}
\expandafter\def\csname sga@target@IV.10.4.2\endcsname{ega4.c04.prop.10-4-2}
\expandafter\def\csname sga@target@IV.10.4.3\endcsname{ega4.c04.statement.10-4-3}
\expandafter\def\csname sga@target@IV.10.4.4\endcsname{ega4.c04.prop.10-4-4}
\expandafter\def\csname sga@target@IV.10.4.5\endcsname{ega4.c04.prop.10-4-5}
\expandafter\def\csname sga@target@IV.10.4.5.1\endcsname{ega4.c04.lem.10-4-5-1}
\expandafter\def\csname sga@target@IV.10.4.6\endcsname{ega4.c04.cor.10-4-6}
\expandafter\def\csname sga@target@IV.10.4.7\endcsname{ega4.c04.cor.10-4-7}
\expandafter\def\csname sga@target@IV.10.4.8\endcsname{ega4.c04.cor.10-4-8}
\expandafter\def\csname sga@target@IV.10.4.9\endcsname{ega4.c04.statement.10-4-9}
\expandafter\def\csname sga@target@IV.10.5.1\endcsname{ega4.c04.prop.10-5-1}
\expandafter\def\csname sga@target@IV.10.5.10\endcsname{ega4.c04.prop.10-5-10}
\expandafter\def\csname sga@target@IV.10.5.2\endcsname{ega4.c04.cor.10-5-2}
\expandafter\def\csname sga@target@IV.10.5.3\endcsname{ega4.c04.cor.10-5-3}
\expandafter\def\csname sga@target@IV.10.5.4\endcsname{ega4.c04.rem.10-5-4}
\expandafter\def\csname sga@target@IV.10.5.5\endcsname{ega4.c04.prop.10-5-5}
\expandafter\def\csname sga@target@IV.10.5.6\endcsname{ega4.c04.rem.10-5-6}
\expandafter\def\csname sga@target@IV.10.5.7\endcsname{ega4.c04.prop.10-5-7}
\expandafter\def\csname sga@target@IV.10.5.8\endcsname{ega4.c04.cor.10-5-8}
\expandafter\def\csname sga@target@IV.10.5.9\endcsname{ega4.c04.cor.10-5-9}
\expandafter\def\csname sga@target@IV.10.6.1\endcsname{ega4.c04.prop.10-6-1}
\expandafter\def\csname sga@target@IV.10.6.1.1\endcsname{ega4.c04.eq.10-6-1-1}
\expandafter\def\csname sga@target@IV.10.6.1.2\endcsname{ega4.c04.eq.10-6-1-2}
\expandafter\def\csname sga@target@IV.10.6.1.3\endcsname{ega4.c04.eq.10-6-1-3}
\expandafter\def\csname sga@target@IV.10.6.1.4\endcsname{ega4.c04.eq.10-6-1-4}
\expandafter\def\csname sga@target@IV.10.6.2\endcsname{ega4.c04.cor.10-6-2}
\expandafter\def\csname sga@target@IV.10.6.3\endcsname{ega4.c04.prop.10-6-3}
\expandafter\def\csname sga@target@IV.10.6.3.1\endcsname{ega4.c04.eq.10-6-3-1}
\expandafter\def\csname sga@target@IV.10.6.3.2\endcsname{ega4.c04.eq.10-6-3-2}
\expandafter\def\csname sga@target@IV.10.6.3.3\endcsname{ega4.c04.eq.10-6-3-3}
\expandafter\def\csname sga@target@IV.10.6.4\endcsname{ega4.c04.cor.10-6-4}
\expandafter\def\csname sga@target@IV.10.6.4.1\endcsname{ega4.c04.eq.10-6-4-1}
\expandafter\def\csname sga@target@IV.10.7.1\endcsname{ega4.c04.statement.10-7-1}
\expandafter\def\csname sga@target@IV.10.7.2\endcsname{ega4.c04.statement.10-7-2}
\expandafter\def\csname sga@target@IV.10.7.3\endcsname{ega4.c04.statement.10-7-3}
\expandafter\def\csname sga@target@IV.10.8.1\endcsname{ega4.c04.def.10-8-1}
\expandafter\def\csname sga@target@IV.10.8.1.1\endcsname{ega4.c04.eq.10-8-1-1}
\expandafter\def\csname sga@target@IV.10.8.1.2\endcsname{ega4.c04.eq.10-8-1-2}
\expandafter\def\csname sga@target@IV.10.8.2\endcsname{ega4.c04.statement.10-8-2}
\expandafter\def\csname sga@target@IV.10.8.3\endcsname{ega4.c04.lem.10-8-3}
\expandafter\def\csname sga@target@IV.10.8.3.1\endcsname{ega4.c04.eq.10-8-3-1}
\expandafter\def\csname sga@target@IV.10.8.4\endcsname{ega4.c04.cor.10-8-4}
\expandafter\def\csname sga@target@IV.10.8.5\endcsname{ega4.c04.prop.10-8-5}
\expandafter\def\csname sga@target@IV.10.8.6\endcsname{ega4.c04.cor.10-8-6}
\expandafter\def\csname sga@target@IV.10.8.7\endcsname{ega4.c04.prop.10-8-7}
\expandafter\def\csname sga@target@IV.10.8.7.1\endcsname{ega4.c04.eq.10-8-7-1}
\expandafter\def\csname sga@target@IV.10.8.8\endcsname{ega4.c04.cor.10-8-8}
\expandafter\def\csname sga@target@IV.10.9.1\endcsname{ega4.c04.statement.10-9-1}
\expandafter\def\csname sga@target@IV.10.9.2\endcsname{ega4.c04.statement.10-9-2}
\expandafter\def\csname sga@target@IV.10.9.3\endcsname{ega4.c04.statement.10-9-3}
\expandafter\def\csname sga@target@IV.10.9.4\endcsname{ega4.c04.prop.10-9-4}
\expandafter\def\csname sga@target@IV.10.9.5\endcsname{ega4.c04.statement.10-9-5}
\expandafter\def\csname sga@target@IV.10.9.6\endcsname{ega4.c04.prop.10-9-6}
\expandafter\def\csname sga@target@IV.11.1.1\endcsname{ega4.c04.thm.11-1-1}
\expandafter\def\csname sga@target@IV.11.1.1.1\endcsname{ega4.c04.sec.11-1-1-1}
\expandafter\def\csname sga@target@IV.11.1.1.2\endcsname{ega4.c04.lem.11-1-1-2}
\expandafter\def\csname sga@target@IV.11.1.2\endcsname{ega4.c04.cor.11-1-2}
\expandafter\def\csname sga@target@IV.11.1.3\endcsname{ega4.c04.cor.11-1-3}
\expandafter\def\csname sga@target@IV.11.1.4\endcsname{ega4.c04.cor.11-1-4}
\expandafter\def\csname sga@target@IV.11.1.5\endcsname{ega4.c04.cor.11-1-5}
\expandafter\def\csname sga@target@IV.11.1.6\endcsname{ega4.c04.cor.11-1-6}
\expandafter\def\csname sga@target@IV.11.10.1\endcsname{ega4.c04.prop.11-10-1}
\expandafter\def\csname sga@target@IV.11.10.1.1\endcsname{ega4.c04.eq.11-10-1-1}
\expandafter\def\csname sga@target@IV.11.10.1.2\endcsname{ega4.c04.eq.11-10-1-2}
\expandafter\def\csname sga@target@IV.11.10.10\endcsname{ega4.c04.prop.11-10-10}
\expandafter\def\csname sga@target@IV.11.10.2\endcsname{ega4.c04.def.11-10-2}
\expandafter\def\csname sga@target@IV.11.10.3\endcsname{ega4.c04.item.11-10-3}
\expandafter\def\csname sga@target@IV.11.10.4\endcsname{ega4.c04.prop.11-10-4}
\expandafter\def\csname sga@target@IV.11.10.5\endcsname{ega4.c04.thm.11-10-5}
\expandafter\def\csname sga@target@IV.11.10.6\endcsname{ega4.c04.prop.11-10-6}
\expandafter\def\csname sga@target@IV.11.10.7\endcsname{ega4.c04.cor.11-10-7}
\expandafter\def\csname sga@target@IV.11.10.8\endcsname{ega4.c04.def.11-10-8}
\expandafter\def\csname sga@target@IV.11.10.9\endcsname{ega4.c04.thm.11-10-9}
\expandafter\def\csname sga@target@IV.11.2.1\endcsname{ega4.c04.statement.11-2-1}
\expandafter\def\csname sga@target@IV.11.2.1.1\endcsname{ega4.c04.eq.11-2-1-1}
\expandafter\def\csname sga@target@IV.11.2.1.2\endcsname{ega4.c04.eq.11-2-1-2}
\expandafter\def\csname sga@target@IV.11.2.1.3\endcsname{ega4.c04.eq.11-2-1-3}
\expandafter\def\csname sga@target@IV.11.2.1.4\endcsname{ega4.c04.eq.11-2-1-4}
\expandafter\def\csname sga@target@IV.11.2.1.5\endcsname{ega4.c04.eq.11-2-1-5}
\expandafter\def\csname sga@target@IV.11.2.10\endcsname{ega4.c04.rem.11-2-10}
\expandafter\def\csname sga@target@IV.11.2.2\endcsname{ega4.c04.statement.11-2-2}
\expandafter\def\csname sga@target@IV.11.2.2.1\endcsname{ega4.c04.eq.11-2-2-1}
\expandafter\def\csname sga@target@IV.11.2.2.2\endcsname{ega4.c04.lem.11-2-2-2}
\expandafter\def\csname sga@target@IV.11.2.3\endcsname{ega4.c04.statement.11-2-3}
\expandafter\def\csname sga@target@IV.11.2.3.1\endcsname{ega4.c04.eq.11-2-3-1}
\expandafter\def\csname sga@target@IV.11.2.3.2\endcsname{ega4.c04.eq.11-2-3-2}
\expandafter\def\csname sga@target@IV.11.2.3.3\endcsname{ega4.c04.eq.11-2-3-3}
\expandafter\def\csname sga@target@IV.11.2.4\endcsname{ega4.c04.lem.11-2-4}
\expandafter\def\csname sga@target@IV.11.2.5\endcsname{ega4.c04.lem.11-2-5}
\expandafter\def\csname sga@target@IV.11.2.6\endcsname{ega4.c04.thm.11-2-6}
\expandafter\def\csname sga@target@IV.11.2.6.1\endcsname{ega4.c04.cor.11-2-6-1}
\expandafter\def\csname sga@target@IV.11.2.7\endcsname{ega4.c04.cor.11-2-7}
\expandafter\def\csname sga@target@IV.11.2.8\endcsname{ega4.c04.prop.11-2-8}
\expandafter\def\csname sga@target@IV.11.2.9\endcsname{ega4.c04.thm.11-2-9}
\expandafter\def\csname sga@target@IV.11.2.9.1\endcsname{ega4.c04.statement.11-2-9-1}
\expandafter\def\csname sga@target@IV.11.2.9.2\endcsname{ega4.c04.lem.11-2-9-2}
\expandafter\def\csname sga@target@IV.11.2.9.3\endcsname{ega4.c04.lem.11-2-9-3}
\expandafter\def\csname sga@target@IV.11.2.9.4\endcsname{ega4.c04.cor.11-2-9-4}
\expandafter\def\csname sga@target@IV.11.2.9.5\endcsname{ega4.c04.lem.11-2-9-5}
\expandafter\def\csname sga@target@IV.11.2.9.6\endcsname{ega4.c04.lem.11-2-9-6}
\expandafter\def\csname sga@target@IV.11.2.9.7\endcsname{ega4.c04.eq.11-2-9-7}
\expandafter\def\csname sga@target@IV.11.2.9.8\endcsname{ega4.c04.eq.11-2-9-8}
\expandafter\def\csname sga@target@IV.11.2.9.9\endcsname{ega4.c04.eq.11-2-9-9}
\expandafter\def\csname sga@target@IV.11.3\endcsname{ega4.c04.subsec.11-3.key-iv-11-3}
\expandafter\def\csname sga@target@IV.11.3.1\endcsname{ega4.c04.thm.11-3-1}
\expandafter\def\csname sga@target@IV.11.3.10\endcsname{ega4.c04.thm.11-3-10}
\expandafter\def\csname sga@target@IV.11.3.10.1\endcsname{ega4.c04.lem.11-3-10-1}
\expandafter\def\csname sga@target@IV.11.3.10.2\endcsname{ega4.c04.lem.11-3-10-2}
\expandafter\def\csname sga@target@IV.11.3.10.3\endcsname{ega4.c04.eq.11-3-10-3}
\expandafter\def\csname sga@target@IV.11.3.10.4\endcsname{ega4.c04.sec.11-3-10-4}
\expandafter\def\csname sga@target@IV.11.3.11\endcsname{ega4.c04.cor.11-3-11}
\expandafter\def\csname sga@target@IV.11.3.12\endcsname{ega4.c04.rem.11-3-12}
\expandafter\def\csname sga@target@IV.11.3.13\endcsname{ega4.c04.prop.11-3-13}
\expandafter\def\csname sga@target@IV.11.3.14\endcsname{ega4.c04.cor.11-3-14}
\expandafter\def\csname sga@target@IV.11.3.15\endcsname{ega4.c04.prop.11-3-15}
\expandafter\def\csname sga@target@IV.11.3.16\endcsname{ega4.c04.prop.11-3-16}
\expandafter\def\csname sga@target@IV.11.3.17\endcsname{ega4.c04.cor.11-3-17}
\expandafter\def\csname sga@target@IV.11.3.2\endcsname{ega4.c04.cor.11-3-2}
\expandafter\def\csname sga@target@IV.11.3.3\endcsname{ega4.c04.cor.11-3-3}
\expandafter\def\csname sga@target@IV.11.3.4\endcsname{ega4.c04.prop.11-3-4}
\expandafter\def\csname sga@target@IV.11.3.5\endcsname{ega4.c04.cor.11-3-5}
\expandafter\def\csname sga@target@IV.11.3.6\endcsname{ega4.c04.prop.11-3-6}
\expandafter\def\csname sga@target@IV.11.3.7\endcsname{ega4.c04.prop.11-3-7}
\expandafter\def\csname sga@target@IV.11.3.8\endcsname{ega4.c04.thm.11-3-8}
\expandafter\def\csname sga@target@IV.11.3.9\endcsname{ega4.c04.prop.11-3-9}
\expandafter\def\csname sga@target@IV.11.3.9.1\endcsname{ega4.c04.lem.11-3-9-1}
\expandafter\def\csname sga@target@IV.11.4.1\endcsname{ega4.c04.thm.11-4-1}
\expandafter\def\csname sga@target@IV.11.4.1.1\endcsname{ega4.c04.lem.11-4-1-1}
\expandafter\def\csname sga@target@IV.11.4.1.2\endcsname{ega4.c04.lem.11-4-1-2}
\expandafter\def\csname sga@target@IV.11.4.10\endcsname{ega4.c04.cor.11-4-10}
\expandafter\def\csname sga@target@IV.11.4.11\endcsname{ega4.c04.prop.11-4-11}
\expandafter\def\csname sga@target@IV.11.4.12\endcsname{ega4.c04.prop.11-4-12}
\expandafter\def\csname sga@target@IV.11.4.13\endcsname{ega4.c04.cor.11-4-13}
\expandafter\def\csname sga@target@IV.11.4.2\endcsname{ega4.c04.rem.11-4-2}
\expandafter\def\csname sga@target@IV.11.4.3\endcsname{ega4.c04.cor.11-4-3}
\expandafter\def\csname sga@target@IV.11.4.4\endcsname{ega4.c04.cor.11-4-4}
\expandafter\def\csname sga@target@IV.11.4.5\endcsname{ega4.c04.prop.11-4-5}
\expandafter\def\csname sga@target@IV.11.4.6\endcsname{ega4.c04.cor.11-4-6}
\expandafter\def\csname sga@target@IV.11.4.7\endcsname{ega4.c04.prop.11-4-7}
\expandafter\def\csname sga@target@IV.11.4.8\endcsname{ega4.c04.rem.11-4-8}
\expandafter\def\csname sga@target@IV.11.4.9\endcsname{ega4.c04.cor.11-4-9}
\expandafter\def\csname sga@target@IV.11.5.1\endcsname{ega4.c04.thm.11-5-1}
\expandafter\def\csname sga@target@IV.11.5.2\endcsname{ega4.c04.cor.11-5-2}
\expandafter\def\csname sga@target@IV.11.5.3\endcsname{ega4.c04.prop.11-5-3}
\expandafter\def\csname sga@target@IV.11.5.4\endcsname{ega4.c04.cor.11-5-4}
\expandafter\def\csname sga@target@IV.11.5.5\endcsname{ega4.c04.prop.11-5-5}
\expandafter\def\csname sga@target@IV.11.5.5.1\endcsname{ega4.c04.lem.11-5-5-1}
\expandafter\def\csname sga@target@IV.11.5.5.2\endcsname{ega4.c04.lem.11-5-5-2}
\expandafter\def\csname sga@target@IV.11.6.1\endcsname{ega4.c04.thm.11-6-1}
\expandafter\def\csname sga@target@IV.11.6.2\endcsname{ega4.c04.cor.11-6-2}
\expandafter\def\csname sga@target@IV.11.7.1\endcsname{ega4.c04.statement.11-7-1}
\expandafter\def\csname sga@target@IV.11.7.2\endcsname{ega4.c04.statement.11-7-2}
\expandafter\def\csname sga@target@IV.11.7.3\endcsname{ega4.c04.statement.11-7-3}
\expandafter\def\csname sga@target@IV.11.7.4\endcsname{ega4.c04.statement.11-7-4}
\expandafter\def\csname sga@target@IV.11.7.5\endcsname{ega4.c04.statement.11-7-5}
\expandafter\def\csname sga@target@IV.11.8.1\endcsname{ega4.c04.thm.11-8-1}
\expandafter\def\csname sga@target@IV.11.8.2\endcsname{ega4.c04.subsec.11-8-2}
\expandafter\def\csname sga@target@IV.11.9.1\endcsname{ega4.c04.statement.11-9-1}
\expandafter\def\csname sga@target@IV.11.9.1.1\endcsname{ega4.c04.statement.11-9-1-1}
\expandafter\def\csname sga@target@IV.11.9.10\endcsname{ega4.c04.thm.11-9-10}
\expandafter\def\csname sga@target@IV.11.9.11\endcsname{ega4.c04.rem.11-9-11}
\expandafter\def\csname sga@target@IV.11.9.12\endcsname{ega4.c04.lem.11-9-12}
\expandafter\def\csname sga@target@IV.11.9.13\endcsname{ega4.c04.prop.11-9-13}
\expandafter\def\csname sga@target@IV.11.9.14\endcsname{ega4.c04.statement.11-9-14}
\expandafter\def\csname sga@target@IV.11.9.15\endcsname{ega4.c04.prop.11-9-15}
\expandafter\def\csname sga@target@IV.11.9.16\endcsname{ega4.c04.thm.11-9-16}
\expandafter\def\csname sga@target@IV.11.9.17\endcsname{ega4.c04.prop.11-9-17}
\expandafter\def\csname sga@target@IV.11.9.17.1\endcsname{ega4.c04.lem.11-9-17-1}
\expandafter\def\csname sga@target@IV.11.9.18\endcsname{ega4.c04.item.11-9-18}
\expandafter\def\csname sga@target@IV.11.9.2\endcsname{ega4.c04.statement.11-9-2}
\expandafter\def\csname sga@target@IV.11.9.3\endcsname{ega4.c04.statement.11-9-3}
\expandafter\def\csname sga@target@IV.11.9.4\endcsname{ega4.c04.statement.11-9-4}
\expandafter\def\csname sga@target@IV.11.9.5\endcsname{ega4.c04.statement.11-9-5}
\expandafter\def\csname sga@target@IV.11.9.6\endcsname{ega4.c04.statement.11-9-6}
\expandafter\def\csname sga@target@IV.11.9.7\endcsname{ega4.c04.prop.11-9-7}
\expandafter\def\csname sga@target@IV.11.9.7.1\endcsname{ega4.c04.lem.11-9-7-1}
\expandafter\def\csname sga@target@IV.11.9.8\endcsname{ega4.c04.subsec.11-9-8}
\expandafter\def\csname sga@target@IV.11.9.9\endcsname{ega4.c04.cor.11-9-9}
\expandafter\def\csname sga@target@IV.12.0.1\endcsname{ega4.c04.statement.12-0-1}
\expandafter\def\csname sga@target@IV.12.0.2\endcsname{ega4.c04.statement.12-0-2}
\expandafter\def\csname sga@target@IV.12.1.1\endcsname{ega4.c04.thm.12-1-1}
\expandafter\def\csname sga@target@IV.12.1.1.1\endcsname{ega4.c04.statement.12-1-1-1}
\expandafter\def\csname sga@target@IV.12.1.1.2\endcsname{ega4.c04.statement.12-1-1-2}
\expandafter\def\csname sga@target@IV.12.1.1.3\endcsname{ega4.c04.statement.12-1-1-3}
\expandafter\def\csname sga@target@IV.12.1.1.4\endcsname{ega4.c04.statement.12-1-1-4}
\expandafter\def\csname sga@target@IV.12.1.1.5\endcsname{ega4.c04.prop.12-1-1-5}
\expandafter\def\csname sga@target@IV.12.1.2\endcsname{ega4.c04.item.12-1-2}
\expandafter\def\csname sga@target@IV.12.1.3\endcsname{ega4.c04.cor.12-1-3}
\expandafter\def\csname sga@target@IV.12.1.4\endcsname{ega4.c04.rem.12-1-4}
\expandafter\def\csname sga@target@IV.12.1.5\endcsname{ega4.c04.cor.12-1-5}
\expandafter\def\csname sga@target@IV.12.1.6\endcsname{ega4.c04.thm.12-1-6}
\expandafter\def\csname sga@target@IV.12.1.7\endcsname{ega4.c04.cor.12-1-7}
\expandafter\def\csname sga@target@IV.12.1.8\endcsname{ega4.c04.item.12-1-8}
\expandafter\def\csname sga@target@IV.12.2.1\endcsname{ega4.c04.thm.12-2-1}
\expandafter\def\csname sga@target@IV.12.2.1.1\endcsname{ega4.c04.statement.12-2-1-1}
\expandafter\def\csname sga@target@IV.12.2.1.2\endcsname{ega4.c04.statement.12-2-1-2}
\expandafter\def\csname sga@target@IV.12.2.1.3\endcsname{ega4.c04.statement.12-2-1-3}
\expandafter\def\csname sga@target@IV.12.2.1.4\endcsname{ega4.c04.lem.12-2-1-4}
\expandafter\def\csname sga@target@IV.12.2.2\endcsname{ega4.c04.cor.12-2-2}
\expandafter\def\csname sga@target@IV.12.2.3\endcsname{ega4.c04.item.12-2-3}
\expandafter\def\csname sga@target@IV.12.2.4\endcsname{ega4.c04.thm.12-2-4}
\expandafter\def\csname sga@target@IV.12.3.1\endcsname{ega4.c04.lem.12-3-1}
\expandafter\def\csname sga@target@IV.12.3.2\endcsname{ega4.c04.prop.12-3-2}
\expandafter\def\csname sga@target@IV.12.3.2.1\endcsname{ega4.c04.eq.12-3-2-1}
\expandafter\def\csname sga@target@IV.12.3.3\endcsname{ega4.c04.prop.12-3-3}
\expandafter\def\csname sga@target@IV.12.3.3.1\endcsname{ega4.c04.lem.12-3-3-1}
\expandafter\def\csname sga@target@IV.12.3.3.2\endcsname{ega4.c04.eq.12-3-3-2}
\expandafter\def\csname sga@target@IV.12.3.3.3\endcsname{ega4.c04.lem.12-3-3-3}
\expandafter\def\csname sga@target@IV.12.3.3.4\endcsname{ega4.c04.eq.12-3-3-4}
\expandafter\def\csname sga@target@IV.12.3.3.5\endcsname{ega4.c04.lem.12-3-3-5}
\expandafter\def\csname sga@target@IV.12.3.4\endcsname{ega4.c04.cor.12-3-4}
\expandafter\def\csname sga@target@IV.13.1.1\endcsname{ega4.c04.lem.13-1-1}
\expandafter\def\csname sga@target@IV.13.1.2\endcsname{ega4.c04.lem.13-1-2}
\expandafter\def\csname sga@target@IV.13.1.3\endcsname{ega4.c04.thm.13-1-3}
\expandafter\def\csname sga@target@IV.13.1.4\endcsname{ega4.c04.cor.13-1-4}
\expandafter\def\csname sga@target@IV.13.1.5\endcsname{ega4.c04.cor.13-1-5}
\expandafter\def\csname sga@target@IV.13.1.6\endcsname{ega4.c04.cor.13-1-6}
\expandafter\def\csname sga@target@IV.13.1.7\endcsname{ega4.c04.prop.13-1-7}
\expandafter\def\csname sga@target@IV.13.2.1\endcsname{ega4.c04.statement.13-2-1}
\expandafter\def\csname sga@target@IV.13.2.10\endcsname{ega4.c04.prop.13-2-10}
\expandafter\def\csname sga@target@IV.13.2.10.1\endcsname{ega4.c04.eq.13-2-10-1}
\expandafter\def\csname sga@target@IV.13.2.10.2\endcsname{ega4.c04.eq.13-2-10-2}
\expandafter\def\csname sga@target@IV.13.2.11\endcsname{ega4.c04.cor.13-2-11}
\expandafter\def\csname sga@target@IV.13.2.12\endcsname{ega4.c04.item.13-2-12}
\expandafter\def\csname sga@target@IV.13.2.2\endcsname{ega4.c04.def.13-2-2}
\expandafter\def\csname sga@target@IV.13.2.3\endcsname{ega4.c04.prop.13-2-3}
\expandafter\def\csname sga@target@IV.13.2.3.1\endcsname{ega4.c04.prop.13-2-3-1}
\expandafter\def\csname sga@target@IV.13.2.4\endcsname{ega4.c04.statement.13-2-4}
\expandafter\def\csname sga@target@IV.13.2.4.1\endcsname{ega4.c04.statement.13-2-4-1}
\expandafter\def\csname sga@target@IV.13.2.4.2\endcsname{ega4.c04.eq.13-2-4-2}
\expandafter\def\csname sga@target@IV.13.2.4.3\endcsname{ega4.c04.eq.13-2-4-3}
\expandafter\def\csname sga@target@IV.13.2.4.4\endcsname{ega4.c04.eq.13-2-4-4}
\expandafter\def\csname sga@target@IV.13.2.4.5\endcsname{ega4.c04.eq.13-2-4-5}
\expandafter\def\csname sga@target@IV.13.2.5\endcsname{ega4.c04.lem.13-2-5}
\expandafter\def\csname sga@target@IV.13.2.6\endcsname{ega4.c04.statement.13-2-6}
\expandafter\def\csname sga@target@IV.13.2.6.1\endcsname{ega4.c04.statement.13-2-6-1}
\expandafter\def\csname sga@target@IV.13.2.7\endcsname{ega4.c04.prop.13-2-7}
\expandafter\def\csname sga@target@IV.13.2.7.1\endcsname{ega4.c04.prop.13-2-7-1}
\expandafter\def\csname sga@target@IV.13.2.8\endcsname{ega4.c04.cor.13-2-8}
\expandafter\def\csname sga@target@IV.13.2.9\endcsname{ega4.c04.cor.13-2-9}
\expandafter\def\csname sga@target@IV.13.2.9.1\endcsname{ega4.c04.eq.13-2-9-1}
\expandafter\def\csname sga@target@IV.13.2.9.2\endcsname{ega4.c04.eq.13-2-9-2}
\expandafter\def\csname sga@target@IV.13.3.1\endcsname{ega4.c04.prop.13-3-1}
\expandafter\def\csname sga@target@IV.13.3.1.1\endcsname{ega4.c04.lem.13-3-1-1}
\expandafter\def\csname sga@target@IV.13.3.2\endcsname{ega4.c04.def.13-3-2}
\expandafter\def\csname sga@target@IV.13.3.3\endcsname{ega4.c04.prop.13-3-3}
\expandafter\def\csname sga@target@IV.13.3.4\endcsname{ega4.c04.cor.13-3-4}
\expandafter\def\csname sga@target@IV.13.3.4.1\endcsname{ega4.c04.cor.13-3-4-1}
\expandafter\def\csname sga@target@IV.13.3.5\endcsname{ega4.c04.cor.13-3-5}
\expandafter\def\csname sga@target@IV.13.3.6\endcsname{ega4.c04.prop.13-3-6}
\expandafter\def\csname sga@target@IV.13.3.6.1\endcsname{ega4.c04.prop.13-3-6-1}
\expandafter\def\csname sga@target@IV.13.3.6.2\endcsname{ega4.c04.eq.13-3-6-2}
\expandafter\def\csname sga@target@IV.13.3.7\endcsname{ega4.c04.prop.13-3-7}
\expandafter\def\csname sga@target@IV.13.3.7.1\endcsname{ega4.c04.prop.13-3-7-1}
\expandafter\def\csname sga@target@IV.13.3.8\endcsname{ega4.c04.prop.13-3-8}
\expandafter\def\csname sga@target@IV.13.3.9\endcsname{ega4.c04.statement.13-3-9}
\expandafter\def\csname sga@target@IV.14.1.1\endcsname{ega4.c04.statement.14-1-1}
\expandafter\def\csname sga@target@IV.14.1.2\endcsname{ega4.c04.prop.14-1-2}
\expandafter\def\csname sga@target@IV.14.1.3\endcsname{ega4.c04.item.14-1-3}
\expandafter\def\csname sga@target@IV.14.1.4\endcsname{ega4.c04.prop.14-1-4}
\expandafter\def\csname sga@target@IV.14.1.5\endcsname{ega4.c04.rem.14-1-5}
\expandafter\def\csname sga@target@IV.14.1.6\endcsname{ega4.c04.statement.14-1-6}
\expandafter\def\csname sga@target@IV.14.2.1\endcsname{ega4.c04.thm.14-2-1}
\expandafter\def\csname sga@target@IV.14.2.1.1\endcsname{ega4.c04.eq.14-2-1-1}
\expandafter\def\csname sga@target@IV.14.2.2\endcsname{ega4.c04.cor.14-2-2}
\expandafter\def\csname sga@target@IV.14.2.3\endcsname{ega4.c04.prop.14-2-3}
\expandafter\def\csname sga@target@IV.14.2.4\endcsname{ega4.c04.cor.14-2-4}
\expandafter\def\csname sga@target@IV.14.2.5\endcsname{ega4.c04.cor.14-2-5}
\expandafter\def\csname sga@target@IV.14.2.6\endcsname{ega4.c04.cor.14-2-6}
\expandafter\def\csname sga@target@IV.14.2.7\endcsname{ega4.c04.rem.14-2-7}
\expandafter\def\csname sga@target@IV.14.3\endcsname{ega4.c04.subsec.14-3.key-iv-14-3}
\expandafter\def\csname sga@target@IV.14.3.1\endcsname{ega4.c04.statement.14-3-1}
\expandafter\def\csname sga@target@IV.14.3.10\endcsname{ega4.c04.lem.14-3-10}
\expandafter\def\csname sga@target@IV.14.3.11\endcsname{ega4.c04.prop.14-3-11}
\expandafter\def\csname sga@target@IV.14.3.12\endcsname{ega4.c04.rem.14-3-12}
\expandafter\def\csname sga@target@IV.14.3.13\endcsname{ega4.c04.prop.14-3-13}
\expandafter\def\csname sga@target@IV.14.3.14\endcsname{ega4.c04.rem.14-3-14}
\expandafter\def\csname sga@target@IV.14.3.2\endcsname{ega4.c04.prop.14-3-2}
\expandafter\def\csname sga@target@IV.14.3.3\endcsname{ega4.c04.def.14-3-3}
\expandafter\def\csname sga@target@IV.14.3.3.1\endcsname{ega4.c04.subsec.14-3-3-1}
\expandafter\def\csname sga@target@IV.14.3.4\endcsname{ega4.c04.prop.14-3-4}
\expandafter\def\csname sga@target@IV.14.3.5\endcsname{ega4.c04.prop.14-3-5}
\expandafter\def\csname sga@target@IV.14.3.6\endcsname{ega4.c04.prop.14-3-6}
\expandafter\def\csname sga@target@IV.14.3.7\endcsname{ega4.c04.cor.14-3-7}
\expandafter\def\csname sga@target@IV.14.3.8\endcsname{ega4.c04.prop.14-3-8}
\expandafter\def\csname sga@target@IV.14.3.9\endcsname{ega4.c04.item.14-3-9}
\expandafter\def\csname sga@target@IV.14.4.1\endcsname{ega4.c04.thm.14-4-1}
\expandafter\def\csname sga@target@IV.14.4.1.1\endcsname{ega4.c04.lem.14-4-1-1}
\expandafter\def\csname sga@target@IV.14.4.1.2\endcsname{ega4.c04.lem.14-4-1-2}
\expandafter\def\csname sga@target@IV.14.4.1.3\endcsname{ega4.c04.lem.14-4-1-3}
\expandafter\def\csname sga@target@IV.14.4.10\endcsname{ega4.c04.item.14-4-10}
\expandafter\def\csname sga@target@IV.14.4.2\endcsname{ega4.c04.cor.14-4-2}
\expandafter\def\csname sga@target@IV.14.4.3\endcsname{ega4.c04.cor.14-4-3}
\expandafter\def\csname sga@target@IV.14.4.4\endcsname{ega4.c04.cor.14-4-4}
\expandafter\def\csname sga@target@IV.14.4.5\endcsname{ega4.c04.rem.14-4-5}
\expandafter\def\csname sga@target@IV.14.4.6\endcsname{ega4.c04.cor.14-4-6}
\expandafter\def\csname sga@target@IV.14.4.7\endcsname{ega4.c04.prop.14-4-7}
\expandafter\def\csname sga@target@IV.14.4.7.1\endcsname{ega4.c04.eq.14-4-7-1}
\expandafter\def\csname sga@target@IV.14.4.7.2\endcsname{ega4.c04.eq.14-4-7-2}
\expandafter\def\csname sga@target@IV.14.4.8\endcsname{ega4.c04.prop.14-4-8}
\expandafter\def\csname sga@target@IV.14.4.8.1\endcsname{ega4.c04.lem.14-4-8-1}
\expandafter\def\csname sga@target@IV.14.4.9\endcsname{ega4.c04.cor.14-4-9}
\expandafter\def\csname sga@target@IV.14.5\endcsname{ega4.c04.subsec.14-5.key-iv-14-5}
\expandafter\def\csname sga@target@IV.14.5.1\endcsname{ega4.c04.lem.14-5-1}
\expandafter\def\csname sga@target@IV.14.5.1.1\endcsname{ega4.c04.eq.14-5-1-1}
\expandafter\def\csname sga@target@IV.14.5.1.2\endcsname{ega4.c04.eq.14-5-1-2}
\expandafter\def\csname sga@target@IV.14.5.1.3\endcsname{ega4.c04.eq.14-5-1-3}
\expandafter\def\csname sga@target@IV.14.5.1.4\endcsname{ega4.c04.eq.14-5-1-4}
\expandafter\def\csname sga@target@IV.14.5.10\endcsname{ega4.c04.prop.14-5-10}
\expandafter\def\csname sga@target@IV.14.5.11\endcsname{ega4.c04.cor.14-5-11}
\expandafter\def\csname sga@target@IV.14.5.11.1\endcsname{ega4.c04.lem.14-5-11-1}
\expandafter\def\csname sga@target@IV.14.5.12\endcsname{ega4.c04.cor.14-5-12}
\expandafter\def\csname sga@target@IV.14.5.2\endcsname{ega4.c04.cor.14-5-2}
\expandafter\def\csname sga@target@IV.14.5.2.1\endcsname{ega4.c04.eq.14-5-2-1}
\expandafter\def\csname sga@target@IV.14.5.3\endcsname{ega4.c04.prop.14-5-3}
\expandafter\def\csname sga@target@IV.14.5.4\endcsname{ega4.c04.cor.14-5-4}
\expandafter\def\csname sga@target@IV.14.5.5\endcsname{ega4.c04.rem.14-5-5}
\expandafter\def\csname sga@target@IV.14.5.6\endcsname{ega4.c04.prop.14-5-6}
\expandafter\def\csname sga@target@IV.14.5.7\endcsname{ega4.c04.rem.14-5-7}
\expandafter\def\csname sga@target@IV.14.5.8\endcsname{ega4.c04.prop.14-5-8}
\expandafter\def\csname sga@target@IV.14.5.9\endcsname{ega4.c04.prop.14-5-9}
\expandafter\def\csname sga@target@IV.15.1.1\endcsname{ega4.c04.prop.15-1-1}
\expandafter\def\csname sga@target@IV.15.1.1-proof\endcsname{ega4.c04.statement.15-1-1-proof}
\expandafter\def\csname sga@target@IV.15.1.1.1\endcsname{ega4.c04.eq.15-1-1-1}
\expandafter\def\csname sga@target@IV.15.1.1.2\endcsname{ega4.c04.eq.15-1-1-2}
\expandafter\def\csname sga@target@IV.15.1.1.3\endcsname{ega4.c04.eq.15-1-1-3}
\expandafter\def\csname sga@target@IV.15.1.1.4\endcsname{ega4.c04.eq.15-1-1-4}
\expandafter\def\csname sga@target@IV.15.1.1.5\endcsname{ega4.c04.eq.15-1-1-5}
\expandafter\def\csname sga@target@IV.15.1.1.6\endcsname{ega4.c04.lem.15-1-1-6}
\expandafter\def\csname sga@target@IV.15.1.2\endcsname{ega4.c04.cor.15-1-2}
\expandafter\def\csname sga@target@IV.15.1.3\endcsname{ega4.c04.cor.15-1-3}
\expandafter\def\csname sga@target@IV.15.1.4\endcsname{ega4.c04.statement.15-1-4}
\expandafter\def\csname sga@target@IV.15.2.1\endcsname{ega4.c04.statement.15-2-1}
\expandafter\def\csname sga@target@IV.15.2.2\endcsname{ega4.c04.thm.15-2-2}
\expandafter\def\csname sga@target@IV.15.2.2-proof\endcsname{ega4.c04.item.15-2-2-proof}
\expandafter\def\csname sga@target@IV.15.2.2.1\endcsname{ega4.c04.lem.15-2-2-1}
\expandafter\def\csname sga@target@IV.15.2.3\endcsname{ega4.c04.cor.15-2-3}
\expandafter\def\csname sga@target@IV.15.2.4\endcsname{ega4.c04.item.15-2-4}
\expandafter\def\csname sga@target@IV.15.3.1\endcsname{ega4.c04.prop.15-3-1}
\expandafter\def\csname sga@target@IV.15.3.1-proof\endcsname{ega4.c04.item.15-3-1-proof}
\expandafter\def\csname sga@target@IV.15.3.2\endcsname{ega4.c04.prop.15-3-2}
\expandafter\def\csname sga@target@IV.15.3.2-proof\endcsname{ega4.c04.item.15-3-2-proof}
\expandafter\def\csname sga@target@IV.15.3.3\endcsname{ega4.c04.cor.15-3-3}
\expandafter\def\csname sga@target@IV.15.3.3-proof\endcsname{ega4.c04.item.15-3-3-proof}
\expandafter\def\csname sga@target@IV.15.3.4\endcsname{ega4.c04.item.15-3-4}
\expandafter\def\csname sga@target@IV.15.4.1\endcsname{ega4.c04.prop.15-4-1}
\expandafter\def\csname sga@target@IV.15.4.2\endcsname{ega4.c04.prop.15-4-2}
\expandafter\def\csname sga@target@IV.15.4.3\endcsname{ega4.c04.prop.15-4-3}
\expandafter\def\csname sga@target@IV.15.4.3-proof\endcsname{ega4.c04.item.15-4-3-proof}
\expandafter\def\csname sga@target@IV.15.5.1\endcsname{ega4.c04.prop.15-5-1}
\expandafter\def\csname sga@target@IV.15.5.10\endcsname{ega4.c04.item.15-5-10}
\expandafter\def\csname sga@target@IV.15.5.2\endcsname{ega4.c04.cor.15-5-2}
\expandafter\def\csname sga@target@IV.15.5.2.1\endcsname{ega4.c04.eq.15-5-2-1}
\expandafter\def\csname sga@target@IV.15.5.2.2\endcsname{ega4.c04.lem.15-5-2-2}
\expandafter\def\csname sga@target@IV.15.5.2.3\endcsname{ega4.c04.eq.15-5-2-3}
\expandafter\def\csname sga@target@IV.15.5.2.4\endcsname{ega4.c04.eq.15-5-2-4}
\expandafter\def\csname sga@target@IV.15.5.3\endcsname{ega4.c04.prop.15-5-3}
\expandafter\def\csname sga@target@IV.15.5.4\endcsname{ega4.c04.cor.15-5-4}
\expandafter\def\csname sga@target@IV.15.5.5\endcsname{ega4.c04.statement.15-5-5}
\expandafter\def\csname sga@target@IV.15.5.6\endcsname{ega4.c04.lem.15-5-6}
\expandafter\def\csname sga@target@IV.15.5.6.1\endcsname{ega4.c04.lem.15-5-6-1}
\expandafter\def\csname sga@target@IV.15.5.7\endcsname{ega4.c04.prop.15-5-7}
\expandafter\def\csname sga@target@IV.15.5.7.1\endcsname{ega4.c04.eq.15-5-7-1}
\expandafter\def\csname sga@target@IV.15.5.8\endcsname{ega4.c04.statement.15-5-8}
\expandafter\def\csname sga@target@IV.15.5.9\endcsname{ega4.c04.prop.15-5-9}
\expandafter\def\csname sga@target@IV.15.6.1\endcsname{ega4.c04.prop.15-6-1}
\expandafter\def\csname sga@target@IV.15.6.1.1\endcsname{ega4.c04.eq.15-6-1-1}
\expandafter\def\csname sga@target@IV.15.6.10\endcsname{ega4.c04.item.15-6-10}
\expandafter\def\csname sga@target@IV.15.6.2\endcsname{ega4.c04.cor.15-6-2}
\expandafter\def\csname sga@target@IV.15.6.3\endcsname{ega4.c04.prop.15-6-3}
\expandafter\def\csname sga@target@IV.15.6.4\endcsname{ega4.c04.prop.15-6-4}
\expandafter\def\csname sga@target@IV.15.6.5\endcsname{ega4.c04.cor.15-6-5}
\expandafter\def\csname sga@target@IV.15.6.6\endcsname{ega4.c04.prop.15-6-6}
\expandafter\def\csname sga@target@IV.15.6.7\endcsname{ega4.c04.cor.15-6-7}
\expandafter\def\csname sga@target@IV.15.6.8\endcsname{ega4.c04.prop.15-6-8}
\expandafter\def\csname sga@target@IV.15.6.9\endcsname{ega4.c04.cor.15-6-9}
\expandafter\def\csname sga@target@IV.15.7.1\endcsname{ega4.c04.def.15-7-1}
\expandafter\def\csname sga@target@IV.15.7.10\endcsname{ega4.c04.cor.15-7-10}
\expandafter\def\csname sga@target@IV.15.7.11\endcsname{ega4.c04.cor.15-7-11}
\expandafter\def\csname sga@target@IV.15.7.11.1\endcsname{ega4.c04.lem.15-7-11-1}
\expandafter\def\csname sga@target@IV.15.7.2\endcsname{ega4.c04.prop.15-7-2}
\expandafter\def\csname sga@target@IV.15.7.3\endcsname{ega4.c04.rem.15-7-3}
\expandafter\def\csname sga@target@IV.15.7.4\endcsname{ega4.c04.cor.15-7-4}
\expandafter\def\csname sga@target@IV.15.7.5\endcsname{ega4.c04.cor.15-7-5}
\expandafter\def\csname sga@target@IV.15.7.6\endcsname{ega4.c04.cor.15-7-6}
\expandafter\def\csname sga@target@IV.15.7.7\endcsname{ega4.c04.item.15-7-7}
\expandafter\def\csname sga@target@IV.15.7.8\endcsname{ega4.c04.statement.15-7-8}
\expandafter\def\csname sga@target@IV.15.7.9\endcsname{ega4.c04.prop.15-7-9}
\expandafter\def\csname sga@target@IV.16.1\endcsname{ega4.c04.subsec.16-1}
\expandafter\def\csname sga@target@IV.16.1.1\endcsname{ega4.c04.statement.16-1-1}
\expandafter\def\csname sga@target@IV.16.1.10\endcsname{ega4.c04.cor.16-1-10}
\expandafter\def\csname sga@target@IV.16.1.11\endcsname{ega4.c04.rem.16-1-11}
\expandafter\def\csname sga@target@IV.16.1.2\endcsname{ega4.c04.def.16-1-2}
\expandafter\def\csname sga@target@IV.16.1.2.1\endcsname{ega4.c04.def.16-1-2-1}
\expandafter\def\csname sga@target@IV.16.1.2.2\endcsname{ega4.c04.eq.16-1-2-2}
\expandafter\def\csname sga@target@IV.16.1.3\endcsname{ega4.c04.statement.16-1-3}
\expandafter\def\csname sga@target@IV.16.1.4\endcsname{ega4.c04.statement.16-1-4}
\expandafter\def\csname sga@target@IV.16.1.5\endcsname{ega4.c04.prop.16-1-5}
\expandafter\def\csname sga@target@IV.16.1.6\endcsname{ega4.c04.cor.16-1-6}
\expandafter\def\csname sga@target@IV.16.1.7\endcsname{ega4.c04.cor.16-1-7}
\expandafter\def\csname sga@target@IV.16.1.8\endcsname{ega4.c04.prop.16-1-8}
\expandafter\def\csname sga@target@IV.16.1.9\endcsname{ega4.c04.prop.16-1-9}
\expandafter\def\csname sga@target@IV.16.10\endcsname{ega4.c04.subsec.16-10}
\expandafter\def\csname sga@target@IV.16.10.1\endcsname{ega4.c04.def.16-10-1}
\expandafter\def\csname sga@target@IV.16.10.2\endcsname{ega4.c04.prop.16-10-2}
\expandafter\def\csname sga@target@IV.16.10.3\endcsname{ega4.c04.prop.16-10-3}
\expandafter\def\csname sga@target@IV.16.10.4\endcsname{ega4.c04.prop.16-10-4}
\expandafter\def\csname sga@target@IV.16.10.5\endcsname{ega4.c04.prop.16-10-5}
\expandafter\def\csname sga@target@IV.16.10.6\endcsname{ega4.c04.rem.16-10-6}
\expandafter\def\csname sga@target@IV.16.11\endcsname{ega4.c04.subsec.16-11}
\expandafter\def\csname sga@target@IV.16.11.1\endcsname{ega4.c04.statement.16-11-1}
\expandafter\def\csname sga@target@IV.16.11.1.1\endcsname{ega4.c04.statement.16-11-1-1}
\expandafter\def\csname sga@target@IV.16.11.1.2\endcsname{ega4.c04.statement.16-11-1-2}
\expandafter\def\csname sga@target@IV.16.11.1.3\endcsname{ega4.c04.eq.16-11-1-3}
\expandafter\def\csname sga@target@IV.16.11.1.4\endcsname{ega4.c04.eq.16-11-1-4}
\expandafter\def\csname sga@target@IV.16.11.1.5\endcsname{ega4.c04.eq.16-11-1-5}
\expandafter\def\csname sga@target@IV.16.11.1.6\endcsname{ega4.c04.eq.16-11-1-6}
\expandafter\def\csname sga@target@IV.16.11.1.7\endcsname{ega4.c04.eq.16-11-1-7}
\expandafter\def\csname sga@target@IV.16.11.2\endcsname{ega4.c04.thm.16-11-2}
\expandafter\def\csname sga@target@IV.16.11.2.1\endcsname{ega4.c04.eq.16-11-2-1}
\expandafter\def\csname sga@target@IV.16.11.2.2\endcsname{ega4.c04.eq.16-11-2-2}
\expandafter\def\csname sga@target@IV.16.11.2.3\endcsname{ega4.c04.eq.16-11-2-3}
\expandafter\def\csname sga@target@IV.16.11.3\endcsname{ega4.c04.statement.16-11-3}
\expandafter\def\csname sga@target@IV.16.12\endcsname{ega4.c04.subsec.16-12}
\expandafter\def\csname sga@target@IV.16.12.1\endcsname{ega4.c04.statement.16-12-1}
\expandafter\def\csname sga@target@IV.16.12.2\endcsname{ega4.c04.thm.16-12-2}
\expandafter\def\csname sga@target@IV.16.12.3\endcsname{ega4.c04.statement.16-12-3}
\expandafter\def\csname sga@target@IV.16.2\endcsname{ega4.c04.subsec.16-2}
\expandafter\def\csname sga@target@IV.16.2.1\endcsname{ega4.c04.statement.16-2-1}
\expandafter\def\csname sga@target@IV.16.2.1.1\endcsname{ega4.c04.statement.16-2-1-1}
\expandafter\def\csname sga@target@IV.16.2.1.2\endcsname{ega4.c04.eq.16-2-1-2}
\expandafter\def\csname sga@target@IV.16.2.1.3\endcsname{ega4.c04.eq.16-2-1-3}
\expandafter\def\csname sga@target@IV.16.2.1.4\endcsname{ega4.c04.eq.16-2-1-4}
\expandafter\def\csname sga@target@IV.16.2.2\endcsname{ega4.c04.prop.16-2-2}
\expandafter\def\csname sga@target@IV.16.2.3\endcsname{ega4.c04.cor.16-2-3}
\expandafter\def\csname sga@target@IV.16.2.3.1\endcsname{ega4.c04.eq.16-2-3-1}
\expandafter\def\csname sga@target@IV.16.2.3.2\endcsname{ega4.c04.eq.16-2-3-2}
\expandafter\def\csname sga@target@IV.16.2.4\endcsname{ega4.c04.cor.16-2-4}
\expandafter\def\csname sga@target@IV.16.2.5\endcsname{ega4.c04.statement.16-2-5}
\expandafter\def\csname sga@target@IV.16.2.6\endcsname{ega4.c04.statement.16-2-6}
\expandafter\def\csname sga@target@IV.16.2.6.1\endcsname{ega4.c04.statement.16-2-6-1}
\expandafter\def\csname sga@target@IV.16.2.6.2\endcsname{ega4.c04.eq.16-2-6-2}
\expandafter\def\csname sga@target@IV.16.2.7\endcsname{ega4.c04.prop.16-2-7}
\expandafter\def\csname sga@target@IV.16.2.7.1\endcsname{ega4.c04.prop.16-2-7-1}
\expandafter\def\csname sga@target@IV.16.3\endcsname{ega4.c04.subsec.16-3}
\expandafter\def\csname sga@target@IV.16.3.1\endcsname{ega4.c04.def.16-3-1}
\expandafter\def\csname sga@target@IV.16.3.1.1\endcsname{ega4.c04.eq.16-3-1-1}
\expandafter\def\csname sga@target@IV.16.3.2\endcsname{ega4.c04.statement.16-3-2}
\expandafter\def\csname sga@target@IV.16.3.3\endcsname{ega4.c04.statement.16-3-3}
\expandafter\def\csname sga@target@IV.16.3.3.1\endcsname{ega4.c04.statement.16-3-3-1}
\expandafter\def\csname sga@target@IV.16.3.4\endcsname{ega4.c04.prop.16-3-4}
\expandafter\def\csname sga@target@IV.16.3.5\endcsname{ega4.c04.statement.16-3-5}
\expandafter\def\csname sga@target@IV.16.3.6\endcsname{ega4.c04.def.16-3-6}
\expandafter\def\csname sga@target@IV.16.3.6.1\endcsname{ega4.c04.eq.16-3-6-1}
\expandafter\def\csname sga@target@IV.16.3.7\endcsname{ega4.c04.statement.16-3-7}
\expandafter\def\csname sga@target@IV.16.3.7.1\endcsname{ega4.c04.eq.16-3-7-1}
\expandafter\def\csname sga@target@IV.16.3.8\endcsname{ega4.c04.prop.16-3-8}
\expandafter\def\csname sga@target@IV.16.3.9\endcsname{ega4.c04.prop.16-3-9}
\expandafter\def\csname sga@target@IV.16.4\endcsname{ega4.c04.subsec.16-4}
\expandafter\def\csname sga@target@IV.16.4.1\endcsname{ega4.c04.statement.16-4-1}
\expandafter\def\csname sga@target@IV.16.4.1.1\endcsname{ega4.c04.statement.16-4-1-1}
\expandafter\def\csname sga@target@IV.16.4.1.2\endcsname{ega4.c04.statement.16-4-1-2}
\expandafter\def\csname sga@target@IV.16.4.1.3\endcsname{ega4.c04.eq.16-4-1-3}
\expandafter\def\csname sga@target@IV.16.4.1.4\endcsname{ega4.c04.eq.16-4-1-4}
\expandafter\def\csname sga@target@IV.16.4.1.5\endcsname{ega4.c04.eq.16-4-1-5}
\expandafter\def\csname sga@target@IV.16.4.10\endcsname{ega4.c04.cor.16-4-10}
\expandafter\def\csname sga@target@IV.16.4.11\endcsname{ega4.c04.prop.16-4-11}
\expandafter\def\csname sga@target@IV.16.4.11.1\endcsname{ega4.c04.prop.16-4-11-1}
\expandafter\def\csname sga@target@IV.16.4.11.2\endcsname{ega4.c04.prop.16-4-11-2}
\expandafter\def\csname sga@target@IV.16.4.11.3\endcsname{ega4.c04.eq.16-4-11-3}
\expandafter\def\csname sga@target@IV.16.4.11.4\endcsname{ega4.c04.eq.16-4-11-4}
\expandafter\def\csname sga@target@IV.16.4.12\endcsname{ega4.c04.cor.16-4-12}
\expandafter\def\csname sga@target@IV.16.4.13\endcsname{ega4.c04.cor.16-4-13}
\expandafter\def\csname sga@target@IV.16.4.14\endcsname{ega4.c04.prop.16-4-14}
\expandafter\def\csname sga@target@IV.16.4.14.1\endcsname{ega4.c04.prop.16-4-14-1}
\expandafter\def\csname sga@target@IV.16.4.15\endcsname{ega4.c04.cor.16-4-15}
\expandafter\def\csname sga@target@IV.16.4.15.1\endcsname{ega4.c04.cor.16-4-15-1}
\expandafter\def\csname sga@target@IV.16.4.15.2\endcsname{ega4.c04.eq.16-4-15-2}
\expandafter\def\csname sga@target@IV.16.4.16\endcsname{ega4.c04.cor.16-4-16}
\expandafter\def\csname sga@target@IV.16.4.16.1\endcsname{ega4.c04.cor.16-4-16-1}
\expandafter\def\csname sga@target@IV.16.4.16.2\endcsname{ega4.c04.eq.16-4-16-2}
\expandafter\def\csname sga@target@IV.16.4.17\endcsname{ega4.c04.cor.16-4-17}
\expandafter\def\csname sga@target@IV.16.4.18\endcsname{ega4.c04.statement.16-4-18}
\expandafter\def\csname sga@target@IV.16.4.18.1\endcsname{ega4.c04.statement.16-4-18-1}
\expandafter\def\csname sga@target@IV.16.4.18.2\endcsname{ega4.c04.statement.16-4-18-2}
\expandafter\def\csname sga@target@IV.16.4.18.3\endcsname{ega4.c04.eq.16-4-18-3}
\expandafter\def\csname sga@target@IV.16.4.19\endcsname{ega4.c04.cor.16-4-19}
\expandafter\def\csname sga@target@IV.16.4.19.1\endcsname{ega4.c04.cor.16-4-19-1}
\expandafter\def\csname sga@target@IV.16.4.2\endcsname{ega4.c04.statement.16-4-2}
\expandafter\def\csname sga@target@IV.16.4.20\endcsname{ega4.c04.prop.16-4-20}
\expandafter\def\csname sga@target@IV.16.4.21\endcsname{ega4.c04.cor.16-4-21}
\expandafter\def\csname sga@target@IV.16.4.21.1\endcsname{ega4.c04.cor.16-4-21-1}
\expandafter\def\csname sga@target@IV.16.4.22\endcsname{ega4.c04.cor.16-4-22}
\expandafter\def\csname sga@target@IV.16.4.23\endcsname{ega4.c04.prop.16-4-23}
\expandafter\def\csname sga@target@IV.16.4.23.1\endcsname{ega4.c04.prop.16-4-23-1}
\expandafter\def\csname sga@target@IV.16.4.24\endcsname{ega4.c04.statement.16-4-24}
\expandafter\def\csname sga@target@IV.16.4.3\endcsname{ega4.c04.statement.16-4-3}
\expandafter\def\csname sga@target@IV.16.4.3.1\endcsname{ega4.c04.statement.16-4-3-1}
\expandafter\def\csname sga@target@IV.16.4.3.2\endcsname{ega4.c04.statement.16-4-3-2}
\expandafter\def\csname sga@target@IV.16.4.3.3\endcsname{ega4.c04.eq.16-4-3-3}
\expandafter\def\csname sga@target@IV.16.4.3.4\endcsname{ega4.c04.eq.16-4-3-4}
\expandafter\def\csname sga@target@IV.16.4.3.5\endcsname{ega4.c04.eq.16-4-3-5}
\expandafter\def\csname sga@target@IV.16.4.3.6\endcsname{ega4.c04.eq.16-4-3-6}
\expandafter\def\csname sga@target@IV.16.4.3.7\endcsname{ega4.c04.eq.16-4-3-7}
\expandafter\def\csname sga@target@IV.16.4.4\endcsname{ega4.c04.statement.16-4-4}
\expandafter\def\csname sga@target@IV.16.4.5\endcsname{ega4.c04.prop.16-4-5}
\expandafter\def\csname sga@target@IV.16.4.6\endcsname{ega4.c04.cor.16-4-6}
\expandafter\def\csname sga@target@IV.16.4.7\endcsname{ega4.c04.statement.16-4-7}
\expandafter\def\csname sga@target@IV.16.4.7.1\endcsname{ega4.c04.statement.16-4-7-1}
\expandafter\def\csname sga@target@IV.16.4.7.2\endcsname{ega4.c04.eq.16-4-7-2}
\expandafter\def\csname sga@target@IV.16.4.7.3\endcsname{ega4.c04.eq.16-4-7-3}
\expandafter\def\csname sga@target@IV.16.4.8\endcsname{ega4.c04.prop.16-4-8}
\expandafter\def\csname sga@target@IV.16.4.8.1\endcsname{ega4.c04.prop.16-4-8-1}
\expandafter\def\csname sga@target@IV.16.4.8.2\endcsname{ega4.c04.eq.16-4-8-2}
\expandafter\def\csname sga@target@IV.16.4.8.3\endcsname{ega4.c04.lem.16-4-8-3}
\expandafter\def\csname sga@target@IV.16.4.8.4\endcsname{ega4.c04.eq.16-4-8-4}
\expandafter\def\csname sga@target@IV.16.4.9\endcsname{ega4.c04.cor.16-4-9}
\expandafter\def\csname sga@target@IV.16.4.9.1\endcsname{ega4.c04.cor.16-4-9-1}
\expandafter\def\csname sga@target@IV.16.5\endcsname{ega4.c04.subsec.16-5}
\expandafter\def\csname sga@target@IV.16.5.1\endcsname{ega4.c04.statement.16-5-1}
\expandafter\def\csname sga@target@IV.16.5.1.1\endcsname{ega4.c04.eq.16-5-1-1}
\expandafter\def\csname sga@target@IV.16.5.1.2\endcsname{ega4.c04.eq.16-5-1-2}
\expandafter\def\csname sga@target@IV.16.5.10\endcsname{ega4.c04.statement.16-5-10}
\expandafter\def\csname sga@target@IV.16.5.10.1\endcsname{ega4.c04.statement.16-5-10-1}
\expandafter\def\csname sga@target@IV.16.5.10.2\endcsname{ega4.c04.statement.16-5-10-2}
\expandafter\def\csname sga@target@IV.16.5.11\endcsname{ega4.c04.prop.16-5-11}
\expandafter\def\csname sga@target@IV.16.5.12\endcsname{ega4.c04.statement.16-5-12}
\expandafter\def\csname sga@target@IV.16.5.12.1\endcsname{ega4.c04.statement.16-5-12-1}
\expandafter\def\csname sga@target@IV.16.5.12.2\endcsname{ega4.c04.eq.16-5-12-2}
\expandafter\def\csname sga@target@IV.16.5.13\endcsname{ega4.c04.statement.16-5-13}
\expandafter\def\csname sga@target@IV.16.5.13.1\endcsname{ega4.c04.statement.16-5-13-1}
\expandafter\def\csname sga@target@IV.16.5.13.2\endcsname{ega4.c04.eq.16-5-13-2}
\expandafter\def\csname sga@target@IV.16.5.13.3\endcsname{ega4.c04.eq.16-5-13-3}
\expandafter\def\csname sga@target@IV.16.5.13.4\endcsname{ega4.c04.eq.16-5-13-4}
\expandafter\def\csname sga@target@IV.16.5.13.5\endcsname{ega4.c04.eq.16-5-13-5}
\expandafter\def\csname sga@target@IV.16.5.14\endcsname{ega4.c04.statement.16-5-14}
\expandafter\def\csname sga@target@IV.16.5.14.1\endcsname{ega4.c04.statement.16-5-14-1}
\expandafter\def\csname sga@target@IV.16.5.14.2\endcsname{ega4.c04.eq.16-5-14-2}
\expandafter\def\csname sga@target@IV.16.5.14.3\endcsname{ega4.c04.eq.16-5-14-3}
\expandafter\def\csname sga@target@IV.16.5.15\endcsname{ega4.c04.statement.16-5-15}
\expandafter\def\csname sga@target@IV.16.5.16\endcsname{ega4.c04.prop.16-5-16}
\expandafter\def\csname sga@target@IV.16.5.17\endcsname{ega4.c04.prop.16-5-17}
\expandafter\def\csname sga@target@IV.16.5.18\endcsname{ega4.c04.cor.16-5-18}
\expandafter\def\csname sga@target@IV.16.5.2\endcsname{ega4.c04.prop.16-5-2}
\expandafter\def\csname sga@target@IV.16.5.3\endcsname{ega4.c04.prop.16-5-3}
\expandafter\def\csname sga@target@IV.16.5.3.1\endcsname{ega4.c04.eq.16-5-3-1}
\expandafter\def\csname sga@target@IV.16.5.4\endcsname{ega4.c04.statement.16-5-4}
\expandafter\def\csname sga@target@IV.16.5.5\endcsname{ega4.c04.cor.16-5-5}
\expandafter\def\csname sga@target@IV.16.5.5.1\endcsname{ega4.c04.cor.16-5-5-1}
\expandafter\def\csname sga@target@IV.16.5.6\endcsname{ega4.c04.cor.16-5-6}
\expandafter\def\csname sga@target@IV.16.5.7\endcsname{ega4.c04.statement.16-5-7}
\expandafter\def\csname sga@target@IV.16.5.7.1\endcsname{ega4.c04.statement.16-5-7-1}
\expandafter\def\csname sga@target@IV.16.5.8\endcsname{ega4.c04.statement.16-5-8}
\expandafter\def\csname sga@target@IV.16.5.8.1\endcsname{ega4.c04.eq.16-5-8-1}
\expandafter\def\csname sga@target@IV.16.5.8.2\endcsname{ega4.c04.eq.16-5-8-2}
\expandafter\def\csname sga@target@IV.16.5.9\endcsname{ega4.c04.statement.16-5-9}
\expandafter\def\csname sga@target@IV.16.5.9.1\endcsname{ega4.c04.statement.16-5-9-1}
\expandafter\def\csname sga@target@IV.16.6\endcsname{ega4.c04.subsec.16-6}
\expandafter\def\csname sga@target@IV.16.6.1\endcsname{ega4.c04.statement.16-6-1}
\expandafter\def\csname sga@target@IV.16.6.1.1\endcsname{ega4.c04.statement.16-6-1-1}
\expandafter\def\csname sga@target@IV.16.6.1.2\endcsname{ega4.c04.statement.16-6-1-2}
\expandafter\def\csname sga@target@IV.16.6.2\endcsname{ega4.c04.thm.16-6-2}
\expandafter\def\csname sga@target@IV.16.6.2.1\endcsname{ega4.c04.eq.16-6-2-1}
\expandafter\def\csname sga@target@IV.16.6.2.2\endcsname{ega4.c04.eq.16-6-2-2}
\expandafter\def\csname sga@target@IV.16.6.2.3\endcsname{ega4.c04.eq.16-6-2-3}
\expandafter\def\csname sga@target@IV.16.6.2.4\endcsname{ega4.c04.eq.16-6-2-4}
\expandafter\def\csname sga@target@IV.16.6.2.5\endcsname{ega4.c04.eq.16-6-2-5}
\expandafter\def\csname sga@target@IV.16.6.3\endcsname{ega4.c04.def.16-6-3}
\expandafter\def\csname sga@target@IV.16.6.4\endcsname{ega4.c04.prop.16-6-4}
\expandafter\def\csname sga@target@IV.16.6.4.1\endcsname{ega4.c04.prop.16-6-4-1}
\expandafter\def\csname sga@target@IV.16.6.5\endcsname{ega4.c04.statement.16-6-5}
\expandafter\def\csname sga@target@IV.16.6.5.1\endcsname{ega4.c04.statement.16-6-5-1}
\expandafter\def\csname sga@target@IV.16.7\endcsname{ega4.c04.subsec.16-7}
\expandafter\def\csname sga@target@IV.16.7.1\endcsname{ega4.c04.statement.16-7-1}
\expandafter\def\csname sga@target@IV.16.7.1.1\endcsname{ega4.c04.eq.16-7-1-1}
\expandafter\def\csname sga@target@IV.16.7.1.2\endcsname{ega4.c04.eq.16-7-1-2}
\expandafter\def\csname sga@target@IV.16.7.10\endcsname{ega4.c04.rem.16-7-10}
\expandafter\def\csname sga@target@IV.16.7.2\endcsname{ega4.c04.statement.16-7-2}
\expandafter\def\csname sga@target@IV.16.7.2.1\endcsname{ega4.c04.statement.16-7-2-1}
\expandafter\def\csname sga@target@IV.16.7.3\endcsname{ega4.c04.prop.16-7-3}
\expandafter\def\csname sga@target@IV.16.7.4\endcsname{ega4.c04.statement.16-7-4}
\expandafter\def\csname sga@target@IV.16.7.4.1\endcsname{ega4.c04.statement.16-7-4-1}
\expandafter\def\csname sga@target@IV.16.7.5\endcsname{ega4.c04.statement.16-7-5}
\expandafter\def\csname sga@target@IV.16.7.5.1\endcsname{ega4.c04.statement.16-7-5-1}
\expandafter\def\csname sga@target@IV.16.7.5.2\endcsname{ega4.c04.statement.16-7-5-2}
\expandafter\def\csname sga@target@IV.16.7.5.3\endcsname{ega4.c04.statement.16-7-5-3}
\expandafter\def\csname sga@target@IV.16.7.5.4\endcsname{ega4.c04.statement.16-7-5-4}
\expandafter\def\csname sga@target@IV.16.7.6\endcsname{ega4.c04.prop.16-7-6}
\expandafter\def\csname sga@target@IV.16.7.7\endcsname{ega4.c04.statement.16-7-7}
\expandafter\def\csname sga@target@IV.16.7.7.1\endcsname{ega4.c04.eq.16-7-7-1}
\expandafter\def\csname sga@target@IV.16.7.7.2\endcsname{ega4.c04.eq.16-7-7-2}
\expandafter\def\csname sga@target@IV.16.7.8\endcsname{ega4.c04.statement.16-7-8}
\expandafter\def\csname sga@target@IV.16.7.8.1\endcsname{ega4.c04.statement.16-7-8-1}
\expandafter\def\csname sga@target@IV.16.7.8.2\endcsname{ega4.c04.eq.16-7-8-2}
\expandafter\def\csname sga@target@IV.16.7.9\endcsname{ega4.c04.statement.16-7-9}
\expandafter\def\csname sga@target@IV.16.7.9.1\endcsname{ega4.c04.statement.16-7-9-1}
\expandafter\def\csname sga@target@IV.16.8\endcsname{ega4.c04.subsec.16-8}
\expandafter\def\csname sga@target@IV.16.8.1\endcsname{ega4.c04.def.16-8-1}
\expandafter\def\csname sga@target@IV.16.8.10\endcsname{ega4.c04.cor.16-8-10}
\expandafter\def\csname sga@target@IV.16.8.11\endcsname{ega4.c04.statement.16-8-11}
\expandafter\def\csname sga@target@IV.16.8.2\endcsname{ega4.c04.statement.16-8-2}
\expandafter\def\csname sga@target@IV.16.8.3\endcsname{ega4.c04.statement.16-8-3}
\expandafter\def\csname sga@target@IV.16.8.3.1\endcsname{ega4.c04.eq.16-8-3-1}
\expandafter\def\csname sga@target@IV.16.8.4\endcsname{ega4.c04.prop.16-8-4}
\expandafter\def\csname sga@target@IV.16.8.4.1\endcsname{ega4.c04.prop.16-8-4-1}
\expandafter\def\csname sga@target@IV.16.8.5\endcsname{ega4.c04.statement.16-8-5}
\expandafter\def\csname sga@target@IV.16.8.5.1\endcsname{ega4.c04.eq.16-8-5-1}
\expandafter\def\csname sga@target@IV.16.8.5.2\endcsname{ega4.c04.eq.16-8-5-2}
\expandafter\def\csname sga@target@IV.16.8.6\endcsname{ega4.c04.prop.16-8-6}
\expandafter\def\csname sga@target@IV.16.8.7\endcsname{ega4.c04.statement.16-8-7}
\expandafter\def\csname sga@target@IV.16.8.7.1\endcsname{ega4.c04.eq.16-8-7-1}
\expandafter\def\csname sga@target@IV.16.8.8\endcsname{ega4.c04.prop.16-8-8}
\expandafter\def\csname sga@target@IV.16.8.8.1\endcsname{ega4.c04.eq.16-8-8-1}
\expandafter\def\csname sga@target@IV.16.8.8.2\endcsname{ega4.c04.eq.16-8-8-2}
\expandafter\def\csname sga@target@IV.16.8.9\endcsname{ega4.c04.prop.16-8-9}
\expandafter\def\csname sga@target@IV.16.8.9.1\endcsname{ega4.c04.lem.16-8-9-1}
\expandafter\def\csname sga@target@IV.16.8.9.2\endcsname{ega4.c04.lem.16-8-9-2}
\expandafter\def\csname sga@target@IV.16.8.9.3\endcsname{ega4.c04.lem.16-8-9-3}
\expandafter\def\csname sga@target@IV.16.8.9.4\endcsname{ega4.c04.eq.16-8-9-4}
\expandafter\def\csname sga@target@IV.16.9\endcsname{ega4.c04.subsec.16-9}
\expandafter\def\csname sga@target@IV.16.9.1\endcsname{ega4.c04.def.16-9-1}
\expandafter\def\csname sga@target@IV.16.9.10\endcsname{ega4.c04.prop.16-9-10}
\expandafter\def\csname sga@target@IV.16.9.11\endcsname{ega4.c04.cor.16-9-11}
\expandafter\def\csname sga@target@IV.16.9.12\endcsname{ega4.c04.statement.16-9-12}
\expandafter\def\csname sga@target@IV.16.9.13\endcsname{ega4.c04.prop.16-9-13}
\expandafter\def\csname sga@target@IV.16.9.13.1\endcsname{ega4.c04.prop.16-9-13-1}
\expandafter\def\csname sga@target@IV.16.9.13.2\endcsname{ega4.c04.lem.16-9-13-2}
\expandafter\def\csname sga@target@IV.16.9.13.3\endcsname{ega4.c04.lem.16-9-13-3}
\expandafter\def\csname sga@target@IV.16.9.13.4\endcsname{ega4.c04.eq.16-9-13-4}
\expandafter\def\csname sga@target@IV.16.9.13.5\endcsname{ega4.c04.cor.16-9-13-5}
\expandafter\def\csname sga@target@IV.16.9.13.6\endcsname{ega4.c04.cor.16-9-13-6}
\expandafter\def\csname sga@target@IV.16.9.13.7\endcsname{ega4.c04.eq.16-9-13-7}
\expandafter\def\csname sga@target@IV.16.9.2\endcsname{ega4.c04.def.16-9-2}
\expandafter\def\csname sga@target@IV.16.9.3\endcsname{ega4.c04.prop.16-9-3}
\expandafter\def\csname sga@target@IV.16.9.4\endcsname{ega4.c04.cor.16-9-4}
\expandafter\def\csname sga@target@IV.16.9.4.1\endcsname{ega4.c04.eq.16-9-4-1}
\expandafter\def\csname sga@target@IV.16.9.5\endcsname{ega4.c04.cor.16-9-5}
\expandafter\def\csname sga@target@IV.16.9.6\endcsname{ega4.c04.statement.16-9-6}
\expandafter\def\csname sga@target@IV.16.9.7\endcsname{ega4.c04.statement.16-9-7}
\expandafter\def\csname sga@target@IV.16.9.8\endcsname{ega4.c04.prop.16-9-8}
\expandafter\def\csname sga@target@IV.16.9.9\endcsname{ega4.c04.cor.16-9-9}
\expandafter\def\csname sga@target@IV.17.1.1\endcsname{ega4.c04.def.17-1-1}
\expandafter\def\csname sga@target@IV.17.1.1.1\endcsname{ega4.c04.def.17-1-1-1}
\expandafter\def\csname sga@target@IV.17.1.2\endcsname{ega4.c04.rem.17-1-2}
\expandafter\def\csname sga@target@IV.17.1.2.1\endcsname{ega4.c04.eq.17-1-2-1}
\expandafter\def\csname sga@target@IV.17.1.3\endcsname{ega4.c04.prop.17-1-3}
\expandafter\def\csname sga@target@IV.17.1.4\endcsname{ega4.c04.prop.17-1-4}
\expandafter\def\csname sga@target@IV.17.1.5\endcsname{ega4.c04.cor.17-1-5}
\expandafter\def\csname sga@target@IV.17.1.6\endcsname{ega4.c04.prop.17-1-6}
\expandafter\def\csname sga@target@IV.17.10.1\endcsname{ega4.c04.def.17-10-1}
\expandafter\def\csname sga@target@IV.17.10.2\endcsname{ega4.c04.prop.17-10-2}
\expandafter\def\csname sga@target@IV.17.10.3\endcsname{ega4.c04.cor.17-10-3}
\expandafter\def\csname sga@target@IV.17.10.3.1\endcsname{ega4.c04.cor.17-10-3-1}
\expandafter\def\csname sga@target@IV.17.10.4\endcsname{ega4.c04.cor.17-10-4}
\expandafter\def\csname sga@target@IV.17.10.4.1\endcsname{ega4.c04.eq.17-10-4-1}
\expandafter\def\csname sga@target@IV.17.11.1\endcsname{ega4.c04.thm.17-11-1}
\expandafter\def\csname sga@target@IV.17.11.1.1\endcsname{ega4.c04.eq.17-11-1-1}
\expandafter\def\csname sga@target@IV.17.11.1.2\endcsname{ega4.c04.eq.17-11-1-2}
\expandafter\def\csname sga@target@IV.17.11.2\endcsname{ega4.c04.cor.17-11-2}
\expandafter\def\csname sga@target@IV.17.11.3\endcsname{ega4.c04.cor.17-11-3}
\expandafter\def\csname sga@target@IV.17.11.4\endcsname{ega4.c04.cor.17-11-4}
\expandafter\def\csname sga@target@IV.17.11.5\endcsname{ega4.c04.prop.17-11-5}
\expandafter\def\csname sga@target@IV.17.12.1\endcsname{ega4.c04.thm.17-12-1}
\expandafter\def\csname sga@target@IV.17.12.2\endcsname{ega4.c04.cor.17-12-2}
\expandafter\def\csname sga@target@IV.17.12.3\endcsname{ega4.c04.cor.17-12-3}
\expandafter\def\csname sga@target@IV.17.12.4\endcsname{ega4.c04.prop.17-12-4}
\expandafter\def\csname sga@target@IV.17.12.5\endcsname{ega4.c04.prop.17-12-5}
\expandafter\def\csname sga@target@IV.17.12.6\endcsname{ega4.c04.cor.17-12-6}
\expandafter\def\csname sga@target@IV.17.12.7\endcsname{ega4.c04.ex.17-12-7}
\expandafter\def\csname sga@target@IV.17.12.8\endcsname{ega4.c04.rem.17-12-8}
\expandafter\def\csname sga@target@IV.17.13\endcsname{ega4.c04.subsec.17-13.key-iv-17-13}
\expandafter\def\csname sga@target@IV.17.13.1\endcsname{ega4.c04.statement.17-13-1}
\expandafter\def\csname sga@target@IV.17.13.1.1\endcsname{ega4.c04.statement.17-13-1-1}
\expandafter\def\csname sga@target@IV.17.13.1.2\endcsname{ega4.c04.statement.17-13-1-2}
\expandafter\def\csname sga@target@IV.17.13.1.3\endcsname{ega4.c04.eq.17-13-1-3}
\expandafter\def\csname sga@target@IV.17.13.10\endcsname{ega4.c04.cor.17-13-10}
\expandafter\def\csname sga@target@IV.17.13.11\endcsname{ega4.c04.statement.17-13-11}
\expandafter\def\csname sga@target@IV.17.13.11.1\endcsname{ega4.c04.statement.17-13-11-1}
\expandafter\def\csname sga@target@IV.17.13.12\endcsname{ega4.c04.cor.17-13-12}
\expandafter\def\csname sga@target@IV.17.13.12.1\endcsname{ega4.c04.cor.17-13-12-1}
\expandafter\def\csname sga@target@IV.17.13.12.2\endcsname{ega4.c04.eq.17-13-12-2}
\expandafter\def\csname sga@target@IV.17.13.2\endcsname{ega4.c04.prop.17-13-2}
\expandafter\def\csname sga@target@IV.17.13.2.1\endcsname{ega4.c04.eq.17-13-2-1}
\expandafter\def\csname sga@target@IV.17.13.2.2\endcsname{ega4.c04.eq.17-13-2-2}
\expandafter\def\csname sga@target@IV.17.13.3\endcsname{ega4.c04.def.17-13-3}
\expandafter\def\csname sga@target@IV.17.13.4\endcsname{ega4.c04.statement.17-13-4}
\expandafter\def\csname sga@target@IV.17.13.5\endcsname{ega4.c04.statement.17-13-5}
\expandafter\def\csname sga@target@IV.17.13.5.1\endcsname{ega4.c04.statement.17-13-5-1}
\expandafter\def\csname sga@target@IV.17.13.5.2\endcsname{ega4.c04.eq.17-13-5-2}
\expandafter\def\csname sga@target@IV.17.13.5.3\endcsname{ega4.c04.eq.17-13-5-3}
\expandafter\def\csname sga@target@IV.17.13.5.4\endcsname{ega4.c04.eq.17-13-5-4}
\expandafter\def\csname sga@target@IV.17.13.6\endcsname{ega4.c04.cor.17-13-6}
\expandafter\def\csname sga@target@IV.17.13.6.1\endcsname{ega4.c04.eq.17-13-6-1}
\expandafter\def\csname sga@target@IV.17.13.6.2\endcsname{ega4.c04.eq.17-13-6-2}
\expandafter\def\csname sga@target@IV.17.13.7\endcsname{ega4.c04.statement.17-13-7}
\expandafter\def\csname sga@target@IV.17.13.7.1\endcsname{ega4.c04.eq.17-13-7-1}
\expandafter\def\csname sga@target@IV.17.13.8\endcsname{ega4.c04.prop.17-13-8}
\expandafter\def\csname sga@target@IV.17.13.8.1\endcsname{ega4.c04.eq.17-13-8-1}
\expandafter\def\csname sga@target@IV.17.13.9\endcsname{ega4.c04.statement.17-13-9}
\expandafter\def\csname sga@target@IV.17.13.9.1\endcsname{ega4.c04.eq.17-13-9-1}
\expandafter\def\csname sga@target@IV.17.13.9.2\endcsname{ega4.c04.eq.17-13-9-2}
\expandafter\def\csname sga@target@IV.17.14.1\endcsname{ega4.c04.prop.17-14-1}
\expandafter\def\csname sga@target@IV.17.14.2\endcsname{ega4.c04.prop.17-14-2}
\expandafter\def\csname sga@target@IV.17.14.3\endcsname{ega4.c04.rem.17-14-3}
\expandafter\def\csname sga@target@IV.17.15.1\endcsname{ega4.c04.prop.17-15-1}
\expandafter\def\csname sga@target@IV.17.15.10\endcsname{ega4.c04.statement.17-15-10}
\expandafter\def\csname sga@target@IV.17.15.11\endcsname{ega4.c04.prop.17-15-11}
\expandafter\def\csname sga@target@IV.17.15.12\endcsname{ega4.c04.prop.17-15-12}
\expandafter\def\csname sga@target@IV.17.15.13\endcsname{ega4.c04.cor.17-15-13}
\expandafter\def\csname sga@target@IV.17.15.14\endcsname{ega4.c04.prop.17-15-14}
\expandafter\def\csname sga@target@IV.17.15.14.1\endcsname{ega4.c04.lem.17-15-14-1}
\expandafter\def\csname sga@target@IV.17.15.14.2\endcsname{ega4.c04.lem.17-15-14-2}
\expandafter\def\csname sga@target@IV.17.15.15\endcsname{ega4.c04.prop.17-15-15}
\expandafter\def\csname sga@target@IV.17.15.2\endcsname{ega4.c04.cor.17-15-2}
\expandafter\def\csname sga@target@IV.17.15.3\endcsname{ega4.c04.prop.17-15-3}
\expandafter\def\csname sga@target@IV.17.15.4\endcsname{ega4.c04.cor.17-15-4}
\expandafter\def\csname sga@target@IV.17.15.5\endcsname{ega4.c04.prop.17-15-5}
\expandafter\def\csname sga@target@IV.17.15.6\endcsname{ega4.c04.cor.17-15-6}
\expandafter\def\csname sga@target@IV.17.15.7\endcsname{ega4.c04.cor.17-15-7}
\expandafter\def\csname sga@target@IV.17.15.8\endcsname{ega4.c04.prop.17-15-8}
\expandafter\def\csname sga@target@IV.17.15.8.1\endcsname{ega4.c04.eq.17-15-8-1}
\expandafter\def\csname sga@target@IV.17.15.8.2\endcsname{ega4.c04.eq.17-15-8-2}
\expandafter\def\csname sga@target@IV.17.15.9\endcsname{ega4.c04.cor.17-15-9}
\expandafter\def\csname sga@target@IV.17.16\endcsname{ega4.c04.subsec.17-16.key-iv-17-16}
\expandafter\def\csname sga@target@IV.17.16.1\endcsname{ega4.c04.prop.17-16-1}
\expandafter\def\csname sga@target@IV.17.16.2\endcsname{ega4.c04.cor.17-16-2}
\expandafter\def\csname sga@target@IV.17.16.3\endcsname{ega4.c04.cor.17-16-3}
\expandafter\def\csname sga@target@IV.17.16.4\endcsname{ega4.c04.prop.17-16-4}
\expandafter\def\csname sga@target@IV.17.16.5\endcsname{ega4.c04.cor.17-16-5}
\expandafter\def\csname sga@target@IV.17.16.6\endcsname{ega4.c04.prop.17-16-6}
\expandafter\def\csname sga@target@IV.17.2.1\endcsname{ega4.c04.prop.17-2-1}
\expandafter\def\csname sga@target@IV.17.2.2\endcsname{ega4.c04.cor.17-2-2}
\expandafter\def\csname sga@target@IV.17.2.3\endcsname{ega4.c04.prop.17-2-3}
\expandafter\def\csname sga@target@IV.17.2.3.1\endcsname{ega4.c04.eq.17-2-3-1}
\expandafter\def\csname sga@target@IV.17.2.4\endcsname{ega4.c04.cor.17-2-4}
\expandafter\def\csname sga@target@IV.17.2.5\endcsname{ega4.c04.prop.17-2-5}
\expandafter\def\csname sga@target@IV.17.2.5.1\endcsname{ega4.c04.prop.17-2-5-1}
\expandafter\def\csname sga@target@IV.17.2.6\endcsname{ega4.c04.prop.17-2-6}
\expandafter\def\csname sga@target@IV.17.3.1\endcsname{ega4.c04.def.17-3-1}
\expandafter\def\csname sga@target@IV.17.3.2\endcsname{ega4.c04.subsec.17-3-2}
\expandafter\def\csname sga@target@IV.17.3.3\endcsname{ega4.c04.prop.17-3-3}
\expandafter\def\csname sga@target@IV.17.3.4\endcsname{ega4.c04.prop.17-3-4}
\expandafter\def\csname sga@target@IV.17.3.5\endcsname{ega4.c04.cor.17-3-5}
\expandafter\def\csname sga@target@IV.17.3.6\endcsname{ega4.c04.prop.17-3-6}
\expandafter\def\csname sga@target@IV.17.3.7\endcsname{ega4.c04.def.17-3-7}
\expandafter\def\csname sga@target@IV.17.3.8\endcsname{ega4.c04.prop.17-3-8}
\expandafter\def\csname sga@target@IV.17.3.9\endcsname{ega4.c04.cor.17-3-9}
\expandafter\def\csname sga@target@IV.17.4.1\endcsname{ega4.c04.thm.17-4-1}
\expandafter\def\csname sga@target@IV.17.4.1.1\endcsname{ega4.c04.eq.17-4-1-1}
\expandafter\def\csname sga@target@IV.17.4.1.2\endcsname{ega4.c04.rem.17-4-1-2}
\expandafter\def\csname sga@target@IV.17.4.10\endcsname{ega4.c04.rem.17-4-10}
\expandafter\def\csname sga@target@IV.17.4.2\endcsname{ega4.c04.cor.17-4-2}
\expandafter\def\csname sga@target@IV.17.4.3\endcsname{ega4.c04.cor.17-4-3}
\expandafter\def\csname sga@target@IV.17.4.4\endcsname{ega4.c04.prop.17-4-4}
\expandafter\def\csname sga@target@IV.17.4.5\endcsname{ega4.c04.statement.17-4-5}
\expandafter\def\csname sga@target@IV.17.4.6\endcsname{ega4.c04.prop.17-4-6}
\expandafter\def\csname sga@target@IV.17.4.7\endcsname{ega4.c04.cor.17-4-7}
\expandafter\def\csname sga@target@IV.17.4.8\endcsname{ega4.c04.cor.17-4-8}
\expandafter\def\csname sga@target@IV.17.4.9\endcsname{ega4.c04.prop.17-4-9}
\expandafter\def\csname sga@target@IV.17.5.1\endcsname{ega4.c04.thm.17-5-1}
\expandafter\def\csname sga@target@IV.17.5.1.1\endcsname{ega4.c04.eq.17-5-1-1}
\expandafter\def\csname sga@target@IV.17.5.1.2\endcsname{ega4.c04.eq.17-5-1-2}
\expandafter\def\csname sga@target@IV.17.5.2\endcsname{ega4.c04.cor.17-5-2}
\expandafter\def\csname sga@target@IV.17.5.3\endcsname{ega4.c04.prop.17-5-3}
\expandafter\def\csname sga@target@IV.17.5.4\endcsname{ega4.c04.rem.17-5-4}
\expandafter\def\csname sga@target@IV.17.5.5\endcsname{ega4.c04.prop.17-5-5}
\expandafter\def\csname sga@target@IV.17.5.6\endcsname{ega4.c04.cor.17-5-6}
\expandafter\def\csname sga@target@IV.17.5.7\endcsname{ega4.c04.prop.17-5-7}
\expandafter\def\csname sga@target@IV.17.5.8\endcsname{ega4.c04.prop.17-5-8}
\expandafter\def\csname sga@target@IV.17.6.1\endcsname{ega4.c04.thm.17-6-1}
\expandafter\def\csname sga@target@IV.17.6.2\endcsname{ega4.c04.cor.17-6-2}
\expandafter\def\csname sga@target@IV.17.6.3\endcsname{ega4.c04.prop.17-6-3}
\expandafter\def\csname sga@target@IV.17.6.4\endcsname{ega4.c04.prop.17-6-4}
\expandafter\def\csname sga@target@IV.17.7\endcsname{ega4.c04.subsec.17-7.key-iv-17-7}
\expandafter\def\csname sga@target@IV.17.7.1\endcsname{ega4.c04.prop.17-7-1}
\expandafter\def\csname sga@target@IV.17.7.10\endcsname{ega4.c04.prop.17-7-10}
\expandafter\def\csname sga@target@IV.17.7.11\endcsname{ega4.c04.prop.17-7-11}
\expandafter\def\csname sga@target@IV.17.7.11.1\endcsname{ega4.c04.lem.17-7-11-1}
\expandafter\def\csname sga@target@IV.17.7.2\endcsname{ega4.c04.cor.17-7-2}
\expandafter\def\csname sga@target@IV.17.7.3\endcsname{ega4.c04.cor.17-7-3}
\expandafter\def\csname sga@target@IV.17.7.4\endcsname{ega4.c04.prop.17-7-4}
\expandafter\def\csname sga@target@IV.17.7.5\endcsname{ega4.c04.lem.17-7-5}
\expandafter\def\csname sga@target@IV.17.7.6\endcsname{ega4.c04.rem.17-7-6}
\expandafter\def\csname sga@target@IV.17.7.7\endcsname{ega4.c04.prop.17-7-7}
\expandafter\def\csname sga@target@IV.17.7.8\endcsname{ega4.c04.prop.17-7-8}
\expandafter\def\csname sga@target@IV.17.7.9\endcsname{ega4.c04.cor.17-7-9}
\expandafter\def\csname sga@target@IV.17.8\endcsname{ega4.c04.subsec.17-8.key-iv-17-8}
\expandafter\def\csname sga@target@IV.17.8.1\endcsname{ega4.c04.prop.17-8-1}
\expandafter\def\csname sga@target@IV.17.8.2\endcsname{ega4.c04.prop.17-8-2}
\expandafter\def\csname sga@target@IV.17.8.3\endcsname{ega4.c04.rem.17-8-3}
\expandafter\def\csname sga@target@IV.17.9\endcsname{ega4.c04.subsec.17-9.key-iv-17-9}
\expandafter\def\csname sga@target@IV.17.9.1\endcsname{ega4.c04.thm.17-9-1}
\expandafter\def\csname sga@target@IV.17.9.2\endcsname{ega4.c04.cor.17-9-2}
\expandafter\def\csname sga@target@IV.17.9.3\endcsname{ega4.c04.cor.17-9-3}
\expandafter\def\csname sga@target@IV.17.9.4\endcsname{ega4.c04.cor.17-9-4}
\expandafter\def\csname sga@target@IV.17.9.5\endcsname{ega4.c04.cor.17-9-5}
\expandafter\def\csname sga@target@IV.17.9.6\endcsname{ega4.c04.prop.17-9-6}
\expandafter\def\csname sga@target@IV.18\endcsname{ega4.c04.sec.18.key-iv-18}
\expandafter\def\csname sga@target@IV.18.1.1\endcsname{ega4.c04.prop.18-1-1}
\expandafter\def\csname sga@target@IV.18.1.1.1\endcsname{ega4.c04.eq.18-1-1-1}
\expandafter\def\csname sga@target@IV.18.1.2\endcsname{ega4.c04.thm.18-1-2}
\expandafter\def\csname sga@target@IV.18.1.3\endcsname{ega4.c04.cor.18-1-3}
\expandafter\def\csname sga@target@IV.18.10.1\endcsname{ega4.c04.thm.18-10-1}
\expandafter\def\csname sga@target@IV.18.10.10\endcsname{ega4.c04.statement.18-10-10}
\expandafter\def\csname sga@target@IV.18.10.11\endcsname{ega4.c04.rem.18-10-11}
\expandafter\def\csname sga@target@IV.18.10.12\endcsname{ega4.c04.cor.18-10-12}
\expandafter\def\csname sga@target@IV.18.10.13\endcsname{ega4.c04.prop.18-10-13}
\expandafter\def\csname sga@target@IV.18.10.14\endcsname{ega4.c04.cor.18-10-14}
\expandafter\def\csname sga@target@IV.18.10.15\endcsname{ega4.c04.statement.18-10-15}
\expandafter\def\csname sga@target@IV.18.10.16\endcsname{ega4.c04.thm.18-10-16}
\expandafter\def\csname sga@target@IV.18.10.17\endcsname{ega4.c04.rem.18-10-17}
\expandafter\def\csname sga@target@IV.18.10.17.1\endcsname{ega4.c04.statement.18-10-17-1}
\expandafter\def\csname sga@target@IV.18.10.17.2\endcsname{ega4.c04.statement.18-10-17-2}
\expandafter\def\csname sga@target@IV.18.10.18\endcsname{ega4.c04.lem.18-10-18}
\expandafter\def\csname sga@target@IV.18.10.19\endcsname{ega4.c04.prop.18-10-19}
\expandafter\def\csname sga@target@IV.18.10.2\endcsname{ega4.c04.rem.18-10-2}
\expandafter\def\csname sga@target@IV.18.10.20\endcsname{ega4.c04.rem.18-10-20}
\expandafter\def\csname sga@target@IV.18.10.21\endcsname{ega4.c04.prop.18-10-21}
\expandafter\def\csname sga@target@IV.18.10.3\endcsname{ega4.c04.cor.18-10-3}
\expandafter\def\csname sga@target@IV.18.10.3.1\endcsname{ega4.c04.rem.18-10-3-1}
\expandafter\def\csname sga@target@IV.18.10.4\endcsname{ega4.c04.cor.18-10-4}
\expandafter\def\csname sga@target@IV.18.10.5\endcsname{ega4.c04.cor.18-10-5}
\expandafter\def\csname sga@target@IV.18.10.6\endcsname{ega4.c04.lem.18-10-6}
\expandafter\def\csname sga@target@IV.18.10.7\endcsname{ega4.c04.prop.18-10-7}
\expandafter\def\csname sga@target@IV.18.10.8\endcsname{ega4.c04.prop.18-10-8}
\expandafter\def\csname sga@target@IV.18.10.9\endcsname{ega4.c04.cor.18-10-9}
\expandafter\def\csname sga@target@IV.18.11.1\endcsname{ega4.c04.lem.18-11-1}
\expandafter\def\csname sga@target@IV.18.11.10\endcsname{ega4.c04.thm.18-11-10}
\expandafter\def\csname sga@target@IV.18.11.11\endcsname{ega4.c04.rem.18-11-11}
\expandafter\def\csname sga@target@IV.18.11.12\endcsname{ega4.c04.cor.18-11-12}
\expandafter\def\csname sga@target@IV.18.11.12.1\endcsname{ega4.c04.lem.18-11-12-1}
\expandafter\def\csname sga@target@IV.18.11.2\endcsname{ega4.c04.lem.18-11-2}
\expandafter\def\csname sga@target@IV.18.11.2.1\endcsname{ega4.c04.eq.18-11-2-1}
\expandafter\def\csname sga@target@IV.18.11.2.2\endcsname{ega4.c04.eq.18-11-2-2}
\expandafter\def\csname sga@target@IV.18.11.3\endcsname{ega4.c04.lem.18-11-3}
\expandafter\def\csname sga@target@IV.18.11.3.1\endcsname{ega4.c04.eq.18-11-3-1}
\expandafter\def\csname sga@target@IV.18.11.3.2\endcsname{ega4.c04.eq.18-11-3-2}
\expandafter\def\csname sga@target@IV.18.11.3.3\endcsname{ega4.c04.lem.18-11-3-3}
\expandafter\def\csname sga@target@IV.18.11.3.4\endcsname{ega4.c04.lem.18-11-3-4}
\expandafter\def\csname sga@target@IV.18.11.3.5\endcsname{ega4.c04.lem.18-11-3-5}
\expandafter\def\csname sga@target@IV.18.11.3.6\endcsname{ega4.c04.eq.18-11-3-6}
\expandafter\def\csname sga@target@IV.18.11.4\endcsname{ega4.c04.lem.18-11-4}
\expandafter\def\csname sga@target@IV.18.11.5\endcsname{ega4.c04.prop.18-11-5}
\expandafter\def\csname sga@target@IV.18.11.5.1\endcsname{ega4.c04.lem.18-11-5-1}
\expandafter\def\csname sga@target@IV.18.11.5.2\endcsname{ega4.c04.lem.18-11-5-2}
\expandafter\def\csname sga@target@IV.18.11.6\endcsname{ega4.c04.rem.18-11-6}
\expandafter\def\csname sga@target@IV.18.11.7\endcsname{ega4.c04.cor.18-11-7}
\expandafter\def\csname sga@target@IV.18.11.8\endcsname{ega4.c04.cor.18-11-8}
\expandafter\def\csname sga@target@IV.18.11.9\endcsname{ega4.c04.cor.18-11-9}
\expandafter\def\csname sga@target@IV.18.12.1\endcsname{ega4.c04.thm.18-12-1}
\expandafter\def\csname sga@target@IV.18.12.10\endcsname{ega4.c04.cor.18-12-10}
\expandafter\def\csname sga@target@IV.18.12.11\endcsname{ega4.c04.cor.18-12-11}
\expandafter\def\csname sga@target@IV.18.12.12\endcsname{ega4.c04.prop.18-12-12}
\expandafter\def\csname sga@target@IV.18.12.13\endcsname{ega4.c04.cor.18-12-13}
\expandafter\def\csname sga@target@IV.18.12.14\endcsname{ega4.c04.rem.18-12-14}
\expandafter\def\csname sga@target@IV.18.12.15\endcsname{ega4.c04.prop.18-12-15}
\expandafter\def\csname sga@target@IV.18.12.16\endcsname{ega4.c04.prop.18-12-16}
\expandafter\def\csname sga@target@IV.18.12.17\endcsname{ega4.c04.cor.18-12-17}
\expandafter\def\csname sga@target@IV.18.12.2\endcsname{ega4.c04.rem.18-12-2}
\expandafter\def\csname sga@target@IV.18.12.3\endcsname{ega4.c04.cor.18-12-3}
\expandafter\def\csname sga@target@IV.18.12.4\endcsname{ega4.c04.cor.18-12-4}
\expandafter\def\csname sga@target@IV.18.12.5\endcsname{ega4.c04.rem.18-12-5}
\expandafter\def\csname sga@target@IV.18.12.6\endcsname{ega4.c04.cor.18-12-6}
\expandafter\def\csname sga@target@IV.18.12.7\endcsname{ega4.c04.prop.18-12-7}
\expandafter\def\csname sga@target@IV.18.12.7.1\endcsname{ega4.c04.lem.18-12-7-1}
\expandafter\def\csname sga@target@IV.18.12.8\endcsname{ega4.c04.prop.18-12-8}
\expandafter\def\csname sga@target@IV.18.12.9\endcsname{ega4.c04.rem.18-12-9}
\expandafter\def\csname sga@target@IV.18.2.1\endcsname{ega4.c04.statement.18-2-1}
\expandafter\def\csname sga@target@IV.18.2.1.1\endcsname{ega4.c04.statement.18-2-1-1}
\expandafter\def\csname sga@target@IV.18.2.1.2\endcsname{ega4.c04.eq.18-2-1-2}
\expandafter\def\csname sga@target@IV.18.2.10\endcsname{ega4.c04.statement.18-2-10}
\expandafter\def\csname sga@target@IV.18.2.2\endcsname{ega4.c04.statement.18-2-2}
\expandafter\def\csname sga@target@IV.18.2.2.1\endcsname{ega4.c04.eq.18-2-2-1}
\expandafter\def\csname sga@target@IV.18.2.2.2\endcsname{ega4.c04.eq.18-2-2-2}
\expandafter\def\csname sga@target@IV.18.2.2.3\endcsname{ega4.c04.eq.18-2-2-3}
\expandafter\def\csname sga@target@IV.18.2.2.4\endcsname{ega4.c04.eq.18-2-2-4}
\expandafter\def\csname sga@target@IV.18.2.3\endcsname{ega4.c04.prop.18-2-3}
\expandafter\def\csname sga@target@IV.18.2.4\endcsname{ega4.c04.cor.18-2-4}
\expandafter\def\csname sga@target@IV.18.2.5\endcsname{ega4.c04.cor.18-2-5}
\expandafter\def\csname sga@target@IV.18.2.6\endcsname{ega4.c04.cor.18-2-6}
\expandafter\def\csname sga@target@IV.18.2.7\endcsname{ega4.c04.item.18-2-7}
\expandafter\def\csname sga@target@IV.18.2.7.1\endcsname{ega4.c04.eq.18-2-7-1}
\expandafter\def\csname sga@target@IV.18.2.7.2\endcsname{ega4.c04.eq.18-2-7-2}
\expandafter\def\csname sga@target@IV.18.2.8\endcsname{ega4.c04.prop.18-2-8}
\expandafter\def\csname sga@target@IV.18.2.9\endcsname{ega4.c04.cor.18-2-9}
\expandafter\def\csname sga@target@IV.18.3.1\endcsname{ega4.c04.prop.18-3-1}
\expandafter\def\csname sga@target@IV.18.3.2\endcsname{ega4.c04.prop.18-3-2}
\expandafter\def\csname sga@target@IV.18.3.2.1\endcsname{ega4.c04.lem.18-3-2-1}
\expandafter\def\csname sga@target@IV.18.3.2.2\endcsname{ega4.c04.eq.18-3-2-2}
\expandafter\def\csname sga@target@IV.18.3.3\endcsname{ega4.c04.statement.18-3-3}
\expandafter\def\csname sga@target@IV.18.3.4\endcsname{ega4.c04.thm.18-3-4}
\expandafter\def\csname sga@target@IV.18.4.1\endcsname{ega4.c04.lem.18-4-1}
\expandafter\def\csname sga@target@IV.18.4.10\endcsname{ega4.c04.prop.18-4-10}
\expandafter\def\csname sga@target@IV.18.4.11\endcsname{ega4.c04.statement.18-4-11}
\expandafter\def\csname sga@target@IV.18.4.12\endcsname{ega4.c04.cor.18-4-12}
\expandafter\def\csname sga@target@IV.18.4.13\endcsname{ega4.c04.cor.18-4-13}
\expandafter\def\csname sga@target@IV.18.4.14\endcsname{ega4.c04.cor.18-4-14}
\expandafter\def\csname sga@target@IV.18.4.2\endcsname{ega4.c04.prop.18-4-2}
\expandafter\def\csname sga@target@IV.18.4.3\endcsname{ega4.c04.cor.18-4-3}
\expandafter\def\csname sga@target@IV.18.4.4\endcsname{ega4.c04.lem.18-4-4}
\expandafter\def\csname sga@target@IV.18.4.5\endcsname{ega4.c04.prop.18-4-5}
\expandafter\def\csname sga@target@IV.18.4.6\endcsname{ega4.c04.thm.18-4-6}
\expandafter\def\csname sga@target@IV.18.4.6.1\endcsname{ega4.c04.lem.18-4-6-1}
\expandafter\def\csname sga@target@IV.18.4.7\endcsname{ega4.c04.cor.18-4-7}
\expandafter\def\csname sga@target@IV.18.4.8\endcsname{ega4.c04.cor.18-4-8}
\expandafter\def\csname sga@target@IV.18.4.9\endcsname{ega4.c04.prop.18-4-9}
\expandafter\def\csname sga@target@IV.18.5.1\endcsname{ega4.c04.statement.18-5-1}
\expandafter\def\csname sga@target@IV.18.5.1.1\endcsname{ega4.c04.eq.18-5-1-1}
\expandafter\def\csname sga@target@IV.18.5.10\endcsname{ega4.c04.prop.18-5-10}
\expandafter\def\csname sga@target@IV.18.5.11\endcsname{ega4.c04.thm.18-5-11}
\expandafter\def\csname sga@target@IV.18.5.12\endcsname{ega4.c04.cor.18-5-12}
\expandafter\def\csname sga@target@IV.18.5.13\endcsname{ega4.c04.rem.18-5-13}
\expandafter\def\csname sga@target@IV.18.5.13.1\endcsname{ega4.c04.lem.18-5-13-1}
\expandafter\def\csname sga@target@IV.18.5.14\endcsname{ega4.c04.prop.18-5-14}
\expandafter\def\csname sga@target@IV.18.5.15\endcsname{ega4.c04.prop.18-5-15}
\expandafter\def\csname sga@target@IV.18.5.16\endcsname{ega4.c04.item.18-5-16}
\expandafter\def\csname sga@target@IV.18.5.17\endcsname{ega4.c04.thm.18-5-17}
\expandafter\def\csname sga@target@IV.18.5.18\endcsname{ega4.c04.rem.18-5-18}
\expandafter\def\csname sga@target@IV.18.5.19\endcsname{ega4.c04.prop.18-5-19}
\expandafter\def\csname sga@target@IV.18.5.2\endcsname{ega4.c04.statement.18-5-2}
\expandafter\def\csname sga@target@IV.18.5.3\endcsname{ega4.c04.lem.18-5-3}
\expandafter\def\csname sga@target@IV.18.5.3.1\endcsname{ega4.c04.lem.18-5-3-1}
\expandafter\def\csname sga@target@IV.18.5.4\endcsname{ega4.c04.prop.18-5-4}
\expandafter\def\csname sga@target@IV.18.5.4.1\endcsname{ega4.c04.eq.18-5-4-1}
\expandafter\def\csname sga@target@IV.18.5.4.2\endcsname{ega4.c04.eq.18-5-4-2}
\expandafter\def\csname sga@target@IV.18.5.4.3\endcsname{ega4.c04.lem.18-5-4-3}
\expandafter\def\csname sga@target@IV.18.5.4.4\endcsname{ega4.c04.rem.18-5-4-4}
\expandafter\def\csname sga@target@IV.18.5.5\endcsname{ega4.c04.def.18-5-5}
\expandafter\def\csname sga@target@IV.18.5.6\endcsname{ega4.c04.prop.18-5-6}
\expandafter\def\csname sga@target@IV.18.5.7\endcsname{ega4.c04.rem.18-5-7}
\expandafter\def\csname sga@target@IV.18.5.8\endcsname{ega4.c04.def.18-5-8}
\expandafter\def\csname sga@target@IV.18.5.9\endcsname{ega4.c04.prop.18-5-9}
\expandafter\def\csname sga@target@IV.18.6.1\endcsname{ega4.c04.statement.18-6-1}
\expandafter\def\csname sga@target@IV.18.6.10\endcsname{ega4.c04.cor.18-6-10}
\expandafter\def\csname sga@target@IV.18.6.11\endcsname{ega4.c04.cor.18-6-11}
\expandafter\def\csname sga@target@IV.18.6.12\endcsname{ega4.c04.prop.18-6-12}
\expandafter\def\csname sga@target@IV.18.6.13\endcsname{ega4.c04.cor.18-6-13}
\expandafter\def\csname sga@target@IV.18.6.14\endcsname{ega4.c04.prop.18-6-14}
\expandafter\def\csname sga@target@IV.18.6.15\endcsname{ega4.c04.cor.18-6-15}
\expandafter\def\csname sga@target@IV.18.6.2\endcsname{ega4.c04.lem.18-6-2}
\expandafter\def\csname sga@target@IV.18.6.3\endcsname{ega4.c04.lem.18-6-3}
\expandafter\def\csname sga@target@IV.18.6.4\endcsname{ega4.c04.statement.18-6-4}
\expandafter\def\csname sga@target@IV.18.6.5\endcsname{ega4.c04.def.18-6-5}
\expandafter\def\csname sga@target@IV.18.6.6\endcsname{ega4.c04.thm.18-6-6}
\expandafter\def\csname sga@target@IV.18.6.7\endcsname{ega4.c04.statement.18-6-7}
\expandafter\def\csname sga@target@IV.18.6.8\endcsname{ega4.c04.prop.18-6-8}
\expandafter\def\csname sga@target@IV.18.6.9\endcsname{ega4.c04.thm.18-6-9}
\expandafter\def\csname sga@target@IV.18.6.9.1\endcsname{ega4.c04.lem.18-6-9-1}
\expandafter\def\csname sga@target@IV.18.6.9.2\endcsname{ega4.c04.lem.18-6-9-2}
\expandafter\def\csname sga@target@IV.18.7.1\endcsname{ega4.c04.statement.18-7-1}
\expandafter\def\csname sga@target@IV.18.7.2\endcsname{ega4.c04.prop.18-7-2}
\expandafter\def\csname sga@target@IV.18.7.3\endcsname{ega4.c04.cor.18-7-3}
\expandafter\def\csname sga@target@IV.18.7.4\endcsname{ega4.c04.cor.18-7-4}
\expandafter\def\csname sga@target@IV.18.7.5\endcsname{ega4.c04.prop.18-7-5}
\expandafter\def\csname sga@target@IV.18.7.5.1\endcsname{ega4.c04.lem.18-7-5-1}
\expandafter\def\csname sga@target@IV.18.7.6\endcsname{ega4.c04.cor.18-7-6}
\expandafter\def\csname sga@target@IV.18.7.7\endcsname{ega4.c04.rem.18-7-7}
\expandafter\def\csname sga@target@IV.18.8.1\endcsname{ega4.c04.prop.18-8-1}
\expandafter\def\csname sga@target@IV.18.8.10\endcsname{ega4.c04.prop.18-8-10}
\expandafter\def\csname sga@target@IV.18.8.11\endcsname{ega4.c04.rem.18-8-11}
\expandafter\def\csname sga@target@IV.18.8.12\endcsname{ega4.c04.prop.18-8-12}
\expandafter\def\csname sga@target@IV.18.8.13\endcsname{ega4.c04.cor.18-8-13}
\expandafter\def\csname sga@target@IV.18.8.14\endcsname{ega4.c04.cor.18-8-14}
\expandafter\def\csname sga@target@IV.18.8.15\endcsname{ega4.c04.prop.18-8-15}
\expandafter\def\csname sga@target@IV.18.8.16\endcsname{ega4.c04.cor.18-8-16}
\expandafter\def\csname sga@target@IV.18.8.17\endcsname{ega4.c04.prop.18-8-17}
\expandafter\def\csname sga@target@IV.18.8.18\endcsname{ega4.c04.prop.18-8-18}
\expandafter\def\csname sga@target@IV.18.8.2\endcsname{ega4.c04.def.18-8-2}
\expandafter\def\csname sga@target@IV.18.8.3\endcsname{ega4.c04.rem.18-8-3}
\expandafter\def\csname sga@target@IV.18.8.4\endcsname{ega4.c04.lem.18-8-4}
\expandafter\def\csname sga@target@IV.18.8.5\endcsname{ega4.c04.lem.18-8-5}
\expandafter\def\csname sga@target@IV.18.8.6\endcsname{ega4.c04.statement.18-8-6}
\expandafter\def\csname sga@target@IV.18.8.6.1\endcsname{ega4.c04.eq.18-8-6-1}
\expandafter\def\csname sga@target@IV.18.8.6.2\endcsname{ega4.c04.eq.18-8-6-2}
\expandafter\def\csname sga@target@IV.18.8.7\endcsname{ega4.c04.def.18-8-7}
\expandafter\def\csname sga@target@IV.18.8.8\endcsname{ega4.c04.prop.18-8-8}
\expandafter\def\csname sga@target@IV.18.8.8.1\endcsname{ega4.c04.statement.18-8-8-1}
\expandafter\def\csname sga@target@IV.18.8.9\endcsname{ega4.c04.statement.18-8-9}
\expandafter\def\csname sga@target@IV.18.9.1\endcsname{ega4.c04.thm.18-9-1}
\expandafter\def\csname sga@target@IV.18.9.10\endcsname{ega4.c04.rem.18-9-10}
\expandafter\def\csname sga@target@IV.18.9.11\endcsname{ega4.c04.cor.18-9-11}
\expandafter\def\csname sga@target@IV.18.9.2\endcsname{ega4.c04.cor.18-9-2}
\expandafter\def\csname sga@target@IV.18.9.3\endcsname{ega4.c04.cor.18-9-3}
\expandafter\def\csname sga@target@IV.18.9.4\endcsname{ega4.c04.cor.18-9-4}
\expandafter\def\csname sga@target@IV.18.9.5\endcsname{ega4.c04.cor.18-9-5}
\expandafter\def\csname sga@target@IV.18.9.5.1\endcsname{ega4.c04.eq.18-9-5-1}
\expandafter\def\csname sga@target@IV.18.9.6\endcsname{ega4.c04.statement.18-9-6}
\expandafter\def\csname sga@target@IV.18.9.7\endcsname{ega4.c04.thm.18-9-7}
\expandafter\def\csname sga@target@IV.18.9.7.1\endcsname{ega4.c04.rem.18-9-7-1}
\expandafter\def\csname sga@target@IV.18.9.7.2\endcsname{ega4.c04.statement.18-9-7-2}
\expandafter\def\csname sga@target@IV.18.9.7.3\endcsname{ega4.c04.lem.18-9-7-3}
\expandafter\def\csname sga@target@IV.18.9.7.4\endcsname{ega4.c04.eq.18-9-7-4}
\expandafter\def\csname sga@target@IV.18.9.7.5\endcsname{ega4.c04.lem.18-9-7-5}
\expandafter\def\csname sga@target@IV.18.9.7.6\endcsname{ega4.c04.rem.18-9-7-6}
\expandafter\def\csname sga@target@IV.18.9.7.7\endcsname{ega4.c04.lem.18-9-7-7}
\expandafter\def\csname sga@target@IV.18.9.7.8\endcsname{ega4.c04.rem.18-9-7-8}
\expandafter\def\csname sga@target@IV.18.9.8\endcsname{ega4.c04.cor.18-9-8}
\expandafter\def\csname sga@target@IV.18.9.9\endcsname{ega4.c04.cor.18-9-9}
\expandafter\def\csname sga@target@IV.19.1.1\endcsname{ega4.c04.prop.19-1-1}
\expandafter\def\csname sga@target@IV.19.1.2\endcsname{ega4.c04.cor.19-1-2}
\expandafter\def\csname sga@target@IV.19.1.3\endcsname{ega4.c04.def.19-1-3}
\expandafter\def\csname sga@target@IV.19.1.4\endcsname{ega4.c04.prop.19-1-4}
\expandafter\def\csname sga@target@IV.19.1.4.1\endcsname{ega4.c04.eq.19-1-4-1}
\expandafter\def\csname sga@target@IV.19.1.5\endcsname{ega4.c04.prop.19-1-5}
\expandafter\def\csname sga@target@IV.19.1.5.1\endcsname{ega4.c04.item.19-1-5-1}
\expandafter\def\csname sga@target@IV.19.1.6\endcsname{ega4.c04.cor.19-1-6}
\expandafter\def\csname sga@target@IV.19.1.6.1\endcsname{ega4.c04.lem.19-1-6-1}
\expandafter\def\csname sga@target@IV.19.1.7\endcsname{ega4.c04.item.19-1-7}
\expandafter\def\csname sga@target@IV.19.1.8\endcsname{ega4.c04.statement.19-1-8}
\expandafter\def\csname sga@target@IV.19.2.1\endcsname{ega4.c04.def.19-2-1}
\expandafter\def\csname sga@target@IV.19.2.2\endcsname{ega4.c04.def.19-2-2}
\expandafter\def\csname sga@target@IV.19.2.3\endcsname{ega4.c04.prop.19-2-3}
\expandafter\def\csname sga@target@IV.19.2.4\endcsname{ega4.c04.prop.19-2-4}
\expandafter\def\csname sga@target@IV.19.2.5\endcsname{ega4.c04.item.19-2-5}
\expandafter\def\csname sga@target@IV.19.2.6\endcsname{ega4.c04.cor.19-2-6}
\expandafter\def\csname sga@target@IV.19.2.7\endcsname{ega4.c04.prop.19-2-7}
\expandafter\def\csname sga@target@IV.19.2.8\endcsname{ega4.c04.prop.19-2-8}
\expandafter\def\csname sga@target@IV.19.2.9\endcsname{ega4.c04.prop.19-2-9}
\expandafter\def\csname sga@target@IV.19.3.1\endcsname{ega4.c04.def.19-3-1}
\expandafter\def\csname sga@target@IV.19.3.10\endcsname{ega4.c04.cor.19-3-10}
\expandafter\def\csname sga@target@IV.19.3.2\endcsname{ega4.c04.prop.19-3-2}
\expandafter\def\csname sga@target@IV.19.3.2.1\endcsname{ega4.c04.lem.19-3-2-1}
\expandafter\def\csname sga@target@IV.19.3.3\endcsname{ega4.c04.cor.19-3-3}
\expandafter\def\csname sga@target@IV.19.3.4\endcsname{ega4.c04.cor.19-3-4}
\expandafter\def\csname sga@target@IV.19.3.5\endcsname{ega4.c04.cor.19-3-5}
\expandafter\def\csname sga@target@IV.19.3.6\endcsname{ega4.c04.def.19-3-6}
\expandafter\def\csname sga@target@IV.19.3.7\endcsname{ega4.c04.prop.19-3-7}
\expandafter\def\csname sga@target@IV.19.3.8\endcsname{ega4.c04.cor.19-3-8}
\expandafter\def\csname sga@target@IV.19.3.9\endcsname{ega4.c04.prop.19-3-9}
\expandafter\def\csname sga@target@IV.19.4.1\endcsname{ega4.c04.statement.19-4-1}
\expandafter\def\csname sga@target@IV.19.4.1.1\endcsname{ega4.c04.subsec.19-4-1-1}
\expandafter\def\csname sga@target@IV.19.4.10\endcsname{ega4.c04.rem.19-4-10}
\expandafter\def\csname sga@target@IV.19.4.11\endcsname{ega4.c04.prop.19-4-11}
\expandafter\def\csname sga@target@IV.19.4.11.1\endcsname{ega4.c04.eq.19-4-11-1}
\expandafter\def\csname sga@target@IV.19.4.2\endcsname{ega4.c04.lem.19-4-2}
\expandafter\def\csname sga@target@IV.19.4.3\endcsname{ega4.c04.cor.19-4-3}
\expandafter\def\csname sga@target@IV.19.4.4\endcsname{ega4.c04.prop.19-4-4}
\expandafter\def\csname sga@target@IV.19.4.5\endcsname{ega4.c04.cor.19-4-5}
\expandafter\def\csname sga@target@IV.19.4.6\endcsname{ega4.c04.prop.19-4-6}
\expandafter\def\csname sga@target@IV.19.4.7\endcsname{ega4.c04.cor.19-4-7}
\expandafter\def\csname sga@target@IV.19.4.8\endcsname{ega4.c04.prop.19-4-8}
\expandafter\def\csname sga@target@IV.19.4.9\endcsname{ega4.c04.cor.19-4-9}
\expandafter\def\csname sga@target@IV.19.5.1\endcsname{ega4.c04.thm.19-5-1}
\expandafter\def\csname sga@target@IV.19.5.1.1\endcsname{ega4.c04.eq.19-5-1-1}
\expandafter\def\csname sga@target@IV.19.5.2\endcsname{ega4.c04.cor.19-5-2}
\expandafter\def\csname sga@target@IV.19.5.3\endcsname{ega4.c04.prop.19-5-3}
\expandafter\def\csname sga@target@IV.19.5.4\endcsname{ega4.c04.statement.19-5-4}
\expandafter\def\csname sga@target@IV.19.5.5\endcsname{ega4.c04.thm.19-5-5}
\expandafter\def\csname sga@target@IV.19.5.5.1\endcsname{ega4.c04.statement.19-5-5-1}
\expandafter\def\csname sga@target@IV.19.5.5.2\endcsname{ega4.c04.statement.19-5-5-2}
\expandafter\def\csname sga@target@IV.19.5.5.3\endcsname{ega4.c04.statement.19-5-5-3}
\expandafter\def\csname sga@target@IV.19.5.5.4\endcsname{ega4.c04.statement.19-5-5-4}
\expandafter\def\csname sga@target@IV.19.5.5.5\endcsname{ega4.c04.statement.19-5-5-5}
\expandafter\def\csname sga@target@IV.19.5.6\endcsname{ega4.c04.statement.19-5-6}
\expandafter\def\csname sga@target@IV.19.6.1\endcsname{ega4.c04.statement.19-6-1}
\expandafter\def\csname sga@target@IV.19.6.1.1\endcsname{ega4.c04.eq.19-6-1-1}
\expandafter\def\csname sga@target@IV.19.6.2\endcsname{ega4.c04.statement.19-6-2}
\expandafter\def\csname sga@target@IV.19.6.2.1\endcsname{ega4.c04.statement.19-6-2-1}
\expandafter\def\csname sga@target@IV.19.6.3\endcsname{ega4.c04.prop.19-6-3}
\expandafter\def\csname sga@target@IV.19.6.4\endcsname{ega4.c04.cor.19-6-4}
\expandafter\def\csname sga@target@IV.19.7\endcsname{ega4.c04.subsec.19-7.key-iv-19-7}
\expandafter\def\csname sga@target@IV.19.7.1\endcsname{ega4.c04.thm.19-7-1}
\expandafter\def\csname sga@target@IV.19.7.1.1\endcsname{ega4.c04.statement.19-7-1-1}
\expandafter\def\csname sga@target@IV.19.7.1.10\endcsname{ega4.c04.statement.19-7-1-10}
\expandafter\def\csname sga@target@IV.19.7.1.11\endcsname{ega4.c04.statement.19-7-1-11}
\expandafter\def\csname sga@target@IV.19.7.1.12\endcsname{ega4.c04.statement.19-7-1-12}
\expandafter\def\csname sga@target@IV.19.7.1.13\endcsname{ega4.c04.statement.19-7-1-13}
\expandafter\def\csname sga@target@IV.19.7.1.14\endcsname{ega4.c04.statement.19-7-1-14}
\expandafter\def\csname sga@target@IV.19.7.1.15\endcsname{ega4.c04.statement.19-7-1-15}
\expandafter\def\csname sga@target@IV.19.7.1.16\endcsname{ega4.c04.statement.19-7-1-16}
\expandafter\def\csname sga@target@IV.19.7.1.17\endcsname{ega4.c04.statement.19-7-1-17}
\expandafter\def\csname sga@target@IV.19.7.1.18\endcsname{ega4.c04.statement.19-7-1-18}
\expandafter\def\csname sga@target@IV.19.7.1.19\endcsname{ega4.c04.statement.19-7-1-19}
\expandafter\def\csname sga@target@IV.19.7.1.2\endcsname{ega4.c04.statement.19-7-1-2}
\expandafter\def\csname sga@target@IV.19.7.1.20\endcsname{ega4.c04.statement.19-7-1-20}
\expandafter\def\csname sga@target@IV.19.7.1.21\endcsname{ega4.c04.statement.19-7-1-21}
\expandafter\def\csname sga@target@IV.19.7.1.22\endcsname{ega4.c04.statement.19-7-1-22}
\expandafter\def\csname sga@target@IV.19.7.1.23\endcsname{ega4.c04.statement.19-7-1-23}
\expandafter\def\csname sga@target@IV.19.7.1.3\endcsname{ega4.c04.statement.19-7-1-3}
\expandafter\def\csname sga@target@IV.19.7.1.5\endcsname{ega4.c04.statement.19-7-1-5}
\expandafter\def\csname sga@target@IV.19.7.1.6\endcsname{ega4.c04.statement.19-7-1-6}
\expandafter\def\csname sga@target@IV.19.7.1.7\endcsname{ega4.c04.lem.19-7-1-7}
\expandafter\def\csname sga@target@IV.19.7.1.8\endcsname{ega4.c04.statement.19-7-1-8}
\expandafter\def\csname sga@target@IV.19.7.1.9\endcsname{ega4.c04.statement.19-7-1-9}
\expandafter\def\csname sga@target@IV.19.7.2\endcsname{ega4.c04.statement.19-7-2}
\expandafter\def\csname sga@target@IV.19.7.3\endcsname{ega4.c04.cor.19-7-3}
\expandafter\def\csname sga@target@IV.19.7.3.1\endcsname{ega4.c04.eq.19-7-3-1}
\expandafter\def\csname sga@target@IV.19.7.4\endcsname{ega4.c04.cor.19-7-4}
\expandafter\def\csname sga@target@IV.19.7.4.1\endcsname{ega4.c04.eq.19-7-4-1}
\expandafter\def\csname sga@target@IV.19.7.5\endcsname{ega4.c04.cor.19-7-5}
\expandafter\def\csname sga@target@IV.19.8.1\endcsname{ega4.c04.prop.19-8-1}
\expandafter\def\csname sga@target@IV.19.8.2\endcsname{ega4.c04.prop.19-8-2}
\expandafter\def\csname sga@target@IV.19.8.3\endcsname{ega4.c04.rem.19-8-3}
\expandafter\def\csname sga@target@IV.19.9.1\endcsname{ega4.c04.statement.19-9-1}
\expandafter\def\csname sga@target@IV.19.9.1.1\endcsname{ega4.c04.eq.19-9-1-1}
\expandafter\def\csname sga@target@IV.19.9.1.2\endcsname{ega4.c04.eq.19-9-1-2}
\expandafter\def\csname sga@target@IV.19.9.1.3\endcsname{ega4.c04.eq.19-9-1-3}
\expandafter\def\csname sga@target@IV.19.9.2\endcsname{ega4.c04.lem.19-9-2}
\expandafter\def\csname sga@target@IV.19.9.3\endcsname{ega4.c04.prop.19-9-3}
\expandafter\def\csname sga@target@IV.19.9.4\endcsname{ega4.c04.cor.19-9-4}
\expandafter\def\csname sga@target@IV.19.9.5\endcsname{ega4.c04.prop.19-9-5}
\expandafter\def\csname sga@target@IV.19.9.6\endcsname{ega4.c04.prop.19-9-6}
\expandafter\def\csname sga@target@IV.19.9.7\endcsname{ega4.c04.cor.19-9-7}
\expandafter\def\csname sga@target@IV.19.9.8\endcsname{ega4.c04.prop.19-9-8}
\expandafter\def\csname sga@target@IV.19.9.9\endcsname{ega4.c04.rem.19-9-9}
\expandafter\def\csname sga@target@IV.2.1.1\endcsname{ega4.c04.statement.2-1-1}
\expandafter\def\csname sga@target@IV.2.1.10\endcsname{ega4.c04.cor.2-1-10}
\expandafter\def\csname sga@target@IV.2.1.11\endcsname{ega4.c04.prop.2-1-11}
\expandafter\def\csname sga@target@IV.2.1.12\endcsname{ega4.c04.prop.2-1-12}
\expandafter\def\csname sga@target@IV.2.1.13\endcsname{ega4.c04.prop.2-1-13}
\expandafter\def\csname sga@target@IV.2.1.14\endcsname{ega4.c04.prop.2-1-14}
\expandafter\def\csname sga@target@IV.2.1.2\endcsname{ega4.c04.prop.2-1-2}
\expandafter\def\csname sga@target@IV.2.1.3\endcsname{ega4.c04.prop.2-1-3}
\expandafter\def\csname sga@target@IV.2.1.4\endcsname{ega4.c04.prop.2-1-4}
\expandafter\def\csname sga@target@IV.2.1.5\endcsname{ega4.c04.prop.2-1-5}
\expandafter\def\csname sga@target@IV.2.1.6\endcsname{ega4.c04.cor.2-1-6}
\expandafter\def\csname sga@target@IV.2.1.7\endcsname{ega4.c04.cor.2-1-7}
\expandafter\def\csname sga@target@IV.2.1.8\endcsname{ega4.c04.prop.2-1-8}
\expandafter\def\csname sga@target@IV.2.1.9\endcsname{ega4.c04.cor.2-1-9}
\expandafter\def\csname sga@target@IV.2.2.1\endcsname{ega4.c04.prop.2-2-1}
\expandafter\def\csname sga@target@IV.2.2.10\endcsname{ega4.c04.prop.2-2-10}
\expandafter\def\csname sga@target@IV.2.2.10.1\endcsname{ega4.c04.eq.2-2-10-1}
\expandafter\def\csname sga@target@IV.2.2.11\endcsname{ega4.c04.cor.2-2-11}
\expandafter\def\csname sga@target@IV.2.2.12\endcsname{ega4.c04.cor.2-2-12}
\expandafter\def\csname sga@target@IV.2.2.13\endcsname{ega4.c04.cor.2-2-13}
\expandafter\def\csname sga@target@IV.2.2.14\endcsname{ega4.c04.prop.2-2-14}
\expandafter\def\csname sga@target@IV.2.2.15\endcsname{ega4.c04.prop.2-2-15}
\expandafter\def\csname sga@target@IV.2.2.16\endcsname{ega4.c04.cor.2-2-16}
\expandafter\def\csname sga@target@IV.2.2.17\endcsname{ega4.c04.prop.2-2-17}
\expandafter\def\csname sga@target@IV.2.2.2\endcsname{ega4.c04.prop.2-2-2}
\expandafter\def\csname sga@target@IV.2.2.2.1\endcsname{ega4.c04.eq.2-2-2-1}
\expandafter\def\csname sga@target@IV.2.2.3\endcsname{ega4.c04.cor.2-2-3}
\expandafter\def\csname sga@target@IV.2.2.4\endcsname{ega4.c04.def.2-2-4}
\expandafter\def\csname sga@target@IV.2.2.5\endcsname{ega4.c04.cor.2-2-5}
\expandafter\def\csname sga@target@IV.2.2.6\endcsname{ega4.c04.cor.2-2-6}
\expandafter\def\csname sga@target@IV.2.2.7\endcsname{ega4.c04.prop.2-2-7}
\expandafter\def\csname sga@target@IV.2.2.8\endcsname{ega4.c04.cor.2-2-8}
\expandafter\def\csname sga@target@IV.2.2.8.1\endcsname{ega4.c04.eq.2-2-8-1}
\expandafter\def\csname sga@target@IV.2.2.9\endcsname{ega4.c04.rem.2-2-9}
\expandafter\def\csname sga@target@IV.2.3.1\endcsname{ega4.c04.lem.2-3-1}
\expandafter\def\csname sga@target@IV.2.3.1.1\endcsname{ega4.c04.eq.2-3-1-1}
\expandafter\def\csname sga@target@IV.2.3.10\endcsname{ega4.c04.thm.2-3-10}
\expandafter\def\csname sga@target@IV.2.3.11\endcsname{ega4.c04.cor.2-3-11}
\expandafter\def\csname sga@target@IV.2.3.12\endcsname{ega4.c04.cor.2-3-12}
\expandafter\def\csname sga@target@IV.2.3.13\endcsname{ega4.c04.cor.2-3-13}
\expandafter\def\csname sga@target@IV.2.3.14\endcsname{ega4.c04.cor.2-3-14}
\expandafter\def\csname sga@target@IV.2.3.15\endcsname{ega4.c04.rem.2-3-15}
\expandafter\def\csname sga@target@IV.2.3.2\endcsname{ega4.c04.prop.2-3-2}
\expandafter\def\csname sga@target@IV.2.3.3\endcsname{ega4.c04.statement.2-3-3}
\expandafter\def\csname sga@target@IV.2.3.4\endcsname{ega4.c04.prop.2-3-4}
\expandafter\def\csname sga@target@IV.2.3.5\endcsname{ega4.c04.cor.2-3-5}
\expandafter\def\csname sga@target@IV.2.3.6\endcsname{ega4.c04.prop.2-3-6}
\expandafter\def\csname sga@target@IV.2.3.7\endcsname{ega4.c04.prop.2-3-7}
\expandafter\def\csname sga@target@IV.2.3.8\endcsname{ega4.c04.cor.2-3-8}
\expandafter\def\csname sga@target@IV.2.3.9\endcsname{ega4.c04.rem.2-3-9}
\expandafter\def\csname sga@target@IV.2.4.1\endcsname{ega4.c04.statement.2-4-1}
\expandafter\def\csname sga@target@IV.2.4.10\endcsname{ega4.c04.cor.2-4-10}
\expandafter\def\csname sga@target@IV.2.4.11\endcsname{ega4.c04.rem.2-4-11}
\expandafter\def\csname sga@target@IV.2.4.2\endcsname{ega4.c04.def.2-4-2}
\expandafter\def\csname sga@target@IV.2.4.3\endcsname{ega4.c04.prop.2-4-3}
\expandafter\def\csname sga@target@IV.2.4.4\endcsname{ega4.c04.prop.2-4-4}
\expandafter\def\csname sga@target@IV.2.4.5\endcsname{ega4.c04.prop.2-4-5}
\expandafter\def\csname sga@target@IV.2.4.6\endcsname{ega4.c04.thm.2-4-6}
\expandafter\def\csname sga@target@IV.2.4.7\endcsname{ega4.c04.cor.2-4-7}
\expandafter\def\csname sga@target@IV.2.4.8\endcsname{ega4.c04.rem.2-4-8}
\expandafter\def\csname sga@target@IV.2.4.9\endcsname{ega4.c04.prop.2-4-9}
\expandafter\def\csname sga@target@IV.2.5.1\endcsname{ega4.c04.prop.2-5-1}
\expandafter\def\csname sga@target@IV.2.5.2\endcsname{ega4.c04.prop.2-5-2}
\expandafter\def\csname sga@target@IV.2.5.3\endcsname{ega4.c04.lem.2-5-3}
\expandafter\def\csname sga@target@IV.2.5.4\endcsname{ega4.c04.rem.2-5-4}
\expandafter\def\csname sga@target@IV.2.5.4.1\endcsname{ega4.c04.lem.2-5-4-1}
\expandafter\def\csname sga@target@IV.2.5.5\endcsname{ega4.c04.prop.2-5-5}
\expandafter\def\csname sga@target@IV.2.5.5.1\endcsname{ega4.c04.statement.2-5-5-1}
\expandafter\def\csname sga@target@IV.2.5.5.2\endcsname{ega4.c04.lem.2-5-5-2}
\expandafter\def\csname sga@target@IV.2.5.5.3\endcsname{ega4.c04.statement.2-5-5-3}
\expandafter\def\csname sga@target@IV.2.5.5.4\endcsname{ega4.c04.rem.2-5-5-4}
\expandafter\def\csname sga@target@IV.2.5.6\endcsname{ega4.c04.lem.2-5-6}
\expandafter\def\csname sga@target@IV.2.5.7\endcsname{ega4.c04.cor.2-5-7}
\expandafter\def\csname sga@target@IV.2.5.8\endcsname{ega4.c04.prop.2-5-8}
\expandafter\def\csname sga@target@IV.2.5.8.1\endcsname{ega4.c04.statement.2-5-8-1}
\expandafter\def\csname sga@target@IV.2.5.8.2\endcsname{ega4.c04.statement.2-5-8-2}
\expandafter\def\csname sga@target@IV.2.5.8.3\endcsname{ega4.c04.statement.2-5-8-3}
\expandafter\def\csname sga@target@IV.2.5.8.4\endcsname{ega4.c04.statement.2-5-8-4}
\expandafter\def\csname sga@target@IV.2.5.8.5\endcsname{ega4.c04.statement.2-5-8-5}
\expandafter\def\csname sga@target@IV.2.5.8.6\endcsname{ega4.c04.statement.2-5-8-6}
\expandafter\def\csname sga@target@IV.2.5.8.7\endcsname{ega4.c04.statement.2-5-8-7}
\expandafter\def\csname sga@target@IV.2.5.8.8\endcsname{ega4.c04.statement.2-5-8-8}
\expandafter\def\csname sga@target@IV.2.5.8.9\endcsname{ega4.c04.statement.2-5-8-9}
\expandafter\def\csname sga@target@IV.2.6.1\endcsname{ega4.c04.prop.2-6-1}
\expandafter\def\csname sga@target@IV.2.6.2\endcsname{ega4.c04.prop.2-6-2}
\expandafter\def\csname sga@target@IV.2.6.3\endcsname{ega4.c04.rem.2-6-3}
\expandafter\def\csname sga@target@IV.2.6.4\endcsname{ega4.c04.cor.2-6-4}
\expandafter\def\csname sga@target@IV.2.7.1\endcsname{ega4.c04.prop.2-7-1}
\expandafter\def\csname sga@target@IV.2.7.1.1\endcsname{ega4.c04.lem.2-7-1-1}
\expandafter\def\csname sga@target@IV.2.7.1.viii\endcsname{ega4.c04.item.2-7-1-viii}
\expandafter\def\csname sga@target@IV.2.7.2\endcsname{ega4.c04.cor.2-7-2}
\expandafter\def\csname sga@target@IV.2.7.3\endcsname{ega4.c04.sec.2-7-3}
\expandafter\def\csname sga@target@IV.2.8.1\endcsname{ega4.c04.prop.2-8-1}
\expandafter\def\csname sga@target@IV.2.8.1.1\endcsname{ega4.c04.lem.2-8-1-1}
\expandafter\def\csname sga@target@IV.2.8.2\endcsname{ega4.c04.cor.2-8-2}
\expandafter\def\csname sga@target@IV.2.8.3\endcsname{ega4.c04.statement.2-8-3}
\expandafter\def\csname sga@target@IV.2.8.4\endcsname{ega4.c04.prop.2-8-4}
\expandafter\def\csname sga@target@IV.2.8.4.1\endcsname{ega4.c04.statement.2-8-4-1}
\expandafter\def\csname sga@target@IV.2.8.5\endcsname{ega4.c04.prop.2-8-5}
\expandafter\def\csname sga@target@IV.2.8.6\endcsname{ega4.c04.cor.2-8-6}
\expandafter\def\csname sga@target@IV.2.8.6.1\endcsname{ega4.c04.statement.2-8-6-1}
\expandafter\def\csname sga@target@IV.20\endcsname{ega4.c04.sec.20}
\expandafter\def\csname sga@target@IV.20.0\endcsname{ega4.c04.subsec.20-0}
\expandafter\def\csname sga@target@IV.20.1\endcsname{ega4.c04.subsec.20-1}
\expandafter\def\csname sga@target@IV.20.1.1\endcsname{ega4.c04.statement.20-1-1}
\expandafter\def\csname sga@target@IV.20.1.1.1\endcsname{ega4.c04.eq.20-1-1-1}
\expandafter\def\csname sga@target@IV.20.1.10\endcsname{ega4.c04.statement.20-1-10}
\expandafter\def\csname sga@target@IV.20.1.11\endcsname{ega4.c04.statement.20-1-11}
\expandafter\def\csname sga@target@IV.20.1.11.1\endcsname{ega4.c04.eq.20-1-11-1}
\expandafter\def\csname sga@target@IV.20.1.11.2\endcsname{ega4.c04.eq.20-1-11-2}
\expandafter\def\csname sga@target@IV.20.1.12\endcsname{ega4.c04.prop.20-1-12}
\expandafter\def\csname sga@target@IV.20.1.13\endcsname{ega4.c04.rem.20-1-13}
\expandafter\def\csname sga@target@IV.20.1.2\endcsname{ega4.c04.statement.20-1-2}
\expandafter\def\csname sga@target@IV.20.1.2.1\endcsname{ega4.c04.statement.20-1-2-1}
\expandafter\def\csname sga@target@IV.20.1.2.2\endcsname{ega4.c04.statement.20-1-2-2}
\expandafter\def\csname sga@target@IV.20.1.3\endcsname{ega4.c04.statement.20-1-3}
\expandafter\def\csname sga@target@IV.20.1.3.1\endcsname{ega4.c04.statement.20-1-3-1}
\expandafter\def\csname sga@target@IV.20.1.4\endcsname{ega4.c04.statement.20-1-4}
\expandafter\def\csname sga@target@IV.20.1.4.1\endcsname{ega4.c04.statement.20-1-4-1}
\expandafter\def\csname sga@target@IV.20.1.5\endcsname{ega4.c04.statement.20-1-5}
\expandafter\def\csname sga@target@IV.20.1.5.1\endcsname{ega4.c04.statement.20-1-5-1}
\expandafter\def\csname sga@target@IV.20.1.6\endcsname{ega4.c04.prop.20-1-6}
\expandafter\def\csname sga@target@IV.20.1.7\endcsname{ega4.c04.statement.20-1-7}
\expandafter\def\csname sga@target@IV.20.1.8\endcsname{ega4.c04.statement.20-1-8}
\expandafter\def\csname sga@target@IV.20.1.9\endcsname{ega4.c04.statement.20-1-9}
\expandafter\def\csname sga@target@IV.20.2\endcsname{ega4.c04.subsec.20-2}
\expandafter\def\csname sga@target@IV.20.2.1\endcsname{ega4.c04.statement.20-2-1}
\expandafter\def\csname sga@target@IV.20.2.10\endcsname{ega4.c04.statement.20-2-10}
\expandafter\def\csname sga@target@IV.20.2.10.1\endcsname{ega4.c04.eq.20-2-10-1}
\expandafter\def\csname sga@target@IV.20.2.10.2\endcsname{ega4.c04.eq.20-2-10-2}
\expandafter\def\csname sga@target@IV.20.2.11\endcsname{ega4.c04.prop.20-2-11}
\expandafter\def\csname sga@target@IV.20.2.11.1\endcsname{ega4.c04.eq.20-2-11-1}
\expandafter\def\csname sga@target@IV.20.2.12\endcsname{ega4.c04.cor.20-2-12}
\expandafter\def\csname sga@target@IV.20.2.13\endcsname{ega4.c04.statement.20-2-13}
\expandafter\def\csname sga@target@IV.20.2.14\endcsname{ega4.c04.statement.20-2-14}
\expandafter\def\csname sga@target@IV.20.2.15\endcsname{ega4.c04.prop.20-2-15}
\expandafter\def\csname sga@target@IV.20.2.16\endcsname{ega4.c04.rem.20-2-16}
\expandafter\def\csname sga@target@IV.20.2.2\endcsname{ega4.c04.statement.20-2-2}
\expandafter\def\csname sga@target@IV.20.2.3\endcsname{ega4.c04.statement.20-2-3}
\expandafter\def\csname sga@target@IV.20.2.3.1\endcsname{ega4.c04.statement.20-2-3-1}
\expandafter\def\csname sga@target@IV.20.2.4\endcsname{ega4.c04.prop.20-2-4}
\expandafter\def\csname sga@target@IV.20.2.5\endcsname{ega4.c04.cor.20-2-5}
\expandafter\def\csname sga@target@IV.20.2.6\endcsname{ega4.c04.cor.20-2-6}
\expandafter\def\csname sga@target@IV.20.2.7\endcsname{ega4.c04.rem.20-2-7}
\expandafter\def\csname sga@target@IV.20.2.8\endcsname{ega4.c04.statement.20-2-8}
\expandafter\def\csname sga@target@IV.20.2.9\endcsname{ega4.c04.lem.20-2-9}
\expandafter\def\csname sga@target@IV.20.3\endcsname{ega4.c04.subsec.20-3}
\expandafter\def\csname sga@target@IV.20.3.1\endcsname{ega4.c04.statement.20-3-1}
\expandafter\def\csname sga@target@IV.20.3.10\endcsname{ega4.c04.lem.20-3-10}
\expandafter\def\csname sga@target@IV.20.3.11\endcsname{ega4.c04.prop.20-3-11}
\expandafter\def\csname sga@target@IV.20.3.11.1\endcsname{ega4.c04.eq.20-3-11-1}
\expandafter\def\csname sga@target@IV.20.3.12\endcsname{ega4.c04.cor.20-3-12}
\expandafter\def\csname sga@target@IV.20.3.13\endcsname{ega4.c04.rem.20-3-13}
\expandafter\def\csname sga@target@IV.20.3.2\endcsname{ega4.c04.statement.20-3-2}
\expandafter\def\csname sga@target@IV.20.3.2.1\endcsname{ega4.c04.eq.20-3-2-1}
\expandafter\def\csname sga@target@IV.20.3.3\endcsname{ega4.c04.statement.20-3-3}
\expandafter\def\csname sga@target@IV.20.3.4\endcsname{ega4.c04.statement.20-3-4}
\expandafter\def\csname sga@target@IV.20.3.5\endcsname{ega4.c04.prop.20-3-5}
\expandafter\def\csname sga@target@IV.20.3.6\endcsname{ega4.c04.statement.20-3-6}
\expandafter\def\csname sga@target@IV.20.3.6.1\endcsname{ega4.c04.eq.20-3-6-1}
\expandafter\def\csname sga@target@IV.20.3.7\endcsname{ega4.c04.prop.20-3-7}
\expandafter\def\csname sga@target@IV.20.3.8\endcsname{ega4.c04.prop.20-3-8}
\expandafter\def\csname sga@target@IV.20.3.8.1\endcsname{ega4.c04.eq.20-3-8-1}
\expandafter\def\csname sga@target@IV.20.3.8.2\endcsname{ega4.c04.lem.20-3-8-2}
\expandafter\def\csname sga@target@IV.20.3.8.3\endcsname{ega4.c04.lem.20-3-8-3}
\expandafter\def\csname sga@target@IV.20.3.8.4\endcsname{ega4.c04.lem.20-3-8-4}
\expandafter\def\csname sga@target@IV.20.3.8.5\endcsname{ega4.c04.sec.20-3-8-5}
\expandafter\def\csname sga@target@IV.20.3.8.6\endcsname{ega4.c04.sec.20-3-8-6}
\expandafter\def\csname sga@target@IV.20.3.8.7\endcsname{ega4.c04.sec.20-3-8-7}
\expandafter\def\csname sga@target@IV.20.3.8.8\endcsname{ega4.c04.lem.20-3-8-8}
\expandafter\def\csname sga@target@IV.20.3.9\endcsname{ega4.c04.cor.20-3-9}
\expandafter\def\csname sga@target@IV.20.4\endcsname{ega4.c04.subsec.20-4}
\expandafter\def\csname sga@target@IV.20.4.1\endcsname{ega4.c04.statement.20-4-1}
\expandafter\def\csname sga@target@IV.20.4.10\endcsname{ega4.c04.cor.20-4-10}
\expandafter\def\csname sga@target@IV.20.4.11\endcsname{ega4.c04.prop.20-4-11}
\expandafter\def\csname sga@target@IV.20.4.12\endcsname{ega4.c04.cor.20-4-12}
\expandafter\def\csname sga@target@IV.20.4.2\endcsname{ega4.c04.statement.20-4-2}
\expandafter\def\csname sga@target@IV.20.4.3\endcsname{ega4.c04.statement.20-4-3}
\expandafter\def\csname sga@target@IV.20.4.4\endcsname{ega4.c04.prop.20-4-4}
\expandafter\def\csname sga@target@IV.20.4.5\endcsname{ega4.c04.prop.20-4-5}
\expandafter\def\csname sga@target@IV.20.4.6\endcsname{ega4.c04.prop.20-4-6}
\expandafter\def\csname sga@target@IV.20.4.7\endcsname{ega4.c04.cor.20-4-7}
\expandafter\def\csname sga@target@IV.20.4.8\endcsname{ega4.c04.rem.20-4-8}
\expandafter\def\csname sga@target@IV.20.4.9\endcsname{ega4.c04.cor.20-4-9}
\expandafter\def\csname sga@target@IV.20.5\endcsname{ega4.c04.subsec.20-5}
\expandafter\def\csname sga@target@IV.20.5.1\endcsname{ega4.c04.statement.20-5-1}
\expandafter\def\csname sga@target@IV.20.5.2\endcsname{ega4.c04.statement.20-5-2}
\expandafter\def\csname sga@target@IV.20.5.2.1\endcsname{ega4.c04.statement.20-5-2-1}
\expandafter\def\csname sga@target@IV.20.5.3\endcsname{ega4.c04.prop.20-5-3}
\expandafter\def\csname sga@target@IV.20.5.4\endcsname{ega4.c04.statement.20-5-4}
\expandafter\def\csname sga@target@IV.20.5.5\endcsname{ega4.c04.statement.20-5-5}
\expandafter\def\csname sga@target@IV.20.5.6\endcsname{ega4.c04.statement.20-5-6}
\expandafter\def\csname sga@target@IV.20.5.7\endcsname{ega4.c04.statement.20-5-7}
\expandafter\def\csname sga@target@IV.20.5.8\endcsname{ega4.c04.statement.20-5-8}
\expandafter\def\csname sga@target@IV.20.6\endcsname{ega4.c04.subsec.20-6}
\expandafter\def\csname sga@target@IV.20.6.1\endcsname{ega4.c04.statement.20-6-1}
\expandafter\def\csname sga@target@IV.20.6.1.1\endcsname{ega4.c04.eq.20-6-1-1}
\expandafter\def\csname sga@target@IV.20.6.2\endcsname{ega4.c04.statement.20-6-2}
\expandafter\def\csname sga@target@IV.20.6.2.1\endcsname{ega4.c04.statement.20-6-2-1}
\expandafter\def\csname sga@target@IV.20.6.2.2\endcsname{ega4.c04.eq.20-6-2-2}
\expandafter\def\csname sga@target@IV.20.6.3\endcsname{ega4.c04.prop.20-6-3}
\expandafter\def\csname sga@target@IV.20.6.4\endcsname{ega4.c04.prop.20-6-4}
\expandafter\def\csname sga@target@IV.20.6.5\endcsname{ega4.c04.statement.20-6-5}
\expandafter\def\csname sga@target@IV.20.6.5.1\endcsname{ega4.c04.eq.20-6-5-1}
\expandafter\def\csname sga@target@IV.20.6.6\endcsname{ega4.c04.prop.20-6-6}
\expandafter\def\csname sga@target@IV.20.6.7\endcsname{ega4.c04.statement.20-6-7}
\expandafter\def\csname sga@target@IV.20.6.8\endcsname{ega4.c04.prop.20-6-8}
\expandafter\def\csname sga@target@IV.20.6.8.1\endcsname{ega4.c04.eq.20-6-8-1}
\expandafter\def\csname sga@target@IV.20.6.9\endcsname{ega4.c04.statement.20-6-9}
\expandafter\def\csname sga@target@IV.20.6.9.1\endcsname{ega4.c04.statement.20-6-9-1}
\expandafter\def\csname sga@target@IV.21.1\endcsname{ega4.c04.subsec.21-1.key-iv-21-1}
\expandafter\def\csname sga@target@IV.21.1.1\endcsname{ega4.c04.statement.21-1-1}
\expandafter\def\csname sga@target@IV.21.1.10\endcsname{ega4.c04.prop.21-1-10}
\expandafter\def\csname sga@target@IV.21.1.2\endcsname{ega4.c04.statement.21-1-2}
\expandafter\def\csname sga@target@IV.21.1.3\endcsname{ega4.c04.statement.21-1-3}
\expandafter\def\csname sga@target@IV.21.1.3.1\endcsname{ega4.c04.statement.21-1-3-1}
\expandafter\def\csname sga@target@IV.21.1.3.2\endcsname{ega4.c04.statement.21-1-3-2}
\expandafter\def\csname sga@target@IV.21.1.3.3\endcsname{ega4.c04.statement.21-1-3-3}
\expandafter\def\csname sga@target@IV.21.1.3.4\endcsname{ega4.c04.statement.21-1-3-4}
\expandafter\def\csname sga@target@IV.21.1.4\endcsname{ega4.c04.statement.21-1-4}
\expandafter\def\csname sga@target@IV.21.1.4.1\endcsname{ega4.c04.eq.21-1-4-1}
\expandafter\def\csname sga@target@IV.21.1.4.2\endcsname{ega4.c04.eq.21-1-4-2}
\expandafter\def\csname sga@target@IV.21.1.4.3\endcsname{ega4.c04.eq.21-1-4-3}
\expandafter\def\csname sga@target@IV.21.1.5\endcsname{ega4.c04.statement.21-1-5}
\expandafter\def\csname sga@target@IV.21.1.5.1\endcsname{ega4.c04.statement.21-1-5-1}
\expandafter\def\csname sga@target@IV.21.1.5.2\endcsname{ega4.c04.eq.21-1-5-2}
\expandafter\def\csname sga@target@IV.21.1.6\endcsname{ega4.c04.def.21-1-6}
\expandafter\def\csname sga@target@IV.21.1.6.1\endcsname{ega4.c04.eq.21-1-6-1}
\expandafter\def\csname sga@target@IV.21.1.6.2\endcsname{ega4.c04.eq.21-1-6-2}
\expandafter\def\csname sga@target@IV.21.1.6.3\endcsname{ega4.c04.eq.21-1-6-3}
\expandafter\def\csname sga@target@IV.21.1.6.4\endcsname{ega4.c04.eq.21-1-6-4}
\expandafter\def\csname sga@target@IV.21.1.7\endcsname{ega4.c04.statement.21-1-7}
\expandafter\def\csname sga@target@IV.21.1.7.1\endcsname{ega4.c04.statement.21-1-7-1}
\expandafter\def\csname sga@target@IV.21.1.8\endcsname{ega4.c04.prop.21-1-8}
\expandafter\def\csname sga@target@IV.21.1.9\endcsname{ega4.c04.cor.21-1-9}
\expandafter\def\csname sga@target@IV.21.10.1\endcsname{ega4.c04.statement.21-10-1}
\expandafter\def\csname sga@target@IV.21.10.1.1\endcsname{ega4.c04.eq.21-10-1-1}
\expandafter\def\csname sga@target@IV.21.10.1.2\endcsname{ega4.c04.eq.21-10-1-2}
\expandafter\def\csname sga@target@IV.21.10.10\endcsname{ega4.c04.statement.21-10-10}
\expandafter\def\csname sga@target@IV.21.10.11\endcsname{ega4.c04.statement.21-10-11}
\expandafter\def\csname sga@target@IV.21.10.11.1\endcsname{ega4.c04.eq.21-10-11-1}
\expandafter\def\csname sga@target@IV.21.10.12\endcsname{ega4.c04.rem.21-10-12}
\expandafter\def\csname sga@target@IV.21.10.13\endcsname{ega4.c04.lem.21-10-13}
\expandafter\def\csname sga@target@IV.21.10.13.1\endcsname{ega4.c04.prop.21-10-13-1}
\expandafter\def\csname sga@target@IV.21.10.13bis\endcsname{ega4.c04.prop.21-10-13bis}
\expandafter\def\csname sga@target@IV.21.10.14\endcsname{ega4.c04.statement.21-10-14}
\expandafter\def\csname sga@target@IV.21.10.14.1\endcsname{ega4.c04.eq.21-10-14-1}
\expandafter\def\csname sga@target@IV.21.10.14.2\endcsname{ega4.c04.eq.21-10-14-2}
\expandafter\def\csname sga@target@IV.21.10.15\endcsname{ega4.c04.prop.21-10-15}
\expandafter\def\csname sga@target@IV.21.10.16\endcsname{ega4.c04.prop.21-10-16}
\expandafter\def\csname sga@target@IV.21.10.16.1\endcsname{ega4.c04.prop.21-10-16-1}
\expandafter\def\csname sga@target@IV.21.10.16.2\endcsname{ega4.c04.eq.21-10-16-2}
\expandafter\def\csname sga@target@IV.21.10.17\endcsname{ega4.c04.prop.21-10-17}
\expandafter\def\csname sga@target@IV.21.10.17.1\endcsname{ega4.c04.prop.21-10-17-1}
\expandafter\def\csname sga@target@IV.21.10.17.10\endcsname{ega4.c04.eq.21-10-17-10}
\expandafter\def\csname sga@target@IV.21.10.17.2\endcsname{ega4.c04.eq.21-10-17-2}
\expandafter\def\csname sga@target@IV.21.10.17.3\endcsname{ega4.c04.lem.21-10-17-3}
\expandafter\def\csname sga@target@IV.21.10.17.4\endcsname{ega4.c04.lem.21-10-17-4}
\expandafter\def\csname sga@target@IV.21.10.17.5\endcsname{ega4.c04.eq.21-10-17-5}
\expandafter\def\csname sga@target@IV.21.10.17.6\endcsname{ega4.c04.eq.21-10-17-6}
\expandafter\def\csname sga@target@IV.21.10.17.7\endcsname{ega4.c04.lem.21-10-17-7}
\expandafter\def\csname sga@target@IV.21.10.17.8\endcsname{ega4.c04.lem.21-10-17-8}
\expandafter\def\csname sga@target@IV.21.10.17.9\endcsname{ega4.c04.eq.21-10-17-9}
\expandafter\def\csname sga@target@IV.21.10.18\endcsname{ega4.c04.prop.21-10-18}
\expandafter\def\csname sga@target@IV.21.10.18.1\endcsname{ega4.c04.prop.21-10-18-1}
\expandafter\def\csname sga@target@IV.21.10.19\endcsname{ega4.c04.cor.21-10-19}
\expandafter\def\csname sga@target@IV.21.10.19.1\endcsname{ega4.c04.cor.21-10-19-1}
\expandafter\def\csname sga@target@IV.21.10.2\endcsname{ega4.c04.statement.21-10-2}
\expandafter\def\csname sga@target@IV.21.10.3\endcsname{ega4.c04.statement.21-10-3}
\expandafter\def\csname sga@target@IV.21.10.3.1\endcsname{ega4.c04.statement.21-10-3-1}
\expandafter\def\csname sga@target@IV.21.10.4\endcsname{ega4.c04.prop.21-10-4}
\expandafter\def\csname sga@target@IV.21.10.4.1\endcsname{ega4.c04.prop.21-10-4-1}
\expandafter\def\csname sga@target@IV.21.10.5\endcsname{ega4.c04.statement.21-10-5}
\expandafter\def\csname sga@target@IV.21.10.6\endcsname{ega4.c04.prop.21-10-6}
\expandafter\def\csname sga@target@IV.21.10.7\endcsname{ega4.c04.rem.21-10-7}
\expandafter\def\csname sga@target@IV.21.10.8\endcsname{ega4.c04.prop.21-10-8}
\expandafter\def\csname sga@target@IV.21.10.9\endcsname{ega4.c04.statement.21-10-9}
\expandafter\def\csname sga@target@IV.21.11\endcsname{ega4.c04.subsec.21-11.key-iv-21-11}
\expandafter\def\csname sga@target@IV.21.11.1\endcsname{ega4.c04.thm.21-11-1}
\expandafter\def\csname sga@target@IV.21.11.1.1\endcsname{ega4.c04.eq.21-11-1-1}
\expandafter\def\csname sga@target@IV.21.11.2\endcsname{ega4.c04.lem.21-11-2}
\expandafter\def\csname sga@target@IV.21.12\endcsname{ega4.c04.subsec.21-12.key-iv-21-12}
\expandafter\def\csname sga@target@IV.21.12.1\endcsname{ega4.c04.statement.21-12-1}
\expandafter\def\csname sga@target@IV.21.12.1.1\endcsname{ega4.c04.eq.21-12-1-1}
\expandafter\def\csname sga@target@IV.21.12.10\endcsname{ega4.c04.prop.21-12-10}
\expandafter\def\csname sga@target@IV.21.12.12\endcsname{ega4.c04.thm.21-12-12}
\expandafter\def\csname sga@target@IV.21.12.13\endcsname{ega4.c04.cor.21-12-13}
\expandafter\def\csname sga@target@IV.21.12.14\endcsname{ega4.c04.item.21-12-14}
\expandafter\def\csname sga@target@IV.21.12.15\endcsname{ega4.c04.prop.21-12-15}
\expandafter\def\csname sga@target@IV.21.12.16\endcsname{ega4.c04.cor.21-12-16}
\expandafter\def\csname sga@target@IV.21.12.17\endcsname{ega4.c04.item.21-12-17}
\expandafter\def\csname sga@target@IV.21.12.2\endcsname{ega4.c04.prop.21-12-2}
\expandafter\def\csname sga@target@IV.21.12.2.1\endcsname{ega4.c04.prop.21-12-2-1}
\expandafter\def\csname sga@target@IV.21.12.3\endcsname{ega4.c04.cor.21-12-3}
\expandafter\def\csname sga@target@IV.21.12.4\endcsname{ega4.c04.eq.21-12-4}
\expandafter\def\csname sga@target@IV.21.12.5\endcsname{ega4.c04.statement.21-12-5}
\expandafter\def\csname sga@target@IV.21.12.5.1\endcsname{ega4.c04.eq.21-12-5-1}
\expandafter\def\csname sga@target@IV.21.12.5.2\endcsname{ega4.c04.eq.21-12-5-2}
\expandafter\def\csname sga@target@IV.21.12.5.3\endcsname{ega4.c04.eq.21-12-5-3}
\expandafter\def\csname sga@target@IV.21.12.6\endcsname{ega4.c04.thm.21-12-6}
\expandafter\def\csname sga@target@IV.21.12.7\endcsname{ega4.c04.cor.21-12-7}
\expandafter\def\csname sga@target@IV.21.12.8\endcsname{ega4.c04.statement.21-12-8}
\expandafter\def\csname sga@target@IV.21.12.9\endcsname{ega4.c04.statement.21-12-9}
\expandafter\def\csname sga@target@IV.21.13\endcsname{ega4.c04.subsec.21-13.key-iv-21-13}
\expandafter\def\csname sga@target@IV.21.13.1\endcsname{ega4.c04.def.21-13-1}
\expandafter\def\csname sga@target@IV.21.13.10\endcsname{ega4.c04.prop.21-13-10}
\expandafter\def\csname sga@target@IV.21.13.11\endcsname{ega4.c04.cor.21-13-11}
\expandafter\def\csname sga@target@IV.21.13.12\endcsname{ega4.c04.prop.21-13-12}
\expandafter\def\csname sga@target@IV.21.13.12.1\endcsname{ega4.c04.statement.21-13-12-1}
\expandafter\def\csname sga@target@IV.21.13.13\endcsname{ega4.c04.rem.21-13-13}
\expandafter\def\csname sga@target@IV.21.13.14\endcsname{ega4.c04.prop.21-13-14}
\expandafter\def\csname sga@target@IV.21.13.15\endcsname{ega4.c04.cor.21-13-15}
\expandafter\def\csname sga@target@IV.21.13.16\endcsname{ega4.c04.rem.21-13-16}
\expandafter\def\csname sga@target@IV.21.13.2\endcsname{ega4.c04.lem.21-13-2}
\expandafter\def\csname sga@target@IV.21.13.2.1\endcsname{ega4.c04.eq.21-13-2-1}
\expandafter\def\csname sga@target@IV.21.13.3\endcsname{ega4.c04.cor.21-13-3}
\expandafter\def\csname sga@target@IV.21.13.4\endcsname{ega4.c04.lem.21-13-4}
\expandafter\def\csname sga@target@IV.21.13.5\endcsname{ega4.c04.prop.21-13-5}
\expandafter\def\csname sga@target@IV.21.13.6\endcsname{ega4.c04.cor.21-13-6}
\expandafter\def\csname sga@target@IV.21.13.7\endcsname{ega4.c04.def.21-13-7}
\expandafter\def\csname sga@target@IV.21.13.8\endcsname{ega4.c04.prop.21-13-8}
\expandafter\def\csname sga@target@IV.21.13.9\endcsname{ega4.c04.statement.21-13-9}
\expandafter\def\csname sga@target@IV.21.13.9.1\endcsname{ega4.c04.eq.21-13-9-1}
\expandafter\def\csname sga@target@IV.21.14\endcsname{ega4.c04.subsec.21-14.key-iv-21-14}
\expandafter\def\csname sga@target@IV.21.14.1\endcsname{ega4.c04.thm.21-14-1}
\expandafter\def\csname sga@target@IV.21.14.2\endcsname{ega4.c04.cor.21-14-2}
\expandafter\def\csname sga@target@IV.21.14.3\endcsname{ega4.c04.prop.21-14-3}
\expandafter\def\csname sga@target@IV.21.14.4\endcsname{ega4.c04.statement.21-14-4}
\expandafter\def\csname sga@target@IV.21.14.4bis\endcsname{ega4.c04.statement.21-14-4bis}
\expandafter\def\csname sga@target@IV.21.15\endcsname{ega4.c04.subsec.21-15.key-iv-21-15}
\expandafter\def\csname sga@target@IV.21.15.1\endcsname{ega4.c04.statement.21-15-1}
\expandafter\def\csname sga@target@IV.21.15.2\endcsname{ega4.c04.def.21-15-2}
\expandafter\def\csname sga@target@IV.21.15.3\endcsname{ega4.c04.statement.21-15-3}
\expandafter\def\csname sga@target@IV.21.15.3.1\endcsname{ega4.c04.eq.21-15-3-1}
\expandafter\def\csname sga@target@IV.21.15.3.2\endcsname{ega4.c04.eq.21-15-3-2}
\expandafter\def\csname sga@target@IV.21.15.3bis\endcsname{ega4.c04.statement.21-15-3bis}
\expandafter\def\csname sga@target@IV.21.15.4\endcsname{ega4.c04.prop.21-15-4}
\expandafter\def\csname sga@target@IV.21.15.5\endcsname{ega4.c04.prop.21-15-5}
\expandafter\def\csname sga@target@IV.21.15.6\endcsname{ega4.c04.prop.21-15-6}
\expandafter\def\csname sga@target@IV.21.15.7\endcsname{ega4.c04.statement.21-15-7}
\expandafter\def\csname sga@target@IV.21.15.8\endcsname{ega4.c04.statement.21-15-8}
\expandafter\def\csname sga@target@IV.21.15.9\endcsname{ega4.c04.statement.21-15-9}
\expandafter\def\csname sga@target@IV.21.15.9.1\endcsname{ega4.c04.eq.21-15-9-1}
\expandafter\def\csname sga@target@IV.21.2.1\endcsname{ega4.c04.statement.21-2-1}
\expandafter\def\csname sga@target@IV.21.2.10\endcsname{ega4.c04.statement.21-2-10}
\expandafter\def\csname sga@target@IV.21.2.10.1\endcsname{ega4.c04.eq.21-2-10-1}
\expandafter\def\csname sga@target@IV.21.2.11\endcsname{ega4.c04.prop.21-2-11}
\expandafter\def\csname sga@target@IV.21.2.11.1\endcsname{ega4.c04.prop.21-2-11-1}
\expandafter\def\csname sga@target@IV.21.2.12\endcsname{ega4.c04.statement.21-2-12}
\expandafter\def\csname sga@target@IV.21.2.2\endcsname{ega4.c04.prop.21-2-2}
\expandafter\def\csname sga@target@IV.21.2.3\endcsname{ega4.c04.cor.21-2-3}
\expandafter\def\csname sga@target@IV.21.2.4\endcsname{ega4.c04.statement.21-2-4}
\expandafter\def\csname sga@target@IV.21.2.5\endcsname{ega4.c04.statement.21-2-5}
\expandafter\def\csname sga@target@IV.21.2.5.1\endcsname{ega4.c04.statement.21-2-5-1}
\expandafter\def\csname sga@target@IV.21.2.5.2\endcsname{ega4.c04.eq.21-2-5-2}
\expandafter\def\csname sga@target@IV.21.2.5.3\endcsname{ega4.c04.eq.21-2-5-3}
\expandafter\def\csname sga@target@IV.21.2.5.4\endcsname{ega4.c04.eq.21-2-5-4}
\expandafter\def\csname sga@target@IV.21.2.6\endcsname{ega4.c04.prop.21-2-6}
\expandafter\def\csname sga@target@IV.21.2.6.1\endcsname{ega4.c04.eq.21-2-6-1}
\expandafter\def\csname sga@target@IV.21.2.7\endcsname{ega4.c04.statement.21-2-7}
\expandafter\def\csname sga@target@IV.21.2.7.1\endcsname{ega4.c04.statement.21-2-7-1}
\expandafter\def\csname sga@target@IV.21.2.7.2\endcsname{ega4.c04.eq.21-2-7-2}
\expandafter\def\csname sga@target@IV.21.2.8\endcsname{ega4.c04.statement.21-2-8}
\expandafter\def\csname sga@target@IV.21.2.8.1\endcsname{ega4.c04.statement.21-2-8-1}
\expandafter\def\csname sga@target@IV.21.2.8.2\endcsname{ega4.c04.statement.21-2-8-2}
\expandafter\def\csname sga@target@IV.21.2.8.3\endcsname{ega4.c04.statement.21-2-8-3}
\expandafter\def\csname sga@target@IV.21.2.9\endcsname{ega4.c04.statement.21-2-9}
\expandafter\def\csname sga@target@IV.21.2.9.1\endcsname{ega4.c04.statement.21-2-9-1}
\expandafter\def\csname sga@target@IV.21.2.9.2\endcsname{ega4.c04.eq.21-2-9-2}
\expandafter\def\csname sga@target@IV.21.2.9.3\endcsname{ega4.c04.eq.21-2-9-3}
\expandafter\def\csname sga@target@IV.21.2.9.4\endcsname{ega4.c04.eq.21-2-9-4}
\expandafter\def\csname sga@target@IV.21.3.1\endcsname{ega4.c04.statement.21-3-1}
\expandafter\def\csname sga@target@IV.21.3.2\endcsname{ega4.c04.statement.21-3-2}
\expandafter\def\csname sga@target@IV.21.3.2.1\endcsname{ega4.c04.statement.21-3-2-1}
\expandafter\def\csname sga@target@IV.21.3.2.2\endcsname{ega4.c04.eq.21-3-2-2}
\expandafter\def\csname sga@target@IV.21.3.2.3\endcsname{ega4.c04.eq.21-3-2-3}
\expandafter\def\csname sga@target@IV.21.3.2.4\endcsname{ega4.c04.eq.21-3-2-4}
\expandafter\def\csname sga@target@IV.21.3.3\endcsname{ega4.c04.prop.21-3-3}
\expandafter\def\csname sga@target@IV.21.3.3.1\endcsname{ega4.c04.eq.21-3-3-1}
\expandafter\def\csname sga@target@IV.21.3.4\endcsname{ega4.c04.prop.21-3-4}
\expandafter\def\csname sga@target@IV.21.3.5\endcsname{ega4.c04.cor.21-3-5}
\expandafter\def\csname sga@target@IV.21.3.6\endcsname{ega4.c04.rem.21-3-6}
\expandafter\def\csname sga@target@IV.21.3.6.1\endcsname{ega4.c04.eq.21-3-6-1}
\expandafter\def\csname sga@target@IV.21.4\endcsname{ega4.c04.subsec.21-4.key-iv-21-4}
\expandafter\def\csname sga@target@IV.21.4.1\endcsname{ega4.c04.statement.21-4-1}
\expandafter\def\csname sga@target@IV.21.4.10\endcsname{ega4.c04.cor.21-4-10}
\expandafter\def\csname sga@target@IV.21.4.11\endcsname{ega4.c04.cor.21-4-11}
\expandafter\def\csname sga@target@IV.21.4.12\endcsname{ega4.c04.rem.21-4-12}
\expandafter\def\csname sga@target@IV.21.4.13\endcsname{ega4.c04.cor.21-4-13}
\expandafter\def\csname sga@target@IV.21.4.2\endcsname{ega4.c04.def.21-4-2}
\expandafter\def\csname sga@target@IV.21.4.2.1\endcsname{ega4.c04.eq.21-4-2-1}
\expandafter\def\csname sga@target@IV.21.4.3\endcsname{ega4.c04.statement.21-4-3}
\expandafter\def\csname sga@target@IV.21.4.4\endcsname{ega4.c04.statement.21-4-4}
\expandafter\def\csname sga@target@IV.21.4.5\endcsname{ega4.c04.prop.21-4-5}
\expandafter\def\csname sga@target@IV.21.4.6\endcsname{ega4.c04.cor.21-4-6}
\expandafter\def\csname sga@target@IV.21.4.6.1\endcsname{ega4.c04.eq.21-4-6-1}
\expandafter\def\csname sga@target@IV.21.4.7\endcsname{ega4.c04.statement.21-4-7}
\expandafter\def\csname sga@target@IV.21.4.8\endcsname{ega4.c04.prop.21-4-8}
\expandafter\def\csname sga@target@IV.21.4.9\endcsname{ega4.c04.prop.21-4-9}
\expandafter\def\csname sga@target@IV.21.5\endcsname{ega4.c04.subsec.21-5.key-iv-21-5}
\expandafter\def\csname sga@target@IV.21.5.1\endcsname{ega4.c04.statement.21-5-1}
\expandafter\def\csname sga@target@IV.21.5.2\endcsname{ega4.c04.lem.21-5-2}
\expandafter\def\csname sga@target@IV.21.5.3\endcsname{ega4.c04.statement.21-5-3}
\expandafter\def\csname sga@target@IV.21.5.3.1\endcsname{ega4.c04.eq.21-5-3-1}
\expandafter\def\csname sga@target@IV.21.5.3.2\endcsname{ega4.c04.eq.21-5-3-2}
\expandafter\def\csname sga@target@IV.21.5.4\endcsname{ega4.c04.statement.21-5-4}
\expandafter\def\csname sga@target@IV.21.5.5\endcsname{ega4.c04.def.21-5-5}
\expandafter\def\csname sga@target@IV.21.5.5.1\endcsname{ega4.c04.eq.21-5-5-1}
\expandafter\def\csname sga@target@IV.21.5.5.2\endcsname{ega4.c04.eq.21-5-5-2}
\expandafter\def\csname sga@target@IV.21.5.5.3\endcsname{ega4.c04.eq.21-5-5-3}
\expandafter\def\csname sga@target@IV.21.5.6\endcsname{ega4.c04.prop.21-5-6}
\expandafter\def\csname sga@target@IV.21.5.6.1\endcsname{ega4.c04.prop.21-5-6-1}
\expandafter\def\csname sga@target@IV.21.5.7\endcsname{ega4.c04.prop.21-5-7}
\expandafter\def\csname sga@target@IV.21.5.7.1\endcsname{ega4.c04.prop.21-5-7-1}
\expandafter\def\csname sga@target@IV.21.5.7.2\endcsname{ega4.c04.lem.21-5-7-2}
\expandafter\def\csname sga@target@IV.21.5.7.3\endcsname{ega4.c04.lem.21-5-7-3}
\expandafter\def\csname sga@target@IV.21.5.7.4\endcsname{ega4.c04.eq.21-5-7-4}
\expandafter\def\csname sga@target@IV.21.5.8\endcsname{ega4.c04.prop.21-5-8}
\expandafter\def\csname sga@target@IV.21.5.8.1\endcsname{ega4.c04.eq.21-5-8-1}
\expandafter\def\csname sga@target@IV.21.6\endcsname{ega4.c04.subsec.21-6.key-iv-21-6}
\expandafter\def\csname sga@target@IV.21.6.1\endcsname{ega4.c04.statement.21-6-1}
\expandafter\def\csname sga@target@IV.21.6.10\endcsname{ega4.c04.cor.21-6-10}
\expandafter\def\csname sga@target@IV.21.6.10.1\endcsname{ega4.c04.eq.21-6-10-1}
\expandafter\def\csname sga@target@IV.21.6.11\endcsname{ega4.c04.cor.21-6-11}
\expandafter\def\csname sga@target@IV.21.6.12\endcsname{ega4.c04.prop.21-6-12}
\expandafter\def\csname sga@target@IV.21.6.12.1\endcsname{ega4.c04.eq.21-6-12-1}
\expandafter\def\csname sga@target@IV.21.6.12.2\endcsname{ega4.c04.eq.21-6-12-2}
\expandafter\def\csname sga@target@IV.21.6.13\endcsname{ega4.c04.cor.21-6-13}
\expandafter\def\csname sga@target@IV.21.6.14\endcsname{ega4.c04.cor.21-6-14}
\expandafter\def\csname sga@target@IV.21.6.2\endcsname{ega4.c04.statement.21-6-2}
\expandafter\def\csname sga@target@IV.21.6.3\endcsname{ega4.c04.statement.21-6-3}
\expandafter\def\csname sga@target@IV.21.6.4\endcsname{ega4.c04.statement.21-6-4}
\expandafter\def\csname sga@target@IV.21.6.4.1\endcsname{ega4.c04.statement.21-6-4-1}
\expandafter\def\csname sga@target@IV.21.6.4.2\endcsname{ega4.c04.eq.21-6-4-2}
\expandafter\def\csname sga@target@IV.21.6.5\endcsname{ega4.c04.statement.21-6-5}
\expandafter\def\csname sga@target@IV.21.6.5.1\endcsname{ega4.c04.eq.21-6-5-1}
\expandafter\def\csname sga@target@IV.21.6.6\endcsname{ega4.c04.prop.21-6-6}
\expandafter\def\csname sga@target@IV.21.6.6.2\endcsname{ega4.c04.statement.21-6-6-2}
\expandafter\def\csname sga@target@IV.21.6.7\endcsname{ega4.c04.statement.21-6-7}
\expandafter\def\csname sga@target@IV.21.6.7.1\endcsname{ega4.c04.eq.21-6-7-1}
\expandafter\def\csname sga@target@IV.21.6.7.2\endcsname{ega4.c04.eq.21-6-7-2}
\expandafter\def\csname sga@target@IV.21.6.7.3\endcsname{ega4.c04.eq.21-6-7-3}
\expandafter\def\csname sga@target@IV.21.6.7.4\endcsname{ega4.c04.eq.21-6-7-4}
\expandafter\def\csname sga@target@IV.21.6.7.5\endcsname{ega4.c04.eq.21-6-7-5}
\expandafter\def\csname sga@target@IV.21.6.8\endcsname{ega4.c04.statement.21-6-8}
\expandafter\def\csname sga@target@IV.21.6.9\endcsname{ega4.c04.thm.21-6-9}
\expandafter\def\csname sga@target@IV.21.6.9.1\endcsname{ega4.c04.eq.21-6-9-1}
\expandafter\def\csname sga@target@IV.21.6.9.2\endcsname{ega4.c04.statement.21-6-9-2}
\expandafter\def\csname sga@target@IV.21.7\endcsname{ega4.c04.subsec.21-7.key-iv-21-7}
\expandafter\def\csname sga@target@IV.21.7.1\endcsname{ega4.c04.statement.21-7-1}
\expandafter\def\csname sga@target@IV.21.7.2\endcsname{ega4.c04.prop.21-7-2}
\expandafter\def\csname sga@target@IV.21.7.2.1\endcsname{ega4.c04.eq.21-7-2-1}
\expandafter\def\csname sga@target@IV.21.7.3\endcsname{ega4.c04.cor.21-7-3}
\expandafter\def\csname sga@target@IV.21.7.3.1\endcsname{ega4.c04.statement.21-7-3-1}
\expandafter\def\csname sga@target@IV.21.7.4\endcsname{ega4.c04.prop.21-7-4}
\expandafter\def\csname sga@target@IV.21.7.5\endcsname{ega4.c04.rem.21-7-5}
\expandafter\def\csname sga@target@IV.21.7.6\endcsname{ega4.c04.cor.21-7-6}
\expandafter\def\csname sga@target@IV.21.7.7\endcsname{ega4.c04.rem.21-7-7}
\expandafter\def\csname sga@target@IV.21.8.1\endcsname{ega4.c04.lem.21-8-1}
\expandafter\def\csname sga@target@IV.21.8.2\endcsname{ega4.c04.cor.21-8-2}
\expandafter\def\csname sga@target@IV.21.8.3\endcsname{ega4.c04.prop.21-8-3}
\expandafter\def\csname sga@target@IV.21.8.3.1\endcsname{ega4.c04.eq.21-8-3-1}
\expandafter\def\csname sga@target@IV.21.8.4\endcsname{ega4.c04.cor.21-8-4}
\expandafter\def\csname sga@target@IV.21.8.5\endcsname{ega4.c04.prop.21-8-5}
\expandafter\def\csname sga@target@IV.21.8.5.1\endcsname{ega4.c04.eq.21-8-5-1}
\expandafter\def\csname sga@target@IV.21.8.5.2\endcsname{ega4.c04.eq.21-8-5-2}
\expandafter\def\csname sga@target@IV.21.8.5.3\endcsname{ega4.c04.eq.21-8-5-3}
\expandafter\def\csname sga@target@IV.21.8.6\endcsname{ega4.c04.rem.21-8-6}
\expandafter\def\csname sga@target@IV.21.9\endcsname{ega4.c04.subsec.21-9.key-iv-21-9}
\expandafter\def\csname sga@target@IV.21.9.1\endcsname{ega4.c04.statement.21-9-1}
\expandafter\def\csname sga@target@IV.21.9.1.1\endcsname{ega4.c04.eq.21-9-1-1}
\expandafter\def\csname sga@target@IV.21.9.1.2\endcsname{ega4.c04.eq.21-9-1-2}
\expandafter\def\csname sga@target@IV.21.9.10\endcsname{ega4.c04.statement.21-9-10}
\expandafter\def\csname sga@target@IV.21.9.11\endcsname{ega4.c04.prop.21-9-11}
\expandafter\def\csname sga@target@IV.21.9.12\endcsname{ega4.c04.cor.21-9-12}
\expandafter\def\csname sga@target@IV.21.9.13\endcsname{ega4.c04.rem.21-9-13}
\expandafter\def\csname sga@target@IV.21.9.2\endcsname{ega4.c04.lem.21-9-2}
\expandafter\def\csname sga@target@IV.21.9.3\endcsname{ega4.c04.rem.21-9-3}
\expandafter\def\csname sga@target@IV.21.9.4\endcsname{ega4.c04.prop.21-9-4}
\expandafter\def\csname sga@target@IV.21.9.4.1\endcsname{ega4.c04.prop.21-9-4-1}
\expandafter\def\csname sga@target@IV.21.9.5\endcsname{ega4.c04.cor.21-9-5}
\expandafter\def\csname sga@target@IV.21.9.6\endcsname{ega4.c04.cor.21-9-6}
\expandafter\def\csname sga@target@IV.21.9.7\endcsname{ega4.c04.cor.21-9-7}
\expandafter\def\csname sga@target@IV.21.9.8\endcsname{ega4.c04.statement.21-9-8}
\expandafter\def\csname sga@target@IV.21.9.9\endcsname{ega4.c04.prop.21-9-9}
\expandafter\def\csname sga@target@IV.3.1.1\endcsname{ega4.c04.def.3-1-1}
\expandafter\def\csname sga@target@IV.3.1.1.1\endcsname{ega4.c04.eq.3-1-1-1}
\expandafter\def\csname sga@target@IV.3.1.10\endcsname{ega4.c04.prop.3-1-10}
\expandafter\def\csname sga@target@IV.3.1.10.1\endcsname{ega4.c04.lem.3-1-10-1}
\expandafter\def\csname sga@target@IV.3.1.10.2\endcsname{ega4.c04.rem.3-1-10-2}
\expandafter\def\csname sga@target@IV.3.1.10.3\endcsname{ega4.c04.eq.3-1-10-3}
\expandafter\def\csname sga@target@IV.3.1.10.4\endcsname{ega4.c04.eq.3-1-10-4}
\expandafter\def\csname sga@target@IV.3.1.10.5\endcsname{ega4.c04.rem.3-1-10-5}
\expandafter\def\csname sga@target@IV.3.1.11\endcsname{ega4.c04.cor.3-1-11}
\expandafter\def\csname sga@target@IV.3.1.12\endcsname{ega4.c04.rem.3-1-12}
\expandafter\def\csname sga@target@IV.3.1.13\endcsname{ega4.c04.prop.3-1-13}
\expandafter\def\csname sga@target@IV.3.1.2\endcsname{ega4.c04.prop.3-1-2}
\expandafter\def\csname sga@target@IV.3.1.3\endcsname{ega4.c04.prop.3-1-3}
\expandafter\def\csname sga@target@IV.3.1.4\endcsname{ega4.c04.cor.3-1-4}
\expandafter\def\csname sga@target@IV.3.1.5\endcsname{ega4.c04.cor.3-1-5}
\expandafter\def\csname sga@target@IV.3.1.6\endcsname{ega4.c04.prop.3-1-6}
\expandafter\def\csname sga@target@IV.3.1.7\endcsname{ega4.c04.prop.3-1-7}
\expandafter\def\csname sga@target@IV.3.1.8\endcsname{ega4.c04.prop.3-1-8}
\expandafter\def\csname sga@target@IV.3.1.9\endcsname{ega4.c04.cor.3-1-9}
\expandafter\def\csname sga@target@IV.3.2.1\endcsname{ega4.c04.prop.3-2-1}
\expandafter\def\csname sga@target@IV.3.2.2\endcsname{ega4.c04.def.3-2-2}
\expandafter\def\csname sga@target@IV.3.2.3\endcsname{ega4.c04.prop.3-2-3}
\expandafter\def\csname sga@target@IV.3.2.4\endcsname{ega4.c04.def.3-2-4}
\expandafter\def\csname sga@target@IV.3.2.5\endcsname{ega4.c04.def.3-2-5}
\expandafter\def\csname sga@target@IV.3.2.6\endcsname{ega4.c04.prop.3-2-6}
\expandafter\def\csname sga@target@IV.3.2.7\endcsname{ega4.c04.cor.3-2-7}
\expandafter\def\csname sga@target@IV.3.2.8\endcsname{ega4.c04.cor.3-2-8}
\expandafter\def\csname sga@target@IV.3.2.9\endcsname{ega4.c04.rem.3-2-9}
\expandafter\def\csname sga@target@IV.3.3.1\endcsname{ega4.c04.prop.3-3-1}
\expandafter\def\csname sga@target@IV.3.3.1.1\endcsname{ega4.c04.prop.3-3-1-1}
\expandafter\def\csname sga@target@IV.3.3.2\endcsname{ega4.c04.cor.3-3-2}
\expandafter\def\csname sga@target@IV.3.3.3\endcsname{ega4.c04.cor.3-3-3}
\expandafter\def\csname sga@target@IV.3.3.4\endcsname{ega4.c04.prop.3-3-4}
\expandafter\def\csname sga@target@IV.3.3.5\endcsname{ega4.c04.cor.3-3-5}
\expandafter\def\csname sga@target@IV.3.3.6\endcsname{ega4.c04.prop.3-3-6}
\expandafter\def\csname sga@target@IV.3.3.6.1\endcsname{ega4.c04.eq.3-3-6-1}
\expandafter\def\csname sga@target@IV.3.3.6.2\endcsname{ega4.c04.eq.3-3-6-2}
\expandafter\def\csname sga@target@IV.3.3.6.3\endcsname{ega4.c04.eq.3-3-6-3}
\expandafter\def\csname sga@target@IV.3.3.6.4\endcsname{ega4.c04.eq.3-3-6-4}
\expandafter\def\csname sga@target@IV.3.3.7\endcsname{ega4.c04.cor.3-3-7}
\expandafter\def\csname sga@target@IV.3.3.7.1\endcsname{ega4.c04.cor.3-3-7-1}
\expandafter\def\csname sga@target@IV.3.3.8\endcsname{ega4.c04.sec.3-3-8}
\expandafter\def\csname sga@target@IV.3.3.9\endcsname{ega4.c04.cor.3-3-9}
\expandafter\def\csname sga@target@IV.3.3.9.1\endcsname{ega4.c04.cor.3-3-9-1}
\expandafter\def\csname sga@target@IV.3.4\endcsname{ega4.c04.subsec.3-4.key-iv-3-4}
\expandafter\def\csname sga@target@IV.3.4.1\endcsname{ega4.c04.prop.3-4-1}
\expandafter\def\csname sga@target@IV.3.4.1.1\endcsname{ega4.c04.eq.3-4-1-1}
\expandafter\def\csname sga@target@IV.3.4.1.2\endcsname{ega4.c04.lem.3-4-1-2}
\expandafter\def\csname sga@target@IV.3.4.1.3\endcsname{ega4.c04.lem.3-4-1-3}
\expandafter\def\csname sga@target@IV.3.4.1.4\endcsname{ega4.c04.lem.3-4-1-4}
\expandafter\def\csname sga@target@IV.3.4.1.5\endcsname{ega4.c04.eq.3-4-1-5}
\expandafter\def\csname sga@target@IV.3.4.2\endcsname{ega4.c04.cor.3-4-2}
\expandafter\def\csname sga@target@IV.3.4.3\endcsname{ega4.c04.prop.3-4-3}
\expandafter\def\csname sga@target@IV.3.4.4\endcsname{ega4.c04.prop.3-4-4}
\expandafter\def\csname sga@target@IV.3.4.5\endcsname{ega4.c04.prop.3-4-5}
\expandafter\def\csname sga@target@IV.3.4.6\endcsname{ega4.c04.prop.3-4-6}
\expandafter\def\csname sga@target@IV.3.4.6.1\endcsname{ega4.c04.lem.3-4-6-1}
\expandafter\def\csname sga@target@IV.3.4.7\endcsname{ega4.c04.cor.3-4-7}
\expandafter\def\csname sga@target@IV.3.4.8\endcsname{ega4.c04.rem.3-4-8}
\expandafter\def\csname sga@target@IV.3.4.9\endcsname{ega4.c04.prop.3-4-9}
\expandafter\def\csname sga@target@IV.4.1.1\endcsname{ega4.c04.def.4-1-1}
\expandafter\def\csname sga@target@IV.4.1.1.1\endcsname{ega4.c04.eq.4-1-1-1}
\expandafter\def\csname sga@target@IV.4.1.1.2\endcsname{ega4.c04.eq.4-1-1-2}
\expandafter\def\csname sga@target@IV.4.1.1.3\endcsname{ega4.c04.eq.4-1-1-3}
\expandafter\def\csname sga@target@IV.4.1.2\endcsname{ega4.c04.thm.4-1-2}
\expandafter\def\csname sga@target@IV.4.1.2.1\endcsname{ega4.c04.cor.4-1-2-1}
\expandafter\def\csname sga@target@IV.4.1.3\endcsname{ega4.c04.rem.4-1-3}
\expandafter\def\csname sga@target@IV.4.1.4\endcsname{ega4.c04.cor.4-1-4}
\expandafter\def\csname sga@target@IV.4.1.5\endcsname{ega4.c04.cor.4-1-5}
\expandafter\def\csname sga@target@IV.4.2.1\endcsname{ega4.c04.prop.4-2-1}
\expandafter\def\csname sga@target@IV.4.2.1.1\endcsname{ega4.c04.eq.4-2-1-1}
\expandafter\def\csname sga@target@IV.4.2.1.2\endcsname{ega4.c04.eq.4-2-1-2}
\expandafter\def\csname sga@target@IV.4.2.1.3\endcsname{ega4.c04.eq.4-2-1-3}
\expandafter\def\csname sga@target@IV.4.2.1.4\endcsname{ega4.c04.rem.4-2-1-4}
\expandafter\def\csname sga@target@IV.4.2.1.5\endcsname{ega4.c04.eq.4-2-1-5}
\expandafter\def\csname sga@target@IV.4.2.10\endcsname{ega4.c04.cor.4-2-10}
\expandafter\def\csname sga@target@IV.4.2.2\endcsname{ega4.c04.cor.4-2-2}
\expandafter\def\csname sga@target@IV.4.2.3\endcsname{ega4.c04.cor.4-2-3}
\expandafter\def\csname sga@target@IV.4.2.4\endcsname{ega4.c04.prop.4-2-4}
\expandafter\def\csname sga@target@IV.4.2.4.i\endcsname{ega4.c04.sec.4-2-4-i}
\expandafter\def\csname sga@target@IV.4.2.4.ii\endcsname{ega4.c04.sec.4-2-4-ii}
\expandafter\def\csname sga@target@IV.4.2.4.iii\endcsname{ega4.c04.sec.4-2-4-iii}
\expandafter\def\csname sga@target@IV.4.2.5\endcsname{ega4.c04.prop.4-2-5}
\expandafter\def\csname sga@target@IV.4.2.6\endcsname{ega4.c04.cor.4-2-6}
\expandafter\def\csname sga@target@IV.4.2.7\endcsname{ega4.c04.prop.4-2-7}
\expandafter\def\csname sga@target@IV.4.2.7.i\endcsname{ega4.c04.sec.4-2-7-i}
\expandafter\def\csname sga@target@IV.4.2.7.ii\endcsname{ega4.c04.sec.4-2-7-ii}
\expandafter\def\csname sga@target@IV.4.2.7.iii\endcsname{ega4.c04.sec.4-2-7-iii}
\expandafter\def\csname sga@target@IV.4.2.7.iv\endcsname{ega4.c04.sec.4-2-7-iv}
\expandafter\def\csname sga@target@IV.4.2.8\endcsname{ega4.c04.cor.4-2-8}
\expandafter\def\csname sga@target@IV.4.2.9\endcsname{ega4.c04.prop.4-2-9}
\expandafter\def\csname sga@target@IV.4.3.1\endcsname{ega4.c04.statement.4-3-1}
\expandafter\def\csname sga@target@IV.4.3.2\endcsname{ega4.c04.prop.4-3-2}
\expandafter\def\csname sga@target@IV.4.3.3\endcsname{ega4.c04.cor.4-3-3}
\expandafter\def\csname sga@target@IV.4.3.4\endcsname{ega4.c04.cor.4-3-4}
\expandafter\def\csname sga@target@IV.4.3.5\endcsname{ega4.c04.prop.4-3-5}
\expandafter\def\csname sga@target@IV.4.3.6\endcsname{ega4.c04.cor.4-3-6}
\expandafter\def\csname sga@target@IV.4.3.7\endcsname{ega4.c04.cor.4-3-7}
\expandafter\def\csname sga@target@IV.4.3.8\endcsname{ega4.c04.cor.4-3-8}
\expandafter\def\csname sga@target@IV.4.3.9\endcsname{ega4.c04.prop.4-3-9}
\expandafter\def\csname sga@target@IV.4.4.1\endcsname{ega4.c04.statement.4-4-1}
\expandafter\def\csname sga@target@IV.4.4.2\endcsname{ega4.c04.lem.4-4-2}
\expandafter\def\csname sga@target@IV.4.4.2.i\endcsname{ega4.c04.sec.4-4-2-i}
\expandafter\def\csname sga@target@IV.4.4.2.ii\endcsname{ega4.c04.sec.4-4-2-ii}
\expandafter\def\csname sga@target@IV.4.4.3\endcsname{ega4.c04.rem.4-4-3}
\expandafter\def\csname sga@target@IV.4.4.4\endcsname{ega4.c04.thm.4-4-4}
\expandafter\def\csname sga@target@IV.4.4.5\endcsname{ega4.c04.cor.4-4-5}
\expandafter\def\csname sga@target@IV.4.4.6\endcsname{ega4.c04.cor.4-4-6}
\expandafter\def\csname sga@target@IV.4.5\endcsname{ega4.c04.subsec.4-5.key-iv-4-5}
\expandafter\def\csname sga@target@IV.4.5.1\endcsname{ega4.c04.prop.4-5-1}
\expandafter\def\csname sga@target@IV.4.5.10\endcsname{ega4.c04.cor.4-5-10}
\expandafter\def\csname sga@target@IV.4.5.11\endcsname{ega4.c04.cor.4-5-11}
\expandafter\def\csname sga@target@IV.4.5.12\endcsname{ega4.c04.sec.4-5-12}
\expandafter\def\csname sga@target@IV.4.5.12.i\endcsname{ega4.c04.sec.4-5-12-i}
\expandafter\def\csname sga@target@IV.4.5.12.ii\endcsname{ega4.c04.sec.4-5-12-ii}
\expandafter\def\csname sga@target@IV.4.5.13\endcsname{ega4.c04.prop.4-5-13}
\expandafter\def\csname sga@target@IV.4.5.13.1\endcsname{ega4.c04.cor.4-5-13-1}
\expandafter\def\csname sga@target@IV.4.5.13.1.i\endcsname{ega4.c04.sec.4-5-13-1-i}
\expandafter\def\csname sga@target@IV.4.5.13.1.ii\endcsname{ega4.c04.sec.4-5-13-1-ii}
\expandafter\def\csname sga@target@IV.4.5.14\endcsname{ega4.c04.cor.4-5-14}
\expandafter\def\csname sga@target@IV.4.5.15\endcsname{ega4.c04.prop.4-5-15}
\expandafter\def\csname sga@target@IV.4.5.15.i\endcsname{ega4.c04.sec.4-5-15-i}
\expandafter\def\csname sga@target@IV.4.5.15.ii\endcsname{ega4.c04.sec.4-5-15-ii}
\expandafter\def\csname sga@target@IV.4.5.16\endcsname{ega4.c04.cor.4-5-16}
\expandafter\def\csname sga@target@IV.4.5.16.i\endcsname{ega4.c04.sec.4-5-16-i}
\expandafter\def\csname sga@target@IV.4.5.16.ii\endcsname{ega4.c04.sec.4-5-16-ii}
\expandafter\def\csname sga@target@IV.4.5.17\endcsname{ega4.c04.cor.4-5-17}
\expandafter\def\csname sga@target@IV.4.5.18\endcsname{ega4.c04.rem.4-5-18}
\expandafter\def\csname sga@target@IV.4.5.19\endcsname{ega4.c04.prop.4-5-19}
\expandafter\def\csname sga@target@IV.4.5.19.1\endcsname{ega4.c04.lem.4-5-19-1}
\expandafter\def\csname sga@target@IV.4.5.19.2\endcsname{ega4.c04.eq.4-5-19-2}
\expandafter\def\csname sga@target@IV.4.5.19.3\endcsname{ega4.c04.cor.4-5-19-3}
\expandafter\def\csname sga@target@IV.4.5.2\endcsname{ega4.c04.def.4-5-2}
\expandafter\def\csname sga@target@IV.4.5.2.a\endcsname{ega4.c04.sec.4-5-2-a}
\expandafter\def\csname sga@target@IV.4.5.2.b\endcsname{ega4.c04.sec.4-5-2-b}
\expandafter\def\csname sga@target@IV.4.5.2.c\endcsname{ega4.c04.sec.4-5-2-c}
\expandafter\def\csname sga@target@IV.4.5.20\endcsname{ega4.c04.rem.4-5-20}
\expandafter\def\csname sga@target@IV.4.5.21\endcsname{ega4.c04.prop.4-5-21}
\expandafter\def\csname sga@target@IV.4.5.3\endcsname{ega4.c04.statement.4-5-3}
\expandafter\def\csname sga@target@IV.4.5.4\endcsname{ega4.c04.prop.4-5-4}
\expandafter\def\csname sga@target@IV.4.5.5\endcsname{ega4.c04.def.4-5-5}
\expandafter\def\csname sga@target@IV.4.5.6\endcsname{ega4.c04.prop.4-5-6}
\expandafter\def\csname sga@target@IV.4.5.6.i\endcsname{ega4.c04.sec.4-5-6-i}
\expandafter\def\csname sga@target@IV.4.5.6.ii\endcsname{ega4.c04.sec.4-5-6-ii}
\expandafter\def\csname sga@target@IV.4.5.7\endcsname{ega4.c04.prop.4-5-7}
\expandafter\def\csname sga@target@IV.4.5.8\endcsname{ega4.c04.cor.4-5-8}
\expandafter\def\csname sga@target@IV.4.5.8.i\endcsname{ega4.c04.sec.4-5-8-i}
\expandafter\def\csname sga@target@IV.4.5.8.ii\endcsname{ega4.c04.sec.4-5-8-ii}
\expandafter\def\csname sga@target@IV.4.5.9\endcsname{ega4.c04.prop.4-5-9}
\expandafter\def\csname sga@target@IV.4.5.9.a\endcsname{ega4.c04.sec.4-5-9-a}
\expandafter\def\csname sga@target@IV.4.5.9.b\endcsname{ega4.c04.sec.4-5-9-b}
\expandafter\def\csname sga@target@IV.4.5.9.c\endcsname{ega4.c04.sec.4-5-9-c}
\expandafter\def\csname sga@target@IV.4.6\endcsname{ega4.c04.subsec.4-6.key-iv-4-6}
\expandafter\def\csname sga@target@IV.4.6.1\endcsname{ega4.c04.prop.4-6-1}
\expandafter\def\csname sga@target@IV.4.6.1.a\endcsname{ega4.c04.sec.4-6-1-a}
\expandafter\def\csname sga@target@IV.4.6.1.b\endcsname{ega4.c04.sec.4-6-1-b}
\expandafter\def\csname sga@target@IV.4.6.1.c\endcsname{ega4.c04.sec.4-6-1-c}
\expandafter\def\csname sga@target@IV.4.6.1.d\endcsname{ega4.c04.sec.4-6-1-d}
\expandafter\def\csname sga@target@IV.4.6.1.e\endcsname{ega4.c04.sec.4-6-1-e}
\expandafter\def\csname sga@target@IV.4.6.10\endcsname{ega4.c04.prop.4-6-10}
\expandafter\def\csname sga@target@IV.4.6.11\endcsname{ega4.c04.cor.4-6-11}
\expandafter\def\csname sga@target@IV.4.6.12\endcsname{ega4.c04.prop.4-6-12}
\expandafter\def\csname sga@target@IV.4.6.12.a\endcsname{ega4.c04.sec.4-6-12-a}
\expandafter\def\csname sga@target@IV.4.6.12.b\endcsname{ega4.c04.sec.4-6-12-b}
\expandafter\def\csname sga@target@IV.4.6.12.c\endcsname{ega4.c04.sec.4-6-12-c}
\expandafter\def\csname sga@target@IV.4.6.12.d\endcsname{ega4.c04.sec.4-6-12-d}
\expandafter\def\csname sga@target@IV.4.6.12.e\endcsname{ega4.c04.sec.4-6-12-e}
\expandafter\def\csname sga@target@IV.4.6.13\endcsname{ega4.c04.cor.4-6-13}
\expandafter\def\csname sga@target@IV.4.6.13.f\endcsname{ega4.c04.sec.4-6-13-f}
\expandafter\def\csname sga@target@IV.4.6.14\endcsname{ega4.c04.statement.4-6-14}
\expandafter\def\csname sga@target@IV.4.6.15\endcsname{ega4.c04.prop.4-6-15}
\expandafter\def\csname sga@target@IV.4.6.15.i\endcsname{ega4.c04.sec.4-6-15-i}
\expandafter\def\csname sga@target@IV.4.6.15.ii\endcsname{ega4.c04.sec.4-6-15-ii}
\expandafter\def\csname sga@target@IV.4.6.16\endcsname{ega4.c04.prop.4-6-16}
\expandafter\def\csname sga@target@IV.4.6.16.a\endcsname{ega4.c04.sec.4-6-16-a}
\expandafter\def\csname sga@target@IV.4.6.16.b\endcsname{ega4.c04.sec.4-6-16-b}
\expandafter\def\csname sga@target@IV.4.6.16.c\endcsname{ega4.c04.sec.4-6-16-c}
\expandafter\def\csname sga@target@IV.4.6.16.d\endcsname{ega4.c04.sec.4-6-16-d}
\expandafter\def\csname sga@target@IV.4.6.17\endcsname{ega4.c04.def.4-6-17}
\expandafter\def\csname sga@target@IV.4.6.18\endcsname{ega4.c04.prop.4-6-18}
\expandafter\def\csname sga@target@IV.4.6.18.a\endcsname{ega4.c04.sec.4-6-18-a}
\expandafter\def\csname sga@target@IV.4.6.18.b\endcsname{ega4.c04.sec.4-6-18-b}
\expandafter\def\csname sga@target@IV.4.6.18.c\endcsname{ega4.c04.sec.4-6-18-c}
\expandafter\def\csname sga@target@IV.4.6.18.d\endcsname{ega4.c04.sec.4-6-18-d}
\expandafter\def\csname sga@target@IV.4.6.19\endcsname{ega4.c04.def.4-6-19}
\expandafter\def\csname sga@target@IV.4.6.2\endcsname{ega4.c04.def.4-6-2}
\expandafter\def\csname sga@target@IV.4.6.20\endcsname{ega4.c04.prop.4-6-20}
\expandafter\def\csname sga@target@IV.4.6.21\endcsname{ega4.c04.prop.4-6-21}
\expandafter\def\csname sga@target@IV.4.6.21.i\endcsname{ega4.c04.sec.4-6-21-i}
\expandafter\def\csname sga@target@IV.4.6.21.ii\endcsname{ega4.c04.sec.4-6-21-ii}
\expandafter\def\csname sga@target@IV.4.6.22\endcsname{ega4.c04.def.4-6-22}
\expandafter\def\csname sga@target@IV.4.6.3\endcsname{ega4.c04.cor.4-6-3}
\expandafter\def\csname sga@target@IV.4.6.4\endcsname{ega4.c04.cor.4-6-4}
\expandafter\def\csname sga@target@IV.4.6.5\endcsname{ega4.c04.prop.4-6-5}
\expandafter\def\csname sga@target@IV.4.6.5.i\endcsname{ega4.c04.sec.4-6-5-i}
\expandafter\def\csname sga@target@IV.4.6.5.ii\endcsname{ega4.c04.sec.4-6-5-ii}
\expandafter\def\csname sga@target@IV.4.6.6\endcsname{ega4.c04.prop.4-6-6}
\expandafter\def\csname sga@target@IV.4.6.7\endcsname{ega4.c04.cor.4-6-7}
\expandafter\def\csname sga@target@IV.4.6.8\endcsname{ega4.c04.cor.4-6-8}
\expandafter\def\csname sga@target@IV.4.6.9\endcsname{ega4.c04.def.4-6-9}
\expandafter\def\csname sga@target@IV.4.7.1\endcsname{ega4.c04.lem.4-7-1}
\expandafter\def\csname sga@target@IV.4.7.1.1\endcsname{ega4.c04.eq.4-7-1-1}
\expandafter\def\csname sga@target@IV.4.7.10\endcsname{ega4.c04.prop.4-7-10}
\expandafter\def\csname sga@target@IV.4.7.11\endcsname{ega4.c04.prop.4-7-11}
\expandafter\def\csname sga@target@IV.4.7.12\endcsname{ega4.c04.def.4-7-12}
\expandafter\def\csname sga@target@IV.4.7.13\endcsname{ega4.c04.prop.4-7-13}
\expandafter\def\csname sga@target@IV.4.7.2\endcsname{ega4.c04.cor.4-7-2}
\expandafter\def\csname sga@target@IV.4.7.2.1\endcsname{ega4.c04.eq.4-7-2-1}
\expandafter\def\csname sga@target@IV.4.7.3\endcsname{ega4.c04.prop.4-7-3}
\expandafter\def\csname sga@target@IV.4.7.3.1\endcsname{ega4.c04.eq.4-7-3-1}
\expandafter\def\csname sga@target@IV.4.7.3.i\endcsname{ega4.c04.sec.4-7-3-i}
\expandafter\def\csname sga@target@IV.4.7.3.ii\endcsname{ega4.c04.sec.4-7-3-ii}
\expandafter\def\csname sga@target@IV.4.7.3.iii\endcsname{ega4.c04.sec.4-7-3-iii}
\expandafter\def\csname sga@target@IV.4.7.4\endcsname{ega4.c04.def.4-7-4}
\expandafter\def\csname sga@target@IV.4.7.5\endcsname{ega4.c04.def.4-7-5}
\expandafter\def\csname sga@target@IV.4.7.6\endcsname{ega4.c04.prop.4-7-6}
\expandafter\def\csname sga@target@IV.4.7.6.1\endcsname{ega4.c04.eq.4-7-6-1}
\expandafter\def\csname sga@target@IV.4.7.7\endcsname{ega4.c04.statement.4-7-7}
\expandafter\def\csname sga@target@IV.4.7.8\endcsname{ega4.c04.prop.4-7-8}
\expandafter\def\csname sga@target@IV.4.7.9\endcsname{ega4.c04.cor.4-7-9}
\expandafter\def\csname sga@target@IV.4.8.1\endcsname{ega4.c04.statement.4-8-1}
\expandafter\def\csname sga@target@IV.4.8.10\endcsname{ega4.c04.cor.4-8-10}
\expandafter\def\csname sga@target@IV.4.8.11\endcsname{ega4.c04.cor.4-8-11}
\expandafter\def\csname sga@target@IV.4.8.12\endcsname{ega4.c04.cor.4-8-12}
\expandafter\def\csname sga@target@IV.4.8.13\endcsname{ega4.c04.prop.4-8-13}
\expandafter\def\csname sga@target@IV.4.8.14\endcsname{ega4.c04.prop.4-8-14}
\expandafter\def\csname sga@target@IV.4.8.15\endcsname{ega4.c04.sec.4-8-15}
\expandafter\def\csname sga@target@IV.4.8.15.i\endcsname{ega4.c04.sec.4-8-15-i}
\expandafter\def\csname sga@target@IV.4.8.15.ii\endcsname{ega4.c04.sec.4-8-15-ii}
\expandafter\def\csname sga@target@IV.4.8.2\endcsname{ega4.c04.statement.4-8-2}
\expandafter\def\csname sga@target@IV.4.8.2.1\endcsname{ega4.c04.eq.4-8-2-1}
\expandafter\def\csname sga@target@IV.4.8.2.2\endcsname{ega4.c04.eq.4-8-2-2}
\expandafter\def\csname sga@target@IV.4.8.2.3\endcsname{ega4.c04.eq.4-8-2-3}
\expandafter\def\csname sga@target@IV.4.8.2.4\endcsname{ega4.c04.eq.4-8-2-4}
\expandafter\def\csname sga@target@IV.4.8.2.5\endcsname{ega4.c04.eq.4-8-2-5}
\expandafter\def\csname sga@target@IV.4.8.2.6\endcsname{ega4.c04.eq.4-8-2-6}
\expandafter\def\csname sga@target@IV.4.8.3\endcsname{ega4.c04.prop.4-8-3}
\expandafter\def\csname sga@target@IV.4.8.4\endcsname{ega4.c04.def.4-8-4}
\expandafter\def\csname sga@target@IV.4.8.5\endcsname{ega4.c04.prop.4-8-5}
\expandafter\def\csname sga@target@IV.4.8.6\endcsname{ega4.c04.lem.4-8-6}
\expandafter\def\csname sga@target@IV.4.8.6.a\endcsname{ega4.c04.sec.4-8-6-a}
\expandafter\def\csname sga@target@IV.4.8.6.b\endcsname{ega4.c04.sec.4-8-6-b}
\expandafter\def\csname sga@target@IV.4.8.6.c\endcsname{ega4.c04.sec.4-8-6-c}
\expandafter\def\csname sga@target@IV.4.8.7\endcsname{ega4.c04.cor.4-8-7}
\expandafter\def\csname sga@target@IV.4.8.8\endcsname{ega4.c04.lem.4-8-8}
\expandafter\def\csname sga@target@IV.4.8.8.i\endcsname{ega4.c04.sec.4-8-8-i}
\expandafter\def\csname sga@target@IV.4.8.8.ii\endcsname{ega4.c04.sec.4-8-8-ii}
\expandafter\def\csname sga@target@IV.4.8.8.iii\endcsname{ega4.c04.sec.4-8-8-iii}
\expandafter\def\csname sga@target@IV.4.8.9\endcsname{ega4.c04.prop.4-8-9}
\expandafter\def\csname sga@target@IV.4.9.1\endcsname{ega4.c04.prop.4-9-1}
\expandafter\def\csname sga@target@IV.4.9.1.a\endcsname{ega4.c04.sec.4-9-1-a}
\expandafter\def\csname sga@target@IV.4.9.1.b\endcsname{ega4.c04.sec.4-9-1-b}
\expandafter\def\csname sga@target@IV.4.9.1.bprime\endcsname{ega4.c04.sec.4-9-1-bprime}
\expandafter\def\csname sga@target@IV.4.9.1.c\endcsname{ega4.c04.sec.4-9-1-c}
\expandafter\def\csname sga@target@IV.4.9.2\endcsname{ega4.c04.cor.4-9-2}
\expandafter\def\csname sga@target@IV.4.9.3\endcsname{ega4.c04.rem.4-9-3}
\expandafter\def\csname sga@target@IV.4.9.4\endcsname{ega4.c04.prop.4-9-4}
\expandafter\def\csname sga@target@IV.4.9.5\endcsname{ega4.c04.cor.4-9-5}
\expandafter\def\csname sga@target@IV.4.9.5.a\endcsname{ega4.c04.sec.4-9-5-a}
\expandafter\def\csname sga@target@IV.4.9.5.b\endcsname{ega4.c04.sec.4-9-5-b}
\expandafter\def\csname sga@target@IV.4.9.6\endcsname{ega4.c04.prop.4-9-6}
\expandafter\def\csname sga@target@IV.4.9.6.a\endcsname{ega4.c04.sec.4-9-6-a}
\expandafter\def\csname sga@target@IV.4.9.6.aprime\endcsname{ega4.c04.sec.4-9-6-aprime}
\expandafter\def\csname sga@target@IV.4.9.6.b\endcsname{ega4.c04.sec.4-9-6-b}
\expandafter\def\csname sga@target@IV.4.9.6.bprime\endcsname{ega4.c04.sec.4-9-6-bprime}
\expandafter\def\csname sga@target@IV.4.9.7\endcsname{ega4.c04.cor.4-9-7}
\expandafter\def\csname sga@target@IV.5.1.1\endcsname{ega4.c04.statement.5-1-1}
\expandafter\def\csname sga@target@IV.5.1.1.1\endcsname{ega4.c04.eq.5-1-1-1}
\expandafter\def\csname sga@target@IV.5.1.1.2\endcsname{ega4.c04.eq.5-1-1-2}
\expandafter\def\csname sga@target@IV.5.1.1.3\endcsname{ega4.c04.eq.5-1-1-3}
\expandafter\def\csname sga@target@IV.5.1.1.4\endcsname{ega4.c04.sec.5-1-1-4}
\expandafter\def\csname sga@target@IV.5.1.1.5\endcsname{ega4.c04.sec.5-1-1-5}
\expandafter\def\csname sga@target@IV.5.1.1.6\endcsname{ega4.c04.sec.5-1-1-6}
\expandafter\def\csname sga@target@IV.5.1.10\endcsname{ega4.c04.prop.5-1-10}
\expandafter\def\csname sga@target@IV.5.1.10.1\endcsname{ega4.c04.lem.5-1-10-1}
\expandafter\def\csname sga@target@IV.5.1.10.i\endcsname{ega4.c04.sec.5-1-10-i}
\expandafter\def\csname sga@target@IV.5.1.10.ii\endcsname{ega4.c04.sec.5-1-10-ii}
\expandafter\def\csname sga@target@IV.5.1.11\endcsname{ega4.c04.cor.5-1-11}
\expandafter\def\csname sga@target@IV.5.1.12\endcsname{ega4.c04.statement.5-1-12}
\expandafter\def\csname sga@target@IV.5.1.12.1\endcsname{ega4.c04.eq.5-1-12-1}
\expandafter\def\csname sga@target@IV.5.1.12.2\endcsname{ega4.c04.eq.5-1-12-2}
\expandafter\def\csname sga@target@IV.5.1.12.3\endcsname{ega4.c04.eq.5-1-12-3}
\expandafter\def\csname sga@target@IV.5.1.13\endcsname{ega4.c04.prop.5-1-13}
\expandafter\def\csname sga@target@IV.5.1.13.1\endcsname{ega4.c04.eq.5-1-13-1}
\expandafter\def\csname sga@target@IV.5.1.2\endcsname{ega4.c04.prop.5-1-2}
\expandafter\def\csname sga@target@IV.5.1.2.1\endcsname{ega4.c04.eq.5-1-2-1}
\expandafter\def\csname sga@target@IV.5.1.3\endcsname{ega4.c04.cor.5-1-3}
\expandafter\def\csname sga@target@IV.5.1.3.1\endcsname{ega4.c04.eq.5-1-3-1}
\expandafter\def\csname sga@target@IV.5.1.3.2\endcsname{ega4.c04.eq.5-1-3-2}
\expandafter\def\csname sga@target@IV.5.1.4\endcsname{ega4.c04.prop.5-1-4}
\expandafter\def\csname sga@target@IV.5.1.4.1\endcsname{ega4.c04.eq.5-1-4-1}
\expandafter\def\csname sga@target@IV.5.1.4.2\endcsname{ega4.c04.eq.5-1-4-2}
\expandafter\def\csname sga@target@IV.5.1.5\endcsname{ega4.c04.cor.5-1-5}
\expandafter\def\csname sga@target@IV.5.1.6\endcsname{ega4.c04.statement.5-1-6}
\expandafter\def\csname sga@target@IV.5.1.7\endcsname{ega4.c04.prop.5-1-7}
\expandafter\def\csname sga@target@IV.5.1.8\endcsname{ega4.c04.cor.5-1-8}
\expandafter\def\csname sga@target@IV.5.1.9\endcsname{ega4.c04.prop.5-1-9}
\expandafter\def\csname sga@target@IV.5.1.9.1\endcsname{ega4.c04.statement.5-1-9-1}
\expandafter\def\csname sga@target@IV.5.10.1\endcsname{ega4.c04.statement.5-10-1}
\expandafter\def\csname sga@target@IV.5.10.1.1\endcsname{ega4.c04.eq.5-10-1-1}
\expandafter\def\csname sga@target@IV.5.10.10\endcsname{ega4.c04.prop.5-10-10}
\expandafter\def\csname sga@target@IV.5.10.11\endcsname{ega4.c04.statement.5-10-11}
\expandafter\def\csname sga@target@IV.5.10.11.1\endcsname{ega4.c04.eq.5-10-11-1}
\expandafter\def\csname sga@target@IV.5.10.12\endcsname{ega4.c04.prop.5-10-12}
\expandafter\def\csname sga@target@IV.5.10.13\endcsname{ega4.c04.statement.5-10-13}
\expandafter\def\csname sga@target@IV.5.10.14\endcsname{ega4.c04.prop.5-10-14}
\expandafter\def\csname sga@target@IV.5.10.15\endcsname{ega4.c04.cor.5-10-15}
\expandafter\def\csname sga@target@IV.5.10.16\endcsname{ega4.c04.prop.5-10-16}
\expandafter\def\csname sga@target@IV.5.10.17\endcsname{ega4.c04.prop.5-10-17}
\expandafter\def\csname sga@target@IV.5.10.2\endcsname{ega4.c04.prop.5-10-2}
\expandafter\def\csname sga@target@IV.5.10.3\endcsname{ega4.c04.cor.5-10-3}
\expandafter\def\csname sga@target@IV.5.10.4\endcsname{ega4.c04.cor.5-10-4}
\expandafter\def\csname sga@target@IV.5.10.5\endcsname{ega4.c04.thm.5-10-5}
\expandafter\def\csname sga@target@IV.5.10.6\endcsname{ega4.c04.cor.5-10-6}
\expandafter\def\csname sga@target@IV.5.10.7\endcsname{ega4.c04.thm.5-10-7}
\expandafter\def\csname sga@target@IV.5.10.8\endcsname{ega4.c04.cor.5-10-8}
\expandafter\def\csname sga@target@IV.5.10.8.1\endcsname{ega4.c04.lem.5-10-8-1}
\expandafter\def\csname sga@target@IV.5.10.9\endcsname{ega4.c04.cor.5-10-9}
\expandafter\def\csname sga@target@IV.5.11.1\endcsname{ega4.c04.prop.5-11-1}
\expandafter\def\csname sga@target@IV.5.11.1.1\endcsname{ega4.c04.lem.5-11-1-1}
\expandafter\def\csname sga@target@IV.5.11.2\endcsname{ega4.c04.cor.5-11-2}
\expandafter\def\csname sga@target@IV.5.11.3\endcsname{ega4.c04.rem.5-11-3}
\expandafter\def\csname sga@target@IV.5.11.4\endcsname{ega4.c04.cor.5-11-4}
\expandafter\def\csname sga@target@IV.5.11.5\endcsname{ega4.c04.cor.5-11-5}
\expandafter\def\csname sga@target@IV.5.11.6\endcsname{ega4.c04.cor.5-11-6}
\expandafter\def\csname sga@target@IV.5.11.7\endcsname{ega4.c04.item.5-11-7}
\expandafter\def\csname sga@target@IV.5.12\endcsname{ega4.c04.subsec.5-12.key-iv-5-12}
\expandafter\def\csname sga@target@IV.5.12.1\endcsname{ega4.c04.prop.5-12-1}
\expandafter\def\csname sga@target@IV.5.12.1.1\endcsname{ega4.c04.eq.5-12-1-1}
\expandafter\def\csname sga@target@IV.5.12.10\endcsname{ega4.c04.cor.5-12-10}
\expandafter\def\csname sga@target@IV.5.12.2\endcsname{ega4.c04.prop.5-12-2}
\expandafter\def\csname sga@target@IV.5.12.3\endcsname{ega4.c04.rem.5-12-3}
\expandafter\def\csname sga@target@IV.5.12.4\endcsname{ega4.c04.cor.5-12-4}
\expandafter\def\csname sga@target@IV.5.12.5\endcsname{ega4.c04.prop.5-12-5}
\expandafter\def\csname sga@target@IV.5.12.6\endcsname{ega4.c04.rem.5-12-6}
\expandafter\def\csname sga@target@IV.5.12.7\endcsname{ega4.c04.cor.5-12-7}
\expandafter\def\csname sga@target@IV.5.12.8\endcsname{ega4.c04.prop.5-12-8}
\expandafter\def\csname sga@target@IV.5.12.9\endcsname{ega4.c04.rem.5-12-9}
\expandafter\def\csname sga@target@IV.5.13.1\endcsname{ega4.c04.prop.5-13-1}
\expandafter\def\csname sga@target@IV.5.13.2\endcsname{ega4.c04.cor.5-13-2}
\expandafter\def\csname sga@target@IV.5.13.3\endcsname{ega4.c04.prop.5-13-3}
\expandafter\def\csname sga@target@IV.5.13.4\endcsname{ega4.c04.cor.5-13-4}
\expandafter\def\csname sga@target@IV.5.13.5\endcsname{ega4.c04.statement.5-13-5}
\expandafter\def\csname sga@target@IV.5.13.6\endcsname{ega4.c04.cor.5-13-6}
\expandafter\def\csname sga@target@IV.5.13.7\endcsname{ega4.c04.prop.5-13-7}
\expandafter\def\csname sga@target@IV.5.13.7.1\endcsname{ega4.c04.lem.5-13-7-1}
\expandafter\def\csname sga@target@IV.5.13.8\endcsname{ega4.c04.cor.5-13-8}
\expandafter\def\csname sga@target@IV.5.2.1\endcsname{ega4.c04.prop.5-2-1}
\expandafter\def\csname sga@target@IV.5.2.1.1\endcsname{ega4.c04.eq.5-2-1-1}
\expandafter\def\csname sga@target@IV.5.2.2\endcsname{ega4.c04.cor.5-2-2}
\expandafter\def\csname sga@target@IV.5.2.3\endcsname{ega4.c04.cor.5-2-3}
\expandafter\def\csname sga@target@IV.5.2.3.1\endcsname{ega4.c04.eq.5-2-3-1}
\expandafter\def\csname sga@target@IV.5.2.4\endcsname{ega4.c04.cor.5-2-4}
\expandafter\def\csname sga@target@IV.5.2.5\endcsname{ega4.c04.sec.5-2-5}
\expandafter\def\csname sga@target@IV.5.2.5.i\endcsname{ega4.c04.sec.5-2-5-i}
\expandafter\def\csname sga@target@IV.5.2.5.ii\endcsname{ega4.c04.sec.5-2-5-ii}
\expandafter\def\csname sga@target@IV.5.3.1\endcsname{ega4.c04.prop.5-3-1}
\expandafter\def\csname sga@target@IV.5.3.1.1\endcsname{ega4.c04.eq.5-3-1-1}
\expandafter\def\csname sga@target@IV.5.3.1.2\endcsname{ega4.c04.eq.5-3-1-2}
\expandafter\def\csname sga@target@IV.5.3.1.3\endcsname{ega4.c04.lem.5-3-1-3}
\expandafter\def\csname sga@target@IV.5.4.1\endcsname{ega4.c04.prop.5-4-1}
\expandafter\def\csname sga@target@IV.5.4.1.i\endcsname{ega4.c04.sec.5-4-1-i}
\expandafter\def\csname sga@target@IV.5.4.1.ii\endcsname{ega4.c04.sec.5-4-1-ii}
\expandafter\def\csname sga@target@IV.5.4.2\endcsname{ega4.c04.cor.5-4-2}
\expandafter\def\csname sga@target@IV.5.4.3\endcsname{ega4.c04.sec.5-4-3}
\expandafter\def\csname sga@target@IV.5.4.3.1\endcsname{ega4.c04.eq.5-4-3-1}
\expandafter\def\csname sga@target@IV.5.4.3.i\endcsname{ega4.c04.sec.5-4-3-i}
\expandafter\def\csname sga@target@IV.5.4.3.ii\endcsname{ega4.c04.sec.5-4-3-ii}
\expandafter\def\csname sga@target@IV.5.5\endcsname{ega4.c04.sec.5-5}
\expandafter\def\csname sga@target@IV.5.5.1\endcsname{ega4.c04.statement.5-5-1}
\expandafter\def\csname sga@target@IV.5.5.1.1\endcsname{ega4.c04.eq.5-5-1-1}
\expandafter\def\csname sga@target@IV.5.5.2\endcsname{ega4.c04.prop.5-5-2}
\expandafter\def\csname sga@target@IV.5.5.2.1\endcsname{ega4.c04.eq.5-5-2-1}
\expandafter\def\csname sga@target@IV.5.5.2.2\endcsname{ega4.c04.eq.5-5-2-2}
\expandafter\def\csname sga@target@IV.5.5.3\endcsname{ega4.c04.prop.5-5-3}
\expandafter\def\csname sga@target@IV.5.5.3.1\endcsname{ega4.c04.eq.5-5-3-1}
\expandafter\def\csname sga@target@IV.5.5.4\endcsname{ega4.c04.cor.5-5-4}
\expandafter\def\csname sga@target@IV.5.5.4.1\endcsname{ega4.c04.eq.5-5-4-1}
\expandafter\def\csname sga@target@IV.5.5.5\endcsname{ega4.c04.cor.5-5-5}
\expandafter\def\csname sga@target@IV.5.5.6\endcsname{ega4.c04.cor.5-5-6}
\expandafter\def\csname sga@target@IV.5.5.6.1\endcsname{ega4.c04.eq.5-5-6-1}
\expandafter\def\csname sga@target@IV.5.5.7\endcsname{ega4.c04.lem.5-5-7}
\expandafter\def\csname sga@target@IV.5.5.7.1\endcsname{ega4.c04.eq.5-5-7-1}
\expandafter\def\csname sga@target@IV.5.5.7.2\endcsname{ega4.c04.statement.5-5-7-2}
\expandafter\def\csname sga@target@IV.5.5.7.3\endcsname{ega4.c04.eq.5-5-7-3}
\expandafter\def\csname sga@target@IV.5.5.7.i\endcsname{ega4.c04.sec.5-5-7-i}
\expandafter\def\csname sga@target@IV.5.5.7.ii\endcsname{ega4.c04.sec.5-5-7-ii}
\expandafter\def\csname sga@target@IV.5.5.8\endcsname{ega4.c04.thm.5-5-8}
\expandafter\def\csname sga@target@IV.5.5.8.1\endcsname{ega4.c04.eq.5-5-8-1}
\expandafter\def\csname sga@target@IV.5.5.8.2\endcsname{ega4.c04.eq.5-5-8-2}
\expandafter\def\csname sga@target@IV.5.5.8.3\endcsname{ega4.c04.eq.5-5-8-3}
\expandafter\def\csname sga@target@IV.5.6\endcsname{ega4.c04.sec.5-6}
\expandafter\def\csname sga@target@IV.5.6.1\endcsname{ega4.c04.prop.5-6-1}
\expandafter\def\csname sga@target@IV.5.6.1.1\endcsname{ega4.c04.eq.5-6-1-1}
\expandafter\def\csname sga@target@IV.5.6.1.2\endcsname{ega4.c04.eq.5-6-1-2}
\expandafter\def\csname sga@target@IV.5.6.1.3\endcsname{ega4.c04.statement.5-6-1-3}
\expandafter\def\csname sga@target@IV.5.6.1.4\endcsname{ega4.c04.statement.5-6-1-4}
\expandafter\def\csname sga@target@IV.5.6.1.5\endcsname{ega4.c04.statement.5-6-1-5}
\expandafter\def\csname sga@target@IV.5.6.1.a\endcsname{ega4.c04.sec.5-6-1-a}
\expandafter\def\csname sga@target@IV.5.6.1.b\endcsname{ega4.c04.sec.5-6-1-b}
\expandafter\def\csname sga@target@IV.5.6.1.c\endcsname{ega4.c04.sec.5-6-1-c}
\expandafter\def\csname sga@target@IV.5.6.10\endcsname{ega4.c04.prop.5-6-10}
\expandafter\def\csname sga@target@IV.5.6.11\endcsname{ega4.c04.ex.5-6-11}
\expandafter\def\csname sga@target@IV.5.6.11.1\endcsname{ega4.c04.lem.5-6-11-1}
\expandafter\def\csname sga@target@IV.5.6.11.1.i\endcsname{ega4.c04.sec.5-6-11-1-i}
\expandafter\def\csname sga@target@IV.5.6.11.1.ii\endcsname{ega4.c04.sec.5-6-11-1-ii}
\expandafter\def\csname sga@target@IV.5.6.2\endcsname{ega4.c04.def.5-6-2}
\expandafter\def\csname sga@target@IV.5.6.3\endcsname{ega4.c04.sec.5-6-3}
\expandafter\def\csname sga@target@IV.5.6.3.i\endcsname{ega4.c04.sec.5-6-3-i}
\expandafter\def\csname sga@target@IV.5.6.3.ii\endcsname{ega4.c04.sec.5-6-3-ii}
\expandafter\def\csname sga@target@IV.5.6.3.iii\endcsname{ega4.c04.sec.5-6-3-iii}
\expandafter\def\csname sga@target@IV.5.6.3.iv\endcsname{ega4.c04.sec.5-6-3-iv}
\expandafter\def\csname sga@target@IV.5.6.4\endcsname{ega4.c04.prop.5-6-4}
\expandafter\def\csname sga@target@IV.5.6.5\endcsname{ega4.c04.prop.5-6-5}
\expandafter\def\csname sga@target@IV.5.6.5.1\endcsname{ega4.c04.eq.5-6-5-1}
\expandafter\def\csname sga@target@IV.5.6.5.2\endcsname{ega4.c04.eq.5-6-5-2}
\expandafter\def\csname sga@target@IV.5.6.5.3\endcsname{ega4.c04.eq.5-6-5-3}
\expandafter\def\csname sga@target@IV.5.6.6\endcsname{ega4.c04.cor.5-6-6}
\expandafter\def\csname sga@target@IV.5.6.6.1\endcsname{ega4.c04.eq.5-6-6-1}
\expandafter\def\csname sga@target@IV.5.6.6.2\endcsname{ega4.c04.eq.5-6-6-2}
\expandafter\def\csname sga@target@IV.5.6.7\endcsname{ega4.c04.cor.5-6-7}
\expandafter\def\csname sga@target@IV.5.6.7.1\endcsname{ega4.c04.eq.5-6-7-1}
\expandafter\def\csname sga@target@IV.5.6.8\endcsname{ega4.c04.cor.5-6-8}
\expandafter\def\csname sga@target@IV.5.6.8.1\endcsname{ega4.c04.eq.5-6-8-1}
\expandafter\def\csname sga@target@IV.5.6.8.2\endcsname{ega4.c04.eq.5-6-8-2}
\expandafter\def\csname sga@target@IV.5.6.9\endcsname{ega4.c04.sec.5-6-9}
\expandafter\def\csname sga@target@IV.5.6.9.i\endcsname{ega4.c04.sec.5-6-9-i}
\expandafter\def\csname sga@target@IV.5.6.9.ii\endcsname{ega4.c04.sec.5-6-9-ii}
\expandafter\def\csname sga@target@IV.5.7.1\endcsname{ega4.c04.def.5-7-1}
\expandafter\def\csname sga@target@IV.5.7.1.1\endcsname{ega4.c04.eq.5-7-1-1}
\expandafter\def\csname sga@target@IV.5.7.10\endcsname{ega4.c04.cor.5-7-10}
\expandafter\def\csname sga@target@IV.5.7.11\endcsname{ega4.c04.cor.5-7-11}
\expandafter\def\csname sga@target@IV.5.7.11.i\endcsname{ega4.c04.sec.5-7-11-i}
\expandafter\def\csname sga@target@IV.5.7.11.ii\endcsname{ega4.c04.sec.5-7-11-ii}
\expandafter\def\csname sga@target@IV.5.7.12\endcsname{ega4.c04.statement.5-7-12}
\expandafter\def\csname sga@target@IV.5.7.13\endcsname{ega4.c04.cor.5-7-13}
\expandafter\def\csname sga@target@IV.5.7.13.1\endcsname{ega4.c04.eq.5-7-13-1}
\expandafter\def\csname sga@target@IV.5.7.2\endcsname{ega4.c04.def.5-7-2}
\expandafter\def\csname sga@target@IV.5.7.2.1\endcsname{ega4.c04.eq.5-7-2-1}
\expandafter\def\csname sga@target@IV.5.7.2.2\endcsname{ega4.c04.eq.5-7-2-2}
\expandafter\def\csname sga@target@IV.5.7.3\endcsname{ega4.c04.sec.5-7-3}
\expandafter\def\csname sga@target@IV.5.7.3.i\endcsname{ega4.c04.sec.5-7-3-i}
\expandafter\def\csname sga@target@IV.5.7.3.ii\endcsname{ega4.c04.sec.5-7-3-ii}
\expandafter\def\csname sga@target@IV.5.7.3.iii\endcsname{ega4.c04.sec.5-7-3-iii}
\expandafter\def\csname sga@target@IV.5.7.3.iv\endcsname{ega4.c04.sec.5-7-3-iv}
\expandafter\def\csname sga@target@IV.5.7.3.v\endcsname{ega4.c04.sec.5-7-3-v}
\expandafter\def\csname sga@target@IV.5.7.3.vi\endcsname{ega4.c04.sec.5-7-3-vi}
\expandafter\def\csname sga@target@IV.5.7.4\endcsname{ega4.c04.prop.5-7-4}
\expandafter\def\csname sga@target@IV.5.7.4.1\endcsname{ega4.c04.eq.5-7-4-1}
\expandafter\def\csname sga@target@IV.5.7.4.2\endcsname{ega4.c04.eq.5-7-4-2}
\expandafter\def\csname sga@target@IV.5.7.4.i\endcsname{ega4.c04.sec.5-7-4-i}
\expandafter\def\csname sga@target@IV.5.7.4.ii\endcsname{ega4.c04.sec.5-7-4-ii}
\expandafter\def\csname sga@target@IV.5.7.5\endcsname{ega4.c04.prop.5-7-5}
\expandafter\def\csname sga@target@IV.5.7.6\endcsname{ega4.c04.cor.5-7-6}
\expandafter\def\csname sga@target@IV.5.7.7\endcsname{ega4.c04.cor.5-7-7}
\expandafter\def\csname sga@target@IV.5.7.8\endcsname{ega4.c04.sec.5-7-8}
\expandafter\def\csname sga@target@IV.5.7.9\endcsname{ega4.c04.prop.5-7-9}
\expandafter\def\csname sga@target@IV.5.7.9.1\endcsname{ega4.c04.eq.5-7-9-1}
\expandafter\def\csname sga@target@IV.5.7.9.2\endcsname{ega4.c04.eq.5-7-9-2}
\expandafter\def\csname sga@target@IV.5.8.1\endcsname{ega4.c04.statement.5-8-1}
\expandafter\def\csname sga@target@IV.5.8.2\endcsname{ega4.c04.def.5-8-2}
\expandafter\def\csname sga@target@IV.5.8.3\endcsname{ega4.c04.prop.5-8-3}
\expandafter\def\csname sga@target@IV.5.8.4\endcsname{ega4.c04.statement.5-8-4}
\expandafter\def\csname sga@target@IV.5.8.5\endcsname{ega4.c04.prop.5-8-5}
\expandafter\def\csname sga@target@IV.5.8.6\endcsname{ega4.c04.thm.5-8-6}
\expandafter\def\csname sga@target@IV.5.8.6.1\endcsname{ega4.c04.eq.5-8-6-1}
\expandafter\def\csname sga@target@IV.5.8.6.i\endcsname{ega4.c04.item.5-8-6-i}
\expandafter\def\csname sga@target@IV.5.8.6.ii\endcsname{ega4.c04.item.5-8-6-ii}
\expandafter\def\csname sga@target@IV.5.9.1\endcsname{ega4.c04.statement.5-9-1}
\expandafter\def\csname sga@target@IV.5.9.1.1\endcsname{ega4.c04.eq.5-9-1-1}
\expandafter\def\csname sga@target@IV.5.9.1.2\endcsname{ega4.c04.eq.5-9-1-2}
\expandafter\def\csname sga@target@IV.5.9.1.3\endcsname{ega4.c04.eq.5-9-1-3}
\expandafter\def\csname sga@target@IV.5.9.10\endcsname{ega4.c04.cor.5-9-10}
\expandafter\def\csname sga@target@IV.5.9.10.i\endcsname{ega4.c04.item.5-9-10-i}
\expandafter\def\csname sga@target@IV.5.9.10.ii\endcsname{ega4.c04.item.5-9-10-ii}
\expandafter\def\csname sga@target@IV.5.9.11\endcsname{ega4.c04.prop.5-9-11}
\expandafter\def\csname sga@target@IV.5.9.12\endcsname{ega4.c04.item.5-9-12}
\expandafter\def\csname sga@target@IV.5.9.12.i\endcsname{ega4.c04.item.5-9-12-i}
\expandafter\def\csname sga@target@IV.5.9.12.ii\endcsname{ega4.c04.item.5-9-12-ii}
\expandafter\def\csname sga@target@IV.5.9.12.iii\endcsname{ega4.c04.item.5-9-12-iii}
\expandafter\def\csname sga@target@IV.5.9.2\endcsname{ega4.c04.prop.5-9-2}
\expandafter\def\csname sga@target@IV.5.9.2.i\endcsname{ega4.c04.item.5-9-2-i}
\expandafter\def\csname sga@target@IV.5.9.2.ii\endcsname{ega4.c04.item.5-9-2-ii}
\expandafter\def\csname sga@target@IV.5.9.3\endcsname{ega4.c04.rem.5-9-3}
\expandafter\def\csname sga@target@IV.5.9.3.1\endcsname{ega4.c04.eq.5-9-3-1}
\expandafter\def\csname sga@target@IV.5.9.4\endcsname{ega4.c04.prop.5-9-4}
\expandafter\def\csname sga@target@IV.5.9.4.1\endcsname{ega4.c04.eq.5-9-4-1}
\expandafter\def\csname sga@target@IV.5.9.5\endcsname{ega4.c04.cor.5-9-5}
\expandafter\def\csname sga@target@IV.5.9.6\endcsname{ega4.c04.cor.5-9-6}
\expandafter\def\csname sga@target@IV.5.9.6.1\endcsname{ega4.c04.eq.5-9-6-1}
\expandafter\def\csname sga@target@IV.5.9.7\endcsname{ega4.c04.statement.5-9-7}
\expandafter\def\csname sga@target@IV.5.9.7.1\endcsname{ega4.c04.eq.5-9-7-1}
\expandafter\def\csname sga@target@IV.5.9.8\endcsname{ega4.c04.prop.5-9-8}
\expandafter\def\csname sga@target@IV.5.9.8.1\endcsname{ega4.c04.eq.5-9-8-1}
\expandafter\def\csname sga@target@IV.5.9.8.a\endcsname{ega4.c04.item.5-9-8-a}
\expandafter\def\csname sga@target@IV.5.9.8.a-prime\endcsname{ega4.c04.item.5-9-8-a-prime}
\expandafter\def\csname sga@target@IV.5.9.8.b\endcsname{ega4.c04.item.5-9-8-b}
\expandafter\def\csname sga@target@IV.5.9.8.b-prime\endcsname{ega4.c04.item.5-9-8-b-prime}
\expandafter\def\csname sga@target@IV.5.9.9\endcsname{ega4.c04.def.5-9-9}
\expandafter\def\csname sga@target@IV.6.1.1\endcsname{ega4.c04.prop.6-1-1}
\expandafter\def\csname sga@target@IV.6.1.1.1\endcsname{ega4.c04.eq.6-1-1-1}
\expandafter\def\csname sga@target@IV.6.1.2\endcsname{ega4.c04.cor.6-1-2}
\expandafter\def\csname sga@target@IV.6.1.2.1\endcsname{ega4.c04.eq.6-1-2-1}
\expandafter\def\csname sga@target@IV.6.1.3\endcsname{ega4.c04.cor.6-1-3}
\expandafter\def\csname sga@target@IV.6.1.4\endcsname{ega4.c04.cor.6-1-4}
\expandafter\def\csname sga@target@IV.6.1.4.1\endcsname{ega4.c04.eq.6-1-4-1}
\expandafter\def\csname sga@target@IV.6.1.5\endcsname{ega4.c04.prop.6-1-5}
\expandafter\def\csname sga@target@IV.6.10.1\endcsname{ega4.c04.statement.6-10-1}
\expandafter\def\csname sga@target@IV.6.10.2\endcsname{ega4.c04.prop.6-10-2}
\expandafter\def\csname sga@target@IV.6.10.3\endcsname{ega4.c04.prop.6-10-3}
\expandafter\def\csname sga@target@IV.6.10.3.1\endcsname{ega4.c04.eq.6-10-3-1}
\expandafter\def\csname sga@target@IV.6.10.3.2\endcsname{ega4.c04.eq.6-10-3-2}
\expandafter\def\csname sga@target@IV.6.10.3.3\endcsname{ega4.c04.eq.6-10-3-3}
\expandafter\def\csname sga@target@IV.6.10.4\endcsname{ega4.c04.cor.6-10-4}
\expandafter\def\csname sga@target@IV.6.10.5\endcsname{ega4.c04.prop.6-10-5}
\expandafter\def\csname sga@target@IV.6.10.5.1\endcsname{ega4.c04.eq.6-10-5-1}
\expandafter\def\csname sga@target@IV.6.10.5.2\endcsname{ega4.c04.eq.6-10-5-2}
\expandafter\def\csname sga@target@IV.6.10.5.3\endcsname{ega4.c04.eq.6-10-5-3}
\expandafter\def\csname sga@target@IV.6.10.5.4\endcsname{ega4.c04.eq.6-10-5-4}
\expandafter\def\csname sga@target@IV.6.10.6\endcsname{ega4.c04.prop.6-10-6}
\expandafter\def\csname sga@target@IV.6.10.6.1\endcsname{ega4.c04.eq.6-10-6-1}
\expandafter\def\csname sga@target@IV.6.10.6.2\endcsname{ega4.c04.eq.6-10-6-2}
\expandafter\def\csname sga@target@IV.6.10.6.3\endcsname{ega4.c04.eq.6-10-6-3}
\expandafter\def\csname sga@target@IV.6.11\endcsname{ega4.c04.subsec.6-11.key-iv-6-11}
\expandafter\def\csname sga@target@IV.6.11.1\endcsname{ega4.c04.lem.6-11-1}
\expandafter\def\csname sga@target@IV.6.11.2\endcsname{ega4.c04.prop.6-11-2}
\expandafter\def\csname sga@target@IV.6.11.2.1\endcsname{ega4.c04.eq.6-11-2-1}
\expandafter\def\csname sga@target@IV.6.11.2.2\endcsname{ega4.c04.eq.6-11-2-2}
\expandafter\def\csname sga@target@IV.6.11.3\endcsname{ega4.c04.cor.6-11-3}
\expandafter\def\csname sga@target@IV.6.11.4\endcsname{ega4.c04.statement.6-11-4}
\expandafter\def\csname sga@target@IV.6.11.5\endcsname{ega4.c04.prop.6-11-5}
\expandafter\def\csname sga@target@IV.6.11.5.1\endcsname{ega4.c04.eq.6-11-5-1}
\expandafter\def\csname sga@target@IV.6.11.6\endcsname{ega4.c04.prop.6-11-6}
\expandafter\def\csname sga@target@IV.6.11.6.1\endcsname{ega4.c04.lem.6-11-6-1}
\expandafter\def\csname sga@target@IV.6.11.7\endcsname{ega4.c04.cor.6-11-7}
\expandafter\def\csname sga@target@IV.6.11.8\endcsname{ega4.c04.prop.6-11-8}
\expandafter\def\csname sga@target@IV.6.11.8.1\endcsname{ega4.c04.eq.6-11-8-1}
\expandafter\def\csname sga@target@IV.6.11.9\endcsname{ega4.c04.statement.6-11-9}
\expandafter\def\csname sga@target@IV.6.12\endcsname{ega4.c04.subsec.6-12.key-iv-6-12}
\expandafter\def\csname sga@target@IV.6.12.1\endcsname{ega4.c04.statement.6-12-1}
\expandafter\def\csname sga@target@IV.6.12.10\endcsname{ega4.c04.prop.6-12-10}
\expandafter\def\csname sga@target@IV.6.12.2\endcsname{ega4.c04.prop.6-12-2}
\expandafter\def\csname sga@target@IV.6.12.3\endcsname{ega4.c04.cor.6-12-3}
\expandafter\def\csname sga@target@IV.6.12.4\endcsname{ega4.c04.thm.6-12-4}
\expandafter\def\csname sga@target@IV.6.12.4.1\endcsname{ega4.c04.lem.6-12-4-1}
\expandafter\def\csname sga@target@IV.6.12.5\endcsname{ega4.c04.cor.6-12-5}
\expandafter\def\csname sga@target@IV.6.12.6\endcsname{ega4.c04.cor.6-12-6}
\expandafter\def\csname sga@target@IV.6.12.7\endcsname{ega4.c04.thm.6-12-7}
\expandafter\def\csname sga@target@IV.6.12.8\endcsname{ega4.c04.cor.6-12-8}
\expandafter\def\csname sga@target@IV.6.12.9\endcsname{ega4.c04.prop.6-12-9}
\expandafter\def\csname sga@target@IV.6.13\endcsname{ega4.c04.subsec.6-13.key-iv-6-13}
\expandafter\def\csname sga@target@IV.6.13.1\endcsname{ega4.c04.statement.6-13-1}
\expandafter\def\csname sga@target@IV.6.13.2\endcsname{ega4.c04.prop.6-13-2}
\expandafter\def\csname sga@target@IV.6.13.3\endcsname{ega4.c04.cor.6-13-3}
\expandafter\def\csname sga@target@IV.6.13.4\endcsname{ega4.c04.prop.6-13-4}
\expandafter\def\csname sga@target@IV.6.13.5\endcsname{ega4.c04.cor.6-13-5}
\expandafter\def\csname sga@target@IV.6.13.6\endcsname{ega4.c04.prop.6-13-6}
\expandafter\def\csname sga@target@IV.6.13.7\endcsname{ega4.c04.prop.6-13-7}
\expandafter\def\csname sga@target@IV.6.13.7.1\endcsname{ega4.c04.lem.6-13-7-1}
\expandafter\def\csname sga@target@IV.6.14.1\endcsname{ega4.c04.prop.6-14-1}
\expandafter\def\csname sga@target@IV.6.14.1.1\endcsname{ega4.c04.lem.6-14-1-1}
\expandafter\def\csname sga@target@IV.6.14.1.2\endcsname{ega4.c04.lem.6-14-1-2}
\expandafter\def\csname sga@target@IV.6.14.1.2.env\endcsname{ega4.c04.statement.6-14-1-2-env}
\expandafter\def\csname sga@target@IV.6.14.1.3\endcsname{ega4.c04.eq.6-14-1-3}
\expandafter\def\csname sga@target@IV.6.14.1.4\endcsname{ega4.c04.lem.6-14-1-4}
\expandafter\def\csname sga@target@IV.6.14.2\endcsname{ega4.c04.cor.6-14-2}
\expandafter\def\csname sga@target@IV.6.14.3\endcsname{ega4.c04.cor.6-14-3}
\expandafter\def\csname sga@target@IV.6.14.4\endcsname{ega4.c04.prop.6-14-4}
\expandafter\def\csname sga@target@IV.6.14.4.1\endcsname{ega4.c04.lem.6-14-4-1}
\expandafter\def\csname sga@target@IV.6.14.5\endcsname{ega4.c04.cor.6-14-5}
\expandafter\def\csname sga@target@IV.6.15.1\endcsname{ega4.c04.statement.6-15-1}
\expandafter\def\csname sga@target@IV.6.15.10\endcsname{ega4.c04.prop.6-15-10}
\expandafter\def\csname sga@target@IV.6.15.11\endcsname{ega4.c04.item.6-15-11}
\expandafter\def\csname sga@target@IV.6.15.2\endcsname{ega4.c04.statement.6-15-2}
\expandafter\def\csname sga@target@IV.6.15.3\endcsname{ega4.c04.statement.6-15-3}
\expandafter\def\csname sga@target@IV.6.15.3.1\endcsname{ega4.c04.lem.6-15-3-1}
\expandafter\def\csname sga@target@IV.6.15.4\endcsname{ega4.c04.statement.6-15-4}
\expandafter\def\csname sga@target@IV.6.15.4.1\endcsname{ega4.c04.lem.6-15-4-1}
\expandafter\def\csname sga@target@IV.6.15.5\endcsname{ega4.c04.prop.6-15-5}
\expandafter\def\csname sga@target@IV.6.15.6\endcsname{ega4.c04.prop.6-15-6}
\expandafter\def\csname sga@target@IV.6.15.7\endcsname{ega4.c04.prop.6-15-7}
\expandafter\def\csname sga@target@IV.6.15.8\endcsname{ega4.c04.lem.6-15-8}
\expandafter\def\csname sga@target@IV.6.15.9\endcsname{ega4.c04.cor.6-15-9}
\expandafter\def\csname sga@target@IV.6.2.1\endcsname{ega4.c04.prop.6-2-1}
\expandafter\def\csname sga@target@IV.6.2.1.1\endcsname{ega4.c04.eq.6-2-1-1}
\expandafter\def\csname sga@target@IV.6.2.1.2\endcsname{ega4.c04.eq.6-2-1-2}
\expandafter\def\csname sga@target@IV.6.2.2\endcsname{ega4.c04.cor.6-2-2}
\expandafter\def\csname sga@target@IV.6.2.3\endcsname{ega4.c04.prop.6-2-3}
\expandafter\def\csname sga@target@IV.6.2.3.1\endcsname{ega4.c04.eq.6-2-3-1}
\expandafter\def\csname sga@target@IV.6.2.4\endcsname{ega4.c04.prop.6-2-4}
\expandafter\def\csname sga@target@IV.6.2.4.1\endcsname{ega4.c04.eq.6-2-4-1}
\expandafter\def\csname sga@target@IV.6.3.1\endcsname{ega4.c04.prop.6-3-1}
\expandafter\def\csname sga@target@IV.6.3.1.1\endcsname{ega4.c04.eq.6-3-1-1}
\expandafter\def\csname sga@target@IV.6.3.2\endcsname{ega4.c04.cor.6-3-2}
\expandafter\def\csname sga@target@IV.6.3.2.1\endcsname{ega4.c04.eq.6-3-2-1}
\expandafter\def\csname sga@target@IV.6.3.3\endcsname{ega4.c04.cor.6-3-3}
\expandafter\def\csname sga@target@IV.6.3.4\endcsname{ega4.c04.cor.6-3-4}
\expandafter\def\csname sga@target@IV.6.3.4.1\endcsname{ega4.c04.eq.6-3-4-1}
\expandafter\def\csname sga@target@IV.6.3.4.2\endcsname{ega4.c04.eq.6-3-4-2}
\expandafter\def\csname sga@target@IV.6.3.5\endcsname{ega4.c04.cor.6-3-5}
\expandafter\def\csname sga@target@IV.6.3.6\endcsname{ega4.c04.cor.6-3-6}
\expandafter\def\csname sga@target@IV.6.3.7\endcsname{ega4.c04.cor.6-3-7}
\expandafter\def\csname sga@target@IV.6.3.8\endcsname{ega4.c04.prop.6-3-8}
\expandafter\def\csname sga@target@IV.6.4.1\endcsname{ega4.c04.prop.6-4-1}
\expandafter\def\csname sga@target@IV.6.4.2\endcsname{ega4.c04.cor.6-4-2}
\expandafter\def\csname sga@target@IV.6.4.3\endcsname{ega4.c04.rem.6-4-3}
\expandafter\def\csname sga@target@IV.6.5.1\endcsname{ega4.c04.prop.6-5-1}
\expandafter\def\csname sga@target@IV.6.5.2\endcsname{ega4.c04.cor.6-5-2}
\expandafter\def\csname sga@target@IV.6.5.3\endcsname{ega4.c04.prop.6-5-3}
\expandafter\def\csname sga@target@IV.6.5.4\endcsname{ega4.c04.cor.6-5-4}
\expandafter\def\csname sga@target@IV.6.5.5\endcsname{ega4.c04.item.6-5-5}
\expandafter\def\csname sga@target@IV.6.6.1\endcsname{ega4.c04.prop.6-6-1}
\expandafter\def\csname sga@target@IV.6.6.2\endcsname{ega4.c04.rem.6-6-2}
\expandafter\def\csname sga@target@IV.6.7.1\endcsname{ega4.c04.prop.6-7-1}
\expandafter\def\csname sga@target@IV.6.7.1.1\endcsname{ega4.c04.lem.6-7-1-1}
\expandafter\def\csname sga@target@IV.6.7.2\endcsname{ega4.c04.cor.6-7-2}
\expandafter\def\csname sga@target@IV.6.7.3\endcsname{ega4.c04.cor.6-7-3}
\expandafter\def\csname sga@target@IV.6.7.4\endcsname{ega4.c04.prop.6-7-4}
\expandafter\def\csname sga@target@IV.6.7.4.1\endcsname{ega4.c04.lem.6-7-4-1}
\expandafter\def\csname sga@target@IV.6.7.5\endcsname{ega4.c04.item.6-7-5}
\expandafter\def\csname sga@target@IV.6.7.6\endcsname{ega4.c04.def.6-7-6}
\expandafter\def\csname sga@target@IV.6.7.7\endcsname{ega4.c04.prop.6-7-7}
\expandafter\def\csname sga@target@IV.6.7.8\endcsname{ega4.c04.cor.6-7-8}
\expandafter\def\csname sga@target@IV.6.8.1\endcsname{ega4.c04.def.6-8-1}
\expandafter\def\csname sga@target@IV.6.8.2\endcsname{ega4.c04.prop.6-8-2}
\expandafter\def\csname sga@target@IV.6.8.3\endcsname{ega4.c04.prop.6-8-3}
\expandafter\def\csname sga@target@IV.6.8.4\endcsname{ega4.c04.item.6-8-4}
\expandafter\def\csname sga@target@IV.6.8.5\endcsname{ega4.c04.prop.6-8-5}
\expandafter\def\csname sga@target@IV.6.8.6\endcsname{ega4.c04.thm.6-8-6}
\expandafter\def\csname sga@target@IV.6.8.7\endcsname{ega4.c04.cor.6-8-7}
\expandafter\def\csname sga@target@IV.6.8.8\endcsname{ega4.c04.rem.6-8-8}
\expandafter\def\csname sga@target@IV.6.9.1\endcsname{ega4.c04.thm.6-9-1}
\expandafter\def\csname sga@target@IV.6.9.2\endcsname{ega4.c04.lem.6-9-2}
\expandafter\def\csname sga@target@IV.6.9.3\endcsname{ega4.c04.cor.6-9-3}
\expandafter\def\csname sga@target@IV.7.1.1\endcsname{ega4.c04.def.7-1-1}
\expandafter\def\csname sga@target@IV.7.1.10\endcsname{ega4.c04.prop.7-1-10}
\expandafter\def\csname sga@target@IV.7.1.11\endcsname{ega4.c04.prop.7-1-11}
\expandafter\def\csname sga@target@IV.7.1.12\endcsname{ega4.c04.item.7-1-12}
\expandafter\def\csname sga@target@IV.7.1.13\endcsname{ega4.c04.cor.7-1-13}
\expandafter\def\csname sga@target@IV.7.1.2\endcsname{ega4.c04.prop.7-1-2}
\expandafter\def\csname sga@target@IV.7.1.3\endcsname{ega4.c04.prop.7-1-3}
\expandafter\def\csname sga@target@IV.7.1.3.1\endcsname{ega4.c04.eq.7-1-3-1}
\expandafter\def\csname sga@target@IV.7.1.3.2\endcsname{ega4.c04.eq.7-1-3-2}
\expandafter\def\csname sga@target@IV.7.1.3.3\endcsname{ega4.c04.eq.7-1-3-3}
\expandafter\def\csname sga@target@IV.7.1.3.4\endcsname{ega4.c04.eq.7-1-3-4}
\expandafter\def\csname sga@target@IV.7.1.3.5\endcsname{ega4.c04.eq.7-1-3-5}
\expandafter\def\csname sga@target@IV.7.1.3.6\endcsname{ega4.c04.eq.7-1-3-6}
\expandafter\def\csname sga@target@IV.7.1.3.7\endcsname{ega4.c04.eq.7-1-3-7}
\expandafter\def\csname sga@target@IV.7.1.4\endcsname{ega4.c04.cor.7-1-4}
\expandafter\def\csname sga@target@IV.7.1.5\endcsname{ega4.c04.cor.7-1-5}
\expandafter\def\csname sga@target@IV.7.1.6\endcsname{ega4.c04.rem.7-1-6}
\expandafter\def\csname sga@target@IV.7.1.7\endcsname{ega4.c04.lem.7-1-7}
\expandafter\def\csname sga@target@IV.7.1.8\endcsname{ega4.c04.thm.7-1-8}
\expandafter\def\csname sga@target@IV.7.1.9\endcsname{ega4.c04.def.7-1-9}
\expandafter\def\csname sga@target@IV.7.2.1\endcsname{ega4.c04.statement.7-2-1}
\expandafter\def\csname sga@target@IV.7.2.1.1\endcsname{ega4.c04.statement.7-2-1-1}
\expandafter\def\csname sga@target@IV.7.2.10\endcsname{ega4.c04.statement.7-2-10}
\expandafter\def\csname sga@target@IV.7.2.2\endcsname{ega4.c04.prop.7-2-2}
\expandafter\def\csname sga@target@IV.7.2.3\endcsname{ega4.c04.prop.7-2-3}
\expandafter\def\csname sga@target@IV.7.2.4\endcsname{ega4.c04.prop.7-2-4}
\expandafter\def\csname sga@target@IV.7.2.5\endcsname{ega4.c04.thm.7-2-5}
\expandafter\def\csname sga@target@IV.7.2.5.1\endcsname{ega4.c04.eq.7-2-5-1}
\expandafter\def\csname sga@target@IV.7.2.6\endcsname{ega4.c04.def.7-2-6}
\expandafter\def\csname sga@target@IV.7.2.7\endcsname{ega4.c04.cor.7-2-7}
\expandafter\def\csname sga@target@IV.7.2.8\endcsname{ega4.c04.cor.7-2-8}
\expandafter\def\csname sga@target@IV.7.2.9\endcsname{ega4.c04.cor.7-2-9}
\expandafter\def\csname sga@target@IV.7.3.1\endcsname{ega4.c04.statement.7-3-1}
\expandafter\def\csname sga@target@IV.7.3.10\endcsname{ega4.c04.statement.7-3-10}
\expandafter\def\csname sga@target@IV.7.3.11\endcsname{ega4.c04.rem.7-3-11}
\expandafter\def\csname sga@target@IV.7.3.12\endcsname{ega4.c04.prop.7-3-12}
\expandafter\def\csname sga@target@IV.7.3.12.1\endcsname{ega4.c04.eq.7-3-12-1}
\expandafter\def\csname sga@target@IV.7.3.13\endcsname{ega4.c04.statement.7-3-13}
\expandafter\def\csname sga@target@IV.7.3.14\endcsname{ega4.c04.prop.7-3-14}
\expandafter\def\csname sga@target@IV.7.3.15\endcsname{ega4.c04.prop.7-3-15}
\expandafter\def\csname sga@target@IV.7.3.16\endcsname{ega4.c04.prop.7-3-16}
\expandafter\def\csname sga@target@IV.7.3.16.1\endcsname{ega4.c04.eq.7-3-16-1}
\expandafter\def\csname sga@target@IV.7.3.16.2\endcsname{ega4.c04.lem.7-3-16-2}
\expandafter\def\csname sga@target@IV.7.3.17\endcsname{ega4.c04.prop.7-3-17}
\expandafter\def\csname sga@target@IV.7.3.18\endcsname{ega4.c04.prop.7-3-18}
\expandafter\def\csname sga@target@IV.7.3.19\endcsname{ega4.c04.statement.7-3-19}
\expandafter\def\csname sga@target@IV.7.3.2\endcsname{ega4.c04.lem.7-3-2}
\expandafter\def\csname sga@target@IV.7.3.3\endcsname{ega4.c04.cor.7-3-3}
\expandafter\def\csname sga@target@IV.7.3.4\endcsname{ega4.c04.statement.7-3-4}
\expandafter\def\csname sga@target@IV.7.3.5\endcsname{ega4.c04.item.7-3-5}
\expandafter\def\csname sga@target@IV.7.3.6\endcsname{ega4.c04.statement.7-3-6}
\expandafter\def\csname sga@target@IV.7.3.7\endcsname{ega4.c04.lem.7-3-7}
\expandafter\def\csname sga@target@IV.7.3.8\endcsname{ega4.c04.statement.7-3-8}
\expandafter\def\csname sga@target@IV.7.3.9\endcsname{ega4.c04.item.7-3-9}
\expandafter\def\csname sga@target@IV.7.4.1\endcsname{ega4.c04.statement.7-4-1}
\expandafter\def\csname sga@target@IV.7.4.2\endcsname{ega4.c04.lem.7-4-2}
\expandafter\def\csname sga@target@IV.7.4.3\endcsname{ega4.c04.cor.7-4-3}
\expandafter\def\csname sga@target@IV.7.4.4\endcsname{ega4.c04.thm.7-4-4}
\expandafter\def\csname sga@target@IV.7.4.4.1\endcsname{ega4.c04.lem.7-4-4-1}
\expandafter\def\csname sga@target@IV.7.4.4.2\endcsname{ega4.c04.lem.7-4-4-2}
\expandafter\def\csname sga@target@IV.7.4.5\endcsname{ega4.c04.cor.7-4-5}
\expandafter\def\csname sga@target@IV.7.4.6\endcsname{ega4.c04.prop.7-4-6}
\expandafter\def\csname sga@target@IV.7.4.7\endcsname{ega4.c04.cor.7-4-7}
\expandafter\def\csname sga@target@IV.7.4.8\endcsname{ega4.c04.rem.7-4-8}
\expandafter\def\csname sga@target@IV.7.5.0\endcsname{ega4.c04.statement.7-5-0}
\expandafter\def\csname sga@target@IV.7.5.1\endcsname{ega4.c04.thm.7-5-1}
\expandafter\def\csname sga@target@IV.7.5.1.1\endcsname{ega4.c04.lem.7-5-1-1}
\expandafter\def\csname sga@target@IV.7.5.2\endcsname{ega4.c04.cor.7-5-2}
\expandafter\def\csname sga@target@IV.7.5.3\endcsname{ega4.c04.item.7-5-3}
\expandafter\def\csname sga@target@IV.7.5.4\endcsname{ega4.c04.item.7-5-4}
\expandafter\def\csname sga@target@IV.7.5.5\endcsname{ega4.c04.lem.7-5-5}
\expandafter\def\csname sga@target@IV.7.5.6\endcsname{ega4.c04.prop.7-5-6}
\expandafter\def\csname sga@target@IV.7.5.7\endcsname{ega4.c04.cor.7-5-7}
\expandafter\def\csname sga@target@IV.7.5.7.1\endcsname{ega4.c04.lem.7-5-7-1}
\expandafter\def\csname sga@target@IV.7.5.8\endcsname{ega4.c04.rem.7-5-8}
\expandafter\def\csname sga@target@IV.7.6.1\endcsname{ega4.c04.prop.7-6-1}
\expandafter\def\csname sga@target@IV.7.6.2\endcsname{ega4.c04.cor.7-6-2}
\expandafter\def\csname sga@target@IV.7.6.3\endcsname{ega4.c04.cor.7-6-3}
\expandafter\def\csname sga@target@IV.7.6.4\endcsname{ega4.c04.thm.7-6-4}
\expandafter\def\csname sga@target@IV.7.6.4.1\endcsname{ega4.c04.lem.7-6-4-1}
\expandafter\def\csname sga@target@IV.7.6.5\endcsname{ega4.c04.cor.7-6-5}
\expandafter\def\csname sga@target@IV.7.6.6\endcsname{ega4.c04.cor.7-6-6}
\expandafter\def\csname sga@target@IV.7.6.7\endcsname{ega4.c04.item.7-6-7}
\expandafter\def\csname sga@target@IV.7.7.1\endcsname{ega4.c04.statement.7-7-1}
\expandafter\def\csname sga@target@IV.7.7.2\endcsname{ega4.c04.thm.7-7-2}
\expandafter\def\csname sga@target@IV.7.7.3\endcsname{ega4.c04.cor.7-7-3}
\expandafter\def\csname sga@target@IV.7.7.4\endcsname{ega4.c04.cor.7-7-4}
\expandafter\def\csname sga@target@IV.7.8.1\endcsname{ega4.c04.statement.7-8-1}
\expandafter\def\csname sga@target@IV.7.8.2\endcsname{ega4.c04.def.7-8-2}
\expandafter\def\csname sga@target@IV.7.8.3\endcsname{ega4.c04.statement.7-8-3}
\expandafter\def\csname sga@target@IV.7.8.3.1\endcsname{ega4.c04.statement.7-8-3-1}
\expandafter\def\csname sga@target@IV.7.8.4\endcsname{ega4.c04.item.7-8-4}
\expandafter\def\csname sga@target@IV.7.8.4.1\endcsname{ega4.c04.statement.7-8-4-1}
\expandafter\def\csname sga@target@IV.7.8.5\endcsname{ega4.c04.statement.7-8-5}
\expandafter\def\csname sga@target@IV.7.8.6\endcsname{ega4.c04.prop.7-8-6}
\expandafter\def\csname sga@target@IV.7.9.1\endcsname{ega4.c04.statement.7-9-1}
\expandafter\def\csname sga@target@IV.7.9.10\endcsname{ega4.c04.item.7-9-10}
\expandafter\def\csname sga@target@IV.7.9.2\endcsname{ega4.c04.statement.7-9-2}
\expandafter\def\csname sga@target@IV.7.9.3\endcsname{ega4.c04.prop.7-9-3}
\expandafter\def\csname sga@target@IV.7.9.3.1\endcsname{ega4.c04.lem.7-9-3-1}
\expandafter\def\csname sga@target@IV.7.9.4\endcsname{ega4.c04.cor.7-9-4}
\expandafter\def\csname sga@target@IV.7.9.5\endcsname{ega4.c04.prop.7-9-5}
\expandafter\def\csname sga@target@IV.7.9.6\endcsname{ega4.c04.rem.7-9-6}
\expandafter\def\csname sga@target@IV.7.9.7\endcsname{ega4.c04.statement.7-9-7}
\expandafter\def\csname sga@target@IV.7.9.8\endcsname{ega4.c04.prop.7-9-8}
\expandafter\def\csname sga@target@IV.7.9.8.1\endcsname{ega4.c04.statement.7-9-8-1}
\expandafter\def\csname sga@target@IV.7.9.9\endcsname{ega4.c04.cor.7-9-9}
\expandafter\def\csname sga@target@IV.8.1.1\endcsname{ega4.c04.statement.8-1-1}
\expandafter\def\csname sga@target@IV.8.1.2\endcsname{ega4.c04.statement.8-1-2}
\expandafter\def\csname sga@target@IV.8.1.3\endcsname{ega4.c04.statement.8-1-3}
\expandafter\def\csname sga@target@IV.8.10.1\endcsname{ega4.c04.prop.8-10-1}
\expandafter\def\csname sga@target@IV.8.10.2\endcsname{ega4.c04.cor.8-10-2}
\expandafter\def\csname sga@target@IV.8.10.3\endcsname{ega4.c04.prop.8-10-3}
\expandafter\def\csname sga@target@IV.8.10.3.1\endcsname{ega4.c04.rem.8-10-3-1}
\expandafter\def\csname sga@target@IV.8.10.4\endcsname{ega4.c04.prop.8-10-4}
\expandafter\def\csname sga@target@IV.8.10.5\endcsname{ega4.c04.thm.8-10-5}
\expandafter\def\csname sga@target@IV.8.10.5.1\endcsname{ega4.c04.prop.8-10-5-1}
\expandafter\def\csname sga@target@IV.8.10.5.2\endcsname{ega4.c04.lem.8-10-5-2}
\expandafter\def\csname sga@target@IV.8.11\endcsname{ega4.c04.subsec.8-11.key-iv-8-11}
\expandafter\def\csname sga@target@IV.8.11.1\endcsname{ega4.c04.thm.8-11-1}
\expandafter\def\csname sga@target@IV.8.11.2\endcsname{ega4.c04.thm.8-11-2}
\expandafter\def\csname sga@target@IV.8.11.3\endcsname{ega4.c04.rem.8-11-3}
\expandafter\def\csname sga@target@IV.8.11.4\endcsname{ega4.c04.statement.8-11-4}
\expandafter\def\csname sga@target@IV.8.11.5\endcsname{ega4.c04.prop.8-11-5}
\expandafter\def\csname sga@target@IV.8.11.5.1\endcsname{ega4.c04.rem.8-11-5-1}
\expandafter\def\csname sga@target@IV.8.11.6\endcsname{ega4.c04.prop.8-11-6}
\expandafter\def\csname sga@target@IV.8.12\endcsname{ega4.c04.subsec.8-12.key-iv-8-12}
\expandafter\def\csname sga@target@IV.8.12.1\endcsname{ega4.c04.lem.8-12-1}
\expandafter\def\csname sga@target@IV.8.12.10\endcsname{ega4.c04.cor.8-12-10}
\expandafter\def\csname sga@target@IV.8.12.10.1\endcsname{ega4.c04.lem.8-12-10-1}
\expandafter\def\csname sga@target@IV.8.12.11\endcsname{ega4.c04.cor.8-12-11}
\expandafter\def\csname sga@target@IV.8.12.12\endcsname{ega4.c04.cor.8-12-12}
\expandafter\def\csname sga@target@IV.8.12.2\endcsname{ega4.c04.lem.8-12-2}
\expandafter\def\csname sga@target@IV.8.12.3\endcsname{ega4.c04.prop.8-12-3}
\expandafter\def\csname sga@target@IV.8.12.4\endcsname{ega4.c04.cor.8-12-4}
\expandafter\def\csname sga@target@IV.8.12.5\endcsname{ega4.c04.statement.8-12-5}
\expandafter\def\csname sga@target@IV.8.12.6\endcsname{ega4.c04.thm.8-12-6}
\expandafter\def\csname sga@target@IV.8.12.6.1\endcsname{ega4.c04.eq.8-12-6-1}
\expandafter\def\csname sga@target@IV.8.12.7\endcsname{ega4.c04.rem.8-12-7}
\expandafter\def\csname sga@target@IV.8.12.8\endcsname{ega4.c04.cor.8-12-8}
\expandafter\def\csname sga@target@IV.8.12.9\endcsname{ega4.c04.cor.8-12-9}
\expandafter\def\csname sga@target@IV.8.13.1\endcsname{ega4.c04.prop.8-13-1}
\expandafter\def\csname sga@target@IV.8.13.1.1\endcsname{ega4.c04.eq.8-13-1-1}
\expandafter\def\csname sga@target@IV.8.13.2\endcsname{ega4.c04.cor.8-13-2}
\expandafter\def\csname sga@target@IV.8.13.2.1\endcsname{ega4.c04.eq.8-13-2-1}
\expandafter\def\csname sga@target@IV.8.13.3\endcsname{ega4.c04.statement.8-13-3}
\expandafter\def\csname sga@target@IV.8.13.4\endcsname{ega4.c04.statement.8-13-4}
\expandafter\def\csname sga@target@IV.8.13.4.1\endcsname{ega4.c04.eq.8-13-4-1}
\expandafter\def\csname sga@target@IV.8.13.4.2\endcsname{ega4.c04.eq.8-13-4-2}
\expandafter\def\csname sga@target@IV.8.13.5\endcsname{ega4.c04.prop.8-13-5}
\expandafter\def\csname sga@target@IV.8.13.6\endcsname{ega4.c04.rem.8-13-6}
\expandafter\def\csname sga@target@IV.8.14.1\endcsname{ega4.c04.statement.8-14-1}
\expandafter\def\csname sga@target@IV.8.14.2\endcsname{ega4.c04.prop.8-14-2}
\expandafter\def\csname sga@target@IV.8.14.2.1\endcsname{ega4.c04.eq.8-14-2-1}
\expandafter\def\csname sga@target@IV.8.14.2.2\endcsname{ega4.c04.cor.8-14-2-2}
\expandafter\def\csname sga@target@IV.8.14.2.3\endcsname{ega4.c04.eq.8-14-2-3}
\expandafter\def\csname sga@target@IV.8.14.2.4\endcsname{ega4.c04.eq.8-14-2-4}
\expandafter\def\csname sga@target@IV.8.14.3\endcsname{ega4.c04.rem.8-14-3}
\expandafter\def\csname sga@target@IV.8.2.1\endcsname{ega4.c04.statement.8-2-1}
\expandafter\def\csname sga@target@IV.8.2.1.1\endcsname{ega4.c04.eq.8-2-1-1}
\expandafter\def\csname sga@target@IV.8.2.10\endcsname{ega4.c04.cor.8-2-10}
\expandafter\def\csname sga@target@IV.8.2.11\endcsname{ega4.c04.cor.8-2-11}
\expandafter\def\csname sga@target@IV.8.2.12\endcsname{ega4.c04.cor.8-2-12}
\expandafter\def\csname sga@target@IV.8.2.12.1\endcsname{ega4.c04.eq.8-2-12-1}
\expandafter\def\csname sga@target@IV.8.2.13\endcsname{ega4.c04.prop.8-2-13}
\expandafter\def\csname sga@target@IV.8.2.14\endcsname{ega4.c04.rem.8-2-14}
\expandafter\def\csname sga@target@IV.8.2.2\endcsname{ega4.c04.statement.8-2-2}
\expandafter\def\csname sga@target@IV.8.2.3\endcsname{ega4.c04.prop.8-2-3}
\expandafter\def\csname sga@target@IV.8.2.3.1\endcsname{ega4.c04.eq.8-2-3-1}
\expandafter\def\csname sga@target@IV.8.2.4\endcsname{ega4.c04.lem.8-2-4}
\expandafter\def\csname sga@target@IV.8.2.5\endcsname{ega4.c04.prop.8-2-5}
\expandafter\def\csname sga@target@IV.8.2.6\endcsname{ega4.c04.lem.8-2-6}
\expandafter\def\csname sga@target@IV.8.2.7\endcsname{ega4.c04.rem.8-2-7}
\expandafter\def\csname sga@target@IV.8.2.8\endcsname{ega4.c04.statement.8-2-8}
\expandafter\def\csname sga@target@IV.8.2.9\endcsname{ega4.c04.prop.8-2-9}
\expandafter\def\csname sga@target@IV.8.3.1\endcsname{ega4.c04.statement.8-3-1}
\expandafter\def\csname sga@target@IV.8.3.10\endcsname{ega4.c04.statement.8-3-10}
\expandafter\def\csname sga@target@IV.8.3.10.1\endcsname{ega4.c04.eq.8-3-10-1}
\expandafter\def\csname sga@target@IV.8.3.11\endcsname{ega4.c04.thm.8-3-11}
\expandafter\def\csname sga@target@IV.8.3.12\endcsname{ega4.c04.cor.8-3-12}
\expandafter\def\csname sga@target@IV.8.3.13\endcsname{ega4.c04.prop.8-3-13}
\expandafter\def\csname sga@target@IV.8.3.2\endcsname{ega4.c04.thm.8-3-2}
\expandafter\def\csname sga@target@IV.8.3.2.1\endcsname{ega4.c04.eq.8-3-2-1}
\expandafter\def\csname sga@target@IV.8.3.3\endcsname{ega4.c04.cor.8-3-3}
\expandafter\def\csname sga@target@IV.8.3.4\endcsname{ega4.c04.cor.8-3-4}
\expandafter\def\csname sga@target@IV.8.3.5\endcsname{ega4.c04.cor.8-3-5}
\expandafter\def\csname sga@target@IV.8.3.6\endcsname{ega4.c04.cor.8-3-6}
\expandafter\def\csname sga@target@IV.8.3.7\endcsname{ega4.c04.cor.8-3-7}
\expandafter\def\csname sga@target@IV.8.3.7.1\endcsname{ega4.c04.eq.8-3-7-1}
\expandafter\def\csname sga@target@IV.8.3.8\endcsname{ega4.c04.prop.8-3-8}
\expandafter\def\csname sga@target@IV.8.3.9\endcsname{ega4.c04.statement.8-3-9}
\expandafter\def\csname sga@target@IV.8.3.9.1\endcsname{ega4.c04.eq.8-3-9-1}
\expandafter\def\csname sga@target@IV.8.3.9.2\endcsname{ega4.c04.eq.8-3-9-2}
\expandafter\def\csname sga@target@IV.8.3.9.3\endcsname{ega4.c04.eq.8-3-9-3}
\expandafter\def\csname sga@target@IV.8.3.9.4\endcsname{ega4.c04.eq.8-3-9-4}
\expandafter\def\csname sga@target@IV.8.3.9.5\endcsname{ega4.c04.eq.8-3-9-5}
\expandafter\def\csname sga@target@IV.8.4.1\endcsname{ega4.c04.prop.8-4-1}
\expandafter\def\csname sga@target@IV.8.4.2\endcsname{ega4.c04.prop.8-4-2}
\expandafter\def\csname sga@target@IV.8.4.2.1\endcsname{ega4.c04.lem.8-4-2-1}
\expandafter\def\csname sga@target@IV.8.4.3\endcsname{ega4.c04.cor.8-4-3}
\expandafter\def\csname sga@target@IV.8.4.4\endcsname{ega4.c04.prop.8-4-4}
\expandafter\def\csname sga@target@IV.8.4.5\endcsname{ega4.c04.cor.8-4-5}
\expandafter\def\csname sga@target@IV.8.4.6\endcsname{ega4.c04.rem.8-4-6}
\expandafter\def\csname sga@target@IV.8.5.1\endcsname{ega4.c04.statement.8-5-1}
\expandafter\def\csname sga@target@IV.8.5.1.1\endcsname{ega4.c04.eq.8-5-1-1}
\expandafter\def\csname sga@target@IV.8.5.1.2\endcsname{ega4.c04.eq.8-5-1-2}
\expandafter\def\csname sga@target@IV.8.5.1.3\endcsname{ega4.c04.eq.8-5-1-3}
\expandafter\def\csname sga@target@IV.8.5.1.4\endcsname{ega4.c04.eq.8-5-1-4}
\expandafter\def\csname sga@target@IV.8.5.10\endcsname{ega4.c04.statement.8-5-10}
\expandafter\def\csname sga@target@IV.8.5.10.1\endcsname{ega4.c04.eq.8-5-10-1}
\expandafter\def\csname sga@target@IV.8.5.10.2\endcsname{ega4.c04.eq.8-5-10-2}
\expandafter\def\csname sga@target@IV.8.5.11\endcsname{ega4.c04.prop.8-5-11}
\expandafter\def\csname sga@target@IV.8.5.2\endcsname{ega4.c04.thm.8-5-2}
\expandafter\def\csname sga@target@IV.8.5.2.1\endcsname{ega4.c04.lem.8-5-2-1}
\expandafter\def\csname sga@target@IV.8.5.2.2\endcsname{ega4.c04.eq.8-5-2-2}
\expandafter\def\csname sga@target@IV.8.5.2.3\endcsname{ega4.c04.eq.8-5-2-3}
\expandafter\def\csname sga@target@IV.8.5.2.4\endcsname{ega4.c04.cor.8-5-2-4}
\expandafter\def\csname sga@target@IV.8.5.2.5\endcsname{ega4.c04.cor.8-5-2-5}
\expandafter\def\csname sga@target@IV.8.5.3\endcsname{ega4.c04.statement.8-5-3}
\expandafter\def\csname sga@target@IV.8.5.4\endcsname{ega4.c04.cor.8-5-4}
\expandafter\def\csname sga@target@IV.8.5.5\endcsname{ega4.c04.prop.8-5-5}
\expandafter\def\csname sga@target@IV.8.5.6\endcsname{ega4.c04.prop.8-5-6}
\expandafter\def\csname sga@target@IV.8.5.7\endcsname{ega4.c04.cor.8-5-7}
\expandafter\def\csname sga@target@IV.8.5.8\endcsname{ega4.c04.cor.8-5-8}
\expandafter\def\csname sga@target@IV.8.5.9\endcsname{ega4.c04.lem.8-5-9}
\expandafter\def\csname sga@target@IV.8.6.1\endcsname{ega4.c04.statement.8-6-1}
\expandafter\def\csname sga@target@IV.8.6.2\endcsname{ega4.c04.statement.8-6-2}
\expandafter\def\csname sga@target@IV.8.6.2.1\endcsname{ega4.c04.eq.8-6-2-1}
\expandafter\def\csname sga@target@IV.8.6.2.2\endcsname{ega4.c04.eq.8-6-2-2}
\expandafter\def\csname sga@target@IV.8.6.2.3\endcsname{ega4.c04.eq.8-6-2-3}
\expandafter\def\csname sga@target@IV.8.6.3\endcsname{ega4.c04.prop.8-6-3}
\expandafter\def\csname sga@target@IV.8.6.3.1\endcsname{ega4.c04.cor.8-6-3-1}
\expandafter\def\csname sga@target@IV.8.6.4\endcsname{ega4.c04.cor.8-6-4}
\expandafter\def\csname sga@target@IV.8.7.1\endcsname{ega4.c04.prop.8-7-1}
\expandafter\def\csname sga@target@IV.8.7.2\endcsname{ega4.c04.prop.8-7-2}
\expandafter\def\csname sga@target@IV.8.7.3\endcsname{ega4.c04.cor.8-7-3}
\expandafter\def\csname sga@target@IV.8.7.4\endcsname{ega4.c04.rem.8-7-4}
\expandafter\def\csname sga@target@IV.8.8.1\endcsname{ega4.c04.statement.8-8-1}
\expandafter\def\csname sga@target@IV.8.8.1.1\endcsname{ega4.c04.eq.8-8-1-1}
\expandafter\def\csname sga@target@IV.8.8.2\endcsname{ega4.c04.thm.8-8-2}
\expandafter\def\csname sga@target@IV.8.8.2.1\endcsname{ega4.c04.lem.8-8-2-1}
\expandafter\def\csname sga@target@IV.8.8.2.2\endcsname{ega4.c04.lem.8-8-2-2}
\expandafter\def\csname sga@target@IV.8.8.2.3\endcsname{ega4.c04.lem.8-8-2-3}
\expandafter\def\csname sga@target@IV.8.8.2.4\endcsname{ega4.c04.cor.8-8-2-4}
\expandafter\def\csname sga@target@IV.8.8.2.5\endcsname{ega4.c04.cor.8-8-2-5}
\expandafter\def\csname sga@target@IV.8.8.3\endcsname{ega4.c04.statement.8-8-3}
\expandafter\def\csname sga@target@IV.8.9.1\endcsname{ega4.c04.prop.8-9-1}
\expandafter\def\csname sga@target@IV.8.9.2\endcsname{ega4.c04.statement.8-9-2}
\expandafter\def\csname sga@target@IV.8.9.3\endcsname{ega4.c04.prop.8-9-3}
\expandafter\def\csname sga@target@IV.8.9.4\endcsname{ega4.c04.prop.8-9-4}
\expandafter\def\csname sga@target@IV.8.9.4.1\endcsname{ega4.c04.lem.8-9-4-1}
\expandafter\def\csname sga@target@IV.8.9.5\endcsname{ega4.c04.cor.8-9-5}
\expandafter\def\csname sga@target@IV.9.1.1\endcsname{ega4.c04.prop.9-1-1}
\expandafter\def\csname sga@target@IV.9.1.2\endcsname{ega4.c04.cor.9-1-2}
\expandafter\def\csname sga@target@IV.9.1.3\endcsname{ega4.c04.rem.9-1-3}
\expandafter\def\csname sga@target@IV.9.1.4\endcsname{ega4.c04.prop.9-1-4}
\expandafter\def\csname sga@target@IV.9.1.5\endcsname{ega4.c04.cor.9-1-5}
\expandafter\def\csname sga@target@IV.9.1.6\endcsname{ega4.c04.prop.9-1-6}
\expandafter\def\csname sga@target@IV.9.2\endcsname{ega4.c04.subsec.9-2.key-iv-9-2}
\expandafter\def\csname sga@target@IV.9.2.1\endcsname{ega4.c04.def.9-2-1}
\expandafter\def\csname sga@target@IV.9.2.2\endcsname{ega4.c04.item.9-2-2}
\expandafter\def\csname sga@target@IV.9.2.3\endcsname{ega4.c04.prop.9-2-3}
\expandafter\def\csname sga@target@IV.9.2.4\endcsname{ega4.c04.cor.9-2-4}
\expandafter\def\csname sga@target@IV.9.2.5\endcsname{ega4.c04.rem.9-2-5}
\expandafter\def\csname sga@target@IV.9.2.6\endcsname{ega4.c04.prop.9-2-6}
\expandafter\def\csname sga@target@IV.9.2.6.1\endcsname{ega4.c04.prop.9-2-6-1}
\expandafter\def\csname sga@target@IV.9.2.6.2\endcsname{ega4.c04.cor.9-2-6-2}
\expandafter\def\csname sga@target@IV.9.3.1\endcsname{ega4.c04.prop.9-3-1}
\expandafter\def\csname sga@target@IV.9.3.2\endcsname{ega4.c04.prop.9-3-2}
\expandafter\def\csname sga@target@IV.9.3.3\endcsname{ega4.c04.prop.9-3-3}
\expandafter\def\csname sga@target@IV.9.3.4\endcsname{ega4.c04.rem.9-3-4}
\expandafter\def\csname sga@target@IV.9.3.5\endcsname{ega4.c04.prop.9-3-5}
\expandafter\def\csname sga@target@IV.9.3.6\endcsname{ega4.c04.statement.9-3-6}
\expandafter\def\csname sga@target@IV.9.4.1\endcsname{ega4.c04.statement.9-4-1}
\expandafter\def\csname sga@target@IV.9.4.2\endcsname{ega4.c04.prop.9-4-2}
\expandafter\def\csname sga@target@IV.9.4.3\endcsname{ega4.c04.cor.9-4-3}
\expandafter\def\csname sga@target@IV.9.4.3.1\endcsname{ega4.c04.eq.9-4-3-1}
\expandafter\def\csname sga@target@IV.9.4.4\endcsname{ega4.c04.prop.9-4-4}
\expandafter\def\csname sga@target@IV.9.4.5\endcsname{ega4.c04.cor.9-4-5}
\expandafter\def\csname sga@target@IV.9.4.6\endcsname{ega4.c04.cor.9-4-6}
\expandafter\def\csname sga@target@IV.9.4.7\endcsname{ega4.c04.prop.9-4-7}
\expandafter\def\csname sga@target@IV.9.4.7.1\endcsname{ega4.c04.lem.9-4-7-1}
\expandafter\def\csname sga@target@IV.9.4.7.2\endcsname{ega4.c04.rem.9-4-7-2}
\expandafter\def\csname sga@target@IV.9.4.8\endcsname{ega4.c04.prop.9-4-8}
\expandafter\def\csname sga@target@IV.9.4.9\endcsname{ega4.c04.rem.9-4-9}
\expandafter\def\csname sga@target@IV.9.5.1\endcsname{ega4.c04.prop.9-5-1}
\expandafter\def\csname sga@target@IV.9.5.2\endcsname{ega4.c04.cor.9-5-2}
\expandafter\def\csname sga@target@IV.9.5.3\endcsname{ega4.c04.prop.9-5-3}
\expandafter\def\csname sga@target@IV.9.5.3.case2\endcsname{ega4.c04.item.9-5-3-case2}
\expandafter\def\csname sga@target@IV.9.5.4\endcsname{ega4.c04.cor.9-5-4}
\expandafter\def\csname sga@target@IV.9.5.5\endcsname{ega4.c04.prop.9-5-5}
\expandafter\def\csname sga@target@IV.9.5.6\endcsname{ega4.c04.cor.9-5-6}
\expandafter\def\csname sga@target@IV.9.5.7\endcsname{ega4.c04.rem.9-5-7}
\expandafter\def\csname sga@target@IV.9.6.1\endcsname{ega4.c04.prop.9-6-1}
\expandafter\def\csname sga@target@IV.9.6.2\endcsname{ega4.c04.prop.9-6-2}
\expandafter\def\csname sga@target@IV.9.6.3\endcsname{ega4.c04.prop.9-6-3}
\expandafter\def\csname sga@target@IV.9.6.3.1\endcsname{ega4.c04.eq.9-6-3-1}
\expandafter\def\csname sga@target@IV.9.6.4\endcsname{ega4.c04.cor.9-6-4}
\expandafter\def\csname sga@target@IV.9.6.5\endcsname{ega4.c04.cor.9-6-5}
\expandafter\def\csname sga@target@IV.9.6.6\endcsname{ega4.c04.sec.9-6-6}
\expandafter\def\csname sga@target@IV.9.6.7\endcsname{ega4.c04.sec.9-6-7}
\expandafter\def\csname sga@target@IV.9.7\endcsname{ega4.c04.subsec.9-7.key-iv-9-7}
\expandafter\def\csname sga@target@IV.9.7.1\endcsname{ega4.c04.lem.9-7-1}
\expandafter\def\csname sga@target@IV.9.7.10\endcsname{ega4.c04.cor.9-7-10}
\expandafter\def\csname sga@target@IV.9.7.11\endcsname{ega4.c04.rem.9-7-11}
\expandafter\def\csname sga@target@IV.9.7.12\endcsname{ega4.c04.prop.9-7-12}
\expandafter\def\csname sga@target@IV.9.7.2\endcsname{ega4.c04.lem.9-7-2}
\expandafter\def\csname sga@target@IV.9.7.3\endcsname{ega4.c04.subsec.9-7-3}
\expandafter\def\csname sga@target@IV.9.7.4\endcsname{ega4.c04.lem.9-7-4}
\expandafter\def\csname sga@target@IV.9.7.5\endcsname{ega4.c04.lem.9-7-5}
\expandafter\def\csname sga@target@IV.9.7.5.1\endcsname{ega4.c04.lem.9-7-5-1}
\expandafter\def\csname sga@target@IV.9.7.6\endcsname{ega4.c04.prop.9-7-6}
\expandafter\def\csname sga@target@IV.9.7.7\endcsname{ega4.c04.thm.9-7-7}
\expandafter\def\csname sga@target@IV.9.7.8\endcsname{ega4.c04.prop.9-7-8}
\expandafter\def\csname sga@target@IV.9.7.9\endcsname{ega4.c04.cor.9-7-9}
\expandafter\def\csname sga@target@IV.9.8\endcsname{ega4.c04.subsec.9-8.key-iv-9-8}
\expandafter\def\csname sga@target@IV.9.8.1\endcsname{ega4.c04.prop.9-8-1}
\expandafter\def\csname sga@target@IV.9.8.1-proof\endcsname{ega4.c04.subsec.9-8-1-proof}
\expandafter\def\csname sga@target@IV.9.8.2\endcsname{ega4.c04.subsec.9-8-2}
\expandafter\def\csname sga@target@IV.9.8.3\endcsname{ega4.c04.prop.9-8-3}
\expandafter\def\csname sga@target@IV.9.8.3-proof\endcsname{ega4.c04.subsec.9-8-3-proof}
\expandafter\def\csname sga@target@IV.9.8.4\endcsname{ega4.c04.cor.9-8-4}
\expandafter\def\csname sga@target@IV.9.8.4-proof\endcsname{ega4.c04.item.9-8-4-proof}
\expandafter\def\csname sga@target@IV.9.8.5\endcsname{ega4.c04.prop.9-8-5}
\expandafter\def\csname sga@target@IV.9.8.5-proof\endcsname{ega4.c04.item.9-8-5-proof}
\expandafter\def\csname sga@target@IV.9.8.6\endcsname{ega4.c04.prop.9-8-6}
\expandafter\def\csname sga@target@IV.9.8.6-proof\endcsname{ega4.c04.item.9-8-6-proof}
\expandafter\def\csname sga@target@IV.9.8.7\endcsname{ega4.c04.cor.9-8-7}
\expandafter\def\csname sga@target@IV.9.8.7-proof\endcsname{ega4.c04.item.9-8-7-proof}
\expandafter\def\csname sga@target@IV.9.8.8\endcsname{ega4.c04.prop.9-8-8}
\expandafter\def\csname sga@target@IV.9.8.8-proof\endcsname{ega4.c04.item.9-8-8-proof}
\expandafter\def\csname sga@target@IV.9.8.9\endcsname{ega4.c04.rem.9-8-9}
\expandafter\def\csname sga@target@IV.9.8.9.1\endcsname{ega4.c04.item.9-8-9-1}
\expandafter\def\csname sga@target@IV.9.9\endcsname{ega4.c04.subsec.9-9.key-iv-9-9}
\expandafter\def\csname sga@target@IV.9.9.1\endcsname{ega4.c04.prop.9-9-1}
\expandafter\def\csname sga@target@IV.9.9.1-proof\endcsname{ega4.c04.item.9-9-1-proof}
\expandafter\def\csname sga@target@IV.9.9.1.1\endcsname{ega4.c04.eq.9-9-1-1}
\expandafter\def\csname sga@target@IV.9.9.2\endcsname{ega4.c04.prop.9-9-2}
\expandafter\def\csname sga@target@IV.9.9.2-proof\endcsname{ega4.c04.item.9-9-2-proof}
\expandafter\def\csname sga@target@IV.9.9.2.1\endcsname{ega4.c04.eq.9-9-2-1}
\expandafter\def\csname sga@target@IV.9.9.2.2\endcsname{ega4.c04.eq.9-9-2-2}
\expandafter\def\csname sga@target@IV.9.9.3\endcsname{ega4.c04.cor.9-9-3}
\expandafter\def\csname sga@target@IV.9.9.3-proof\endcsname{ega4.c04.statement.9-9-3-proof}
\expandafter\def\csname sga@target@IV.9.9.4\endcsname{ega4.c04.prop.9-9-4}
\expandafter\def\csname sga@target@IV.9.9.4-proof\endcsname{ega4.c04.item.9-9-4-proof}
\expandafter\def\csname sga@target@IV.9.9.5\endcsname{ega4.c04.cor.9-9-5}
\expandafter\def\csname sga@target@IV.9.9.5-proof\endcsname{ega4.c04.item.9-9-5-proof}
\expandafter\def\csname sga@target@IV.9.9.6\endcsname{ega4.c04.prop.9-9-6}
\expandafter\def\csname sga@target@IV.9.9.6-proof\endcsname{ega4.c04.item.9-9-6-proof}
\expandafter\def\csname sga@target@section:IV.1\endcsname{ega4.c04.sec.1}
\expandafter\def\csname sga@target@section:IV.10\endcsname{ega4.c04.sec.10}
\expandafter\def\csname sga@target@section:IV.11\endcsname{ega4.c04.sec.11}
\expandafter\def\csname sga@target@section:IV.12\endcsname{ega4.c04.sec.12}
\expandafter\def\csname sga@target@section:IV.13\endcsname{ega4.c04.sec.13}
\expandafter\def\csname sga@target@section:IV.14\endcsname{ega4.c04.sec.14}
\expandafter\def\csname sga@target@section:IV.15\endcsname{ega4.c04.sec.15}
\expandafter\def\csname sga@target@section:IV.16\endcsname{ega4.c04.sec.16}
\expandafter\def\csname sga@target@section:IV.17\endcsname{ega4.c04.sec.17}
\expandafter\def\csname sga@target@section:IV.18\endcsname{ega4.c04.sec.18}
\expandafter\def\csname sga@target@section:IV.19\endcsname{ega4.c04.sec.19}
\expandafter\def\csname sga@target@section:IV.2\endcsname{ega4.c04.sec.2}
\expandafter\def\csname sga@target@section:IV.21\endcsname{ega4.c04.sec.21}
\expandafter\def\csname sga@target@section:IV.3\endcsname{ega4.c04.sec.3}
\expandafter\def\csname sga@target@section:IV.4\endcsname{ega4.c04.sec.4}
\expandafter\def\csname sga@target@section:IV.5\endcsname{ega4.c04.sec.5}
\expandafter\def\csname sga@target@section:IV.6\endcsname{ega4.c04.sec.6}
\expandafter\def\csname sga@target@section:IV.7\endcsname{ega4.c04.sec.7}
\expandafter\def\csname sga@target@section:IV.8\endcsname{ega4.c04.sec.8}
\expandafter\def\csname sga@target@section:IV.9\endcsname{ega4.c04.sec.9}
\expandafter\def\csname sga@target@subsection:IV.1.1\endcsname{ega4.c04.subsec.1-1}
\expandafter\def\csname sga@target@subsection:IV.1.10\endcsname{ega4.c04.subsec.1-10}
\expandafter\def\csname sga@target@subsection:IV.1.2\endcsname{ega4.c04.subsec.1-2}
\expandafter\def\csname sga@target@subsection:IV.1.3\endcsname{ega4.c04.subsec.1-3}
\expandafter\def\csname sga@target@subsection:IV.1.4\endcsname{ega4.c04.subsec.1-4}
\expandafter\def\csname sga@target@subsection:IV.1.5\endcsname{ega4.c04.subsec.1-5}
\expandafter\def\csname sga@target@subsection:IV.1.6\endcsname{ega4.c04.subsec.1-6}
\expandafter\def\csname sga@target@subsection:IV.1.7\endcsname{ega4.c04.subsec.1-7}
\expandafter\def\csname sga@target@subsection:IV.1.8\endcsname{ega4.c04.subsec.1-8}
\expandafter\def\csname sga@target@subsection:IV.1.9\endcsname{ega4.c04.subsec.1-9}
\expandafter\def\csname sga@target@subsection:IV.10.1\endcsname{ega4.c04.subsec.10-1}
\expandafter\def\csname sga@target@subsection:IV.10.10\endcsname{ega4.c04.subsec.10-10}
\expandafter\def\csname sga@target@subsection:IV.10.2\endcsname{ega4.c04.subsec.10-2}
\expandafter\def\csname sga@target@subsection:IV.10.3\endcsname{ega4.c04.subsec.10-3}
\expandafter\def\csname sga@target@subsection:IV.10.4\endcsname{ega4.c04.subsec.10-4}
\expandafter\def\csname sga@target@subsection:IV.10.5\endcsname{ega4.c04.subsec.10-5}
\expandafter\def\csname sga@target@subsection:IV.10.6\endcsname{ega4.c04.subsec.10-6}
\expandafter\def\csname sga@target@subsection:IV.10.7\endcsname{ega4.c04.subsec.10-7}
\expandafter\def\csname sga@target@subsection:IV.10.8\endcsname{ega4.c04.subsec.10-8}
\expandafter\def\csname sga@target@subsection:IV.10.9\endcsname{ega4.c04.subsec.10-9}
\expandafter\def\csname sga@target@subsection:IV.11.1\endcsname{ega4.c04.subsec.11-1}
\expandafter\def\csname sga@target@subsection:IV.11.10\endcsname{ega4.c04.subsec.11-10}
\expandafter\def\csname sga@target@subsection:IV.11.2\endcsname{ega4.c04.subsec.11-2}
\expandafter\def\csname sga@target@subsection:IV.11.3\endcsname{ega4.c04.subsec.11-3}
\expandafter\def\csname sga@target@subsection:IV.11.4\endcsname{ega4.c04.subsec.11-4}
\expandafter\def\csname sga@target@subsection:IV.11.5\endcsname{ega4.c04.subsec.11-5}
\expandafter\def\csname sga@target@subsection:IV.11.6\endcsname{ega4.c04.subsec.11-6}
\expandafter\def\csname sga@target@subsection:IV.11.7\endcsname{ega4.c04.subsec.11-7}
\expandafter\def\csname sga@target@subsection:IV.11.8\endcsname{ega4.c04.subsec.11-8}
\expandafter\def\csname sga@target@subsection:IV.11.9\endcsname{ega4.c04.subsec.11-9}
\expandafter\def\csname sga@target@subsection:IV.12.0\endcsname{ega4.c04.subsec.12-0}
\expandafter\def\csname sga@target@subsection:IV.12.1\endcsname{ega4.c04.subsec.12-1}
\expandafter\def\csname sga@target@subsection:IV.12.2\endcsname{ega4.c04.subsec.12-2}
\expandafter\def\csname sga@target@subsection:IV.12.3\endcsname{ega4.c04.subsec.12-3}
\expandafter\def\csname sga@target@subsection:IV.13.1\endcsname{ega4.c04.subsec.13-1}
\expandafter\def\csname sga@target@subsection:IV.13.2\endcsname{ega4.c04.subsec.13-2}
\expandafter\def\csname sga@target@subsection:IV.13.3\endcsname{ega4.c04.subsec.13-3}
\expandafter\def\csname sga@target@subsection:IV.14.1\endcsname{ega4.c04.subsec.14-1}
\expandafter\def\csname sga@target@subsection:IV.14.2\endcsname{ega4.c04.subsec.14-2}
\expandafter\def\csname sga@target@subsection:IV.14.3\endcsname{ega4.c04.subsec.14-3}
\expandafter\def\csname sga@target@subsection:IV.14.4\endcsname{ega4.c04.subsec.14-4}
\expandafter\def\csname sga@target@subsection:IV.14.5\endcsname{ega4.c04.subsec.14-5}
\expandafter\def\csname sga@target@subsection:IV.15.1\endcsname{ega4.c04.subsec.15-1}
\expandafter\def\csname sga@target@subsection:IV.15.2\endcsname{ega4.c04.subsec.15-2}
\expandafter\def\csname sga@target@subsection:IV.15.3\endcsname{ega4.c04.subsec.15-3}
\expandafter\def\csname sga@target@subsection:IV.15.4\endcsname{ega4.c04.subsec.15-4}
\expandafter\def\csname sga@target@subsection:IV.15.5\endcsname{ega4.c04.subsec.15-5}
\expandafter\def\csname sga@target@subsection:IV.15.6\endcsname{ega4.c04.subsec.15-6}
\expandafter\def\csname sga@target@subsection:IV.15.7\endcsname{ega4.c04.subsec.15-7}
\expandafter\def\csname sga@target@subsection:IV.17.1\endcsname{ega4.c04.subsec.17-1}
\expandafter\def\csname sga@target@subsection:IV.17.10\endcsname{ega4.c04.subsec.17-10}
\expandafter\def\csname sga@target@subsection:IV.17.11\endcsname{ega4.c04.subsec.17-11}
\expandafter\def\csname sga@target@subsection:IV.17.12\endcsname{ega4.c04.subsec.17-12}
\expandafter\def\csname sga@target@subsection:IV.17.13\endcsname{ega4.c04.subsec.17-13}
\expandafter\def\csname sga@target@subsection:IV.17.14\endcsname{ega4.c04.subsec.17-14}
\expandafter\def\csname sga@target@subsection:IV.17.15\endcsname{ega4.c04.subsec.17-15}
\expandafter\def\csname sga@target@subsection:IV.17.16\endcsname{ega4.c04.subsec.17-16}
\expandafter\def\csname sga@target@subsection:IV.17.2\endcsname{ega4.c04.subsec.17-2}
\expandafter\def\csname sga@target@subsection:IV.17.3\endcsname{ega4.c04.subsec.17-3}
\expandafter\def\csname sga@target@subsection:IV.17.4\endcsname{ega4.c04.subsec.17-4}
\expandafter\def\csname sga@target@subsection:IV.17.5\endcsname{ega4.c04.subsec.17-5}
\expandafter\def\csname sga@target@subsection:IV.17.6\endcsname{ega4.c04.subsec.17-6}
\expandafter\def\csname sga@target@subsection:IV.17.7\endcsname{ega4.c04.subsec.17-7}
\expandafter\def\csname sga@target@subsection:IV.17.8\endcsname{ega4.c04.subsec.17-8}
\expandafter\def\csname sga@target@subsection:IV.17.9\endcsname{ega4.c04.subsec.17-9}
\expandafter\def\csname sga@target@subsection:IV.18.1\endcsname{ega4.c04.subsec.18-1}
\expandafter\def\csname sga@target@subsection:IV.18.10\endcsname{ega4.c04.subsec.18-10}
\expandafter\def\csname sga@target@subsection:IV.18.11\endcsname{ega4.c04.subsec.18-11}
\expandafter\def\csname sga@target@subsection:IV.18.12\endcsname{ega4.c04.subsec.18-12}
\expandafter\def\csname sga@target@subsection:IV.18.2\endcsname{ega4.c04.subsec.18-2}
\expandafter\def\csname sga@target@subsection:IV.18.3\endcsname{ega4.c04.subsec.18-3}
\expandafter\def\csname sga@target@subsection:IV.18.4\endcsname{ega4.c04.subsec.18-4}
\expandafter\def\csname sga@target@subsection:IV.18.5\endcsname{ega4.c04.subsec.18-5}
\expandafter\def\csname sga@target@subsection:IV.18.6\endcsname{ega4.c04.subsec.18-6}
\expandafter\def\csname sga@target@subsection:IV.18.7\endcsname{ega4.c04.subsec.18-7}
\expandafter\def\csname sga@target@subsection:IV.18.8\endcsname{ega4.c04.subsec.18-8}
\expandafter\def\csname sga@target@subsection:IV.18.9\endcsname{ega4.c04.subsec.18-9}
\expandafter\def\csname sga@target@subsection:IV.19.1\endcsname{ega4.c04.subsec.19-1}
\expandafter\def\csname sga@target@subsection:IV.19.2\endcsname{ega4.c04.subsec.19-2}
\expandafter\def\csname sga@target@subsection:IV.19.3\endcsname{ega4.c04.subsec.19-3}
\expandafter\def\csname sga@target@subsection:IV.19.4\endcsname{ega4.c04.subsec.19-4}
\expandafter\def\csname sga@target@subsection:IV.19.5\endcsname{ega4.c04.subsec.19-5}
\expandafter\def\csname sga@target@subsection:IV.19.6\endcsname{ega4.c04.subsec.19-6}
\expandafter\def\csname sga@target@subsection:IV.19.7\endcsname{ega4.c04.subsec.19-7}
\expandafter\def\csname sga@target@subsection:IV.19.8\endcsname{ega4.c04.subsec.19-8}
\expandafter\def\csname sga@target@subsection:IV.19.9\endcsname{ega4.c04.subsec.19-9}
\expandafter\def\csname sga@target@subsection:IV.2.1\endcsname{ega4.c04.subsec.2-1}
\expandafter\def\csname sga@target@subsection:IV.2.2\endcsname{ega4.c04.subsec.2-2}
\expandafter\def\csname sga@target@subsection:IV.2.3\endcsname{ega4.c04.subsec.2-3}
\expandafter\def\csname sga@target@subsection:IV.2.4\endcsname{ega4.c04.subsec.2-4}
\expandafter\def\csname sga@target@subsection:IV.2.5\endcsname{ega4.c04.subsec.2-5}
\expandafter\def\csname sga@target@subsection:IV.2.6\endcsname{ega4.c04.subsec.2-6}
\expandafter\def\csname sga@target@subsection:IV.2.7\endcsname{ega4.c04.subsec.2-7}
\expandafter\def\csname sga@target@subsection:IV.2.8\endcsname{ega4.c04.subsec.2-8}
\expandafter\def\csname sga@target@subsection:IV.21.1\endcsname{ega4.c04.subsec.21-1}
\expandafter\def\csname sga@target@subsection:IV.21.10\endcsname{ega4.c04.subsec.21-10}
\expandafter\def\csname sga@target@subsection:IV.21.11\endcsname{ega4.c04.subsec.21-11}
\expandafter\def\csname sga@target@subsection:IV.21.12\endcsname{ega4.c04.subsec.21-12}
\expandafter\def\csname sga@target@subsection:IV.21.13\endcsname{ega4.c04.subsec.21-13}
\expandafter\def\csname sga@target@subsection:IV.21.14\endcsname{ega4.c04.subsec.21-14}
\expandafter\def\csname sga@target@subsection:IV.21.15\endcsname{ega4.c04.subsec.21-15}
\expandafter\def\csname sga@target@subsection:IV.21.2\endcsname{ega4.c04.subsec.21-2}
\expandafter\def\csname sga@target@subsection:IV.21.3\endcsname{ega4.c04.subsec.21-3}
\expandafter\def\csname sga@target@subsection:IV.21.4\endcsname{ega4.c04.subsec.21-4}
\expandafter\def\csname sga@target@subsection:IV.21.5\endcsname{ega4.c04.subsec.21-5}
\expandafter\def\csname sga@target@subsection:IV.21.6\endcsname{ega4.c04.subsec.21-6}
\expandafter\def\csname sga@target@subsection:IV.21.7\endcsname{ega4.c04.subsec.21-7}
\expandafter\def\csname sga@target@subsection:IV.21.8\endcsname{ega4.c04.subsec.21-8}
\expandafter\def\csname sga@target@subsection:IV.21.9\endcsname{ega4.c04.subsec.21-9}
\expandafter\def\csname sga@target@subsection:IV.3.1\endcsname{ega4.c04.subsec.3-1}
\expandafter\def\csname sga@target@subsection:IV.3.2\endcsname{ega4.c04.subsec.3-2}
\expandafter\def\csname sga@target@subsection:IV.3.3\endcsname{ega4.c04.subsec.3-3}
\expandafter\def\csname sga@target@subsection:IV.3.4\endcsname{ega4.c04.subsec.3-4}
\expandafter\def\csname sga@target@subsection:IV.4.1\endcsname{ega4.c04.subsec.4-1}
\expandafter\def\csname sga@target@subsection:IV.4.2\endcsname{ega4.c04.subsec.4-2}
\expandafter\def\csname sga@target@subsection:IV.4.3\endcsname{ega4.c04.subsec.4-3}
\expandafter\def\csname sga@target@subsection:IV.4.4\endcsname{ega4.c04.subsec.4-4}
\expandafter\def\csname sga@target@subsection:IV.4.5\endcsname{ega4.c04.subsec.4-5}
\expandafter\def\csname sga@target@subsection:IV.4.6\endcsname{ega4.c04.subsec.4-6}
\expandafter\def\csname sga@target@subsection:IV.4.7\endcsname{ega4.c04.subsec.4-7}
\expandafter\def\csname sga@target@subsection:IV.4.8\endcsname{ega4.c04.subsec.4-8}
\expandafter\def\csname sga@target@subsection:IV.4.9\endcsname{ega4.c04.subsec.4-9}
\expandafter\def\csname sga@target@subsection:IV.5.1\endcsname{ega4.c04.subsec.5-1}
\expandafter\def\csname sga@target@subsection:IV.5.10\endcsname{ega4.c04.subsec.5-10}
\expandafter\def\csname sga@target@subsection:IV.5.11\endcsname{ega4.c04.subsec.5-11}
\expandafter\def\csname sga@target@subsection:IV.5.12\endcsname{ega4.c04.subsec.5-12}
\expandafter\def\csname sga@target@subsection:IV.5.13\endcsname{ega4.c04.subsec.5-13}
\expandafter\def\csname sga@target@subsection:IV.5.2\endcsname{ega4.c04.subsec.5-2}
\expandafter\def\csname sga@target@subsection:IV.5.3\endcsname{ega4.c04.subsec.5-3}
\expandafter\def\csname sga@target@subsection:IV.5.4\endcsname{ega4.c04.subsec.5-4}
\expandafter\def\csname sga@target@subsection:IV.5.5\endcsname{ega4.c04.subsec.5-5}
\expandafter\def\csname sga@target@subsection:IV.5.6\endcsname{ega4.c04.subsec.5-6}
\expandafter\def\csname sga@target@subsection:IV.5.7\endcsname{ega4.c04.subsec.5-7}
\expandafter\def\csname sga@target@subsection:IV.5.8\endcsname{ega4.c04.subsec.5-8}
\expandafter\def\csname sga@target@subsection:IV.5.9\endcsname{ega4.c04.subsec.5-9}
\expandafter\def\csname sga@target@subsection:IV.6.1\endcsname{ega4.c04.subsec.6-1}
\expandafter\def\csname sga@target@subsection:IV.6.10\endcsname{ega4.c04.subsec.6-10}
\expandafter\def\csname sga@target@subsection:IV.6.11\endcsname{ega4.c04.subsec.6-11}
\expandafter\def\csname sga@target@subsection:IV.6.12\endcsname{ega4.c04.subsec.6-12}
\expandafter\def\csname sga@target@subsection:IV.6.13\endcsname{ega4.c04.subsec.6-13}
\expandafter\def\csname sga@target@subsection:IV.6.14\endcsname{ega4.c04.subsec.6-14}
\expandafter\def\csname sga@target@subsection:IV.6.15\endcsname{ega4.c04.subsec.6-15}
\expandafter\def\csname sga@target@subsection:IV.6.2\endcsname{ega4.c04.subsec.6-2}
\expandafter\def\csname sga@target@subsection:IV.6.3\endcsname{ega4.c04.subsec.6-3}
\expandafter\def\csname sga@target@subsection:IV.6.4\endcsname{ega4.c04.subsec.6-4}
\expandafter\def\csname sga@target@subsection:IV.6.5\endcsname{ega4.c04.subsec.6-5}
\expandafter\def\csname sga@target@subsection:IV.6.6\endcsname{ega4.c04.subsec.6-6}
\expandafter\def\csname sga@target@subsection:IV.6.7\endcsname{ega4.c04.subsec.6-7}
\expandafter\def\csname sga@target@subsection:IV.6.8\endcsname{ega4.c04.subsec.6-8}
\expandafter\def\csname sga@target@subsection:IV.6.9\endcsname{ega4.c04.subsec.6-9}
\expandafter\def\csname sga@target@subsection:IV.7.1\endcsname{ega4.c04.subsec.7-1}
\expandafter\def\csname sga@target@subsection:IV.7.2\endcsname{ega4.c04.subsec.7-2}
\expandafter\def\csname sga@target@subsection:IV.7.3\endcsname{ega4.c04.subsec.7-3}
\expandafter\def\csname sga@target@subsection:IV.7.4\endcsname{ega4.c04.subsec.7-4}
\expandafter\def\csname sga@target@subsection:IV.7.5\endcsname{ega4.c04.subsec.7-5}
\expandafter\def\csname sga@target@subsection:IV.7.6\endcsname{ega4.c04.subsec.7-6}
\expandafter\def\csname sga@target@subsection:IV.7.7\endcsname{ega4.c04.subsec.7-7}
\expandafter\def\csname sga@target@subsection:IV.7.8\endcsname{ega4.c04.subsec.7-8}
\expandafter\def\csname sga@target@subsection:IV.7.9\endcsname{ega4.c04.subsec.7-9}
\expandafter\def\csname sga@target@subsection:IV.8.1\endcsname{ega4.c04.subsec.8-1}
\expandafter\def\csname sga@target@subsection:IV.8.10\endcsname{ega4.c04.subsec.8-10}
\expandafter\def\csname sga@target@subsection:IV.8.11\endcsname{ega4.c04.subsec.8-11}
\expandafter\def\csname sga@target@subsection:IV.8.12\endcsname{ega4.c04.subsec.8-12}
\expandafter\def\csname sga@target@subsection:IV.8.13\endcsname{ega4.c04.subsec.8-13}
\expandafter\def\csname sga@target@subsection:IV.8.14\endcsname{ega4.c04.subsec.8-14}
\expandafter\def\csname sga@target@subsection:IV.8.2\endcsname{ega4.c04.subsec.8-2}
\expandafter\def\csname sga@target@subsection:IV.8.3\endcsname{ega4.c04.subsec.8-3}
\expandafter\def\csname sga@target@subsection:IV.8.4\endcsname{ega4.c04.subsec.8-4}
\expandafter\def\csname sga@target@subsection:IV.8.5\endcsname{ega4.c04.subsec.8-5}
\expandafter\def\csname sga@target@subsection:IV.8.6\endcsname{ega4.c04.subsec.8-6}
\expandafter\def\csname sga@target@subsection:IV.8.7\endcsname{ega4.c04.subsec.8-7}
\expandafter\def\csname sga@target@subsection:IV.8.8\endcsname{ega4.c04.subsec.8-8}
\expandafter\def\csname sga@target@subsection:IV.8.9\endcsname{ega4.c04.subsec.8-9}
\expandafter\def\csname sga@target@subsection:IV.9.1\endcsname{ega4.c04.subsec.9-1}
\expandafter\def\csname sga@target@subsection:IV.9.2\endcsname{ega4.c04.subsec.9-2}
\expandafter\def\csname sga@target@subsection:IV.9.3\endcsname{ega4.c04.subsec.9-3}
\expandafter\def\csname sga@target@subsection:IV.9.4\endcsname{ega4.c04.subsec.9-4}
\expandafter\def\csname sga@target@subsection:IV.9.5\endcsname{ega4.c04.subsec.9-5}
\expandafter\def\csname sga@target@subsection:IV.9.6\endcsname{ega4.c04.subsec.9-6}
\expandafter\def\csname sga@target@subsection:IV.9.7\endcsname{ega4.c04.subsec.9-7}
\expandafter\def\csname sga@target@subsection:IV.9.8\endcsname{ega4.c04.subsec.9-8}
\expandafter\def\csname sga@target@subsection:IV.9.9\endcsname{ega4.c04.subsec.9-9}
\expandafter\def\csname sga@bib@I-4\endcsname{ega4.reader.bib.i-4}
\expandafter\def\csname sga@bib@IV-32\endcsname{ega4.reader.bib.iv-32}
\expandafter\def\csname sga@bib@IV-34\endcsname{ega4.reader.bib.iv-34}
\expandafter\def\csname sga@bib@IV-37\endcsname{ega4.reader.bib.iv-37}
\expandafter\def\csname sga@bib@IV-41\endcsname{ega4.reader.bib.iv-41}
\expandafter\def\csname sga@bib@IV-42\endcsname{ega4.reader.bib.iv-42}
\expandafter\def\csname sga@bib@IV-43\endcsname{ega4.reader.bib.iv-43}
\expandafter\def\csname sga@bib@IV-44\endcsname{ega4.reader.bib.iv-44}
\expandafter\def\csname sga@bib@SGA1\endcsname{ega4.reader.bib.sga1}
\let\sgaOriginalLabel\label
\renewcommand{\label}[1]{%
  \sgaOriginalLabel{#1}%
  \ifcsname sga@target@#1\endcsname
    \expandafter\sgaStableTarget\expandafter{\csname sga@target@#1\endcsname}%
  \fi
}
\let\sgaOriginalBibitem\bibitem
\newcommand{\sgaBibTarget}[1]{%
  \ifcsname sga@bib@#1\endcsname
    \expandafter\sgaStableTarget\expandafter{\csname sga@bib@#1\endcsname}%
  \fi
}
\newcommand{\sgaBibitemPlain}[1]{%
  \sgaOriginalBibitem{#1}\sgaBibTarget{#1}}
\newcommand{\sgaBibitemOptional}[2][]{%
  \sgaOriginalBibitem[#1]{#2}\sgaBibTarget{#2}}
\renewcommand{\bibitem}{%
  \@ifnextchar[\sgaBibitemOptional\sgaBibitemPlain}
\makeatother
% END NATIVE INPUT EGA4_reference_v2_registry
% BEGIN NATIVE INPUT EGA4_reference_v2_links
% Generated stable-link dispatch for the EGA IV reader.
% Stable destinations are emitted by EGA4_reference_v2_registry.tex.
\makeatletter
\expandafter\def\csname sga@target@IV.1.1.1\endcsname{ega4.c04.def.1-1-1}
\expandafter\def\csname sga@target@IV.1.1.2\endcsname{ega4.c04.prop.1-1-2}
\expandafter\def\csname sga@target@IV.1.1.3\endcsname{ega4.c04.prop.1-1-3}
\expandafter\def\csname sga@target@IV.1.1.4\endcsname{ega4.c04.cor.1-1-4}
\expandafter\def\csname sga@target@IV.1.1.5\endcsname{ega4.c04.prop.1-1-5}
\expandafter\def\csname sga@target@IV.1.1.6\endcsname{ega4.c04.prop.1-1-6}
\expandafter\def\csname sga@target@IV.1.10.1\endcsname{ega4.c04.thm.1-10-1}
\expandafter\def\csname sga@target@IV.1.10.2\endcsname{ega4.c04.cor.1-10-2}
\expandafter\def\csname sga@target@IV.1.10.3\endcsname{ega4.c04.prop.1-10-3}
\expandafter\def\csname sga@target@IV.1.10.4\endcsname{ega4.c04.cor.1-10-4}
\expandafter\def\csname sga@target@IV.1.2.1\endcsname{ega4.c04.def.1-2-1}
\expandafter\def\csname sga@target@IV.1.2.2\endcsname{ega4.c04.prop.1-2-2}
\expandafter\def\csname sga@target@IV.1.2.3\endcsname{ega4.c04.cor.1-2-3}
\expandafter\def\csname sga@target@IV.1.2.4\endcsname{ega4.c04.prop.1-2-4}
\expandafter\def\csname sga@target@IV.1.2.4.1\endcsname{ega4.c04.eq.1-2-4-1}
\expandafter\def\csname sga@target@IV.1.2.5\endcsname{ega4.c04.prop.1-2-5}
\expandafter\def\csname sga@target@IV.1.2.6\endcsname{ega4.c04.prop.1-2-6}
\expandafter\def\csname sga@target@IV.1.2.7\endcsname{ega4.c04.prop.1-2-7}
\expandafter\def\csname sga@target@IV.1.2.8\endcsname{ega4.c04.cor.1-2-8}
\expandafter\def\csname sga@target@IV.1.2.9\endcsname{ega4.c04.prop.1-2-9}
\expandafter\def\csname sga@target@IV.1.3.1\endcsname{ega4.c04.statement.1-3-1}
\expandafter\def\csname sga@target@IV.1.3.10\endcsname{ega4.c04.statement.1-3-10}
\expandafter\def\csname sga@target@IV.1.3.11\endcsname{ega4.c04.prop.1-3-11}
\expandafter\def\csname sga@target@IV.1.3.2\endcsname{ega4.c04.def.1-3-2}
\expandafter\def\csname sga@target@IV.1.3.3\endcsname{ega4.c04.prop.1-3-3}
\expandafter\def\csname sga@target@IV.1.3.4\endcsname{ega4.c04.prop.1-3-4}
\expandafter\def\csname sga@target@IV.1.3.5\endcsname{ega4.c04.cor.1-3-5}
\expandafter\def\csname sga@target@IV.1.3.6\endcsname{ega4.c04.prop.1-3-6}
\expandafter\def\csname sga@target@IV.1.3.7\endcsname{ega4.c04.prop.1-3-7}
\expandafter\def\csname sga@target@IV.1.3.8\endcsname{ega4.c04.statement.1-3-8}
\expandafter\def\csname sga@target@IV.1.3.9\endcsname{ega4.c04.prop.1-3-9}
\expandafter\def\csname sga@target@IV.1.4.1\endcsname{ega4.c04.statement.1-4-1}
\expandafter\def\csname sga@target@IV.1.4.2\endcsname{ega4.c04.def.1-4-2}
\expandafter\def\csname sga@target@IV.1.4.3\endcsname{ega4.c04.prop.1-4-3}
\expandafter\def\csname sga@target@IV.1.4.3.1\endcsname{ega4.c04.cor.1-4-3-1}
\expandafter\def\csname sga@target@IV.1.4.4\endcsname{ega4.c04.prop.1-4-4}
\expandafter\def\csname sga@target@IV.1.4.4.1\endcsname{ega4.c04.lem.1-4-4-1}
\expandafter\def\csname sga@target@IV.1.4.5\endcsname{ega4.c04.cor.1-4-5}
\expandafter\def\csname sga@target@IV.1.4.6\endcsname{ega4.c04.prop.1-4-6}
\expandafter\def\csname sga@target@IV.1.4.7\endcsname{ega4.c04.prop.1-4-7}
\expandafter\def\csname sga@target@IV.1.4.7.1\endcsname{ega4.c04.lem.1-4-7-1}
\expandafter\def\csname sga@target@IV.1.5.1\endcsname{ega4.c04.prop.1-5-1}
\expandafter\def\csname sga@target@IV.1.5.2\endcsname{ega4.c04.def.1-5-2}
\expandafter\def\csname sga@target@IV.1.5.3\endcsname{ega4.c04.prop.1-5-3}
\expandafter\def\csname sga@target@IV.1.5.4\endcsname{ega4.c04.prop.1-5-4}
\expandafter\def\csname sga@target@IV.1.5.5\endcsname{ega4.c04.cor.1-5-5}
\expandafter\def\csname sga@target@IV.1.5.6\endcsname{ega4.c04.cor.1-5-6}
\expandafter\def\csname sga@target@IV.1.5.7\endcsname{ega4.c04.prop.1-5-7}
\expandafter\def\csname sga@target@IV.1.6.1\endcsname{ega4.c04.def.1-6-1}
\expandafter\def\csname sga@target@IV.1.6.2\endcsname{ega4.c04.prop.1-6-2}
\expandafter\def\csname sga@target@IV.1.6.3\endcsname{ega4.c04.cor.1-6-3}
\expandafter\def\csname sga@target@IV.1.6.4\endcsname{ega4.c04.rem.1-6-4}
\expandafter\def\csname sga@target@IV.1.6.5\endcsname{ega4.c04.prop.1-6-5}
\expandafter\def\csname sga@target@IV.1.7.1\endcsname{ega4.c04.statement.1-7-1}
\expandafter\def\csname sga@target@IV.1.7.10\endcsname{ega4.c04.statement.1-7-10}
\expandafter\def\csname sga@target@IV.1.7.11\endcsname{ega4.c04.statement.1-7-11}
\expandafter\def\csname sga@target@IV.1.7.12\endcsname{ega4.c04.statement.1-7-12}
\expandafter\def\csname sga@target@IV.1.7.13\endcsname{ega4.c04.statement.1-7-13}
\expandafter\def\csname sga@target@IV.1.7.14\endcsname{ega4.c04.statement.1-7-14}
\expandafter\def\csname sga@target@IV.1.7.15\endcsname{ega4.c04.statement.1-7-15}
\expandafter\def\csname sga@target@IV.1.7.16\endcsname{ega4.c04.statement.1-7-16}
\expandafter\def\csname sga@target@IV.1.7.17\endcsname{ega4.c04.statement.1-7-17}
\expandafter\def\csname sga@target@IV.1.7.18\endcsname{ega4.c04.statement.1-7-18}
\expandafter\def\csname sga@target@IV.1.7.19\endcsname{ega4.c04.statement.1-7-19}
\expandafter\def\csname sga@target@IV.1.7.2\endcsname{ega4.c04.statement.1-7-2}
\expandafter\def\csname sga@target@IV.1.7.20\endcsname{ega4.c04.statement.1-7-20}
\expandafter\def\csname sga@target@IV.1.7.21\endcsname{ega4.c04.statement.1-7-21}
\expandafter\def\csname sga@target@IV.1.7.22\endcsname{ega4.c04.statement.1-7-22}
\expandafter\def\csname sga@target@IV.1.7.3\endcsname{ega4.c04.statement.1-7-3}
\expandafter\def\csname sga@target@IV.1.7.4\endcsname{ega4.c04.statement.1-7-4}
\expandafter\def\csname sga@target@IV.1.7.5\endcsname{ega4.c04.statement.1-7-5}
\expandafter\def\csname sga@target@IV.1.7.6\endcsname{ega4.c04.statement.1-7-6}
\expandafter\def\csname sga@target@IV.1.7.7\endcsname{ega4.c04.statement.1-7-7}
\expandafter\def\csname sga@target@IV.1.7.8\endcsname{ega4.c04.statement.1-7-8}
\expandafter\def\csname sga@target@IV.1.7.9\endcsname{ega4.c04.statement.1-7-9}
\expandafter\def\csname sga@target@IV.1.8.1\endcsname{ega4.c04.statement.1-8-1}
\expandafter\def\csname sga@target@IV.1.8.2\endcsname{ega4.c04.prop.1-8-2}
\expandafter\def\csname sga@target@IV.1.8.3\endcsname{ega4.c04.lem.1-8-3}
\expandafter\def\csname sga@target@IV.1.8.4\endcsname{ega4.c04.thm.1-8-4}
\expandafter\def\csname sga@target@IV.1.8.4.1\endcsname{ega4.c04.lem.1-8-4-1}
\expandafter\def\csname sga@target@IV.1.8.4.2\endcsname{ega4.c04.lem.1-8-4-2}
\expandafter\def\csname sga@target@IV.1.8.5\endcsname{ega4.c04.cor.1-8-5}
\expandafter\def\csname sga@target@IV.1.8.5.1\endcsname{ega4.c04.lem.1-8-5-1}
\expandafter\def\csname sga@target@IV.1.8.6\endcsname{ega4.c04.cor.1-8-6}
\expandafter\def\csname sga@target@IV.1.8.7\endcsname{ega4.c04.prop.1-8-7}
\expandafter\def\csname sga@target@IV.1.8.7.1\endcsname{ega4.c04.lem.1-8-7-1}
\expandafter\def\csname sga@target@IV.1.9.1\endcsname{ega4.c04.lem.1-9-1}
\expandafter\def\csname sga@target@IV.1.9.10\endcsname{ega4.c04.prop.1-9-10}
\expandafter\def\csname sga@target@IV.1.9.11\endcsname{ega4.c04.prop.1-9-11}
\expandafter\def\csname sga@target@IV.1.9.12\endcsname{ega4.c04.cor.1-9-12}
\expandafter\def\csname sga@target@IV.1.9.13\endcsname{ega4.c04.statement.1-9-13}
\expandafter\def\csname sga@target@IV.1.9.14\endcsname{ega4.c04.prop.1-9-14}
\expandafter\def\csname sga@target@IV.1.9.15\endcsname{ega4.c04.prop.1-9-15}
\expandafter\def\csname sga@target@IV.1.9.16\endcsname{ega4.c04.statement.1-9-16}
\expandafter\def\csname sga@target@IV.1.9.2\endcsname{ega4.c04.prop.1-9-2}
\expandafter\def\csname sga@target@IV.1.9.2.1\endcsname{ega4.c04.statement.1-9-2-1}
\expandafter\def\csname sga@target@IV.1.9.2.i\endcsname{ega4.c04.statement.1-9-2-i}
\expandafter\def\csname sga@target@IV.1.9.2.ii\endcsname{ega4.c04.statement.1-9-2-ii}
\expandafter\def\csname sga@target@IV.1.9.3\endcsname{ega4.c04.prop.1-9-3}
\expandafter\def\csname sga@target@IV.1.9.3.1\endcsname{ega4.c04.lem.1-9-3-1}
\expandafter\def\csname sga@target@IV.1.9.3.2\endcsname{ega4.c04.lem.1-9-3-2}
\expandafter\def\csname sga@target@IV.1.9.4\endcsname{ega4.c04.def.1-9-4}
\expandafter\def\csname sga@target@IV.1.9.5\endcsname{ega4.c04.prop.1-9-5}
\expandafter\def\csname sga@target@IV.1.9.6\endcsname{ega4.c04.statement.1-9-6}
\expandafter\def\csname sga@target@IV.1.9.7\endcsname{ega4.c04.thm.1-9-7}
\expandafter\def\csname sga@target@IV.1.9.8\endcsname{ega4.c04.cor.1-9-8}
\expandafter\def\csname sga@target@IV.1.9.9\endcsname{ega4.c04.cor.1-9-9}
\expandafter\def\csname sga@target@IV.10.1.1\endcsname{ega4.c04.statement.10-1-1}
\expandafter\def\csname sga@target@IV.10.1.2\endcsname{ega4.c04.prop.10-1-2}
\expandafter\def\csname sga@target@IV.10.1.3\endcsname{ega4.c04.def.10-1-3}
\expandafter\def\csname sga@target@IV.10.1.4\endcsname{ega4.c04.cor.10-1-4}
\expandafter\def\csname sga@target@IV.10.10.1\endcsname{ega4.c04.statement.10-10-1}
\expandafter\def\csname sga@target@IV.10.10.2\endcsname{ega4.c04.statement.10-10-2}
\expandafter\def\csname sga@target@IV.10.10.3\endcsname{ega4.c04.statement.10-10-3}
\expandafter\def\csname sga@target@IV.10.10.4\endcsname{ega4.c04.statement.10-10-4}
\expandafter\def\csname sga@target@IV.10.10.5\endcsname{ega4.c04.statement.10-10-5}
\expandafter\def\csname sga@target@IV.10.2.1\endcsname{ega4.c04.prop.10-2-1}
\expandafter\def\csname sga@target@IV.10.2.2\endcsname{ega4.c04.def.10-2-2}
\expandafter\def\csname sga@target@IV.10.2.3\endcsname{ega4.c04.cor.10-2-3}
\expandafter\def\csname sga@target@IV.10.2.4\endcsname{ega4.c04.cor.10-2-4}
\expandafter\def\csname sga@target@IV.10.2.5\endcsname{ega4.c04.cor.10-2-5}
\expandafter\def\csname sga@target@IV.10.2.6\endcsname{ega4.c04.cor.10-2-6}
\expandafter\def\csname sga@target@IV.10.2.7\endcsname{ega4.c04.prop.10-2-7}
\expandafter\def\csname sga@target@IV.10.2.8\endcsname{ega4.c04.item.10-2-8}
\expandafter\def\csname sga@target@IV.10.2.8.1\endcsname{ega4.c04.eq.10-2-8-1}
\expandafter\def\csname sga@target@IV.10.3.1\endcsname{ega4.c04.def.10-3-1}
\expandafter\def\csname sga@target@IV.10.3.2\endcsname{ega4.c04.prop.10-3-2}
\expandafter\def\csname sga@target@IV.10.3.3\endcsname{ega4.c04.prop.10-3-3}
\expandafter\def\csname sga@target@IV.10.4.1\endcsname{ega4.c04.def.10-4-1}
\expandafter\def\csname sga@target@IV.10.4.10\endcsname{ega4.c04.statement.10-4-10}
\expandafter\def\csname sga@target@IV.10.4.11\endcsname{ega4.c04.statement.10-4-11}
\expandafter\def\csname sga@target@IV.10.4.11.1\endcsname{ega4.c04.lem.10-4-11-1}
\expandafter\def\csname sga@target@IV.10.4.12\endcsname{ega4.c04.cor.10-4-12}
\expandafter\def\csname sga@target@IV.10.4.12.1\endcsname{ega4.c04.sec.10-4-12-1}
\expandafter\def\csname sga@target@IV.10.4.12.1.errata\endcsname{ega4.c04.eq.10-4-12-1-errata}
\expandafter\def\csname sga@target@IV.10.4.13\endcsname{ega4.c04.cor.10-4-13}
\expandafter\def\csname sga@target@IV.10.4.14\endcsname{ega4.c04.rem.10-4-14}
\expandafter\def\csname sga@target@IV.10.4.2\endcsname{ega4.c04.prop.10-4-2}
\expandafter\def\csname sga@target@IV.10.4.3\endcsname{ega4.c04.statement.10-4-3}
\expandafter\def\csname sga@target@IV.10.4.4\endcsname{ega4.c04.prop.10-4-4}
\expandafter\def\csname sga@target@IV.10.4.5\endcsname{ega4.c04.prop.10-4-5}
\expandafter\def\csname sga@target@IV.10.4.5.1\endcsname{ega4.c04.lem.10-4-5-1}
\expandafter\def\csname sga@target@IV.10.4.6\endcsname{ega4.c04.cor.10-4-6}
\expandafter\def\csname sga@target@IV.10.4.7\endcsname{ega4.c04.cor.10-4-7}
\expandafter\def\csname sga@target@IV.10.4.8\endcsname{ega4.c04.cor.10-4-8}
\expandafter\def\csname sga@target@IV.10.4.9\endcsname{ega4.c04.statement.10-4-9}
\expandafter\def\csname sga@target@IV.10.5.1\endcsname{ega4.c04.prop.10-5-1}
\expandafter\def\csname sga@target@IV.10.5.10\endcsname{ega4.c04.prop.10-5-10}
\expandafter\def\csname sga@target@IV.10.5.2\endcsname{ega4.c04.cor.10-5-2}
\expandafter\def\csname sga@target@IV.10.5.3\endcsname{ega4.c04.cor.10-5-3}
\expandafter\def\csname sga@target@IV.10.5.4\endcsname{ega4.c04.rem.10-5-4}
\expandafter\def\csname sga@target@IV.10.5.5\endcsname{ega4.c04.prop.10-5-5}
\expandafter\def\csname sga@target@IV.10.5.6\endcsname{ega4.c04.rem.10-5-6}
\expandafter\def\csname sga@target@IV.10.5.7\endcsname{ega4.c04.prop.10-5-7}
\expandafter\def\csname sga@target@IV.10.5.8\endcsname{ega4.c04.cor.10-5-8}
\expandafter\def\csname sga@target@IV.10.5.9\endcsname{ega4.c04.cor.10-5-9}
\expandafter\def\csname sga@target@IV.10.6.1\endcsname{ega4.c04.prop.10-6-1}
\expandafter\def\csname sga@target@IV.10.6.1.1\endcsname{ega4.c04.eq.10-6-1-1}
\expandafter\def\csname sga@target@IV.10.6.1.2\endcsname{ega4.c04.eq.10-6-1-2}
\expandafter\def\csname sga@target@IV.10.6.1.3\endcsname{ega4.c04.eq.10-6-1-3}
\expandafter\def\csname sga@target@IV.10.6.1.4\endcsname{ega4.c04.eq.10-6-1-4}
\expandafter\def\csname sga@target@IV.10.6.2\endcsname{ega4.c04.cor.10-6-2}
\expandafter\def\csname sga@target@IV.10.6.3\endcsname{ega4.c04.prop.10-6-3}
\expandafter\def\csname sga@target@IV.10.6.3.1\endcsname{ega4.c04.eq.10-6-3-1}
\expandafter\def\csname sga@target@IV.10.6.3.2\endcsname{ega4.c04.eq.10-6-3-2}
\expandafter\def\csname sga@target@IV.10.6.3.3\endcsname{ega4.c04.eq.10-6-3-3}
\expandafter\def\csname sga@target@IV.10.6.4\endcsname{ega4.c04.cor.10-6-4}
\expandafter\def\csname sga@target@IV.10.6.4.1\endcsname{ega4.c04.eq.10-6-4-1}
\expandafter\def\csname sga@target@IV.10.7.1\endcsname{ega4.c04.statement.10-7-1}
\expandafter\def\csname sga@target@IV.10.7.2\endcsname{ega4.c04.statement.10-7-2}
\expandafter\def\csname sga@target@IV.10.7.3\endcsname{ega4.c04.statement.10-7-3}
\expandafter\def\csname sga@target@IV.10.8.1\endcsname{ega4.c04.def.10-8-1}
\expandafter\def\csname sga@target@IV.10.8.1.1\endcsname{ega4.c04.eq.10-8-1-1}
\expandafter\def\csname sga@target@IV.10.8.1.2\endcsname{ega4.c04.eq.10-8-1-2}
\expandafter\def\csname sga@target@IV.10.8.2\endcsname{ega4.c04.statement.10-8-2}
\expandafter\def\csname sga@target@IV.10.8.3\endcsname{ega4.c04.lem.10-8-3}
\expandafter\def\csname sga@target@IV.10.8.3.1\endcsname{ega4.c04.eq.10-8-3-1}
\expandafter\def\csname sga@target@IV.10.8.4\endcsname{ega4.c04.cor.10-8-4}
\expandafter\def\csname sga@target@IV.10.8.5\endcsname{ega4.c04.prop.10-8-5}
\expandafter\def\csname sga@target@IV.10.8.6\endcsname{ega4.c04.cor.10-8-6}
\expandafter\def\csname sga@target@IV.10.8.7\endcsname{ega4.c04.prop.10-8-7}
\expandafter\def\csname sga@target@IV.10.8.7.1\endcsname{ega4.c04.eq.10-8-7-1}
\expandafter\def\csname sga@target@IV.10.8.8\endcsname{ega4.c04.cor.10-8-8}
\expandafter\def\csname sga@target@IV.10.9.1\endcsname{ega4.c04.statement.10-9-1}
\expandafter\def\csname sga@target@IV.10.9.2\endcsname{ega4.c04.statement.10-9-2}
\expandafter\def\csname sga@target@IV.10.9.3\endcsname{ega4.c04.statement.10-9-3}
\expandafter\def\csname sga@target@IV.10.9.4\endcsname{ega4.c04.prop.10-9-4}
\expandafter\def\csname sga@target@IV.10.9.5\endcsname{ega4.c04.statement.10-9-5}
\expandafter\def\csname sga@target@IV.10.9.6\endcsname{ega4.c04.prop.10-9-6}
\expandafter\def\csname sga@target@IV.11.1.1\endcsname{ega4.c04.thm.11-1-1}
\expandafter\def\csname sga@target@IV.11.1.1.1\endcsname{ega4.c04.sec.11-1-1-1}
\expandafter\def\csname sga@target@IV.11.1.1.2\endcsname{ega4.c04.lem.11-1-1-2}
\expandafter\def\csname sga@target@IV.11.1.2\endcsname{ega4.c04.cor.11-1-2}
\expandafter\def\csname sga@target@IV.11.1.3\endcsname{ega4.c04.cor.11-1-3}
\expandafter\def\csname sga@target@IV.11.1.4\endcsname{ega4.c04.cor.11-1-4}
\expandafter\def\csname sga@target@IV.11.1.5\endcsname{ega4.c04.cor.11-1-5}
\expandafter\def\csname sga@target@IV.11.1.6\endcsname{ega4.c04.cor.11-1-6}
\expandafter\def\csname sga@target@IV.11.10.1\endcsname{ega4.c04.prop.11-10-1}
\expandafter\def\csname sga@target@IV.11.10.1.1\endcsname{ega4.c04.eq.11-10-1-1}
\expandafter\def\csname sga@target@IV.11.10.1.2\endcsname{ega4.c04.eq.11-10-1-2}
\expandafter\def\csname sga@target@IV.11.10.10\endcsname{ega4.c04.prop.11-10-10}
\expandafter\def\csname sga@target@IV.11.10.2\endcsname{ega4.c04.def.11-10-2}
\expandafter\def\csname sga@target@IV.11.10.3\endcsname{ega4.c04.item.11-10-3}
\expandafter\def\csname sga@target@IV.11.10.4\endcsname{ega4.c04.prop.11-10-4}
\expandafter\def\csname sga@target@IV.11.10.5\endcsname{ega4.c04.thm.11-10-5}
\expandafter\def\csname sga@target@IV.11.10.6\endcsname{ega4.c04.prop.11-10-6}
\expandafter\def\csname sga@target@IV.11.10.7\endcsname{ega4.c04.cor.11-10-7}
\expandafter\def\csname sga@target@IV.11.10.8\endcsname{ega4.c04.def.11-10-8}
\expandafter\def\csname sga@target@IV.11.10.9\endcsname{ega4.c04.thm.11-10-9}
\expandafter\def\csname sga@target@IV.11.2.1\endcsname{ega4.c04.statement.11-2-1}
\expandafter\def\csname sga@target@IV.11.2.1.1\endcsname{ega4.c04.eq.11-2-1-1}
\expandafter\def\csname sga@target@IV.11.2.1.2\endcsname{ega4.c04.eq.11-2-1-2}
\expandafter\def\csname sga@target@IV.11.2.1.3\endcsname{ega4.c04.eq.11-2-1-3}
\expandafter\def\csname sga@target@IV.11.2.1.4\endcsname{ega4.c04.eq.11-2-1-4}
\expandafter\def\csname sga@target@IV.11.2.1.5\endcsname{ega4.c04.eq.11-2-1-5}
\expandafter\def\csname sga@target@IV.11.2.10\endcsname{ega4.c04.rem.11-2-10}
\expandafter\def\csname sga@target@IV.11.2.2\endcsname{ega4.c04.statement.11-2-2}
\expandafter\def\csname sga@target@IV.11.2.2.1\endcsname{ega4.c04.eq.11-2-2-1}
\expandafter\def\csname sga@target@IV.11.2.2.2\endcsname{ega4.c04.lem.11-2-2-2}
\expandafter\def\csname sga@target@IV.11.2.3\endcsname{ega4.c04.statement.11-2-3}
\expandafter\def\csname sga@target@IV.11.2.3.1\endcsname{ega4.c04.eq.11-2-3-1}
\expandafter\def\csname sga@target@IV.11.2.3.2\endcsname{ega4.c04.eq.11-2-3-2}
\expandafter\def\csname sga@target@IV.11.2.3.3\endcsname{ega4.c04.eq.11-2-3-3}
\expandafter\def\csname sga@target@IV.11.2.4\endcsname{ega4.c04.lem.11-2-4}
\expandafter\def\csname sga@target@IV.11.2.5\endcsname{ega4.c04.lem.11-2-5}
\expandafter\def\csname sga@target@IV.11.2.6\endcsname{ega4.c04.thm.11-2-6}
\expandafter\def\csname sga@target@IV.11.2.6.1\endcsname{ega4.c04.cor.11-2-6-1}
\expandafter\def\csname sga@target@IV.11.2.7\endcsname{ega4.c04.cor.11-2-7}
\expandafter\def\csname sga@target@IV.11.2.8\endcsname{ega4.c04.prop.11-2-8}
\expandafter\def\csname sga@target@IV.11.2.9\endcsname{ega4.c04.thm.11-2-9}
\expandafter\def\csname sga@target@IV.11.2.9.1\endcsname{ega4.c04.statement.11-2-9-1}
\expandafter\def\csname sga@target@IV.11.2.9.2\endcsname{ega4.c04.lem.11-2-9-2}
\expandafter\def\csname sga@target@IV.11.2.9.3\endcsname{ega4.c04.lem.11-2-9-3}
\expandafter\def\csname sga@target@IV.11.2.9.4\endcsname{ega4.c04.cor.11-2-9-4}
\expandafter\def\csname sga@target@IV.11.2.9.5\endcsname{ega4.c04.lem.11-2-9-5}
\expandafter\def\csname sga@target@IV.11.2.9.6\endcsname{ega4.c04.lem.11-2-9-6}
\expandafter\def\csname sga@target@IV.11.2.9.7\endcsname{ega4.c04.eq.11-2-9-7}
\expandafter\def\csname sga@target@IV.11.2.9.8\endcsname{ega4.c04.eq.11-2-9-8}
\expandafter\def\csname sga@target@IV.11.2.9.9\endcsname{ega4.c04.eq.11-2-9-9}
\expandafter\def\csname sga@target@IV.11.3\endcsname{ega4.c04.subsec.11-3.key-iv-11-3}
\expandafter\def\csname sga@target@IV.11.3.1\endcsname{ega4.c04.thm.11-3-1}
\expandafter\def\csname sga@target@IV.11.3.10\endcsname{ega4.c04.thm.11-3-10}
\expandafter\def\csname sga@target@IV.11.3.10.1\endcsname{ega4.c04.lem.11-3-10-1}
\expandafter\def\csname sga@target@IV.11.3.10.2\endcsname{ega4.c04.lem.11-3-10-2}
\expandafter\def\csname sga@target@IV.11.3.10.3\endcsname{ega4.c04.eq.11-3-10-3}
\expandafter\def\csname sga@target@IV.11.3.10.4\endcsname{ega4.c04.sec.11-3-10-4}
\expandafter\def\csname sga@target@IV.11.3.11\endcsname{ega4.c04.cor.11-3-11}
\expandafter\def\csname sga@target@IV.11.3.12\endcsname{ega4.c04.rem.11-3-12}
\expandafter\def\csname sga@target@IV.11.3.13\endcsname{ega4.c04.prop.11-3-13}
\expandafter\def\csname sga@target@IV.11.3.14\endcsname{ega4.c04.cor.11-3-14}
\expandafter\def\csname sga@target@IV.11.3.15\endcsname{ega4.c04.prop.11-3-15}
\expandafter\def\csname sga@target@IV.11.3.16\endcsname{ega4.c04.prop.11-3-16}
\expandafter\def\csname sga@target@IV.11.3.17\endcsname{ega4.c04.cor.11-3-17}
\expandafter\def\csname sga@target@IV.11.3.2\endcsname{ega4.c04.cor.11-3-2}
\expandafter\def\csname sga@target@IV.11.3.3\endcsname{ega4.c04.cor.11-3-3}
\expandafter\def\csname sga@target@IV.11.3.4\endcsname{ega4.c04.prop.11-3-4}
\expandafter\def\csname sga@target@IV.11.3.5\endcsname{ega4.c04.cor.11-3-5}
\expandafter\def\csname sga@target@IV.11.3.6\endcsname{ega4.c04.prop.11-3-6}
\expandafter\def\csname sga@target@IV.11.3.7\endcsname{ega4.c04.prop.11-3-7}
\expandafter\def\csname sga@target@IV.11.3.8\endcsname{ega4.c04.thm.11-3-8}
\expandafter\def\csname sga@target@IV.11.3.9\endcsname{ega4.c04.prop.11-3-9}
\expandafter\def\csname sga@target@IV.11.3.9.1\endcsname{ega4.c04.lem.11-3-9-1}
\expandafter\def\csname sga@target@IV.11.4.1\endcsname{ega4.c04.thm.11-4-1}
\expandafter\def\csname sga@target@IV.11.4.1.1\endcsname{ega4.c04.lem.11-4-1-1}
\expandafter\def\csname sga@target@IV.11.4.1.2\endcsname{ega4.c04.lem.11-4-1-2}
\expandafter\def\csname sga@target@IV.11.4.10\endcsname{ega4.c04.cor.11-4-10}
\expandafter\def\csname sga@target@IV.11.4.11\endcsname{ega4.c04.prop.11-4-11}
\expandafter\def\csname sga@target@IV.11.4.12\endcsname{ega4.c04.prop.11-4-12}
\expandafter\def\csname sga@target@IV.11.4.13\endcsname{ega4.c04.cor.11-4-13}
\expandafter\def\csname sga@target@IV.11.4.2\endcsname{ega4.c04.rem.11-4-2}
\expandafter\def\csname sga@target@IV.11.4.3\endcsname{ega4.c04.cor.11-4-3}
\expandafter\def\csname sga@target@IV.11.4.4\endcsname{ega4.c04.cor.11-4-4}
\expandafter\def\csname sga@target@IV.11.4.5\endcsname{ega4.c04.prop.11-4-5}
\expandafter\def\csname sga@target@IV.11.4.6\endcsname{ega4.c04.cor.11-4-6}
\expandafter\def\csname sga@target@IV.11.4.7\endcsname{ega4.c04.prop.11-4-7}
\expandafter\def\csname sga@target@IV.11.4.8\endcsname{ega4.c04.rem.11-4-8}
\expandafter\def\csname sga@target@IV.11.4.9\endcsname{ega4.c04.cor.11-4-9}
\expandafter\def\csname sga@target@IV.11.5.1\endcsname{ega4.c04.thm.11-5-1}
\expandafter\def\csname sga@target@IV.11.5.2\endcsname{ega4.c04.cor.11-5-2}
\expandafter\def\csname sga@target@IV.11.5.3\endcsname{ega4.c04.prop.11-5-3}
\expandafter\def\csname sga@target@IV.11.5.4\endcsname{ega4.c04.cor.11-5-4}
\expandafter\def\csname sga@target@IV.11.5.5\endcsname{ega4.c04.prop.11-5-5}
\expandafter\def\csname sga@target@IV.11.5.5.1\endcsname{ega4.c04.lem.11-5-5-1}
\expandafter\def\csname sga@target@IV.11.5.5.2\endcsname{ega4.c04.lem.11-5-5-2}
\expandafter\def\csname sga@target@IV.11.6.1\endcsname{ega4.c04.thm.11-6-1}
\expandafter\def\csname sga@target@IV.11.6.2\endcsname{ega4.c04.cor.11-6-2}
\expandafter\def\csname sga@target@IV.11.7.1\endcsname{ega4.c04.statement.11-7-1}
\expandafter\def\csname sga@target@IV.11.7.2\endcsname{ega4.c04.statement.11-7-2}
\expandafter\def\csname sga@target@IV.11.7.3\endcsname{ega4.c04.statement.11-7-3}
\expandafter\def\csname sga@target@IV.11.7.4\endcsname{ega4.c04.statement.11-7-4}
\expandafter\def\csname sga@target@IV.11.7.5\endcsname{ega4.c04.statement.11-7-5}
\expandafter\def\csname sga@target@IV.11.8.1\endcsname{ega4.c04.thm.11-8-1}
\expandafter\def\csname sga@target@IV.11.8.2\endcsname{ega4.c04.subsec.11-8-2}
\expandafter\def\csname sga@target@IV.11.9.1\endcsname{ega4.c04.statement.11-9-1}
\expandafter\def\csname sga@target@IV.11.9.1.1\endcsname{ega4.c04.statement.11-9-1-1}
\expandafter\def\csname sga@target@IV.11.9.10\endcsname{ega4.c04.thm.11-9-10}
\expandafter\def\csname sga@target@IV.11.9.11\endcsname{ega4.c04.rem.11-9-11}
\expandafter\def\csname sga@target@IV.11.9.12\endcsname{ega4.c04.lem.11-9-12}
\expandafter\def\csname sga@target@IV.11.9.13\endcsname{ega4.c04.prop.11-9-13}
\expandafter\def\csname sga@target@IV.11.9.14\endcsname{ega4.c04.statement.11-9-14}
\expandafter\def\csname sga@target@IV.11.9.15\endcsname{ega4.c04.prop.11-9-15}
\expandafter\def\csname sga@target@IV.11.9.16\endcsname{ega4.c04.thm.11-9-16}
\expandafter\def\csname sga@target@IV.11.9.17\endcsname{ega4.c04.prop.11-9-17}
\expandafter\def\csname sga@target@IV.11.9.17.1\endcsname{ega4.c04.lem.11-9-17-1}
\expandafter\def\csname sga@target@IV.11.9.18\endcsname{ega4.c04.item.11-9-18}
\expandafter\def\csname sga@target@IV.11.9.2\endcsname{ega4.c04.statement.11-9-2}
\expandafter\def\csname sga@target@IV.11.9.3\endcsname{ega4.c04.statement.11-9-3}
\expandafter\def\csname sga@target@IV.11.9.4\endcsname{ega4.c04.statement.11-9-4}
\expandafter\def\csname sga@target@IV.11.9.5\endcsname{ega4.c04.statement.11-9-5}
\expandafter\def\csname sga@target@IV.11.9.6\endcsname{ega4.c04.statement.11-9-6}
\expandafter\def\csname sga@target@IV.11.9.7\endcsname{ega4.c04.prop.11-9-7}
\expandafter\def\csname sga@target@IV.11.9.7.1\endcsname{ega4.c04.lem.11-9-7-1}
\expandafter\def\csname sga@target@IV.11.9.8\endcsname{ega4.c04.subsec.11-9-8}
\expandafter\def\csname sga@target@IV.11.9.9\endcsname{ega4.c04.cor.11-9-9}
\expandafter\def\csname sga@target@IV.12.0.1\endcsname{ega4.c04.statement.12-0-1}
\expandafter\def\csname sga@target@IV.12.0.2\endcsname{ega4.c04.statement.12-0-2}
\expandafter\def\csname sga@target@IV.12.1.1\endcsname{ega4.c04.thm.12-1-1}
\expandafter\def\csname sga@target@IV.12.1.1.1\endcsname{ega4.c04.statement.12-1-1-1}
\expandafter\def\csname sga@target@IV.12.1.1.2\endcsname{ega4.c04.statement.12-1-1-2}
\expandafter\def\csname sga@target@IV.12.1.1.3\endcsname{ega4.c04.statement.12-1-1-3}
\expandafter\def\csname sga@target@IV.12.1.1.4\endcsname{ega4.c04.statement.12-1-1-4}
\expandafter\def\csname sga@target@IV.12.1.1.5\endcsname{ega4.c04.prop.12-1-1-5}
\expandafter\def\csname sga@target@IV.12.1.2\endcsname{ega4.c04.item.12-1-2}
\expandafter\def\csname sga@target@IV.12.1.3\endcsname{ega4.c04.cor.12-1-3}
\expandafter\def\csname sga@target@IV.12.1.4\endcsname{ega4.c04.rem.12-1-4}
\expandafter\def\csname sga@target@IV.12.1.5\endcsname{ega4.c04.cor.12-1-5}
\expandafter\def\csname sga@target@IV.12.1.6\endcsname{ega4.c04.thm.12-1-6}
\expandafter\def\csname sga@target@IV.12.1.7\endcsname{ega4.c04.cor.12-1-7}
\expandafter\def\csname sga@target@IV.12.1.8\endcsname{ega4.c04.item.12-1-8}
\expandafter\def\csname sga@target@IV.12.2.1\endcsname{ega4.c04.thm.12-2-1}
\expandafter\def\csname sga@target@IV.12.2.1.1\endcsname{ega4.c04.statement.12-2-1-1}
\expandafter\def\csname sga@target@IV.12.2.1.2\endcsname{ega4.c04.statement.12-2-1-2}
\expandafter\def\csname sga@target@IV.12.2.1.3\endcsname{ega4.c04.statement.12-2-1-3}
\expandafter\def\csname sga@target@IV.12.2.1.4\endcsname{ega4.c04.lem.12-2-1-4}
\expandafter\def\csname sga@target@IV.12.2.2\endcsname{ega4.c04.cor.12-2-2}
\expandafter\def\csname sga@target@IV.12.2.3\endcsname{ega4.c04.item.12-2-3}
\expandafter\def\csname sga@target@IV.12.2.4\endcsname{ega4.c04.thm.12-2-4}
\expandafter\def\csname sga@target@IV.12.3.1\endcsname{ega4.c04.lem.12-3-1}
\expandafter\def\csname sga@target@IV.12.3.2\endcsname{ega4.c04.prop.12-3-2}
\expandafter\def\csname sga@target@IV.12.3.2.1\endcsname{ega4.c04.eq.12-3-2-1}
\expandafter\def\csname sga@target@IV.12.3.3\endcsname{ega4.c04.prop.12-3-3}
\expandafter\def\csname sga@target@IV.12.3.3.1\endcsname{ega4.c04.lem.12-3-3-1}
\expandafter\def\csname sga@target@IV.12.3.3.2\endcsname{ega4.c04.eq.12-3-3-2}
\expandafter\def\csname sga@target@IV.12.3.3.3\endcsname{ega4.c04.lem.12-3-3-3}
\expandafter\def\csname sga@target@IV.12.3.3.4\endcsname{ega4.c04.eq.12-3-3-4}
\expandafter\def\csname sga@target@IV.12.3.3.5\endcsname{ega4.c04.lem.12-3-3-5}
\expandafter\def\csname sga@target@IV.12.3.4\endcsname{ega4.c04.cor.12-3-4}
\expandafter\def\csname sga@target@IV.13.1.1\endcsname{ega4.c04.lem.13-1-1}
\expandafter\def\csname sga@target@IV.13.1.2\endcsname{ega4.c04.lem.13-1-2}
\expandafter\def\csname sga@target@IV.13.1.3\endcsname{ega4.c04.thm.13-1-3}
\expandafter\def\csname sga@target@IV.13.1.4\endcsname{ega4.c04.cor.13-1-4}
\expandafter\def\csname sga@target@IV.13.1.5\endcsname{ega4.c04.cor.13-1-5}
\expandafter\def\csname sga@target@IV.13.1.6\endcsname{ega4.c04.cor.13-1-6}
\expandafter\def\csname sga@target@IV.13.1.7\endcsname{ega4.c04.prop.13-1-7}
\expandafter\def\csname sga@target@IV.13.2.1\endcsname{ega4.c04.statement.13-2-1}
\expandafter\def\csname sga@target@IV.13.2.10\endcsname{ega4.c04.prop.13-2-10}
\expandafter\def\csname sga@target@IV.13.2.10.1\endcsname{ega4.c04.eq.13-2-10-1}
\expandafter\def\csname sga@target@IV.13.2.10.2\endcsname{ega4.c04.eq.13-2-10-2}
\expandafter\def\csname sga@target@IV.13.2.11\endcsname{ega4.c04.cor.13-2-11}
\expandafter\def\csname sga@target@IV.13.2.12\endcsname{ega4.c04.item.13-2-12}
\expandafter\def\csname sga@target@IV.13.2.2\endcsname{ega4.c04.def.13-2-2}
\expandafter\def\csname sga@target@IV.13.2.3\endcsname{ega4.c04.prop.13-2-3}
\expandafter\def\csname sga@target@IV.13.2.3.1\endcsname{ega4.c04.prop.13-2-3-1}
\expandafter\def\csname sga@target@IV.13.2.4\endcsname{ega4.c04.statement.13-2-4}
\expandafter\def\csname sga@target@IV.13.2.4.1\endcsname{ega4.c04.statement.13-2-4-1}
\expandafter\def\csname sga@target@IV.13.2.4.2\endcsname{ega4.c04.eq.13-2-4-2}
\expandafter\def\csname sga@target@IV.13.2.4.3\endcsname{ega4.c04.eq.13-2-4-3}
\expandafter\def\csname sga@target@IV.13.2.4.4\endcsname{ega4.c04.eq.13-2-4-4}
\expandafter\def\csname sga@target@IV.13.2.4.5\endcsname{ega4.c04.eq.13-2-4-5}
\expandafter\def\csname sga@target@IV.13.2.5\endcsname{ega4.c04.lem.13-2-5}
\expandafter\def\csname sga@target@IV.13.2.6\endcsname{ega4.c04.statement.13-2-6}
\expandafter\def\csname sga@target@IV.13.2.6.1\endcsname{ega4.c04.statement.13-2-6-1}
\expandafter\def\csname sga@target@IV.13.2.7\endcsname{ega4.c04.prop.13-2-7}
\expandafter\def\csname sga@target@IV.13.2.7.1\endcsname{ega4.c04.prop.13-2-7-1}
\expandafter\def\csname sga@target@IV.13.2.8\endcsname{ega4.c04.cor.13-2-8}
\expandafter\def\csname sga@target@IV.13.2.9\endcsname{ega4.c04.cor.13-2-9}
\expandafter\def\csname sga@target@IV.13.2.9.1\endcsname{ega4.c04.eq.13-2-9-1}
\expandafter\def\csname sga@target@IV.13.2.9.2\endcsname{ega4.c04.eq.13-2-9-2}
\expandafter\def\csname sga@target@IV.13.3.1\endcsname{ega4.c04.prop.13-3-1}
\expandafter\def\csname sga@target@IV.13.3.1.1\endcsname{ega4.c04.lem.13-3-1-1}
\expandafter\def\csname sga@target@IV.13.3.2\endcsname{ega4.c04.def.13-3-2}
\expandafter\def\csname sga@target@IV.13.3.3\endcsname{ega4.c04.prop.13-3-3}
\expandafter\def\csname sga@target@IV.13.3.4\endcsname{ega4.c04.cor.13-3-4}
\expandafter\def\csname sga@target@IV.13.3.4.1\endcsname{ega4.c04.cor.13-3-4-1}
\expandafter\def\csname sga@target@IV.13.3.5\endcsname{ega4.c04.cor.13-3-5}
\expandafter\def\csname sga@target@IV.13.3.6\endcsname{ega4.c04.prop.13-3-6}
\expandafter\def\csname sga@target@IV.13.3.6.1\endcsname{ega4.c04.prop.13-3-6-1}
\expandafter\def\csname sga@target@IV.13.3.6.2\endcsname{ega4.c04.eq.13-3-6-2}
\expandafter\def\csname sga@target@IV.13.3.7\endcsname{ega4.c04.prop.13-3-7}
\expandafter\def\csname sga@target@IV.13.3.7.1\endcsname{ega4.c04.prop.13-3-7-1}
\expandafter\def\csname sga@target@IV.13.3.8\endcsname{ega4.c04.prop.13-3-8}
\expandafter\def\csname sga@target@IV.13.3.9\endcsname{ega4.c04.statement.13-3-9}
\expandafter\def\csname sga@target@IV.14.1.1\endcsname{ega4.c04.statement.14-1-1}
\expandafter\def\csname sga@target@IV.14.1.2\endcsname{ega4.c04.prop.14-1-2}
\expandafter\def\csname sga@target@IV.14.1.3\endcsname{ega4.c04.item.14-1-3}
\expandafter\def\csname sga@target@IV.14.1.4\endcsname{ega4.c04.prop.14-1-4}
\expandafter\def\csname sga@target@IV.14.1.5\endcsname{ega4.c04.rem.14-1-5}
\expandafter\def\csname sga@target@IV.14.1.6\endcsname{ega4.c04.statement.14-1-6}
\expandafter\def\csname sga@target@IV.14.2.1\endcsname{ega4.c04.thm.14-2-1}
\expandafter\def\csname sga@target@IV.14.2.1.1\endcsname{ega4.c04.eq.14-2-1-1}
\expandafter\def\csname sga@target@IV.14.2.2\endcsname{ega4.c04.cor.14-2-2}
\expandafter\def\csname sga@target@IV.14.2.3\endcsname{ega4.c04.prop.14-2-3}
\expandafter\def\csname sga@target@IV.14.2.4\endcsname{ega4.c04.cor.14-2-4}
\expandafter\def\csname sga@target@IV.14.2.5\endcsname{ega4.c04.cor.14-2-5}
\expandafter\def\csname sga@target@IV.14.2.6\endcsname{ega4.c04.cor.14-2-6}
\expandafter\def\csname sga@target@IV.14.2.7\endcsname{ega4.c04.rem.14-2-7}
\expandafter\def\csname sga@target@IV.14.3\endcsname{ega4.c04.subsec.14-3.key-iv-14-3}
\expandafter\def\csname sga@target@IV.14.3.1\endcsname{ega4.c04.statement.14-3-1}
\expandafter\def\csname sga@target@IV.14.3.10\endcsname{ega4.c04.lem.14-3-10}
\expandafter\def\csname sga@target@IV.14.3.11\endcsname{ega4.c04.prop.14-3-11}
\expandafter\def\csname sga@target@IV.14.3.12\endcsname{ega4.c04.rem.14-3-12}
\expandafter\def\csname sga@target@IV.14.3.13\endcsname{ega4.c04.prop.14-3-13}
\expandafter\def\csname sga@target@IV.14.3.14\endcsname{ega4.c04.rem.14-3-14}
\expandafter\def\csname sga@target@IV.14.3.2\endcsname{ega4.c04.prop.14-3-2}
\expandafter\def\csname sga@target@IV.14.3.3\endcsname{ega4.c04.def.14-3-3}
\expandafter\def\csname sga@target@IV.14.3.3.1\endcsname{ega4.c04.subsec.14-3-3-1}
\expandafter\def\csname sga@target@IV.14.3.4\endcsname{ega4.c04.prop.14-3-4}
\expandafter\def\csname sga@target@IV.14.3.5\endcsname{ega4.c04.prop.14-3-5}
\expandafter\def\csname sga@target@IV.14.3.6\endcsname{ega4.c04.prop.14-3-6}
\expandafter\def\csname sga@target@IV.14.3.7\endcsname{ega4.c04.cor.14-3-7}
\expandafter\def\csname sga@target@IV.14.3.8\endcsname{ega4.c04.prop.14-3-8}
\expandafter\def\csname sga@target@IV.14.3.9\endcsname{ega4.c04.item.14-3-9}
\expandafter\def\csname sga@target@IV.14.4.1\endcsname{ega4.c04.thm.14-4-1}
\expandafter\def\csname sga@target@IV.14.4.1.1\endcsname{ega4.c04.lem.14-4-1-1}
\expandafter\def\csname sga@target@IV.14.4.1.2\endcsname{ega4.c04.lem.14-4-1-2}
\expandafter\def\csname sga@target@IV.14.4.1.3\endcsname{ega4.c04.lem.14-4-1-3}
\expandafter\def\csname sga@target@IV.14.4.10\endcsname{ega4.c04.item.14-4-10}
\expandafter\def\csname sga@target@IV.14.4.2\endcsname{ega4.c04.cor.14-4-2}
\expandafter\def\csname sga@target@IV.14.4.3\endcsname{ega4.c04.cor.14-4-3}
\expandafter\def\csname sga@target@IV.14.4.4\endcsname{ega4.c04.cor.14-4-4}
\expandafter\def\csname sga@target@IV.14.4.5\endcsname{ega4.c04.rem.14-4-5}
\expandafter\def\csname sga@target@IV.14.4.6\endcsname{ega4.c04.cor.14-4-6}
\expandafter\def\csname sga@target@IV.14.4.7\endcsname{ega4.c04.prop.14-4-7}
\expandafter\def\csname sga@target@IV.14.4.7.1\endcsname{ega4.c04.eq.14-4-7-1}
\expandafter\def\csname sga@target@IV.14.4.7.2\endcsname{ega4.c04.eq.14-4-7-2}
\expandafter\def\csname sga@target@IV.14.4.8\endcsname{ega4.c04.prop.14-4-8}
\expandafter\def\csname sga@target@IV.14.4.8.1\endcsname{ega4.c04.lem.14-4-8-1}
\expandafter\def\csname sga@target@IV.14.4.9\endcsname{ega4.c04.cor.14-4-9}
\expandafter\def\csname sga@target@IV.14.5\endcsname{ega4.c04.subsec.14-5.key-iv-14-5}
\expandafter\def\csname sga@target@IV.14.5.1\endcsname{ega4.c04.lem.14-5-1}
\expandafter\def\csname sga@target@IV.14.5.1.1\endcsname{ega4.c04.eq.14-5-1-1}
\expandafter\def\csname sga@target@IV.14.5.1.2\endcsname{ega4.c04.eq.14-5-1-2}
\expandafter\def\csname sga@target@IV.14.5.1.3\endcsname{ega4.c04.eq.14-5-1-3}
\expandafter\def\csname sga@target@IV.14.5.1.4\endcsname{ega4.c04.eq.14-5-1-4}
\expandafter\def\csname sga@target@IV.14.5.10\endcsname{ega4.c04.prop.14-5-10}
\expandafter\def\csname sga@target@IV.14.5.11\endcsname{ega4.c04.cor.14-5-11}
\expandafter\def\csname sga@target@IV.14.5.11.1\endcsname{ega4.c04.lem.14-5-11-1}
\expandafter\def\csname sga@target@IV.14.5.12\endcsname{ega4.c04.cor.14-5-12}
\expandafter\def\csname sga@target@IV.14.5.2\endcsname{ega4.c04.cor.14-5-2}
\expandafter\def\csname sga@target@IV.14.5.2.1\endcsname{ega4.c04.eq.14-5-2-1}
\expandafter\def\csname sga@target@IV.14.5.3\endcsname{ega4.c04.prop.14-5-3}
\expandafter\def\csname sga@target@IV.14.5.4\endcsname{ega4.c04.cor.14-5-4}
\expandafter\def\csname sga@target@IV.14.5.5\endcsname{ega4.c04.rem.14-5-5}
\expandafter\def\csname sga@target@IV.14.5.6\endcsname{ega4.c04.prop.14-5-6}
\expandafter\def\csname sga@target@IV.14.5.7\endcsname{ega4.c04.rem.14-5-7}
\expandafter\def\csname sga@target@IV.14.5.8\endcsname{ega4.c04.prop.14-5-8}
\expandafter\def\csname sga@target@IV.14.5.9\endcsname{ega4.c04.prop.14-5-9}
\expandafter\def\csname sga@target@IV.15.1.1\endcsname{ega4.c04.prop.15-1-1}
\expandafter\def\csname sga@target@IV.15.1.1-proof\endcsname{ega4.c04.statement.15-1-1-proof}
\expandafter\def\csname sga@target@IV.15.1.1.1\endcsname{ega4.c04.eq.15-1-1-1}
\expandafter\def\csname sga@target@IV.15.1.1.2\endcsname{ega4.c04.eq.15-1-1-2}
\expandafter\def\csname sga@target@IV.15.1.1.3\endcsname{ega4.c04.eq.15-1-1-3}
\expandafter\def\csname sga@target@IV.15.1.1.4\endcsname{ega4.c04.eq.15-1-1-4}
\expandafter\def\csname sga@target@IV.15.1.1.5\endcsname{ega4.c04.eq.15-1-1-5}
\expandafter\def\csname sga@target@IV.15.1.1.6\endcsname{ega4.c04.lem.15-1-1-6}
\expandafter\def\csname sga@target@IV.15.1.2\endcsname{ega4.c04.cor.15-1-2}
\expandafter\def\csname sga@target@IV.15.1.3\endcsname{ega4.c04.cor.15-1-3}
\expandafter\def\csname sga@target@IV.15.1.4\endcsname{ega4.c04.statement.15-1-4}
\expandafter\def\csname sga@target@IV.15.2.1\endcsname{ega4.c04.statement.15-2-1}
\expandafter\def\csname sga@target@IV.15.2.2\endcsname{ega4.c04.thm.15-2-2}
\expandafter\def\csname sga@target@IV.15.2.2-proof\endcsname{ega4.c04.item.15-2-2-proof}
\expandafter\def\csname sga@target@IV.15.2.2.1\endcsname{ega4.c04.lem.15-2-2-1}
\expandafter\def\csname sga@target@IV.15.2.3\endcsname{ega4.c04.cor.15-2-3}
\expandafter\def\csname sga@target@IV.15.2.4\endcsname{ega4.c04.item.15-2-4}
\expandafter\def\csname sga@target@IV.15.3.1\endcsname{ega4.c04.prop.15-3-1}
\expandafter\def\csname sga@target@IV.15.3.1-proof\endcsname{ega4.c04.item.15-3-1-proof}
\expandafter\def\csname sga@target@IV.15.3.2\endcsname{ega4.c04.prop.15-3-2}
\expandafter\def\csname sga@target@IV.15.3.2-proof\endcsname{ega4.c04.item.15-3-2-proof}
\expandafter\def\csname sga@target@IV.15.3.3\endcsname{ega4.c04.cor.15-3-3}
\expandafter\def\csname sga@target@IV.15.3.3-proof\endcsname{ega4.c04.item.15-3-3-proof}
\expandafter\def\csname sga@target@IV.15.3.4\endcsname{ega4.c04.item.15-3-4}
\expandafter\def\csname sga@target@IV.15.4.1\endcsname{ega4.c04.prop.15-4-1}
\expandafter\def\csname sga@target@IV.15.4.2\endcsname{ega4.c04.prop.15-4-2}
\expandafter\def\csname sga@target@IV.15.4.3\endcsname{ega4.c04.prop.15-4-3}
\expandafter\def\csname sga@target@IV.15.4.3-proof\endcsname{ega4.c04.item.15-4-3-proof}
\expandafter\def\csname sga@target@IV.15.5.1\endcsname{ega4.c04.prop.15-5-1}
\expandafter\def\csname sga@target@IV.15.5.10\endcsname{ega4.c04.item.15-5-10}
\expandafter\def\csname sga@target@IV.15.5.2\endcsname{ega4.c04.cor.15-5-2}
\expandafter\def\csname sga@target@IV.15.5.2.1\endcsname{ega4.c04.eq.15-5-2-1}
\expandafter\def\csname sga@target@IV.15.5.2.2\endcsname{ega4.c04.lem.15-5-2-2}
\expandafter\def\csname sga@target@IV.15.5.2.3\endcsname{ega4.c04.eq.15-5-2-3}
\expandafter\def\csname sga@target@IV.15.5.2.4\endcsname{ega4.c04.eq.15-5-2-4}
\expandafter\def\csname sga@target@IV.15.5.3\endcsname{ega4.c04.prop.15-5-3}
\expandafter\def\csname sga@target@IV.15.5.4\endcsname{ega4.c04.cor.15-5-4}
\expandafter\def\csname sga@target@IV.15.5.5\endcsname{ega4.c04.statement.15-5-5}
\expandafter\def\csname sga@target@IV.15.5.6\endcsname{ega4.c04.lem.15-5-6}
\expandafter\def\csname sga@target@IV.15.5.6.1\endcsname{ega4.c04.lem.15-5-6-1}
\expandafter\def\csname sga@target@IV.15.5.7\endcsname{ega4.c04.prop.15-5-7}
\expandafter\def\csname sga@target@IV.15.5.7.1\endcsname{ega4.c04.eq.15-5-7-1}
\expandafter\def\csname sga@target@IV.15.5.8\endcsname{ega4.c04.statement.15-5-8}
\expandafter\def\csname sga@target@IV.15.5.9\endcsname{ega4.c04.prop.15-5-9}
\expandafter\def\csname sga@target@IV.15.6.1\endcsname{ega4.c04.prop.15-6-1}
\expandafter\def\csname sga@target@IV.15.6.1.1\endcsname{ega4.c04.eq.15-6-1-1}
\expandafter\def\csname sga@target@IV.15.6.10\endcsname{ega4.c04.item.15-6-10}
\expandafter\def\csname sga@target@IV.15.6.2\endcsname{ega4.c04.cor.15-6-2}
\expandafter\def\csname sga@target@IV.15.6.3\endcsname{ega4.c04.prop.15-6-3}
\expandafter\def\csname sga@target@IV.15.6.4\endcsname{ega4.c04.prop.15-6-4}
\expandafter\def\csname sga@target@IV.15.6.5\endcsname{ega4.c04.cor.15-6-5}
\expandafter\def\csname sga@target@IV.15.6.6\endcsname{ega4.c04.prop.15-6-6}
\expandafter\def\csname sga@target@IV.15.6.7\endcsname{ega4.c04.cor.15-6-7}
\expandafter\def\csname sga@target@IV.15.6.8\endcsname{ega4.c04.prop.15-6-8}
\expandafter\def\csname sga@target@IV.15.6.9\endcsname{ega4.c04.cor.15-6-9}
\expandafter\def\csname sga@target@IV.15.7.1\endcsname{ega4.c04.def.15-7-1}
\expandafter\def\csname sga@target@IV.15.7.10\endcsname{ega4.c04.cor.15-7-10}
\expandafter\def\csname sga@target@IV.15.7.11\endcsname{ega4.c04.cor.15-7-11}
\expandafter\def\csname sga@target@IV.15.7.11.1\endcsname{ega4.c04.lem.15-7-11-1}
\expandafter\def\csname sga@target@IV.15.7.2\endcsname{ega4.c04.prop.15-7-2}
\expandafter\def\csname sga@target@IV.15.7.3\endcsname{ega4.c04.rem.15-7-3}
\expandafter\def\csname sga@target@IV.15.7.4\endcsname{ega4.c04.cor.15-7-4}
\expandafter\def\csname sga@target@IV.15.7.5\endcsname{ega4.c04.cor.15-7-5}
\expandafter\def\csname sga@target@IV.15.7.6\endcsname{ega4.c04.cor.15-7-6}
\expandafter\def\csname sga@target@IV.15.7.7\endcsname{ega4.c04.item.15-7-7}
\expandafter\def\csname sga@target@IV.15.7.8\endcsname{ega4.c04.statement.15-7-8}
\expandafter\def\csname sga@target@IV.15.7.9\endcsname{ega4.c04.prop.15-7-9}
\expandafter\def\csname sga@target@IV.16.1\endcsname{ega4.c04.subsec.16-1}
\expandafter\def\csname sga@target@IV.16.1.1\endcsname{ega4.c04.statement.16-1-1}
\expandafter\def\csname sga@target@IV.16.1.10\endcsname{ega4.c04.cor.16-1-10}
\expandafter\def\csname sga@target@IV.16.1.11\endcsname{ega4.c04.rem.16-1-11}
\expandafter\def\csname sga@target@IV.16.1.2\endcsname{ega4.c04.def.16-1-2}
\expandafter\def\csname sga@target@IV.16.1.2.1\endcsname{ega4.c04.def.16-1-2-1}
\expandafter\def\csname sga@target@IV.16.1.2.2\endcsname{ega4.c04.eq.16-1-2-2}
\expandafter\def\csname sga@target@IV.16.1.3\endcsname{ega4.c04.statement.16-1-3}
\expandafter\def\csname sga@target@IV.16.1.4\endcsname{ega4.c04.statement.16-1-4}
\expandafter\def\csname sga@target@IV.16.1.5\endcsname{ega4.c04.prop.16-1-5}
\expandafter\def\csname sga@target@IV.16.1.6\endcsname{ega4.c04.cor.16-1-6}
\expandafter\def\csname sga@target@IV.16.1.7\endcsname{ega4.c04.cor.16-1-7}
\expandafter\def\csname sga@target@IV.16.1.8\endcsname{ega4.c04.prop.16-1-8}
\expandafter\def\csname sga@target@IV.16.1.9\endcsname{ega4.c04.prop.16-1-9}
\expandafter\def\csname sga@target@IV.16.10\endcsname{ega4.c04.subsec.16-10}
\expandafter\def\csname sga@target@IV.16.10.1\endcsname{ega4.c04.def.16-10-1}
\expandafter\def\csname sga@target@IV.16.10.2\endcsname{ega4.c04.prop.16-10-2}
\expandafter\def\csname sga@target@IV.16.10.3\endcsname{ega4.c04.prop.16-10-3}
\expandafter\def\csname sga@target@IV.16.10.4\endcsname{ega4.c04.prop.16-10-4}
\expandafter\def\csname sga@target@IV.16.10.5\endcsname{ega4.c04.prop.16-10-5}
\expandafter\def\csname sga@target@IV.16.10.6\endcsname{ega4.c04.rem.16-10-6}
\expandafter\def\csname sga@target@IV.16.11\endcsname{ega4.c04.subsec.16-11}
\expandafter\def\csname sga@target@IV.16.11.1\endcsname{ega4.c04.statement.16-11-1}
\expandafter\def\csname sga@target@IV.16.11.1.1\endcsname{ega4.c04.statement.16-11-1-1}
\expandafter\def\csname sga@target@IV.16.11.1.2\endcsname{ega4.c04.statement.16-11-1-2}
\expandafter\def\csname sga@target@IV.16.11.1.3\endcsname{ega4.c04.eq.16-11-1-3}
\expandafter\def\csname sga@target@IV.16.11.1.4\endcsname{ega4.c04.eq.16-11-1-4}
\expandafter\def\csname sga@target@IV.16.11.1.5\endcsname{ega4.c04.eq.16-11-1-5}
\expandafter\def\csname sga@target@IV.16.11.1.6\endcsname{ega4.c04.eq.16-11-1-6}
\expandafter\def\csname sga@target@IV.16.11.1.7\endcsname{ega4.c04.eq.16-11-1-7}
\expandafter\def\csname sga@target@IV.16.11.2\endcsname{ega4.c04.thm.16-11-2}
\expandafter\def\csname sga@target@IV.16.11.2.1\endcsname{ega4.c04.eq.16-11-2-1}
\expandafter\def\csname sga@target@IV.16.11.2.2\endcsname{ega4.c04.eq.16-11-2-2}
\expandafter\def\csname sga@target@IV.16.11.2.3\endcsname{ega4.c04.eq.16-11-2-3}
\expandafter\def\csname sga@target@IV.16.11.3\endcsname{ega4.c04.statement.16-11-3}
\expandafter\def\csname sga@target@IV.16.12\endcsname{ega4.c04.subsec.16-12}
\expandafter\def\csname sga@target@IV.16.12.1\endcsname{ega4.c04.statement.16-12-1}
\expandafter\def\csname sga@target@IV.16.12.2\endcsname{ega4.c04.thm.16-12-2}
\expandafter\def\csname sga@target@IV.16.12.3\endcsname{ega4.c04.statement.16-12-3}
\expandafter\def\csname sga@target@IV.16.2\endcsname{ega4.c04.subsec.16-2}
\expandafter\def\csname sga@target@IV.16.2.1\endcsname{ega4.c04.statement.16-2-1}
\expandafter\def\csname sga@target@IV.16.2.1.1\endcsname{ega4.c04.statement.16-2-1-1}
\expandafter\def\csname sga@target@IV.16.2.1.2\endcsname{ega4.c04.eq.16-2-1-2}
\expandafter\def\csname sga@target@IV.16.2.1.3\endcsname{ega4.c04.eq.16-2-1-3}
\expandafter\def\csname sga@target@IV.16.2.1.4\endcsname{ega4.c04.eq.16-2-1-4}
\expandafter\def\csname sga@target@IV.16.2.2\endcsname{ega4.c04.prop.16-2-2}
\expandafter\def\csname sga@target@IV.16.2.3\endcsname{ega4.c04.cor.16-2-3}
\expandafter\def\csname sga@target@IV.16.2.3.1\endcsname{ega4.c04.eq.16-2-3-1}
\expandafter\def\csname sga@target@IV.16.2.3.2\endcsname{ega4.c04.eq.16-2-3-2}
\expandafter\def\csname sga@target@IV.16.2.4\endcsname{ega4.c04.cor.16-2-4}
\expandafter\def\csname sga@target@IV.16.2.5\endcsname{ega4.c04.statement.16-2-5}
\expandafter\def\csname sga@target@IV.16.2.6\endcsname{ega4.c04.statement.16-2-6}
\expandafter\def\csname sga@target@IV.16.2.6.1\endcsname{ega4.c04.statement.16-2-6-1}
\expandafter\def\csname sga@target@IV.16.2.6.2\endcsname{ega4.c04.eq.16-2-6-2}
\expandafter\def\csname sga@target@IV.16.2.7\endcsname{ega4.c04.prop.16-2-7}
\expandafter\def\csname sga@target@IV.16.2.7.1\endcsname{ega4.c04.prop.16-2-7-1}
\expandafter\def\csname sga@target@IV.16.3\endcsname{ega4.c04.subsec.16-3}
\expandafter\def\csname sga@target@IV.16.3.1\endcsname{ega4.c04.def.16-3-1}
\expandafter\def\csname sga@target@IV.16.3.1.1\endcsname{ega4.c04.eq.16-3-1-1}
\expandafter\def\csname sga@target@IV.16.3.2\endcsname{ega4.c04.statement.16-3-2}
\expandafter\def\csname sga@target@IV.16.3.3\endcsname{ega4.c04.statement.16-3-3}
\expandafter\def\csname sga@target@IV.16.3.3.1\endcsname{ega4.c04.statement.16-3-3-1}
\expandafter\def\csname sga@target@IV.16.3.4\endcsname{ega4.c04.prop.16-3-4}
\expandafter\def\csname sga@target@IV.16.3.5\endcsname{ega4.c04.statement.16-3-5}
\expandafter\def\csname sga@target@IV.16.3.6\endcsname{ega4.c04.def.16-3-6}
\expandafter\def\csname sga@target@IV.16.3.6.1\endcsname{ega4.c04.eq.16-3-6-1}
\expandafter\def\csname sga@target@IV.16.3.7\endcsname{ega4.c04.statement.16-3-7}
\expandafter\def\csname sga@target@IV.16.3.7.1\endcsname{ega4.c04.eq.16-3-7-1}
\expandafter\def\csname sga@target@IV.16.3.8\endcsname{ega4.c04.prop.16-3-8}
\expandafter\def\csname sga@target@IV.16.3.9\endcsname{ega4.c04.prop.16-3-9}
\expandafter\def\csname sga@target@IV.16.4\endcsname{ega4.c04.subsec.16-4}
\expandafter\def\csname sga@target@IV.16.4.1\endcsname{ega4.c04.statement.16-4-1}
\expandafter\def\csname sga@target@IV.16.4.1.1\endcsname{ega4.c04.statement.16-4-1-1}
\expandafter\def\csname sga@target@IV.16.4.1.2\endcsname{ega4.c04.statement.16-4-1-2}
\expandafter\def\csname sga@target@IV.16.4.1.3\endcsname{ega4.c04.eq.16-4-1-3}
\expandafter\def\csname sga@target@IV.16.4.1.4\endcsname{ega4.c04.eq.16-4-1-4}
\expandafter\def\csname sga@target@IV.16.4.1.5\endcsname{ega4.c04.eq.16-4-1-5}
\expandafter\def\csname sga@target@IV.16.4.10\endcsname{ega4.c04.cor.16-4-10}
\expandafter\def\csname sga@target@IV.16.4.11\endcsname{ega4.c04.prop.16-4-11}
\expandafter\def\csname sga@target@IV.16.4.11.1\endcsname{ega4.c04.prop.16-4-11-1}
\expandafter\def\csname sga@target@IV.16.4.11.2\endcsname{ega4.c04.prop.16-4-11-2}
\expandafter\def\csname sga@target@IV.16.4.11.3\endcsname{ega4.c04.eq.16-4-11-3}
\expandafter\def\csname sga@target@IV.16.4.11.4\endcsname{ega4.c04.eq.16-4-11-4}
\expandafter\def\csname sga@target@IV.16.4.12\endcsname{ega4.c04.cor.16-4-12}
\expandafter\def\csname sga@target@IV.16.4.13\endcsname{ega4.c04.cor.16-4-13}
\expandafter\def\csname sga@target@IV.16.4.14\endcsname{ega4.c04.prop.16-4-14}
\expandafter\def\csname sga@target@IV.16.4.14.1\endcsname{ega4.c04.prop.16-4-14-1}
\expandafter\def\csname sga@target@IV.16.4.15\endcsname{ega4.c04.cor.16-4-15}
\expandafter\def\csname sga@target@IV.16.4.15.1\endcsname{ega4.c04.cor.16-4-15-1}
\expandafter\def\csname sga@target@IV.16.4.15.2\endcsname{ega4.c04.eq.16-4-15-2}
\expandafter\def\csname sga@target@IV.16.4.16\endcsname{ega4.c04.cor.16-4-16}
\expandafter\def\csname sga@target@IV.16.4.16.1\endcsname{ega4.c04.cor.16-4-16-1}
\expandafter\def\csname sga@target@IV.16.4.16.2\endcsname{ega4.c04.eq.16-4-16-2}
\expandafter\def\csname sga@target@IV.16.4.17\endcsname{ega4.c04.cor.16-4-17}
\expandafter\def\csname sga@target@IV.16.4.18\endcsname{ega4.c04.statement.16-4-18}
\expandafter\def\csname sga@target@IV.16.4.18.1\endcsname{ega4.c04.statement.16-4-18-1}
\expandafter\def\csname sga@target@IV.16.4.18.2\endcsname{ega4.c04.statement.16-4-18-2}
\expandafter\def\csname sga@target@IV.16.4.18.3\endcsname{ega4.c04.eq.16-4-18-3}
\expandafter\def\csname sga@target@IV.16.4.19\endcsname{ega4.c04.cor.16-4-19}
\expandafter\def\csname sga@target@IV.16.4.19.1\endcsname{ega4.c04.cor.16-4-19-1}
\expandafter\def\csname sga@target@IV.16.4.2\endcsname{ega4.c04.statement.16-4-2}
\expandafter\def\csname sga@target@IV.16.4.20\endcsname{ega4.c04.prop.16-4-20}
\expandafter\def\csname sga@target@IV.16.4.21\endcsname{ega4.c04.cor.16-4-21}
\expandafter\def\csname sga@target@IV.16.4.21.1\endcsname{ega4.c04.cor.16-4-21-1}
\expandafter\def\csname sga@target@IV.16.4.22\endcsname{ega4.c04.cor.16-4-22}
\expandafter\def\csname sga@target@IV.16.4.23\endcsname{ega4.c04.prop.16-4-23}
\expandafter\def\csname sga@target@IV.16.4.23.1\endcsname{ega4.c04.prop.16-4-23-1}
\expandafter\def\csname sga@target@IV.16.4.24\endcsname{ega4.c04.statement.16-4-24}
\expandafter\def\csname sga@target@IV.16.4.3\endcsname{ega4.c04.statement.16-4-3}
\expandafter\def\csname sga@target@IV.16.4.3.1\endcsname{ega4.c04.statement.16-4-3-1}
\expandafter\def\csname sga@target@IV.16.4.3.2\endcsname{ega4.c04.statement.16-4-3-2}
\expandafter\def\csname sga@target@IV.16.4.3.3\endcsname{ega4.c04.eq.16-4-3-3}
\expandafter\def\csname sga@target@IV.16.4.3.4\endcsname{ega4.c04.eq.16-4-3-4}
\expandafter\def\csname sga@target@IV.16.4.3.5\endcsname{ega4.c04.eq.16-4-3-5}
\expandafter\def\csname sga@target@IV.16.4.3.6\endcsname{ega4.c04.eq.16-4-3-6}
\expandafter\def\csname sga@target@IV.16.4.3.7\endcsname{ega4.c04.eq.16-4-3-7}
\expandafter\def\csname sga@target@IV.16.4.4\endcsname{ega4.c04.statement.16-4-4}
\expandafter\def\csname sga@target@IV.16.4.5\endcsname{ega4.c04.prop.16-4-5}
\expandafter\def\csname sga@target@IV.16.4.6\endcsname{ega4.c04.cor.16-4-6}
\expandafter\def\csname sga@target@IV.16.4.7\endcsname{ega4.c04.statement.16-4-7}
\expandafter\def\csname sga@target@IV.16.4.7.1\endcsname{ega4.c04.statement.16-4-7-1}
\expandafter\def\csname sga@target@IV.16.4.7.2\endcsname{ega4.c04.eq.16-4-7-2}
\expandafter\def\csname sga@target@IV.16.4.7.3\endcsname{ega4.c04.eq.16-4-7-3}
\expandafter\def\csname sga@target@IV.16.4.8\endcsname{ega4.c04.prop.16-4-8}
\expandafter\def\csname sga@target@IV.16.4.8.1\endcsname{ega4.c04.prop.16-4-8-1}
\expandafter\def\csname sga@target@IV.16.4.8.2\endcsname{ega4.c04.eq.16-4-8-2}
\expandafter\def\csname sga@target@IV.16.4.8.3\endcsname{ega4.c04.lem.16-4-8-3}
\expandafter\def\csname sga@target@IV.16.4.8.4\endcsname{ega4.c04.eq.16-4-8-4}
\expandafter\def\csname sga@target@IV.16.4.9\endcsname{ega4.c04.cor.16-4-9}
\expandafter\def\csname sga@target@IV.16.4.9.1\endcsname{ega4.c04.cor.16-4-9-1}
\expandafter\def\csname sga@target@IV.16.5\endcsname{ega4.c04.subsec.16-5}
\expandafter\def\csname sga@target@IV.16.5.1\endcsname{ega4.c04.statement.16-5-1}
\expandafter\def\csname sga@target@IV.16.5.1.1\endcsname{ega4.c04.eq.16-5-1-1}
\expandafter\def\csname sga@target@IV.16.5.1.2\endcsname{ega4.c04.eq.16-5-1-2}
\expandafter\def\csname sga@target@IV.16.5.10\endcsname{ega4.c04.statement.16-5-10}
\expandafter\def\csname sga@target@IV.16.5.10.1\endcsname{ega4.c04.statement.16-5-10-1}
\expandafter\def\csname sga@target@IV.16.5.10.2\endcsname{ega4.c04.statement.16-5-10-2}
\expandafter\def\csname sga@target@IV.16.5.11\endcsname{ega4.c04.prop.16-5-11}
\expandafter\def\csname sga@target@IV.16.5.12\endcsname{ega4.c04.statement.16-5-12}
\expandafter\def\csname sga@target@IV.16.5.12.1\endcsname{ega4.c04.statement.16-5-12-1}
\expandafter\def\csname sga@target@IV.16.5.12.2\endcsname{ega4.c04.eq.16-5-12-2}
\expandafter\def\csname sga@target@IV.16.5.13\endcsname{ega4.c04.statement.16-5-13}
\expandafter\def\csname sga@target@IV.16.5.13.1\endcsname{ega4.c04.statement.16-5-13-1}
\expandafter\def\csname sga@target@IV.16.5.13.2\endcsname{ega4.c04.eq.16-5-13-2}
\expandafter\def\csname sga@target@IV.16.5.13.3\endcsname{ega4.c04.eq.16-5-13-3}
\expandafter\def\csname sga@target@IV.16.5.13.4\endcsname{ega4.c04.eq.16-5-13-4}
\expandafter\def\csname sga@target@IV.16.5.13.5\endcsname{ega4.c04.eq.16-5-13-5}
\expandafter\def\csname sga@target@IV.16.5.14\endcsname{ega4.c04.statement.16-5-14}
\expandafter\def\csname sga@target@IV.16.5.14.1\endcsname{ega4.c04.statement.16-5-14-1}
\expandafter\def\csname sga@target@IV.16.5.14.2\endcsname{ega4.c04.eq.16-5-14-2}
\expandafter\def\csname sga@target@IV.16.5.14.3\endcsname{ega4.c04.eq.16-5-14-3}
\expandafter\def\csname sga@target@IV.16.5.15\endcsname{ega4.c04.statement.16-5-15}
\expandafter\def\csname sga@target@IV.16.5.16\endcsname{ega4.c04.prop.16-5-16}
\expandafter\def\csname sga@target@IV.16.5.17\endcsname{ega4.c04.prop.16-5-17}
\expandafter\def\csname sga@target@IV.16.5.18\endcsname{ega4.c04.cor.16-5-18}
\expandafter\def\csname sga@target@IV.16.5.2\endcsname{ega4.c04.prop.16-5-2}
\expandafter\def\csname sga@target@IV.16.5.3\endcsname{ega4.c04.prop.16-5-3}
\expandafter\def\csname sga@target@IV.16.5.3.1\endcsname{ega4.c04.eq.16-5-3-1}
\expandafter\def\csname sga@target@IV.16.5.4\endcsname{ega4.c04.statement.16-5-4}
\expandafter\def\csname sga@target@IV.16.5.5\endcsname{ega4.c04.cor.16-5-5}
\expandafter\def\csname sga@target@IV.16.5.5.1\endcsname{ega4.c04.cor.16-5-5-1}
\expandafter\def\csname sga@target@IV.16.5.6\endcsname{ega4.c04.cor.16-5-6}
\expandafter\def\csname sga@target@IV.16.5.7\endcsname{ega4.c04.statement.16-5-7}
\expandafter\def\csname sga@target@IV.16.5.7.1\endcsname{ega4.c04.statement.16-5-7-1}
\expandafter\def\csname sga@target@IV.16.5.8\endcsname{ega4.c04.statement.16-5-8}
\expandafter\def\csname sga@target@IV.16.5.8.1\endcsname{ega4.c04.eq.16-5-8-1}
\expandafter\def\csname sga@target@IV.16.5.8.2\endcsname{ega4.c04.eq.16-5-8-2}
\expandafter\def\csname sga@target@IV.16.5.9\endcsname{ega4.c04.statement.16-5-9}
\expandafter\def\csname sga@target@IV.16.5.9.1\endcsname{ega4.c04.statement.16-5-9-1}
\expandafter\def\csname sga@target@IV.16.6\endcsname{ega4.c04.subsec.16-6}
\expandafter\def\csname sga@target@IV.16.6.1\endcsname{ega4.c04.statement.16-6-1}
\expandafter\def\csname sga@target@IV.16.6.1.1\endcsname{ega4.c04.statement.16-6-1-1}
\expandafter\def\csname sga@target@IV.16.6.1.2\endcsname{ega4.c04.statement.16-6-1-2}
\expandafter\def\csname sga@target@IV.16.6.2\endcsname{ega4.c04.thm.16-6-2}
\expandafter\def\csname sga@target@IV.16.6.2.1\endcsname{ega4.c04.eq.16-6-2-1}
\expandafter\def\csname sga@target@IV.16.6.2.2\endcsname{ega4.c04.eq.16-6-2-2}
\expandafter\def\csname sga@target@IV.16.6.2.3\endcsname{ega4.c04.eq.16-6-2-3}
\expandafter\def\csname sga@target@IV.16.6.2.4\endcsname{ega4.c04.eq.16-6-2-4}
\expandafter\def\csname sga@target@IV.16.6.2.5\endcsname{ega4.c04.eq.16-6-2-5}
\expandafter\def\csname sga@target@IV.16.6.3\endcsname{ega4.c04.def.16-6-3}
\expandafter\def\csname sga@target@IV.16.6.4\endcsname{ega4.c04.prop.16-6-4}
\expandafter\def\csname sga@target@IV.16.6.4.1\endcsname{ega4.c04.prop.16-6-4-1}
\expandafter\def\csname sga@target@IV.16.6.5\endcsname{ega4.c04.statement.16-6-5}
\expandafter\def\csname sga@target@IV.16.6.5.1\endcsname{ega4.c04.statement.16-6-5-1}
\expandafter\def\csname sga@target@IV.16.7\endcsname{ega4.c04.subsec.16-7}
\expandafter\def\csname sga@target@IV.16.7.1\endcsname{ega4.c04.statement.16-7-1}
\expandafter\def\csname sga@target@IV.16.7.1.1\endcsname{ega4.c04.eq.16-7-1-1}
\expandafter\def\csname sga@target@IV.16.7.1.2\endcsname{ega4.c04.eq.16-7-1-2}
\expandafter\def\csname sga@target@IV.16.7.10\endcsname{ega4.c04.rem.16-7-10}
\expandafter\def\csname sga@target@IV.16.7.2\endcsname{ega4.c04.statement.16-7-2}
\expandafter\def\csname sga@target@IV.16.7.2.1\endcsname{ega4.c04.statement.16-7-2-1}
\expandafter\def\csname sga@target@IV.16.7.3\endcsname{ega4.c04.prop.16-7-3}
\expandafter\def\csname sga@target@IV.16.7.4\endcsname{ega4.c04.statement.16-7-4}
\expandafter\def\csname sga@target@IV.16.7.4.1\endcsname{ega4.c04.statement.16-7-4-1}
\expandafter\def\csname sga@target@IV.16.7.5\endcsname{ega4.c04.statement.16-7-5}
\expandafter\def\csname sga@target@IV.16.7.5.1\endcsname{ega4.c04.statement.16-7-5-1}
\expandafter\def\csname sga@target@IV.16.7.5.2\endcsname{ega4.c04.statement.16-7-5-2}
\expandafter\def\csname sga@target@IV.16.7.5.3\endcsname{ega4.c04.statement.16-7-5-3}
\expandafter\def\csname sga@target@IV.16.7.5.4\endcsname{ega4.c04.statement.16-7-5-4}
\expandafter\def\csname sga@target@IV.16.7.6\endcsname{ega4.c04.prop.16-7-6}
\expandafter\def\csname sga@target@IV.16.7.7\endcsname{ega4.c04.statement.16-7-7}
\expandafter\def\csname sga@target@IV.16.7.7.1\endcsname{ega4.c04.eq.16-7-7-1}
\expandafter\def\csname sga@target@IV.16.7.7.2\endcsname{ega4.c04.eq.16-7-7-2}
\expandafter\def\csname sga@target@IV.16.7.8\endcsname{ega4.c04.statement.16-7-8}
\expandafter\def\csname sga@target@IV.16.7.8.1\endcsname{ega4.c04.statement.16-7-8-1}
\expandafter\def\csname sga@target@IV.16.7.8.2\endcsname{ega4.c04.eq.16-7-8-2}
\expandafter\def\csname sga@target@IV.16.7.9\endcsname{ega4.c04.statement.16-7-9}
\expandafter\def\csname sga@target@IV.16.7.9.1\endcsname{ega4.c04.statement.16-7-9-1}
\expandafter\def\csname sga@target@IV.16.8\endcsname{ega4.c04.subsec.16-8}
\expandafter\def\csname sga@target@IV.16.8.1\endcsname{ega4.c04.def.16-8-1}
\expandafter\def\csname sga@target@IV.16.8.10\endcsname{ega4.c04.cor.16-8-10}
\expandafter\def\csname sga@target@IV.16.8.11\endcsname{ega4.c04.statement.16-8-11}
\expandafter\def\csname sga@target@IV.16.8.2\endcsname{ega4.c04.statement.16-8-2}
\expandafter\def\csname sga@target@IV.16.8.3\endcsname{ega4.c04.statement.16-8-3}
\expandafter\def\csname sga@target@IV.16.8.3.1\endcsname{ega4.c04.eq.16-8-3-1}
\expandafter\def\csname sga@target@IV.16.8.4\endcsname{ega4.c04.prop.16-8-4}
\expandafter\def\csname sga@target@IV.16.8.4.1\endcsname{ega4.c04.prop.16-8-4-1}
\expandafter\def\csname sga@target@IV.16.8.5\endcsname{ega4.c04.statement.16-8-5}
\expandafter\def\csname sga@target@IV.16.8.5.1\endcsname{ega4.c04.eq.16-8-5-1}
\expandafter\def\csname sga@target@IV.16.8.5.2\endcsname{ega4.c04.eq.16-8-5-2}
\expandafter\def\csname sga@target@IV.16.8.6\endcsname{ega4.c04.prop.16-8-6}
\expandafter\def\csname sga@target@IV.16.8.7\endcsname{ega4.c04.statement.16-8-7}
\expandafter\def\csname sga@target@IV.16.8.7.1\endcsname{ega4.c04.eq.16-8-7-1}
\expandafter\def\csname sga@target@IV.16.8.8\endcsname{ega4.c04.prop.16-8-8}
\expandafter\def\csname sga@target@IV.16.8.8.1\endcsname{ega4.c04.eq.16-8-8-1}
\expandafter\def\csname sga@target@IV.16.8.8.2\endcsname{ega4.c04.eq.16-8-8-2}
\expandafter\def\csname sga@target@IV.16.8.9\endcsname{ega4.c04.prop.16-8-9}
\expandafter\def\csname sga@target@IV.16.8.9.1\endcsname{ega4.c04.lem.16-8-9-1}
\expandafter\def\csname sga@target@IV.16.8.9.2\endcsname{ega4.c04.lem.16-8-9-2}
\expandafter\def\csname sga@target@IV.16.8.9.3\endcsname{ega4.c04.lem.16-8-9-3}
\expandafter\def\csname sga@target@IV.16.8.9.4\endcsname{ega4.c04.eq.16-8-9-4}
\expandafter\def\csname sga@target@IV.16.9\endcsname{ega4.c04.subsec.16-9}
\expandafter\def\csname sga@target@IV.16.9.1\endcsname{ega4.c04.def.16-9-1}
\expandafter\def\csname sga@target@IV.16.9.10\endcsname{ega4.c04.prop.16-9-10}
\expandafter\def\csname sga@target@IV.16.9.11\endcsname{ega4.c04.cor.16-9-11}
\expandafter\def\csname sga@target@IV.16.9.12\endcsname{ega4.c04.statement.16-9-12}
\expandafter\def\csname sga@target@IV.16.9.13\endcsname{ega4.c04.prop.16-9-13}
\expandafter\def\csname sga@target@IV.16.9.13.1\endcsname{ega4.c04.prop.16-9-13-1}
\expandafter\def\csname sga@target@IV.16.9.13.2\endcsname{ega4.c04.lem.16-9-13-2}
\expandafter\def\csname sga@target@IV.16.9.13.3\endcsname{ega4.c04.lem.16-9-13-3}
\expandafter\def\csname sga@target@IV.16.9.13.4\endcsname{ega4.c04.eq.16-9-13-4}
\expandafter\def\csname sga@target@IV.16.9.13.5\endcsname{ega4.c04.cor.16-9-13-5}
\expandafter\def\csname sga@target@IV.16.9.13.6\endcsname{ega4.c04.cor.16-9-13-6}
\expandafter\def\csname sga@target@IV.16.9.13.7\endcsname{ega4.c04.eq.16-9-13-7}
\expandafter\def\csname sga@target@IV.16.9.2\endcsname{ega4.c04.def.16-9-2}
\expandafter\def\csname sga@target@IV.16.9.3\endcsname{ega4.c04.prop.16-9-3}
\expandafter\def\csname sga@target@IV.16.9.4\endcsname{ega4.c04.cor.16-9-4}
\expandafter\def\csname sga@target@IV.16.9.4.1\endcsname{ega4.c04.eq.16-9-4-1}
\expandafter\def\csname sga@target@IV.16.9.5\endcsname{ega4.c04.cor.16-9-5}
\expandafter\def\csname sga@target@IV.16.9.6\endcsname{ega4.c04.statement.16-9-6}
\expandafter\def\csname sga@target@IV.16.9.7\endcsname{ega4.c04.statement.16-9-7}
\expandafter\def\csname sga@target@IV.16.9.8\endcsname{ega4.c04.prop.16-9-8}
\expandafter\def\csname sga@target@IV.16.9.9\endcsname{ega4.c04.cor.16-9-9}
\expandafter\def\csname sga@target@IV.17.1.1\endcsname{ega4.c04.def.17-1-1}
\expandafter\def\csname sga@target@IV.17.1.1.1\endcsname{ega4.c04.def.17-1-1-1}
\expandafter\def\csname sga@target@IV.17.1.2\endcsname{ega4.c04.rem.17-1-2}
\expandafter\def\csname sga@target@IV.17.1.2.1\endcsname{ega4.c04.eq.17-1-2-1}
\expandafter\def\csname sga@target@IV.17.1.3\endcsname{ega4.c04.prop.17-1-3}
\expandafter\def\csname sga@target@IV.17.1.4\endcsname{ega4.c04.prop.17-1-4}
\expandafter\def\csname sga@target@IV.17.1.5\endcsname{ega4.c04.cor.17-1-5}
\expandafter\def\csname sga@target@IV.17.1.6\endcsname{ega4.c04.prop.17-1-6}
\expandafter\def\csname sga@target@IV.17.10.1\endcsname{ega4.c04.def.17-10-1}
\expandafter\def\csname sga@target@IV.17.10.2\endcsname{ega4.c04.prop.17-10-2}
\expandafter\def\csname sga@target@IV.17.10.3\endcsname{ega4.c04.cor.17-10-3}
\expandafter\def\csname sga@target@IV.17.10.3.1\endcsname{ega4.c04.cor.17-10-3-1}
\expandafter\def\csname sga@target@IV.17.10.4\endcsname{ega4.c04.cor.17-10-4}
\expandafter\def\csname sga@target@IV.17.10.4.1\endcsname{ega4.c04.eq.17-10-4-1}
\expandafter\def\csname sga@target@IV.17.11.1\endcsname{ega4.c04.thm.17-11-1}
\expandafter\def\csname sga@target@IV.17.11.1.1\endcsname{ega4.c04.eq.17-11-1-1}
\expandafter\def\csname sga@target@IV.17.11.1.2\endcsname{ega4.c04.eq.17-11-1-2}
\expandafter\def\csname sga@target@IV.17.11.2\endcsname{ega4.c04.cor.17-11-2}
\expandafter\def\csname sga@target@IV.17.11.3\endcsname{ega4.c04.cor.17-11-3}
\expandafter\def\csname sga@target@IV.17.11.4\endcsname{ega4.c04.cor.17-11-4}
\expandafter\def\csname sga@target@IV.17.11.5\endcsname{ega4.c04.prop.17-11-5}
\expandafter\def\csname sga@target@IV.17.12.1\endcsname{ega4.c04.thm.17-12-1}
\expandafter\def\csname sga@target@IV.17.12.2\endcsname{ega4.c04.cor.17-12-2}
\expandafter\def\csname sga@target@IV.17.12.3\endcsname{ega4.c04.cor.17-12-3}
\expandafter\def\csname sga@target@IV.17.12.4\endcsname{ega4.c04.prop.17-12-4}
\expandafter\def\csname sga@target@IV.17.12.5\endcsname{ega4.c04.prop.17-12-5}
\expandafter\def\csname sga@target@IV.17.12.6\endcsname{ega4.c04.cor.17-12-6}
\expandafter\def\csname sga@target@IV.17.12.7\endcsname{ega4.c04.ex.17-12-7}
\expandafter\def\csname sga@target@IV.17.12.8\endcsname{ega4.c04.rem.17-12-8}
\expandafter\def\csname sga@target@IV.17.13\endcsname{ega4.c04.subsec.17-13.key-iv-17-13}
\expandafter\def\csname sga@target@IV.17.13.1\endcsname{ega4.c04.statement.17-13-1}
\expandafter\def\csname sga@target@IV.17.13.1.1\endcsname{ega4.c04.statement.17-13-1-1}
\expandafter\def\csname sga@target@IV.17.13.1.2\endcsname{ega4.c04.statement.17-13-1-2}
\expandafter\def\csname sga@target@IV.17.13.1.3\endcsname{ega4.c04.eq.17-13-1-3}
\expandafter\def\csname sga@target@IV.17.13.10\endcsname{ega4.c04.cor.17-13-10}
\expandafter\def\csname sga@target@IV.17.13.11\endcsname{ega4.c04.statement.17-13-11}
\expandafter\def\csname sga@target@IV.17.13.11.1\endcsname{ega4.c04.statement.17-13-11-1}
\expandafter\def\csname sga@target@IV.17.13.12\endcsname{ega4.c04.cor.17-13-12}
\expandafter\def\csname sga@target@IV.17.13.12.1\endcsname{ega4.c04.cor.17-13-12-1}
\expandafter\def\csname sga@target@IV.17.13.12.2\endcsname{ega4.c04.eq.17-13-12-2}
\expandafter\def\csname sga@target@IV.17.13.2\endcsname{ega4.c04.prop.17-13-2}
\expandafter\def\csname sga@target@IV.17.13.2.1\endcsname{ega4.c04.eq.17-13-2-1}
\expandafter\def\csname sga@target@IV.17.13.2.2\endcsname{ega4.c04.eq.17-13-2-2}
\expandafter\def\csname sga@target@IV.17.13.3\endcsname{ega4.c04.def.17-13-3}
\expandafter\def\csname sga@target@IV.17.13.4\endcsname{ega4.c04.statement.17-13-4}
\expandafter\def\csname sga@target@IV.17.13.5\endcsname{ega4.c04.statement.17-13-5}
\expandafter\def\csname sga@target@IV.17.13.5.1\endcsname{ega4.c04.statement.17-13-5-1}
\expandafter\def\csname sga@target@IV.17.13.5.2\endcsname{ega4.c04.eq.17-13-5-2}
\expandafter\def\csname sga@target@IV.17.13.5.3\endcsname{ega4.c04.eq.17-13-5-3}
\expandafter\def\csname sga@target@IV.17.13.5.4\endcsname{ega4.c04.eq.17-13-5-4}
\expandafter\def\csname sga@target@IV.17.13.6\endcsname{ega4.c04.cor.17-13-6}
\expandafter\def\csname sga@target@IV.17.13.6.1\endcsname{ega4.c04.eq.17-13-6-1}
\expandafter\def\csname sga@target@IV.17.13.6.2\endcsname{ega4.c04.eq.17-13-6-2}
\expandafter\def\csname sga@target@IV.17.13.7\endcsname{ega4.c04.statement.17-13-7}
\expandafter\def\csname sga@target@IV.17.13.7.1\endcsname{ega4.c04.eq.17-13-7-1}
\expandafter\def\csname sga@target@IV.17.13.8\endcsname{ega4.c04.prop.17-13-8}
\expandafter\def\csname sga@target@IV.17.13.8.1\endcsname{ega4.c04.eq.17-13-8-1}
\expandafter\def\csname sga@target@IV.17.13.9\endcsname{ega4.c04.statement.17-13-9}
\expandafter\def\csname sga@target@IV.17.13.9.1\endcsname{ega4.c04.eq.17-13-9-1}
\expandafter\def\csname sga@target@IV.17.13.9.2\endcsname{ega4.c04.eq.17-13-9-2}
\expandafter\def\csname sga@target@IV.17.14.1\endcsname{ega4.c04.prop.17-14-1}
\expandafter\def\csname sga@target@IV.17.14.2\endcsname{ega4.c04.prop.17-14-2}
\expandafter\def\csname sga@target@IV.17.14.3\endcsname{ega4.c04.rem.17-14-3}
\expandafter\def\csname sga@target@IV.17.15.1\endcsname{ega4.c04.prop.17-15-1}
\expandafter\def\csname sga@target@IV.17.15.10\endcsname{ega4.c04.statement.17-15-10}
\expandafter\def\csname sga@target@IV.17.15.11\endcsname{ega4.c04.prop.17-15-11}
\expandafter\def\csname sga@target@IV.17.15.12\endcsname{ega4.c04.prop.17-15-12}
\expandafter\def\csname sga@target@IV.17.15.13\endcsname{ega4.c04.cor.17-15-13}
\expandafter\def\csname sga@target@IV.17.15.14\endcsname{ega4.c04.prop.17-15-14}
\expandafter\def\csname sga@target@IV.17.15.14.1\endcsname{ega4.c04.lem.17-15-14-1}
\expandafter\def\csname sga@target@IV.17.15.14.2\endcsname{ega4.c04.lem.17-15-14-2}
\expandafter\def\csname sga@target@IV.17.15.15\endcsname{ega4.c04.prop.17-15-15}
\expandafter\def\csname sga@target@IV.17.15.2\endcsname{ega4.c04.cor.17-15-2}
\expandafter\def\csname sga@target@IV.17.15.3\endcsname{ega4.c04.prop.17-15-3}
\expandafter\def\csname sga@target@IV.17.15.4\endcsname{ega4.c04.cor.17-15-4}
\expandafter\def\csname sga@target@IV.17.15.5\endcsname{ega4.c04.prop.17-15-5}
\expandafter\def\csname sga@target@IV.17.15.6\endcsname{ega4.c04.cor.17-15-6}
\expandafter\def\csname sga@target@IV.17.15.7\endcsname{ega4.c04.cor.17-15-7}
\expandafter\def\csname sga@target@IV.17.15.8\endcsname{ega4.c04.prop.17-15-8}
\expandafter\def\csname sga@target@IV.17.15.8.1\endcsname{ega4.c04.eq.17-15-8-1}
\expandafter\def\csname sga@target@IV.17.15.8.2\endcsname{ega4.c04.eq.17-15-8-2}
\expandafter\def\csname sga@target@IV.17.15.9\endcsname{ega4.c04.cor.17-15-9}
\expandafter\def\csname sga@target@IV.17.16\endcsname{ega4.c04.subsec.17-16.key-iv-17-16}
\expandafter\def\csname sga@target@IV.17.16.1\endcsname{ega4.c04.prop.17-16-1}
\expandafter\def\csname sga@target@IV.17.16.2\endcsname{ega4.c04.cor.17-16-2}
\expandafter\def\csname sga@target@IV.17.16.3\endcsname{ega4.c04.cor.17-16-3}
\expandafter\def\csname sga@target@IV.17.16.4\endcsname{ega4.c04.prop.17-16-4}
\expandafter\def\csname sga@target@IV.17.16.5\endcsname{ega4.c04.cor.17-16-5}
\expandafter\def\csname sga@target@IV.17.16.6\endcsname{ega4.c04.prop.17-16-6}
\expandafter\def\csname sga@target@IV.17.2.1\endcsname{ega4.c04.prop.17-2-1}
\expandafter\def\csname sga@target@IV.17.2.2\endcsname{ega4.c04.cor.17-2-2}
\expandafter\def\csname sga@target@IV.17.2.3\endcsname{ega4.c04.prop.17-2-3}
\expandafter\def\csname sga@target@IV.17.2.3.1\endcsname{ega4.c04.eq.17-2-3-1}
\expandafter\def\csname sga@target@IV.17.2.4\endcsname{ega4.c04.cor.17-2-4}
\expandafter\def\csname sga@target@IV.17.2.5\endcsname{ega4.c04.prop.17-2-5}
\expandafter\def\csname sga@target@IV.17.2.5.1\endcsname{ega4.c04.prop.17-2-5-1}
\expandafter\def\csname sga@target@IV.17.2.6\endcsname{ega4.c04.prop.17-2-6}
\expandafter\def\csname sga@target@IV.17.3.1\endcsname{ega4.c04.def.17-3-1}
\expandafter\def\csname sga@target@IV.17.3.2\endcsname{ega4.c04.subsec.17-3-2}
\expandafter\def\csname sga@target@IV.17.3.3\endcsname{ega4.c04.prop.17-3-3}
\expandafter\def\csname sga@target@IV.17.3.4\endcsname{ega4.c04.prop.17-3-4}
\expandafter\def\csname sga@target@IV.17.3.5\endcsname{ega4.c04.cor.17-3-5}
\expandafter\def\csname sga@target@IV.17.3.6\endcsname{ega4.c04.prop.17-3-6}
\expandafter\def\csname sga@target@IV.17.3.7\endcsname{ega4.c04.def.17-3-7}
\expandafter\def\csname sga@target@IV.17.3.8\endcsname{ega4.c04.prop.17-3-8}
\expandafter\def\csname sga@target@IV.17.3.9\endcsname{ega4.c04.cor.17-3-9}
\expandafter\def\csname sga@target@IV.17.4.1\endcsname{ega4.c04.thm.17-4-1}
\expandafter\def\csname sga@target@IV.17.4.1.1\endcsname{ega4.c04.eq.17-4-1-1}
\expandafter\def\csname sga@target@IV.17.4.1.2\endcsname{ega4.c04.rem.17-4-1-2}
\expandafter\def\csname sga@target@IV.17.4.10\endcsname{ega4.c04.rem.17-4-10}
\expandafter\def\csname sga@target@IV.17.4.2\endcsname{ega4.c04.cor.17-4-2}
\expandafter\def\csname sga@target@IV.17.4.3\endcsname{ega4.c04.cor.17-4-3}
\expandafter\def\csname sga@target@IV.17.4.4\endcsname{ega4.c04.prop.17-4-4}
\expandafter\def\csname sga@target@IV.17.4.5\endcsname{ega4.c04.statement.17-4-5}
\expandafter\def\csname sga@target@IV.17.4.6\endcsname{ega4.c04.prop.17-4-6}
\expandafter\def\csname sga@target@IV.17.4.7\endcsname{ega4.c04.cor.17-4-7}
\expandafter\def\csname sga@target@IV.17.4.8\endcsname{ega4.c04.cor.17-4-8}
\expandafter\def\csname sga@target@IV.17.4.9\endcsname{ega4.c04.prop.17-4-9}
\expandafter\def\csname sga@target@IV.17.5.1\endcsname{ega4.c04.thm.17-5-1}
\expandafter\def\csname sga@target@IV.17.5.1.1\endcsname{ega4.c04.eq.17-5-1-1}
\expandafter\def\csname sga@target@IV.17.5.1.2\endcsname{ega4.c04.eq.17-5-1-2}
\expandafter\def\csname sga@target@IV.17.5.2\endcsname{ega4.c04.cor.17-5-2}
\expandafter\def\csname sga@target@IV.17.5.3\endcsname{ega4.c04.prop.17-5-3}
\expandafter\def\csname sga@target@IV.17.5.4\endcsname{ega4.c04.rem.17-5-4}
\expandafter\def\csname sga@target@IV.17.5.5\endcsname{ega4.c04.prop.17-5-5}
\expandafter\def\csname sga@target@IV.17.5.6\endcsname{ega4.c04.cor.17-5-6}
\expandafter\def\csname sga@target@IV.17.5.7\endcsname{ega4.c04.prop.17-5-7}
\expandafter\def\csname sga@target@IV.17.5.8\endcsname{ega4.c04.prop.17-5-8}
\expandafter\def\csname sga@target@IV.17.6.1\endcsname{ega4.c04.thm.17-6-1}
\expandafter\def\csname sga@target@IV.17.6.2\endcsname{ega4.c04.cor.17-6-2}
\expandafter\def\csname sga@target@IV.17.6.3\endcsname{ega4.c04.prop.17-6-3}
\expandafter\def\csname sga@target@IV.17.6.4\endcsname{ega4.c04.prop.17-6-4}
\expandafter\def\csname sga@target@IV.17.7\endcsname{ega4.c04.subsec.17-7.key-iv-17-7}
\expandafter\def\csname sga@target@IV.17.7.1\endcsname{ega4.c04.prop.17-7-1}
\expandafter\def\csname sga@target@IV.17.7.10\endcsname{ega4.c04.prop.17-7-10}
\expandafter\def\csname sga@target@IV.17.7.11\endcsname{ega4.c04.prop.17-7-11}
\expandafter\def\csname sga@target@IV.17.7.11.1\endcsname{ega4.c04.lem.17-7-11-1}
\expandafter\def\csname sga@target@IV.17.7.2\endcsname{ega4.c04.cor.17-7-2}
\expandafter\def\csname sga@target@IV.17.7.3\endcsname{ega4.c04.cor.17-7-3}
\expandafter\def\csname sga@target@IV.17.7.4\endcsname{ega4.c04.prop.17-7-4}
\expandafter\def\csname sga@target@IV.17.7.5\endcsname{ega4.c04.lem.17-7-5}
\expandafter\def\csname sga@target@IV.17.7.6\endcsname{ega4.c04.rem.17-7-6}
\expandafter\def\csname sga@target@IV.17.7.7\endcsname{ega4.c04.prop.17-7-7}
\expandafter\def\csname sga@target@IV.17.7.8\endcsname{ega4.c04.prop.17-7-8}
\expandafter\def\csname sga@target@IV.17.7.9\endcsname{ega4.c04.cor.17-7-9}
\expandafter\def\csname sga@target@IV.17.8\endcsname{ega4.c04.subsec.17-8.key-iv-17-8}
\expandafter\def\csname sga@target@IV.17.8.1\endcsname{ega4.c04.prop.17-8-1}
\expandafter\def\csname sga@target@IV.17.8.2\endcsname{ega4.c04.prop.17-8-2}
\expandafter\def\csname sga@target@IV.17.8.3\endcsname{ega4.c04.rem.17-8-3}
\expandafter\def\csname sga@target@IV.17.9\endcsname{ega4.c04.subsec.17-9.key-iv-17-9}
\expandafter\def\csname sga@target@IV.17.9.1\endcsname{ega4.c04.thm.17-9-1}
\expandafter\def\csname sga@target@IV.17.9.2\endcsname{ega4.c04.cor.17-9-2}
\expandafter\def\csname sga@target@IV.17.9.3\endcsname{ega4.c04.cor.17-9-3}
\expandafter\def\csname sga@target@IV.17.9.4\endcsname{ega4.c04.cor.17-9-4}
\expandafter\def\csname sga@target@IV.17.9.5\endcsname{ega4.c04.cor.17-9-5}
\expandafter\def\csname sga@target@IV.17.9.6\endcsname{ega4.c04.prop.17-9-6}
\expandafter\def\csname sga@target@IV.18\endcsname{ega4.c04.sec.18.key-iv-18}
\expandafter\def\csname sga@target@IV.18.1.1\endcsname{ega4.c04.prop.18-1-1}
\expandafter\def\csname sga@target@IV.18.1.1.1\endcsname{ega4.c04.eq.18-1-1-1}
\expandafter\def\csname sga@target@IV.18.1.2\endcsname{ega4.c04.thm.18-1-2}
\expandafter\def\csname sga@target@IV.18.1.3\endcsname{ega4.c04.cor.18-1-3}
\expandafter\def\csname sga@target@IV.18.10.1\endcsname{ega4.c04.thm.18-10-1}
\expandafter\def\csname sga@target@IV.18.10.10\endcsname{ega4.c04.statement.18-10-10}
\expandafter\def\csname sga@target@IV.18.10.11\endcsname{ega4.c04.rem.18-10-11}
\expandafter\def\csname sga@target@IV.18.10.12\endcsname{ega4.c04.cor.18-10-12}
\expandafter\def\csname sga@target@IV.18.10.13\endcsname{ega4.c04.prop.18-10-13}
\expandafter\def\csname sga@target@IV.18.10.14\endcsname{ega4.c04.cor.18-10-14}
\expandafter\def\csname sga@target@IV.18.10.15\endcsname{ega4.c04.statement.18-10-15}
\expandafter\def\csname sga@target@IV.18.10.16\endcsname{ega4.c04.thm.18-10-16}
\expandafter\def\csname sga@target@IV.18.10.17\endcsname{ega4.c04.rem.18-10-17}
\expandafter\def\csname sga@target@IV.18.10.17.1\endcsname{ega4.c04.statement.18-10-17-1}
\expandafter\def\csname sga@target@IV.18.10.17.2\endcsname{ega4.c04.statement.18-10-17-2}
\expandafter\def\csname sga@target@IV.18.10.18\endcsname{ega4.c04.lem.18-10-18}
\expandafter\def\csname sga@target@IV.18.10.19\endcsname{ega4.c04.prop.18-10-19}
\expandafter\def\csname sga@target@IV.18.10.2\endcsname{ega4.c04.rem.18-10-2}
\expandafter\def\csname sga@target@IV.18.10.20\endcsname{ega4.c04.rem.18-10-20}
\expandafter\def\csname sga@target@IV.18.10.21\endcsname{ega4.c04.prop.18-10-21}
\expandafter\def\csname sga@target@IV.18.10.3\endcsname{ega4.c04.cor.18-10-3}
\expandafter\def\csname sga@target@IV.18.10.3.1\endcsname{ega4.c04.rem.18-10-3-1}
\expandafter\def\csname sga@target@IV.18.10.4\endcsname{ega4.c04.cor.18-10-4}
\expandafter\def\csname sga@target@IV.18.10.5\endcsname{ega4.c04.cor.18-10-5}
\expandafter\def\csname sga@target@IV.18.10.6\endcsname{ega4.c04.lem.18-10-6}
\expandafter\def\csname sga@target@IV.18.10.7\endcsname{ega4.c04.prop.18-10-7}
\expandafter\def\csname sga@target@IV.18.10.8\endcsname{ega4.c04.prop.18-10-8}
\expandafter\def\csname sga@target@IV.18.10.9\endcsname{ega4.c04.cor.18-10-9}
\expandafter\def\csname sga@target@IV.18.11.1\endcsname{ega4.c04.lem.18-11-1}
\expandafter\def\csname sga@target@IV.18.11.10\endcsname{ega4.c04.thm.18-11-10}
\expandafter\def\csname sga@target@IV.18.11.11\endcsname{ega4.c04.rem.18-11-11}
\expandafter\def\csname sga@target@IV.18.11.12\endcsname{ega4.c04.cor.18-11-12}
\expandafter\def\csname sga@target@IV.18.11.12.1\endcsname{ega4.c04.lem.18-11-12-1}
\expandafter\def\csname sga@target@IV.18.11.2\endcsname{ega4.c04.lem.18-11-2}
\expandafter\def\csname sga@target@IV.18.11.2.1\endcsname{ega4.c04.eq.18-11-2-1}
\expandafter\def\csname sga@target@IV.18.11.2.2\endcsname{ega4.c04.eq.18-11-2-2}
\expandafter\def\csname sga@target@IV.18.11.3\endcsname{ega4.c04.lem.18-11-3}
\expandafter\def\csname sga@target@IV.18.11.3.1\endcsname{ega4.c04.eq.18-11-3-1}
\expandafter\def\csname sga@target@IV.18.11.3.2\endcsname{ega4.c04.eq.18-11-3-2}
\expandafter\def\csname sga@target@IV.18.11.3.3\endcsname{ega4.c04.lem.18-11-3-3}
\expandafter\def\csname sga@target@IV.18.11.3.4\endcsname{ega4.c04.lem.18-11-3-4}
\expandafter\def\csname sga@target@IV.18.11.3.5\endcsname{ega4.c04.lem.18-11-3-5}
\expandafter\def\csname sga@target@IV.18.11.3.6\endcsname{ega4.c04.eq.18-11-3-6}
\expandafter\def\csname sga@target@IV.18.11.4\endcsname{ega4.c04.lem.18-11-4}
\expandafter\def\csname sga@target@IV.18.11.5\endcsname{ega4.c04.prop.18-11-5}
\expandafter\def\csname sga@target@IV.18.11.5.1\endcsname{ega4.c04.lem.18-11-5-1}
\expandafter\def\csname sga@target@IV.18.11.5.2\endcsname{ega4.c04.lem.18-11-5-2}
\expandafter\def\csname sga@target@IV.18.11.6\endcsname{ega4.c04.rem.18-11-6}
\expandafter\def\csname sga@target@IV.18.11.7\endcsname{ega4.c04.cor.18-11-7}
\expandafter\def\csname sga@target@IV.18.11.8\endcsname{ega4.c04.cor.18-11-8}
\expandafter\def\csname sga@target@IV.18.11.9\endcsname{ega4.c04.cor.18-11-9}
\expandafter\def\csname sga@target@IV.18.12.1\endcsname{ega4.c04.thm.18-12-1}
\expandafter\def\csname sga@target@IV.18.12.10\endcsname{ega4.c04.cor.18-12-10}
\expandafter\def\csname sga@target@IV.18.12.11\endcsname{ega4.c04.cor.18-12-11}
\expandafter\def\csname sga@target@IV.18.12.12\endcsname{ega4.c04.prop.18-12-12}
\expandafter\def\csname sga@target@IV.18.12.13\endcsname{ega4.c04.cor.18-12-13}
\expandafter\def\csname sga@target@IV.18.12.14\endcsname{ega4.c04.rem.18-12-14}
\expandafter\def\csname sga@target@IV.18.12.15\endcsname{ega4.c04.prop.18-12-15}
\expandafter\def\csname sga@target@IV.18.12.16\endcsname{ega4.c04.prop.18-12-16}
\expandafter\def\csname sga@target@IV.18.12.17\endcsname{ega4.c04.cor.18-12-17}
\expandafter\def\csname sga@target@IV.18.12.2\endcsname{ega4.c04.rem.18-12-2}
\expandafter\def\csname sga@target@IV.18.12.3\endcsname{ega4.c04.cor.18-12-3}
\expandafter\def\csname sga@target@IV.18.12.4\endcsname{ega4.c04.cor.18-12-4}
\expandafter\def\csname sga@target@IV.18.12.5\endcsname{ega4.c04.rem.18-12-5}
\expandafter\def\csname sga@target@IV.18.12.6\endcsname{ega4.c04.cor.18-12-6}
\expandafter\def\csname sga@target@IV.18.12.7\endcsname{ega4.c04.prop.18-12-7}
\expandafter\def\csname sga@target@IV.18.12.7.1\endcsname{ega4.c04.lem.18-12-7-1}
\expandafter\def\csname sga@target@IV.18.12.8\endcsname{ega4.c04.prop.18-12-8}
\expandafter\def\csname sga@target@IV.18.12.9\endcsname{ega4.c04.rem.18-12-9}
\expandafter\def\csname sga@target@IV.18.2.1\endcsname{ega4.c04.statement.18-2-1}
\expandafter\def\csname sga@target@IV.18.2.1.1\endcsname{ega4.c04.statement.18-2-1-1}
\expandafter\def\csname sga@target@IV.18.2.1.2\endcsname{ega4.c04.eq.18-2-1-2}
\expandafter\def\csname sga@target@IV.18.2.10\endcsname{ega4.c04.statement.18-2-10}
\expandafter\def\csname sga@target@IV.18.2.2\endcsname{ega4.c04.statement.18-2-2}
\expandafter\def\csname sga@target@IV.18.2.2.1\endcsname{ega4.c04.eq.18-2-2-1}
\expandafter\def\csname sga@target@IV.18.2.2.2\endcsname{ega4.c04.eq.18-2-2-2}
\expandafter\def\csname sga@target@IV.18.2.2.3\endcsname{ega4.c04.eq.18-2-2-3}
\expandafter\def\csname sga@target@IV.18.2.2.4\endcsname{ega4.c04.eq.18-2-2-4}
\expandafter\def\csname sga@target@IV.18.2.3\endcsname{ega4.c04.prop.18-2-3}
\expandafter\def\csname sga@target@IV.18.2.4\endcsname{ega4.c04.cor.18-2-4}
\expandafter\def\csname sga@target@IV.18.2.5\endcsname{ega4.c04.cor.18-2-5}
\expandafter\def\csname sga@target@IV.18.2.6\endcsname{ega4.c04.cor.18-2-6}
\expandafter\def\csname sga@target@IV.18.2.7\endcsname{ega4.c04.item.18-2-7}
\expandafter\def\csname sga@target@IV.18.2.7.1\endcsname{ega4.c04.eq.18-2-7-1}
\expandafter\def\csname sga@target@IV.18.2.7.2\endcsname{ega4.c04.eq.18-2-7-2}
\expandafter\def\csname sga@target@IV.18.2.8\endcsname{ega4.c04.prop.18-2-8}
\expandafter\def\csname sga@target@IV.18.2.9\endcsname{ega4.c04.cor.18-2-9}
\expandafter\def\csname sga@target@IV.18.3.1\endcsname{ega4.c04.prop.18-3-1}
\expandafter\def\csname sga@target@IV.18.3.2\endcsname{ega4.c04.prop.18-3-2}
\expandafter\def\csname sga@target@IV.18.3.2.1\endcsname{ega4.c04.lem.18-3-2-1}
\expandafter\def\csname sga@target@IV.18.3.2.2\endcsname{ega4.c04.eq.18-3-2-2}
\expandafter\def\csname sga@target@IV.18.3.3\endcsname{ega4.c04.statement.18-3-3}
\expandafter\def\csname sga@target@IV.18.3.4\endcsname{ega4.c04.thm.18-3-4}
\expandafter\def\csname sga@target@IV.18.4.1\endcsname{ega4.c04.lem.18-4-1}
\expandafter\def\csname sga@target@IV.18.4.10\endcsname{ega4.c04.prop.18-4-10}
\expandafter\def\csname sga@target@IV.18.4.11\endcsname{ega4.c04.statement.18-4-11}
\expandafter\def\csname sga@target@IV.18.4.12\endcsname{ega4.c04.cor.18-4-12}
\expandafter\def\csname sga@target@IV.18.4.13\endcsname{ega4.c04.cor.18-4-13}
\expandafter\def\csname sga@target@IV.18.4.14\endcsname{ega4.c04.cor.18-4-14}
\expandafter\def\csname sga@target@IV.18.4.2\endcsname{ega4.c04.prop.18-4-2}
\expandafter\def\csname sga@target@IV.18.4.3\endcsname{ega4.c04.cor.18-4-3}
\expandafter\def\csname sga@target@IV.18.4.4\endcsname{ega4.c04.lem.18-4-4}
\expandafter\def\csname sga@target@IV.18.4.5\endcsname{ega4.c04.prop.18-4-5}
\expandafter\def\csname sga@target@IV.18.4.6\endcsname{ega4.c04.thm.18-4-6}
\expandafter\def\csname sga@target@IV.18.4.6.1\endcsname{ega4.c04.lem.18-4-6-1}
\expandafter\def\csname sga@target@IV.18.4.7\endcsname{ega4.c04.cor.18-4-7}
\expandafter\def\csname sga@target@IV.18.4.8\endcsname{ega4.c04.cor.18-4-8}
\expandafter\def\csname sga@target@IV.18.4.9\endcsname{ega4.c04.prop.18-4-9}
\expandafter\def\csname sga@target@IV.18.5.1\endcsname{ega4.c04.statement.18-5-1}
\expandafter\def\csname sga@target@IV.18.5.1.1\endcsname{ega4.c04.eq.18-5-1-1}
\expandafter\def\csname sga@target@IV.18.5.10\endcsname{ega4.c04.prop.18-5-10}
\expandafter\def\csname sga@target@IV.18.5.11\endcsname{ega4.c04.thm.18-5-11}
\expandafter\def\csname sga@target@IV.18.5.12\endcsname{ega4.c04.cor.18-5-12}
\expandafter\def\csname sga@target@IV.18.5.13\endcsname{ega4.c04.rem.18-5-13}
\expandafter\def\csname sga@target@IV.18.5.13.1\endcsname{ega4.c04.lem.18-5-13-1}
\expandafter\def\csname sga@target@IV.18.5.14\endcsname{ega4.c04.prop.18-5-14}
\expandafter\def\csname sga@target@IV.18.5.15\endcsname{ega4.c04.prop.18-5-15}
\expandafter\def\csname sga@target@IV.18.5.16\endcsname{ega4.c04.item.18-5-16}
\expandafter\def\csname sga@target@IV.18.5.17\endcsname{ega4.c04.thm.18-5-17}
\expandafter\def\csname sga@target@IV.18.5.18\endcsname{ega4.c04.rem.18-5-18}
\expandafter\def\csname sga@target@IV.18.5.19\endcsname{ega4.c04.prop.18-5-19}
\expandafter\def\csname sga@target@IV.18.5.2\endcsname{ega4.c04.statement.18-5-2}
\expandafter\def\csname sga@target@IV.18.5.3\endcsname{ega4.c04.lem.18-5-3}
\expandafter\def\csname sga@target@IV.18.5.3.1\endcsname{ega4.c04.lem.18-5-3-1}
\expandafter\def\csname sga@target@IV.18.5.4\endcsname{ega4.c04.prop.18-5-4}
\expandafter\def\csname sga@target@IV.18.5.4.1\endcsname{ega4.c04.eq.18-5-4-1}
\expandafter\def\csname sga@target@IV.18.5.4.2\endcsname{ega4.c04.eq.18-5-4-2}
\expandafter\def\csname sga@target@IV.18.5.4.3\endcsname{ega4.c04.lem.18-5-4-3}
\expandafter\def\csname sga@target@IV.18.5.4.4\endcsname{ega4.c04.rem.18-5-4-4}
\expandafter\def\csname sga@target@IV.18.5.5\endcsname{ega4.c04.def.18-5-5}
\expandafter\def\csname sga@target@IV.18.5.6\endcsname{ega4.c04.prop.18-5-6}
\expandafter\def\csname sga@target@IV.18.5.7\endcsname{ega4.c04.rem.18-5-7}
\expandafter\def\csname sga@target@IV.18.5.8\endcsname{ega4.c04.def.18-5-8}
\expandafter\def\csname sga@target@IV.18.5.9\endcsname{ega4.c04.prop.18-5-9}
\expandafter\def\csname sga@target@IV.18.6.1\endcsname{ega4.c04.statement.18-6-1}
\expandafter\def\csname sga@target@IV.18.6.10\endcsname{ega4.c04.cor.18-6-10}
\expandafter\def\csname sga@target@IV.18.6.11\endcsname{ega4.c04.cor.18-6-11}
\expandafter\def\csname sga@target@IV.18.6.12\endcsname{ega4.c04.prop.18-6-12}
\expandafter\def\csname sga@target@IV.18.6.13\endcsname{ega4.c04.cor.18-6-13}
\expandafter\def\csname sga@target@IV.18.6.14\endcsname{ega4.c04.prop.18-6-14}
\expandafter\def\csname sga@target@IV.18.6.15\endcsname{ega4.c04.cor.18-6-15}
\expandafter\def\csname sga@target@IV.18.6.2\endcsname{ega4.c04.lem.18-6-2}
\expandafter\def\csname sga@target@IV.18.6.3\endcsname{ega4.c04.lem.18-6-3}
\expandafter\def\csname sga@target@IV.18.6.4\endcsname{ega4.c04.statement.18-6-4}
\expandafter\def\csname sga@target@IV.18.6.5\endcsname{ega4.c04.def.18-6-5}
\expandafter\def\csname sga@target@IV.18.6.6\endcsname{ega4.c04.thm.18-6-6}
\expandafter\def\csname sga@target@IV.18.6.7\endcsname{ega4.c04.statement.18-6-7}
\expandafter\def\csname sga@target@IV.18.6.8\endcsname{ega4.c04.prop.18-6-8}
\expandafter\def\csname sga@target@IV.18.6.9\endcsname{ega4.c04.thm.18-6-9}
\expandafter\def\csname sga@target@IV.18.6.9.1\endcsname{ega4.c04.lem.18-6-9-1}
\expandafter\def\csname sga@target@IV.18.6.9.2\endcsname{ega4.c04.lem.18-6-9-2}
\expandafter\def\csname sga@target@IV.18.7.1\endcsname{ega4.c04.statement.18-7-1}
\expandafter\def\csname sga@target@IV.18.7.2\endcsname{ega4.c04.prop.18-7-2}
\expandafter\def\csname sga@target@IV.18.7.3\endcsname{ega4.c04.cor.18-7-3}
\expandafter\def\csname sga@target@IV.18.7.4\endcsname{ega4.c04.cor.18-7-4}
\expandafter\def\csname sga@target@IV.18.7.5\endcsname{ega4.c04.prop.18-7-5}
\expandafter\def\csname sga@target@IV.18.7.5.1\endcsname{ega4.c04.lem.18-7-5-1}
\expandafter\def\csname sga@target@IV.18.7.6\endcsname{ega4.c04.cor.18-7-6}
\expandafter\def\csname sga@target@IV.18.7.7\endcsname{ega4.c04.rem.18-7-7}
\expandafter\def\csname sga@target@IV.18.8.1\endcsname{ega4.c04.prop.18-8-1}
\expandafter\def\csname sga@target@IV.18.8.10\endcsname{ega4.c04.prop.18-8-10}
\expandafter\def\csname sga@target@IV.18.8.11\endcsname{ega4.c04.rem.18-8-11}
\expandafter\def\csname sga@target@IV.18.8.12\endcsname{ega4.c04.prop.18-8-12}
\expandafter\def\csname sga@target@IV.18.8.13\endcsname{ega4.c04.cor.18-8-13}
\expandafter\def\csname sga@target@IV.18.8.14\endcsname{ega4.c04.cor.18-8-14}
\expandafter\def\csname sga@target@IV.18.8.15\endcsname{ega4.c04.prop.18-8-15}
\expandafter\def\csname sga@target@IV.18.8.16\endcsname{ega4.c04.cor.18-8-16}
\expandafter\def\csname sga@target@IV.18.8.17\endcsname{ega4.c04.prop.18-8-17}
\expandafter\def\csname sga@target@IV.18.8.18\endcsname{ega4.c04.prop.18-8-18}
\expandafter\def\csname sga@target@IV.18.8.2\endcsname{ega4.c04.def.18-8-2}
\expandafter\def\csname sga@target@IV.18.8.3\endcsname{ega4.c04.rem.18-8-3}
\expandafter\def\csname sga@target@IV.18.8.4\endcsname{ega4.c04.lem.18-8-4}
\expandafter\def\csname sga@target@IV.18.8.5\endcsname{ega4.c04.lem.18-8-5}
\expandafter\def\csname sga@target@IV.18.8.6\endcsname{ega4.c04.statement.18-8-6}
\expandafter\def\csname sga@target@IV.18.8.6.1\endcsname{ega4.c04.eq.18-8-6-1}
\expandafter\def\csname sga@target@IV.18.8.6.2\endcsname{ega4.c04.eq.18-8-6-2}
\expandafter\def\csname sga@target@IV.18.8.7\endcsname{ega4.c04.def.18-8-7}
\expandafter\def\csname sga@target@IV.18.8.8\endcsname{ega4.c04.prop.18-8-8}
\expandafter\def\csname sga@target@IV.18.8.8.1\endcsname{ega4.c04.statement.18-8-8-1}
\expandafter\def\csname sga@target@IV.18.8.9\endcsname{ega4.c04.statement.18-8-9}
\expandafter\def\csname sga@target@IV.18.9.1\endcsname{ega4.c04.thm.18-9-1}
\expandafter\def\csname sga@target@IV.18.9.10\endcsname{ega4.c04.rem.18-9-10}
\expandafter\def\csname sga@target@IV.18.9.11\endcsname{ega4.c04.cor.18-9-11}
\expandafter\def\csname sga@target@IV.18.9.2\endcsname{ega4.c04.cor.18-9-2}
\expandafter\def\csname sga@target@IV.18.9.3\endcsname{ega4.c04.cor.18-9-3}
\expandafter\def\csname sga@target@IV.18.9.4\endcsname{ega4.c04.cor.18-9-4}
\expandafter\def\csname sga@target@IV.18.9.5\endcsname{ega4.c04.cor.18-9-5}
\expandafter\def\csname sga@target@IV.18.9.5.1\endcsname{ega4.c04.eq.18-9-5-1}
\expandafter\def\csname sga@target@IV.18.9.6\endcsname{ega4.c04.statement.18-9-6}
\expandafter\def\csname sga@target@IV.18.9.7\endcsname{ega4.c04.thm.18-9-7}
\expandafter\def\csname sga@target@IV.18.9.7.1\endcsname{ega4.c04.rem.18-9-7-1}
\expandafter\def\csname sga@target@IV.18.9.7.2\endcsname{ega4.c04.statement.18-9-7-2}
\expandafter\def\csname sga@target@IV.18.9.7.3\endcsname{ega4.c04.lem.18-9-7-3}
\expandafter\def\csname sga@target@IV.18.9.7.4\endcsname{ega4.c04.eq.18-9-7-4}
\expandafter\def\csname sga@target@IV.18.9.7.5\endcsname{ega4.c04.lem.18-9-7-5}
\expandafter\def\csname sga@target@IV.18.9.7.6\endcsname{ega4.c04.rem.18-9-7-6}
\expandafter\def\csname sga@target@IV.18.9.7.7\endcsname{ega4.c04.lem.18-9-7-7}
\expandafter\def\csname sga@target@IV.18.9.7.8\endcsname{ega4.c04.rem.18-9-7-8}
\expandafter\def\csname sga@target@IV.18.9.8\endcsname{ega4.c04.cor.18-9-8}
\expandafter\def\csname sga@target@IV.18.9.9\endcsname{ega4.c04.cor.18-9-9}
\expandafter\def\csname sga@target@IV.19.1.1\endcsname{ega4.c04.prop.19-1-1}
\expandafter\def\csname sga@target@IV.19.1.2\endcsname{ega4.c04.cor.19-1-2}
\expandafter\def\csname sga@target@IV.19.1.3\endcsname{ega4.c04.def.19-1-3}
\expandafter\def\csname sga@target@IV.19.1.4\endcsname{ega4.c04.prop.19-1-4}
\expandafter\def\csname sga@target@IV.19.1.4.1\endcsname{ega4.c04.eq.19-1-4-1}
\expandafter\def\csname sga@target@IV.19.1.5\endcsname{ega4.c04.prop.19-1-5}
\expandafter\def\csname sga@target@IV.19.1.5.1\endcsname{ega4.c04.item.19-1-5-1}
\expandafter\def\csname sga@target@IV.19.1.6\endcsname{ega4.c04.cor.19-1-6}
\expandafter\def\csname sga@target@IV.19.1.6.1\endcsname{ega4.c04.lem.19-1-6-1}
\expandafter\def\csname sga@target@IV.19.1.7\endcsname{ega4.c04.item.19-1-7}
\expandafter\def\csname sga@target@IV.19.1.8\endcsname{ega4.c04.statement.19-1-8}
\expandafter\def\csname sga@target@IV.19.2.1\endcsname{ega4.c04.def.19-2-1}
\expandafter\def\csname sga@target@IV.19.2.2\endcsname{ega4.c04.def.19-2-2}
\expandafter\def\csname sga@target@IV.19.2.3\endcsname{ega4.c04.prop.19-2-3}
\expandafter\def\csname sga@target@IV.19.2.4\endcsname{ega4.c04.prop.19-2-4}
\expandafter\def\csname sga@target@IV.19.2.5\endcsname{ega4.c04.item.19-2-5}
\expandafter\def\csname sga@target@IV.19.2.6\endcsname{ega4.c04.cor.19-2-6}
\expandafter\def\csname sga@target@IV.19.2.7\endcsname{ega4.c04.prop.19-2-7}
\expandafter\def\csname sga@target@IV.19.2.8\endcsname{ega4.c04.prop.19-2-8}
\expandafter\def\csname sga@target@IV.19.2.9\endcsname{ega4.c04.prop.19-2-9}
\expandafter\def\csname sga@target@IV.19.3.1\endcsname{ega4.c04.def.19-3-1}
\expandafter\def\csname sga@target@IV.19.3.10\endcsname{ega4.c04.cor.19-3-10}
\expandafter\def\csname sga@target@IV.19.3.2\endcsname{ega4.c04.prop.19-3-2}
\expandafter\def\csname sga@target@IV.19.3.2.1\endcsname{ega4.c04.lem.19-3-2-1}
\expandafter\def\csname sga@target@IV.19.3.3\endcsname{ega4.c04.cor.19-3-3}
\expandafter\def\csname sga@target@IV.19.3.4\endcsname{ega4.c04.cor.19-3-4}
\expandafter\def\csname sga@target@IV.19.3.5\endcsname{ega4.c04.cor.19-3-5}
\expandafter\def\csname sga@target@IV.19.3.6\endcsname{ega4.c04.def.19-3-6}
\expandafter\def\csname sga@target@IV.19.3.7\endcsname{ega4.c04.prop.19-3-7}
\expandafter\def\csname sga@target@IV.19.3.8\endcsname{ega4.c04.cor.19-3-8}
\expandafter\def\csname sga@target@IV.19.3.9\endcsname{ega4.c04.prop.19-3-9}
\expandafter\def\csname sga@target@IV.19.4.1\endcsname{ega4.c04.statement.19-4-1}
\expandafter\def\csname sga@target@IV.19.4.1.1\endcsname{ega4.c04.subsec.19-4-1-1}
\expandafter\def\csname sga@target@IV.19.4.10\endcsname{ega4.c04.rem.19-4-10}
\expandafter\def\csname sga@target@IV.19.4.11\endcsname{ega4.c04.prop.19-4-11}
\expandafter\def\csname sga@target@IV.19.4.11.1\endcsname{ega4.c04.eq.19-4-11-1}
\expandafter\def\csname sga@target@IV.19.4.2\endcsname{ega4.c04.lem.19-4-2}
\expandafter\def\csname sga@target@IV.19.4.3\endcsname{ega4.c04.cor.19-4-3}
\expandafter\def\csname sga@target@IV.19.4.4\endcsname{ega4.c04.prop.19-4-4}
\expandafter\def\csname sga@target@IV.19.4.5\endcsname{ega4.c04.cor.19-4-5}
\expandafter\def\csname sga@target@IV.19.4.6\endcsname{ega4.c04.prop.19-4-6}
\expandafter\def\csname sga@target@IV.19.4.7\endcsname{ega4.c04.cor.19-4-7}
\expandafter\def\csname sga@target@IV.19.4.8\endcsname{ega4.c04.prop.19-4-8}
\expandafter\def\csname sga@target@IV.19.4.9\endcsname{ega4.c04.cor.19-4-9}
\expandafter\def\csname sga@target@IV.19.5.1\endcsname{ega4.c04.thm.19-5-1}
\expandafter\def\csname sga@target@IV.19.5.1.1\endcsname{ega4.c04.eq.19-5-1-1}
\expandafter\def\csname sga@target@IV.19.5.2\endcsname{ega4.c04.cor.19-5-2}
\expandafter\def\csname sga@target@IV.19.5.3\endcsname{ega4.c04.prop.19-5-3}
\expandafter\def\csname sga@target@IV.19.5.4\endcsname{ega4.c04.statement.19-5-4}
\expandafter\def\csname sga@target@IV.19.5.5\endcsname{ega4.c04.thm.19-5-5}
\expandafter\def\csname sga@target@IV.19.5.5.1\endcsname{ega4.c04.statement.19-5-5-1}
\expandafter\def\csname sga@target@IV.19.5.5.2\endcsname{ega4.c04.statement.19-5-5-2}
\expandafter\def\csname sga@target@IV.19.5.5.3\endcsname{ega4.c04.statement.19-5-5-3}
\expandafter\def\csname sga@target@IV.19.5.5.4\endcsname{ega4.c04.statement.19-5-5-4}
\expandafter\def\csname sga@target@IV.19.5.5.5\endcsname{ega4.c04.statement.19-5-5-5}
\expandafter\def\csname sga@target@IV.19.5.6\endcsname{ega4.c04.statement.19-5-6}
\expandafter\def\csname sga@target@IV.19.6.1\endcsname{ega4.c04.statement.19-6-1}
\expandafter\def\csname sga@target@IV.19.6.1.1\endcsname{ega4.c04.eq.19-6-1-1}
\expandafter\def\csname sga@target@IV.19.6.2\endcsname{ega4.c04.statement.19-6-2}
\expandafter\def\csname sga@target@IV.19.6.2.1\endcsname{ega4.c04.statement.19-6-2-1}
\expandafter\def\csname sga@target@IV.19.6.3\endcsname{ega4.c04.prop.19-6-3}
\expandafter\def\csname sga@target@IV.19.6.4\endcsname{ega4.c04.cor.19-6-4}
\expandafter\def\csname sga@target@IV.19.7\endcsname{ega4.c04.subsec.19-7.key-iv-19-7}
\expandafter\def\csname sga@target@IV.19.7.1\endcsname{ega4.c04.thm.19-7-1}
\expandafter\def\csname sga@target@IV.19.7.1.1\endcsname{ega4.c04.statement.19-7-1-1}
\expandafter\def\csname sga@target@IV.19.7.1.10\endcsname{ega4.c04.statement.19-7-1-10}
\expandafter\def\csname sga@target@IV.19.7.1.11\endcsname{ega4.c04.statement.19-7-1-11}
\expandafter\def\csname sga@target@IV.19.7.1.12\endcsname{ega4.c04.statement.19-7-1-12}
\expandafter\def\csname sga@target@IV.19.7.1.13\endcsname{ega4.c04.statement.19-7-1-13}
\expandafter\def\csname sga@target@IV.19.7.1.14\endcsname{ega4.c04.statement.19-7-1-14}
\expandafter\def\csname sga@target@IV.19.7.1.15\endcsname{ega4.c04.statement.19-7-1-15}
\expandafter\def\csname sga@target@IV.19.7.1.16\endcsname{ega4.c04.statement.19-7-1-16}
\expandafter\def\csname sga@target@IV.19.7.1.17\endcsname{ega4.c04.statement.19-7-1-17}
\expandafter\def\csname sga@target@IV.19.7.1.18\endcsname{ega4.c04.statement.19-7-1-18}
\expandafter\def\csname sga@target@IV.19.7.1.19\endcsname{ega4.c04.statement.19-7-1-19}
\expandafter\def\csname sga@target@IV.19.7.1.2\endcsname{ega4.c04.statement.19-7-1-2}
\expandafter\def\csname sga@target@IV.19.7.1.20\endcsname{ega4.c04.statement.19-7-1-20}
\expandafter\def\csname sga@target@IV.19.7.1.21\endcsname{ega4.c04.statement.19-7-1-21}
\expandafter\def\csname sga@target@IV.19.7.1.22\endcsname{ega4.c04.statement.19-7-1-22}
\expandafter\def\csname sga@target@IV.19.7.1.23\endcsname{ega4.c04.statement.19-7-1-23}
\expandafter\def\csname sga@target@IV.19.7.1.3\endcsname{ega4.c04.statement.19-7-1-3}
\expandafter\def\csname sga@target@IV.19.7.1.5\endcsname{ega4.c04.statement.19-7-1-5}
\expandafter\def\csname sga@target@IV.19.7.1.6\endcsname{ega4.c04.statement.19-7-1-6}
\expandafter\def\csname sga@target@IV.19.7.1.7\endcsname{ega4.c04.lem.19-7-1-7}
\expandafter\def\csname sga@target@IV.19.7.1.8\endcsname{ega4.c04.statement.19-7-1-8}
\expandafter\def\csname sga@target@IV.19.7.1.9\endcsname{ega4.c04.statement.19-7-1-9}
\expandafter\def\csname sga@target@IV.19.7.2\endcsname{ega4.c04.statement.19-7-2}
\expandafter\def\csname sga@target@IV.19.7.3\endcsname{ega4.c04.cor.19-7-3}
\expandafter\def\csname sga@target@IV.19.7.3.1\endcsname{ega4.c04.eq.19-7-3-1}
\expandafter\def\csname sga@target@IV.19.7.4\endcsname{ega4.c04.cor.19-7-4}
\expandafter\def\csname sga@target@IV.19.7.4.1\endcsname{ega4.c04.eq.19-7-4-1}
\expandafter\def\csname sga@target@IV.19.7.5\endcsname{ega4.c04.cor.19-7-5}
\expandafter\def\csname sga@target@IV.19.8.1\endcsname{ega4.c04.prop.19-8-1}
\expandafter\def\csname sga@target@IV.19.8.2\endcsname{ega4.c04.prop.19-8-2}
\expandafter\def\csname sga@target@IV.19.8.3\endcsname{ega4.c04.rem.19-8-3}
\expandafter\def\csname sga@target@IV.19.9.1\endcsname{ega4.c04.statement.19-9-1}
\expandafter\def\csname sga@target@IV.19.9.1.1\endcsname{ega4.c04.eq.19-9-1-1}
\expandafter\def\csname sga@target@IV.19.9.1.2\endcsname{ega4.c04.eq.19-9-1-2}
\expandafter\def\csname sga@target@IV.19.9.1.3\endcsname{ega4.c04.eq.19-9-1-3}
\expandafter\def\csname sga@target@IV.19.9.2\endcsname{ega4.c04.lem.19-9-2}
\expandafter\def\csname sga@target@IV.19.9.3\endcsname{ega4.c04.prop.19-9-3}
\expandafter\def\csname sga@target@IV.19.9.4\endcsname{ega4.c04.cor.19-9-4}
\expandafter\def\csname sga@target@IV.19.9.5\endcsname{ega4.c04.prop.19-9-5}
\expandafter\def\csname sga@target@IV.19.9.6\endcsname{ega4.c04.prop.19-9-6}
\expandafter\def\csname sga@target@IV.19.9.7\endcsname{ega4.c04.cor.19-9-7}
\expandafter\def\csname sga@target@IV.19.9.8\endcsname{ega4.c04.prop.19-9-8}
\expandafter\def\csname sga@target@IV.19.9.9\endcsname{ega4.c04.rem.19-9-9}
\expandafter\def\csname sga@target@IV.2.1.1\endcsname{ega4.c04.statement.2-1-1}
\expandafter\def\csname sga@target@IV.2.1.10\endcsname{ega4.c04.cor.2-1-10}
\expandafter\def\csname sga@target@IV.2.1.11\endcsname{ega4.c04.prop.2-1-11}
\expandafter\def\csname sga@target@IV.2.1.12\endcsname{ega4.c04.prop.2-1-12}
\expandafter\def\csname sga@target@IV.2.1.13\endcsname{ega4.c04.prop.2-1-13}
\expandafter\def\csname sga@target@IV.2.1.14\endcsname{ega4.c04.prop.2-1-14}
\expandafter\def\csname sga@target@IV.2.1.2\endcsname{ega4.c04.prop.2-1-2}
\expandafter\def\csname sga@target@IV.2.1.3\endcsname{ega4.c04.prop.2-1-3}
\expandafter\def\csname sga@target@IV.2.1.4\endcsname{ega4.c04.prop.2-1-4}
\expandafter\def\csname sga@target@IV.2.1.5\endcsname{ega4.c04.prop.2-1-5}
\expandafter\def\csname sga@target@IV.2.1.6\endcsname{ega4.c04.cor.2-1-6}
\expandafter\def\csname sga@target@IV.2.1.7\endcsname{ega4.c04.cor.2-1-7}
\expandafter\def\csname sga@target@IV.2.1.8\endcsname{ega4.c04.prop.2-1-8}
\expandafter\def\csname sga@target@IV.2.1.9\endcsname{ega4.c04.cor.2-1-9}
\expandafter\def\csname sga@target@IV.2.2.1\endcsname{ega4.c04.prop.2-2-1}
\expandafter\def\csname sga@target@IV.2.2.10\endcsname{ega4.c04.prop.2-2-10}
\expandafter\def\csname sga@target@IV.2.2.10.1\endcsname{ega4.c04.eq.2-2-10-1}
\expandafter\def\csname sga@target@IV.2.2.11\endcsname{ega4.c04.cor.2-2-11}
\expandafter\def\csname sga@target@IV.2.2.12\endcsname{ega4.c04.cor.2-2-12}
\expandafter\def\csname sga@target@IV.2.2.13\endcsname{ega4.c04.cor.2-2-13}
\expandafter\def\csname sga@target@IV.2.2.14\endcsname{ega4.c04.prop.2-2-14}
\expandafter\def\csname sga@target@IV.2.2.15\endcsname{ega4.c04.prop.2-2-15}
\expandafter\def\csname sga@target@IV.2.2.16\endcsname{ega4.c04.cor.2-2-16}
\expandafter\def\csname sga@target@IV.2.2.17\endcsname{ega4.c04.prop.2-2-17}
\expandafter\def\csname sga@target@IV.2.2.2\endcsname{ega4.c04.prop.2-2-2}
\expandafter\def\csname sga@target@IV.2.2.2.1\endcsname{ega4.c04.eq.2-2-2-1}
\expandafter\def\csname sga@target@IV.2.2.3\endcsname{ega4.c04.cor.2-2-3}
\expandafter\def\csname sga@target@IV.2.2.4\endcsname{ega4.c04.def.2-2-4}
\expandafter\def\csname sga@target@IV.2.2.5\endcsname{ega4.c04.cor.2-2-5}
\expandafter\def\csname sga@target@IV.2.2.6\endcsname{ega4.c04.cor.2-2-6}
\expandafter\def\csname sga@target@IV.2.2.7\endcsname{ega4.c04.prop.2-2-7}
\expandafter\def\csname sga@target@IV.2.2.8\endcsname{ega4.c04.cor.2-2-8}
\expandafter\def\csname sga@target@IV.2.2.8.1\endcsname{ega4.c04.eq.2-2-8-1}
\expandafter\def\csname sga@target@IV.2.2.9\endcsname{ega4.c04.rem.2-2-9}
\expandafter\def\csname sga@target@IV.2.3.1\endcsname{ega4.c04.lem.2-3-1}
\expandafter\def\csname sga@target@IV.2.3.1.1\endcsname{ega4.c04.eq.2-3-1-1}
\expandafter\def\csname sga@target@IV.2.3.10\endcsname{ega4.c04.thm.2-3-10}
\expandafter\def\csname sga@target@IV.2.3.11\endcsname{ega4.c04.cor.2-3-11}
\expandafter\def\csname sga@target@IV.2.3.12\endcsname{ega4.c04.cor.2-3-12}
\expandafter\def\csname sga@target@IV.2.3.13\endcsname{ega4.c04.cor.2-3-13}
\expandafter\def\csname sga@target@IV.2.3.14\endcsname{ega4.c04.cor.2-3-14}
\expandafter\def\csname sga@target@IV.2.3.15\endcsname{ega4.c04.rem.2-3-15}
\expandafter\def\csname sga@target@IV.2.3.2\endcsname{ega4.c04.prop.2-3-2}
\expandafter\def\csname sga@target@IV.2.3.3\endcsname{ega4.c04.statement.2-3-3}
\expandafter\def\csname sga@target@IV.2.3.4\endcsname{ega4.c04.prop.2-3-4}
\expandafter\def\csname sga@target@IV.2.3.5\endcsname{ega4.c04.cor.2-3-5}
\expandafter\def\csname sga@target@IV.2.3.6\endcsname{ega4.c04.prop.2-3-6}
\expandafter\def\csname sga@target@IV.2.3.7\endcsname{ega4.c04.prop.2-3-7}
\expandafter\def\csname sga@target@IV.2.3.8\endcsname{ega4.c04.cor.2-3-8}
\expandafter\def\csname sga@target@IV.2.3.9\endcsname{ega4.c04.rem.2-3-9}
\expandafter\def\csname sga@target@IV.2.4.1\endcsname{ega4.c04.statement.2-4-1}
\expandafter\def\csname sga@target@IV.2.4.10\endcsname{ega4.c04.cor.2-4-10}
\expandafter\def\csname sga@target@IV.2.4.11\endcsname{ega4.c04.rem.2-4-11}
\expandafter\def\csname sga@target@IV.2.4.2\endcsname{ega4.c04.def.2-4-2}
\expandafter\def\csname sga@target@IV.2.4.3\endcsname{ega4.c04.prop.2-4-3}
\expandafter\def\csname sga@target@IV.2.4.4\endcsname{ega4.c04.prop.2-4-4}
\expandafter\def\csname sga@target@IV.2.4.5\endcsname{ega4.c04.prop.2-4-5}
\expandafter\def\csname sga@target@IV.2.4.6\endcsname{ega4.c04.thm.2-4-6}
\expandafter\def\csname sga@target@IV.2.4.7\endcsname{ega4.c04.cor.2-4-7}
\expandafter\def\csname sga@target@IV.2.4.8\endcsname{ega4.c04.rem.2-4-8}
\expandafter\def\csname sga@target@IV.2.4.9\endcsname{ega4.c04.prop.2-4-9}
\expandafter\def\csname sga@target@IV.2.5.1\endcsname{ega4.c04.prop.2-5-1}
\expandafter\def\csname sga@target@IV.2.5.2\endcsname{ega4.c04.prop.2-5-2}
\expandafter\def\csname sga@target@IV.2.5.3\endcsname{ega4.c04.lem.2-5-3}
\expandafter\def\csname sga@target@IV.2.5.4\endcsname{ega4.c04.rem.2-5-4}
\expandafter\def\csname sga@target@IV.2.5.4.1\endcsname{ega4.c04.lem.2-5-4-1}
\expandafter\def\csname sga@target@IV.2.5.5\endcsname{ega4.c04.prop.2-5-5}
\expandafter\def\csname sga@target@IV.2.5.5.1\endcsname{ega4.c04.statement.2-5-5-1}
\expandafter\def\csname sga@target@IV.2.5.5.2\endcsname{ega4.c04.lem.2-5-5-2}
\expandafter\def\csname sga@target@IV.2.5.5.3\endcsname{ega4.c04.statement.2-5-5-3}
\expandafter\def\csname sga@target@IV.2.5.5.4\endcsname{ega4.c04.rem.2-5-5-4}
\expandafter\def\csname sga@target@IV.2.5.6\endcsname{ega4.c04.lem.2-5-6}
\expandafter\def\csname sga@target@IV.2.5.7\endcsname{ega4.c04.cor.2-5-7}
\expandafter\def\csname sga@target@IV.2.5.8\endcsname{ega4.c04.prop.2-5-8}
\expandafter\def\csname sga@target@IV.2.5.8.1\endcsname{ega4.c04.statement.2-5-8-1}
\expandafter\def\csname sga@target@IV.2.5.8.2\endcsname{ega4.c04.statement.2-5-8-2}
\expandafter\def\csname sga@target@IV.2.5.8.3\endcsname{ega4.c04.statement.2-5-8-3}
\expandafter\def\csname sga@target@IV.2.5.8.4\endcsname{ega4.c04.statement.2-5-8-4}
\expandafter\def\csname sga@target@IV.2.5.8.5\endcsname{ega4.c04.statement.2-5-8-5}
\expandafter\def\csname sga@target@IV.2.5.8.6\endcsname{ega4.c04.statement.2-5-8-6}
\expandafter\def\csname sga@target@IV.2.5.8.7\endcsname{ega4.c04.statement.2-5-8-7}
\expandafter\def\csname sga@target@IV.2.5.8.8\endcsname{ega4.c04.statement.2-5-8-8}
\expandafter\def\csname sga@target@IV.2.5.8.9\endcsname{ega4.c04.statement.2-5-8-9}
\expandafter\def\csname sga@target@IV.2.6.1\endcsname{ega4.c04.prop.2-6-1}
\expandafter\def\csname sga@target@IV.2.6.2\endcsname{ega4.c04.prop.2-6-2}
\expandafter\def\csname sga@target@IV.2.6.3\endcsname{ega4.c04.rem.2-6-3}
\expandafter\def\csname sga@target@IV.2.6.4\endcsname{ega4.c04.cor.2-6-4}
\expandafter\def\csname sga@target@IV.2.7.1\endcsname{ega4.c04.prop.2-7-1}
\expandafter\def\csname sga@target@IV.2.7.1.1\endcsname{ega4.c04.lem.2-7-1-1}
\expandafter\def\csname sga@target@IV.2.7.1.viii\endcsname{ega4.c04.item.2-7-1-viii}
\expandafter\def\csname sga@target@IV.2.7.2\endcsname{ega4.c04.cor.2-7-2}
\expandafter\def\csname sga@target@IV.2.7.3\endcsname{ega4.c04.sec.2-7-3}
\expandafter\def\csname sga@target@IV.2.8.1\endcsname{ega4.c04.prop.2-8-1}
\expandafter\def\csname sga@target@IV.2.8.1.1\endcsname{ega4.c04.lem.2-8-1-1}
\expandafter\def\csname sga@target@IV.2.8.2\endcsname{ega4.c04.cor.2-8-2}
\expandafter\def\csname sga@target@IV.2.8.3\endcsname{ega4.c04.statement.2-8-3}
\expandafter\def\csname sga@target@IV.2.8.4\endcsname{ega4.c04.prop.2-8-4}
\expandafter\def\csname sga@target@IV.2.8.4.1\endcsname{ega4.c04.statement.2-8-4-1}
\expandafter\def\csname sga@target@IV.2.8.5\endcsname{ega4.c04.prop.2-8-5}
\expandafter\def\csname sga@target@IV.2.8.6\endcsname{ega4.c04.cor.2-8-6}
\expandafter\def\csname sga@target@IV.2.8.6.1\endcsname{ega4.c04.statement.2-8-6-1}
\expandafter\def\csname sga@target@IV.20\endcsname{ega4.c04.sec.20}
\expandafter\def\csname sga@target@IV.20.0\endcsname{ega4.c04.subsec.20-0}
\expandafter\def\csname sga@target@IV.20.1\endcsname{ega4.c04.subsec.20-1}
\expandafter\def\csname sga@target@IV.20.1.1\endcsname{ega4.c04.statement.20-1-1}
\expandafter\def\csname sga@target@IV.20.1.1.1\endcsname{ega4.c04.eq.20-1-1-1}
\expandafter\def\csname sga@target@IV.20.1.10\endcsname{ega4.c04.statement.20-1-10}
\expandafter\def\csname sga@target@IV.20.1.11\endcsname{ega4.c04.statement.20-1-11}
\expandafter\def\csname sga@target@IV.20.1.11.1\endcsname{ega4.c04.eq.20-1-11-1}
\expandafter\def\csname sga@target@IV.20.1.11.2\endcsname{ega4.c04.eq.20-1-11-2}
\expandafter\def\csname sga@target@IV.20.1.12\endcsname{ega4.c04.prop.20-1-12}
\expandafter\def\csname sga@target@IV.20.1.13\endcsname{ega4.c04.rem.20-1-13}
\expandafter\def\csname sga@target@IV.20.1.2\endcsname{ega4.c04.statement.20-1-2}
\expandafter\def\csname sga@target@IV.20.1.2.1\endcsname{ega4.c04.statement.20-1-2-1}
\expandafter\def\csname sga@target@IV.20.1.2.2\endcsname{ega4.c04.statement.20-1-2-2}
\expandafter\def\csname sga@target@IV.20.1.3\endcsname{ega4.c04.statement.20-1-3}
\expandafter\def\csname sga@target@IV.20.1.3.1\endcsname{ega4.c04.statement.20-1-3-1}
\expandafter\def\csname sga@target@IV.20.1.4\endcsname{ega4.c04.statement.20-1-4}
\expandafter\def\csname sga@target@IV.20.1.4.1\endcsname{ega4.c04.statement.20-1-4-1}
\expandafter\def\csname sga@target@IV.20.1.5\endcsname{ega4.c04.statement.20-1-5}
\expandafter\def\csname sga@target@IV.20.1.5.1\endcsname{ega4.c04.statement.20-1-5-1}
\expandafter\def\csname sga@target@IV.20.1.6\endcsname{ega4.c04.prop.20-1-6}
\expandafter\def\csname sga@target@IV.20.1.7\endcsname{ega4.c04.statement.20-1-7}
\expandafter\def\csname sga@target@IV.20.1.8\endcsname{ega4.c04.statement.20-1-8}
\expandafter\def\csname sga@target@IV.20.1.9\endcsname{ega4.c04.statement.20-1-9}
\expandafter\def\csname sga@target@IV.20.2\endcsname{ega4.c04.subsec.20-2}
\expandafter\def\csname sga@target@IV.20.2.1\endcsname{ega4.c04.statement.20-2-1}
\expandafter\def\csname sga@target@IV.20.2.10\endcsname{ega4.c04.statement.20-2-10}
\expandafter\def\csname sga@target@IV.20.2.10.1\endcsname{ega4.c04.eq.20-2-10-1}
\expandafter\def\csname sga@target@IV.20.2.10.2\endcsname{ega4.c04.eq.20-2-10-2}
\expandafter\def\csname sga@target@IV.20.2.11\endcsname{ega4.c04.prop.20-2-11}
\expandafter\def\csname sga@target@IV.20.2.11.1\endcsname{ega4.c04.eq.20-2-11-1}
\expandafter\def\csname sga@target@IV.20.2.12\endcsname{ega4.c04.cor.20-2-12}
\expandafter\def\csname sga@target@IV.20.2.13\endcsname{ega4.c04.statement.20-2-13}
\expandafter\def\csname sga@target@IV.20.2.14\endcsname{ega4.c04.statement.20-2-14}
\expandafter\def\csname sga@target@IV.20.2.15\endcsname{ega4.c04.prop.20-2-15}
\expandafter\def\csname sga@target@IV.20.2.16\endcsname{ega4.c04.rem.20-2-16}
\expandafter\def\csname sga@target@IV.20.2.2\endcsname{ega4.c04.statement.20-2-2}
\expandafter\def\csname sga@target@IV.20.2.3\endcsname{ega4.c04.statement.20-2-3}
\expandafter\def\csname sga@target@IV.20.2.3.1\endcsname{ega4.c04.statement.20-2-3-1}
\expandafter\def\csname sga@target@IV.20.2.4\endcsname{ega4.c04.prop.20-2-4}
\expandafter\def\csname sga@target@IV.20.2.5\endcsname{ega4.c04.cor.20-2-5}
\expandafter\def\csname sga@target@IV.20.2.6\endcsname{ega4.c04.cor.20-2-6}
\expandafter\def\csname sga@target@IV.20.2.7\endcsname{ega4.c04.rem.20-2-7}
\expandafter\def\csname sga@target@IV.20.2.8\endcsname{ega4.c04.statement.20-2-8}
\expandafter\def\csname sga@target@IV.20.2.9\endcsname{ega4.c04.lem.20-2-9}
\expandafter\def\csname sga@target@IV.20.3\endcsname{ega4.c04.subsec.20-3}
\expandafter\def\csname sga@target@IV.20.3.1\endcsname{ega4.c04.statement.20-3-1}
\expandafter\def\csname sga@target@IV.20.3.10\endcsname{ega4.c04.lem.20-3-10}
\expandafter\def\csname sga@target@IV.20.3.11\endcsname{ega4.c04.prop.20-3-11}
\expandafter\def\csname sga@target@IV.20.3.11.1\endcsname{ega4.c04.eq.20-3-11-1}
\expandafter\def\csname sga@target@IV.20.3.12\endcsname{ega4.c04.cor.20-3-12}
\expandafter\def\csname sga@target@IV.20.3.13\endcsname{ega4.c04.rem.20-3-13}
\expandafter\def\csname sga@target@IV.20.3.2\endcsname{ega4.c04.statement.20-3-2}
\expandafter\def\csname sga@target@IV.20.3.2.1\endcsname{ega4.c04.eq.20-3-2-1}
\expandafter\def\csname sga@target@IV.20.3.3\endcsname{ega4.c04.statement.20-3-3}
\expandafter\def\csname sga@target@IV.20.3.4\endcsname{ega4.c04.statement.20-3-4}
\expandafter\def\csname sga@target@IV.20.3.5\endcsname{ega4.c04.prop.20-3-5}
\expandafter\def\csname sga@target@IV.20.3.6\endcsname{ega4.c04.statement.20-3-6}
\expandafter\def\csname sga@target@IV.20.3.6.1\endcsname{ega4.c04.eq.20-3-6-1}
\expandafter\def\csname sga@target@IV.20.3.7\endcsname{ega4.c04.prop.20-3-7}
\expandafter\def\csname sga@target@IV.20.3.8\endcsname{ega4.c04.prop.20-3-8}
\expandafter\def\csname sga@target@IV.20.3.8.1\endcsname{ega4.c04.eq.20-3-8-1}
\expandafter\def\csname sga@target@IV.20.3.8.2\endcsname{ega4.c04.lem.20-3-8-2}
\expandafter\def\csname sga@target@IV.20.3.8.3\endcsname{ega4.c04.lem.20-3-8-3}
\expandafter\def\csname sga@target@IV.20.3.8.4\endcsname{ega4.c04.lem.20-3-8-4}
\expandafter\def\csname sga@target@IV.20.3.8.5\endcsname{ega4.c04.sec.20-3-8-5}
\expandafter\def\csname sga@target@IV.20.3.8.6\endcsname{ega4.c04.sec.20-3-8-6}
\expandafter\def\csname sga@target@IV.20.3.8.7\endcsname{ega4.c04.sec.20-3-8-7}
\expandafter\def\csname sga@target@IV.20.3.8.8\endcsname{ega4.c04.lem.20-3-8-8}
\expandafter\def\csname sga@target@IV.20.3.9\endcsname{ega4.c04.cor.20-3-9}
\expandafter\def\csname sga@target@IV.20.4\endcsname{ega4.c04.subsec.20-4}
\expandafter\def\csname sga@target@IV.20.4.1\endcsname{ega4.c04.statement.20-4-1}
\expandafter\def\csname sga@target@IV.20.4.10\endcsname{ega4.c04.cor.20-4-10}
\expandafter\def\csname sga@target@IV.20.4.11\endcsname{ega4.c04.prop.20-4-11}
\expandafter\def\csname sga@target@IV.20.4.12\endcsname{ega4.c04.cor.20-4-12}
\expandafter\def\csname sga@target@IV.20.4.2\endcsname{ega4.c04.statement.20-4-2}
\expandafter\def\csname sga@target@IV.20.4.3\endcsname{ega4.c04.statement.20-4-3}
\expandafter\def\csname sga@target@IV.20.4.4\endcsname{ega4.c04.prop.20-4-4}
\expandafter\def\csname sga@target@IV.20.4.5\endcsname{ega4.c04.prop.20-4-5}
\expandafter\def\csname sga@target@IV.20.4.6\endcsname{ega4.c04.prop.20-4-6}
\expandafter\def\csname sga@target@IV.20.4.7\endcsname{ega4.c04.cor.20-4-7}
\expandafter\def\csname sga@target@IV.20.4.8\endcsname{ega4.c04.rem.20-4-8}
\expandafter\def\csname sga@target@IV.20.4.9\endcsname{ega4.c04.cor.20-4-9}
\expandafter\def\csname sga@target@IV.20.5\endcsname{ega4.c04.subsec.20-5}
\expandafter\def\csname sga@target@IV.20.5.1\endcsname{ega4.c04.statement.20-5-1}
\expandafter\def\csname sga@target@IV.20.5.2\endcsname{ega4.c04.statement.20-5-2}
\expandafter\def\csname sga@target@IV.20.5.2.1\endcsname{ega4.c04.statement.20-5-2-1}
\expandafter\def\csname sga@target@IV.20.5.3\endcsname{ega4.c04.prop.20-5-3}
\expandafter\def\csname sga@target@IV.20.5.4\endcsname{ega4.c04.statement.20-5-4}
\expandafter\def\csname sga@target@IV.20.5.5\endcsname{ega4.c04.statement.20-5-5}
\expandafter\def\csname sga@target@IV.20.5.6\endcsname{ega4.c04.statement.20-5-6}
\expandafter\def\csname sga@target@IV.20.5.7\endcsname{ega4.c04.statement.20-5-7}
\expandafter\def\csname sga@target@IV.20.5.8\endcsname{ega4.c04.statement.20-5-8}
\expandafter\def\csname sga@target@IV.20.6\endcsname{ega4.c04.subsec.20-6}
\expandafter\def\csname sga@target@IV.20.6.1\endcsname{ega4.c04.statement.20-6-1}
\expandafter\def\csname sga@target@IV.20.6.1.1\endcsname{ega4.c04.eq.20-6-1-1}
\expandafter\def\csname sga@target@IV.20.6.2\endcsname{ega4.c04.statement.20-6-2}
\expandafter\def\csname sga@target@IV.20.6.2.1\endcsname{ega4.c04.statement.20-6-2-1}
\expandafter\def\csname sga@target@IV.20.6.2.2\endcsname{ega4.c04.eq.20-6-2-2}
\expandafter\def\csname sga@target@IV.20.6.3\endcsname{ega4.c04.prop.20-6-3}
\expandafter\def\csname sga@target@IV.20.6.4\endcsname{ega4.c04.prop.20-6-4}
\expandafter\def\csname sga@target@IV.20.6.5\endcsname{ega4.c04.statement.20-6-5}
\expandafter\def\csname sga@target@IV.20.6.5.1\endcsname{ega4.c04.eq.20-6-5-1}
\expandafter\def\csname sga@target@IV.20.6.6\endcsname{ega4.c04.prop.20-6-6}
\expandafter\def\csname sga@target@IV.20.6.7\endcsname{ega4.c04.statement.20-6-7}
\expandafter\def\csname sga@target@IV.20.6.8\endcsname{ega4.c04.prop.20-6-8}
\expandafter\def\csname sga@target@IV.20.6.8.1\endcsname{ega4.c04.eq.20-6-8-1}
\expandafter\def\csname sga@target@IV.20.6.9\endcsname{ega4.c04.statement.20-6-9}
\expandafter\def\csname sga@target@IV.20.6.9.1\endcsname{ega4.c04.statement.20-6-9-1}
\expandafter\def\csname sga@target@IV.21.1\endcsname{ega4.c04.subsec.21-1.key-iv-21-1}
\expandafter\def\csname sga@target@IV.21.1.1\endcsname{ega4.c04.statement.21-1-1}
\expandafter\def\csname sga@target@IV.21.1.10\endcsname{ega4.c04.prop.21-1-10}
\expandafter\def\csname sga@target@IV.21.1.2\endcsname{ega4.c04.statement.21-1-2}
\expandafter\def\csname sga@target@IV.21.1.3\endcsname{ega4.c04.statement.21-1-3}
\expandafter\def\csname sga@target@IV.21.1.3.1\endcsname{ega4.c04.statement.21-1-3-1}
\expandafter\def\csname sga@target@IV.21.1.3.2\endcsname{ega4.c04.statement.21-1-3-2}
\expandafter\def\csname sga@target@IV.21.1.3.3\endcsname{ega4.c04.statement.21-1-3-3}
\expandafter\def\csname sga@target@IV.21.1.3.4\endcsname{ega4.c04.statement.21-1-3-4}
\expandafter\def\csname sga@target@IV.21.1.4\endcsname{ega4.c04.statement.21-1-4}
\expandafter\def\csname sga@target@IV.21.1.4.1\endcsname{ega4.c04.eq.21-1-4-1}
\expandafter\def\csname sga@target@IV.21.1.4.2\endcsname{ega4.c04.eq.21-1-4-2}
\expandafter\def\csname sga@target@IV.21.1.4.3\endcsname{ega4.c04.eq.21-1-4-3}
\expandafter\def\csname sga@target@IV.21.1.5\endcsname{ega4.c04.statement.21-1-5}
\expandafter\def\csname sga@target@IV.21.1.5.1\endcsname{ega4.c04.statement.21-1-5-1}
\expandafter\def\csname sga@target@IV.21.1.5.2\endcsname{ega4.c04.eq.21-1-5-2}
\expandafter\def\csname sga@target@IV.21.1.6\endcsname{ega4.c04.def.21-1-6}
\expandafter\def\csname sga@target@IV.21.1.6.1\endcsname{ega4.c04.eq.21-1-6-1}
\expandafter\def\csname sga@target@IV.21.1.6.2\endcsname{ega4.c04.eq.21-1-6-2}
\expandafter\def\csname sga@target@IV.21.1.6.3\endcsname{ega4.c04.eq.21-1-6-3}
\expandafter\def\csname sga@target@IV.21.1.6.4\endcsname{ega4.c04.eq.21-1-6-4}
\expandafter\def\csname sga@target@IV.21.1.7\endcsname{ega4.c04.statement.21-1-7}
\expandafter\def\csname sga@target@IV.21.1.7.1\endcsname{ega4.c04.statement.21-1-7-1}
\expandafter\def\csname sga@target@IV.21.1.8\endcsname{ega4.c04.prop.21-1-8}
\expandafter\def\csname sga@target@IV.21.1.9\endcsname{ega4.c04.cor.21-1-9}
\expandafter\def\csname sga@target@IV.21.10.1\endcsname{ega4.c04.statement.21-10-1}
\expandafter\def\csname sga@target@IV.21.10.1.1\endcsname{ega4.c04.eq.21-10-1-1}
\expandafter\def\csname sga@target@IV.21.10.1.2\endcsname{ega4.c04.eq.21-10-1-2}
\expandafter\def\csname sga@target@IV.21.10.10\endcsname{ega4.c04.statement.21-10-10}
\expandafter\def\csname sga@target@IV.21.10.11\endcsname{ega4.c04.statement.21-10-11}
\expandafter\def\csname sga@target@IV.21.10.11.1\endcsname{ega4.c04.eq.21-10-11-1}
\expandafter\def\csname sga@target@IV.21.10.12\endcsname{ega4.c04.rem.21-10-12}
\expandafter\def\csname sga@target@IV.21.10.13\endcsname{ega4.c04.lem.21-10-13}
\expandafter\def\csname sga@target@IV.21.10.13.1\endcsname{ega4.c04.prop.21-10-13-1}
\expandafter\def\csname sga@target@IV.21.10.13bis\endcsname{ega4.c04.prop.21-10-13bis}
\expandafter\def\csname sga@target@IV.21.10.14\endcsname{ega4.c04.statement.21-10-14}
\expandafter\def\csname sga@target@IV.21.10.14.1\endcsname{ega4.c04.eq.21-10-14-1}
\expandafter\def\csname sga@target@IV.21.10.14.2\endcsname{ega4.c04.eq.21-10-14-2}
\expandafter\def\csname sga@target@IV.21.10.15\endcsname{ega4.c04.prop.21-10-15}
\expandafter\def\csname sga@target@IV.21.10.16\endcsname{ega4.c04.prop.21-10-16}
\expandafter\def\csname sga@target@IV.21.10.16.1\endcsname{ega4.c04.prop.21-10-16-1}
\expandafter\def\csname sga@target@IV.21.10.16.2\endcsname{ega4.c04.eq.21-10-16-2}
\expandafter\def\csname sga@target@IV.21.10.17\endcsname{ega4.c04.prop.21-10-17}
\expandafter\def\csname sga@target@IV.21.10.17.1\endcsname{ega4.c04.prop.21-10-17-1}
\expandafter\def\csname sga@target@IV.21.10.17.10\endcsname{ega4.c04.eq.21-10-17-10}
\expandafter\def\csname sga@target@IV.21.10.17.2\endcsname{ega4.c04.eq.21-10-17-2}
\expandafter\def\csname sga@target@IV.21.10.17.3\endcsname{ega4.c04.lem.21-10-17-3}
\expandafter\def\csname sga@target@IV.21.10.17.4\endcsname{ega4.c04.lem.21-10-17-4}
\expandafter\def\csname sga@target@IV.21.10.17.5\endcsname{ega4.c04.eq.21-10-17-5}
\expandafter\def\csname sga@target@IV.21.10.17.6\endcsname{ega4.c04.eq.21-10-17-6}
\expandafter\def\csname sga@target@IV.21.10.17.7\endcsname{ega4.c04.lem.21-10-17-7}
\expandafter\def\csname sga@target@IV.21.10.17.8\endcsname{ega4.c04.lem.21-10-17-8}
\expandafter\def\csname sga@target@IV.21.10.17.9\endcsname{ega4.c04.eq.21-10-17-9}
\expandafter\def\csname sga@target@IV.21.10.18\endcsname{ega4.c04.prop.21-10-18}
\expandafter\def\csname sga@target@IV.21.10.18.1\endcsname{ega4.c04.prop.21-10-18-1}
\expandafter\def\csname sga@target@IV.21.10.19\endcsname{ega4.c04.cor.21-10-19}
\expandafter\def\csname sga@target@IV.21.10.19.1\endcsname{ega4.c04.cor.21-10-19-1}
\expandafter\def\csname sga@target@IV.21.10.2\endcsname{ega4.c04.statement.21-10-2}
\expandafter\def\csname sga@target@IV.21.10.3\endcsname{ega4.c04.statement.21-10-3}
\expandafter\def\csname sga@target@IV.21.10.3.1\endcsname{ega4.c04.statement.21-10-3-1}
\expandafter\def\csname sga@target@IV.21.10.4\endcsname{ega4.c04.prop.21-10-4}
\expandafter\def\csname sga@target@IV.21.10.4.1\endcsname{ega4.c04.prop.21-10-4-1}
\expandafter\def\csname sga@target@IV.21.10.5\endcsname{ega4.c04.statement.21-10-5}
\expandafter\def\csname sga@target@IV.21.10.6\endcsname{ega4.c04.prop.21-10-6}
\expandafter\def\csname sga@target@IV.21.10.7\endcsname{ega4.c04.rem.21-10-7}
\expandafter\def\csname sga@target@IV.21.10.8\endcsname{ega4.c04.prop.21-10-8}
\expandafter\def\csname sga@target@IV.21.10.9\endcsname{ega4.c04.statement.21-10-9}
\expandafter\def\csname sga@target@IV.21.11\endcsname{ega4.c04.subsec.21-11.key-iv-21-11}
\expandafter\def\csname sga@target@IV.21.11.1\endcsname{ega4.c04.thm.21-11-1}
\expandafter\def\csname sga@target@IV.21.11.1.1\endcsname{ega4.c04.eq.21-11-1-1}
\expandafter\def\csname sga@target@IV.21.11.2\endcsname{ega4.c04.lem.21-11-2}
\expandafter\def\csname sga@target@IV.21.12\endcsname{ega4.c04.subsec.21-12.key-iv-21-12}
\expandafter\def\csname sga@target@IV.21.12.1\endcsname{ega4.c04.statement.21-12-1}
\expandafter\def\csname sga@target@IV.21.12.1.1\endcsname{ega4.c04.eq.21-12-1-1}
\expandafter\def\csname sga@target@IV.21.12.10\endcsname{ega4.c04.prop.21-12-10}
\expandafter\def\csname sga@target@IV.21.12.12\endcsname{ega4.c04.thm.21-12-12}
\expandafter\def\csname sga@target@IV.21.12.13\endcsname{ega4.c04.cor.21-12-13}
\expandafter\def\csname sga@target@IV.21.12.14\endcsname{ega4.c04.item.21-12-14}
\expandafter\def\csname sga@target@IV.21.12.15\endcsname{ega4.c04.prop.21-12-15}
\expandafter\def\csname sga@target@IV.21.12.16\endcsname{ega4.c04.cor.21-12-16}
\expandafter\def\csname sga@target@IV.21.12.17\endcsname{ega4.c04.item.21-12-17}
\expandafter\def\csname sga@target@IV.21.12.2\endcsname{ega4.c04.prop.21-12-2}
\expandafter\def\csname sga@target@IV.21.12.2.1\endcsname{ega4.c04.prop.21-12-2-1}
\expandafter\def\csname sga@target@IV.21.12.3\endcsname{ega4.c04.cor.21-12-3}
\expandafter\def\csname sga@target@IV.21.12.4\endcsname{ega4.c04.eq.21-12-4}
\expandafter\def\csname sga@target@IV.21.12.5\endcsname{ega4.c04.statement.21-12-5}
\expandafter\def\csname sga@target@IV.21.12.5.1\endcsname{ega4.c04.eq.21-12-5-1}
\expandafter\def\csname sga@target@IV.21.12.5.2\endcsname{ega4.c04.eq.21-12-5-2}
\expandafter\def\csname sga@target@IV.21.12.5.3\endcsname{ega4.c04.eq.21-12-5-3}
\expandafter\def\csname sga@target@IV.21.12.6\endcsname{ega4.c04.thm.21-12-6}
\expandafter\def\csname sga@target@IV.21.12.7\endcsname{ega4.c04.cor.21-12-7}
\expandafter\def\csname sga@target@IV.21.12.8\endcsname{ega4.c04.statement.21-12-8}
\expandafter\def\csname sga@target@IV.21.12.9\endcsname{ega4.c04.statement.21-12-9}
\expandafter\def\csname sga@target@IV.21.13\endcsname{ega4.c04.subsec.21-13.key-iv-21-13}
\expandafter\def\csname sga@target@IV.21.13.1\endcsname{ega4.c04.def.21-13-1}
\expandafter\def\csname sga@target@IV.21.13.10\endcsname{ega4.c04.prop.21-13-10}
\expandafter\def\csname sga@target@IV.21.13.11\endcsname{ega4.c04.cor.21-13-11}
\expandafter\def\csname sga@target@IV.21.13.12\endcsname{ega4.c04.prop.21-13-12}
\expandafter\def\csname sga@target@IV.21.13.12.1\endcsname{ega4.c04.statement.21-13-12-1}
\expandafter\def\csname sga@target@IV.21.13.13\endcsname{ega4.c04.rem.21-13-13}
\expandafter\def\csname sga@target@IV.21.13.14\endcsname{ega4.c04.prop.21-13-14}
\expandafter\def\csname sga@target@IV.21.13.15\endcsname{ega4.c04.cor.21-13-15}
\expandafter\def\csname sga@target@IV.21.13.16\endcsname{ega4.c04.rem.21-13-16}
\expandafter\def\csname sga@target@IV.21.13.2\endcsname{ega4.c04.lem.21-13-2}
\expandafter\def\csname sga@target@IV.21.13.2.1\endcsname{ega4.c04.eq.21-13-2-1}
\expandafter\def\csname sga@target@IV.21.13.3\endcsname{ega4.c04.cor.21-13-3}
\expandafter\def\csname sga@target@IV.21.13.4\endcsname{ega4.c04.lem.21-13-4}
\expandafter\def\csname sga@target@IV.21.13.5\endcsname{ega4.c04.prop.21-13-5}
\expandafter\def\csname sga@target@IV.21.13.6\endcsname{ega4.c04.cor.21-13-6}
\expandafter\def\csname sga@target@IV.21.13.7\endcsname{ega4.c04.def.21-13-7}
\expandafter\def\csname sga@target@IV.21.13.8\endcsname{ega4.c04.prop.21-13-8}
\expandafter\def\csname sga@target@IV.21.13.9\endcsname{ega4.c04.statement.21-13-9}
\expandafter\def\csname sga@target@IV.21.13.9.1\endcsname{ega4.c04.eq.21-13-9-1}
\expandafter\def\csname sga@target@IV.21.14\endcsname{ega4.c04.subsec.21-14.key-iv-21-14}
\expandafter\def\csname sga@target@IV.21.14.1\endcsname{ega4.c04.thm.21-14-1}
\expandafter\def\csname sga@target@IV.21.14.2\endcsname{ega4.c04.cor.21-14-2}
\expandafter\def\csname sga@target@IV.21.14.3\endcsname{ega4.c04.prop.21-14-3}
\expandafter\def\csname sga@target@IV.21.14.4\endcsname{ega4.c04.statement.21-14-4}
\expandafter\def\csname sga@target@IV.21.14.4bis\endcsname{ega4.c04.statement.21-14-4bis}
\expandafter\def\csname sga@target@IV.21.15\endcsname{ega4.c04.subsec.21-15.key-iv-21-15}
\expandafter\def\csname sga@target@IV.21.15.1\endcsname{ega4.c04.statement.21-15-1}
\expandafter\def\csname sga@target@IV.21.15.2\endcsname{ega4.c04.def.21-15-2}
\expandafter\def\csname sga@target@IV.21.15.3\endcsname{ega4.c04.statement.21-15-3}
\expandafter\def\csname sga@target@IV.21.15.3.1\endcsname{ega4.c04.eq.21-15-3-1}
\expandafter\def\csname sga@target@IV.21.15.3.2\endcsname{ega4.c04.eq.21-15-3-2}
\expandafter\def\csname sga@target@IV.21.15.3bis\endcsname{ega4.c04.statement.21-15-3bis}
\expandafter\def\csname sga@target@IV.21.15.4\endcsname{ega4.c04.prop.21-15-4}
\expandafter\def\csname sga@target@IV.21.15.5\endcsname{ega4.c04.prop.21-15-5}
\expandafter\def\csname sga@target@IV.21.15.6\endcsname{ega4.c04.prop.21-15-6}
\expandafter\def\csname sga@target@IV.21.15.7\endcsname{ega4.c04.statement.21-15-7}
\expandafter\def\csname sga@target@IV.21.15.8\endcsname{ega4.c04.statement.21-15-8}
\expandafter\def\csname sga@target@IV.21.15.9\endcsname{ega4.c04.statement.21-15-9}
\expandafter\def\csname sga@target@IV.21.15.9.1\endcsname{ega4.c04.eq.21-15-9-1}
\expandafter\def\csname sga@target@IV.21.2.1\endcsname{ega4.c04.statement.21-2-1}
\expandafter\def\csname sga@target@IV.21.2.10\endcsname{ega4.c04.statement.21-2-10}
\expandafter\def\csname sga@target@IV.21.2.10.1\endcsname{ega4.c04.eq.21-2-10-1}
\expandafter\def\csname sga@target@IV.21.2.11\endcsname{ega4.c04.prop.21-2-11}
\expandafter\def\csname sga@target@IV.21.2.11.1\endcsname{ega4.c04.prop.21-2-11-1}
\expandafter\def\csname sga@target@IV.21.2.12\endcsname{ega4.c04.statement.21-2-12}
\expandafter\def\csname sga@target@IV.21.2.2\endcsname{ega4.c04.prop.21-2-2}
\expandafter\def\csname sga@target@IV.21.2.3\endcsname{ega4.c04.cor.21-2-3}
\expandafter\def\csname sga@target@IV.21.2.4\endcsname{ega4.c04.statement.21-2-4}
\expandafter\def\csname sga@target@IV.21.2.5\endcsname{ega4.c04.statement.21-2-5}
\expandafter\def\csname sga@target@IV.21.2.5.1\endcsname{ega4.c04.statement.21-2-5-1}
\expandafter\def\csname sga@target@IV.21.2.5.2\endcsname{ega4.c04.eq.21-2-5-2}
\expandafter\def\csname sga@target@IV.21.2.5.3\endcsname{ega4.c04.eq.21-2-5-3}
\expandafter\def\csname sga@target@IV.21.2.5.4\endcsname{ega4.c04.eq.21-2-5-4}
\expandafter\def\csname sga@target@IV.21.2.6\endcsname{ega4.c04.prop.21-2-6}
\expandafter\def\csname sga@target@IV.21.2.6.1\endcsname{ega4.c04.eq.21-2-6-1}
\expandafter\def\csname sga@target@IV.21.2.7\endcsname{ega4.c04.statement.21-2-7}
\expandafter\def\csname sga@target@IV.21.2.7.1\endcsname{ega4.c04.statement.21-2-7-1}
\expandafter\def\csname sga@target@IV.21.2.7.2\endcsname{ega4.c04.eq.21-2-7-2}
\expandafter\def\csname sga@target@IV.21.2.8\endcsname{ega4.c04.statement.21-2-8}
\expandafter\def\csname sga@target@IV.21.2.8.1\endcsname{ega4.c04.statement.21-2-8-1}
\expandafter\def\csname sga@target@IV.21.2.8.2\endcsname{ega4.c04.statement.21-2-8-2}
\expandafter\def\csname sga@target@IV.21.2.8.3\endcsname{ega4.c04.statement.21-2-8-3}
\expandafter\def\csname sga@target@IV.21.2.9\endcsname{ega4.c04.statement.21-2-9}
\expandafter\def\csname sga@target@IV.21.2.9.1\endcsname{ega4.c04.statement.21-2-9-1}
\expandafter\def\csname sga@target@IV.21.2.9.2\endcsname{ega4.c04.eq.21-2-9-2}
\expandafter\def\csname sga@target@IV.21.2.9.3\endcsname{ega4.c04.eq.21-2-9-3}
\expandafter\def\csname sga@target@IV.21.2.9.4\endcsname{ega4.c04.eq.21-2-9-4}
\expandafter\def\csname sga@target@IV.21.3.1\endcsname{ega4.c04.statement.21-3-1}
\expandafter\def\csname sga@target@IV.21.3.2\endcsname{ega4.c04.statement.21-3-2}
\expandafter\def\csname sga@target@IV.21.3.2.1\endcsname{ega4.c04.statement.21-3-2-1}
\expandafter\def\csname sga@target@IV.21.3.2.2\endcsname{ega4.c04.eq.21-3-2-2}
\expandafter\def\csname sga@target@IV.21.3.2.3\endcsname{ega4.c04.eq.21-3-2-3}
\expandafter\def\csname sga@target@IV.21.3.2.4\endcsname{ega4.c04.eq.21-3-2-4}
\expandafter\def\csname sga@target@IV.21.3.3\endcsname{ega4.c04.prop.21-3-3}
\expandafter\def\csname sga@target@IV.21.3.3.1\endcsname{ega4.c04.eq.21-3-3-1}
\expandafter\def\csname sga@target@IV.21.3.4\endcsname{ega4.c04.prop.21-3-4}
\expandafter\def\csname sga@target@IV.21.3.5\endcsname{ega4.c04.cor.21-3-5}
\expandafter\def\csname sga@target@IV.21.3.6\endcsname{ega4.c04.rem.21-3-6}
\expandafter\def\csname sga@target@IV.21.3.6.1\endcsname{ega4.c04.eq.21-3-6-1}
\expandafter\def\csname sga@target@IV.21.4\endcsname{ega4.c04.subsec.21-4.key-iv-21-4}
\expandafter\def\csname sga@target@IV.21.4.1\endcsname{ega4.c04.statement.21-4-1}
\expandafter\def\csname sga@target@IV.21.4.10\endcsname{ega4.c04.cor.21-4-10}
\expandafter\def\csname sga@target@IV.21.4.11\endcsname{ega4.c04.cor.21-4-11}
\expandafter\def\csname sga@target@IV.21.4.12\endcsname{ega4.c04.rem.21-4-12}
\expandafter\def\csname sga@target@IV.21.4.13\endcsname{ega4.c04.cor.21-4-13}
\expandafter\def\csname sga@target@IV.21.4.2\endcsname{ega4.c04.def.21-4-2}
\expandafter\def\csname sga@target@IV.21.4.2.1\endcsname{ega4.c04.eq.21-4-2-1}
\expandafter\def\csname sga@target@IV.21.4.3\endcsname{ega4.c04.statement.21-4-3}
\expandafter\def\csname sga@target@IV.21.4.4\endcsname{ega4.c04.statement.21-4-4}
\expandafter\def\csname sga@target@IV.21.4.5\endcsname{ega4.c04.prop.21-4-5}
\expandafter\def\csname sga@target@IV.21.4.6\endcsname{ega4.c04.cor.21-4-6}
\expandafter\def\csname sga@target@IV.21.4.6.1\endcsname{ega4.c04.eq.21-4-6-1}
\expandafter\def\csname sga@target@IV.21.4.7\endcsname{ega4.c04.statement.21-4-7}
\expandafter\def\csname sga@target@IV.21.4.8\endcsname{ega4.c04.prop.21-4-8}
\expandafter\def\csname sga@target@IV.21.4.9\endcsname{ega4.c04.prop.21-4-9}
\expandafter\def\csname sga@target@IV.21.5\endcsname{ega4.c04.subsec.21-5.key-iv-21-5}
\expandafter\def\csname sga@target@IV.21.5.1\endcsname{ega4.c04.statement.21-5-1}
\expandafter\def\csname sga@target@IV.21.5.2\endcsname{ega4.c04.lem.21-5-2}
\expandafter\def\csname sga@target@IV.21.5.3\endcsname{ega4.c04.statement.21-5-3}
\expandafter\def\csname sga@target@IV.21.5.3.1\endcsname{ega4.c04.eq.21-5-3-1}
\expandafter\def\csname sga@target@IV.21.5.3.2\endcsname{ega4.c04.eq.21-5-3-2}
\expandafter\def\csname sga@target@IV.21.5.4\endcsname{ega4.c04.statement.21-5-4}
\expandafter\def\csname sga@target@IV.21.5.5\endcsname{ega4.c04.def.21-5-5}
\expandafter\def\csname sga@target@IV.21.5.5.1\endcsname{ega4.c04.eq.21-5-5-1}
\expandafter\def\csname sga@target@IV.21.5.5.2\endcsname{ega4.c04.eq.21-5-5-2}
\expandafter\def\csname sga@target@IV.21.5.5.3\endcsname{ega4.c04.eq.21-5-5-3}
\expandafter\def\csname sga@target@IV.21.5.6\endcsname{ega4.c04.prop.21-5-6}
\expandafter\def\csname sga@target@IV.21.5.6.1\endcsname{ega4.c04.prop.21-5-6-1}
\expandafter\def\csname sga@target@IV.21.5.7\endcsname{ega4.c04.prop.21-5-7}
\expandafter\def\csname sga@target@IV.21.5.7.1\endcsname{ega4.c04.prop.21-5-7-1}
\expandafter\def\csname sga@target@IV.21.5.7.2\endcsname{ega4.c04.lem.21-5-7-2}
\expandafter\def\csname sga@target@IV.21.5.7.3\endcsname{ega4.c04.lem.21-5-7-3}
\expandafter\def\csname sga@target@IV.21.5.7.4\endcsname{ega4.c04.eq.21-5-7-4}
\expandafter\def\csname sga@target@IV.21.5.8\endcsname{ega4.c04.prop.21-5-8}
\expandafter\def\csname sga@target@IV.21.5.8.1\endcsname{ega4.c04.eq.21-5-8-1}
\expandafter\def\csname sga@target@IV.21.6\endcsname{ega4.c04.subsec.21-6.key-iv-21-6}
\expandafter\def\csname sga@target@IV.21.6.1\endcsname{ega4.c04.statement.21-6-1}
\expandafter\def\csname sga@target@IV.21.6.10\endcsname{ega4.c04.cor.21-6-10}
\expandafter\def\csname sga@target@IV.21.6.10.1\endcsname{ega4.c04.eq.21-6-10-1}
\expandafter\def\csname sga@target@IV.21.6.11\endcsname{ega4.c04.cor.21-6-11}
\expandafter\def\csname sga@target@IV.21.6.12\endcsname{ega4.c04.prop.21-6-12}
\expandafter\def\csname sga@target@IV.21.6.12.1\endcsname{ega4.c04.eq.21-6-12-1}
\expandafter\def\csname sga@target@IV.21.6.12.2\endcsname{ega4.c04.eq.21-6-12-2}
\expandafter\def\csname sga@target@IV.21.6.13\endcsname{ega4.c04.cor.21-6-13}
\expandafter\def\csname sga@target@IV.21.6.14\endcsname{ega4.c04.cor.21-6-14}
\expandafter\def\csname sga@target@IV.21.6.2\endcsname{ega4.c04.statement.21-6-2}
\expandafter\def\csname sga@target@IV.21.6.3\endcsname{ega4.c04.statement.21-6-3}
\expandafter\def\csname sga@target@IV.21.6.4\endcsname{ega4.c04.statement.21-6-4}
\expandafter\def\csname sga@target@IV.21.6.4.1\endcsname{ega4.c04.statement.21-6-4-1}
\expandafter\def\csname sga@target@IV.21.6.4.2\endcsname{ega4.c04.eq.21-6-4-2}
\expandafter\def\csname sga@target@IV.21.6.5\endcsname{ega4.c04.statement.21-6-5}
\expandafter\def\csname sga@target@IV.21.6.5.1\endcsname{ega4.c04.eq.21-6-5-1}
\expandafter\def\csname sga@target@IV.21.6.6\endcsname{ega4.c04.prop.21-6-6}
\expandafter\def\csname sga@target@IV.21.6.6.2\endcsname{ega4.c04.statement.21-6-6-2}
\expandafter\def\csname sga@target@IV.21.6.7\endcsname{ega4.c04.statement.21-6-7}
\expandafter\def\csname sga@target@IV.21.6.7.1\endcsname{ega4.c04.eq.21-6-7-1}
\expandafter\def\csname sga@target@IV.21.6.7.2\endcsname{ega4.c04.eq.21-6-7-2}
\expandafter\def\csname sga@target@IV.21.6.7.3\endcsname{ega4.c04.eq.21-6-7-3}
\expandafter\def\csname sga@target@IV.21.6.7.4\endcsname{ega4.c04.eq.21-6-7-4}
\expandafter\def\csname sga@target@IV.21.6.7.5\endcsname{ega4.c04.eq.21-6-7-5}
\expandafter\def\csname sga@target@IV.21.6.8\endcsname{ega4.c04.statement.21-6-8}
\expandafter\def\csname sga@target@IV.21.6.9\endcsname{ega4.c04.thm.21-6-9}
\expandafter\def\csname sga@target@IV.21.6.9.1\endcsname{ega4.c04.eq.21-6-9-1}
\expandafter\def\csname sga@target@IV.21.6.9.2\endcsname{ega4.c04.statement.21-6-9-2}
\expandafter\def\csname sga@target@IV.21.7\endcsname{ega4.c04.subsec.21-7.key-iv-21-7}
\expandafter\def\csname sga@target@IV.21.7.1\endcsname{ega4.c04.statement.21-7-1}
\expandafter\def\csname sga@target@IV.21.7.2\endcsname{ega4.c04.prop.21-7-2}
\expandafter\def\csname sga@target@IV.21.7.2.1\endcsname{ega4.c04.eq.21-7-2-1}
\expandafter\def\csname sga@target@IV.21.7.3\endcsname{ega4.c04.cor.21-7-3}
\expandafter\def\csname sga@target@IV.21.7.3.1\endcsname{ega4.c04.statement.21-7-3-1}
\expandafter\def\csname sga@target@IV.21.7.4\endcsname{ega4.c04.prop.21-7-4}
\expandafter\def\csname sga@target@IV.21.7.5\endcsname{ega4.c04.rem.21-7-5}
\expandafter\def\csname sga@target@IV.21.7.6\endcsname{ega4.c04.cor.21-7-6}
\expandafter\def\csname sga@target@IV.21.7.7\endcsname{ega4.c04.rem.21-7-7}
\expandafter\def\csname sga@target@IV.21.8.1\endcsname{ega4.c04.lem.21-8-1}
\expandafter\def\csname sga@target@IV.21.8.2\endcsname{ega4.c04.cor.21-8-2}
\expandafter\def\csname sga@target@IV.21.8.3\endcsname{ega4.c04.prop.21-8-3}
\expandafter\def\csname sga@target@IV.21.8.3.1\endcsname{ega4.c04.eq.21-8-3-1}
\expandafter\def\csname sga@target@IV.21.8.4\endcsname{ega4.c04.cor.21-8-4}
\expandafter\def\csname sga@target@IV.21.8.5\endcsname{ega4.c04.prop.21-8-5}
\expandafter\def\csname sga@target@IV.21.8.5.1\endcsname{ega4.c04.eq.21-8-5-1}
\expandafter\def\csname sga@target@IV.21.8.5.2\endcsname{ega4.c04.eq.21-8-5-2}
\expandafter\def\csname sga@target@IV.21.8.5.3\endcsname{ega4.c04.eq.21-8-5-3}
\expandafter\def\csname sga@target@IV.21.8.6\endcsname{ega4.c04.rem.21-8-6}
\expandafter\def\csname sga@target@IV.21.9\endcsname{ega4.c04.subsec.21-9.key-iv-21-9}
\expandafter\def\csname sga@target@IV.21.9.1\endcsname{ega4.c04.statement.21-9-1}
\expandafter\def\csname sga@target@IV.21.9.1.1\endcsname{ega4.c04.eq.21-9-1-1}
\expandafter\def\csname sga@target@IV.21.9.1.2\endcsname{ega4.c04.eq.21-9-1-2}
\expandafter\def\csname sga@target@IV.21.9.10\endcsname{ega4.c04.statement.21-9-10}
\expandafter\def\csname sga@target@IV.21.9.11\endcsname{ega4.c04.prop.21-9-11}
\expandafter\def\csname sga@target@IV.21.9.12\endcsname{ega4.c04.cor.21-9-12}
\expandafter\def\csname sga@target@IV.21.9.13\endcsname{ega4.c04.rem.21-9-13}
\expandafter\def\csname sga@target@IV.21.9.2\endcsname{ega4.c04.lem.21-9-2}
\expandafter\def\csname sga@target@IV.21.9.3\endcsname{ega4.c04.rem.21-9-3}
\expandafter\def\csname sga@target@IV.21.9.4\endcsname{ega4.c04.prop.21-9-4}
\expandafter\def\csname sga@target@IV.21.9.4.1\endcsname{ega4.c04.prop.21-9-4-1}
\expandafter\def\csname sga@target@IV.21.9.5\endcsname{ega4.c04.cor.21-9-5}
\expandafter\def\csname sga@target@IV.21.9.6\endcsname{ega4.c04.cor.21-9-6}
\expandafter\def\csname sga@target@IV.21.9.7\endcsname{ega4.c04.cor.21-9-7}
\expandafter\def\csname sga@target@IV.21.9.8\endcsname{ega4.c04.statement.21-9-8}
\expandafter\def\csname sga@target@IV.21.9.9\endcsname{ega4.c04.prop.21-9-9}
\expandafter\def\csname sga@target@IV.3.1.1\endcsname{ega4.c04.def.3-1-1}
\expandafter\def\csname sga@target@IV.3.1.1.1\endcsname{ega4.c04.eq.3-1-1-1}
\expandafter\def\csname sga@target@IV.3.1.10\endcsname{ega4.c04.prop.3-1-10}
\expandafter\def\csname sga@target@IV.3.1.10.1\endcsname{ega4.c04.lem.3-1-10-1}
\expandafter\def\csname sga@target@IV.3.1.10.2\endcsname{ega4.c04.rem.3-1-10-2}
\expandafter\def\csname sga@target@IV.3.1.10.3\endcsname{ega4.c04.eq.3-1-10-3}
\expandafter\def\csname sga@target@IV.3.1.10.4\endcsname{ega4.c04.eq.3-1-10-4}
\expandafter\def\csname sga@target@IV.3.1.10.5\endcsname{ega4.c04.rem.3-1-10-5}
\expandafter\def\csname sga@target@IV.3.1.11\endcsname{ega4.c04.cor.3-1-11}
\expandafter\def\csname sga@target@IV.3.1.12\endcsname{ega4.c04.rem.3-1-12}
\expandafter\def\csname sga@target@IV.3.1.13\endcsname{ega4.c04.prop.3-1-13}
\expandafter\def\csname sga@target@IV.3.1.2\endcsname{ega4.c04.prop.3-1-2}
\expandafter\def\csname sga@target@IV.3.1.3\endcsname{ega4.c04.prop.3-1-3}
\expandafter\def\csname sga@target@IV.3.1.4\endcsname{ega4.c04.cor.3-1-4}
\expandafter\def\csname sga@target@IV.3.1.5\endcsname{ega4.c04.cor.3-1-5}
\expandafter\def\csname sga@target@IV.3.1.6\endcsname{ega4.c04.prop.3-1-6}
\expandafter\def\csname sga@target@IV.3.1.7\endcsname{ega4.c04.prop.3-1-7}
\expandafter\def\csname sga@target@IV.3.1.8\endcsname{ega4.c04.prop.3-1-8}
\expandafter\def\csname sga@target@IV.3.1.9\endcsname{ega4.c04.cor.3-1-9}
\expandafter\def\csname sga@target@IV.3.2.1\endcsname{ega4.c04.prop.3-2-1}
\expandafter\def\csname sga@target@IV.3.2.2\endcsname{ega4.c04.def.3-2-2}
\expandafter\def\csname sga@target@IV.3.2.3\endcsname{ega4.c04.prop.3-2-3}
\expandafter\def\csname sga@target@IV.3.2.4\endcsname{ega4.c04.def.3-2-4}
\expandafter\def\csname sga@target@IV.3.2.5\endcsname{ega4.c04.def.3-2-5}
\expandafter\def\csname sga@target@IV.3.2.6\endcsname{ega4.c04.prop.3-2-6}
\expandafter\def\csname sga@target@IV.3.2.7\endcsname{ega4.c04.cor.3-2-7}
\expandafter\def\csname sga@target@IV.3.2.8\endcsname{ega4.c04.cor.3-2-8}
\expandafter\def\csname sga@target@IV.3.2.9\endcsname{ega4.c04.rem.3-2-9}
\expandafter\def\csname sga@target@IV.3.3.1\endcsname{ega4.c04.prop.3-3-1}
\expandafter\def\csname sga@target@IV.3.3.1.1\endcsname{ega4.c04.prop.3-3-1-1}
\expandafter\def\csname sga@target@IV.3.3.2\endcsname{ega4.c04.cor.3-3-2}
\expandafter\def\csname sga@target@IV.3.3.3\endcsname{ega4.c04.cor.3-3-3}
\expandafter\def\csname sga@target@IV.3.3.4\endcsname{ega4.c04.prop.3-3-4}
\expandafter\def\csname sga@target@IV.3.3.5\endcsname{ega4.c04.cor.3-3-5}
\expandafter\def\csname sga@target@IV.3.3.6\endcsname{ega4.c04.prop.3-3-6}
\expandafter\def\csname sga@target@IV.3.3.6.1\endcsname{ega4.c04.eq.3-3-6-1}
\expandafter\def\csname sga@target@IV.3.3.6.2\endcsname{ega4.c04.eq.3-3-6-2}
\expandafter\def\csname sga@target@IV.3.3.6.3\endcsname{ega4.c04.eq.3-3-6-3}
\expandafter\def\csname sga@target@IV.3.3.6.4\endcsname{ega4.c04.eq.3-3-6-4}
\expandafter\def\csname sga@target@IV.3.3.7\endcsname{ega4.c04.cor.3-3-7}
\expandafter\def\csname sga@target@IV.3.3.7.1\endcsname{ega4.c04.cor.3-3-7-1}
\expandafter\def\csname sga@target@IV.3.3.8\endcsname{ega4.c04.sec.3-3-8}
\expandafter\def\csname sga@target@IV.3.3.9\endcsname{ega4.c04.cor.3-3-9}
\expandafter\def\csname sga@target@IV.3.3.9.1\endcsname{ega4.c04.cor.3-3-9-1}
\expandafter\def\csname sga@target@IV.3.4\endcsname{ega4.c04.subsec.3-4.key-iv-3-4}
\expandafter\def\csname sga@target@IV.3.4.1\endcsname{ega4.c04.prop.3-4-1}
\expandafter\def\csname sga@target@IV.3.4.1.1\endcsname{ega4.c04.eq.3-4-1-1}
\expandafter\def\csname sga@target@IV.3.4.1.2\endcsname{ega4.c04.lem.3-4-1-2}
\expandafter\def\csname sga@target@IV.3.4.1.3\endcsname{ega4.c04.lem.3-4-1-3}
\expandafter\def\csname sga@target@IV.3.4.1.4\endcsname{ega4.c04.lem.3-4-1-4}
\expandafter\def\csname sga@target@IV.3.4.1.5\endcsname{ega4.c04.eq.3-4-1-5}
\expandafter\def\csname sga@target@IV.3.4.2\endcsname{ega4.c04.cor.3-4-2}
\expandafter\def\csname sga@target@IV.3.4.3\endcsname{ega4.c04.prop.3-4-3}
\expandafter\def\csname sga@target@IV.3.4.4\endcsname{ega4.c04.prop.3-4-4}
\expandafter\def\csname sga@target@IV.3.4.5\endcsname{ega4.c04.prop.3-4-5}
\expandafter\def\csname sga@target@IV.3.4.6\endcsname{ega4.c04.prop.3-4-6}
\expandafter\def\csname sga@target@IV.3.4.6.1\endcsname{ega4.c04.lem.3-4-6-1}
\expandafter\def\csname sga@target@IV.3.4.7\endcsname{ega4.c04.cor.3-4-7}
\expandafter\def\csname sga@target@IV.3.4.8\endcsname{ega4.c04.rem.3-4-8}
\expandafter\def\csname sga@target@IV.3.4.9\endcsname{ega4.c04.prop.3-4-9}
\expandafter\def\csname sga@target@IV.4.1.1\endcsname{ega4.c04.def.4-1-1}
\expandafter\def\csname sga@target@IV.4.1.1.1\endcsname{ega4.c04.eq.4-1-1-1}
\expandafter\def\csname sga@target@IV.4.1.1.2\endcsname{ega4.c04.eq.4-1-1-2}
\expandafter\def\csname sga@target@IV.4.1.1.3\endcsname{ega4.c04.eq.4-1-1-3}
\expandafter\def\csname sga@target@IV.4.1.2\endcsname{ega4.c04.thm.4-1-2}
\expandafter\def\csname sga@target@IV.4.1.2.1\endcsname{ega4.c04.cor.4-1-2-1}
\expandafter\def\csname sga@target@IV.4.1.3\endcsname{ega4.c04.rem.4-1-3}
\expandafter\def\csname sga@target@IV.4.1.4\endcsname{ega4.c04.cor.4-1-4}
\expandafter\def\csname sga@target@IV.4.1.5\endcsname{ega4.c04.cor.4-1-5}
\expandafter\def\csname sga@target@IV.4.2.1\endcsname{ega4.c04.prop.4-2-1}
\expandafter\def\csname sga@target@IV.4.2.1.1\endcsname{ega4.c04.eq.4-2-1-1}
\expandafter\def\csname sga@target@IV.4.2.1.2\endcsname{ega4.c04.eq.4-2-1-2}
\expandafter\def\csname sga@target@IV.4.2.1.3\endcsname{ega4.c04.eq.4-2-1-3}
\expandafter\def\csname sga@target@IV.4.2.1.4\endcsname{ega4.c04.rem.4-2-1-4}
\expandafter\def\csname sga@target@IV.4.2.1.5\endcsname{ega4.c04.eq.4-2-1-5}
\expandafter\def\csname sga@target@IV.4.2.10\endcsname{ega4.c04.cor.4-2-10}
\expandafter\def\csname sga@target@IV.4.2.2\endcsname{ega4.c04.cor.4-2-2}
\expandafter\def\csname sga@target@IV.4.2.3\endcsname{ega4.c04.cor.4-2-3}
\expandafter\def\csname sga@target@IV.4.2.4\endcsname{ega4.c04.prop.4-2-4}
\expandafter\def\csname sga@target@IV.4.2.4.i\endcsname{ega4.c04.sec.4-2-4-i}
\expandafter\def\csname sga@target@IV.4.2.4.ii\endcsname{ega4.c04.sec.4-2-4-ii}
\expandafter\def\csname sga@target@IV.4.2.4.iii\endcsname{ega4.c04.sec.4-2-4-iii}
\expandafter\def\csname sga@target@IV.4.2.5\endcsname{ega4.c04.prop.4-2-5}
\expandafter\def\csname sga@target@IV.4.2.6\endcsname{ega4.c04.cor.4-2-6}
\expandafter\def\csname sga@target@IV.4.2.7\endcsname{ega4.c04.prop.4-2-7}
\expandafter\def\csname sga@target@IV.4.2.7.i\endcsname{ega4.c04.sec.4-2-7-i}
\expandafter\def\csname sga@target@IV.4.2.7.ii\endcsname{ega4.c04.sec.4-2-7-ii}
\expandafter\def\csname sga@target@IV.4.2.7.iii\endcsname{ega4.c04.sec.4-2-7-iii}
\expandafter\def\csname sga@target@IV.4.2.7.iv\endcsname{ega4.c04.sec.4-2-7-iv}
\expandafter\def\csname sga@target@IV.4.2.8\endcsname{ega4.c04.cor.4-2-8}
\expandafter\def\csname sga@target@IV.4.2.9\endcsname{ega4.c04.prop.4-2-9}
\expandafter\def\csname sga@target@IV.4.3.1\endcsname{ega4.c04.statement.4-3-1}
\expandafter\def\csname sga@target@IV.4.3.2\endcsname{ega4.c04.prop.4-3-2}
\expandafter\def\csname sga@target@IV.4.3.3\endcsname{ega4.c04.cor.4-3-3}
\expandafter\def\csname sga@target@IV.4.3.4\endcsname{ega4.c04.cor.4-3-4}
\expandafter\def\csname sga@target@IV.4.3.5\endcsname{ega4.c04.prop.4-3-5}
\expandafter\def\csname sga@target@IV.4.3.6\endcsname{ega4.c04.cor.4-3-6}
\expandafter\def\csname sga@target@IV.4.3.7\endcsname{ega4.c04.cor.4-3-7}
\expandafter\def\csname sga@target@IV.4.3.8\endcsname{ega4.c04.cor.4-3-8}
\expandafter\def\csname sga@target@IV.4.3.9\endcsname{ega4.c04.prop.4-3-9}
\expandafter\def\csname sga@target@IV.4.4.1\endcsname{ega4.c04.statement.4-4-1}
\expandafter\def\csname sga@target@IV.4.4.2\endcsname{ega4.c04.lem.4-4-2}
\expandafter\def\csname sga@target@IV.4.4.2.i\endcsname{ega4.c04.sec.4-4-2-i}
\expandafter\def\csname sga@target@IV.4.4.2.ii\endcsname{ega4.c04.sec.4-4-2-ii}
\expandafter\def\csname sga@target@IV.4.4.3\endcsname{ega4.c04.rem.4-4-3}
\expandafter\def\csname sga@target@IV.4.4.4\endcsname{ega4.c04.thm.4-4-4}
\expandafter\def\csname sga@target@IV.4.4.5\endcsname{ega4.c04.cor.4-4-5}
\expandafter\def\csname sga@target@IV.4.4.6\endcsname{ega4.c04.cor.4-4-6}
\expandafter\def\csname sga@target@IV.4.5\endcsname{ega4.c04.subsec.4-5.key-iv-4-5}
\expandafter\def\csname sga@target@IV.4.5.1\endcsname{ega4.c04.prop.4-5-1}
\expandafter\def\csname sga@target@IV.4.5.10\endcsname{ega4.c04.cor.4-5-10}
\expandafter\def\csname sga@target@IV.4.5.11\endcsname{ega4.c04.cor.4-5-11}
\expandafter\def\csname sga@target@IV.4.5.12\endcsname{ega4.c04.sec.4-5-12}
\expandafter\def\csname sga@target@IV.4.5.12.i\endcsname{ega4.c04.sec.4-5-12-i}
\expandafter\def\csname sga@target@IV.4.5.12.ii\endcsname{ega4.c04.sec.4-5-12-ii}
\expandafter\def\csname sga@target@IV.4.5.13\endcsname{ega4.c04.prop.4-5-13}
\expandafter\def\csname sga@target@IV.4.5.13.1\endcsname{ega4.c04.cor.4-5-13-1}
\expandafter\def\csname sga@target@IV.4.5.13.1.i\endcsname{ega4.c04.sec.4-5-13-1-i}
\expandafter\def\csname sga@target@IV.4.5.13.1.ii\endcsname{ega4.c04.sec.4-5-13-1-ii}
\expandafter\def\csname sga@target@IV.4.5.14\endcsname{ega4.c04.cor.4-5-14}
\expandafter\def\csname sga@target@IV.4.5.15\endcsname{ega4.c04.prop.4-5-15}
\expandafter\def\csname sga@target@IV.4.5.15.i\endcsname{ega4.c04.sec.4-5-15-i}
\expandafter\def\csname sga@target@IV.4.5.15.ii\endcsname{ega4.c04.sec.4-5-15-ii}
\expandafter\def\csname sga@target@IV.4.5.16\endcsname{ega4.c04.cor.4-5-16}
\expandafter\def\csname sga@target@IV.4.5.16.i\endcsname{ega4.c04.sec.4-5-16-i}
\expandafter\def\csname sga@target@IV.4.5.16.ii\endcsname{ega4.c04.sec.4-5-16-ii}
\expandafter\def\csname sga@target@IV.4.5.17\endcsname{ega4.c04.cor.4-5-17}
\expandafter\def\csname sga@target@IV.4.5.18\endcsname{ega4.c04.rem.4-5-18}
\expandafter\def\csname sga@target@IV.4.5.19\endcsname{ega4.c04.prop.4-5-19}
\expandafter\def\csname sga@target@IV.4.5.19.1\endcsname{ega4.c04.lem.4-5-19-1}
\expandafter\def\csname sga@target@IV.4.5.19.2\endcsname{ega4.c04.eq.4-5-19-2}
\expandafter\def\csname sga@target@IV.4.5.19.3\endcsname{ega4.c04.cor.4-5-19-3}
\expandafter\def\csname sga@target@IV.4.5.2\endcsname{ega4.c04.def.4-5-2}
\expandafter\def\csname sga@target@IV.4.5.2.a\endcsname{ega4.c04.sec.4-5-2-a}
\expandafter\def\csname sga@target@IV.4.5.2.b\endcsname{ega4.c04.sec.4-5-2-b}
\expandafter\def\csname sga@target@IV.4.5.2.c\endcsname{ega4.c04.sec.4-5-2-c}
\expandafter\def\csname sga@target@IV.4.5.20\endcsname{ega4.c04.rem.4-5-20}
\expandafter\def\csname sga@target@IV.4.5.21\endcsname{ega4.c04.prop.4-5-21}
\expandafter\def\csname sga@target@IV.4.5.3\endcsname{ega4.c04.statement.4-5-3}
\expandafter\def\csname sga@target@IV.4.5.4\endcsname{ega4.c04.prop.4-5-4}
\expandafter\def\csname sga@target@IV.4.5.5\endcsname{ega4.c04.def.4-5-5}
\expandafter\def\csname sga@target@IV.4.5.6\endcsname{ega4.c04.prop.4-5-6}
\expandafter\def\csname sga@target@IV.4.5.6.i\endcsname{ega4.c04.sec.4-5-6-i}
\expandafter\def\csname sga@target@IV.4.5.6.ii\endcsname{ega4.c04.sec.4-5-6-ii}
\expandafter\def\csname sga@target@IV.4.5.7\endcsname{ega4.c04.prop.4-5-7}
\expandafter\def\csname sga@target@IV.4.5.8\endcsname{ega4.c04.cor.4-5-8}
\expandafter\def\csname sga@target@IV.4.5.8.i\endcsname{ega4.c04.sec.4-5-8-i}
\expandafter\def\csname sga@target@IV.4.5.8.ii\endcsname{ega4.c04.sec.4-5-8-ii}
\expandafter\def\csname sga@target@IV.4.5.9\endcsname{ega4.c04.prop.4-5-9}
\expandafter\def\csname sga@target@IV.4.5.9.a\endcsname{ega4.c04.sec.4-5-9-a}
\expandafter\def\csname sga@target@IV.4.5.9.b\endcsname{ega4.c04.sec.4-5-9-b}
\expandafter\def\csname sga@target@IV.4.5.9.c\endcsname{ega4.c04.sec.4-5-9-c}
\expandafter\def\csname sga@target@IV.4.6\endcsname{ega4.c04.subsec.4-6.key-iv-4-6}
\expandafter\def\csname sga@target@IV.4.6.1\endcsname{ega4.c04.prop.4-6-1}
\expandafter\def\csname sga@target@IV.4.6.1.a\endcsname{ega4.c04.sec.4-6-1-a}
\expandafter\def\csname sga@target@IV.4.6.1.b\endcsname{ega4.c04.sec.4-6-1-b}
\expandafter\def\csname sga@target@IV.4.6.1.c\endcsname{ega4.c04.sec.4-6-1-c}
\expandafter\def\csname sga@target@IV.4.6.1.d\endcsname{ega4.c04.sec.4-6-1-d}
\expandafter\def\csname sga@target@IV.4.6.1.e\endcsname{ega4.c04.sec.4-6-1-e}
\expandafter\def\csname sga@target@IV.4.6.10\endcsname{ega4.c04.prop.4-6-10}
\expandafter\def\csname sga@target@IV.4.6.11\endcsname{ega4.c04.cor.4-6-11}
\expandafter\def\csname sga@target@IV.4.6.12\endcsname{ega4.c04.prop.4-6-12}
\expandafter\def\csname sga@target@IV.4.6.12.a\endcsname{ega4.c04.sec.4-6-12-a}
\expandafter\def\csname sga@target@IV.4.6.12.b\endcsname{ega4.c04.sec.4-6-12-b}
\expandafter\def\csname sga@target@IV.4.6.12.c\endcsname{ega4.c04.sec.4-6-12-c}
\expandafter\def\csname sga@target@IV.4.6.12.d\endcsname{ega4.c04.sec.4-6-12-d}
\expandafter\def\csname sga@target@IV.4.6.12.e\endcsname{ega4.c04.sec.4-6-12-e}
\expandafter\def\csname sga@target@IV.4.6.13\endcsname{ega4.c04.cor.4-6-13}
\expandafter\def\csname sga@target@IV.4.6.13.f\endcsname{ega4.c04.sec.4-6-13-f}
\expandafter\def\csname sga@target@IV.4.6.14\endcsname{ega4.c04.statement.4-6-14}
\expandafter\def\csname sga@target@IV.4.6.15\endcsname{ega4.c04.prop.4-6-15}
\expandafter\def\csname sga@target@IV.4.6.15.i\endcsname{ega4.c04.sec.4-6-15-i}
\expandafter\def\csname sga@target@IV.4.6.15.ii\endcsname{ega4.c04.sec.4-6-15-ii}
\expandafter\def\csname sga@target@IV.4.6.16\endcsname{ega4.c04.prop.4-6-16}
\expandafter\def\csname sga@target@IV.4.6.16.a\endcsname{ega4.c04.sec.4-6-16-a}
\expandafter\def\csname sga@target@IV.4.6.16.b\endcsname{ega4.c04.sec.4-6-16-b}
\expandafter\def\csname sga@target@IV.4.6.16.c\endcsname{ega4.c04.sec.4-6-16-c}
\expandafter\def\csname sga@target@IV.4.6.16.d\endcsname{ega4.c04.sec.4-6-16-d}
\expandafter\def\csname sga@target@IV.4.6.17\endcsname{ega4.c04.def.4-6-17}
\expandafter\def\csname sga@target@IV.4.6.18\endcsname{ega4.c04.prop.4-6-18}
\expandafter\def\csname sga@target@IV.4.6.18.a\endcsname{ega4.c04.sec.4-6-18-a}
\expandafter\def\csname sga@target@IV.4.6.18.b\endcsname{ega4.c04.sec.4-6-18-b}
\expandafter\def\csname sga@target@IV.4.6.18.c\endcsname{ega4.c04.sec.4-6-18-c}
\expandafter\def\csname sga@target@IV.4.6.18.d\endcsname{ega4.c04.sec.4-6-18-d}
\expandafter\def\csname sga@target@IV.4.6.19\endcsname{ega4.c04.def.4-6-19}
\expandafter\def\csname sga@target@IV.4.6.2\endcsname{ega4.c04.def.4-6-2}
\expandafter\def\csname sga@target@IV.4.6.20\endcsname{ega4.c04.prop.4-6-20}
\expandafter\def\csname sga@target@IV.4.6.21\endcsname{ega4.c04.prop.4-6-21}
\expandafter\def\csname sga@target@IV.4.6.21.i\endcsname{ega4.c04.sec.4-6-21-i}
\expandafter\def\csname sga@target@IV.4.6.21.ii\endcsname{ega4.c04.sec.4-6-21-ii}
\expandafter\def\csname sga@target@IV.4.6.22\endcsname{ega4.c04.def.4-6-22}
\expandafter\def\csname sga@target@IV.4.6.3\endcsname{ega4.c04.cor.4-6-3}
\expandafter\def\csname sga@target@IV.4.6.4\endcsname{ega4.c04.cor.4-6-4}
\expandafter\def\csname sga@target@IV.4.6.5\endcsname{ega4.c04.prop.4-6-5}
\expandafter\def\csname sga@target@IV.4.6.5.i\endcsname{ega4.c04.sec.4-6-5-i}
\expandafter\def\csname sga@target@IV.4.6.5.ii\endcsname{ega4.c04.sec.4-6-5-ii}
\expandafter\def\csname sga@target@IV.4.6.6\endcsname{ega4.c04.prop.4-6-6}
\expandafter\def\csname sga@target@IV.4.6.7\endcsname{ega4.c04.cor.4-6-7}
\expandafter\def\csname sga@target@IV.4.6.8\endcsname{ega4.c04.cor.4-6-8}
\expandafter\def\csname sga@target@IV.4.6.9\endcsname{ega4.c04.def.4-6-9}
\expandafter\def\csname sga@target@IV.4.7.1\endcsname{ega4.c04.lem.4-7-1}
\expandafter\def\csname sga@target@IV.4.7.1.1\endcsname{ega4.c04.eq.4-7-1-1}
\expandafter\def\csname sga@target@IV.4.7.10\endcsname{ega4.c04.prop.4-7-10}
\expandafter\def\csname sga@target@IV.4.7.11\endcsname{ega4.c04.prop.4-7-11}
\expandafter\def\csname sga@target@IV.4.7.12\endcsname{ega4.c04.def.4-7-12}
\expandafter\def\csname sga@target@IV.4.7.13\endcsname{ega4.c04.prop.4-7-13}
\expandafter\def\csname sga@target@IV.4.7.2\endcsname{ega4.c04.cor.4-7-2}
\expandafter\def\csname sga@target@IV.4.7.2.1\endcsname{ega4.c04.eq.4-7-2-1}
\expandafter\def\csname sga@target@IV.4.7.3\endcsname{ega4.c04.prop.4-7-3}
\expandafter\def\csname sga@target@IV.4.7.3.1\endcsname{ega4.c04.eq.4-7-3-1}
\expandafter\def\csname sga@target@IV.4.7.3.i\endcsname{ega4.c04.sec.4-7-3-i}
\expandafter\def\csname sga@target@IV.4.7.3.ii\endcsname{ega4.c04.sec.4-7-3-ii}
\expandafter\def\csname sga@target@IV.4.7.3.iii\endcsname{ega4.c04.sec.4-7-3-iii}
\expandafter\def\csname sga@target@IV.4.7.4\endcsname{ega4.c04.def.4-7-4}
\expandafter\def\csname sga@target@IV.4.7.5\endcsname{ega4.c04.def.4-7-5}
\expandafter\def\csname sga@target@IV.4.7.6\endcsname{ega4.c04.prop.4-7-6}
\expandafter\def\csname sga@target@IV.4.7.6.1\endcsname{ega4.c04.eq.4-7-6-1}
\expandafter\def\csname sga@target@IV.4.7.7\endcsname{ega4.c04.statement.4-7-7}
\expandafter\def\csname sga@target@IV.4.7.8\endcsname{ega4.c04.prop.4-7-8}
\expandafter\def\csname sga@target@IV.4.7.9\endcsname{ega4.c04.cor.4-7-9}
\expandafter\def\csname sga@target@IV.4.8.1\endcsname{ega4.c04.statement.4-8-1}
\expandafter\def\csname sga@target@IV.4.8.10\endcsname{ega4.c04.cor.4-8-10}
\expandafter\def\csname sga@target@IV.4.8.11\endcsname{ega4.c04.cor.4-8-11}
\expandafter\def\csname sga@target@IV.4.8.12\endcsname{ega4.c04.cor.4-8-12}
\expandafter\def\csname sga@target@IV.4.8.13\endcsname{ega4.c04.prop.4-8-13}
\expandafter\def\csname sga@target@IV.4.8.14\endcsname{ega4.c04.prop.4-8-14}
\expandafter\def\csname sga@target@IV.4.8.15\endcsname{ega4.c04.sec.4-8-15}
\expandafter\def\csname sga@target@IV.4.8.15.i\endcsname{ega4.c04.sec.4-8-15-i}
\expandafter\def\csname sga@target@IV.4.8.15.ii\endcsname{ega4.c04.sec.4-8-15-ii}
\expandafter\def\csname sga@target@IV.4.8.2\endcsname{ega4.c04.statement.4-8-2}
\expandafter\def\csname sga@target@IV.4.8.2.1\endcsname{ega4.c04.eq.4-8-2-1}
\expandafter\def\csname sga@target@IV.4.8.2.2\endcsname{ega4.c04.eq.4-8-2-2}
\expandafter\def\csname sga@target@IV.4.8.2.3\endcsname{ega4.c04.eq.4-8-2-3}
\expandafter\def\csname sga@target@IV.4.8.2.4\endcsname{ega4.c04.eq.4-8-2-4}
\expandafter\def\csname sga@target@IV.4.8.2.5\endcsname{ega4.c04.eq.4-8-2-5}
\expandafter\def\csname sga@target@IV.4.8.2.6\endcsname{ega4.c04.eq.4-8-2-6}
\expandafter\def\csname sga@target@IV.4.8.3\endcsname{ega4.c04.prop.4-8-3}
\expandafter\def\csname sga@target@IV.4.8.4\endcsname{ega4.c04.def.4-8-4}
\expandafter\def\csname sga@target@IV.4.8.5\endcsname{ega4.c04.prop.4-8-5}
\expandafter\def\csname sga@target@IV.4.8.6\endcsname{ega4.c04.lem.4-8-6}
\expandafter\def\csname sga@target@IV.4.8.6.a\endcsname{ega4.c04.sec.4-8-6-a}
\expandafter\def\csname sga@target@IV.4.8.6.b\endcsname{ega4.c04.sec.4-8-6-b}
\expandafter\def\csname sga@target@IV.4.8.6.c\endcsname{ega4.c04.sec.4-8-6-c}
\expandafter\def\csname sga@target@IV.4.8.7\endcsname{ega4.c04.cor.4-8-7}
\expandafter\def\csname sga@target@IV.4.8.8\endcsname{ega4.c04.lem.4-8-8}
\expandafter\def\csname sga@target@IV.4.8.8.i\endcsname{ega4.c04.sec.4-8-8-i}
\expandafter\def\csname sga@target@IV.4.8.8.ii\endcsname{ega4.c04.sec.4-8-8-ii}
\expandafter\def\csname sga@target@IV.4.8.8.iii\endcsname{ega4.c04.sec.4-8-8-iii}
\expandafter\def\csname sga@target@IV.4.8.9\endcsname{ega4.c04.prop.4-8-9}
\expandafter\def\csname sga@target@IV.4.9.1\endcsname{ega4.c04.prop.4-9-1}
\expandafter\def\csname sga@target@IV.4.9.1.a\endcsname{ega4.c04.sec.4-9-1-a}
\expandafter\def\csname sga@target@IV.4.9.1.b\endcsname{ega4.c04.sec.4-9-1-b}
\expandafter\def\csname sga@target@IV.4.9.1.bprime\endcsname{ega4.c04.sec.4-9-1-bprime}
\expandafter\def\csname sga@target@IV.4.9.1.c\endcsname{ega4.c04.sec.4-9-1-c}
\expandafter\def\csname sga@target@IV.4.9.2\endcsname{ega4.c04.cor.4-9-2}
\expandafter\def\csname sga@target@IV.4.9.3\endcsname{ega4.c04.rem.4-9-3}
\expandafter\def\csname sga@target@IV.4.9.4\endcsname{ega4.c04.prop.4-9-4}
\expandafter\def\csname sga@target@IV.4.9.5\endcsname{ega4.c04.cor.4-9-5}
\expandafter\def\csname sga@target@IV.4.9.5.a\endcsname{ega4.c04.sec.4-9-5-a}
\expandafter\def\csname sga@target@IV.4.9.5.b\endcsname{ega4.c04.sec.4-9-5-b}
\expandafter\def\csname sga@target@IV.4.9.6\endcsname{ega4.c04.prop.4-9-6}
\expandafter\def\csname sga@target@IV.4.9.6.a\endcsname{ega4.c04.sec.4-9-6-a}
\expandafter\def\csname sga@target@IV.4.9.6.aprime\endcsname{ega4.c04.sec.4-9-6-aprime}
\expandafter\def\csname sga@target@IV.4.9.6.b\endcsname{ega4.c04.sec.4-9-6-b}
\expandafter\def\csname sga@target@IV.4.9.6.bprime\endcsname{ega4.c04.sec.4-9-6-bprime}
\expandafter\def\csname sga@target@IV.4.9.7\endcsname{ega4.c04.cor.4-9-7}
\expandafter\def\csname sga@target@IV.5.1.1\endcsname{ega4.c04.statement.5-1-1}
\expandafter\def\csname sga@target@IV.5.1.1.1\endcsname{ega4.c04.eq.5-1-1-1}
\expandafter\def\csname sga@target@IV.5.1.1.2\endcsname{ega4.c04.eq.5-1-1-2}
\expandafter\def\csname sga@target@IV.5.1.1.3\endcsname{ega4.c04.eq.5-1-1-3}
\expandafter\def\csname sga@target@IV.5.1.1.4\endcsname{ega4.c04.sec.5-1-1-4}
\expandafter\def\csname sga@target@IV.5.1.1.5\endcsname{ega4.c04.sec.5-1-1-5}
\expandafter\def\csname sga@target@IV.5.1.1.6\endcsname{ega4.c04.sec.5-1-1-6}
\expandafter\def\csname sga@target@IV.5.1.10\endcsname{ega4.c04.prop.5-1-10}
\expandafter\def\csname sga@target@IV.5.1.10.1\endcsname{ega4.c04.lem.5-1-10-1}
\expandafter\def\csname sga@target@IV.5.1.10.i\endcsname{ega4.c04.sec.5-1-10-i}
\expandafter\def\csname sga@target@IV.5.1.10.ii\endcsname{ega4.c04.sec.5-1-10-ii}
\expandafter\def\csname sga@target@IV.5.1.11\endcsname{ega4.c04.cor.5-1-11}
\expandafter\def\csname sga@target@IV.5.1.12\endcsname{ega4.c04.statement.5-1-12}
\expandafter\def\csname sga@target@IV.5.1.12.1\endcsname{ega4.c04.eq.5-1-12-1}
\expandafter\def\csname sga@target@IV.5.1.12.2\endcsname{ega4.c04.eq.5-1-12-2}
\expandafter\def\csname sga@target@IV.5.1.12.3\endcsname{ega4.c04.eq.5-1-12-3}
\expandafter\def\csname sga@target@IV.5.1.13\endcsname{ega4.c04.prop.5-1-13}
\expandafter\def\csname sga@target@IV.5.1.13.1\endcsname{ega4.c04.eq.5-1-13-1}
\expandafter\def\csname sga@target@IV.5.1.2\endcsname{ega4.c04.prop.5-1-2}
\expandafter\def\csname sga@target@IV.5.1.2.1\endcsname{ega4.c04.eq.5-1-2-1}
\expandafter\def\csname sga@target@IV.5.1.3\endcsname{ega4.c04.cor.5-1-3}
\expandafter\def\csname sga@target@IV.5.1.3.1\endcsname{ega4.c04.eq.5-1-3-1}
\expandafter\def\csname sga@target@IV.5.1.3.2\endcsname{ega4.c04.eq.5-1-3-2}
\expandafter\def\csname sga@target@IV.5.1.4\endcsname{ega4.c04.prop.5-1-4}
\expandafter\def\csname sga@target@IV.5.1.4.1\endcsname{ega4.c04.eq.5-1-4-1}
\expandafter\def\csname sga@target@IV.5.1.4.2\endcsname{ega4.c04.eq.5-1-4-2}
\expandafter\def\csname sga@target@IV.5.1.5\endcsname{ega4.c04.cor.5-1-5}
\expandafter\def\csname sga@target@IV.5.1.6\endcsname{ega4.c04.statement.5-1-6}
\expandafter\def\csname sga@target@IV.5.1.7\endcsname{ega4.c04.prop.5-1-7}
\expandafter\def\csname sga@target@IV.5.1.8\endcsname{ega4.c04.cor.5-1-8}
\expandafter\def\csname sga@target@IV.5.1.9\endcsname{ega4.c04.prop.5-1-9}
\expandafter\def\csname sga@target@IV.5.1.9.1\endcsname{ega4.c04.statement.5-1-9-1}
\expandafter\def\csname sga@target@IV.5.10.1\endcsname{ega4.c04.statement.5-10-1}
\expandafter\def\csname sga@target@IV.5.10.1.1\endcsname{ega4.c04.eq.5-10-1-1}
\expandafter\def\csname sga@target@IV.5.10.10\endcsname{ega4.c04.prop.5-10-10}
\expandafter\def\csname sga@target@IV.5.10.11\endcsname{ega4.c04.statement.5-10-11}
\expandafter\def\csname sga@target@IV.5.10.11.1\endcsname{ega4.c04.eq.5-10-11-1}
\expandafter\def\csname sga@target@IV.5.10.12\endcsname{ega4.c04.prop.5-10-12}
\expandafter\def\csname sga@target@IV.5.10.13\endcsname{ega4.c04.statement.5-10-13}
\expandafter\def\csname sga@target@IV.5.10.14\endcsname{ega4.c04.prop.5-10-14}
\expandafter\def\csname sga@target@IV.5.10.15\endcsname{ega4.c04.cor.5-10-15}
\expandafter\def\csname sga@target@IV.5.10.16\endcsname{ega4.c04.prop.5-10-16}
\expandafter\def\csname sga@target@IV.5.10.17\endcsname{ega4.c04.prop.5-10-17}
\expandafter\def\csname sga@target@IV.5.10.2\endcsname{ega4.c04.prop.5-10-2}
\expandafter\def\csname sga@target@IV.5.10.3\endcsname{ega4.c04.cor.5-10-3}
\expandafter\def\csname sga@target@IV.5.10.4\endcsname{ega4.c04.cor.5-10-4}
\expandafter\def\csname sga@target@IV.5.10.5\endcsname{ega4.c04.thm.5-10-5}
\expandafter\def\csname sga@target@IV.5.10.6\endcsname{ega4.c04.cor.5-10-6}
\expandafter\def\csname sga@target@IV.5.10.7\endcsname{ega4.c04.thm.5-10-7}
\expandafter\def\csname sga@target@IV.5.10.8\endcsname{ega4.c04.cor.5-10-8}
\expandafter\def\csname sga@target@IV.5.10.8.1\endcsname{ega4.c04.lem.5-10-8-1}
\expandafter\def\csname sga@target@IV.5.10.9\endcsname{ega4.c04.cor.5-10-9}
\expandafter\def\csname sga@target@IV.5.11.1\endcsname{ega4.c04.prop.5-11-1}
\expandafter\def\csname sga@target@IV.5.11.1.1\endcsname{ega4.c04.lem.5-11-1-1}
\expandafter\def\csname sga@target@IV.5.11.2\endcsname{ega4.c04.cor.5-11-2}
\expandafter\def\csname sga@target@IV.5.11.3\endcsname{ega4.c04.rem.5-11-3}
\expandafter\def\csname sga@target@IV.5.11.4\endcsname{ega4.c04.cor.5-11-4}
\expandafter\def\csname sga@target@IV.5.11.5\endcsname{ega4.c04.cor.5-11-5}
\expandafter\def\csname sga@target@IV.5.11.6\endcsname{ega4.c04.cor.5-11-6}
\expandafter\def\csname sga@target@IV.5.11.7\endcsname{ega4.c04.item.5-11-7}
\expandafter\def\csname sga@target@IV.5.12\endcsname{ega4.c04.subsec.5-12.key-iv-5-12}
\expandafter\def\csname sga@target@IV.5.12.1\endcsname{ega4.c04.prop.5-12-1}
\expandafter\def\csname sga@target@IV.5.12.1.1\endcsname{ega4.c04.eq.5-12-1-1}
\expandafter\def\csname sga@target@IV.5.12.10\endcsname{ega4.c04.cor.5-12-10}
\expandafter\def\csname sga@target@IV.5.12.2\endcsname{ega4.c04.prop.5-12-2}
\expandafter\def\csname sga@target@IV.5.12.3\endcsname{ega4.c04.rem.5-12-3}
\expandafter\def\csname sga@target@IV.5.12.4\endcsname{ega4.c04.cor.5-12-4}
\expandafter\def\csname sga@target@IV.5.12.5\endcsname{ega4.c04.prop.5-12-5}
\expandafter\def\csname sga@target@IV.5.12.6\endcsname{ega4.c04.rem.5-12-6}
\expandafter\def\csname sga@target@IV.5.12.7\endcsname{ega4.c04.cor.5-12-7}
\expandafter\def\csname sga@target@IV.5.12.8\endcsname{ega4.c04.prop.5-12-8}
\expandafter\def\csname sga@target@IV.5.12.9\endcsname{ega4.c04.rem.5-12-9}
\expandafter\def\csname sga@target@IV.5.13.1\endcsname{ega4.c04.prop.5-13-1}
\expandafter\def\csname sga@target@IV.5.13.2\endcsname{ega4.c04.cor.5-13-2}
\expandafter\def\csname sga@target@IV.5.13.3\endcsname{ega4.c04.prop.5-13-3}
\expandafter\def\csname sga@target@IV.5.13.4\endcsname{ega4.c04.cor.5-13-4}
\expandafter\def\csname sga@target@IV.5.13.5\endcsname{ega4.c04.statement.5-13-5}
\expandafter\def\csname sga@target@IV.5.13.6\endcsname{ega4.c04.cor.5-13-6}
\expandafter\def\csname sga@target@IV.5.13.7\endcsname{ega4.c04.prop.5-13-7}
\expandafter\def\csname sga@target@IV.5.13.7.1\endcsname{ega4.c04.lem.5-13-7-1}
\expandafter\def\csname sga@target@IV.5.13.8\endcsname{ega4.c04.cor.5-13-8}
\expandafter\def\csname sga@target@IV.5.2.1\endcsname{ega4.c04.prop.5-2-1}
\expandafter\def\csname sga@target@IV.5.2.1.1\endcsname{ega4.c04.eq.5-2-1-1}
\expandafter\def\csname sga@target@IV.5.2.2\endcsname{ega4.c04.cor.5-2-2}
\expandafter\def\csname sga@target@IV.5.2.3\endcsname{ega4.c04.cor.5-2-3}
\expandafter\def\csname sga@target@IV.5.2.3.1\endcsname{ega4.c04.eq.5-2-3-1}
\expandafter\def\csname sga@target@IV.5.2.4\endcsname{ega4.c04.cor.5-2-4}
\expandafter\def\csname sga@target@IV.5.2.5\endcsname{ega4.c04.sec.5-2-5}
\expandafter\def\csname sga@target@IV.5.2.5.i\endcsname{ega4.c04.sec.5-2-5-i}
\expandafter\def\csname sga@target@IV.5.2.5.ii\endcsname{ega4.c04.sec.5-2-5-ii}
\expandafter\def\csname sga@target@IV.5.3.1\endcsname{ega4.c04.prop.5-3-1}
\expandafter\def\csname sga@target@IV.5.3.1.1\endcsname{ega4.c04.eq.5-3-1-1}
\expandafter\def\csname sga@target@IV.5.3.1.2\endcsname{ega4.c04.eq.5-3-1-2}
\expandafter\def\csname sga@target@IV.5.3.1.3\endcsname{ega4.c04.lem.5-3-1-3}
\expandafter\def\csname sga@target@IV.5.4.1\endcsname{ega4.c04.prop.5-4-1}
\expandafter\def\csname sga@target@IV.5.4.1.i\endcsname{ega4.c04.sec.5-4-1-i}
\expandafter\def\csname sga@target@IV.5.4.1.ii\endcsname{ega4.c04.sec.5-4-1-ii}
\expandafter\def\csname sga@target@IV.5.4.2\endcsname{ega4.c04.cor.5-4-2}
\expandafter\def\csname sga@target@IV.5.4.3\endcsname{ega4.c04.sec.5-4-3}
\expandafter\def\csname sga@target@IV.5.4.3.1\endcsname{ega4.c04.eq.5-4-3-1}
\expandafter\def\csname sga@target@IV.5.4.3.i\endcsname{ega4.c04.sec.5-4-3-i}
\expandafter\def\csname sga@target@IV.5.4.3.ii\endcsname{ega4.c04.sec.5-4-3-ii}
\expandafter\def\csname sga@target@IV.5.5\endcsname{ega4.c04.sec.5-5}
\expandafter\def\csname sga@target@IV.5.5.1\endcsname{ega4.c04.statement.5-5-1}
\expandafter\def\csname sga@target@IV.5.5.1.1\endcsname{ega4.c04.eq.5-5-1-1}
\expandafter\def\csname sga@target@IV.5.5.2\endcsname{ega4.c04.prop.5-5-2}
\expandafter\def\csname sga@target@IV.5.5.2.1\endcsname{ega4.c04.eq.5-5-2-1}
\expandafter\def\csname sga@target@IV.5.5.2.2\endcsname{ega4.c04.eq.5-5-2-2}
\expandafter\def\csname sga@target@IV.5.5.3\endcsname{ega4.c04.prop.5-5-3}
\expandafter\def\csname sga@target@IV.5.5.3.1\endcsname{ega4.c04.eq.5-5-3-1}
\expandafter\def\csname sga@target@IV.5.5.4\endcsname{ega4.c04.cor.5-5-4}
\expandafter\def\csname sga@target@IV.5.5.4.1\endcsname{ega4.c04.eq.5-5-4-1}
\expandafter\def\csname sga@target@IV.5.5.5\endcsname{ega4.c04.cor.5-5-5}
\expandafter\def\csname sga@target@IV.5.5.6\endcsname{ega4.c04.cor.5-5-6}
\expandafter\def\csname sga@target@IV.5.5.6.1\endcsname{ega4.c04.eq.5-5-6-1}
\expandafter\def\csname sga@target@IV.5.5.7\endcsname{ega4.c04.lem.5-5-7}
\expandafter\def\csname sga@target@IV.5.5.7.1\endcsname{ega4.c04.eq.5-5-7-1}
\expandafter\def\csname sga@target@IV.5.5.7.2\endcsname{ega4.c04.statement.5-5-7-2}
\expandafter\def\csname sga@target@IV.5.5.7.3\endcsname{ega4.c04.eq.5-5-7-3}
\expandafter\def\csname sga@target@IV.5.5.7.i\endcsname{ega4.c04.sec.5-5-7-i}
\expandafter\def\csname sga@target@IV.5.5.7.ii\endcsname{ega4.c04.sec.5-5-7-ii}
\expandafter\def\csname sga@target@IV.5.5.8\endcsname{ega4.c04.thm.5-5-8}
\expandafter\def\csname sga@target@IV.5.5.8.1\endcsname{ega4.c04.eq.5-5-8-1}
\expandafter\def\csname sga@target@IV.5.5.8.2\endcsname{ega4.c04.eq.5-5-8-2}
\expandafter\def\csname sga@target@IV.5.5.8.3\endcsname{ega4.c04.eq.5-5-8-3}
\expandafter\def\csname sga@target@IV.5.6\endcsname{ega4.c04.sec.5-6}
\expandafter\def\csname sga@target@IV.5.6.1\endcsname{ega4.c04.prop.5-6-1}
\expandafter\def\csname sga@target@IV.5.6.1.1\endcsname{ega4.c04.eq.5-6-1-1}
\expandafter\def\csname sga@target@IV.5.6.1.2\endcsname{ega4.c04.eq.5-6-1-2}
\expandafter\def\csname sga@target@IV.5.6.1.3\endcsname{ega4.c04.statement.5-6-1-3}
\expandafter\def\csname sga@target@IV.5.6.1.4\endcsname{ega4.c04.statement.5-6-1-4}
\expandafter\def\csname sga@target@IV.5.6.1.5\endcsname{ega4.c04.statement.5-6-1-5}
\expandafter\def\csname sga@target@IV.5.6.1.a\endcsname{ega4.c04.sec.5-6-1-a}
\expandafter\def\csname sga@target@IV.5.6.1.b\endcsname{ega4.c04.sec.5-6-1-b}
\expandafter\def\csname sga@target@IV.5.6.1.c\endcsname{ega4.c04.sec.5-6-1-c}
\expandafter\def\csname sga@target@IV.5.6.10\endcsname{ega4.c04.prop.5-6-10}
\expandafter\def\csname sga@target@IV.5.6.11\endcsname{ega4.c04.ex.5-6-11}
\expandafter\def\csname sga@target@IV.5.6.11.1\endcsname{ega4.c04.lem.5-6-11-1}
\expandafter\def\csname sga@target@IV.5.6.11.1.i\endcsname{ega4.c04.sec.5-6-11-1-i}
\expandafter\def\csname sga@target@IV.5.6.11.1.ii\endcsname{ega4.c04.sec.5-6-11-1-ii}
\expandafter\def\csname sga@target@IV.5.6.2\endcsname{ega4.c04.def.5-6-2}
\expandafter\def\csname sga@target@IV.5.6.3\endcsname{ega4.c04.sec.5-6-3}
\expandafter\def\csname sga@target@IV.5.6.3.i\endcsname{ega4.c04.sec.5-6-3-i}
\expandafter\def\csname sga@target@IV.5.6.3.ii\endcsname{ega4.c04.sec.5-6-3-ii}
\expandafter\def\csname sga@target@IV.5.6.3.iii\endcsname{ega4.c04.sec.5-6-3-iii}
\expandafter\def\csname sga@target@IV.5.6.3.iv\endcsname{ega4.c04.sec.5-6-3-iv}
\expandafter\def\csname sga@target@IV.5.6.4\endcsname{ega4.c04.prop.5-6-4}
\expandafter\def\csname sga@target@IV.5.6.5\endcsname{ega4.c04.prop.5-6-5}
\expandafter\def\csname sga@target@IV.5.6.5.1\endcsname{ega4.c04.eq.5-6-5-1}
\expandafter\def\csname sga@target@IV.5.6.5.2\endcsname{ega4.c04.eq.5-6-5-2}
\expandafter\def\csname sga@target@IV.5.6.5.3\endcsname{ega4.c04.eq.5-6-5-3}
\expandafter\def\csname sga@target@IV.5.6.6\endcsname{ega4.c04.cor.5-6-6}
\expandafter\def\csname sga@target@IV.5.6.6.1\endcsname{ega4.c04.eq.5-6-6-1}
\expandafter\def\csname sga@target@IV.5.6.6.2\endcsname{ega4.c04.eq.5-6-6-2}
\expandafter\def\csname sga@target@IV.5.6.7\endcsname{ega4.c04.cor.5-6-7}
\expandafter\def\csname sga@target@IV.5.6.7.1\endcsname{ega4.c04.eq.5-6-7-1}
\expandafter\def\csname sga@target@IV.5.6.8\endcsname{ega4.c04.cor.5-6-8}
\expandafter\def\csname sga@target@IV.5.6.8.1\endcsname{ega4.c04.eq.5-6-8-1}
\expandafter\def\csname sga@target@IV.5.6.8.2\endcsname{ega4.c04.eq.5-6-8-2}
\expandafter\def\csname sga@target@IV.5.6.9\endcsname{ega4.c04.sec.5-6-9}
\expandafter\def\csname sga@target@IV.5.6.9.i\endcsname{ega4.c04.sec.5-6-9-i}
\expandafter\def\csname sga@target@IV.5.6.9.ii\endcsname{ega4.c04.sec.5-6-9-ii}
\expandafter\def\csname sga@target@IV.5.7.1\endcsname{ega4.c04.def.5-7-1}
\expandafter\def\csname sga@target@IV.5.7.1.1\endcsname{ega4.c04.eq.5-7-1-1}
\expandafter\def\csname sga@target@IV.5.7.10\endcsname{ega4.c04.cor.5-7-10}
\expandafter\def\csname sga@target@IV.5.7.11\endcsname{ega4.c04.cor.5-7-11}
\expandafter\def\csname sga@target@IV.5.7.11.i\endcsname{ega4.c04.sec.5-7-11-i}
\expandafter\def\csname sga@target@IV.5.7.11.ii\endcsname{ega4.c04.sec.5-7-11-ii}
\expandafter\def\csname sga@target@IV.5.7.12\endcsname{ega4.c04.statement.5-7-12}
\expandafter\def\csname sga@target@IV.5.7.13\endcsname{ega4.c04.cor.5-7-13}
\expandafter\def\csname sga@target@IV.5.7.13.1\endcsname{ega4.c04.eq.5-7-13-1}
\expandafter\def\csname sga@target@IV.5.7.2\endcsname{ega4.c04.def.5-7-2}
\expandafter\def\csname sga@target@IV.5.7.2.1\endcsname{ega4.c04.eq.5-7-2-1}
\expandafter\def\csname sga@target@IV.5.7.2.2\endcsname{ega4.c04.eq.5-7-2-2}
\expandafter\def\csname sga@target@IV.5.7.3\endcsname{ega4.c04.sec.5-7-3}
\expandafter\def\csname sga@target@IV.5.7.3.i\endcsname{ega4.c04.sec.5-7-3-i}
\expandafter\def\csname sga@target@IV.5.7.3.ii\endcsname{ega4.c04.sec.5-7-3-ii}
\expandafter\def\csname sga@target@IV.5.7.3.iii\endcsname{ega4.c04.sec.5-7-3-iii}
\expandafter\def\csname sga@target@IV.5.7.3.iv\endcsname{ega4.c04.sec.5-7-3-iv}
\expandafter\def\csname sga@target@IV.5.7.3.v\endcsname{ega4.c04.sec.5-7-3-v}
\expandafter\def\csname sga@target@IV.5.7.3.vi\endcsname{ega4.c04.sec.5-7-3-vi}
\expandafter\def\csname sga@target@IV.5.7.4\endcsname{ega4.c04.prop.5-7-4}
\expandafter\def\csname sga@target@IV.5.7.4.1\endcsname{ega4.c04.eq.5-7-4-1}
\expandafter\def\csname sga@target@IV.5.7.4.2\endcsname{ega4.c04.eq.5-7-4-2}
\expandafter\def\csname sga@target@IV.5.7.4.i\endcsname{ega4.c04.sec.5-7-4-i}
\expandafter\def\csname sga@target@IV.5.7.4.ii\endcsname{ega4.c04.sec.5-7-4-ii}
\expandafter\def\csname sga@target@IV.5.7.5\endcsname{ega4.c04.prop.5-7-5}
\expandafter\def\csname sga@target@IV.5.7.6\endcsname{ega4.c04.cor.5-7-6}
\expandafter\def\csname sga@target@IV.5.7.7\endcsname{ega4.c04.cor.5-7-7}
\expandafter\def\csname sga@target@IV.5.7.8\endcsname{ega4.c04.sec.5-7-8}
\expandafter\def\csname sga@target@IV.5.7.9\endcsname{ega4.c04.prop.5-7-9}
\expandafter\def\csname sga@target@IV.5.7.9.1\endcsname{ega4.c04.eq.5-7-9-1}
\expandafter\def\csname sga@target@IV.5.7.9.2\endcsname{ega4.c04.eq.5-7-9-2}
\expandafter\def\csname sga@target@IV.5.8.1\endcsname{ega4.c04.statement.5-8-1}
\expandafter\def\csname sga@target@IV.5.8.2\endcsname{ega4.c04.def.5-8-2}
\expandafter\def\csname sga@target@IV.5.8.3\endcsname{ega4.c04.prop.5-8-3}
\expandafter\def\csname sga@target@IV.5.8.4\endcsname{ega4.c04.statement.5-8-4}
\expandafter\def\csname sga@target@IV.5.8.5\endcsname{ega4.c04.prop.5-8-5}
\expandafter\def\csname sga@target@IV.5.8.6\endcsname{ega4.c04.thm.5-8-6}
\expandafter\def\csname sga@target@IV.5.8.6.1\endcsname{ega4.c04.eq.5-8-6-1}
\expandafter\def\csname sga@target@IV.5.8.6.i\endcsname{ega4.c04.item.5-8-6-i}
\expandafter\def\csname sga@target@IV.5.8.6.ii\endcsname{ega4.c04.item.5-8-6-ii}
\expandafter\def\csname sga@target@IV.5.9.1\endcsname{ega4.c04.statement.5-9-1}
\expandafter\def\csname sga@target@IV.5.9.1.1\endcsname{ega4.c04.eq.5-9-1-1}
\expandafter\def\csname sga@target@IV.5.9.1.2\endcsname{ega4.c04.eq.5-9-1-2}
\expandafter\def\csname sga@target@IV.5.9.1.3\endcsname{ega4.c04.eq.5-9-1-3}
\expandafter\def\csname sga@target@IV.5.9.10\endcsname{ega4.c04.cor.5-9-10}
\expandafter\def\csname sga@target@IV.5.9.10.i\endcsname{ega4.c04.item.5-9-10-i}
\expandafter\def\csname sga@target@IV.5.9.10.ii\endcsname{ega4.c04.item.5-9-10-ii}
\expandafter\def\csname sga@target@IV.5.9.11\endcsname{ega4.c04.prop.5-9-11}
\expandafter\def\csname sga@target@IV.5.9.12\endcsname{ega4.c04.item.5-9-12}
\expandafter\def\csname sga@target@IV.5.9.12.i\endcsname{ega4.c04.item.5-9-12-i}
\expandafter\def\csname sga@target@IV.5.9.12.ii\endcsname{ega4.c04.item.5-9-12-ii}
\expandafter\def\csname sga@target@IV.5.9.12.iii\endcsname{ega4.c04.item.5-9-12-iii}
\expandafter\def\csname sga@target@IV.5.9.2\endcsname{ega4.c04.prop.5-9-2}
\expandafter\def\csname sga@target@IV.5.9.2.i\endcsname{ega4.c04.item.5-9-2-i}
\expandafter\def\csname sga@target@IV.5.9.2.ii\endcsname{ega4.c04.item.5-9-2-ii}
\expandafter\def\csname sga@target@IV.5.9.3\endcsname{ega4.c04.rem.5-9-3}
\expandafter\def\csname sga@target@IV.5.9.3.1\endcsname{ega4.c04.eq.5-9-3-1}
\expandafter\def\csname sga@target@IV.5.9.4\endcsname{ega4.c04.prop.5-9-4}
\expandafter\def\csname sga@target@IV.5.9.4.1\endcsname{ega4.c04.eq.5-9-4-1}
\expandafter\def\csname sga@target@IV.5.9.5\endcsname{ega4.c04.cor.5-9-5}
\expandafter\def\csname sga@target@IV.5.9.6\endcsname{ega4.c04.cor.5-9-6}
\expandafter\def\csname sga@target@IV.5.9.6.1\endcsname{ega4.c04.eq.5-9-6-1}
\expandafter\def\csname sga@target@IV.5.9.7\endcsname{ega4.c04.statement.5-9-7}
\expandafter\def\csname sga@target@IV.5.9.7.1\endcsname{ega4.c04.eq.5-9-7-1}
\expandafter\def\csname sga@target@IV.5.9.8\endcsname{ega4.c04.prop.5-9-8}
\expandafter\def\csname sga@target@IV.5.9.8.1\endcsname{ega4.c04.eq.5-9-8-1}
\expandafter\def\csname sga@target@IV.5.9.8.a\endcsname{ega4.c04.item.5-9-8-a}
\expandafter\def\csname sga@target@IV.5.9.8.a-prime\endcsname{ega4.c04.item.5-9-8-a-prime}
\expandafter\def\csname sga@target@IV.5.9.8.b\endcsname{ega4.c04.item.5-9-8-b}
\expandafter\def\csname sga@target@IV.5.9.8.b-prime\endcsname{ega4.c04.item.5-9-8-b-prime}
\expandafter\def\csname sga@target@IV.5.9.9\endcsname{ega4.c04.def.5-9-9}
\expandafter\def\csname sga@target@IV.6.1.1\endcsname{ega4.c04.prop.6-1-1}
\expandafter\def\csname sga@target@IV.6.1.1.1\endcsname{ega4.c04.eq.6-1-1-1}
\expandafter\def\csname sga@target@IV.6.1.2\endcsname{ega4.c04.cor.6-1-2}
\expandafter\def\csname sga@target@IV.6.1.2.1\endcsname{ega4.c04.eq.6-1-2-1}
\expandafter\def\csname sga@target@IV.6.1.3\endcsname{ega4.c04.cor.6-1-3}
\expandafter\def\csname sga@target@IV.6.1.4\endcsname{ega4.c04.cor.6-1-4}
\expandafter\def\csname sga@target@IV.6.1.4.1\endcsname{ega4.c04.eq.6-1-4-1}
\expandafter\def\csname sga@target@IV.6.1.5\endcsname{ega4.c04.prop.6-1-5}
\expandafter\def\csname sga@target@IV.6.10.1\endcsname{ega4.c04.statement.6-10-1}
\expandafter\def\csname sga@target@IV.6.10.2\endcsname{ega4.c04.prop.6-10-2}
\expandafter\def\csname sga@target@IV.6.10.3\endcsname{ega4.c04.prop.6-10-3}
\expandafter\def\csname sga@target@IV.6.10.3.1\endcsname{ega4.c04.eq.6-10-3-1}
\expandafter\def\csname sga@target@IV.6.10.3.2\endcsname{ega4.c04.eq.6-10-3-2}
\expandafter\def\csname sga@target@IV.6.10.3.3\endcsname{ega4.c04.eq.6-10-3-3}
\expandafter\def\csname sga@target@IV.6.10.4\endcsname{ega4.c04.cor.6-10-4}
\expandafter\def\csname sga@target@IV.6.10.5\endcsname{ega4.c04.prop.6-10-5}
\expandafter\def\csname sga@target@IV.6.10.5.1\endcsname{ega4.c04.eq.6-10-5-1}
\expandafter\def\csname sga@target@IV.6.10.5.2\endcsname{ega4.c04.eq.6-10-5-2}
\expandafter\def\csname sga@target@IV.6.10.5.3\endcsname{ega4.c04.eq.6-10-5-3}
\expandafter\def\csname sga@target@IV.6.10.5.4\endcsname{ega4.c04.eq.6-10-5-4}
\expandafter\def\csname sga@target@IV.6.10.6\endcsname{ega4.c04.prop.6-10-6}
\expandafter\def\csname sga@target@IV.6.10.6.1\endcsname{ega4.c04.eq.6-10-6-1}
\expandafter\def\csname sga@target@IV.6.10.6.2\endcsname{ega4.c04.eq.6-10-6-2}
\expandafter\def\csname sga@target@IV.6.10.6.3\endcsname{ega4.c04.eq.6-10-6-3}
\expandafter\def\csname sga@target@IV.6.11\endcsname{ega4.c04.subsec.6-11.key-iv-6-11}
\expandafter\def\csname sga@target@IV.6.11.1\endcsname{ega4.c04.lem.6-11-1}
\expandafter\def\csname sga@target@IV.6.11.2\endcsname{ega4.c04.prop.6-11-2}
\expandafter\def\csname sga@target@IV.6.11.2.1\endcsname{ega4.c04.eq.6-11-2-1}
\expandafter\def\csname sga@target@IV.6.11.2.2\endcsname{ega4.c04.eq.6-11-2-2}
\expandafter\def\csname sga@target@IV.6.11.3\endcsname{ega4.c04.cor.6-11-3}
\expandafter\def\csname sga@target@IV.6.11.4\endcsname{ega4.c04.statement.6-11-4}
\expandafter\def\csname sga@target@IV.6.11.5\endcsname{ega4.c04.prop.6-11-5}
\expandafter\def\csname sga@target@IV.6.11.5.1\endcsname{ega4.c04.eq.6-11-5-1}
\expandafter\def\csname sga@target@IV.6.11.6\endcsname{ega4.c04.prop.6-11-6}
\expandafter\def\csname sga@target@IV.6.11.6.1\endcsname{ega4.c04.lem.6-11-6-1}
\expandafter\def\csname sga@target@IV.6.11.7\endcsname{ega4.c04.cor.6-11-7}
\expandafter\def\csname sga@target@IV.6.11.8\endcsname{ega4.c04.prop.6-11-8}
\expandafter\def\csname sga@target@IV.6.11.8.1\endcsname{ega4.c04.eq.6-11-8-1}
\expandafter\def\csname sga@target@IV.6.11.9\endcsname{ega4.c04.statement.6-11-9}
\expandafter\def\csname sga@target@IV.6.12\endcsname{ega4.c04.subsec.6-12.key-iv-6-12}
\expandafter\def\csname sga@target@IV.6.12.1\endcsname{ega4.c04.statement.6-12-1}
\expandafter\def\csname sga@target@IV.6.12.10\endcsname{ega4.c04.prop.6-12-10}
\expandafter\def\csname sga@target@IV.6.12.2\endcsname{ega4.c04.prop.6-12-2}
\expandafter\def\csname sga@target@IV.6.12.3\endcsname{ega4.c04.cor.6-12-3}
\expandafter\def\csname sga@target@IV.6.12.4\endcsname{ega4.c04.thm.6-12-4}
\expandafter\def\csname sga@target@IV.6.12.4.1\endcsname{ega4.c04.lem.6-12-4-1}
\expandafter\def\csname sga@target@IV.6.12.5\endcsname{ega4.c04.cor.6-12-5}
\expandafter\def\csname sga@target@IV.6.12.6\endcsname{ega4.c04.cor.6-12-6}
\expandafter\def\csname sga@target@IV.6.12.7\endcsname{ega4.c04.thm.6-12-7}
\expandafter\def\csname sga@target@IV.6.12.8\endcsname{ega4.c04.cor.6-12-8}
\expandafter\def\csname sga@target@IV.6.12.9\endcsname{ega4.c04.prop.6-12-9}
\expandafter\def\csname sga@target@IV.6.13\endcsname{ega4.c04.subsec.6-13.key-iv-6-13}
\expandafter\def\csname sga@target@IV.6.13.1\endcsname{ega4.c04.statement.6-13-1}
\expandafter\def\csname sga@target@IV.6.13.2\endcsname{ega4.c04.prop.6-13-2}
\expandafter\def\csname sga@target@IV.6.13.3\endcsname{ega4.c04.cor.6-13-3}
\expandafter\def\csname sga@target@IV.6.13.4\endcsname{ega4.c04.prop.6-13-4}
\expandafter\def\csname sga@target@IV.6.13.5\endcsname{ega4.c04.cor.6-13-5}
\expandafter\def\csname sga@target@IV.6.13.6\endcsname{ega4.c04.prop.6-13-6}
\expandafter\def\csname sga@target@IV.6.13.7\endcsname{ega4.c04.prop.6-13-7}
\expandafter\def\csname sga@target@IV.6.13.7.1\endcsname{ega4.c04.lem.6-13-7-1}
\expandafter\def\csname sga@target@IV.6.14.1\endcsname{ega4.c04.prop.6-14-1}
\expandafter\def\csname sga@target@IV.6.14.1.1\endcsname{ega4.c04.lem.6-14-1-1}
\expandafter\def\csname sga@target@IV.6.14.1.2\endcsname{ega4.c04.lem.6-14-1-2}
\expandafter\def\csname sga@target@IV.6.14.1.2.env\endcsname{ega4.c04.statement.6-14-1-2-env}
\expandafter\def\csname sga@target@IV.6.14.1.3\endcsname{ega4.c04.eq.6-14-1-3}
\expandafter\def\csname sga@target@IV.6.14.1.4\endcsname{ega4.c04.lem.6-14-1-4}
\expandafter\def\csname sga@target@IV.6.14.2\endcsname{ega4.c04.cor.6-14-2}
\expandafter\def\csname sga@target@IV.6.14.3\endcsname{ega4.c04.cor.6-14-3}
\expandafter\def\csname sga@target@IV.6.14.4\endcsname{ega4.c04.prop.6-14-4}
\expandafter\def\csname sga@target@IV.6.14.4.1\endcsname{ega4.c04.lem.6-14-4-1}
\expandafter\def\csname sga@target@IV.6.14.5\endcsname{ega4.c04.cor.6-14-5}
\expandafter\def\csname sga@target@IV.6.15.1\endcsname{ega4.c04.statement.6-15-1}
\expandafter\def\csname sga@target@IV.6.15.10\endcsname{ega4.c04.prop.6-15-10}
\expandafter\def\csname sga@target@IV.6.15.11\endcsname{ega4.c04.item.6-15-11}
\expandafter\def\csname sga@target@IV.6.15.2\endcsname{ega4.c04.statement.6-15-2}
\expandafter\def\csname sga@target@IV.6.15.3\endcsname{ega4.c04.statement.6-15-3}
\expandafter\def\csname sga@target@IV.6.15.3.1\endcsname{ega4.c04.lem.6-15-3-1}
\expandafter\def\csname sga@target@IV.6.15.4\endcsname{ega4.c04.statement.6-15-4}
\expandafter\def\csname sga@target@IV.6.15.4.1\endcsname{ega4.c04.lem.6-15-4-1}
\expandafter\def\csname sga@target@IV.6.15.5\endcsname{ega4.c04.prop.6-15-5}
\expandafter\def\csname sga@target@IV.6.15.6\endcsname{ega4.c04.prop.6-15-6}
\expandafter\def\csname sga@target@IV.6.15.7\endcsname{ega4.c04.prop.6-15-7}
\expandafter\def\csname sga@target@IV.6.15.8\endcsname{ega4.c04.lem.6-15-8}
\expandafter\def\csname sga@target@IV.6.15.9\endcsname{ega4.c04.cor.6-15-9}
\expandafter\def\csname sga@target@IV.6.2.1\endcsname{ega4.c04.prop.6-2-1}
\expandafter\def\csname sga@target@IV.6.2.1.1\endcsname{ega4.c04.eq.6-2-1-1}
\expandafter\def\csname sga@target@IV.6.2.1.2\endcsname{ega4.c04.eq.6-2-1-2}
\expandafter\def\csname sga@target@IV.6.2.2\endcsname{ega4.c04.cor.6-2-2}
\expandafter\def\csname sga@target@IV.6.2.3\endcsname{ega4.c04.prop.6-2-3}
\expandafter\def\csname sga@target@IV.6.2.3.1\endcsname{ega4.c04.eq.6-2-3-1}
\expandafter\def\csname sga@target@IV.6.2.4\endcsname{ega4.c04.prop.6-2-4}
\expandafter\def\csname sga@target@IV.6.2.4.1\endcsname{ega4.c04.eq.6-2-4-1}
\expandafter\def\csname sga@target@IV.6.3.1\endcsname{ega4.c04.prop.6-3-1}
\expandafter\def\csname sga@target@IV.6.3.1.1\endcsname{ega4.c04.eq.6-3-1-1}
\expandafter\def\csname sga@target@IV.6.3.2\endcsname{ega4.c04.cor.6-3-2}
\expandafter\def\csname sga@target@IV.6.3.2.1\endcsname{ega4.c04.eq.6-3-2-1}
\expandafter\def\csname sga@target@IV.6.3.3\endcsname{ega4.c04.cor.6-3-3}
\expandafter\def\csname sga@target@IV.6.3.4\endcsname{ega4.c04.cor.6-3-4}
\expandafter\def\csname sga@target@IV.6.3.4.1\endcsname{ega4.c04.eq.6-3-4-1}
\expandafter\def\csname sga@target@IV.6.3.4.2\endcsname{ega4.c04.eq.6-3-4-2}
\expandafter\def\csname sga@target@IV.6.3.5\endcsname{ega4.c04.cor.6-3-5}
\expandafter\def\csname sga@target@IV.6.3.6\endcsname{ega4.c04.cor.6-3-6}
\expandafter\def\csname sga@target@IV.6.3.7\endcsname{ega4.c04.cor.6-3-7}
\expandafter\def\csname sga@target@IV.6.3.8\endcsname{ega4.c04.prop.6-3-8}
\expandafter\def\csname sga@target@IV.6.4.1\endcsname{ega4.c04.prop.6-4-1}
\expandafter\def\csname sga@target@IV.6.4.2\endcsname{ega4.c04.cor.6-4-2}
\expandafter\def\csname sga@target@IV.6.4.3\endcsname{ega4.c04.rem.6-4-3}
\expandafter\def\csname sga@target@IV.6.5.1\endcsname{ega4.c04.prop.6-5-1}
\expandafter\def\csname sga@target@IV.6.5.2\endcsname{ega4.c04.cor.6-5-2}
\expandafter\def\csname sga@target@IV.6.5.3\endcsname{ega4.c04.prop.6-5-3}
\expandafter\def\csname sga@target@IV.6.5.4\endcsname{ega4.c04.cor.6-5-4}
\expandafter\def\csname sga@target@IV.6.5.5\endcsname{ega4.c04.item.6-5-5}
\expandafter\def\csname sga@target@IV.6.6.1\endcsname{ega4.c04.prop.6-6-1}
\expandafter\def\csname sga@target@IV.6.6.2\endcsname{ega4.c04.rem.6-6-2}
\expandafter\def\csname sga@target@IV.6.7.1\endcsname{ega4.c04.prop.6-7-1}
\expandafter\def\csname sga@target@IV.6.7.1.1\endcsname{ega4.c04.lem.6-7-1-1}
\expandafter\def\csname sga@target@IV.6.7.2\endcsname{ega4.c04.cor.6-7-2}
\expandafter\def\csname sga@target@IV.6.7.3\endcsname{ega4.c04.cor.6-7-3}
\expandafter\def\csname sga@target@IV.6.7.4\endcsname{ega4.c04.prop.6-7-4}
\expandafter\def\csname sga@target@IV.6.7.4.1\endcsname{ega4.c04.lem.6-7-4-1}
\expandafter\def\csname sga@target@IV.6.7.5\endcsname{ega4.c04.item.6-7-5}
\expandafter\def\csname sga@target@IV.6.7.6\endcsname{ega4.c04.def.6-7-6}
\expandafter\def\csname sga@target@IV.6.7.7\endcsname{ega4.c04.prop.6-7-7}
\expandafter\def\csname sga@target@IV.6.7.8\endcsname{ega4.c04.cor.6-7-8}
\expandafter\def\csname sga@target@IV.6.8.1\endcsname{ega4.c04.def.6-8-1}
\expandafter\def\csname sga@target@IV.6.8.2\endcsname{ega4.c04.prop.6-8-2}
\expandafter\def\csname sga@target@IV.6.8.3\endcsname{ega4.c04.prop.6-8-3}
\expandafter\def\csname sga@target@IV.6.8.4\endcsname{ega4.c04.item.6-8-4}
\expandafter\def\csname sga@target@IV.6.8.5\endcsname{ega4.c04.prop.6-8-5}
\expandafter\def\csname sga@target@IV.6.8.6\endcsname{ega4.c04.thm.6-8-6}
\expandafter\def\csname sga@target@IV.6.8.7\endcsname{ega4.c04.cor.6-8-7}
\expandafter\def\csname sga@target@IV.6.8.8\endcsname{ega4.c04.rem.6-8-8}
\expandafter\def\csname sga@target@IV.6.9.1\endcsname{ega4.c04.thm.6-9-1}
\expandafter\def\csname sga@target@IV.6.9.2\endcsname{ega4.c04.lem.6-9-2}
\expandafter\def\csname sga@target@IV.6.9.3\endcsname{ega4.c04.cor.6-9-3}
\expandafter\def\csname sga@target@IV.7.1.1\endcsname{ega4.c04.def.7-1-1}
\expandafter\def\csname sga@target@IV.7.1.10\endcsname{ega4.c04.prop.7-1-10}
\expandafter\def\csname sga@target@IV.7.1.11\endcsname{ega4.c04.prop.7-1-11}
\expandafter\def\csname sga@target@IV.7.1.12\endcsname{ega4.c04.item.7-1-12}
\expandafter\def\csname sga@target@IV.7.1.13\endcsname{ega4.c04.cor.7-1-13}
\expandafter\def\csname sga@target@IV.7.1.2\endcsname{ega4.c04.prop.7-1-2}
\expandafter\def\csname sga@target@IV.7.1.3\endcsname{ega4.c04.prop.7-1-3}
\expandafter\def\csname sga@target@IV.7.1.3.1\endcsname{ega4.c04.eq.7-1-3-1}
\expandafter\def\csname sga@target@IV.7.1.3.2\endcsname{ega4.c04.eq.7-1-3-2}
\expandafter\def\csname sga@target@IV.7.1.3.3\endcsname{ega4.c04.eq.7-1-3-3}
\expandafter\def\csname sga@target@IV.7.1.3.4\endcsname{ega4.c04.eq.7-1-3-4}
\expandafter\def\csname sga@target@IV.7.1.3.5\endcsname{ega4.c04.eq.7-1-3-5}
\expandafter\def\csname sga@target@IV.7.1.3.6\endcsname{ega4.c04.eq.7-1-3-6}
\expandafter\def\csname sga@target@IV.7.1.3.7\endcsname{ega4.c04.eq.7-1-3-7}
\expandafter\def\csname sga@target@IV.7.1.4\endcsname{ega4.c04.cor.7-1-4}
\expandafter\def\csname sga@target@IV.7.1.5\endcsname{ega4.c04.cor.7-1-5}
\expandafter\def\csname sga@target@IV.7.1.6\endcsname{ega4.c04.rem.7-1-6}
\expandafter\def\csname sga@target@IV.7.1.7\endcsname{ega4.c04.lem.7-1-7}
\expandafter\def\csname sga@target@IV.7.1.8\endcsname{ega4.c04.thm.7-1-8}
\expandafter\def\csname sga@target@IV.7.1.9\endcsname{ega4.c04.def.7-1-9}
\expandafter\def\csname sga@target@IV.7.2.1\endcsname{ega4.c04.statement.7-2-1}
\expandafter\def\csname sga@target@IV.7.2.1.1\endcsname{ega4.c04.statement.7-2-1-1}
\expandafter\def\csname sga@target@IV.7.2.10\endcsname{ega4.c04.statement.7-2-10}
\expandafter\def\csname sga@target@IV.7.2.2\endcsname{ega4.c04.prop.7-2-2}
\expandafter\def\csname sga@target@IV.7.2.3\endcsname{ega4.c04.prop.7-2-3}
\expandafter\def\csname sga@target@IV.7.2.4\endcsname{ega4.c04.prop.7-2-4}
\expandafter\def\csname sga@target@IV.7.2.5\endcsname{ega4.c04.thm.7-2-5}
\expandafter\def\csname sga@target@IV.7.2.5.1\endcsname{ega4.c04.eq.7-2-5-1}
\expandafter\def\csname sga@target@IV.7.2.6\endcsname{ega4.c04.def.7-2-6}
\expandafter\def\csname sga@target@IV.7.2.7\endcsname{ega4.c04.cor.7-2-7}
\expandafter\def\csname sga@target@IV.7.2.8\endcsname{ega4.c04.cor.7-2-8}
\expandafter\def\csname sga@target@IV.7.2.9\endcsname{ega4.c04.cor.7-2-9}
\expandafter\def\csname sga@target@IV.7.3.1\endcsname{ega4.c04.statement.7-3-1}
\expandafter\def\csname sga@target@IV.7.3.10\endcsname{ega4.c04.statement.7-3-10}
\expandafter\def\csname sga@target@IV.7.3.11\endcsname{ega4.c04.rem.7-3-11}
\expandafter\def\csname sga@target@IV.7.3.12\endcsname{ega4.c04.prop.7-3-12}
\expandafter\def\csname sga@target@IV.7.3.12.1\endcsname{ega4.c04.eq.7-3-12-1}
\expandafter\def\csname sga@target@IV.7.3.13\endcsname{ega4.c04.statement.7-3-13}
\expandafter\def\csname sga@target@IV.7.3.14\endcsname{ega4.c04.prop.7-3-14}
\expandafter\def\csname sga@target@IV.7.3.15\endcsname{ega4.c04.prop.7-3-15}
\expandafter\def\csname sga@target@IV.7.3.16\endcsname{ega4.c04.prop.7-3-16}
\expandafter\def\csname sga@target@IV.7.3.16.1\endcsname{ega4.c04.eq.7-3-16-1}
\expandafter\def\csname sga@target@IV.7.3.16.2\endcsname{ega4.c04.lem.7-3-16-2}
\expandafter\def\csname sga@target@IV.7.3.17\endcsname{ega4.c04.prop.7-3-17}
\expandafter\def\csname sga@target@IV.7.3.18\endcsname{ega4.c04.prop.7-3-18}
\expandafter\def\csname sga@target@IV.7.3.19\endcsname{ega4.c04.statement.7-3-19}
\expandafter\def\csname sga@target@IV.7.3.2\endcsname{ega4.c04.lem.7-3-2}
\expandafter\def\csname sga@target@IV.7.3.3\endcsname{ega4.c04.cor.7-3-3}
\expandafter\def\csname sga@target@IV.7.3.4\endcsname{ega4.c04.statement.7-3-4}
\expandafter\def\csname sga@target@IV.7.3.5\endcsname{ega4.c04.item.7-3-5}
\expandafter\def\csname sga@target@IV.7.3.6\endcsname{ega4.c04.statement.7-3-6}
\expandafter\def\csname sga@target@IV.7.3.7\endcsname{ega4.c04.lem.7-3-7}
\expandafter\def\csname sga@target@IV.7.3.8\endcsname{ega4.c04.statement.7-3-8}
\expandafter\def\csname sga@target@IV.7.3.9\endcsname{ega4.c04.item.7-3-9}
\expandafter\def\csname sga@target@IV.7.4.1\endcsname{ega4.c04.statement.7-4-1}
\expandafter\def\csname sga@target@IV.7.4.2\endcsname{ega4.c04.lem.7-4-2}
\expandafter\def\csname sga@target@IV.7.4.3\endcsname{ega4.c04.cor.7-4-3}
\expandafter\def\csname sga@target@IV.7.4.4\endcsname{ega4.c04.thm.7-4-4}
\expandafter\def\csname sga@target@IV.7.4.4.1\endcsname{ega4.c04.lem.7-4-4-1}
\expandafter\def\csname sga@target@IV.7.4.4.2\endcsname{ega4.c04.lem.7-4-4-2}
\expandafter\def\csname sga@target@IV.7.4.5\endcsname{ega4.c04.cor.7-4-5}
\expandafter\def\csname sga@target@IV.7.4.6\endcsname{ega4.c04.prop.7-4-6}
\expandafter\def\csname sga@target@IV.7.4.7\endcsname{ega4.c04.cor.7-4-7}
\expandafter\def\csname sga@target@IV.7.4.8\endcsname{ega4.c04.rem.7-4-8}
\expandafter\def\csname sga@target@IV.7.5.0\endcsname{ega4.c04.statement.7-5-0}
\expandafter\def\csname sga@target@IV.7.5.1\endcsname{ega4.c04.thm.7-5-1}
\expandafter\def\csname sga@target@IV.7.5.1.1\endcsname{ega4.c04.lem.7-5-1-1}
\expandafter\def\csname sga@target@IV.7.5.2\endcsname{ega4.c04.cor.7-5-2}
\expandafter\def\csname sga@target@IV.7.5.3\endcsname{ega4.c04.item.7-5-3}
\expandafter\def\csname sga@target@IV.7.5.4\endcsname{ega4.c04.item.7-5-4}
\expandafter\def\csname sga@target@IV.7.5.5\endcsname{ega4.c04.lem.7-5-5}
\expandafter\def\csname sga@target@IV.7.5.6\endcsname{ega4.c04.prop.7-5-6}
\expandafter\def\csname sga@target@IV.7.5.7\endcsname{ega4.c04.cor.7-5-7}
\expandafter\def\csname sga@target@IV.7.5.7.1\endcsname{ega4.c04.lem.7-5-7-1}
\expandafter\def\csname sga@target@IV.7.5.8\endcsname{ega4.c04.rem.7-5-8}
\expandafter\def\csname sga@target@IV.7.6.1\endcsname{ega4.c04.prop.7-6-1}
\expandafter\def\csname sga@target@IV.7.6.2\endcsname{ega4.c04.cor.7-6-2}
\expandafter\def\csname sga@target@IV.7.6.3\endcsname{ega4.c04.cor.7-6-3}
\expandafter\def\csname sga@target@IV.7.6.4\endcsname{ega4.c04.thm.7-6-4}
\expandafter\def\csname sga@target@IV.7.6.4.1\endcsname{ega4.c04.lem.7-6-4-1}
\expandafter\def\csname sga@target@IV.7.6.5\endcsname{ega4.c04.cor.7-6-5}
\expandafter\def\csname sga@target@IV.7.6.6\endcsname{ega4.c04.cor.7-6-6}
\expandafter\def\csname sga@target@IV.7.6.7\endcsname{ega4.c04.item.7-6-7}
\expandafter\def\csname sga@target@IV.7.7.1\endcsname{ega4.c04.statement.7-7-1}
\expandafter\def\csname sga@target@IV.7.7.2\endcsname{ega4.c04.thm.7-7-2}
\expandafter\def\csname sga@target@IV.7.7.3\endcsname{ega4.c04.cor.7-7-3}
\expandafter\def\csname sga@target@IV.7.7.4\endcsname{ega4.c04.cor.7-7-4}
\expandafter\def\csname sga@target@IV.7.8.1\endcsname{ega4.c04.statement.7-8-1}
\expandafter\def\csname sga@target@IV.7.8.2\endcsname{ega4.c04.def.7-8-2}
\expandafter\def\csname sga@target@IV.7.8.3\endcsname{ega4.c04.statement.7-8-3}
\expandafter\def\csname sga@target@IV.7.8.3.1\endcsname{ega4.c04.statement.7-8-3-1}
\expandafter\def\csname sga@target@IV.7.8.4\endcsname{ega4.c04.item.7-8-4}
\expandafter\def\csname sga@target@IV.7.8.4.1\endcsname{ega4.c04.statement.7-8-4-1}
\expandafter\def\csname sga@target@IV.7.8.5\endcsname{ega4.c04.statement.7-8-5}
\expandafter\def\csname sga@target@IV.7.8.6\endcsname{ega4.c04.prop.7-8-6}
\expandafter\def\csname sga@target@IV.7.9.1\endcsname{ega4.c04.statement.7-9-1}
\expandafter\def\csname sga@target@IV.7.9.10\endcsname{ega4.c04.item.7-9-10}
\expandafter\def\csname sga@target@IV.7.9.2\endcsname{ega4.c04.statement.7-9-2}
\expandafter\def\csname sga@target@IV.7.9.3\endcsname{ega4.c04.prop.7-9-3}
\expandafter\def\csname sga@target@IV.7.9.3.1\endcsname{ega4.c04.lem.7-9-3-1}
\expandafter\def\csname sga@target@IV.7.9.4\endcsname{ega4.c04.cor.7-9-4}
\expandafter\def\csname sga@target@IV.7.9.5\endcsname{ega4.c04.prop.7-9-5}
\expandafter\def\csname sga@target@IV.7.9.6\endcsname{ega4.c04.rem.7-9-6}
\expandafter\def\csname sga@target@IV.7.9.7\endcsname{ega4.c04.statement.7-9-7}
\expandafter\def\csname sga@target@IV.7.9.8\endcsname{ega4.c04.prop.7-9-8}
\expandafter\def\csname sga@target@IV.7.9.8.1\endcsname{ega4.c04.statement.7-9-8-1}
\expandafter\def\csname sga@target@IV.7.9.9\endcsname{ega4.c04.cor.7-9-9}
\expandafter\def\csname sga@target@IV.8.1.1\endcsname{ega4.c04.statement.8-1-1}
\expandafter\def\csname sga@target@IV.8.1.2\endcsname{ega4.c04.statement.8-1-2}
\expandafter\def\csname sga@target@IV.8.1.3\endcsname{ega4.c04.statement.8-1-3}
\expandafter\def\csname sga@target@IV.8.10.1\endcsname{ega4.c04.prop.8-10-1}
\expandafter\def\csname sga@target@IV.8.10.2\endcsname{ega4.c04.cor.8-10-2}
\expandafter\def\csname sga@target@IV.8.10.3\endcsname{ega4.c04.prop.8-10-3}
\expandafter\def\csname sga@target@IV.8.10.3.1\endcsname{ega4.c04.rem.8-10-3-1}
\expandafter\def\csname sga@target@IV.8.10.4\endcsname{ega4.c04.prop.8-10-4}
\expandafter\def\csname sga@target@IV.8.10.5\endcsname{ega4.c04.thm.8-10-5}
\expandafter\def\csname sga@target@IV.8.10.5.1\endcsname{ega4.c04.prop.8-10-5-1}
\expandafter\def\csname sga@target@IV.8.10.5.2\endcsname{ega4.c04.lem.8-10-5-2}
\expandafter\def\csname sga@target@IV.8.11\endcsname{ega4.c04.subsec.8-11.key-iv-8-11}
\expandafter\def\csname sga@target@IV.8.11.1\endcsname{ega4.c04.thm.8-11-1}
\expandafter\def\csname sga@target@IV.8.11.2\endcsname{ega4.c04.thm.8-11-2}
\expandafter\def\csname sga@target@IV.8.11.3\endcsname{ega4.c04.rem.8-11-3}
\expandafter\def\csname sga@target@IV.8.11.4\endcsname{ega4.c04.statement.8-11-4}
\expandafter\def\csname sga@target@IV.8.11.5\endcsname{ega4.c04.prop.8-11-5}
\expandafter\def\csname sga@target@IV.8.11.5.1\endcsname{ega4.c04.rem.8-11-5-1}
\expandafter\def\csname sga@target@IV.8.11.6\endcsname{ega4.c04.prop.8-11-6}
\expandafter\def\csname sga@target@IV.8.12\endcsname{ega4.c04.subsec.8-12.key-iv-8-12}
\expandafter\def\csname sga@target@IV.8.12.1\endcsname{ega4.c04.lem.8-12-1}
\expandafter\def\csname sga@target@IV.8.12.10\endcsname{ega4.c04.cor.8-12-10}
\expandafter\def\csname sga@target@IV.8.12.10.1\endcsname{ega4.c04.lem.8-12-10-1}
\expandafter\def\csname sga@target@IV.8.12.11\endcsname{ega4.c04.cor.8-12-11}
\expandafter\def\csname sga@target@IV.8.12.12\endcsname{ega4.c04.cor.8-12-12}
\expandafter\def\csname sga@target@IV.8.12.2\endcsname{ega4.c04.lem.8-12-2}
\expandafter\def\csname sga@target@IV.8.12.3\endcsname{ega4.c04.prop.8-12-3}
\expandafter\def\csname sga@target@IV.8.12.4\endcsname{ega4.c04.cor.8-12-4}
\expandafter\def\csname sga@target@IV.8.12.5\endcsname{ega4.c04.statement.8-12-5}
\expandafter\def\csname sga@target@IV.8.12.6\endcsname{ega4.c04.thm.8-12-6}
\expandafter\def\csname sga@target@IV.8.12.6.1\endcsname{ega4.c04.eq.8-12-6-1}
\expandafter\def\csname sga@target@IV.8.12.7\endcsname{ega4.c04.rem.8-12-7}
\expandafter\def\csname sga@target@IV.8.12.8\endcsname{ega4.c04.cor.8-12-8}
\expandafter\def\csname sga@target@IV.8.12.9\endcsname{ega4.c04.cor.8-12-9}
\expandafter\def\csname sga@target@IV.8.13.1\endcsname{ega4.c04.prop.8-13-1}
\expandafter\def\csname sga@target@IV.8.13.1.1\endcsname{ega4.c04.eq.8-13-1-1}
\expandafter\def\csname sga@target@IV.8.13.2\endcsname{ega4.c04.cor.8-13-2}
\expandafter\def\csname sga@target@IV.8.13.2.1\endcsname{ega4.c04.eq.8-13-2-1}
\expandafter\def\csname sga@target@IV.8.13.3\endcsname{ega4.c04.statement.8-13-3}
\expandafter\def\csname sga@target@IV.8.13.4\endcsname{ega4.c04.statement.8-13-4}
\expandafter\def\csname sga@target@IV.8.13.4.1\endcsname{ega4.c04.eq.8-13-4-1}
\expandafter\def\csname sga@target@IV.8.13.4.2\endcsname{ega4.c04.eq.8-13-4-2}
\expandafter\def\csname sga@target@IV.8.13.5\endcsname{ega4.c04.prop.8-13-5}
\expandafter\def\csname sga@target@IV.8.13.6\endcsname{ega4.c04.rem.8-13-6}
\expandafter\def\csname sga@target@IV.8.14.1\endcsname{ega4.c04.statement.8-14-1}
\expandafter\def\csname sga@target@IV.8.14.2\endcsname{ega4.c04.prop.8-14-2}
\expandafter\def\csname sga@target@IV.8.14.2.1\endcsname{ega4.c04.eq.8-14-2-1}
\expandafter\def\csname sga@target@IV.8.14.2.2\endcsname{ega4.c04.cor.8-14-2-2}
\expandafter\def\csname sga@target@IV.8.14.2.3\endcsname{ega4.c04.eq.8-14-2-3}
\expandafter\def\csname sga@target@IV.8.14.2.4\endcsname{ega4.c04.eq.8-14-2-4}
\expandafter\def\csname sga@target@IV.8.14.3\endcsname{ega4.c04.rem.8-14-3}
\expandafter\def\csname sga@target@IV.8.2.1\endcsname{ega4.c04.statement.8-2-1}
\expandafter\def\csname sga@target@IV.8.2.1.1\endcsname{ega4.c04.eq.8-2-1-1}
\expandafter\def\csname sga@target@IV.8.2.10\endcsname{ega4.c04.cor.8-2-10}
\expandafter\def\csname sga@target@IV.8.2.11\endcsname{ega4.c04.cor.8-2-11}
\expandafter\def\csname sga@target@IV.8.2.12\endcsname{ega4.c04.cor.8-2-12}
\expandafter\def\csname sga@target@IV.8.2.12.1\endcsname{ega4.c04.eq.8-2-12-1}
\expandafter\def\csname sga@target@IV.8.2.13\endcsname{ega4.c04.prop.8-2-13}
\expandafter\def\csname sga@target@IV.8.2.14\endcsname{ega4.c04.rem.8-2-14}
\expandafter\def\csname sga@target@IV.8.2.2\endcsname{ega4.c04.statement.8-2-2}
\expandafter\def\csname sga@target@IV.8.2.3\endcsname{ega4.c04.prop.8-2-3}
\expandafter\def\csname sga@target@IV.8.2.3.1\endcsname{ega4.c04.eq.8-2-3-1}
\expandafter\def\csname sga@target@IV.8.2.4\endcsname{ega4.c04.lem.8-2-4}
\expandafter\def\csname sga@target@IV.8.2.5\endcsname{ega4.c04.prop.8-2-5}
\expandafter\def\csname sga@target@IV.8.2.6\endcsname{ega4.c04.lem.8-2-6}
\expandafter\def\csname sga@target@IV.8.2.7\endcsname{ega4.c04.rem.8-2-7}
\expandafter\def\csname sga@target@IV.8.2.8\endcsname{ega4.c04.statement.8-2-8}
\expandafter\def\csname sga@target@IV.8.2.9\endcsname{ega4.c04.prop.8-2-9}
\expandafter\def\csname sga@target@IV.8.3.1\endcsname{ega4.c04.statement.8-3-1}
\expandafter\def\csname sga@target@IV.8.3.10\endcsname{ega4.c04.statement.8-3-10}
\expandafter\def\csname sga@target@IV.8.3.10.1\endcsname{ega4.c04.eq.8-3-10-1}
\expandafter\def\csname sga@target@IV.8.3.11\endcsname{ega4.c04.thm.8-3-11}
\expandafter\def\csname sga@target@IV.8.3.12\endcsname{ega4.c04.cor.8-3-12}
\expandafter\def\csname sga@target@IV.8.3.13\endcsname{ega4.c04.prop.8-3-13}
\expandafter\def\csname sga@target@IV.8.3.2\endcsname{ega4.c04.thm.8-3-2}
\expandafter\def\csname sga@target@IV.8.3.2.1\endcsname{ega4.c04.eq.8-3-2-1}
\expandafter\def\csname sga@target@IV.8.3.3\endcsname{ega4.c04.cor.8-3-3}
\expandafter\def\csname sga@target@IV.8.3.4\endcsname{ega4.c04.cor.8-3-4}
\expandafter\def\csname sga@target@IV.8.3.5\endcsname{ega4.c04.cor.8-3-5}
\expandafter\def\csname sga@target@IV.8.3.6\endcsname{ega4.c04.cor.8-3-6}
\expandafter\def\csname sga@target@IV.8.3.7\endcsname{ega4.c04.cor.8-3-7}
\expandafter\def\csname sga@target@IV.8.3.7.1\endcsname{ega4.c04.eq.8-3-7-1}
\expandafter\def\csname sga@target@IV.8.3.8\endcsname{ega4.c04.prop.8-3-8}
\expandafter\def\csname sga@target@IV.8.3.9\endcsname{ega4.c04.statement.8-3-9}
\expandafter\def\csname sga@target@IV.8.3.9.1\endcsname{ega4.c04.eq.8-3-9-1}
\expandafter\def\csname sga@target@IV.8.3.9.2\endcsname{ega4.c04.eq.8-3-9-2}
\expandafter\def\csname sga@target@IV.8.3.9.3\endcsname{ega4.c04.eq.8-3-9-3}
\expandafter\def\csname sga@target@IV.8.3.9.4\endcsname{ega4.c04.eq.8-3-9-4}
\expandafter\def\csname sga@target@IV.8.3.9.5\endcsname{ega4.c04.eq.8-3-9-5}
\expandafter\def\csname sga@target@IV.8.4.1\endcsname{ega4.c04.prop.8-4-1}
\expandafter\def\csname sga@target@IV.8.4.2\endcsname{ega4.c04.prop.8-4-2}
\expandafter\def\csname sga@target@IV.8.4.2.1\endcsname{ega4.c04.lem.8-4-2-1}
\expandafter\def\csname sga@target@IV.8.4.3\endcsname{ega4.c04.cor.8-4-3}
\expandafter\def\csname sga@target@IV.8.4.4\endcsname{ega4.c04.prop.8-4-4}
\expandafter\def\csname sga@target@IV.8.4.5\endcsname{ega4.c04.cor.8-4-5}
\expandafter\def\csname sga@target@IV.8.4.6\endcsname{ega4.c04.rem.8-4-6}
\expandafter\def\csname sga@target@IV.8.5.1\endcsname{ega4.c04.statement.8-5-1}
\expandafter\def\csname sga@target@IV.8.5.1.1\endcsname{ega4.c04.eq.8-5-1-1}
\expandafter\def\csname sga@target@IV.8.5.1.2\endcsname{ega4.c04.eq.8-5-1-2}
\expandafter\def\csname sga@target@IV.8.5.1.3\endcsname{ega4.c04.eq.8-5-1-3}
\expandafter\def\csname sga@target@IV.8.5.1.4\endcsname{ega4.c04.eq.8-5-1-4}
\expandafter\def\csname sga@target@IV.8.5.10\endcsname{ega4.c04.statement.8-5-10}
\expandafter\def\csname sga@target@IV.8.5.10.1\endcsname{ega4.c04.eq.8-5-10-1}
\expandafter\def\csname sga@target@IV.8.5.10.2\endcsname{ega4.c04.eq.8-5-10-2}
\expandafter\def\csname sga@target@IV.8.5.11\endcsname{ega4.c04.prop.8-5-11}
\expandafter\def\csname sga@target@IV.8.5.2\endcsname{ega4.c04.thm.8-5-2}
\expandafter\def\csname sga@target@IV.8.5.2.1\endcsname{ega4.c04.lem.8-5-2-1}
\expandafter\def\csname sga@target@IV.8.5.2.2\endcsname{ega4.c04.eq.8-5-2-2}
\expandafter\def\csname sga@target@IV.8.5.2.3\endcsname{ega4.c04.eq.8-5-2-3}
\expandafter\def\csname sga@target@IV.8.5.2.4\endcsname{ega4.c04.cor.8-5-2-4}
\expandafter\def\csname sga@target@IV.8.5.2.5\endcsname{ega4.c04.cor.8-5-2-5}
\expandafter\def\csname sga@target@IV.8.5.3\endcsname{ega4.c04.statement.8-5-3}
\expandafter\def\csname sga@target@IV.8.5.4\endcsname{ega4.c04.cor.8-5-4}
\expandafter\def\csname sga@target@IV.8.5.5\endcsname{ega4.c04.prop.8-5-5}
\expandafter\def\csname sga@target@IV.8.5.6\endcsname{ega4.c04.prop.8-5-6}
\expandafter\def\csname sga@target@IV.8.5.7\endcsname{ega4.c04.cor.8-5-7}
\expandafter\def\csname sga@target@IV.8.5.8\endcsname{ega4.c04.cor.8-5-8}
\expandafter\def\csname sga@target@IV.8.5.9\endcsname{ega4.c04.lem.8-5-9}
\expandafter\def\csname sga@target@IV.8.6.1\endcsname{ega4.c04.statement.8-6-1}
\expandafter\def\csname sga@target@IV.8.6.2\endcsname{ega4.c04.statement.8-6-2}
\expandafter\def\csname sga@target@IV.8.6.2.1\endcsname{ega4.c04.eq.8-6-2-1}
\expandafter\def\csname sga@target@IV.8.6.2.2\endcsname{ega4.c04.eq.8-6-2-2}
\expandafter\def\csname sga@target@IV.8.6.2.3\endcsname{ega4.c04.eq.8-6-2-3}
\expandafter\def\csname sga@target@IV.8.6.3\endcsname{ega4.c04.prop.8-6-3}
\expandafter\def\csname sga@target@IV.8.6.3.1\endcsname{ega4.c04.cor.8-6-3-1}
\expandafter\def\csname sga@target@IV.8.6.4\endcsname{ega4.c04.cor.8-6-4}
\expandafter\def\csname sga@target@IV.8.7.1\endcsname{ega4.c04.prop.8-7-1}
\expandafter\def\csname sga@target@IV.8.7.2\endcsname{ega4.c04.prop.8-7-2}
\expandafter\def\csname sga@target@IV.8.7.3\endcsname{ega4.c04.cor.8-7-3}
\expandafter\def\csname sga@target@IV.8.7.4\endcsname{ega4.c04.rem.8-7-4}
\expandafter\def\csname sga@target@IV.8.8.1\endcsname{ega4.c04.statement.8-8-1}
\expandafter\def\csname sga@target@IV.8.8.1.1\endcsname{ega4.c04.eq.8-8-1-1}
\expandafter\def\csname sga@target@IV.8.8.2\endcsname{ega4.c04.thm.8-8-2}
\expandafter\def\csname sga@target@IV.8.8.2.1\endcsname{ega4.c04.lem.8-8-2-1}
\expandafter\def\csname sga@target@IV.8.8.2.2\endcsname{ega4.c04.lem.8-8-2-2}
\expandafter\def\csname sga@target@IV.8.8.2.3\endcsname{ega4.c04.lem.8-8-2-3}
\expandafter\def\csname sga@target@IV.8.8.2.4\endcsname{ega4.c04.cor.8-8-2-4}
\expandafter\def\csname sga@target@IV.8.8.2.5\endcsname{ega4.c04.cor.8-8-2-5}
\expandafter\def\csname sga@target@IV.8.8.3\endcsname{ega4.c04.statement.8-8-3}
\expandafter\def\csname sga@target@IV.8.9.1\endcsname{ega4.c04.prop.8-9-1}
\expandafter\def\csname sga@target@IV.8.9.2\endcsname{ega4.c04.statement.8-9-2}
\expandafter\def\csname sga@target@IV.8.9.3\endcsname{ega4.c04.prop.8-9-3}
\expandafter\def\csname sga@target@IV.8.9.4\endcsname{ega4.c04.prop.8-9-4}
\expandafter\def\csname sga@target@IV.8.9.4.1\endcsname{ega4.c04.lem.8-9-4-1}
\expandafter\def\csname sga@target@IV.8.9.5\endcsname{ega4.c04.cor.8-9-5}
\expandafter\def\csname sga@target@IV.9.1.1\endcsname{ega4.c04.prop.9-1-1}
\expandafter\def\csname sga@target@IV.9.1.2\endcsname{ega4.c04.cor.9-1-2}
\expandafter\def\csname sga@target@IV.9.1.3\endcsname{ega4.c04.rem.9-1-3}
\expandafter\def\csname sga@target@IV.9.1.4\endcsname{ega4.c04.prop.9-1-4}
\expandafter\def\csname sga@target@IV.9.1.5\endcsname{ega4.c04.cor.9-1-5}
\expandafter\def\csname sga@target@IV.9.1.6\endcsname{ega4.c04.prop.9-1-6}
\expandafter\def\csname sga@target@IV.9.2\endcsname{ega4.c04.subsec.9-2.key-iv-9-2}
\expandafter\def\csname sga@target@IV.9.2.1\endcsname{ega4.c04.def.9-2-1}
\expandafter\def\csname sga@target@IV.9.2.2\endcsname{ega4.c04.item.9-2-2}
\expandafter\def\csname sga@target@IV.9.2.3\endcsname{ega4.c04.prop.9-2-3}
\expandafter\def\csname sga@target@IV.9.2.4\endcsname{ega4.c04.cor.9-2-4}
\expandafter\def\csname sga@target@IV.9.2.5\endcsname{ega4.c04.rem.9-2-5}
\expandafter\def\csname sga@target@IV.9.2.6\endcsname{ega4.c04.prop.9-2-6}
\expandafter\def\csname sga@target@IV.9.2.6.1\endcsname{ega4.c04.prop.9-2-6-1}
\expandafter\def\csname sga@target@IV.9.2.6.2\endcsname{ega4.c04.cor.9-2-6-2}
\expandafter\def\csname sga@target@IV.9.3.1\endcsname{ega4.c04.prop.9-3-1}
\expandafter\def\csname sga@target@IV.9.3.2\endcsname{ega4.c04.prop.9-3-2}
\expandafter\def\csname sga@target@IV.9.3.3\endcsname{ega4.c04.prop.9-3-3}
\expandafter\def\csname sga@target@IV.9.3.4\endcsname{ega4.c04.rem.9-3-4}
\expandafter\def\csname sga@target@IV.9.3.5\endcsname{ega4.c04.prop.9-3-5}
\expandafter\def\csname sga@target@IV.9.3.6\endcsname{ega4.c04.statement.9-3-6}
\expandafter\def\csname sga@target@IV.9.4.1\endcsname{ega4.c04.statement.9-4-1}
\expandafter\def\csname sga@target@IV.9.4.2\endcsname{ega4.c04.prop.9-4-2}
\expandafter\def\csname sga@target@IV.9.4.3\endcsname{ega4.c04.cor.9-4-3}
\expandafter\def\csname sga@target@IV.9.4.3.1\endcsname{ega4.c04.eq.9-4-3-1}
\expandafter\def\csname sga@target@IV.9.4.4\endcsname{ega4.c04.prop.9-4-4}
\expandafter\def\csname sga@target@IV.9.4.5\endcsname{ega4.c04.cor.9-4-5}
\expandafter\def\csname sga@target@IV.9.4.6\endcsname{ega4.c04.cor.9-4-6}
\expandafter\def\csname sga@target@IV.9.4.7\endcsname{ega4.c04.prop.9-4-7}
\expandafter\def\csname sga@target@IV.9.4.7.1\endcsname{ega4.c04.lem.9-4-7-1}
\expandafter\def\csname sga@target@IV.9.4.7.2\endcsname{ega4.c04.rem.9-4-7-2}
\expandafter\def\csname sga@target@IV.9.4.8\endcsname{ega4.c04.prop.9-4-8}
\expandafter\def\csname sga@target@IV.9.4.9\endcsname{ega4.c04.rem.9-4-9}
\expandafter\def\csname sga@target@IV.9.5.1\endcsname{ega4.c04.prop.9-5-1}
\expandafter\def\csname sga@target@IV.9.5.2\endcsname{ega4.c04.cor.9-5-2}
\expandafter\def\csname sga@target@IV.9.5.3\endcsname{ega4.c04.prop.9-5-3}
\expandafter\def\csname sga@target@IV.9.5.3.case2\endcsname{ega4.c04.item.9-5-3-case2}
\expandafter\def\csname sga@target@IV.9.5.4\endcsname{ega4.c04.cor.9-5-4}
\expandafter\def\csname sga@target@IV.9.5.5\endcsname{ega4.c04.prop.9-5-5}
\expandafter\def\csname sga@target@IV.9.5.6\endcsname{ega4.c04.cor.9-5-6}
\expandafter\def\csname sga@target@IV.9.5.7\endcsname{ega4.c04.rem.9-5-7}
\expandafter\def\csname sga@target@IV.9.6.1\endcsname{ega4.c04.prop.9-6-1}
\expandafter\def\csname sga@target@IV.9.6.2\endcsname{ega4.c04.prop.9-6-2}
\expandafter\def\csname sga@target@IV.9.6.3\endcsname{ega4.c04.prop.9-6-3}
\expandafter\def\csname sga@target@IV.9.6.3.1\endcsname{ega4.c04.eq.9-6-3-1}
\expandafter\def\csname sga@target@IV.9.6.4\endcsname{ega4.c04.cor.9-6-4}
\expandafter\def\csname sga@target@IV.9.6.5\endcsname{ega4.c04.cor.9-6-5}
\expandafter\def\csname sga@target@IV.9.6.6\endcsname{ega4.c04.sec.9-6-6}
\expandafter\def\csname sga@target@IV.9.6.7\endcsname{ega4.c04.sec.9-6-7}
\expandafter\def\csname sga@target@IV.9.7\endcsname{ega4.c04.subsec.9-7.key-iv-9-7}
\expandafter\def\csname sga@target@IV.9.7.1\endcsname{ega4.c04.lem.9-7-1}
\expandafter\def\csname sga@target@IV.9.7.10\endcsname{ega4.c04.cor.9-7-10}
\expandafter\def\csname sga@target@IV.9.7.11\endcsname{ega4.c04.rem.9-7-11}
\expandafter\def\csname sga@target@IV.9.7.12\endcsname{ega4.c04.prop.9-7-12}
\expandafter\def\csname sga@target@IV.9.7.2\endcsname{ega4.c04.lem.9-7-2}
\expandafter\def\csname sga@target@IV.9.7.3\endcsname{ega4.c04.subsec.9-7-3}
\expandafter\def\csname sga@target@IV.9.7.4\endcsname{ega4.c04.lem.9-7-4}
\expandafter\def\csname sga@target@IV.9.7.5\endcsname{ega4.c04.lem.9-7-5}
\expandafter\def\csname sga@target@IV.9.7.5.1\endcsname{ega4.c04.lem.9-7-5-1}
\expandafter\def\csname sga@target@IV.9.7.6\endcsname{ega4.c04.prop.9-7-6}
\expandafter\def\csname sga@target@IV.9.7.7\endcsname{ega4.c04.thm.9-7-7}
\expandafter\def\csname sga@target@IV.9.7.8\endcsname{ega4.c04.prop.9-7-8}
\expandafter\def\csname sga@target@IV.9.7.9\endcsname{ega4.c04.cor.9-7-9}
\expandafter\def\csname sga@target@IV.9.8\endcsname{ega4.c04.subsec.9-8.key-iv-9-8}
\expandafter\def\csname sga@target@IV.9.8.1\endcsname{ega4.c04.prop.9-8-1}
\expandafter\def\csname sga@target@IV.9.8.1-proof\endcsname{ega4.c04.subsec.9-8-1-proof}
\expandafter\def\csname sga@target@IV.9.8.2\endcsname{ega4.c04.subsec.9-8-2}
\expandafter\def\csname sga@target@IV.9.8.3\endcsname{ega4.c04.prop.9-8-3}
\expandafter\def\csname sga@target@IV.9.8.3-proof\endcsname{ega4.c04.subsec.9-8-3-proof}
\expandafter\def\csname sga@target@IV.9.8.4\endcsname{ega4.c04.cor.9-8-4}
\expandafter\def\csname sga@target@IV.9.8.4-proof\endcsname{ega4.c04.item.9-8-4-proof}
\expandafter\def\csname sga@target@IV.9.8.5\endcsname{ega4.c04.prop.9-8-5}
\expandafter\def\csname sga@target@IV.9.8.5-proof\endcsname{ega4.c04.item.9-8-5-proof}
\expandafter\def\csname sga@target@IV.9.8.6\endcsname{ega4.c04.prop.9-8-6}
\expandafter\def\csname sga@target@IV.9.8.6-proof\endcsname{ega4.c04.item.9-8-6-proof}
\expandafter\def\csname sga@target@IV.9.8.7\endcsname{ega4.c04.cor.9-8-7}
\expandafter\def\csname sga@target@IV.9.8.7-proof\endcsname{ega4.c04.item.9-8-7-proof}
\expandafter\def\csname sga@target@IV.9.8.8\endcsname{ega4.c04.prop.9-8-8}
\expandafter\def\csname sga@target@IV.9.8.8-proof\endcsname{ega4.c04.item.9-8-8-proof}
\expandafter\def\csname sga@target@IV.9.8.9\endcsname{ega4.c04.rem.9-8-9}
\expandafter\def\csname sga@target@IV.9.8.9.1\endcsname{ega4.c04.item.9-8-9-1}
\expandafter\def\csname sga@target@IV.9.9\endcsname{ega4.c04.subsec.9-9.key-iv-9-9}
\expandafter\def\csname sga@target@IV.9.9.1\endcsname{ega4.c04.prop.9-9-1}
\expandafter\def\csname sga@target@IV.9.9.1-proof\endcsname{ega4.c04.item.9-9-1-proof}
\expandafter\def\csname sga@target@IV.9.9.1.1\endcsname{ega4.c04.eq.9-9-1-1}
\expandafter\def\csname sga@target@IV.9.9.2\endcsname{ega4.c04.prop.9-9-2}
\expandafter\def\csname sga@target@IV.9.9.2-proof\endcsname{ega4.c04.item.9-9-2-proof}
\expandafter\def\csname sga@target@IV.9.9.2.1\endcsname{ega4.c04.eq.9-9-2-1}
\expandafter\def\csname sga@target@IV.9.9.2.2\endcsname{ega4.c04.eq.9-9-2-2}
\expandafter\def\csname sga@target@IV.9.9.3\endcsname{ega4.c04.cor.9-9-3}
\expandafter\def\csname sga@target@IV.9.9.3-proof\endcsname{ega4.c04.statement.9-9-3-proof}
\expandafter\def\csname sga@target@IV.9.9.4\endcsname{ega4.c04.prop.9-9-4}
\expandafter\def\csname sga@target@IV.9.9.4-proof\endcsname{ega4.c04.item.9-9-4-proof}
\expandafter\def\csname sga@target@IV.9.9.5\endcsname{ega4.c04.cor.9-9-5}
\expandafter\def\csname sga@target@IV.9.9.5-proof\endcsname{ega4.c04.item.9-9-5-proof}
\expandafter\def\csname sga@target@IV.9.9.6\endcsname{ega4.c04.prop.9-9-6}
\expandafter\def\csname sga@target@IV.9.9.6-proof\endcsname{ega4.c04.item.9-9-6-proof}
\expandafter\def\csname sga@target@section:IV.1\endcsname{ega4.c04.sec.1}
\expandafter\def\csname sga@target@section:IV.10\endcsname{ega4.c04.sec.10}
\expandafter\def\csname sga@target@section:IV.11\endcsname{ega4.c04.sec.11}
\expandafter\def\csname sga@target@section:IV.12\endcsname{ega4.c04.sec.12}
\expandafter\def\csname sga@target@section:IV.13\endcsname{ega4.c04.sec.13}
\expandafter\def\csname sga@target@section:IV.14\endcsname{ega4.c04.sec.14}
\expandafter\def\csname sga@target@section:IV.15\endcsname{ega4.c04.sec.15}
\expandafter\def\csname sga@target@section:IV.16\endcsname{ega4.c04.sec.16}
\expandafter\def\csname sga@target@section:IV.17\endcsname{ega4.c04.sec.17}
\expandafter\def\csname sga@target@section:IV.18\endcsname{ega4.c04.sec.18}
\expandafter\def\csname sga@target@section:IV.19\endcsname{ega4.c04.sec.19}
\expandafter\def\csname sga@target@section:IV.2\endcsname{ega4.c04.sec.2}
\expandafter\def\csname sga@target@section:IV.21\endcsname{ega4.c04.sec.21}
\expandafter\def\csname sga@target@section:IV.3\endcsname{ega4.c04.sec.3}
\expandafter\def\csname sga@target@section:IV.4\endcsname{ega4.c04.sec.4}
\expandafter\def\csname sga@target@section:IV.5\endcsname{ega4.c04.sec.5}
\expandafter\def\csname sga@target@section:IV.6\endcsname{ega4.c04.sec.6}
\expandafter\def\csname sga@target@section:IV.7\endcsname{ega4.c04.sec.7}
\expandafter\def\csname sga@target@section:IV.8\endcsname{ega4.c04.sec.8}
\expandafter\def\csname sga@target@section:IV.9\endcsname{ega4.c04.sec.9}
\expandafter\def\csname sga@target@subsection:IV.1.1\endcsname{ega4.c04.subsec.1-1}
\expandafter\def\csname sga@target@subsection:IV.1.10\endcsname{ega4.c04.subsec.1-10}
\expandafter\def\csname sga@target@subsection:IV.1.2\endcsname{ega4.c04.subsec.1-2}
\expandafter\def\csname sga@target@subsection:IV.1.3\endcsname{ega4.c04.subsec.1-3}
\expandafter\def\csname sga@target@subsection:IV.1.4\endcsname{ega4.c04.subsec.1-4}
\expandafter\def\csname sga@target@subsection:IV.1.5\endcsname{ega4.c04.subsec.1-5}
\expandafter\def\csname sga@target@subsection:IV.1.6\endcsname{ega4.c04.subsec.1-6}
\expandafter\def\csname sga@target@subsection:IV.1.7\endcsname{ega4.c04.subsec.1-7}
\expandafter\def\csname sga@target@subsection:IV.1.8\endcsname{ega4.c04.subsec.1-8}
\expandafter\def\csname sga@target@subsection:IV.1.9\endcsname{ega4.c04.subsec.1-9}
\expandafter\def\csname sga@target@subsection:IV.10.1\endcsname{ega4.c04.subsec.10-1}
\expandafter\def\csname sga@target@subsection:IV.10.10\endcsname{ega4.c04.subsec.10-10}
\expandafter\def\csname sga@target@subsection:IV.10.2\endcsname{ega4.c04.subsec.10-2}
\expandafter\def\csname sga@target@subsection:IV.10.3\endcsname{ega4.c04.subsec.10-3}
\expandafter\def\csname sga@target@subsection:IV.10.4\endcsname{ega4.c04.subsec.10-4}
\expandafter\def\csname sga@target@subsection:IV.10.5\endcsname{ega4.c04.subsec.10-5}
\expandafter\def\csname sga@target@subsection:IV.10.6\endcsname{ega4.c04.subsec.10-6}
\expandafter\def\csname sga@target@subsection:IV.10.7\endcsname{ega4.c04.subsec.10-7}
\expandafter\def\csname sga@target@subsection:IV.10.8\endcsname{ega4.c04.subsec.10-8}
\expandafter\def\csname sga@target@subsection:IV.10.9\endcsname{ega4.c04.subsec.10-9}
\expandafter\def\csname sga@target@subsection:IV.11.1\endcsname{ega4.c04.subsec.11-1}
\expandafter\def\csname sga@target@subsection:IV.11.10\endcsname{ega4.c04.subsec.11-10}
\expandafter\def\csname sga@target@subsection:IV.11.2\endcsname{ega4.c04.subsec.11-2}
\expandafter\def\csname sga@target@subsection:IV.11.3\endcsname{ega4.c04.subsec.11-3}
\expandafter\def\csname sga@target@subsection:IV.11.4\endcsname{ega4.c04.subsec.11-4}
\expandafter\def\csname sga@target@subsection:IV.11.5\endcsname{ega4.c04.subsec.11-5}
\expandafter\def\csname sga@target@subsection:IV.11.6\endcsname{ega4.c04.subsec.11-6}
\expandafter\def\csname sga@target@subsection:IV.11.7\endcsname{ega4.c04.subsec.11-7}
\expandafter\def\csname sga@target@subsection:IV.11.8\endcsname{ega4.c04.subsec.11-8}
\expandafter\def\csname sga@target@subsection:IV.11.9\endcsname{ega4.c04.subsec.11-9}
\expandafter\def\csname sga@target@subsection:IV.12.0\endcsname{ega4.c04.subsec.12-0}
\expandafter\def\csname sga@target@subsection:IV.12.1\endcsname{ega4.c04.subsec.12-1}
\expandafter\def\csname sga@target@subsection:IV.12.2\endcsname{ega4.c04.subsec.12-2}
\expandafter\def\csname sga@target@subsection:IV.12.3\endcsname{ega4.c04.subsec.12-3}
\expandafter\def\csname sga@target@subsection:IV.13.1\endcsname{ega4.c04.subsec.13-1}
\expandafter\def\csname sga@target@subsection:IV.13.2\endcsname{ega4.c04.subsec.13-2}
\expandafter\def\csname sga@target@subsection:IV.13.3\endcsname{ega4.c04.subsec.13-3}
\expandafter\def\csname sga@target@subsection:IV.14.1\endcsname{ega4.c04.subsec.14-1}
\expandafter\def\csname sga@target@subsection:IV.14.2\endcsname{ega4.c04.subsec.14-2}
\expandafter\def\csname sga@target@subsection:IV.14.3\endcsname{ega4.c04.subsec.14-3}
\expandafter\def\csname sga@target@subsection:IV.14.4\endcsname{ega4.c04.subsec.14-4}
\expandafter\def\csname sga@target@subsection:IV.14.5\endcsname{ega4.c04.subsec.14-5}
\expandafter\def\csname sga@target@subsection:IV.15.1\endcsname{ega4.c04.subsec.15-1}
\expandafter\def\csname sga@target@subsection:IV.15.2\endcsname{ega4.c04.subsec.15-2}
\expandafter\def\csname sga@target@subsection:IV.15.3\endcsname{ega4.c04.subsec.15-3}
\expandafter\def\csname sga@target@subsection:IV.15.4\endcsname{ega4.c04.subsec.15-4}
\expandafter\def\csname sga@target@subsection:IV.15.5\endcsname{ega4.c04.subsec.15-5}
\expandafter\def\csname sga@target@subsection:IV.15.6\endcsname{ega4.c04.subsec.15-6}
\expandafter\def\csname sga@target@subsection:IV.15.7\endcsname{ega4.c04.subsec.15-7}
\expandafter\def\csname sga@target@subsection:IV.17.1\endcsname{ega4.c04.subsec.17-1}
\expandafter\def\csname sga@target@subsection:IV.17.10\endcsname{ega4.c04.subsec.17-10}
\expandafter\def\csname sga@target@subsection:IV.17.11\endcsname{ega4.c04.subsec.17-11}
\expandafter\def\csname sga@target@subsection:IV.17.12\endcsname{ega4.c04.subsec.17-12}
\expandafter\def\csname sga@target@subsection:IV.17.13\endcsname{ega4.c04.subsec.17-13}
\expandafter\def\csname sga@target@subsection:IV.17.14\endcsname{ega4.c04.subsec.17-14}
\expandafter\def\csname sga@target@subsection:IV.17.15\endcsname{ega4.c04.subsec.17-15}
\expandafter\def\csname sga@target@subsection:IV.17.16\endcsname{ega4.c04.subsec.17-16}
\expandafter\def\csname sga@target@subsection:IV.17.2\endcsname{ega4.c04.subsec.17-2}
\expandafter\def\csname sga@target@subsection:IV.17.3\endcsname{ega4.c04.subsec.17-3}
\expandafter\def\csname sga@target@subsection:IV.17.4\endcsname{ega4.c04.subsec.17-4}
\expandafter\def\csname sga@target@subsection:IV.17.5\endcsname{ega4.c04.subsec.17-5}
\expandafter\def\csname sga@target@subsection:IV.17.6\endcsname{ega4.c04.subsec.17-6}
\expandafter\def\csname sga@target@subsection:IV.17.7\endcsname{ega4.c04.subsec.17-7}
\expandafter\def\csname sga@target@subsection:IV.17.8\endcsname{ega4.c04.subsec.17-8}
\expandafter\def\csname sga@target@subsection:IV.17.9\endcsname{ega4.c04.subsec.17-9}
\expandafter\def\csname sga@target@subsection:IV.18.1\endcsname{ega4.c04.subsec.18-1}
\expandafter\def\csname sga@target@subsection:IV.18.10\endcsname{ega4.c04.subsec.18-10}
\expandafter\def\csname sga@target@subsection:IV.18.11\endcsname{ega4.c04.subsec.18-11}
\expandafter\def\csname sga@target@subsection:IV.18.12\endcsname{ega4.c04.subsec.18-12}
\expandafter\def\csname sga@target@subsection:IV.18.2\endcsname{ega4.c04.subsec.18-2}
\expandafter\def\csname sga@target@subsection:IV.18.3\endcsname{ega4.c04.subsec.18-3}
\expandafter\def\csname sga@target@subsection:IV.18.4\endcsname{ega4.c04.subsec.18-4}
\expandafter\def\csname sga@target@subsection:IV.18.5\endcsname{ega4.c04.subsec.18-5}
\expandafter\def\csname sga@target@subsection:IV.18.6\endcsname{ega4.c04.subsec.18-6}
\expandafter\def\csname sga@target@subsection:IV.18.7\endcsname{ega4.c04.subsec.18-7}
\expandafter\def\csname sga@target@subsection:IV.18.8\endcsname{ega4.c04.subsec.18-8}
\expandafter\def\csname sga@target@subsection:IV.18.9\endcsname{ega4.c04.subsec.18-9}
\expandafter\def\csname sga@target@subsection:IV.19.1\endcsname{ega4.c04.subsec.19-1}
\expandafter\def\csname sga@target@subsection:IV.19.2\endcsname{ega4.c04.subsec.19-2}
\expandafter\def\csname sga@target@subsection:IV.19.3\endcsname{ega4.c04.subsec.19-3}
\expandafter\def\csname sga@target@subsection:IV.19.4\endcsname{ega4.c04.subsec.19-4}
\expandafter\def\csname sga@target@subsection:IV.19.5\endcsname{ega4.c04.subsec.19-5}
\expandafter\def\csname sga@target@subsection:IV.19.6\endcsname{ega4.c04.subsec.19-6}
\expandafter\def\csname sga@target@subsection:IV.19.7\endcsname{ega4.c04.subsec.19-7}
\expandafter\def\csname sga@target@subsection:IV.19.8\endcsname{ega4.c04.subsec.19-8}
\expandafter\def\csname sga@target@subsection:IV.19.9\endcsname{ega4.c04.subsec.19-9}
\expandafter\def\csname sga@target@subsection:IV.2.1\endcsname{ega4.c04.subsec.2-1}
\expandafter\def\csname sga@target@subsection:IV.2.2\endcsname{ega4.c04.subsec.2-2}
\expandafter\def\csname sga@target@subsection:IV.2.3\endcsname{ega4.c04.subsec.2-3}
\expandafter\def\csname sga@target@subsection:IV.2.4\endcsname{ega4.c04.subsec.2-4}
\expandafter\def\csname sga@target@subsection:IV.2.5\endcsname{ega4.c04.subsec.2-5}
\expandafter\def\csname sga@target@subsection:IV.2.6\endcsname{ega4.c04.subsec.2-6}
\expandafter\def\csname sga@target@subsection:IV.2.7\endcsname{ega4.c04.subsec.2-7}
\expandafter\def\csname sga@target@subsection:IV.2.8\endcsname{ega4.c04.subsec.2-8}
\expandafter\def\csname sga@target@subsection:IV.21.1\endcsname{ega4.c04.subsec.21-1}
\expandafter\def\csname sga@target@subsection:IV.21.10\endcsname{ega4.c04.subsec.21-10}
\expandafter\def\csname sga@target@subsection:IV.21.11\endcsname{ega4.c04.subsec.21-11}
\expandafter\def\csname sga@target@subsection:IV.21.12\endcsname{ega4.c04.subsec.21-12}
\expandafter\def\csname sga@target@subsection:IV.21.13\endcsname{ega4.c04.subsec.21-13}
\expandafter\def\csname sga@target@subsection:IV.21.14\endcsname{ega4.c04.subsec.21-14}
\expandafter\def\csname sga@target@subsection:IV.21.15\endcsname{ega4.c04.subsec.21-15}
\expandafter\def\csname sga@target@subsection:IV.21.2\endcsname{ega4.c04.subsec.21-2}
\expandafter\def\csname sga@target@subsection:IV.21.3\endcsname{ega4.c04.subsec.21-3}
\expandafter\def\csname sga@target@subsection:IV.21.4\endcsname{ega4.c04.subsec.21-4}
\expandafter\def\csname sga@target@subsection:IV.21.5\endcsname{ega4.c04.subsec.21-5}
\expandafter\def\csname sga@target@subsection:IV.21.6\endcsname{ega4.c04.subsec.21-6}
\expandafter\def\csname sga@target@subsection:IV.21.7\endcsname{ega4.c04.subsec.21-7}
\expandafter\def\csname sga@target@subsection:IV.21.8\endcsname{ega4.c04.subsec.21-8}
\expandafter\def\csname sga@target@subsection:IV.21.9\endcsname{ega4.c04.subsec.21-9}
\expandafter\def\csname sga@target@subsection:IV.3.1\endcsname{ega4.c04.subsec.3-1}
\expandafter\def\csname sga@target@subsection:IV.3.2\endcsname{ega4.c04.subsec.3-2}
\expandafter\def\csname sga@target@subsection:IV.3.3\endcsname{ega4.c04.subsec.3-3}
\expandafter\def\csname sga@target@subsection:IV.3.4\endcsname{ega4.c04.subsec.3-4}
\expandafter\def\csname sga@target@subsection:IV.4.1\endcsname{ega4.c04.subsec.4-1}
\expandafter\def\csname sga@target@subsection:IV.4.2\endcsname{ega4.c04.subsec.4-2}
\expandafter\def\csname sga@target@subsection:IV.4.3\endcsname{ega4.c04.subsec.4-3}
\expandafter\def\csname sga@target@subsection:IV.4.4\endcsname{ega4.c04.subsec.4-4}
\expandafter\def\csname sga@target@subsection:IV.4.5\endcsname{ega4.c04.subsec.4-5}
\expandafter\def\csname sga@target@subsection:IV.4.6\endcsname{ega4.c04.subsec.4-6}
\expandafter\def\csname sga@target@subsection:IV.4.7\endcsname{ega4.c04.subsec.4-7}
\expandafter\def\csname sga@target@subsection:IV.4.8\endcsname{ega4.c04.subsec.4-8}
\expandafter\def\csname sga@target@subsection:IV.4.9\endcsname{ega4.c04.subsec.4-9}
\expandafter\def\csname sga@target@subsection:IV.5.1\endcsname{ega4.c04.subsec.5-1}
\expandafter\def\csname sga@target@subsection:IV.5.10\endcsname{ega4.c04.subsec.5-10}
\expandafter\def\csname sga@target@subsection:IV.5.11\endcsname{ega4.c04.subsec.5-11}
\expandafter\def\csname sga@target@subsection:IV.5.12\endcsname{ega4.c04.subsec.5-12}
\expandafter\def\csname sga@target@subsection:IV.5.13\endcsname{ega4.c04.subsec.5-13}
\expandafter\def\csname sga@target@subsection:IV.5.2\endcsname{ega4.c04.subsec.5-2}
\expandafter\def\csname sga@target@subsection:IV.5.3\endcsname{ega4.c04.subsec.5-3}
\expandafter\def\csname sga@target@subsection:IV.5.4\endcsname{ega4.c04.subsec.5-4}
\expandafter\def\csname sga@target@subsection:IV.5.5\endcsname{ega4.c04.subsec.5-5}
\expandafter\def\csname sga@target@subsection:IV.5.6\endcsname{ega4.c04.subsec.5-6}
\expandafter\def\csname sga@target@subsection:IV.5.7\endcsname{ega4.c04.subsec.5-7}
\expandafter\def\csname sga@target@subsection:IV.5.8\endcsname{ega4.c04.subsec.5-8}
\expandafter\def\csname sga@target@subsection:IV.5.9\endcsname{ega4.c04.subsec.5-9}
\expandafter\def\csname sga@target@subsection:IV.6.1\endcsname{ega4.c04.subsec.6-1}
\expandafter\def\csname sga@target@subsection:IV.6.10\endcsname{ega4.c04.subsec.6-10}
\expandafter\def\csname sga@target@subsection:IV.6.11\endcsname{ega4.c04.subsec.6-11}
\expandafter\def\csname sga@target@subsection:IV.6.12\endcsname{ega4.c04.subsec.6-12}
\expandafter\def\csname sga@target@subsection:IV.6.13\endcsname{ega4.c04.subsec.6-13}
\expandafter\def\csname sga@target@subsection:IV.6.14\endcsname{ega4.c04.subsec.6-14}
\expandafter\def\csname sga@target@subsection:IV.6.15\endcsname{ega4.c04.subsec.6-15}
\expandafter\def\csname sga@target@subsection:IV.6.2\endcsname{ega4.c04.subsec.6-2}
\expandafter\def\csname sga@target@subsection:IV.6.3\endcsname{ega4.c04.subsec.6-3}
\expandafter\def\csname sga@target@subsection:IV.6.4\endcsname{ega4.c04.subsec.6-4}
\expandafter\def\csname sga@target@subsection:IV.6.5\endcsname{ega4.c04.subsec.6-5}
\expandafter\def\csname sga@target@subsection:IV.6.6\endcsname{ega4.c04.subsec.6-6}
\expandafter\def\csname sga@target@subsection:IV.6.7\endcsname{ega4.c04.subsec.6-7}
\expandafter\def\csname sga@target@subsection:IV.6.8\endcsname{ega4.c04.subsec.6-8}
\expandafter\def\csname sga@target@subsection:IV.6.9\endcsname{ega4.c04.subsec.6-9}
\expandafter\def\csname sga@target@subsection:IV.7.1\endcsname{ega4.c04.subsec.7-1}
\expandafter\def\csname sga@target@subsection:IV.7.2\endcsname{ega4.c04.subsec.7-2}
\expandafter\def\csname sga@target@subsection:IV.7.3\endcsname{ega4.c04.subsec.7-3}
\expandafter\def\csname sga@target@subsection:IV.7.4\endcsname{ega4.c04.subsec.7-4}
\expandafter\def\csname sga@target@subsection:IV.7.5\endcsname{ega4.c04.subsec.7-5}
\expandafter\def\csname sga@target@subsection:IV.7.6\endcsname{ega4.c04.subsec.7-6}
\expandafter\def\csname sga@target@subsection:IV.7.7\endcsname{ega4.c04.subsec.7-7}
\expandafter\def\csname sga@target@subsection:IV.7.8\endcsname{ega4.c04.subsec.7-8}
\expandafter\def\csname sga@target@subsection:IV.7.9\endcsname{ega4.c04.subsec.7-9}
\expandafter\def\csname sga@target@subsection:IV.8.1\endcsname{ega4.c04.subsec.8-1}
\expandafter\def\csname sga@target@subsection:IV.8.10\endcsname{ega4.c04.subsec.8-10}
\expandafter\def\csname sga@target@subsection:IV.8.11\endcsname{ega4.c04.subsec.8-11}
\expandafter\def\csname sga@target@subsection:IV.8.12\endcsname{ega4.c04.subsec.8-12}
\expandafter\def\csname sga@target@subsection:IV.8.13\endcsname{ega4.c04.subsec.8-13}
\expandafter\def\csname sga@target@subsection:IV.8.14\endcsname{ega4.c04.subsec.8-14}
\expandafter\def\csname sga@target@subsection:IV.8.2\endcsname{ega4.c04.subsec.8-2}
\expandafter\def\csname sga@target@subsection:IV.8.3\endcsname{ega4.c04.subsec.8-3}
\expandafter\def\csname sga@target@subsection:IV.8.4\endcsname{ega4.c04.subsec.8-4}
\expandafter\def\csname sga@target@subsection:IV.8.5\endcsname{ega4.c04.subsec.8-5}
\expandafter\def\csname sga@target@subsection:IV.8.6\endcsname{ega4.c04.subsec.8-6}
\expandafter\def\csname sga@target@subsection:IV.8.7\endcsname{ega4.c04.subsec.8-7}
\expandafter\def\csname sga@target@subsection:IV.8.8\endcsname{ega4.c04.subsec.8-8}
\expandafter\def\csname sga@target@subsection:IV.8.9\endcsname{ega4.c04.subsec.8-9}
\expandafter\def\csname sga@target@subsection:IV.9.1\endcsname{ega4.c04.subsec.9-1}
\expandafter\def\csname sga@target@subsection:IV.9.2\endcsname{ega4.c04.subsec.9-2}
\expandafter\def\csname sga@target@subsection:IV.9.3\endcsname{ega4.c04.subsec.9-3}
\expandafter\def\csname sga@target@subsection:IV.9.4\endcsname{ega4.c04.subsec.9-4}
\expandafter\def\csname sga@target@subsection:IV.9.5\endcsname{ega4.c04.subsec.9-5}
\expandafter\def\csname sga@target@subsection:IV.9.6\endcsname{ega4.c04.subsec.9-6}
\expandafter\def\csname sga@target@subsection:IV.9.7\endcsname{ega4.c04.subsec.9-7}
\expandafter\def\csname sga@target@subsection:IV.9.8\endcsname{ega4.c04.subsec.9-8}
\expandafter\def\csname sga@target@subsection:IV.9.9\endcsname{ega4.c04.subsec.9-9}
\let\sgaOriginalRef\ref
\let\sgaOriginalPageref\pageref
\renewcommand{\ref}[1]{%
  \ifcsname sga@target@#1\endcsname
    \hyperlink{\csname sga@target@#1\endcsname}{\sgaOriginalRef*{#1}}%
  \else
    \sgaOriginalRef{#1}%
  \fi
}
\renewcommand{\pageref}[1]{%
  \ifcsname sga@target@#1\endcsname
    \hyperlink{\csname sga@target@#1\endcsname}{\sgaOriginalPageref*{#1}}%
  \else
    \sgaOriginalPageref{#1}%
  \fi
}
\renewcommand{\eqref}[1]{%
  \ifcsname sga@target@#1\endcsname
    \hyperlink{\csname sga@target@#1\endcsname}{\textup{\tagform@{\sgaOriginalRef*{#1}}}}%
  \else
    \textup{\tagform@{\sgaOriginalRef{#1}}}%
  \fi
}
% Stable semantic routing for the reader's prose reference helpers.
\newcommand{\sgaStableOrOriginalLink}[2]{%
  \ifcsname sga@target@#1\endcsname
    \hyperlink{\csname sga@target@#1\endcsname}{#2}%
  \else
    \@ifundefined{r@#1}{#2}{\hyperref[#1]{#2}}%
  \fi
}
\RenewDocumentCommand{\sref}{o m o}{%
  \IfNoValueTF{#1}%
    {\IfNoValueTF{#3}%
      {\sgaStableOrOriginalLink{#2}{\normalfont{(\StrBehind{#2}{.})}}}%
      {\sgaStableOrOriginalLink{#2}{\normalfont{(\StrBehind{#2}{.},~{#3})}}}}%
    {\IfNoValueTF{#3}%
      {\sgaStableOrOriginalLink{#2}{\normalfont{(\textbf{#1},~\StrBehind{#2}{.})}}}%
      {\sgaStableOrOriginalLink{#2}{\normalfont{(\textbf{#1},~\StrBehind{#2}{.},~{#3})}}}}%
}
\renewcommand{\eref}[1]{%
  \sgaStableOrOriginalLink{#1}{\normalfont{(\StrBehind{#1}{.})}}}
\makeatother
% END NATIVE INPUT EGA4_reference_v2_links
\setlength{\emergencystretch}{3em}

\hypersetup{
  pdftitle={Elements of Algebraic Geometry IV, Part Two -- English Edition},
  pdfauthor={Alexandre Grothendieck and Jean Dieudonne; English translation},
  pdfsubject={Standalone English reader of EGA IV, Part Two}
}

\begin{document}

\title{Elements of Algebraic Geometry IV\\Part Two}
\author{Alexandre Grothendieck and Jean Dieudonn\'e}
\date{}
\maketitle

% BEGIN SCHOLARLY TRANSLATION ATTRIBUTION

\begingroup
\small
\begin{quote}
\textbf{Translation sources and contributors.}
The original work is by Alexandre Grothendieck and Jean Dieudonn\'e.
The EGA English edition series draws on the human community translation
\texttt{github.com/ryankeleti/ega}, source version \texttt{3d7afe28}.
Its named contributors are Mateo Carmona, Tim Holzschuh, Tim Hosgood,
Ryan Keleti, Thiago Solovera \& Nery, and Wayne Yang.
This volume belongs to the wider EGA English edition series; the community translation's contribution to that series is described in the accompanying attribution notice.
The additional translation is AI-translated; subsequent corrections, editing,
typesetting and assembly are AI-produced. The publisher did not translate
the text. Recorded AI work includes OpenAI Codex,
GPT-6 Astra, Ultra.
The accompanying \texttt{ATTRIBUTION.md} gives source links, contribution
scope and rights information. Attribution does not assert endorsement
or a licence grant from the upstream contributors.

\medskip
\textbf{Rights and availability.}
Rights in the original works and in any protected contributions to these
derivative translations and transcriptions remain with their respective
rights holders. This project does not dedicate those works or contributions
to the public domain or grant a licence over them. Existing third-party
notices and any permissions already granted remain applicable.

These editions are available free of charge for non-commercial, purely
educational purposes, on the presumption that the respective rights holders
would not object to these translations and transcriptions being freely
available alongside the French originals on NUMDAM
(\texttt{www.numdam.org/journals/PMIHES/}). This presumption is not a statement
that permission has been obtained; free access to the originals does not
itself grant permission to translate, adapt or republish them.

Rights holders who object to the availability of their material may open
an issue at \texttt{github.com/KokunoYumeto/ega-en/issues}. The publisher
will take down the affected material at the rights holder's request.
\end{quote}
\endgroup
% END SCHOLARLY TRANSLATION ATTRIBUTION


\begin{center}
\large Standalone English reader
\end{center}

\tableofcontents
% BEGIN NATIVE INPUT EGA_IV2_standalone_reference_fallbacks_r52
% Standalone-only routing for references whose targets are outside EGA IV-2.
% Mathematical source files remain unchanged. All fallback destinations land
% on this notice; the cumulative reader retains statement-level routing.
\clearpage
\phantomsection
\section*{Cross-volume references}

This part-only reader preserves the citation text and link affordances of the
maintained edition. References whose source statements are not reproduced in
Part Two resolve to this notice. Statement-level cross-volume navigation is
available in the cumulative EGA reader.

% One inherited raw hyperref target is not represented in the EGA IV stable
% registry; provide its non-linking standalone destination explicitly.
\phantomsection
\label{0.16.4.6.i}

\makeatletter
\newcommand{\EGAIVTwoRegisterExternalTargets}[1]{%
  \@for\EGAIVTwoExternalTarget:=#1\do{%
    \edef\EGAIVTwoResolvedTarget{\EGAIVTwoExternalTarget}%
    \expandafter\sgaStableTarget\expandafter{\EGAIVTwoResolvedTarget}%
  }%
}
\makeatother

\EGAIVTwoRegisterExternalTargets{%
ega4.c04.cor.1-1-4,ega4.c04.cor.1-10-4,ega4.c04.cor.1-8-5,ega4.c04.cor.1-9-12,ega4.c04.cor.11-3-11,ega4.c04.cor.8-4-5,%
ega4.c04.def.1-4-2,ega4.c04.def.1-5-2,ega4.c04.def.1-6-1,ega4.c04.lem.13-1-1,%
ega4.c04.prop.1-1-2,ega4.c04.prop.1-1-5,ega4.c04.prop.1-1-6,ega4.c04.prop.1-10-3,ega4.c04.prop.1-2-2,%
ega4.c04.prop.1-2-5,ega4.c04.prop.1-3-11,ega4.c04.prop.1-3-4,ega4.c04.prop.1-3-6,ega4.c04.prop.1-3-9,%
ega4.c04.prop.1-4-3,ega4.c04.prop.1-4-4,ega4.c04.prop.1-4-6,ega4.c04.prop.1-5-4,ega4.c04.prop.10-6-1,%
ega4.c04.prop.1-9-5,ega4.c04.def.1-9-4,ega4.c04.statement.1-7-21,%
ega4.c04.prop.14-3-2,ega4.c04.prop.17-6-3,ega4.c04.prop.8-9-3,ega4.c04.rem.12-1-4,ega4.c04.rem.18-8-11,%
ega4.c04.statement.1-3-10,ega4.c04.statement.1-3-8,ega4.c04.statement.1-7-4,ega4.c04.statement.1-7-8,%
ega4.c04.statement.18-9-6,ega4.c04.subsec.17-3-2,ega4.c04.thm.11-3-10,ega4.c04.thm.17-5-1%
}
% END NATIVE INPUT EGA_IV2_standalone_reference_fallbacks_r52

% EGA IV, Part Two contains sections 2--7 of Chapter IV.
\setcounter{section}{1}
% BEGIN NATIVE INPUT ega4-2.source_aligned_r1
\section{Base change and flatness}
\label{section:IV.2}

\oldpage[IV-2]{5}

This section (unlike \textsection~6) only exceptionally makes use of Noetherian techniques.
The n\textsuperscript{os}~1 and~2 are essentially nothing more than translations of elementary properties of flatness in commutative Algebra (cf.~Bourbaki, \emph{Alg.~comm.}, chap.~I\textsuperscript{er}), and are included here for the convenience of references.
The following numbers are primarily devoted to properties of ``descent'' by flat or faithfully flat morphisms: if $g:Y'\to Y$ is such a morphism, we want to be able to assert that a subset of~$Y$, or an $\sh{O}_Y$-Module, or a morphism $X\to Y$, has a certain property, when we know that its \emph{inverse image} by $g$ possesses this property.
We limit ourselves here to properties that do not make use of the general technique of ``descent'', which will be developed in Chapter~V.

\subsection{Flat Modules on preschemes}
\label{subsection:IV.2.1}

\begin{env}[2.1.1]
\label{IV.2.1.1}
Let $f:X\to Y$ be a morphism of preschemes, $\sh{F}$ an $\sh{O}_X$-Module;
recall~\sref[0\textsubscript{I}]{0.6.7.1} that $\sh{F}$ is said to be \emph{$f$-flat} (or \emph{$Y$-flat}) \emph{at a point} $x\in X$ if $\sh{F}_x$ is a flat $\sh{O}_{f(x)}$-module, \emph{$f$-flat} (or \emph{$Y$-flat}) if it is $f$-flat at every point $x\in X$, and finally that the morphism $f$ is said to be \emph{flat at the point} $x\in X$ (resp.\ \emph{flat}) if $\sh{O}_X$ is $f$-flat at the point $x$ (resp.\ $f$-flat).
When $f=1_X$, we will simply say that an $\sh{O}_X$-Module $\sh{F}$ is \emph{flat at the point} $x$ (resp.\ \emph{flat}) if it is $X$-flat at this point (resp.\ at every point $x\in X$), that is to say if $\sh{F}_x$ is a flat $\sh{O}_x$-module (resp.\ if this is the case for every $x\in X$).
Recall that we have shown~\sref[III]{III.1.4.15.1} the following property:
\end{env}

\begin{proposition}[2.1.2]
\label{IV.2.1.2}
Let $A$, $B$ be two rings, $\varphi:A\to B$ a ring homomorphism, $X=\Spec(B)$, $Y=\Spec(A)$, $f:X\to Y$ the corresponding morphism, and $M$ a $B$-module.
For $\sh{F}=\widetilde{M}$ to be $f$-flat, it is necessary and sufficient that $M$ be a flat $A$-module.
\end{proposition}

\begin{proposition}[2.1.3]
\label{IV.2.1.3}
Let $f:X\to Y$ be a morphism of preschemes, $\sh{F}$ a quasi-coherent $\sh{O}_X$-Module. The following conditions are equivalent:
\begin{enumerate}
  \item[\emph{a)}] For every base change $g:Y'\to Y$, if we set $X'=X\times_Y Y'$, the functor $\sh{G}'\mapsto\sh{F}\otimes_{\sh{O}_Y}\sh{G}'$ from the category of quasi-coherent $\sh{O}_{Y'}$-Modules to that of quasi-coherent $\sh{O}_{X'}$-Modules is exact.
  \item[\emph{a$'$)}] Condition~\emph{a)} is satisfied for all canonical morphisms $g:\Spec(\sh{O}_y)\to Y$~\sref[I]{I.2.4.1}, where $y$ ranges over $Y$.
  \item[\emph{b)}] $\sh{F}$ is $f$-flat.
\end{enumerate}
\end{proposition}

\begin{proof}
The questions being local on $X$ and $Y$, we can restrict to the case where $Y=\Spec(B)$, $X=\Spec(A)$, $\sh{F}=\widetilde{M}$, where $M$ is an $A$-module.
It is clear that \emph{a)} implies \emph{a$'$)};
condition~\emph{a$'$)} implies that for every $x\in X$, the functor $N\mapsto M_\mathfrak{n}\otimes_{B_\mathfrak{n}}N$ is exact in the category of $B_\mathfrak{n}$-modules, $\mathfrak{n}$ being the ideal $\mathfrak{j}_{f(x)}$ of $B$; this means that $M_\mathfrak{n}$ is a flat $B_\mathfrak{n}$-module, and it follows from~\sref[0\textsubscript{I}]{0.6.3.3} and from~\sref{IV.2.1.2} that $\sh{F}$ is $f$-flat.
Finally, to see that \emph{b)} implies \emph{a)}, we can also restrict to the case where $Y'=\Spec(A')$ is affine and $\sh{G}'=\widetilde{N'}$, where $N'$ is an $A'$-module; the conclusion then follows from~\sref{IV.2.1.2} and from the definition of flatness, since $(M\otimes_B N')^{\sim}=\sh{F}\otimes_{\sh{O}_Y}\sh{G}'$.
\end{proof}

\begin{proposition}[2.1.4]
\label{IV.2.1.4}
\oldpage[IV-2]{6}
Let $f:X\to Y$, $g:Y'\to Y$ be two morphisms of preschemes, $\sh{F}$ a quasi-coherent $\sh{O}_X$-Module; set $X'=X\times_Y Y'$, $f'=f_{(Y')}:X'\to Y'$, $\sh{F}'=\sh{F}\otimes_{\sh{O}_Y}\sh{O}_{Y'}$, and let $g':X'\to X$ be the canonical projection.
Let $x'$ be a point of $X'$, and set $x=g'(x')$, $y'=f'(x')$, and $y=g(y')=f(x)$.
If $\sh{F}$ is $f$-flat at the point $x$, $\sh{F}'$ is $f'$-flat at the point $x'$; in particular, if $\sh{F}$ is $f$-flat, $\sh{F}'$ is $f'$-flat; if $f$ is flat, $f'$ is flat.
\end{proposition}

\begin{proof}
It suffices to prove the first assertion; applying three times~\sref[I]{I.3.6.5}, as well as~\sref[I]{I.2.4.4}, we can reduce to the case where $Y=\Spec(\sh{O}_y)$, $X=\Spec(\sh{O}_x)$, $Y'=\Spec(\sh{O}_{y'})$, $\sh{F}=\widetilde{M}$, where $M=\sh{F}_x$.
The hypothesis and~\sref{IV.2.1.2} then imply that $\sh{F}$ is $f$-flat; in other words, we are reduced to proving a special case of the second assertion, and the latter follows immediately from~\sref{IV.2.1.3}.
\end{proof}

\begin{proposition}[2.1.5]
\label{IV.2.1.5}
Consider a commutative diagram of morphisms of preschemes
\[
\xymatrix{
  X \ar[d]_f & X' \ar[l]_{g'} \ar[d]^{f'} \\
  Y & Y' \ar[l]^g \ar[d]^h \\
  & Z
}
\]
where $X'=X\times_Y Y'$ and $f'=f_{(Y')}$.
Let $x'$ be a point of $X'$, and set $x=g'(x')$, $y'=f'(x')$, $y=f(x)=g(y')$, $z=h(y')$.
Let $\sh{F}$ be a quasi-coherent $\sh{O}_X$-Module $f$-flat at the point $x$ \textnormal{(}resp.\ $f$-flat\textnormal{)}, and $\sh{G}'$ a quasi-coherent $\sh{O}_{Y'}$-Module $h$-flat at the point $y'$ \textnormal{(}resp.\ $h$-flat\textnormal{)}; then $\sh{F}\otimes_{\sh{O}_{Y'}}\sh{G}'$ is a quasi-coherent $\sh{O}_{X'}$-Module $(h\circ f')$-flat at the point $x'$ \textnormal{(}resp.\ $(h\circ f')$-flat\textnormal{)}.
\end{proposition}

\begin{proof}
As in~\sref{IV.2.1.4}, we reduce to the case where $X=\Spec(\sh{O}_x)$, $Y=\Spec(\sh{O}_y)$,
\oldpage[IV-2]{7}
$Y'=\Spec(\sh{O}_{y'})$ and $Z=\Spec(\sh{O}_z)$, and it suffices then to prove that $\sh{F}\otimes_{\sh{O}_{Y'}}\sh{G}'$ is $(h\circ f')$-flat.
Taking into account~\sref{IV.2.1.2}, the proposition follows from Bourbaki, \emph{Alg.~comm.}, chap.~I\textsuperscript{er}, \textsection~2, n\textsuperscript{o}~7, prop.~8.
\end{proof}

\begin{corollary}[2.1.6]
\label{IV.2.1.6}
Let $f:X\to Y$, $g:Y\to Z$ be two morphisms of preschemes, $\sh{F}$ an $\sh{O}_X$-Module. If $\sh{F}$ is $f$-flat at the point $x\in X$ and if $g$ is flat at the point $f(x)$, $\sh{F}$ is $(g\circ f)$-flat at the point $x$.
In particular, if $f$ and $g$ are flat morphisms, so is $g\circ f$.
\end{corollary}

\begin{proof}
This is the special case of~\sref{IV.2.1.5} with $Y'=Y$, $\sh{G}'=\sh{O}_{Y'}$.
\end{proof}

\begin{corollary}[2.1.7]
\label{IV.2.1.7}
If $f:X\to X'$, $g:Y\to Y'$ are two flat $S$-morphisms, $f\times_S g:X\times_S Y\to X'\times_S Y'$ is a flat morphism.
\end{corollary}

\begin{proof}
This follows from~\sref{IV.2.1.4} and~\sref{IV.2.1.6} (cf.~\sref[I]{I.3.5.1}).
\end{proof}

\begin{proposition}[2.1.8]
\label{IV.2.1.8}
Let $f:X\to Y$ be a morphism of preschemes,
\[
  0 \to \sh{F}' \to \sh{F} \to \sh{F}'' \to 0
\]
an exact sequence of quasi-coherent $\sh{O}_X$-Modules, such that $\sh{F}''$ is $Y$-flat.
\begin{enumerate}[(i)]
  \item For every morphism $g:Y'\to Y$ and every quasi-coherent $\sh{O}_{Y'}$-Module $\sh{G}'$, the sequence
    $0\to\sh{F}'\otimes_{\sh{O}_Y}\sh{G}'\to\sh{F}\otimes_{\sh{O}_Y}\sh{G}'\to\sh{F}''\otimes_{\sh{O}_Y}\sh{G}'\to 0$ of $\sh{O}_{X'}$-Modules \textnormal{(}where $X'=X\times_Y Y'$\textnormal{)} is exact.
  \item For $\sh{F}$ to be $Y$-flat, it is necessary and sufficient that $\sh{F}'$ be so.
\end{enumerate}
\end{proposition}

\begin{proof}
We can evidently suppose $X$, $Y$, $Y'$ affine; the conclusion then follows from~\sref{IV.2.1.2} and from~\sref[0\textsubscript{I}]{0.6.1.2}.
\end{proof}

\begin{corollary}[2.1.9]
\label{IV.2.1.9}
Let $\sh{L}^\bullet$ be a complex of quasi-coherent $\sh{O}_X$-Modules, $i$ an index such that if $d^i:\sh{L}^i\to\sh{L}^{i+1}$ is the derivation operator, $\sh{B}^{i+1}(\sh{L}^\bullet)=\operatorname{Im}(d^i)$ and $\sh{Z}^{i+1}(\sh{L}^\bullet)=\operatorname{Coker}(d^i)$ are $Y$-flat.
Then, with the notations of~\sref{IV.2.1.8}, the canonical homomorphism
\[
  \sh{H}^i(\sh{L}^\bullet)\otimes_{\sh{O}_Y}\sh{G}' \to \sh{H}^i(\sh{L}^\bullet\otimes_{\sh{O}_Y}\sh{G}')
\]
is bijective.
\end{corollary}

\begin{proof}
As the tensor product is right exact, we have
\[
  \sh{Z}^{i+1}(\sh{L}^\bullet)\otimes_{\sh{O}_Y}\sh{G}' = \operatorname{Coker}(d^i\otimes 1) = \sh{Z}^{i+1}(\sh{L}^\bullet\otimes_{\sh{O}_Y}\sh{G}')
\]
and $\sh{Z}^{i}(\sh{L}^\bullet)\otimes_{\sh{O}_Y}\sh{G}' = \sh{Z}^{i}(\sh{L}^\bullet\otimes_{\sh{O}_Y}\sh{G}')$.
Moreover, in the exact sequence
\[
  0\to\sh{B}^{i+1}(\sh{L}^\bullet)\to\sh{L}^{i+1}\to\sh{Z}^{i+1}(\sh{L}^\bullet)\to 0
\]
$\sh{Z}^{i+1}(\sh{L}^\bullet)$ is $Y$-flat, so it follows from~\sref{IV.2.1.8}[(i)] that the exact sequence
\[
  0\to\sh{B}^{i+1}(\sh{L}^\bullet)\otimes_{\sh{O}_Y}\sh{G}'\to\sh{L}^{i+1}\otimes_{\sh{O}_Y}\sh{G}'\to\sh{Z}^{i+1}(\sh{L}^\bullet\otimes_{\sh{O}_Y}\sh{G}')\to 0
\]
whence $\sh{B}^{i+1}(\sh{L}^\bullet)\otimes_{\sh{O}_Y}\sh{G}'=\operatorname{Im}(d^i\otimes 1)=\sh{B}^{i+1}(\sh{L}^\bullet\otimes_{\sh{O}_Y}\sh{G}')$.
Then, as in the exact sequence
\[
  0\to\sh{H}^i(\sh{L}^\bullet)\to\sh{Z}^i(\sh{L}^\bullet)\to\sh{B}^{i+1}(\sh{L}^\bullet)\to 0
\]
$\sh{B}^{i+1}(\sh{L}^\bullet)$ is $Y$-flat, it follows from~\sref{IV.2.1.8}[(i)] and from what precedes that the sequence
\[
  0\to\sh{H}^i(\sh{L}^\bullet)\otimes_{\sh{O}_Y}\sh{G}'\to\sh{Z}^i(\sh{L}^\bullet\otimes_{\sh{O}_Y}\sh{G}')\to\sh{B}^{i+1}(\sh{L}^\bullet\otimes_{\sh{O}_Y}\sh{G}')\to 0
\]
is exact, which proves the corollary.
\end{proof}

\begin{corollary}[2.1.10]
\label{IV.2.1.10}
\oldpage[IV-2]{8}
Let $f:X\to Y$ be a morphism of preschemes, $\sh{F}$ a quasi-coherent $\sh{O}_X$-Module which is $Y$-flat, $\sh{L}_\bullet=(\sh{L}_i)$ a left resolution of $\sh{F}$ formed of quasi-coherent $\sh{O}_X$-Modules which are $Y$-flat.
Then, for every morphism $g:Y'\to Y$ and every quasi-coherent $\sh{O}_{Y'}$-Module $\sh{G}'$, the complex $\sh{L}_\bullet\otimes_{\sh{O}_Y}\sh{G}'=(\sh{L}_i\otimes_{\sh{O}_Y}\sh{G}')_{i\geq 0}$ is a left resolution of $\sh{F}\otimes_{\sh{O}_Y}\sh{G}'$.
Moreover, if $\sh{Z}_i(\sh{L}_\bullet)=\operatorname{Ker}(\sh{L}_i\to\sh{L}_{i-1})$, the $\sh{Z}_i(\sh{L}_\bullet)$ are $Y$-flat, and we have $\sh{Z}_i(\sh{L}_\bullet)\otimes_{\sh{O}_Y}\sh{G}'=\sh{Z}_i(\sh{L}_\bullet\otimes_{\sh{O}_Y}\sh{G}')=\operatorname{Ker}(\sh{L}_i\otimes_{\sh{O}_Y}\sh{G}'\to\sh{L}_{i-1}\otimes_{\sh{O}_Y}\sh{G}')$.
Setting $\sh{R}_i=\operatorname{Im}(\sh{L}_{i+1}\to\sh{L}_i)=\sh{Z}_i(\sh{L}_\bullet)$; we then have the exact sequences
\[
  0\leftarrow\sh{F}\leftarrow\sh{L}_0\leftarrow\sh{R}_0\leftarrow 0
\]
\[
  0\leftarrow\sh{R}_i\leftarrow\sh{L}_{i+1}\leftarrow\sh{R}_{i+1}\leftarrow 0 \qquad (i\geq 0)
\]
and since $\sh{F}$ and $\sh{L}_i$ are $Y$-flat, we deduce from~\sref{IV.2.1.8}[(ii)] by induction that all the $\sh{R}_i$ are also $Y$-flat; using~\sref{IV.2.1.8}[(i)], we thus have the exact sequences
\[
  0\leftarrow\sh{F}\otimes_{\sh{O}_Y}\sh{G}'\leftarrow\sh{L}_0\otimes_{\sh{O}_Y}\sh{G}'\leftarrow\sh{R}_0\otimes_{\sh{O}_Y}\sh{G}'\leftarrow 0
\]
\[
  0\leftarrow\sh{R}_i\otimes_{\sh{O}_Y}\sh{G}'\leftarrow\sh{L}_{i+1}\otimes_{\sh{O}_Y}\sh{G}'\leftarrow\sh{R}_{i+1}\otimes_{\sh{O}_Y}\sh{G}'\leftarrow 0 \qquad (i\geq 0)
\]
which prove the corollary.
\end{corollary}

\begin{proposition}[2.1.11]
\label{IV.2.1.11}
Let $f:X\to Y$ be a flat morphism, $\sh{F}$ a quasi-coherent $\sh{O}_Y$-Module of finite presentation.
If $\sh{J}$ is the Ideal of $\sh{O}_Y$ annihilator of $\sh{F}$, $f^*(\sh{J})$ is the Ideal of $\sh{O}_X$ annihilator of $f^*(\sh{F})$.
\end{proposition}

\begin{proof}
We have indeed by definition an exact sequence~\sref[0\textsubscript{I}]{0.5.3.7}
\[
  0\to\sh{J}\to\sh{O}_Y\to\shHom_{\sh{O}_Y}(\sh{F},\,\sh{F})
\]
whence, since $f$ is flat, an exact sequence
\[
  0\to f^*(\sh{J})\to\sh{O}_X\to f^*(\shHom_{\sh{O}_Y}(\sh{F},\,\sh{F}))
\]
and since $\sh{F}$ is an $\sh{O}_Y$-Module of finite presentation, $f^*(\shHom_{\sh{O}_Y}(\sh{F},\,\sh{F}))$ canonically identifies with $\shHom_{\sh{O}_X}(f^*(\sh{F}),f^*(\sh{F}))$~\sref[0\textsubscript{I}]{0.6.7.6}, whence the conclusion.
\end{proof}

\begin{proposition}[2.1.12]
\label{IV.2.1.12}
Let $X$ be a prescheme, $\sh{F}$ an $\sh{O}_X$-Module of finite presentation, $x$ a point of $X$.
The following conditions are equivalent:
\begin{enumerate}
  \item[\emph{a)}] $\sh{F}_x$ is a flat $\sh{O}_x$-module.
  \item[\emph{b)}] There exists an open neighbourhood $U$ of $x$ such that $\sh{F}|U$ is a locally free $(\sh{O}_X|U)$-Module.
\end{enumerate}
\end{proposition}

\begin{proof}
Indeed, $\sh{F}_x$ is a module of finite presentation over $\sh{O}_x$, and $\sh{O}_x$ is a local ring; it amounts to the same thing to say that $\sh{F}_x$ is a flat $\sh{O}_x$-module or a free $\sh{O}_x$-module (Bourbaki, \emph{Alg.~comm.}, chap.~II, \textsection~3, n\textsuperscript{o}~2, cor.~2 of prop.~5); whence the conclusion, taking into account~\sref[0\textsubscript{I}]{0.5.2.7}.
We note that the proposition is valid for any \emph{ringed space with local rings}.
\end{proof}

\begin{proposition}[2.1.13]
\label{IV.2.1.13}
Let $f:X\to Y$ be a morphism of preschemes. If $f$ is flat at a point $x\in X$ and if the ring $\sh{O}_x$ is reduced \textnormal{(}resp.\ integral and integrally closed\textnormal{)}, then the ring $\sh{O}_{f(x)}$
\oldpage[IV-2]{9}
is reduced \textnormal{(}resp.\ integral and integrally closed\textnormal{)}.
If $f$ is faithfully flat and if $X$ is reduced \textnormal{(}resp.\ normal\textnormal{)}, then $Y$ is reduced \textnormal{(}resp.\ normal\textnormal{)}.
\end{proposition}

\begin{proof}
Set $\sh{O}_{f(x)}=A$, $\sh{O}_x=B$.
If $B$ is a flat $A$-module, it is also a faithfully flat $A$-module~\sref[0\textsubscript{I}]{0.6.6.2}, so $A$ identifies with a subring of $B$; if $B$ is reduced, so is $A$.
Suppose now that $B$ is integral and integrally closed, and let $L$ be its field of fractions; then $A\subset B$ is integral; denote by $K\subset L$ its field of fractions.
The hypothesis implies that $B\cap K=A$ (Bourbaki, \emph{Alg.~comm.}, chap.~I, \textsection~3, n\textsuperscript{o}~5, prop.~10).
If then $t\in K$ is integral over $A$, it is also integral over $B$, so belongs to $B$ by hypothesis, and thus $t\in A$, which proves that $A$ is integrally closed.
\end{proof}

\begin{proposition}[2.1.14]
\label{IV.2.1.14}
Let $f:X\to Y$ be a faithfully flat morphism. If $X$ is locally integral and if the space underlying $Y$ is locally Noetherian, then $Y$ is locally integral.
\end{proposition}

\begin{proof}
Indeed, every $y\in Y$ is of the form $f(x)$ for some $x\in X$ and by hypothesis $\sh{O}_y$ identifies with a subring of $\sh{O}_x$~\sref[0\textsubscript{I}]{0.6.6.1}; as $\sh{O}_x$ is integral, so is $\sh{O}_y$ and this proves the proposition~\sref[I]{I.5.1.4}.
\end{proof}

\subsection{Faithfully flat Modules on preschemes}
\label{subsection:IV.2.2}

\begin{proposition}[2.2.1]
\label{IV.2.2.1}
Let $f:X\to Y$ be a morphism of preschemes, $\sh{F}$ a quasi-coherent $\sh{O}_X$-Module. The following conditions are equivalent:
\begin{enumerate}
  \item[\emph{a)}] For every base change $g:Y'\to Y$, if we set $X'=X\times_Y Y'$, the functor $\sh{G}'\mapsto\sh{F}\otimes_{\sh{O}_Y}\sh{G}'$, from the category of quasi-coherent $\sh{O}_{Y'}$-Modules to that of quasi-coherent $\sh{O}_{X'}$-Modules, is exact and faithful.
  \item[\emph{a$'$)}] Condition~\emph{a)} is satisfied for all canonical morphisms $g:\Spec(\sh{O}_y)\to Y$~\sref[I]{I.2.4.1}, where $y$ ranges over $Y$.
  \item[\emph{a$''$)}] Condition~\emph{a)} is satisfied for all canonical immersions $Y'\to Y$, where $Y'$ ranges over the set of affine open sub-preschemes of $Y$.
  \item[\emph{b)}] $\sh{F}$ is $Y$-flat and, for every $y\in Y$, if we denote by $X_y$ the fibre $f^{-1}(y)$, the $\sh{O}_{X_y}$-Module $\sh{F}_y=\sh{F}\otimes_{\sh{O}_Y}k(y)$ is nonzero.
\end{enumerate}
\end{proposition}

\begin{proof}
It is clear that \emph{a)} implies \emph{a$'$)} and \emph{a$''$)}; condition~\emph{a$'$)} implies first that $\sh{F}$ is $Y$-flat~\sref{IV.2.1.3}; on the other hand it implies that for every $y\in Y$, the functor $N\mapsto\sh{F}\otimes_{\sh{O}_y}\widetilde{N}$ is faithful in the category of $\sh{O}_y$-modules; taking in particular $N=k(y)$, we obtain the second assertion of~\emph{b)}.
To show that \emph{b)} implies \emph{a)}, we can restrict to the case where $Y$ is affine, since the question is local over~$Y$.
Likewise, to prove that \emph{a$''$)} implies \emph{a)}, we are reduced to proving that, when $Y$ is affine, the fact that $\sh{G}\mapsto\sh{F}\otimes_{\sh{O}_Y}\sh{G}$ is an exact and faithful functor implies condition~\emph{a)}.
In other words, we are reduced to proving the following more precise proposition:
\end{proof}

\begin{proposition}[2.2.2]
\label{IV.2.2.2}
\oldpage[IV-2]{10}
Let $Y=\Spec(A)$ be an affine scheme, $f:X\to Y$ a morphism of preschemes, $\sh{F}$ a quasi-coherent $\sh{O}_X$-Module.
Condition~\emph{a)} of~\sref{IV.2.2.1} is equivalent to each of the following:
\begin{enumerate}
  \item[\emph{b$'$)}] $\sh{F}$ is $Y$-flat and, for every closed point $y$ of $Y$, we have $\sh{F}_y\neq 0$.
  \item[\emph{c)}] The functor $\sh{G}\mapsto\sh{F}\otimes_{\sh{O}_Y}\sh{G}$ from the category of quasi-coherent $\sh{O}_Y$-Modules to that of quasi-coherent $\sh{O}_X$-Modules is exact and faithful.
\end{enumerate}
\end{proposition}

\begin{proof}
If \emph{b$'$)} holds and $y$ is a closed point of $Y$, there is at least one $x\in f^{-1}(y)$ such that $(\sh{F}_y)_x\neq 0$; let $U=\Spec(B)$ be an affine open neighbourhood of $x$, and write $\sh{F}|_U=\widetilde{M}$, where $M$ is a $B$-module.
Then \emph{b$'$)} implies that $M/\mathfrak{j}_yM$ is nonzero and hence, since $M$ is a flat $A$-module~\sref{IV.2.1.2}, that $M\otimes_A A_y$ is a faithfully flat $A_y$-module~\sref[0\textsubscript{I}]{0.6.4.5}.
The relation $(\sh{F}\otimes_{\sh{O}_Y}\sh{G})\otimes_{\sh{O}_Y}\sh{O}_y=0$ for a closed point $y$ of $Y$ implies
\[
  (\sh{F}\otimes_{\sh{O}_Y}\sh{O}_y)\otimes_{\sh{O}_y}\sh{G}_y=0,
\]
and therefore $\sh{G}_y=0$.
But if $\sh{G}_y=0$ for every closed point $y$ of $Y$, we conclude that $\sh{G}=0$: indeed, if $\sh{G}=\widetilde{N}$, the annihilator of every element of $N$ is contained in no maximal ideal of $A$, and hence is all of~$A$.
The relation $\sh{F}\otimes_{\sh{O}_Y}\sh{G}=0$ implies therefore $\sh{G}=0$, in other words the functor $\sh{G}\mapsto\sh{F}\otimes_{\sh{O}_Y}\sh{G}$ is faithful; we know moreover that this functor is exact~\sref{IV.2.1.3}, which shows that \emph{b$'$)} implies \emph{c)}.

Finally, to see that \emph{c)} implies \emph{a)}, we can restrict to the case where $Y'$ is also affine.
The morphism $g:Y'\to Y$ is then affine, and so is the projection $g':X'\to X$~\sref[II]{II.1.5.5}.
Moreover, the functor $\sh{H}'\mapsto g'_*(\sh{H}')$ is exact in the category of quasi-coherent $\sh{O}_{X'}$-Modules~\sref[I]{I.1.6.4}; if $g'_*(\sh{H}')=0$, then $\sh{H}'=0$, so this functor is also faithful.
Consequently, to show that $\sh{G}'\mapsto\sh{F}\otimes_{\sh{O}_Y}\sh{G}'$ is exact and faithful, it suffices to show that the functor
\[
  \sh{G}'\longmapsto
  g'_*\bigl(g'^*(\sh{F})\otimes_{\sh{O}_{X'}}f'^*(\sh{G}')\bigr)
\]
is exact and faithful.
The fact that $g$ is affine implies that we have a canonical isomorphism

\begin{equation}
\tag{2.2.2.1}
\label{IV.2.2.2.1}
  \sh{F}\otimes_{\sh{O}_X}f^*(g_*(\sh{G}'))\xrightarrow{\sim} g'_*(g'^*(\sh{F})\otimes_{\sh{O}_{X'}}f'^*(\sh{G}')).
\end{equation}

Indeed, we know~\sref[II]{II.1.5.2} that we have a canonical isomorphism
\[
  f^*(g_*(\sh{G}'))\xrightarrow{\sim} g'_*(f'^*(\sh{G}')),
\]
and on the other hand a canonical homomorphism $\sh{F}\to g'_*(g'^*(\sh{F}))$~\sref[0\textsubscript{I}]{0.4.4.3.2}; composing the homomorphism $\sh{F}\otimes_{\sh{O}_X}f^*(g_*(\sh{G}'))\to g'_*(g'^*(\sh{F}))\otimes_{\sh{O}_X}g'_*(f'^*(\sh{G}'))$ with the canonical homomorphism~\sref[0\textsubscript{I}]{0.4.2.2.1}, we deduce the homomorphism~\sref{IV.2.2.2.1}; the verification that it is an isomorphism is immediate by reducing to the case where $X$ is affine.
This being so, the functor $\sh{G}'\mapsto g_*(\sh{G}')$ is exact and faithful, and by hypothesis the same is true of the functor $\sh{G}\mapsto\sh{F}\otimes_{\sh{O}_Y}\sh{G}=\sh{F}\otimes_{\sh{O}_X}f^*(\sh{G})$ from quasi-coherent $\sh{O}_Y$-Modules to quasi-coherent $\sh{O}_X$-Modules.
Their composite is therefore exact and faithful, which completes the proof of~\sref{IV.2.2.1} and~\sref{IV.2.2.2}.
\end{proof}

\begin{corollary}[2.2.3]
\label{IV.2.2.3}
Let $X=\Spec(B)$, $Y=\Spec(A)$ be two affine schemes, $f:X\to Y$ a morphism, $\sh{F}=\widetilde{M}$ a quasi-coherent $\sh{O}_X$-Module.
For $\sh{F}$ to satisfy the equivalent conditions of~\sref{IV.2.2.1} \textnormal{(}or~\sref{IV.2.2.2}\textnormal{)}, it is necessary and sufficient that the $A$-module $M$ be faithfully flat.
\end{corollary}

\begin{proof}
Indeed, condition~\emph{c)} of~\sref{IV.2.2.2} means that the functor $N\mapsto M\otimes_A N$ from the category of $A$-modules to that of $B$-modules is exact and faithful, and the conclusion follows from~\sref[0\textsubscript{I}]{0.6.4.1}.
\end{proof}

\begin{definition}[2.2.4]
\label{IV.2.2.4}
When the equivalent conditions of~\sref{IV.2.2.1} are satisfied, we say that the quasi-coherent $\sh{O}_X$-Module $\sh{F}$ is \emph{faithfully flat relative to} $f$ (or to~$Y$).
\end{definition}

\oldpage[IV-2]{11}
We note that this notion is local over $Y$, but \emph{not over} $X$; in particular we can have $\sh{F}_x=0$ for certain $x\in X$, in other words $\operatorname{Supp}(\sh{F})$ is not necessarily equal to~$X$.
However, it follows immediately from criterion~\sref{IV.2.2.1}\emph{(b)} that for every $y\in Y$, there exists at least one $x\in f^{-1}(y)$ such that $(\sh{F}_y)_x=\sh{F}_x\otimes_{\sh{O}_y}k(y)\neq 0$, and \emph{a fortiori} $\sh{F}_x\neq 0$; in other terms:

\begin{corollary}[2.2.5]
\label{IV.2.2.5}
If $\sh{F}$ is a quasi-coherent $\sh{O}_X$-Module, faithfully flat relative to $f$, then
\[
  f(\operatorname{Supp}(\sh{F}))=Y.
\]
In particular, $f$ is a surjective morphism.
\end{corollary}

This result admits a partial converse:

\begin{corollary}[2.2.6]
\label{IV.2.2.6}
Let $\sh{F}$ be a quasi-coherent $\sh{O}_X$-Module of finite type. For $\sh{F}$ to be faithfully flat relative to $f$, it is necessary and sufficient that $\sh{F}$ be $f$-flat and that $f(\operatorname{Supp}(\sh{F}))=Y$.
\end{corollary}

\begin{proof}
Indeed, by~\sref[I]{I.9.1.13} and~\sref[I]{I.3.6.1}, we have $\operatorname{Supp}(\sh{F}_y)=f^{-1}(y)\cap\operatorname{Supp}(\sh{F})$, so in this case criterion~\sref{IV.2.2.1}\emph{(b)} is precisely the condition of the corollary.
In particular, the $\sh{O}_X$-Module $\sh{O}_X$ is faithfully flat relative to $f$ if and only if $f$ is flat and surjective, that is to say if and only if the \emph{morphism $f$ is faithfully flat}~\sref[0\textsubscript{I}]{0.6.7.8}.
\end{proof}

Let us spell out the properties intervening in definition~\sref{IV.2.2.4}:

\begin{proposition}[2.2.7]
\label{IV.2.2.7}
Let $f:X\to Y$ be a morphism of preschemes, $\sh{F}$ a quasi-coherent $\sh{O}_X$-Module faithfully flat relative to $f$.
Then, for a sequence $\sh{G}'\to\sh{G}\to\sh{G}''$ of quasi-coherent $\sh{O}_Y$-Modules to be exact, it is necessary and sufficient that the corresponding sequence $\sh{F}\otimes_{\sh{O}_Y}\sh{G}'\to\sh{F}\otimes_{\sh{O}_Y}\sh{G}\to\sh{F}\otimes_{\sh{O}_Y}\sh{G}''$ be so.
In particular, for a homomorphism $u:\sh{G}\to\sh{G}'$ of quasi-coherent $\sh{O}_Y$-Modules to be injective \textnormal{(}resp.\ surjective, bijective, zero\textnormal{)}, it is necessary and sufficient that $1_{\sh{F}}\otimes f^*(u):\sh{F}\otimes_{\sh{O}_Y}\sh{G}\to\sh{F}\otimes_{\sh{O}_Y}\sh{G}'$ be so.
For a quasi-coherent $\sh{O}_Y$-Module $\sh{G}$ to be zero, it is necessary and sufficient that $\sh{F}\otimes_{\sh{O}_Y}\sh{G}=0$.
For a quasi-coherent $\sh{O}_Y$-Module $\sh{G}$, the map $\sh{G}'\mapsto\sh{F}\otimes_{\sh{O}_Y}\sh{G}'$ from the set of quasi-coherent $\sh{O}_Y$-submodules of $\sh{G}$ to the set of quasi-coherent $\sh{O}_X$-submodules of $\sh{F}\otimes_{\sh{O}_Y}\sh{G}$ is injective.
\end{proposition}

\begin{proof}
To prove the last assertion, that is to say that for two quasi-coherent sub-$\sh{O}_Y$-Modules $\sh{G}'$, $\sh{G}''$ of $\sh{G}$, the relation $\sh{F}\otimes_{\sh{O}_Y}\sh{G}'=\sh{F}\otimes_{\sh{O}_Y}\sh{G}''$ implies $\sh{G}'=\sh{G}''$, we can reduce to the case where $\sh{G}'\subset\sh{G}''$ (by replacing $\sh{G}''$ by $\sh{G}'+\sh{G}''$), and it then suffices to apply the second assertion of the statement to the injection $u:\sh{G}'\to\sh{G}''$.
\end{proof}

\begin{corollary}[2.2.8]
\label{IV.2.2.8}
Let $f:X\to Y$ be a faithfully flat morphism, $\sh{G}$ a quasi-coherent $\sh{O}_Y$-Module.
The canonical map
\begin{equation}
\tag{2.2.8.1}
\label{IV.2.2.8.1}
  \Gamma(Y,\,\sh{G}) \to \Gamma(X,\,f^*(\sh{G}))
\end{equation}
is injective.
\end{corollary}

\begin{proof}
Indeed, there are canonical identifications~\sref[0\textsubscript{I}]{0.5.1.1}
\begin{align*}
  \Gamma(Y,\,\sh{G})
    &\simeq \operatorname{Hom}_{\sh{O}_Y}(\sh{O}_Y,\,\sh{G}),\\
  \Gamma(X,\,f^*(\sh{G}))
    &\simeq \operatorname{Hom}_{\sh{O}_X}(\sh{O}_X,\,f^*(\sh{G})).
\end{align*}
By virtue of~\sref{IV.2.2.1} and~\sref{IV.2.2.4}, the hypothesis implies that the functor $\sh{G}\mapsto f^*(\sh{G})$ is exact and faithful in the category of quasi-coherent $\sh{O}_Y$-Modules, and thus a homomorphism $u:\sh{O}_Y\to\sh{G}$ is zero if and only if $f^*(u):f^*(\sh{O}_Y)=\sh{O}_X\to f^*(\sh{G})$ is zero.
\end{proof}

\begin{remark}[2.2.9]
\label{IV.2.2.9}
The assertions of~\sref{IV.2.2.7} and~\sref{IV.2.2.8} are still true when the $\sh{O}_Y$-Modules $\sh{G}'$, $\sh{G}$, $\sh{G}''$ appearing there are arbitrary (not necessarily quasi-coherent).
Indeed, for every $y\in Y$, there exists $x\in f^{-1}(y)$ such that $\sh{F}_x$ is a flat $\sh{O}_y$-module
\oldpage[IV-2]{12}
faithfully flat, and thus the functor $\sh{G}_y\mapsto\sh{G}_y\otimes_{\sh{O}_y}\sh{F}_x$ is faithful; since moreover for \emph{every} $x\in f^{-1}(y)$ the functor $\sh{G}_y\mapsto\sh{G}_y\otimes_{\sh{O}_y}\sh{F}_x$ is exact, we immediately deduce the conclusion.
\end{remark}

\begin{proposition}[2.2.10]
\label{IV.2.2.10}
Let $f:X\to Y$, $g:Y'\to Y$, $h:Y'\to Z$ be three morphisms of preschemes, $X'=X\times_Y Y'$, $\sh{F}$ a quasi-coherent $\sh{O}_X$-Module faithfully flat relative to $Y$, $\sh{G}'$ a quasi-coherent $\sh{O}_{Y'}$-Module.
For $\sh{G}'$ to be flat \textnormal{(}resp.\ faithfully flat\textnormal{)} relative to $Z$, it is necessary and sufficient that $\sh{F}\otimes_{\sh{O}_Y}\sh{G}'$ be a flat \textnormal{(}resp.\ faithfully flat\textnormal{)} $\sh{O}_{X'}$-Module relative to $Z$.
\end{proposition}

\begin{proof}
We know already that if $\sh{G}'$ is $Z$-flat, so is $\sh{F}\otimes_{\sh{O}_Y}\sh{G}'$~\sref{IV.2.1.5}.
Consider a base change $Z''\to Z$ and set $X''=X'\times_Z Z''$; if $\sh{G}'$ is faithfully flat relative to $Z$, the functor
\begin{equation}
\tag{2.2.10.1}
\label{IV.2.2.10.1}
\sh{H}''\mapsto\sh{H}''\otimes_Z\sh{G}'\to(\sh{H}''\otimes_Z\sh{G}')\otimes_{\sh{O}_Y}\sh{F}=\sh{H}''\otimes_Z(\sh{G}'\otimes_{\sh{O}_Y}\sh{F})
\end{equation}
from the category of $\sh{O}_{Z''}$-Modules to that of $\sh{O}_{X''}$-Modules is composed of two exact and faithful functors, so this functor is also exact and faithful.
Conversely, if this composite functor is exact \textnormal{(}resp.\ exact and faithful\textnormal{)}, the same is true of $\sh{H}''\mapsto\sh{H}''\otimes_Z\sh{G}'$, since the functor $\sh{M}'\mapsto\sh{M}'\otimes_{\sh{O}_Y}\sh{F}$ from quasi-coherent $\sh{O}_{Y'}$-Modules to quasi-coherent $\sh{O}_{X'}$-Modules is exact and faithful by hypothesis.
\end{proof}

\begin{corollary}[2.2.11]
\label{IV.2.2.11}
\medskip\noindent
\begin{enumerate}
  \item[\textnormal{(i)}] Let $f:X\to Y$, $g:Y'\to Y$ be two morphisms, $X'=X\times_Y Y'$, $\sh{F}$ a quasi-coherent $\sh{O}_X$-Module. If $\sh{F}$ is faithfully flat relative to $Y$, $\sh{F}'=\sh{F}\otimes_{\sh{O}_Y}\sh{O}_{Y'}$ is faithfully flat relative to $Y'$.
  \item[\textnormal{(ii)}] Let $f:X\to Y$, $g:Y\to Z$ be two morphisms, $\sh{F}$ a quasi-coherent $\sh{O}_X$-Module faithfully flat relative to $Y$.
    For $g$ to be a faithfully flat morphism, it is necessary and sufficient that $\sh{F}$ be faithfully flat relative to $Z$.
  \item[\textnormal{(iii)}] Let $f:X\to Y$, $g:Y\to Z$ be two morphisms, $\sh{G}$ a quasi-coherent $\sh{O}_Y$-Module.
    Suppose the morphism $f$ is faithfully flat.
    For $\sh{G}$ to be flat \textnormal{(}resp.\ faithfully flat\textnormal{)} relative to $Z$, it is necessary and sufficient that $f^*(\sh{G})$ be flat \textnormal{(}resp.\ faithfully flat\textnormal{)} relative to $Z$.
  \item[\textnormal{(iv)}] Let $f:X\to Y$, $g:Y\to Z$ be two morphisms, $x$ a point of $X$.
    Suppose $f$ is flat at the point $x$. For $g\circ f$ to be flat at the point $x$, it is necessary and sufficient that $g$ be flat at the point $f(x)$.
\end{enumerate}
\end{corollary}

\begin{proof}
To prove~(i), we apply~\sref{IV.2.2.10}, replacing $Z$ by $Y'$ and $\sh{G}'$ by $\sh{O}_{Y'}$.
To prove~(ii), we apply~\sref{IV.2.2.10} taking $Y'\to Y$ the identity and replacing $\sh{G}'$ by $\sh{O}_Y$.
To prove~(iii), we apply~\sref{IV.2.2.10}, again taking $Y'\to Y$ to be the identity, and replacing $\sh{F}$ by $\sh{O}_X$ and $\sh{G}'$ by $\sh{G}$.
Finally, (iv) is deduced from~(ii) by replacing $X$ by $\Spec(\sh{O}_x)$, $\sh{F}$ by $\sh{O}_X$, $Y$ by $\Spec(\sh{O}_{f(x)})$ and $Z$ by $\Spec(\sh{O}_{g(f(x))})$, and taking into account~\sref[0\textsubscript{I}]{0.6.6.2}.
\end{proof}

\begin{corollary}[2.2.12]
\label{IV.2.2.12}
Let $Y$ be an affine scheme, $f:X\to Y$ a quasi-compact morphism,
\oldpage[IV-2]{13}
$\sh{F}$ a quasi-coherent $\sh{O}_X$-Module.
If $\sh{F}$ is faithfully flat relative to $f$, there exists an affine scheme $X'$ and a surjective local isomorphism $g:X'\to X$ such that $g^*(\sh{F})$ is faithfully flat relative to $f\circ g$.
\end{corollary}

\begin{proof}
Indeed, $X$ is quasi-compact, hence a finite union of affine open subsets $X_i$; it suffices to take for $X'$ the affine scheme given by the disjoint sum of the preschemes induced on the open subsets $X_i$, and for $g:X'\to X$ the canonical morphism.
It is clear that $g$ is a surjective local isomorphism, hence faithfully flat, and the hypothesis implies that $g^*(\sh{F})$ is faithfully flat relative to $f\circ g$, by virtue of~\sref{IV.2.2.11}\textnormal{(iii)}.
\end{proof}

\begin{corollary}[2.2.13]
\label{IV.2.2.13}
\medskip\noindent
\begin{enumerate}
  \item[\textnormal{(i)}] Let $f:X\to Y$, $g:Y'\to Y$ be two morphisms, $X'=X\times_Y Y'$, $f'=f_{(Y')}:X'\to Y'$.
    If $f$ is a faithfully flat morphism, so is $f'$.
  \item[\textnormal{(ii)}] If $f:X\to X'$, $g:Y\to Y'$ are two faithfully flat $S$-morphisms,
    \[
      f\times_S g: X\times_S Y \to X'\times_S Y'
    \]
    is faithfully flat.
  \item[\textnormal{(iii)}] Let $f:X\to Y$, $g:Y\to Z$ be two morphisms such that $f$ is faithfully flat.
    For $g$ to be a flat \textnormal{(}resp.\ faithfully flat\textnormal{)} morphism, it is necessary and sufficient that $g\circ f$ be a flat \textnormal{(}resp.\ faithfully flat\textnormal{)} morphism.
\end{enumerate}
\end{corollary}

\begin{proposition}[2.2.14]
\label{IV.2.2.14}
Let $f:X\to Y$ be a faithfully flat morphism and quasi-compact.
If $X$ is locally Noetherian, so is $Y$.
\end{proposition}

\begin{proof}
The question being local on $Y$, we can suppose $Y=\Spec(A)$ affine; as $f$ is quasi-compact, it follows from~\sref{IV.2.2.12} that we can also restrict to the case where $X=\Spec(B)$ is affine.
Then $B$ is a Noetherian ring by hypothesis, and a faithfully flat $A$-module~\sref{IV.2.2.3}; so $A$ is Noetherian~\sref[0\textsubscript{I}]{0.6.5.2}.
\end{proof}

\begin{proposition}[2.2.15]
\label{IV.2.2.15}
Let $f:X\to Y$ be a faithfully flat morphism.
If $\mathfrak{S}(X)$ and $\mathfrak{S}(Y)$ denote respectively the sets of sub-preschemes of $X$ and of $Y$, the map $Z\mapsto f^{-1}(Z)$ of $\mathfrak{S}(Y)$ to $\mathfrak{S}(X)$ is injective.
\end{proposition}

\begin{proof}
Since $f$ is surjective, on the underlying sets of a sub-prescheme $Z$ of $Y$ we have $f(f^{-1}(Z))=Z$.
On the other hand, if $U$ is an open subset of $Y$ containing $Z$ and in which $Z$ is closed, then $f^{-1}(U)$ is open in $X$, $f^{-1}(Z)$ is closed in $f^{-1}(U)$, and the restriction $f^{-1}(U)\to U$ of $f$ is a faithfully flat morphism.
We can therefore restrict to considering only the \emph{closed} sub-preschemes of $Y$.
Now, if $Z$ is a closed sub-prescheme of $Y$ corresponding to a quasi-coherent Ideal $\sh{J}$ of $\sh{O}_Y$, we know that $f^{-1}(Z)$ corresponds to the quasi-coherent Ideal $f^*(\sh{J})\sh{O}_X$ of $\sh{O}_X$~\sref[I]{I.4.4.5}, and since $f$ is flat, $f^*(\sh{J})\sh{O}_X$ identifies with $f^*(\sh{J})$.
But the map $\sh{J}\mapsto f^*(\sh{J})$ from the set of quasi-coherent Ideals of $\sh{O}_Y$ to the set of quasi-coherent Ideals of $\sh{O}_X$ is injective~\sref{IV.2.2.7}, whence the conclusion.
\end{proof}

\begin{corollary}[2.2.16]
\label{IV.2.2.16}
Let $X$, $Y$ be two $S$-preschemes; if $S'\to S$ is a faithfully flat morphism, the map $f\mapsto f_{(S')}$ from $\operatorname{Hom}_S(X,Y)$ to $\operatorname{Hom}_{S'}(X_{(S')},Y_{(S')})$ is injective.
\end{corollary}

\begin{proof}
We have $X_{(S')}\times_{S'}Y_{(S')}=(X\times_S Y)_{(S')}$~\sref[I]{I.3.3.10}, so the morphism projection $X_{(S')}\times_{S'}Y_{(S')}\to X\times_S Y$ is faithfully flat~\sref{IV.2.2.13}.
The elements of $\operatorname{Hom}_S(X,Y)$ correspond bijectively to the sub-preschemes of $X\times_S Y$ that are their graphs~\sref[I]{I.5.3.11}; if $Z_f$ is the graph of $f\in\operatorname{Hom}_S(X,Y)$, then $Z_{f_{(S')}}=g^{-1}(Z_f)$~\sref[I]{I.5.3.12}.
It suffices therefore to apply~\sref{IV.2.2.15}.
\end{proof}

\begin{proposition}[2.2.17]
\label{IV.2.2.17}
Let $A$ be a ring, $B$ an $A$-algebra such that $B$ is a faithfully flat $A$-module and of finite presentation.
Then the structural homomorphism $\varphi:A\to B$ is an isomorphism of the $A$-module $A$ onto a direct summand of the $A$-module $B$.
If $A$ is a local ring, $B$ is a free $A$-module and there exists a basis of this module containing the identity element of $B$.
\end{proposition}

\begin{proof}
\oldpage[IV-2]{14}
By virtue of Bourbaki, \emph{Alg.~comm.}, chap.~II, \textsection~3, n\textsuperscript{o}~3, prop.~12, it suffices to prove the proposition when $A$ is a local ring; we know then (\emph{loc.~cit.}, n\textsuperscript{o}~2, cor.~2 of prop.~5) that $B$ is an $A$-module free of finite type, and the conclusion follows from \emph{loc.~cit.}, prop.~5.
\end{proof}

\subsection{Topological properties of flat morphisms}
\label{subsection:IV.2.3}

\begin{lemma}[2.3.1]
\label{IV.2.3.1}
Let $f:X\to Y$ be a quasi-compact and quasi-separated morphism, $g:Y'\to Y$ a flat morphism; set $X'=X\times_Y Y'$, $f'=f_{(Y')}:X'\to Y'$.
Then, for every quasi-coherent $\sh{O}_X$-Module $\sh{F}$, the canonical homomorphism
\begin{equation}
\tag{2.3.1.1}
\label{IV.2.3.1.1}
  g^*(f_*(\sh{F})) \to f'_*(\sh{F}\otimes_{\sh{O}_X}\sh{O}_{X'})
\end{equation}
is bijective.
\end{lemma}

\begin{proof}
This is a special case of~\sref[III]{III.1.4.15} (improved in~{\hyperlink{ega4.c04.statement.1-7-21}{(1.7.21)}}).
\end{proof}

\begin{proposition}[2.3.2]
\label{IV.2.3.2}
Let $S$ be a prescheme, $f:X\to Y$ a quasi-compact and quasi-separated $S$-morphism, and let $Z$ be the sub-prescheme of $Y$ which is the closed image of $X$ under $f$~\textnormal{(}\sref[I]{I.9.5.3} and~{\hyperlink{ega4.c04.statement.1-7-8}{(1.7.8)}}\textnormal{)}.
Let $j:Z\to Y$ be the canonical injection, so that $f=j\circ g$, where $g:X\to Z$ is a morphism \textnormal{(}\emph{loc.~cit.}\textnormal{)}.
Let $h:S'\to S$ be a flat morphism, and set $f'=f_{(S')}:X_{(S')}\to Y_{(S')}$.
Then $j'=j_{(S')}:Z_{(S')}\to Y_{(S')}$ identifies with the canonical injection of the sub-prescheme of $Y_{(S')}$ which is the closed image of $X_{(S')}$ under $f'$.
\end{proposition}

\begin{proof}
Since the morphism $Y_{(S')}\to Y$ is flat~\sref{IV.2.1.4}, we may reduce to the case $S=Y$~\sref[I]{I.3.3.11}.
If $f=(\psi,\theta)$, then $Z$ is the closed sub-prescheme of $S$ defined by the quasi-coherent Ideal $\sh{J}\subset\sh{O}_S$ which is the kernel of
\[
  \theta:\sh{O}_S\to f_*(\sh{O}_X)
\]
\sref[I]{I.9.5.2}.
Since $h$ is flat, the quasi-coherent Ideal $\sh{J}\sh{O}_{S'}$ identifies with the kernel of
\[
  h^*(\theta):\sh{O}_{S'}\to h^*(f_*(\sh{O}_X)).
\]
If $f'=(\psi',\theta')$, then $\theta':\sh{O}_{S'}\to f'_*(\sh{O}_{X'})$ is the composite of $h^*(\theta)$ with the canonical homomorphism
\[
  h^*(f_*(\sh{O}_X))\to f'_*(\sh{O}_{X'})
\]
of~\sref{IV.2.3.1.1}.
The conclusion follows from~\sref{IV.2.3.1} and~\sref[I]{I.9.5.2}, since $f'$ is quasi-compact and quasi-separated~({\hyperlink{ega4.c04.prop.1-1-2}{1.1.2}} and {\hyperlink{ega4.c04.prop.1-2-2}{1.2.2}}).
\end{proof}

\begin{env}[2.3.3]
\label{IV.2.3.3}
We will say that a morphism $f:X\to Y$ is \emph{quasi-flat} if there exists a quasi-coherent $\sh{O}_X$-Module $\sh{F}$ \emph{of finite type} which is $f$-flat and whose support is equal to $X$.
We will say that $f$ is \emph{quasi-faithfully flat} if it is quasi-flat and surjective.
Every flat (resp.\ faithfully flat) morphism is quasi-flat (resp.\ quasi-faithfully flat), for then $\sh{F}=\sh{O}_X$ satisfies the preceding conditions.
\end{env}

It follows immediately from~\sref{IV.2.1.4} and from~\sref[I]{I.9.1.13}, that if $f$ is quasi-flat, then, for every morphism $Y'\to Y$, the morphism $f_{(Y')}:X\times_Y Y'\to Y'$ is quasi-flat.
Similarly (taking into account~\sref[I]{I.3.5.2}), if $f$ is quasi-faithfully flat, so is $f_{(Y')}$.

\begin{proposition}[2.3.4]
\label{IV.2.3.4}
Let $f:X\to Y$ be a quasi-flat morphism~\sref{IV.2.3.3}. Then $f$ possesses the following properties \textnormal{(}which are equivalent by virtue of~{\hyperlink{ega4.c04.cor.1-10-4}{(1.10.4)}}\textnormal{)}:
\begin{enumerate}
  \item[(i)] For every $x\in X$ and every generization $y'$ of $y=f(x)$, there exists a generization $x'$ of $x$ such that $f(x')=y'$.
  \item[(ii)] For every $x\in X$, the image under $f$ of $\Spec(\sh{O}_{X,x})$ is $\Spec(\sh{O}_{Y,y})$.
\oldpage[IV-2]{15}
  \item[\emph{(iii)}] For every closed irreducible part $Y'$ of $Y$, every irreducible component of $f^{-1}(Y')$ dominates $Y'$.
\end{enumerate}
\end{proposition}

\begin{proof}
It suffices, for example, to prove~(ii).
By hypothesis, there is an $\sh{O}_X$-Module quasi-coherent of finite type $\sh{F}$, which is $f$-flat and such that $\operatorname{Supp}(\sh{F})=X$.
For every $x\in X$, $\sh{F}_x$ is then an $\sh{O}_x$-module of finite type, nonzero, and moreover $\sh{F}_x$ is an $\sh{O}_y$-module \emph{flat}, for the homomorphism $\rho:\sh{O}_y\to\sh{O}_x$.
As the latter is local and $\sh{F}_x\neq 0$, the Nakayama lemma shows that $\sh{F}_x\otimes_{\sh{O}_y}k(y)\neq 0$, so $\sh{F}_x$ is a faithfully flat $\sh{O}_y$-module~\sref[0]{0.6.4.1}.
It follows that for every prime ideal $\mathfrak{q}$ of $\sh{O}_y$, there exists a prime ideal $\mathfrak{p}$ of $\sh{O}_x$ such that $\mathfrak{q}=\rho^{-1}(\mathfrak{p})$~\sref[0]{0.6.5.1}, which proves~(ii).
\end{proof}

\begin{corollary}[2.3.5]
\label{IV.2.3.5}
Let $f:X\to Y$ be a morphism satisfying the equivalent conditions~(i), (ii), (iii) of~\sref{IV.2.3.4} \textnormal{(}in particular a quasi-flat or open morphism~{\hyperlink{ega4.c04.cor.1-10-4}{(1.10.4)}}\textnormal{)}.
\begin{enumerate}
  \item[(i)] Let $Z$, $Z'$ be two closed irreducible subsets of $Y$ such that $Z\subset Z'$, and let $\mathrm{T}$ be an irreducible component of $f^{-1}(Z)$; then there exists an irreducible component $\mathrm{T}'$ of $f^{-1}(Z')$ containing $\mathrm{T}$ \textnormal{(}and dominating $Z'$\textnormal{)}.
  \item[(ii)] For every irreducible component $\mathrm{T}$ of $X$, $\overline{f(\mathrm{T})}$ is an irreducible component of $Y$.
  \item[(iii)] Suppose $Y$ irreducible, and, if $y$ is its generic point, suppose that $f^{-1}(y)$ is irreducible. Then $X$ is irreducible.
\end{enumerate}
\end{corollary}

\begin{proof}
(i) It suffices to apply~\sref{IV.2.3.4},~(i) by taking for $x$ the generic point of $\mathrm{T}$ ($y=f(x)$ is then the generic point of $Z$) and for $y'$ the generic point of $Z'$.

(ii) It is clear that $\overline{f(\mathrm{T})}=Z$ is irreducible, and by virtue of~(i), $Z$ cannot be strictly contained in a closed irreducible part of $Y$.

(iii) By virtue of~(ii), every irreducible component of $X$ dominates $Y$, and hence meets $f^{-1}(y)$, where $y$ is the generic point of $Y$.
The conclusion then follows from~\sref[0]{0.2.1.8}.
\end{proof}

\begin{proposition}[2.3.6]
\label{IV.2.3.6}
Let $Y$ be a prescheme whose set of irreducible components is locally finite \textnormal{(}which is the case if the underlying space is locally Noetherian~(cf.~\sref[I]{I.6.1.9})\textnormal{)}.
\begin{enumerate}
  \item[(i)] Every closed subset $W$ of $Y$ stable by generization~\sref[0]{0.2.1.2} is open. In particular, every connected component of $Y$ is open.
  \item[(ii)] Let $f:X\to Y$ be a closed morphism satisfying in addition the equivalent conditions~(i), (ii), (iii) of~\sref{IV.2.3.4} \textnormal{(}which will be the case if $f$ is quasi-flat or open\textnormal{)}. Then the image under $f$ of every connected component $C$ of $X$ is a connected component of $Y$.
  \item[(iii)] Let $f:X\to Y$ be a morphism satisfying the equivalent conditions~(i), (ii), (iii) of~\sref{IV.2.3.4}, and such that moreover for every $y\in Y$, the set $f^{-1}(y)$ is finite \textnormal{(}which will be the case if $f$ is quasi-finite or radicial\textnormal{)}. Then the set of irreducible components of $X$ is locally finite.
\end{enumerate}
\end{proposition}

\begin{proof}
\oldpage[IV-2]{16}
(i) If $y\in W$, the generic points $\eta_i$ ($1\leq i\leq m$) of the irreducible components $Y_i$ of $Y$ containing $y$ belong to $W$ by hypothesis.
Since $W$ is closed, these components themselves are contained in $W$; and since there is by hypothesis a neighbourhood $U$ of $y$ which is the union of the $U\cap Y_i$~\sref[0]{0.2.1.6}, we have $U\subset W$, so $W$ is open.
Since for every $y\in Y$, $\overline{\{y\}}$ is connected, a connected component of $Y$ is stable by generization, so the second assertion follows immediately from the first.

(ii) Since $C$ is closed in $X$, $f(C)$ is closed in $Y$ by hypothesis.
Moreover, since $C$ is stable by generization, the hypothesis on $f$ implies that $f(C)$ is stable by generization; so, by virtue of~(i), $f(C)$ is open and closed, and since it is connected, it is a connected component of $Y$.

(iii) Let $x\in X$.
By hypothesis, there is an open neighbourhood $U$ of $y=f(x)$ which meets only finitely many irreducible components $Y_i$ of $Y$ ($1\leq i\leq n$); let $y_i$ be the generic point of $Y_i$.
For every irreducible component $Z$ of $X$ meeting $f^{-1}(U)$, the generic point $z$ of $Z$ is necessarily in one of the sets $f^{-1}(y_i)$~\sref{IV.2.3.4}.
Since each of these sets is finite by hypothesis, this proves our assertion.
\end{proof}

\begin{proposition}[2.3.7]
\label{IV.2.3.7}
Let $f:X\to Y$ be a morphism, $g:Y'\to Y$ a quasi-flat morphism, $X'=X\times_Y Y'$, $f'=f_{(Y')}:X'\to Y'$.
\begin{enumerate}
  \item[(i)] If $f$ is quasi-compact and dominant, so is $f'$.
  \item[(ii)] If every irreducible component of $X$ dominates an irreducible component of $Y$, then every irreducible component of $X'$ dominates an irreducible component of $Y'$.
\end{enumerate}
\end{proposition}

\begin{proof}
Denote by $g':X'\to X$ the canonical projection, which is a quasi-flat morphism~\sref{IV.2.3.3}.

(i) We know already~\sref{IV.1.1.2} that $f'$ is quasi-compact; moreover, if $y'$ is a maximal point of $Y'$, $y=g(y')$ is a maximal point of $Y$, as follows from~\sref{IV.2.3.4},~(iii).
By hypothesis~\sref{IV.1.1.5}, there exists $x\in X$ such that $f(x)=y$; so~\sref[I]{I.3.4.7}, there exists $x'\in X'$ such that $f'(x')=y'$.

(ii) Let $x'$ be a maximal point of $X'$, and let $x=g'(x')$; it follows from~\sref{IV.2.3.4},~(ii) that $x$ is a maximal point of $X$, and by hypothesis $y=f(x)$ is a maximal point of $Y$.
Set $y'=f'(x')$, so that $g(y')=y$; it remains to show that $y'$ is a maximal point of $Y'$.
By virtue of~\sref[I]{I.5.1.7} and~\sref{IV.2.3.3}, we may reduce to the case where $X$ and $Y$ are reduced, so that $\sh{O}_x$ and $\sh{O}_y$ are fields.
By virtue of~\sref[I]{I.3.6.5}, applied twice, and~\sref[I]{I.2.4.4}, we may further reduce to the case where $Y=\Spec(\sh{O}_y)$ and $Y'=\Spec(\sh{O}_{y'})$.
Then $f$ is flat because $\sh{O}_y$ is a field~\sref{IV.2.1.2}, and the same is true of $f'$~\sref{IV.2.1.4}; hence~\sref{IV.2.3.4},~(ii) shows that $y'=f'(x')$ is a maximal point of $Y'$.
\end{proof}

\begin{corollary}[2.3.8]
\label{IV.2.3.8}
\emph{(Zariski).}
Let $A$, $B$ be two Noetherian local rings, $\mathfrak{m}$, $\mathfrak{n}$ their maximal ideals respectively, $\varphi:A\to B$ a local homomorphism. Suppose the following hypotheses satisfied:
\begin{enumerate}
  \item[$1^\circ$] $B$ is an $A$-algebra essentially of finite type~{\hyperlink{ega4.c04.statement.1-3-8}{(1.3.8)}}.
  \item[$2^\circ$] The $\mathfrak{m}$-adic completion $\widehat{A}$ of $A$ for the $\mathfrak{m}$-adic topology is integral.
  \item[$3^\circ$] $\varphi$ is injective.
\end{enumerate}
Then the $\mathfrak{m}$-adic topology of $A$ is induced by the $\mathfrak{n}$-adic topology of $B$.
\end{corollary}

\begin{proof}
\oldpage[IV-2]{17}
Set $B'=B\otimes_A\widehat{A}$; by hypothesis $1^\circ$, $B'$ is of the form $S^{-1}(C\otimes_A\widehat{A})$, where $C$ is an $A$-algebra of finite type and $S$ is a multiplicative subset of $C$; hence $B'$ is Noetherian.
As $A$ identifies with a subring of $\widehat{A}$~\sref[0]{0.7.3.5}, $A$ is integral by hypothesis $2^\circ$.
Hypothesis $3^\circ$ implies that there exists a prime ideal $\mathfrak{q}$ of $B$ inducing the zero ideal of $A$~\sref[0]{0.1.5.8}; consequently, the local homomorphism $A\to B/\mathfrak{q}$ is injective.
We can therefore restrict to proving the conclusion of~\sref{IV.2.3.8} under the added hypothesis that $B$ is a local ring which is \emph{integral}.
Apply~\sref{IV.2.3.7},~(ii) to $Y=\Spec(A)$, $X=\Spec(B)$, $Y'=\Spec(\widehat{A})$, and $X'=\Spec(B')$.
Since $Y'\to Y$ is flat and the integral scheme $X$ dominates the integral scheme $Y$, every irreducible component of $X'$ dominates $Y'$.
If $y$, $x$, and $y'$ are the closed points of $Y$, $X$, and $Y'$ respectively, there is a point $x'\in X'$ (in fact unique) lying above $x$ and $y'$~\sref[I]{I.3.4.9}, and $\Spec(\sh{O}_{x'})$ therefore dominates $\Spec(\sh{O}_{y'})$.
Consequently, there is a commutative diagram of local homomorphisms of Noetherian local rings
\[
\xymatrix{
  B=\sh{O}_x \ar[r] & \sh{O}_{x'}\\
  A=\sh{O}_y \ar[u]^{\varphi} \ar[r]_{u} & \sh{O}_{y'}=\widehat{A} \ar[u]_{v}
}
\]
in which $u$ and $v$ are injective~\sref[I]{I.1.2.7}.
Identify $A$ and $\widehat{A}$ with subrings of $\sh{O}_{x'}$, and let $\mathfrak{r}$ be the maximal ideal of $\sh{O}_{x'}$.
Then $\bigcap_k(\mathfrak{r}^k\cap\widehat{A})=0$~\sref[0]{0.7.3.5}; since $\widehat{A}$ is complete and these ideals are open in $\widehat{A}$, it follows (Bourbaki, \emph{Alg.~comm.}, chap.~III, \textsection~2, n\textsuperscript{o}~7, prop.~8) that the topology of $\widehat{A}$ is induced by the $\mathfrak{r}$-preadic topology of $\sh{O}_{x'}$.
\emph{A fortiori}, the same is true of the topology of $A$~\sref[0]{0.7.3.5}.
Moreover, $\mathfrak{n}^k\cap A\subset\mathfrak{r}^k\cap A$, so the $\mathfrak{n}$-preadic topology of $B$ induces on $A$ a topology finer than the $\mathfrak{m}$-preadic topology; but $\mathfrak{m}^k\subset\mathfrak{n}^k\cap A$, so these two topologies are identical.
\end{proof}

\begin{remark}[2.3.9]
\label{IV.2.3.9}
We will see further on~({\hyperlink{ega4.c04.statement.7-8-3}{7.8.3}}, (vii)) that for the Noetherian local rings $A$ most commonly used in Algebraic Geometry, the hypothesis that $A$ is integral and integrally closed implies that the same is true of $\widehat{A}$.
This is why, in Algebraic Geometry over a base field, one generally states~\sref{IV.2.3.8} under the hypothesis that $A$ is \emph{integral and integrally closed}.
\end{remark}

\begin{theorem}[2.3.10]
\label{IV.2.3.10}
Let $f:X\to Y$ be a quasi-flat morphism~\sref{IV.2.3.3}.
Then, for every pro-constructible part~{\hyperlink{ega4.c04.def.1-9-4}{(1.9.4)}} $Z$ of $Y$, we have $\overline{f^{-1}(Z)}=f^{-1}(\overline{Z})$.
\end{theorem}

\begin{proof}
Since $f$ is continuous, we have $\overline{f^{-1}(Z)}\subset f^{-1}(\overline{Z})$ and it amounts to proving that for every $x\in X$ such that $f(x)\in\overline{Z}$, $x$ is adherent to $f^{-1}(Z)$;
it is clear that the question is local on $Y$, so we can suppose $Y$ affine.
By virtue of the hypothesis, there exists an affine scheme $Y'$ and a morphism $g:Y'\to Y$ such that $g(Y')=Z$~({\hyperlink{ega4.c04.prop.1-9-5}{1.9.5}},~(ix)).
Let $Y_1$ be the closed image of $Y'$ under $g$~\sref[I]{I.9.5.3}, and let $X_1$ be the closed sub-prescheme $f^{-1}(Y_1)$ of $X$; if $f_1:X_1\to Y_1$ is the restriction of $f$ to $X_1$, we know that $f_1$ is quasi-flat~\sref{IV.2.3.3}; we can therefore replace $X$, $Y$ by $X_1$, $Y_1$ respectively, in other words suppose that $g$ is \emph{dominant}.
Set then $X'=X\times_Y Y'$, and let $f'$ and $g'$ be the projections of $X'$ into $Y'$ and $X$ respectively, so that we have the commutative diagram
\[
\xymatrix{
  X \ar[d]_f & X' \ar[l]_{g'} \ar[d]^{f'} \\
  Y & Y' \ar[l]^g
}
\]
Since $f$ is quasi-flat and $g$ is quasi-compact and dominant, it follows from~\sref{IV.2.3.7},~(i) (with the roles of $f$ and $g$ interchanged) that $g'$ is dominant, which proves the theorem.
\end{proof}

\begin{corollary}[2.3.11]
\label{IV.2.3.11}
Let $f$ be a quasi-flat and quasi-compact morphism, $F$ a closed part of $X$ such that $F=f^{-1}(f(F))$; then $F=f^{-1}(\overline{f(F)})$.
\end{corollary}

\begin{proof}
\oldpage[IV-2]{18}
Let $Y'$ be the reduced sub-prescheme of $X$ having $F$ as underlying space~\sref[I]{I.5.2.1}, and let $j:Y'\to X$ be the canonical injection.
Then $f\circ j$ is quasi-compact~{\hyperlink{ega4.c04.prop.1-1-2}{(1.1.2)}}, so $Z=f(F)$ is pro-constructible in $Y$~({\hyperlink{ega4.c04.prop.1-9-5}{1.9.5}},~(vii)), and the corollary follows from the fact that $F$ is closed.

We can further write the result of~\sref{IV.2.3.11} in the form $F=f^{-1}(f(X)\cap\overline{f(F)})$.
In other words, the closed subsets of the subspace $f(X)$ of $Y$ are the parts $Z\subset f(X)$ such that $f^{-1}(Z)$ is closed in $X$; this also means that \emph{the topology induced by $Y$ on $f(X)$ is the quotient of that of $X$ by the equivalence relation defined by $f$}.
In particular:
\end{proof}

\begin{corollary}[2.3.12]
\label{IV.2.3.12}
Let $f:X\to Y$ be a quasi-faithfully flat morphism~\sref{IV.2.3.3} and quasi-compact.
Then the topology of $Y$ is the quotient of that of $X$ by the equivalence relation defined by $f$, in other words, for $Z\subset Y$ to be open \textnormal{(}resp.\ closed\textnormal{)} in $Y$, it is necessary and sufficient that $f^{-1}(Z)$ be open \textnormal{(}resp.\ closed\textnormal{)} in $X$.
\end{corollary}

\begin{proof}
Indeed, we then have $f(X)=Y$.
\end{proof}

\begin{corollary}[2.3.13]
\label{IV.2.3.13}
Let $X$, $Y$ be two $S$-preschemes, $f:X\to Y$ an $S$-morphism faithfully flat and quasi-compact.
For $Y$ to be separated over $S$, it is necessary and sufficient that the canonical immersion $j:X\times_Y X\to X\times_S X$~\sref[I]{I.5.3.10} be \emph{closed}.
\end{corollary}

\begin{proof}
Let us note for this that we have the commutative diagram~\sref[I]{I.5.3.5}
\[
\xymatrix{
  X\times_Y X \ar[r]^j \ar[d]_{\pi} & X\times_S X \ar[d]^{f\times_S f}\\
  Y \ar[r]_{\Delta_Y} & Y\times_S Y
}
\]
identifying $X\times_Y X$ with the product of the $(Y\times_S Y)$-preschemes $Y$ and $X\times_S X$.
Since $f$ is surjective, so are $\pi$ and $f\times_S f$; hence the inverse image of the diagonal $\Delta_Y(Y)$ under $f\times_S f$ is the image $j(X\times_Y X)$~\sref[I]{I.3.4.8}.
Since $f\times_S f$ is faithfully flat and quasi-compact~({\hyperlink{ega4.c04.prop.1-1-2}{1.1.2}} and {\hyperlink{ega4.c04.cor.2-2-13}{2.2.13}}), it suffices to apply~\sref{IV.2.3.12} to this morphism.
\end{proof}

\begin{corollary}[2.3.14]
\label{IV.2.3.14}
Let $f:X\to Y$ be a quasi-faithfully flat and quasi-compact morphism, and let $Z$ be a part of $Y$.
For $Z$ to be a locally closed part pro-constructible of $Y$, it is necessary and sufficient that $f^{-1}(Z)$ be a locally closed part pro-constructible of $X$.
\end{corollary}

\begin{proof}
We already know~{\hyperlink{ega4.c04.cor.1-9-12}{(1.9.12)}} that, for $Z$ to be a pro-constructible part of $Y$, it is necessary and sufficient that $f^{-1}(Z)$ be a pro-constructible part of $X$.
The condition is evidently necessary.
To prove that it is sufficient, consider the closed reduced sub-prescheme $Y_1$ of $Y$ having $\overline{Z}$ as underlying space and let $X_1$ be its inverse image, so that $f^{-1}(\overline{Z})=\overline{f^{-1}(Z)}$ by virtue of~\sref{IV.2.3.10}.
As $f^{-1}(Z)$ is locally closed in $X$, it is open in $\overline{f^{-1}(Z)}$, so in $X_1$; now, the morphism $f_1:X_1\to Y_1$ deduced from $f$ by restriction is quasi-faithfully flat and quasi-compact~({\hyperlink{ega4.c04.prop.1-1-2}{1.1.2}} and {\hyperlink{ega4.c04.statement.2-3-3}{2.3.3}}); we deduce from~\sref{IV.2.3.12} that $Z$ is open in $\overline{Z}$, which shows that $Z$ is locally closed in $Y$.
\end{proof}

\begin{remark}[2.3.15]
\label{IV.2.3.15}
It does not suffice, in~\sref{IV.2.3.12}, to suppose only that $f$ is \emph{faithfully flat}.
For example, take for $Y$ a reduced irreducible algebraic curve~\sref[II]{II.7.4.2}, and let
\oldpage[IV-2]{19}
$X$ be the prescheme sum $\coprod_{y\in Y}\Spec(\sh{O}_y)$~\sref[I]{I.3.1}.
Let $f:X\to Y$ be the canonical morphism, which is faithfully flat~\sref[I]{I.2.4.2}.
If $\eta$ denotes the generic point of $Y$, then $Z=\{\eta\}$ is not open in $Y$~\sref[II]{II.7.4.3}, whereas $f^{-1}(Z)=\Spec(\sh{O}_\eta)$ is open in $X$, since $\sh{O}_\eta$ is a field.
\end{remark}

\subsection{Universally open morphisms and flat morphisms}
\label{subsection:IV.2.4}

\begin{env}[2.4.1]
\label{IV.2.4.1}
We have already defined~\sref[II]{II.5.4.9} the notion of \emph{universally closed} morphism; in the same manner, we give the following definitions:
\end{env}

\begin{definition}[2.4.2]
\label{IV.2.4.2}
We say that a morphism of preschemes $f:X\to Y$ is \emph{universally open} \textnormal{(}resp.\ \emph{universally bicontinuous}, resp.\ a \emph{universal homeomorphism}\textnormal{)} if, for every morphism $g:Y'\to Y$, the morphism $f_{(Y')}:X\times_Y Y'\to Y'$ is open \textnormal{(}resp.\ a homeomorphism onto its image, resp.\ a homeomorphism onto $Y'$\textnormal{)}.
\end{definition}

We will see further on~\sref{IV.14.3.2} that when $Y$ is locally Noetherian, the definition of universally open morphism given here is equivalent to the definition~\sref[III]{III.4.3.9} for morphisms of finite type; the reader can verify that we do not use the latter definition before \textsection~14.

\begin{proposition}[2.4.3]
\label{IV.2.4.3}
\medskip\noindent
\begin{enumerate}
  \item[(i)] An immersion \textnormal{(}resp.\ an open immersion, resp.\ a closed immersion\textnormal{)} is universally bicontinuous \textnormal{(}resp.\ universally open, resp.\ universally closed\textnormal{)}.
  \item[(ii)] The composite of two universally open morphisms \textnormal{(}resp.\ universally closed morphisms, resp.\ universally bicontinuous morphisms, resp.\ universal homeomorphisms\textnormal{)} has the same property.
  \item[(iii)] If $f:X\to Y$ is a universally open \textnormal{(}resp.\ universally closed, resp.\ universally bicontinuous, resp.\ a universal homeomorphism\textnormal{)} $S$-morphism, then so is $f_{(S')}:X_{(S')}\to Y_{(S')}$ after every base change $S'\to S$.
  \item[(iv)] If the $S$-morphisms $f:X\to X'$ and $g:Y\to Y'$ are both universally open \textnormal{(}resp.\ universally closed, resp.\ universally bicontinuous, resp.\ universal homeomorphisms\textnormal{)}, then so is
  \[
    f\times_S g:X\times_S Y\longrightarrow X'\times_S Y'.
  \]
  \item[(v)] Let $f:X\to Y$ and $g:Y\to Z$ be morphisms, with $f$ surjective. If $g\circ f$ is universally open \textnormal{(}resp.\ universally closed, resp.\ universally bicontinuous, resp.\ a universal homeomorphism\textnormal{)}, then so is $g$.
  \item[(vi)] A morphism $f$ is universally open \textnormal{(}resp.\ universally closed, resp.\ universally bicontinuous, resp.\ a universal homeomorphism\textnormal{)} if and only if $f_{\mathrm{red}}$ has the same property.
  \item[(vii)] Let $(U_\alpha)$ be an open covering of $Y$. A morphism $f:X\to Y$ is universally open \textnormal{(}resp.\ universally closed, resp.\ universally bicontinuous, resp.\ a universal homeomorphism\textnormal{)} if and only if, for every $\alpha$, its restriction $f^{-1}(U_\alpha)\to U_\alpha$ has the same property.
\end{enumerate}
\end{proposition}

\begin{proof}
Assertion~(i) follows from~\sref[I]{I.4.3.2}.
Assertion~(ii) follows immediately from the definitions. The same is true of~(iii), after reducing to the case $Y=S$ and $Y'=S'$, as permitted by~\sref[I]{I.3.3.11}; assertion~(iv) follows from~(ii),~(iii), and~\sref[I]{I.3.5.1}.

To prove~(v), note that after every morphism $Z'\to Z$, the base change
\[
  f_{(Z')}:X_{(Z')}\longrightarrow Y_{(Z')}
\]
is surjective~\sref[I]{I.3.5.2}. It is therefore enough to use the following
purely topological fact: if $f$ is surjective and $g\circ f$ is open
\textnormal{(}resp.\ closed, resp.\ a
\oldpage[IV-2]{20}
homeomorphism onto its image, resp.\ a bijective homeomorphism\textnormal{)},
then $g$ has the corresponding property. In the open and closed cases this
is Bourbaki, \emph{Top.~g\'{e}n.}, chap.~I\textsuperscript{er},
3\textsuperscript{e}~\'{e}d., \textsection~5, n\textsuperscript{o}~1,
prop.~1. In the other two cases we may replace $Z$ by $g(Y)$ and thus
suppose that $g(f(X))=g(Y)=Z$. Then $g\circ f$ is a homeomorphism of $X$
onto $Z$; since $f$ is surjective, $g$ is bijective, and since the preceding
case shows that $g$ is open, $g$ is a homeomorphism of $Y$ onto $Z$.
To prove~(vi), note that a morphism $g$ is open \textnormal{(}resp.\ closed, resp.\ a homeomorphism onto its image, resp.\ a bijective homeomorphism\textnormal{)} if and only if $g_{\mathrm{red}}$ has the same property.
On the other hand, for every morphism $Y''\to Y$, we have
\[
  \bigl(X_{\mathrm{red}}\times_{Y_{\mathrm{red}}}Y''_{\mathrm{red}}\bigr)_{\mathrm{red}}
  =(X\times_Y Y'')_{\mathrm{red}}
\]
by~\sref[I]{I.5.1.8}; hence the preceding remark shows that if $f_{\mathrm{red}}$ is universally open \textnormal{(}resp.\ universally closed, resp.\ universally bicontinuous, resp.\ a universal homeomorphism\textnormal{)}, then so is $f$.
The converse is proved in the same manner, noting here that for every morphism $Y''\to Y_{\mathrm{red}}$, we have $(X_{\mathrm{red}}\times_{Y_{\mathrm{red}}}Y'')_{\mathrm{red}}=(X\times_Y Y'')_{\mathrm{red}}$~\sref[I]{I.5.1.8}.

Finally, the necessity of condition~(vii) follows immediately from~(iii).
Conversely, suppose the condition satisfied, and let $g:Y'\to Y$ be a morphism; then the $g^{-1}(U_\alpha)=U'_\alpha$ form an open covering of $Y'$ and if one notes $f_\alpha$ the restriction $f^{-1}(U_\alpha)\to U_\alpha$ of $f$, and $f'$ the morphism $f_{(Y')}$, the restriction $f'^{-1}(U'_\alpha)\to U'_\alpha$ is nothing other than $(f_\alpha)_{(U'_\alpha)}$.
We can therefore restrict to proving that $f'$ is open \textnormal{(}resp.\ closed, resp.\ a homeomorphism onto its image, resp.\ a homeomorphism onto $Y'$\textnormal{)}, which is immediate.
\end{proof}

\begin{proposition}[2.4.4]
\label{IV.2.4.4}
A universally bicontinuous morphism $f:X\to Y$ is radicial \textnormal{(}and hence separated~\sref[I]{I.7.7.1}\textnormal{)}.
\end{proposition}

\begin{proof}
Indeed, $f$ is universally injective by hypothesis~\sref[I]{I.3.5.11}.
\end{proof}

\begin{proposition}[2.4.5]
\label{IV.2.4.5}
\medskip\noindent
\begin{enumerate}
  \item[(i)] An integral, surjective, and radicial morphism $f:X\to Y$ is a universal homeomorphism.
  \item[(ii)] Conversely, suppose that $Y$ is locally Noetherian. If a morphism of finite type $f:X\to Y$ is a universal homeomorphism, then it is finite, surjective, and radicial.
\end{enumerate}
\end{proposition}

\begin{proof}
(i) The three properties of $f$ are preserved by base change~\sref[I]{I.3.5.2},~\sref[I]{I.3.5.7}, and~\sref[II]{II.6.1.5}. Since an integral morphism is closed~\sref[II]{II.6.1.10}, $f$ is a homeomorphism of $X$ onto $Y$.

(ii) The morphism $f$ is of finite type and universally closed by hypothesis, and it is separated by~\sref{IV.2.4.4}; hence it is proper~\sref[II]{II.5.4.1}. For every $y\in Y$, the fiber $f^{-1}(y)$ consists of a single point, so~\sref[III]{III.4.4.2} implies that $f$ is finite. It is plainly surjective, and it is radicial because it is universally injective~\sref[I]{I.3.5.11}.
\end{proof}

\begin{theorem}[2.4.6]
\label{IV.2.4.6}
Let $f:X\to Y$ be a quasi-flat morphism~\sref{IV.2.3.3} and locally of finite presentation.
Then $f$ is universally open.
In particular, a flat morphism locally of finite presentation is universally open.
\end{theorem}

\begin{proof}
We know that for every base change $Y'\to Y$, $f_{(Y')}$ is quasi-flat~\sref{IV.2.3.3} and locally of finite presentation~\sref{IV.1.4.3}[iii].
It suffices therefore to prove that $f$ is an
\oldpage[IV-2]{21}
open morphism.
But this follows from the criterion~\sref{IV.1.10.4} for morphisms locally of finite presentation, conditions~\emph{b)}, \emph{b$'$)}, and \emph{c)} of~\sref{IV.1.10.4} being nothing other than conditions~(i), (ii), (iii) of~\sref{IV.2.3.4}.
\end{proof}

\begin{corollary}[2.4.7]
\label{IV.2.4.7}
For every prescheme $Y$, the structural morphism
\[
  \mathbf{V}_Y^n\longrightarrow Y,
  \qquad
  \mathbf{V}_Y^n
  =Y\times_{\Spec(\mathbf{Z})}\Spec\bigl(\mathbf{Z}[T_1,\ldots,T_n]\bigr),
\]
where $\mathbf{V}_Y^n$ is also denoted by $Y[T_1,\ldots,T_n]$, is universally open.
\end{corollary}

\begin{proof}
Indeed, for $Y=\Spec(A)$, $Y[T_1,\ldots,T_n]=\Spec(A[T_1,\ldots,T_n])$, and $A[T_1,\ldots,T_n]$ is a free $A$-algebra of finite presentation.
\end{proof}

\begin{remark}[2.4.8]
\label{IV.2.4.8}
\medskip\noindent
\begin{enumerate}
  \item[(i)] A faithfully flat and quasi-compact morphism $f:X\to Y$ is \emph{not necessarily open}, even when $X$ and $Y$ are Noetherian.
    For example, let $Y$ be a reduced irreducible algebraic curve~\sref[II]{II.7.4.2} with generic point $y$, let $X=Y\coprod\Spec(k(y))$, and let $f$ be the canonical morphism. It is clear that $f$ is flat and surjective, hence faithfully flat, and it is quasi-compact. But the image under $f$ of the open subset $\Spec(k(y))$ of $X$ is $\{y\}$, which is not open in $Y$~\sref[II]{II.7.4.3}.
  \item[(ii)] For every prescheme $X$, the canonical morphism $f:X_{\mathrm{red}}\to X$ is a closed immersion and a \emph{universal homeomorphism} by~\sref{IV.2.4.3},~(vi). But when $X$ is locally Noetherian, $f$ is flat only if $X$ is reduced, in which case $f=1_X$~\sref{IV.2.2.17}.
\end{enumerate}
\end{remark}

\begin{proposition}[2.4.9]
\label{IV.2.4.9}
Let $Y$ be a discrete prescheme. Then every morphism $f:X\to Y$ is universally open.
\end{proposition}

\begin{proof}
The question is local on $Y$~\sref{IV.2.4.3},~(vii), so we may assume that the underlying space of $Y$ consists of a single point. Replacing $f$ by $f_{\mathrm{red}}$~\sref{IV.2.4.3},~(vi), we may further assume that $Y$ is the spectrum of a \emph{field} $k$. For every base change $Y'\to Y$, the inverse images in $X'=X\times_Y Y'$ of the affine open subsets of $X$ cover $X'$, so we may also suppose that $X=\Spec(B)$ is affine, with $B$ a $k$-algebra.
It suffices then to prove that for every $k$-algebra $A'$, if we set $B'=B\otimes_k A'$, the image of every open subset of the form $U=D(t)$ ($t\in B'$) is open in $\Spec(A')$.
Now $B$ is the filtered direct limit of the increasing family $(B_\alpha)$ of its sub-$k$-algebras \emph{of finite type}. Since filtered direct limits commute with tensor products, $B'$ is the direct limit of the $A'$-algebras $B_\alpha\otimes_k A'=B'_\alpha$.
There exists $\alpha$ such that $t$ is the image in $B'$ of an element $t_\alpha\in B'_\alpha$, so $D(t)$ is the inverse image under the canonical morphism $u_\alpha:\Spec(B')\to\Spec(B'_\alpha)$ of the open subset $U_\alpha=D(t_\alpha)$~\sref[I]{I.1.2.2.2}.
But since $k$ is a field, $B_\alpha$ is a $k$-algebra of finite presentation and a $k$-module flat, hence $B'_\alpha$ is an $A'$-algebra of finite presentation and a flat $A'$-module; we thus conclude from~\sref{IV.2.4.6} that the image of $U_\alpha$ by $f'_\alpha:\Spec(B'_\alpha)\to\Spec(A')$ is \emph{open} in $\Spec(A')$.
It remains to prove that $f'(U)=f'_\alpha(U_\alpha)$. The inclusion $f'(U)\subset f'_\alpha(U_\alpha)$ is clear, so it remains to show that, for every point $y'\in f'_\alpha(U_\alpha)$, the intersection $V=U\cap f'^{-1}(y')$ is nonempty.
Now $V=u_\alpha^{-1}(V_\alpha)$, where $V_\alpha=U_\alpha\cap f_\alpha'^{-1}(y')$; in other words, $V_\alpha$ is a non-empty open subset of the prescheme $\Spec(B'_\alpha\otimes_{A'}k(y'))=\Spec(B_\alpha\otimes_k k(y'))$ and $V$ is its inverse image by the morphism $v:\Spec(B\otimes_k k(y'))\to\Spec(B_\alpha\otimes_k k(y'))$.
  Since $B_\alpha$ is a subalgebra of $B$ and $k(y')$ is flat over the field
  $k$, the homomorphism
  $B_\alpha\otimes_k k(y')\to B\otimes_k k(y')$ is \emph{injective}.
  \oldpage[IV-2]{22}
  Thus $v$ is dominant~\sref[I]{I.1.2.7}, which shows that $V$ is nonempty.
\end{proof}

\begin{corollary}[2.4.10]
\label{IV.2.4.10}
Let $k$ be a field, $X$, $Y$ two $k$-preschemes; then the morphism projection $X\times_k Y\to X$ is universally open.
In particular, for every extension $K$ of $k$, the morphism projection $X_{(K)}\to X$ is universally open.
\end{corollary}

\begin{proof}
It suffices to apply~\sref{IV.2.4.9} to the structural morphism $Y\to\Spec(k)$.
\end{proof}

\begin{remark}[2.4.11]
\label{IV.2.4.11}
If $f:X\to Y$ is an \emph{open} morphism, then for every subset $E$ of $Y$ we have $f^{-1}(\overline{E})=\overline{f^{-1}(E)}$ (Bourbaki, \emph{Top.~g\'{e}n.}, chap.~I, 3\textsuperscript{e}~\'{e}d., \textsection~5, n\textsuperscript{o}~4, prop.~7).
This remark applies for example when $f$ is a morphism \emph{flat} locally of finite presentation~\sref{IV.2.4.6}, or a morphism projection $X\times_k Y\to X$ where $X$, $Y$ are preschemes over a field $k$~\sref{IV.2.4.10}, and generalises then~\sref{IV.2.3.10}.
\end{remark}

\subsection{Permanence of properties of Modules by faithfully flat descent}
\label{subsection:IV.2.5}

\begin{proposition}[2.5.1]
\label{IV.2.5.1}
Let $f:X\to Y$ be a morphism, let $\sh{F}$ be a quasi-coherent
$\sh{O}_X$-Module, and let $g:Y'\to Y$ be faithfully flat. Set
\[
  X'=X\times_Y Y',\qquad
  f'=f_{(Y')}:X'\to Y',
\]
let $g':X'\to X$ be the canonical projection, and set
$\sh{F}'=g'^*(\sh{F})$. Then $\sh{F}$ is flat \textnormal{(}resp.\
faithfully flat\textnormal{)} relative to $f$ if and only if $\sh{F}'$ is
flat \textnormal{(}resp.\ faithfully flat\textnormal{)} relative to $f'$.
\end{proposition}

\begin{proof}
The direct implication follows from~\sref{IV.2.1.5} in the flat case and
from~\sref{IV.2.2.11}[i] in the faithfully flat case. Conversely, suppose
that $\sh{F}'$ is flat \textnormal{(}resp.\ faithfully flat\textnormal{)}
relative to $f'$. Since $g$ is faithfully flat, $\sh{F}'$ is then flat
relative to $g\circ f'$ in the first case, by~\sref{IV.2.1.6}, and
faithfully flat relative to $g\circ f'$ in the second, by
~\sref{IV.2.2.11}[ii]. But $g':X'\to X$ is faithfully
flat~\sref{IV.2.2.13}, and
\[
  g\circ f'=f\circ g'.
\]
Thus~\sref{IV.2.2.11}[iii] shows that $\sh{F}$ is flat
\textnormal{(}resp.\ faithfully flat\textnormal{)} relative to $f$.
\end{proof}

\begin{proposition}[2.5.2]
\label{IV.2.5.2}
Let $f:X'\to X$ be a morphism faithfully flat and quasi-compact,
$\sh{F}$ a quasi-coherent $\sh{O}_X$-Module, and
$\sh{F}'=f^*(\sh{F})$. Consider, for a quasi-coherent Module, the property
of being:
\begin{enumerate}
  \item[(i)] of finite type;
  \item[(ii)] of finite presentation;
  \item[(iii)] locally free of finite type;
  \item[(iv)] locally free of rank $n$.
\end{enumerate}
If $\mathbf{P}$ denotes any one of these properties, then $\sh{F}$ has
property $\mathbf{P}$ if and only if $\sh{F}'$ does.
\end{proposition}

\begin{proof}
For a quasi-coherent $\sh{O}_X$-Module to be locally free of finite type,
it is necessary and sufficient that it be flat over $X$ and of finite
presentation (Bourbaki, \emph{Alg.\ comm.}, chap.~II, \S~5, no.~2,
cor.~2 of th.~1, taking~\sref{IV.2.1.2} into account). By
\sref{IV.2.5.1}, applied with the identity morphism for $f$, $\sh{F}$ is
flat over $X$ if and only if $\sh{F}'$ is flat over $X'$. Thus case~(iii)
follows from cases~(i) and~(ii). The same reduction proves~(iv), since
$f^*(\sh{O}_X^n)=\sh{O}_{X'}^n$: if $\sh{F}$ and $\sh{F}'$ are locally
free of finite type and $x=f(x')$, their ranks at $x$ and $x'$ are equal,
and $f$ is surjective.

\oldpage[IV-2]{23}

To treat cases~(i) and~(ii), note that the question is local on $X$, so we
may suppose that $X$ is affine. By~\sref{IV.2.2.12}, there is an affine
scheme $X''$ and a faithfully flat local isomorphism $g:X''\to X'$.
Consequently, it is equivalent for $f^*(\sh{F})$ or
$g^*f^*(\sh{F})$ to have property $\mathbf{P}$. We are therefore reduced
to
\[
  X=\Spec(\mathrm{A}),\qquad X'=\Spec(\mathrm{A}').
\]
In view of~\sref{IV.2.2.3} and~\sref[II]{II.6.1.4.1}, it remains only to
prove the following lemma.
\end{proof}

\begin{lemma}[2.5.3]
\label{IV.2.5.3}
Let $\mathrm{A}$ be a ring, let $\mathrm{A}'$ be a faithfully flat
$\mathrm{A}$-algebra, and let $\mathrm{M}$ be an $\mathrm{A}$-module.
Set $\mathrm{M}'=\mathrm{M}\otimes_{\mathrm{A}}\mathrm{A}'$. Then
$\mathrm{M}$ is of finite type \textnormal{(}resp.\ of finite
presentation\textnormal{)} if and only if $\mathrm{M}'$ is.
\end{lemma}

\begin{proof}
For the proof, see Bourbaki, \emph{Alg.\ comm.}, chap.~I\textsuperscript{er}, \S~3, no.~6, prop.~11.
\end{proof}

\begin{remark}[2.5.4]
\label{IV.2.5.4}
The assertions of~\sref{IV.2.5.2} for properties~(i) and~(ii) remain
valid if $f$ is assumed only \emph{quasi-faithfully flat}
\sref{IV.2.3.3} and quasi-compact. Indeed, one is reduced (Bourbaki,
\emph{loc.\ cit.}) to the following lemma.
\end{remark}

\begin{lemma}[2.5.4.1]
\label{IV.2.5.4.1}
Let $\rho:\mathrm{A}\to\mathrm{A}'$ be a homomorphism of rings such that
the corresponding morphism
$f:\Spec(\mathrm{A}')\to\Spec(\mathrm{A})$ is surjective. Suppose that
there is an $\mathrm{A}'$-module $\mathrm{N}'$ of finite type which is
flat over $\mathrm{A}$ and has support $\Spec(\mathrm{A}')$. If
$u:\mathrm{P}\to\mathrm{Q}$ is a homomorphism of $\mathrm{A}$-modules
such that
\[
  u\otimes 1:
  \mathrm{P}\otimes_{\mathrm{A}}\mathrm{A}'
  \longrightarrow
  \mathrm{Q}\otimes_{\mathrm{A}}\mathrm{A}'
\]
is surjective, then $u$ is surjective.
\end{lemma}

\begin{proof}
Tensoring the assumed surjection with $\mathrm{N}'$ first shows that
\[
  u\otimes 1_{\mathrm{N}'}:
  (\mathrm{P}\otimes_{\mathrm{A}}\mathrm{A}')
    \otimes_{\mathrm{A}'}\mathrm{N}'
  \longrightarrow
  (\mathrm{Q}\otimes_{\mathrm{A}}\mathrm{A}')
    \otimes_{\mathrm{A}'}\mathrm{N}'
\]
is surjective. Let $\mathfrak{q}$ be a prime ideal of $\mathrm{A}'$ and
put $\mathfrak{p}=\rho^{-1}(\mathfrak{q})$. After localization, the
corresponding map is
\[
  u_{\mathfrak{p}}\otimes 1:
  \mathrm{P}_{\mathfrak{p}}
    \otimes_{\mathrm{A}_{\mathfrak{p}}}\mathrm{N}'_{\mathfrak{q}}
  \longrightarrow
  \mathrm{Q}_{\mathfrak{p}}
    \otimes_{\mathrm{A}_{\mathfrak{p}}}\mathrm{N}'_{\mathfrak{q}}
\]
\sref[0]{0.1.5.4}. By the support hypothesis
$\mathrm{N}'_{\mathfrak{q}}\neq0$, and
$\mathrm{N}'_{\mathfrak{q}}$ is flat over
$\mathrm{A}_{\mathfrak{p}}$~\sref[0]{0.6.3.1}. Nakayama's lemma gives
$\mathfrak{q}\mathrm{N}'_{\mathfrak{q}}\neq
\mathrm{N}'_{\mathfrak{q}}$, hence a fortiori
$\mathfrak{p}\mathrm{N}'_{\mathfrak{q}}\neq
\mathrm{N}'_{\mathfrak{q}}$. Thus $\mathrm{N}'_{\mathfrak{q}}$ is
faithfully flat over $\mathrm{A}_{\mathfrak{p}}$
\sref[0]{0.6.4.1}, and $u_{\mathfrak{p}}$ is surjective. Since $f$ is
surjective, every $\mathfrak{p}\in\Spec(\mathrm{A})$ occurs in this way;
therefore $u$ is surjective (Bourbaki, \emph{Alg.\ comm.}, chap.~II,
\S~3, no.~3, th.~1).
\end{proof}

\begin{proposition}[2.5.5]
\label{IV.2.5.5}
Let $f:X\to Y$ be a morphism, let $\sh{F}$ be a quasi-coherent
$\sh{O}_X$-Module of finite type which is flat relative to $f$, and let
$\sh{G}$ be a quasi-coherent $\sh{O}_Y$-Module of finite type. For
$y\in Y$, set
\[
  \sh{F}_y=\sh{F}\otimes_{\sh{O}_Y}k(y).
\]
A point $x\in X$ is a maximal point of
$\operatorname{Supp}(\sh{F}\otimes_Y\sh{G})$ if and only if
$y=f(x)$ is a maximal point of $\operatorname{Supp}(\sh{G})$ and $x$
is a maximal point of $\operatorname{Supp}(\sh{F}_y)$ in $f^{-1}(y)$.
When this holds, we have
\end{proposition}

\begin{env}[2.5.5.1]
\label{IV.2.5.5.1}
\[
  \operatorname{length}\bigl((\sh{F}\otimes_Y\sh{G})_x\bigr)
  =
  \operatorname{length}(\sh{G}_y)\,
  \operatorname{length}\bigl((\sh{F}_y)_x\bigr).
\]
\end{env}

\begin{proof}
It is clear that
\[
  f\bigl(\operatorname{Supp}(\sh{F}\otimes_Y\sh{G})\bigr)
  \subset \operatorname{Supp}(\sh{G})
\]
\sref[0]{0.5.2.2}; hence the image under $f$ of every irreducible
component of $\operatorname{Supp}(\sh{F}\otimes_Y\sh{G})$ is contained
in an irreducible component of $\operatorname{Supp}(\sh{G})$.

  We may reduce to the case $\operatorname{Supp}(\sh{F})=X$. Indeed
  \textnormal{(}\textbf{Err}\textsubscript{\textbf{III}}, 30\textnormal{)},
  since $\sh{F}$ is of finite type, there is a closed subprescheme
  $X'\subset X$ with underlying space $\operatorname{Supp}(\sh{F})$ and a
  quasi-coherent $\sh{O}_{X'}$-Module $\sh{H}$ of finite type such that,
  for the canonical injection $j:X'\to X$,
  \[
    \sh{F}=j_*(\sh{H}).
  \]
  Set $f'=f\circ j$. Then $\sh{H}$ is flat relative to $f'$, and
  \[
    \operatorname{Supp}(\sh{F}\otimes_Y\sh{G})
    =
    \operatorname{Supp}(\sh{H}\otimes_Y\sh{G}).
  \]

  Suppose therefore that $\operatorname{Supp}(\sh{F})=X$. If $Z$ is an
  irreducible component of $\operatorname{Supp}(\sh{G})$, then
  \[
    f^{-1}(Z)\subset
    \operatorname{Supp}(\sh{F}\otimes_Y\sh{G})
  \]
  \sref[I]{I.9.1.13}; and~\sref{IV.2.3.4} shows that every irreducible
  component of $f^{-1}(Z)$ dominates $Z$. Thus, if $x$ is the generic
  point of an irreducible component of
  $\operatorname{Supp}(\sh{F}\otimes_Y\sh{G})$ contained in $f^{-1}(Z)$,
  then $y=f(x)$ is the generic point of $Z$. Moreover,
  by~\sref[0]{0.2.1.8}, $x$ is the generic point of an irreducible
  component of
  \[
    f^{-1}(y)=\operatorname{Supp}(\sh{F}_y)
  \]
  \sref[I]{I.9.1.13}; conversely, every generic point of one of these
  components is the generic point of an irreducible component of
  $f^{-1}(Z)$.

\oldpage[IV-2]{24}

It remains to prove~\sref{IV.2.5.5.1}. We have
\[
  (\sh{F}\otimes_Y\sh{G})_x
  =
  \sh{F}_x\otimes_{\sh{O}_{Y,y}}\sh{G}_y
\]
\sref[I]{I.9.1.12} and
$(\sh{F}_y)_x=\sh{F}_x/\mathfrak{m}_y\sh{F}_x$; hence we are reduced to
the following lemma.
\end{proof}

\begin{lemma}[2.5.5.2]
\label{IV.2.5.5.2}
Let $\mathrm{A}$ and $\mathrm{B}$ be local rings,
$\rho:\mathrm{A}\to\mathrm{B}$ a local homomorphism, and
$\mathfrak{m}$ the maximal ideal of $\mathrm{A}$. Let $\mathrm{M}$ be an
$\mathrm{A}$-module and $\mathrm{N}$ a $\mathrm{B}$-module which is
faithfully flat over $\mathrm{A}$ and such that
$\mathrm{N}/\mathfrak{m}\mathrm{N}$ has finite length as a
$\mathrm{B}$-module. Then
\end{lemma}

\begin{env}[2.5.5.3]
\label{IV.2.5.5.3}
\[
  \operatorname{length}_{\mathrm{B}}(\mathrm{M}\otimes_{\mathrm{A}}\mathrm{N}) = \operatorname{length}_{\mathrm{A}}(\mathrm{M})\cdot\operatorname{length}_{\mathrm{B}}(\mathrm{N}/\mathfrak{m}\mathrm{N}).
\]
\end{env}

\begin{proof}
If $\mathrm{M}$ has infinite length, then so does
$\mathrm{M}\otimes_{\mathrm{A}}\mathrm{N}$. Indeed, every strictly
increasing chain of $n$ submodules $\mathrm{M}_i$ of $\mathrm{M}$
\textnormal{(}$1\leq i\leq n$\textnormal{)} yields pairwise distinct
sub-$\mathrm{B}$-modules
$\mathrm{M}_i\otimes_{\mathrm{A}}\mathrm{N}$ of
$\mathrm{M}\otimes_{\mathrm{A}}\mathrm{N}$. Since
$\mathrm{N}\neq\mathfrak{m}\mathrm{N}$, formula~\sref{IV.2.5.5.3}
holds in this case.

Suppose that $\mathrm{M}$ has finite length. If $\mathrm{M}=0$, both
sides of~\sref{IV.2.5.5.3} vanish, so assume $\mathrm{M}\neq0$. Each
$\mathfrak{m}^k\mathrm{M}$ has finite length, and if
$\mathfrak{m}^k\mathrm{M}\neq0$, Nakayama's lemma gives
$\mathfrak{m}^{k+1}\mathrm{M}\neq\mathfrak{m}^k\mathrm{M}$. Hence some
$r$ satisfies $\mathfrak{m}^r\mathrm{M}=0$. The modules
$\mathfrak{m}^k\mathrm{M}\otimes_{\mathrm{A}}\mathrm{N}$ identify with
sub-$\mathrm{B}$-modules of
$\mathrm{M}\otimes_{\mathrm{A}}\mathrm{N}$, and
\[
 \frac{\mathfrak{m}^k\mathrm{M}\otimes_{\mathrm{A}}\mathrm{N}}
      {\mathfrak{m}^{k+1}\mathrm{M}\otimes_{\mathrm{A}}\mathrm{N}}
 \simeq
 \frac{\mathfrak{m}^k\mathrm{M}}{\mathfrak{m}^{k+1}\mathrm{M}}
 \otimes_{\mathrm{A}/\mathfrak{m}}
 \frac{\mathrm{N}}{\mathfrak{m}\mathrm{N}}
\]
for $0\leq k\leq r-1$. Its length as a $\mathrm{B}$-module is the
product of
$\operatorname{length}_{\mathrm{B}}(\mathrm{N}/\mathfrak{m}\mathrm{N})$
and the dimension of the $(\mathrm{A}/\mathfrak{m})$-vector space
$\mathfrak{m}^k\mathrm{M}/\mathfrak{m}^{k+1}\mathrm{M}$; the latter
dimension is the length of that quotient as an $\mathrm{A}$-module.
Summing over $0\leq k\leq r-1$ gives~\sref{IV.2.5.5.3}.
\end{proof}

\begin{remark}[2.5.5.4]
\label{IV.2.5.5.4}
When $\mathrm{N}$ is an $\mathrm{A}$-module \emph{of finite type}, it is
faithfully flat over $\mathrm{A}$ if and only if it is flat and
$\mathrm{N}\neq0$; indeed, Nakayama's lemma then shows that
$\mathfrak{m}\mathrm{N}\neq\mathrm{N}$.
\end{remark}

\begin{lemma}[2.5.6]
\label{IV.2.5.6}
Let $\mathrm{B}$ be a not necessarily commutative ring, let $\mathrm{V}$
and $\mathrm{W}$ be isomorphic left $\mathrm{B}$-modules, set
\[
  \mathrm{C}=\operatorname{End}_{\mathrm{B}}(\mathrm{V}),\qquad
  \mathrm{M}=\operatorname{Hom}_{\mathrm{B}}(\mathrm{V},\mathrm{W}),
\]
and equip $\mathrm{M}$ with its canonical right
$\mathrm{C}$-module structure. Then $\mathrm{M}$ is isomorphic to the
right regular module $\mathrm{C}_d$. Moreover, for $u\in\mathrm{M}$,
the following conditions are equivalent:
\begin{enumerate}
  \item[\emph{a)}] $\{u\}$ is a basis of the $\mathrm{C}$-module
  $\mathrm{M}$;
  \item[\emph{b)}] $u$ is an isomorphism from $\mathrm{V}$ onto
  $\mathrm{W}$.
\end{enumerate}
\end{lemma}

\begin{proof}
If $u$ is an isomorphism from $\mathrm{V}$ onto $\mathrm{W}$, then
$v\mapsto u\circ v$ is a bijection from $\mathrm{C}$ to $\mathrm{M}$,
so \emph{b)} implies \emph{a)}. Conversely, suppose that $\{u\}$ is a
basis of $\mathrm{M}$. By hypothesis there is an isomorphism
$u':\mathrm{V}\to\mathrm{W}$, and $\{u'\}$ is also a basis of
$\mathrm{M}$. Hence $u=u'\circ w$ for some invertible $w\in\mathrm{C}$,
that is, for some automorphism $w$ of $\mathrm{V}$. Thus $u$ is an
isomorphism.
\end{proof}

\begin{corollary}[2.5.7]
\label{IV.2.5.7}
Under the hypotheses of~\sref{IV.2.5.6}, suppose in addition that one of
the following conditions holds:
\begin{enumerate}
  \item[(i)] $\mathrm{V}$ and $\mathrm{W}$ are Noetherian
  $\mathrm{B}$-modules;
  \item[(ii)] $\mathrm{V}$ and $\mathrm{W}$ are modules of finite
  presentation over a commutative subring of $\mathrm{B}$.
\end{enumerate}
Then conditions~\emph{a)} and~\emph{b)} of~\sref{IV.2.5.6} are also
equivalent to:
\begin{enumerate}
  \item[\emph{a}$'$)] $u$ is a generator of the $\mathrm{C}$-module
  $\mathrm{M}$;
  \item[\emph{b}$'$)] $u$ is an epimorphism of $\mathrm{V}$ onto $\mathrm{W}$.
\end{enumerate}
\end{corollary}

\begin{proof}
An epimorphism from an $\mathrm{A}$-module $\mathrm{E}$ onto itself is
bijective in either of the following cases: \(1^\circ\) $\mathrm{E}$ is
Noetherian (Bourbaki, \emph{Alg.}, chap.~VIII, \S~2, no.~2, lemma~3);
\(2^\circ\) $\mathrm{A}$ is commutative and $\mathrm{E}$ is of finite
presentation~\sref{IV.8.9.3}\footnote{The reader may verify that
\sref{IV.2.5.7} and \sref{IV.2.5.8} are not used before \S~9.}.
  Therefore \emph{b)} and \emph{b}$'$) are equivalent. If $u$ generates
  $\mathrm{M}$,
  \oldpage[IV-2]{25}
  and $\{u'\}$ is a basis, there is $v\in\mathrm{C}$ such
  that $u'=u\circ v$, so $u$ is surjective. Hence \emph{a}$'$) implies
  \emph{b}$'$); and \emph{a)} plainly implies \emph{a}$'$). This proves
  the corollary.
\end{proof}

\begin{proposition}[2.5.8]
\label{IV.2.5.8}
Let $\mathrm{A}$ be a commutative semilocal ring, let $\mathrm{B}$ be an
$\mathrm{A}$-algebra \textnormal{(}not necessarily
commutative\textnormal{)}, and let $\mathrm{V}$ and $\mathrm{W}$ be
$\mathrm{B}$-modules. Let $\mathrm{A}'$ be a commutative
$\mathrm{A}$-algebra which is faithfully flat as an $\mathrm{A}$-module,
and set
\[
  \mathrm{B}'=\mathrm{B}\otimes_{\mathrm{A}}\mathrm{A}',\qquad
  \mathrm{V}'=\mathrm{V}\otimes_{\mathrm{A}}\mathrm{A}',\qquad
  \mathrm{W}'=\mathrm{W}\otimes_{\mathrm{A}}\mathrm{A}'.
\]
Thus $\mathrm{B}'$ is an $\mathrm{A}'$-algebra and $\mathrm{V}'$,
$\mathrm{W}'$ are $\mathrm{B}'$-modules. Suppose in addition that one of
the following conditions holds:
\begin{enumerate}
  \item[(i)] $\mathrm{A}$ and $\mathrm{A}'$ are Noetherian, and
  $\mathrm{V}$ and $\mathrm{W}$ are $\mathrm{A}$-modules of finite type;
  \item[(ii)] $\mathrm{B}$ is an $\mathrm{A}$-module of finite type, and
  $\mathrm{V}$ is both a projective $\mathrm{B}$-module of finite type
  and an $\mathrm{A}$-module of finite presentation.
\end{enumerate}
If $\mathrm{V}'$ and $\mathrm{W}'$ are isomorphic as
$\mathrm{B}'$-modules, then $\mathrm{V}$ and $\mathrm{W}$ are isomorphic
as $\mathrm{B}$-modules.
\end{proposition}

\begin{proof}
In case~(ii), $\mathrm{W}'$, being isomorphic to $\mathrm{V}'$, is an
$\mathrm{A}'$-module of finite type. It follows that $\mathrm{W}$ is an
$\mathrm{A}$-module of finite type (Bourbaki, \emph{Alg.\ comm.},
chap.~I\textsuperscript{er}, \S~3, no.~6, prop.~11). Thus in both cases
$\mathrm{V}$ and $\mathrm{W}$ are $\mathrm{A}$-modules of finite type.

\begin{env}[2.5.8.1]
\label{IV.2.5.8.1}
\emph{Under either hypothesis} (i) \emph{or} (ii),
$\operatorname{Hom}_{\mathrm{B}}(\mathrm{V},\mathrm{W})$
\emph{is an $\mathrm{A}$-module of finite type.}
\end{env}

This is immediate in case~(i), because
$\operatorname{Hom}_{\mathrm{A}}(\mathrm{V},\mathrm{W})$ is then an
$\mathrm{A}$-module of finite type and
$\operatorname{Hom}_{\mathrm{B}}(\mathrm{V},\mathrm{W})$ is an
$\mathrm{A}$-submodule. In case~(ii), $\mathrm{V}$ is a direct factor
of a free $\mathrm{B}$-module $\mathrm{B}^n$. Consequently,
$\operatorname{Hom}_{\mathrm{B}}(\mathrm{V},\mathrm{W})$ is a direct
factor of
\[
  \operatorname{Hom}_{\mathrm{B}}(\mathrm{B}^n,\mathrm{W})
  =\mathrm{W}^n.
\]
It is therefore of finite type over $\mathrm{A}$.

Set
\[
  \mathrm{C}=\operatorname{End}_{\mathrm{B}}(\mathrm{V}),\qquad \mathrm{M}=\operatorname{Hom}_{\mathrm{B}}(\mathrm{V},\mathrm{W}).
\]
These are $\mathrm{A}$-modules of finite type in both cases. Under
either hypothesis, the canonical homomorphism

\begin{env}[2.5.8.2]
\label{IV.2.5.8.2}
\[
  \operatorname{Hom}_{\mathrm{A}}(\mathrm{V},\mathrm{W})\otimes_{\mathrm{A}}\mathrm{A}' \to \operatorname{Hom}_{\mathrm{A}'}(\mathrm{V}',\mathrm{W}')
\]
\end{env}

is bijective (Bourbaki, \emph{Alg.\ comm.}, chap.~II, \S~2, no.~10,
prop.~11). Since $\mathrm{A}'$ is flat over $\mathrm{A}$,
\[
  \operatorname{Hom}_{\mathrm{B}}(\mathrm{V},\mathrm{W})
  \otimes_{\mathrm{A}}\mathrm{A}'
\]
identifies canonically with a sub-$\mathrm{A}'$-module of
$\operatorname{Hom}_{\mathrm{A}}(\mathrm{V},\mathrm{W})
\otimes_{\mathrm{A}}\mathrm{A}'$. Its image under~\sref{IV.2.5.8.2}
lies in $\operatorname{Hom}_{\mathrm{B}'}(\mathrm{V}',\mathrm{W}')$:
if $u\in\operatorname{Hom}_{\mathrm{B}}(\mathrm{V},\mathrm{W})$ and
$a'\in\mathrm{A}'$, the image of $u\otimes a'$ is the homomorphism
$u':\mathrm{V}'\to\mathrm{W}'$ satisfying
$u'(x\otimes1)=u(x)\otimes a'$, and
\[
  u'\bigl((b\otimes1)(x\otimes1)\bigr)
  =u(bx)\otimes a'
  =bu(x)\otimes a'
  =(b\otimes1)\bigl(u(x)\otimes a'\bigr).
\]

\begin{env}[2.5.8.3]
\label{IV.2.5.8.3}
\emph{Under either hypothesis} (i) \emph{or} (ii), \emph{the
homomorphism}
\end{env}

\begin{env}[2.5.8.4]
\label{IV.2.5.8.4}
\[
  \operatorname{Hom}_{\mathrm{B}}(\mathrm{V},\mathrm{W})\otimes_{\mathrm{A}}\mathrm{A}' \to \operatorname{Hom}_{\mathrm{B}'}(\mathrm{V}',\mathrm{W}')
\]
\end{env}

\emph{is bijective.}

For $b\in\mathrm{B}$, let $h(b)$ \textnormal{(}resp.\ $h'(b)$\textnormal{)}
denote the $\mathrm{A}$-endomorphism $x\mapsto bx$ of $\mathrm{V}$
\textnormal{(}resp.\ $\mathrm{W}$\textnormal{)}. Let
$(b_\alpha)_{\alpha\in\mathrm{I}}$ be a system of generators of the
$\mathrm{A}$-algebra $\mathrm{B}$. The map
\[
  u\mapsto (h'(b_\alpha)\circ u - u\circ h(b_\alpha))_{\alpha\in\mathrm{I}}
\]
from $\operatorname{Hom}_{\mathrm{A}}(\mathrm{V},\mathrm{W})$ to
$(\operatorname{Hom}_{\mathrm{A}}(\mathrm{V},\mathrm{W}))^{\mathrm{I}}$
is $\mathrm{A}$-linear, and its kernel is precisely
$\operatorname{Hom}_{\mathrm{B}}(\mathrm{V},\mathrm{W})$. Thus
\[
  0 \to \operatorname{Hom}_{\mathrm{B}}(\mathrm{V},\mathrm{W}) \to \operatorname{Hom}_{\mathrm{A}}(\mathrm{V},\mathrm{W}) \to (\operatorname{Hom}_{\mathrm{A}}(\mathrm{V},\mathrm{W}))^{\mathrm{I}}.
\]

The same reasoning applies after replacing $\mathrm{A}$, $\mathrm{B}$,
$\mathrm{V}$, $\mathrm{W}$ by $\mathrm{A}'$, $\mathrm{B}'$,
$\mathrm{V}'$, $\mathrm{W}'$. Moreover, there is a commutative diagram

\begin{env}[2.5.8.5]
\label{IV.2.5.8.5}
\[
\xymatrix{
  0 \ar[r] & \operatorname{Hom}_{\mathrm{B}}(\mathrm{V},\mathrm{W})\otimes_{\mathrm{A}}\mathrm{A}' \ar[r] \ar[d]^{r} & \operatorname{Hom}_{\mathrm{A}}(\mathrm{V},\mathrm{W})\otimes_{\mathrm{A}}\mathrm{A}' \ar[r] \ar[d]^{s} & (\operatorname{Hom}_{\mathrm{A}}(\mathrm{V},\mathrm{W}))^{\mathrm{I}}\otimes_{\mathrm{A}}\mathrm{A}' \ar[d]^{t} \\
  0 \ar[r] & \operatorname{Hom}_{\mathrm{B}'}(\mathrm{V}',\mathrm{W}') \ar[r] & \operatorname{Hom}_{\mathrm{A}'}(\mathrm{V}',\mathrm{W}') \ar[r] & (\operatorname{Hom}_{\mathrm{A}'}(\mathrm{V}',\mathrm{W}'))^{\mathrm{I}}
}
\]
\end{env}

Here $r$ is~\sref{IV.2.5.8.4}, $s$ is~\sref{IV.2.5.8.2}, and $t$ is
the composite
\[
  (\operatorname{Hom}_{\mathrm{A}}(\mathrm{V},\mathrm{W}))^{\mathrm{I}}\otimes_{\mathrm{A}}\mathrm{A}' \xrightarrow{w} (\operatorname{Hom}_{\mathrm{A}}(\mathrm{V},\mathrm{W})\otimes_{\mathrm{A}}\mathrm{A}')^{\mathrm{I}} \xrightarrow{s^{\mathrm{I}}} (\operatorname{Hom}_{\mathrm{A}'}(\mathrm{V}',\mathrm{W}'))^{\mathrm{I}}
\]

where $w$ is the canonical homomorphism (Bourbaki, \emph{Alg.},
chap.~II, 3\textsuperscript{e}~ed., \S~3, no.~7). Since $\mathrm{A}'$
is flat over $\mathrm{A}$, the rows of~\sref{IV.2.5.8.5} are exact.

\oldpage[IV-2]{26}

The map $s$ is an isomorphism, hence so is $s^{\mathrm{I}}$. In case~(ii)
we may take $\mathrm{I}$ finite, and then $w$ is bijective (Bourbaki,
\emph{loc.\ cit.}, prop.~7). In case~(i), let $\mathrm{B}_1$
\textnormal{(}resp.\ $\mathrm{B}_2$\textnormal{)} be the image of
$\mathrm{B}$ under $h$ \textnormal{(}resp.\ $h'$\textnormal{)} in
$\operatorname{End}_{\mathrm{A}}(\mathrm{V})$
\textnormal{(}resp.\
$\operatorname{End}_{\mathrm{A}}(\mathrm{W})$\textnormal{)}. Both
images are $\mathrm{A}$-modules of finite type, so here too
$\mathrm{I}$ may be chosen finite. Thus $t$, and consequently $r$, is
bijective in both cases.

\begin{env}[2.5.8.6]
\label{IV.2.5.8.6}
It follows from~\sref{IV.2.5.8.4} that, if
\[
  \mathrm{C}'=\mathrm{C}\otimes_{\mathrm{A}}\mathrm{A}',\qquad \mathrm{M}'=\mathrm{M}\otimes_{\mathrm{A}}\mathrm{A}'
\]
then there are canonical bijections
\[
  \mathrm{C}' \xrightarrow{\;\sim\;} \operatorname{End}_{\mathrm{B}'}(\mathrm{V}'),\qquad \mathrm{M}' \xrightarrow{\;\sim\;} \operatorname{Hom}_{\mathrm{B}'}(\mathrm{V}',\mathrm{W}')
\]
the first an isomorphism of $\mathrm{A}'$-algebras and the second,
through the first, an isomorphism of right $\mathrm{C}'$-modules.
\end{env}

\begin{env}[2.5.8.7]
\label{IV.2.5.8.7}
\emph{Reduction to the case $\mathrm{V}=\mathrm{B}_s$.}
The hypothesis that $\mathrm{V}'$ and $\mathrm{W}'$ are isomorphic
$\mathrm{B}'$-modules implies that $\mathrm{C}'_d$ and $\mathrm{M}'$
are isomorphic right $\mathrm{C}'$-modules~\sref{IV.2.5.6}. To prove
\sref{IV.2.5.8}, it is enough to find $u\in\mathrm{M}$ which generates
$\mathrm{M}$ as a right $\mathrm{C}$-module. Indeed,
$u'=u\otimes1$ then generates $\mathrm{M}'$ as a right
$\mathrm{C}'$-module. In case~(i), $\mathrm{V}'$ and $\mathrm{W}'$ are
Noetherian $\mathrm{A}'$-modules; in case~(ii), they are isomorphic
$\mathrm{A}'$-modules of finite presentation. Corollary~\sref{IV.2.5.7}
therefore shows that $u'$ is a $\mathrm{B}'$-isomorphism
$\mathrm{V}'\to\mathrm{W}'$. Since $\mathrm{A}'$ is faithfully flat
over $\mathrm{A}$, $u$ is bijective~\sref[0]{0.6.4.1}.

In case~(i), $\mathrm{C}$ and $\mathrm{M}$ are $\mathrm{A}$-modules of
finite type. In case~(ii), $\mathrm{C}$ is a direct factor of
$\mathrm{V}^n$ by~\sref{IV.2.5.8.1}, hence a projective
$\mathrm{A}$-module of finite type. After changing notation, we are
therefore reduced to proving~\sref{IV.2.5.8} when
$\mathrm{V}=\mathrm{B}_s$; it is enough to prove that the
$\mathrm{B}$-module $\mathrm{W}$ has a generator. In this case
$\mathrm{B}$ is an $\mathrm{A}$-module of finite type.
\end{env}

\begin{env}[2.5.8.8]
\label{IV.2.5.8.8}
\emph{Reduction to the case where $\mathrm{A}$ and $\mathrm{A}'$ are
fields and $\mathrm{B}$ is a simple $\mathrm{A}$-algebra with centre
$\mathrm{A}$.}

Let $\mathfrak{r}$ be the radical of the semilocal ring $\mathrm{A}$.
It is enough to prove that
$\mathrm{W}/\mathfrak{r}\mathrm{W}$ is a cyclic
$(\mathrm{B}/\mathfrak{r}\mathrm{B})$-module. Indeed, a surjection
\[
  \mathrm{B}_s/\mathfrak{r}\mathrm{B}_s
  \longrightarrow \mathrm{W}/\mathfrak{r}\mathrm{W}
\]
composes with $\mathrm{B}_s\to
\mathrm{B}_s/\mathfrak{r}\mathrm{B}_s$ and, since $\mathrm{B}_s$ is
free as a $\mathrm{B}$-module, lifts to a map
$u:\mathrm{B}_s\to\mathrm{W}$. The displayed surjection is
$u\otimes1$ after tensoring over $\mathrm{A}$ with
$\mathrm{A}/\mathfrak{r}$. Since $\mathrm{W}$ is of finite type over
$\mathrm{A}$, Nakayama's lemma shows that $u$ is surjective (Bourbaki,
\emph{Alg.\ comm.}, chap.~II, \S~3, no.~2, cor.~1 of prop.~4).

Put
\[
\begin{aligned}
  \mathrm{A}_1&=\mathrm{A}/\mathfrak{r},&
  \mathrm{A}'_1&=\mathrm{A}'\otimes_{\mathrm{A}}\mathrm{A}_1
                    =\mathrm{A}'/\mathfrak{r}\mathrm{A}',\\
  \mathrm{B}_1&=\mathrm{B}/\mathfrak{r}\mathrm{B}
                    =\mathrm{B}\otimes_{\mathrm{A}}\mathrm{A}_1,&
  \mathrm{W}_1&=\mathrm{W}/\mathfrak{r}\mathrm{W}
                    =\mathrm{W}\otimes_{\mathrm{A}}\mathrm{A}_1.
\end{aligned}
\]
The hypotheses remain valid after replacing $\mathrm{A}$, $\mathrm{A}'$,
$\mathrm{B}$, $\mathrm{V}=\mathrm{B}_s$, and $\mathrm{W}$ by
$\mathrm{A}_1$, $\mathrm{A}'_1$, $\mathrm{B}_1$,
$\mathrm{V}_1=(\mathrm{B}_1)_s$, and $\mathrm{W}_1$. Moreover,
$\mathrm{A}'_1$ is faithfully flat over $\mathrm{A}_1$, and the
corresponding base changes of $\mathrm{V}_1$ and $\mathrm{W}_1$ remain
isomorphic. We may therefore suppose that $\mathrm{A}$ is a finite
product of fields.

The ring $\mathrm{B}$ is then Artinian. Let $\mathfrak{N}$ be its
radical. It is enough, by the same Nakayama argument, to prove that
$\mathrm{W}/\mathfrak{N}\mathrm{W}$ is cyclic over
$\mathrm{B}/\mathfrak{N}$. Since
\[
\begin{aligned}
  \mathrm{B}'/\mathfrak{N}\mathrm{B}'
    &=(\mathrm{B}/\mathfrak{N})\otimes_{\mathrm{A}}\mathrm{A}',\\
  \mathrm{W}'/\mathfrak{N}\mathrm{W}'
    &=(\mathrm{W}/\mathfrak{N}\mathrm{W})
      \otimes_{\mathrm{A}}\mathrm{A}',
\end{aligned}
\]
and the latter is isomorphic to the right regular module over the
former, we may further suppose that $\mathfrak{N}=0$, that is, that
$\mathrm{B}$ is a semisimple $\mathrm{A}$-algebra.

Write $\mathrm{A}=\prod_{i=1}^n k_i$. Then
$\mathrm{A}'=\prod_{i=1}^n\mathrm{A}'_i$, where each
$\mathrm{A}'_i$ is a $k_i$-algebra and is annihilated by the other
factors $k_j$. Faithful flatness implies
$\mathrm{A}'_i\neq0$. Choose a field quotient
$k'_i$ of each $\mathrm{A}'_i$ and set
$\mathrm{A}''=\prod_i k'_i$. Then $\mathrm{A}''$ is faithfully flat
over $\mathrm{A}$ and is a quotient of $\mathrm{A}'$. After the further
base change,
\[
  \mathrm{B}''=\mathrm{B}\otimes_{\mathrm{A}}\mathrm{A}'',
  \qquad
  \mathrm{W}''=\mathrm{W}\otimes_{\mathrm{A}}\mathrm{A}'',
\]
the module $\mathrm{W}''$ remains isomorphic to
$\mathrm{B}''_s$. Hence we may also suppose that $\mathrm{A}'$ is a
finite product of fields.

Let $\mathrm{Z}$ be the centre of $\mathrm{B}$. It is a finite product
of fields and an $\mathrm{A}$-module of finite type; $\mathrm{B}$ and
$\mathrm{W}$ are finite, hence projective, $\mathrm{Z}$-modules. Put
$\mathrm{Z}'=\mathrm{Z}\otimes_{\mathrm{A}}\mathrm{A}'$. Then
\[
  \mathrm{B}'=\mathrm{B}\otimes_{\mathrm{Z}}\mathrm{Z}',
  \qquad
  \mathrm{W}'=\mathrm{W}\otimes_{\mathrm{Z}}\mathrm{Z}',
\]
and $\mathrm{Z}'$ is faithfully flat over $\mathrm{Z}$. We may replace
$\mathrm{A}$ by $\mathrm{Z}$ and thus suppose that $\mathrm{A}$ is the
centre of $\mathrm{B}$.

If $k_i$ are the field factors of $\mathrm{A}$, then
$\mathrm{B}=\prod_i\mathrm{B}_i$, with $\mathrm{B}_i$ a simple ring of
centre $k_i$, and
$\mathrm{W}=\bigoplus_i\mathrm{W}_i$, where the other factors
$\mathrm{B}_j$ annihilate $\mathrm{W}_i$. We may suppose likewise that
$\mathrm{A}'=\prod_i k'_i$, where $k'_i/k_i$ is a field extension.
For every $i$,
\[
  (\mathrm{B}_i\otimes_{k_i}k'_i)_s
  \quad\text{and}\quad
  \mathrm{W}_i\otimes_{k_i}k'_i
\]
are isomorphic modules over $\mathrm{B}_i\otimes_{k_i}k'_i$.
It therefore suffices to treat $n=1$: $\mathrm{A}$ and $\mathrm{A}'$
are fields, and $\mathrm{B}$ is a simple algebra with centre
$\mathrm{A}$.
\end{env}

\begin{env}[2.5.8.9]
\label{IV.2.5.8.9}
\emph{End of the proof.} Every $\mathrm{B}$-module is a direct sum of
modules isomorphic to a minimal ideal of $\mathrm{B}$ (Bourbaki,
\emph{Alg.}, chap.~VIII, \S~5, props.~6 and~8). Thus two
$\mathrm{B}$-modules of finite rank over $\mathrm{A}$ are isomorphic
if and only if they have the same rank over $\mathrm{A}$. By hypothesis,
\[
  [\mathrm{W}':\mathrm{A}']=[\mathrm{B}'_s:\mathrm{A}'].
\]
But
\[
  [\mathrm{W}':\mathrm{A}']=[\mathrm{W}:\mathrm{A}],
  \qquad
  [\mathrm{B}'_s:\mathrm{A}']=[\mathrm{B}_s:\mathrm{A}],
\]
so $[\mathrm{W}:\mathrm{A}]=[\mathrm{B}_s:\mathrm{A}]$. This completes
the proof.
\end{env}
\end{proof}

\oldpage[IV-2]{27}

\subsection{Permanence of set-theoretic and topological properties of morphisms
by faithfully flat descent}
\label{subsection:IV.2.6}

\begin{proposition}[2.6.1]
\phantomsection
\label{IV.2.6.1}
Let $f:X\to Y$ be an $\mathrm{S}$-morphism of $\mathrm{S}$-preschemes,
and let $g:\mathrm{S}'\to\mathrm{S}$ be a surjective morphism. Set
\[
  X'=X_{(\mathrm{S}')},\qquad
  Y'=Y_{(\mathrm{S}')},\qquad
  f'=f_{(\mathrm{S}')}:X'\to Y'.
\]
Consider, for a morphism, the property of being:
\begin{enumerate}
  \item[(i)] surjective;
  \item[(ii)] injective;
  \item[(iii)] with finite fibres \textnormal{(}as sets\textnormal{)};
  \item[(iv)] bijective;
  \item[(v)] radicial.
\end{enumerate}
If $\mathbf{P}$ denotes any one of these properties and $f'$ has
property $\mathbf{P}$, then so does $f$.
\end{proposition}

\begin{proof}
The projection $Y'\to Y$ is itself surjective~\sref[I]{I.3.5.2};
hence, by~\sref[I]{I.3.3.11}, we may reduce to the case
$Y=\mathrm{S}$ and $Y'=\mathrm{S}'$. For $y\in Y$
\textnormal{(}resp.\ $y'\in Y'$\textnormal{)}, denote by $X_y$
\textnormal{(}resp.\ $X'_{y'}$\textnormal{)} the fibre prescheme
$f^{-1}(y)$ \textnormal{(}resp.\ $f'^{-1}(y')$\textnormal{)}
\sref[I]{I.3.6.2}. If $y'\in Y'$ and $y=g(y')$, there is a canonical
isomorphism
\[
  X'_{y'}\xrightarrow{\;\sim\;}
  X_y\otimes_{k(y)}k(y')
\]
\sref[I]{I.3.6.4}. Since
$\Spec(k(y'))\to\Spec(k(y))$ is surjective, so is the projection
$X'_{y'}\to X_y$~\sref[I]{I.3.5.2}. Therefore, if $X'_{y'}$ is
nonempty \textnormal{(}resp.\ has at most one point, resp.\ is a finite
set\textnormal{)}, the same is true of $X_y$. Since $g:Y'\to Y$ is
surjective, this proves the assertion in cases~(i), (ii), and~(iii);
case~(iv) follows from~(i) and~(ii).

Finally, for~(v), it is enough to show that if $f'$ is universally
injective~\sref[I]{I.3.5.11}, then so is $f$. Let $Y_1\to Y$ be any
morphism and set
\[
  X_1=X\times_Y Y_1,\qquad f_1=f_{(Y_1)}.
\]
Also set
\[
\begin{aligned}
  Y'_1&=Y_1\times_Y Y',\\
  X'_1&=X'\times_{Y'}Y'_1=X_1\times_{Y_1}Y'_1,\\
  f'_1&=f'_{(Y'_1)}=(f_1)_{(Y'_1)}:X'_1\to Y'_1.
\end{aligned}
\]
The projection $g_1=g_{(Y_1)}:Y'_1\to Y_1$ is
surjective~\sref[I]{I.3.5.2}; and, since $f'$ is universally injective,
$f'_1$ is injective. Case~(ii) then shows that $f_1$ is injective,
which proves the assertion.
\end{proof}

\begin{proposition}[2.6.2]
\phantomsection
\label{IV.2.6.2}
With the notation of~\sref{IV.2.6.1}, suppose that
$g:\mathrm{S}'\to\mathrm{S}$ is faithfully flat and quasi-compact.
Consider, for a morphism, the property of being:
\begin{enumerate}
  \item[(i)] open;
  \item[(ii)] closed;
  \item[(iii)] quasi-compact and a homeomorphism onto its image
    subspace;
  \item[(iv)] a homeomorphism.
\end{enumerate}
If $\mathbf{P}$ denotes any one of these properties and $f'$ has
property $\mathbf{P}$, then so does $f$.
\end{proposition}

\oldpage[IV-2]{28}

\begin{proof}
The morphism $Y'\to Y$ is faithfully flat and
quasi-compact~\sref{IV.2.2.13} and~\sref{IV.1.1.2}; hence, by
\sref[I]{I.3.3.11}, we may reduce to the case $Y=\mathrm{S}$,
$Y'=\mathrm{S}'$. Let $g':X'\to X$ be the projection. For every subset
$M$ of $X$ we have
\[
  g^{-1}\bigl(f(M)\bigr)=f'\bigl(g'^{-1}(M)\bigr)
\]
\sref[I]{I.3.4.8}. If $f'$ is open \textnormal{(}resp.\ closed\textnormal{)},
then, for every open \textnormal{(}resp.\ closed\textnormal{)} subset
$M$ of $X$, the subset $f'(g'^{-1}(M))$ is open
\textnormal{(}resp.\ closed\textnormal{)} in $Y'$. Since $g$ is
faithfully flat and quasi-compact, \sref{IV.2.3.12} implies that $f(M)$
is open \textnormal{(}resp.\ closed\textnormal{)} in $Y$. This proves
cases~(i) and~(ii).

Let us prove~(iii), which also yields~(iv) in view of
\sref{IV.2.6.1}[(iv)]. By~\sref{IV.2.6.1}[(ii)], $f$ is injective; it
remains to prove that $f$, considered as a map from $X$ onto $f(X)$,
is quasi-compact and open. Since $f'$ is quasi-compact, so is
$f$~\sref{IV.1.1.4}. It is therefore enough to prove, for every closed
subset $Z$ of $X$, that
\[
  Z=f^{-1}\bigl(\overline{f(Z)}\bigr).
\]
Since $g'$ is surjective, this is equivalent to
\[
\begin{aligned}
  g'^{-1}(Z)
    &=g'^{-1}\!\left(f^{-1}\bigl(\overline{f(Z)}\bigr)\right)\\
    &=f'^{-1}\!\left(g^{-1}\bigl(\overline{f(Z)}\bigr)\right).
\end{aligned}
\]
The composite of the canonical injection $Z\to X$ with $f$ is
quasi-compact, where $Z$ denotes here the reduced closed subprescheme of
$X$ with underlying space $Z$. Applying~\sref{IV.2.3.10} to the subset
$f(Z)$ of $Y$, the image of the morphism $f|_Z$, gives
\[
  g^{-1}\bigl(\overline{f(Z)}\bigr)
  =
  \overline{g^{-1}\bigl(f(Z)\bigr)}.
\]
If $Z'=g'^{-1}(Z)$, the required formula is consequently equivalent to
\[
  Z'=f'^{-1}\bigl(\overline{f'(Z')}\bigr),
\]
which follows from the hypothesis that $f'$ is a homeomorphism from
$X'$ onto $f'(X')$.
\end{proof}

\begin{remark}[2.6.3]
\phantomsection
\label{IV.2.6.3}
In cases~(i) and~(ii), the conclusions of~\sref{IV.2.6.2} remain valid
if $g$ is assumed only \emph{quasi-faithfully flat}~\sref{IV.2.3.3}
and quasi-compact. Indeed, by~\sref{IV.2.1.4}, \sref[I]{I.3.5.2},
and~\sref[I]{I.9.1.13.1}, we may again reduce to the case
$Y=\mathrm{S}$, $Y'=\mathrm{S}'$, and the conclusion follows from
\sref{IV.2.3.12}. In cases~(iii) and~(iv), the conclusions remain valid
if $g$ is assumed only quasi-faithfully flat, provided that $f$ is
already assumed quasi-compact; the proof then uses only
\sref{IV.2.3.10} and the surjectivity of $g$.

Finally, if $g$ is faithfully flat and locally of finite presentation,
or if $g$ is surjective and $\mathrm{S}$ is discrete, the conclusion of
\sref{IV.2.6.2} remains valid when $\mathbf{P}$ is the property
\begin{enumerate}
  \item[(iii \emph{bis})] of being a homeomorphism onto its image
    subspace;
\end{enumerate}
as follows from the proof of~\sref{IV.2.6.2} and
Remark~\sref{IV.2.4.11}.
\end{remark}

\begin{corollary}[2.6.4]
\phantomsection
\label{IV.2.6.4}
With the notation of~\sref{IV.2.6.1}, suppose that
$g:\mathrm{S}'\to\mathrm{S}$ is faithfully flat and quasi-compact.
Consider, for a morphism, the property of being:
\begin{enumerate}
  \item[(i)] universally open;
  \item[(ii)] universally closed;
  \item[(iii)] quasi-compact and universally bicontinuous;
  \item[(iv)] a universal homeomorphism;
  \item[(v)] quasi-compact;
  \item[(vi)] quasi-compact and dominant.
\end{enumerate}
If $\mathbf{P}$ denotes any one of these properties, then $f$ has
property $\mathbf{P}$ if and only if $f'$ does.
\end{corollary}

\begin{proof}
Properties~(v) and~(vi) are only given for the record, being consequences
of~\sref{IV.1.1.4}, \sref{IV.1.1.6}, and~\sref{IV.2.3.7}. For the
other properties, necessity follows from~\sref{IV.2.4.3}. Conversely,
suppose, for example, that $f'$ is universally open, and let
$Y_1\to Y$ be any morphism. Set
\[
  X_1=X\times_Y Y_1,\qquad f_1=f_{(Y_1)}.
\]
Also set
\[
\begin{aligned}
  Y'_1&=Y_1\times_Y Y',\\
  X'_1&=X'\times_{Y'}Y'_1=X_1\times_{Y_1}Y'_1,\\
  f'_1&=f'_{(Y'_1)}=(f_1)_{(Y'_1)}:X'_1\to Y'_1.
\end{aligned}
\]
The projection $g_1=g_{(Y_1)}:Y'_1\to Y_1$ is faithfully flat and
quasi-compact~\sref{IV.2.2.13} and~\sref{IV.1.1.2}. Since $f'$ is
universally open, $f'_1$ is open, and~\sref{IV.2.6.2} then shows that
$f_1$ is open. Hence $f$ is universally open. The same argument applies
in the other cases.

One may again replace ``faithfully flat'' by ``quasi-faithfully flat''.
Moreover, when $g$ is locally of finite presentation, or when $g$ is
surjective and $\mathrm{S}$ is discrete, property~(iii) may be replaced
by
\begin{enumerate}
  \item[(iii \emph{bis})] universally bicontinuous.
\end{enumerate}
\end{proof}

\oldpage[IV-2]{29}

\subsection{Permanence of various properties of morphisms under faithfully flat descent}
\label{subsection:IV.2.7}

\begin{proposition}[2.7.1]
\phantomsection
\label{IV.2.7.1}
Let $f:X\to Y$ be an $\mathrm{S}$-morphism of $\mathrm{S}$-preschemes,
and let $g:\mathrm{S}'\to\mathrm{S}$ be faithfully flat and
quasi-compact. Set
\[
  X'=X_{(\mathrm{S}')},\qquad
  Y'=Y_{(\mathrm{S}')},\qquad
  f'=f_{(\mathrm{S}')}:X'\to Y'.
\]
Consider, for a morphism, the property of being:
\begin{enumerate}
  \item[(i)] separated;
  \item[(ii)] quasi-separated;
  \item[(iii)] locally of finite type;
  \item[(iv)] locally of finite presentation;
  \item[(v)] of finite type;
  \item[(vi)] of finite presentation;
  \item[(vii)] proper;
  \item[(viii)] \phantomsection\label{IV.2.7.1.viii}an isomorphism;
  \item[(ix)] a monomorphism;
  \item[(x)] an open immersion;
  \item[(xi)] a quasi-compact immersion;
  \item[(xii)] a closed immersion;
  \item[(xiii)] affine;
  \item[(xiv)] quasi-affine;
  \item[(xv)] finite;
  \item[(xvi)] quasi-finite;
  \item[(xvii)] integral.
\end{enumerate}
If $\mathbf{P}$ denotes any one of these properties, then $f$ has
property $\mathbf{P}$ if and only if $f'$ does.
\end{proposition}

\begin{proof}
It was proved in Chapters~I and~II and in Chapter~IV, \S~1, that if
$f$ has one of the properties $\mathbf{P}$ above, then so does $f'$
\textnormal{(}without any hypothesis on $g:\mathrm{S}'\to\mathrm{S}$\textnormal{)}.
It remains to prove the converse. Since the projection $Y'\to Y$ is

\oldpage[IV-2]{30}

a faithfully flat and quasi-compact morphism~\sref{IV.2.2.13}
and~\sref{IV.1.1.2}, we may, by~\sref[I]{I.3.3.11}, reduce to the
case $\mathrm{S}=Y$, $\mathrm{S}'=Y'$.

\textnormal{(i)} To say that $f$ is separated means that its diagonal
$\Delta_f:X\to X\times_Y X$ is closed. Since
$\Delta_{f'}=(\Delta_f)_{(Y')}$~\sref[I]{I.5.3.4}, if
$\Delta_{f'}$ is closed, then so is $\Delta_f$ by~\sref{IV.2.6.2};
hence $f$ is separated.

\textnormal{(ii)} This was already proved under weaker
hypotheses~\sref{IV.1.2.5}.

\textnormal{(iii)} and~\textnormal{(iv)}. The question is plainly local
on $X$ and $Y$ and, in view of~\sref{IV.2.2.12}, it is therefore enough
to prove the following lemma.

\begin{lemma}[2.7.1.1]
\phantomsection
\label{IV.2.7.1.1}
Let $A$ be a ring, $B$ an $A$-algebra, and $A'$ an $A$-algebra that
is faithfully flat as an $A$-module; put $B'=B\otimes_A A'$. Then $B$
is an $A$-algebra of finite type \textnormal{(}resp.\ of finite
presentation\textnormal{)} if and only if $B'$ is an $A'$-algebra of
finite type \textnormal{(}resp.\ of finite presentation\textnormal{)}.
\end{lemma}

Necessity is already known without any hypothesis on
$A'$~\sref{IV.1.3.4}, \sref{IV.1.3.6}, \sref{IV.1.4.3},
and~\sref{IV.1.4.6}. Suppose that $B'$ is an $A'$-algebra of finite
type. Let $(B_\alpha)_{\alpha\in I}$ be the filtered increasing family
of finite-type $A$-subalgebras of $B$. Then
\[
  B=\varinjlim_\alpha B_\alpha,\qquad
  B'=\varinjlim_\alpha(B_\alpha\otimes_A A'),
\]
because tensor product commutes with inductive limits. If
$(x'_i)$ is a finite set of generators of the $A'$-algebra $B'$, there
is an index $\alpha$ such that every $x'_i$ belongs to the subalgebra
$B_\alpha\otimes_A A'$ of $B'$. Hence
$B'=B_\alpha\otimes_A A'$, and faithful flatness gives
$B=B_\alpha$~\sref[0]{0.6.4.1}.

Now suppose that $B'$ is an $A'$-algebra of finite presentation. What
precedes shows that $B$ is an $A$-algebra of finite type. Thus there
are a polynomial $A$-algebra $C=A[T_1,\ldots,T_m]$ and a surjective
$A$-homomorphism of algebras $C\to B$. Let $\mathfrak{J}$ be its
kernel. We have an exact sequence
\[
  0\to\mathfrak{J}\to C\to B\to0,
\]
and, since $A'$ is flat over $A$, an exact sequence
\[
  0\to\mathfrak{J}'\to C'\to B'\to0,
\]
where
\[
  C'=C\otimes_A A'=A'[T_1,\ldots,T_m],
  \qquad
  \mathfrak{J}'=\mathfrak{J}\otimes_A A'
\]
and $\mathfrak{J}'$ is identified with an ideal of $C'$. Since $B'$ is
of finite presentation over $A'$, the $C'$-module $\mathfrak{J}'$ is
of finite type~\sref{IV.1.4.4}. But
$\mathfrak{J}'=\mathfrak{J}\otimes_C C'$, and $C'$ is faithfully flat
over $C$~\sref{IV.2.2.13} and~\sref{IV.2.2.3}. It follows that
$\mathfrak{J}$ is a $C$-module of finite type
\textnormal{(}Bourbaki, \emph{Alg\`ebre commutative}, Chap.~I, \S~3,
no.~6, Prop.~11\textnormal{)}. Hence $B$ is an $A$-algebra of finite
presentation.

\textnormal{(v)} This follows from~\textnormal{(iii)} and
\sref{IV.2.6.4}[(v)], by~\sref{IV.1.5.2}.

\textnormal{(vi)} This follows similarly from~\textnormal{(iv)},
\textnormal{(v)}, and~\textnormal{(ii)}, by~\sref{IV.1.6.1}.

\textnormal{(vii)} This follows from~\textnormal{(i)},
\textnormal{(v)}, and~\sref{IV.2.6.4}[(ii)], by~\sref[II]{II.5.4.1}.

\textnormal{(viii)} Since $f'$ is an isomorphism, it is a universal
homeomorphism; hence so is $f$~\sref{IV.2.6.4}. In particular, $f$ is
quasi-compact and separated~\sref{IV.2.4.4}. Write
$f=(\psi,\theta)$, where $\psi$ is therefore a homeomorphism. It remains
to prove that
\[
  \theta:\sh{O}_Y\to f_*(\sh{O}_X)
\]
is an isomorphism of $\sh{O}_Y$-modules. If
$f'=(\psi',\theta')$, then
\[
  \theta':\sh{O}_{Y'}\to f'_*(\sh{O}_{X'})
\]
is the composite of $g^*(\theta)$ with the canonical homomorphism
\[
  g^*\bigl(f_*(\sh{O}_X)\bigr)\longrightarrow f'_*(\sh{O}_{X'})
\]
\sref{IV.2.3.2}. The latter is bijective by~\sref{IV.2.3.1}. Thus,
if $\theta'$ is bijective, then so is $g^*(\theta)$; since $g$ is
faithfully flat, $\theta$ is bijective~\sref{IV.2.2.7}. This
proves~\textnormal{(viii)}.

\textnormal{(ix)} This follows from~\textnormal{(viii)},
\sref[I]{I.5.3.4}, and~\sref[I]{I.5.3.8}, which reduces
monomorphisms to isomorphisms.

\textnormal{(x)} If $f'$ is an open immersion, then $f'(X')$ is open
in $Y'$, and
\[
  f'(X')=g^{-1}\bigl(f(X)\bigr)
\]
\sref[I]{I.3.4.8}. It follows from~\sref{IV.2.3.12} that $f(X)$ is
open. Replacing $Y$ \textnormal{(}resp.\ $Y'$\textnormal{)} by the
subprescheme induced on $f(X)$ \textnormal{(}resp.\ $f'(X')$\textnormal{)},
and using~\sref{IV.1.1.2} and~\sref{IV.2.2.13}, we reduce to the case
where $f'$ is an isomorphism. Then $f$ is an isomorphism
by~\textnormal{(viii)}, proving~\textnormal{(x)}.

\oldpage[IV-2]{31}

\textnormal{(xi)} If $f'$ is a quasi-compact immersion, then it is a
quasi-compact and quasi-separated morphism~\sref{IV.1.2.2}; hence so
is $f$, by~\textnormal{(ii)} and~\sref{IV.2.6.4}[(v)]. Let $Z$ be the
closed subprescheme of $Y$ that is the closed image of $X$ under
$f$~\sref{IV.1.7.8}, and write $f=j\circ q$, where $j:Z\to Y$ is the
canonical injection. Then $f'=j'\circ q'$, with
$j'=j_{(Y')}$ and $q'=q_{(Y')}$. Moreover, $j'$ identifies with the
canonical injection $Z'\to Y'$ of the closed image $Z'$ of $X'$ under
$f'$~\sref{IV.2.3.2}. The hypothesis on $f'$ now says that $q'$ is an
open immersion~\sref[I]{I.9.5.10}; hence so is $q$,
by~\textnormal{(x)}, and therefore $f$ is an immersion.

\textnormal{(xii)} To say that $f$ \textnormal{(}resp.\ $f'$\textnormal{)}
is a closed immersion is to say that it is a quasi-compact immersion
and a closed morphism. Thus~\textnormal{(xii)} follows
from~\textnormal{(xi)} and~\sref{IV.2.6.2}[(ii)].

\textnormal{(xiii)} and~\textnormal{(xiv)}. Suppose that $f'$ is affine
\textnormal{(}resp.\ quasi-affine\textnormal{)}. Then $f'$ is
quasi-compact and quasi-separated~\sref[II]{II.5.1.1}, and so is $f$
by~\textnormal{(ii)} and~\sref{IV.2.6.4}[(v)]. Put
\[
  \sh{A}=f_*(\sh{O}_X),\qquad
  \sh{A}'=f'_*(\sh{O}_{X'}).
\]
By~\sref{IV.2.3.1}, the canonical homomorphism of
$\sh{O}_{Y'}$-algebras $g^*(\sh{A})\to\sh{A}'$ is bijective.
Consequently, if $h:Z=\Spec(\sh{A})\to Y$ is the structure morphism,
then the structure morphism $h':Z'=\Spec(\sh{A}')\to Y'$ identifies
with $h_{(Y')}$~\sref[II]{II.1.5.2}. Let $u:X\to Z$
\textnormal{(}resp.\ $u':X'\to Z'$\textnormal{)} be the canonical
$Y$-morphism \textnormal{(}resp.\ $Y'$-morphism\textnormal{)}
corresponding to the identity homomorphism of $\sh{A}$
\textnormal{(}resp.\ $\sh{A}'$\textnormal{)}~\sref[II]{II.1.2.7}.
We have a commutative diagram
\[
\xymatrix{
  X \ar[d]_u & X' \ar[l] \ar[d]^{u'} \\
  Z \ar[d]_h & Z' \ar[l]_{g'} \ar[d]^{h'} \\
  Y & Y' \ar[l]^g
}
\]
and $h'\circ u'=f'$. It follows from~\sref[II]{II.1.2.7} that
$u'=u_{(Z')}$. The projection $g':Z'\to Z$ is faithfully flat and
quasi-compact~\sref{IV.1.1.2} and~\sref{IV.2.2.13}. The hypothesis on
$f'$ says that $u'$ is an isomorphism
\textnormal{(}resp.\ an open immersion\textnormal{)}
\sref[II]{II.5.1.6}; by~\textnormal{(viii)}
\textnormal{(}resp.\ (x)\textnormal{)}, $u$ is an
isomorphism \textnormal{(}resp.\ an open immersion\textnormal{)}. This
proves~\textnormal{(xiii)} \textnormal{(}resp.\
\textnormal{(xiv)}\textnormal{)}.

\oldpage[IV-2]{32}

\textnormal{(xv)} If $f'$ is finite, it is affine; hence $f$ is affine
by~\textnormal{(xiii)}. With the notation used there, $\sh{A}'$ is an
$\sh{O}_{Y'}$-module of finite type and is isomorphic to
$g^*(\sh{A})$. It follows from~\sref{IV.2.5.2} that $\sh{A}$ is an
$\sh{O}_Y$-module of finite type, and therefore $f$ is finite.

\textnormal{(xvi)} To say that $f$ is quasi-finite is to say that $f$
is of finite type and that $f^{-1}(y)$ is finite for every
$y\in Y$~\sref[II]{II.6.2.2} and~\sref[I]{I.6.4.4}. The conclusion
therefore follows from~\textnormal{(v)} and~\textnormal{(xv)}, applied
to the fibres.

\textnormal{(xvii)} As in~\textnormal{(xv)}, one first sees that $f$ is
affine. We may reduce to the case
\[
  Y=\Spec(A),\qquad Y'=\Spec(A'),\qquad
  X=\Spec(B),\qquad X'=\Spec(B'),
\]
where $B'=B\otimes_A A'$. The ring $B$ is the inductive limit of its
finite-type $A$-subalgebras $B_\alpha$; hence
\[
  B'=\varinjlim_\alpha B'_\alpha,\qquad
  B'_\alpha=B_\alpha\otimes_A A',
\]
and each $B'_\alpha$ is a finite-type $A'$-algebra. By hypothesis,
$B'$ is integral over $A'$, so each $B'_\alpha$ is an $A'$-module of
finite type. Hence $B_\alpha$ is an $A$-module of finite
type~\sref{IV.2.5.2}, and is therefore integral over $A$. Since every
element of $B$ lies in some $B_\alpha$, the algebra $B$ is integral
over $A$.
\end{proof}

\begin{corollary}[2.7.2]
\phantomsection
\label{IV.2.7.2}
Under the hypotheses and notation of~\sref{IV.2.7.1}, suppose that $f$
is quasi-compact. Let $\sh{L}$ be an invertible $\sh{O}_X$-module and,
if $g':X'\to X$ is the projection, let
$\sh{L}'=g'^*(\sh{L})$ be its inverse image. Then $\sh{L}$ is ample
\textnormal{(}resp.\ very ample\textnormal{)} for $f$ if and only if
$\sh{L}'$ is ample \textnormal{(}resp.\ very ample\textnormal{)} for
$f'$.
\end{corollary}

\begin{proof}
Necessity holds without any hypothesis on
$g:\mathrm{S}'\to\mathrm{S}$~\sref[II]{II.4.4.10}
and~\sref[II]{II.4.6.13}. To prove sufficiency, we may, as
in~\sref{IV.2.7.1}, reduce to the case $\mathrm{S}=Y$,
$\mathrm{S}'=Y'$. The hypothesis on $\sh{L}'$ implies that $f'$ is
quasi-compact and separated~\sref[II]{II.4.6.1}; hence so is $f$,
by~\sref{IV.2.6.4}[(v)] and~\sref{IV.2.7.1}[(i)]. Put
\[
  \sh{E}=f_*(\sh{L}),\qquad
  \sh{E}'=f'_*(\sh{L}').
\]
By~\sref{IV.2.3.1}, the canonical homomorphism
\[
  u:g^*(\sh{E})\longrightarrow\sh{E}'
\]
is bijective. Suppose first that $\sh{L}'$ is very ample for $f'$. The
canonical homomorphism
$\sigma':f'^*(\sh{E}')\to\sh{L}'$ is then surjective, and
\[
  r'=r_{\sh{L}',\sigma'}:X'\longrightarrow\mathbf{P}(\sh{E}')
\]
is a necessarily quasi-compact immersion~\sref[II]{II.4.4.4}[(b)]
and~\sref{IV.1.1.2}[(v)]. If
\[
  h:\mathbf{P}(\sh{E})\to Y,\qquad
  h':\mathbf{P}(\sh{E}')\to Y'
\]
are the structure morphisms, the bijectivity of $u$ implies that $h'$
identifies with $h_{(Y')}$~\sref[II]{II.4.1.3}. On the other hand,
the projection $g':X'\to X$ is faithfully flat~\sref{IV.2.2.13}, and
$f\circ g'=g\circ f'$. Let
$\sigma:f^*(\sh{E})\to\sh{L}$ be the canonical homomorphism. One
readily verifies that $g'^*(\sigma)$ is the composite
\[
  f'^*\!\left(g^*\!\left(f_*(\sh{L})\right)\right)
    \xrightarrow{\,f'^*(u)\,}
  f'^*\!\left(f'_*(\sh{L}')\right)
    \xrightarrow{\,\sigma'\,}
  \sh{L}' .
\]
This may be checked, for example, by reducing to the case where $Y$
and $Y'$ are affine. Since $\sigma'$ is surjective and $f'^*(u)$ is
bijective, $g'^*(\sigma)$ is surjective; hence $\sigma$ is surjective
by~\sref{IV.2.2.7}. Thus the morphism
\[
  r=r_{\sh{L},\sigma}:X\longrightarrow\mathbf{P}(\sh{E})
\]
is everywhere defined~\sref[II]{II.3.7.4}. Put
$P=\mathbf{P}(\sh{E})$ and $P'=\mathbf{P}(\sh{E}')$, and let
$g'':P'\to P$ be the projection. Then $r'$ identifies with
$r_{(P')}$~\sref[II]{II.4.2.10}, while $g''$ is faithfully flat and
quasi-compact~\sref{IV.1.1.2} and~\sref{IV.2.2.13}.
By~\sref{IV.2.7.1}[(xi)], $r$ is an immersion, and therefore
$\sh{L}$ is very ample~\sref[II]{II.4.4.4}[(b)].

Now suppose that $\sh{L}'$ is ample for $f'$. To prove that $\sh{L}$
is ample for $f$, we may reduce to the case where $Y$ is
affine~\sref[II]{II.4.6.4}; by~\sref{IV.2.2.12}
and~\sref[II]{II.4.6.13}, we may also suppose that $Y'$ is affine.
Then $X$ and $X'$ are quasi-compact schemes, and we may apply the
criterion of~\sref[II]{II.4.6.8}[(c)]. Let $\sh{F}$ be a
quasi-coherent $\sh{O}_X$-module of finite type, and let
\[
  \sigma:
  f^*f_*(\sh{F}\otimes\sh{L}^{\otimes n})
  \longrightarrow
  \sh{F}\otimes\sh{L}^{\otimes n}
\]
be the canonical homomorphism. As above, $g'^*(\sigma)$ is the
composite
\[
  f'^*\!\left(g^*f_*(\sh{F}\otimes\sh{L}^{\otimes n})\right)
    \xrightarrow{\,f'^*(u)\,}
  f'^*f'_*(\sh{F}'\otimes\sh{L}'^{\otimes n})
    \xrightarrow{\,\sigma'\,}
  \sh{F}'\otimes\sh{L}'^{\otimes n},
\]
where $\sh{F}'=g'^*(\sh{F})$, using~\sref[0]{0.4.3.3.1}, and where
$u$ is the canonical homomorphism
\[
  g^*f_*(\sh{F}\otimes\sh{L}^{\otimes n})
  \longrightarrow
  f'_*(\sh{F}'\otimes\sh{L}'^{\otimes n}).
\]
The homomorphism $u$ is bijective for every $n$ by~\sref{IV.2.3.1}.
Moreover, $\sh{F}'$ is quasi-coherent and of finite type, so the
ampleness of $\sh{L}'$ for $f'$ gives an $n_0$ such that $\sigma'$ is
surjective for $n\geq n_0$. Thus $g'^*(\sigma)$ is surjective for
$n\geq n_0$; since $g'$ is faithfully flat, $\sigma$ is surjective for
these values of $n$~\sref{IV.2.2.7}. This completes the proof.
\end{proof}

\oldpage[IV-2]{33}

\begin{remarks}[2.7.3]
\phantomsection
\label{IV.2.7.3}
\begin{enumerate}
  \item[(i)] It follows from~\sref{IV.2.6.1}, \sref{IV.2.6.4},
  and~\sref{IV.2.5.4.1} that the conclusions of~\sref{IV.2.7.1}
  remain valid in cases~(i), (iii), (v), (vii), and~(xvi) when $g$ is
  assumed only quasi-compact and \emph{quasi-faithfully
  flat}~\sref{IV.2.3.3}. We already noted that case~(ii)
  of~\sref{IV.2.7.1} is valid under the sole assumptions that $g$ is
  surjective and quasi-compact.

  \item[(ii)] Under the hypotheses and notation of~\sref{IV.2.7.1},
  it can happen that $f$ is proper and $f'$ is projective while $f$ is
  not quasi-projective. Indeed, Hironaka~\cite{IV-34} gave an example
  of a proper nonprojective morphism $f:X\to Y$, where $X$ and $Y$ are
  regular algebraic schemes~\sref[0]{0.4.1.4} over the same field
  $k$, with $Y$ projective over $k$. Moreover, $Y$ is the union of two
  affine open subsets $Y_i$ \textnormal{(}$i=1,2$\textnormal{)} such
  that
  \[
    f_i:X\times_Y Y_i\longrightarrow Y_i
  \]
  is projective for $i=1,2$. Let $Y'=Y_1\amalg Y_2$ be their
  coproduct as preschemes. The canonical morphism $g:Y'\to Y$, whose
  restrictions are the canonical injections of $Y_1$ and $Y_2$ into
  $Y$, is faithfully flat and is quasi-compact by~\sref[I]{I.5.5.10}.
  Yet $f':X\times_Y Y'\to Y'$, whose restriction to each $Y_i$ is
  $f_i$, is projective~\sref[II]{II.5.5.6}, whereas $f$ is not.
  Consequently there exists an invertible $\sh{O}_{X'}$-module
  $\sh{L}'$ that is $f'$-ample but is not of the form
  $g'^*(\sh{L})$ for an invertible $\sh{O}_X$-module $\sh{L}$, where
  $g':X'\to X$ is the projection, by~\sref{IV.2.7.2}.

  \item[(iii)] Under the hypotheses of~\sref{IV.2.7.1}, it can happen
  that $f'$ is a local isomorphism while $f$ is not a local immersion.
  Indeed, let $k$ be a field, let $\bar{k}$ be an algebraic closure of
  $k$, and let $K\neq k$ be a finite separable extension of $k$. The
  structure morphism $f:X\to Y$, where $X=\Spec(K)$ and
  $Y=\Spec(k)$, is not a local immersion. But if
  $Y'=\Spec(\bar{k})$, then $Y'\to Y$ is faithfully flat and
  quasi-compact, and $f'$ is a local isomorphism, since
  $X'=X\times_Y Y'$ is a coproduct of finitely many schemes isomorphic
  to $Y'$.
\end{enumerate}
\end{remarks}

\subsection{Preschemes over a regular base of dimension
\texorpdfstring{$1$}{1}; closure of a closed subprescheme of the generic
fibre}
\label{subsection:IV.2.8}

\begin{proposition}[2.8.1]
\phantomsection
\label{IV.2.8.1}
Let $Y$ be a locally Noetherian, regular, irreducible prescheme of
dimension $1$, with generic point $\eta$. Let $f:X\to Y$ be a
morphism, let $X_\eta=f^{-1}(\eta)$ be the fibre at the generic point,

\oldpage[IV-2]{34}

and let $i:X_\eta\to X$ be the canonical morphism. Let $\sh{F}$ be a
quasi-coherent $\sh{O}_X$-module, put
$\sh{F}_\eta=i^*(\sh{F})$, let $\sh{G}'$ be an
$\sh{O}_{X_\eta}$-module quotient of $\sh{F}_\eta$, and let
$\overline{\sh{G}'}$ be the image of $\sh{F}$ under the composite
homomorphism~\sref[0]{0.4.4.3.2}
\[
  \sh{F}
    \xrightarrow{\rho_{\sh{F}}}
  i_*\bigl(i^*(\sh{F})\bigr)
    \longrightarrow
  i_*(\sh{G}').
\]

Then $\overline{\sh{G}'}$ is a quasi-coherent
$\sh{O}_X$-module quotient of $\sh{F}$, flat over $Y$, such that
\[
  i^*\bigl(\overline{\sh{G}'}\bigr)=\sh{G}',
\]
and it is the unique $\sh{O}_X$-module quotient of $\sh{F}$ having
these properties.
\end{proposition}

\begin{proof}
Since $i$ is quasi-compact and quasi-separated
\sref{IV.1.1.2} and~\sref{IV.1.2.2}, it follows
from~\sref{IV.1.7.4} that, for every quasi-coherent
$\sh{O}_{X_\eta}$-module $\sh{H}'$, the $\sh{O}_X$-module
$i_*(\sh{H}')$ is quasi-coherent. Moreover, for every open subset
$U$ of $X$, by definition~\sref[0]{0.3.4.1},
\[
  i_*(\sh{H}')|U
    =
  (i|U\cap X_\eta)_*
    \bigl(\sh{H}'|U\cap X_\eta\bigr).
\]
If the proposition is proved when $X$ and $Y$ are affine, the general
case follows by gluing, in view of the uniqueness assertion in the
affine case. In other words, we are reduced to proving the following
lemma.

\begin{lemma}[2.8.1.1]
\phantomsection
\label{IV.2.8.1.1}
Let $A$ be a Noetherian regular integral ring of dimension $1$
\sref[0]{0.17.3.6}, let $K$ be its field of fractions, let $M$ be an
$A$-module, and let $N'$ be a $K$-module quotient of
\[
  M_{(K)}=M\otimes_A K
\]
by a sub-$K$-module $P'$. Let $\overline{N'}$ be the image of $M$
under the composite homomorphism
\[
  M\longrightarrow M_{(K)}\longrightarrow N'.
\]
Then $\overline{N'}$ is a flat $A$-module, and it is the unique
quotient module $N$ of $M$ that is flat over $A$ and for which the
kernel of the surjective homomorphism
\[
  M_{(K)}\longrightarrow N_{(K)}
\]
is equal to $P'$.
\end{lemma}

For every maximal ideal $\mathfrak{m}$ of $A$, the ring
$A_\mathfrak{m}$ is a regular local ring of dimension $1$, hence a
discrete valuation ring. Therefore an $A$-module $N$ is flat if and
only if it is torsion-free~\sref[0]{0.6.3.4}. Since $N'$ is a
$K$-vector space, it is torsion-free as an $A$-module, and the same is
therefore true of its submodule $\overline{N'}$. Moreover, one
immediately verifies that $\overline{N'}_{(K)}$ identifies with $N'$.

Conversely, let $N$ be an $A$-module quotient of $M$ having the
properties in the statement. Since $N$ is flat over $A$, the
canonical homomorphism
\[
  N\longrightarrow N_{(K)}=N\otimes_A K
\]
is injective. Since $N_{(K)}$ identifies with $N'$, the conclusion
follows from the commutativity of the diagram
\[
\xymatrix{
  M \ar[r] \ar[d] & N \ar[d] \\
  M_{(K)} \ar[r] & N_{(K)}.
}
\]
\end{proof}

\begin{corollary}[2.8.2]
\phantomsection
\label{IV.2.8.2}
Under the conditions of~\sref{IV.2.8.1}, the module $\sh{F}$ is flat
over $Y$ if and only if the canonical homomorphism
\[
  \sh{F}\longrightarrow i_*(\sh{F}_\eta)
    =i_*\bigl(i^*(\sh{F})\bigr)
\]
is injective.
\end{corollary}

\begin{env}[2.8.3]
\phantomsection
\label{IV.2.8.3}
The formation of the $\sh{O}_X$-module $\overline{\sh{G}'}$ is
\emph{functorial} in $\sh{F}$ and $\sh{G}'$. More precisely, let
$\sh{F}_1$ and $\sh{F}_2$ be two quasi-coherent
$\sh{O}_X$-modules, let
$u:\sh{F}_1\to\sh{F}_2$ be an $\sh{O}_X$-homomorphism, let
$\sh{G}'_i$ be an $\sh{O}_{X_\eta}$-module quotient of
$(\sh{F}_i)_\eta$ for $i=1,2$, and let
$v:\sh{G}'_1\to\sh{G}'_2$ be a homomorphism making the diagram
\[
\xymatrix{
  (\sh{F}_1)_\eta \ar[r]^{i^*(u)} \ar[d]
    &(\sh{F}_2)_\eta \ar[d] \\
  \sh{G}'_1 \ar[r]_v
    &\sh{G}'_2
}
\]

\oldpage[IV-2]{35}

commutative. Such a homomorphism $v$, when it exists, is uniquely
determined by this property. Then the diagram
\[
\xymatrix{
  \sh{F}_1 \ar[r]^u \ar[d]
    &\sh{F}_2 \ar[d] \\
  i_*(\sh{G}'_1) \ar[r]_{i_*(v)}
    &i_*(\sh{G}'_2)
}
\]
is commutative, and consequently there is a unique homomorphism
\[
  w:\overline{\sh{G}'_1}\longrightarrow\overline{\sh{G}'_2}
\]
making the diagram
\[
\xymatrix{
  \sh{F}_1 \ar[r]^u \ar[d]
    &\sh{F}_2 \ar[d] \\
  \overline{\sh{G}'_1} \ar[r]_w
    &\overline{\sh{G}'_2}
}
\]
commutative.
\end{env}

\begin{proposition}[2.8.4]
\phantomsection
\label{IV.2.8.4}
Suppose that $Y$ satisfies the hypotheses of~\sref{IV.2.8.1}. Let
$X_1$ and $X_2$ be two $Y$-preschemes, let $\sh{F}_i$ be a
quasi-coherent $\sh{O}_{X_i}$-module, and let $\sh{G}'_i$ be an
$\sh{O}_{(X_i)_\eta}$-module quotient of $(\sh{F}_i)_\eta$, for
$i=1,2$. Then
\end{proposition}

\begin{env}[2.8.4.1]
\phantomsection
\label{IV.2.8.4.1}
\[
  \overline{\sh{G}'_1}\otimes_{\sh{O}_Y}\overline{\sh{G}'_2}
    =
  \overline{\sh{G}'_1\otimes_{k(\eta)}\sh{G}'_2}.
\]
\end{env}

\begin{proof}
Put $X=X_1\times_Y X_2$. The left-hand side
of~\sref{IV.2.8.4.1} is a quasi-coherent $\sh{O}_X$-module flat over
$Y$~\sref[0]{0.6.2.1}; its inverse image on $X_\eta$ is
\[
  \sh{G}'_1\otimes_{k(\eta)}\sh{G}'_2
\]
\sref[I]{I.9.1.5}, and it is a quotient of
$\sh{F}_1\otimes_{\sh{O}_Y}\sh{F}_2$. The conclusion therefore
follows from the uniqueness assertion of~\sref{IV.2.8.1}.
\end{proof}

\begin{proposition}[2.8.5]
\phantomsection
\label{IV.2.8.5}
Suppose that $X$ and $Y$ satisfy the hypotheses
of~\sref{IV.2.8.1}, and let $Z'$ be a closed subprescheme of
$X_\eta$. Then there exists a unique closed subprescheme
$\overline{Z'}$ of $X$ that is flat over $Y$ and satisfies
\[
  i^{-1}(\overline{Z'})=Z'.
\]
\end{proposition}

\begin{proof}
Let $\sh{J}'$ be the quasi-coherent ideal of $\sh{O}_{X_\eta}$
defining $Z'$. Apply~\sref{IV.2.8.1} with
\[
  \sh{F}=\sh{O}_X,\qquad
  \sh{G}'=\sh{O}_{X_\eta}/\sh{J}'.
\]
If $\overline{\sh{G}'}=\sh{O}_X/\sh{J}$, then
\[
  \sh{J}'=i^*(\sh{J})\sh{O}_{X_\eta},
\]
and hence
$i^{-1}(\overline{Z'})=Z'$~\sref[I]{I.4.4.5}.
\end{proof}

The prescheme $\overline{Z'}$ is the \emph{closed image} of $Z'$ under
the composite morphism
\[
  Z'\longrightarrow X_\eta\xrightarrow{i}X,
\]
where the first arrow is the canonical injection~\sref[I]{I.9.5.3}.
Its underlying space is the closure of $Z'$ \emph{in} $X$
\sref[I]{I.9.5.4}, which justifies the notation. We also say that
$\overline{Z'}$ is the closure subprescheme of $X$ of $Z'$.

\begin{corollary}[2.8.6]
\phantomsection
\label{IV.2.8.6}
Let $X_1$ and $X_2$ be two $Y$-preschemes, and let $Z'_i$ be a closed
subprescheme of $(X_i)_\eta$, for $i=1,2$. Then
\end{corollary}

\begin{env}[2.8.6.1]
\phantomsection
\label{IV.2.8.6.1}
\[
  \overline{Z'_1}\times_Y\overline{Z'_2}
    =
  \overline{Z'_1\times_{k(\eta)}Z'_2}.
\]
\end{env}

\begin{proof}
This follows from~\sref{IV.2.8.4} and~\sref{IV.2.8.5}.
\end{proof}
% END NATIVE INPUT ega4-2.source_aligned_r1
% BEGIN NATIVE INPUT ega4-3.source_aligned_r1
\section{Associated prime cycles and primary decompositions}
\label{section:IV.3}

\oldpage[IV-2]{36}

In this section, we mainly give the translation of results on modules presented in Bourbaki, \emph{Alg. comm.}, chap.~IV, which we follow very closely.
The notions that follow seem to be of interest only in the case of \emph{locally Noetherian} preschemes.

\subsection{Associated prime cycles of a module}
\label{subsection:IV.3.1}

\begin{definition}[3.1.1]
\phantomsection
\label{IV.3.1.1}
Let $X$ be a prescheme and $\sh{F}$ a quasi-coherent $\sh{O}_X$-module.
We say that a point $x\in X$ is \emph{associated} to $\sh{F}$ if the maximal ideal $\mathfrak{m}_x$ of $\sh{O}_x$ is associated to the $\sh{O}_x$-module $\sh{F}_x$ (in other words, if $\mathfrak{m}_x$ is the annihilator of an element of $\sh{F}_x$).
We denote by $\operatorname{Ass}(\sh{F})$ the set of $x\in X$ associated to $\sh{F}$.
We say that a closed irreducible subset $Z$ of $X$ is an \emph{associated prime cycle of $\sh{F}$} if its generic point is associated to $\sh{F}$.
When $\sh{F}=\sh{O}_X$, we also say that the points (resp.\ \emph{prime cycles}) associated to $\sh{F}$ are associated to the \emph{prescheme} $X$.
\end{definition}

We say that an associated prime cycle of $\sh{F}$ (resp.\ $X$) is \emph{embedded} if it is contained in another associated prime cycle of $\sh{F}$ (resp.\ $X$) (in other words, if it is not \emph{maximal} in the set of these cycles).

If $X$ is locally Noetherian, the \emph{irreducible components} of $X$ are none other than the \emph{maximal} (or \emph{non-embedded}) prime cycles associated to $X$.

It is clear that if $x\in\operatorname{Ass}(\sh{F})$, then $\sh{F}_x\neq 0$, in other words
\[
\label{IV.3.1.1.1}
\operatorname{Ass}(\sh{F})\subset\operatorname{Supp}(\sh{F}).
\tag{3.1.1.1}
\]

If $x\in X$ is associated to $\sh{F}$, it is evidently associated to $\sh{F}|U$ for every open neighbourhood $U$ of $x$, and conversely, if it is associated to $\sh{F}|U$ for one of these neighbourhoods, it is associated to $\sh{F}$.

We note finally that for a quasi-coherent $\sh{O}_X$-module $\sh{F}$, the existence of embedded associated prime cycles is a \emph{local} question, for if $y$ and $z$ are two points of $\operatorname{Ass}(\sh{F})$ such that $y\in\overline{\{z\}}$, every neighbourhood of $y$ contains $z$.

\begin{proposition}[3.1.2]
\phantomsection
\label{IV.3.1.2}
Let $A$ be a Noetherian ring, $M$ an $A$-module, $X=\operatorname{Spec}(A)$, $\sh{F}=\widetilde{M}$; for a point $x\in X$ to be associated to $\sh{F}$, it is necessary and sufficient that the prime ideal $\mathfrak{j}_x$ of $A$ be associated to the module $M$ (in other words, be the annihilator of an element $f\in M$).
\end{proposition}

\begin{proof}
This follows from the definition \sref{IV.3.1.1} and from Bourbaki, \emph{loc.\ cit.}, \textsection~1, no~2, cor.\ of prop.~5, applied to $S=A-\mathfrak{j}_x$.
\end{proof}

\begin{proposition}[3.1.3]
\phantomsection
\label{IV.3.1.3}
Let $X$ be a locally Noetherian prescheme, $\sh{F}$ a quasi-coherent $\sh{O}_X$-module, $x$ a point of $X$, and $Z$ the reduced closed sub-prescheme of $X$ having $\overline{\{x\}}$ as underlying space \sref[I]{I.5.2.1}.
The following conditions are equivalent:
\begin{enumerate}
  \item[\emph{a)}] $x\in\operatorname{Ass}(\sh{F})$.
\oldpage[IV-2]{37}
  \item[\emph{b)}] There exist an open neighbourhood $U$ of $x$ and a section $f\in\Gamma(U,\sh{F})$ such that $U\cap Z$ is equal to $\operatorname{Supp}(\sh{O}_U\cdot f)$.
  \item[\emph{b$'$)}] There exist an open neighbourhood $U$ of $x$ and a section $f\in\Gamma(U,\sh{F})$ such that $U\cap Z$ is an irreducible component of $\operatorname{Supp}(\sh{O}_U\cdot f)$.
  \item[\emph{c)}] There exist an open neighbourhood $U$ of $x$ and a submodule of $\sh{F}|U$ isomorphic to $\sh{O}_Z|U$ ($\sh{O}_Z$ being identified with a quotient of $\sh{O}_X$).
  \item[\emph{c$'$)}] There exist an open neighbourhood $U$ of $x$ and a coherent submodule $\sh{G}$ of $\sh{F}|U$ such that $U\cap Z$ is an irreducible component of $\operatorname{Supp}(\sh{G})$.
\end{enumerate}
\end{proposition}

\begin{proof}
It is clear that \emph{c)} implies \emph{b)}, for it suffices to take for $f$ the element of $\Gamma(U,\sh{F})$ which corresponds to the unit section of $\sh{O}_Z|U$.
As $U\cap Z$ is irreducible \sref[0]{0.2.1.6}, \emph{b)} implies \emph{b$'$)}, and \emph{b$'$)} implies \emph{c$'$)} since $\sh{O}_X$ is coherent \sref[0]{0.5.3.4}.
To see that \emph{c$'$)} implies \emph{a)}, we can reduce to the case where $U=X=\operatorname{Spec}(A)$ is affine, $A$ being Noetherian, and where $\sh{F}=\widetilde{M}$, where $M$ is an $A$-module, and $\sh{G}=\widetilde{N}$, where $N$ is a sub-module of $M$, of finite type.
We then know that the minimal elements of $\operatorname{Supp}(\sh{G})$ are the maximal points of $V(\operatorname{Ann}(N))$ \sref[0]{0.1.7.4}, and these are also the minimal elements of $\operatorname{Ass}(N)$ (Bourbaki, \emph{Alg.\ comm.}, chap.~IV, \textsection~1, no~3, cor.~1 of prop.~7);
as $\operatorname{Ass}(N)\subset\operatorname{Ass}(M)=\operatorname{Ass}(\sh{F})$, we see that \emph{c$'$)} implies \emph{a)}.
Finally, \emph{a)} implies \emph{c)} by \sref{IV.3.1.2}, taking $X$ affine again, $\sh{F}=\widetilde{M}$, and $Z$ defined by the ideal $\mathfrak{j}_xA$ (with the notation of \sref{IV.3.1.2}).
\end{proof}

\begin{corollary}[3.1.4]
\phantomsection
\label{IV.3.1.4}
Let $X$ be a locally Noetherian prescheme and $\sh{F}$ a coherent $\sh{O}_X$-module.
Then the maximal prime cycles associated to $\sh{F}$ are the irreducible components of $\operatorname{Supp}(\sh{F})$, and the generic points of these components are the points $x\in X$ such that $\sh{F}_x$ is an $\sh{O}_x$-module of finite nonzero length.
\end{corollary}

\begin{proof}
Indeed, if $x$ is the generic point of an irreducible component $Z$ of $\operatorname{Supp}(\sh{F})$, the equivalence of \emph{a)} and \emph{c$'$)} in \sref{IV.3.1.3} shows that $x$ belongs to $\operatorname{Ass}(\sh{F})$, and $Z$ is an associated prime cycle of $\sh{F}$, necessarily maximal by \sref{IV.3.1.1.1}; the converse follows trivially from \sref{IV.3.1.1.1}. Finally, the last assertion is local and follows from Bourbaki, \emph{Alg.\ comm.}, chap.~IV, \textsection~2, no~5, cor.~2 of prop.~7.
\end{proof}

\begin{corollary}[3.1.5]
\phantomsection
\label{IV.3.1.5}
Let $X$ be a locally Noetherian prescheme and $\sh{F}$ a quasi-coherent $\sh{O}_X$-module.
For $\sh{F}=0$, it is necessary and sufficient that $\operatorname{Ass}(\sh{F})=\varnothing$.
\end{corollary}

\begin{proof}
The question being local, we reduce to the case where $X$ is affine, and the conclusion follows immediately from \sref{IV.3.1.2} and from Bourbaki, \emph{Alg.\ comm.}, chap.~IV, \textsection~1, no~1, cor.~1 of prop.~2.
\end{proof}

\begin{proposition}[3.1.6]
\phantomsection
\label{IV.3.1.6}
Let $X$ be a locally Noetherian prescheme and $\sh{F}$ a coherent $\sh{O}_X$-module; then $\operatorname{Ass}(\sh{F})$ is \emph{locally finite} (that is to say, every point of $X$ admits a neighbourhood whose intersection with $\operatorname{Ass}(\sh{F})$ is finite).
\end{proposition}

\begin{proof}
It suffices to consider the case where $X$ is affine, hence Noetherian, and then the proposition follows from \sref{IV.3.1.2} and from Bourbaki, \emph{Alg.\ comm.}, chap.~IV, \textsection~1, no~4, cor.\ of th.~2.
\end{proof}

\begin{proposition}[3.1.7]
\phantomsection
\label{IV.3.1.7}
Let $X$ be a prescheme.
\begin{enumerate}
  \item[(i)] Let $0\to\sh{F}'\to\sh{F}\to\sh{F}''\to 0$ be an exact sequence of quasi-coherent $\sh{O}_X$-modules.
  Then $\operatorname{Ass}(\sh{F}')\subset\operatorname{Ass}(\sh{F})\subset\operatorname{Ass}(\sh{F}')\cup\operatorname{Ass}(\sh{F}'')$.
  \item[(ii)] Let $\sh{F}$ be a quasi-coherent $\sh{O}_X$-module, $(\sh{F}_\alpha)$ a family of quasi-coherent $\sh{O}_X$-submodules of $\sh{F}$ such that $\sh{F}$ is the union of the $\sh{F}_\alpha$.
  Then $\operatorname{Ass}(\sh{F})=\bigcup_\alpha\operatorname{Ass}(\sh{F}_\alpha)$.
  \item[(iii)] For every family $(\sh{F}_\alpha)$ of quasi-coherent $\sh{O}_X$-modules, we have $\operatorname{Ass}(\bigoplus_\alpha\sh{F}_\alpha)=\bigcup_\alpha\operatorname{Ass}(\sh{F}_\alpha)$.
\end{enumerate}
\end{proposition}

\begin{proof}
We immediately reduce to the corresponding propositions for modules (Bourbaki, \emph{loc.\ cit.}, \textsection~1, no~1, formula~(1), prop.~3 and cor.~1 of prop.~3).
\end{proof}

\oldpage[IV-2]{38}

\begin{proposition}[3.1.8]
\phantomsection
\label{IV.3.1.8}
Let $X$ be a locally Noetherian prescheme, $\sh{F}$ a quasi-coherent $\sh{O}_X$-module, $U$ an open subset of $X$, and $\sh{J}$ a coherent ideal of $\sh{O}_X$ defining a closed sub-prescheme of $X$ having $X-U$ as underlying space.
The following conditions are equivalent:
\begin{enumerate}
  \item[\emph{a)}] $\operatorname{Ass}(\sh{F})\subset U$.
  \item[\emph{b)}] For every affine open subset $V$ of $X$, every section of $\sh{F}$ over $V$ whose restriction to $V\cap U$ is zero is itself zero.
  \item[\emph{c)}] The canonical homomorphism $\sh{F}\to\shHom_{\sh{O}_X}(\sh{J},\sh{F})$ is injective.
\end{enumerate}
\end{proposition}

\begin{proof}
The question being local, we can suppose $X=\operatorname{Spec}(A)$ affine, $A$ being a Noetherian ring, $\sh{F}=\widetilde{M}$, where $M$ is an $A$-module, and $\sh{J}=\widetilde{\mathfrak{J}}$, where $\mathfrak{J}$ is an ideal of $A$.
The homomorphism $M\to\operatorname{Hom}_A(\mathfrak{J},M)$ sends $m\in M$ to the homomorphism $x\mapsto xm$ of $\mathfrak{J}$ into $M$; to say that it is not injective means that there exists $m\neq 0$ in $M$ such that $\mathfrak{J} m=0$.

Let us first show that \emph{c)} implies \emph{a)}. If $\operatorname{Ass}(\sh{F})$ met $X-U$, there would be a prime ideal $\mathfrak{p}\in\operatorname{Ass}(M)$ containing $\mathfrak{J}$, hence an element $m\neq 0$ of $M$ such that $\mathfrak{p}m=0$, and therefore $\mathfrak{J}m=0$, contradicting \emph{c)}.

Next, \emph{b)} implies \emph{c)}. Indeed, if $\mathfrak{J}m=0$ for some $m\neq 0$ in $M$, then for every prime ideal $\mathfrak{q}$ not containing $\mathfrak{J}$ there exists $a\in\mathfrak{J}$ with $a\notin\mathfrak{q}$. The relation $am=0$ then implies that the canonical image of $m$ in $M_\mathfrak{q}$ is zero. Thus $m$ would define a nonzero section of $\sh{F}$ over $X$ whose restriction to $U$ is zero, contrary to \emph{b)}.

Finally, \emph{a)} implies \emph{b)}. The canonical homomorphism
\[
M\longrightarrow\prod_{\mathfrak{p}\in\operatorname{Ass}(M)}M_\mathfrak{p}
\]
is injective: if $N$ is its kernel, then $\operatorname{Ass}(N)\subset\operatorname{Ass}(M)$; if there existed $\mathfrak{p}\in\operatorname{Ass}(N)$, there would be an element $n\in N$ whose annihilator is $\mathfrak{p}$, while the definition of $N$ would give an $s\notin\mathfrak{p}$ such that $sn=0$, a contradiction. Hence $\operatorname{Ass}(N)=\varnothing$, and therefore $N=0$. Under \emph{a)}, a section $m\in M$ whose restriction to $U$ is zero has zero image in $M_\mathfrak{p}$ for every $\mathfrak{p}\in\operatorname{Ass}(M)$; consequently $m=0$.
\end{proof}

\begin{corollary}[3.1.9]
\phantomsection
\label{IV.3.1.9}
Let $X$ be a locally Noetherian prescheme, $\sh{F}$ a quasi-coherent $\sh{O}_X$-module, and $f$ a section of $\sh{O}_X$ over $X$.
For $f$ to be $\sh{F}$-regular \sref[0]{0.15.2.2}, it is necessary and sufficient that $\operatorname{Ass}(\sh{F})\subset X_f$.
\end{corollary}

\begin{proof}
Indeed, it is immediate that to say that $f$ is $\sh{F}$-regular means that the canonical homomorphism $\sh{F}\to\shHom_{\sh{O}_X}(f\sh{O}_X,\sh{F})$ is injective, and it suffices to apply \sref{IV.3.1.8} to the ideal $\sh{J}=f\sh{O}_X$.
\end{proof}

\begin{proposition}[3.1.10]
\phantomsection
\label{IV.3.1.10}
Let $X$ and $Y$ be locally Noetherian preschemes and $f:X\to Y$ an integral morphism.
Then, for every quasi-coherent $\sh{O}_X$-module $\sh{F}$, we have $f(\operatorname{Ass}(\sh{F}))=\operatorname{Ass}(f_*(\sh{F}))$.
\end{proposition}

\begin{proof}
The question being local on $Y$, and the morphism $f$ being affine, we are immediately reduced to the case where $Y$ is affine, in other words, to the following lemma.
\end{proof}

\oldpage[IV-2]{39}

\begin{lemma}[3.1.10.1]
\phantomsection
\label{IV.3.1.10.1}
Let $A$ and $B$ be two Noetherian rings, $\rho:A\to B$ a ring homomorphism making $B$ an integral $A$-algebra, and $M$ a $B$-module.
Then the prime ideals $\mathfrak{p}\in\operatorname{Ass}(M_{[\rho]})$ are the inverse images by $\rho$ of the prime ideals $\mathfrak{q}\in\operatorname{Ass}(M)$.
\end{lemma}

\begin{proof}
Indeed, if $\mathfrak{q}\in\operatorname{Ass}(M)$, then $\mathfrak{q}$ is the annihilator in $B$ of an element $x\in M$, and hence $\rho^{-1}(\mathfrak{q})$ is the annihilator in $A$ of $x$. Conversely, suppose $\mathfrak{p}\in\operatorname{Ass}(M_{[\rho]})$, so that $\mathfrak{p}$ is the inverse image under $\rho$ of the annihilator $\mathfrak{b}$ in $B$ of an element $x\in M$. It follows from the first Cohen--Seidenberg theorem that there exists a prime ideal $\mathfrak{q}$ of $B$ containing $\mathfrak{b}$ and whose inverse image is $\mathfrak{p}$ (Bourbaki, \emph{Alg.\ comm.}, chap.~V, \textsection~2, no~1, cor.~2 of th.~1). By considering a prime ideal minimal among those contained in $\mathfrak{q}$ and containing $\mathfrak{b}$, we may evidently suppose that $\mathfrak{q}$ itself is one of these minimal primes. Since $B.x\subset M$ is isomorphic to $B/\mathfrak{b}$, we then have $\mathfrak{q}\in\operatorname{Ass}(B/\mathfrak{b})\subset\operatorname{Ass}(M)$ (Bourbaki, \emph{Alg.\ comm.}, chap.~IV, \textsection~1, no~4, th.~2).
\end{proof}

\begin{corollary}[3.1.11]
\phantomsection
\label{IV.3.1.11}
Under the hypotheses of \sref{IV.3.1.10}, for $\sh{F}$ to have no embedded associated prime cycle, it suffices that the same be true of $f_*(\sh{F})$.
\end{corollary}

\begin{proof}
Suppose indeed that $f_*(\sh{F})$ has no embedded associated prime cycle.
Note that if $A$ is an integral algebra over a field $k$, all prime ideals of $A$ are maximal (Bourbaki, \emph{Alg.\ comm.}, chap.~V, \textsection~2, no~1, prop.~1); it follows from \sref[I]{I.6.2.2} that the fibres of $f$ are \emph{discrete} spaces. If $x$ and $x'$ are two distinct points of $\operatorname{Ass}(\sh{F})$, neither can be adherent to the other when $f(x)=f(x')$; and if $f(x)\neq f(x')$, then \sref{IV.3.1.10} and the hypothesis imply that neither of the two points $f(x)$ and $f(x')$ can be adherent to the other, so the same is true of $x$ and $x'$.
\end{proof}

\begin{remark}[3.1.12]
\phantomsection
\label{IV.3.1.12}
Under the hypotheses of \sref{IV.3.1.10}, it may happen, on the contrary, that $\sh{F}$ has no embedded associated prime cycle whereas $f_*(\sh{F})$ does.
For example, take $Y=\operatorname{Spec}(k[\mathrm{T}])$, where $k$ is a field (an ``affine line''), and let $X$ be the disjoint union of $X_1=Y$ and $X_2=\operatorname{Spec}(k)$, the morphism $X_2\to Y$ corresponding to the canonical homomorphism $k[\mathrm{T}]\to k[\mathrm{T}]/\mathfrak{m}$, where $\mathfrak{m}$ is the maximal ideal $(\mathrm{T})$.
It is clear that the morphism $f:X\to Y$ is finite. If one takes $\sh{F}=\sh{O}_X$, then $\sh{F}$ has no embedded associated prime cycle, but $f_*(\sh{F})=\widetilde{M}$, where $M$ is the direct sum of the $k[\mathrm{T}]$-modules $k[\mathrm{T}]$ and $k$; hence $\operatorname{Ass}(M)$ consists of the generic point $(0)$ of $Y$ and the point $\mathfrak{m}$.
\end{remark}

\begin{proposition}[3.1.13]
\phantomsection
\label{IV.3.1.13}
Let $X$ be a locally Noetherian prescheme, $U$ an open subset of $X$, and $i:U\to X$ the canonical inclusion.
For every quasi-coherent $\sh{O}_U$-module $\sh{F}$, we have $\operatorname{Ass}(i_*(\sh{F}))=\operatorname{Ass}(\sh{F})$.
\end{proposition}

\begin{proof}
Recall that $i_*(\sh{F})$ is a quasi-coherent $\sh{O}_X$-module \sref[I]{I.2.2} and \sref[I]{I.7.4}; as $i_*(\sh{F})|U=\sh{F}$, we have $\operatorname{Ass}(i_*(\sh{F}))\cap U=\operatorname{Ass}(\sh{F})$, and it therefore remains to prove that $\operatorname{Ass}(i_*(\sh{F}))\subset U$.
But by \sref{IV.3.1.8}, this relation means that for every affine open subset $V$ of $X$, every section of $i_*(\sh{F})$ over $V$ which is zero on $U\cap V$, is zero, a condition which is trivially satisfied since $\Gamma(V,i_*(\sh{F}))=\Gamma(U\cap V,\sh{F})=\Gamma(U\cap V,i_*(\sh{F}))$.
\end{proof}

\oldpage[IV-2]{40}

\subsection{Irredundant decompositions}
\label{subsection:IV.3.2}

\phantomsection
\begin{proposition}[3.2.1]
\label{IV.3.2.1}
Let $X$ be a locally Noetherian prescheme and $U$ a dense open subset of $X$.
The following conditions are equivalent:
\begin{enumerate}
  \item[\emph{a)}] $X$ is reduced.
  \item[\emph{b)}] The induced sub-prescheme on $U$ is reduced, and $X$ has no embedded associated prime cycle.
  \item[\emph{c)}] $X$ has no embedded associated prime cycle and, for every generic point $x$ of an irreducible component of $X$, one has $\operatorname{length}(\sh{O}_x)=1$.
\end{enumerate}
The prime cycles associated with $X$ are then precisely the irreducible components of $X$.
\end{proposition}

\begin{proof}
It is clear that if $X$ is reduced, then so is the sub-prescheme induced on $U$.
Moreover, since the existence of embedded associated prime cycles is a local question, we may restrict to the case where $X=\operatorname{Spec}(A)$ is affine and $A$ is Noetherian.
If $A$ is reduced, its minimal prime ideals form a reduced primary decomposition of $(0)$ (Bourbaki, \emph{Alg.\ comm.}, chap.~IV, \textsection~2, no~5, prop.~10) and are precisely the elements of $\operatorname{Ass}(A)$; hence $A$ has no embedded associated prime ideals.  Thus \emph{a)} implies \emph{b)}.
It is immediate that \emph{b)} implies \emph{c)}, since a generic point $x$ of an irreducible component belongs to $U$, and consequently $\sh{O}_x$ is a field.
Finally, \emph{c)} implies \emph{a)}.  Indeed, let $\sh{N}$ be the nilradical of $\sh{O}_X$, which is a coherent ideal.  By hypothesis, $\operatorname{Supp}(\sh{N})$ contains none of the generic points of the irreducible components of $X$.  If $\operatorname{Supp}(\sh{N})$ were nonempty and $x$ were one of its maximal points, criterion \sref{IV.3.1.3}, \emph{c$'$)}, would give $x\in\operatorname{Ass}(\sh{O}_X)$; then $\overline{\{x\}}$ would be an embedded associated prime cycle of $X$, contrary to the hypothesis.  Hence $\sh{N}=0$.
\end{proof}

\phantomsection
\begin{definition}[3.2.2]
\label{IV.3.2.2}
Let $X$ be a locally Noetherian prescheme and $\sh{F}$ a coherent $\sh{O}_X$-module.
We say that $\sh{F}$ is \emph{reduced} if it satisfies the following two conditions: 1\textsuperscript{o} $\sh{F}$ has no embedded associated prime cycle; 2\textsuperscript{o} for every maximal point $x$ of $\operatorname{Supp}(\sh{F})$, one has $\operatorname{length}(\sh{F}_x)=1$.
\end{definition}

Condition 1\textsuperscript{o} means that the prime cycles associated with $\sh{F}$ are the irreducible components of $\operatorname{Supp}(\sh{F})$ \sref{IV.3.1.4}, and condition 2\textsuperscript{o} means that for every generic point $x$ of such a component one has $\operatorname{length}(\sh{F}_x)=1$.

For an affine scheme $X$, this definition gives in particular the notion of a \emph{reduced module} over a Noetherian ring $A$: a finitely generated $A$-module $M$ is said to be \emph{reduced} if it has no embedded associated prime ideals and if, for every $\mathfrak{p}\in\operatorname{Ass}(M)$, one has $\operatorname{length}_{A_\mathfrak{p}}(M_\mathfrak{p})=1$.
Returning to a locally Noetherian prescheme $X$ and a coherent $\sh{O}_X$-module $\sh{F}$, we say that $\sh{F}$ is \emph{reduced at a point} $x\in X$ if $\sh{F}_x$ is a reduced $\sh{O}_x$-module.  Equivalently, on the local scheme $\operatorname{Spec}(\sh{O}_x)$, the module associated with $\sh{F}_x$ is reduced.  Thus $x$ belongs to no embedded associated prime cycle of $\sh{F}$ and $\operatorname{length}(\sh{F}_z)=1$ \emph{for every maximal point $z$ of $\operatorname{Supp}(\sh{F})$ such that $x\in\overline{\{z\}}$}.
It is clear that if $\sh{F}$ is a reduced coherent $\sh{O}_X$-module, then it is reduced at every point of $X$.  Conversely, if $\sh{F}$ is \emph{reduced at a point} $x$, there exists an open neighbourhood $U$ of $x$ such that $\sh{F}|U$ is a reduced $\sh{O}_U$-module: it suffices to take $U$ meeting no embedded associated prime cycle of $\sh{F}$ (such a neighbourhood exists because these cycles form a locally finite family of closed subsets of $X$).
To say that $\sh{O}_X$ is reduced at $x$ is equivalent to saying that $X$ is \emph{reduced at the point} $x$.

\oldpage[IV-2]{41}

\phantomsection
\begin{proposition}[3.2.3]
\label{IV.3.2.3}
Let $X$ be a locally Noetherian prescheme, $U$ an open subset of $X$, and $\sh{F}$ a coherent $\sh{O}_X$-module such that $U\cap\operatorname{Supp}(\sh{F})$ is dense in $\operatorname{Supp}(\sh{F})$.
The following conditions are equivalent:
\begin{enumerate}
  \item[\emph{a)}] $\sh{F}$ is reduced.
  \item[\emph{b)}] $\sh{F}|U$ is reduced, and $\sh{F}$ has no embedded associated prime cycle.
  \item[\emph{c)}] There exist a reduced closed sub-prescheme $X'$ of $X$ and a coherent torsion-free $\sh{O}_{X'}$-module $\sh{F}'$ of rank $1$ on every irreducible component of $X'$ such that, if $j:X'\to X$ is the canonical inclusion, then $j_*(\sh{F}')=\sh{F}$.
\end{enumerate}
Moreover, when these conditions hold, $X'$ is defined by the ideal $\sh{J}$ of $\sh{O}_X$ that annihilates $\sh{F}$.
\end{proposition}

\begin{proof}
To see that \emph{c)} implies \emph{a)}, we may, by \sref{IV.3.1.3}, restrict to the case where $X'=X$ is integral, with generic point $x$.  We may further suppose that $X=\operatorname{Spec}(A)$ is affine and $\sh{F}=\widetilde{M}$, where $A$ is integral and $M$ is a torsion-free $A$-module of rank $1$.  The annihilator of every element of $M$ is then zero, so $\operatorname{Ass}(\sh{F})=\{x\}$ and $\sh{F}_x$ is isomorphic to $\sh{O}_x$, the field of fractions of $A$; hence the conditions of Definition \sref{IV.3.2.2} are satisfied.
Since the existence of embedded associated prime cycles is a local question, if $\sh{F}$ has no such cycles and $U\cap\operatorname{Supp}(\sh{F})$ is dense in $\operatorname{Supp}(\sh{F})$, then $\operatorname{Ass}(\sh{F})=\operatorname{Ass}(\sh{F}|U)$.  Thus \emph{a)} and \emph{b)} are equivalent.
If \emph{a)} holds, take for $X'$ the reduced closed sub-prescheme of $X$ whose underlying space is $\operatorname{Supp}(\sh{F})$ \sref[I]{I.5.2.1}, and put $\sh{F}'=j^*(\sh{F})$.  A point $x$ of $\operatorname{Ass}(\sh{F})$ is necessarily a maximal point of $X'$, and since $\operatorname{length}(\sh{F}_x)=1$, $\sh{F}'_x$ is isomorphic to $\sh{F}_x$, hence to the residue field $\boldsymbol{k}(x)=\sh{O}_{X',x}$.  This proves that $\sh{F}'$ is torsion-free and of rank $1$ \sref[I]{I.7.4.6} and \sref[I]{I.7.4.2}.
Finally, the last assertion is immediate, since for every $y\in X'$ the annihilator of the $\sh{O}_{X',y}$-module $\sh{F}'_y$ is zero.
\end{proof}

\phantomsection
\begin{definition}[3.2.4]
\label{IV.3.2.4}
Let $X$ be a locally Noetherian prescheme and $\sh{F}$ a quasi-coherent $\sh{O}_X$-module.
We say that $\sh{F}$ is \emph{irredundant} if $\operatorname{Ass}(\sh{F})$ consists of a single element $x$; if $\sh{F}$ is of finite type, we say that $\sh{F}$ is \emph{integral} if, in addition, $\sh{F}$ is reduced (in other words, if $\operatorname{length}(\sh{F}_x)=1$).
We say that a quasi-coherent $\sh{O}_X$-submodule $\sh{G}$ of a quasi-coherent $\sh{O}_X$-module $\sh{F}$ is \emph{primary in $\sh{F}$} if $\sh{F}/\sh{G}$ is irredundant.
\end{definition}

For an affine scheme $X$, this definition gives in particular the notion of an \emph{integral module} over a Noetherian ring $A$: an $A$-module $M$ is said to be \emph{integral} if it is of finite type, if $M$ is primary (that is, if $\operatorname{Ass}(M)$ consists of a single prime ideal $\mathfrak{p}$), and if, in addition, $\operatorname{length}_{A_\mathfrak{p}}(M_\mathfrak{p})=1$.
Returning to an arbitrary locally Noetherian prescheme $X$ and a coherent $\sh{O}_X$-module $\sh{F}$, we say that $\sh{F}$ is \emph{integral at a point} $x\in X$ if $\sh{F}_x$ is an integral $\sh{O}_x$-module.  Equivalently, $x$ belongs to only one associated prime cycle of $\sh{F}$ (necessarily a nonembedded one), and $\operatorname{length}(\sh{F}_z)=1$ at its generic point $z$.
It is clear that if $\sh{F}$ is an integral coherent $\sh{O}_X$-module, then it is integral at every point of $X$.  Conversely, if $\sh{F}$ is \emph{integral at a point} $x$, there exists an open neighbourhood $U$ of $x$ such that $\sh{F}|U$ is an integral $\sh{O}_U$-module: it suffices to take $U$ such that $\sh{F}|U$ is a reduced $\sh{O}_U$-module \sref{IV.3.2.2}.

\oldpage[IV-2]{42}

We say that a locally Noetherian prescheme $X$ is \emph{irredundant} if $\sh{O}_X$ is irredundant (which implies that $X$ is \emph{irreducible}).  For $X$ to be \emph{integral}, it is necessary and sufficient that the $\sh{O}_X$-module $\sh{O}_X$ be integral \sref[I]{I.2.1.8} and \sref{IV.3.2.1}.
If $\sh{O}_X$ is integral at a point $x$, that is, if the ring $\sh{O}_x$ is integral, we say that $X$ is \emph{integral at the point} $x$.
We say that a closed sub-prescheme $Y$ of $X$ is \emph{primary in $X$} if the ideal $\sh{J}$ of $\sh{O}_X$ defining $Y$ is primary in $\sh{O}_X$.

\phantomsection
\begin{definition}[3.2.5]
\label{IV.3.2.5}
Let $X$ be a locally Noetherian prescheme and $\sh{F}$ a coherent $\sh{O}_X$-module.
An \emph{irredundant decomposition} of $\sh{F}$ is a family $(\sh{F}_\alpha)_{\alpha\in I}$ of quotient $\sh{O}_X$-modules of $\sh{F}$ such that the $\sh{F}_\alpha$ are irredundant, the family $(\operatorname{Supp}(\sh{F}_\alpha))_{\alpha\in I}$ is locally finite, and the canonical homomorphism $\sh{F}\to\bigoplus_{\alpha\in I}\sh{F}_\alpha$ is injective.
We say that such a decomposition is \emph{reduced} if the sets $\operatorname{Ass}(\sh{F}_\alpha)$ are pairwise distinct and there is no proper subset $J\subsetneq I$ for which the subfamily $(\sh{F}_\alpha)_{\alpha\in J}$ is also an irredundant decomposition of $\sh{F}$.
\end{definition}

If $(\sh{F}_\alpha)_{\alpha\in I}$ is an irredundant decomposition (resp.\ a reduced irredundant decomposition) of $\sh{F}$ and $\sh{F}_\alpha=\sh{F}/\sh{G}_\alpha$, we also say that the family $(\sh{G}_\alpha)_{\alpha\in I}$ of $\sh{O}_X$-submodules of $\sh{F}$ is a \emph{primary decomposition of $0$ in $\sh{F}$}.  The injectivity of the canonical homomorphism $\sh{F}\to\bigoplus_{\alpha\in I}\sh{F}_\alpha$ is equivalent to the condition $\bigcap_{\alpha\in I}\sh{G}_\alpha=0$.

If $(\sh{F}_\alpha)$ is an irredundant decomposition of $\sh{F}$, to say that it is reduced is equivalent to saying that the sets $\operatorname{Ass}(\sh{F}_\alpha)$ are pairwise distinct and contained in $\operatorname{Ass}(\sh{F})$.  If $\operatorname{Ass}(\sh{F}_\alpha)=\{x_\alpha\}$ for every $\alpha\in I$, then $\alpha\mapsto x_\alpha$ is a bijection from $I$ onto $\operatorname{Ass}(\sh{F})$.  These properties are local and therefore follow from Bourbaki, \emph{Alg.\ comm.}, chap.~IV, \textsection~2, no~3, prop.~4.

\phantomsection
\begin{proposition}[3.2.6]
\label{IV.3.2.6}
Let $X$ be a locally Noetherian prescheme and $\sh{F}$ a coherent $\sh{O}_X$-module.
Then there exists a reduced irredundant decomposition $(\sh{F}^{(x)})_{x\in\operatorname{Ass}(\sh{F})}$ formed of coherent $\sh{O}_X$-modules such that, for every $x\in\operatorname{Ass}(\sh{F})$, one has $\operatorname{Ass}(\sh{F}^{(x)})=\{x\}$.
For every $x\in\operatorname{Ass}(\sh{F})$ such that $\overline{\{x\}}$ is not embedded, $\sh{F}^{(x)}$ is uniquely determined as the image of the canonical homomorphism $\sh{F}\to i_*(i^*(\sh{F}))$, where $i$ is the canonical morphism $\operatorname{Spec}(\boldsymbol{k}(x))\to X$.
\end{proposition}

\begin{proof}
For every $x\in\operatorname{Ass}(\sh{F})$, let $U$ be an affine open neighbourhood of $x$, with ring $A$, and write $\sh{F}|U=\widetilde{M}$, where $M$ is a finitely generated $A$-module.
We know (Bourbaki, \emph{Alg.\ comm.}, chap.~IV, \textsection~1, no~1, prop.~4) that there exists a submodule $N$ of $M$ such that, on setting $P=M/N$, one has $\operatorname{Ass}(P)=\{x\}$ and $\operatorname{Ass}(N)=\operatorname{Ass}(M)-\{x\}$.
Let $\sh{G}=\widetilde{P}$, a quasi-coherent $\sh{O}_U$-module, and let $j:U\to X$ be the canonical inclusion.  Let $u:j^*(\sh{F})\to\sh{G}$ be the surjective homomorphism corresponding to $M\to P$.  It induces a homomorphism $j_*(u):j_*(j^*(\sh{F}))\to j_*(\sh{G})$, and hence, by composition, a homomorphism
\[
v:\sh{F}\xrightarrow{\rho_\sh{F}}j_*(j^*(\sh{F}))\xrightarrow{j_*(u)}j_*(\sh{G})
\]
whose restriction to $U$ is $u$.  We denote by $\sh{F}^{(x)}$ the image of $\sh{F}$ under this homomorphism; it is a coherent $\sh{O}_X$-module \sref[I]{I.6.1.1}.
By \sref{IV.3.1.13}, one has $\operatorname{Ass}(j_*(\sh{G}))=\{x\}$, and therefore, by \sref{IV.3.1.7}, $\operatorname{Ass}(\sh{F}^{(x)})=\{x\}$, since $\sh{F}^{(x)}\neq 0$.

Let $\sh{N}^{(x)}=\operatorname{Ker}(v)$.  Then $\sh{N}^{(x)}|U=\widetilde{N}$, and hence $x\notin\operatorname{Ass}(\sh{N}^{(x)})$.
It follows that the homomorphism $\sh{F}\to\bigoplus_{x\in\operatorname{Ass}(\sh{F})}\sh{F}^{(x)}$ is injective.  Indeed, its kernel $\sh{H}$ is contained in every $\sh{N}^{(x)}$, so $\operatorname{Ass}(\sh{H})$ is contained in the intersection of the sets $\operatorname{Ass}(\sh{N}^{(x)})$, which is empty; consequently, by \sref{IV.3.1.5}, $\sh{H}=0$.
Since $\operatorname{Ass}(\sh{F})$ is locally finite \sref{IV.3.1.6}, it is clear that $(\sh{F}^{(x)})_{x\in\operatorname{Ass}(\sh{F})}$ is a reduced irredundant decomposition of $\sh{F}$ satisfying the conditions of the statement.
Finally, when $\overline{\{x\}}$ is not embedded, the characterization of $\sh{F}^{(x)}$ follows from Bourbaki, \emph{Alg.\ comm.}, chap.~IV, \textsection~2, no~3, prop.~5, since the question is local, together with \sref[I]{I.1.6.7}.
\end{proof}

\oldpage[IV-2]{43}

\phantomsection
\begin{corollary}[3.2.7]
\label{IV.3.2.7}
Under the hypotheses of \sref{IV.3.2.6}, if $\sh{F}$ has no embedded associated prime cycle, there exists a unique reduced irredundant decomposition of $\sh{F}$.
\end{corollary}

\phantomsection
\begin{corollary}[3.2.8]
\label{IV.3.2.8}
Let $X$ be a Noetherian prescheme and $\sh{F}$ a coherent $\sh{O}_X$-module.
There exists a finite filtration $(\sh{F}_i)_{0\leqslant i\leqslant n}$ of $\sh{F}$ such that $\sh{F}_0=\sh{F}$ and $\sh{F}_n=0$, formed of coherent $\sh{O}_X$-modules, and such that the quotients $\sh{F}_i/\sh{F}_{i+1}$ are either zero or irredundant, with $\operatorname{Ass}(\sh{F}_i/\sh{F}_{i+1})\subset\operatorname{Ass}(\sh{F})$.
\end{corollary}

\begin{proof}
Indeed, by \sref{IV.3.2.6}, $\sh{F}$ is isomorphic to an $\sh{O}_X$-submodule of a finite direct sum $\bigoplus_{j=1}^n\sh{G}_j$, where the $\sh{G}_j$ are irredundant and coherent.  Since every quasi-coherent $\sh{O}_X$-submodule of $\sh{G}_j$ is either zero or irredundant \sref{IV.3.1.7}, the modules
\[
\sh{F}_i=\sh{F}\cap\left(\bigoplus_{j=1}^{n-i}\sh{G}_j\right)
\]
satisfy the requirements, since $\sh{F}_i/\sh{F}_{i+1}$ is isomorphic to a coherent $\sh{O}_X$-submodule of $\sh{G}_{n-i}$.
\end{proof}

\subsection{Relations with flatness}
\label{subsection:IV.3.3}

\phantomsection
\begin{proposition}[3.3.1]
\label{IV.3.3.1}
Let $f:X\to Y$ be a morphism, $\sh{F}$ a quasi-coherent $\sh{O}_X$-module that is $f$-flat, and $\sh{E}$ a quasi-coherent $\sh{O}_Y$-module.
If, for every $y\in Y$, we set $\sh{F}_y=\sh{F}\otimes_{\sh{O}_Y}\boldsymbol{k}(y)$, then
\[
\label{IV.3.3.1.1}
\operatorname{Ass}(\sh{E}\otimes_{\sh{O}_Y}\sh{F})\supset\bigcup_{y\in\operatorname{Ass}(\sh{E})}\operatorname{Ass}(\sh{F}_y)
\tag{3.3.1.1}
\]
and the two sides are equal if $Y$ is locally Noetherian.
\end{proposition}

\begin{proof}
(Of course, $\sh{F}_y$ is a sheaf on the fibre $f^{-1}(y)$, and this fibre is identified with a subspace of $X$ \sref[I]{I.3.6.1}.)
The question being local on $X$ and on $Y$, we reduce to the case where $X$ and $Y$ are affine, and the proposition is then proved in Bourbaki, \emph{Alg.\ comm.}, chap.~IV, \textsection~2, no~6, th.~2.
\end{proof}

\phantomsection
\begin{corollary}[3.3.2]
\label{IV.3.3.2}
Let $Y$ be a locally Noetherian prescheme with no embedded associated prime cycles, $f:X\to Y$ a morphism, and $\sh{F}$ a quasi-coherent $\sh{O}_X$-module that is $f$-flat.
Then, for every $x\in\operatorname{Ass}(\sh{F})$, $f(x)$ is a maximal point of $Y$.
\end{corollary}

\begin{proof}
It suffices to apply \sref{IV.3.3.1} with $\sh{E}=\sh{O}_Y$, since $\operatorname{Ass}(\sh{O}_Y)$ is by hypothesis the set of maximal points of $Y$.
\end{proof}

\phantomsection
\begin{corollary}[3.3.3]
\label{IV.3.3.3}
Under the hypotheses of \sref{IV.3.3.1}, suppose in addition that $X$ and $Y$ are locally Noetherian and that $\sh{E}$ and $\sh{F}$ are coherent.
Then the following conditions are equivalent:
\begin{enumerate}
  \item[\emph{a)}] $\sh{E}\otimes_Y\sh{F}$ has no embedded associated prime cycle.
  \item[\emph{b)}] For every point $y\in\operatorname{Ass}(\sh{E})\cap f(\operatorname{Supp}(\sh{F}))$, $\overline{\{y\}}$ is a nonembedded associated prime cycle of $\sh{E}$ and $\sh{F}_y$ has no embedded associated prime cycle.
\end{enumerate}
\end{corollary}

\begin{proof}
\oldpage[IV-2]{44}
Suppose \emph{a)} is satisfied.
The hypotheses imply that $\sh{E}\otimes_Y\sh{F}$ is a coherent $\sh{O}_X$-module \sref[0]{0.5.3.11} and \sref[0]{0.5.3.5}; its associated prime cycles are therefore the irreducible components of $\operatorname{Supp}(\sh{E}\otimes_Y\sh{F})$ \sref{IV.3.1.4}, and for every maximal point $x$ of $\operatorname{Supp}(\sh{E}\otimes_Y\sh{F})$, $f(x)=y$ is a maximal point of $\operatorname{Supp}(\sh{E})$ and $x$ is a maximal point of $\operatorname{Supp}(\sh{F}_y)$ \sref{IV.2.5.5}.
As, by virtue of \sref{IV.3.3.1} and the fact that $y\in f(\operatorname{Supp}(\sh{F}))$ implies $\sh{F}_y\neq 0$ \sref[I]{I.9.1.13}, every point of $\operatorname{Ass}(\sh{E})\cap f(\operatorname{Supp}(\sh{F}))$ is the image by $f$ of a maximal point of $\operatorname{Supp}(\sh{E}\otimes_Y\sh{F})$, we see that condition \emph{b)} is satisfied.

Conversely, suppose \emph{b)} is satisfied, and let us show that if $z$, $z'$ are two distinct points of $\operatorname{Ass}(\sh{E}\otimes_Y\sh{F})$, neither of the two can be adherent to the other.
In the first place, if $f(z)=f(z')=y$, we have $y\in\operatorname{Ass}(\sh{E})$ and $z$ and $z'$ belong to $\operatorname{Ass}(\sh{F}_y)$ by \sref{IV.3.3.1}, whence $y\in f(\operatorname{Supp}(\sh{F}))$; as by hypothesis neither of the two points $z$, $z'$ is adherent to the other \emph{within} $f^{-1}(y)$, neither can be adherent to the other \emph{within} $X$.
If $y=f(z)$ and $y'=f(z')$ are distinct, they belong to $\operatorname{Ass}(\sh{E})\cap f(\operatorname{Supp}(\sh{F}))$, hence neither of the two can be adherent to the other in $Y$; it follows from the continuity of $f$ that neither of the points $z$, $z'$ can be adherent to the other in $X$.
\end{proof}

\phantomsection
\begin{proposition}[3.3.4]
\label{IV.3.3.4}
Let $X$ and $Y$ be locally Noetherian preschemes, $f:X\to Y$ a morphism, $\sh{E}$ a coherent $\sh{O}_Y$-module, and $\sh{F}$ a coherent $\sh{O}_X$-module that is $f$-flat.
Then the following conditions are equivalent:
\begin{enumerate}
  \item[\emph{a)}] $\sh{E}\otimes_Y\sh{F}$ is reduced \sref{IV.3.2.2}.
  \item[\emph{b)}] For every point $y\in\operatorname{Ass}(\sh{E})\cap f(\operatorname{Supp}(\sh{F}))$, $\overline{\{y\}}$ is a nonembedded associated prime cycle of $\sh{E}$, $\operatorname{length}(\sh{E}_y)=1$, and $\sh{F}_y$ is reduced.
\end{enumerate}
\end{proposition}

\begin{proof}
Suppose \emph{a)} is satisfied.
We already know from \sref{IV.3.3.3} that, for every $y\in\operatorname{Ass}(\sh{E})$ lying in $f(\operatorname{Supp}(\sh{F}))$, $\overline{\{y\}}$ is a nonembedded associated prime cycle of $\sh{E}$ and $\sh{F}_y$ has no embedded associated prime cycle.
Furthermore, by \sref{IV.2.5.5}, for every $x\in\operatorname{Ass}(\sh{E}\otimes_Y\sh{F})\cap f^{-1}(y)$, we have
\[
1=\operatorname{length}((\sh{E}\otimes_Y\sh{F})_x)
=\operatorname{length}(\sh{E}_y)\operatorname{length}((\sh{F}_y)_x),
\]
hence $\operatorname{length}(\sh{E}_y)=\operatorname{length}((\sh{F}_y)_x)=1$, which proves \emph{b)}.

Conversely, suppose \emph{b)} is satisfied; we already know that every point $x\in\operatorname{Ass}(\sh{E}\otimes_Y\sh{F})$ is a maximal point of $\operatorname{Supp}(\sh{E}\otimes_Y\sh{F})$, that $y=f(x)$ is a maximal point of $\operatorname{Supp}(\sh{E})$ and $x$ a maximal point of $\operatorname{Supp}(\sh{F}_y)$ \sref{IV.3.3.1} and \sref{IV.3.3.3}; in addition it follows from the hypothesis and from \sref{IV.2.5.5} that $\operatorname{length}((\sh{E}\otimes_Y\sh{F})_x)=1$, which proves \emph{a)}.
\end{proof}

\phantomsection
\begin{corollary}[3.3.5]
\label{IV.3.3.5}
Let $X$ and $Y$ be locally Noetherian preschemes and let $f:X\to Y$ be a flat morphism.
If $Y$ is reduced at the points of $f(X)$ and $f^{-1}(y)$ is a reduced $\boldsymbol{k}(y)$-prescheme for every $y\in f(X)$, then $X$ is reduced.
\end{corollary}

\begin{proof}
Since the nilradical $\sh{N}_Y$ is coherent, the set of points where $Y$ is reduced is open \sref[0]{0.5.2.2}, and we may restrict to the case where $Y$ is reduced.
It suffices then to apply \sref{IV.3.3.4} with $\sh{E}=\sh{O}_Y$ and $\sh{F}=\sh{O}_X$.
\end{proof}

\phantomsection
\begin{proposition}[3.3.6]
\label{IV.3.3.6}
Let $f:X\to S$ and $g:Y\to S$ be two morphisms, $\sh{F}$ a quasi-coherent $\sh{O}_X$-module, and $\sh{G}$ a quasi-coherent $\sh{O}_Y$-module.
Suppose that: 1\textsuperscript{o} $\sh{G}$ is $g$-flat; 2\textsuperscript{o} $X$ is locally Noetherian, and for every $s\in S$, $g^{-1}(s)$ is locally Noetherian (which will hold if $Y$ is also locally Noetherian).
Let $Z=X\times_S Y$; for every pair $(x,y)$ such that $x\in X$, $y\in Y$,
\oldpage[IV-2]{45}
and $f(x)=g(y)=s$, let $\mathrm{T}_{x,y}$ be the prescheme $\operatorname{Spec}(\boldsymbol{k}(x)\otimes_{\boldsymbol{k}(s)}\boldsymbol{k}(y))$, and let $\mathrm{I}_{x,y}$ be the image of $\operatorname{Ass}(\sh{O}_{\mathrm{T}_{x,y}})$ by the canonical monomorphism $\mathrm{T}_{x,y}\to Z$ \sref[I]{I.3.4.9}.
Then
\[
\label{IV.3.3.6.1}
\operatorname{Ass}(\sh{F}\otimes_S\sh{G})=\bigcup_{x\in\operatorname{Ass}(\sh{F})}\bigg(\bigcup_{y\in\operatorname{Ass}(\sh{G}_{f(x)})}\mathrm{I}_{x,y}\bigg)
\tag{3.3.6.1}
\]
where for all $s\in S$, $\sh{G}_s=\sh{G}\otimes_{\sh{O}_S}\boldsymbol{k}(s)$.
\end{proposition}

\begin{proof}
Let $p:Z\to X$, $q:Z\to Y$ be the canonical projections, so that we have the commutative diagram
\[
\xymatrix{
X\ar[d]_{f} & Z\ar[l]_{p}\ar[d]^{q}\\
S & Y\ar[l]^{g}
}
\]
Set $\sh{G}'=q^*(\sh{G})$, so that $\sh{F}\otimes_S\sh{G}=\sh{F}\otimes_X\sh{G}'$; as $\sh{G}'$ is $p$-flat \sref{IV.2.1.4}, it follows from \sref{IV.3.3.1} that
\[
\label{IV.3.3.6.2}
\operatorname{Ass}(\sh{F}\otimes_S\sh{G})=\bigcup_{x\in\operatorname{Ass}(\sh{F})}\operatorname{Ass}(\sh{G}'_x)
\tag{3.3.6.2}
\]
with $\sh{G}'_x=\sh{G}'\otimes_{\sh{O}_X}\boldsymbol{k}(x)$.
If $s=f(x)$, we have $\sh{G}'_x=\sh{G}_s\otimes_{\boldsymbol{k}(s)}\boldsymbol{k}(x)$, and $p^{-1}(x)=g^{-1}(s)\otimes_{\boldsymbol{k}(s)}\boldsymbol{k}(x)$; furthermore, as the field $\boldsymbol{k}(x)$ is a flat $\boldsymbol{k}(s)$-module, the morphism $p^{-1}(x)\to g^{-1}(s)$ is flat; applying \sref{IV.3.3.1} to this morphism, we get
\[
\label{IV.3.3.6.3}
\operatorname{Ass}(\sh{G}'_x)=\bigcup_{y\in\operatorname{Ass}(\sh{G}_s)}\operatorname{Ass}(\sh{O}_{\mathrm{T}_{x,y}})
\tag{3.3.6.3}
\]
whence the proposition.
\end{proof}

We note that if, in the statement, we suppress hypothesis 2\textsuperscript{o}, we can still conclude, by virtue of \sref{IV.3.3.1}, the relation
\[
\label{IV.3.3.6.4}
\operatorname{Ass}(\sh{F}\otimes_S\sh{G})\supset\bigcup_{x\in\operatorname{Ass}(\sh{F})}\bigg(\bigcup_{y\in\operatorname{Ass}(\sh{G}_{f(x)})}\mathrm{I}_{x,y}\bigg).
\tag{3.3.6.4}
\]

\phantomsection
\begin{corollary}[3.3.7]
\label{IV.3.3.7}
Under the hypotheses of \sref{IV.3.3.6}, suppose in addition that $S$ is also locally Noetherian and that $f(\operatorname{Ass}(\sh{F}))\subset\operatorname{Ass}(\sh{O}_S)$.
Then
\[
\label{IV.3.3.7.1}
\operatorname{Ass}(\sh{F}\otimes_S\sh{G})=\bigcup_{(x,y)\in C}\mathrm{I}_{x,y}
\tag{3.3.7.1}
\]
where $C$ is the set of pairs $(x,y)$ such that $x\in\operatorname{Ass}(\sh{F})$, $y\in\operatorname{Ass}(\sh{G})$ and $f(x)=g(y)$.
\end{corollary}

\begin{proof}
Since $\sh{G}$ is $g$-flat, it follows from \sref{IV.3.3.1} that the relation
``$s\in\operatorname{Ass}(\sh{O}_S)$ and $y\in\operatorname{Ass}(\sh{G}_s)$''
is equivalent to ``$y\in\operatorname{Ass}(\sh{G})$''; the conclusion follows from \sref{IV.3.3.6.1}.
\end{proof}

\phantomsection
\begin{remarks}[3.3.8]
\label{IV.3.3.8}
We shall see later \sref{IV.4.2.2} that under the hypotheses of \sref{IV.3.3.6}, $\mathrm{T}_{x,y}$ is a prescheme \emph{with no embedded associated prime cycle}; it will follow that if $\sh{F}$ and the $\sh{G}_s$ have no embedded associated prime cycles, the same is true of $\sh{F}\otimes_S\sh{G}$.
\end{remarks}

\phantomsection
\begin{corollary}[3.3.9]
\label{IV.3.3.9}
Under the conditions of \sref{IV.3.3.7}, we have
\[
\label{IV.3.3.9.1}
q(\operatorname{Ass}(\sh{F}\otimes_S\sh{G}))\subset\operatorname{Ass}(\sh{G})
\tag{3.3.9.1}
\]
(where $q:X\times_S Y\to Y$ is the canonical projection).
\end{corollary}

\begin{proof}
Indeed, if $(x,y)\in Z$, we have $q(\mathrm{I}_{x,y})=\{y\}\subset\operatorname{Ass}(\sh{G})$.
\end{proof}

\oldpage[IV-2]{46}

\subsection{Properties of the sheaves \texorpdfstring{$\sh{F}/t\sh{F}$}{F/tF}}
\label{subsection:IV.3.4}
\label{IV.3.4}

\begin{proposition}[3.4.1]
\phantomsection
\label{IV.3.4.1}
Let $X$ be a locally Noetherian prescheme, $t$ a section of $\sh{O}_X$ over $X$, and $Y$ the closed sub-prescheme of $X$ defined by the ideal $t\sh{O}_X$.
Let $\sh{F}$ be a coherent $\sh{O}_X$-module, let $S$ be the reduced closed sub-prescheme of $X$ having $\operatorname{Supp}(\sh{F})$ as underlying space, and let $(S_i)$ be the family of reduced closed sub-preschemes of $X$ whose underlying spaces are the irreducible components of $S$; denote by $s_i$ the generic point of $S_i$.
Finally, let $Z$ be an irreducible component of $\operatorname{Supp}(\sh{F}/t\sh{F})=S\cap Y$, $z$ its generic point.
\begin{enumerate}
  \item[\emph{(i)}] For every $i$ such that $Z\subset S_i$, $Z$ is an irreducible component of $S_i\cap Y$.
  \item[\emph{(ii)}] If $Z$ is not equal to any of the $S_i$, then
\[
\label{IV.3.4.1.1}
\operatorname{length}((\sh{F}/t\sh{F})_z)\geqslant\sum_i\operatorname{length}(\sh{F}_{s_i})
\tag{3.4.1.1}
\]
where the sum on the right-hand side is extended to all $i$ such that $Z\subset S_i$.
  \item[\emph{(iii)}] Suppose that $Z$ is not equal to any of the $S_i$.
  For the two sides of \sref{IV.3.4.1.1} to be equal, it is necessary and sufficient that the following two conditions be satisfied:
  \begin{enumerate}
    \item[\emph{$\alpha$)}] $t_z$ is $\sh{F}_z$-regular \sref[0]{0.15.1.4}.
    \item[\emph{$\beta$)}] For every $i$ such that $Z\subset S_i$, the canonical image of the germ $t_z$ in $\sh{O}_{S_i,z}$ generates the maximal ideal of this ring (which implies that $\sh{O}_{S_i,z}$ is a discrete valuation ring and that the image of $t_z$ is a uniformizer of this ring).
  \end{enumerate}
\end{enumerate}
\end{proposition}

\begin{proof}
\emph{(i)} If $j:Y\to X$ is the canonical injection, then
\[
\sh{F}/t\sh{F}
  =\sh{F}\otimes_{\sh{O}_X}\sh{O}_Y
  =j^*\sh{F},
\]
and hence
\[
\operatorname{Supp}(\sh{F}/t\sh{F})
  =j^{-1}(S)=S\cap Y
\]
by \sref[I]{I.9.1.13}, which proves the assertion.

\emph{(ii) and (iii)} The points $s_i$ such that $Z\subset S_i$ are precisely those belonging to $\operatorname{Spec}(\sh{O}_z)$.
Thus, to prove \emph{(ii)} and \emph{(iii)}, we may replace $X$ by $\operatorname{Spec}(\sh{O}_z)$.
If $M=\sh{F}_z$, we may therefore suppose that $\sh{F}=\widetilde{M}$, and then $\sh{F}_{s_i}=M_{\mathfrak p_i}$, where the $\mathfrak p_i$ are the minimal prime ideals of $\sh{O}_z$.
Moreover, since $M$ is an $\sh{O}_z$-module of finite type, we have $S=S'_{\mathrm{red}}$, where
\[
S'=\operatorname{Spec}(\sh{O}_z/\mathfrak a)
\]
and $\mathfrak a$ is the annihilator of $M$ in $\sh{O}_z$ \sref[0]{0.1.7.4}.
The two sides of \sref{IV.3.4.1.1} have the same values whether $M$ is regarded as an $\sh{O}_z$-module or as an $(\sh{O}_z/\mathfrak a)$-module.
We may therefore finally replace $X$ by $S'=\operatorname{Spec}(A)$, where $A$ is a Noetherian local ring and $M$ is a faithful $A$-module.
Since $Z$ is closed in $S$, the hypothesis that $Z\neq S_i$ for every $i$ means that $s_i\notin Z$, and hence that $\dim(A)>0$.
Finally, to say that $z$ is the generic point of $Z$, an irreducible component of $S_i\cap V(t)$, means that $A/tA$ has dimension $0$ (in other words, that it is a local Artinian ring).
We are thus reduced to the following lemma.
\end{proof}

\begin{lemma}[3.4.1.2]
\phantomsection
\label{IV.3.4.1.2}
Let $A$ be a Noetherian local ring of dimension $>0$, let $\mathfrak p_i$ be the minimal prime ideals of $A$, let $\mathfrak m$ be its maximal ideal, and let $t\in\mathfrak m$ be such that $A/tA$ is Artinian.
Then, for every $A$-module of finite type $M$,
\[
\label{IV.3.4.1.3}
\operatorname{length}(M/tM)
  \geqslant
  \sum_i\operatorname{length}(M_{\mathfrak p_i}).
\tag{3.4.1.3}
\]
Moreover, equality holds if and only if the following two conditions are satisfied:
\begin{enumerate}
  \item[\emph{$\alpha$)}] $t$ is $M$-regular;
  \item[\emph{$\beta$)}] for every $i$ such that $M_{\mathfrak p_i}\neq 0$, the image of $t$ in $A/\mathfrak p_i$ generates the maximal ideal of that ring (which implies that $A/\mathfrak p_i$ is a discrete valuation ring).
\end{enumerate}
\end{lemma}

\oldpage[IV-2]{47}

\begin{proof}
Since $A$ does not have dimension $0$ and $A/tA$ is Artinian, necessarily $\dim(A)=1$ \sref[0]{0.16.3.4}, and $t\notin\mathfrak p_i$ for every $i$.
The principal ideal $(t)$ is therefore an ideal of definition of $A$, and hence contains a power of the maximal ideal $\mathfrak m$.
Let $N$ be the submodule of $M$ consisting of the elements annihilated by a power of $t$ (or, equivalently in view of what we have just proved, by a power of $\mathfrak m$), and put $P=M/N$.
Then $t$ is $P$-regular: if $tx\in N$ for some $x\in M$, then $t^k(tx)=0$ for some integer $k$, and hence $x\in N$.
We shall use the following elementary lemma.

\begin{lemma}[3.4.1.4]
\phantomsection
\label{IV.3.4.1.4}
Let $A$ be a ring and
\[
0\longrightarrow M'\longrightarrow M\longrightarrow M''\longrightarrow 0
\]
an exact sequence of $A$-modules.
If $t\in A$ is $M''$-regular, then the sequence
\[
0\longrightarrow M'/tM'
 \longrightarrow M/tM
 \longrightarrow M''/tM''
 \longrightarrow 0
\]
is exact.
\end{lemma}

\begin{proof}
Since $M/tM=M\otimes_A(A/tA)$, it suffices to prove exactness at $M'/tM'$.
Suppose that the image $x\in M$ of an element $x'\in M'$ is of the form $x=ty$, with $y\in M$.
For the images $x'',y''$ of $x,y$ in $M''$, we have $x''=ty''$.
But $x''=0$, so the hypothesis gives $y''=0$.
Thus $y$ is the image of an element $y'\in M'$, and the relation $x=ty$ gives $x'=ty'$ because $M'\to M$ is injective.
\end{proof}

It follows from this lemma that
\[
\label{IV.3.4.1.5}
\operatorname{length}(M/tM)
  =
  \operatorname{length}(N/tN)
  +
  \operatorname{length}(P/tP).
\tag{3.4.1.5}
\]
On the other hand, $N_{\mathfrak p_i}=0$ for every $i$, since $t\notin\mathfrak p_i$, and hence $M_{\mathfrak p_i}=P_{\mathfrak p_i}$.
Thus, to prove \sref{IV.3.4.1.3}, it suffices to prove it with $P$ in place of $M$.
Furthermore, if the two sides of \sref{IV.3.4.1.3} are equal, then the same inequality for $P$, together with \sref{IV.3.4.1.5}, gives $\operatorname{length}(N/tN)=0$.
Consequently $N/tN=0$, and Nakayama's lemma gives $N=0$, since $N$ is of finite type.
But $N=0$ means that $t$ is $M$-regular.
We may therefore reduce to the case $M=P$, that is, suppose from the outset that $t$ is $M$-regular.
This implies that $\mathfrak m\notin\operatorname{Ass}(M)$, since $t$ cannot annihilate a nonzero element of $M$.
Because $\dim(A)=1$, it follows that
\[
\operatorname{Ass}(M)\subset\{\mathfrak p_i\}_i.
\]

We now argue by induction on
\[
n=\sum_i\operatorname{length}(M_{\mathfrak p_i}).
\]
If $n=0$, then $M_{\mathfrak p_i}=0$ for every $i$, and hence $M=0$, since none of the $\mathfrak p_i$ belongs to $\operatorname{Ass}(M)$.
Both sides of \sref{IV.3.4.1.3} are then zero, and condition \emph{$\beta$)} is vacuous.
If $n>0$, the argument at the beginning of the proof of \sref{IV.3.4.1} allows us, moreover, to suppose that $M$ is faithful.
It follows that $M_{\mathfrak p_i}\neq 0$ for every $i$ (Bourbaki, \emph{Alg.\ comm.}, chap.~II, \textsection~2, no~2, cor.~2 of prop.~4), and consequently
\[
\operatorname{Ass}(M)=\{\mathfrak p_i\}_i.
\]

\oldpage[IV-2]{48}

Suppose first that $n=1$.
There is then only one minimal prime ideal $\mathfrak p$ of $A$, and the assertion that $M_{\mathfrak p}$ has length $1$ means that it is isomorphic, as an $A_{\mathfrak p}$-module, to the residue field
\[
k=A_{\mathfrak p}/\mathfrak pA_{\mathfrak p}.
\]
Thus $M_{\mathfrak p}$ is annihilated by $\mathfrak pA_{\mathfrak p}$, and consequently $\mathfrak p$ is the annihilator of $M$ (Bourbaki, \emph{Alg.\ comm.}, chap.~II, \textsection~2, no~4, formula~(9)).
Since $M$ is faithful, $\mathfrak p=0$, and hence $A$ is integral.
Because $M\neq 0$, Nakayama's lemma gives $M/tM\neq 0$, so that $\operatorname{length}(M/tM)\geqslant 1$; this proves \sref{IV.3.4.1.3} in this case.
Moreover, if $\operatorname{length}(M/tM)=1$, then $M$ is cyclic (Bourbaki, \emph{Alg.\ comm.}, chap.~II, \textsection~3, no~2, cor.~2 of prop.~4), and hence isomorphic to a quotient $A/\mathfrak b$.
Since $M$ is faithful, necessarily $\mathfrak b=0$, and $M$ is isomorphic to $A$.
As $\operatorname{length}(A/tA)=1$, the ideal $tA$ is the maximal ideal $\mathfrak m$.
Since $A$ is a Noetherian integral local ring, it follows that $A$ is a discrete valuation ring (Bourbaki, \emph{Alg.\ comm.}, chap.~VI, \textsection~3, no~6, prop.~9), with uniformizer $t$.
Conversely, if $A$ is a discrete valuation ring, $t$ is its uniformizer, $\operatorname{length}(M_{\mathfrak p})=1$, and $t$ is $M$-regular, then $M$ is torsion-free and therefore isomorphic to a submodule of $A$ (as $M$ is of finite type), hence to $A$ itself.
Thus
\[
\operatorname{length}(M/tM)=\operatorname{length}(A/tA)=1.
\]

Suppose now that $n\geqslant 2$.
There exists an exact sequence
\[
0\longrightarrow M'\longrightarrow M\longrightarrow M''\longrightarrow 0
\]
with $M'\neq 0$, $M''\neq 0$, and
\[
\operatorname{Ass}(M)=\operatorname{Ass}(M')\cup\operatorname{Ass}(M'').
\]
Indeed, if $\operatorname{Ass}(M)$ has more than one element, this follows from Bourbaki, \emph{Alg.\ comm.}, chap.~IV, \textsection~1, no~1, prop.~4.
If, on the contrary, $\operatorname{Ass}(M)$ consists of a single prime ideal, that ideal is necessarily the unique minimal prime ideal $\mathfrak p$ of $A$.
The hypothesis then gives $\operatorname{length}(M_{\mathfrak p})\geqslant 2$, and it is enough to take for $M'$ the inverse image of a submodule of $M_{\mathfrak p}$ distinct from both $0$ and $M_{\mathfrak p}$.
Since $t$ is $M$-regular, it belongs to none of the prime ideals in $\operatorname{Ass}(M)$ (Bourbaki, \emph{Alg.\ comm.}, chap.~IV, \textsection~1, no~1, cor.~2 of prop.~2).
For the same reason, $t$ is both $M'$-regular and $M''$-regular.
By \sref{IV.3.4.1.4}, the sequence
\[
0\longrightarrow M'/tM'
 \longrightarrow M/tM
 \longrightarrow M''/tM''
 \longrightarrow 0
\]
is exact.
For every $i$, so is the sequence
\[
0\longrightarrow M'_{\mathfrak p_i}
 \longrightarrow M_{\mathfrak p_i}
 \longrightarrow M''_{\mathfrak p_i}
 \longrightarrow 0.
\]
Consequently
\[
\begin{aligned}
\operatorname{length}(M/tM)
  &=
  \operatorname{length}(M'/tM')
  +
  \operatorname{length}(M''/tM''),\\
\operatorname{length}(M_{\mathfrak p_i})
  &=
  \operatorname{length}(M'_{\mathfrak p_i})
  +
  \operatorname{length}(M''_{\mathfrak p_i}).
\end{aligned}
\]
The induction hypothesis now proves \sref{IV.3.4.1.3}.
Moreover, the two sides can be equal only if the analogous inequalities for $M'$ and $M''$ are both equalities.
By the induction hypothesis, this is equivalent to condition \emph{$\beta$)} for those $\mathfrak p_i$ such that $M'_{\mathfrak p_i}\neq 0$ or $M''_{\mathfrak p_i}\neq 0$; these are precisely the ideals for which $M_{\mathfrak p_i}\neq 0$.
\end{proof}

\oldpage[IV-2]{49}

\begin{corollary}[3.4.2]
\phantomsection
\label{IV.3.4.2}
Under the general hypotheses of \sref{IV.3.4.1}, suppose that $Z$ is not equal to any of the $S_i$ and that $\operatorname{length}((\sh{F}/t\sh{F})_z)=1$.
Then there is a unique $S_i$ containing $Z$, and for this value of $i$, $\operatorname{length}(\sh{F}_{s_i})=1$; moreover, $\sh{O}_{S_i,z}$ is a discrete valuation ring with uniformizer $t_z$, and $t_z$ is $\sh{F}_z$-regular.
\end{corollary}

\begin{proof}
This follows from \sref{IV.3.4.1}, the two sides of \sref{IV.3.4.1.1} being then equal.
\end{proof}

\begin{proposition}[3.4.3]
\phantomsection
\label{IV.3.4.3}
Let $X$ be a locally Noetherian prescheme, $t$ a section of $\sh{O}_X$ over $X$, and $Y$ the closed sub-prescheme of $X$ defined by the ideal $t\sh{O}_X$.
Let $\sh{F}$ be a coherent $\sh{O}_X$-module, let $T$ be an associated prime cycle of $\sh{F}$, let $T'$ be an irreducible component of $T\cap Y$, and let $x$ be the generic point of $T'$.
Suppose that $t_x$ is $\sh{F}_x$-regular; then $x\in\operatorname{Ass}(\sh{F}/t\sh{F})$.
\end{proposition}

\begin{proof}
As in the proof of \sref{IV.3.4.1}, we may reduce to the case $X=\operatorname{Spec}(\sh{O}_x)$.
Taking \sref{IV.3.1.2} into account, the proposition is then a consequence of \sref[0]{0.16.4.6.3}.
\end{proof}

\begin{proposition}[3.4.4]
\phantomsection
\label{IV.3.4.4}
Let $X$ be a locally Noetherian prescheme, $t$ a section of $\sh{O}_X$ over $X$, and $Y$ the closed sub-prescheme of $X$ defined by the ideal $t\sh{O}_X$.
Let $\sh{F}$ be a coherent $\sh{O}_X$-module, and let $(S_i)$ be the family of irreducible components of $\operatorname{Supp}(\sh{F})$.
Let $y$ be a point of $Y$ such that $t_y$ is $\sh{F}_y$-regular and no embedded associated prime cycle of $\sh{F}/t\sh{F}$ contains $y$.
Then the irreducible components of $\operatorname{Supp}(\sh{F}/t\sh{F})$ containing $y$ are precisely the irreducible components of the $S_i\cap Y$ containing $y$, and the associated prime cycles of $\sh{F}$ containing $y$ are non-embedded.
\end{proposition}

\begin{proof}
Let us first prove the last assertion.
Let $T\supset T_1$ be two associated prime cycles of $\sh{F}$ containing $y$.
If $x$ is the generic point of an irreducible component of $T\cap Y$ containing $y$, then $x$ is a generization of $y$, and hence belongs to every neighbourhood of $y$.
The hypothesis that $t_y$ is $\sh{F}_y$-regular therefore implies that $t_x$ is $\sh{F}_x$-regular \sref[0]{0.15.2.4}, so \sref{IV.3.4.3} gives $x\in\operatorname{Ass}(\sh{F}/t\sh{F})$.
Let $x_1$ be the generic point of an irreducible component of $T_1\cap Y$ containing $y$, and let $x$ be the generic point of an irreducible component of $T\cap Y$ containing $x_1$.
The preceding argument shows that $x_1$ and $x$ both belong to $\operatorname{Ass}(\sh{F}/t\sh{F})$.
Since $x_1\in\overline{\{x\}}$, the hypothesis implies $x_1=x$.

Again denote by $T$ and $T_1$ the integral closed sub-preschemes of $X$ having these spaces as their respective underlying spaces, and set
\[
A=\sh{O}_{T,x},
\qquad
A_1=\sh{O}_{T_1,x}.
\]
Then $A_1=A/\mathfrak p$ for a prime ideal $\mathfrak p$ of $A$.
By the definitions of $x$ and $x_1$, the rings $A/tA$ and $A_1/tA_1$ are Artinian.
On the other hand, we saw above that $t_x$ is $\sh{F}_x$-regular, so $x\notin\operatorname{Ass}(\sh{F})$, and consequently $A_1$ is not Artinian.
Thus $\dim A=\dim A_1=1$.
This implies $\mathfrak p=0$ and $A=A_1$ \sref[0]{0.16.1.2.2}.
Since $\operatorname{Spec}(A)$ and $\operatorname{Spec}(A_1)$ are respectively dense in $T$ and $T_1$, it follows that $T=T_1$.

The $S_i$ containing $y$ are therefore all the associated prime cycles of $\sh{F}$ containing $y$.
If $x$ is the generic point of an irreducible component of $S_i\cap Y$ containing $y$, then \sref[0]{0.15.2.4} again shows that $t_x$ is $\sh{F}_x$-regular, and \sref{IV.3.4.3} gives $x\in\operatorname{Ass}(\sh{F}/t\sh{F})$.
This proves the first assertion of \sref{IV.3.4.4}.
\end{proof}

\begin{proposition}[3.4.5]
\phantomsection
\label{IV.3.4.5}
Let $X$ be a locally Noetherian prescheme, $t$ a section of $\sh{O}_X$ over $X$, and $Y$ the closed sub-prescheme of $X$ defined by the ideal $t\sh{O}_X$.
Let $\sh{F}$ be a coherent $\sh{O}_X$-module and $y$ a point of $Y$.
\oldpage[IV-2]{50}
Suppose that $t_y$ is $\sh{F}_y$-regular and that $\sh{F}/t\sh{F}$ is integral at $y$ \sref{IV.3.2.4}.
Then $\sh{F}$ is integral at the point $y$.
\end{proposition}

\begin{proof}
Taking into account \sref{IV.3.4.4}, it suffices to prove that $y$ is contained in a \emph{unique} irreducible component of $\operatorname{Supp}(\sh{F})$, and that if $s$ is the generic point of this component, $\operatorname{length}(\sh{F}_s)=1$.
By hypothesis, $y$ belongs to only one irreducible component of $\operatorname{Supp}(\sh{F}/t\sh{F})$, and if $z$ is the generic point of this component, then $\operatorname{length}((\sh{F}/t\sh{F})_z)=1$.
The conclusion follows from \sref{IV.3.4.2}.
\end{proof}

\begin{proposition}[3.4.6]
\phantomsection
\label{IV.3.4.6}
The hypotheses being those of \sref{IV.3.4.1}, let $x$ be a point of $Y$.
Suppose that $Y$ does not contain any of the $S_i$ containing $x$, and that $\sh{F}/t\sh{F}$ is reduced at the point $x$ \sref{IV.3.2.2}.
Then $t_x$ is $\sh{F}_x$-regular and $\sh{F}$ is reduced at the point $x$.
Moreover, if $z_j$ is the generic point of an irreducible component of $\operatorname{Supp}(\sh{F}/t\sh{F})$ containing $x$, then $z_j$ is contained in a unique $S_i$, and $\sh{O}_{S_i,z_j}$ is a discrete valuation ring with uniformizer $t_{z_j}$.
\end{proposition}

\begin{proof}
The assertion that $t_x$ is $\sh{F}_x$-regular follows by applying the following lemma to the ring $\sh{O}_x$.

\begin{lemma}[3.4.6.1]
\phantomsection
\label{IV.3.4.6.1}
Let $A$ be a Noetherian ring, $M$ an $A$-module of finite type, $\mathfrak p_i$ the minimal elements of $\operatorname{Supp}(M)$, and $t\in A$.
Suppose that $t$ belongs to none of the $\mathfrak p_i$ and that $M/tM$ is a reduced $A$-module \sref{IV.3.2.2}.
Then $t$ is $M$-regular.
\end{lemma}

\begin{proof}
Every prime ideal $\mathfrak p\in\operatorname{Supp}(M)$ contains one of the $\mathfrak p_i$.
Since $t$ belongs to none of the $\mathfrak p_i$, multiplication by $t$ on $M_{\mathfrak p}$ is not nilpotent (Bourbaki, \emph{Alg.\ comm.}, chap.~IV, \textsection~1, no~4, cor.\ of prop.~9).
Let $N$ be the submodule of $M$ consisting of the elements annihilated by a power of $t$, and set $P=M/N$.
We shall prove that $N=0$.
Since $t$ is $P$-regular, \sref{IV.3.4.1.4} gives an exact sequence
\[
0\longrightarrow N/tN
 \longrightarrow M/tM
 \longrightarrow P/tP
 \longrightarrow 0.
\]
Since $N$ is of finite type, it is annihilated by a power of $t$, and it therefore suffices to prove that $N/tN=0$.
As $N/tN$ is a submodule of $M/tM$, it suffices to prove that $(N/tN)_{\mathfrak p}=0$ for every $\mathfrak p\in\operatorname{Ass}(M/tM)$, or equivalently that
\[
u_{\mathfrak p}:(M/tM)_{\mathfrak p}\longrightarrow(P/tP)_{\mathfrak p}
\]
is bijective for every $\mathfrak p\in\operatorname{Ass}(M/tM)$.
Now $(P/tP)_{\mathfrak p}\neq 0$.
Indeed,
\[
\mathfrak p\in\operatorname{Supp}(M/tM)
  =\operatorname{Supp}(M)\cap V(t)
\]
by \sref[0]{0.1.7.5}, so the image of $t$ in $A_{\mathfrak p}$ lies in the maximal ideal $\mathfrak pA_{\mathfrak p}$.
The equality $P_{\mathfrak p}=tP_{\mathfrak p}$ would therefore imply $P_{\mathfrak p}=0$ by Nakayama's lemma.
We would then have $M_{\mathfrak p}=N_{\mathfrak p}$, and multiplication by $t$ on $M_{\mathfrak p}$ would be nilpotent, contradicting the observation above because $\mathfrak p\in\operatorname{Supp}(M)$.
Finally, since $M/tM$ is reduced, $\operatorname{length}((M/tM)_{\mathfrak p})=1$.
As $(P/tP)_{\mathfrak p}\neq 0$, the map $u_{\mathfrak p}$ is necessarily bijective.
This proves the lemma.
\end{proof}

By hypothesis, no embedded associated prime cycle of $\sh{F}/t\sh{F}$ contains $x$.
It follows from \sref{IV.3.4.4} that no embedded associated prime cycle of $\sh{F}$ contains $x$.
On the other hand, applying \sref{IV.3.4.2} to an irreducible component of $\operatorname{Supp}(\sh{F}/t\sh{F})$ containing $x$, we find that $\operatorname{length}(\sh{F}_{s_i})=1$ for every $S_i$ containing $x$.
This completes the proof that $\sh{F}_x$ is reduced.
The final assertions also follow from \sref{IV.3.4.2}.
\end{proof}

\begin{corollary}[3.4.7]
\phantomsection
\label{IV.3.4.7}
Let $A$ be a Noetherian local ring, $\mathfrak m$ its maximal ideal, $M$ an $A$-module of finite type, and $(x_i)_{1\leqslant i\leqslant k}$ a family of elements of $\mathfrak m$ forming part of a system of parameters for $M$ \sref[0]{0.16.3.6}.
\oldpage[IV-2]{51}
If the $A$-module
\[
N=M\Big/\Big(\sum_{i=1}^k x_iM\Big)
\]
is integral \sref{IV.3.2.4}, then $M$ is integral and the sequence $(x_i)_{1\leqslant i\leqslant k}$ is $M$-regular.
\end{corollary}

\begin{proof}
Induction on $k$ immediately reduces us to the case $k=1$; write $x$ for $x_1$.
The hypothesis that $x$ forms part of a system of parameters for $M$ implies
\[
\dim(N)=\dim(M)-1
\]
by \sref[0]{0.16.3.7}.
Put $n=\dim(N)$.
There is therefore a minimal element $\mathfrak p$ of $\operatorname{Supp}(M)$ such that
\[
\dim(M/\mathfrak pM)=n+1
\]
by \sref[0]{0.16.3.4}; for every integer $j>0$ one also has
\[
\dim(M/\mathfrak p^jM)=n+1
\]
by \sref[0]{0.16.3.5}.
Moreover, $x$ forms part of a system of parameters for $M/\mathfrak p^jM$ \sref[0]{0.16.3.5}.
Set
\[
M'=M/\mathfrak p^jM,
\qquad
N'=M'/xM';
\]
then $\dim(N')=n$.

There is an evident surjective homomorphism $v:N\to N'$.
We claim that it is bijective.
Indeed, if $P=\operatorname{Ker}(v)$, then $\operatorname{Ass}(P)\subset\operatorname{Ass}(N)$.
Since $N$ is integral, if $P\neq 0$, then $\operatorname{Ass}(P)$ and $\operatorname{Ass}(N)$ would both consist of the unique point $\mathfrak q$ of
\[
\operatorname{Ass}(N)=\operatorname{Supp}(N).
\]
But $\dim(N')=\dim(N)$ implies $N'_{\mathfrak q}\neq 0$, and $\operatorname{length}(N_{\mathfrak q})=1$ implies $\operatorname{length}(N'_{\mathfrak q})=1$, since $N_{\mathfrak q}\to N'_{\mathfrak q}$ is surjective.
Thus $P_{\mathfrak q}=0$, a contradiction.
Hence $v$ is bijective.

It follows that $N'$, being isomorphic to $N$, is integral.
Moreover, the support of $N'$, which is the intersection of $\operatorname{Supp}(M')$ with $V(x)$, cannot contain $\operatorname{Supp}(M')$, and the latter is irreducible by construction.
Since $M'/xM'$ is integral, and hence reduced, \sref{IV.3.4.6} shows that $x$ is $M'$-regular.
Consequently the kernel of multiplication by $x$ on $M$ is contained in
\[
\mathfrak p^jM\subset\mathfrak m^jM
\]
for every integer $j$, and this kernel is therefore zero by \sref[0]{0.7.3.5}.
Thus $x$ is $M$-regular.
Applying \sref{IV.3.4.5}, we conclude that $M$ is integral.
\end{proof}

\begin{remark}[3.4.8]
\phantomsection
\label{IV.3.4.8}
The analogue of \sref{IV.3.4.7} obtained by replacing ``integral'' by ``reduced'' need not be true.
For example, let $C=K[X,Y,Z]$ be the polynomial ring over a field $K$, let
\[
B=C/\mathfrak p\mathfrak q,
\qquad
\mathfrak p=CZ,
\qquad
\mathfrak q=CX^2+CY,
\]
and let $A$ be the local ring of $B$ corresponding to the maximal ideal that is the image of $CX+CY+CZ$ in $B$.
If $x,z$ are the canonical images of $X,Z$ in $A$, then $xz\neq 0$ but $x^2z^2=0$.
On the other hand, $A/xA$ is isomorphic to $K[Y,Z]/(YZ)$, so $\dim(A/xA)=1$ whereas $\dim(A)=2$.
Thus $x$ belongs to a system of parameters of $A$ \sref[0]{0.16.3.4}, and $A/xA$ is reduced, but $A$ is not.
\end{remark}

\begin{proposition}[3.4.9]
\phantomsection
\label{IV.3.4.9}
Let $A$ be a Noetherian ring, $M$ an $A$-module of finite type, and $f\in A$ an $M$-regular element such that $M/fM$ has no embedded associated prime ideals.
If $\mathfrak p_i$ $(1\leqslant i\leqslant m)$ are the prime ideals associated to $M/fM$, then, for every integer $n>0$, $f^nM$ is the intersection of the inverse images of the $f^nM_{\mathfrak p_i}$ under the canonical maps $M\to M_{\mathfrak p_i}$ $(1\leqslant i\leqslant m)$.
\end{proposition}

\begin{proof}
It suffices to prove that the $\mathfrak p_i$ are also the prime ideals associated to $M/f^nM$, for then the $\mathfrak p_i$-saturations of $f^nM$ are the components of the necessarily unique reduced primary decomposition of $f^nM$ in $M$.
Now, we have
\[
\operatorname{Ass}(f^{n-1}M/f^nM)
\subset
\operatorname{Ass}(M/f^nM)
\subset
\operatorname{Ass}(M/f^{n-1}M)
\cup
\operatorname{Ass}(f^{n-1}M/f^nM)
\]
by \sref{IV.3.1.7}.
Since $f$ is $M$-regular, $f^{n-1}M/f^nM$ is isomorphic to $M/fM$, and the result follows by induction on $n$.
\end{proof}
% END NATIVE INPUT ega4-3.source_aligned_r1
% BEGIN NATIVE INPUT ega4-4.source_aligned_r1
% BEGIN NATIVE INPUT components/01_section41_dimension_of_algebraic_preschemes
\section{Change of base field for algebraic preschemes}
\label{section:IV.4}

\oldpage[IV-2]{52}

\subsection{Dimension of algebraic preschemes}
\label{subsection:IV.4.1}

We shall develop in \S~5 the general theory of the dimension of preschemes.
The theory of the dimension of algebraic preschemes, however, can be developed more elementarily and also has many special features; we therefore give here a rapid account independent of the general theory.

If $K$ is a field and $L$ an extension of $K$, we write $\operatorname{tr.deg}_K L$ for the transcendence degree of $L$ over $K$.

\begin{definition}[4.1.1]
\phantomsection
\label{IV.4.1.1}
Let $k$ be a field and $X$ a prescheme locally of finite type over $k$.
The \emph{dimension} of $X$ is the number
\begin{equation}
\dim(X)=\sup_x\operatorname{tr.deg}_k\kappa(x),
\tag{4.1.1.1}
\label{IV.4.1.1.1}
\end{equation}
where $x$ runs through the set of maximal points of $X$.
\end{definition}

We shall see in \sref{IV.5.2.2} that Definition \sref{IV.4.1.1} depends only apparently on the base field $k$ over which $X$ is locally of finite type, and that the number $\dim(X)$ so defined agrees with the dimension of the underlying space of $X$ \sref[0]{0.14.1.2}.
Evidently,
\[
\dim(X)=\dim(X_{\mathrm{red}}).
\]

Each $\kappa(x)$ is an extension of finite type of $k$ \sref[I]{I.6.3.3}, hence has finite transcendence degree over $k$.
If $X$ is of finite type over $k$ and nonempty, it is Noetherian \sref[I]{I.6.3.7}, hence has only finitely many irreducible components; consequently $\dim(X)$ is finite and nonnegative.
Definition \sref{IV.4.1.1} gives
\[
\dim(\emptyset)=-\infty.
\]

If $(X_\alpha)$ denotes the family of reduced closed sub-preschemes of $X$ whose underlying spaces are the irreducible components of $X$ \sref[I]{I.5.2.1}, then clearly
\begin{equation}
\dim(X)=\sup_\alpha\dim(X_\alpha).
\tag{4.1.1.2}
\label{IV.4.1.1.2}
\end{equation}
The calculation of the dimension is thus reduced to the case of integral preschemes locally of finite type over $k$.
Finally,
\begin{equation}
\dim(X)=\dim(U)
\tag{4.1.1.3}
\label{IV.4.1.1.3}
\end{equation}
for every induced sub-prescheme on an everywhere dense open subset $U$ of $X$.
The notion of dimension is thereby reduced to the case of an affine scheme of finite type over $k$.

\begin{theorem}[4.1.2]
\phantomsection
\label{IV.4.1.2}
Let $X$ and $Y$ be two preschemes locally of finite type over a field $k$, and let $f:X\to Y$ be a $k$-morphism.

\oldpage[IV-2]{53}

\begin{enumerate}[(i)]
\item If $f$ is quasi-compact and dominant, then
\[
\dim(Y)\leq\dim(X).
\]
\item If $f$ is quasi-finite, then
\[
\dim(X)\leq\dim(Y).
\]
\item Suppose that $X$ is of finite type over $k$.
For $\dim(X)\geq n$ (respectively $\dim(X)\leq n$, respectively $\dim(X)=n$), it is necessary and sufficient that there exist an open subset $U$ dense in $X$ and a $k$-morphism
\[
g:U\longrightarrow\mathbf{V}_k^n
=\Spec\bigl(k[T_1,\ldots,T_n]\bigr)
\]
that is surjective (respectively finite, respectively finite and surjective).
If $X$ is affine, one may take $U=X$.
\end{enumerate}
\end{theorem}

\begin{proof}
\emph{(i)} Let $y$ be a maximal point of $Y$.
The fibre $f^{-1}(y)$ contains a maximal point $x$ of $X$ \sref{IV.1.1.5}; hence $\kappa(x)$ is an extension of $\kappa(y)$, and therefore
\[
\operatorname{tr.deg}_k\kappa(y)
\leq
\operatorname{tr.deg}_k\kappa(x)
\leq
\dim(X).
\]
This proves (i).

\emph{(ii)} If $x$ is a maximal point of $X$, then $\kappa(x)$ is a finite extension of $\kappa(f(x))$ \sref[II]{II.6.2.2}, and hence has the same transcendence degree over $k$.
On considering the reduced closed sub-prescheme of $Y$ whose underlying space is $\overline{\{f(x)\}}$, we are reduced to the following corollary.
\end{proof}

\begin{corollary}[4.1.2.1]
\phantomsection
\label{IV.4.1.2.1}
Let $Y$ be a $k$-prescheme locally of finite type.
For every sub-prescheme $Z$ of $Y$, one has
\[
\dim(Z)\leq\dim(Y).
\]
Suppose in addition that all the irreducible components of $Y$ have the same dimension.
Then $Z$ is rare in $Y$ if and only if
\[
\dim(Z)<\dim(Y).
\]
\end{corollary}

\begin{proof}
Let $z$ be a maximal point of $Z$.
By considering one of the reduced closed sub-preschemes of $Y$ whose underlying space is an irreducible component of $Y$ containing $z$, we reduce to the case in which $Y$ is integral, with generic point $y$.
Then, by taking in $Y$ an affine open neighbourhood of $z$, we reduce to the case in which $Y$ is affine and of finite type over $k$.
Write
\[
Y=\Spec(A),
\qquad
Z=\Spec(A/\mathfrak a),
\]
where $A$ is an integral ring and $\mathfrak a$ is an ideal of $A$ distinct from $A$.

By the normalization lemma (Bourbaki, \emph{Alg.\ comm.}, chap.~V, \S~3, no.~1, th.~1), there exists a finite sequence $(x_i)_{1\leq i\leq n}$ of elements of $A$, algebraically independent over $k$, such that $A$ is integral over
\[
B=k[x_1,\ldots,x_n]
\]
and $\mathfrak a\cap B$ is generated by a subfamily $(x_i)_{1\leq i\leq p}$, possibly empty, of $(x_i)_{1\leq i\leq n}$.
Let
\[
g:Y\longrightarrow\Spec(B)=\mathbf{V}_k^n
\]
be the finite dominant morphism corresponding to the canonical injection $B\to A$.
Since
\[
C=B/(\mathfrak a\cap B)
\cong k[x_{p+1},\ldots,x_n],
\]
$g$ induces on $Z$ a finite dominant morphism
\[
Z\longrightarrow\mathbf{V}_k^{n-p}.
\]
Thus
\[
\operatorname{tr.deg}_k\kappa(y)=n,
\qquad
\operatorname{tr.deg}_k\kappa(z)=n-p,
\]
because $\kappa(y)$ and $\kappa(z)$ are finite extensions of $\kappa(g(y))$ and $\kappa(g(z))$, respectively.
This proves the first assertion.

If $p=0$, then necessarily $z=y$, since $y$ is the only point of $Y$ whose image under $g$ is the generic point of $\mathbf{V}_k^n$; in this case $Z$ contains a nonempty open subset of $Y$.
If $p>0$, then $z\neq y$, and $\overline{\{z\}}$ is rare in $\overline{\{y\}}$.
This completes the proof.
\end{proof}

\begin{proof}[Proof of \sref{IV.4.1.2}, continued]
The corollary proves (ii).

\emph{(iii)} The sufficiency of the stated conditions follows immediately from (i), (ii), and \sref{IV.1.5.4}(v).
To prove necessity, choose a dense open subset $U$ of $X$ that is the union of pairwise disjoint irreducible affine open subsets $U_i$ $(1\leq i\leq m)$, each containing a maximal point of $X$ \sref[0]{0.2.1.6}.
We may further suppose that $X$ is reduced; if $X$ is affine, we may of course take $U=X$.

\oldpage[IV-2]{54}

The same argument as in \sref{IV.4.1.2.1} shows that, if $\dim(U_i)=n_i$, there exists a $k$-morphism
\[
g_i:U_i\longrightarrow\mathbf{V}_k^{n_i}
\]
that is finite and dominant, hence surjective \sref[II]{II.6.1.10}.
Set
\[
n=\dim(X)=\sup_i n_i.
\]
For every $n_i$ there is a morphism
\[
h_i:\mathbf{V}_k^{n_i}\longrightarrow\mathbf{V}_k^n
\]
which is the closed immersion corresponding to the canonical homomorphism
\[
k[T_1,\ldots,T_n]\longrightarrow k[T_1,\ldots,T_{n_i}]
\]
and is the identity when $n_i=n$.
Since $U$ is the disjoint union of the $U_i$, define $g$ on each $U_i$ to be $h_i\circ g_i$.
Then $g$ is finite and surjective.

When $n'\geq n$, composing $g$ with the canonical closed immersion
\[
\mathbf{V}_k^n\longrightarrow\mathbf{V}_k^{n'}
\]
gives a finite morphism $U\to\mathbf{V}_k^{n'}$.
When $n'\leq n$, compose $g$ with the canonical morphism
\[
p:\mathbf{V}_k^n\longrightarrow\mathbf{V}_k^{n'}
\]
corresponding to the canonical injection
\[
k[T_1,\ldots,T_{n'}]\longrightarrow k[T_1,\ldots,T_n].
\]
The morphism $p$ is faithfully flat, hence surjective.
\end{proof}

\begin{remark}[4.1.3]
\phantomsection
\label{IV.4.1.3}
Corollary \sref{IV.4.1.2.1} shows that in formula \sref{IV.4.1.1.1} one may let $x$ run through the set of all points of $X$.
\end{remark}

\begin{corollary}[4.1.4]
\phantomsection
\label{IV.4.1.4}
Let $X$ be a prescheme locally of finite type over a field $k$, and let $K$ be an extension of $k$.
Then
\[
\dim(X\otimes_k K)=\dim(X).
\]
\end{corollary}

\begin{proof}
It is enough to consider the case in which $X$ is of finite type over $k$.
The morphism
\[
u:X\otimes_k K\longrightarrow X
\]
is faithfully flat \sref{IV.2.2.13}(i).
Thus, if $U$ is an everywhere dense open subset of $X$, then
\[
u^{-1}(U)=U\otimes_k K
\]
is dense in $X\otimes_k K$ \sref{IV.2.3.10}.
If
\[
g:U\longrightarrow\mathbf{V}_k^n
\]
is finite and surjective, then so is
\[
g_{(K)}:U\otimes_k K\longrightarrow\mathbf{V}_K^n
\]
\sref[I]{I.3.5.2} and \sref[II]{II.6.1.5}.
The corollary follows.
\end{proof}

\begin{corollary}[4.1.5]
\phantomsection
\label{IV.4.1.5}
Let $X$ and $Y$ be two preschemes locally of finite type over a field $k$.
Then
\[
\dim(X\times_k Y)=\dim(X)+\dim(Y).
\]
\end{corollary}

\begin{proof}
It is enough to prove the following.
Let $U$ and $V$ be affine neighbourhoods of points $x\in X$ and $y\in Y$ such that
\[
\operatorname{tr.deg}_k\kappa(x)=m=\dim(U),
\qquad
\operatorname{tr.deg}_k\kappa(y)=n=\dim(V).
\]
Then
\[
\dim(U\times_k V)=m+n.
\]
Equivalently, by \sref{IV.4.1.2}, we may suppose that there exist finite surjective $k$-morphisms
\[
f:X\longrightarrow\mathbf{V}_k^m,
\qquad
g:Y\longrightarrow\mathbf{V}_k^n.
\]
The product morphism
\[
f\times g:X\times_k Y\longrightarrow\mathbf{V}_k^{m+n}
\]
is finite and surjective \sref[I]{I.3.5.2} and \sref[II]{II.6.1.5}, and the result follows.
\end{proof}
% END NATIVE INPUT components/01_section41_dimension_of_algebraic_preschemes
% BEGIN NATIVE INPUT components/02_section42_associated_prime_cycles
\subsection{Associated prime cycles on algebraic preschemes}
\label{subsection:IV.4.2}

\begin{proposition}[4.2.1]
\phantomsection
\label{IV.4.2.1}
Let $K$ and $L$ be two extensions of a field $k$ such that
$K\otimes_k L$ is Noetherian.
Then the prime ideals associated with $K\otimes_k L$ are minimal, and if
$E$ is the residue field of the local ring at such an ideal, then
\begin{equation}
\operatorname{tr.deg}_K E=\operatorname{tr.deg}_k L,
\qquad
\operatorname{tr.deg}_L E=\operatorname{tr.deg}_k K,
\tag{4.2.1.1}
\label{IV.4.2.1.1}
\end{equation}
whence
\begin{equation}
\operatorname{tr.deg}_k E
=\operatorname{tr.deg}_k K+\operatorname{tr.deg}_k L.
\tag{4.2.1.2}
\label{IV.4.2.1.2}
\end{equation}
\end{proposition}

\begin{proof}
It is known that $K$ is an algebraic extension of a purely transcendental
extension
\[
K'=k(\mathbf t),
\qquad
\mathbf t=(t_\alpha)_{\alpha\in I},
\]
where $(t_\alpha)_{\alpha\in I}$ is a family of indeterminates.
The ring
\[
k[\mathbf t]\otimes_k L=L[\mathbf t]
\]
is integral, and hence so is $K'\otimes_k L$, which is a ring of fractions
of it; the field of fractions of $K'\otimes_k L$ is $L(\mathbf t)$.
There is a commutative diagram of canonical homomorphisms

\oldpage[IV-2]{55}

\begin{equation}
\xymatrix@C=3.5em@R=2.8em{
K \ar[r] &
K\otimes_k L \ar[r] &
K\otimes_{K'}L(\mathbf t)
\\
K' \ar[u] \ar[r] &
K'\otimes_k L \ar[u] \ar[r] &
L(\mathbf t) \ar[u]
\\
k \ar[u] \ar[r] &
L \ar[u] &
}
\tag{4.2.1.3}
\label{IV.4.2.1.3}
\end{equation}

Since $K$ is faithfully flat over $K'$,
\[
K\otimes_k L
=K\otimes_{K'}(K'\otimes_k L)
\]
is faithfully flat over $K'\otimes_k L$, and therefore
$K'\otimes_k L$ is Noetherian
\sref[0\textsubscript{I}]{0.6.5.2}.
Moreover, $K'\otimes_k L$ identifies with a subring of
$K\otimes_k L$.
The contraction to $K'\otimes_k L$ of a prime ideal $\mathfrak p$
associated with $K\otimes_k L$ is the zero ideal, since no nonzero element
of $K'\otimes_k L$ is a zero-divisor in $K\otimes_k L$
\sref[0\textsubscript{I}]{0.6.3.4} and
Bourbaki, \emph{Alg.\ comm.}, chap.~IV, \S~1, no.~1,
cor.~3 to prop.~2.

Since $K$ is algebraic over $K'$, the ring $K\otimes_k L$ is integral
over $K'\otimes_k L$, and the prime ideals of $K\otimes_k L$ lying over
the zero ideal of $K'\otimes_k L$ are necessarily pairwise incomparable
(Bourbaki, \emph{Alg.\ comm.}, chap.~V, \S~2, no.~1,
cor.~1 to prop.~1).
This proves the first assertion
(Bourbaki, \emph{Alg.\ comm.}, chap.~IV, \S~1, no.~3,
cor.~1 to prop.~7).

Furthermore, the residue field $E$ of $\mathfrak p$ is algebraic over the
residue field of the zero ideal of $K'\otimes_k L$, namely
$L(\mathbf t)$.
Consequently
\[
\operatorname{tr.deg}_L E
=\operatorname{tr.deg}_L L(\mathbf t)
=\operatorname{Card}(I)
=\operatorname{tr.deg}_k K.
\]
This gives the second equality in \sref{IV.4.2.1.1}; interchanging the
roles of $K$ and $L$ gives the first, and
\sref{IV.4.2.1.2} follows.
\end{proof}

\begin{corollary}[4.2.2]
\phantomsection
\label{IV.4.2.2}
Under the hypotheses of \sref{IV.3.3.6}, if the preschemes
$T_{x,y}$ are locally Noetherian, they have no immersed associated
prime cycle.
\end{corollary}

\begin{corollary}[4.2.3]
\phantomsection
\label{IV.4.2.3}
Under the hypotheses of \sref{IV.3.3.6}
(respectively \sref{IV.3.3.7}), if the $T_{x,y}$ are locally Noetherian
and if $\sh F$ and the $\sh G_s$ $(s\in S)$
(respectively $\sh F$ and $\sh G$) have no immersed associated prime
cycle, then neither does
\[
\sh F\otimes_{\sh S}\sh G.
\]
\end{corollary}

This follows from \sref{IV.4.2.2}, \sref{IV.3.3.2}, and the proof of
\sref{IV.3.3.6}.
In particular, since every prescheme over a field $k$ is flat over $k$,
we have thereby proved assertion~{\hyperlink{ega4.c04.sec.4-2-4-i}{(i)}} of the
following proposition.

\begin{proposition}[4.2.4]
\phantomsection
\label{IV.4.2.4}
Let $k$ be a field, and let $X$ and $Y$ be two locally Noetherian
$k$-preschemes such that $X\times_k Y$ is locally Noetherian.
Suppose in addition that $X$ and $Y$ are integral.
Then:
\begin{enumerate}[(i)]
\item
\phantomsection
\label{IV.4.2.4.i}
$X\times_k Y$ has no immersed associated prime cycle.
Each irreducible component of $X\times_k Y$ dominates $X$ and $Y$, and
the set of these components is in bijective correspondence with the set
of irreducible components of
\[
\Spec\bigl(\mathrm R(X)\otimes_k\mathrm R(Y)\bigr),
\]
or equivalently with the set of minimal prime ideals of
$\mathrm R(X)\otimes_k\mathrm R(Y)$, where $\mathrm R(X)$ and
$\mathrm R(Y)$ are the fields of rational functions of $X$ and $Y$,
respectively.

\item
\phantomsection
\label{IV.4.2.4.ii}
If a maximal point $z$ of $X\times_k Y$ is identified with a minimal
prime ideal $\mathfrak p$ of
$\mathrm R(X)\otimes_k\mathrm R(Y)$, then the local ring
\[
\sh O_{X\times_k Y,z}
\]
is isomorphic to the ring of fractions
\[
\bigl(\mathrm R(X)\otimes_k\mathrm R(Y)\bigr)_{\mathfrak p}.
\]
In particular, if $\mathrm R(X)$ or $\mathrm R(Y)$ is separable over
$k$, then $X\times_k Y$ is reduced.

\oldpage[IV-2]{56}

\item
\phantomsection
\label{IV.4.2.4.iii}
If, in addition, $X$ and $Y$ are locally of finite type over $k$, then
every irreducible component of $X\times_k Y$ has dimension
\[
\dim(X)+\dim(Y).
\]
\end{enumerate}
\end{proposition}

\begin{proof}
Assertion~{\hyperlink{ega4.c04.sec.4-2-4-iii}{(iii)}} follows from
\sref{IV.4.2.1.2}, taking account of
{\hyperlink{ega4.c04.sec.4-2-4-i}{(i)}} and {\hyperlink{ega4.c04.sec.4-2-4-ii}{(ii)}}.

To prove {\hyperlink{ega4.c04.sec.4-2-4-ii}{(ii)}}, we may plainly reduce to the case
in which
\[
X=\Spec(A),
\qquad
Y=\Spec(B)
\]
are affine, with $A$ and $B$ integral rings having respective fields of
fractions
\[
K=\mathrm R(X),
\qquad
L=\mathrm R(Y).
\]
Assertion~{\hyperlink{ega4.c04.sec.4-2-4-i}{(i)}} shows that every minimal prime ideal
$\mathfrak q$ of $A\otimes_k B$ is the contraction to $A\otimes_k B$
of a minimal prime ideal $\mathfrak p$ of $K\otimes_k L$.
Since $K\otimes_k L$ is a ring of fractions of $A\otimes_k B$, the
isomorphism
\[
(A\otimes_k B)_{\mathfrak q}
\simeq
(K\otimes_k L)_{\mathfrak p}
\]
follows from \sref[0\textsubscript{I}]{0.1.2.5}.

Finally, if, for example, $\mathrm R(Y)$ is separable over $k$, then
$\mathrm R(X)\otimes_k\mathrm R(Y)$ is reduced
(Bourbaki, \emph{Alg.}, chap.~VIII, \S~7, no.~3, th.~1).
The local rings at the maximal points of $X\times_k Y$ are therefore
also reduced, and it follows that $X\times_k Y$ is reduced
\sref{IV.3.2.1}.
\end{proof}

\begin{proposition}[4.2.5]
\phantomsection
\label{IV.4.2.5}
Let $k$ be a field, let $X$ and $Y$ be locally Noetherian
$k$-preschemes, and let $\sh F$ (respectively $\sh G$) be a
quasi-coherent $\sh O_X$-Module (respectively $\sh O_Y$-Module).
Let $(Z'_\lambda)$ (respectively $(Z''_\mu)$) be the family of prime
cycles associated with $\sh F$ (respectively $\sh G$), and denote again
by $Z'_\lambda$ (respectively $Z''_\mu$) the reduced sub-prescheme of
$X$ (respectively $Y$) having $Z'_\lambda$ (respectively $Z''_\mu$) as
its underlying space.
If $Z'_\lambda\times_k Z''_\mu$ is locally Noetherian, then its
irreducible components $Z_{\lambda\mu\nu}$ dominate $Z'_\lambda$ and
$Z''_\mu$, and $(Z_{\lambda\mu\nu})$ is the family of distinct prime
cycles associated with $\sh F\otimes_k\sh G$.
\end{proposition}

\begin{proof}
It suffices to apply \sref{IV.4.2.4} to the product
\[
\Spec\bigl(\kappa(x)\bigr)\times_k\Spec\bigl(\kappa(y)\bigr).
\]
\end{proof}

\begin{corollary}[4.2.6]
\phantomsection
\label{IV.4.2.6}
Let $k$ be a field, and let $X$ and $Y$ be two locally Noetherian
$k$-preschemes such that $X\times_k Y$ is locally Noetherian.
Let $(Z'_\lambda)$ (respectively $(Z''_\mu)$) be the family of reduced
sub-preschemes of $X$ (respectively $Y$) whose underlying spaces are
the irreducible components of $X$ (respectively $Y$).
Then the irreducible components $Z_{\lambda\mu\nu}$ of
$Z'_\lambda\times_k Z''_\mu$ dominate $Z'_\lambda$ and $Z''_\mu$, and
$(Z_{\lambda\mu\nu})$ is the family of irreducible components of
$X\times_k Y$.
\end{corollary}

\begin{proof}
We may reduce to the case in which $X$ and $Y$ are reduced
\sref[I]{I.5.1.8}.
The irreducible components of $X$ (respectively $Y$) are then the prime
cycles associated with $\sh O_X$ (respectively $\sh O_Y$)
\sref{IV.3.2.1}.
Apply \sref{IV.4.2.5} with
\[
\sh F=\sh O_X,
\qquad
\sh G=\sh O_Y,
\]
noting that, by definition,
\[
\sh O_X\otimes_k\sh O_Y=\sh O_{X\times_k Y}.
\]
The corollary follows because $\sh O_X\otimes_k\sh O_Y$ has no immersed
associated prime cycles, since by hypothesis the same is true of
$\sh O_X$ and $\sh O_Y$ \sref{IV.4.2.3}.
\end{proof}

We apply the preceding results to the case in which $Y$ is the spectrum
of an extension $K$ of $k$.

\begin{proposition}[4.2.7]
\phantomsection
\label{IV.4.2.7}
Let $k$ be a field, let $X$ be a $k$-prescheme, and let $K$ be an
extension of $k$ such that $X\otimes_k K$ is locally Noetherian.
Let $\sh F$ be a quasi-coherent $\sh O_X$-Module, let $x'$ be a point
of $X\otimes_k K$, and let $x$ be its image in $X$.
\begin{enumerate}[(i)]
\item
\phantomsection
\label{IV.4.2.7.i}
Let $(Z_\lambda)$ be the family of reduced sub-preschemes of $X$ whose
underlying spaces are the prime cycles associated with $\sh F$.
Then the irreducible components $Z_{\lambda\mu}$ of
$Z_\lambda\otimes_k K$ are the prime cycles associated with
$\sh F\otimes_k K$, and $Z_{\lambda\mu}$ dominates $Z_\lambda$.
Moreover, $Z_{\lambda\mu}$ is immersed if and only if $Z_\lambda$ is
immersed.

\item
\phantomsection
\label{IV.4.2.7.ii}
The point $x$ belongs to an immersed prime cycle associated with
$\sh F$ if and only if $x'$ belongs to an immersed prime cycle
associated with $\sh F\otimes_k K$.
The Module $\sh F$ has no immersed associated prime cycle if and only if
the same is true of $\sh F\otimes_k K$.

\oldpage[IV-2]{57}

\item
\phantomsection
\label{IV.4.2.7.iii}
The indices $\lambda$ for which $x\in Z_\lambda$ are precisely those for
which there exists an index $\mu$ such that
$x'\in Z_{\lambda\mu}$.
In particular, if $x'$ belongs to only one prime cycle associated with
$\sh F\otimes_k K$, then $x$ belongs to only one prime cycle associated
with $\sh F$.

\item
\phantomsection
\label{IV.4.2.7.iv}
If $X$ is locally of finite type over $k$, then
\[
\dim(Z_{\lambda\mu})=\dim(Z_\lambda).
\]
\end{enumerate}
\end{proposition}

\begin{proof}
The hypothesis implies that $X$ itself is locally Noetherian
\sref{IV.2.2.13} and \sref{IV.2.2.14}.
Assertion~{\hyperlink{ega4.c04.sec.4-2-7-i}{(i)}} follows from
\sref{IV.4.2.5} and the proof of \sref{IV.4.2.3}, with
$\sh G=\sh O_Y$ and $Y=\Spec(K)$.
Assertions~{\hyperlink{ega4.c04.sec.4-2-7-ii}{(ii)}} and
{\hyperlink{ega4.c04.sec.4-2-7-iii}{(iii)}} follow from
{\hyperlink{ega4.c04.sec.4-2-7-i}{(i)}} and \sref{IV.2.3.5}.
Finally, {\hyperlink{ega4.c04.sec.4-2-7-iv}{(iv)}} is a special case of
{\hyperlink{ega4.c04.sec.4-2-4-iii}{\sref{IV.4.2.4}, (iii)}}.
\end{proof}

\begin{corollary}[4.2.8]
\phantomsection
\label{IV.4.2.8}
If $X$ is locally of finite type over $k$, then the set of dimensions
of the prime cycles associated with $\sh F$ is the same as that of the
prime cycles associated with $\sh F\otimes_k K$.
The set of dimensions of the irreducible components is likewise the
same for $X$ and $X\otimes_k K$.
\end{corollary}

\begin{proposition}[4.2.9]
\phantomsection
\label{IV.4.2.9}
Suppose that the hypotheses of \sref{IV.4.2.5} hold, and suppose in
addition that $\sh F$ and $\sh G$ are coherent.
Let $(\sh F_\lambda)_{\lambda\in L}$ and
$(\sh G_\mu)_{\mu\in M}$ be irredundant decompositions of $\sh F$ and
$\sh G$, respectively.
For every pair $(\lambda,\mu)\in L\times M$, let
\[
(\sh H_{\lambda\mu\nu})_{\nu\in S(\lambda,\mu)}
\]
be a reduced irredundant decomposition of
$\sh F_\lambda\otimes_k\sh G_\mu$, where
\[
S(\lambda,\mu)
=\operatorname{Ass}(\sh F_\lambda\otimes_k\sh G_\mu)
\]
\sref{IV.3.2.5}.
Then $(\sh H_{\lambda\mu\nu})$, over all triples
$(\lambda,\mu,\nu)$, is an irredundant decomposition of
$\sh F\otimes_k\sh G$.
It is reduced if $(\sh F_\lambda)$ and $(\sh G_\mu)$ are reduced.
\end{proposition}

\begin{proof}
The Module $\sh F\otimes_k\sh G$ and each
$\sh F_\lambda\otimes_k\sh G_\mu$ are coherent, and each
$\sh F_\lambda\otimes_k\sh G_\mu$ identifies with a quotient of
$\sh F\otimes_k\sh G$.
By definition, each $\sh H_{\lambda\mu\nu}$ identifies with a coherent
quotient of $\sh F_\lambda\otimes_k\sh G_\mu$, and hence also with a
coherent quotient of $\sh F\otimes_k\sh G$.

The family of supports of the $\sh H_{\lambda\mu\nu}$ is locally finite.
Indeed,
\[
\Supp(\sh F_\lambda\otimes_k\sh G_\mu)
\subset
\bigl|\Supp(\sh F_\lambda)\times_k\Supp(\sh G_\mu)\bigr|,
\]
where $\Supp(\sh F_\lambda)$ and $\Supp(\sh G_\mu)$ denote the
corresponding reduced closed sub-preschemes.
For fixed $\lambda$ and $\mu$, the family of supports of the
$\sh H_{\lambda\mu\nu}$ is locally finite by hypothesis; the assertion
therefore follows from the local finiteness of the families
$(\Supp(\sh F_\lambda))$ and $(\Supp(\sh G_\mu))$.

It remains to show that the canonical homomorphism
\[
\sh F\otimes_k\sh G
\longrightarrow
\bigoplus_{\lambda,\mu,\nu}\sh H_{\lambda\mu\nu}
\]
is injective.
It is the composite
\[
\begin{aligned}
\sh F\otimes_k\sh G
&\longrightarrow
\left(\bigoplus_\lambda\sh F_\lambda\right)
\otimes_k
\left(\bigoplus_\mu\sh G_\mu\right)
\\
&\simeq
\bigoplus_{\lambda,\mu}
(\sh F_\lambda\otimes_k\sh G_\mu)
\longrightarrow
\bigoplus_{\lambda,\mu,\nu}\sh H_{\lambda\mu\nu}.
\end{aligned}
\]
The last homomorphism is injective by definition, and the first is
injective because it is the tensor product of the canonical injections
\[
\sh F\longrightarrow\bigoplus_\lambda\sh F_\lambda,
\qquad
\sh G\longrightarrow\bigoplus_\mu\sh G_\mu,
\]
and $\sh F$ and $\sh G$ are flat over $k$.

Finally, if $(\sh F_\lambda)$ and $(\sh G_\mu)$ are reduced, we may
suppose that
\[
L=\operatorname{Ass}(\sh F),
\qquad
M=\operatorname{Ass}(\sh G).
\]
The $\sh H_{\lambda\mu\nu}$ then form a reduced decomposition because,
by \sref{IV.4.2.2}, the sets
\[
\operatorname{Ass}(\sh F_\lambda\otimes_k\sh G_\mu)
\]
partition
$\operatorname{Ass}(\sh F\otimes_k\sh G)$
\sref{IV.3.2.5}.
\end{proof}

\begin{corollary}[4.2.10]
\phantomsection
\label{IV.4.2.10}
Under the hypotheses of \sref{IV.4.2.7}, suppose in addition that
$\sh F$ is coherent, and let
$(\sh F_\lambda)_{\lambda\in L}$ be an irredundant decomposition of
$\sh F$.
For every $\lambda\in L$, let
\[
(\sh F_{\lambda\mu})_{
  \mu\in\operatorname{Ass}(\sh F_\lambda\otimes_k K)}
\]
be a reduced irredundant decomposition of
$\sh F_\lambda\otimes_k K$.
Then $(\sh F_{\lambda\mu})$ is an irredundant decomposition of
$\sh F\otimes_k K$.
It is reduced if $(\sh F_\lambda)$ is reduced.
\end{corollary}

\begin{proof}
Apply \sref{IV.4.2.9} with
\[
Y=\Spec(K),
\qquad
\sh G=\sh O_Y=K,
\]
and with $M$ consisting of a single element.
\end{proof}
% END NATIVE INPUT components/02_section42_associated_prime_cycles
% BEGIN NATIVE INPUT components/03_section43_tensor_products_of_fields
\oldpage[IV-2]{58}

\subsection{Reminders on tensor products of fields}
\label{subsection:IV.4.3}

For the reader's convenience, we recall here certain properties of
tensor products of fields that we shall use in the following numbers;
for the proofs, we refer to [1].

\begin{env}[4.3.1]
\phantomsection
\label{IV.4.3.1}
Recall that an extension $L$ of a field $k$ is called
\emph{primary} if the largest separable algebraic extension of $k$
contained in $L$ is $k$ itself.
\end{env}

\begin{proposition}[4.3.2]
\phantomsection
\label{IV.4.3.2}
Let $K$ and $L$ be two extensions of a field $k$.
If $L$ is a primary extension of $k$, then
$\Spec(L\otimes_k K)$ is irreducible and, if $\zeta$ is its generic
point, $\kappa(\zeta)$ is a primary extension of $K$.
Conversely, if $\Spec(L\otimes_k K)$ is irreducible for every finite
separable extension $K$ of $k$, then $L$ is a primary extension of $k$.
\end{proposition}

\begin{proof}
See [1], pp.~14-03 to 14-06.
\end{proof}

\begin{corollary}[4.3.3]
\phantomsection
\label{IV.4.3.3}
If $k$ is separably closed (that is, if its algebraic closure is
radicial over $k$), then $\Spec(L\otimes_k K)$ is irreducible for every
two extensions $K$ and $L$ of $k$; conversely, if this conclusion holds
for every such pair, then $k$ is separably closed.
\end{corollary}

\begin{proof}
This follows immediately from \sref{IV.4.3.2}.
\end{proof}

\begin{corollary}[4.3.4]
\phantomsection
\label{IV.4.3.4}
Let $L$ be an extension of a field $k$, and let $L_s$ be the largest
separable algebraic extension of $k$ contained in $L$ (sometimes called
the separable algebraic closure of $k$ in $L$).
Suppose that $L_s$ is of finite degree over $k$, and let $K$ be a
Galois extension of $k$ containing $L_s$.
Then $\Spec(L\otimes_k K)$ has $[L_s:k]$ irreducible components and is
their union, and the residue fields at their generic points are primary
extensions of $K$.
\end{corollary}

\begin{proof}
Indeed,
\[
L\otimes_k K
=L\otimes_{L_s}(L_s\otimes_k K),
\]
and $L_s\otimes_k K$ is isomorphic to a product of $[L_s:k]$ fields,
each isomorphic to $K$
(Bourbaki, \emph{Alg.}, chap.~VIII, \S~7, no.~3,
cor.~2 to th.~1).
Thus $L\otimes_k K$ is isomorphic to a product of $[L_s:k]$ rings,
each isomorphic to $L\otimes_{L_s}K$.
Since $L$ is a primary extension of $L_s$, the conclusion follows from
\sref{IV.4.3.2}.
\end{proof}

\begin{proposition}[4.3.5]
\phantomsection
\label{IV.4.3.5}
Let $K$ and $L$ be two extensions of a field $k$.
If $L$ is a separable extension of $k$, then the ring
$L\otimes_k K$ is reduced.
Conversely, if $L\otimes_k K$ is reduced for every finite radicial
extension $K$ of $k$, then $L$ is a separable extension of $k$.
\end{proposition}

\begin{proof}
See Bourbaki, \emph{Alg.}, chap.~VIII, \S~7, no.~3, th.~1.
\end{proof}

\begin{corollary}[4.3.6]
\phantomsection
\label{IV.4.3.6}
If $k$ is perfect, then $K\otimes_k L$ is reduced for every two
extensions $K$ and $L$ of $k$; conversely, if this conclusion holds for
every such pair, then $k$ is perfect.
\end{corollary}

\begin{proof}
This follows immediately from \sref{IV.4.3.5}.
\end{proof}

\begin{corollary}[4.3.7]
\phantomsection
\label{IV.4.3.7}
Let $L$ be a separable extension of $k$, and let $K$ be an extension of
$k$ such that either $K$ or $L$ is finite over $k$.
Then the residue fields of the semilocal ring $L\otimes_k K$ are
separable extensions of $K$.
\end{corollary}

\begin{proof}
The ring $L\otimes_k K$ is the direct product of its residue fields
$E_i$.
For every extension $K'$ of $K$, the ring $L\otimes_k K'$ is therefore
isomorphic to the direct product of the rings $E_i\otimes_K K'$.
Since $L\otimes_k K'$ is reduced, the fields $E_i$ are separable over
$K$ by \sref{IV.4.3.5}.
\end{proof}

\begin{corollary}[4.3.8]
\phantomsection
\label{IV.4.3.8}
If $K$ and $L$ are two finite separable extensions of $k$, then
$K\otimes_k L$ is a product of finitely many separable extensions of
$k$.
\end{corollary}

\oldpage[IV-2]{59}

\begin{proposition}[4.3.9]
\phantomsection
\label{IV.4.3.9}
If $k$ is algebraically closed, then the ring $L\otimes_k K$ is an
integral domain for every two extensions $K$ and $L$ of $k$; conversely,
if this conclusion holds for every such pair, then $k$ is algebraically
closed.
\end{proposition}

\begin{proof}
This follows from \sref{IV.4.3.3} and \sref{IV.4.3.6}, since a field
that is both perfect and separably closed is algebraically closed.
\end{proof}
% END NATIVE INPUT components/03_section43_tensor_products_of_fields
% BEGIN NATIVE INPUT components/04_section44_irreducible_connected_over_algebraically_closed
\subsection{Irreducible preschemes and connected preschemes over an algebraically closed field}
\label{subsection:IV.4.4}

\begin{env}[4.4.1]
\phantomsection
\label{IV.4.4.1}
Let $k$ be a field, let $X$ be a $k$-prescheme, and let $K$ be an
extension of $k$.
Since the morphism
\[
\Spec(K)\longrightarrow\Spec(k)
\]
is faithfully flat, quasi-compact, and universally open
\sref{IV.2.4.9}, the same is true of the projection
\[
p:X\otimes_k K\longrightarrow X
\]
\sref{IV.2.2.13}.
It therefore follows from \sref{IV.2.3.5} that every irreducible
component of $X\otimes_k K$ dominates an irreducible component of $X$,
and, since $p$ is surjective, every irreducible component of $X$ is
dominated by an irreducible component of $X\otimes_k K$.
Thus $p$ induces a surjective map from the set of irreducible components
of $X\otimes_k K$ to the set of irreducible components of $X$.

Likewise, because $p$ is continuous and surjective, every connected
component of $X$ is the image under $p$ of a union of connected
components of $X\otimes_k K$.
Hence there is a surjective map from the set of connected components of
$X\otimes_k K$ to the set of connected components of $X$.
It follows that every irreducible (resp.\ connected) component $Z'$ of
$X\otimes_k K$ is contained in a unique set of the form $p^{-1}(Z)$,
where $Z$ is an irreducible (resp.\ connected) component of $X$.
For every sub-prescheme of $X$ having $Z$ as its underlying space (and
again denoted by $Z$), the space $Z'$ is an irreducible
(resp.\ connected) component of $Z\otimes_k K$.

In particular, suppose that $X\otimes_k K$ has finitely many
irreducible (resp.\ connected) components, say $n$ (resp.\ $n'$).
This is the case when $X$ is of finite type over $k$, since then
$X\otimes_k K$ is of finite type over $K$ and hence Noetherian.
The number of irreducible (resp.\ connected) components of $X$ is then
at most $n$ (resp.\ at most $n'$).
Equality with $n$ (resp.\ $n'$) means that, for every reduced
sub-prescheme $Z$ of $X$ whose underlying space is an irreducible
(resp.\ connected) component of $X$, the prescheme $Z\otimes_k K$ is
irreducible (resp.\ connected) \sref[I]{I.5.1.8}.
Consequently, the number of irreducible (resp.\ connected) components of
$X\otimes_k K$ is independent of $K$ if and only if, for every
irreducible (resp.\ connected) component $Z$ of $X$, the prescheme
$Z\otimes_k K$ is irreducible (resp.\ connected) for every extension
$K$ of $k$.

In particular, if $X\otimes_k K$ is irreducible (resp.\ connected), then
so is $X$, as also follows directly from the surjectivity of $p$.
Similarly, if $X\otimes_k K$ is reduced (resp.\ integral), then so is
$X$, since $p$ is faithfully flat \sref{IV.2.1.13}.

In this number and the next, we shall examine more closely what can be
said about the irreducible (resp.\ connected) components of
$X\otimes_k K$ as $K$ varies.
The preceding observations already show that the number of these
components is an increasing function of $K$.
\end{env}

\begin{lemma}[4.4.2]
\phantomsection
\label{IV.4.4.2}
Let $X'$ and $X$ be two topological spaces, and let
$f:X'\to X$ be a continuous map.
Suppose that the following conditions hold:
\begin{enumerate}
\item[\rm(i)]
\phantomsection
\label{IV.4.4.2.i}
$f$ is open and surjective (resp.\ the topological space $X$ is
canonically identified with the quotient of $X'$ by the equivalence
relation defined by $f$);

\oldpage[IV-2]{60}

\item[\rm(ii)]
\phantomsection
\label{IV.4.4.2.ii}
for every $x\in X$, the fibre $f^{-1}(x)$ is irreducible
(resp.\ connected).
\end{enumerate}
Then $X'$ is irreducible (resp.\ connected) if and only if $X$ is.
\end{lemma}

\begin{proof}
The assertion concerning connectedness is Bourbaki,
\emph{Top.\ g\'en.}, chap.~I, 3rd ed., \S~11, no.~3, prop.~7.
We prove the assertion concerning irreducibility.
The condition is plainly necessary because $f$ is surjective; let us
prove that it is sufficient.

Let $X'_1$ and $X'_2$ be two closed subsets of $X'$ such that
$X'=X'_1\cup X'_2$, and, for $i=1,2$, let $X_i$ be the set of
$x\in X$ such that $f^{-1}(x)\subset X'_i$.
Then
\[
X-X_i=f(X'-X'_i),
\]
and, since $f$ is open, $X_i$ is a closed subset of $X$.
For every $x\in X$, the fibre $f^{-1}(x)$ is irreducible by hypothesis
and is the union of the two closed subsets
\[
X'_1\cap f^{-1}(x)
\quad\text{and}\quad
X'_2\cap f^{-1}(x).
\]
One of them must equal $f^{-1}(x)$, so $x\in X_1$ or $x\in X_2$.
Thus $X=X_1\cup X_2$.
Since $X$ is irreducible, $X=X_1$ or $X=X_2$, and therefore
$X'=X'_1$ or $X'=X'_2$.
\end{proof}

\begin{remark}[4.4.3]
\phantomsection
\label{IV.4.4.3}
If the topology of $X$ is not the quotient of that of $X'$ by the
equivalence relation defined by $f$, it can happen that $X$ and all the
fibres $f^{-1}(x)$ are irreducible without $X'$ being connected.
For example, take for $X$ an irreducible scheme---the affine line
$\Spec(k[T])$ over an algebraically closed field $k$, say---choose a
closed point $x_0$ of $X$, and take for $X'$ the topological sum of
$\{x_0\}$ and the open subspace $X-\{x_0\}$ of $X$.

Likewise, if the hypothesis that $f$ is open is replaced by the
hypothesis that $f$ is closed, or even proper, then it can happen that
$X$ and all the fibres $f^{-1}(x)$ are irreducible without $X'$ being
irreducible, even though this replacement does ensure that the topology
of $X$ is the quotient of that of $X'$ by the relation defined by $f$.
Again take $X$ to be the affine line over $k$, and put
\[
S=X\times_k\mathbf P,
\]
where $\mathbf P$ is the projective line over $k$.
Let $x_0$ be a closed point of $X$, let $t_0$ be a closed point of
$\mathbf P$, and let
\[
p:S\longrightarrow X,
\qquad
q:S\longrightarrow\mathbf P
\]
be the projections.
Let $X'$ be the reduced sub-prescheme whose underlying space is the
closed set
\[
p^{-1}(x_0)\cup q^{-1}(t_0),
\]
and let $f$ be the restriction of $p$ to $X'$.
Then $f$ is proper, but $X'$ is not irreducible.
\end{remark}

\begin{theorem}[4.4.4]
\phantomsection
\label{IV.4.4.4}
Let $k$ be an algebraically closed field and let $X$ be a
$k$-prescheme.
If $X$ is irreducible (resp.\ connected), then $X\otimes_k K$ is
irreducible (resp.\ connected) for every extension $K$ of $k$.
\end{theorem}

\begin{proof}
We saw in \sref{IV.4.4.1} that
\[
p:X\otimes_k K\longrightarrow X
\]
is faithfully flat and open.
To apply \sref{IV.4.4.2}, it therefore suffices to verify that every
fibre $p^{-1}(x)$ is irreducible.
The $\kappa(x)$-prescheme $p^{-1}(x)$ is isomorphic to
\[
\Spec\bigl(\kappa(x)\otimes_k K\bigr)
\]
\sref[I]{I.3.6.2}, and is therefore integral by the hypothesis on $k$
and \sref{IV.4.3.9}.
\end{proof}

\begin{corollary}[4.4.5]
\phantomsection
\label{IV.4.4.5}
Let $k$ be an algebraically closed field, let $X$ be a $k$-prescheme,
let $K$ be an extension of $k$, and let
\[
p:X\otimes_k K\longrightarrow X
\]
be the canonical projection.
If $Z$ is an irreducible (resp.\ connected) subset of $X$, then
$p^{-1}(Z)$ is irreducible (resp.\ connected).
In particular, if $X_0$ is an irreducible component
(resp.\ the connected component) of $X$ containing $Z$, then
$p^{-1}(X_0)$ is an irreducible (resp.\ connected) component of
$X\otimes_k K$ containing $p^{-1}(Z)$.
\end{corollary}

\begin{proof}
The second assertion follows from \sref{IV.4.4.1} and the first assertion
applied after replacing $Z$ by $X_0$.
To prove the first, let $X'$ be the underlying space of
$X\otimes_k K$, and put $Z'=p^{-1}(Z)$.

\oldpage[IV-2]{61}

The equivalence relation $R$ defined by $p$ on $X'$ is open.
The induced relation $R_{Z'}$ on the saturated subset $Z'$ is therefore
also open, and $Z$ identifies with the quotient space $Z'/R_{Z'}$
(Bourbaki, \emph{Top.\ g\'en.}, chap.~I, 3rd ed., \S~5, no.~2,
prop.~4).
We may thus apply \sref{IV.4.4.2} to $Z$ and $Z'$, proving the first
assertion of \sref{IV.4.4.5}.
\end{proof}

\begin{corollary}[4.4.6]
\phantomsection
\label{IV.4.4.6}
Under the hypotheses and notation of \sref{IV.4.4.5}, the map
\[
Z\longmapsto p^{-1}(Z)
\]
is a bijection from the set of irreducible (resp.\ connected)
components of $X$ to the set of irreducible (resp.\ connected)
components of $X\otimes_k K$.
Its inverse is
\[
Z'\longmapsto p(Z').
\]
\end{corollary}
% END NATIVE INPUT components/04_section44_irreducible_connected_over_algebraically_closed
% BEGIN NATIVE INPUT components/05_section45_geometric_irreducibility_connectedness
\subsection{Geometrically irreducible and geometrically connected preschemes}
\label{subsection:IV.4.5}
\label{IV.4.5}

\begin{proposition}[4.5.1]
\phantomsection
\label{IV.4.5.1}
Let $k$ be a field, let $X$ be a $k$-prescheme, and let $\Omega$ be an
algebraically closed extension of $k$.
Let $n$ (resp.\ $n'$) be the cardinality of the set of irreducible
components (resp.\ connected components) of $X\otimes_k\Omega$.
This cardinality is independent of the chosen algebraically closed
extension $\Omega$.
Moreover, for every extension $K$ of $k$, the cardinality of the set of
irreducible components (resp.\ connected components) of $X\otimes_k K$
is at most $n$ (resp.\ at most $n'$).
\end{proposition}

\begin{proof}
Two algebraically closed extensions $\Omega$ and $\Omega'$ of $k$ may
always be regarded as subextensions of a common extension $E$ of $k$.
Applying \sref{IV.4.4.6} to $X\otimes_k\Omega$ and
$X\otimes_k\Omega'$, and to the common extensions $E/\Omega$ and
$E/\Omega'$, proves the first assertion.
For the second, take $\Omega$ to be an algebraically closed extension of
$K$ and use \sref{IV.4.4.1}.
\end{proof}

\begin{definition}[4.5.2]
\phantomsection
\label{IV.4.5.2}
The cardinality $n$ (resp.\ $n'$) of \sref{IV.4.5.1} is called the
\emph{geometric number of irreducible components}
(resp.\ \emph{of connected components}) of $X$ relative to $k$.
If $n=1$ (resp.\ $n'\leq 1$), the $k$-prescheme $X$ is said to be
\emph{geometrically irreducible} (resp.\ \emph{geometrically
connected}).
\end{definition}

This recovers the definition given in the special case
\sref[III]{III.4.3.4}.
By \sref{IV.4.5.1}, the following properties are equivalent:
\begin{enumerate}
\item[\rm(a)]
\phantomsection
\label{IV.4.5.2.a}
$X$ is geometrically irreducible (resp.\ geometrically connected);

\item[\rm(b)]
\phantomsection
\label{IV.4.5.2.b}
for one algebraically closed extension $\Omega$ of $k$,
$X\otimes_k\Omega$ is irreducible (resp.\ connected);

\item[\rm(c)]
\phantomsection
\label{IV.4.5.2.c}
for every extension $K$ of $k$, $X\otimes_k K$ is irreducible
(resp.\ connected).
\end{enumerate}

\begin{env}[4.5.3]
\phantomsection
\label{IV.4.5.3}
Let $X$ be a $k$-prescheme and let $Z$ be a locally closed subset of
$X$.
If any sub-prescheme of $X$ having $Z$ as underlying space is again
denoted by $Z$ \sref[I]{I.5.2.1}, it follows from \sref[I]{I.5.1.8}
that, for an extension $K$ of $k$, the number of irreducible
(resp.\ connected) components of $Z\otimes_k K$ is independent of the
chosen sub-prescheme with underlying space $Z$.
The \emph{geometric number of irreducible components}
(resp.\ \emph{of connected components}) of $Z$ is therefore defined to
be that of any sub-prescheme of $X$ having $Z$ as underlying space.
\end{env}

\begin{proposition}[4.5.4]
\phantomsection
\label{IV.4.5.4}
Let $X$ and $Y$ be two $k$-preschemes, and let $f:X\to Y$ be a
surjective $k$-morphism.
If $X$ is geometrically irreducible (resp.\ geometrically connected),
then so is $Y$.
\end{proposition}

\oldpage[IV-2]{62}

\begin{proof}
For every extension $K$ of $k$, the morphism
\[
f_K:X\otimes_k K\longrightarrow Y\otimes_k K
\]
is surjective \sref[I]{I.3.5.2}.
Since $X\otimes_k K$ is irreducible (resp.\ connected) by hypothesis,
the same is true of $Y\otimes_k K$.
\end{proof}

\begin{definition}[4.5.5]
\phantomsection
\label{IV.4.5.5}
A morphism $f:X\to Y$ of preschemes is said to be
\emph{irreducible} (resp.\ \emph{connected}) if, for every $y\in Y$,
the $\kappa(y)$-prescheme $f^{-1}(y)$ is geometrically irreducible
(resp.\ geometrically connected).
\end{definition}

\begin{proposition}[4.5.6]
\phantomsection
\label{IV.4.5.6}
\begin{enumerate}
\item[\rm(i)]
\phantomsection
\label{IV.4.5.6.i}
Let $X$ be a $k$-prescheme and let $K$ be an extension of $k$.
The geometric number of irreducible components (resp.\ connected
components) of the $K$-prescheme $X\otimes_k K$, relative to $K$, is
equal to the geometric number of irreducible components
(resp.\ connected components) of $X$, relative to $k$.
In particular, $X\otimes_k K$ is geometrically irreducible
(resp.\ geometrically connected) as a $K$-prescheme if and only if
$X$ is so as a $k$-prescheme.

\item[\rm(ii)]
\phantomsection
\label{IV.4.5.6.ii}
Let $f:X\to Y$ and $g:Y'\to Y$ be two morphisms, and put
\[
X'=X_{(Y')}=X\times_Y Y',
\qquad
f'=f_{(Y')}:X'\longrightarrow Y'.
\]
If $f$ is irreducible (resp.\ connected), then so is $f'$.
The converse holds if $g$ is surjective.
\end{enumerate}
\end{proposition}

\begin{proof}
For (i), if $\Omega$ is an algebraically closed extension of $K$, then
\[
X\otimes_k\Omega
  =(X\otimes_k K)\otimes_K\Omega,
\]
which proves the assertion.

For (ii), for every $y'\in Y'$, put $y=g(y')$.
Then
\[
f'^{-1}(y')
  =f^{-1}(y)\otimes_{\kappa(y)}\kappa(y')
\]
\sref[I]{I.3.6.4}, so both assertions follow from (i).
\end{proof}

\begin{proposition}[4.5.7]
\phantomsection
\label{IV.4.5.7}
Let $f:X\to Y$ be a surjective irreducible (resp.\ connected)
morphism, let $g:Y'\to Y$ be a morphism, and put
$X'=X\times_Y Y'$.
Suppose in addition that $Y'$ is irreducible (resp.\ connected), and
that $f$ is universally open (resp.\ flat and quasi-compact, or
universally open, or universally closed).
Then $X'$ is irreducible (resp.\ connected).
\end{proposition}

\begin{proof}
All the properties of $f$ under consideration are stable under base
change, so it is enough to treat the case $Y'=Y$.
The result then follows from \sref{IV.4.4.2}, taking
\sref{IV.2.3.12} into account.
\end{proof}

\begin{corollary}[4.5.8]
\phantomsection
\label{IV.4.5.8}
Let $X$ and $Y$ be two $k$-preschemes.
\begin{enumerate}
\item[\rm(i)]
\phantomsection
\label{IV.4.5.8.i}
If $X$ is geometrically irreducible (resp.\ geometrically connected)
and $Y$ is irreducible (resp.\ connected), then $X\times_k Y$ is
irreducible (resp.\ connected).

\item[\rm(ii)]
\phantomsection
\label{IV.4.5.8.ii}
If $X$ and $Y$ are geometrically irreducible
(resp.\ geometrically connected), then so is $X\times_k Y$.
\end{enumerate}
\end{corollary}

\begin{proof}
For (i), the structure morphism
$X\to\Spec(k)$ is surjective and universally open
\sref{IV.2.4.9}, and it is also irreducible (resp.\ connected).
It is therefore enough to apply \sref{IV.4.5.7}.

For (ii), if $\Omega$ is an algebraically closed extension of $k$, then
\[
(X\times_k Y)\otimes_k\Omega
  =(X\otimes_k\Omega)\times_\Omega(Y\otimes_k\Omega)
\]
\sref[I]{I.3.3.10}.
By hypothesis and \sref{IV.4.5.6}, the two factors on the right are
geometrically irreducible (resp.\ geometrically connected)
$\Omega$-preschemes, so (i) applies.
\end{proof}

\begin{proposition}[4.5.9]
\phantomsection
\label{IV.4.5.9}
Let $X$ be a $k$-prescheme.  The following conditions are equivalent:
\begin{enumerate}
\item[\rm(a)]
\phantomsection
\label{IV.4.5.9.a}
$X$ is geometrically irreducible; in other words,
$X\otimes_k K$ is irreducible for every extension $K$ of $k$;

\item[\rm(b)]
\phantomsection
\label{IV.4.5.9.b}
$X\otimes_k K$ is irreducible for every finite separable extension
$K$ of $k$;

\item[\rm(c)]
\phantomsection
\label{IV.4.5.9.c}
$X$ is irreducible and, if $x$ is its generic point,
$\kappa(x)$ is a primary extension of $k$.
\end{enumerate}
\end{proposition}

\begin{proof}
It is clear that (a) implies (b).
To see that (b) implies (c), let $K$ be a finite separable extension
of $k$ and let
\[
p:X\otimes_k K\longrightarrow X
\]
be the canonical projection.
The maximal points of $X\otimes_k K$ are those of the fibre
\[
p^{-1}(x)=\Spec\bigl(\kappa(x)\otimes_k K\bigr)
\]
\sref{IV.2.3.4} and $(\mathbf{0}_{\mathrm I},2.1.8)$.
Thus, saying that this fibre is irreducible for every finite separable
extension $K$ of $k$ is equivalent, by \sref{IV.4.3.2}, to saying that
$\kappa(x)$ is a primary extension of $k$.
Conversely, if (c) holds, the same argument, again using
\sref{IV.4.3.2}, shows that $X\otimes_k K$ is irreducible for every
extension $K$ of $k$.
Hence (c) implies (a).
\end{proof}

\oldpage[IV-2]{63}

\begin{corollary}[4.5.10]
\phantomsection
\label{IV.4.5.10}
Let $X$ be an irreducible $k$-prescheme, let $x$ be its generic point,
let $k'$ be the separable algebraic closure of $k$ in $\kappa(x)$,
and let $k''$ be a Galois extension of $k$, of finite degree or not,
containing $k'$.
Suppose that $k'$ is finite over $k$.
Then the irreducible components of $X\otimes_k k''$ are geometrically
irreducible; their number is $[k':k]$, and this is also the geometric
number of irreducible components of $X$.
\end{corollary}

\begin{proof}
As in \sref{IV.4.5.9}, it is enough to consider the maximal points of
\[
\Spec\bigl(\kappa(x)\otimes_k k''\bigr).
\]
By \sref{IV.4.3.4}, there are $[k':k]$ of them, and their residue
fields, which are also the residue fields at the maximal points of
$X\otimes_k k''$, are primary extensions of $k''$.
The result follows from \sref{IV.4.5.9}.
\end{proof}

\begin{corollary}[4.5.11]
\phantomsection
\label{IV.4.5.11}
Let $X$ be a $k$-prescheme.
Suppose that $X$ has only finitely many maximal points
$x_i$ $(1\leq i\leq r)$ and that, for each $i$, the separable
algebraic closure $k'_i$ of $k$ in $\kappa(x_i)$ has finite degree
over $k$.
Then there is a finite separable extension $L$ of $k$ such that the
irreducible components of $X\otimes_k L$ are geometrically
irreducible.
Their number, equal to
\[
\sum_i[k'_i:k],
\]
is the geometric number of irreducible components of $X$.
\end{corollary}

\begin{proof}
The maximal points of $X\otimes_k L$ are those of the various fibres
\[
\Spec\bigl(\kappa(x_i)\otimes_k L\bigr),
\]
by \sref{IV.2.3.4} and $(\mathbf{0}_{\mathrm I},2.1.8)$, applied to
the irreducible components of $X$ and their inverse images in
$X\otimes_k L$.
It therefore suffices to take for $L$ a finite Galois extension of
$k$ containing all the $k'_i$ and to apply \sref{IV.4.5.10}.
\end{proof}

Notice that \sref{IV.4.5.11} applies in particular to every
$k$-prescheme $X$ of finite type over $k$.
Indeed, $X$ is then Noetherian, so it has only finitely many
irreducible components; moreover, each $\kappa(x_i)$ is an extension
of finite type of $k$, and therefore so is the algebraic closure of
$k$ in $\kappa(x_i)$ ([1], p.~6-06, Lemma~5).

\begin{remarks}[4.5.12]
\phantomsection
\label{IV.4.5.12}
\begin{enumerate}
\item[\rm(i)]
\phantomsection
\label{IV.4.5.12.i}
The notions defined in \sref{IV.4.5.2} depend on the base field $k$.
For brevity, when it is necessary to name the base field, we may say
``$k$-irreducible'' (resp.\ ``$k$-connected'') instead of
``geometrically irreducible (resp.\ connected) relative to $k$''.
This terminology is natural in the context of schemes, but is exactly
opposite to Weil's.
For Weil, ``$k$-irreducible'' (resp.\ ``$k$-connected'',
``$k$-normal'', and so on) designates an intrinsic property of a
prescheme $X$, independent of the field $k$ chosen as base field
(that is, independent of the chosen morphism
$X\to\Spec(k)$); we express these properties simply by saying that
$X$ is irreducible (resp.\ connected, normal, and so on).
On the other hand, Weil says ``absolutely irreducible'',
``absolutely connected'', ``absolutely normal'', and so on, for the
corresponding notions relative to $X\otimes_k\Omega$, where $\Omega$
is a suitable algebraically closed extension of $k$.
This terminology is explained by his different point of view:
an ``algebraic variety'' $V$ is first given as a geometric object over

\oldpage[IV-2]{64}

an algebraically closed field $\Omega$; the datum of a smaller
``field of definition'' $k$ (that is, a scheme over $k$ which yields
$V$ after extension of the base field to $\Omega$) is an additional
and, to some extent, secondary structure.
Thus his use of ``absolute'' or ``intrinsic'' on the one hand and
``relative to $k$'' on the other is opposite to ours.\footnote{Zariski's
point of view is already closer to ours, and his terminology is not
generally in conflict with that introduced here.}
We shall therefore avoid the qualifier ``absolute'', which may cause
confusion; whenever we use the abbreviated terminology introduced
above, in conflict with established terminology, we shall refer to
\sref{IV.4.5.12} to avoid ambiguity.

\item[\rm(ii)]
\phantomsection
\label{IV.4.5.12.ii}
Proposition \sref{IV.4.5.9} gives a ``birational'' criterion---one
depending only on the residue field at the generic point---for a
$k$-prescheme $X$ to be geometrically irreducible.
There is no analogous criterion for geometric connectedness.
For example, take
\[
X=\Spec\bigl(\mathbf{R}[[T,U]]/(T^2+U^2)\bigr),
\]
where $T$ and $U$ are indeterminates.
Then $X$ is an integral $\mathbf{R}$-scheme, and its field of rational
functions $\kappa(x)$ is $\mathbf{C}((T))$, as is readily verified.
The scheme $X\otimes_{\mathbf R}\mathbf C$ is the union of two
irreducible components, the two ``isotropic lines'', which have a
point in common: the maximal ideal in
$\mathbf C[[T,U]]/(T^2+U^2)$ induced by the maximal ideal
$(T)+(U)$ of $\mathbf C[[T,U]]$.
Thus $X\otimes_{\mathbf R}\mathbf C$ is connected.
Let $X'$ be the open complement of the closed point of $X$
corresponding to the maximal ideal $(T)+(U)$ of
$\mathbf R[[T,U]]$.
Then $X'\otimes_{\mathbf R}\mathbf C$ is not connected, although $X$
and $X'$ have the same field of rational functions.
\end{enumerate}
\end{remarks}

\begin{proposition}[4.5.13]
\phantomsection
\label{IV.4.5.13}
Let $X$ and $Y$ be two $k$-preschemes and let
$f:Y\to X$ be a $k$-morphism.
Suppose that $Y$ is nonempty and geometrically connected, and that
$X$ is connected.
Then $X$ is geometrically connected.
\end{proposition}

\begin{proof}
Let $\Omega$ be an algebraic closure of $k$, put
\[
X'=X\otimes_k\Omega,\qquad
Y'=Y\otimes_k\Omega,
\]
and let $f'=f_{(\Omega)}:Y'\to X'$.
It is enough to prove that $X'$ is connected.
Let $U'$ be a nonempty open-and-closed subset of $X'$.
The projections
\[
p:X'\longrightarrow X,\qquad q:Y'\longrightarrow Y
\]
are open \sref{IV.2.4.10}, and they are closed because $\Omega/k$ is
algebraic \sref[II]{II.6.1.10}.
Consequently $U=p(U')$ is a nonempty open-and-closed subset of $X$,
and hence $U=X$.
Since $Y$ is nonempty, it follows from \sref[I]{I.3.4.7} that
\[
V'=f'^{-1}(U')
\]
is nonempty.
It is also open and closed in $Y'$, which is connected by hypothesis,
so $V'=Y'$.
We conclude that $U'=X'$: otherwise, applying the same argument to
the open-and-closed subset $X'\setminus U'$ would give
$f'^{-1}(X'\setminus U')=Y'$, which is impossible because $Y'$ is
nonempty.
\end{proof}

\begin{corollary}[4.5.13.1]
\phantomsection
\label{IV.4.5.13.1}
\begin{enumerate}
\item[\rm(i)]
\phantomsection
\label{IV.4.5.13.1.i}
Let $X$ and $Y$ be two $k$-preschemes and let $f:Y\to X$ be a
$k$-morphism.
If $Y$ is nonempty and geometrically connected, then the connected
component $X_0$ of $X$ containing $f(Y)$ is geometrically connected.

\item[\rm(ii)]
\phantomsection
\label{IV.4.5.13.1.ii}
Let $X$ be a $k$-prescheme.
If $Y$ is a geometrically irreducible irreducible component of $X$,
then the connected component $X_0$ of $X$ containing $Y$ is
geometrically connected.
\end{enumerate}
\end{corollary}

\begin{proof}
For (i), we may suppose that $Y$ is reduced, so that $f$ factors as
\[
Y\xrightarrow{g}X_0\xrightarrow{j}X,
\]
where $X_0$ denotes the reduced sub-prescheme of $X$ whose underlying
space is the connected component $X_0$ \sref[I]{I.5.2.2}.
Apply \sref{IV.4.5.13} to $g$.

For (ii), choose a prescheme structure on the underlying space $Y$.
Since a geometrically irreducible prescheme is in particular
geometrically connected, (i) applies to the canonical injection
$Y\to X$.
\end{proof}

\oldpage[IV-2]{65}

\begin{corollary}[4.5.14]
\phantomsection
\label{IV.4.5.14}
Let $X$ be a $k$-prescheme.
If $x\in X$ is such that $\kappa(x)$ is a primary extension of $k$
(in particular, if $x$ is a rational point of $X$), then the
connected component of $x$ in $X$ is geometrically connected.
\end{corollary}

\begin{proof}
Apply \sref{IV.4.5.13.1} to
$Y=\Spec(\kappa(x))$ and to the canonical morphism $Y\to X$.
By \sref{IV.4.5.9}, the prescheme $Y$ is geometrically irreducible.
\end{proof}

\begin{proposition}[4.5.15]
\phantomsection
\label{IV.4.5.15}
Let $X$ be a $k$-prescheme, let $x\in X$, and let $k'$ be the
separable algebraic closure of $k$ in $\kappa(x)$.
Suppose that:
\begin{enumerate}
\item[\rm(i)]
\phantomsection
\label{IV.4.5.15.i}
$X$ is connected;

\item[\rm(ii)]
\phantomsection
\label{IV.4.5.15.ii}
$k'$ is a finite extension of $k$.
\end{enumerate}
Then the geometric number of connected components of $X$ is at most
$[k':k]$.
Moreover, if $k''$ is a finite Galois extension of $k$ containing
$k'$, the connected components of $X\otimes_k k''$ are geometrically
connected.
\end{proposition}

Condition (ii) is satisfied in particular when $X$ is locally of
finite type over $k$.

\begin{proof}
By (ii), there are finite Galois extensions $k''$ of $k$ containing
$k'$.
Put $X''=X\otimes_k k''$.
The projection $p:X''\to X$ is faithfully flat and finite, hence
closed \sref[II]{II.6.1.10}.
It follows from \sref{IV.2.3.6},~(ii), that the image under $p$ of
every connected component $X''_\alpha$ of $X''$ is all of $X$.
Thus $X''_\alpha$ contains a point
$x''_\alpha\in p^{-1}(x)$.

The number of points of $p^{-1}(x)$ is the number of irreducible
components of
\[
\Spec\bigl(\kappa(x)\otimes_k k''\bigr)
\]
\sref[I]{I.3.4.9}, namely $[k':k]$ by \sref{IV.4.3.4}.
Therefore the number of connected components of $X''$ is at most
$[k':k]$.
For every $x''_\alpha\in p^{-1}(x)$, the field
$\kappa(x''_\alpha)$ is a primary extension of $k''$, by
\sref{IV.4.3.5} and \sref[I]{I.3.4.9}.
Corollary \sref{IV.4.5.14}, applied to $x''_\alpha$ and $X''$, now
shows that every connected component of $X''$ is geometrically
connected.
\end{proof}

\begin{corollary}[4.5.16]
\phantomsection
\label{IV.4.5.16}
Let $X$ be a $k$-prescheme, let $(X_\alpha)$ be the family of its
connected components, and choose $x_\alpha\in X_\alpha$ for every
$\alpha$.
Suppose that:
\begin{enumerate}
\item[\rm(i)]
\phantomsection
\label{IV.4.5.16.i}
the family $(X_\alpha)$ is finite;

\item[\rm(ii)]
\phantomsection
\label{IV.4.5.16.ii}
for every $\alpha$, the separable algebraic closure $k'_\alpha$ of
$k$ in $\kappa(x_\alpha)$ is a finite extension of $k$.
\end{enumerate}
Then the geometric number of connected components of $X$ is at most
\[
\sum_\alpha[k'_\alpha:k],
\]
and there exists a finite separable extension $k''$ of $k$ such that
all the connected components of $X\otimes_k k''$ are geometrically
connected.
\end{corollary}

\begin{proof}
By (i), the connected components of $X$ are open.
Apply \sref{IV.4.5.15} to each sub-prescheme induced on the open
subset $X_\alpha$.
To obtain one extension $k''$ with the asserted property, take a
finite Galois extension of $k$ containing all the $k'_\alpha$.
\end{proof}

\oldpage[IV-2]{66}

\begin{corollary}[4.5.17]
\phantomsection
\label{IV.4.5.17}
Suppose that the $k$-prescheme $X$ contains a point $x$ such that the
separable algebraic closure of $k$ in $\kappa(x)$ is finite over $k$.
Then $X$ is geometrically connected if and only if
$X\otimes_k K$ is connected for every finite separable extension
$K$ of $k$.
\end{corollary}

\begin{remark}[4.5.18]
\phantomsection
\label{IV.4.5.18}
We shall see in \S~8 \sref{IV.8.4.5} that the conclusion of
\sref{IV.4.5.17} remains valid if, instead of the hypothesis in its
statement, one assumes that $X$ is quasi-compact.
\end{remark}

\begin{proposition}[4.5.19]
\phantomsection
\label{IV.4.5.19}
Let $X$ be a $k$-prescheme, let $Z$ be a subset of $X$, let $k'$ be
an algebraically closed extension of $k$, put
\[
X'=X\otimes_k k',
\]
and let $p:X'\to X$ be the canonical projection.
Suppose that $Z'=p^{-1}(Z)$ is contained in a single irreducible
component $X'_0$ of $X'$.
Then
\[
X'_0=p^{-1}(X_0),
\]
where $X_0=p(X'_0)$ is an irreducible component of $X$ containing
$Z$; moreover, $X_0$ is geometrically irreducible.
If, in addition, $X'_0$ is the only irreducible component of $X'$
which meets $Z'$, then $X_0$ is the only irreducible component of $X$
which meets $Z$.
\end{proposition}

To prove that $X'_0=p^{-1}(X_0)$, we shall use the following lemma.

\begin{lemma}[4.5.19.1]
\phantomsection
\label{IV.4.5.19.1}
Let $T$ and $T'$ be preschemes, let $p:T'\to T$ be a morphism, put
\[
T''=T'\times_T T',
\]
and let $p_1,p_2:T''\to T'$ be the canonical projections.
A subset $U'$ of $T'$ is of the form $p^{-1}(U)$, for some subset
$U\subset T$, if and only if
\[
p_1^{-1}(U')=p_2^{-1}(U').
\]
\end{lemma}

\begin{proof}
The structure morphism $q:T''\to T$ is, by definition, both
$p\circ p_1$ and $p\circ p_2$, so the condition is necessary.
Conversely, suppose the condition holds.
Let $u'\in U'$ and let $t'\in p^{-1}(p(u'))$.
Since $p(t')=p(u')$, there is a point $t''\in T''$ such that
\[
p_1(t'')=u',\qquad p_2(t'')=t'
\]
\sref[I]{I.3.4.7}.
Thus
$t''\in p_1^{-1}(U')=p_2^{-1}(U')$, and consequently
$t'=p_2(t'')\in U'$.
This proves the lemma.
\end{proof}

\begin{proof}[Proof of {\rm \sref{IV.4.5.19}}]
Form the fibre product
\[
X''=X'\times_X X'
\]
relative to $p:X'\to X$, and let $p_1,p_2:X''\to X'$ be the
canonical projections.
By \sref{IV.4.5.19.1}, it is enough to prove that
\[
p_1^{-1}(X'_0)=p_2^{-1}(X'_0).
\]
Put
\[
S=\Spec(k),\qquad S'=\Spec(k'),\qquad
S''=S'\times_S S'=\Spec(k'\otimes_k k').
\]
Then $X''$ may also be written $X'\times_{S'}S''$.
If $\varphi'':X''\to S''$ is the structure morphism, it is enough to
show, for every $s''\in S''$, that
\[
p_1^{-1}(X'_0)\cap\varphi''^{-1}(s'')
 =
p_2^{-1}(X'_0)\cap\varphi''^{-1}(s'').
\]
Put
\[
T=\Spec(\kappa(s'')),\qquad
U=X''\times_{S''}T=\varphi''^{-1}(s''),
\]
and let $v:U\to X''$ be the canonical projection.
We must prove that
\[
U_1=v^{-1}\bigl(p_1^{-1}(X'_0)\bigr),
\qquad
U_2=v^{-1}\bigl(p_2^{-1}(X'_0)\bigr),
\]
which are irreducible components of $U$ by \sref{IV.4.4.5}, are
equal.
The relevant morphisms form the commutative diagram
\[
\label{IV.4.5.19.2}
\tag{4.5.19.2}
\xymatrix@C=3.2em@R=3.0em{
X\ar[d]
  & X'\ar[l]_-p\ar[d]
  & X''\ar@<.55ex>[l]^-{p_2}\ar@<-.55ex>[l]_-{p_1}
       \ar[d]^-{\varphi''}
  & U\ar[l]_-v\ar[d]
\\
S
  & S'\ar[l]
  & S''\ar@<.55ex>[l]\ar@<-.55ex>[l]
  & T\ar[l].
}
\]
By construction, $U_1$ and $U_2$ both contain $w^{-1}(Z)$, where
\[
w=p\circ p_1\circ v=p\circ p_2\circ v.
\]

\oldpage[IV-2]{67}

There is a field $K$ which is a common extension of $k'$ and
$\kappa(s'')$.
If $P=\Spec(K)$, the diagram
\[
\xymatrix@C=3.5em@R=3em{
S'\ar[d] & P\ar[l]\ar[d]\\
S        & T\ar[l]
}
\]
is commutative.
Putting $Q=X\otimes_k K$, the diagram
\[
\xymatrix@C=3.5em@R=3em{
X'\ar[d]_-p & Q\ar[l]_-q\ar[d]^-r\\
X           & U\ar[l]^-w
}
\]
is likewise commutative.
The two irreducible components $r^{-1}(U_1)$ and $r^{-1}(U_2)$ of
$Q$ \sref{IV.4.4.5} therefore contain $q^{-1}(Z')$.
By \sref{IV.4.4.5}, their images
\[
q\bigl(r^{-1}(U_1)\bigr)
\quad\hbox{and}\quad
q\bigl(r^{-1}(U_2)\bigr)
\]
are irreducible components of $X'$ containing $Z'$.
Both are therefore $X'_0$.
Again by \sref{IV.4.4.5},
$r^{-1}(U_1)=r^{-1}(U_2)$, and hence $U_1=U_2$.

Since $p$ is surjective and $X'_0$ dominates an irreducible component
of $X$, the subset $X_0=p(X'_0)$ is an irreducible component of $X$
containing $Z$.
Since $p^{-1}(X_0)=X'_0$ is an irreducible component of $X'$, the
component $X_0$ is geometrically irreducible.

For the last assertion, let $X_1$ be another irreducible component of
$X$ meeting $Z$.
Choose $z\in Z\cap X_1$.
There is an irreducible component of $X'$ which meets $Z'$ and
dominates $X_1$ \sref{IV.2.3.5}.
By hypothesis $X'_0$ is the only such component, so necessarily
$X_1=X_0$.
\end{proof}

\begin{remark}[4.5.20]
\phantomsection
\label{IV.4.5.20}
The fact that $Z'$ is contained in a single irreducible component
$X'_0$ of $X'$ does \emph{not} imply that
$X_0=p(X'_0)$ is the only irreducible component of $X$ containing
$Z$.
Indeed, take $k=\mathbf R$, $k'=\mathbf C$, and let $X$ be the
$k$-scheme obtained by gluing\footnote{The technique of gluing
preschemes will be developed in detail in Chapter~V.  For the case
considered here, which concerns algebraic curves, see [38],
pp.~68--71.} the following two schemes at a nongeneric point of each:
\[
X_1=\Spec\bigl(\mathbf R[S,T]/(S^2+T^2+1)\bigr),
\qquad
X_2=\Spec\bigl(\mathbf C[T]\bigr);
\]
the local rings of these two curves at each of their closed points
have the same residue field, isomorphic to $\mathbf C$.
The schemes $X_1$ and $X_2$ identify with the two irreducible
components of $X$.
If $x$ is their common point, then $p^{-1}(x)$ consists of two
distinct points $y'$ and $z'$; the inverse image $p^{-1}(X_1)$ is
irreducible, while $p^{-1}(X_2)$ is the union of two irreducible
components $Y'$ and $Z'$ of $X'$, with $y'\in Y'$ and $z'\in Z'$.
Thus $p^{-1}(x)$ is contained in only one irreducible component of
$X'$.
\end{remark}

\begin{proposition}[4.5.21]
\phantomsection
\label{IV.4.5.21}
Let $k$ be a separably closed field and let $k'$ be an algebraic
closure of $k$.
For every $k$-prescheme $X$, the canonical projection
\[
X\otimes_k k'\longrightarrow X
\]
is a universal homeomorphism.

\oldpage[IV-2]{68}

In particular, the irreducible components (resp.\ connected
components) of $X$ are geometrically irreducible
(resp.\ geometrically connected).
\end{proposition}

\begin{proof}
By definition, $k'/k$ is radicial.
The morphism $\Spec(k')\to\Spec(k)$ is integral, surjective, and
radicial, so the first assertion follows from
\sref{IV.2.4.5},~(i).
The second follows from the first and \sref{IV.4.5.1}.
\end{proof}
% END NATIVE INPUT components/05_section45_geometric_irreducibility_connectedness
% BEGIN NATIVE INPUT components/06_section46_geometrically_reduced_preschemes
\subsection{Geometrically reduced algebraic preschemes}
\label{subsection:IV.4.6}
\label{IV.4.6}

\begin{proposition}[4.6.1]
\phantomsection
\label{IV.4.6.1}
Let $k$ be a field, let $X$ be a $k$-prescheme, and let $\Omega$ be a
perfect extension of $k$.
The following conditions are equivalent:
\begin{enumerate}
\item[\rm(a)]
\phantomsection
\label{IV.4.6.1.a}
for every reduced $k$-prescheme $S$, the prescheme $X\times_k S$ is
reduced;

\item[\rm(b)]
\phantomsection
\label{IV.4.6.1.b}
for every extension $K$ of $k$, the prescheme $X\otimes_k K$ is
reduced;

\item[\rm(c)]
\phantomsection
\label{IV.4.6.1.c}
the prescheme $X\otimes_k\Omega$ is reduced;

\item[\rm(d)]
\phantomsection
\label{IV.4.6.1.d}
for every finite radicial extension $k'$ of $k$, the prescheme
$X\otimes_k k'$ is reduced;

\item[\rm(e)]
\phantomsection
\label{IV.4.6.1.e}
$X$ is reduced and, for every irreducible component $X_\alpha$ of
$X$, with generic point $x_\alpha$, the field $\kappa(x_\alpha)$ is a
separable extension of $k$.
\end{enumerate}
\end{proposition}

\begin{proof}
The implications \emph{(a)}$\Rightarrow$\emph{(b)}
$\Rightarrow$\emph{(c)} are immediate.
Since $k'$ may be regarded as a subextension of $\Omega$,
\sref{IV.4.4.1} shows that \emph{(c)} implies \emph{(d)}.

Suppose \emph{(d)}.
Taking $k'=k$ first shows that $X$ is reduced.
To prove that $\kappa(x_\alpha)$ is separable over $k$, it suffices to
show that $\kappa(x_\alpha)\otimes_k k'$ is reduced for every finite
radicial extension $k'$ of $k$ \sref{IV.4.3.5}.
We may reduce to the case $X=\Spec(A)$, where $A$ is a reduced
$k$-algebra.
Then $A\otimes_k k'$ is reduced by hypothesis \sref[I]{I.5.1.4}, and
$\kappa(x_\alpha)\otimes_k k'$ is a ring of fractions of this ring;
hence it is reduced $(\mathbf{0}_I,\,1.2.8)$.
Thus \emph{(d)} implies \emph{(e)}.

Finally, suppose \emph{(e)}.
We may reduce to the affine case
$X=\Spec(A)$ and $S=\Spec(B)$, where $A$ and $B$ are $k$-algebras.
The fields $\kappa(x_\alpha)$ are the fields of fractions of the
quotient rings $A/\mathfrak p_\alpha$, where
$\mathfrak p_\alpha=\mathfrak j_{x_\alpha}$ are the minimal prime
ideals of $A$.
Since $\bigcap_\alpha\mathfrak p_\alpha=0$, the algebra $A$ embeds in
$\prod_\alpha A/\mathfrak p_\alpha$, and hence in
$\prod_\alpha\kappa(x_\alpha)$.
Likewise, $B$ may be regarded as a subalgebra of a product
$\prod_\beta L_\beta$ of extensions of $k$.
It follows that $A\otimes_k B$ embeds in
\[
 \left(\prod_\alpha\kappa(x_\alpha)\right)
 \otimes_k
 \left(\prod_\beta L_\beta\right),
\]
and this tensor product itself embeds in
\[
 \prod_{\alpha,\beta}
 \bigl(\kappa(x_\alpha)\otimes_k L_\beta\bigr)
\]
(Bourbaki, \emph{Alg.}, chap.~II, 3rd ed., \S~7,
n\textsuperscript{o}~7, prop.~15).
By hypothesis the rings
$\kappa(x_\alpha)\otimes_k L_\beta$ are reduced
\sref{IV.4.3.5}; hence $A\otimes_k B$ is reduced.
Thus \emph{(e)} implies \emph{(a)}.
\end{proof}

\begin{definition}[4.6.2]
\phantomsection
\label{IV.4.6.2}
If the equivalent conditions \emph{(a)}--\emph{(e)} of
\sref{IV.4.6.1} hold, $X$ is said to be \emph{separable}
(or \emph{geometrically reduced}, or \emph{universally reduced}) over
$k$.
A $k$-prescheme $X$ is said to be \emph{geometrically integral} over
$k$ if, for every extension $K$ of $k$, the prescheme
$X\otimes_k K$ is integral.
By \sref{IV.4.6.1}, this is equivalent to saying that $X$ is separable
and geometrically irreducible over $k$.
\end{definition}

A commutative $k$-algebra $A$ is said to be \emph{separable} if
$\Spec(A)$ is separable over $k$; equivalently, for every extension
$K$ of $k$, the ring $A_{(K)}=A\otimes_k K$ is reduced.
This definition agrees with that of Bourbaki, \emph{Alg.}, chap.~VIII,
\oldpage[IV-2]{69}
\S~7, n\textsuperscript{o}~5, def.~1, when $A$ has finite rank over
$k$ (loc.\ cit., corollary to prop.~7), but not in general: a
$k$-algebra may have nonzero radical even when it is integral.

\begin{corollary}[4.6.3]
\phantomsection
\label{IV.4.6.3}
Let $X$ be an integral $k$-prescheme.
Then $X$ is geometrically reduced (resp.\ geometrically integral) over
$k$ if and only if its field of rational functions $\mathrm{R}(X)$ is
a separable (resp.\ separable and primary) extension of $k$.
\end{corollary}

\begin{proof}
This follows immediately from \sref{IV.4.5.9} and
\sref{IV.4.6.1}.
\end{proof}

\begin{corollary}[4.6.4]
\phantomsection
\label{IV.4.6.4}
Let $X$ be a reduced $k$-prescheme.
For every separable extension $k'$ of $k$, the prescheme
$X'=X\otimes_k k'$ is reduced.
\end{corollary}

\begin{proof}
Apply the equivalence of \emph{(a)} and \emph{(e)} in
\sref{IV.4.6.1}, replacing $X$ there by $\Spec(k')$ and $S$ by $X$.
\end{proof}

\begin{proposition}[4.6.5]
\phantomsection
\label{IV.4.6.5}
\begin{enumerate}
\item[\rm(i)]
\phantomsection
\label{IV.4.6.5.i}
Let $X$ be a $k$-prescheme and let $K$ be an extension of $k$.
Then $X$ is geometrically reduced (resp.\ geometrically integral) over
$k$ if and only if $X\otimes_k K$ is geometrically reduced
(resp.\ geometrically integral) over $K$.

\item[\rm(ii)]
\phantomsection
\label{IV.4.6.5.ii}
Let $X$ and $Y$ be two $k$-preschemes.
If $X$ and $Y$ are geometrically reduced
(resp.\ geometrically integral) over $k$, then so is $X\times_k Y$.
\end{enumerate}
\end{proposition}

\begin{proof}
Assertion~\emph{(i)} is an immediate consequence of the definitions
and \sref{IV.4.4.1}.
For the reduced assertion in~\emph{(ii)}, let $\Omega$ be an
algebraically closed extension of $k$.
Then
\[
 (X\times_kY)\otimes_k\Omega
 =X\times_k(Y\otimes_k\Omega).
\]
Since $Y\otimes_k\Omega$ is reduced, the right-hand side is reduced by
\sref{IV.4.6.1.a}.
The remainder of~\emph{(ii)} follows from this and
\sref{IV.4.5.8.ii}.
\end{proof}

\begin{proposition}[4.6.6]
\phantomsection
\label{IV.4.6.6}
Let $X$ be a $k$-prescheme of finite type.
There exists a finite radicial extension $k'$ of $k$ such that
$(X_{(k')})_{\mathrm{red}}$ is geometrically reduced over $k'$.
\end{proposition}

\begin{proof}
Since $X$ is covered by finitely many affine open subsets $U_i$, we
may reduce to the affine case.
Indeed, suppose for each $i$ that $k_i$ is a finite radicial extension
of $k$ such that
$(U_i\otimes_k k_i)_{\mathrm{red}}$ is geometrically reduced over
$k_i$.
We may choose a common finite radicial extension $k'$ of $k$ containing
all the $k_i$.
By \sref[I]{I.5.1.8},
\[
 (U_i\otimes_k k')_{\mathrm{red}}
 \simeq
 (U_i\otimes_k k_i)_{\mathrm{red}}\otimes_{k_i}k',
\]
and this is geometrically reduced over $k'$.
Thus the same holds for $(X\otimes_k k')_{\mathrm{red}}$.

Now let $X=\Spec(A)$, where $A$ is a $k$-algebra of finite type, and
let $\Omega$ be an algebraic closure of $k$.
Write $\mathfrak N'$ for the nilradical of $A\otimes_k\Omega$.
Since $A\otimes_k\Omega$ is Noetherian, $\mathfrak N'$ is generated by
finitely many elements
\[
 y_i=\sum_j x_{ij}\otimes\xi_{ij},
 \qquad x_{ij}\in A,\quad \xi_{ij}\in\Omega.
\]
Let $K$ be a finite subextension of $\Omega$ containing all the
$\xi_{ij}$, and let $\mathfrak N$ be the ideal of $A\otimes_kK$
generated by the $y_i$, regarding $A\otimes_kK$ as a subring of
$A\otimes_k\Omega$.
The $y_i$ are nilpotent in $A\otimes_kK$.
Moreover,
$\mathfrak N\otimes_K\Omega=\mathfrak N'$, and hence
\[
 \mathfrak N'\cap(A\otimes_kK)=\mathfrak N.
\]
Thus $\mathfrak N$ contains the nilradical of $A\otimes_kK$ and is
therefore equal to it.
Since
\[
 (A\otimes_k\Omega)/\mathfrak N'
 =
 \bigl((A\otimes_kK)/\mathfrak N\bigr)\otimes_K\Omega,
\]
the prescheme
$(X\otimes_kK)_{\mathrm{red}}\otimes_K\Omega$ is reduced.
Consequently $(X_{(K)})_{\mathrm{red}}$ is geometrically reduced over
$K$ by \sref{IV.4.6.1}.

Replacing $K$ by the quasi-Galois extension of $k$ that it generates
in $\Omega$, we may further suppose, by \sref[I]{I.5.1.8}, that $K$ is
quasi-Galois over $k$.
It is then a separable extension of a finite radicial extension $k'$
of $k$ (Bourbaki, \emph{Alg.}, chap.~V, \S~10,
n\textsuperscript{o}~9, prop.~14).
\oldpage[IV-2]{70}
By \sref[I]{I.5.1.8},
$(X_{(k')})_{\mathrm{red}}\otimes_{k'}K$ is isomorphic to
$(X_{(K)})_{\mathrm{red}}$, and hence is geometrically reduced.
It follows from \sref{IV.4.6.5.i} that
$(X_{(k')})_{\mathrm{red}}$ is geometrically reduced.
\end{proof}

\begin{corollary}[4.6.7]
\phantomsection
\label{IV.4.6.7}
If $K$ is an extension of finite type of $k$, there exists a finite
radicial extension $k'$ of $k$ such that the residue fields of the
semilocal ring $K\otimes_k k'$ are separable over $k'$.
\end{corollary}

\begin{proof}
Apply \sref{IV.4.6.6} to $X=\Spec(A)$, where $A$ is a $k$-algebra of
finite type having $K$ as its field of fractions.
\end{proof}

\begin{corollary}[4.6.8]
\phantomsection
\label{IV.4.6.8}
Let $X$ be a $k$-prescheme of finite type.
There exists a finite extension $k'$ of $k$ such that
$(X_{(k')})_{\mathrm{red}}$ is geometrically reduced over $k'$, the
irreducible components of $X_{(k')}$ are geometrically irreducible,
and its connected components are geometrically connected.
\end{corollary}

\begin{proof}
This follows immediately from \sref{IV.4.6.6},
\sref{IV.4.5.10}, and \sref{IV.4.5.15}.
\end{proof}

\begin{definition}[4.6.9]
\phantomsection
\label{IV.4.6.9}
Let $k$ be a field, let $X$ be a $k$-prescheme, and let $x\in X$.
We say that $X$ is \emph{geometrically reduced}
(or \emph{separable}) (resp.\ \emph{geometrically punctually
integral}) at $x$ over $k$ if, for every extension $k'$ of $k$ and
every point $x'$ of $X'=X\otimes_k k'$ above $x$, the prescheme $X'$
is reduced (resp.\ integral) at $x'$; that is,
$\sh O_{X',x'}$ is reduced (resp.\ integral).

We say that $X$ is \emph{geometrically punctually integral} if it is
geometrically punctually integral at all of its points.
Equivalently, for every extension $k'$ of $k$, all local rings of
$X\otimes_k k'$ are integral; in that case $X\otimes_k k'$ is also
said to be \emph{punctually integral}.
\end{definition}

By \sref{IV.4.6.2}, $X$ is geometrically reduced over $k$ if and only
if it is geometrically reduced over $k$ at every point $x\in X$.

\begin{proposition}[4.6.10]
\phantomsection
\label{IV.4.6.10}
Let $k$ be a field, let $X$ be a $k$-prescheme, let $k'$ be an
extension of $k$, and let $x'$ be a point of
$X'=X\otimes_k k'$ with image $x$ in $X$.
Then $X$ is geometrically reduced
(resp.\ geometrically punctually integral) at $x$ over $k$ if and only
if $X'$ has the corresponding property at $x'$ over $k'$.
In particular, $X$ is geometrically punctually integral over $k$ if
and only if $X'$ is geometrically punctually integral over $k'$.
\end{proposition}

\begin{proof}
It is enough to prove the first assertion.
Necessity is immediate.
Suppose conversely that $X'$ is geometrically reduced
(resp.\ geometrically punctually integral) at $x'$ over $k'$.
Let $k''$ be an extension of $k$ and let $x''$ be a point of
$X''=X\otimes_k k''$ above $x$.
There is a point $z$ of $X'\times_X X''$ projecting to $x'$ and $x''$
\sref[I]{I.3.4.7}.
Let $s$ be the image of $z$ in
$\Spec(k'\otimes_k k'')$ \sref[I]{I.3.4.9}, put
$K=\kappa(s)$, and put $Z=X\otimes_kK$.
Then $K$ may be regarded as a composite extension of $k'$ and $k''$,
and $z$ as a point of $Z$ whose images in $X'$ and $X''$ are,
respectively, $x'$ and $x''$.
By hypothesis, $\sh O_{Z,z}$ is reduced (resp.\ integral).
Since $\sh O_{Z,z}$ is a faithfully flat
$\sh O_{X'',x''}$-module, it follows that $\sh O_{X'',x''}$ is
reduced (resp.\ integral) $(\mathbf{0}_I,\,6.5.1)$.
\end{proof}

\begin{corollary}[4.6.11]
\phantomsection
\label{IV.4.6.11}
Suppose that $k$ is perfect (resp.\ algebraically closed).
Then $X$ is geometrically reduced
(resp.\ geometrically punctually integral) at $x$ over $k$ if and
only if $\sh O_{X,x}$ is reduced (resp.\ integral).
\end{corollary}

\begin{proof}
Necessity is immediate.
Conversely, for every extension $k'$ of $k$, the prescheme
$\Spec(\sh O_{X,x})\otimes_k k'$ is reduced
\sref{IV.4.6.4} (resp.\ integral \sref{IV.4.4.4}).
\oldpage[IV-2]{71}
For every point $x'$ of $X'=X\otimes_k k'$ above $x$, the local ring
$\sh O_{X',x'}$ is isomorphic to a local ring of
$\Spec(\sh O_{X,x})\otimes_k k'$, and is therefore reduced
(resp.\ integral).
\end{proof}

\begin{proposition}[4.6.12]
\phantomsection
\label{IV.4.6.12}
Let $k$ be a field, let $X$ be a $k$-prescheme, let $x\in X$, and let
$\Omega$ be a perfect extension of $k$.
The following conditions are equivalent:
\begin{enumerate}
\item[\rm(a)]
\phantomsection
\label{IV.4.6.12.a}
$X$ is geometrically reduced over $k$ at $x$; that is, for every
extension $k'$ of $k$ and every point $x'$ of
$X'=X\otimes_k k'$ above $x$, the local ring
$\sh O_{X',x'}$ is reduced;

\item[\rm(b)]
\phantomsection
\label{IV.4.6.12.b}
the prescheme $X\otimes_k\Omega$ is reduced at one point above $x$;

\item[\rm(c)]
\phantomsection
\label{IV.4.6.12.c}
for every finite radicial extension $k'$ of $k$, the prescheme
$X'=X\otimes_k k'$ is reduced at the unique point $x'$ above $x$;

\item[\rm(d)]
\phantomsection
\label{IV.4.6.12.d}
the prescheme $\Spec(\sh O_{X,x})$ is geometrically reduced over $k$;

\item[\rm(e)]
\phantomsection
\label{IV.4.6.12.e}
the local ring $\sh O_{X,x}$ is reduced and, for every irreducible
component $Z$ of $X$ containing $x$, with generic point $z$, the field
$\kappa(z)$ is a separable extension of $k$.
\end{enumerate}
\end{proposition}

\begin{proof}
The implications
\emph{(d)}$\Rightarrow$\emph{(a)}$\Rightarrow$\emph{(b)}
$\Rightarrow$\emph{(c)} are immediate, using for the last implication
\sref{IV.4.6.10}, \sref{IV.4.6.11}, and the fact that every radicial
extension of $k$ is isomorphic to a subextension of $\Omega$.

Suppose \emph{(c)}.
By \sref{IV.4.6.1}, it suffices to show that
$\sh O_{X,x}\otimes_k k'$ is reduced for every finite radicial
extension $k'$ of $k$.
This ring is the local ring of $X'=X\otimes_k k'$ at the unique point
$x'$ above $x$ \sref[I]{I.3.5.8} \sref[I]{I.3.5.7}, and is reduced by
\emph{(c)}.
Thus \emph{(c)} implies \emph{(d)}.

Finally, the equivalence of \emph{(d)} and \emph{(e)} follows from the
equivalence of \emph{(b)} and \emph{(e)} in \sref{IV.4.6.1}, since the
irreducible components of $X$ containing $x$ correspond bijectively
to the irreducible components of $\Spec(\sh O_{X,x})$
\sref[I]{I.2.4.2}.
\end{proof}

\begin{corollary}[4.6.13]
\phantomsection
\label{IV.4.6.13}
Under the hypotheses of \sref{IV.4.6.12}, suppose in addition that
$X$ is locally Noetherian.
Then conditions \emph{(a)}--\emph{(e)} of \sref{IV.4.6.12} are also
equivalent to:
\begin{enumerate}
\item[\rm(f)]
\phantomsection
\label{IV.4.6.13.f}
there exists an open neighbourhood $U$ of $x$ that is geometrically
reduced over $k$.
\end{enumerate}
\end{corollary}

\begin{proof}
Condition~\emph{(f)} plainly implies that $X$ is geometrically reduced
at $x$ over $k$.
Conversely, suppose that $X$ is geometrically reduced at $x$.
Since $\sh O_{X,x}$ is reduced, there is a reduced open neighbourhood
$U$ of $x$ in $X$ \sref[I]{I.6.1.13}.
After shrinking $U$, we may further suppose that it meets no
irreducible component of $X$ that does not contain $x$.
Condition~\emph{(e)} of \sref{IV.4.6.12}, together with
condition~\emph{(e)} of \sref{IV.4.6.1}, then shows that $U$ is
geometrically reduced over $k$.
\end{proof}

It follows from \sref{IV.4.6.13} that, when $X$ is locally Noetherian,
the set of points $x\in X$ at which $X$ is geometrically reduced over
$k$ is open; it is the largest open subset of $X$ that is
geometrically reduced over $k$.

\begin{env}[4.6.14]
\phantomsection
\label{IV.4.6.14}
It follows from \sref{IV.4.6.10} and \sref{IV.4.6.11} that, if
$\Omega$ is an algebraically closed extension of $k$, saying that $X$
is geometrically punctually integral at $x$ is equivalent to saying
that $X\otimes_k\Omega$ is integral at one point above $x$.

For $X$ to be geometrically punctually integral at $x$, it is
necessary that $X$ be integral at $x$ and that, if $z$ is the generic
point of the unique irreducible component of $X$ containing $x$, the
field $\kappa(z)$ be a separable extension of $k$.
This follows from \sref{IV.4.6.10}.
These conditions are not sufficient, as is shown by
\sref{IV.4.5.12.ii}.
A sufficient but not necessary condition is that $X$ be geometrically
reduced at $x$ and belong to only one irreducible component of $X$,
and that this component be geometrically irreducible.
\oldpage[IV-2]{72}
If in addition $X$ is locally Noetherian, then $x$ has an open
neighbourhood in $X$ that is geometrically integral.
\end{env}

Recall that a locally Noetherian prescheme that is punctually integral
is locally integral \sref[I]{I.6.1.13}.
Thus, if a $k$-prescheme $X$ locally of finite type over $k$ is
geometrically punctually integral over $k$, then
$X\otimes_k k'$ is locally integral for every extension $k'$ of $k$.
In this case $X$ is also said to be
\emph{geometrically locally integral}.

\begin{proposition}[4.6.15]
\phantomsection
\label{IV.4.6.15}
\begin{enumerate}
\item[\rm(i)]
\phantomsection
\label{IV.4.6.15.i}
Let $k$ be a field and let $X$ be a $k$-prescheme of finite type.
Then $X$ is geometrically punctually integral if and only if it is
geometrically reduced and its geometric number of irreducible
components equals its geometric number of connected components.

\item[\rm(ii)]
\phantomsection
\label{IV.4.6.15.ii}
Let $X$ be a prescheme locally of finite type over $k$.
If $X$ is geometrically punctually integral at $x$, then $x$ has an
open neighbourhood $U$ that is geometrically punctually integral.
Equivalently, the set of points $x\in X$ at which $X$ is geometrically
punctually integral is open in $X$.
\end{enumerate}
\end{proposition}

\begin{proof}
\begin{enumerate}
\item[\rm(i)]
Let $\Omega$ be an algebraically closed extension of $k$.
The prescheme $X_{(\Omega)}$ is punctually integral
(equivalently, locally integral, since it is Noetherian) if and only
if it is reduced and the number of its connected components equals
the number of its irreducible components \sref[I]{I.6.1.10}.
Now apply \sref{IV.4.5.1}, \sref{IV.4.6.1}, and the first observation
in \sref{IV.4.6.14}.

\item[\rm(ii)]
The question is local on $X$, so we may reduce to the case in which
$X$ is of finite type over $k$.
By \sref{IV.4.6.13}, $x$ has a geometrically reduced open
neighbourhood; hence we may suppose that $X$ itself is geometrically
reduced.
By \sref{IV.4.6.8}, there is a finite extension $k'$ of $k$ such that
$X'=X\otimes_k k'$ is reduced and its irreducible components are
geometrically irreducible.
Let $p:X'\to X$ be the canonical projection, which is finite and
surjective, and let $x'_j$ ($1\leq j\leq r$) be the points of
$p^{-1}(x)$.
By hypothesis $X'$ is integral at each $x'_j$, so each $x'_j$ has an
integral open neighbourhood $V'_j$ in $X'$.
Since $p$ is closed \sref[II]{II.6.1.10}, there is an open
neighbourhood $U$ of $x$ such that $p^{-1}(U)$ is contained in
$\bigcup_jV'_j$.
Thus $p^{-1}(U)$ is integral at every point, and hence locally
integral.
Its irreducible components are consequently open and pairwise
disjoint; since they are geometrically irreducible
\sref{IV.4.5.9}, assertion~\emph{(i)} shows that
$p^{-1}(U)=U\otimes_k k'$ is geometrically punctually integral.
By definition the same is true of $U$.
\end{enumerate}
\end{proof}

\begin{proposition}[4.6.16]
\phantomsection
\label{IV.4.6.16}
Let $k$ be a field, let $X$ be a locally Noetherian $k$-prescheme,
and let $\sh F$ be a coherent $\sh O_X$-Module.
The following conditions are equivalent:
\begin{enumerate}
\item[\rm(a)]
\phantomsection
\label{IV.4.6.16.a}
for every extension of finite type $k'$ of $k$, if
$X'=X\otimes_k k'$, then
$\sh F'=\sh F\otimes_k k'$ is a reduced
$\sh O_{X'}$-Module \sref{IV.3.2.2};

\item[\rm(b)]
\phantomsection
\label{IV.4.6.16.b}
for every finite radicial extension $k'$ of $k$, the Module $\sh F'$
is a reduced $\sh O_{X'}$-Module;

\item[\rm(c)]
\phantomsection
\label{IV.4.6.16.c}
$\sh F$ is reduced and, if $\sh J$ is the Ideal of $\sh O_X$
annihilating $\sh F$, the closed sub-prescheme defined by $\sh J$ is
geometrically reduced over $k$.
\end{enumerate}
If, moreover, $X$ is locally of finite type over $k$, these conditions
are also equivalent to:
\oldpage[IV-2]{73}
\begin{enumerate}
\item[\rm(d)]
\phantomsection
\label{IV.4.6.16.d}
for every extension $k'$ of $k$
(or for one algebraically closed extension $k'$ of $k$), the Module
$\sh F'$ is a reduced $\sh O_{X'}$-Module.
\end{enumerate}
\end{proposition}

\begin{proof}
It is clear that \emph{(a)} implies \emph{(b)}.
Condition~\emph{(b)} first implies that $\sh F$ is reduced.
By \sref{IV.3.2.3}, this means the following.
If $Y$ is the closed sub-prescheme of $X$ defined by the coherent
Ideal $\sh J$ annihilating $\sh F$, then $Y$ is reduced and there is a
coherent torsion-free $\sh O_Y$-Module $\sh G$, of rank~$1$ on every
irreducible component of $Y$, such that
$j_*(\sh G)=\sh F$, where $j:Y\to X$ is the canonical immersion.

For every extension $k'$ of $k$, the annihilator of $\sh F'$ is
$\sh J'=\sh J\otimes_k k'$, since the projection $X'\to X$ is flat
\sref{IV.2.1.11}.
The closed sub-prescheme of $X'$ defined by $\sh J'$ is
$Y'=Y\otimes_k k'$.
If $j':Y'\to X'$ is the canonical immersion and
$\sh G'=\sh G\otimes_k k'$, then
$\sh F'=j'_*(\sh G')$.
For every extension $k'$ for which $X'$ is locally Noetherian, the
Module $\sh F'$ has no embedded associated prime cycle
\sref{IV.4.2.7}.
Moreover, every maximal point $y'$ of $Y'$ lies above a maximal point
$y$ of $Y$; since
$\sh G_y\simeq\sh O_{Y,y}$, one has
$\sh G'_{y'}\simeq\sh O_{Y',y'}$.
Consequently, by \sref{IV.3.2.3}, when $X'$ is locally Noetherian,
$\sh F'$ is reduced if and only if $Y'$ is reduced.

It now follows from conditions~\emph{(d)} and~\emph{(b)} of
\sref{IV.4.6.1} that \emph{(b)} implies \emph{(c)} and
\emph{(c)} implies \emph{(a)}.
In~\emph{(a)}, the finite-type hypothesis on $k'/k$ serves only to
ensure that $X'$ is locally Noetherian.
When $X$ is locally of finite type over $k$, the prescheme $X'$ is
locally Noetherian for every extension $k'$ of $k$; this proves the
equivalence of~\emph{(d)} with the other conditions in that case.
\end{proof}

\begin{definition}[4.6.17]
\phantomsection
\label{IV.4.6.17}
When $\sh F$ satisfies the equivalent conditions of
\sref{IV.4.6.16}, it is said to be \emph{geometrically reduced} over
$k$, or \emph{separable} over $k$.
\end{definition}

\begin{proposition}[4.6.18]
\phantomsection
\label{IV.4.6.18}
Let $k$ be a field, let $X$ be a locally Noetherian $k$-prescheme,
and let $\sh F$ be a coherent $\sh O_X$-Module.
The following conditions are equivalent:
\begin{enumerate}
\item[\rm(a)]
\phantomsection
\label{IV.4.6.18.a}
for every extension of finite type $k'$ of $k$, if
$X'=X\otimes_k k'$, then
$\sh F'=\sh F\otimes_k k'$ is an integral
$\sh O_{X'}$-Module \sref{IV.3.2.4};

\item[\rm(b)]
\phantomsection
\label{IV.4.6.18.b}
for every finite extension $k'$ of $k$, the Module $\sh F'$ is an
integral $\sh O_{X'}$-Module;

\item[\rm(c)]
\phantomsection
\label{IV.4.6.18.c}
$\sh F$ is reduced (or integral) and, if $\sh J$ is the Ideal of
$\sh O_X$ annihilating $\sh F$, the closed sub-prescheme defined by
$\sh J$ is geometrically integral over $k$.
\end{enumerate}
If, moreover, $X$ is locally of finite type over $k$, these conditions
are also equivalent to:
\begin{enumerate}
\item[\rm(d)]
\phantomsection
\label{IV.4.6.18.d}
for every extension $k'$ of $k$
(or for one algebraically closed extension $k'$ of $k$), the Module
$\sh F'$ is an integral $\sh O_{X'}$-Module.
\end{enumerate}
\end{proposition}

\begin{proof}
Conditions~\emph{(a)}, \emph{(b)}, and \emph{(c)} imply that $\sh F$
is integral.
Hence $\operatorname{Ass}(\sh F)$ consists of one point $x$, the
generic point of $\operatorname{Supp}(\sh F)$.
For every extension $k'$ of $k$ for which $X'$ is locally
Noetherian, the prime cycles associated to $\sh F'$ are therefore
nonembedded and are the maximal points of
$\operatorname{Supp}(\sh F')$, all lying above $x$.
By \sref{IV.4.6.16}, each of conditions~\emph{(a)}, \emph{(b)}, and
\emph{(c)} also implies that $\sh F$ is geometrically reduced.
It is immediate that \emph{(a)} implies \emph{(b)}, and
\emph{(b)} implies that $\operatorname{Supp}(\sh F)$ is geometrically
irreducible \sref{IV.4.5.9}.
Thus \emph{(b)} implies \emph{(c)}, because the closed sub-prescheme
$Y$ defined by $\sh J$ is geometrically reduced and geometrically
irreducible, and hence geometrically integral.
\oldpage[IV-2]{74}
Conversely, if $Y$ is geometrically integral, then for every extension
$k'$ of $k$, the support $\operatorname{Supp}(\sh F')$ is
irreducible.
Consequently, whenever $X'$ is locally Noetherian, $\sh F'$ is
integral.
This proves that \emph{(c)} implies \emph{(a)}, and also that
\emph{(c)} implies \emph{(d)} when $X$ is locally of finite type over
$k$.
\end{proof}

\begin{definition}[4.6.19]
\phantomsection
\label{IV.4.6.19}
When $\sh F$ satisfies the equivalent conditions of
\sref{IV.4.6.18}, it is said to be \emph{geometrically integral} over
$k$.
\end{definition}

\begin{proposition}[4.6.20]
\phantomsection
\label{IV.4.6.20}
Let $k$ be a field, let $X$ be a locally Noetherian $k$-prescheme,
let $\sh F$ be a coherent $\sh O_X$-Module, and let $K$ be an
extension of $k$ such that $X\otimes_kK$ is locally Noetherian.
If $\sh F$ is geometrically reduced
(resp.\ geometrically integral) over $k$, then
$\sh F\otimes_kK$ is geometrically reduced
(resp.\ geometrically integral) over $K$.
\end{proposition}

\begin{proof}
For every finite extension $K'$ of $K$,
\[
 (X\otimes_kK)\otimes_KK'=X\otimes_kK'
\]
is locally Noetherian.
The proofs of \sref{IV.4.6.16} and \sref{IV.4.6.18} show that the
hypothesis on $\sh F$ implies that $\sh F\otimes_kK'$ is reduced
(resp.\ integral), which proves the assertion.
\end{proof}

\begin{proposition}[4.6.21]
\phantomsection
\label{IV.4.6.21}
Let $k$ be a field, let $X$ and $Y$ be two locally Noetherian
$k$-preschemes such that $X\times_kY$ is locally Noetherian, let
$\sh F$ be a coherent $\sh O_X$-Module, and let $\sh G$ be a coherent
$\sh O_Y$-Module.
\begin{enumerate}
\item[\rm(i)]
\phantomsection
\label{IV.4.6.21.i}
If $\sh F$ is geometrically reduced
(resp.\ geometrically integral) and $\sh G$ is reduced
(resp.\ integral), then $\sh F\otimes_k\sh G$ is reduced
(resp.\ integral).

\item[\rm(ii)]
\phantomsection
\label{IV.4.6.21.ii}
If $\sh F$ and $\sh G$ are both geometrically reduced
(resp.\ geometrically integral), then so is
$\sh F\otimes_k\sh G$.
\end{enumerate}
\end{proposition}

\begin{proof}
\begin{enumerate}
\item[\rm(i)]
Using \sref{IV.3.2.3}, \sref{IV.4.6.16}, and
\sref{IV.4.6.18}, we may reduce to the case in which
\[
 \operatorname{Supp}(\sh F)=X,\qquad
 \operatorname{Supp}(\sh G)=Y,
\]
$X$ is geometrically reduced (resp.\ geometrically integral), $Y$ is
reduced (resp.\ integral), $\sh F$ and $\sh G$ have no embedded
associated prime cycles, and
\[
 \sh F_x\simeq\sh O_x,\qquad
 \sh G_y\simeq\sh O_y
\]
at every maximal point $x$ of $X$ and $y$ of $Y$.
Then $\sh F\otimes_k\sh G$ has no embedded associated prime cycle
\sref{IV.4.2.3}, and its associated prime cycles are exactly the
irreducible components of $X\times_kY$ \sref{IV.4.2.5}.
We may therefore reduce further to the case in which $X$ and $Y$ are
irreducible.
If $X$ is geometrically reduced and $Y$ is reduced, then
$X\times_kY$ is reduced \sref{IV.4.2.4.ii}.
If $X$ is geometrically integral and $Y$ is integral, then
$X\times_kY$ is moreover irreducible \sref{IV.4.5.8.i}, and hence
integral.
Finally, every maximal point $z$ of $X\times_kY$ lies above a maximal
point $x$ of $X$ and a maximal point $y$ of $Y$.
The hypotheses on $\sh F$ and $\sh G$ therefore imply, by
\sref[I]{I.9.1.12}, that
$(\sh F\otimes_k\sh G)_z\simeq\sh O_z$.

\item[\rm(ii)]
The argument is analogous.
After reducing to
$X=\operatorname{Supp}(\sh F)$ and
$Y=\operatorname{Supp}(\sh G)$, the prescheme $X\times_kY$ is
geometrically reduced (resp.\ geometrically integral) by
\sref{IV.4.6.5.ii}.
\end{enumerate}
\end{proof}

The definitions in \sref{IV.4.6.17} and \sref{IV.4.6.19} localize as
follows.

\begin{definition}[4.6.22]
\phantomsection
\label{IV.4.6.22}
Let $k$ be a field, let $X$ be a locally Noetherian $k$-prescheme,
and let $\sh F$ be a coherent $\sh O_X$-Module.
We say that $\sh F$ is \emph{geometrically reduced}, or
\emph{separable} (resp.\ \emph{geometrically punctually integral}),
at a point $x\in X$ if, for every finite radicial
(resp.\ finite) extension $k'$ of $k$, the Module
$\sh F\otimes_k k'$ is reduced (resp.\ integral) at every point $x'$
of $X\otimes_k k'$ above $x$; see \sref{IV.3.2.2} and
\sref{IV.3.2.4}.
\end{definition}

\oldpage[IV-2]{75}
If $\sh F$ is reduced at $x$, there is an open neighbourhood $U$ of
$x$ such that $\sh F|U$ is reduced \sref{IV.3.2.2}.
Arguing as in \sref{IV.4.6.16}, one sees that $\sh F$ is
geometrically reduced at $x$ if and only if $\sh F$ is reduced at $x$
and the closed sub-prescheme $Y$ defined by the annihilator $\sh J$ of
$\sh F$ is geometrically reduced at $x$.
It follows that $x$ has an open neighbourhood $U$ such that
$\sh F|U$ is geometrically reduced \sref{IV.4.6.13}.
The arguments of \sref{IV.4.6.16} and \sref{IV.4.6.18} likewise show
that, if $X$ is locally of finite type over $k$, then $\sh F$ is
geometrically punctually integral at $x$ if and only if $\sh F$ is
reduced at $x$ and $Y$ is geometrically punctually integral at $x$.
% END NATIVE INPUT components/06_section46_geometrically_reduced_preschemes
% BEGIN NATIVE INPUT components/07_section47_multiplicities_primary_decomposition
\subsection{Multiplicities in the primary decomposition on an algebraic prescheme}
\phantomsection
\label{subsection:IV.4.7}

\begin{lemma}[4.7.1]
\phantomsection
\label{IV.4.7.1}
Let $A$ and $B$ be two local rings, let $\mathfrak m$ be the maximal
ideal of $A$, and let $\rho:A\to B$ be a local homomorphism.
Suppose that $B$ is flat as an $A$-module and that $B/\mathfrak mB$
has finite length as a $B$-module.
Then an $A$-module $M$ has finite length if and only if $M_{(B)}$ has
finite length as a $B$-module, and
\begin{equation}
\operatorname{length}_B(M_{(B)})
=\operatorname{length}_A(M)\,
 \operatorname{length}_B(B/\mathfrak mB).
\tag{4.7.1.1}
\label{IV.4.7.1.1}
\end{equation}
\end{lemma}

\begin{proof}
This is a special case of \sref{IV.2.5.5.2}.
\end{proof}

\begin{corollary}[4.7.2]
\phantomsection
\label{IV.4.7.2}
Let $A$ and $B$ be two Noetherian rings, and let
$\rho:A\to B$ be a ring homomorphism such that $B$ is flat as an
$A$-module.
Let $\mathfrak q$ be a minimal prime ideal of $B$, put
$\mathfrak p=\rho^{-1}(\mathfrak q)$, and let $M$ be an $A$-module
of finite type.
Then $\mathfrak p$ is a minimal prime ideal of $A$,
$M_{\mathfrak p}$ has finite length as an $A_{\mathfrak p}$-module,
and
\begin{equation}
\operatorname{length}_{B_{\mathfrak q}}
 \bigl((M_{(B)})_{\mathfrak q}\bigr)
=
\operatorname{length}_{A_{\mathfrak p}}(M_{\mathfrak p})\,
\operatorname{length}_{B_{\mathfrak q}}
 \bigl(B_{\mathfrak q}/\mathfrak pB_{\mathfrak q}\bigr).
\tag{4.7.2.1}
\label{IV.4.7.2.1}
\end{equation}
\end{corollary}

\begin{proof}
The ring $B_{\mathfrak q}$ is flat over $A_{\mathfrak p}$
$(\mathbf 0_I,\,6.3.2)$, and the homomorphism
$A_{\mathfrak p}\to B_{\mathfrak q}$ is local; hence
$B_{\mathfrak q}$ is faithfully flat over $A_{\mathfrak p}$
$(\mathbf 0_I,\,6.6.2)$.
It follows from \sref{IV.2.3.5} that $\mathfrak p$ is minimal in $A$.
The ring $A_{\mathfrak p}$ is therefore Artinian
(Bourbaki, \emph{Alg.\ comm.}, chap.~IV, \S~2,
n\textsuperscript{o}~5, prop.~9), so $M_{\mathfrak p}$ has finite
length over $A_{\mathfrak p}$.
Formula \sref{IV.4.7.2.1} is now a special case of
\sref{IV.4.7.1.1}.
\end{proof}

\begin{proposition}[4.7.3]
\phantomsection
\label{IV.4.7.3}
Let $k$ be a field of characteristic exponent $p$, let $X$ be an
integral locally Noetherian prescheme, and let
$K=\mathrm R(X)$ be the field of rational functions on $X$.
Let $k'$ be an extension of $k$, let $(X'_\alpha)$ be the family of
irreducible components of $X'=X_{(k')}$, and let $x'_\alpha$ be the
generic point of $X'_\alpha$.
\begin{enumerate}
\item[\rm(i)]
\phantomsection
\label{IV.4.7.3.i}
Suppose that $X'$ is locally Noetherian, as is the case if $X$ is
locally of finite type over $k$ or if $k'$ is an extension of finite
type of $k$.
Let $k''$ be any extension of $k'$, set
$X''=X\otimes_k k''$, and let $(X''_\beta)$ be the family of
irreducible components of $X''$, with generic points $x''_\beta$.
If $x''_\beta$ lies above $x'_\alpha$, set
\[
 Z''=(X'_{\mathrm{red}})\otimes_{k'}k'',
\]
Then
\begin{equation}
\operatorname{length}(\sh O_{X'',x''_\beta})
=
\operatorname{length}(\sh O_{X',x'_\alpha})\,
\operatorname{length}(\sh O_{Z'',x''_\beta}).
\tag{4.7.3.1}
\label{IV.4.7.3.1}
\end{equation}
In particular, if $k''$ is a separable extension of $k'$ such that
$Z''$ is locally Noetherian, then
\[
\operatorname{length}(\sh O_{X'',x''_\beta})
=\operatorname{length}(\sh O_{X',x'_\alpha}).
\]

\oldpage[IV-2]{76}
\item[\rm(ii)]
\phantomsection
\label{IV.4.7.3.ii}
If $X$ is locally of finite type over $k$, or if the extension $k'/k$
is of finite type, then each number
\[
\operatorname{length}(\sh O_{X',x'_\alpha})
\]
is a power of $p$.

\item[\rm(iii)]
\phantomsection
\label{IV.4.7.3.iii}
Suppose that $X$ is locally of finite type over $k$.
If $k'$ is perfect, all the numbers
$\operatorname{length}(\sh O_{X',x'_\alpha})$ are equal and depend
only on the extension $K/k$, not on the particular perfect extension
$k'/k$ under consideration.
Moreover, there exists a finite radicial extension $k_1$ of $k$ such
that, if $x_1$ is the generic point of
$X_1=X\otimes_k k_1$, which is irreducible, then
\[
\operatorname{length}(\sh O_{X',x'_\alpha})
=\operatorname{length}(\sh O_{X_1,x_1})
\quad\text{for every }\alpha.
\]
\end{enumerate}
\end{proposition}

\begin{proof}
\begin{enumerate}
\item[\rm(i)]
Recall from \sref{IV.4.2.4} that the finitely many irreducible
components $X'_\alpha$ correspond bijectively to the minimal prime
ideals $\mathfrak p_\alpha$ of $K\otimes_k k'$, and that
$\sh O_{X',x'_\alpha}$ is an Artinian ring isomorphic to
$(K\otimes_k k')_{\mathfrak p_\alpha}$.
These rings therefore depend only on the extensions $K/k$ and $k'/k$;
in particular, with $k'$ fixed, they are the same for every integral
$k$-prescheme having $K$ as its field of rational functions.

Formula \sref{IV.4.7.3.1} follows from \sref{IV.4.7.2.1}.
Indeed, apply that formula to the Noetherian ring $A$ of an affine
neighbourhood of $x'_\alpha$ in $X'$ and to the ring $B$ of a
sufficiently small affine neighbourhood of $x''_\beta$ in $X''$,
where $\mathfrak p$ and $\mathfrak q$ are the minimal prime ideals
corresponding respectively to $x'_\alpha$ and $x''_\beta$, and take
$M=A$.
Here $A/\mathfrak p$ is the ring of an affine neighbourhood of
$x'_\alpha$ in $X'_{\mathrm{red}}$, while
$B=A\otimes_{k'}k''$.
If $\mathfrak q'=\mathfrak q/\mathfrak pB$, a minimal prime ideal of
$B/\mathfrak pB$, then
\[
\sh O_{Z'',x''_\beta}
=(B/\mathfrak pB)_{\mathfrak q'}
=B_{\mathfrak q}/\mathfrak pB_{\mathfrak q},
\]
which proves \sref{IV.4.7.3.1}.
The last assertion of \emph{(i)} follows because $Z''$ is reduced in
this case \sref{IV.4.2.4}, and hence
$\mathfrak pB_{\mathfrak q}=\mathfrak qB_{\mathfrak q}$.

\item[\rm(ii)]
If $k'$ is of finite type over $k$, there is a purely transcendental
extension $k'_0$ of $k$ such that $k'$ is finite over $k'_0$.
Since $k'_0/k$ is separable, \emph{(i)}, applied with $k'$ and $k''$
replaced by $k$ and $k'_0$, shows that we may replace $X$ by
$X\otimes_k k'_0$; thus we may assume that $k'/k$ is finite.
There is then a finite quasi-Galois extension $k''/k$ containing
$k'$, and \sref{IV.4.7.3.1} shows that it suffices to prove the
assertion for $k''$.
But $k''$ is separable over a finite radicial extension $k_1$ of $k$,
so \emph{(i)} reduces us to the case $k'=k_1$ radicial.
The ring $K\otimes_k k_1$ is then Artinian and has a unique minimal
prime ideal \sref{IV.4.3.2}.
If $\overline K$ is an algebraic closure of $K$, the length of
$K\otimes_k k_1$ divides the length of the Artinian ring
$\overline K\otimes_k k_1$ by \sref{IV.4.7.1.1}.
The Artinian $\overline K$-algebra
$\overline K\otimes_k k_1$ has a unique prime ideal, necessarily
maximal, and its quotient by that ideal is a finite extension of
$\overline K$, hence is $\overline K$.
Its length is therefore its rank over $\overline K$, namely
$[k_1:k]$, a power of $p$.

Suppose next that $X$ is locally of finite type over $k$.
Then $X''$ is locally Noetherian for every extension $k''$ of $k'$.
Taking $k''$ algebraically closed, \sref{IV.4.7.3.1} reduces us to
the case in which $k'$ is algebraically closed.
Since $k'$ is then separable over the algebraic closure $\overline k$
of $k$, another application of \emph{(i)} reduces us to
$k'=\overline k$.

\oldpage[IV-2]{77}

After replacing $X$ by an affine open of finite type over $k$, there
is a finite radicial extension $k_1/k$ such that
$(X\otimes_k k_1)_{\mathrm{red}}$ is separable over $k_1$
\sref{IV.4.6.6}.
Consequently,
$(X\otimes_k k_1)_{\mathrm{red}}\otimes_{k_1}\overline k$ is reduced
\sref{IV.4.6.4}.
By \sref{IV.4.7.3.1}, this again reduces us to the case in which
$k'$ is a finite radicial extension of $k$, and the preceding
argument applies.

\item[\rm(iii)]
Any two extensions $k'$ and $k''$ of $k$ embed in a common
algebraically closed extension of $k$.
To show that the set of numbers
$\operatorname{length}(\sh O_{X',x'_\alpha})$ for perfect $k'$ does
not depend on $k'$, it is therefore enough to compare them with those
obtained after an algebraically closed extension $k''/k'$; the
assertion then follows from \emph{(i)}, since $k''/k'$ is separable.
To prove that all these numbers are equal, we may thus assume
$k'=\overline k$.
For the finite radicial extension $k_1/k$ chosen in \emph{(ii)}, all
the points $x'_\alpha$ lie above the unique generic point $x_1$ of
$X_1=X\otimes_k k_1$.
The equality
\[
\operatorname{length}(\sh O_{X',x'_\alpha})
=\operatorname{length}(\sh O_{X_1,x_1})
\]
for every $\alpha$ follows once more from \sref{IV.4.7.3.1} and the
choice of $k_1$.
\end{enumerate}
\end{proof}

\begin{definition}[4.7.4]
\phantomsection
\label{IV.4.7.4}
Let $k$ be a field, let $X$ be an integral $k$-prescheme, and let $K$
be the field of rational functions on $X$.
The \emph{radicial multiplicity} of $K$ (or of $X$) over $k$ is the
supremum of the lengths of the Artinian rings $K\otimes_k k_1$, as
$k_1$ ranges over the finite radicial extensions of $k$.
The \emph{separable multiplicity} of $K$ (or of $X$) over $k$ is the
geometric number of irreducible components of $X$ if this number is
finite, and $+\infty$ otherwise.
Finally, the \emph{total multiplicity} of $K$ (or of $X$) over $k$ is
the product of its radicial and separable multiplicities.
\end{definition}

If $p$ is the characteristic exponent of $k$ and the radicial
multiplicity of $K$ over $k$ is finite, then it is a power
$q=p^f$ of $p$, by \sref{IV.4.7.3.ii}.
When $p\ne1$, the well-defined integer $f$ is called the
\emph{inseparability exponent} of $K$ over $k$; when $p=1$, the
radicial multiplicity is always $1$.

To say that the radicial (resp.\ separable, resp.\ total)
multiplicity of $X$ over $k$ is $1$ means that $X$ is geometrically
reduced (resp.\ geometrically irreducible, resp.\ geometrically
integral) over $k$.

When $X$ is locally of finite type, its radicial multiplicity over
$k$ is also the common length of the local rings $A_{\mathfrak q}$ at
the minimal prime ideals of the total ring of fractions $A$ of
$K\otimes_k k'$, for every perfect extension $k'/k$.
Its separable multiplicity is the number of minimal prime ideals of
$A=K\otimes_k\Omega$, for every algebraically closed extension
$\Omega/k$.
Finally, its total multiplicity is the sum of the lengths of the local
rings $A_{\mathfrak q}$ of $A=K\otimes_k\Omega$ at its minimal prime
ideals.
In this case all three numbers are finite.

\begin{definition}[4.7.5]
\phantomsection
\label{IV.4.7.5}
Let $k$ be a field, let $X$ be a $k$-prescheme locally of finite type
over $k$, let $\sh F$ be a coherent $\sh O_X$-Module, and let $x\in X$
be such that $\sh F_x$ has finite length as an $\sh O_{X,x}$-module.
The \emph{geometric length} of $\sh F$ at $x$ relative to $k$, or the
\emph{radicial multiplicity of $x$ for $\sh F$}, is
\[
\lambda_x(\sh F)
=\operatorname{length}_{\sh O_{X,x}}(\sh F_x)\,
 \mu_r(\kappa(x)\mid k),
\]
where $\mu_r(\kappa(x)\mid k)$ denotes the radicial multiplicity of
$\kappa(x)$ over $k$.
\end{definition}

When $X$ is integral and $x$ is its generic point,
$\operatorname{length}(\sh O_{X,x})=1$.
Thus the radicial multiplicity of $x$ for $\sh O_X$ is precisely the
radicial multiplicity of $X$ over $k$ defined in
\sref{IV.4.7.4}.

\oldpage[IV-2]{78}

\begin{proposition}[4.7.6]
\phantomsection
\label{IV.4.7.6}
Let
\[
0\longrightarrow\sh F'\longrightarrow\sh F
\longrightarrow\sh F''\longrightarrow0
\]
be an exact sequence of coherent $\sh O_X$-Modules.
At a point $x\in X$, the geometric length $\lambda_x(\sh F)$ is
defined if and only if both $\lambda_x(\sh F')$ and
$\lambda_x(\sh F'')$ are defined; in that case
\begin{equation}
\lambda_x(\sh F)=\lambda_x(\sh F')+\lambda_x(\sh F'').
\tag{4.7.6.1}
\label{IV.4.7.6.1}
\end{equation}
\end{proposition}

\begin{proof}
This follows immediately from the definition: the length of
$\sh F_x$ is finite if and only if the lengths of $\sh F'_x$ and
$\sh F''_x$ are finite, and then
\[
\operatorname{length}(\sh F_x)
=\operatorname{length}(\sh F'_x)+
 \operatorname{length}(\sh F''_x).
\]
\end{proof}

\begin{env}[4.7.7]
\phantomsection
\label{IV.4.7.7}
Under the hypotheses of \sref{IV.4.7.5}, the only points $x\in X$ for
which $\sh F_x$ is a nonzero $\sh O_{X,x}$-module of finite length
are, by \sref{IV.3.1.2}, the maximal points of
$\operatorname{Supp}(\sh F)$.
This follows from Bourbaki, \emph{Alg.\ comm.}, chap.~IV, \S~2,
n\textsuperscript{o}~5, cor.~2 to prop.~7, since the question is local
on $X$.
There is a closed sub-prescheme $Y$ with underlying space
$\operatorname{Supp}(\sh F)$ and a coherent $\sh O_Y$-Module $\sh G$
such that $\sh F=j_*(\sh G)$, where $j:Y\to X$ is the canonical
injection.
If $x$ is a maximal point of $Y$, then the geometric lengths of
$\sh F$ and $\sh G$ at $x$ are equal, because
$\sh O_{X,x}$ and $\sh O_{Y,x}$ have the same residue field.
This gives the following characterization of geometric length.
\end{env}

\begin{proposition}[4.7.8]
\phantomsection
\label{IV.4.7.8}
Under the hypotheses of \sref{IV.4.7.5}, let $x$ be a maximal point
of $\operatorname{Supp}(\sh F)$.
Let $k'$ be a perfect extension of $k$, and set
$X'=X\otimes_k k'$ and $\sh F'=\sh F\otimes_k k'$.
Then $\lambda_x(\sh F)$ equals
$\operatorname{length}_{\sh O_{X',x'}}(\sh F'_{x'})$ at every
maximal point $x'$ of $\operatorname{Supp}(\sh F')$ above $x$.
Moreover, there exists a finite radicial extension $k_1$ of $k$ such
that, on setting $X_1=X\otimes_k k_1$ and
$\sh F_1=\sh F\otimes_k k_1$, the number $\lambda_x(\sh F)$ equals
$\operatorname{length}_{\sh O_{X_1,x_1}}((\sh F_1)_{x_1})$ at every
maximal point $x_1$ of $\operatorname{Supp}(\sh F_1)$ above $x$.
\end{proposition}

\begin{proof}
The remark in \sref{IV.4.7.7} reduces us to the case in which $x$ is
a maximal point of $X$.
Since the question is local on $X$, suppose
$X=\Spec(A)$ is affine Noetherian and $\sh F=\widetilde M$, where
$M$ is an $A$-module of finite type.
Put $B=A\otimes_k k'$, and let $\mathfrak p$ and $\mathfrak q$ be the
minimal prime ideals of $A$ and $B$ corresponding respectively to
$x$ and $x'$.
Formula \sref{IV.4.7.2.1} gives
\[
\operatorname{length}_{B_{\mathfrak q}}
 \bigl((M_{(B)})_{\mathfrak q}\bigr)
=
\operatorname{length}_{A_{\mathfrak p}}(M_{\mathfrak p})\,
\operatorname{length}_{B_{\mathfrak q}}
 \bigl(B_{\mathfrak q}/\mathfrak pB_{\mathfrak q}\bigr).
\]
Since
$\kappa(x)=A_{\mathfrak p}/\mathfrak pA_{\mathfrak p}$, the argument
of \sref{IV.4.7.3} shows that the last factor is the length of
$(B/\mathfrak pB)_{\mathfrak q'}$, where
$\mathfrak q'=\mathfrak q/\mathfrak pB$.
This is a local ring of
$(A/\mathfrak p)\otimes_k k'$ at a minimal prime ideal, and therefore
also a local ring of the total ring of fractions of
$\kappa(x)\otimes_k k'$.
By definition its length is $\mu_r(\kappa(x)\mid k)$, proving the
first assertion.
The second follows in the same way by replacing $k'$ with $k_1$ and
using \sref{IV.4.7.3.iii}.
\end{proof}

When $x$ is a maximal point of $\operatorname{Supp}(\sh F)$, the
geometric length of $\sh F$ at $x$ is also called the
\emph{radicial multiplicity of the maximal prime cycle}
$\overline{\{x\}}$ associated with $\sh F$.

\begin{corollary}[4.7.9]
\phantomsection
\label{IV.4.7.9}
Under the hypotheses of \sref{IV.4.7.5}, let $k'$ be an extension of
$k$, set $X'=X\otimes_k k'$ and $\sh F'=\sh F\otimes_k k'$, and let
$Z$ be an irreducible component of $\operatorname{Supp}(\sh F)$.
If $Z'$ is an irreducible component of
$\operatorname{Supp}(\sh F')$ dominating $Z$
\sref{IV.4.2.7}, then the radicial multiplicities of $Z$ with respect
to $\sh F$ and of $Z'$ with respect to $\sh F'$ are equal.
\end{corollary}

\oldpage[IV-2]{79}

\begin{proof}
Choose a perfect extension $k''$ of $k'$, set
$X''=X\otimes_k k''$ and $\sh F''=\sh F\otimes_k k''$, and choose a
maximal point $x''$ of $\operatorname{Supp}(\sh F'')$ above the
generic point of $Z'$.
It follows from \sref{IV.4.7.8} that
$\operatorname{length}(\sh F''_{x''})$ equals both multiplicities.
\end{proof}

\begin{proposition}[4.7.10]
\phantomsection
\label{IV.4.7.10}
Let $k$ be a field, let $X$ be a $k$-prescheme locally of finite type
over $k$, let $\sh F$ be a coherent $\sh O_X$-Module, let $x\in X$,
and let $\xi_i$ be the generic points of the irreducible components
of $\operatorname{Supp}(\sh F)$ that contain $x$.
Then $\sh F$ is geometrically reduced at $x$
\sref{IV.4.6.22} if and only if $x$ belongs to no embedded associated
prime cycle of $\sh F$ and
$\lambda_{\xi_i}(\sh F)=1$ for every $i$.
\end{proposition}

\begin{proof}
First suppose the stated conditions hold.
Let $k'/k$ be finite, let $x'$ be a point of
$X'=X\otimes_k k'$ above $x$, and put
$\sh F'=\sh F\otimes_k k'$.
By \sref{IV.4.2.7}, $x'$ belongs to no embedded associated prime cycle
of $\sh F'$.
Moreover, \sref{IV.4.7.8} shows that $\sh F'$ has geometric length
$1$ at every generic point $\xi'$ of an irreducible component of
$\operatorname{Supp}(\sh F')$ containing $x'$; by
\sref{IV.2.3.4}, such a point lies above one of the $\xi_i$.
A fortiori,
$\operatorname{length}(\sh F'_{\xi'})=1$, so $\sh F'$ is reduced at
$x'$.

Conversely, \sref{IV.4.2.7} first shows that if, for one finite
extension $k'/k$, the Module $\sh F'$ is reduced at every point $x'$
above $x$, then $x$ belongs to no embedded associated prime cycle of
$\sh F$.
For each $\xi_i$, take for $k'$ the field $k_1$ supplied by
\sref{IV.4.7.8}.
The hypothesis that $\sh F$ is geometrically reduced at $x$ then
implies $\lambda_{\xi_i}(\sh F)=1$.
\end{proof}

\begin{proposition}[4.7.11]
\phantomsection
\label{IV.4.7.11}
Let $k$ be a field, let $X$ be a prescheme locally of finite type
over $k$, let $\sh F$ be a coherent $\sh O_X$-Module, let $k'$ be an
extension of $k$, set $X'=X\otimes_k k'$ and
$\sh F'=\sh F\otimes_k k'$, and let $x'\in X'$ lie above $x\in X$.
Then $\sh F$ is geometrically reduced
(resp.\ geometrically punctually integral) at $x$ if and only if
$\sh F'$ is geometrically reduced
(resp.\ geometrically punctually integral) at $x'$.
\end{proposition}

\begin{proof}
Every irreducible component of $\operatorname{Supp}(\sh F')$
containing $x'$ dominates an irreducible component of
$\operatorname{Supp}(\sh F)$ containing $x$, and conversely every
such component downstairs is dominated by one containing $x'$
(\sref{IV.4.2.7.i} and \sref{IV.2.3.4}).
The assertion concerning geometric reducedness therefore follows
from \sref{IV.4.2.7.ii}, \sref{IV.4.7.9}, and
\sref{IV.4.7.10}.

Let $Y$ be the closed sub-prescheme of $X$ defined by the annihilator
Ideal $\sh J$ of $\sh F$.
Then $Y'=Y\otimes_k k'$ is the closed sub-prescheme of $X'$ defined by
the annihilator Ideal
$\sh J'=\sh J\otimes_k k'$ of $\sh F'$.
To say that $\sh F$ (resp.\ $\sh F'$) is geometrically punctually
integral at $x$ (resp.\ $x'$) is to say that $\sh F$ (resp.\
$\sh F'$) is geometrically reduced at $x$ (resp.\ $x'$) and that
$Y$ (resp.\ $Y'$) is geometrically punctually integral at $x$
(resp.\ $x'$).
By the assertion just proved, either of the two geometric
pointwise-integrality hypotheses implies that $Y$ has a separable
neighbourhood of $x$.
The conclusion follows from \sref{IV.4.6.10}.
\end{proof}

\begin{definition}[4.7.12]
\phantomsection
\label{IV.4.7.12}
Let $k$ be a field, let $X$ be a $k$-prescheme locally of finite type
over $k$, let $\sh F$ be a coherent $\sh O_X$-Module, and let $x$ be
a maximal point of $\operatorname{Supp}(\sh F)$.
The

\oldpage[IV-2]{80}

\emph{total multiplicity} of $x$ (or of the prime cycle
$\overline{\{x\}}$) for $\sh F$, relative to $k$, is the product of
the radicial multiplicity of $x$ for $\sh F$ and the separable
multiplicity of $\kappa(x)$ over $k$.
\end{definition}

Equivalently, the total multiplicity of $x$ for $\sh F$ is the
product of
$\operatorname{length}_{\sh O_{X,x}}(\sh F_x)$ and the total
multiplicity of $\kappa(x)$ over $k$ \sref{IV.4.7.4}.
With the notation of \sref{IV.4.7.11}, the total multiplicity of $x$
for $\sh F$ is the sum of the total multiplicities, for $\sh F'$, of
the maximal points $x'_i$ of $\operatorname{Supp}(\sh F')$ above $x$.
Moreover, there exists a finite extension $k'/k$ for which the total
multiplicity of $x$ for $\sh F$ is
\[
\sum_i\operatorname{length}(\sh F'_{x'_i}).
\]

\begin{proposition}[4.7.13]
\phantomsection
\label{IV.4.7.13}
Under the hypotheses and notation of \sref{IV.4.7.10}, it is
sufficient, for $\sh F$ to be geometrically punctually integral at
$x$, that $x$ belong to only one associated prime cycle of $\sh F$
(necessarily non-embedded), and that the total multiplicity for
$\sh F$ of the generic point $\xi$ of that cycle be $1$.
\end{proposition}

\begin{proof}
Let $k'/k$ be finite, let $x'$ be a point of
$X'=X\otimes_k k'$ above $x$, and put
$\sh F'=\sh F\otimes_k k'$.
The point $x'$ can belong to only one associated prime cycle of
$\sh F'$, since by hypothesis there is only one maximal point of
$\operatorname{Supp}(\sh F')$ above $\xi$.
Moreover, $\sh F'$ is geometrically reduced at $x'$ by
\sref{IV.4.7.10}.
The conclusion follows.
\end{proof}
% END NATIVE INPUT components/07_section47_multiplicities_primary_decomposition
% BEGIN NATIVE INPUT components/08_section48_fields_of_definition
\subsection{Fields of definition}
\phantomsection
\label{subsection:IV.4.8}

\begin{env}[4.8.1]
\phantomsection
\label{IV.4.8.1}
Given a prescheme $X$, throughout this section we denote by
$\mathfrak S(X)$ the set of sub-preschemes of $X$, and, as usual, by
$\mathfrak P(X)$ the set of subsets of the underlying space of $X$.
For every quasi-coherent $\sh O_X$-Module $\sh F$, we denote by
$\Phi(\sh F)$ the set of quasi-coherent sub-$\sh O_X$-Modules of
$\sh F$.
\end{env}

\begin{env}[4.8.2]
\phantomsection
\label{IV.4.8.2}
Let $k$ be a field, let $X$ and $Y$ be two $k$-preschemes, let
$\sh F$ and $\sh G$ be two quasi-coherent $\sh O_X$-Modules, and let
$K'\subset K$ be two extensions of $k$.
Corresponding to the morphism $\Spec(K)\to\Spec(K')$, there are
canonical maps
\begin{align}
\Phi(\sh F\otimes_kK')
  &\longrightarrow \Phi(\sh F\otimes_kK),
\tag{4.8.2.1}
\label{IV.4.8.2.1}\\
\operatorname{Hom}(\sh F\otimes_kK',\sh G\otimes_kK')
  &\longrightarrow
   \operatorname{Hom}(\sh F\otimes_kK,\sh G\otimes_kK),
\tag{4.8.2.2}
\label{IV.4.8.2.2}\\
\mathfrak S(X_{(K')})
  &\longrightarrow \mathfrak S(X_{(K)}),
\tag{4.8.2.3}
\label{IV.4.8.2.3}\\
\operatorname{Hom}_{K'}(X_{(K')},Y_{(K')})
  &\longrightarrow \operatorname{Hom}_{K}(X_{(K)},Y_{(K)}),
\tag{4.8.2.4}
\label{IV.4.8.2.4}\\
\mathfrak P(X_{(K')})
  &\longrightarrow \mathfrak P(X_{(K)}).
\tag{4.8.2.5}
\label{IV.4.8.2.5}
\end{align}
If $p_X:X_{(K)}\to X_{(K')}$ is the canonical projection, then
\sref{IV.4.8.2.5} sends $M$ to $p_X^{-1}(M)$, and similarly
\sref{IV.4.8.2.3} sends $Z$ to
$p_X^{-1}(Z)=Z\otimes_{K'}K$ \sref[I]{I.4.4.1}.
Map \sref{IV.4.8.2.1} sends
\[
\sh H\longmapsto
p_X^*(\sh H)
=
\sh H\otimes_{\sh O_{X_{(K')}}}\sh O_{X_{(K)}};
\]
\sref{IV.4.8.2.2} sends $u$ to $p_X^*(u)$; and
\sref{IV.4.8.2.4} sends $f$ to $f_{(K)}$.

If $\varepsilon_{K,K'}$ denotes any one of these canonical maps, and if
$K''\subset K'\subset K$ are three extensions of $k$, then
\begin{equation}
\varepsilon_{K,K''}=\varepsilon_{K,K'}\circ\varepsilon_{K',K''}.
\tag{4.8.2.6}
\label{IV.4.8.2.6}
\end{equation}
\end{env}

\oldpage[IV-2]{81}

\begin{proposition}[4.8.3]
\phantomsection
\label{IV.4.8.3}
The canonical maps \sref{IV.4.8.2.1}--\sref{IV.4.8.2.5} are
injective.
\end{proposition}

\begin{proof}
The morphism $p_X$ is faithfully flat and quasi-compact.
The injectivity of \sref{IV.4.8.2.5} follows from the surjectivity of
$p_X$.
The injectivity of \sref{IV.4.8.2.1} follows from
\sref{IV.2.2.2}, that of \sref{IV.4.8.2.2} from
\sref{IV.2.2.7}, that of \sref{IV.4.8.2.3} from
\sref{IV.2.2.15}, and that of \sref{IV.4.8.2.4} from
\sref{IV.2.2.16}.
\end{proof}

\begin{definition}[4.8.4]
\phantomsection
\label{IV.4.8.4}
With the notation of \sref{IV.4.8.2}, a quasi-coherent
sub-$\sh O_{X_{(K)}}$-Module of $\sh F\otimes_kK$
(resp.\ a homomorphism
$\sh F\otimes_kK\to\sh G\otimes_kK$, a sub-prescheme of $X_{(K)}$,
a $K$-morphism $X_{(K)}\to Y_{(K)}$, or a subset of $X_{(K)}$)
is said to be \emph{defined over $K'$} if it belongs to the image of
\sref{IV.4.8.2.1}
(resp.\ \sref{IV.4.8.2.2}, \sref{IV.4.8.2.3},
\sref{IV.4.8.2.4}, or \sref{IV.4.8.2.5}).
The field $K'$ is then called a \emph{field of definition} of the
object in question.
\end{definition}

Clearly, $K$ itself is a field of definition for every object in any
of the five target sets in \sref{IV.4.8.2}.
One should note, however, that for a $K$-prescheme $Z$, to say that a
quasi-coherent $\sh O_Z$-Module $\sh H$ is ``defined over a subfield
$K'$ of $K$'' has no meaning unless $Z$ has been \emph{given} in the
form $X_{(K)}$ for a prescheme $X$ over a subfield $k$ of $K$, and
unless $\sh H$ is a sub-$\sh O_Z$-Module of an $\sh O_Z$-Module that
has been \emph{given} in the form $\sh F\otimes_kK$, where $\sh F$ is
a quasi-coherent $\sh O_X$-Module.
Analogous remarks apply to the other notions.

It follows from \sref{IV.4.8.2.6} that if an object in one of the
five target sets in \sref{IV.4.8.2} is defined over $K'$, then it is
also defined over every intermediate field $K''$ with
$K'\subset K''\subset K$.
Finally, if $K_1/K$ is an extension, then an object in a source set
of $\varepsilon_{K_1,K}$ is defined over $K'$ if and only if its image under
$\varepsilon_{K_1,K}$ is defined over $K'$; this follows from
\[
\varepsilon_{K_1,K'}=\varepsilon_{K_1,K}\circ\varepsilon_{K,K'}.
\]

\begin{proposition}[4.8.5]
\phantomsection
\label{IV.4.8.5}
With the notation of \sref{IV.4.8.2}, let $(X_\alpha)$ be an open
covering of $X$.
An element
\[
\sh H\in\Phi(\sh F\otimes_kK)
\quad
\text{(resp. }u\in\operatorname{Hom}(\sh F\otimes_kK,\sh G\otimes_kK),
\ Z\in\mathfrak S(X_{(K)}),\
\text{or }M\in\mathfrak P(X_{(K)})\text{)}
\]
is defined over $K'$ if and only if, for every $\alpha$, its
restriction
\[
\sh H|_{(X_\alpha)_{(K)}}
\quad
\text{(resp. }u|_{(X_\alpha)_{(K)}},\
 Z\cap(X_\alpha)_{(K)},\
\text{or }M\cap(X_\alpha)_{(K)}\text{)}
\]
is defined over $K'$.
\end{proposition}

\begin{proof}
This follows immediately from the injectivity of
\sref{IV.4.8.2.1}, \sref{IV.4.8.2.2}, \sref{IV.4.8.2.3}, and
\sref{IV.4.8.2.5}.
For example, suppose that for every $\alpha$ there is a
quasi-coherent sub-$\sh O_{(X_\alpha)_{(K')}}$-Module
$\sh H'_\alpha$ such that
\[
\sh H|_{(X_\alpha)_{(K)}}=\sh H'_\alpha\otimes_{K'}K.
\]
For any $\alpha,\beta$, injectivity after restriction to
$(X_\alpha\cap X_\beta)_{(K')}$ gives
\[
\sh H'_\alpha|_{(X_\alpha\cap X_\beta)_{(K')}}
=
\sh H'_\beta|_{(X_\alpha\cap X_\beta)_{(K')}}.
\]
The $\sh H'_\alpha$ therefore glue to a quasi-coherent
sub-$\sh O_{X_{(K')}}$-Module $\sh H'$ of
$\sh F\otimes_kK'$ such that $\sh H=\sh H'\otimes_{K'}K$.
The other cases are treated in the same way.
\end{proof}

\begin{lemma}[4.8.6]
\phantomsection
\label{IV.4.8.6}
Let $k$ be a field, let $K/k$ be an extension, let $A$ be a
$k$-algebra, let $M$ be an $A$-module, let $\overline N$ be an
$A_{(K)}$-submodule of $M_{(K)}$, and let $K'$ be an intermediate
field of $K/k$.
The following conditions are equivalent:
\begin{enumerate}
\item[\rm(a)]
\phantomsection
\label{IV.4.8.6.a}
$\overline N=N'\otimes_{K'}K$, where $N'$ is an
$A_{(K')}$-submodule of $M_{(K')}$;

\item[\rm(b)]
\phantomsection
\label{IV.4.8.6.b}
if $X=\Spec(A)$, the quasi-coherent
$\sh O_{X_{(K)}}$-Module $\widetilde{\overline N}$ on
$X_{(K)}=\Spec(A_{(K)})$ is defined over $K'$;

\item[\rm(c)]
\phantomsection
\label{IV.4.8.6.c}
$\overline N=N'\otimes_{K'}K$, where $N'$ is a $K'$-vector
subspace of $M_{(K')}$.
\end{enumerate}
\end{lemma}

\begin{proof}
\oldpage[IV-2]{82}
The equivalence of \emph{(a)} and \emph{(b)} follows immediately from
\sref[I]{I.1.6.5}, and \emph{(a)} plainly implies \emph{(c)}.
Conversely, if \emph{(c)} holds, then, up to the canonical
identification,
\[
N'=\overline N\cap M_{(K')}.
\]
Since $\overline N$ and $M_{(K')}$ are $A_{(K')}$-submodules, so is
$N'$.
\end{proof}

\begin{corollary}[4.8.7]
\phantomsection
\label{IV.4.8.7}
Under the hypotheses of \sref{IV.4.8.6}, there is a smallest subfield
$K'$ of $K$ containing $k$ for which the equivalent conditions of
\sref{IV.4.8.6} hold.
\end{corollary}

\begin{proof}
It is enough to use condition \emph{(c)}, for which the assertion is
Bourbaki, \emph{Alg.}, chap.~II, 3rd ed., \S~8,
n\textsuperscript{o}~6, prop.~6.
\end{proof}

\begin{lemma}[4.8.8]
\phantomsection
\label{IV.4.8.8}
With the notation of \sref{IV.4.8.2}:
\begin{enumerate}
\item[\rm(i)]
\phantomsection
\label{IV.4.8.8.i}
Let
$\overline u:\sh F\otimes_kK\to\sh G\otimes_kK$ be an
$\sh O_{X_{(K)}}$-homomorphism, and let
\[
\overline{\sh H}\subset
(\sh F\otimes_kK)\oplus(\sh G\otimes_kK)
\]
be its graph.
Then $\overline u$ is defined over $K'$ if and only if
$\overline{\sh H}$ is.

\item[\rm(ii)]
\phantomsection
\label{IV.4.8.8.ii}
Let $\overline Z$ be a closed sub-prescheme of $X_{(K)}$, and let
$\overline{\sh J}$ be the quasi-coherent Ideal of
$\sh O_{X_{(K)}}$ defining $\overline Z$.
Then $\overline Z$ is defined over $K'$ if and only if
$\overline{\sh J}$ is.

\item[\rm(iii)]
\phantomsection
\label{IV.4.8.8.iii}
Let $\overline f:X_{(K)}\to Y_{(K)}$ be a $K$-morphism, and let
$\overline Z$ be its graph, regarded as a sub-prescheme of
$X_{(K)}\times_KY_{(K)}$ \sref[I]{I.5.3.11}.
Then $\overline f$ is defined over $K'$ if and only if
$\overline Z$ is.
\end{enumerate}
\end{lemma}

\begin{proof}
Assertion \emph{(ii)} follows immediately from \sref[I]{I.4.4.5}.
Necessity in \emph{(iii)} is clear.
Conversely, suppose
$\overline Z=Z'_{(K)}$, where $Z'$ is a sub-prescheme of
$(X\times_kY)_{(K')}$.
If $p':Z'\to X_{(K')}$ is the restriction of the first projection,
then $p'_{(K)}$ is an isomorphism, and hence so is $p'$
\sref{IV.2.7.1.viii}.
Thus $Z'$ is the graph of a $K'$-morphism
$f':X_{(K')}\to Y_{(K')}$ \sref[I]{I.5.3.11}, and
$\overline f=f'_{(K)}$ \sref[I]{I.5.3.12}.

Finally, to prove \emph{(i)}, \sref{IV.4.8.5} permits us to assume
$X=\Spec(A)$ affine, with $\sh F=\widetilde M$,
$\sh G=\widetilde N$, and
$\overline u=\widetilde{\overline v}$ for an
$A_{(K)}$-homomorphism
$\overline v:M_{(K)}\to N_{(K)}$.
Suppose the graph of $\overline v$ is
$P'_{(K)}$, where $P'$ is an $A_{(K')}$-submodule of
$M_{(K')}\oplus N_{(K')}$.
If $p':P'\to M_{(K')}$ is the restriction of the first projection,
then $p'\otimes1_K$ is an isomorphism of $A_{(K)}$-modules.
Faithful flatness $(\mathbf 0_I,\,6.4.1)$ implies that $p'$ is an
isomorphism of $A_{(K')}$-modules.
Hence $P'$ is the graph of an $A_{(K')}$-homomorphism
$v':M_{(K')}\to N_{(K')}$, and plainly
$\overline v=v'\otimes1_K$.
\end{proof}

\begin{proposition}[4.8.9]
\phantomsection
\label{IV.4.8.9}
Let $k$ be a field, let $K/k$ be an extension, let $X$ be a
$k$-prescheme, let $\sh F$ be a quasi-coherent $\sh O_X$-Module, and
let $\overline{\sh H}$ be a quasi-coherent
sub-$\sh O_{X_{(K)}}$-Module of
$\sh F\otimes_kK$.
Then $\overline{\sh H}$ has a smallest field of definition.
\end{proposition}

\begin{proof}
When $X$ is affine, this is precisely \sref{IV.4.8.7}.
In general, let $(X_\alpha)$ be an affine open covering of $X$.
For every $\alpha$, there is a smallest subfield $K'_\alpha$ of $K$
containing $k$ over which
$\overline{\sh H}|_{(X_\alpha)_{(K)}}$ is defined.
By \sref{IV.4.8.5}, the subfield of $K$ generated by the
$K'_\alpha$ is the smallest field of definition of $\overline{\sh H}$.
\end{proof}

\begin{corollary}[4.8.10]
\phantomsection
\label{IV.4.8.10}
With the notation of \sref{IV.4.8.9}, let $\sh G$ be a second
quasi-coherent $\sh O_X$-Module and let
\[
\overline u:\sh F\otimes_kK\longrightarrow\sh G\otimes_kK
\]
be an $\sh O_{X_{(K)}}$-homomorphism.
Then $\overline u$ has a smallest field of definition.
\end{corollary}

\begin{proof}
By \sref{IV.4.8.8.i}, such a field is also the smallest field of
definition of the graph of $\overline u$, whose existence follows
from \sref{IV.4.8.9}.
\end{proof}

\oldpage[IV-2]{83}

\begin{corollary}[4.8.11]
\phantomsection
\label{IV.4.8.11}
Let $k$ be a field, let $X$ be a $k$-prescheme, let $K/k$ be an
extension, and let $\overline Z$ be a closed sub-prescheme of
$X_{(K)}$.
Then $\overline Z$ has a smallest field of definition.
\end{corollary}

\begin{proof}
Let $\overline{\sh J}$ be the quasi-coherent Ideal of
$\sh O_{X_{(K)}}$ defining $\overline Z$.
By \sref{IV.4.8.8.ii}, a field defines $\overline Z$ if and only if
it defines $\overline{\sh J}$, and the assertion follows from
\sref{IV.4.8.9}.
\end{proof}

\begin{corollary}[4.8.12]
\phantomsection
\label{IV.4.8.12}
Let $k$ be a field, let $X$ and $Y$ be two $k$-preschemes with $Y$
separated, let $K/k$ be an extension, and let
$\overline f:X_{(K)}\to Y_{(K)}$ be a $K$-morphism.
Then $\overline f$ has a smallest field of definition.
\end{corollary}

\begin{proof}
By \sref{IV.4.8.8.iii}, the existence of such a field is equivalent
to the existence of a smallest field of definition for the graph
$\overline Z$ of $\overline f$.
Since $\overline Z$ is a closed sub-prescheme of
$(X\times_kY)_{(K)}$ \sref[I]{I.5.4.3}, the assertion follows from
\sref{IV.4.8.11}.
\end{proof}

\begin{proposition}[4.8.13]
\phantomsection
\label{IV.4.8.13}
Under the hypotheses of \sref{IV.4.8.11}
(resp.\ \sref{IV.4.8.9}, \sref{IV.4.8.10}, or
\sref{IV.4.8.12}), suppose in addition that $X$ is of finite type
over $k$
(resp.\ that $X$ is of finite type over $k$ and $\sh F$ is coherent;
that $X$ is of finite type over $k$ and both $\sh F$ and $\sh G$ are
coherent; or that $X$ and $Y$ are of finite type over $k$).
Then the smallest field of definition of $\overline Z$
(resp.\ $\overline{\sh H}$, $\overline u$, or $\overline f$) is an extension of
finite type of $k$.
\end{proposition}

\begin{proof}
It is enough to treat \sref{IV.4.8.9}, and we may assume that $X$ is
affine because it is a finite union of affine opens of finite type
over $k$.
With the notation of \sref{IV.4.8.6}, it remains to show that if $A$
is a $k$-algebra of finite type and $M$ an $A$-module of finite type,
then one of the fields $K'$ satisfying the equivalent conditions of
\sref{IV.4.8.6} may be chosen of finite type over $k$.

Let $(m_i)_{1\le i\le n}$ generate $M$ over $A$.
Since $A$ is Noetherian, so is $A_{(K)}$, and the
$A_{(K)}$-submodule $\overline N$ of $M_{(K)}$ has a finite set of
generators $n_j$ $(1\le j\le r)$.
Write
\[
n_j=\sum_i(m_i\otimes1)b_{ij},
\qquad b_{ij}\in A_{(K)},
\]
and then
\[
b_{ij}=\sum_h a_{ijh}\otimes c_{ijh},
\qquad a_{ijh}\in A,\quad c_{ijh}\in K.
\]
The extension $K'$ of $k$ generated by the finitely many
$c_{ijh}$ has the required property.
\end{proof}

\begin{proposition}[4.8.14]
\phantomsection
\label{IV.4.8.14}
Let $k$ be a field, let $X$ be a $k$-prescheme of finite type, and
let $K/k$ be an extension containing a perfect extension of $k$.
Then the smallest field of definition of the closed sub-prescheme
$(X_{(K)})_{\mathrm{red}}$ of $X_{(K)}$ is a finite radicial
extension of $k$.
\end{proposition}

\begin{proof}
If $p$ is the characteristic exponent of $k$, the hypothesis implies
that $k'=k^{p^{-\infty}}$ is contained in $K$.
The prescheme $(X_{(k')})_{\mathrm{red}}$ is geometrically reduced
over $k'$ \sref{IV.4.6.1}.
Thus
\[
(X_{(k')})_{\mathrm{red}}\otimes_{k'}K
\]
is reduced and is consequently equal to
$(X_{(K)})_{\mathrm{red}}$ \sref[I]{I.5.1.8}.
In other words, $k'$ is a field of definition of
$(X_{(K)})_{\mathrm{red}}$.
Since $k'/k$ is algebraic, the conclusion follows from
\sref{IV.4.8.13}.
\end{proof}

\begin{remarks}[4.8.15]
\phantomsection
\label{IV.4.8.15}
\begin{enumerate}
\item[\rm(i)]
\phantomsection
\label{IV.4.8.15.i}
The conclusion of \sref{IV.4.8.14} need not hold if $K$ is not
assumed to contain a perfect extension of $k$.
For example, let $k_0$

\oldpage[IV-2]{84}

be a field of characteristic $p>0$, let
$k=k_0(s,t)$ be the field of rational functions in two
indeterminates, let
\[
L=k(s^{1/p},t^{1/p}),
\qquad
X=\Spec(L),
\]
let $u$ be a third indeterminate, and put
\[
K=k\bigl(u,s^{1/p}+ut^{1/p}\bigr).
\]
One readily verifies that $k$ is algebraically closed in $K$ and that
$L\otimes_kK$ has a nonzero nilradical, so $X_{(K)}$ is not reduced.
But $(X_{(K)})_{\mathrm{red}}$ is not defined over $k$, while $K$
contains no intermediate extension of finite degree over $k$ other
than $k$ itself.
Hence the smallest field of definition of
$(X_{(K)})_{\mathrm{red}}$ cannot be finite over $k$.

\item[\rm(ii)]
\phantomsection
\label{IV.4.8.15.ii}
Given an object in one of the target sets in \sref{IV.4.8.2} and an
extension $K_1/K$, the object has a smallest field of definition if
and only if its image under $\varepsilon_{K_1,K}$ does; the two smallest
subfields are necessarily the same, as observed in
\sref{IV.4.8.4}.
\end{enumerate}
\end{remarks}
% END NATIVE INPUT components/08_section48_fields_of_definition
% BEGIN NATIVE INPUT components/09_section49_fields_of_definition_of_subsets
\subsection{Fields of definition of subsets of a prescheme}
\phantomsection
\label{subsection:IV.4.9}

\begin{proposition}[4.9.1]
\phantomsection
\label{IV.4.9.1}
Let $k$ be a field, let $X$ be a $k$-prescheme, let $K/k$ be an
extension, let $\overline T$ be a closed subset of $X_{(K)}$, and let
$K'$ be an intermediate field of $K/k$.
The following conditions are equivalent:
\begin{enumerate}
\item[\rm(a)]
\phantomsection
\label{IV.4.9.1.a}
$\overline T$ is defined over $K'$.
\item[\rm(b)]
\phantomsection
\label{IV.4.9.1.b}
The open subset $\overline U=X_{(K)}-\overline T$ of $X_{(K)}$ is
defined over $K'$.
\item[\rm(b$'$)]
\phantomsection
\label{IV.4.9.1.bprime}
The sub-prescheme of $X_{(K)}$ induced on the open subset
$\overline U$ is defined over $K'$.
\item[\rm(c)]
\phantomsection
\label{IV.4.9.1.c}
There exists a closed sub-prescheme $\overline Z$ of $X_{(K)}$,
defined over $K'$, whose underlying space is $\overline T$.
\end{enumerate}
\end{proposition}

\begin{proof}
Let $g:X_{(K)}\to X_{(K')}$ be the canonical projection.
To say that $\overline T$ is defined over $K'$ means that there is a
closed subset $T'$ of $X_{(K')}$ such that
$\overline T=g^{-1}(T')$.
Since
\[
g^{-1}(X_{(K')}-T')=X_{(K)}-g^{-1}(T')
\]
for every subset $T'$ of $X_{(K')}$, conditions
\sref{IV.4.9.1.a}, \sref{IV.4.9.1.b}, and
\sref{IV.4.9.1.bprime} are equivalent.

If $\overline Z$, whose underlying space is $\overline T$, is defined
over $K'$, then
\[
\overline Z=Z'\otimes_{K'}K
\]
for a closed sub-prescheme $Z'$ of $X_{(K')}$.
If $T'$ is the underlying space of $Z'$, then
$\overline T=g^{-1}(T')$ \sref[I]{I.4.4.1}; thus
\sref{IV.4.9.1.c} implies \sref{IV.4.9.1.a}.
Conversely, if \sref{IV.4.9.1.a} holds, choose a closed sub-prescheme
$Z'$ of $X_{(K')}$ whose underlying space is $T'$
\sref[I]{I.5.2.1}, and take
$\overline Z=Z'\otimes_{K'}K$.
\end{proof}

\begin{corollary}[4.9.2]
\phantomsection
\label{IV.4.9.2}
With the notation of \sref{IV.4.9.1}, suppose that $K'$ is a field
of definition of $\overline T$.
Every intermediate field $K''$ with $k\subset K''\subset K'$ for
which $K'/K''$ is radicial is also a field of definition of
$\overline T$.
\end{corollary}

\begin{proof}
The canonical projection
\[
g:X_{(K')}\longrightarrow X_{(K'')}
\]
is integral, surjective, and radicial, hence is a homeomorphism
\sref{IV.2.4.5}.
\end{proof}

\begin{remark}[4.9.3]
\phantomsection
\label{IV.4.9.3}
It follows from \sref{IV.4.9.2} that a closed subset
$\overline T$ of $X_{(K)}$ need not have a smallest field of
definition.
For example, let $K=k(t)$, where $t$ is an indeterminate and $k$ is
a perfect field of characteristic $p>0$.
A closed subset $\overline F$ of $X_{(K)}$ can have a smallest field
of definition only if it is already defined over $k$.
Indeed, for every subfield $K'$ of $K$ containing $k$ with
$K'\ne k$, one has $(K')^p\ne K'$, the field $(K')^p$ contains $k$,
and $K'/(K')^p$ is radicial.
Thus, if $K'$ is a field of definition of $\overline F$, then so is
$(K')^p$.
Nevertheless, one has the following result.
\end{remark}

\oldpage[IV-2]{85}

\begin{proposition}[4.9.4]
\phantomsection
\label{IV.4.9.4}
With the notation of \sref{IV.4.9.1}, suppose that $K/K'$ is
separable.
Then $K'$ is a field of definition of $\overline T$ if and only if
it is a field of definition of the reduced sub-prescheme of
$X_{(K)}$ whose underlying space is $\overline T$.
\end{proposition}

\begin{proof}
The condition is sufficient by \sref{IV.4.9.1}.
Conversely, suppose that
$\overline T=g^{-1}(T')$, where $T'$ is a closed subset of
$X_{(K')}$, and let $Z'$ be the reduced sub-prescheme of $X_{(K')}$
whose underlying space is $T'$.
Then
\[
\overline Z=Z'\otimes_{K'}K=\Spec(K)\times_{K'}Z'
\]
is reduced by \sref{IV.4.6.1.e} and the separability of $K/K'$,
and its underlying space is $\overline T$.
\end{proof}

\begin{corollary}[4.9.5]
\phantomsection
\label{IV.4.9.5}
Suppose that one of the following hypotheses holds:
\begin{enumerate}
\item[\rm(a)]
\phantomsection
\label{IV.4.9.5.a}
$k$ has characteristic zero.
\item[\rm(b)]
\phantomsection
\label{IV.4.9.5.b}
$K/k$ is algebraic and $X$ is of finite type over $k$.
\end{enumerate}
Then every closed subset $\overline T$ of $X_{(K)}$ has a smallest
field of definition.
This field is a separable extension of $k$, and under
\sref{IV.4.9.5.b} it is finite over $k$.
\end{corollary}

\begin{proof}
Under \sref{IV.4.9.5.a}, $K$ is separable over each of its
subfields.
It follows from \sref{IV.4.9.4} that the fields of definition of
$\overline T$ are exactly those of the reduced sub-prescheme of
$X_{(K)}$ whose underlying space is $\overline T$.
The assertion therefore follows from \sref{IV.4.8.11}.

Under \sref{IV.4.9.5.b}, the extension $K$ is radicial over a
separable extension $K_1$ of $k$.
For every intermediate field $K'$ of $K/k$, the extension
$K'/(K'\cap K_1)$ is radicial.
By \sref{IV.4.9.2}, $K'$ is a field of definition of
$\overline T$ if and only if $K'\cap K_1$ is.
This reduces the question to the case $K=K_1$, hence to the case of
a separable algebraic extension of $k$, and the argument finishes as
under \sref{IV.4.9.5.a}.
\end{proof}

\begin{proposition}[4.9.6]
\phantomsection
\label{IV.4.9.6}
Let $k$ be a field, let $X$ be a $k$-prescheme, let $K/k$ be an
extension, let $\overline U$ be a subset of $X_{(K)}$ that is both
open and closed, and let $\overline V$ be its complement.
For an intermediate field $K'$ of $K/k$, the following conditions
are equivalent:
\begin{enumerate}
\item[\rm(a)]
\phantomsection
\label{IV.4.9.6.a}
$K'$ is a field of definition of the subset $\overline U$ of
$X_{(K)}$.
\item[\rm(a$'$)]
\phantomsection
\label{IV.4.9.6.aprime}
$K'$ is a field of definition of the closed sub-prescheme of
$X_{(K)}$ induced on the open subset $\overline U$.
\item[\rm(b)]
\phantomsection
\label{IV.4.9.6.b}
$K'$ is a field of definition of the subset $\overline V$ of
$X_{(K)}$.
\item[\rm(b$'$)]
\phantomsection
\label{IV.4.9.6.bprime}
$K'$ is a field of definition of the closed sub-prescheme of
$X_{(K)}$ induced on the open subset $\overline V$.
\end{enumerate}
\end{proposition}

\begin{proof}
It is clear that
{\hyperlink{ega4.c04.sec.4-9-6-aprime}{(a$'$)}} implies
{\hyperlink{ega4.c04.sec.4-9-6-a}{(a)}}, and
{\hyperlink{ega4.c04.sec.4-9-6-a}{(a)}} implies
{\hyperlink{ega4.c04.sec.4-9-6-bprime}{(b$'$)}} by \sref{IV.4.9.1}.
Since $\overline U$ and $\overline V$ play symmetric roles,
{\hyperlink{ega4.c04.sec.4-9-6-bprime}{(b$'$)}} implies
{\hyperlink{ega4.c04.sec.4-9-6-b}{(b)}}, and
{\hyperlink{ega4.c04.sec.4-9-6-b}{(b)}} implies
{\hyperlink{ega4.c04.sec.4-9-6-aprime}{(a$'$)}}.
\end{proof}

\begin{corollary}[4.9.7]
\phantomsection
\label{IV.4.9.7}
Under the hypotheses of \sref{IV.4.9.6}, there is a smallest field
of definition $K'$ for the subset $\overline U$ of $X_{(K)}$.
It is also the smallest field of definition of the subset
$\overline V$ and of the closed sub-preschemes of $X_{(K)}$ induced
on the open subsets $\overline U$ and $\overline V$.
If, in addition, $X$ is of finite type over $k$, then $K'/k$ is
finite and separable.
\end{corollary}

\begin{proof}
In view of \sref{IV.4.9.6} and \sref{IV.4.8.11}, only the last
assertion remains to be proved.
By \sref{IV.4.8.15.ii}, we may reduce to the case in which $K$ is
algebraically closed.
Let $\overline k$ be the algebraic closure of $k$ contained in $K$,
and let $W_i$ $(1\le i\le m)$ be the connected components of
$X_{(\overline k)}$.

\oldpage[IV-2]{86}
If
\[
p:X_{(K)}\longrightarrow X_{(\overline k)}
\]
is the canonical projection, then the $p^{-1}(W_i)$ are the
connected components of $X_{(K)}$ \sref{IV.4.4.6}.
Thus $\overline U$ is a union of some of these components and is
therefore defined over $\overline k$.
The conclusion follows from \sref{IV.4.9.5.b}.
\end{proof}
% END NATIVE INPUT components/09_section49_fields_of_definition_of_subsets
% END NATIVE INPUT ega4-4.source_aligned_r1
% BEGIN NATIVE INPUT ega4-5.source_aligned_r1
% BEGIN NATIVE INPUT components/01_section51_dimension_of_preschemes
\section{Dimension, depth, and regularity in locally Noetherian preschemes}
\phantomsection
\label{section:IV.5}

\oldpage[IV-2]{86}

This section merely restates, in geometric language and with various
technical complements, the notions and results of commutative algebra
set out in Chapter~0, \S\S~16 and~17.

\subsection{Dimension of preschemes}
\phantomsection
\label{subsection:IV.5.1}

\begin{env}[5.1.1]
\phantomsection
\label{IV.5.1.1}
Let $A$ be a ring and let $\mathfrak J$ be an ideal of $A$.
Recall \sref[0]{0.16.1} that the definitions in the theory of the
dimension of a ring \sref[0]{0.16.1.1} and those in the theory of the
combinatorial dimension of topological spaces \sref[0]{0.14.1.2} give
the following relations:
\begin{align}
\dim(\Spec(A))&=\dim(A),
\tag{5.1.1.1}\label{IV.5.1.1.1}\\
\dim(\mathrm V(\mathfrak J))&=\dim(A/\mathfrak J),
\tag{5.1.1.2}\label{IV.5.1.1.2}\\
\operatorname{codim}(\mathrm V(\mathfrak J),\Spec(A))
  &=\operatorname{ht}(\mathfrak J).
\tag{5.1.1.3}\label{IV.5.1.1.3}
\end{align}

\begin{enumerate}
\item[\rm(5.1.1.4)]
\phantomsection
\label{IV.5.1.1.4}
$\Spec(A)$ is a catenary space if and only if $A$ is a catenary ring.
\item[\rm(5.1.1.5)]
\phantomsection
\label{IV.5.1.1.5}
$\Spec(A)$ is equidimensional if and only if $A$ is equidimensional,
if and only if the minimal prime ideals $\mathfrak p_\alpha$ of $A$
are such that the rings $A/\mathfrak p_\alpha$ all have the same
dimension.
\item[\rm(5.1.1.6)]
\phantomsection
\label{IV.5.1.1.6}
$\Spec(A)$ is equicodimensional if and only if $A$ is
equicodimensional, if and only if all maximal ideals of $A$ have the
same height.
\end{enumerate}

Recall that a Noetherian ring $A$ is said to be
\emph{biequidimensional} if $\Spec(A)$ is biequidimensional; that is,
if $\Spec(A)$ is Noetherian and $A$ is equidimensional,
equicodimensional, catenary, and of finite dimension.
\end{env}

\begin{proposition}[5.1.2]
\phantomsection
\label{IV.5.1.2}
Let $X$ be a prescheme, let $Y$ be an irreducible closed subset of
$X$, and let $y$ be the generic point of $Y$. Then
\begin{equation}
\operatorname{codim}(Y,X)=\dim(\sh O_{X,y}).
\tag{5.1.2.1}\label{IV.5.1.2.1}
\end{equation}
\end{proposition}

\begin{proof}
Indeed, the irreducible closed subsets of $X$ containing $y$ are
canonically in bijection with the irreducible closed subsets of
$\Spec(\sh O_{X,y})$ \sref[I]{I.2.4.2}, and hence with the prime
ideals of $\sh O_{X,y}$; formula \eqref{IV.5.1.2.1} follows from the
definitions.

In particular, this recovers the fact that the generic points of the
irreducible components of $X$ are the only points $x\in X$ such that
$\dim(\sh O_x)=0$ \sref[I]{I.1.1.14}.
\end{proof}

\oldpage[IV-2]{87}

\begin{corollary}[5.1.3]
\phantomsection
\label{IV.5.1.3}
For every closed subset $Y$ of a prescheme $X$, one has
\begin{equation}
\operatorname{codim}(Y,X)=\inf_{y\in Y}\dim(\sh O_{X,y}).
\tag{5.1.3.1}\label{IV.5.1.3.1}
\end{equation}
Moreover, if $X$ is locally Noetherian, then for every $x\in Y$,
\begin{equation}
\operatorname{codim}_x(Y,X)
=\inf_{\substack{y\in Y\\x\in\overline{\{y\}}}}
  \dim(\sh O_{X,y}).
\tag{5.1.3.2}\label{IV.5.1.3.2}
\end{equation}
\end{corollary}

\begin{proof}
Formula \eqref{IV.5.1.3.2} follows from
\eqref{IV.5.1.3.1} and \sref[0]{0.14.2.6}.
This corollary allows us to \emph{define}, for any subset $Y$ of a
prescheme $X$, the \emph{codimension}
$\operatorname{codim}(Y,X)$ of $Y$ in $X$ to be the right-hand side
of \eqref{IV.5.1.3.1}.
\end{proof}

\begin{proposition}[5.1.4]
\phantomsection
\label{IV.5.1.4}
For every prescheme $X$, one has
\begin{equation}
\dim(X)=\sup_{x\in X}\dim(\sh O_x).
\tag{5.1.4.1}\label{IV.5.1.4.1}
\end{equation}
If every irreducible closed subset of $X$ contains a closed point,
then one also has
\begin{equation}
\dim(X)=\sup_{x\in F}\dim(\sh O_x),
\tag{5.1.4.2}\label{IV.5.1.4.2}
\end{equation}
where $F$ is the set of closed points of $X$.
\end{proposition}

\begin{proof}
It suffices to observe, by \sref[I]{I.2.4.2}, that the chains of
irreducible closed subsets of $X$ correspond bijectively to the
chains of irreducible closed subsets of a local scheme of $X$.
When every irreducible closed subset of $X$ contains a closed point,
one may plainly restrict to the local schemes at closed points of
$X$.

Every irreducible closed subset of $X$ contains a closed point when
$X$ is quasi-compact \sref[0]{0.2.1.3}; we shall see below
\sref{IV.5.1.11} that the same is true when $X$ is locally
Noetherian.
\end{proof}

\begin{corollary}[5.1.5]
\phantomsection
\label{IV.5.1.5}
A prescheme $X$ is catenary if and only if the local ring $\sh O_x$
is catenary for every $x\in X$.
If, in addition, every irreducible closed subset of $X$ contains a
closed point, then for $X$ to be catenary it is enough that
$\sh O_x$ be catenary for every closed point $x$ of $X$.
\end{corollary}

\begin{proof}
The argument is the same as in \sref{IV.5.1.4}, since it is a matter
of comparing chains of irreducible closed subsets with the same
endpoints.
\end{proof}

\begin{env}[5.1.6]
\phantomsection
\label{IV.5.1.6}
We now examine the more special properties tied to Noetherian
hypotheses.

Recall \sref[0]{0.16.2.3} that a Noetherian local ring $A\ne0$ has
finite dimension, equal to the minimum number of generators of an
ideal of definition of $A$.
For every prime ideal $\mathfrak p$ of $A$, its height, equal to
$\dim(A_{\mathfrak p})$, is therefore also finite.
Together with \eqref{IV.5.1.2.1}, \eqref{IV.5.1.3.1}, and
\eqref{IV.5.1.4.1}, these facts show the following.
\end{env}

\begin{proposition}[5.1.7]
\phantomsection
\label{IV.5.1.7}
For every nonempty closed subset $Y$ of a locally Noetherian
prescheme $X$, the number $\operatorname{codim}(Y,X)$ is finite.
If $X$ is Noetherian and affine and $Y$ is an irreducible closed
subset of $X$, then $\operatorname{codim}(Y,X)$ is the minimum
number of sections $s_i$ of $\sh O_X$ over $X$ such that $Y$ is an
irreducible component of the set of points $x\in X$ for which
$s_i(x)=0$ for every $i$.
\end{proposition}

\begin{corollary}[5.1.8]
\phantomsection
\label{IV.5.1.8}
Let $X$ be a locally Noetherian prescheme, let $\sh L$ be an
invertible $\sh O_X$-module, and let $f$ be a section of $\sh L$ over
$X$.
Every irreducible component of the set $Z$ of points $x\in X$ such
that $f(x)=0$ has codimension at most $1$ in $X$; it has codimension
$1$ if $Z$ contains no irreducible component of $X$.
\end{corollary}

\begin{proof}
We may restrict to the case $\sh L=\sh O_X$.
If $y$ is a generic point of an irreducible component $Y$ of $Z$,
the ideal $(f_y)$ of $\sh O_{X,y}$ must be such that
$\sh O_{X,y}/(f_y)$ has only one prime ideal.
Thus $f_y$ generates an ideal of definition of the Noetherian local
ring $\sh O_{X,y}$, and
$\operatorname{codim}(Y,X)\leq1$ by \sref{IV.5.1.7}.
If $Z$ contains no irreducible component of $X$, then
$\operatorname{codim}(Y,X)$ cannot be $0$ by
\sref[0]{0.14.2.1}.
\end{proof}

\oldpage[IV-2]{88}

\begin{proposition}[5.1.9]
\phantomsection
\label{IV.5.1.9}
Let $X$ be a locally Noetherian prescheme and let $Y$ be a closed
subset of $X$.
If $x\in Y$ and $\sh O_{X,x}$ is a catenary ring, then
\begin{equation}
\begin{split}
\operatorname{codim}_x(Y,X)
&=\dim(\sh O_{X,x})-\operatorname{codim}(\overline{\{x\}},Y)\\
&=\dim(\sh O_{X,x})-\dim(\sh O_{Y,x}).
\end{split}
\tag{5.1.9.1}\label{IV.5.1.9.1}
\end{equation}
\end{proposition}

\begin{proof}
Let $Y_i$ ($1\leq i\leq n$) be the irreducible components of $Y$
containing $x$; there are only finitely many because $Y$ is locally
Noetherian.
Let $y_i$ be the generic point of $Y_i$.
Put $A=\sh O_{X,x}$, and let $\mathfrak p_i$ be the prime ideal of
$A$ corresponding to $Y_i\cap\Spec(A)$.
The irreducible closed subsets of $Y$ containing $x$ correspond
bijectively to the prime ideals of $A$ containing one of the
$\mathfrak p_i$, so
\[
\dim(\sh O_{Y,x})=\sup_i\dim(A/\mathfrak p_i).
\]
On the other hand, $\sh O_{X,y_i}=A_{\mathfrak p_i}$, and the
hypothesis on $A$ gives
\[
\dim(A)=\dim(A/\mathfrak p_i)+\dim(A_{\mathfrak p_i})
\]
by \sref[0]{0.16.1.4}.
The conclusion follows from
\[
\operatorname{codim}_x(Y,X)
  =\inf_i\dim(\sh O_{X,y_i})
\]
using \eqref{IV.5.1.2.1} and \sref[0]{0.14.2.6}.
\end{proof}

\begin{proposition}[5.1.10]
\phantomsection
\label{IV.5.1.10}
\begin{enumerate}[{\rm(i)}]
\item
\phantomsection
\label{IV.5.1.10.i}
In every prescheme $X$, every nonempty locally constructible subset
$E$ contains a point $x$ such that $\{x\}$ is locally closed in $X$;
equivalently, $x$ is isolated in $\overline{\{x\}}$.
\item
\phantomsection
\label{IV.5.1.10.ii}
Let $X$ be a locally Noetherian prescheme and let $x\in X$ be such
that $\{x\}$ is locally closed in $X$.
Then $\dim(\overline{\{x\}})\leq1$; consequently, every point
$y\ne x$ of $\overline{\{x\}}$ is closed in $X$.
\end{enumerate}
\end{proposition}

\begin{proof}
\sref{IV.5.1.10.i} is a special case of the following lemma.
\end{proof}

\begin{lemma}[5.1.10.1]
\phantomsection
\label{IV.5.1.10.1}
Let $X$ be a topological space with the following property: for every
nonempty locally closed subset $Z$ of $X$, there exist a subset
$Z'$ of $Z$, locally closed in $X$ (or, equivalently, in $Z$), and
a point $x\in Z'$ that is closed in $Z'$.
Then every nonempty locally closed subset $Z$ of $X$ contains a point
$x$ isolated in $\overline{\{x\}}$.
\end{lemma}

\begin{proof}
Let $Z'\subset Z$ be a locally closed subset of $Z$ containing a
point $x$ that is closed in $Z'$.
There is an open neighbourhood $U$ of $x$ in $X$ such that $x$ is
closed in $U$; hence
\[
U\cap\overline{\{x\}}=\{x\},
\]
which says precisely that $x$ is isolated in $\overline{\{x\}}$.

The lemma applies to the underlying space of every prescheme $X$.
Indeed, $Z$ is then itself the underlying space of a prescheme
\sref[I]{I.5.2.1}, and one may take for $Z'$ an affine open subset of
$Z$, which is a quasi-compact Kolmogorov space and therefore contains
a closed point \sref[0]{0.2.1.3}.
\end{proof}

\oldpage[IV-2]{89}

\begin{proof}[Proof of \sref{IV.5.1.10.ii}]
Let $Z$ be the reduced sub-prescheme of $X$ whose underlying space is
$\overline{\{x\}}$.
The hypothesis implies that $\{x\}$ is open in $Z$.
Thus, for every $z\in Z$, the generic point $x$ of
$\Spec(\sh O_{Z,z})$ is isolated in $\Spec(\sh O_{Z,z})$.
The ring $A=\sh O_{Z,z}$ is a Noetherian local integral domain, and
the hypothesis implies that there exists $f\in A$ such that $A_f$ is
a field.
It follows from \sref[0]{0.16.3.3} that $\dim(A)\leq1$.
Since $\dim(\sh O_{Z,z})\leq1$ for every $z\in Z$, one has
$\dim(Z)\leq1$.
\end{proof}

\begin{corollary}[5.1.11]
\phantomsection
\label{IV.5.1.11}
If $X$ is a locally Noetherian prescheme, every nonempty closed
subset of $X$ contains a closed point.
\end{corollary}

\begin{proof}
Every closed subset of $X$ is constructible
\sref[0]{0.9.1.1} and \sref[0]{0.9.1.5}, so it suffices to apply
\sref{IV.5.1.10}.
\end{proof}

\begin{env}[5.1.12]
\phantomsection
\label{IV.5.1.12}
Let $X$ be a prescheme, let $\sh F$ be a quasi-coherent
$\sh O_X$-module of finite type, and let
$S=\operatorname{Supp}(\sh F)$ be its support, which is closed in
$X$ \sref[0]{0.5.2.2}.
For every $x\in X$, regard $\operatorname{Supp}(\sh F_x)$ as a closed
subset of the local scheme $\Spec(\sh O_x)$.
By definition \sref[0]{0.16.1.7},
\[
\dim(\sh F_x)=\dim(\operatorname{Supp}(\sh F_x)).
\]
But
\[
\operatorname{Supp}(\sh F_x)
  =S\cap\Spec(\sh O_{X,x})
  =\Spec(\sh O_{S,x}),
\]
where in the last expression $S$ denotes any closed sub-prescheme of
$X$ whose underlying space is $S$.
The irreducible components of $\Spec(\sh O_{S,x})$ are the
intersections of $\Spec(\sh O_{X,x})$ with the irreducible components
of $S$ containing $x$, and correspond bijectively to the latter.
Hence, by \eqref{IV.5.1.1.1},
\begin{equation}
\dim(\sh F_x)=\dim(\sh O_{S,x}).
\tag{5.1.12.1}\label{IV.5.1.12.1}
\end{equation}
It also follows from \eqref{IV.5.1.2.1} that, if $X$ is locally
Noetherian,
\begin{equation}
\dim(\sh F_x)=\operatorname{codim}(\overline{\{x\}},S).
\tag{5.1.12.2}\label{IV.5.1.12.2}
\end{equation}

We say that $\sh F$ is \emph{equidimensional at the point} $x\in X$
if $\sh F_x$ is an equidimensional $\sh O_{X,x}$-module; that is
\sref[0]{0.16.1.7}, if $\operatorname{Supp}(\sh F_x)$ is
equidimensional as a closed subset of $\Spec(\sh O_{X,x})$.
Equivalently, the ring $\sh O_{S,x}$ is equidimensional.

The \emph{dimension} of $\sh F$, denoted $\dim(\sh F)$, is the
dimension of its support $\operatorname{Supp}(\sh F)$.
Equations \eqref{IV.5.1.4.1} and \eqref{IV.5.1.12.1} give
\begin{equation}
\dim(\sh F)=\sup_{x\in X}\dim(\sh F_x).
\tag{5.1.12.3}\label{IV.5.1.12.3}
\end{equation}

If $X=\Spec(A)$ is affine and $\sh F=\widetilde M$, where $M$ is a
finite $A$-module, then
$\dim(\sh F)=\dim(M)$ by \sref[0]{0.16.1.7} and
\eqref{IV.5.1.4.1}.
\end{env}

\begin{proposition}[5.1.13]
\phantomsection
\label{IV.5.1.13}
Let $X$ be a prescheme, let $\sh F$ be a quasi-coherent
$\sh O_X$-module of finite type, let $x\in X$, and let $x'$ be a
generization of $x$ in $X$. Then
\begin{equation}
\dim(\sh F_{x'})\leq\dim(\sh F_x).
\tag{5.1.13.1}\label{IV.5.1.13.1}
\end{equation}
\end{proposition}

\begin{proof}
This follows immediately from \eqref{IV.5.1.12.1} and the
definitions.
\end{proof}
% END NATIVE INPUT components/01_section51_dimension_of_preschemes
% BEGIN NATIVE INPUT components/02_section52_dimension_of_algebraic_preschemes
\subsection{Dimension of an algebraic prescheme}
\phantomsection
\label{subsection:IV.5.2}

\oldpage[IV-2]{90}

\begin{proposition}[5.2.1]
\phantomsection
\label{IV.5.2.1}
Let $k$ be a field, let $X$ be an irreducible prescheme locally of
finite type over $k$, and let $\xi$ be its generic point.
Then $X$ is biequidimensional and
\[
\dim(X)=\operatorname{tr.deg}_k\kappa(\xi).
\]
\end{proposition}

\begin{proof}
Every local ring $\sh O_x$ is a local ring of a finite-type
$k$-algebra, hence is a quotient of a regular local ring
\sref[0]{0.17.3.9}.
Such rings are catenary \sref[0]{0.16.5.12}; consequently $X$ is a
catenary space \sref{IV.5.1.5}.
Since $X$ is irreducible, it follows from \sref{IV.5.1.1} that it is
enough to prove, for every closed point $x\in X$, that
\begin{equation}
\dim(\sh O_x)=\operatorname{tr.deg}_k\kappa(\xi).
\tag{5.2.1.1}\label{IV.5.2.1.1}
\end{equation}

We may plainly suppose that $X$ is reduced and affine, hence integral,
with ring $A$, where $A$ is a finite-type $k$-algebra.
Put
\[
n=\operatorname{tr.deg}_k\kappa(\xi),
\qquad K=\kappa(\xi)=\operatorname{Frac}(A).
\]
There exists a $k$-subalgebra
\[
B=k[t_1,\ldots,t_n]\subset A,
\]
where the $t_i$ are algebraically independent over $k$, such that $A$
is a finite $B$-algebra
(Bourbaki, \emph{Alg\`ebre commutative}, chap.~V, \S~3,
n\textsuperscript{o}~1, th.~1).
Let $\mathfrak m=\mathfrak j_x$, which is maximal by hypothesis.
Then $\mathfrak n=B\cap\mathfrak m$ is a maximal ideal of $B$
(ibid., chap.~V, \S~2, n\textsuperscript{o}~1, prop.~1), and
$A_{\mathfrak m}$ is a local ring of the finite
$B_{\mathfrak n}$-algebra $S^{-1}A$, where
$S=B-\mathfrak n$.
Since $B_{\mathfrak n}$ is integrally closed and $S^{-1}A$ is
integral,
\[
\dim(A_{\mathfrak m})=\dim(B_{\mathfrak n})
\]
by \sref[0]{0.16.1.6}.

We may therefore reduce to $A=k[t_1,\ldots,t_n]$.
The residue field $k'=\kappa(x)$ is then a finite extension of $k$
\sref[I]{I.6.4.2}.
Put $A'=k'[t_1,\ldots,t_n]$.
By \sref[I]{I.2.4.6} and \sref[I]{I.3.3.14}, there exists a maximal
ideal $\mathfrak m'$ of $A'$ lying above $\mathfrak m$ such that $k'$
is the residue field of $A'_{\mathfrak m'}$.
Since $A'$ is integral and finite over $A$, the preceding argument
gives
\[
\dim(A'_{\mathfrak m'})=\dim(A_{\mathfrak m}),
\]
so we are reduced to the case $A/\mathfrak m=k$.

In that case $\mathfrak m$ is generated by the polynomials
$t_i-a_i$, where $a_i\in k$ is the image of $t_i$ in
$A/\mathfrak m$.
Replacing $t_i$ by $t_i-a_i$, we are finally reduced to
\[
\mathfrak m=\sum_{i=1}^n At_i.
\]
The completion of $A_{\mathfrak m}$ is then the formal power-series
ring $k[[t_1,\ldots,t_n]]$.
It has the same dimension as $A_{\mathfrak m}$
\sref[0]{0.16.2.4}, while
$\dim(k[[t_1,\ldots,t_n]])=n$ by \sref[0]{0.17.1.4}.
This proves \eqref{IV.5.2.1.1}.
\end{proof}

\begin{corollary}[5.2.2]
\phantomsection
\label{IV.5.2.2}
For a prescheme $X$ locally of finite type over a field $k$,
$\dim(X)$ agrees with the number defined in \sref{IV.4.1.1}.
\end{corollary}

\begin{proof}
Let $X_\alpha$ be the reduced sub-preschemes whose underlying spaces
are the irreducible components of $X$.
Then
\[
\dim(X)=\sup_\alpha\dim(X_\alpha)
\]
by \sref[0]{0.14.1.2.1}, and it suffices to apply
\sref{IV.5.2.1} to the $X_\alpha$.
\end{proof}

\begin{corollary}[5.2.3]
\phantomsection
\label{IV.5.2.3}
Let $X$ be a prescheme locally of finite type over a field $k$, and
let $x\in X$. Then
\begin{equation}
\dim_x(X)=\dim(\sh O_x)+\operatorname{tr.deg}_k\kappa(x).
\tag{5.2.3.1}\label{IV.5.2.3.1}
\end{equation}
\end{corollary}

\begin{proof}
There is an open neighbourhood $U$ of $x$ in $X$ such that
$\dim_x(X)=\dim(U)$ \sref[0]{0.14.1.4.1}.
We may suppose that the irreducible components of $U$ are the
$X_i\cap U$, where the $X_i$ are the irreducible components of $X$
containing $x$.

\oldpage[IV-2]{91}

Since $U\cap X_i$ is dense in $X_i$, formula \sref{IV.4.1.1.3}
gives $\dim(X_i)=\dim(U\cap X_i)$, and hence
\[
\dim_x(X)=\sup_i\dim(X_i).
\]
The minimal prime ideals of $\sh O_x$ correspond to the generic
points of the $X_i$, so
\[
\dim(\sh O_{X,x})=\sup_i\dim(\sh O_{X_i,x})
\]
by \sref[0]{0.16.1.1.1}.
We are therefore reduced to the case where $X$ is irreducible.
Since $X$ is biequidimensional by \sref{IV.5.2.1},
\[
\dim(X)
  =\dim(\overline{\{x\}})
   +\operatorname{codim}(\overline{\{x\}},X)
\]
by \sref[0]{0.14.3.5.1}.
Finally,
\[
\dim(\overline{\{x\}})
  =\operatorname{tr.deg}_k\kappa(x)
\]
by \sref{IV.5.2.1}, while
\[
\operatorname{codim}(\overline{\{x\}},X)=\dim(\sh O_x)
\]
by \sref{IV.5.1.2}.
\end{proof}

\begin{corollary}[5.2.4]
\phantomsection
\label{IV.5.2.4}
Let $X$ be a prescheme locally of finite type over a field $k$, let
$\sh L$ be an invertible $\sh O_X$-module, and let $f$ be a section
of $\sh L$ over $X$ such that the set
\[
Y=\{x\in X\mid f(x)=0\}
\]
is rare in $X$.
Then
\[
\dim(Y)\leq\dim(X)-1.
\]
Equality holds if $Y$ meets every irreducible component of $X$ having
maximum dimension.
\end{corollary}

\begin{proof}
If $(X_\alpha)$ is the family of irreducible components of $X$, then
\[
\dim(Y)=\sup_\alpha\dim(Y\cap X_\alpha).
\]
We are therefore reduced to the case where $X$ is irreducible,
because $Y\cap X_\alpha$ is rare in $X_\alpha$ and every
$X_\alpha$ has nonempty interior in $X$.
We may also suppose that $Y\ne\varnothing$.
For every maximal point $x$ of $Y$, biequidimensionality gives
\[
\dim(\overline{\{x\}})
  =\dim(X)-\operatorname{codim}(\overline{\{x\}},X).
\]
Since $Y$ is rare in $X$,
\[
\operatorname{codim}(\overline{\{x\}},X)=1
\]
by \sref{IV.5.1.8}, and the result follows.
\end{proof}

\begin{remarks}[5.2.5]
\phantomsection
\label{IV.5.2.5}
\begin{enumerate}[{\rm(i)}]
\item
\phantomsection
\label{IV.5.2.5.i}
Unlike an algebraic $k$-prescheme, a locally Noetherian prescheme
$X$ need not be catenary; see \sref{IV.5.6.11}.
It is catenary, however, if all its local rings $\sh O_x$ are
quotients of regular local rings \sref[0]{0.16.5.12}, and in
particular if $X$ is regular \sref[I]{I.4.1.4}.

Nevertheless, even an integral scheme $X=\Spec(A)$ with $A$ regular
need not be biequidimensional.
In other words, the closed points of $X$ need not all have the same
codimension \sref[0]{0.14.3.3}.
For example, let $B$ be a discrete valuation ring, let
$\mathfrak m=B\pi$ be its maximal ideal, let
$k=B/\mathfrak m$ be its residue field, let $K$ be its fraction
field, and put $A=B[T]$.
The ring $A$ has maximal ideals of height $2$, such as
$(\pi)+(T)$, but also maximal ideals of height $1$, such as the
principal ideal $(\pi T-1)$.
Indeed, a nonzero principal prime ideal has height $1$
\sref{IV.5.1.8}, while
\[
A/(\pi T-1)\simeq B_\pi=K
\]
by \sref[0]{0.1.2.3}; thus $(\pi T-1)$ is maximal and has height $1$.

\item
\phantomsection
\label{IV.5.2.5.ii}
If $X$ is locally of finite type over a field, then
$\dim(U)=\dim(X)$ for every dense open subset $U$ of $X$
\sref{IV.4.1.1.3}.
This does not extend to regular Noetherian preschemes, even when they
are biequidimensional.
For example, let $B$ be a discrete valuation ring with maximal ideal
$\mathfrak m$.
Then $X=\Spec(B)$ has two points, $(0)$ and $\mathfrak m$, the latter
being its only closed point.
One has $\dim(X)=1$, whereas the open subset $U=\{(0)\}$ has
dimension $0$ (see \S~10).\footnote{The printed French text has
$X=\Spec(A)$ here; the ring just specified, and the asserted
two-point spectrum, require $X=\Spec(B)$.}
\end{enumerate}
\end{remarks}
% END NATIVE INPUT components/02_section52_dimension_of_algebraic_preschemes
% BEGIN NATIVE INPUT components/03_section53_support_dimension_and_hilbert_polynomial
\subsection{Dimension of the support of a module and Hilbert polynomial}
\phantomsection
\label{subsection:IV.5.3}

\oldpage[IV-2]{92}

This section uses the results of Chapter~III; it will not be used in
the remainder of this chapter.

\begin{proposition}[5.3.1]
\phantomsection
\label{IV.5.3.1}
Let $A$ be an Artinian local ring, let $X$ be a projective scheme over
$\Spec(A)$, let $\sh L$ be an invertible $\sh O_X$-module that is very
ample relative to $A$, and let $\sh F$ be a nonzero coherent
$\sh O_X$-module. Put
\[
\sh F(n)=\sh F\otimes_{\sh O_X}\sh L^{\otimes n}
\qquad(n\in\bb Z).
\]
Then the degree of the Hilbert polynomial
\[
P(n)=\chi_A(\sh F(n))
\]
of $\sh F$ relative to $A$ \sref[III]{III.2.5.3} equals
$\dim(\operatorname{Supp}(\sh F))$.
\end{proposition}

\begin{proof}
We argue by induction on
\[
d=\dim(\operatorname{Supp}(\sh F)).
\]
There is a closed subprescheme $Y$ of $X$ whose underlying space is
$\operatorname{Supp}(\sh F)$, together with a coherent
$\sh O_Y$-module $\sh G$ such that $\sh F=j_*(\sh G)$, where
$j:Y\to X$ is the canonical injection
(\textbf{Err}\textsubscript{\textbf{III}},~30).
The Hilbert polynomials of $\sh F$ and $\sh G$ are plainly the same,
so we may restrict to the case
$X=\operatorname{Supp}(\sh F)$.

Suppose first that $d=0$.
Every point of $X$ is then closed, and $X$ is an Artinian scheme
\sref[I]{I.6.2.2}; hence $\sh F(n)=\sh F$ for every integer $n$.
Consequently, by \sref[III]{III.2.5.3},
\[
\chi_A(\sh F(n))
  =\operatorname{length}_A\bigl(\Gamma(X,\sh F)\bigr)
\]
for all sufficiently large $n$, and the Hilbert polynomial therefore
has degree $0$.

Suppose now that $d>0$, and put
$Z=\operatorname{Ass}(\sh F)$, a finite set \sref{IV.3.1.6}.
There are an integer $m>0$ and a section
$f\in\Gamma(X,\sh L^{\otimes m})$ such that $X_f$ is a neighbourhood
of $Z$ \sref[II]{II.4.5.4}.
Multiplication by $f$ is an injective homomorphism
\[
\mu_f:\sh F\longrightarrow\sh F(m)
\]
by \sref{IV.3.1.8}; thus there is an exact sequence
\begin{equation}
0\longrightarrow\sh F
 \xrightarrow{\mu_f}\sh F(m)
 \longrightarrow\sh G
 \longrightarrow0,
\tag{5.3.1.1}\label{IV.5.3.1.1}
\end{equation}
where $\sh G$ is coherent.
By Nakayama's lemma, the points
$x\in\operatorname{Supp}(\sh G)$ are exactly those for which
$f(x)=0$.
We shall deduce that
\begin{equation}
\dim(\operatorname{Supp}(\sh G))=d-1.
\tag{5.3.1.2}\label{IV.5.3.1.2}
\end{equation}
This follows from the next lemma.
\end{proof}

\begin{lemma}[5.3.1.3]
\phantomsection
\label{IV.5.3.1.3}
Let $A$ be an Artinian local ring, let $X$ be a projective scheme over
$\Spec(A)$, and let $\sh L$ be an invertible $\sh O_X$-module that is
ample relative to $A$.
For every section $g$ of $\sh L$ over $X$, the set
\[
X\setminus X_g=\{x\in X\mid g(x)=0\}
\]
meets every irreducible component of $X$ having nonzero dimension.
\end{lemma}

\begin{proof}
Indeed, $X_g$ is an affine open subset of $X$
\sref[II]{II.5.5.7}.
If it contains an irreducible component $X'$ of $X$, the reduced
closed subprescheme of $X$ whose underlying space is $X'$ is both
projective and affine over $A$, hence finite over $A$
\sref[III]{III.4.4.2}.
It is therefore an Artinian scheme and has dimension $0$.
\end{proof}

\begin{proof}[Proof of \emph{{\hyperlink{ega4.c04.prop.5-3-1}{(5.3.1)}}}, continued]
Since $Z$ contains the maximal points of $X$, the open subset $X_f$
is dense; hence $\operatorname{Supp}(\sh G)$ is rare in $X$.
Equation \eqref{IV.5.3.1.2} now follows from the lemma and
\sref{IV.5.2.4}.

\oldpage[IV-2]{93}

From the exact sequence \eqref{IV.5.3.1.1}, for every $n$ we obtain
the exact sequence
\[
0\longrightarrow\sh F(n)
 \longrightarrow\sh F(n+m)
 \longrightarrow\sh G(n)
 \longrightarrow0.
\]
For all sufficiently large $n$, there is therefore also an exact
sequence \sref[III]{III.2.2.3}
\[
0\longrightarrow\Gamma(X,\sh F(n))
 \longrightarrow\Gamma(X,\sh F(n+m))
 \longrightarrow\Gamma(X,\sh G(n))
 \longrightarrow0.
\]
Taking \sref[III]{III.2.5.3} into account, we have, for all
sufficiently large $n$,
\[
\chi_A(\sh G(n))
  =\chi_A(\sh F(n+m))-\chi_A(\sh F(n)).
\]
By \eqref{IV.5.3.1.2} and the induction hypothesis, the polynomial
$\chi_A(\sh G(n))$ has degree $d-1$.
The preceding relation therefore implies that
$\chi_A(\sh F(n))$ has degree $d$.
\end{proof}
% END NATIVE INPUT components/03_section53_support_dimension_and_hilbert_polynomial
% BEGIN NATIVE INPUT components/04_section54_dimension_of_morphism_image
\subsection{Dimension of the image of a morphism}
\phantomsection
\label{subsection:IV.5.4}

\begin{proposition}[5.4.1]
\phantomsection
\label{IV.5.4.1}
Let $X$ and $Y$ be locally Noetherian preschemes, and let
$f:X\to Y$ be a morphism.
\begin{enumerate}[{\rm(i)}]
\item
\phantomsection
\label{IV.5.4.1.i}
If $f$ is quasi-finite, then
\[
\dim(X)\leq\dim(\overline{f(X)})\leq\dim(Y).
\]
\item
\phantomsection
\label{IV.5.4.1.ii}
If $f$ is surjective and open (respectively, surjective and closed),
then
\[
\dim(X)\geq\dim(Y).
\]
\end{enumerate}
\end{proposition}

\begin{proof}
\begin{enumerate}[{\rm(i)}]
\item
We may replace $f$ by $f_{\mathrm{red}}$ \sref[II]{II.6.2.4}, and
therefore suppose that $X$ and $Y$ are reduced.
Let $Z$ be the reduced closed subprescheme of $Y$ whose underlying
space is $\overline{f(X)}$.
Then $f=j\circ g$, where $j:Z\to Y$ is the canonical injection and
$g:X\to Z$ is quasi-finite
\sref[I]{I.5.2.2}, \sref[II]{II.6.2.4}.
We may thus restrict to the case in which $f(X)$ is dense in $Y$.

For every $x\in X$, the ring $\sh O_x$ is then a quasi-finite
$\sh O_{f(x)}$-module \sref[II]{II.6.2.2}; consequently
\[
\mathfrak m_{f(x)}\sh O_x
\]
is an ideal of definition of $\sh O_x$ \sref[0]{0.7.4.4}.
If $A\to B$ is a local homomorphism of Noetherian local rings, with
maximal ideal $\mathfrak m$ of $A$, and if $\mathfrak mB$ is an ideal
of definition of $B$, then
\[
\dim(B)\leq\dim(A)
\]
by \sref[0]{0.16.3.10}.
Part~(i) now follows from \sref{IV.5.1.4}.

\item
By the definition of dimension \sref[0]{0.14.1.2}, it is enough to
prove the following.
For every sequence $(y_i)_{0\leq i\leq n}$ of distinct points of $Y$
such that
\[
y_i\in\overline{\{y_{i+1}\}}
\qquad(0\leq i\leq n-1),
\]
there is a sequence $(x_i)_{0\leq i\leq n}$ of points of $X$ such
that
\[
x_i\in\overline{\{x_{i+1}\}}
\qquad(0\leq i\leq n-1)
\]
and $f(x_i)=y_i$ for every $i$.

Suppose first that $f$ is surjective and open.
We construct the $x_i$ by induction on $i$.
Surjectivity gives an $x_0\in X$ such that $f(x_0)=y_0$.
Suppose the $x_i$ have been chosen for $i\leq m$ so that
$f(x_i)=y_i$ for $i\leq m$ and
$x_i\in\overline{\{x_{i+1}\}}$ for $i<m$.
Since $f$ is open and $y_{m+1}$ is a generization of $y_m$, there is
an $x_{m+1}\in f^{-1}(y_{m+1})$ that is a generization of $x_m$
\sref{IV.1.10.3}.
The induction therefore continues.

\oldpage[IV-2]{94}

Suppose now that $f$ is surjective and closed.
We construct the $x_i$ by descending induction on $i$.
Surjectivity gives an $x_n$ such that $f(x_n)=y_n$.
Suppose the $x_i$ have been chosen with the required properties for
$i>m$.
Because $f$ is closed,
\[
f\bigl(\overline{\{x_{m+1}\}}\bigr)
  =\overline{\{f(x_{m+1})\}}
  =\overline{\{y_{m+1}\}}
\]
(Bourbaki, \emph{Topologie g\'en\'erale}, chap.~I,
3\textsuperscript{rd} ed., \S~5, n\textsuperscript{o}~4, prop.~9).
There is therefore an
$x_m\in\overline{\{x_{m+1}\}}$ such that $f(x_m)=y_m$, and the
descending induction continues.
\end{enumerate}
\end{proof}

\begin{corollary}[5.4.2]
\phantomsection
\label{IV.5.4.2}
Let $X$ and $Y$ be locally Noetherian preschemes, and let
$f:X\to Y$ be a finite morphism, hence a closed morphism.
Then
\[
\dim(X)=\dim(\overline{f(X)}).
\]
If $f$ is moreover surjective, then $\dim(X)=\dim(Y)$.
\end{corollary}

\begin{remarks}[5.4.3]
\phantomsection
\label{IV.5.4.3}
\begin{enumerate}[{\rm(i)}]
\item
\phantomsection
\label{IV.5.4.3.i}
We saw in \sref{IV.4.1.2} that if $X$ and $Y$ are preschemes locally
of finite type over a field $k$, and $f$ is a $k$-morphism, the
inequality
\[
\dim(Y)\leq\dim(X)
\]
already holds when $f$ is quasi-compact and dominant.
By contrast, if one assumes only that $X$ and $Y$ are locally
Noetherian, it is possible to have $\dim(Y)>\dim(X)$ even when $f$ is
of finite type, bijective, and a local immersion.

For example, let $A$ be a discrete valuation ring, let $K$ be its
fraction field, and let $k$ be its residue field.
Put $Y=\Spec(A)$, and denote by $a$ the generic point of $Y$ and by
$b$ its closed point; thus $\kappa(a)=K$ and $\kappa(b)=k$.
Let
\[
X=\Spec(K)\amalg\Spec(k),
\]
and let $f:X\to Y$ agree on the two summands with the respective
canonical morphisms \sref[I]{I.2.4.5}.
The morphism $f$ is plainly a bijective local immersion, since
$\{a\}$ is open in $Y$.
It is also of finite type: if $\pi$ is a uniformizer of $A$, then
$K=A[\pi^{-1}]$, so $K$ is a finitely generated $A$-algebra.
Nevertheless,
\[
\dim(X)=0
\qquad\text{and}\qquad
\dim(Y)=1.
\]

\item
\phantomsection
\label{IV.5.4.3.ii}
Let $A\subset B$ be Noetherian rings, with $B$ a finite
$A$-algebra.
Corollary~\sref{IV.5.4.2} again gives
\[
\dim(B)=\dim(A)
\]
\sref[0]{0.16.1.5}.
Suppose in addition that $A$ is a Noetherian local ring.
Then $B$ is a Noetherian semilocal ring
(Bourbaki, \emph{Alg\`ebre commutative}, chap.~IV, \S~2,
n\textsuperscript{o}~5, cor.~3 of prop.~9).
If $\mathfrak n_i$ $(1\leq i\leq r)$ are the maximal ideals of $B$,
then
\begin{equation}
\dim(A)=\sup_i\dim(B_{\mathfrak n_i}).
\tag{5.4.3.1}\label{IV.5.4.3.1}
\end{equation}

Under these hypotheses one need not have
\[
\dim(B_{\mathfrak n_i})=\dim(A)
\]
for every $i$; replacing $B$ by the direct product of $B$ with the
residue field $k$ of $A$ gives an immediate example.
There are even examples in which $A$ and $B$ are integral domains of
dimension $2$ and some of the $B_{\mathfrak n_i}$ have dimension
strictly less than $\dim(A)$ \sref{IV.5.6.11}.
We shall see later, in \sref{IV.5.6.4} and \sref{IV.5.6.10}, that
this last phenomenon cannot occur when $A$ is a quotient of a
regular local ring.
\end{enumerate}
\end{remarks}
% END NATIVE INPUT components/04_section54_dimension_of_morphism_image
% BEGIN NATIVE INPUT components/05_section55_dimension_formula_finite_type
\subsection{Dimension formula for a morphism of finite type}
\phantomsection
\label{subsection:IV.5.5}
\label{IV.5.5}

\begin{env}[5.5.1]
\phantomsection
\label{IV.5.5.1}

\oldpage[IV-2]{94}

Recall \sref[0]{0.16.3.9} that if $A$ and $B$ are Noetherian local
rings, $k$ is the residue field of $A$, and
$\varphi:A\to B$ is a local homomorphism, then
\begin{equation}
\dim(B)\leq\dim(A)+\dim(B\otimes_A k).
\tag{5.5.1.1}\label{IV.5.5.1.1}
\end{equation}
\end{env}

We deduce the following.

\begin{proposition}[5.5.2]
\phantomsection
\label{IV.5.5.2}
Let $X$ and $Y$ be locally Noetherian preschemes, let
$f:X\to Y$ be a morphism, let $x\in X$, and put $y=f(x)$.
Then
\begin{equation}
\dim(\sh O_x)
 \leq\dim(\sh O_y)
   +\dim\bigl(\sh O_x\otimes_{\sh O_y}\kappa(y)\bigr).
\tag{5.5.2.1}\label{IV.5.5.2.1}
\end{equation}
In particular, if $x$ is a maximal point of the fibre $f^{-1}(y)$,
then
\begin{equation}
\dim(\sh O_x)\leq\dim(\sh O_y),
\tag{5.5.2.2}\label{IV.5.5.2.2}
\end{equation}
because
\[
\sh O_x\otimes_{\sh O_y}\kappa(y)
  =\sh O_x/\mathfrak m_y\sh O_x
\]
is the local ring of $x$ in the prescheme $f^{-1}(y)$ and therefore
has dimension $0$.
\end{proposition}

\oldpage[IV-2]{95}

We shall now obtain a formula more precise than
\eqref{IV.5.5.2.1} when $f$ is assumed to be of finite type.

\begin{proposition}[5.5.3]
\phantomsection
\label{IV.5.5.3}
Let $A$ be a Noetherian ring, let $T$ be an indeterminate, and let
$\mathfrak p$ be a prime ideal of $A$.
Then
\[
\mathfrak p'=\mathfrak pA[T]
\]
is a prime ideal of $B=A[T]$, and
$\mathfrak p'\cap A=\mathfrak p$.
There are infinitely many prime ideals of $B$, distinct from
$\mathfrak p'$, whose intersection with $A$ is $\mathfrak p$; these
prime ideals are pairwise incomparable.
Moreover, if $\mathfrak q$ is any such ideal, then
\begin{equation}
\dim(B_{\mathfrak q})
  =\dim(B_{\mathfrak p'})+1
  =\dim(A_{\mathfrak p})+1.
\tag{5.5.3.1}\label{IV.5.5.3.1}
\end{equation}
\end{proposition}

\begin{proof}
For the first assertions, replacing $A$ by $A/\mathfrak p$ and using
\[
A[T]/\mathfrak pA[T]=(A/\mathfrak p)[T]
\]
reduces us immediately to the case $\mathfrak p=0$.
The prime ideals of $B=A[T]$ whose intersection with $A$ is $0$ are
exactly those disjoint from the multiplicative set
$S=A-\{0\}$ of the domain $A$.
There is an inclusion-preserving bijection from this set of prime
ideals to the prime ideals of
\[
S^{-1}A[T]=K[T],
\]
where $K$ is the fraction field of $A$
(Bourbaki, \emph{Alg\`ebre commutative}, chap.~II, \S~2,
n\textsuperscript{o}~5, prop.~11).
The assertions follow.

Furthermore, by \eqref{IV.5.5.1.1},\footnote{The printed French text
cites \emph{(5.5.1.2)}, which does not exist.  The displayed
inequality used here is \emph{{\hyperlink{ega4.c04.eq.5-5-1-1}{(5.5.1.1)}}}.}
\[
\dim(B_{\mathfrak q})
 \leq\dim(A_{\mathfrak p})
   +\dim(B_{\mathfrak q}/\mathfrak pB_{\mathfrak q}).
\]
Let $k$ be the fraction field of $A/\mathfrak p$.
The ring $B_{\mathfrak q}/\mathfrak pB_{\mathfrak q}$ is canonically
identified with
\[
(k[T])_{\overline{\mathfrak q}},
\]
where $\overline{\mathfrak q}$ is a nonzero prime ideal of $k[T]$.
It is therefore a discrete valuation ring and has dimension $1$.
Thus
\[
\dim(B_{\mathfrak q})\leq\dim(A_{\mathfrak p})+1.
\]

Let
\[
\mathfrak p=\mathfrak p_0
 \supset\mathfrak p_1
 \supset\cdots
 \supset\mathfrak p_m
\]
be a chain of prime ideals of $A$ of maximum length.
The ideals $\mathfrak p_jB$ $(0\leq j\leq m)$ are prime in $B$,
pairwise distinct, and contained in $\mathfrak q$.
Hence
\[
\dim(B_{\mathfrak q})
 \geq m+1
 =\dim(A_{\mathfrak p})+1.
\]
It follows that
\[
\operatorname{ht}(\mathfrak q)
  =\operatorname{ht}(\mathfrak p)+1.
\]
Since $\mathfrak q\neq\mathfrak p'$, we also have
\[
\operatorname{ht}(\mathfrak q)
 \geq\operatorname{ht}(\mathfrak p')+1
 \geq\operatorname{ht}(\mathfrak p)+1.
\]
All these inequalities are equalities, and
\eqref{IV.5.5.3.1} follows.
\end{proof}

\begin{corollary}[5.5.4]
\phantomsection
\label{IV.5.5.4}
For every Noetherian ring $A$,
\begin{equation}
\dim(A[T_1,\ldots,T_r])=\dim(A)+r,
\tag{5.5.4.1}\label{IV.5.5.4.1}
\end{equation}
where the $T_i$ are indeterminates.
\end{corollary}

By contrast, there are examples of non-Noetherian local rings $A$
such that
\[
\dim(A)=1
\qquad\text{and}\qquad
\dim(A[T])=3
\]
[30].

\begin{corollary}[5.5.5]
\phantomsection
\label{IV.5.5.5}
For every locally Noetherian prescheme $X$, the dimension of
\[
X[T_1,\ldots,T_r]
  =X\otimes_{\bb Z}\bb Z[T_1,\ldots,T_r],
\]
where the $T_i$ are indeterminates, is $\dim(X)+r$.
\end{corollary}

\begin{proof}
This follows from \sref{IV.5.5.4} and \sref[0]{0.14.1.7}.
\end{proof}

\begin{corollary}[5.5.6]
\phantomsection
\label{IV.5.5.6}
Under the hypotheses of \sref{IV.5.5.3}, let $\mathfrak q$ be a
prime ideal of $B=A[T]$ such that
$\mathfrak q\cap A=\mathfrak p$.
If $k$ and $k'$ are the residue fields of $A_{\mathfrak p}$ and
$B_{\mathfrak q}$, respectively, then
\begin{equation}
\dim(A_{\mathfrak p})+1
 =\dim(B_{\mathfrak q})+\operatorname{tr.deg}_k k'.
\tag{5.5.6.1}\label{IV.5.5.6.1}
\end{equation}
\end{corollary}

\begin{proof}
\oldpage[IV-2]{96}

If $\mathfrak q=\mathfrak pB$, then
\[
\dim(B_{\mathfrak q})=\dim(A_{\mathfrak p})
\]
by \eqref{IV.5.5.3.1}, while $k'=k(T)$.
Thus \eqref{IV.5.5.6.1} holds.
If $\mathfrak q\neq\mathfrak pB$, then
\[
\dim(B_{\mathfrak q})=\dim(A_{\mathfrak p})+1,
\]
and $\mathfrak q$ corresponds to a nonzero prime ideal of $k[T]$.
Consequently $k'$ is algebraic over $k$, and
\eqref{IV.5.5.6.1} again holds.
\end{proof}

\begin{lemma}[5.5.7]
\phantomsection
\label{IV.5.5.7}
Let $A$ be a Noetherian local domain, let $\mathfrak m$ be its
maximal ideal, and let $k$ be its residue field.
\begin{enumerate}[{\rm(i)}]
\item
\phantomsection
\label{IV.5.5.7.i}
Let $B$ be a domain containing $A$ such that $B=A[x]$ for some
$x\in B$, and let $\mathfrak q$ be a prime ideal of $B$ such that
$\mathfrak q\cap A=\mathfrak m$.
Then
\begin{equation}
\dim(A)+\operatorname{tr.deg}_A B
 \geq\dim(B_{\mathfrak q})+\operatorname{tr.deg}_k k',
\tag{5.5.7.1}\label{IV.5.5.7.1}
\end{equation}
where $k'$ is the residue field of $B_{\mathfrak q}$ and
$\operatorname{tr.deg}_A B$ denotes the transcendence degree of the
fraction field of $B$ over that of $A$.

\item
\phantomsection
\label{IV.5.5.7.ii}
Suppose that for every maximal ideal $\mathfrak n$ of $A[T]$ such
that $\mathfrak n\cap A=\mathfrak m$, the ring
$(A[T])_{\mathfrak n}$ is catenary.
Then the two sides of \eqref{IV.5.5.7.1} are equal.
\end{enumerate}
\end{lemma}

\begin{proof}
\begin{enumerate}[{\rm(i)}]
\item
If $x$ is transcendental over the fraction field of $A$, then
$\operatorname{tr.deg}_A B=1$, and the two sides of
\eqref{IV.5.5.7.1} are equal by \sref{IV.5.5.6}.

Suppose instead that $x$ is algebraic over the fraction field of
$A$.
Then
\[
B=A[T]/\mathfrak p,
\]
where $\mathfrak p$ is a nonzero prime ideal of $A[T]$ such that
$\mathfrak p\cap A=0$, because $B$ contains $A$.
Thus
\[
\operatorname{ht}(\mathfrak p)=1
\]
by \sref{IV.5.5.3}.
The ideal $\mathfrak q$ is of the form
$\mathfrak n/\mathfrak p$, where
$\mathfrak n\supset\mathfrak p$ is a prime ideal of $A[T]$ satisfying
$\mathfrak n\cap A=\mathfrak m$, and
\[
B_{\mathfrak q}
 =(A[T])_{\mathfrak n}/\mathfrak p(A[T])_{\mathfrak n}.
\]
Formula \sref[0]{0.16.1.4.1}, applied to
\[
X=\Spec((A[T])_{\mathfrak n})
\quad\text{and}\quad
Y=\Spec(B_{\mathfrak q}),
\]
gives
\begin{equation}
\begin{split}
\dim((A[T])_{\mathfrak n})
 &\geq
 \operatorname{ht}\bigl(\mathfrak p(A[T])_{\mathfrak n}\bigr)
   +\dim(B_{\mathfrak q}) \\
 &=1+\dim(B_{\mathfrak q}).
\end{split}
\tag{5.5.7.2}\label{IV.5.5.7.2}
\end{equation}
Indeed,
\[
\operatorname{ht}\bigl(\mathfrak p(A[T])_{\mathfrak n}\bigr)
 =\operatorname{ht}(\mathfrak p)=1
\]
by the bijection between prime ideals of $A[T]$ contained in
$\mathfrak n$ and prime ideals of $(A[T])_{\mathfrak n}$.
Finally, \eqref{IV.5.5.6.1} gives
\begin{equation}
\dim((A[T])_{\mathfrak n})
 =\dim(A)+1-\operatorname{tr.deg}_k k',
\tag{5.5.7.3}\label{IV.5.5.7.3}
\end{equation}
because $B_{\mathfrak q}$ and $(A[T])_{\mathfrak n}$ have the same
residue field.
Here $\operatorname{tr.deg}_A B=0$, and
\eqref{IV.5.5.7.1} follows.

\item
If $(A[T])_{\mathfrak n}$ is catenary, the two sides of
\eqref{IV.5.5.7.2} are equal \sref[0]{0.16.1.4}; hence the two sides
of \eqref{IV.5.5.7.1} are equal as well.
\end{enumerate}
\end{proof}

\begin{theorem}[5.5.8 (dimension formula)]
\phantomsection
\label{IV.5.5.8}
Let $A$ be a Noetherian local domain, let $B$ be a domain containing
$A$ and finitely generated as an $A$-algebra, and let $\mathfrak q$
be a prime ideal of $B$ such that
$\mathfrak q\cap A$ is the maximal ideal $\mathfrak m$ of $A$.
Let $k$ and $k'$ be the residue fields of $A$ and
$B_{\mathfrak q}$, respectively.
Then
\begin{equation}
\dim(A)+\operatorname{tr.deg}_A B
 \geq\dim(B_{\mathfrak q})+\operatorname{tr.deg}_k k'.
\tag{5.5.8.1}\label{IV.5.5.8.1}
\end{equation}
Moreover, the two sides are equal if, for every finitely generated
$A$-subalgebra $A'$ of $B[T]$ and every maximal ideal
$\mathfrak m'$ of $A'$ such that
$\mathfrak m'\cap A=\mathfrak m$, the local ring
$A'_{\mathfrak m'}$ is catenary.
\end{theorem}

\begin{proof}
\oldpage[IV-2]{97}

Write
\[
B=A[x_1,\ldots,x_n]
\]
and argue by induction on $n$.
Put
\[
C=A[x_1,\ldots,x_{n-1}],
\qquad
\mathfrak r=\mathfrak q\cap C.
\]
The ring $C_{\mathfrak r}$ is a Noetherian local domain.
Let
\[
S=C-\mathfrak r,\qquad
B'=S^{-1}B,\qquad
\mathfrak q'=S^{-1}\mathfrak q.
\]
Then
\[
B_{\mathfrak q}=B'_{\mathfrak q'},
\qquad
B'=C_{\mathfrak r}[x_n],
\qquad
\mathfrak q'\cap C_{\mathfrak r}
  =\mathfrak rC_{\mathfrak r}.
\]
Moreover, the fraction fields of $B'$ and $C_{\mathfrak r}$ are
those of $B$ and $C$, respectively.
If $k_1$ is the residue field of $C_{\mathfrak r}$, then
\eqref{IV.5.5.7.1} gives
\begin{equation}
\dim(C_{\mathfrak r})+\operatorname{tr.deg}_C B
 \geq\dim(B_{\mathfrak q})+\operatorname{tr.deg}_{k_1}k'.
\tag{5.5.8.2}\label{IV.5.5.8.2}
\end{equation}
The induction hypothesis gives
\begin{equation}
\dim(A)+\operatorname{tr.deg}_A C
 \geq\dim(C_{\mathfrak r})+\operatorname{tr.deg}_k k_1.
\tag{5.5.8.3}\label{IV.5.5.8.3}
\end{equation}
Adding \eqref{IV.5.5.8.2} and \eqref{IV.5.5.8.3} term by term yields
\eqref{IV.5.5.8.1}.

For the second assertion, first note that every finitely generated
$A$-subalgebra of $C[T]$ is also a finitely generated
$A$-subalgebra of $B[T]$.
By induction on $n$, we may therefore suppose that the two sides of
\eqref{IV.5.5.8.3} are equal.

By \sref{IV.5.5.7}, to prove equality in
\eqref{IV.5.5.8.2}, it is enough to show that whenever
$\mathfrak n$ is a maximal ideal of $C_{\mathfrak r}[T]$ satisfying
\[
\mathfrak n\cap C_{\mathfrak r}
  =\mathfrak rC_{\mathfrak r},
\]
the ring $(C_{\mathfrak r}[T])_{\mathfrak n}$ is catenary.
We have
\[
C_{\mathfrak r}[T]=S^{-1}C[T].
\]
Thus $\mathfrak n=S^{-1}\mathfrak n'$ for a prime ideal
$\mathfrak n'$ of $C[T]$ such that
$\mathfrak n'\cap C=\mathfrak r$, and
\[
(C_{\mathfrak r}[T])_{\mathfrak n}
  =(C[T])_{\mathfrak n'}.
\]
Choose a maximal ideal $\mathfrak m'$ of $C[T]$ containing
$\mathfrak n'$.
Then $\mathfrak m'\cap A=\mathfrak m$, and
$(C[T])_{\mathfrak m'}$ is catenary by hypothesis.
Since $(C[T])_{\mathfrak n'}$ is a localization of this ring, it too
is catenary \sref[0]{0.16.1.4}.
\end{proof}
% END NATIVE INPUT components/05_section55_dimension_formula_finite_type
% BEGIN NATIVE INPUT components/06_section56_dimension_formula_universally_catenary
\subsection{Dimension formula and universally catenary rings}
\phantomsection
\label{subsection:IV.5.6}
\label{IV.5.6}

\oldpage[IV-2]{97}

\begin{proposition}[5.6.1]
\phantomsection
\label{IV.5.6.1}
Let $A$ be a Noetherian ring.  The following properties are
equivalent:
\begin{enumerate}[{\rm(a)}]
\item
\phantomsection
\label{IV.5.6.1.a}
Every polynomial ring $A[T_1,\ldots,T_n]$ is catenary.

\item
\phantomsection
\label{IV.5.6.1.b}
Every $A$-algebra of finite type is catenary.

\item
\phantomsection
\label{IV.5.6.1.c}
The ring $A$ is catenary, and the following dimension formula holds.
Let $B'$ be any integral local $A$-algebra essentially of finite type
\sref{IV.1.3.8}; let $\mathfrak q$ be its maximal ideal and
$\mathfrak p$ the inverse image of $\mathfrak q$ in $A$; let $A'$
be the image of $A_{\mathfrak p}$ in $B'$; and let $k$ and $k'$ be
the residue fields of $A'$ and $B'$, respectively.  Then
\begin{equation}
\dim(A')+\operatorname{tr.deg}_{A'}B'
 =\dim(B')+\operatorname{tr.deg}_k k'.
\tag{5.6.1.1}\label{IV.5.6.1.1}
\end{equation}
\end{enumerate}
\end{proposition}

\begin{proof}
Suppose first that condition
{\hyperlink{ega4.c04.sec.5-6-1-b}{\textup{(b)}}} holds.  Then $B'$ is a local
$A'$-algebra essentially of finite type \sref{IV.1.3.10}; hence
\[
B'=B''_{\mathfrak q''},
\]
where $B''$ is a finitely generated $A'$-subalgebra of $B'$ and
$\mathfrak q''$ is a prime ideal of $B''$ lying over the maximal
ideal of $A'$.  Moreover, $B'$ and $B''$ have the same fraction
field, so
\[
\operatorname{tr.deg}_{A'}B'
 =\operatorname{tr.deg}_{A'}B''.
\]
To prove \eqref{IV.5.6.1.1}, it is enough to show that every
finitely generated $A'$-subalgebra $A_1$ of $B'[T]$ is catenary.
Indeed, the two sides of \eqref{IV.5.5.8.1}, applied to $A'$, $B''$,
and $\mathfrak q''$, are then equal.  But
Condition {\hyperlink{ega4.c04.sec.5-6-1-b}{\textup{(b)}}} implies that every
$A$-algebra essentially of
finite type is catenary \sref[0]{0.16.1.4}, and $A_1$ is such an
$A$-algebra \sref{IV.1.3.9}.  Thus (b) implies (c).

It is immediate that (b) implies (a).  Conversely, (a) implies (b),
because every $A$-algebra of finite type is a quotient of a
polynomial algebra \sref[0]{0.16.1.4}.

It remains to prove that (c) implies (b).  Since every quotient by a
prime ideal of an $A$-algebra of finite type is an integral
$A$-algebra of finite type, it is enough to prove the following.
Let $B$ be an integral $A$-algebra of finite type, and let
$\mathfrak q'\subset\mathfrak q$ be prime ideals of $B$.  Then
\begin{equation}
\dim(B_{\mathfrak q}/\mathfrak q'B_{\mathfrak q})
 +\dim(B_{\mathfrak q'})
 =\dim(B_{\mathfrak q}).
\tag{5.6.1.2}\label{IV.5.6.1.2}
\end{equation}
By \sref[0]{0.16.1.4.2}, this is precisely the required catenarity
condition.

\oldpage[IV-2]{98}

Let $\mathfrak p$ and $\mathfrak p'$ be the inverse images in $A$ of
$\mathfrak q$ and $\mathfrak q'$, respectively, and let
$\mathfrak n$ be the kernel of
\[
A_{\mathfrak p}\longrightarrow B_{\mathfrak q}.
\]
The image of $A_{\mathfrak p}$ in $B_{\mathfrak q}$ is isomorphic to
$A_{\mathfrak p}/\mathfrak n$.  Applying
\eqref{IV.5.6.1.1} to this image and to $B_{\mathfrak q}$ gives
\begin{equation}
\begin{split}
\dim(A_{\mathfrak p}/\mathfrak n)
 &+\operatorname{tr.deg}_{A_{\mathfrak p}/\mathfrak n}
       B_{\mathfrak q}\\
 &=\dim(B_{\mathfrak q})
   +\operatorname{tr.deg}_{\kappa(\mathfrak p)}
       \kappa(\mathfrak q).
\end{split}
\tag{5.6.1.3}\label{IV.5.6.1.3}
\end{equation}

The kernel of
$A_{\mathfrak p'}\to B_{\mathfrak q'}$ is
$\mathfrak nA_{\mathfrak p'}$.  Applying
\eqref{IV.5.6.1.1} again gives
\begin{equation}
\begin{split}
\dim(A_{\mathfrak p'}/\mathfrak nA_{\mathfrak p'})
 &+\operatorname{tr.deg}_{
      A_{\mathfrak p'}/\mathfrak nA_{\mathfrak p'}}
      B_{\mathfrak q'}\\
 &=\dim(B_{\mathfrak q'})
   +\operatorname{tr.deg}_{\kappa(\mathfrak p')}
      \kappa(\mathfrak q').
\end{split}
\tag{5.6.1.4}\label{IV.5.6.1.4}
\end{equation}

Finally, $B/\mathfrak q'$ is an integral $A$-algebra of finite type.
The inverse image of $\mathfrak q/\mathfrak q'$ in $A$ is
$\mathfrak p$, and the kernel of
\[
A_{\mathfrak p}\longrightarrow
 B_{\mathfrak q}/\mathfrak q'B_{\mathfrak q}
\]
is $\mathfrak p'A_{\mathfrak p}$.  Hence
\begin{equation}
\begin{split}
\dim(A_{\mathfrak p}/\mathfrak p'A_{\mathfrak p})
 &+\operatorname{tr.deg}_{\kappa(\mathfrak p')}
   (B_{\mathfrak q}/\mathfrak q'B_{\mathfrak q})\\
 &=\dim(B_{\mathfrak q}/\mathfrak q'B_{\mathfrak q})
   +\operatorname{tr.deg}_{\kappa(\mathfrak p)}
      \kappa(\mathfrak q).
\end{split}
\tag{5.6.1.5}\label{IV.5.6.1.5}
\end{equation}

Add \eqref{IV.5.6.1.4} and \eqref{IV.5.6.1.5} term by term.
The fields $\kappa(\mathfrak p')$ and
$\kappa(\mathfrak q')$ are the fraction fields of
$A_{\mathfrak p}/\mathfrak p'A_{\mathfrak p}$ and
$B_{\mathfrak q}/\mathfrak q'B_{\mathfrak q}$, respectively.
Furthermore,
$A_{\mathfrak p'}/\mathfrak nA_{\mathfrak p'}$ and
$A_{\mathfrak p}/\mathfrak n$ have the same fraction field, as do
$B_{\mathfrak q'}$ and $B_{\mathfrak q}$.  Since $A$ is catenary,
\sref[0]{0.16.1.4.2} gives
\[
\dim(A_{\mathfrak p}/\mathfrak p'A_{\mathfrak p})
 +\dim(A_{\mathfrak p'}/\mathfrak nA_{\mathfrak p'})
 =\dim(A_{\mathfrak p}/\mathfrak n).
\]
Taking these observations and \eqref{IV.5.6.1.3} into account yields
\eqref{IV.5.6.1.2}.
\end{proof}

\begin{definition}[5.6.2]
\phantomsection
\label{IV.5.6.2}
A Noetherian ring is said to be \emph{universally catenary} if it
satisfies the equivalent conditions of \sref{IV.5.6.1}.
\end{definition}

\begin{remarks}[5.6.3]
\phantomsection
\label{IV.5.6.3}
\begin{enumerate}[{\rm(i)}]
\item
\phantomsection
\label{IV.5.6.3.i}
If $A$ is universally catenary, then so is $S^{-1}A$ for every
multiplicative subset $S$ of $A$.  This follows at once from
{\hyperlink{ega4.c04.sec.5-6-1-a}{5.6.1\textup{(a)}}} and the fact that every ring
of fractions of a
catenary ring is catenary.

Conversely, suppose that $A_{\mathfrak p}$ is universally catenary
for every prime ideal $\mathfrak p$ of $A$.  Then $A$ is universally
catenary.  Indeed, apply
{\hyperlink{ega4.c04.sec.5-6-1-c}{5.6.1\textup{(c)}}}; if
$S=A-\mathfrak p$, then $S^{-1}B$ is an
$A_{\mathfrak p}$-algebra of finite type and
\[
B_{\mathfrak q}=(S^{-1}B)_{S^{-1}\mathfrak q}.
\]

\item
\phantomsection
\label{IV.5.6.3.ii}
A locally Noetherian prescheme $X$ is said to be
\emph{universally catenary} if $\sh O_x$ is universally catenary for
every $x\in X$.  By (i), a Noetherian ring $A$ is universally
catenary if and only if $\Spec(A)$ is universally catenary.

\item
\phantomsection
\label{IV.5.6.3.iii}
The criterion {\hyperlink{ega4.c04.sec.5-6-1-c}{5.6.1\textup{(c)}}} involves only
the images of $A$ in
integral $A$-algebras of finite type, hence only integral quotient
rings of $A$.

\oldpage[IV-2]{99}

It follows that if
$\mathfrak p_i$ $(1\leq i\leq n)$ are the minimal prime ideals of
$A$, then $A$ is universally catenary if and only if every
$A/\mathfrak p_i$ is universally catenary.  Likewise, a locally
Noetherian prescheme $X$ is universally catenary if and only if
$X_{\mathrm{red}}$ is universally catenary.

\item
\phantomsection
\label{IV.5.6.3.iv}
The criterion {\hyperlink{ega4.c04.sec.5-6-1-b}{5.6.1\textup{(b)}}} and (i) show
that every $A$-algebra
essentially of finite type over a universally catenary ring $A$ is
universally catenary \sref{IV.1.3.8}.
\end{enumerate}
\end{remarks}

\begin{proposition}[5.6.4]
\phantomsection
\label{IV.5.6.4}
Every quotient ring of a regular ring is universally catenary.
\end{proposition}

\begin{proof}
By {\hyperlink{ega4.c04.sec.5-6-3-iv}{5.6.3\textup{(iv)}}}, it is enough to show
that every regular ring
$A$ is universally catenary.  Each polynomial ring
$A[T_1,\ldots,T_n]$ is regular \sref[0]{0.17.3.7}, hence catenary
\sref[0]{0.16.5.12}, and the assertion follows from
{\hyperlink{ega4.c04.sec.5-6-1-a}{5.6.1\textup{(a)}}}.
\end{proof}

In particular, Cohen's theorem \sref[0]{0.19.8.8} implies that every
complete Noetherian local ring is universally catenary.  Likewise,
every algebra of finite type over a field is universally catenary,
and every prescheme locally of finite type over a field is
universally catenary.

We shall see in \sref{IV.6.3.7} that ``regular ring'' may be replaced
by ``Cohen--Macaulay ring'' in \sref{IV.5.6.4}.

\begin{proposition}[5.6.5]
\phantomsection
\label{IV.5.6.5}
Let $Y$ be an irreducible locally Noetherian prescheme, let $X$ be
an irreducible prescheme, and let $f:X\to Y$ be a dominant morphism
locally of finite type.  Let $\xi$ and $\eta$ be the generic points
of $X$ and $Y$, respectively, and put
\[
e=\dim(f^{-1}(\eta))
  =\operatorname{tr.deg}_{\kappa(\eta)}\kappa(\xi),
\]
the dimension of the generic fibre
\sref[0]{0.2.1.8}, \sref{IV.4.1.1}.
For every $x\in X$, with $y=f(x)$, one has
\begin{equation}
e+\dim(\sh O_y)
 \geq \operatorname{tr.deg}_{\kappa(y)}\kappa(x)
       +\dim(\sh O_x)
\tag{5.6.5.1}\label{IV.5.6.5.1}
\end{equation}
and
\begin{equation}
\dim(\sh O_x)
 \leq \dim(\sh O_y)
      +\dim(\sh O_x\otimes_{\sh O_y}\kappa(y))
      -\delta(x),
\tag{5.6.5.2}\label{IV.5.6.5.2}
\end{equation}
where
\[
\delta(x)=\dim_x(f^{-1}(y))-e.
\]
If $Y$ is universally catenary, equality holds in
\eqref{IV.5.6.5.1} and \eqref{IV.5.6.5.2}.  If, in addition, $x$ is
closed in $f^{-1}(y)$, then
\begin{equation}
\dim(\sh O_x)=\dim(\sh O_y)+e.
\tag{5.6.5.3}\label{IV.5.6.5.3}
\end{equation}
\end{proposition}

\begin{proof}
We may restrict to the case in which $X$ and $Y$ are affine, and
hence $f$ is of finite type, and then replace both preschemes by
their reductions \sref{IV.1.5.4}.  Thus $X$ and $Y$ may be assumed
integral.

Since $f$ is dominant, $\sh O_y$ identifies with a subring of
$\sh O_x$, and their fraction fields are
$\kappa(\eta)$ and $\kappa(\xi)$, respectively
\sref[I]{I.8.2.7}.  Moreover, $\sh O_x$ is the localization at a
prime ideal $\mathfrak q$ of an integral $\sh O_y$-algebra $B$ of
finite type \sref[I]{I.3.6.5}.  Applying
\eqref{IV.5.5.8.1} to $A=\sh O_y$, $B$, and $\mathfrak q$ gives
\eqref{IV.5.6.5.1}.

Furthermore,
\[
\sh O_x\otimes_{\sh O_y}\kappa(y)
 =\sh O_x/\mathfrak m_y\sh O_x
\]
is the local ring of the fibre $f^{-1}(y)$ at $x$
\sref[I]{I.3.6.1}.  Since the fibre is of finite type over
$\kappa(y)$, \eqref{IV.5.2.3.1} gives
\[
\dim(\sh O_x\otimes_{\sh O_y}\kappa(y))
 =\dim_x(f^{-1}(y))
  -\operatorname{tr.deg}_{\kappa(y)}\kappa(x).
\]
Thus \eqref{IV.5.6.5.2} is simply another form of
\eqref{IV.5.6.5.1}.

\oldpage[IV-2]{100}

When $Y$ is universally catenary, equality in
\eqref{IV.5.6.5.1} is exactly \eqref{IV.5.6.1.1}, applied to
$A=\sh O_y$ and $B'=\sh O_x$.  Finally, if $x$ is closed in the
fibre, then in general
\[
\operatorname{tr.deg}_{\kappa(y)}\kappa(x)
 =\dim(\overline{\{x\}}\cap f^{-1}(y));
\]
this follows from \sref{IV.5.2.1}, applied to the reduced closed
sub-prescheme of $f^{-1}(y)$ whose underlying space is the displayed
subspace.  Equation \eqref{IV.5.6.5.3} follows.
\end{proof}

We shall prove in \sref{IV.13.1.1} that under the hypotheses of
\sref{IV.5.6.5} one always has $\delta(x)\geq0$.  Thus in this case
\eqref{IV.5.6.5.2} sharpens \eqref{IV.5.5.2.1}.

\begin{corollary}[5.6.6]
\phantomsection
\label{IV.5.6.6}
Under the general hypotheses of \sref{IV.5.6.5},
\begin{equation}
\dim(X)\leq\dim(Y)+e.
\tag{5.6.6.1}\label{IV.5.6.6.1}
\end{equation}
If $Y$ is universally catenary, equality holds in
\eqref{IV.5.6.6.1} if and only if
\begin{equation}
\dim(Y)=\sup_{y\in f(X)}\dim(\sh O_y).
\tag{5.6.6.2}\label{IV.5.6.6.2}
\end{equation}
In particular, equality holds in \eqref{IV.5.6.6.1} when $Y$ is
locally of finite type over a field.
\end{corollary}

\begin{proof}
Apply \eqref{IV.5.6.5.1} at a closed point $x$ of $X$, and use
\sref{IV.5.1.4.2} and \sref{IV.5.1.11}; this gives
\eqref{IV.5.6.6.1}.  Conversely, every nonempty fibre $f^{-1}(y)$
contains a point closed in that fibre \sref{IV.5.1.11}.  If $Y$ is
universally catenary, \eqref{IV.5.6.5.3} holds at such a point.
Taking suprema in \eqref{IV.5.6.5.3} as $x$ ranges over $X$, and
observing that every point closed in $X$ is \emph{a fortiori} closed
in $f^{-1}(f(x))$, yields the second assertion by
\sref{IV.5.1.4.1} and \sref{IV.5.1.4.2}.

For the last assertion, choose an affine open neighbourhood $U$ of
the generic point of $X$ such that $f|_U$ is of finite type.  Then
$f(U)$ is constructible \sref{IV.1.8.5} and dense in $Y$, and
therefore contains a nonempty open, hence dense, subset of $Y$
\sref[0]{0.9.2.2}.  Since $Y$ is locally of finite type over a
field, it is universally catenary \sref{IV.5.6.4}; moreover, the two
sides of \eqref{IV.5.6.6.2} are equal by
\sref{IV.5.2.2} and \sref{IV.4.1.1.3}.

Condition \eqref{IV.5.6.6.2} is trivially satisfied when $f$ is
surjective.
\end{proof}

\begin{corollary}[5.6.7]
\phantomsection
\label{IV.5.6.7}
Let $Y$ be a locally Noetherian prescheme, let $f:X\to Y$ be a
morphism locally of finite type, and let $n$ be an integer.  If
\[
\dim(f^{-1}(y))\leq n
\]
for every $y\in Y$, then
\begin{equation}
\dim(X)\leq\dim(Y)+n.
\tag{5.6.7.1}\label{IV.5.6.7.1}
\end{equation}
\end{corollary}

\begin{proof}
By \sref[0]{0.14.1.7} and \sref{IV.1.5.4}, we may assume that $X$
and $Y$ are affine and reduced, so that $f$ is of finite type.
Let $X_i$ $(1\leq i\leq m)$ be the integral closed sub-preschemes
whose underlying spaces are the irreducible components of $X$.
Then
\[
\dim(X)=\sup_i\dim(X_i).
\]
Let $Z_i$ be the reduced closed sub-prescheme of $Y$ whose
underlying space is $\overline{f(X_i)}$.  It is integral
\sref[0]{0.2.1.5}, and the restriction of $f$ factors as
\[
X_i\xrightarrow{g_i}Z_i\xrightarrow{j_i}Y,
\]
where $j_i$ is the canonical injection \sref[I]{I.5.2.2}, and $g_i$
is dominant and of finite type \sref{IV.1.5.4}.

\oldpage[IV-2]{101}

We have $\dim(Z_i)\leq\dim(Y)$.  If $\zeta_i$ is the generic point
of $Z_i$, then
\[
\dim(g_i^{-1}(\zeta_i))\leq n
\]
by hypothesis.  Thus it is enough to show
\[
\dim(X_i)\leq\dim(Z_i)+n
\]
for every $i$, and this is \eqref{IV.5.6.6.1}.
\end{proof}

\begin{corollary}[5.6.8]
\phantomsection
\label{IV.5.6.8}
Let $Y$ be a locally Noetherian prescheme, let $X$ be irreducible,
let $f:X\to Y$ be a morphism locally of finite type, and let
$n\geq0$ be an integer.  Suppose that $Y$ is universally catenary
and that
\[
\dim(f^{-1}(y))\geq n
\quad\text{for every }y\in Y.
\]
Then
\begin{equation}
\dim(X)\geq\dim(Y)+n.
\tag{5.6.8.1}\label{IV.5.6.8.1}
\end{equation}
If instead every fibre has dimension exactly $n$, then
\begin{equation}
\dim(X)=\dim(Y)+n.
\tag{5.6.8.2}\label{IV.5.6.8.2}
\end{equation}
\end{corollary}

\begin{proof}
Since $n\geq0$, either hypothesis implies that $f$ is surjective,
and therefore $Y$ is irreducible.  If $\eta$ is the generic point of
$Y$, then $\dim(f^{-1}(\eta))\geq n$, respectively
$\dim(f^{-1}(\eta))=n$.  The conclusion follows from
\sref{IV.5.6.6}.
\end{proof}

\begin{remarks}[5.6.9]
\phantomsection
\label{IV.5.6.9}
\begin{enumerate}[{\rm(i)}]
\item
\phantomsection
\label{IV.5.6.9.i}
Even if $Y$ is regular, $X$ is irreducible, and $f:X\to Y$ is
dominant and of finite type, the two sides of
\eqref{IV.5.6.6.1} need not be equal.  Take
$Y=\Spec(A)$ for a discrete valuation ring $A$, take
$X=\Spec(K)$ for its fraction field $K$, and let $f$ be the
canonical morphism.

\item
\phantomsection
\label{IV.5.6.9.ii}
Example {\hyperlink{ega4.c04.sec.5-4-3-i}{5.4.3\textup{(i)}}} shows that the
hypothesis that $X$ is
irreducible cannot be omitted from \sref{IV.5.6.6} or
\sref{IV.5.6.8}, even when all the other hypotheses hold.  We shall
see in \sref{IV.10.6.1} that such phenomena cannot occur under
additional hypotheses on $Y$, satisfied for example when $Y$ is
locally of finite type over a field or over a Dedekind ring having
infinitely many prime ideals.
\end{enumerate}
\end{remarks}

\begin{proposition}[5.6.10]
\phantomsection
\label{IV.5.6.10}
Let $A$ be an integral universally catenary Noetherian local ring,
and let $B$ be an integral finite $A$-algebra containing $A$.  For
every maximal ideal $\mathfrak n$ of $B$,
\[
\dim(B_{\mathfrak n})=\dim(A).
\]
\end{proposition}

\begin{proof}
We have $\operatorname{tr.deg}_A B=0$, and the residue field of
$B_{\mathfrak n}$ is algebraic over the residue field of
$A$.\footnote{The printed French text says ``the fraction field of
$A$''.  Here the field required by the dimension formula is the
residue field of $A$.}  Every maximal ideal of $B$ lies over the
maximal ideal of $A$, so the assertion follows from
\eqref{IV.5.6.1.1}.
\end{proof}

\begin{example}[5.6.11]
\phantomsection
\label{IV.5.6.11}
Let $A$ be an integral, integrally closed Noetherian local ring.
Then \sref[0]{0.16.1.6} shows that the conclusion of
\sref{IV.5.6.10} holds for every integral finite $A$-algebra $B$
containing $A$.  We now construct, by contrast, an integral
catenary Noetherian local ring $A$ for which that conclusion fails.

Use the construction of {\hyperlink{ega4.c04.sec.5-2-5-i}{5.2.5\textup{(i)}}}.
Let $k_0$ be a field,
let
\[
k=k_0(S_i)_{i\in\bb N}
\]
be a purely transcendental extension of infinite transcendence
degree, and let $V=k[S]_{(S)}$ be the discrete valuation ring
obtained by localizing $k[S]$ at $(S)$.  Put $E=V[T]$.  In $E$, the
maximal ideal
\[
\mathfrak m=(S)+(T)
\]
has height $2$, whereas the maximal ideal
\[
\mathfrak m'=(ST-1)
\]
has height $1$.

\oldpage[IV-2]{102}

The corresponding residue fields are
\[
\kappa(\mathfrak m)=k,
\qquad
\kappa(\mathfrak m')=k(S),
\]
the latter being the fraction field of $V$.  By the choice of $k$,
these fields are isomorphic.  Let
\[
\varepsilon:E\longrightarrow E/\mathfrak m,
\qquad
\varepsilon':E\longrightarrow E/\mathfrak m'
\]
be the canonical homomorphisms, choose an isomorphism
\[
\sigma:\kappa(\mathfrak m)\xrightarrow{\sim}\kappa(\mathfrak m'),
\]
and let $C$ be the subring of $E$ formed by the elements $x$ such
that
\[
\varepsilon'(x)=\sigma(\varepsilon(x)).
\]
This is a special case of the ``gluing procedures'' to be studied
systematically in Chapter~V.

It is immediate that
\[
\mathfrak n=\mathfrak m\mathfrak m'
 =\mathfrak m\cap\mathfrak m'
\]
is a maximal ideal of $C$, and that $C/\mathfrak n$ identifies with
the subfield of
$(E/\mathfrak m)\times(E/\mathfrak m')$ formed by the pairs
$(z,\sigma(z))$.  Moreover,
\[
E=C+C(ST-1),
\]
so $E$ is a finite $C$-algebra; it is plainly the integral closure
of $C$.  We next prove that $C$ is Noetherian.

\begin{lemma}[5.6.11.1]
\phantomsection
\label{IV.5.6.11.1}
Let $R$ be a ring, let $S$ be a subring, and let
\[
\mathfrak K=\operatorname{Ann}_S(R/S)
\]
be the conductor of $S$ in $R$, that is, the largest ideal of $S$
which is also an ideal of $R$.
\begin{enumerate}[{\rm(i)}]
\item
\phantomsection
\label{IV.5.6.11.1.i}
For every ideal $\mathfrak J\subset\mathfrak K$ of $R$, there is a
strictly increasing map from the set of ideals $\mathfrak a$ of $S$
such that $R\mathfrak a=\mathfrak J$ to the set of
$(S/\mathfrak K)$-submodules of
$\mathfrak J/\mathfrak K\mathfrak J$.

\item
\phantomsection
\label{IV.5.6.11.1.ii}
If $S/\mathfrak K$ and $R$ are Noetherian and $R$ is a finite
$S$-module, then every increasing sequence of ideals of $S$
contained in $\mathfrak K$ is stationary.
\end{enumerate}
\end{lemma}

\begin{proof}
For (i), if $R\mathfrak a=\mathfrak J$, then
\[
\mathfrak K\mathfrak J
 =\mathfrak K(R\mathfrak a)
 =\mathfrak K\mathfrak a
 \subset\mathfrak a.
\]
Thus $\mathfrak a\mapsto
\mathfrak a/\mathfrak K\mathfrak J$ gives the asserted strictly
increasing map.

For (ii), let $(\mathfrak L_r)$ be an increasing sequence of ideals
of $S$ contained in $\mathfrak K$.  Since $R$ is Noetherian, the
sequence $(R\mathfrak L_r)$ is stationary; after discarding finitely
many terms, suppose that all its members are the same ideal
$\mathfrak J$ of $R$.  The module
$\mathfrak J/\mathfrak K\mathfrak J$ is finite over $R$, hence over
$S$, and therefore finite over the Noetherian ring
$S/\mathfrak K$.  It is consequently Noetherian, and (i) shows that
$(\mathfrak L_r)$ is stationary.
\end{proof}

Apply the lemma with $R=E$ and $S=C$.  The ideal $\mathfrak n$ is
the conductor of $C$ in $E$.  For every ideal $\mathfrak a$ of $C$,
\[
\mathfrak a/(\mathfrak n\cap\mathfrak a)
 \simeq(\mathfrak a+\mathfrak n)/\mathfrak n
\]
is a submodule of the simple $C$-module $C/\mathfrak n$.  It is
therefore enough to prove that every ideal
$\mathfrak a\cap\mathfrak n$ is finitely generated; this follows
from \sref{IV.5.6.11.1}.  Thus $C$ is Noetherian.

By \sref[0]{0.16.1.5},
\[
\dim(C)=\dim(E)=2.
\]
Set $A=C_{\mathfrak n}$ and $U=C-\mathfrak n$.  Then
\[
B=U^{-1}E
\]
is the integral closure of $A$ and is a finite $A$-module.  The
ideals $\mathfrak m$ and $\mathfrak m'$ are the only prime ideals of
$E$ containing $\mathfrak n$.  Hence $B$ is semilocal, and its local
rings at its two maximal ideals are isomorphic to
$E_{\mathfrak m}$ and $E_{\mathfrak m'}$, of dimensions $2$ and
$1$, respectively.  Thus $\dim(B)=2$ and, by
\sref[0]{0.16.1.5}, $\dim(A)=2$.

Since $A$ is an integral local ring of dimension $2$, it is
catenary: every prime ideal other than $0$ and the maximal ideal has
height $1$.  Nevertheless it does not satisfy the conclusion of
\sref{IV.5.6.10}, and therefore is not universally catenary.

The ideal $\mathfrak n$ has height $2$ in $C$.  If $\mathfrak p$ is
a height-$1$ prime ideal of $C$, there is a unique prime ideal
$\mathfrak p'$ of $E$ above $\mathfrak p$; it has height $1$, and
\[
C_{\mathfrak p}=E_{\mathfrak p'}.
\]
Existence follows from integrality, and the Cohen--Seidenberg theorem
shows that every prime ideal of $E$ above $\mathfrak p$ has height
$1$.

\oldpage[IV-2]{103}

Moreover, $E$ is a finite $C$-algebra and
$C_{\mathfrak p}$ is integrally closed
(Bourbaki, \emph{Alg\`ebre commutative}, chap.~V, \S~1,
n\textsuperscript{o}~5, cor.~5 of prop.~16).  This proves uniqueness
and the displayed equality.

It would be interesting to know whether every integral Noetherian
local ring satisfying the conclusion of \sref{IV.5.6.10} is
universally catenary.  Nagata asserted this [33], but his proof does
not appear to be complete.
\end{example}
% END NATIVE INPUT components/06_section56_dimension_formula_universally_catenary
% BEGIN NATIVE INPUT components/07_section57_depth_and_property_sk
\subsection{Depth and property \texorpdfstring{$(S_k)$}{(Sk)}}
\phantomsection
\label{subsection:IV.5.7}

\oldpage[IV-2]{103}

\begin{definition}[5.7.1]
\phantomsection
\label{IV.5.7.1}
Let $X$ be a locally Noetherian prescheme and let $\sh F$ be a
coherent $\sh O_X$-module.  The \emph{depth} (respectively,
\emph{codepth}) of $\sh F$ at a point $x\in X$ is the number
$\operatorname{prof}(\sh F_x)$ (respectively,
$\operatorname{coprof}(\sh F_x)$)
\sref[0]{0.16.4.5}, \sref[0]{0.16.4.9}.  The \emph{codepth} of
$\sh F$ is the number
\begin{equation}
\operatorname{coprof}(\sh F)
 =\sup_{x\in X}\operatorname{coprof}(\sh F_x).
\tag{5.7.1.1}\label{IV.5.7.1.1}
\end{equation}

The module $\sh F$ is said to be a
\emph{Cohen--Macaulay $\sh O_X$-module at $x$} if $\sh F_x$ is a
Cohen--Macaulay $\sh O_x$-module, that is
\sref[0]{0.16.5.1}, if
$\operatorname{coprof}(\sh F_x)=0$.  It is said to be a
\emph{Cohen--Macaulay $\sh O_X$-module} if it is so at every point,
or equivalently if $\operatorname{coprof}(\sh F)=0$.  A point
$x\in X$ such that $\sh O_x$ is a Cohen--Macaulay ring is also
called a \emph{Cohen--Macaulay point} of $X$.

The \emph{codepth of $X$}, denoted by $\operatorname{coprof}(X)$,
is the number $\operatorname{coprof}(\sh O_X)$.  The prescheme $X$
is said to be \emph{Cohen--Macaulay} if $\sh O_X$ is a
Cohen--Macaulay $\sh O_X$-module, or equivalently if
$\operatorname{coprof}(X)=0$.  Every locally Noetherian prescheme of
dimension $0$ is evidently Cohen--Macaulay.  To say that
$\Spec(A)$ is Cohen--Macaulay means that $A$ is a Cohen--Macaulay
ring \sref[0]{0.16.5.13}.
\end{definition}

\begin{definition}[5.7.2]
\phantomsection
\label{IV.5.7.2}
Let $X$ be a locally Noetherian prescheme, let $\sh F$ be a coherent
$\sh O_X$-module, and let $k$ be an integer, positive or negative.
The module $\sh F$ is said to satisfy property $(S_k)$ if, for every
$x\in X$,
\begin{equation}
\operatorname{prof}(\sh F_x)
 \geq \inf\{k,\dim(\sh F_x)\}.
\tag{5.7.2.1}\label{IV.5.7.2.1}
\end{equation}
It is said to satisfy property $(S_k)$ \emph{at a point $x\in X$}
if, for every generization $x'$ of $x$ in $X$,
\begin{equation}
\operatorname{prof}(\sh F_{x'})
 \geq \inf\{k,\dim(\sh F_{x'})\}.
\tag{5.7.2.2}\label{IV.5.7.2.2}
\end{equation}

The prescheme $X$ is said to satisfy property $(S_k)$
(respectively, property $(S_k)$ at $x$) if $\sh O_X$ does so
(respectively, does so at $x$).

The module $\sh F$ satisfies $(S_k)$ if and only if it satisfies
$(S_k)$ at every point of $X$.  If $U$ is an open subset of $X$ and
$\sh F$ satisfies $(S_k)$, then $\sh F|U$ satisfies $(S_k)$.
Conversely, if $(U_\alpha)$ is an open covering of $X$ and
$\sh F|U_\alpha$ satisfies $(S_k)$ for every $\alpha$, then
$\sh F$ satisfies $(S_k)$.
\end{definition}

\begin{remarks}[5.7.3]
\phantomsection
\label{IV.5.7.3}
\begin{enumerate}[{\rm(i)}]
\item
\phantomsection
\label{IV.5.7.3.i}
Recall that
$\operatorname{prof}(\sh F_x)\leq\dim(\sh F_x)$
\sref[0]{0.16.4.5.1} whenever $\sh F_x\neq0$.  Thus $\sh F$
satisfies $(S_k)$ precisely when
$\operatorname{prof}(\sh F_x)\geq k$ except at points $x$ for which
$\dim(\sh F_x)<k$, while at those latter points
\[
\dim(\sh F_x)=\operatorname{prof}(\sh F_x).
\]
Equivalently \sref[0]{0.16.5.1}, $\sh F$ is Cohen--Macaulay at those
points.  At points where $\dim(\sh F_x)=k$, one likewise has
$\operatorname{prof}(\sh F_x)=k$.

\oldpage[IV-2]{104}

Hence $\sh F$ is also Cohen--Macaulay at those points.  To say that
$\sh F$ satisfies $(S_k)$ for every $k$ is therefore the same as to
say that $\sh F$ is a Cohen--Macaulay $\sh O_X$-module.  If
$k'\geq k$, then $(S_{k'})$ implies $(S_k)$; if $k\leq0$, every
coherent $\sh O_X$-module satisfies $(S_k)$.

\item
\phantomsection
\label{IV.5.7.3.ii}
To verify \eqref{IV.5.7.2.1}, it is enough to consider
$x\in\operatorname{Supp}(\sh F)$; otherwise
$\dim(\sh F_x)=-\infty$ \sref[0]{0.14.1.2}.

\item
\phantomsection
\label{IV.5.7.3.iii}
If $X=\Spec(A)$, where $A$ is Noetherian, and
$\sh F=\widetilde M$, where $M$ is a finite $A$-module, then $M$ is
said to satisfy $(S_k)$ when $\sh F$ does.

For an arbitrary locally Noetherian prescheme $X$, saying that
$\sh F$ satisfies $(S_k)$ at $x$ means the following.  Put
$Y=\Spec(\sh O_x)$; then the $\sh O_Y$-module
$\widetilde{\sh F_x}$ satisfies $(S_k)$.  Notice that
\eqref{IV.5.7.2.1} does not in general imply
\eqref{IV.5.7.2.2} for every generization $x'$ of $x$, since there
is no general inequality between
$\operatorname{prof}(\sh F_x)$ and
$\operatorname{prof}(\sh F_{x'})$ \sref[0]{0.16.4.6}.

\item
\phantomsection
\label{IV.5.7.3.iv}
It follows immediately from the definition that if $\sh F$
satisfies $(S_k)$ at $x$, then it also satisfies $(S_k)$ at every
generization of $x$.

\item
\phantomsection
\label{IV.5.7.3.v}
Property $(S_k)$ is especially important for $k=1$ and $k=2$.
Serre introduced it for $k=2$ in order to express his normality
criterion; see \sref{IV.5.8.6}.\footnote{The printed French text points to
\textup{(5.8.5)}; Serre's normality criterion is Theorem
\textup{(5.8.6)}.}

\item
\phantomsection
\label{IV.5.7.3.vi}
Let $Y$ be a closed sub-prescheme of $X$, let
$j:Y\to X$ be the canonical immersion, and let $\sh G$ be a coherent
$\sh O_Y$-module.  For every $x\in Y$,
\[
\dim(\sh G_x)=\dim((j_*\sh G)_x),\qquad
\operatorname{prof}(\sh G_x)
 =\operatorname{prof}((j_*\sh G)_x),
\]
and consequently
\[
\operatorname{coprof}(\sh G_x)
 =\operatorname{coprof}((j_*\sh G)_x).
\]
Thus $\sh G$ satisfies $(S_k)$ if and only if $j_*\sh G$ does.
\end{enumerate}
\end{remarks}

\begin{proposition}[5.7.4]
\phantomsection
\label{IV.5.7.4}
Let $X$ be a locally Noetherian prescheme and let $\sh F$ be a
coherent $\sh O_X$-module.  Put
$S=\operatorname{Supp}(\sh F)$, and for every $n\geq0$ let $Z_n$ be
the set of points $z\in X$ such that
$\operatorname{coprof}(\sh F_z)>n$; thus $Z_n\subset S$.
\begin{enumerate}[{\rm(i)}]
\item
\phantomsection
\label{IV.5.7.4.i}
The module $\sh F$ satisfies $(S_k)$ if and only if, for every
$n\geq0$,
\begin{equation}
\operatorname{codim}(Z_n,S)>n+k.
\tag{5.7.4.1}\label{IV.5.7.4.1}
\end{equation}

\item
\phantomsection
\label{IV.5.7.4.ii}
Suppose in addition that every $Z_n$ is closed in $X$.  Then
$\sh F$ satisfies $(S_k)$ at a point $x\in S$ if and only if, for
every $n\geq0$,
\begin{equation}
\operatorname{codim}_x(Z_n,S)>n+k.
\tag{5.7.4.2}\label{IV.5.7.4.2}
\end{equation}
\end{enumerate}
\end{proposition}

\begin{proof}
\begin{enumerate}[{\rm(i)}]
\item
By definition \sref{IV.5.1.3},
\[
\operatorname{codim}(Z_n,S)
 =\inf_{z\in Z_n}\dim(\sh O_{S,z}),
\]
and \eqref{IV.5.7.4.1} therefore says
\sref{IV.5.1.12.2} that, for every $z\in X$ and every $n\geq0$,
\[
\dim(\sh F_z)-\operatorname{prof}(\sh F_z)\geq n+1
\]
implies
\[
\dim(\sh F_z)\geq n+k+1.
\]
If $a=\dim(\sh F_z)$ and
$b=\operatorname{prof}(\sh F_z)$, then $b\leq a$; requiring, for
every $n\geq0$, that $b\leq a-n-1$ imply $k\leq a-n-1$ is
equivalent to $b\geq\inf\{k,a\}$.  This proves (i).

\oldpage[IV-2]{105}

\item
The argument is the same as in
{\hyperlink{ega4.c04.sec.5-7-4-i}{\textup{(i)}}}, except that one restricts to
the points $z$ which are generizations of $x$ and uses
\sref{IV.5.1.3.2}.
\end{enumerate}
\end{proof}

\begin{proposition}[5.7.5]
\phantomsection
\label{IV.5.7.5}
Let $X$ be a locally Noetherian prescheme and let $\sh F$ be a
coherent $\sh O_X$-module.  The module $\sh F$ satisfies $(S_k)$,
where $k\geq1$, if and only if the following condition holds:
for every integer $r$ with $0\leq r<k$, every open subset $U$ of
$X$, and every $(\sh F|U)$-regular sequence
$(f_i)_{1\leq i\leq r}$ of sections of $\sh O_X$ over $U$, the
module
\[
(\sh F|U)\Big/\Big(\sum_{i=1}^{r}f_i(\sh F|U)\Big)
\]
has no embedded associated prime cycle.
\end{proposition}

\begin{proof}
For $k=1$, the assertion says that $\sh F$ satisfies $(S_1)$ if and
only if $\sh F$ has no embedded associated prime cycle.  Indeed,
$(S_1)$ says that
$\operatorname{prof}(\sh F_x)\geq1$ at precisely the points
$x\in\operatorname{Supp}(\sh F)$ for which
$\dim(\sh F_x)>0$; at the other points of the support,
$\dim(\sh F_x)=\operatorname{prof}(\sh F_x)=0$.
Now $\operatorname{prof}(\sh F_x)\geq1$ means that the maximal ideal
$\mathfrak m_x$ is not associated to $\sh F_x$
\hyperref[0.16.4.6.i]{\textup{(0, 16.4.6(i))}}, or equivalently that $x$ is not associated to
$\sh F$ \sref{IV.3.1.1}.  If $S'$ is a sub-prescheme of $X$ whose
underlying space is $\operatorname{Supp}(\sh F)$, then
\[
\dim(\sh F_x)=\dim(\sh O_{S',x})
\]
\sref{IV.5.1.12.1}; saying that this dimension is positive is the
same as saying that $x$ is not a maximal point of
$\operatorname{Supp}(\sh F)$ \sref{IV.5.1.2}.  This proves the
assertion for $k=1$.

Let $k>1$, and first suppose that $\sh F$ satisfies $(S_k)$.  By
induction on $k$ and the case just proved, it is enough to show that
if $f$ is an $\sh F$-regular section of $\sh O_X$, then
$\sh F/f\sh F$ satisfies $(S_{k-1})$.  At each $x\in X$, the element
$f_x$ is $\sh F_x$-regular.  If it is invertible, the quotient stalk
is zero.  Otherwise $f_x\in\mathfrak m_x$, and
\[
\operatorname{prof}(\sh F_x/f_x\sh F_x)
 =\operatorname{prof}(\sh F_x)-1,\qquad
\dim(\sh F_x/f_x\sh F_x)=\dim(\sh F_x)-1
\]
\hyperref[0.16.4.6.i]{\textup{(0, 16.4.6(i))}}, \sref[0]{0.16.3.4}; the assertion follows.
Thus the condition is necessary.

Conversely, suppose the stated condition holds.  We argue by
induction on $k$.  Let $x\in X$, and first suppose that
$\dim(\sh F_x)\geq k$.  The induction hypothesis gives $(S_{k-1})$,
so $\operatorname{prof}(\sh F_x)\geq k-1$.  By
\sref[0]{0.15.2.4}, after replacing $X$ by an open neighbourhood
$U$ of $x$, there is an $(\sh F|U)$-regular sequence
$(f_i)_{1\leq i\leq k-1}$ with
$(f_i)_x\in\mathfrak m_x$.  Put
\[
\sh G=(\sh F|U)\Big/\Big(
 \sum_{i=1}^{k-1}f_i(\sh F|U)\Big).
\]
By hypothesis $\sh G$ has no embedded associated prime cycle, while
\[
\dim(\sh G_x)=\dim(\sh F_x)-(k-1)\geq1
\]
\sref[0]{0.16.3.4}.  Hence $x$ is not associated to $\sh G$, so
$\operatorname{prof}(\sh G_x)\geq1$ \sref[0]{0.16.4.6}; since
\[
\operatorname{prof}(\sh G_x)
 =\operatorname{prof}(\sh F_x)-(k-1),
\]
we obtain $\operatorname{prof}(\sh F_x)\geq k$.

If instead $\dim(\sh F_x)=r<k$, the induction hypothesis gives
\[
\operatorname{prof}(\sh F_x)
 \geq\inf\{k-1,r\}=r.
\]
Thus $\sh F$ satisfies $(S_k)$, completing the proof.
\end{proof}

\begin{corollary}[5.7.6]
\phantomsection
\label{IV.5.7.6}
Suppose that $\sh F$ satisfies $(S_k)$, where $k\geq1$.  If
$(f_i)_{1\leq i\leq r}$ is an $\sh F$-regular sequence of sections
of $\sh O_X$ over $X$ and $r<k$, then
\[
\sh F\Big/\Big(\sum_{i=1}^{r}f_i\sh F\Big)
\]
satisfies $(S_{k-r})$.
\end{corollary}

\begin{proof}
This follows immediately from \sref{IV.5.7.5}.
\end{proof}

\begin{corollary}[5.7.7]
\phantomsection
\label{IV.5.7.7}
The module $\sh F$ satisfies $(S_2)$ if and only if $\sh F$ has no
embedded associated prime cycle and, for every open subset $U$ of
$X$ and every $(\sh F|U)$-regular section $f$ of $\sh O_X$ over
$U$, the module $(\sh F|U)/f(\sh F|U)$ has no embedded associated
prime cycle.
\end{corollary}

\begin{proof}
This follows immediately from \sref{IV.5.7.5}.
\end{proof}

\oldpage[IV-2]{106}

\begin{remarks}[5.7.8]
\phantomsection
\label{IV.5.7.8}
Let $X$ be a locally Noetherian prescheme of dimension $1$.  The
following conditions are equivalent: $X$ is Cohen--Macaulay; $X$
satisfies $(S_1)$; $X$ satisfies $(S_n)$ for some $n\geq1$.  This
follows from Definitions \sref{IV.5.7.1} and \sref{IV.5.7.2}.
By \sref{IV.5.7.5}, these conditions are also equivalent to the
absence of an embedded associated prime cycle.  In particular,
every reduced locally Noetherian prescheme of dimension $1$ is
Cohen--Macaulay.
\end{remarks}

\begin{proposition}[5.7.9]
\phantomsection
\label{IV.5.7.9}
Let $A$ and $B$ be Noetherian rings, let $\rho:A\to B$ be a ring
homomorphism, and let $M$ be a $B$-module such that the
$A$-module $M_{[\rho]}$ is finite.  Let $\mathfrak p$ be a prime
ideal of $A$.  The prime ideals of $B$ lying over $\mathfrak p$ and
belonging to $\operatorname{Supp}(M)$ are finite in number.  If
$(\mathfrak q_i)_{1\leq i\leq n}$ is the family of these ideals,
then
\begin{align}
\dim_{A_{\mathfrak p}}(M_{\mathfrak p})
 &=\sup_i\dim_{B_{\mathfrak q_i}}(M_{\mathfrak q_i}),
\tag{5.7.9.1}\label{IV.5.7.9.1}\\
\operatorname{prof}_{A_{\mathfrak p}}(M_{\mathfrak p})
 &=\inf_i\operatorname{prof}_{B_{\mathfrak q_i}}
   (M_{\mathfrak q_i}).
\tag{5.7.9.2}\label{IV.5.7.9.2}
\end{align}
\end{proposition}

\begin{proof}
Let $\mathfrak b$ be the annihilator of $M$.  The ring
$B/\mathfrak b$ identifies with an $A$-submodule of
$\operatorname{End}_A(M_{[\rho]})$, and is therefore finite over
$A$.  Hence only finitely many prime ideals of $B$ lie over
$\mathfrak p$ and contain $\mathfrak b$; these are precisely the
ones belonging to $\operatorname{Supp}(M)$.

Replacing $B$ by $B/\mathfrak b$ does not change the right-hand
sides of \eqref{IV.5.7.9.1} and \eqref{IV.5.7.9.2}
\sref[0]{0.16.1.9}, \sref[0]{0.16.4.8}; we may therefore suppose
that $B$ is finite over $A$.  Put $S=A-\mathfrak p$ and
$B'=S^{-1}B$.  Then $B'$ is a Noetherian semilocal ring whose
maximal ideals are $S^{-1}\mathfrak q_i$, and
\[
\dim_{A_{\mathfrak p}}(M_{\mathfrak p})
 =\dim_{B'}(M_{\mathfrak p})
\]
\sref[0]{0.16.1.9}.  Formula \eqref{IV.5.7.9.1} is now the
semilocal dimension formula \sref[0]{0.16.1.7.4}.

For \eqref{IV.5.7.9.2}, as in \sref[0]{0.16.4.8}, it is enough to
treat the case
$\operatorname{prof}_{A_{\mathfrak p}}(M_{\mathfrak p})=0$.  The
same argument shows that $M_{\mathfrak p}$ contains a
$B'$-submodule of finite length, and hence a simple
$B'$-submodule.  Such a module is isomorphic to the residue field
of one of the local rings $B_{\mathfrak q_i}$.  Thus
\[
\operatorname{prof}_{B_{\mathfrak q_i}}(M_{\mathfrak q_i})=0
\]
for at least one $i$ \sref[0]{0.16.4.6}, proving the formula.
\end{proof}

\begin{corollary}[5.7.10]
\phantomsection
\label{IV.5.7.10}
Assume the hypotheses of \sref{IV.5.7.9}, and suppose in addition
that $A$ is local.  The $A$-module $M_{[\rho]}$ is Cohen--Macaulay
if and only if every $M_{\mathfrak q_i}$, for
$\mathfrak q_i$ lying over the maximal ideal of $A$, is a
Cohen--Macaulay $B_{\mathfrak q_i}$-module and all the numbers
$\dim_{B_{\mathfrak q_i}}(M_{\mathfrak q_i})$ are equal.
\end{corollary}

\begin{proof}
By \eqref{IV.5.7.9.1} and \eqref{IV.5.7.9.2}, these conditions are
equivalent to
\[
\dim_A(M_{[\rho]})
 =\operatorname{prof}_A(M_{[\rho]}).
\]
\end{proof}

\begin{corollary}[5.7.11]
\phantomsection
\label{IV.5.7.11}
Assume the hypotheses of \sref{IV.5.7.9}.
\begin{enumerate}[{\rm(i)}]
\item
\phantomsection
\label{IV.5.7.11.i}
If $M_{[\rho]}$ satisfies $(S_k)$, then $M$ satisfies $(S_k)$.

\item
\phantomsection
\label{IV.5.7.11.ii}
Suppose that for every pair of prime ideals
$\mathfrak q,\mathfrak q'$ of $B$ satisfying
$\rho^{-1}(\mathfrak q)=\rho^{-1}(\mathfrak q')$, one has
\[
\dim_{B_{\mathfrak q}}(M_{\mathfrak q})
 =\dim_{B_{\mathfrak q'}}(M_{\mathfrak q'}).
\]
If $M$ satisfies $(S_k)$, then $M_{[\rho]}$ satisfies $(S_k)$.
\end{enumerate}
\end{corollary}

\oldpage[IV-2]{107}

\begin{proof}
This follows immediately from \eqref{IV.5.7.9.1},
\eqref{IV.5.7.9.2}, and the definition of $(S_k)$.
\end{proof}

\begin{env}[5.7.12]
\phantomsection
\label{IV.5.7.12}
In accordance with \sref{IV.5.7.1}, if $A$ is any Noetherian ring
and $M$ is a finite $A$-module, define
\[
\operatorname{coprof}_A(M)
 =\operatorname{coprof}(\widetilde M)
 =\sup_{x\in X}\operatorname{coprof}_{A_x}(M_x),
\qquad X=\Spec(A).
\]
We shall see in \sref{IV.6.11.5} that this definition agrees with
that of \sref[0]{0.16.4.9} when $A$ is a Noetherian local ring.
\end{env}

\begin{corollary}[5.7.13]
\phantomsection
\label{IV.5.7.13}
Let $A$ and $B$ be Noetherian rings, let $\rho:A\to B$ be a ring
homomorphism, and let $M$ be a $B$-module such that $M_{[\rho]}$ is
a finite $A$-module.  Then
\begin{equation}
\operatorname{coprof}_A(M_{[\rho]})
 \geq \operatorname{coprof}_B(M).
\tag{5.7.13.1}\label{IV.5.7.13.1}
\end{equation}
\end{corollary}

\begin{proof}
This follows from \sref{IV.5.7.12} and
\eqref{IV.5.7.9.1}, \eqref{IV.5.7.9.2}.
\end{proof}
% END NATIVE INPUT components/07_section57_depth_and_property_sk
% BEGIN NATIVE INPUT components/08_section58_regular_preschemes_serre_criterion
\subsection{Regular preschemes and property
\texorpdfstring{$(\mathrm{R}_k)$}{(Rk)}. Serre's normality criterion}
\label{subsection:IV.5.8}

\begin{env}[5.8.1]
\label{IV.5.8.1}
Recall \sref[0]{0.4.1.4} that a ringed space $(X,\sh O_X)$ is said
to be \emph{regular} at a point $x\in X$ if $\sh O_x$ is a regular
local ring.  When dealing with preschemes, we shall use this
terminology in this chapter only when $X$ is locally Noetherian.
\end{env}

\begin{definition}[5.8.2]
\label{IV.5.8.2}
A locally Noetherian prescheme $X$ is said to be \emph{regular in
codimension $\leq k$}, or to satisfy property $(\mathrm R_k)$, if the
set of points where $X$ is not regular has codimension $>k$.  In other
words \sref{IV.5.1.3}, for every $x\in X$,
\[
  \dim(\sh O_x)\leq k
  \quad\Longrightarrow\quad
  \sh O_x\ \text{is regular}.
\]
\end{definition}

To say that $X$ is regular means that $X$ satisfies $(\mathrm R_k)$
for every $k$.

If $X=\Spec(A)$, where $A$ is a Noetherian ring, we say that $A$
satisfies $(\mathrm R_k)$ if $X$ does; to say that $X$ is regular is
to say that the ring $A$ is regular \sref[0]{0.17.3.6}.  For an
arbitrary locally Noetherian prescheme $X$, we say that $X$ satisfies
$(\mathrm R_k)$ \emph{at a point} $x\in X$ if the local ring
$\sh O_x$ satisfies $(\mathrm R_k)$.  Thus, for every generization
$x'$ of $x$ in $X$, the relation
$\dim(\sh O_{x'})\leq k$ must imply that $\sh O_{x'}$ is a regular
local ring.  By \sref[0]{0.17.3.6}, saying that $X$ is regular at
$x$ is equivalent to saying that $X$ satisfies $(\mathrm R_n)$ at
$x$ for every $n\geq0$.

\begin{proposition}[5.8.3]
\label{IV.5.8.3}
Let $k$ be a field and let $X$ be a prescheme locally of finite type
over $k$.  For every $x\in X$, there is an open neighbourhood of $x$
in $X$ which is isomorphic to a subscheme of a regular $k$-scheme.
\end{proposition}

\begin{proof}
There is an affine open neighbourhood $U$ of $x$ isomorphic to a
$k$-scheme $\Spec(A)$, where $A$ is a $k$-algebra of finite type.
The algebra $A$ is therefore isomorphic to a quotient of a polynomial
algebra $k[T_1,\ldots,T_n]$, and the latter is a regular ring
\sref[0]{0.17.3.7}.
\end{proof}

\begin{env}[5.8.4]
\label{IV.5.8.4}
By \sref{IV.5.8.2}, saying that $X$ satisfies $(\mathrm R_0)$ means
that, at every maximal point $x$ of $X$, the ring $\sh O_x$ is a
field \sref[0]{0.17.1.4}; in other words, $X$ is reduced at $x$.
The set $U$ of points where $X$ is reduced is open, because the
nilradical of $X$ is a coherent $\sh O_X$-module
\sref[I]{I.6.1.1}, \sref[0]{0.5.2.2}.  Consequently, saying that
$X$ satisfies $(\mathrm R_0)$ is equivalent to saying that $U$ is
everywhere dense.
\end{env}

\oldpage[IV-2]{108}

\begin{proposition}[5.8.5]
\label{IV.5.8.5}
A locally Noetherian prescheme $X$ is reduced if and only if it
satisfies $(\mathrm S_1)$ and $(\mathrm R_0)$.
\end{proposition}

\begin{proof}
Taking \sref{IV.5.7.5} into account, this follows from the preceding
remark and \sref{IV.3.2.1}.
\end{proof}

\begin{theorem}[5.8.6]
\label{IV.5.8.6}
\emph{(Serre's criterion).}
Let $X$ be a locally Noetherian prescheme.  Then $X$ is normal if and
only if it satisfies $(\mathrm S_2)$ and $(\mathrm R_1)$.  Equivalently,
for every $x\in X$, the following conditions hold:
\begin{enumerate}[{\rm(i)}]
  \item\label{IV.5.8.6.i}
    If $\dim(\sh O_x)\leq1$, then $\sh O_x$ is regular; that is, it is
    a field or a discrete valuation ring \sref[0]{0.17.1.4}.
  \item\label{IV.5.8.6.ii}
    If $\dim(\sh O_x)\geq2$, then
    $\operatorname{prof}(\sh O_x)\geq2$.
\end{enumerate}
\end{theorem}

\begin{proof}
The conditions are necessary.  Indeed, saying that $X$ is normal
means that, for every $x\in X$, the ring $\sh O_x$ is an integrally
closed Noetherian local ring.  If $\dim(\sh O_x)=0$, then $\sh O_x$
is a field because it is a domain; if $\dim(\sh O_x)=1$, then
$\sh O_x$ is a discrete valuation ring by \sref[II]{II.7.1.6}.
Furthermore, for every nonzero element $f_x\in\sh O_x$, the associated
prime ideals of $\sh O_x/f_x\sh O_x$ are non-embedded
(Bourbaki, \emph{Alg.\ comm.}, chap.~VII, \S~1, no.~4, prop.~8).
Hence $\sh O_x$ satisfies $(\mathrm S_2)$ by \sref{IV.5.7.7}.

Conversely, suppose that the two conditions hold.  It follows first
from \sref{IV.5.8.5} that $X$ is reduced.  Since the question is
local, we may suppose that $X=\Spec(A)$, where $A$ is a reduced
Noetherian ring \sref[I]{I.5.1.4}.  If $R$ is the total ring of
fractions of $A$, then $R$ is a direct product of finitely many
fields; taking \sref[II]{II.6.3.6} into account, it is enough to prove
that $A$ is integrally closed in $R$.

Let $h=f/g\in R$ be integral over $A$, where $f,g\in A$ and $g$ is a
non-zero-divisor.  There is a relation
\begin{equation}
\label{IV.5.8.6.1}
  f^n+\sum_{i=1}^{n}a_i f^{\,n-i}g^i=0,
  \qquad a_i\in A\quad(1\leq i\leq n).
\tag{5.8.6.1}
\end{equation}
Let $\mathfrak p$ be a prime ideal of $A$ such that
$\dim(A_{\mathfrak p})=1$, and write $f_{\mathfrak p}$ and
$g_{\mathfrak p}$ for the images of $f$ and $g$ in
$A_{\mathfrak p}$.  It follows from \eqref{IV.5.8.6.1} that
$f_{\mathfrak p}/g_{\mathfrak p}$ is integral over
$A_{\mathfrak p}$.  This quotient belongs to the total ring of
fractions of $A_{\mathfrak p}$, since flatness implies that
$g_{\mathfrak p}$ is still a non-zero-divisor
\sref[0]{0.5.3.1}.  But $A_{\mathfrak p}$ is regular and hence
integrally closed, so
$f_{\mathfrak p}/g_{\mathfrak p}\in A_{\mathfrak p}$.  Equivalently,
\[
  (fA)_{\mathfrak p}\subset (gA)_{\mathfrak p}.
\]

Since $g$ is a non-zero-divisor in $A$, condition $(\mathrm S_2)$
and \sref{IV.5.7.7} imply that $A/gA$ has only non-embedded
associated prime ideals $\mathfrak p_i$ $(1\leq i\leq n)$.  The ideal
$gA$ is therefore an intersection of primary ideals
$\mathfrak q_i$ corresponding to the $\mathfrak p_i$; moreover, each
$\mathfrak q_i$ is the inverse image in $A$, under
$A\to A_{\mathfrak p_i}$, of $(gA)_{\mathfrak p_i}$
(Bourbaki, \emph{Alg.\ comm.}, chap.~IV, \S~2, no.~3, prop.~5).
By the Hauptidealsatz \sref[0]{0.16.3.2},
$\dim(A_{\mathfrak p_i})=1$ for every $i$, so the preceding argument
gives
\[
  (fA)_{\mathfrak p_i}\subset(gA)_{\mathfrak p_i}
  \qquad(1\leq i\leq n).
\]
Now $fA$ is contained in the intersection of the inverse images of
the $(fA)_{\mathfrak p_i}$.  It is therefore contained in the
intersection of the inverse images of the $(gA)_{\mathfrak p_i}$,
which is $gA$.  Thus $fA\subset gA$, and hence $f/g\in A$.
\end{proof}

\oldpage[IV-2]{109}
% END NATIVE INPUT components/08_section58_regular_preschemes_serre_criterion
% BEGIN NATIVE INPUT components/09_section59_z_pure_z_closed_modules
\subsection{\texorpdfstring{$Z$}{Z}-pure and
\texorpdfstring{$Z$}{Z}-closed modules}
\label{subsection:IV.5.9}

Some of the notions and results in this section and the next are
special cases of notions and results developed in Chapter~III in the
theory of local cohomology.  For the reader's convenience, we give an
independent account here.

\begin{env}[5.9.1]
\label{IV.5.9.1}
Let $X$ be a locally Noetherian prescheme and let $Z$ be a subset of
$X$ \emph{stable under specialization}.  This means that the closure
of every finite subset $M$ of $Z$ is contained in $Z$.  It follows
that $Z$ is the union of a filtered increasing family $(Z_\alpha)$ of
closed subsets of $X$; conversely, such a union is plainly stable
under specialization.

Put $U_\alpha=X-Z_\alpha$.  Then $X-Z$ is the intersection of the
filtered decreasing family of open subsets $U_\alpha$.  Let
$i_\alpha:U_\alpha\to X$ be the canonical injection, and, whenever
$U_\alpha\supset U_\beta$, let
$i_{\alpha\beta}:U_\beta\to U_\alpha$ be the canonical injection, so
that $i_\beta=i_\alpha\circ i_{\alpha\beta}$.

Let $\sh F$ be an $\sh O_X$-module, not necessarily quasi-coherent.
We have
\[
 (i_\beta)_*(\sh F|U_\beta)
 = (i_\alpha)_*((i_{\alpha\beta})_*(\sh F|U_\beta)).
\]
The canonical homomorphism \sref[0]{0.4.4.3.2}
\[
  \sh F|U_\alpha\longrightarrow
  (i_{\alpha\beta})_*(\sh F|U_\beta)
\]
therefore gives, after applying $(i_\alpha)_*$, a homomorphism
\[
 \rho_{\beta\alpha}:
 (i_\alpha)_*(\sh F|U_\alpha)\longrightarrow
 (i_\beta)_*(\sh F|U_\beta).
\]
For $U_\alpha\supset U_\beta\supset U_\gamma$, one immediately
checks that
$\rho_{\gamma\alpha}=\rho_{\gamma\beta}\circ\rho_{\beta\alpha}$.
Thus the modules $(i_\alpha)_*(\sh F|U_\alpha)$ form an inductive
system under the homomorphisms $\rho_{\beta\alpha}$.  We put
\end{env}

\begin{equation}
\label{IV.5.9.1.1}
 \sh H^0_{X/Z}(\sh F)
 =\varinjlim_\alpha(i_\alpha)_*(\sh F|U_\alpha).
\tag{5.9.1.1}
\end{equation}

This $\sh O_X$-module does not depend on the increasing family
$(Z_\alpha)$ of closed subsets having union $Z$.  Indeed, let $V$ be
a Noetherian open subset of $X$.  In the category of
$\sh O_V$-modules, the functor
$\sh G\mapsto\Gamma(V,\sh G)$ commutes with inductive limits
\sref[G]{G.II.3.10.1}; hence \eqref{IV.5.9.1.1} gives
\begin{equation}
\label{IV.5.9.1.2}
 \Gamma(V,\sh H^0_{X/Z}(\sh F))
 =\varinjlim_\alpha\Gamma(V\cap U_\alpha,\sh F).
\tag{5.9.1.2}
\end{equation}

Let $(Z'_\lambda)$ be another filtered increasing family of closed
subsets of $X$ with union $Z$.  For fixed $\alpha$, the set
$V\cap Z_\alpha$ is the union of the
$V\cap Z_\alpha\cap Z'_\lambda$.  Now $V\cap Z_\alpha$ is locally
closed in $X$, and every closed irreducible subset of it has a
generic point.  Since the $V\cap Z_\alpha\cap Z'_\lambda$ form a
filtered increasing family of closed subsets of $V\cap Z_\alpha$,
there is a $\lambda$ such that
\[
 V\cap Z_\alpha\cap Z'_\lambda=V\cap Z_\alpha
 \qquad\text{\sref[0]{0.9.2.4}}.
\]
Equivalently, $V\cap Z_\alpha\subset V\cap Z'_\lambda$.  The two
filtered decreasing families $V\cap U_\alpha$ and
$V\cap U'_\lambda$, where $U'_\lambda=X-Z'_\lambda$, are therefore
cofinal.  The asserted independence follows from
\eqref{IV.5.9.1.2}.

\oldpage[IV-2]{110}

The subset $Z$ need not be constructible.  For example, take
$X=\Spec(A)$, where $A$ is a Noetherian domain having infinitely many
maximal ideals, and let $Z$ be the complement of the generic point of
$X$.

If $Z$ is closed and $i:X-Z\to X$ is the canonical injection, then
\begin{equation}
\label{IV.5.9.1.3}
 \sh H^0_{X/Z}(\sh F)=i_*(\sh F|X-Z).
\tag{5.9.1.3}
\end{equation}
In particular, when $Z=\varnothing$,
$\sh H^0_{X/Z}(\sh F)=\sh F$.

\begin{proposition}[5.9.2]
\label{IV.5.9.2}
\begin{enumerate}[{\rm(i)}]
  \item\label{IV.5.9.2.i}
    The functor $\sh F\mapsto\sh H^0_{X/Z}(\sh F)$ is left exact.
  \item\label{IV.5.9.2.ii}
    If $\sh F$ is quasi-coherent, so is
    $\sh H^0_{X/Z}(\sh F)$.
\end{enumerate}
\end{proposition}

\begin{proof}
Assertion~(i) follows from \eqref{IV.5.9.1.1}, from the left
exactness of $(i_\alpha)_*$, and from the fact that inductive limits
preserve exactness in the category of $\sh O_X$-modules.\ 
Assertion~(ii) follows from \sref[I]{I.9.2.2} and from the fact that
an inductive limit of quasi-coherent $\sh O_X$-modules is
quasi-coherent \sref[I]{I.1.3.9}.
\end{proof}

\begin{remark}[5.9.3]
\label{IV.5.9.3}
If $\sh F$ is an $\sh O_X$-algebra, then
$\sh H^0_{X/Z}(\sh F)$ is one as well \sref[0]{0.4.2.4}.  In
particular, $\sh H^0_{X/Z}(\sh O_X)$ is a quasi-coherent
$\sh O_X$-algebra, and, for every $\sh O_X$-module $\sh F$,
$\sh H^0_{X/Z}(\sh F)$ is a module over that algebra; it is
quasi-coherent when $\sh F$ is quasi-coherent \sref[I]{I.9.6.1}.

More specifically, suppose that $X=\Spec(A)$, where $A$ is a
Noetherian domain.  Then $\sh H^0_{X/Z}(\sh O_X)$ is the
$\sh O_X$-algebra $\widetilde B$, where
\begin{equation}
\label{IV.5.9.3.1}
 B=\bigcap_{\mathfrak p\in X-Z}A_{\mathfrak p}.
\tag{5.9.3.1}
\end{equation}
This follows at once from \eqref{IV.5.9.1.2} and
\sref[I]{I.8.2.1.1}.
\end{remark}

\begin{proposition}[5.9.4]
\label{IV.5.9.4}
Let $X$ be a locally Noetherian prescheme, let $Z\subset X$ be stable
under specialization, let $X'$ be a locally Noetherian prescheme,
and let $f:X'\to X$ be flat.  Then $Z'=f^{-1}(Z)$ is stable under
specialization, and, for every quasi-coherent $\sh O_X$-module
$\sh F$, there is a canonical isomorphism
\begin{equation}
\label{IV.5.9.4.1}
 f^*(\sh H^0_{X/Z}(\sh F))
 \xrightarrow{\sim}
 \sh H^0_{X'/Z'}(f^*\sh F).
\tag{5.9.4.1}
\end{equation}
\end{proposition}

\begin{proof}
With the notation of \sref{IV.5.9.1},
$Z'_\alpha=f^{-1}(Z_\alpha)$ is closed in $X'$, and $Z'$ is the
union of the $Z'_\alpha$.  The base change of $i_\alpha$ to $X'$ is
the canonical injection
$i'_\alpha:U'_\alpha\to X'$, where
$U'_\alpha=X'-Z'_\alpha=f^{-1}(U_\alpha)$.
Since $f$ is flat, the canonical homomorphism
\[
 f^*((i_\alpha)_*(\sh F|U_\alpha))
 \longrightarrow
 (i'_\alpha)_*(f^*\sh F|U'_\alpha)
\]
is bijective \sref{IV.2.3.1}.  For $\alpha\leq\beta$, the diagram
\[
\xymatrix{
 f^*((i_\alpha)_*(\sh F|U_\alpha)) \ar[r]^-{\sim} \ar[d]
   & (i'_\alpha)_*(f^*\sh F|U'_\alpha) \ar[d] \\
 f^*((i_\beta)_*(\sh F|U_\beta)) \ar[r]^-{\sim}
   & (i'_\beta)_*(f^*\sh F|U'_\beta)
}
\]
is commutative.  Passing to the limit gives a canonical isomorphism
\[
 \varinjlim_\alpha f^*((i_\alpha)_*(\sh F|U_\alpha))
 \xrightarrow{\sim}
 \sh H^0_{X'/Z'}(f^*\sh F).
\]

\oldpage[IV-2]{111}

Since $f^*$ commutes with inductive limits \sref[0]{0.4.3.2}, this
is precisely \eqref{IV.5.9.4.1}.
\end{proof}

\begin{corollary}[5.9.5]
\label{IV.5.9.5}
Under the hypotheses of \sref{IV.5.9.4}, if
$\sh H^0_{X/Z}(\sh F)$ is coherent, then
$\sh H^0_{X'/Z'}(f^*\sh F)$ is coherent.  The converse holds when
$f$ is faithfully flat and quasi-compact.
\end{corollary}

\begin{proof}
The first assertion follows from \eqref{IV.5.9.4.1} and
\sref[0]{0.5.3.11}.  The second amounts to saying that if
$\sh H^0_{X'/Z'}(f^*\sh F)$ is of finite type, then so is
$\sh H^0_{X/Z}(\sh F)$; this follows from \eqref{IV.5.9.4.1} and
\sref{IV.2.5.2}.
\end{proof}

\begin{corollary}[5.9.6]
\label{IV.5.9.6}
Let $X$ be a locally Noetherian prescheme, let $Z\subset X$ be stable
under specialization, and let $\sh F$ be a quasi-coherent
$\sh O_X$-module.  For every $x\in X$, put
$X_x=\Spec(\sh O_x)$ and $Z_x=Z\cap X_x$.  There is a functorial
isomorphism
\begin{equation}
\label{IV.5.9.6.1}
 (\sh H^0_{X/Z}(\sh F))_x
 \xrightarrow{\sim}
 \sh H^0_{X_x/Z_x}(\widetilde{\sh F}_x).
\tag{5.9.6.1}
\end{equation}
\end{corollary}

\begin{proof}
Apply \sref{IV.5.9.4} to the canonical flat morphism $X_x\to X$ and
use \sref[I]{I.1.6.5}.
\end{proof}

\begin{env}[5.9.7]
\label{IV.5.9.7}
With the notation of \sref{IV.5.9.1}, for every $\alpha$ there is a
canonical functorial homomorphism
\[
 \sh F\longrightarrow(i_\alpha)_*(\sh F|U_\alpha)
 \qquad\text{\sref[0]{0.4.4.3.2}}.
\]
These homomorphisms form an inductive system, and passage to the
limit gives a canonical functorial homomorphism
\begin{equation}
\label{IV.5.9.7.1}
 \rho_{X/Z}:\sh F\longrightarrow\sh H^0_{X/Z}(\sh F).
\tag{5.9.7.1}
\end{equation}
\end{env}

\begin{proposition}[5.9.8]
\label{IV.5.9.8}
Let $X$ be a locally Noetherian prescheme, let $Z\subset X$ be stable
under specialization, and let $\sh F$ be an $\sh O_X$-module.  The
following conditions are equivalent:
\begin{enumerate}[{\rm a)}]
  \item\label{IV.5.9.8.a}
    The homomorphism $\rho_{X/Z}$ in \eqref{IV.5.9.7.1} is injective
    (respectively bijective).
  \item\label{IV.5.9.8.b}
    For every Noetherian open subset $V$ of $X$, the homomorphism
    \[
      (\rho_{X/Z})_V:\Gamma(V,\sh F)\longrightarrow
      \varinjlim_\alpha\Gamma(V\cap U_\alpha,\sh F)
    \]
    is injective (respectively bijective).
\end{enumerate}
They are also equivalent to:
\begin{enumerate}
  \item[{\rm a')}] \label{IV.5.9.8.a-prime}
    For every closed subset $T\subset Z$ of $X$, the canonical
    homomorphism \sref[0]{0.4.4.3.2}
    \[
      \sh F\longrightarrow i_*(\sh F|X-T),
    \]
    where $i:X-T\to X$ is the canonical injection, is injective
    (respectively bijective).
  \item[{\rm b')}] \label{IV.5.9.8.b-prime}
    For every closed subset $T\subset Z$ and every Noetherian open
    subset $V\subset X$, the restriction homomorphism
    \[
      \Gamma(V,\sh F)\longrightarrow
      \Gamma(V\cap(X-T),\sh F)
    \]
    is injective (respectively bijective).
\end{enumerate}
\end{proposition}

\begin{proof}
Taking \eqref{IV.5.9.1.2} into account, the equivalence of a) and b),
and likewise that of a') and b'), follows from the definition of the
functor $\Gamma$ and its left exactness.

For every closed subset $T\subset Z$, the homomorphism
$(\rho_{X/Z})_V$ is the composite
\begin{equation}
\label{IV.5.9.8.1}
 \Gamma(V,\sh F)\longrightarrow
 \Gamma(V\cap(X-T),\sh F)\longrightarrow
 \varinjlim_\alpha\Gamma(V\cap U_\alpha,\sh F).
\tag{5.9.8.1}
\end{equation}
Thus b) implies b') in the injective case, while b') implies b)
directly from the definition of an inductive limit.

\oldpage[IV-2]{112}

It remains to show that, if $\rho_{X/Z}$ is bijective, then so is
$\Gamma(V,\sh F)\to\Gamma(V\cap(X-T),\sh F)$.  By
\eqref{IV.5.9.8.1}, it is enough to observe that whenever
$U'\subset U$ are open subsets of $V$ containing $V\cap(X-Z)$, the
restriction homomorphism
$\Gamma(U,\sh F)\to\Gamma(U',\sh F)$ is injective.  This follows from
the injectivity of $\rho_{X/Z}$ by replacing $V$ with $U$ and
$V\cap(X-T)$ with $U'$.
\end{proof}

\begin{definition}[5.9.9]
\label{IV.5.9.9}
Under the hypotheses of \sref{IV.5.9.8}, the module $\sh F$ is called
\emph{$Z$-pure} (respectively \emph{$Z$-closed}) if $\rho_{X/Z}$ is
injective (respectively bijective).
\end{definition}

If $X=\Spec(A)$ is affine and $\sh F=\widetilde M$, where $M$ is an
$A$-module, we call $M$ $Z$-pure (respectively $Z$-closed) when
$\sh F$ is.

We say that $\sh F$ is \emph{$Z$-pure} (respectively
\emph{$Z$-closed}) \emph{at a point} $x\in X$ if, with the notation
of \sref{IV.5.9.6}, $\widetilde{\sh F}_x$ is $Z_x$-pure
(respectively $Z_x$-closed).  By \sref{IV.5.9.6}, this is equivalent
to saying that
\[
 \sh F_x\longrightarrow(\sh H^0_{X/Z}(\sh F))_x
\]
is injective (respectively bijective).  At every point $x\in X-Z$,
$\sh F$ is $Z$-closed by \sref{IV.5.9.8}.

\begin{corollary}[5.9.10]
\label{IV.5.9.10}
\begin{enumerate}[{\rm(i)}]
  \item\label{IV.5.9.10.i}
    Let $(V_\lambda)$ be an open covering of $X$.  Then $\sh F$ is
    $Z$-pure (respectively $Z$-closed) if and only if, for every
    $\lambda$, $\sh F|V_\lambda$ is $(Z\cap V_\lambda)$-pure
    (respectively $(Z\cap V_\lambda)$-closed).
  \item\label{IV.5.9.10.ii}
    Let $Z'\subset Z$ be stable under specialization.  If $\sh F$ is
    $Z$-pure (respectively $Z$-closed), then it is $Z'$-pure
    (respectively $Z'$-closed).
\end{enumerate}
\end{corollary}

\begin{proof}
This follows immediately from condition b') of \sref{IV.5.9.8}.
\end{proof}

\begin{proposition}[5.9.11]
\label{IV.5.9.11}
Under the hypotheses of \sref{IV.5.9.8}, the supports of
$\ker(\rho_{X/Z})$ and $\operatorname{coker}(\rho_{X/Z})$ are
contained in $Z$, and $\sh H^0_{X/Z}(\sh F)$ is $Z$-closed.

Furthermore, if $u:\sh F\to\sh F'$ is an $\sh O_X$-module
homomorphism and $\sh F'$ is $Z$-closed, then $u$ factors uniquely as
\[
 \sh F\xrightarrow{\rho_{X/Z}}\sh H^0_{X/Z}(\sh F)
 \xrightarrow{v}\sh F'.
\]
If, in addition, the supports of $\ker(u)$ and $\operatorname{coker}(u)$
are contained in $Z$, then $v$ is an isomorphism.
\end{proposition}

\begin{proof}
The first assertion says that, for every $x\in X-Z$,
\[
 \ker(\rho_{X/Z})_x=\operatorname{coker}(\rho_{X/Z})_x=0.
\]
Equivalently, $(\rho_{X/Z})_x$ is bijective, or $\sh F$ is
$Z$-closed at $x$, as noted after \sref{IV.5.9.9}.

To prove that $\sh H^0_{X/Z}(\sh F)$ is $Z$-closed, it is enough to
show, for every Noetherian open $V\subset X$, that
\[
 \varinjlim_\beta
 \Gamma(V\cap U_\beta,\sh H^0_{X/Z}(\sh F))
 =\Gamma(V,\sh H^0_{X/Z}(\sh F)).
\]
By definition,
\[
 \Gamma(V\cap U_\beta,\sh H^0_{X/Z}(\sh F))
 =\varinjlim_\alpha
 \Gamma(V\cap U_\alpha\cap U_\beta,\sh F).
\]
The double family $(U_\alpha\cap U_\beta)$ is filtered decreasing, so
the assertion follows from \sref{IV.5.9.1} and the theorem on double
inductive limits.

\oldpage[IV-2]{113}

We now prove the second part.  For every $\alpha$,
$(i_\alpha)_*(u)$ is the unique homomorphism making the diagram
\[
\xymatrix{
\sh F \ar[r]^{u} \ar[d]
  & \sh F' \ar[d] \\
 (i_\alpha)_*(\sh F|U_\alpha)
   \ar[r]_{(i_\alpha)_*(u)}
   & (i_\alpha)_*(\sh F'|U_\alpha)
}
\]
commute.  Hence there is a unique homomorphism $w$ making all diagrams
\[
\xymatrix{
 (i_\alpha)_*(\sh F|U_\alpha)
   \ar[r]^{(i_\alpha)_*(u)} \ar[d]
   & (i_\alpha)_*(\sh F'|U_\alpha) \ar[d] \\
 \sh H^0_{X/Z}(\sh F) \ar[r]_{w}
   & \sh H^0_{X/Z}(\sh F')
}
\]
commute.  Since $\sh F'$ is $Z$-closed, it identifies canonically
with $\sh H^0_{X/Z}(\sh F')$.  This proves existence and uniqueness
of $v$.

Suppose finally that the supports of $\ker(u)$ and
$\operatorname{coker}(u)$ are contained in $Z$.  We show that, for
every Noetherian open $V\subset X$, the homomorphism
\[
 \Gamma(V,\sh H^0_{X/Z}(\sh F))\longrightarrow\Gamma(V,\sh F')
\]
is an isomorphism.  If a section
$t\in\Gamma(V,\sh H^0_{X/Z}(\sh F))$ maps to zero, then, for some
$\alpha$, it is represented by a section on $V\cap U_\alpha$.
By the hypothesis on $u$, its germ vanishes at every point of
$V\cap(X-Z)$.  It therefore vanishes on an open neighbourhood of
$V\cap(X-Z)$, and so represents zero in
$\Gamma(V,\sh H^0_{X/Z}(\sh F))$.

Conversely, let $s'\in\Gamma(V,\sh F')$.  For every
$x\in V\cap(X-Z)$, there is an open neighbourhood $W^{(x)}$ of $x$
and a section $s^{(x)}$ of $\sh F$ on $W^{(x)}$ whose image under
$u$ is $s'|W^{(x)}$.  Hence $s'|W^{(x)}$ is also the image under $v$
of a section $t^{(x)}$ of $\sh H^0_{X/Z}(\sh F)$.  Since $v$ is
injective, these sections agree on overlaps.  They therefore glue to
a section $t$ over an open subset $U\subset V$ containing
$V\cap(X-Z)$.  As $\sh H^0_{X/Z}(\sh F)$ is $Z$-closed, $t$ extends
uniquely over $V$; its image under $v$ agrees with $s'$ on $U$, and
hence on $V$ by the same argument.

The module $\sh H^0_{X/Z}(\sh F)$ is called the \emph{$Z$-closure}
of $\sh F$.
\end{proof}

\oldpage[IV-2]{114}

\begin{remarks}[5.9.12]
\label{IV.5.9.12}
\begin{enumerate}[{\rm(i)}]
  \item\label{IV.5.9.12.i}
    Let $\mathcal C(X)$ be the category of $\sh O_X$-modules, and let
    $\mathcal C_Z(X)$ be its full subcategory of modules whose support
    is contained in $Z$.  This is a localizing subcategory in the
    sense of Gabriel, and $\sh H^0_{X/Z}$ is precisely Gabriel's
    localization functor (cf.~\textup{[27]}); this gives another proof of
    \sref{IV.5.9.11}.  When $Z$ is closed, the functor
    $i^*:\mathcal C(X)\to\mathcal C(X-Z)$, where
    $i:X-Z\to X$, defines an equivalence
    \[
      \mathcal C(X)/\mathcal C_Z(X)\simeq\mathcal C(X-Z).
    \]
  \item\label{IV.5.9.12.ii}
    By \sref{IV.5.9.11},
    $\sh H^0_{X/Z}(\sh F)=0$ if and only if
    $\operatorname{Supp}(\sh F)\subset Z$.  Indeed, the vanishing
    implies that $\ker(\rho_{X/Z})=\sh F$.  Conversely, if
    $\operatorname{Supp}(\sh F)\subset Z$, apply the second part of
    \sref{IV.5.9.11} to the unique map $u:\sh F\to0$; the
    corresponding map $v:\sh H^0_{X/Z}(\sh F)\to0$ is an
    isomorphism.
  \item\label{IV.5.9.12.iii}
    The preceding constructions make sense for every locally
    Noetherian ringed space in which every closed irreducible subset
    has a unique generic point.  In particular, on a locally
    Noetherian prescheme they apply to arbitrary sheaves of abelian
    groups, regarded as modules over the sheaf associated with the
    constant presheaf $\mathbf Z$.  For every $x\in X$ one still has
    the canonical isomorphism \eqref{IV.5.9.6.1}, in which
    $\widetilde{\sh F}_x$ denotes the sheaf induced by $\sh F$ on the
    subspace $X_x$ of $X$.  A direct proof follows from
    \eqref{IV.5.9.1.2} and the theorem on double inductive limits.
\end{enumerate}
\end{remarks}
% END NATIVE INPUT components/09_section59_z_pure_z_closed_modules
% BEGIN NATIVE INPUT components/10a_section510_s2_z_closure_through_corollary5109
\subsection{Property \texorpdfstring{$(\mathrm{S}_2)$}{(S2)} and
\texorpdfstring{$Z$}{Z}-closure}
\label{subsection:IV.5.10}

\begin{env}[5.10.1]
\label{IV.5.10.1}
Let $X$ be a locally Noetherian prescheme and let $\sh F$ be a
coherent $\sh O_X$-module.  For every subset $T$ of $X$, put
\end{env}

\begin{equation}
\label{IV.5.10.1.1}
\operatorname{prof}_T(\sh F)
 =\inf_{z\in T}\operatorname{prof}(\sh F_z).
\tag{5.10.1.1}
\end{equation}

\begin{proposition}[5.10.2]
\label{IV.5.10.2}
Let $X$ be a locally Noetherian prescheme, let $Z$ be a subset of
$X$ stable under specialization, and let $\sh F$ be a
quasi-coherent $\sh O_X$-module.  The following conditions are
equivalent:
\begin{enumerate}[{\rm a)}]
  \item $\sh F$ is $Z$-pure.
  \item $\operatorname{Ass}(\sh F)$ does not meet $Z$.
\end{enumerate}
If, in addition, $\sh F$ is coherent, these conditions are also
equivalent to:
\begin{enumerate}
  \item[\rm c)] $\operatorname{prof}_Z(\sh F)\geq1$.
\end{enumerate}
\end{proposition}

\begin{proof}
To say that $\sh F$ is $Z$-pure means that, for every Noetherian
open subset $V$ of $X$ and every open subset $U\supset X-Z$, the
restriction homomorphism
\[
  \Gamma(V,\sh F)\longrightarrow\Gamma(V\cap U,\sh F)
\]
is injective \sref{IV.5.9.8}.  By \sref{IV.3.1.8}, this is
equivalent to
$V\cap\operatorname{Ass}(\sh F)\subset U$, which proves the
equivalence of {\rm a)} and {\rm b)}.  Further, to say that
$x\in\operatorname{Ass}(\sh F)$ means that no element of
$\mathfrak m_x$ is $\sh F_x$-regular \sref{IV.3.1.2}; when
$\sh F$ is coherent, this is equivalent to
$\operatorname{prof}(\sh F_x)=0$.  The equivalence of {\rm b)} and
{\rm c)} follows immediately.
\end{proof}

\begin{corollary}[5.10.3]
\label{IV.5.10.3}
Let
\[
  0\longrightarrow\sh F'\longrightarrow\sh F
   \longrightarrow\sh F''\longrightarrow0
\]
be an exact sequence of quasi-coherent $\sh O_X$-modules.  If
$\sh F$ is $Z$-pure, then so is $\sh F'$.  Conversely, if $\sh F'$
and $\sh F''$ are $Z$-pure, then so is $\sh F$.
\end{corollary}

\oldpage[IV-2]{115}

\begin{proof}
This follows from the form \sref{IV.5.10.2}, {\rm b)}, of the
condition for a quasi-coherent $\sh O_X$-module to be $Z$-pure,
and from \sref{IV.3.1.7}.
\end{proof}

\begin{corollary}[5.10.4]
\label{IV.5.10.4}
Suppose that $\sh F$ is coherent.  For $\sh F$ to be $Z$-pure at
a point $x\in X$, it is necessary and sufficient that
\[
  \operatorname{prof}_{Z_x}(\widetilde{\sh F}_x)\geq1,
\]
with the notation of \sref{IV.5.9.6}.
\end{corollary}

\begin{proof}
This follows immediately from \sref{IV.5.10.2} and
\sref{IV.5.9.6}.
\end{proof}

\begin{theorem}[5.10.5]
\label{IV.5.10.5}
Let $X$ be a locally Noetherian prescheme, let $Z$ be a subset of
$X$ stable under specialization, and let $\sh F$ be a coherent
$\sh O_X$-module.  For $\sh F$ to be $Z$-closed, it is necessary
and sufficient that
\[
  \operatorname{prof}_Z(\sh F)\geq2.
\]
\end{theorem}

\begin{proof}
By \sref{IV.5.10.2}, we may restrict ourselves to the case in
which $\sh F$ is $Z$-pure and
$\operatorname{prof}_Z(\sh F)\geq1$.  Moreover, to say that
$\operatorname{prof}_Z(\sh F)\geq2$ is equivalent to saying that
\[
  \operatorname{prof}_{Z_\alpha}(\sh F)\geq2
\]
for every closed subset $Z_\alpha$ of $Z$; similarly, by
\sref{IV.5.9.8}, to say that $\sh F$ is $Z$-closed is equivalent
to saying that it is $Z_\alpha$-closed for every $\alpha$.  We
may therefore first reduce to the case in which $Z$ is closed.

The question is local.  For each $x\in Z$, it is enough to prove
the theorem for $\sh F|U$, where $U$ is an affine open
neighbourhood of $x$; we may thus suppose that $X=U$ is affine.
Then $\operatorname{Ass}(\sh F)$ is finite \sref{IV.3.1.6}, and,
since $\operatorname{Ass}(\sh F)\subset X-Z$, there is a section
$f$ of $\sh O_X$ over $X$ such that
\[
  \operatorname{Ass}(\sh F)\subset X_f\subset X-Z
  \qquad\text{\sref[II]{II.4.5.4}}.
\]
It follows that $f$ is $\sh F$-regular \sref{IV.3.1.9}, and for
every $y\in Z$ one has $f_y\in\mathfrak m_y$, hence
\[
  \operatorname{prof}(\sh F_y)
   =1+\operatorname{prof}(\sh F_y/f_y\sh F_y)
  \qquad\text{\sref[0]{0.16.4.6}}.
\]
Thus $\operatorname{prof}_Z(\sh F)\geq2$ is equivalent to
$\operatorname{prof}_Z(\sh F/f\sh F)\geq1$, or, by
\sref{IV.5.10.2}, to $\sh F/f\sh F$ being $Z$-pure.  It remains
to show that this last property is equivalent to $\sh F$ being
$Z$-closed.

Consider the exact sequence
\[
  0\longrightarrow\sh F\xrightarrow{\,f\,}\sh F
   \longrightarrow\sh F/f\sh F\longrightarrow0,
\]
where multiplication by $f$ is injective by hypothesis.  Put
$W=X-Z$.  We have the commutative diagram
\[
\xymatrix{
  0 \ar[r] &
  \Gamma(X,\sh F) \ar[r]^{f} \ar[d] &
  \Gamma(X,\sh F) \ar[r] \ar[d] &
  \Gamma(X,\sh F/f\sh F) \ar[r] \ar[d] &
  0 \\
  0 \ar[r] &
  \Gamma(W,\sh F) \ar[r]^{f} &
  \Gamma(W,\sh F) \ar[r] &
  \Gamma(W,\sh F/f\sh F) &
}
\]
whose two rows are exact, since $X$ is affine.  If the
restriction homomorphism
\[
  \Gamma(X,\sh F)\longrightarrow\Gamma(W,\sh F)
\]
is bijective, the diagram shows that
\[
  \Gamma(X,\sh F/f\sh F)
   \longrightarrow\Gamma(W,\sh F/f\sh F)
\]
is injective; hence $\sh F/f\sh F$ is $Z$-pure
\sref{IV.5.9.8}.

Conversely, suppose that $\sh F/f\sh F$ is $Z$-pure, and let $s$
be a section of $\sh F$ over $W$.  Since $X_f\subset W$, there is
an integer $n>0$ such that $f^n(s|X_f)$ extends to a section $t$
of $\sh F$ over $X$ \sref[I]{I.1.4.1}.  The restrictions of $t$
and $f^ns$ to $X_f$ are equal; the relation
$\operatorname{Ass}(\sh F)\subset X_f$ therefore implies, by
\sref{IV.5.10.2}, that $t|W=f^ns$.  Since $f$ is
$\sh F$-regular, it is enough to show that $t=f^nt'$ for some
$t'\in\Gamma(X,\sh F)$.

\oldpage[IV-2]{116}

This is equivalent to the image of $t$ in
$\Gamma(X,\sh F/f^n\sh F)$ being zero.  For every $k$,
$f^k\sh F/f^{k+1}\sh F$ is isomorphic to $\sh F/f\sh F$, hence is
$Z$-pure.  By induction on $n$, using \sref{IV.5.10.3}, it follows
that $\sh F/f^n\sh F$ is $Z$-pure.  But the image of
$t|W=f^ns$ in $\Gamma(W,\sh F/f^n\sh F)$ is zero; the assertion
follows from $Z$-purity.
\end{proof}

\begin{corollary}[5.10.6]
\label{IV.5.10.6}
Let $\sh F$ be a coherent $\sh O_X$-module.  For $\sh F$ to be
$Z$-closed at a point $x$, it is necessary and sufficient that
\[
  \operatorname{prof}_{Z_x}(\widetilde{\sh F}_x)\geq2.
\]
\end{corollary}

\begin{proof}
This follows from \sref{IV.5.9.6} and \sref{IV.5.10.5}.
\end{proof}

\begin{theorem}[5.10.7]
\label{IV.5.10.7}
\emph{(Hartshorne).}
Let $X$ be a locally Noetherian prescheme and let $Y$ be a closed
subset of $X$.  Suppose that
$\operatorname{prof}(\sh O_{X,y})\geq2$ for every $y\in Y$.
Then, for every connected component $C$ of $X$,
$C-(C\cap Y)$ is connected.
\end{theorem}

\begin{proof}
We may restrict ourselves to the case in which $X$ is connected.
By \sref{IV.5.10.5}, the canonical homomorphism
\[
  \sh O_X\longrightarrow i_*(\sh O_X|X-Y),
\]
where $i:X-Y\to X$ is the canonical injection, is bijective.
Consequently the restriction homomorphism
\[
  \Gamma(X,\sh O_X)\longrightarrow\Gamma(X-Y,\sh O_X)
\]
is also bijective.  It remains to apply
\sref[III]{III.7.8.6.1}.
\end{proof}

\begin{corollary}[5.10.8]
\label{IV.5.10.8}
Let $X$ be a locally Noetherian prescheme and let $d$ be an
integer such that, for every $x\in X$,
\[
  \dim(\sh O_x)\geq d
  \quad\Longrightarrow\quad
  \operatorname{prof}(\sh O_x)\geq2.
\]
Suppose that $X$ is connected.  If $X'$ and $X''$ are two
distinct irreducible components of $X$, there is a sequence
$(X_i)_{0\leq i\leq n}$ of irreducible components of $X$ such
that $X_0=X'$, $X_n=X''$, and
\[
  \codim(X_{i-1}\cap X_i,X)\leq d-1
  \qquad(1\leq i\leq n).
\]
In this case one says that $X$ is
\emph{connected in codimension $\leq d-1$}.
\end{corollary}

\begin{proof}
If $Y$ is a closed subset of $X$ such that
$\codim(Y,X)\geq d$, then
$\dim(\sh O_{X,y})\geq d$ for every $y\in Y$
\sref{IV.5.1.3}; hence
$\operatorname{prof}(\sh O_{X,y})\geq2$ for every $y\in Y$, and
\sref{IV.5.10.7} shows that $X-Y$ is connected.

For $\codim(Y,X)\geq d$, it is necessary and sufficient that,
for every $y\in Y$, there be an open neighbourhood $V$ of $y$
such that $\codim(Y\cap V,V)\geq d$
\sref[0]{0.14.2.3}.  Let $\mathfrak Y$ be the set of closed
subsets $Y$ of $X$ of codimension at least $d$.  The union of two
members of $\mathfrak Y$ belongs to $\mathfrak Y$
\sref[0]{0.14.2.5}, and every closed subset of a member of
$\mathfrak Y$ also belongs to $\mathfrak Y$.  In other words,
$\mathfrak Y$ is an \emph{antifilter} of closed subsets of $X$.
The corollary now follows from the following topological lemma.
\end{proof}

\begin{lemma}[5.10.8.1]
\label{IV.5.10.8.1}
Let $X$ be a connected locally Noetherian topological space and
let $\mathfrak Y$ be an antifilter of closed subsets of $X$.
Suppose that a closed subset $Y$ belongs to $\mathfrak Y$ whenever,
for every $y\in Y$, there exist an open neighbourhood $V$ of $y$
in $X$ and a $Y_y\in\mathfrak Y$ such that
$V\cap Y=V\cap Y_y$.  The following conditions are equivalent:
\begin{enumerate}[{\rm a)}]
  \item For every $Y\in\mathfrak Y$, the space $X-Y$ is connected.
  \item If $X'$ and $X''$ are two distinct irreducible components
  of $X$, there is a sequence $(X_i)_{0\leq i\leq n}$ of
  irreducible components of $X$ such that $X_0=X'$, $X_n=X''$,
  and
  \[
    X_{i-1}\cap X_i\notin\mathfrak Y
    \qquad(1\leq i\leq n).
  \]
\end{enumerate}
\end{lemma}

\begin{proof}
Suppose {\rm b)} holds and let $Y\in\mathfrak Y$.  Put
$U=X-Y$.  If $U'$ and $U''$ are two distinct irreducible
components of $U$, there are irreducible components $X'$ and
$X''$ of $X$ such that
$X'\cap U=U'$ and $X''\cap U=U''$ \sref[0]{0.2.1.6}.

\oldpage[IV-2]{117}

Choose a sequence $(X_i)$ with the property in {\rm b)}, and put
$U_i=X_i\cap U$ for $0\leq i\leq n$.  Each $U_i$ is an
irreducible component of $U$ \sref[0]{0.2.1.6}, and
$U_i\cap U_{i-1}\ne\varnothing$ for $1\leq i\leq n$: otherwise
$X_i\cap X_{i-1}\subset Y$, so
$X_i\cap X_{i-1}\in\mathfrak Y$, contrary to the choice of the
$X_i$.  Thus $U$ is connected.

Conversely, suppose {\rm a)} holds.  Let $Y$ be the union of the
sets $X_\alpha\cap X_\beta$, where $(X_\alpha,X_\beta)$ ranges
over the pairs of distinct irreducible components of $X$ for
which $X_\alpha\cap X_\beta\in\mathfrak Y$.  Every $y\in Y$ has
an open neighbourhood $V$ meeting only finitely many irreducible
components of $X$.  It follows both that $Y$ is closed and that
$V\cap Y$ is the intersection of $V$ with a member of
$\mathfrak Y$.  By hypothesis, $Y\in\mathfrak Y$; hence
$U=X-Y$ is connected.  Moreover, $Y$ is nowhere dense in $X$.

Let $X'$ and $X''$ be distinct irreducible components of $X$, and
let $U'$ and $U''$ be their respective traces on $U$.  They are
distinct irreducible components of $U$ \sref[0]{0.2.1.6}.
The union of those irreducible components $W$ of $U$ for which
there is a sequence $(U_i)_{0\leq i\leq n}$ of irreducible
components of $U$ satisfying
\[
  U_0=U',\qquad
  U_{i-1}\cap U_i\ne\varnothing\ (1\leq i\leq n),\qquad
  U_n=W,
\]
is both open and closed in $U$, because $U$ is locally
Noetherian and its irreducible components form a locally finite
family of closed subsets.  Since $U$ is connected, there is such
a sequence ending at $U''$.  Let $X_i$ be the irreducible
component of $X$ such that $X_i\cap U=U_i$
\sref[0]{0.2.1.6}.  Since $U_{i-1}\ne U_i$, one has
$X_{i-1}\ne X_i$.  If
$X_{i-1}\cap X_i\in\mathfrak Y$ for some $i$, then
$X_{i-1}\cap X_i\subset Y$, contrary to
$U_{i-1}\cap U_i\ne\varnothing$.  This proves {\rm b)}.
\end{proof}

The hypothesis on $X$ in \sref{IV.5.10.8} is satisfied when $X$
is a Cohen--Macaulay prescheme and $d\geq2$.

\begin{corollary}[5.10.9]
\label{IV.5.10.9}
If a Noetherian local ring $A$ satisfies $(\mathrm S_2)$ and is
catenary, then it is equidimensional.
\end{corollary}

\begin{proof}
The hypothesis of \sref{IV.5.10.8} is satisfied by
$X=\Spec(A)$ with $d=2$.  By that corollary, to show that all
irreducible components of $X$ have the same dimension, it is
enough to consider two such components $X'$ and $X''$ with
\[
  \codim(X'\cap X'',X)=1.
\]
There is then an irreducible component $Z$ of $X'\cap X''$ such
that $\codim(Z,X)=1$.  Since
\[
  1=\codim(Z,X)\geq\codim(Z,X')\geq1,
\]
one has $\codim(Z,X')=1$; similarly
$\codim(Z,X'')=1$.  Since $X$ is catenary, this implies
$\dim(X')=\dim(X'')$.
\end{proof}
% END NATIVE INPUT components/10a_section510_s2_z_closure_through_corollary5109
% BEGIN NATIVE INPUT components/10b_section510_z_closure_proposition51010_through_51012
\begin{proposition}[5.10.10]
\label{IV.5.10.10}
Let $X$ be a locally Noetherian prescheme, let $Z$ be a subset of
$X$ stable under specialization, and let $\sh F$ be a coherent
$\sh O_X$-module.  Suppose that the $\sh O_X$-module
$\sh H^0_{X/Z}(\sh F)$ is coherent.  Then:
\begin{enumerate}[{\rm(i)}]
  \item
  \[
    \operatorname{prof}_Z\bigl(\sh H^0_{X/Z}(\sh F)\bigr)\geq2.
  \]
  \item For every
  $x\in\operatorname{Ass}(\sh F)\cap(X-Z)$,
  \[
    \codim\bigl(Z\cap\overline{\{x\}},\overline{\{x\}}\bigr)
    \geq2.
  \]
  \item The set $U$ of points $x\in X$ such that
  \[
    \operatorname{prof}_{Z_x}(\widetilde{\sh F}_x)\geq2
  \]
  (with the notation of \sref{IV.5.9.6}) is open in $X$;
  moreover, $X-U\subset Z$, and $U$ is the largest open subset of
  $X$ such that $\sh F|U$ is $(Z\cap U)$-closed.
\end{enumerate}
\end{proposition}

\oldpage[IV-2]{118}

\begin{proof}
Put, for brevity,
\[
  \sh F'=\sh H^0_{X/Z}(\sh F).
\]
By \sref{IV.5.9.11}, $\sh F'$ is $Z$-closed, so assertion
{\rm(i)} follows from \sref{IV.5.10.5}.

Let $x\in\operatorname{Ass}(\sh F)\cap(X-Z)$.  The restrictions of
$\sh F$ and $\sh F'$ to $X-Z$ are canonically isomorphic
\sref{IV.5.9.11}; hence $x\in\operatorname{Ass}(\sh F')$.
For $y\in Z\cap\overline{\{x\}}$, let $\mathfrak p$ be the prime
ideal of $\sh O_y$ corresponding to $x$.  It is associated to the
$\sh O_y$-module $\sh F'_y$.  Consequently, by {\rm(i)} and
\sref[0]{0.16.4.6.2},
\[
  2\leq\operatorname{prof}(\sh F'_y)
   \leq\dim(\sh O_y/\mathfrak p)
   =\codim\bigl(\overline{\{y\}},\overline{\{x\}}\bigr),
\]
which proves {\rm(ii)}.

Finally, by \sref{IV.5.10.5}, $U$ is the set of points $x$ for
which $\widetilde{\sh F}_x$ is $Z_x$-closed; by
\sref{IV.5.9.6}, it is equivalently the set of points at which
the canonical homomorphism
\[
  \sh F_x\longrightarrow\sh F'_x
\]
is bijective.  Thus $X-U$ is the union of the supports of
$\operatorname{Ker}(\rho_{X/Z})$ and
$\operatorname{Coker}(\rho_{X/Z})$.  These are coherent
$\sh O_X$-modules by the hypothesis and \sref[0]{0.5.3.4}, so
their supports are closed \sref[0]{0.5.2.2}.  Hence $U$ is open.
It is plainly the largest open subset on which $\sh F$ is
$(Z\cap U)$-closed, and $X-U\subset Z$ follows from
\sref{IV.5.9.11}.
\end{proof}

We shall see in \sref{IV.5.11.1} that, in the most important
cases, assertion {\rm(ii)} conversely implies that
$\sh H^0_{X/Z}(\sh F)$ is coherent.

\begin{env}[5.10.11]
\label{IV.5.10.11}
Let $X$ be a locally Noetherian prescheme and let $Z$ be a subset
of $X$ stable under specialization.  We saw in \sref{IV.5.9.3}
that
\[
  \mathscr A=\sh H^0_{X/Z}(\sh O_X)
\]
is a quasi-coherent $\sh O_X$-algebra.  The $X$-scheme
\[
  X'=\Spec(\mathscr A)
\]
\sref[II]{II.1.3.1} is called the \emph{$Z$-closure} of $X$.

For every $\sh O_X$-module $\sh F$,
$\sh H^0_{X/Z}(\sh F)$ is an $\mathscr A$-module, and it is
quasi-coherent when $\sh F$ is quasi-coherent.  In that case there
is a unique quasi-coherent $\sh O_{X'}$-module $\sh F'$ such that,
if $g:X'\to X$ denotes the structural morphism, then
\end{env}

\begin{equation}
\label{IV.5.10.11.1}
  \sh H^0_{X/Z}(\sh F)=g_*(\sh F').
\tag{5.10.11.1}
\end{equation}

\begin{proposition}[5.10.12]
\label{IV.5.10.12}
Use the notation of \sref{IV.5.10.11}.
\begin{enumerate}[{\rm(i)}]
  \item Let $x\in X$.  The morphism obtained from $g$ by
  localization,
  \[
    X'\times_X X_x\longrightarrow
    X_x=\Spec(\sh O_{X,x}),
  \]
  is an isomorphism if and only if $\sh O_X$ is $Z$-closed at
  $x$.  This condition holds for every $x\in X-Z$.

  \item Put $Z'=g^{-1}(Z)$ and suppose that $X'$ is locally
  Noetherian.  Then $\sh F'$ is $Z'$-closed.  If, in addition,
  $\sh F'$ is a coherent $\sh O_{X'}$-module, then
  \[
    \operatorname{prof}_{Z'}(\sh F')\geq2.
  \]

  \item Suppose that both
  $\sh H^0_{X/Z}(\sh O_X)$ and
  $\sh H^0_{X/Z}(\sh F)$ are coherent.  Then
  $g:X'\to X$ is finite.  The set $U$ of points $x\in X$ such
  that
  \[
    \operatorname{prof}_{Z_x}(\widetilde{\sh O}_x)\geq2
    \quad\text{and}\quad
    \operatorname{prof}_{Z_x}(\widetilde{\sh F}_x)\geq2
  \]
  is open in $X$ and satisfies $X-U\subset Z$.  Moreover, $U$ is
  the largest open subset of $X$ such that
  $g^{-1}(U)\to U$ is an isomorphism and the restriction to $U$
  of the canonical $g$-morphism from $\sh F$ to $\sh F'$ is an
  isomorphism.
\end{enumerate}
\end{proposition}

\begin{proof}
Assertion {\rm(i)} follows from the definitions, and {\rm(iii)}
is an immediate consequence of \sref{IV.5.10.10}, {\rm(iii)}.

For {\rm(ii)}, let $V$ be an open subset of $X$ containing
$X-Z$, put $V'=g^{-1}(V)$, and let
$i:V\to X$ and $i':V'\to X'$ be the canonical injections.

\oldpage[IV-2]{119}

The canonical homomorphism
\[
  \rho_{X'/Z'}:\sh F'\longrightarrow i'_*(\sh F'|V')
\]
has the property that $g_*(\rho_{X'/Z'})$ is the canonical
homomorphism
\[
  \rho_{X/Z}:\sh H^0_{X/Z}(\sh F)
   \longrightarrow
   i_*\bigl(\sh H^0_{X/Z}(\sh F)|V\bigr),
\]
by \sref[II]{II.1.4.2} and \sref{IV.5.10.11.1}.
Since $\sh H^0_{X/Z}(\sh F)$ is $Z$-closed
\sref{IV.5.9.11}, $\rho_{X/Z}$ is an isomorphism, and therefore
so is $\rho_{X'/Z'}$.  The open subset $X'-Z'$ is the
intersection of the filtered family of inverse images
$g^{-1}(V)$, as $V$ ranges over the open subsets of $X$
containing $X-Z$.  Since $X'$ is locally Noetherian,
\sref{IV.5.9.1} now shows that $\sh F'$ is $Z'$-closed.
The final assertion follows from \sref{IV.5.10.5}.
\end{proof}
% END NATIVE INPUT components/10b_section510_z_closure_proposition51010_through_51012
% BEGIN NATIVE INPUT components/10c_section510_s2_closure_applications_through_51017
\begin{env}[5.10.13]
\label{IV.5.10.13}
We shall now apply the preceding results when $Z$ is one of the
sets $Z^{(n)}(X)$ (or simply $Z^{(n)}$), defined as the set of
points $x\in X$ such that
\[
  \dim(\sh O_x)\geq n.
\]
It is clear that $Z^{(n)}$ is stable under specialization.  A
closed subset $T$ of $X$ is contained in $Z^{(n)}$ if and only if
\[
  \codim(T,X)\geq n.
\]
We shall be concerned here with the case $n=2$.
\end{env}

\begin{proposition}[5.10.14]
\label{IV.5.10.14}
Let $X$ be a locally Noetherian prescheme, and let $\sh F$ be a
coherent $\sh O_X$-module whose support is $X$.
\begin{enumerate}[{\rm(i)}]
  \item The module $\sh F$ satisfies $(\mathrm S_1)$ if and only
  if it is $Z^{(1)}$-pure.
  \item The module $\sh F$ satisfies $(\mathrm S_2)$ if and only
  if it is $Z^{(2)}$-closed and $Z^{(1)}$-pure, or equivalently,
  if it is $Z^{(2)}$-closed and has no associated prime cycle of
  codimension $1$.
\end{enumerate}
\end{proposition}

\begin{proof}
To say that $\sh F$ has property $(\mathrm S_1)$ means that it
has no embedded associated prime cycle \sref{IV.5.7.5}, or
equivalently that
\[
  \dim(\sh F_x)=0
  \qquad\text{for every }x\in\operatorname{Ass}(\sh F)
\]
\sref{IV.3.1.4}.  Since $\operatorname{Supp}(\sh F)=X$, this is
equivalent to $\dim(\sh O_x)=0$ \sref{IV.5.1.12.1}; in other
words, $\operatorname{Ass}(\sh F)$ does not meet $Z^{(1)}$.
Assertion {\rm(i)} therefore follows from \sref{IV.5.10.2}.

To say that $\sh F$ is $Z^{(2)}$-closed means that
\[
  \operatorname{prof}_{Z^{(2)}}(\sh F)\geq2,
\]
or equivalently that, for every $x\in X$,
\[
  \dim(\sh O_x)\geq2
  \quad\Longrightarrow\quad
  \operatorname{prof}(\sh F_x)\geq2.
\]
Thus $(\mathrm S_2)$ implies that $\sh F$ is $Z^{(2)}$-closed.
It also implies $(\mathrm S_1)$, so $\sh F$ has no embedded
associated prime cycle \sref{IV.5.7.5}; since its support is
$X$, this says that all its associated prime cycles have
codimension $0$.

Conversely, suppose that $\sh F$ is $Z^{(2)}$-closed and has no
associated prime cycle of codimension $1$.  It remains only to
consider a point $x$ for which $\dim(\sh O_x)=1$.  By hypothesis
$x\notin\operatorname{Ass}(\sh F)$, and hence
$\operatorname{prof}(\sh F_x)>0$.  Since
$\dim(\sh F_x)=\dim(\sh O_x)=1$, it follows that
$\operatorname{prof}(\sh F_x)=1$.  Therefore $\sh F$ satisfies
$(\mathrm S_2)$.  Finally, if $\sh F$ is $Z^{(1)}$-pure, it
satisfies $(\mathrm S_1)$ by {\rm(i)}; as observed above, all its
associated prime cycles then have codimension $0$, so the same
argument applies.

Notice that a module may be $Z^{(2)}$-closed without satisfying
$(\mathrm S_1)$.  For example, if $\dim X=1$, then
$Z^{(2)}=\varnothing$ and every $\sh O_X$-module is
$Z^{(2)}$-closed, even though a module may have embedded
associated prime cycles.

We recall that Chapter~III gives, by means of local cohomology, a
characterization of $(\mathrm S_n)$ for every $n\geq1$ that
generalizes \sref{IV.5.10.14}.
\end{proof}

\begin{corollary}[5.10.15]
\label{IV.5.10.15}
Let $X$ be a locally Noetherian prescheme, and let $\sh F$ be a
coherent $\sh O_X$-module with support $X$.  Suppose that $\sh F$
has no associated prime cycle of codimension $1$ and that
\[
  \sh F'=\sh H^0_{X/Z^{(2)}}(\sh F)
\]
is coherent.  Then:
\begin{enumerate}[{\rm(i)}]
  \item $\sh F'$ satisfies $(\mathrm S_2)$.

\oldpage[IV-2]{120}

  \item The set $U$ of points $x\in X$ at which $\sh F$
  satisfies $(\mathrm S_2)$ \sref{IV.5.7.2} is open in $X$, and
  \[
    \codim(X-U,X)\geq2.
  \]
\end{enumerate}
\end{corollary}

\begin{proof}
By \sref{IV.5.9.11}, $\sh F'$ is $Z^{(2)}$-closed.  Moreover,
$\operatorname{Supp}(\sh F')=X$: every maximal point of $X$
lies outside $Z^{(2)}$, and at such a point
$\sh F'_x=\sh F_x\neq0$.  Thus the support of $\sh F'$ is dense
in $X$; it is also closed because $\sh F'$ is coherent, and
hence it is all of $X$.

It remains to see that $\sh F'$ has no associated prime cycle of
codimension $1$.  If $x\in\operatorname{Ass}(\sh F')$ and
\[
  \dim(\sh F'_x)=\dim(\sh O_x)=1,
\]
then $x\notin Z^{(2)}$, so $\sh F'_x=\sh F_x$ and consequently
$x\in\operatorname{Ass}(\sh F)$, contrary to the hypothesis.
Assertion {\rm(i)} now follows from \sref{IV.5.10.14}.

For {\rm(ii)}, one has
\[
  Z^{(n)}(X_x)=Z^{(n)}(X)\cap X_x
\]
with the notation of \sref{IV.5.9.6}, by \sref[I]{I.2.4.2}.
The absence of associated prime cycles of codimension $1$ for
$\sh F$ implies the same property for
$\widetilde{\sh F}_x$.  By \sref{IV.5.10.14},
$\widetilde{\sh F}_x$ satisfies $(\mathrm S_2)$ if and only if
it is $Z^{(2)}(X_x)$-closed.  The assertion therefore follows
from \sref{IV.5.10.6} and \sref{IV.5.10.10}, {\rm(iii)}.
\end{proof}

\begin{proposition}[5.10.16]
\label{IV.5.10.16}
Let $X$ be a locally Noetherian prescheme, let
\[
  X'=\Spec\bigl(\sh H^0_{X/Z^{(2)}}(\sh O_X)\bigr)
\]
be its $Z^{(2)}$-closure, and let $g:X'\to X$ be the structural
morphism.  Suppose that $X$ has no associated prime cycle of
codimension $1$.
\begin{enumerate}[{\rm(i)}]
  \item At a point $x\in X$, the prescheme $X$ satisfies
  $(\mathrm S_2)$ if and only if the morphism
  \[
    X'_x\longrightarrow X_x
  \]
  deduced from $g$ (with the notation of \sref{IV.5.10.12}) is
  an isomorphism.  This condition always holds if
  \[
    \codim\bigl(\overline{\{x\}},X\bigr)\leq1.
  \]
  \item Suppose in addition that $g$ is finite (see
  \sref{IV.5.11.2} for sufficient conditions).  Then the set
  $U$ of points at which $X$ satisfies $(\mathrm S_2)$ is open,
  \[
    \codim(X-U,X)\geq2,
  \]
  and $U$ is the largest open subset of $X$ for which
  $g^{-1}(U)\to U$ is an isomorphism.
  \item Under the hypotheses of {\rm(ii)}, $X'$ satisfies
  $(\mathrm S_2)$.  Moreover, for every $x'\in X'$ such that
  \[
    \codim\bigl(\overline{\{x'\}},X'\bigr)\leq1,
  \]
  the point $x=g(x')$ satisfies
  \[
    \codim\bigl(\overline{\{x\}},X\bigr)
    =
    \codim\bigl(\overline{\{x'\}},X'\bigr).
  \]
  \item Under the hypotheses of {\rm(ii)}, let $\sh F$ be a
  coherent $\sh O_X$-module with support $X$, and suppose that
  $\sh H^0_{X/Z^{(2)}}(\sh F)$ is coherent.  The
  $\sh O_{X'}$-module $\sh F'$ defined by
  \[
    g_*(\sh F')=\sh H^0_{X/Z^{(2)}}(\sh F)
  \]
  is coherent and satisfies $(\mathrm S_2)$, and its support is
  a union of irreducible components of $X'$.
\end{enumerate}
\end{proposition}

\begin{proof}
Assertions {\rm(i)} and {\rm(ii)} have already been proved in
substance: {\rm(i)} follows from \sref{IV.5.10.12}, {\rm(i)}, and
\sref{IV.5.10.14}, while {\rm(ii)} is the special case of
\sref{IV.5.10.15}, {\rm(ii)}, obtained by taking
$\sh F=\sh O_X$.

We prove {\rm(iii)}.  Put $x=g(x')$.  Since $g$ is finite, so is
$X'_x\to X_x$, and therefore
\[
  \dim(\sh O_{X',x'})
  \leq \dim(X'_x)
  \leq \dim(X_x)
  =\dim(\sh O_{X,x})
\]
by \sref{IV.5.4.1}.  Suppose first that
$\dim(\sh O_{X',x'})\leq1$.  Then
$\dim(\sh O_{X,x})\leq1$ as well: otherwise $x\in Z^{(2)}$ and
\sref{IV.5.10.12}, {\rm(ii)}, applied to $\sh O_X$, would give
\[
  \operatorname{prof}(\sh O_{X',x'})\geq2,
\]
which is impossible.  Hence $x\notin Z^{(2)}$, and
\sref{IV.5.10.12}, {\rm(i)}, gives an isomorphism
\[
  \sh O_{X,x}\simeq\sh O_{X',x'}.
\]
In particular their dimensions are equal.  Since $X$ has no
associated prime cycle of codimension $1$, the equality
$\dim(\sh O_{X,x})=1$ implies
$x\notin\operatorname{Ass}(\sh O_X)$, so
\[
  \operatorname{prof}(\sh O_{X,x})
  =\operatorname{prof}(\sh O_{X',x'})=1.
\]
If instead $\dim(\sh O_{X',x'})\geq2$, then
$\dim(\sh O_{X,x})\geq2$, so $x\in Z^{(2)}$ and
\sref{IV.5.10.12}, {\rm(ii)}, gives
\[
  \operatorname{prof}(\sh O_{X',x'})\geq2.
\]
This proves all the assertions in {\rm(iii)}.

\oldpage[IV-2]{121}

For {\rm(iv)}, replace $\sh O_X$ by $\sh F$ in the preceding
argument.  Since $g$ is finite and $g_*(\sh F')$ is coherent,
$\sh F'$ is coherent.  The same argument proves that $\sh F'$
satisfies $(\mathrm S_2)$ and that, if
$\dim(\sh F'_{x'})\leq1$, then $\sh F_x$ and
$\sh F'_{x'}$ are di-isomorphic.  In particular,
\[
  \dim(\sh F'_{x'})=0
  \quad\Longrightarrow\quad
  \dim(\sh F_x)=0
  \quad\Longrightarrow\quad
  \dim(\sh O_{X,x})=0,
\]
because $\operatorname{Supp}(\sh F)=X$, and hence also
$\dim(\sh O_{X',x'})=0$.  Every irreducible component of
$\operatorname{Supp}(\sh F')$ is therefore an irreducible
component of $X'$.  Finally, $\operatorname{Supp}(\sh F')$ is
closed because $\sh F'$ is coherent.
\end{proof}

\begin{proposition}[5.10.17]
\label{IV.5.10.17}
Let $A$ be a Noetherian integral domain, and denote by $A^{(1)}$
the intersection of the local rings $A_{\mathfrak p}$, where
$\mathfrak p$ ranges over the prime ideals of $A$ of height $1$.
Suppose that $A^{(1)}$ is a finite $A$-algebra.  Then:
\begin{enumerate}[{\rm(i)}]
  \item The ring $A^{(1)}$ satisfies $(\mathrm S_2)$.
  \item The set $U$ of primes $\mathfrak p\in\Spec(A)$ for which
  the canonical homomorphism
  \[
    A_{\mathfrak p}\longrightarrow(A^{(1)})_{\mathfrak p}
  \]
  is bijective is precisely the set of primes for which
  $A_{\mathfrak p}$ satisfies $(\mathrm S_2)$.  The set $U$ is
  open in $X=\Spec(A)$ and
  \[
    \codim(X-U,X)\geq2.
  \]
  \item For every multiplicative subset $S$ of $A$,
  $(S^{-1}A)^{(1)}$ is a finite $S^{-1}A$-algebra.
  \item Let $B$ be a finite integral $A$-algebra that is an
  integral domain and contains $A$.  Then $B^{(1)}$ is a finite
  $B$-algebra.  Moreover, if $\mathfrak q$ is a prime ideal of
  $B$ of height $1$, then $\mathfrak q\cap A$ has height $1$ in
  $A$.
\end{enumerate}
\end{proposition}

\begin{proof}
Formula \sref{IV.5.9.3.1} shows that
\[
  X'=\Spec(A^{(1)})
\]
is the $Z^{(2)}$-closure of $X=\Spec(A)$.  Since $A$ has no
embedded associated prime ideals, assertions {\rm(i)} and
{\rm(ii)} are special cases of \sref{IV.5.10.16},
{\rm(i)--(iii)}.

For {\rm(iii)}, observe that
\[
  (S^{-1}A)^{(1)}=S^{-1}A^{(1)},
\]
which is a special case of \sref{IV.5.9.4}.  Indeed, $S^{-1}A$
is flat over $A$, its prime ideals are the ideals
$S^{-1}\mathfrak p$ with
$\mathfrak p\in\Spec(A)$ disjoint from $S$, and
\[
  \operatorname{ht}(S^{-1}\mathfrak p)
  =\operatorname{ht}(\mathfrak p).
\]
Since $A^{(1)}$ is finite over $A$, its localization
$S^{-1}A^{(1)}$ is finite over $S^{-1}A$.

To prove {\rm(iv)}, put $Y=\Spec(B)$ and let $f:Y\to X$ be the
structural morphism.  Since $f$ is finite, \sref{IV.5.4.1}
gives, for every $y\in Y$,
\[
  \dim(\sh O_{Y,y})\leq\dim(\sh O_{X,f(y)}).
\]
Consequently, if
\[
  T=f^{-1}\bigl(Z^{(2)}(X)\bigr),
\]
then $T\supset Z^{(2)}(Y)$.

We claim that
\[
  \sh G=
  \sh H^0_{X/Z^{(2)}(X)}\bigl(f_*(\sh O_Y)\bigr)
\]
is coherent.  The module $f_*(\sh O_Y)$ corresponds to $B$
viewed as an $A$-module.  Since $B$ is a finite integral
$A$-algebra and an integral domain, its field of fractions is
finite over the field of fractions of $A$; hence $B$ is
contained in a finite free $A$-module.  By
\sref{IV.5.9.2}, {\rm(i)}, $\sh G$ is therefore a
quasi-coherent $\sh O_X$-submodule of
\[
  \bigl(\sh H^0_{X/Z^{(2)}(X)}(\sh O_X)\bigr)^n
\]
for a suitable $n$.  The latter module is coherent by
hypothesis, and hence so is $\sh G$.

By the definition \sref{IV.5.9.1.2}, there is an isomorphism
\[
  \sh G\simeq
  f_*\bigl(\sh H^0_{Y/T}(\sh O_Y)\bigr).
\]
It follows, in particular, that
$\sh H^0_{Y/T}(\sh O_Y)$ is a coherent $\sh O_Y$-module.
Applying \sref{IV.5.10.10}, {\rm(ii)}, to $\sh O_Y$ and the
generic point of $Y$, we obtain
\[
  \codim(T,Y)\geq2,
\]
that is, $T\subset Z^{(2)}(Y)$.  Thus
$T=Z^{(2)}(Y)$.  Formula \sref{IV.5.9.3.1} now identifies the
corresponding algebra with $B^{(1)}$, proving that $B^{(1)}$ is
finite over $B$.  Finally, if $\mathfrak q$ has height $1$, then
$\mathfrak q\notin T$, so
$\mathfrak q\cap A\notin Z^{(2)}(X)$.  Incomparability for the
integral extension $A\subset B$ excludes
$\mathfrak q\cap A=(0)$; hence
\[
  \operatorname{ht}(\mathfrak q\cap A)=1.
\]
\end{proof}
% END NATIVE INPUT components/10c_section510_s2_closure_applications_through_51017
% BEGIN NATIVE INPUT components/11a_section511_coherence_criteria_through_corollary5112
\oldpage[IV-2]{122}

\subsection{Coherence criteria for the modules
\texorpdfstring{$\sh H^0_{X/Z}(\sh F)$}{H0(X/Z,F)}}
\label{subsection:IV.5.11}

\begin{proposition}[5.11.1]
\label{IV.5.11.1}
Let $X$ be a locally Noetherian prescheme, let $Z$ be a subset of
$X$ stable under specialization, and let $\sh F$ be a coherent
$\sh O_X$-module.  Denote by $(x_\alpha)$ the family of points of
\[
  \operatorname{Ass}(\sh F)\cap(X-Z),
\]
and, for every $\alpha$, let $Y_\alpha$ be the reduced closed
subprescheme of $X$ whose underlying space is
$\overline{\{x_\alpha\}}$, and put
$Z_\alpha=Z\cap\overline{\{x_\alpha\}}$.  The following two
conditions are equivalent:
\begin{enumerate}[{\rm a)}]
  \item $\sh H^0_{X/Z}(\sh F)$ is a coherent $\sh O_X$-module.
  \item For every $\alpha$,
  $\sh H^0_{Y_\alpha/Z_\alpha}(\sh O_{Y_\alpha})$ is a coherent
  $\sh O_{Y_\alpha}$-module.
\end{enumerate}
These conditions imply:
\begin{enumerate}
  \item[{\rm c)}] For every $\alpha$,
  $\codim(Z_\alpha,Y_\alpha)\geq2$.
\end{enumerate}
Moreover, conditions {\rm a)}, {\rm b)}, and {\rm c)} are
equivalent if, in addition, either of the following properties
holds:
\begin{enumerate}[{\rm(i)}]
  \item Every point of $X$ has an open neighbourhood isomorphic
  to a subscheme of a regular scheme; one also says in this case
  that $X$ is \emph{locally immersible in a regular scheme}.
  \item For every $\alpha$, $Y_\alpha$ is universally catenary
  \sref{IV.5.6.2}, and its normalization
  $\widetilde Y_\alpha$ \sref[II]{II.6.3.8} is finite over
  $Y_\alpha$.
\end{enumerate}
\end{proposition}

\noindent\emph{Proof.}\quad
All the properties under consideration are local on $X$.  We may
therefore restrict ourselves to the case in which
$X=\Spec(A)$ is affine, $A$ is Noetherian, and
$\sh F=\widetilde M$, where $M$ is a finite $A$-module.  For
every $\alpha$, let $h_\alpha:Y_\alpha\to X$ be the canonical
injection.  Then
\[
  (h_\alpha)_*(\sh O_{Y_\alpha})=\sh G_\alpha
\]
is the $\sh O_X$-module corresponding to the quotient
$A/\mathfrak j_{x_\alpha}$; by the definition of
$\operatorname{Ass}(\sh F)$, this quotient is isomorphic to an
$A$-submodule of $M$.  Since
$\sh H^0_{X/Z}(\sh G_\alpha)$ is a quasi-coherent
$\sh O_X$-submodule of $\sh H^0_{X/Z}(\sh F)$
\sref{IV.5.9.2}, condition {\rm a)} implies that it is coherent.
On the other hand, the definition \sref{IV.5.9.1.2} gives
\[
  \sh H^0_{X/Z}(\sh G_\alpha)
  \simeq
  (h_\alpha)_*
  \bigl(\sh H^0_{Y_\alpha/Z_\alpha}(\sh O_{Y_\alpha})\bigr).
\]
Thus {\rm a)} implies {\rm b)}.

To prove the converse, it is enough to find a finite filtration
$(\sh F_i)_{0\leq i\leq n}$ by coherent $\sh O_X$-modules, with
$\sh F_0=\sh F$ and $\sh F_n=0$, such that
\[
  \sh H^0_{X/Z}(\sh F_i/\sh F_{i+1})
\]
is coherent for every $i$.  Indeed, the assertion then follows by
descending induction from the exact sequences
\[
  0\longrightarrow\sh F_{i+1}\longrightarrow\sh F_i
  \longrightarrow\sh F_i/\sh F_{i+1},
\]
the left exactness of $\sh H^0_{X/Z}$ \sref{IV.5.9.2}, and
\sref[0I]{0.5.3.3} and \sref[I]{I.6.1.1}.  In view of
\sref{IV.3.2.8}, it is therefore enough to treat the case in which
$\sh F$ is irredundant, that is,
$\operatorname{Ass}(M)=\{\mathfrak p\}$ consists of one element.
We use the following lemma.

\begin{lemma}[5.11.1.1]
\label{IV.5.11.1.1}
Let $A$ be a Noetherian ring and let $M$ be a finite $A$-module
such that $\operatorname{Ass}(M)=\{\mathfrak p\}$.  There is a
finite filtration $(M_i)_{0\leq i\leq r}$ of $M$, with
$M_0=M$ and $M_r=0$, such that every quotient
$M_i/M_{i+1}$ is isomorphic to a submodule of $A/\mathfrak p$.
\end{lemma}

\begin{proof}
The canonical homomorphism
$M\to M_{\mathfrak p}=N$ is injective
(Bourbaki, \emph{Alg\`ebre commutative}, Chapter~IV,
\S1, no.~2, Proposition~6).  Put
\[
  B=A_{\mathfrak p},
  \qquad
  \mathfrak m=\mathfrak pA_{\mathfrak p}.
\]
Then $\mathfrak m$ is the maximal ideal of $B$ and
$\operatorname{Ass}(N)=\{\mathfrak m\}$
(\emph{loc.\ cit.}, Proposition~5).

\oldpage[IV-2]{123}

Since $N$ is a finite
$B$-module, there is an integer $s$ such that
$\mathfrak m^sN=0$.  The filtration by the powers
$\mathfrak m^jN$ may be refined to a finite filtration
$(N_i)$ of $N$ whose successive quotients are isomorphic to
\[
  B/\mathfrak m=\operatorname{Frac}(A/\mathfrak p).
\]
Intersecting it with $M$ gives $M_i=M\cap N_i$.  Each quotient
$M_i/M_{i+1}$ is then isomorphic to a finite
$(A/\mathfrak p)$-submodule of
$\operatorname{Frac}(A/\mathfrak p)$, and every such module is
isomorphic to a submodule of $A/\mathfrak p$.
\end{proof}

\begin{proof}[Completion of the proof of {\hyperlink{ega4.c04.prop.5-11-1}{Proposition~5.11.1}}]
The lemma and the preceding filtration argument reduce the proof,
under condition {\rm b)}, to showing that
$\sh H^0_{X/Z}(\sh G)$ is coherent when
\[
  \sh G=(A/\mathfrak p)^{\sim}
\]
and $\mathfrak p$ is an associated prime of $M$.  If
$\mathfrak p=\mathfrak j_y$ with $y\in Z$, the support of
$\sh G$ is a closed subset contained in $Z$, because $Z$ is
stable under specialization.  Hence
$\sh H^0_{X/Z}(\sh G)=0$ by \sref{IV.5.9.1.2}.  Otherwise,
$\mathfrak p=\mathfrak j_{x_\alpha}$ for some $\alpha$, and
\[
  \sh G=(h_\alpha)_*(\sh O_{Y_\alpha}).
\]
The displayed isomorphism above and condition {\rm b)} then show
that $\sh H^0_{X/Z}(\sh G)$ is coherent.  This proves the
equivalence of {\rm a)} and {\rm b)}.

Condition {\rm a)} implies {\rm c)} by
\sref{IV.5.10.10}[ii].  It remains to prove, under either
hypothesis {\rm(i)} or {\rm(ii)}, that {\rm c)} implies {\rm b)}.
If $X$ satisfies {\rm(i)}, each $Y_\alpha$ does as well.
Consequently, in view of \sref{IV.5.9.3.1}, the required
assertion follows from {\hyperlink{ega4.c04.cor.5-11-2}{Corollary~5.11.2}} below.
\end{proof}

\begin{corollary}[5.11.2]
\label{IV.5.11.2}
Let $A$ be a Noetherian domain satisfying either of the following
hypotheses:
\begin{enumerate}[{\rm(i)}]
  \item $A$ is a quotient of a regular Noetherian ring.
  \item $A$ is universally catenary, and its integral closure
  $A'$ is a finite $A$-algebra.
\end{enumerate}
Then
\[
  A^{(1)}
  =\bigcap_{\substack{\mathfrak p\in\Spec(A)\\
                      \operatorname{ht}(\mathfrak p)=1}}
    A_{\mathfrak p}
\]
is a finite $A$-algebra.
\end{corollary}

\begin{proof}
Suppose first that {\rm(i)} holds, and put $X=\Spec(A)$.  The set
$U$ of points at which $X$ satisfies $(\mathrm S_2)$ is open by
\sref{IV.6.11.2}.\footnote{The reader may verify that
{\hyperlink{ega4.c04.cor.5-11-2}{Corollary~5.11.2}} is not used in the proof of
\sref{IV.6.11.2}.}
Put $Z=X-U$.  Then $\codim(Z,X)\geq2$.  Indeed, if
$x\in X$ and $\dim(\sh O_x)\leq1$, then
$\dim(\sh O_{x'})\leq1$ for every generization $x'$ of $x$.
Since $\sh O_{x'}$ is a domain, it has positive depth whenever
its dimension is positive, and therefore $X$ satisfies
$(\mathrm S_2)$ at $x$.

Thus $Z\subset Z^{(2)}(X)$.  If $j:U\to X$ is the canonical
injection, \sref{IV.5.9.1.2} gives
\[
  \sh H^0_{X/Z^{(2)}(X)}(\sh O_X)
  \simeq
  j_*\bigl(
    \sh H^0_{U/Z^{(2)}(U)}(\sh O_U)
  \bigr).
\]
The prescheme $U$ satisfies $(\mathrm S_2)$, so
\sref{IV.5.10.14} identifies the module inside the parentheses
with $\sh O_U$.  Since $\codim(Z,X)\geq2$, Chapter~III,
\S9 shows that $j_*(\sh O_U)$ is coherent.  Finally,
\sref{IV.5.9.3.1} identifies the corresponding $A$-module with
$A^{(1)}$, proving the assertion under {\rm(i)}.

Now suppose that {\rm(ii)} holds.  The ring $A'$ is a Noetherian
normal domain, hence it is the intersection of its local rings
$A'_{\mathfrak p'}$ as $\mathfrak p'$ ranges over the
height-one prime ideals of $A'$
(Bourbaki, \emph{Alg\`ebre commutative}, Chapter~VII,
\S1, no.~6, Theorem~4).  For such a prime $\mathfrak p'$, put
\[
  \mathfrak p=\mathfrak p'\cap A,
  \qquad
  S=A-\mathfrak p.
\]

\oldpage[IV-2]{124}

The ring $A'_{\mathfrak p'}$ is the localization of the finite
$A_{\mathfrak p}$-algebra $S^{-1}A'$ at
$S^{-1}\mathfrak p'$.  The latter prime lies above the maximal
ideal $\mathfrak pA_{\mathfrak p}$ and is therefore maximal.
Moreover, $A_{\mathfrak p}$ is universally catenary by
\sref{IV.5.6.3}[i].  Formula \sref{IV.5.6.10} consequently gives
\[
  \dim(A_{\mathfrak p})
  =\dim(A'_{\mathfrak p'})
  =1.
\]
It follows that $A^{(1)}\subset A'$.  Since $A'$ is a finite
$A$-module and $A$ is Noetherian, its submodule $A^{(1)}$ is
finite as well.
\end{proof}
% END NATIVE INPUT components/11a_section511_coherence_criteria_through_corollary5112
% BEGIN NATIVE INPUT components/11b_section511_corollaries_and_remarks
\begin{remark}[5.11.3]
\label{IV.5.11.3}
The conclusion that $A^{(1)}$ is a finite $A$-algebra need no
longer hold if, in hypothesis {\rm(ii)} of
\sref{IV.5.11.2}, one assumes only that $A$ is catenary.

Indeed, consider the catenary local ring $A$ constructed in
\sref{IV.5.6.11}; its integral closure $A'$ (denoted there by
$B$) is a finite $A$-algebra.  If $A^{(1)}$ were also finite over
$A$, then, being contained in the fraction field of $A$, it would
be contained in $A'$.  On the other hand, in the notation of
\sref{IV.5.6.11}, every height-one prime ideal of $A$ is of the
form $\mathfrak pA$, where $\mathfrak p$ is a height-one prime of
$C$, and
\[
  A_{\mathfrak pA}=C_{\mathfrak p}.
\]
Moreover, $C_{\mathfrak p}=E_{\mathfrak p'}$, where
$\mathfrak p'$ is the unique prime of $E$ above $\mathfrak p$.
Thus $A_{\mathfrak pA}$ is integrally closed and consequently
contains $A'$.  By definition this gives
$A'\subset A^{(1)}$, so finiteness of $A^{(1)}$ would force
$A^{(1)}=A'$.  This is impossible: one of the two prime ideals of
$A'$ above the maximal ideal $\mathfrak nA$ of $A$ has height
$1$, whereas $\mathfrak nA$ has height $2$, contradicting
\sref{IV.5.10.17}[iv].
\end{remark}

\begin{corollary}[5.11.4]
\label{IV.5.11.4}
Let $X$ be a locally Noetherian prescheme, let $Z$ be a closed
subset of $X$, put $U=X-Z$, and let $i:U\to X$ be the canonical
injection.  Let $\sh F$ be a coherent $\sh O_U$-module.  A
necessary condition for $i_*(\sh F)$ to be a coherent
$\sh O_X$-module is that
\[
  \codim\bigl(\overline{\{x\}}\cap Z,\overline{\{x\}}\bigr)
  \geq2
  \qquad
  \text{for every }x\in\operatorname{Ass}(\sh F).
\]
This condition is sufficient in either of the following cases:
\begin{enumerate}[{\rm(i)}]
  \item $X$ is locally immersible in a regular scheme.
  \item $X$ is universally catenary and universally Japanese;
  by definition, the latter means that every point $x\in X$ has
  an affine open neighbourhood whose ring is universally
  Japanese \sref[0]{0.23.1.1}.
\end{enumerate}
\end{corollary}

\begin{proof}
By \sref[I]{I.9.4.7}, there is a coherent
$\sh O_X$-submodule $\sh G$ of $i_*(\sh F)$ such that
$\sh G|U=\sh F$.  Plainly,
\[
  \operatorname{Ass}(\sh G)
  \subset\operatorname{Ass}(\sh F)
  \subset\operatorname{Ass}\bigl(i_*(\sh F)\bigr),
\]
and
\[
  \operatorname{Ass}\bigl(i_*(\sh F)\bigr)
  =\operatorname{Ass}(\sh F)
\]
by \sref{IV.3.1.13}.  Hence
$\operatorname{Ass}(\sh G)=\operatorname{Ass}(\sh F)$.
It remains to apply \sref{IV.5.11.1} to $\sh G$, observing that
\[
  i_*(\sh F)=\sh H^0_{X/Z}(\sh G),
\]
and that, when $X$ is universally catenary and universally
Japanese, hypothesis {\rm(ii)} of \sref{IV.5.11.1} holds by
definition.
\end{proof}

\begin{corollary}[5.11.5]
\label{IV.5.11.5}
Let $X$ be a locally Noetherian prescheme, let $Z$ be a subset
of $X$ stable under specialization, and let $\sh F$ be a coherent
$\sh O_X$-module such that
$\operatorname{Ass}(\sh F)\subset X-Z$.  Consider the conditions:
\begin{enumerate}[{\rm a)}]
  \item $\sh H^0_{X/Z}(\sh F)$ is a coherent
  $\sh O_X$-module.
  \item[{\rm d)}] For every subset $T$ of $Z$ closed in $X$,
  $i_*(\sh F|X-T)$ is coherent, where
  $i:X-T\to X$ is the canonical injection.  It is enough here

\oldpage[IV-2]{125}

  to consider an increasing filtered family of such closed
  subsets whose union is $Z$.
\end{enumerate}
Then {\rm a)} implies {\rm d)}.  If $X$ satisfies either
hypothesis {\rm(i)} or {\rm(ii)} of \sref{IV.5.11.1}, the two
conditions are equivalent.
\end{corollary}

\begin{proof}
The hypothesis and \sref{IV.3.1.13} give
\[
  \operatorname{Ass}\bigl(i_*(\sh F|X-T)\bigr)
  =\operatorname{Ass}(\sh F).
\]
It follows from \sref{IV.5.10.2} that the canonical maps
\[
  i_*(\sh F|X-T)\longrightarrow\sh H^0_{X/Z}(\sh F)
\]
are injective.  Condition {\rm a)} therefore implies {\rm d)}.

Conversely, condition {\rm d)} and \sref{IV.5.11.1} imply, in
the notation of that proposition, that
\[
  \codim(T\cap Y_\alpha,Y_\alpha)\geq2.
\]
Since $Z$ is the union of its subsets closed in $X$, it follows
that
\[
  \codim(Z\cap Y_\alpha,Y_\alpha)\geq2.
\]
The last assertion now follows from \sref{IV.5.11.1}.
\end{proof}

\begin{corollary}[5.11.6]
\label{IV.5.11.6}
Let $A$ be a Noetherian ring and put $X=\Spec(A)$.  Consider the
following properties:
\begin{enumerate}[{\rm a)}]
  \item For every quotient domain $B$ of $A$, the ring $B^{(1)}$
  \sref{IV.5.11.2} is a finite $B$-algebra.
  \item For every coherent $\sh O_X$-module $\sh F$ and every
  subset $Z$ of $X$ stable under specialization such that
  \[
    \codim\bigl(\overline{\{x\}}\cap Z,\overline{\{x\}}\bigr)
    \geq2
  \]
  for every
  $x\in\operatorname{Ass}(\sh F)\cap(X-Z)$, the module
  $\sh H^0_{X/Z}(\sh F)$ is coherent.
  \item For every closed subset $T$ of $X$ and every coherent
  $\sh O_U$-module $\sh G$, where $U=X-T$, such that
  \[
    \codim\bigl(\overline{\{x\}}\cap T,\overline{\{x\}}\bigr)
    \geq2
    \qquad
    (x\in\operatorname{Ass}(\sh G)),
  \]
  the module $i_*(\sh G)$ is coherent, where $i:U\to X$ is the
  canonical injection.
  \item For every quotient domain $B$ of $A$ and every ideal
  $\mathfrak J$ of $B$ of height at least $2$, the ring
  \[
    \bigcap_{\mathfrak p\not\supseteq\mathfrak J}
      B_{\mathfrak p}
  \]
  is a finite $B$-algebra.
\end{enumerate}
Then
\[
  {\rm a)}\Longleftrightarrow{\rm b)}
  \Longrightarrow{\rm c)}
  \Longleftrightarrow{\rm d)}.
\]
Moreover, all four properties hold in either of the following
cases:
\begin{enumerate}[{\rm(i)}]
  \item $A$ is a quotient of a regular ring.
  \item $A$ is universally catenary and universally Japanese.
\end{enumerate}
\end{corollary}

\begin{proof}
Let $B=A/\mathfrak q$, where $\mathfrak q$ is a prime ideal of
$A$, so that $Y=\Spec(B)=V(\mathfrak q)$ is a closed subset of
$X$.  Put $Z=Z^{(2)}(Y)$, viewed as a subset of $X$ stable under
specialization.  Since $B$ is a domain,
$\operatorname{Ass}(A/\mathfrak q)$ consists only of the generic
point $\mathfrak q$ of $Y$.  If {\rm b)} holds, apply it to
\[
  \sh F=(A/\mathfrak q)^{\sim}
\]
and to $Z$.  Formula \sref{IV.5.9.3.1} then shows that $B^{(1)}$
is a finite $A$-module, and hence a finite $B$-module.  Thus
{\rm b)} implies {\rm a)}.

Conversely, suppose {\rm a)} holds.  Let $\sh F$ and $Z$ satisfy
the assumptions in {\rm b)}, and use the notation of
\sref{IV.5.11.1}.  Each
$Y_\alpha=\Spec(B_\alpha)$, where $B_\alpha$ is a quotient
domain of $A$.  Since
$Z_\alpha\subset Z^{(2)}(Y_\alpha)$, condition {\rm a)},
\sref{IV.5.10.2}, and the fact that
$\operatorname{Ass}(\sh O_{Y_\alpha})$ consists of the generic
point show that
\[
  \sh H^0_{Y_\alpha/Z_\alpha}(\sh O_{Y_\alpha})
\]
is coherent.  Property {\rm b)} follows from
\sref{IV.5.11.1}.

\oldpage[IV-2]{126}

To prove that {\rm c)} implies {\rm d)}, argue as above with
$\sh G=(A/\mathfrak q)^{\sim}|U$ and $T=V(\mathfrak J)$.
Conversely, {\rm d)} implies {\rm c)} by applying the same
componentwise argument and \sref{IV.5.11.1}.  Property {\rm c)}
is plainly a special case of {\rm b)}.  Finally,
\sref{IV.5.11.2}, together with the definitions of universally
catenary and universally Japanese rings, proves all four
properties under either hypothesis {\rm(i)} or {\rm(ii)}.
\end{proof}

\begin{remarks}[5.11.7]
\label{IV.5.11.7}
\begin{enumerate}[{\rm(i)}]
  \item It is not known whether condition {\rm d)} in
  \sref{IV.5.11.5} implies condition {\rm a)} without an
  additional hypothesis on $X$.  We shall see in
  \sref{IV.7.2.4} that it does when $X$ is a local scheme.
  Likewise, when $A$ is a Noetherian local ring, the four
  properties {\rm a)}, {\rm b)}, {\rm c)}, and {\rm d)} of
  \sref{IV.5.11.6} are equivalent \sref{IV.7.2.4}.  It is not
  known whether this result extends to every Noetherian ring.

  \item If $A$ satisfies property {\rm a)} of
  \sref{IV.5.11.6}, then so does every localization $S^{-1}A$
  and every finite $A$-algebra $C$.  Indeed, every quotient
  domain of $S^{-1}A$ has the form
  \[
    S^{-1}(A/\mathfrak q),
  \]
  where $\mathfrak q$ is a prime ideal of $A$ disjoint from
  $S$.  On the other hand, if $\mathfrak r$ is a prime ideal of
  $C$ and $\mathfrak p$ is its inverse image in $A$, then
  $C/\mathfrak r$ is a finite integral
  $(A/\mathfrak p)$-algebra containing $A/\mathfrak p$.  The
  assertions therefore follow from \sref{IV.5.10.17}[iii] and
  \sref{IV.5.10.17}[iv].
\end{enumerate}
\end{remarks}
% END NATIVE INPUT components/11b_section511_corollaries_and_remarks
% BEGIN NATIVE INPUT components/12a_section512_quotient_by_regular_element_through_5127
\subsection{Relations between the properties of a Noetherian local
ring \texorpdfstring{$A$}{A} and of a quotient ring
\texorpdfstring{$A/tA$}{A/tA}}
\label{subsection:IV.5.12}
\label{IV.5.12}

We have already seen in \sref{IV.3.4} relations between the
properties of $A$ and those of $A/tA$ concerning associated prime
ideals, as well as the properties of being a domain or being
reduced.  In this subsection we give further relations between the
properties of these rings, involving the notions of dimension and
depth.

\begin{proposition}[5.12.1]
\label{IV.5.12.1}
Let $A$ be a Noetherian local ring, put $X=\Spec(A)$, and let $t$
be an element of $A$ belonging to a system of parameters
\sref[0]{0.16.3.6}.  Let $X_0$ be the closed subspace $V(t)$ of
$X$, and let $X_1=X-X_0$.  Let $Z$ be a subset of $X$ stable
under specialization.  Suppose that $X$ is biequidimensional
(as is the case when $A$ is equidimensional and is a quotient of
a regular local ring \sref[0]{0.16.5.12}).  If
$Z_0=Z\cap X_0$ and $Z_1=Z\cap X_1$, then
\begin{equation}
\tag{5.12.1.1}
\label{IV.5.12.1.1}
  \codim(Z_0,X_0)\leq\codim(Z,X)\leq\codim(Z_1,X_1).
\end{equation}
\end{proposition}

\begin{proof}
The second inequality follows from the definition
\sref{IV.5.1.3}.  For the first, it is enough to treat the case
in which $Z$ is closed.  Indeed, by hypothesis $Z$ is a union of
closed subsets $Z_\alpha$, and inequalities
\[
  \codim(X_0\cap Z_\alpha,X_0)\leq\codim(Z_\alpha,X)
\]
for all $\alpha$ give
\[
\begin{split}
  \codim(Z_0,X_0)
  &=\inf_\alpha\codim(X_0\cap Z_\alpha,X_0)\\
  &\leq\inf_\alpha\codim(Z_\alpha,X)
   =\codim(Z,X).
\end{split}
\]
Suppose therefore that $Z$ is closed.  Then
\[
  \codim(Z,X)=\dim(X)-\dim(Z)
\]
by \sref[0]{0.14.3.5.1}.

\oldpage[IV-2]{127}

Moreover, $X_0$ is catenary and equidimensional, and
$\dim(X_0)=\dim(X)-1$ by \sref[0]{0.16.3.4}; hence
\[
  \codim(Z_0,X_0)
  =\dim(X)-1-\dim(Z_0).
\]
Since $\dim(Z_0)\geq\dim(Z)-1$ \sref[0]{0.16.3.4}, the first
inequality in \eqref{IV.5.12.1.1} follows.
\end{proof}

\begin{proposition}[5.12.2]
\label{IV.5.12.2}
Let $A$ be a catenary Noetherian local ring, let $M$ be a finite
$A$-module, let $t$ be an $M$-regular element belonging to the
maximal ideal $\mathfrak m$ of $A$, and let $k\geq1$.  If $M/tM$
is equidimensional and satisfies $(\mathrm S_k)$, then the same is
true of $M$.
\end{proposition}

\begin{proof}
In view of $(\mathbf{Err}_{\mathrm{III}},30)$ and
\sref{IV.5.7.3}[vi], we may reduce to the case
\[
  \operatorname{Supp}(M)=\Spec(A)=X.
\]
Put $X_0=V(t)=\Spec(A/tA)$.  Since $M/tM$ satisfies
$(\mathrm S_1)$, it has no embedded associated prime ideals
\sref{IV.5.7.5}.  Applying \sref{IV.3.4.4} at the closed point of
$X$ (and of every closed subprescheme of $X$), we see that the
irreducible components of $\operatorname{Supp}(M/tM)$ are exactly
the irreducible components of $Y_i\cap X_0$, where
$Y_i$ $(1\leq i\leq r)$ are the irreducible components of $X$.

Because $t$ is $M$-regular, $V(t)$ contains no maximal point of
$X$ \sref{IV.3.1.8}.  Every irreducible component of
$Y_i\cap X_0$ therefore has codimension $1$ in $Y_i$
\sref{IV.5.1.8}.  Since $X$ is catenary and all irreducible
components of $\operatorname{Supp}(M/tM)$ have the same
dimension, the $Y_i$ all have the same dimension.  Thus $X$ is
equidimensional, and hence biequidimensional because $A$ is local
and catenary.

It remains to prove $(\mathrm S_k)$.  By the criterion
\sref{IV.5.7.4}, with its notation, we must show that
\[
  \codim(Z_n,X)>n+k
\]
for every integer $n\geq0$.  The hypothesis on $M/tM$ and the same
criterion give
\[
  \codim(Z_n\cap X_0,X_0)>n+k.
\]
Every specialization of a point of $Z_n$ again belongs to $Z_n$,
as will be proved in \sref{IV.6.11.5}.\footnote{The reader may
verify that {\hyperlink{ega4.c04.prop.5-12-2}{Proposition~5.12.2}} is not used in the proof of
\sref{IV.6.11.5}.}  Finally, $t$ belongs to a system of
parameters for $M$ \sref[0]{0.16.4.1}, and hence also to one for
$A$, because the annihilator of $M$ is nilpotent under the
hypothesis $\operatorname{Supp}(M)=\Spec(A)$.  The required
inequality now follows from \sref{IV.5.12.1}.
\end{proof}

\begin{remark}[5.12.3]
\label{IV.5.12.3}
If in {\hyperlink{ega4.c04.prop.5-12-2}{Proposition~5.12.2}} one assumes only that $M/tM$ is
equidimensional, it need not follow that $M$ is
equidimensional.  For example, let $k$ be a field, let
$B=k[T,U,V]$, let $C$ be the local ring of $B$ at the maximal
ideal $BT+BU+BV$, and put
\[
  \mathfrak p=CU,\qquad
  \mathfrak q=CV+C(T+U),\qquad
  A=C/\mathfrak p\mathfrak q.
\]
Geometrically, $A$ is the local ring, at their intersection
point $x$, of the closed subscheme of affine $3$-space formed by
a plane and a line meeting that plane at $x$.  Let $t,u,v$ be the
images of $T,U,V$ in $A$.  The element $t$ is not a zero divisor
in $A$, and
\[
  A/tA\simeq C_0/\mathfrak p_0\mathfrak q_0,
\]
where $C_0$ is the local ring of $B_0=k[U,V]$ at the maximal
ideal $B_0U+B_0V$, while
\[
  \mathfrak p_0=C_0U,\qquad
  \mathfrak q_0=C_0U+C_0V.
\]
  Thus $\mathfrak q_0$ is the maximal ideal of $C_0$.

\oldpage[IV-2]{128}

  It is
immediate that $\Spec(A/tA)$ is irreducible of dimension $1$,
that $A/tA$ does not satisfy $(\mathrm S_1)$, and that
$\Spec(A)$ is not equidimensional.
\end{remark}

\begin{corollary}[5.12.4]
\label{IV.5.12.4}
Let $A$ be a catenary Noetherian local ring, let $t$ be a regular
element belonging to the maximal ideal of $A$, and let $k\geq1$.
If $A/tA$ satisfies $(\mathrm S_k)$, then so does $A$.
\end{corollary}

\begin{proof}
For $k=1$ this follows from \sref{IV.3.4.4}, in view of the
interpretation of $(\mathrm S_1)$ in \sref{IV.5.7.5}.  For
$k\geq2$, Hartshorne's criterion \sref{IV.5.10.9} shows that
$A/tA$ is equidimensional, and {\hyperlink{ega4.c04.prop.5-12-2}{Proposition~5.12.2}} applies.
\end{proof}

\begin{proposition}[5.12.5]
\label{IV.5.12.5}
Let $A$ be a catenary Noetherian local ring, let $t$ be a regular
element of $A$ belonging to its maximal ideal, and let $k\geq0$.
If $A/tA$ is reduced and equidimensional and satisfies
$(\mathrm R_k)$, then $A$ has these same three properties.
\end{proposition}

\begin{proof}
We already know from \sref{IV.3.4.6} that $A$ is reduced, and
{\hyperlink{ega4.c04.prop.5-12-2}{Proposition~5.12.2}}, applied with $k=1$, shows that $A$ is
equidimensional.  Put
\[
  X=\Spec(A),\qquad X_0=V(t)=\Spec(A/tA),
\]
and let $Z$ be the set of points at which $X$ is not regular.
Every specialization of a point of $Z$ belongs to $Z$
\sref[0]{0.17.3.2}.  On the other hand, if $x\in X_0$ and $X_0$
is regular at $x$, then $X$ is regular at $x$
\sref[0]{0.17.1.8}.  Consequently the nonregular locus $Z'_0$
of $X_0$ contains $Z_0=Z\cap X_0$.  The hypothesis gives
\[
  k<\codim(Z'_0,X_0)\leq\codim(Z_0,X_0).
\]
Since $A$ is equidimensional and catenary, {\hyperlink{ega4.c04.prop.5-12-1}{Proposition~5.12.1}}
gives
\[
  \codim(Z_0,X_0)\leq\codim(Z,X).
\]
Thus $X$ satisfies $(\mathrm R_k)$.
\end{proof}

\begin{remark}[5.12.6]
\label{IV.5.12.6}
If in {\hyperlink{ega4.c04.prop.5-12-5}{Proposition~5.12.5}} one assumes only that $A/tA$ is reduced
and satisfies $(\mathrm R_k)$, it need not follow that $A$
satisfies $(\mathrm R_k)$.  For example, let $k$ be a field and
let $P(U,V)\in k[U,V]$ be irreducible, with the curve $\Gamma$
defined by $(P)$ in $\Spec(k[U,V])$ singular at the maximal ideal
$(U)+(V)$.  One may take, for example,
\[
  P(U,V)=U(U^2+V^2)+(U^2-V^2),
\]
which defines a cubic with a double point.  Let
$B=k[T,U,V,W]$, let $C$ be the local ring of $B$ at the maximal
ideal $BT+BU+BV+BW$, and put
\[
  \mathfrak p=CW+CP(U-T,V),\qquad \mathfrak q=CU.
\]
Consider the local ring $A=C/\mathfrak p\mathfrak q$.
Geometrically, the corresponding closed subscheme $X$ of affine
$4$-space is the union of a hyperplane $H$ and a
$2$-dimensional ``cylinder'' $L$ with base $\Gamma$; the
``singular line'' $Y$ of $L$ is not contained in $H$, and $A$ is
the local ring of $X$ at the point $x=Y\cap H$.

Let $t$ be the image of $T$ in $A$.  The scheme
$\Spec(A/tA)$ has $x$ as its only singular point; its local ring
there is $A/tA$, which is reduced and has dimension $2$.  By
contrast, the generic point $y$ of $Y$, defined by the image in
$A$ of
\[
  C(U-T)+CV+CW,
\]
is a singular point of $X$, and $\sh O_{X,y}$ has dimension $1$.
Thus $A/tA$ is reduced and satisfies $(\mathrm R_1)$, whereas
$A$ does not satisfy $(\mathrm R_1)$.
\end{remark}

\oldpage[IV-2]{129}

\begin{corollary}[5.12.7]
\label{IV.5.12.7}
Let $A$ be a catenary Noetherian local ring and let $t$ be a
regular element of $A$ belonging to its maximal ideal.  If
$A/tA$ is a domain and is integrally closed, then so is $A$.
\end{corollary}

\begin{proof}
By Serre's criterion \sref{IV.5.8.6} and the fact that $A$ is
local, it is enough to prove that $\Spec(A)$ satisfies
$(\mathrm S_2)$ and $(\mathrm R_1)$.  The hypothesis on $A/tA$
implies $(\mathrm R_1)$ for $A$ by {\hyperlink{ega4.c04.prop.5-12-5}{Proposition~5.12.5}} and
$(\mathrm S_2)$ for $A$ by {\hyperlink{ega4.c04.cor.5-12-4}{Corollary~5.12.4}}.
\end{proof}
% END NATIVE INPUT components/12a_section512_quotient_by_regular_element_through_5127
% BEGIN NATIVE INPUT components/12b_section512_hironaka_lemma_and_regular_sequences
\begin{proposition}[5.12.8]
\label{IV.5.12.8}
{\normalfont(Hironaka's lemma).}
Let $A$ be a reduced, equidimensional, catenary Noetherian local
ring.  Suppose in addition that, for every minimal prime ideal
$\mathfrak q_i$ of $A$, the ring
$B_i=A/\mathfrak q_i$ has the property that $B_i^{(1)}$
(notation of \sref{IV.5.10.17}) is a finite $B_i$-algebra.  This
last condition holds, for example, when $A$ has either property
{\rm(i)} or property {\rm(ii)} of \sref{IV.5.11.6}.

Let $t\in A$ be a non-zero-divisor satisfying the following
conditions:
\begin{enumerate}[{\rm(i)}]
  \item The $A$-module $A/tA$ has a unique nonembedded associated
  prime ideal $\mathfrak p$.  (It necessarily has height $1$ by
  the Hauptidealsatz \sref[0]{0.16.3.2}, but it is not assumed to
  be the only associated prime of $A/tA$.)
  \item The ideal $\mathfrak pA_{\mathfrak p}$ is generated by
  $t/1$.
  \item The ring $A/\mathfrak p$ is integrally closed.
\end{enumerate}
Then $A$ is a domain and is integrally closed, and
$\mathfrak p=tA$.
\end{proposition}

\begin{proof}
Put $X=\Spec(A)$ and $Z=V(tA)$.  Condition {\rm(i)} says that
$Z$ is irreducible, with generic point corresponding to
$\mathfrak p$.  By {\rm(ii)}, the maximal ideal of the
Noetherian local ring $A_{\mathfrak p}$ is generated by one
element.  Thus $A_{\mathfrak p}$ is a discrete valuation ring
with uniformizer $t/1$ (Bourbaki, \emph{Alg\`ebre commutative},
chap.~VI, \S3, no.~6, prop.~9), and
$A_{\mathfrak p}/tA_{\mathfrak p}$, its residue field, has length
$1$ as an $A_{\mathfrak p}$-module.

It follows from \sref{IV.3.4.2} that $Z$ is contained in only one
irreducible component, say $X_i$, of $X$.  For every other
irreducible component $X_j$ one consequently has
\[
  \dim(X_j\cap Z)<\dim(Z)=\dim(X_i\cap Z).
\]
Since $t$ is not a zero divisor, it belongs to none of the
minimal prime ideals $\mathfrak q_j$, and hence
\[
  \codim(X_i\cap Z,X_i)
  =\codim(X_j\cap Z,X_j)=1
\]
for every $j\ne i$ \sref{IV.5.1.8}.  But $A$ is
biequidimensional; the displayed strict inequality would
therefore force $\dim(X_i)\ne\dim(X_j)$
\sref[0]{0.14.3.3.1}, a contradiction.  Hence $A$ has only one
minimal prime.  Since it is reduced, $A$ is a domain.

Let $A^{(1)}$ have the meaning of \sref{IV.5.10.17}.  It is a
finite $A$-algebra and therefore a Noetherian semilocal ring.
Every nonembedded associated prime $\mathfrak r$ of
$A^{(1)}/tA^{(1)}$ has height $1$ by the Hauptidealsatz.  Its
contraction $\mathfrak r\cap A$ also has height $1$
\sref{IV.5.10.17}[iv], and, since it contains $tA$, condition
{\rm(i)} forces
\[
  \mathfrak r\cap A=\mathfrak p.
\]

Let $A'$ be the integral closure of $A$ and put
$S=A-\mathfrak p$.  The integral closure of $A_{\mathfrak p}$ is
$S^{-1}A'$ (Bourbaki, \emph{Alg\`ebre commutative}, chap.~V,
\S1, no.~5, prop.~17); it equals $A_{\mathfrak p}$ because this
ring is a discrete valuation ring.  Consequently there is only
one prime $\mathfrak p'$ of $A'$ above $\mathfrak p$
(\emph{loc.\ cit.}, \S2, no.~1, lemma~1), and therefore only one
prime $\mathfrak p^{(1)}$ of $A^{(1)}$ above $\mathfrak p$.
Moreover,
\[
  A_{\mathfrak p}
  =A'_{\mathfrak p'}
  =A^{(1)}_{\mathfrak p^{(1)}}.
\]

\oldpage[IV-2]{130}

The ring $A^{(1)}$ satisfies $(\mathrm S_2)$
\sref{IV.5.10.17}[i].  Since $t$ is a non-zero-divisor on
$A^{(1)}$, the quotient
\[
  A^{(1)}/tA^{(1)}
\]
has no embedded associated primes \sref{IV.5.7.7}.  Thus
$\mathfrak p^{(1)}$ is its only associated prime, so
$tA^{(1)}$ is $\mathfrak p^{(1)}$-primary.  Equivalently,
$tA^{(1)}$ is the inverse image in $A^{(1)}$ of
\[
  tA^{(1)}_{\mathfrak p^{(1)}}=tA_{\mathfrak p},
\]
which is the maximal ideal of $A_{\mathfrak p}$ by {\rm(ii)}.
It follows that
\[
  tA^{(1)}=\mathfrak p^{(1)}
\]
is prime in $A^{(1)}$.

The domain $A^{(1)}/\mathfrak p^{(1)}$ is finite over
$A/\mathfrak p$ and has the same fraction field, namely the
residue field of
$A^{(1)}_{\mathfrak p^{(1)}}=A_{\mathfrak p}$.  Condition
{\rm(iii)} therefore gives
\[
  A^{(1)}/\mathfrak p^{(1)}=A/\mathfrak p.
\]
Hence
\[
  A^{(1)}=A+\mathfrak p^{(1)}=A+tA^{(1)}.
\]
Since $t$ lies in the Jacobson radical of the Noetherian
semilocal ring $A^{(1)}$, Nakayama's lemma yields
$A^{(1)}=A$.  Thus
\[
  \mathfrak p^{(1)}=\mathfrak p=tA.
\]
Finally, $A$ is catenary and
$A/tA=A/\mathfrak p$ is a domain and integrally closed by
{\rm(iii)}.  {\hyperlink{ega4.c04.cor.5-12-7}{Corollary~5.12.7}} shows that $A$ is integrally
closed.
\end{proof}

\begin{remark}[5.12.9]
\label{IV.5.12.9}
The conclusion of Hironaka's lemma fails without the hypothesis
that $A$ is equidimensional.  Let $k$ be a field, put
\[
  B=k[X,Y,Z],\qquad
  \mathfrak r_1=BZ,\qquad
  \mathfrak r_2=BX+BY,\qquad
  C=B/\mathfrak r_1\mathfrak r_2.
\]
Let $\mathfrak m=BX+BY+BZ$, let
$\mathfrak n=\mathfrak m/\mathfrak r_1\mathfrak r_2$ be its
image in $C$, and let $A=C_{\mathfrak n}$.  Finally, let $t_0$
be the image of $Y+Z$ in $C$, and let $t$ be its image in $A$.
Thus $\Spec(C)$ is the union of a plane $P$ and a line $D$ not
contained in $P$, while $\Spec(C/t_0C)$ is a line $D'$ contained
in $P$ and passing through $D\cap P$.

The ring $A$ is reduced and catenary, since it is a quotient of a
regular ring, and its minimal prime ideals $\mathfrak q_1$ and
$\mathfrak q_2$ are the images of $\mathfrak r_1$ and
$\mathfrak r_2$.  The $A$-module $A/tA$ has a unique
nonembedded associated prime $\mathfrak p$, the image of
$BY+BZ$; the ideal $\mathfrak pA_{\mathfrak p}$ is generated by
$t/1$, and $A/\mathfrak p\simeq k[X]$.  Nevertheless, $A$ is
not a domain.
\end{remark}

\begin{corollary}[5.12.10]
\label{IV.5.12.10}
Let $A$ be a Noetherian local domain satisfying one of the
following hypotheses:
\begin{enumerate}[{\rm a)}]
  \item $A$ is a quotient of a regular ring;
  \item $A$ is universally catenary and universally Japanese.
\end{enumerate}
Let $(x_i)_{1\leq i\leq n}$ be a sequence of $n$ elements of
$A$, and put
\[
  \mathfrak J=x_1A+\cdots+x_nA.
\]
Suppose that:
\begin{enumerate}[{\rm(i)}]
  \item $A/\mathfrak J$ has a unique nonembedded associated
  prime $\mathfrak p$, and $\mathfrak p$ has height $n$;
  \item the maximal ideal $\mathfrak pA_{\mathfrak p}$ is
  generated by the $x_i/1$ (so $A_{\mathfrak p}$ is a regular
  local ring of dimension $n$);
  \item $A/\mathfrak p$ is integrally closed.
\end{enumerate}
Then, for every $i$ with $0\leq i\leq n$, the quotient
\[
  A/(x_1A+\cdots+x_iA)
\]
is integrally closed and has dimension $\dim(A)-i$.  In
particular, $\mathfrak J$ is prime and equals $\mathfrak p$, the
ring $A$ is integrally closed, and $(x_i)_{1\leq i\leq n}$ is an
$A$-regular sequence.
\end{corollary}

\begin{proof}
We argue by induction on $n$.  The assertion is trivial for
$n=0$ and is Hironaka's lemma for $n=1$, so suppose $n\geq2$.
Put
\[
  \mathfrak J'=x_1A+\cdots+x_{n-1}A,
\]
and let $\mathfrak q$ be a prime minimal among those containing
$\mathfrak J'$.  Then
$\operatorname{ht}(\mathfrak q)\leq n-1$
\sref[0]{0.16.3.1}.  We first show that
$\mathfrak q\subset\mathfrak p$.  If $\mathfrak p'$ is minimal
among the primes containing $\mathfrak q+x_nA$, then
$\mathfrak p'/\mathfrak q$ has height at most $1$ by the
Hauptidealsatz \sref[0]{0.16.3.2}; because $A$ is catenary,
$\operatorname{ht}(\mathfrak p')\leq n$.  But $\mathfrak p'$
contains $\mathfrak J$, and the unique minimal prime of
$A/\mathfrak J$ is $\mathfrak p$.  Hence
$\mathfrak p'=\mathfrak p$, and $\mathfrak q\subset\mathfrak p$.

\oldpage[IV-2]{131}

The prime $\mathfrak q$ is therefore induced from a prime of
$A_{\mathfrak p}$ minimal over
$\mathfrak J'A_{\mathfrak p}$.  By {\rm(ii)} and
\sref[0]{0.17.1.7}, the ideal
$\mathfrak J'A_{\mathfrak p}$ is prime.  Consequently
\[
  \mathfrak q=A\cap\mathfrak J'A_{\mathfrak p}
\]
is the unique nonembedded associated prime of
$A/\mathfrak J'$, and has height $n-1$.  Moreover,
\[
  \mathfrak qA_{\mathfrak q}
  =\mathfrak J'A_{\mathfrak q},
\]
and, since
$A_{\mathfrak q}=(A_{\mathfrak p})_{\mathfrak qA_{\mathfrak p}}$,
the maximal ideal of $A_{\mathfrak q}$ is generated by the
images of $x_1,\ldots,x_{n-1}$.  Thus the shortened sequence
already satisfies conditions {\rm(i)} and {\rm(ii)}.

We verify condition {\rm(iii)} for it.  Put
$B=A/\mathfrak q$, let $\mathfrak n=\mathfrak p/\mathfrak q$,
and let $t$ be the image of $x_n$ in $B$.  Property {\rm a)}
(respectively {\rm b)}) passes from $A$ to $B$
\sref{IV.5.6.3}[iv], and $t\ne0$ because
$x_n\notin\mathfrak q$.  A prime of $B$ minimal over $t$ comes
from a prime of $A$ minimal over $\mathfrak q+x_nA$; the only
such prime is $\mathfrak p$.  Hence $\mathfrak n$ is the unique
nonembedded associated prime of $B/tB$.  Furthermore,
\[
  B_{\mathfrak n}
    =A_{\mathfrak p}/\mathfrak qA_{\mathfrak p},
  \qquad
  \mathfrak nB_{\mathfrak n}
    =\mathfrak pA_{\mathfrak p}/\mathfrak qA_{\mathfrak p}.
\]
Since
$\mathfrak qA_{\mathfrak p}=\mathfrak J'A_{\mathfrak p}$
and $\mathfrak pA_{\mathfrak p}$ is generated by
$x_1,\ldots,x_n$, the ideal
$\mathfrak nB_{\mathfrak n}$ is generated by $t$.  Finally,
$B/\mathfrak n=A/\mathfrak p$ is integrally closed.  Hironaka's
lemma therefore gives
\[
  B=A/\mathfrak q\ \text{integrally closed},
  \qquad
  \mathfrak p=\mathfrak q+x_nA.
\]

The induction hypothesis applied to
$x_1,\ldots,x_{n-1}$ now shows that
$A/(x_1A+\cdots+x_iA)$ is integrally closed and has dimension
$\dim(A)-i$ for $0\leq i\leq n-1$, and that
$\mathfrak q=\mathfrak J'$.  Hence
\[
  \mathfrak p=\mathfrak J'+x_nA=\mathfrak J,
\]
so $A/\mathfrak J=A/\mathfrak p$ is integrally closed.  Finally,
$\dim(A_{\mathfrak p})=n$, and biequidimensionality of $A$
gives
\[
  \dim(A/\mathfrak p)
  =\dim(A)-\dim(A_{\mathfrak p})
  =\dim(A)-n.
\]
This completes the induction.
\end{proof}
% END NATIVE INPUT components/12b_section512_hironaka_lemma_and_regular_sequences
% BEGIN NATIVE INPUT components/13a_section513_direct_limits_through_integral_closure
\subsection{Properties preserved under passage to a direct limit}
\label{subsection:IV.5.13}

\oldpage[IV-2]{131}

Throughout this subsection,
$(A_\alpha,\varphi_{\beta\alpha})$ denotes a direct system of
rings whose index set is filtered and preordered.  We put
\[
  A=\varinjlim_\alpha A_\alpha
\]
and denote by $\varphi_\alpha:A_\alpha\to A$ the canonical
homomorphism.

\begin{proposition}[5.13.1]
\label{IV.5.13.1}
\begin{enumerate}[{\rm(i)}]
  \item For every index $\alpha$, let $\mathfrak a_\alpha$ be an
  ideal of $A_\alpha$ such that
  \[
    \varphi_{\beta\alpha}(\mathfrak a_\alpha)
    \subset\mathfrak a_\beta
    \qquad(\alpha\leq\beta).
  \]
  Then the $\mathfrak a_\alpha$, with the restrictions of the
  $\varphi_{\beta\alpha}$, form a direct system;
  $\mathfrak a=\varinjlim\mathfrak a_\alpha$ is canonically
  identified with an ideal of $A$, and
  \[
    A/\mathfrak a
    \simeq
    \varinjlim_\alpha(A_\alpha/\mathfrak a_\alpha).
  \]
  If, in addition,
  \[
    \mathfrak a_\alpha
    =\varphi_{\beta\alpha}^{-1}(\mathfrak a_\beta)
    \qquad(\alpha\leq\beta),
  \]
  then
  \[
    \mathfrak a_\alpha=\varphi_\alpha^{-1}(\mathfrak a)
    \qquad\text{for every }\alpha.
  \]

  \item Conversely, for every ideal $\mathfrak a$ of $A$, if
  \[
    \mathfrak a_\alpha=\varphi_\alpha^{-1}(\mathfrak a)
    \qquad\text{for every }\alpha,
  \]
  then the $\mathfrak a_\alpha$ form a direct system and
  \[
    \mathfrak a=\varinjlim_\alpha\mathfrak a_\alpha.
  \]
\end{enumerate}
\end{proposition}

\begin{proof}
Assertion {\rm(ii)} is plainly a special case of {\rm(i)}.
In {\rm(i)}, the canonical identification of
$\mathfrak a$ with an ideal of $A$, and that of
$A/\mathfrak a$ with
$\varinjlim(A_\alpha/\mathfrak a_\alpha)$, follow from the
exactness of the direct-limit functor in the category of abelian
groups.

Moreover, $\mathfrak a$ is the union of the increasing family of
images $\varphi_\alpha(\mathfrak a_\alpha)$.  Thus, if
$x_\alpha\in A_\alpha$ and
$\varphi_\alpha(x_\alpha)\in\mathfrak a$, there are
$\beta\geq\alpha$ and $y_\beta\in\mathfrak a_\beta$ such that
\[
  \varphi_\alpha(x_\alpha)=\varphi_\beta(y_\beta).
\]
For some $\gamma\geq\beta$ one then has
\[
  \varphi_{\gamma\alpha}(x_\alpha)
  =\varphi_{\gamma\beta}(y_\beta)\in\mathfrak a_\gamma.
\]
Under the additional inverse-image hypothesis this implies
$x_\alpha\in\mathfrak a_\alpha$, and proves the last assertion of
{\rm(i)}.
\end{proof}

\begin{corollary}[5.13.2]
\label{IV.5.13.2}
Let $\mathfrak N_\alpha$ be the nilradical of $A_\alpha$.  The
nilradical of $A$ is canonically identified with
$\varinjlim\mathfrak N_\alpha$, and
\[
  A_{\mathrm{red}}
  \simeq
  \varinjlim_\alpha(A_\alpha)_{\mathrm{red}}.
\]
In particular, if every $A_\alpha$ is reduced, then $A$ is
reduced.
\end{corollary}

\oldpage[IV-2]{132}

\begin{proof}
It is clear that
\[
  \varphi_{\beta\alpha}(\mathfrak N_\alpha)
  \subset\mathfrak N_\beta
  \qquad(\alpha\leq\beta),
\]
and that $\varinjlim\mathfrak N_\alpha$ is a nil ideal of $A$,
hence is contained in the nilradical $\mathfrak N$ of $A$.
Conversely, if $x\in\mathfrak N$, then $x^h=0$ for some integer
$h$.  Write $x=\varphi_\alpha(x_\alpha)$ with
$x_\alpha\in A_\alpha$.  The relation
$\varphi_\alpha(x_\alpha^h)=0$ implies that for some
$\beta\geq\alpha$,
\[
  \varphi_{\beta\alpha}(x_\alpha^h)=0.
\]
Putting $x_\beta=\varphi_{\beta\alpha}(x_\alpha)$, we obtain
$x_\beta^h=0$, hence $x_\beta\in\mathfrak N_\beta$, while
$x=\varphi_\beta(x_\beta)$.  This proves the assertion, in view
of \sref{IV.5.13.1}.
\end{proof}

\begin{proposition}[5.13.3]
\label{IV.5.13.3}
\begin{enumerate}[{\rm(i)}]
  \item For every index $\alpha$, let $\mathfrak p_\alpha$ be a
  prime ideal of $A_\alpha$ such that
  \[
    \varphi_{\beta\alpha}(\mathfrak p_\alpha)
    \subset\mathfrak p_\beta
    \qquad(\alpha\leq\beta).
  \]
  Then $\mathfrak p=\varinjlim\mathfrak p_\alpha$ is a prime
  ideal of $A$.  In particular, if all the $A_\alpha$ are
  domains, then $A$ is a domain.

  \item If, in addition,
  \[
    \mathfrak p_\alpha
    =\varphi_{\beta\alpha}^{-1}(\mathfrak p_\beta)
    \qquad(\alpha\leq\beta),
  \]
  then there is a canonical identification
  \[
    A_{\mathfrak p}
    \simeq
    \varinjlim_\alpha(A_\alpha)_{\mathfrak p_\alpha}.
  \]

  \item Suppose that $A$ is local, with maximal ideal
  $\mathfrak p$, and that the homomorphisms
  $\varphi_\alpha:A_\alpha\to A$ are injective.  If
  $\mathfrak p_\alpha=\varphi_\alpha^{-1}(\mathfrak p)$, then
  \[
    A
    \simeq
    \varinjlim_\alpha(A_\alpha)_{\mathfrak p_\alpha},
  \]
  and every canonical homomorphism
  $(A_\alpha)_{\mathfrak p_\alpha}\to A$ is injective.
\end{enumerate}
\end{proposition}

\begin{proof}
The two assertions in {\rm(i)} are equivalent, since
\[
  A/\mathfrak p
  \simeq
  \varinjlim_\alpha(A_\alpha/\mathfrak p_\alpha).
\]
Suppose therefore that all the $A_\alpha$ are domains, and let
$x,y\in A$ satisfy $xy=0$.  For some $\alpha$ one may write
\[
  x=\varphi_\alpha(x_\alpha),
  \qquad
  y=\varphi_\alpha(y_\alpha),
\]
with $x_\alpha,y_\alpha\in A_\alpha$ and
$\varphi_\alpha(x_\alpha y_\alpha)=0$.  Hence there is
$\beta\geq\alpha$ such that, on putting
\[
  x_\beta=\varphi_{\beta\alpha}(x_\alpha),
  \qquad
  y_\beta=\varphi_{\beta\alpha}(y_\alpha),
\]
one has $x_\beta y_\beta=0$.  Since $A_\beta$ is a domain,
$x_\beta=0$ or $y_\beta=0$, and therefore $x=0$ or $y=0$.

For {\rm(ii)}, the hypothesis says that the multiplicative
subsets $A_\alpha-\mathfrak p_\alpha$ form a direct system whose
direct limit is $A-\mathfrak p$.  The local rings
$(A_\alpha)_{\mathfrak p_\alpha}$ therefore form a direct system
with local transition maps, and the canonical homomorphisms to
$A_{\mathfrak p}$ induce a surjection
\[
  \psi:
  \varinjlim_\alpha(A_\alpha)_{\mathfrak p_\alpha}
  \longrightarrow A_{\mathfrak p}.
\]
To prove injectivity, let $x_\alpha\in A_\alpha$ and
$s_\alpha\in A_\alpha-\mathfrak p_\alpha$, and suppose that
\[
  \frac{\varphi_\alpha(x_\alpha)}
       {\varphi_\alpha(s_\alpha)}
  =0
  \quad\text{in }A_{\mathfrak p}.
\]
There is $s'\in A-\mathfrak p$ such that
$s'\varphi_\alpha(x_\alpha)=0$.  After replacing $\alpha$ by a
larger index, we may write
$s'=\varphi_\alpha(s'_\alpha)$ with
$s'_\alpha\in A_\alpha-\mathfrak p_\alpha$.  Then
$\varphi_\alpha(s'_\alpha x_\alpha)=0$, so for some
$\beta\geq\alpha$,
\[
  \varphi_{\beta\alpha}(s'_\alpha x_\alpha)=0
  \quad\text{in }A_\beta.
\]
It follows that the image of $x_\alpha/s_\alpha$ in the direct
limit is zero.

Finally, the first assertion of {\rm(iii)} follows from {\rm(ii)}
and $A_{\mathfrak p}=A$.  Every element of
$A_\alpha-\mathfrak p_\alpha$ maps to a unit of $A$, and the
injectivity of
$(A_\alpha)_{\mathfrak p_\alpha}\to A$ follows from that of
$\varphi_\alpha$.
\end{proof}

\begin{corollary}[5.13.4]
\label{IV.5.13.4}
Suppose that all the $A_\alpha$ are domains and that all the
transition homomorphisms $\varphi_{\beta\alpha}$ are injective.
If $A'_\alpha$ denotes the integral closure of $A_\alpha$, then
\[
  A'=\varinjlim_\alpha A'_\alpha
\]
is the integral closure of $A$.  In particular, if every
$A_\alpha$ is integrally closed, then $A$ is integrally closed.

If the $A_\alpha$ are unibranch local rings
(respectively, geometrically unibranch local rings)
\sref[0]{0.23.2.1}, and the $\varphi_{\beta\alpha}$ are injective
local homomorphisms, then $A$ is unibranch
(respectively, geometrically unibranch).
\end{corollary}

\begin{proof}
Let $K_\alpha$ be the fraction field of $A_\alpha$.  By
\sref{IV.5.13.3}[ii], applied to $\mathfrak p_\alpha=(0)$,
\[
  K=\varinjlim_\alpha K_\alpha
\]
is the fraction field of $A$.  Thus
$A'=\varinjlim A'_\alpha$ is identified with

\oldpage[IV-2]{133}

a subring of $K$ which is plainly integral over $A$.

Conversely, let $z\in K$ satisfy an equation of integral
dependence
\[
  z^n+\sum_{i=1}^{n}c_i z^{\,n-i}=0,
  \qquad c_i\in A.
\]
There are an index $\alpha$, an element $z_\alpha\in K_\alpha$,
and elements $c_{i\alpha}\in A_\alpha$ whose images are
$z$ and the $c_i$, respectively, and such that the same equation
holds in $K_\alpha$.  Hence $z_\alpha\in A'_\alpha$, so
$z\in A'$.  This also proves at once that a direct limit of the
integrally closed rings occurring here is integrally closed.

Suppose now that the $A_\alpha$ are local and that the transition
maps are local.  Then $A$ is local and, if $k_\alpha$ denotes the
residue field of $A_\alpha$,
\[
  k=\varinjlim_\alpha k_\alpha
\]
is the residue field of $A$ \sref[0]{0.10.3.1.3}.  If every
$A'_\alpha$ is local, then $A'$ is local, and hence $A$ is
unibranch.  If, moreover, the residue field $k'_\alpha$ of
$A'_\alpha$ is a radicial extension of $k_\alpha$ for every
$\alpha$, then
\[
  k'=\varinjlim_\alpha k'_\alpha,
\]
the residue field of $A'$, is a radicial extension of $k$.
Thus $A$ is geometrically unibranch.
\end{proof}
% END NATIVE INPUT components/13a_section513_direct_limits_through_integral_closure
% BEGIN NATIVE INPUT components/13b_section513_normality_regularity_and_Rk
\begin{env}[5.13.5]
\label{IV.5.13.5}
A ring $A$ is called \emph{normal} if $\Spec(A)$ is normal; in
other words, if for every prime ideal $\mathfrak p$ of $A$, the
ring $A_{\mathfrak p}$ is a domain and is integrally closed.

This implies that $A$ is reduced and that $\mathfrak p$ can
contain only one minimal prime ideal $\mathfrak q$ of $A$; in
other words, $\mathfrak p$ belongs to only one irreducible
component of $\Spec(A)$.  Moreover, $\mathfrak q$ is the inverse
image of $(0)$ under the canonical homomorphism
$A\to A_{\mathfrak p}$, and hence
\[
  A_{\mathfrak p}
  \simeq
  (A/\mathfrak q)_{\mathfrak p/\mathfrak q}.
\]
It follows that the domain $A/\mathfrak q$ is integrally closed,
being the intersection of integrally closed rings
\sref[I]{I.8.2.1.1}.

If $A$ is Noetherian and normal, it is the direct product of
finitely many integrally closed domains \sref[I]{I.5.1.4}.
\end{env}

\begin{corollary}[5.13.6]
\label{IV.5.13.6}
Suppose that all the $A_\alpha$ are normal and that, for
$\alpha\leq\beta$, every irreducible component of
$\Spec(A_\beta)$ dominates an irreducible component of
$\Spec(A_\alpha)$.  Then
\[
  A=\varinjlim_\alpha A_\alpha
\]
is normal.
\end{corollary}

\begin{proof}
Let $\mathfrak p$ be a prime ideal of $A$, and put
\[
  \mathfrak p_\alpha=\varphi_\alpha^{-1}(\mathfrak p)
\]
for every $\alpha$.  By hypothesis, $\mathfrak p_\alpha$
contains a unique minimal prime ideal $\mathfrak q_\alpha$ of
$A_\alpha$.  If $\alpha\leq\beta$, then
$\varphi_{\beta\alpha}^{-1}(\mathfrak q_\beta)$ is a minimal
prime ideal of $A_\alpha$ contained in $\mathfrak p_\alpha$,
hence equals $\mathfrak q_\alpha$.

Put
\[
  B_\alpha=A_\alpha/\mathfrak q_\alpha.
\]
The domains $B_\alpha$ form a direct system under the injective
homomorphisms
\[
  \varphi'_{\beta\alpha}:B_\alpha\longrightarrow B_\beta
\]
induced by the $\varphi_{\beta\alpha}$.  Each $B_\alpha$ is
integrally closed, so
\[
  B=\varinjlim_\alpha B_\alpha
\]
is integrally closed by \sref{IV.5.13.4}.  If
\[
  \mathfrak p'_\alpha
  =\mathfrak p_\alpha/\mathfrak q_\alpha,
\]
then
\[
  (B_\alpha)_{\mathfrak p'_\alpha}
  =(A_\alpha)_{\mathfrak p_\alpha}.
\]
Consequently,
\[
  \varinjlim_\alpha(B_\alpha)_{\mathfrak p'_\alpha}
  =
  \varinjlim_\alpha(A_\alpha)_{\mathfrak p_\alpha}
  =
  A_{\mathfrak p}
\]
is a localization of $B$ at a prime ideal
\sref{IV.5.13.3}.  It is therefore a domain and is integrally
closed, proving the corollary.
\end{proof}

\begin{proposition}[5.13.7]
\label{IV.5.13.7}
Suppose that every $A_\alpha$ is regular and that, for
$\alpha\leq\beta$, the homomorphism $\varphi_{\beta\alpha}$
makes $A_\beta$ a flat $A_\alpha$-module.  Suppose in addition
that
\[
  A=\varinjlim_\alpha A_\alpha
\]
is Noetherian.  Then $A$ is regular.
\end{proposition}

\begin{proof}
By definition, it is enough to prove that the Noetherian local
ring $A_{\mathfrak p}$ is regular for every prime ideal
$\mathfrak p$ of $A$.  Put
\[
  \mathfrak p_\alpha
  =\varphi_\alpha^{-1}(\mathfrak p).
\]
By \sref{IV.5.13.3},
\[
  \mathfrak p=\varinjlim_\alpha\mathfrak p_\alpha,
  \qquad
  A_{\mathfrak p}
  =\varinjlim_\alpha(A_\alpha)_{\mathfrak p_\alpha}.
\]
The rings $(A_\alpha)_{\mathfrak p_\alpha}$ are regular, and for
$\alpha\leq\beta$ the ring
$(A_\beta)_{\mathfrak p_\beta}$ is flat over
$(A_\alpha)_{\mathfrak p_\alpha}$ \sref[0]{0.6.3.2}.  We are
therefore reduced to the case in which all the $A_\alpha$, as
well as $A$, are local.

To prove that $A$ is regular, it is enough to show that for any
two finite $A$-modules $M$ and $N$, there is an integer $i_0$
such that

\oldpage[IV-2]{134}

\[
  \operatorname{Tor}^{A}_i(M,N)=0
  \qquad(i\geq i_0)
\]
\sref[0]{0.17.3.1} and \sref[0]{0.17.2.6}.  We use the following
lemma, in which the rings $A_\alpha$ are again arbitrary.
\end{proof}

\begin{lemma}[5.13.7.1]
\label{IV.5.13.7.1}
For every finitely presented $A$-module $M$, there are an index
$\alpha$ and a finitely presented $A_\alpha$-module $M_\alpha$
such that
\[
  M\simeq M_\alpha\otimes_{A_\alpha}A.
\]
\end{lemma}

\begin{proof}
We may write $M=A^r/P$, where $P$ is a finitely generated
$A$-submodule of $A^r$.  Since
\[
  A^r=\varinjlim_\alpha A_\alpha^r
\]
and $P$ is finitely generated, there are an index $\alpha$ and a
finitely generated $A_\alpha$-submodule
$P_\alpha\subset A_\alpha^r$ whose canonical image generates
$P$.  Right exactness of the tensor product then gives
\[
  M
  \simeq
  (A_\alpha^r/P_\alpha)\otimes_{A_\alpha}A.
\]
\end{proof}

\begin{proof}[Completion of the proof of {\hyperlink{ega4.c04.prop.5-13-7}{Proposition~5.13.7}}]
Since $A$ is Noetherian, the finite modules $M$ and $N$ are
finitely presented.  After increasing $\alpha$ if necessary, the
lemma gives
\[
  M=M_\alpha\otimes_{A_\alpha}A,
  \qquad
  N=N_\alpha\otimes_{A_\alpha}A.
\]
The $A_\alpha$-module $A$ is flat \sref[0]{0.6.1.2}, and base
change for Tor gives
\[
  \operatorname{Tor}^{A}_i(M,N)
  \simeq
  \operatorname{Tor}^{A_\alpha}_i(M_\alpha,N_\alpha)
  \otimes_{A_\alpha}A
\]
\sref[III]{III.6.3.9}.  Since $A_\alpha$ is regular, the
right-hand side vanishes for all sufficiently large $i$.  This
proves the proposition.
\end{proof}

\begin{corollary}[5.13.8]
\label{IV.5.13.8}
Let $k\geq0$.  Suppose that the $A_\alpha$ are Noetherian and
satisfy $(\mathrm R_k)$.  Suppose also that, for
$\alpha\leq\beta$, the homomorphism $\varphi_{\beta\alpha}$
makes $A_\beta$ a flat $A_\alpha$-module, that
$A=\varinjlim A_\alpha$ is Noetherian, and that for every
$\alpha$ the induced map
\[
  {}^{a}\varphi_\alpha:
  \Spec(A)\longrightarrow\Spec(A_\alpha)
\]
is a homeomorphism.  Then $A$ satisfies $(\mathrm R_k)$.
\end{corollary}

\begin{proof}
Let $\mathfrak p$ be a prime ideal of $A$, and put
\[
  \mathfrak p_\alpha
  =\varphi_\alpha^{-1}(\mathfrak p).
\]
Then
\[
  \mathfrak p=\varinjlim_\alpha\mathfrak p_\alpha,
  \qquad
  A_{\mathfrak p}
  =\varinjlim_\alpha(A_\alpha)_{\mathfrak p_\alpha}
\]
by \sref{IV.5.13.3}.  It follows from
\sref[I]{I.2.4.2} and the hypothesis that
\[
  \Spec(A_{\mathfrak p})
  \longrightarrow
  \Spec\bigl((A_\alpha)_{\mathfrak p_\alpha}\bigr)
\]
is a surjective homeomorphism for every $\alpha$.  Hence
\[
  \dim(A_{\mathfrak p})
  =
  \dim\bigl((A_\alpha)_{\mathfrak p_\alpha}\bigr)
  \qquad\text{for every }\alpha.
\]
If $\dim(A_{\mathfrak p})\leq k$, then each
$(A_\alpha)_{\mathfrak p_\alpha}$ is regular.  Moreover, for
$\alpha\leq\beta$,
$(A_\beta)_{\mathfrak p_\beta}$ is flat over
$(A_\alpha)_{\mathfrak p_\alpha}$ \sref[0]{0.6.3.2}.
Proposition~\sref{IV.5.13.7} therefore shows that
$A_{\mathfrak p}$ is regular.  Thus $A$ satisfies
$(\mathrm R_k)$.
\end{proof}
% END NATIVE INPUT components/13b_section513_normality_regularity_and_Rk
% END NATIVE INPUT ega4-5.source_aligned_r1
% BEGIN NATIVE INPUT source_aligned/ega4-6
\section{Flat morphisms of locally Noetherian preschemes}
\label{section:IV.6}

\oldpage[IV-2]{134}
Let $X$, $Y$ be two locally Noetherian preschemes, $f:X\to Y$ a morphism.
For every $y\in Y$, the fibre $f^{-1}(y)=X\times_Y\Spec(\boldsymbol{k}(y))$ is also a locally Noetherian prescheme, since it suffices to see this when $Y=\Spec(\mathrm{A})$ and $X=\Spec(\mathrm{B})$ are spectra of Noetherian rings, and then $\mathrm{B}\otimes_\mathrm{A}\boldsymbol{k}(y)$ is a ring of fractions of the quotient ring $\mathrm{B}/\mathfrak{I}_y\mathrm{B}$, hence is Noetherian.
We propose first of all in this paragraph to relate the properties of $X$, of $Y$, and of the fibres $f^{-1}(y)$ (where $y$ ranges over $Y$), under the hypothesis that the morphism $f$ is \emph{flat}.
The questions treated will reduce to the study of relations between the Noetherian local rings $\sh{O}_y$, $\sh{O}_x$, and $\sh{O}_x\otimes_{\sh{O}_y}\boldsymbol{k}(y)$, the homomorphism $\sh{O}_y\to\sh{O}_x$ being local and making $\sh{O}_x$ a flat $\sh{O}_y$-module.
In n\textsuperscript{os}~{\hyperlink{ega4.c04.subsec.6-11}{6.11}} to~{\hyperlink{ega4.c04.subsec.6-13}{6.13}} (rather separate from the rest of the paragraph, by their ``absolute'' rather than ``relative'' nature), we apply certain preceding results to find
\oldpage[IV-2]{135}
criteria allowing us to affirm that the \emph{singular locus} (or certain analogous sets) of certain preschemes are closed sets; these criteria will play an important role in~\textsection\,7.

\subsection{Flatness and dimension}
\label{subsection:IV.6.1}

\begin{proposition}[6.1.1]
\label{IV.6.1.1}
Let $\mathrm{A}$, $\mathrm{B}$ be two local Noetherian rings, $\mathfrak{m}$ the maximal ideal of $\mathrm{A}$, $k=\mathrm{A}/\mathfrak{m}$ its residue field, $\varphi:\mathrm{A}\to\mathrm{B}$ a local homomorphism.
Suppose that for every prime ideal $\mathfrak{p}$ of $\mathrm{A}$, distinct from $\mathfrak{m}$, and for every minimal element $\mathfrak{q}$ of the set of prime ideals of $\mathrm{B}$ containing $\mathfrak{p}\mathrm{B}$, $\varphi^{-1}(\mathfrak{q})$ is distinct from $\mathfrak{m}$ (in other words, no irreducible component of $\Spec(\mathrm{B}/\mathfrak{p}\mathrm{B})$ is contained in the inverse image of the closed point $\mathfrak{m}$ of $\Spec(\mathrm{A})$ in $\Spec(\mathrm{B})$).
Then we have
\begin{equation}
  \dim(\mathrm{B})=\dim(\mathrm{A})+\dim(\mathrm{B}\otimes_\mathrm{A} k).
  \tag{6.1.1.1}
\label{IV.6.1.1.1}
\end{equation}
\end{proposition}

\begin{proof}
We argue by induction on $n=\dim(\mathrm{A})$; the assertion is trivial for $n=0$, $\mathfrak{m}\mathrm{B}$ being then contained in the nilradical of $\mathrm{B}$.
We may therefore suppose $n>0$.
Let $\mathfrak{q}_i$ ($1\leqslant i\leqslant m$) be the minimal prime ideals of $\mathrm{B}$, and let $\mathfrak{p}_i=\varphi^{-1}(\mathfrak{q}_i)$; for every $i$, we have $\mathfrak{p}_i\neq\mathfrak{m}$; otherwise (since $n>0$) there would exist a prime ideal $\mathfrak{p}\neq\mathfrak{m}$ contained in $\mathfrak{p}_i$ and as $\mathfrak{q}_i$ is minimal among the prime ideals of $\mathrm{B}$ containing $\mathfrak{p}\mathrm{B}$, we would arrive at a contradiction with the hypothesis.
Therefore $\mathfrak{m}$ is distinct from the union of the $\mathfrak{p}_i$ and of the minimal prime ideals $\mathfrak{p}'_j$ ($1\leqslant j\leqslant r$) of $\mathrm{A}$ (Bourbaki, \emph{Alg.\ comm.}, chap.~II, \textsection\,1, n\textsuperscript{o}\,1, prop.~2) and there exists an $a\in\mathfrak{m}$ not belonging to any of the $\mathfrak{p}_i$ nor any of the $\mathfrak{p}'_j$.
Set $\mathrm{A}'=\mathrm{A}/a\mathrm{A}$, $\mathrm{B}'=\mathrm{B}/a\mathrm{B}$; we have \sref[0]{0.16.3.4}
\[
  \dim(\mathrm{A}')=\dim(\mathrm{A})-1,\quad
  \dim(\mathrm{B}')=\dim(\mathrm{B})-1
\]
by construction of $a$, and on the other hand $\mathrm{B}'\otimes_{\mathrm{A}'}k=\mathrm{B}\otimes_\mathrm{A} k$, hence
\[
  \dim(\mathrm{B}'\otimes_{\mathrm{A}'}k)=\dim(\mathrm{B}\otimes_\mathrm{A} k)
\]
and it suffices to prove \sref{IV.6.1.1} for $\mathrm{A}'$ and $\mathrm{B}'$; now, by virtue of the bijective correspondence between ideals of $\mathrm{A}$ (resp.\ $\mathrm{B}$) containing $a$ and ideals of $\mathrm{A}'$ (resp.\ $\mathrm{B}'$), the hypotheses of the statement are also valid for $\mathrm{A}'$ and $\mathrm{B}'$; we may therefore apply the induction hypothesis, which completes the proof.
\end{proof}

\begin{corollary}[6.1.2]
\label{IV.6.1.2}
Let $\mathrm{A}$, $\mathrm{B}$ be two local Noetherian rings, $k$ the residue field of $\mathrm{A}$, $\varphi:\mathrm{A}\to\mathrm{B}$ a local homomorphism, $\mathrm{M}\neq 0$ an $\mathrm{A}$-module of finite type, $\mathrm{N}\neq 0$ a $\mathrm{B}$-module of finite type.
If $\mathrm{N}$ is a flat $\mathrm{A}$-module \textnormal{(}resp.\ if $\mathrm{B}$ is a flat $\mathrm{A}$-module\textnormal{)}, then we have
\begin{equation}
  \dim_\mathrm{B}(\mathrm{M}\otimes_\mathrm{A}\mathrm{N})=\dim_\mathrm{A}(\mathrm{M})+\dim_{\mathrm{B}\otimes_\mathrm{A} k}(\mathrm{N}\otimes_\mathrm{A} k)
  \label{IV.6.1.2.1}
\tag{6.1.2.1}
\end{equation}
\textnormal{(}resp.\ \sref{IV.6.1.1}\textnormal{)}.
\end{corollary}

\begin{proof}
It suffices to prove the assertion concerning $\mathrm{N}$; on the other hand, if $\mathfrak{b}$ is the annihilator of $\mathrm{N}$, we may replace $\mathrm{B}$ by $\mathrm{B}/\mathfrak{b}$, and therefore suppose that $\operatorname{Supp}(\mathrm{N})=\Spec(\mathrm{B})$; the hypothesis then means that the morphism $f:\Spec(\mathrm{B})\to\Spec(\mathrm{A})$ corresponding to $\varphi$ is \emph{quasi-flat} \sref{IV.2.3.3}; we conclude \sref{IV.2.3.4} that for every prime ideal $\mathfrak{p}\neq\mathfrak{m}$ of $\mathrm{A}$,
\oldpage[IV-2]{136}
every irreducible component of $f^{-1}(\mathrm{V}(\mathfrak{p}))$ \emph{dominates} $\mathrm{V}(\mathfrak{p})$, and therefore the condition of \sref{IV.6.1.1} is satisfied, whence the conclusion.
\end{proof}

\begin{corollary}[6.1.3]
\label{IV.6.1.3}
Let $\mathrm{A}$, $\mathrm{B}$ be two local Noetherian rings, $\mathfrak{m}$ the maximal ideal of $\mathrm{A}$, $k$ its residue field, $\varphi:\mathrm{A}\to\mathrm{B}$ a local homomorphism.
Suppose that $\dim(\mathrm{B}\otimes_\mathrm{A} k)=0$ (that is to say \sref[0]{0.16.2.1} that $\mathfrak{m}\mathrm{B}$ is an \emph{ideal of definition} of $\mathrm{B}$).
Then we have $\dim(\mathrm{B})\leqslant\dim(\mathrm{A})$.
If in addition there exists a $\mathrm{B}$-module of finite type $\mathrm{N}$ that is a flat $\mathrm{A}$-module and has support equal to $\Spec(\mathrm{B})$ (which holds in particular if $\mathrm{B}$ is a flat $\mathrm{A}$-module), then $\dim(\mathrm{B})=\dim(\mathrm{A})$.
\end{corollary}

\begin{proof}
The first assertion follows from \sref{IV.5.5.1} and the second from \sref{IV.6.1.2}.
\end{proof}

\begin{corollary}[6.1.4]
\label{IV.6.1.4}
Let $X$, $Y$ be two locally Noetherian preschemes, $f:X\to Y$ a surjective morphism.
For every closed subset $Z$ of $Y$, we have
\begin{equation}
  \operatorname{codim}(f^{-1}(Z),X)\leqslant\operatorname{codim}(Z,Y)
  \label{IV.6.1.4.1}
\tag{6.1.4.1}
\end{equation}
Moreover, if $f$ is quasi-flat \sref{IV.2.3.3}, the two sides of \sref{IV.6.1.4.1} are equal.
\end{corollary}

\begin{proof}
Indeed, if $y$ is a maximal point of $Z$ and $x$ a maximal point of the fibre $f^{-1}(y)$, we have $\dim(\sh{O}_{X,x})\leqslant\dim(\sh{O}_{Y,y})$ by virtue of \sref{IV.5.5.2}; the inequality \sref{IV.6.1.4.1} then follows from \sref{IV.5.1.2.1} and from \sref{IV.5.1.3.1} and from the fact that $f$ is surjective.
If in addition $f$ is quasi-flat, we know \sref{IV.2.3.4} that every irreducible component of $f^{-1}(Z)$ dominates an irreducible component of $Z$; the maximal points of $f^{-1}(Z)$ are therefore the maximal points of the fibres $f^{-1}(y)$, where $y$ ranges over the set of maximal points of $Z$ \sref[0]{0.2.1.8}, and for each of these maximal points $x$ we have $\dim(\sh{O}_x\otimes_{\sh{O}_y}\boldsymbol{k}(y))=0$; as in addition $\sh{O}_x$ is a flat $\sh{O}_y$-module, we have $\dim(\sh{O}_x)=\dim(\sh{O}_y)$ by virtue of \sref{IV.6.1.1}, whence the equality in \sref{IV.6.1.4.1}, taking into account that $f$ is surjective.
\end{proof}

The proposition \sref{IV.6.1.1} admits the following partial converse:

\begin{proposition}[6.1.5]
\label{IV.6.1.5}
Let $\mathrm{A}$, $\mathrm{B}$ be two local Noetherian rings, $k$ the residue field of $\mathrm{A}$, $\varphi:\mathrm{A}\to\mathrm{B}$ a local homomorphism, $\mathrm{M}$ a $\mathrm{B}$-module of finite type.
Suppose that:
\begin{enumerate}[label=\upshape\arabic{enumi}\textsuperscript{o}]
  \item $\mathrm{A}$ is a regular ring.
  \item $\mathrm{M}$ is a Cohen-Macaulay $\mathrm{B}$-module.
  \item $\dim_\mathrm{B}(\mathrm{M})=\dim(\mathrm{A})+\dim_{\mathrm{B}\otimes_\mathrm{A} k}(\mathrm{M}\otimes_\mathrm{A} k)$.
\end{enumerate}
Then $\mathrm{M}$ is a flat $\mathrm{A}$-module.
\end{proposition}

\begin{proof}
We argue by induction on $n=\dim(\mathrm{A})$.
If $n=0$, $\mathrm{A}$ is a field since it is regular \sref[0]{0.17.1.4} and the assertion is trivial.
Suppose $n>0$; let $\mathfrak{m}$ be the maximal ideal of $\mathrm{A}$, and let $x\in\mathfrak{m}$ be an element not belonging to $\mathfrak{m}^2$; we know then \sref[0]{0.17.1.8} that $\mathrm{A}'=\mathrm{A}/x\mathrm{A}$ is regular and $\dim(\mathrm{A}')=\dim(\mathrm{A})-1$.
Set
\[
  \mathrm{B}'=\mathrm{B}/x\mathrm{B},\quad \mathrm{M}'=\mathrm{M}/x\mathrm{M}=\mathrm{M}\otimes_\mathrm{A}\mathrm{A}';
\]
then
\[
  \mathrm{B}'\otimes_{\mathrm{A}'}k=\mathrm{B}\otimes_\mathrm{A} k,\quad \mathrm{M}'\otimes_{\mathrm{A}'}k=\mathrm{M}\otimes_\mathrm{A} k
\]
and by virtue of \sref[0]{0.16.3.9.2} we have
\[
  \dim_{\mathrm{B}'}(\mathrm{M}')\leqslant\dim(\mathrm{A}')+\dim_{\mathrm{B}'\otimes_{\mathrm{A}'}k}(\mathrm{M}'\otimes_{\mathrm{A}'} k);
\]
\oldpage[IV-2]{137}
we conclude therefore from \sref[0]{0.16.3.4} that we have
\[
  \dim_\mathrm{B}(\mathrm{M})\leqslant\dim_{\mathrm{B}'}(\mathrm{M}')+1\leqslant(\dim(\mathrm{A}')+\dim_{\mathrm{B}\otimes_\mathrm{A} k}(\mathrm{M}\otimes_\mathrm{A} k))+1=\dim(\mathrm{A})+\dim_{\mathrm{B}\otimes_\mathrm{A} k}(\mathrm{M}\otimes_\mathrm{A} k).
\]

Since by hypothesis the extreme terms of these inequalities are equal, we have necessarily: (i) $\dim_{\mathrm{B}'}(\mathrm{M}')=\dim_\mathrm{B}(\mathrm{M})-1$, and as $\mathrm{M}$ is by hypothesis a Cohen-Macaulay $\mathrm{B}$-module, this means that $x$ is M-\emph{regular} \sref[0]{0.16.1.9} and \sref[0]{0.16.5.5}; (ii) $\dim_{\mathrm{B}'}(\mathrm{M}')=\dim(\mathrm{A}')+\dim_{\mathrm{B}'\otimes_{\mathrm{A}'}k}(\mathrm{M}'\otimes_{\mathrm{A}'}k)$; as $\mathrm{M}'$ is a Cohen-Macaulay $\mathrm{B}'$-module by virtue of (i) and of \sref[0]{0.16.1.9} and \sref[0]{0.16.5.5} and $\mathrm{A}'$ is regular, the induction hypothesis proves that $\mathrm{M}'$ is a flat $\mathrm{A}'$-module; we deduce then from \sref[0]{0.10.2.7} that $\mathrm{M}$ is a flat $\mathrm{A}$-module, since $x$ is M-regular by~(i).
\end{proof}

\subsection{Flatness and projective dimension}
\label{subsection:IV.6.2}

\begin{proposition}[6.2.1]
\label{IV.6.2.1}
\medskip\noindent
\begin{enumerate}[label=(\roman*)]
  \item Let $\mathrm{A}$, $\mathrm{B}$ be two rings, $\varphi:\mathrm{A}\to\mathrm{B}$ a homomorphism such that $\mathrm{B}$ is a flat $\mathrm{A}$-module.
Then, for every $\mathrm{A}$-module $\mathrm{M}$, we have
\begin{equation}
  \dim.\operatorname{proj}_\mathrm{B}(\mathrm{M}\otimes_\mathrm{A}\mathrm{B})\leqslant\dim.\operatorname{proj}_\mathrm{A}(\mathrm{M}).
  \label{IV.6.2.1.1}
\tag{6.2.1.1}
\end{equation}

  \item Suppose in addition that $\mathrm{A}$ is a Noetherian ring, $\mathrm{B}$ a faithfully flat $\mathrm{A}$-module, and $\mathrm{M}$ an $\mathrm{A}$-module of finite type; then
\begin{equation}
  \dim.\operatorname{proj}_\mathrm{B}(\mathrm{M}\otimes_\mathrm{A}\mathrm{B})=\dim.\operatorname{proj}_\mathrm{A}(\mathrm{M}).
  \label{IV.6.2.1.2}
\tag{6.2.1.2}
\end{equation}
\end{enumerate}
\end{proposition}

\begin{proof}
\medskip\noindent
\begin{enumerate}[label=(\roman*)]
  \item We may restrict to the case where $n=\dim.\operatorname{proj}_\mathrm{A}(\mathrm{M})$ is finite; there therefore exists a left resolution
\[
  0\to\mathrm{P}_n\to\mathrm{P}_{n-1}\to\ldots\to\mathrm{P}_0\to\mathrm{M}\to 0
\]
of $\mathrm{M}$ by projective $\mathrm{A}$-modules; as $\mathrm{B}$ is a flat $\mathrm{A}$-module, the sequence
\[
  0\to\mathrm{P}_n\otimes_\mathrm{A}\mathrm{B}\to\mathrm{P}_{n-1}\otimes_\mathrm{A}\mathrm{B}\to\ldots\to\mathrm{P}_0\otimes_\mathrm{A}\mathrm{B}\to\mathrm{M}\otimes_\mathrm{A}\mathrm{B}\to 0
\]
is exact; on the other hand, the $\mathrm{P}_i\otimes_\mathrm{A}\mathrm{B}$ are projective $\mathrm{B}$-modules; whence the conclusion.

  \item Suppose that $\dim.\operatorname{proj}_\mathrm{B}(\mathrm{M}\otimes_\mathrm{A}\mathrm{B})=m$, and consider an exact sequence
\[
  0\to\mathrm{R}\to\mathrm{P}_{m-1}\to\mathrm{P}_{m-2}\to\ldots\to\mathrm{P}_0\to\mathrm{M}\to 0
\]
where the $\mathrm{P}_i$ ($0\leqslant i\leqslant m-1$) are projective $\mathrm{A}$-modules of finite type; as $\mathrm{A}$ is Noetherian, $\mathrm{R}$ is also an $\mathrm{A}$-module of finite type.
As $\mathrm{B}$ is a flat $\mathrm{A}$-module, we also have an exact sequence
\[
  0\to\mathrm{R}\otimes_\mathrm{A}\mathrm{B}\to\mathrm{P}_{m-1}\otimes_\mathrm{A}\mathrm{B}\to\ldots\to\mathrm{P}_0\otimes_\mathrm{A}\mathrm{B}\to\mathrm{M}\otimes_\mathrm{A}\mathrm{B}\to 0
\]
and the hypothesis on $\mathrm{M}\otimes_\mathrm{A}\mathrm{B}$ implies that $\mathrm{R}\otimes_\mathrm{A}\mathrm{B}$ is a projective $\mathrm{B}$-module \sref[0]{0.17.2.1}.
As $\mathrm{B}$ is a faithfully flat $\mathrm{A}$-module, we conclude that $\mathrm{R}$ is a projective $\mathrm{A}$-module (Bourbaki, \emph{Alg.\ comm.}, chap.~I\textsuperscript{er}, \textsection\,3, n\textsuperscript{o}\,6, prop.~12); hence $\dim.\operatorname{proj}_\mathrm{A}(\mathrm{M})\leqslant m$, which completes the proof.
\end{enumerate}
\end{proof}

\begin{corollary}[6.2.2]
\label{IV.6.2.2}
Let $f:X\to Y$ be a flat morphism of preschemes, $\sh{E}$ a quasi-coherent $\sh{O}_Y$-Module.
If $\dim.\operatorname{proj}(\sh{E})\leqslant n$ \sref[0]{0.17.2.14}, then $\dim.\operatorname{proj}(f^*(\sh{E}))\leqslant n$.
\end{corollary}

\oldpage[IV-2]{138}
\begin{proposition}[6.2.3]
\label{IV.6.2.3}
Let $\mathrm{A}$, $\mathrm{B}$ be two local Noetherian rings, $k$ the residue field of $\mathrm{A}$, $\varphi:\mathrm{A}\to\mathrm{B}$ a local homomorphism, $\mathrm{N}$ a $\mathrm{B}$-module of finite type.
Suppose that $\mathrm{B}$ and $\mathrm{N}$ are flat $\mathrm{A}$-modules.
Then we have
\begin{equation}
  \dim.\operatorname{proj}_\mathrm{B}(\mathrm{N})=\dim.\operatorname{proj}_{\mathrm{B}\otimes_\mathrm{A} k}(\mathrm{N}\otimes_\mathrm{A} k).
  \label{IV.6.2.3.1}
\tag{6.2.3.1}
\end{equation}
\end{proposition}

\begin{proof}
Consider a left resolution of $\mathrm{N}$ by free $\mathrm{B}$-modules of finite type
\[
  \ldots\to\mathrm{L}_i\to\mathrm{L}_{i-1}\to\ldots\to\mathrm{L}_0\to\mathrm{N}\to 0
\]
As the $\mathrm{L}_i$ and $\mathrm{N}$ are flat $\mathrm{A}$-modules ($\mathrm{B}$ being a flat $\mathrm{A}$-module) it follows from \sref{IV.2.1.10} that the $\mathrm{Z}_i(\mathrm{L}_\bullet)$ are flat $\mathrm{A}$-modules, that $\mathrm{L}_\bullet\otimes_\mathrm{A} k$ is a resolution of $\mathrm{N}\otimes_\mathrm{A} k$ and that $\mathrm{Z}_i(\mathrm{L}_\bullet\otimes_\mathrm{A} k)=\mathrm{Z}_i(\mathrm{L}_\bullet)\otimes_\mathrm{A} k$ for all $i$.
Note that since $\mathrm{B}$ is Noetherian, the $\mathrm{Z}_i(\mathrm{L}_\bullet)$ are $\mathrm{B}$-modules of finite type, hence the $\mathrm{Z}_i(\mathrm{L}_\bullet\otimes_\mathrm{A} k)$ are ($\mathrm{B}\otimes_\mathrm{A} k$)-modules of finite type.
Now, to say that a $\mathrm{B}$-module (resp. a ($\mathrm{B}\otimes_\mathrm{A} k$)-module) of finite type is flat is equivalent to saying that it is \emph{free} or even \emph{projective} \sref[0]{0.10.1.3}.
As the $\mathrm{Z}_i(\mathrm{L}_\bullet)$ are flat $\mathrm{A}$-modules, it follows on the other hand from \sref[0]{0.10.2.5} (where we set $\mathrm{C}=\mathrm{B}$) that, for $\mathrm{Z}_i(\mathrm{L}_\bullet)$ to be a flat $\mathrm{B}$-module, it is necessary and sufficient that $\mathrm{Z}_i(\mathrm{L}_\bullet)\otimes_\mathrm{A} k$ be a ($\mathrm{B}\otimes_\mathrm{A} k$)-module that is flat.
The smallest integer $n$ such that $\mathrm{Z}_{n-1}(\mathrm{L}_\bullet)$ is a free $\mathrm{B}$-module is therefore also the smallest integer $n$ such that $\mathrm{Z}_{n-1}(\mathrm{L}_\bullet\otimes_\mathrm{A} k)$ is a free ($\mathrm{B}\otimes_\mathrm{A} k$)-module, which proves the proposition \sref[0]{0.17.2.1}.
\end{proof}

\subsection{Flatness and depth}
\label{subsection:IV.6.3}

\begin{proposition}[6.3.1]
\label{IV.6.3.1}
Let $\mathrm{A}$, $\mathrm{B}$ be two local Noetherian rings, $k$ the residue field of $\mathrm{A}$, $\varphi:\mathrm{A}\to\mathrm{B}$ a local homomorphism, $\mathrm{M}$ an $\mathrm{A}$-module of finite type, $\mathrm{N}$ a $\mathrm{B}$-module of finite type.
If $\mathrm{N}$ is a flat $\mathrm{A}$-module $\neq 0$, then we have
\begin{equation}
  \operatorname{prof}_\mathrm{B}(\mathrm{M}\otimes_\mathrm{A}\mathrm{N})=\operatorname{prof}_\mathrm{A}(\mathrm{M})+\operatorname{prof}_{\mathrm{B}\otimes_\mathrm{A} k}(\mathrm{N}\otimes_\mathrm{A} k).
  \tag{6.3.1.1}
\label{IV.6.3.1.1}
\end{equation}
\end{proposition}

\begin{proof}
We may restrict to the case where $\mathrm{M}\neq 0$, without which both sides of \sref{IV.6.3.1} are all equal to $+\infty$.
We will argue by induction on the integer $n$ equal to the second member of \sref{IV.6.3.1.1} (which is finite by virtue of the hypotheses, of \sref[0]{0.16.4.6.2} and of the lemma of Nakayama applied to $\mathrm{N}$).
Suppose first $n=0$; then, if $\mathfrak{m}$ and $\mathfrak{n}$ denote the maximal ideals of $\mathrm{A}$ and $\mathrm{B}$ respectively, $\mathfrak{m}$ is a prime ideal associated to $\mathrm{M}$ and $\mathfrak{n}/\mathfrak{m}\mathrm{B}$ a prime ideal of $\mathrm{B}/\mathfrak{m}\mathrm{B}=\mathrm{B}\otimes_\mathrm{A} k$ associated to $\mathrm{N}\otimes_\mathrm{A} k$.
We conclude then from \sref{IV.3.3.1} that $\mathfrak{n}$ is a prime ideal associated to $\mathrm{M}\otimes_\mathrm{A}\mathrm{N}$; hence, by \sref[0]{0.16.4.6}, (i), $\operatorname{prof}_\mathrm{B}(\mathrm{M}\otimes_\mathrm{A}\mathrm{N})=0$.
Suppose therefore $n>0$, and distinguish two cases:

\emph{a)} Suppose $\operatorname{prof}_\mathrm{A}(\mathrm{M})>0$.
Let $x\in\mathfrak{m}$ be an M-regular element, and set
\[
  \mathrm{A}'=\mathrm{A}/x\mathrm{A},\quad\mathrm{B}'=\mathrm{B}/x\mathrm{B},\quad\mathrm{M}'=\mathrm{M}/x\mathrm{M},\quad\mathrm{N}'=\mathrm{N}/x\mathrm{N};
\]
then
\[
  \mathrm{B}'\otimes_{\mathrm{A}'}k=\mathrm{B}\otimes_\mathrm{A} k,\quad \mathrm{N}'\otimes_{\mathrm{A}'}k=\mathrm{N}\otimes_\mathrm{A} k
\]
since $x\in\mathfrak{m}$, and
\[
  \mathrm{M}'\otimes_{\mathrm{A}'}\mathrm{N}'=(\mathrm{M}\otimes_\mathrm{A}\mathrm{N})/x(\mathrm{M}\otimes_\mathrm{A}\mathrm{N});
\]
\oldpage[IV-2]{139}
in addition, as $\mathrm{N}$ is a flat $\mathrm{A}$-module, the hypothesis that $x$ is M-regular implies that $x$ is also ($\mathrm{M}\otimes_\mathrm{A}\mathrm{N}$)-regular.
We thus have \sref[0]{0.16.4.6}, (ii), and \sref[0]{0.16.4.8}
\[
  \operatorname{prof}_{\mathrm{A}'}(\mathrm{M}')=\operatorname{prof}_\mathrm{A}(\mathrm{M})-1,\quad \operatorname{prof}_{\mathrm{B}'}(\mathrm{M}'\otimes_{\mathrm{A}'}\mathrm{N}')=\operatorname{prof}_\mathrm{B}(\mathrm{M}\otimes_\mathrm{A}\mathrm{N})-1
\]
and
\[
  \operatorname{prof}_{\mathrm{B}'\otimes_{\mathrm{A}'}k}(\mathrm{N}'\otimes_{\mathrm{A}'}k)=\operatorname{prof}_{\mathrm{B}\otimes_\mathrm{A} k}(\mathrm{N}\otimes_\mathrm{A} k).
\]

The equality \sref{IV.6.3.1.1} is therefore a consequence of the same relation for $\mathrm{A}'$, $\mathrm{B}'$, $\mathrm{M}'$ and $\mathrm{N}'$; but as $\mathrm{N}$ is a flat $\mathrm{A}$-module, $\mathrm{N}'=\mathrm{N}\otimes_\mathrm{A}\mathrm{A}'$ is a flat $\mathrm{A}'$-module \sref[0]{0.6.2.1}; we may therefore apply the induction hypothesis, which proves \sref{IV.6.3.1.1} in this case.

\emph{b)} Suppose $\operatorname{prof}_{\mathrm{B}\otimes_\mathrm{A} k}(\mathrm{N}\otimes_\mathrm{A} k)>0$.
Consider an element $y\in\mathfrak{n}$ which is ($\mathrm{N}\otimes_\mathrm{A} k$)-regular; it follows from \sref[0]{0.10.2.4} that $y$ is N-regular and that we then have
\[
  \mathrm{N}'=\mathrm{N}/y\mathrm{N}
\]
is a flat $\mathrm{A}$-module, since $\mathrm{N}$ is supposed to be a flat $\mathrm{A}$-module.
Applying \sref[0]{0.6.1.2} to the exact sequence of $\mathrm{A}$-modules
\[
  0\to\mathrm{N}\xrightarrow{y}\mathrm{N}\to\mathrm{N}'\to 0
\]
we conclude that we have isomorphisms
\[
  (\mathrm{M}\otimes_\mathrm{A}\mathrm{N})/y(\mathrm{M}\otimes_\mathrm{A}\mathrm{N})\xrightarrow{\sim}\mathrm{M}\otimes_\mathrm{A}\mathrm{N}' \quad\text{and}\quad \mathrm{N}'\otimes_\mathrm{A} k\xrightarrow{\sim}(\mathrm{N}\otimes_\mathrm{A} k)/y(\mathrm{N}\otimes_\mathrm{A} k)
\]
and in addition that $y$ is ($\mathrm{M}\otimes_\mathrm{A}\mathrm{N}$)-regular.
We thus have
\[
  \operatorname{prof}_{\mathrm{B}\otimes_\mathrm{A} k}(\mathrm{N}'\otimes_\mathrm{A} k)=\operatorname{prof}_{\mathrm{B}\otimes_\mathrm{A} k}(\mathrm{N}\otimes_\mathrm{A} k)-1
\]
and
\[
  \operatorname{prof}_\mathrm{B}(\mathrm{M}\otimes_\mathrm{A}\mathrm{N}')=\operatorname{prof}_\mathrm{B}(\mathrm{M}\otimes_\mathrm{A}\mathrm{N})-1;
\]
the induction hypothesis shows that the relation \sref{IV.6.3.1.1} is valid for $\mathrm{A}$, $\mathrm{B}$, $\mathrm{M}$ and $\mathrm{N}'$, and from what precedes, we deduce \sref{IV.6.3.1.1} for $\mathrm{A}$, $\mathrm{B}$, $\mathrm{M}$ and $\mathrm{N}$.
\end{proof}

\begin{corollary}[6.3.2]
\label{IV.6.3.2}
Under the hypotheses of \sref{IV.6.3.1}, suppose in addition $\mathrm{M}\neq 0$ and $\mathrm{N}\neq 0$; then we have
\begin{equation}
  \operatorname{coprof}_\mathrm{B}(\mathrm{M}\otimes_\mathrm{A}\mathrm{N})=\operatorname{coprof}_\mathrm{A}(\mathrm{M})+\operatorname{coprof}_{\mathrm{B}\otimes_\mathrm{A} k}(\mathrm{N}\otimes_\mathrm{A} k).
  \label{IV.6.3.2.1}
\tag{6.3.2.1}
\end{equation}
\end{corollary}

\begin{proof}
This follows immediately from \sref{IV.6.3.1.1} and \sref{IV.6.1.2} and from the definition of the codepth \sref[0]{0.16.4.9}.
\end{proof}

\begin{corollary}[6.3.3]
\label{IV.6.3.3}
Under the hypotheses of \sref{IV.6.3.1}, suppose in addition $\mathrm{M}\neq 0$ and $\mathrm{N}\neq 0$; for $\mathrm{M}\otimes_\mathrm{A}\mathrm{N}$ to be a Cohen-Macaulay $\mathrm{B}$-module, it is necessary and sufficient that $\mathrm{M}$ be a Cohen-Macaulay $\mathrm{A}$-module and that $\mathrm{N}\otimes_\mathrm{A} k$ be a ($\mathrm{B}\otimes_\mathrm{A} k$)-module of Cohen-Macaulay.
\end{corollary}

\begin{proof}
This follows from the corollary \sref{IV.6.3.2} and from the definition of Cohen-Macaulay modules by the fact that their codepth is zero \sref[0]{0.16.5.1}, taking into account that the condition $\mathrm{N}\neq 0$ is equivalent to $\mathrm{N}\otimes_\mathrm{A} k\neq 0$ by virtue of the lemma of Nakayama.
\end{proof}

\begin{corollary}[6.3.4]
\label{IV.6.3.4}
Under the hypotheses of \sref{IV.6.3.1}, suppose in addition that $\mathrm{N}\otimes_\mathrm{A} k$ is a $\mathrm{B}$-module of finite length; then we have
\begin{equation}
  \operatorname{prof}_\mathrm{B}(\mathrm{M}\otimes_\mathrm{A}\mathrm{N})=\operatorname{prof}_\mathrm{A}(\mathrm{M}).
  \label{IV.6.3.4.1}
\tag{6.3.4.1}
\end{equation}
\oldpage[IV-2]{140}
If in addition $\mathrm{M}\neq 0$ and $\mathrm{N}\neq 0$, then we have
\begin{equation}
  \operatorname{coprof}_\mathrm{B}(\mathrm{M}\otimes_\mathrm{A}\mathrm{N})=\operatorname{coprof}_\mathrm{A}(\mathrm{M}).
  \label{IV.6.3.4.2}
\tag{6.3.4.2}
\end{equation}
\end{corollary}

\begin{proof}
It amounts to the same to say that $\mathrm{N}\otimes_\mathrm{A} k$ is a $\mathrm{B}$-module of finite length or a ($\mathrm{B}\otimes_\mathrm{A} k$)-module of finite length, and we know \sref[0]{0.16.2.3} that modules of finite length $\neq 0$ are of dimension~$0$ and depth~$0$.

The corollary \sref{IV.6.3.4} is applicable in particular when $\mathrm{B}\otimes_\mathrm{A} k$ is an Artinian ring, that is to say when $\mathfrak{m}\mathrm{B}$ is an \emph{ideal of definition} of the ring $\mathrm{B}$.
\end{proof}

\begin{corollary}[6.3.5]
\label{IV.6.3.5}
Let $X$, $Y$ be two locally Noetherian preschemes, $f:X\to Y$ a flat morphism.
\begin{enumerate}[label=(\roman*)]
  \item If, at a point $x\in X$, $\operatorname{coprof}(\sh{O}_x)\leqslant n$, then, setting $y=f(x)$, we have $\operatorname{coprof}(\sh{O}_y)\leqslant n$ and $\operatorname{coprof}(\sh{O}_x\otimes_{\sh{O}_y}\boldsymbol{k}(y))\leqslant n$; in particular, if $\sh{O}_x$ is a Cohen-Macaulay ring, the same is true of $\sh{O}_y$ and of $\sh{O}_x\otimes_{\sh{O}_y}\boldsymbol{k}(y)$.
  \item Conversely, suppose that $\sh{O}_x\otimes_{\sh{O}_y}\boldsymbol{k}(y)$ is a Cohen-Macaulay ring.
Then, if $\operatorname{coprof}(\sh{O}_y)\leqslant n$ \textnormal{(}resp.\ if $\sh{O}_y$ is a Cohen-Macaulay ring\textnormal{)}, we have $\operatorname{coprof}(\sh{O}_x)\leqslant n$ \textnormal{(}resp.\ $\sh{O}_x$ is a Cohen-Macaulay ring\textnormal{)}.
\end{enumerate}
\end{corollary}

\begin{proof}
It suffices to apply \sref{IV.6.3.2} for $\mathrm{M}=\sh{O}_y$ and $\mathrm{N}=\mathrm{B}=\sh{O}_x$.
\end{proof}

\begin{corollary}[6.3.6]
\label{IV.6.3.6}
Let $\mathrm{A}$ be a Cohen-Macaulay ring.
Then $\mathrm{A}'=\mathrm{A}[\mathrm{T}_1,\ldots,\mathrm{T}_n]$ ($\mathrm{T}_i$ indeterminates) is a Cohen-Macaulay ring.
\end{corollary}

\begin{proof}
Indeed, if we set $Y=\Spec(\mathrm{A})$, $X=\Spec(\mathrm{A}')$ and if $f:X\to Y$ is the canonical morphism, $f$ is \emph{flat} (since $\mathrm{A}'$ is a free $\mathrm{A}$-module \sref{IV.2.1.2}) and for every $y\in Y$, $\boldsymbol{k}(y)[\mathrm{T}_1,\ldots,\mathrm{T}_n]$ is a regular ring \sref[0]{0.17.3.7}, hence \emph{a fortiori} a Cohen-Macaulay ring; it therefore suffices to apply \sref{IV.6.3.5}.
\end{proof}

\begin{corollary}[6.3.7]
\label{IV.6.3.7}
Every quotient ring of a Cohen-Macaulay ring is universally catenary.
\end{corollary}

\begin{proof}
This follows from \sref{IV.6.3.6}, from \sref{IV.5.6.1} and from \sref[0]{0.16.5.12}.
\end{proof}

\begin{proposition}[6.3.8]
\label{IV.6.3.8}
Let $\mathrm{A}$ be a local Noetherian ring; suppose that there exists a Cohen-Macaulay $\mathrm{A}$-module $\mathrm{M}$ of finite type whose support is equal to $\Spec(\mathrm{A})$.
Let $\mathrm{B}$ be a quotient ring of $\mathrm{A}$, and let $f:\Spec(\hat{\mathrm{B}})\to\Spec(\mathrm{B})$ be the canonical morphism; then, for every $x\in\Spec(\mathrm{B})$, $f^{-1}(x)$ is a Cohen-Macaulay prescheme.
\end{proposition}

\begin{proof}
As $\mathrm{B}$ is an $\mathrm{A}$-module of finite type, one has $\hat{\mathrm{B}}=\mathrm{B}\otimes_\mathrm{A}\hat{\mathrm{A}}$ \sref[0]{0.7.3.3}; hence $f$ is the restriction to $\Spec(\hat{\mathrm{B}})$ of the canonical morphism $\Spec(\hat{\mathrm{A}})\to\Spec(\mathrm{A})$, and the fibres of these two morphisms at a point of $\Spec(\mathrm{B})$ are the same.
We are therefore reduced to proving the proposition when $\mathrm{B}=\mathrm{A}$.
Let $\mathfrak{p}=\mathfrak{j}_x$ be a prime ideal of $\mathrm{A}$. Since $\mathfrak{p}\in\operatorname{Supp}(\mathrm{M})$ by hypothesis, there exists \sref[0]{0.16.5.6} a $\mathrm{M}$-regular sequence $(t_i)_{1\leqslant i\leqslant r}$ of elements of $\mathfrak{p}$ such that
\[
 \mathrm{N}=\mathrm{M}\big/\Bigl(\sum_{i=1}^r t_i\mathrm{M}\Bigr)
\]
is a Cohen--Macaulay $\mathrm{A}$-module, $\mathfrak{p}$ is a minimal associated prime ideal of $\mathrm{N}$, and
\[
 \dim(\mathrm{N})=\dim(\mathrm{M}/\mathfrak{p}\mathrm{M})=\dim(\mathrm{A}/\mathfrak{p})
\]
\sref[0]{0.16.5.9}. The same argument as at the beginning shows that one may replace $\mathrm{M}$ by $\mathrm{N}$ and $\mathrm{A}$ by $\mathrm{A}/(\sum_{i=1}^r t_i\mathrm{A})$; consequently one may suppose that $\mathfrak{p}$ is a minimal prime ideal of $\mathrm{A}$.
\oldpage[IV-2]{141}
Set, to simplify, $\mathrm{A}'=\hat{\mathrm{A}}$, $\mathrm{M}'=\hat{\mathrm{M}}=\mathrm{M}\otimes_\mathrm{A}\mathrm{A}'$ \sref[0]{0.7.3.3}; we know \sref[0]{0.16.5.2} that $\mathrm{M}'$ is a Cohen-Macaulay $\mathrm{A}'$-module, which implies that for every prime ideal $\mathfrak{p}'$ of $\mathrm{A}'$, $\mathrm{M}'_{\mathfrak{p}'}$ is a Cohen-Macaulay $\mathrm{A}'_{\mathfrak{p}'}$-module \sref[0]{0.16.5.10}.
Taking into account \sref[I]{I.3.6.5}, one sees that if we set $\mathrm{A}''=\mathrm{A}'\otimes_\mathrm{A}\mathrm{A}_\mathfrak{p}$ and $\mathrm{M}''=\mathrm{M}'\otimes_\mathrm{A}\mathrm{A}_\mathfrak{p}$, $\mathrm{M}''$ is a Cohen-Macaulay $\mathrm{A}''$-module.
For every prime ideal $\mathfrak{p}''$ of $\mathrm{A}''$ above $\mathfrak{p}$, $\mathrm{M}''_{\mathfrak{p}''}=\mathrm{M}_\mathfrak{p}\otimes_{\mathrm{A}_\mathfrak{p}}\mathrm{A}''_{\mathfrak{p}''}$ is a Cohen-Macaulay $\mathrm{A}''_{\mathfrak{p}''}$-module \sref[0]{0.16.5.10}; on the other hand, $\mathrm{M}_\mathfrak{p}$ is by hypothesis a Cohen-Macaulay $\mathrm{A}_\mathfrak{p}$-module and $\mathrm{A}''_{\mathfrak{p}''}$ is a flat $\mathrm{A}_\mathfrak{p}$-module since $\mathrm{A}'$ is a flat $\mathrm{A}$-module \sref[0]{0.7.3.3} and \sref{IV.6.3.2}.
We conclude from \sref{IV.6.3.3} that, if $k$ is the residue field of $\mathrm{A}_\mathfrak{p}$, $k\otimes_{\mathrm{A}_\mathfrak{p}}\mathrm{A}''_{\mathfrak{p}''}$ is a Cohen-Macaulay ring.
But since $\mathrm{A}_\mathfrak{p}$ is of dimension~$0$, the prime ideals of $\mathrm{A}'$ correspond bijectively to those of $k\otimes_{\mathrm{A}_\mathfrak{p}}\mathrm{A}''$ \sref[I]{I.3.5.7}, and if $\mathfrak{q}''$ is the prime ideal of $k\otimes_{\mathrm{A}_\mathfrak{p}}\mathrm{A}''$ corresponding to $\mathfrak{p}''$, the local rings $(k\otimes_{\mathrm{A}_\mathfrak{p}}\mathrm{A}'')_{\mathfrak{q}''}$ and $k\otimes_{\mathrm{A}_\mathfrak{p}}\mathrm{A}''_{\mathfrak{p}''}$ are isomorphic.
Therefore \sref[0]{0.16.5.13}, the ring $k\otimes_{\mathrm{A}_\mathfrak{p}}\mathrm{A}''$ is a Cohen-Macaulay ring.
\end{proof}

\subsection{Flatness and property \texorpdfstring{$(\mathrm{S}_k)$}{(Sk)}}
\label{subsection:IV.6.4}

\begin{proposition}[6.4.1]
\label{IV.6.4.1}
Let $X$, $Y$ be two locally Noetherian preschemes, $f:X\to Y$ a morphism, $\sh{E}$ a coherent $\sh{O}_Y$-Module, $\sh{F}$ a coherent $\sh{O}_X$-Module and $f$-flat.
\begin{enumerate}[label=(\roman*)]
  \item Let $x\in\operatorname{Supp}(\sh{F})$ be such that $\sh{E}\otimes_{\sh{O}_Y}\sh{F}$ has property $(\mathrm{S}_k)$ at the point $x$; then $\sh{E}$ has property $(\mathrm{S}_k)$ at the point $y=f(x)$.
  \item Suppose that for every $y\in f(\operatorname{Supp}(\sh{F}))$, $\sh{F}_y=\sh{F}\otimes_{\sh{O}_y}\boldsymbol{k}(y)$ has property $(\mathrm{S}_k)$; then, if for a point $y\in f(\operatorname{Supp}(\sh{F}))$, $\sh{E}$ has property $(\mathrm{S}_k)$ at the point $y$, $\sh{E}\otimes_{\sh{O}_Y}\sh{F}$ has property $(\mathrm{S}_k)$ at every point of $f^{-1}(y)$.
\end{enumerate}
\end{proposition}

\begin{proof}
\medskip\noindent
\begin{enumerate}[label=(\roman*)]
  \item We know (\textbf{Err}\textsubscript{\textbf{III}},~30) that there exists a closed sub-prescheme $X'$ of $X$ having for underlying space $\operatorname{Supp}(\sh{F})$ and a coherent $\sh{O}_{X'}$-Module $\sh{F}'$ such that $\sh{F}=j_*(\sh{F}')$, where $j$ is the canonical injection $X'\hookrightarrow X$; we may replace $X$ by $X'$ and $\sh{F}$ by $\sh{F}'$, in other words suppose that $\operatorname{Supp}(\sh{F})=X$.
Let $y'$ be a generization of $y$ in $Y$; we know \sref{IV.2.3.4} that there is a generization $x'$ of $x$ in $X$ such that $y'=f(x')$; we may in addition suppose that $x'$ is the generic point of an irreducible component of $f^{-1}(y')$; we then have $\dim(\sh{F}_{x'}\otimes_{\sh{O}_{y'}}\boldsymbol{k}(y'))=0$.
By virtue of \sref{IV.6.1.2}, we then have, setting $\sh{G}=\sh{E}\otimes_{\sh{O}_Y}\sh{F}$,
\[
  \dim(\sh{G}_{x'})=\dim(\sh{E}_{y'})
\]
and by virtue of \sref{IV.6.3.1}
\[
  \operatorname{prof}(\sh{G}_{x'})=\operatorname{prof}(\sh{E}_{y'})
\]
taking into account that the depth is always at most equal to the dimension.
By hypothesis, we have
\[
  \operatorname{prof}(\sh{G}_{x'})\geqslant\inf(k,\dim(\sh{G}_{x'}))
\]
hence
\[
  \operatorname{prof}(\sh{E}_{y'})\geqslant\inf(k,\dim(\sh{E}_{y'}))
\]
which proves the first assertion.

\oldpage[IV-2]{142}
  \item As for every generization $x'$ of $x\in f^{-1}(y)$, $y'=f(x')$ is a generization of $y$, we may restrict to verifying that if $x\in\operatorname{Supp}(\sh{G})$ and $f(x)=y$, we have $\operatorname{prof}(\sh{G}_x)\geqslant\inf(k,\dim(\sh{G}_x))$; it follows from \sref{IV.6.1.2} and \sref{IV.6.3.1} that we have
\[
  \dim(\sh{G}_x)=\dim(\sh{E}_y)+\dim(\sh{F}_x\otimes_{\sh{O}_y}\boldsymbol{k}(y))
\]
\[
  \operatorname{prof}(\sh{G}_x)=\operatorname{prof}(\sh{E}_y)+\operatorname{prof}(\sh{F}_x\otimes_{\sh{O}_y}\boldsymbol{k}(y)).
\]

By hypothesis, we have
\[
  \operatorname{prof}(\sh{E}_y)\geqslant\inf(k,\dim(\sh{E}_y))
\]
\[
  \operatorname{prof}(\sh{F}_x\otimes_{\sh{O}_y}\boldsymbol{k}(y))\geqslant\inf(k,\dim(\sh{F}_x\otimes_{\sh{O}_y}\boldsymbol{k}(y)))
\]
whence, adding term by term,
\begin{align*}
  \operatorname{prof}(\sh{G}_x)&\geqslant\inf(k,\dim(\sh{E}_y))+\inf(k,\dim(\sh{F}_x\otimes_{\sh{O}_y}\boldsymbol{k}(y)))\geqslant{}\\
  &\geqslant\inf(k,\dim(\sh{E}_y)+\dim(\sh{F}_x\otimes_{\sh{O}_y}\boldsymbol{k}(y)))=\inf(k,\dim(\sh{G}_x))
\end{align*}
which proves~(ii).
\end{enumerate}
\end{proof}

\begin{corollary}[6.4.2]
\label{IV.6.4.2}
Let $X$, $Y$ be two locally Noetherian preschemes, $f:X\to Y$ a flat morphism, $\sh{E}$ a coherent $\sh{O}_Y$-Module.
Suppose that for every $y\in Y$, the prescheme $f^{-1}(y)$ has property $(\mathrm{S}_k)$.
Then, for $f^*(\sh{E})$ to have property $(\mathrm{S}_k)$ at a point $x$, it is necessary and sufficient that $\sh{E}$ have property $(\mathrm{S}_k)$ at the point $f(x)$.
\end{corollary}

\begin{remark}[6.4.3]
\label{IV.6.4.3}
Let $\mathrm{A}$, $\mathrm{B}$ be two local Noetherian rings, $k$ the residue field of $\mathrm{A}$, $\varphi:\mathrm{A}\to\mathrm{B}$ a local homomorphism, $\mathrm{M}$ a $\mathrm{B}$-module of finite type that is a flat $\mathrm{A}$-module, and suppose in addition that the $\mathrm{A}$-module $\mathrm{A}$ and the ($\mathrm{B}\otimes_\mathrm{A} k$)-module $\mathrm{M}\otimes_\mathrm{A} k$ have property $(\mathrm{S}_n)$; can one then conclude that the $\mathrm{B}$-module $\mathrm{M}$ has property $(\mathrm{S}_n)$?
We do not know even when $n=1$, $\mathrm{M}=\mathrm{B}$ and $\mathrm{B}=\hat{\mathrm{A}}$, which is to say we do not know whether in general the \emph{completion} of a local Noetherian ring satisfying $(\mathrm{S}_n)$ also satisfies $(\mathrm{S}_n)$, even for $n=1$.
Setting $Y=\Spec(\mathrm{A})$, $X=\Spec(\mathrm{B})$ and $\sh{F}=\widetilde{\mathrm{M}}$, it would suffice, by virtue of \sref{IV.6.4.1}, (ii), to show that for \emph{every} $y\in Y$, $\sh{F}_y$ has property $(\mathrm{S}_k)$ (and not only when $y$ is the closed point of $Y$); it would also suffice to show that the set $U$ of $x\in X$ such that $\sh{F}_x\otimes_{\sh{O}_{f(x)}}\boldsymbol{k}(f(x))$ satisfies $(\mathrm{S}_k)$ is \emph{open} in $X$ (since it contains the closed point of $X$ by hypothesis).
We will show further on \sref{IV.12.1.4} that $U$ is open when $\mathrm{B}$ is a local ring of a finite type $\mathrm{A}$-algebra and $\mathrm{M}=\mathrm{B}$.
On the other hand, for $\mathrm{B}=\hat{\mathrm{A}}$ and $\mathrm{M}=\mathrm{B}$, we have seen in \sref{IV.6.3.8} that the preschemes $f^{-1}(y)$ are Cohen-Macaulay preschemes (in other words, satisfy $(\mathrm{S}_n)$ for \emph{every} $n$) when $\mathrm{A}$ is a quotient of a local Cohen-Macaulay ring.
We conclude that when $\mathrm{A}$ is a quotient of a local Cohen-Macaulay ring (or, more generally, of a local ring satisfying the hypotheses of \sref{IV.6.3.8}), \emph{for $\mathrm{A}$ to satisfy} $(\mathrm{S}_n)$, \emph{it is necessary and sufficient that its completion} $\hat{\mathrm{A}}$ \emph{satisfy} $(\mathrm{S}_n)$.
It would remain to see whether this property holds without restriction on $\mathrm{A}$.
\end{remark}

\subsection{Flatness and property \texorpdfstring{$(\mathrm{R}_k)$}{(Rk)}}
\label{subsection:IV.6.5}
\oldpage[IV-2]{143}

\begin{proposition}[6.5.1]
\label{IV.6.5.1}
Let $\mathrm{A}$, $\mathrm{B}$ be two local Noetherian rings, $k$ the residue field of $\mathrm{A}$, $\varphi:\mathrm{A}\to\mathrm{B}$ a local homomorphism, for which $\mathrm{B}$ is a flat $\mathrm{A}$-module.
Then:
\begin{enumerate}[label=(\roman*)]
  \item If $\mathrm{B}$ is regular, $\mathrm{A}$ is regular.
  \item If $\mathrm{A}$ and $\mathrm{B}\otimes_\mathrm{A} k$ are regular, $\mathrm{B}$ is regular.
\end{enumerate}
This proposition is given here for the record, having already been proved in \sref[0]{0.17.3.3}.
\end{proposition}

\begin{corollary}[6.5.2]
\label{IV.6.5.2}
Let $X$, $Y$ be two locally Noetherian preschemes, $f:X\to Y$ a flat morphism.
\begin{enumerate}[label=(\roman*)]
  \item If $X$ is regular at a point $x$, $Y$ is regular at the point $f(x)$.
  \item If for a $y\in f(X)$, $Y$ is regular at the point $y$ and if $f^{-1}(y)$ is a regular prescheme at a point $x$, then $X$ is regular at the point $x$.
\end{enumerate}
\end{corollary}

\begin{proposition}[6.5.3]
\label{IV.6.5.3}
Let $X$, $Y$ be two locally Noetherian preschemes, $f:X\to Y$ a flat morphism.
\begin{enumerate}[label=(\roman*)]
  \item If $X$ has property $(\mathrm{R}_k)$ at a point $x$ \textnormal{(}resp.\ has property $(\mathrm{R}_k)$\textnormal{)}, $Y$ has property $(\mathrm{R}_k)$ at the point $f(x)$ \textnormal{(}resp.\ $Y$ has property $(\mathrm{R}_k)$ at all points of $f(X)$\textnormal{)}.
  \item Suppose that for every $y\in f(X)$, the prescheme $f^{-1}(y)$ has property $(\mathrm{R}_k)$; then, if for a point $y\in f(X)$, $Y$ has property $(\mathrm{R}_k)$ at the point $y$, $X$ has property $(\mathrm{R}_k)$ at every point of $f^{-1}(y)$.
\end{enumerate}
\end{proposition}

\begin{proof}
\medskip\noindent
\begin{enumerate}[label=(\roman*)]
  \item Set $y=f(x)$ and let $y'$ be a generization of $y$; as in the proof of \sref{IV.6.4.1}, there is a generic point $x'$ of an irreducible component of $f^{-1}(y')$ which is a generization of $x$; hence $\dim(\sh{O}_{x'}\otimes_{\sh{O}_{y'}}\boldsymbol{k}(y'))=0$ and by virtue of the hypothesis and of \sref{IV.6.1.2}, $\dim(\sh{O}_{x'})=\dim(\sh{O}_{y'})$.
The hypothesis implies, provided that $\dim(\sh{O}_{x'})>k$, that $\sh{O}_{x'}$ is a regular ring, and then $\sh{O}_{y'}$ is regular by virtue of \sref{IV.6.5.1}.

  \item As for every generization $x'$ of $x\in f^{-1}(y)$, $y'=f(x')$ is a generization of $y$, we may restrict to verifying that if $\dim(\sh{O}_x)\leqslant k$, then $\sh{O}_x$ is a regular ring.
Now, by virtue of \sref{IV.6.1.2}
\[
  \dim(\sh{O}_x)=\dim(\sh{O}_y)+\dim(\sh{O}_x\otimes_{\sh{O}_y}\boldsymbol{k}(y))
\]
hence if $\dim(\sh{O}_x)\leqslant k$, we have \emph{a fortiori} $\dim(\sh{O}_y)\leqslant k$ and $\dim(\sh{O}_x\otimes_{\sh{O}_y}\boldsymbol{k}(y))\leqslant k$ and the hypothesis implies that $\sh{O}_y$ and $\sh{O}_x\otimes_{\sh{O}_y}\boldsymbol{k}(y)$ are regular rings.
We conclude then from \sref{IV.6.5.1} that $\sh{O}_x$ is a regular ring.
\end{enumerate}
\end{proof}

\begin{corollary}[6.5.4]
\label{IV.6.5.4}
Let $X$, $Y$ be two locally Noetherian preschemes, $f:X\to Y$ a flat morphism.
\begin{enumerate}[label=(\roman*)]
  \item If $X$ is normal at a point $x$, $Y$ is normal at the point $f(x)$.
  \item If, for every $y\in f(X)$, $f^{-1}(y)$ is a normal prescheme and if $Y$ is normal at a point $y\in f(X)$, then $X$ is a normal prescheme at every point of $f^{-1}(y)$.
\end{enumerate}
\end{corollary}

\begin{proof}
This follows immediately from the normality criterion of Serre \sref{IV.5.8.6} and from \sref{IV.6.4.1} applied for $k=2$ and \sref{IV.6.5.3} applied for $k=1$.
\end{proof}

\begin{remarks}[6.5.5]
\label{IV.6.5.5}
\oldpage[IV-2]{144}
\medskip\noindent
\begin{enumerate}[label=(\roman*)]
  \item Let $\mathrm{A}$, $\mathrm{B}$ be two local Noetherian rings, $\varphi:\mathrm{A}\to\mathrm{B}$ a local homomorphism such that $\mathrm{B}$ is a flat $\mathrm{A}$-module.
Let $k$ be the residue field of $\mathrm{A}$, and suppose that the two rings $\mathrm{A}$ and $\mathrm{B}\otimes_\mathrm{A} k$ satisfy property $(\mathrm{R}_i)$ \sref{IV.5.8.2}; then \emph{it does not necessarily follow} that $\mathrm{B}$ satisfies $(\mathrm{R}_i)$, even in the particular case where $i=0$ or $i=1$ and where $\mathrm{B}$ is the \emph{completion} $\hat{\mathrm{A}}$ of $\mathrm{A}$.
Nagata has in fact given an example where $\mathrm{A}$ is \emph{normal} (hence satisfies $(\mathrm{R}_1)$) but where $\hat{\mathrm{A}}$ is not even reduced (hence does not satisfy $(\mathrm{R}_0)$)~[30].
We cannot apply the proposition \sref{IV.6.5.3} to this case because the fibres $f^{-1}(y)$ do not necessarily satisfy property $(\mathrm{R}_i)$ at the points distinct from the closed point of $Y=\Spec(\mathrm{A})$.
We will show further on \sref{IV.7.8.3}, (v)) that such phenomena do not arise for the local Noetherian rings that one encounters most often in applications.

  \item The property for a prescheme of being \emph{integral} does not at all behave like the properties we have just examined and the preceding ones: it can happen that $f:X\to Y$ is a flat morphism of finite type, that all the fibres $f^{-1}(y)$ are geometrically regular (and even geometrically regular \sref{IV.6.7.6}) and that $Y$ is integral, without $X$ being even locally integral.
For example, let $k$ be an algebraically closed field of characteristic~$0$, and let $\mathrm{A}$ be the integral ring $k[\mathrm{U},\mathrm{V}]/(\mathrm{U}\mathrm{V}-(\mathrm{U}+\mathrm{V})^3)$ (whose spectrum is therefore a ``cubic with a double point''); $\mathrm{A}$ is not integrally closed, and if $u$, $v$ are the canonical images of $\mathrm{U}$, $\mathrm{V}$ in $\mathrm{A}$, the integral closure of $\mathrm{A}$ is the ring $\mathrm{C}=\mathrm{A}[t]$, with $t=(u-v)/(u+v)$, which satisfies the equation $t^2=1+u+v$, whence $u=\frac{1}{2}(t^3+t^2-t)$, $v=\frac{1}{2}(-t^3+t^2+t)$ and therefore $\mathrm{C}=k[t]$, isomorphic to a polynomial ring in one indeterminate over $k$.
If $\mathfrak{m}=\mathrm{A} u+\mathrm{A} v$, the maximal ideal of $\mathrm{A}$ (corresponding to the ``double point'' of the cubic), there are two maximal ideals $\mathfrak{n}_1$, $\mathfrak{n}_2$ of $\mathrm{C}$, generated respectively by $t-1$ and $t+1$.
Let $\mathrm{B}$ be the subring of the product $\mathrm{C}\times\mathrm{C}$ formed by the pairs of polynomials $(f,g)$ such that $f(1)=g(-1)$ and $f(-1)=g(1)$ ($\Spec(\mathrm{B})$ is the prescheme obtained by ``gluing'' two copies of $\Spec(\mathrm{C})$, the point $\mathfrak{n}_1$ (resp.\ $\mathfrak{n}_2$) of one of the two copies being ``glued'' to the point $\mathfrak{n}_2$ (resp.\ $\mathfrak{n}_1$) of the other; cf.\ chap.~V, where this operation will be discussed in a general way).
There are therefore two maximal ideals $\mathfrak{r}_1$, $\mathfrak{r}_2$ of $\mathrm{B}$ above the maximal ideal $\mathfrak{m}$ of $\mathrm{A}$.
Moreover, as the ``gluing'' process commutes with localisation and completion, one easily verifies that the canonical homomorphisms $\hat{\mathrm{A}}_\mathfrak{m}\to\hat{\mathrm{B}}_{\mathfrak{r}_1}$ and $\hat{\mathrm{A}}_\mathfrak{m}\to\hat{\mathrm{B}}_{\mathfrak{r}_2}$ are bijective; consequently $\mathrm{B}_{\mathfrak{r}_1}$ and $\mathrm{B}_{\mathfrak{r}_2}$ are flat $\mathrm{A}_\mathfrak{m}$-modules (Bourbaki, \emph{Alg.\ comm.}, chap.~III, \textsection\,3, n\textsuperscript{o}\,5, prop.~10).
For every other maximal ideal $\mathfrak{p}$ of $\mathrm{A}$, there are two maximal ideals $\mathfrak{q}_1$, $\mathfrak{q}_2$ of $\mathrm{B}$ above $\mathfrak{p}$, and the homomorphisms $\mathrm{A}_\mathfrak{p}\to\mathrm{B}_{\mathfrak{q}_1}$ and $\mathrm{A}_\mathfrak{p}\to\mathrm{B}_{\mathfrak{q}_2}$ are bijective.
We see thus that the morphism $\Spec(\mathrm{B})\to\Spec(\mathrm{A})$ is flat and finite, and that all its fibres are geometrically regular (it is even \emph{\'etale}, as we will see later \sref{IV.17.6.3}); however it is immediate that $\mathrm{B}$ is not integral.
\end{enumerate}
\end{remarks}

\subsection{Properties of transitivity}
\label{subsection:IV.6.6}
\oldpage[IV-2]{145}

\begin{proposition}[6.6.1]
\label{IV.6.6.1}
For a locally Noetherian prescheme $\mathrm{Z}$, let us denote by $\boldsymbol{P}(\mathrm{Z})$ any one of the following properties:
\begin{enumerate}[label=\emph{\alph*)}]
  \item $\mathrm{Z}$ is a Cohen-Macaulay prescheme.
  \item $\mathrm{Z}$ satisfies $(\mathrm{S}_n)$.
  \item $\mathrm{Z}$ is regular.
  \item $\mathrm{Z}$ satisfies $(\mathrm{R}_n)$.
  \item $\mathrm{Z}$ is normal.
  \item $\mathrm{Z}$ is reduced.
\end{enumerate}
Let then $\mathrm{X}$, $\mathrm{Y}$, $\mathrm{Z}$ be three locally Noetherian preschemes, $f:\mathrm{X}\to\mathrm{Y}$, $g:\mathrm{Y}\to\mathrm{Z}$ two morphisms.
\begin{enumerate}[label=(\roman*)]
  \item Suppose that $f$ and $g$ are flat and that for every $y\in f(\mathrm{X})$ \textnormal{(}resp.\ every $z\in g(\mathrm{Y})$\textnormal{)}, the prescheme $f^{-1}(y)$ \textnormal{(}resp.\ $g^{-1}(z)$\textnormal{)} satisfies $\boldsymbol{P}$. Then $h=g\circ f$ is flat and, for every $z\in h(\mathrm{X})$, $h^{-1}(z)$ satisfies $\boldsymbol{P}$.

  \item Suppose the following conditions are satisfied: $f$ is faithfully flat, $h=g\circ f$ is flat, for every $y\in\mathrm{Y}$, the prescheme $f^{-1}(y)$ satisfies $\boldsymbol{P}$, and for every $z\in g(\mathrm{Y})$, the prescheme $h^{-1}(z)$ satisfies $\boldsymbol{P}$. Then $g$ is flat, and for every $z\in g(\mathrm{Y})$, $g^{-1}(z)$ satisfies $\boldsymbol{P}$.
\end{enumerate}
\end{proposition}

\begin{proof}
(i) We already know \sref[0]{0.2.1.6} that $h$ is flat; furthermore, for every $z\in\mathrm{Z}$, $f_z=f\otimes\mathbf{1}_{\boldsymbol{k}(z)}:h^{-1}(z)\to g^{-1}(z)$ is flat \sref[0]{0.2.1.4}, and for every $y\in g^{-1}(z)$, $f_z^{-1}(y)$ is isomorphic to $f^{-1}(y)$ \sref[I]{I.3.6.4}. The conclusion follows then respectively from \sref{IV.6.3.5},~(ii), \sref{IV.6.4.2}, \sref{IV.6.5.2},~(ii), \sref{IV.6.5.3},~(ii), \sref{IV.6.5.4},~(ii) and \sref{IV.3.3.5},~(ii).

(ii) We already know that $g$ is flat \sref{IV.2.2.13}; in addition, for every $z\in\mathrm{Z}$, $f_z$ is faithfully flat \sref{IV.2.2.13}. The conclusion follows then respectively from \sref{IV.6.3.5},~(i), \sref{IV.6.4.2}, \sref{IV.6.5.2},~(i), \sref{IV.6.5.3},~(i), \sref{IV.6.5.4},~(ii) and \sref{IV.2.1.13}.
\end{proof}

\begin{remark}[6.6.2]
\label{IV.6.6.2}
Suppose $h$ flat, $f$ faithfully flat, and suppose that for every $z\in\mathrm{Z}$, $h^{-1}(z)$ is of coprofundity $\leqslant n$ \sref{IV.5.7.1}; then it follows from the reasoning of \sref{IV.6.6.1},~(ii) and \sref{IV.6.3.2} that $g^{-1}(z)$ is of coprofundity $\leqslant n$ for every $z\in g(\mathrm{Y})$ and that for every $y\in\mathrm{Y}$, $f^{-1}(y)$ is of coprofundity $\leqslant n$.
\end{remark}

\subsection{Application to base changes in algebraic preschemes}
\label{subsection:IV.6.7}

\begin{proposition}[6.7.1]
\label{IV.6.7.1}
Let $k$ be a field, $\mathrm{X}$ a $k$-prescheme locally Noetherian, $\sh{F}$ an $\sh{O}_\mathrm{X}$-Module coherent. Let $k'$ be an extension of $k$; set $\mathrm{X}'=\mathrm{X}\otimes_k k'$, $\sh{F}'=\sh{F}\otimes_k k'$ and let $p:\mathrm{X}'\to\mathrm{X}$ be the canonical projection. We suppose either that $\mathrm{X}$ is locally of finite type over $k$, or that $k'$ is an extension of finite type of $k$, so that in both cases $\mathrm{X}'$ is locally Noetherian. Let $x'$ be a point of $\mathrm{X}'$, $x=p(x')$. Then:
\begin{enumerate}[label=(\roman*)]
  \item We have $\operatorname{coprof}(\sh{F}_x)=\operatorname{coprof}(\sh{F}'_{x'})$; in particular, for $\sh{F}_x$ to be an $\sh{O}_x$-module of Cohen-Macaulay, it is necessary and sufficient that $\sh{F}'_{x'}$ be an $\sh{O}_{x'}$-module of Cohen-Macaulay.
  \item For $\sh{F}$ to satisfy the property $(\mathrm{S}_n)$ at the point $x$, it is necessary and sufficient that $\sh{F}'$ satisfy $(\mathrm{S}_n)$ at the point $x'$.
\end{enumerate}
\end{proposition}

\begin{proof}
\oldpage[IV-2]{146}
We know that $p$ is faithfully flat \sref{IV.2.2.13}; the two assertions are then consequences respectively of \sref{IV.6.3.2} and \sref{IV.6.4.2}, provided that one proves that $p^{-1}(x)=\operatorname{Spec}(\boldsymbol{k}(x)\otimes_k k')$ is a Cohen-Macaulay prescheme; as one of the two fields $\boldsymbol{k}(x)$, $k'$ is an extension of finite type of $k$ by hypothesis, we are reduced to proving the following:
\end{proof}

\begin{lemma}[6.7.1.1]
\label{IV.6.7.1.1}
Let $\mathrm{K}$, $\mathrm{L}$ be two extensions of a field $k$, one of which is of finite type (so that the ring $\mathrm{K}\otimes_k\mathrm{L}$ is Noetherian). Then $\mathrm{K}\otimes_k\mathrm{L}$ is a Cohen-Macaulay ring.
\end{lemma}

\begin{proof}
Suppose for example that $\mathrm{L}$ is an extension of finite type of $k$, so that $\mathrm{L}$ is extension of a pure extension $k(\mathbf{t})$ of $k$ ($\mathbf{t}$ being a system $(t_i)_{1\leqslant i\leqslant m}$ of indeterminates). We then set $\mathrm{A}=\mathrm{K}\otimes_k k(\mathbf{t})$, $\mathrm{B}=\mathrm{K}\otimes_k\mathrm{L}$, $\mathrm{B}$ is an $\mathrm{A}$-module flat \sref[0]{0.6.2.1} and of finite type; the morphism $h:\operatorname{Spec}(\mathrm{B})\to\operatorname{Spec}(\mathrm{A})$ is finite, and as every Artinian prescheme is of Cohen-Macaulay, to prove that $\mathrm{B}$ is a Cohen-Macaulay ring, it suffices to prove that $\mathrm{A}$ is a Cohen-Macaulay ring \sref{IV.6.3.5}. Now, if $\mathrm{S}$ is the set of nonzero elements of $k[\mathbf{t}]$, then $\mathrm{A}=\mathrm{S}^{-1}\mathrm{A}'$, where $\mathrm{A}'=\mathrm{K}\otimes_k k[\mathbf{t}]=\mathrm{K}[\mathbf{t}]$; by \sref[0]{0.16.5.13}, it suffices to prove that $\mathrm{A}'$ is regular \sref[0]{0.17.3.7}.
\end{proof}

\begin{corollary}[6.7.2]
\label{IV.6.7.2}
For $\sh{O}_x$ to be a Cohen--Macaulay ring (resp. for $\mathrm{X}$ to satisfy $(\mathrm{S}_n)$ at the point $x$), it is necessary and sufficient that $\sh{O}_{x'}$ be a Cohen--Macaulay ring (resp. that $\mathrm{X}'$ satisfy $(\mathrm{S}_n)$ at the point $x'$).
\end{corollary}

\begin{corollary}[6.7.3]
\label{IV.6.7.3}
Let $\mathrm{X}$ and $\mathrm{Y}$ be two locally Noetherian $k$-preschemes, at least one of which is locally of finite type over $k$. Let $\sh{F}$ (resp. $\sh{G}$) be a coherent $\sh{O}_{\mathrm{X}}$-Module (resp. a coherent $\sh{O}_{\mathrm{Y}}$-Module). If $\sh{F}$ and $\sh{G}$ satisfy property $(\mathrm{S}_n)$, then so does $\sh{F}\otimes_k\sh{G}$.
\end{corollary}

\begin{proof}
The hypothesis implies that $\mathrm{X}\times_k\mathrm{Y}$ is locally Noetherian; let $p:\mathrm{X}\times_k\mathrm{Y}\to\mathrm{X}$ and $q:\mathrm{X}\times_k\mathrm{Y}\to\mathrm{Y}$ be the canonical projections, which are flat morphisms. Suppose, for example, that $\mathrm{X}$ is locally of finite type over $k$; we may write $\sh{F}\otimes_k\sh{G}=p^*(\sh{F})\otimes_k\sh{G}$, and since $\sh{F}$ is flat relative to the structural morphism $\mathrm{X}\to\operatorname{Spec}(k)$, $p^*(\sh{F})$ is $q$-flat \sref{IV.2.1.4}. Applying the criterion \sref{IV.6.4.1},~(ii) to the morphism $q$, it suffices to see that, for every $y\in\mathrm{Y}$, $p^*(\sh{F})\otimes_{\sh{O}_{\mathrm{Y}}}\boldsymbol{k}(y)$ has property $(\mathrm{S}_n)$; but
$p^*(\sh{F})\otimes_{\sh{O}_{\mathrm{Y}}}\boldsymbol{k}(y)=\sh{F}\otimes_k\boldsymbol{k}(y)$, and since $\mathrm{X}$ is locally of finite type over $k$, the conclusion follows from \sref{IV.6.7.1},~(ii).
\end{proof}

\begin{proposition}[6.7.4]
\label{IV.6.7.4}
For a locally Noetherian prescheme $\mathrm{Z}$ and a point $z\in\mathrm{Z}$, let $\boldsymbol{P}(\mathrm{Z},z)$ be one of the following properties:
\begin{enumerate}[label=\alph*)]
  \item $\mathrm{Z}$ is regular at the point $z$.
  \item $\mathrm{Z}$ satisfies $(\mathrm{R}_n)$ at the point $z$.
  \item $\mathrm{Z}$ is normal at the point $z$.
  \item $\mathrm{Z}$ is reduced at the point $z$.
\end{enumerate}
Under the hypotheses and with the notation of \sref{IV.6.7.1}, if $\boldsymbol{P}(\mathrm{X}',x')$ is true, then $\boldsymbol{P}(\mathrm{X},x)$ is true. The converse is true if $k'$ is a separable extension of $k$.
\end{proposition}

\begin{proof}
The case where $\boldsymbol{P}$ is property d) has been included only for the record, having already been treated in \sref{IV.2.1.13} and \sref{IV.4.6.1}. The same reasoning as at the beginning of the proof of \sref{IV.6.7.1} shows that the first assertion follows respectively from \sref{IV.6.5.2},~(i), \sref{IV.6.5.3},~(i), and \sref{IV.6.5.4},~(i); similarly, the second assertion follows from \sref{IV.6.5.2},~(ii), \sref{IV.6.5.3},~(ii), and \sref{IV.6.5.4},~(ii),

\oldpage[IV-2]{147}

provided that one proves that $p^{-1}(x)$ is a regular prescheme; in other words, we are reduced to proving the following lemma.
\end{proof}

\begin{lemma}[6.7.4.1]
\label{IV.6.7.4.1}
Let $\mathrm{K}$ and $\mathrm{L}$ be two extensions of a field $k$, one of which is of finite type. If $\mathrm{L}$ is a separable extension of $k$, then the ring $\mathrm{K}\otimes_k\mathrm{L}$ is regular.
\end{lemma}

\begin{proof}
We distinguish two cases.

A) $\mathrm{L}$ is an extension of finite type of $k$. With the notation of \sref{IV.6.7.1.1}, we may then suppose that $\mathrm{L}$ is a finite separable extension of $k(\mathbf{t})$. For every $s\in\operatorname{Spec}(\mathrm{A})$, the ring $\boldsymbol{k}(s)\otimes_{k(\mathbf{t})}\mathrm{L}$ is a direct sum of finitely many fields, finite separable extensions of $\boldsymbol{k}(s)$, and is therefore regular; it follows from \sref{IV.6.5.2},~(ii) that it suffices to prove that $\mathrm{A}$ is regular. Since $\mathrm{A}=\mathrm{S}^{-1}\mathrm{A}'$, it suffices to prove that $\mathrm{A}'$ is regular \sref[0]{0.17.3.6}, and this was proved in the proof of \sref{IV.6.7.1.1}.

B) Let $(\mathrm{L}_\lambda)$ be the filtered family of subextensions of $\mathrm{L}$ that are of finite type over $k$. By A), each ring $\mathrm{C}_\lambda=\mathrm{K}\otimes_k\mathrm{L}_\lambda$ is regular; moreover, if $\mathrm{L}_\lambda\subset\mathrm{L}_\mu$, then $\mathrm{L}_\mu$ is a flat $\mathrm{L}_\lambda$-module, so $\mathrm{C}_\mu$ is a flat $\mathrm{C}_\lambda$-module \sref[0]{0.6.2.1}. Since $\mathrm{C}=\mathrm{K}\otimes_k\mathrm{L}=\varinjlim(\mathrm{K}\otimes_k\mathrm{L}_\lambda)$ is Noetherian, we may apply \sref{IV.5.13.7}, and $\mathrm{C}$ is therefore regular.
\end{proof}

\begin{remarks}[6.7.5]
\label{IV.6.7.5}
(i) In the proof of the fact that $\boldsymbol{P}(\mathrm{X}, x)$ implies $\boldsymbol{P}(\mathrm{X}', x')$, one cannot dispense with the hypothesis that $k'$ is a \emph{separable} extension of $k$. We have indeed already seen \sref{IV.4.6.1} when $\boldsymbol{P}$ is the property $d)$; but even when $\mathrm{X}$ is \emph{geometrically integral} \sref{IV.4.6.2}, it can happen that $\mathrm{X}$ is \emph{regular} without $\mathrm{X}'$ being \emph{normal}.

Take, for example, for $\mathrm{X}$ a normal algebraic curve over $k$ \sref[II]{II.7.4.2}. Since the local rings of $\mathrm{X}$ are integrally closed and of dimension~$1$, they are discrete valuation rings and hence regular \sref[II]{II.7.1.6}; consequently $\mathrm{X}$ is a regular $k$-scheme (and \emph{a fortiori} satisfies $(\mathrm{R}_n)$ for every $n\geqslant0$). To say that $\mathrm{X}$ is geometrically integral then means that its field $\mathrm{K}$ of rational functions is a separable and primary extension of $k$ \sref{IV.4.6.3}.

Now let $k$ be a nonperfect field of characteristic $p>2$, and let $a\in k$ be an element not belonging to $k^p$. Let $\mathrm{B}=k[\mathrm{S},\mathrm{T}]$ be the polynomial ring in two indeterminates. The polynomial
\[
 \mathrm{P}(\mathrm{S},\mathrm{T})=\mathrm{T}^2-\mathrm{S}^p+a
\]
is irreducible in $\mathrm{B}$, since one checks at once that $\mathrm{S}^p-a$ is not a square in $k[\mathrm{S}]$; hence $\mathrm{X}=\operatorname{Spec}(\mathrm{A})$, where $\mathrm{A}=\mathrm{B}/\mathrm{PB}$, is an irreducible affine curve over $k$. To see that $\mathrm{X}$ is regular, it suffices to show that it is normal \sref[II]{II.7.4.5}. But $\mathrm{A}=k[\mathrm{S}][t]$, where $t$ is a root of $\mathrm{P}$ considered as a polynomial in $\mathrm{T}$ over $k[\mathrm{S}]$, so the field $\mathrm{K}=\mathrm{R}(\mathrm{X})$ of rational functions of $\mathrm{X}$ is the field of fractions $k(\mathrm{S})[\mathrm{T}]/(\mathrm{P})$ of $\mathrm{A}$, a quadratic extension of $k(\mathrm{S})$, hence separable over $k(\mathrm{S})$ and \emph{a fortiori} over $k$. Since $2$ is invertible in $k$, one checks at once that $\mathrm{A}$ is the integral closure of $k[\mathrm{S}]$ in $\mathrm{K}$; thus $\mathrm{A}$ is integrally closed, and $\mathrm{X}$ is regular. Moreover, if an element $f+tg$ of $\mathrm{K}$ (with $f,g\in k(\mathrm{S})$) is algebraic over $k$, then so are its norm and trace over $k(\mathrm{S})$; since $k(\mathrm{S})$ is a purely transcendental extension of $k$, it follows readily that $f\in k$ and $g=0$. Thus $k$ is algebraically closed in $\mathrm{K}$, and \emph{a fortiori} $\mathrm{K}$ is a primary extension

\oldpage[IV-2]{148}

of $k$ \sref{IV.4.3.1}. Nevertheless, if $k'=k(a^{1/p})$, then $\mathrm{X}'=\mathrm{X}\otimes_k k'$ is not normal: in $k'[\mathrm{S}]$ one has
\[
 \mathrm{S}^p-a=(\mathrm{S}-a^{1/p})^p,
\]
and $\mathrm{X}'$ is therefore isomorphic to $\operatorname{Spec}(\mathrm{A}')$, where $\mathrm{A}'=k'[\mathrm{S}][t']$ and $t'$ is a root of $\mathrm{T}^2-\mathrm{S}^p$. The ring $\mathrm{A}'$ is not integrally closed: the element $t''=t'/\mathrm{S}$ of its field of fractions $\mathrm{K}'=k'(\mathrm{S})[t']$ satisfies the integral-dependence equation
\[
 (t'')^2-\mathrm{S}^{p-2}=0
\]
over $\mathrm{A}'$, but does not belong to $\mathrm{A}'$. In classical terminology, the ``genus'' over $k'$ of the field $\mathrm{K}'$ of rational functions of $\mathrm{X}'$ is \emph{strictly smaller} than that of $\mathrm{K}$ over $k$.

(ii) As recalled in (i), it follows from \sref{IV.4.6.1} that when $\mathrm{X}$ is not a geometrically reduced algebraic prescheme over $k$, there exist finite radicial extensions $k'$ of $k$ such that $\mathrm{X}'=\mathrm{X}\otimes_k k'$ is not reduced. It is interesting to give an example where $\mathrm{X}$ is a regular scheme over $k$, such that $k$ is \emph{algebraically closed} in the field $\mathrm{K}$ of rational functions of $\mathrm{X}$. Let $k$ be a field of characteristic $p>0$, in which there exist two elements $a$, $b$ forming a $p$-free family over $k^p$. Denoting again by $\mathrm{B}$ the ring $k[\mathrm{S},\mathrm{T}]$, consider this time the polynomial $\mathrm{P}(\mathrm{S},\mathrm{T})=\mathrm{T}^p-a\mathrm{S}^p-b$; as $a\mathrm{S}^p+b$ is not a $p$-th power in $k(\mathrm{S})$, $\mathrm{P}$ is irreducible as a polynomial of $k(\mathrm{S})[\mathrm{T}]$, and the scheme $\mathrm{X}_0=\operatorname{Spec}(\mathrm{B}/\mathrm{PB})$ is therefore an integral affine curve over $k$, whose field of rational functions is $\mathrm{K}=k(\mathrm{S})[t]$, where $t$ is a root of $\mathrm{P}$; let us show that $k$ is \emph{algebraically closed} in $\mathrm{K}$. Suppose indeed that $\mathrm{K}$ contains an element $z$ algebraic over $k$ and not in $k$; one would then also have $z\notin k(\mathrm{S})$, so $[k(\mathrm{S})[z]:k(\mathrm{S})]>1$; as $[\mathrm{K}:k(\mathrm{S})]=p$, and as $\mathrm{K}$ is radicial over $k(\mathrm{S})$, one would have $z^p\in k(\mathrm{S})$, so $c=z^p\in k$ since $z$ is algebraic over $k$; but then one would have $t^p=a\mathrm{S}^p+b\in k^p(\mathrm{S}^p)(c)$ hence $a$ and $b$ would belong to $k^p(c)$, which is absurd. So $\mathrm{X}$ is the normalisation of the curve $\mathrm{X}_0$ in $\mathrm{K}$, which is therefore a normal curve (and hence \emph{regular}) over $k$. If $k'=k(a^{1/p},b^{1/p})$, it is clear that $a\mathrm{S}^p+b$ is a $p$-th power in $k'(\mathrm{S})$, hence $\mathrm{K}\otimes_k k'$ is \emph{not reduced}, nor \emph{a fortiori} the scheme $\mathrm{X}\otimes_k k'$.
\end{remarks}

\begin{definition}[6.7.6]
\label{IV.6.7.6}
Let $k$ be a field, $\mathrm{X}$ a $k$-prescheme locally Noetherian. We say that $\mathrm{X}$ is \emph{geometrically regular} at a point $x$ (resp.\ has the property $(\mathrm{R}_n)$ \emph{geometric}) at a point $x$, resp.\ is geometrically normal at a point $x$) if, \emph{for every} finite extension $k'$ of $k$, $\mathrm{X}'=\mathrm{X}_{(k')}=\mathrm{X}\otimes_k k'$ is regular (resp.\ has the property $(\mathrm{R}_n)$, resp.\ is normal) at every point $x'$ of $\mathrm{X}'$ whose projection in $\mathrm{X}$ is $x$. We say that $\mathrm{X}$ is \emph{geometrically regular} (resp.\ has the property $(\mathrm{R}_n)$ \emph{geometric}, or equivalently is geometrically regular in codimension $\leqslant n$), resp.\ is \emph{geometrically normal}) if $\mathrm{X}$ is geometrically regular (resp.\ has the property $(\mathrm{R}_n)$ geometric, resp.\ is geometrically normal) at every point.

We say that an algebra $\mathrm{A}$ over $k$ is a \emph{geometrically regular ring} (resp.\ geometrically normal, resp.\ geometrically reduced, resp.\ having the property $(\mathrm{R}_n)$ geometric) if $\operatorname{Spec}(\mathrm{A})$ has the same property.

If $\mathrm{X}=\operatorname{Spec}(\mathrm{K})$, where $\mathrm{K}$ is an extension of $k$, then saying that $\mathrm{X}$ is geometrically regular, or geometrically normal, or \emph{geometrically reduced}, amounts to saying that $\mathrm{K}$ is a \emph{separable} extension of $k$: this follows from \sref{IV.4.6.1} and the fact that if $\mathrm{K}$ is a separable extension of $k$ and $k'$ a finite extension of $k$, $\mathrm{K}\otimes_k k'$ is composed of a direct product of a finite number of fields (Bourbaki, \emph{Alg.}, chap.~VIII, \S~7, no~3, cor.~1 of th.~1).
\end{definition}

\oldpage[IV-2]{149}

\begin{proposition}[6.7.7]
\label{IV.6.7.7}
Let $k$ be a field, $\mathrm{X}$ a $k$-prescheme locally Noetherian, $x$ a point of $\mathrm{X}$; denote by $Q(k')$ the relation ``$k'$ is an extension of $k$, and $\boldsymbol{P}(\mathrm{X}_{(k')},x')$ is true for every point $x'$ of $\mathrm{X}_{(k')}$ whose projection in $\mathrm{X}$ is $x$'', where $\boldsymbol{P}(\mathrm{Z},z)$ is one of the properties \emph{a)}, \emph{b)}, \emph{c)} of the statement \sref{IV.6.7.4}. Then the following properties are equivalent:
\begin{enumerate}[label=\emph{\alph*)}]
  \item $Q(k')$ is true for every finite extension $k'$ of $k$.
  \item $Q(k')$ is true for every finite radicial extension $k'$ of $k$.
  \item $Q(k')$ is true for every extension of finite type $k'$ of $k$.
\end{enumerate}
Suppose in addition that $\mathrm{X}$ is locally of finite type over $k$; then the three preceding properties are also equivalent to the following:
\begin{enumerate}[label=\emph{\alph*)}]
\setcounter{enumi}{3}
  \item $Q(k')$ is true for every extension $k'$ of $k$.
  \item $Q(k')$ is true for a perfect extension $k'$ of $k$.
  \item $Q(k')$ is true for every extension $k'=k^{p^{-s}}$, where $p$ is the characteristic exponent of $k$, and $s>0$.
\end{enumerate}
\end{proposition}

\begin{proof}
To prove the equivalence of $a)$, $b)$ and $c)$, it clearly suffices to show that $b)$ implies $c)$. Let $\mathrm{K}$ be an extension of finite type of $k$, which is then the field of fractions of a $k$-algebra of finite type $\mathrm{A}$; set $\mathrm{Y}=\operatorname{Spec}(\mathrm{A})$. We know \sref{IV.4.6.6} that there exists a finite and radicial extension $k'$ of $k$ such that $\mathrm{Y}'=(\mathrm{Y}\otimes_k k')_{\mathrm{red}}$ is separable over $k'$; if $\eta'$ is the generic point of $\mathrm{Y}'$, $\mathrm{K}'=\boldsymbol{k}(\eta')$ is a separable extension of $k'$ by \sref{IV.4.6.1}. This being so, $\boldsymbol{P}(\mathrm{X}_{(k')},x')$ is true by hypothesis for every point $x'$ of $\mathrm{X}_{(k')}$ above $x$, so it follows from \sref{IV.6.7.4} that $\boldsymbol{P}(\mathrm{X}_{(\mathrm{K}')},x'')$ is true for every point $x''$ of $\mathrm{X}_{(\mathrm{K}')}$ above $x$, since $\mathrm{X}_{(\mathrm{K}')}=(\mathrm{X}_{(k')})_{(\mathrm{K}')}$, $x''$ is above a point $x'$ of $\mathrm{X}_{(k')}$, itself above $x$, and $\mathrm{K}'$ is separable over $k'$. But we also have $\mathrm{X}_{(\mathrm{K}')}=(\mathrm{X}_{(\mathrm{K})})_{(\mathrm{K}')}$ and for every $x_0\in\mathrm{X}_{(\mathrm{K})}$ above $x$, there exists $x''\in\mathrm{X}_{(\mathrm{K}')}$ above $x_0$; it therefore follows from \sref{IV.6.7.4} that $\boldsymbol{P}(\mathrm{X}_{(\mathrm{K})},x_0)$ is true.

Suppose now that $\mathrm{X}$ is locally of finite type over $k$, so that for \emph{every} extension $\mathrm{K}$ of $k$, $\mathrm{X}_{(\mathrm{K})}$ is locally Noetherian. As a radicial extension of $k$ is isomorphic to a sub-extension of a perfect extension of $k$, it follows immediately from \sref{IV.6.7.4} that $e)$ implies $f)$; also, every finite radicial extension of $k$ is contained in an extension $k^{p^{-s}}\subset k^{p^{-\infty}}$, so $f)$ implies $b)$; $d)$ trivially implies $e)$, and finally $e)$ implies $d)$. Indeed, if $\mathrm{K}$ is any extension \emph{whatever} of $k$, there exists an extension $\mathrm{K}'$ of $k$ containing (up to $k$-isomorphisms) $k'$ and $\mathrm{K}$; since $\mathrm{K}'$ is a separable extension of $k'$, one deduces from \sref{IV.6.7.4} that $Q(k')$ is true, then that $Q(\mathrm{K})$ is true, reasoning as in the first part of the proof. It therefore suffices to prove that $b)$ \emph{implies} $e)$. We shall take $k'=k^{p^{-\infty}}$ in $e)$.

The question being local on $\mathrm{X}$, one can furthermore limit oneself to the case where $\mathrm{X}$ is affine and of finite type over $k$, by replacing $\mathrm{X}$ by a neighbourhood of $x$; let then $\mathrm{X}=\operatorname{Spec}(\mathrm{B})$, where $\mathrm{B}$ is a $k$-algebra of finite type; furthermore, by definition of the property $\boldsymbol{P}$ in the cases considered, one can also replace $\mathrm{X}$ by the local scheme of $\mathrm{X}$ at the point $x$, hence suppose that $\mathrm{B}$ is a \emph{local} Noetherian ring. Set $\mathrm{B}'=\mathrm{B}\otimes_k k'$; $k'$ is the inductive limit of its finite sub-extensions $k'_\lambda$ and if one sets $\mathrm{B}'_\lambda=\mathrm{B}\otimes_k k'_\lambda$, the morphism $f_\lambda:\operatorname{Spec}(\mathrm{B}')\to\operatorname{Spec}(\mathrm{B}'_\lambda)$ is a \emph{homeomorphism} of $\operatorname{Spec}(\mathrm{B}')$ onto $\operatorname{Spec}(\mathrm{B}'_\lambda)$ for every $\lambda$;

\oldpage[IV-2]{150}

indeed $f_\lambda$ is bijective by virtue of \sref[I]{I.3.5.2}, \sref{I.3.5.7} and \sref{I.3.5.8}; on the other hand, it is closed by \sref[II]{II.6.1.10}. One can now apply \sref{IV.5.13.6} which proves (by virtue of the hypothesis $b)$) that $b)$ implies $e)$ when $\boldsymbol{P}(\mathrm{Z},z)$ is the property $(\mathrm{R}_n)$ at the point $z$. The case where $\boldsymbol{P}$ is the property of being regular (i.e.\ of having the property $(\mathrm{R}_n)$ for every $n$) follows trivially. Finally, taking into account Serre's normality criterion \sref{IV.5.8.6}, $b)$ implies $e)$ further when $\boldsymbol{P}(\mathrm{Z},z)$ is the property of being normal at the point $z$, since this is equivalent to saying that $\mathrm{Z}$ has at the point $z$ the properties $(\mathrm{S}_2)$ and $(\mathrm{R}_1)$: one applies what precedes for $n=1$, and \sref{IV.6.7.2},~(ii) for $n=2$.
\end{proof}

\begin{corollary}[6.7.8]
\label{IV.6.7.8}
Let $k$ be a field; for a $k$-prescheme locally Noetherian $\mathrm{X}$ and a point $x\in\mathrm{X}$, denote by $\boldsymbol{P}(\mathrm{X},x)$ any of the following properties:
\begin{enumerate}[label=(\roman*)]
  \item $\operatorname{coprof}(\sh{O}_x)\leqslant n$.
  \item $\sh{O}_x$ is a Cohen-Macaulay ring.
  \item $\mathrm{X}$ satisfies the property $(\mathrm{S}_n)$ at the point $x$.
  \item $\mathrm{X}$ is geometrically regular at the point $x$.
  \item $\mathrm{X}$ has the property $(\mathrm{R}_n)$ geometric at the point $x$.
  \item $\mathrm{X}$ is geometrically normal at the point $x$.
  \item $\mathrm{X}$ is geometrically reduced (i.e.\ separable) at the point $x$.
\end{enumerate}
Let $k'$ be an extension of $k$; one supposes either that $k'$ is of finite type, or that $\mathrm{X}$ is locally of finite type over $k$, so that $\mathrm{X}'=\mathrm{X}\otimes_k k'$ is locally Noetherian. Let $x'$ be a point of $\mathrm{X}'$ whose projection in $\mathrm{X}$ is $x$. Then, for $\boldsymbol{P}(\mathrm{X},x)$ to be true, it is necessary and sufficient that $\boldsymbol{P}(\mathrm{X}',x')$ be so.
\end{corollary}

\begin{proof}
This has already been seen for property (vii) \sref{IV.4.6.11}, and for properties (i), (ii), and (iii) \sref{IV.6.7.1}. For (iv), (v), and (vi), this follows from \sref{IV.6.7.7}: the necessity follows from the equivalence of criteria $a)$ and $c)$, and from the equivalence of $c)$ and $d)$ when $\mathrm{X}$ is locally of finite type over $k$. To prove sufficiency, let $k''$ be a finite radicial extension of $k$; one may regard $k'$ and $k''$ as subextensions of an extension $\mathrm{K}$ of $k$. Set $\mathrm{X}''=\mathrm{X}\otimes_k k''$ and $\mathrm{X}_0=\mathrm{X}\otimes_k\mathrm{K}$, and note that, since $k''$ is radicial over $k$, there is only \emph{one} point $x''$ of $\mathrm{X}''$ above $x$ \sref[I]{I.3.5.7} and \sref{I.3.5.8}. Let $x_0$ be any point of $\mathrm{X}_0$ above $x'$. If $\boldsymbol{P}(\mathrm{X}',x')$ is true, then so is $\boldsymbol{P}(\mathrm{X}_0,x_0)$ by \sref{IV.6.7.7},~$c)$ and $d)$ (if $k'$ is of finite type over $k$, one may suppose the same of $\mathrm{K}$; if $\mathrm{X}$ is locally of finite type over $k$, then $\mathrm{X}'$ is locally of finite type over $k'$). It follows from \sref{IV.6.7.4} that $Q(k'')$ is true (with the notation of \sref{IV.6.7.7}), and hence $\boldsymbol{P}(\mathrm{X},x)$ is true by \sref{IV.6.7.7},~$b)$.
\end{proof}

\subsection{Regular, normal, reduced, smooth morphisms}
\label{subsection:IV.6.8}

\begin{definition}[6.8.1]
\label{IV.6.8.1}
Let $\mathrm{X}$, $\mathrm{Y}$ be two preschemes, $f:\mathrm{X}\to\mathrm{Y}$ a morphism such that the fibre $f^{-1}(y)$ is a locally Noetherian prescheme for every $y\in\mathrm{Y}$, $x$ a point of $\mathrm{X}$. We say that $f$ is respectively a morphism:
\begin{enumerate}[label=(\roman*)]
  \item of coprofundity $\leqslant n$ at the point $x$;
  \item of Cohen-Macaulay at the point $x$;
\oldpage[IV-2]{151}
  \item $(\mathrm{S}_n)$ at the point $x$;
  \item regular at the point $x$;
  \item $(\mathrm{R}_n)$ at the point $x$;
  \item normal at the point $x$;
  \item reduced at the point $x$;
\end{enumerate}
if $f$ is flat at the point $x$, and if in addition the corresponding property $\boldsymbol{P}(f^{-1}(f(x)),x)$ \textnormal{(}notation of \sref{IV.6.7.8}\textnormal{)} is true.

We say that $f$ is \emph{smooth} at the point $x$ if it is locally of finite presentation in a neighbourhood of $x$ and if it is regular at the point $x$. We say respectively that $f$ is a morphism: of coprofundity $\leqslant n$, of Cohen-Macaulay, $(\mathrm{S}_n)$, regular, $(\mathrm{R}_n)$, normal, reduced, smooth, if it has the corresponding property at every point of $\mathrm{X}$.
\end{definition}

\begin{proposition}[6.8.2]
\label{IV.6.8.2}
Let $\mathrm{X}$, $\mathrm{Y}$ be two locally Noetherian preschemes, $f:\mathrm{X}\to\mathrm{Y}$ a morphism, $x$ a point of $\mathrm{X}$. Denote by $\boldsymbol{M}(f,x)$ one of the properties (i) to (vii) of the definition \sref{IV.6.8.1}, or the property of being smooth at the point $x$. Let $\mathrm{Y}'$ be a locally Noetherian prescheme, $g:\mathrm{Y}'\to\mathrm{Y}$ a morphism, $\mathrm{X}'=\mathrm{X}\times_\mathrm{Y}\mathrm{Y}'$, $f'=f_{(\mathrm{Y}')}:\mathrm{X}'\to\mathrm{Y}'$. Suppose that $f$ or $g$ is locally of finite type. Then, for every $x'\in\mathrm{X}'$ above $x$, the property $\boldsymbol{M}(f,x)$ implies $\boldsymbol{M}(f',x')$.
\end{proposition}

\begin{proof}
Set $y=f(x)$, $y'=f'(x')$; by transitivity of fibres \sref[I]{I.3.6.4}, we have $f'^{-1}(y')=f^{-1}(y)\otimes_{\boldsymbol{k}(y)}\boldsymbol{k}(y')$. Since either $f^{-1}(y)$ is locally of finite type over $\boldsymbol{k}(y)$, or $\boldsymbol{k}(y')$ is of finite type over $\boldsymbol{k}(y)$ \sref[I]{I.6.4.11}, it follows from \sref{IV.6.7.8} that the properties $\boldsymbol{P}(f^{-1}(y),x)$ and $\boldsymbol{P}(f'^{-1}(y'),x')$ are equivalent; furthermore, if $f$ is flat at the point $x$, $f'$ is flat at the point $x'$ \sref{IV.2.1.4}, which proves the proposition, taking into account \sref{IV.1.4.3},~(iii).
\end{proof}

\begin{proposition}[6.8.3]
\label{IV.6.8.3}
For a morphism $f$ of locally Noetherian preschemes, let $\boldsymbol{M}(f)$ be one of the following properties: being Cohen-Macaulay, $(\mathrm{S}_n)$, regular, $(\mathrm{R}_n)$, normal, reduced.
\begin{enumerate}[label=(\roman*)]
  \item Let $\mathrm{X}$, $\mathrm{Y}$, $\mathrm{Z}$ be three locally Noetherian preschemes, $f:\mathrm{X}\to\mathrm{Y}$, $g:\mathrm{Y}\to\mathrm{Z}$ two morphisms. If $\boldsymbol{M}(f)$ and $\boldsymbol{M}(g)$ are true, $\boldsymbol{M}(g\circ f)$ is true.
  \item Conversely, if $f$ is surjective and if $\boldsymbol{M}(f)$ and $\boldsymbol{M}(g\circ f)$ are true, then $\boldsymbol{M}(g)$ is true.
  \item Let $\mathrm{X}$, $\mathrm{Y}$, $\mathrm{Y}'$ be three locally Noetherian preschemes, $f:\mathrm{X}\to\mathrm{Y}$, $h:\mathrm{Y}'\to\mathrm{Y}$ two morphisms; set $\mathrm{X}'=\mathrm{X}\times_\mathrm{Y}\mathrm{Y}'$, $f'=f_{(\mathrm{Y}')}:\mathrm{X}'\to\mathrm{Y}'$. Suppose that $f$ or $h$ is locally of finite type. If $\boldsymbol{M}(f)$ is true, then so is $\boldsymbol{M}(f')$; the converse is true when $h$ is faithfully flat.
\end{enumerate}
The conclusions of (i) and (iii) are also true when $\boldsymbol{M}$ is the property of being smooth and (in (iii)) $h$ is quasi-compact.
\end{proposition}

\begin{proof}
(i) We already know that if $f$ and $g$ are flat, so is $u=g\circ f$, and that if $f$ and $g\circ f$ are flat and $f$ is surjective, then $g$ is flat \sref{IV.2.1.6} and \sref{IV.2.2.13}. On the other hand, for every $z\in\mathrm{Z}$, the morphism $f_z=f\otimes\mathbf{1}_{\boldsymbol{k}(z)}:u^{-1}(z)\to g^{-1}(z)$ is flat (resp.\ faithfully flat if $f$ is) and for every $y\in g^{-1}(z)$, $f_z^{-1}(y)$ is isomorphic to $f^{-1}(y)$ \sref[I]{I.3.6.4}. If $\boldsymbol{M}(f)$ and $\boldsymbol{M}(g)$ are true, $\boldsymbol{M}(g\circ f)$ is then true for the cases where $\boldsymbol{M}$ is the property of being Cohen-Macaulay or $(\mathrm{S}_n)$, by virtue of \sref{IV.6.6.1}. On the other hand, let $k'$ be a finite extension of $\boldsymbol{k}(z)$; setting $\mathrm{Y}'_z=g^{-1}(z)\otimes_{\boldsymbol{k}(z)}k'$, $\mathrm{X}'_z=u^{-1}(z)\otimes_{\boldsymbol{k}(z)}k'$, and $f'_z=f_z\otimes\mathbf{1}_{k'}:\mathrm{X}'_z\to\mathrm{Y}'_z$; the morphism $f'_z$

\oldpage[IV-2]{152}

is flat (resp.\ faithfully flat) and for every $y'\in\mathrm{Y}'_z$, the fibre $f'^{-1}_z(y')$ is isomorphic to $f^{-1}(y)\otimes_{\boldsymbol{k}(y)}\boldsymbol{k}(y')$, denoting by $y$ the image of $y'$ in $g^{-1}(z)$ \sref[I]{I.3.6.4}. When $\boldsymbol{M}$ is the property of being regular, $(\mathrm{R}_n)$, normal or reduced, the hypothesis that $\boldsymbol{M}(f)$ and $\boldsymbol{M}(g)$ are true implies that $\mathrm{Y}'_z$ and each of the fibres $f'^{-1}_z(y')$ for $y'\in\mathrm{Y}'_z$ have the corresponding property among the properties $c)$, $d)$, $e)$, $f)$ of \sref{IV.6.6.1}; one deduces from \sref{IV.6.6.1},~(i) that $\mathrm{X}'_z$ has the same property, hence $\boldsymbol{M}(g\circ f)$ is true.

(ii) Conversely, the hypothesis that $\boldsymbol{M}(g\circ f)$ and $\boldsymbol{M}(f)$ are true and that $f$ is surjective implies that $\mathrm{Y}'_z$ has the corresponding property by virtue of \sref{IV.6.6.1},~(ii), since $f'_z$ is surjective for every $z$; hence $\boldsymbol{M}(g)$ is true.

(iii) The first assertion follows immediately from \sref{IV.6.8.2}. On the other hand, if $h$ is faithfully flat and $f'$ flat, $f$ is flat \sref{IV.2.4.1}; since, with the notations of \sref{IV.6.8.2}, the properties $\boldsymbol{P}(f^{-1}(y),x)$ and $\boldsymbol{P}(f'^{-1}(y'),x')$ are equivalent \sref{IV.6.7.8}, one sees that $\boldsymbol{M}(f)$ and $\boldsymbol{M}(f')$ are then equivalent.

Finally, the last assertion of the proposition follows from \sref{IV.1.4.3},~(iii) and \sref{IV.2.7.1},~(iv).
\end{proof}

\begin{remarks}[6.8.4]
\label{IV.6.8.4}
(i) If $f$ is faithfully flat, $g$ flat and if $g\circ f$ is of coprofundity $\leqslant n$, then $g$ is of coprofundity $\leqslant n$, as follows from the reasoning of \sref{IV.6.6.2}.

(ii) When, in \sref{IV.6.8.1}, one takes for $\mathrm{Y}$ the spectrum of a field $k$, the notions (iv), (v) and (vi) reduce to those defined in \sref{IV.6.7.5}. It is clear that these latter are relative to the base field $k$. The definition \sref{IV.6.8.1} thus leads to saying that a prescheme $\mathrm{X}$ is ``regular (resp.\ $(\mathrm{R}_n)$, resp.\ normal) over $k$'' instead of saying that it is ``geometrically regular (resp.\ $(\mathrm{R}_n)$, resp.\ normal) relative to $k$''; one will avoid confusing this notion with the property of being regular (resp.\ $(\mathrm{R}_n)$, resp.\ normal) which is independent of $k$. One can make in this connection the same remarks as in \sref{IV.4.5.12}.
\end{remarks}

\begin{proposition}[6.8.5]
\label{IV.6.8.5}
Let $k$ be a field, $\mathrm{X}$, $\mathrm{Y}$ two $k$-preschemes locally Noetherian, at least one of which is locally of finite type over $k$. For a $k$-prescheme $\mathrm{Z}$, denote by $\boldsymbol{P}(\mathrm{Z})$ one of the properties $c)$ to $f)$ of \sref{IV.6.6.1}; the property ``$\boldsymbol{P}(\mathrm{Z})$ geometric'' is then defined in \sref{IV.6.7.6} (resp.\ \sref{IV.4.6.1}). Then:
\begin{enumerate}[label=(\roman*)]
  \item If $\mathrm{X}$ has the property $\boldsymbol{P}$, and $\mathrm{Y}$ the property $\boldsymbol{P}$ geometric, $\mathrm{X}\times_k\mathrm{Y}$ has the property $\boldsymbol{P}$.
  \item If $\mathrm{X}$ and $\mathrm{Y}$ have the property $\boldsymbol{P}$ geometric, so does $\mathrm{X}\times_k\mathrm{Y}$.
\end{enumerate}
\end{proposition}

\begin{proof}
Indeed, let $f:\mathrm{X}\times_k\mathrm{Y}\to\mathrm{X}$, $g:\mathrm{X}\to\operatorname{Spec}(k)$ be the structural morphisms, which are flat by \sref{IV.2.1.2} and \sref{IV.2.1.4}. The hypothesis that $\mathrm{Y}$ has the property $\boldsymbol{P}$ geometric implies by virtue of \sref{IV.6.8.2} that $\boldsymbol{M}(f)$ is true, $\boldsymbol{M}$ being the property \sref{IV.6.8.3} which corresponds to $\boldsymbol{P}$; under the hypothesis (ii), $\boldsymbol{M}(g)$ is also true, hence the assertion (ii) follows from \sref{IV.6.8.3},~(i). As for the assertion (i), it follows directly from \sref{IV.6.6.1},~(i).
\end{proof}

\begin{theorem}[6.8.6]
\label{IV.6.8.6}
Let $\mathrm{X}$, $\mathrm{Y}$ be two locally Noetherian preschemes, $f:\mathrm{X}\to\mathrm{Y}$ a morphism locally of finite type, $x$ a point of $\mathrm{X}$, $y=f(x)$. The following properties are equivalent:
\begin{enumerate}[label=\emph{\alph*)}]
  \item $f$ is smooth at the point $x$.
  \item $f$ is regular at the point $x$.
\oldpage[IV-2]{153}
  \item $\sh{O}_x$ is a formally smooth $\sh{O}_y$-algebra \sref[0]{0.19.3.1} for the adic topologies on $\sh{O}_x$ and $\sh{O}_y$.
  \item[\emph{c$'$)}] $\sh{O}_x$ is a formally smooth $\sh{O}_y$-algebra \sref[0]{0.19.3.1} for the discrete topologies on $\sh{O}_x$ and $\sh{O}_y$.
\end{enumerate}
\end{theorem}

\begin{proof}
The equivalence of $a)$ and $b)$ follows from the definitions \sref{IV.6.8.1} and from the fact that for locally Noetherian preschemes, morphisms locally of finite type are locally of finite presentation \sref{IV.1.4.2}.

Secondly, for $\sh{O}_x$ to be a formally smooth $\sh{O}_y$-algebra for the adic topologies, it is necessary and sufficient that $\sh{O}_x$ be flat over $\sh{O}_y$ and that $\sh{O}_x\otimes_{\sh{O}_y}\boldsymbol{k}(y)$ be formally smooth over $\boldsymbol{k}(y)$ for its adic topology \sref[0]{0.19.7.1}. But for $\sh{O}_x\otimes_{\sh{O}_y}\boldsymbol{k}(y)$ to be a formally smooth $\boldsymbol{k}(y)$-algebra for its adic topology, it is necessary and sufficient that it be a geometrically regular $\boldsymbol{k}(y)$-algebra \sref[0]{0.19.6.6}; this proves the equivalence of $b)$ and $c)$. Finally, to prove the equivalence of $c)$ and $c')$, one may reduce to the affine case $\mathrm{Y}=\operatorname{Spec}(\mathrm{A})$, $\mathrm{X}=\operatorname{Spec}(\mathrm{C})$, where $\mathrm{A}$ is Noetherian and $\mathrm{C}$ is a finite-type Noetherian $\mathrm{A}$-algebra, say $\mathrm{C}=\mathrm{A}[\mathrm{T}_1,\ldots,\mathrm{T}_n]/\mathfrak{J}$. Since $\mathfrak{J}/\mathfrak{J}^2$ is a finitely presented $\mathrm{C}$-module, the equivalence follows from \sref[0]{0.22.6.4} applied to $\mathrm{A}$, $\mathrm{B}=\mathrm{A}[\mathrm{T}_1,\ldots,\mathrm{T}_n]$ and $\mathrm{C}=\mathrm{B}/\mathfrak{J}$.
\end{proof}

\begin{corollary}[6.8.7]
\label{IV.6.8.7}
Let $\mathrm{X}$, $\mathrm{Y}$ be two locally Noetherian preschemes, $f:\mathrm{X}\to\mathrm{Y}$ a morphism locally of finite type. Then the set of points $x\in\mathrm{X}$ where $f$ is smooth \textnormal{(}or regular\textnormal{)} is open in $\mathrm{X}$.
\end{corollary}

\begin{proof}
Indeed it follows from \sref[0]{0.22.6.5} that the set of $x\in\mathrm{X}$ satisfying the condition $c')$ of \sref{IV.6.8.6} is open in $\mathrm{X}$, and we conclude by \sref{IV.6.8.6}.
\end{proof}

\begin{remark}[6.8.8]
\label{IV.6.8.8}
In \sref{IV.17.5.1}, we shall show that the equivalence of $b)$ and $c')$ in \sref{IV.6.8.6}, as well as the corollary \sref{IV.6.8.7}, remain valid without the Noetherian hypothesis on $\mathrm{X}$ and $\mathrm{Y}$, provided that one restricts to morphisms \emph{locally of finite presentation}.
\end{remark}

\subsection{The theorem of generic flatness}
\label{subsection:IV.6.9}

\begin{theorem}[6.9.1]
\label{IV.6.9.1}
Let $\mathrm{Y}$ be a locally Noetherian and integral prescheme, $u:\mathrm{X}\to\mathrm{Y}$ a morphism of finite type, $\sh{F}$ a coherent $\sh{O}_\mathrm{X}$-Module. There then exists a non-empty open subset $\mathrm{U}$ of $\mathrm{Y}$ such that $\sh{F}|u^{-1}(\mathrm{U})$ is flat over $\mathrm{U}$.
\end{theorem}

\begin{proof}
One can clearly limit oneself to the case where $\mathrm{Y}=\operatorname{Spec}(\mathrm{A})$ is affine; then $\mathrm{X}$ is a finite union of affine opens $\mathrm{X}_i$ of finite type over $\mathrm{Y}$; if, for each $i$, there exists a non-empty open $\mathrm{U}_i$ in $\mathrm{Y}$ such that $\sh{F}|(\mathrm{X}_i\cap u^{-1}(\mathrm{U}_i))$ is flat over $\mathrm{U}_i$, it is clear that taking for $\mathrm{U}$ the intersection of the $\mathrm{U}_i$, we shall have answered the question; one can therefore suppose that $\mathrm{X}=\operatorname{Spec}(\mathrm{B})$, where $\mathrm{B}$ is an $\mathrm{A}$-algebra of finite type. We then have $\sh{F}=\widetilde{\mathrm{M}}$, where $\mathrm{M}$ is a $\mathrm{B}$-module of finite type; the theorem will follow from the following:
\end{proof}

\begin{lemma}[6.9.2]
\label{IV.6.9.2}
Let $\mathrm{A}$ be an integral Noetherian ring, $\mathrm{B}$ a finite-type $\mathrm{A}$-algebra, and $\mathrm{M}$ a finite $\mathrm{B}$-module. Then there exists $f\neq 0$ in $\mathrm{A}$ such that $\mathrm{M}_f$ is a free $\mathrm{A}_f$-module.
\end{lemma}

\begin{proof}
\oldpage[IV-2]{154}
Denote by $\mathrm{K}$ the field of fractions of $\mathrm{A}$; then $\mathrm{B}\otimes_\mathrm{A}\mathrm{K}$ is an algebra of finite type over $\mathrm{K}$ and $\mathrm{M}\otimes_\mathrm{A}\mathrm{K}$ is a $(\mathrm{B}\otimes_\mathrm{A}\mathrm{K})$-module of finite type. We shall reason by induction on the \emph{dimension} $n$ of the support of $\mathrm{M}\otimes_\mathrm{A}\mathrm{K}$, which is either $-\infty$ or finite and $\geqslant 0$. Suppose first $n=-\infty$, i.e.\ $\mathrm{M}\otimes_\mathrm{A}\mathrm{K}=0$; if $(m_i)_{1\leqslant i\leqslant r}$ is a system of generators of the $\mathrm{B}$-module $\mathrm{M}$, there therefore exists $f\neq 0$ in $\mathrm{A}$ such that $fm_i=0$ for $1\leqslant i\leqslant r$; hence $\mathrm{M}_f=0$ and the lemma is true in this case.

Suppose now $n\geqslant 0$. We know that there exists a composition series $\mathrm{M}=\mathrm{M}_1\supset\mathrm{M}_2\supset\ldots\supset\mathrm{M}_q=0$ of the $\mathrm{B}$-module $\mathrm{M}$ such that each quotient $\mathrm{N}_i=\mathrm{M}_i/\mathrm{M}_{i+1}$ is isomorphic to a $\mathrm{B}$-module of the form $\mathrm{B}/\mathfrak{p}_i$, where $\mathfrak{p}_i$ is a prime ideal of $\mathrm{B}$ (Bourbaki, \emph{Alg.\ comm.}, chap.~IV, \S~1, no~4, th.~1). If the theorem is true for each $\mathrm{N}_i$, there exists for each $i$ an $f_i\neq 0$ in $\mathrm{A}$ such that $(\mathrm{N}_i)_{f_i}$ is free over $\mathrm{A}_{f_i}$; setting $f=f_1f_2\cdots f_{q-1}$, it follows that $(\mathrm{N}_i)_f$ is a free $\mathrm{A}_f$-module for $1\leqslant i\leqslant q-1$. But $(\mathrm{N}_i)_f=(\mathrm{M}_i)_f/(\mathrm{M}_{i+1})_f$ \sref[0]{0.1.3.2}, and since an extension of free modules is free, we deduce that $\mathrm{M}_f$ is a free $\mathrm{A}_f$-module.

Replacing $\mathrm{B}$ by $\mathrm{B}/\mathfrak{p}$, where $\mathfrak{p}$ is a prime ideal of $\mathrm{B}$, one can therefore reduce to the case $\mathrm{M}=\mathrm{B}$ with $\mathrm{B}$ integral. We then know (Bourbaki, \emph{Alg.\ comm.}, chap.~V, \S~3, no~1, cor.~1 of th.~1) that there exist $g\neq 0$ in $\mathrm{A}$ and elements $t_i$ ($1\leqslant i\leqslant m$) of $\mathrm{B}$, algebraically independent over $\mathrm{A}$, such that $\mathrm{B}_g$ is integral over $\mathrm{A}_g[t_1,\ldots,t_m]$. One can replace $\mathrm{A}$ by $\mathrm{A}_g$, $\mathrm{B}$ by $\mathrm{B}_g$ and suppose that $\mathrm{B}$ is integral over $\mathrm{C}=\mathrm{A}[t_1,\ldots,t_m]$, hence a $\mathrm{C}$-module of finite type and \emph{without torsion}. We know furthermore \sref{IV.4.1.2} that the dimension of $\operatorname{Spec}(\mathrm{B}\otimes_\mathrm{A}\mathrm{K})$ is equal to $m$, hence $m=n$.

This being so, if $h$ is the rank of the $\mathrm{C}$-module without torsion $\mathrm{B}$, there exists a short exact sequence of $\mathrm{C}$-modules
\[
0\to\mathrm{C}^h\to\mathrm{B}\to\mathrm{M}'\to 0
\]
where $\mathrm{M}'$ is a torsion $\mathrm{C}$-module of finite type; the support of $\mathrm{M}'$ does not therefore contain the generic point of $\operatorname{Spec}(\mathrm{C})$ \sref[I]{I.7.4.6} and hence the support of $\mathrm{M}'\otimes_\mathrm{A}\mathrm{K}$ does not contain the generic point of $\operatorname{Spec}(\mathrm{C}\otimes_\mathrm{A}\mathrm{K})$ \sref[I]{I.9.1.13.1}; one concludes \sref{IV.4.1.2.1} that its dimension is $<n$. By the induction hypothesis, there exists $f\neq 0$ in $\mathrm{A}$ such that $\mathrm{M}'_f$ is a free $\mathrm{A}_f$-module; moreover $\mathrm{C}^h_f$ is a free $\mathrm{A}_f$-module, and therefore so is $\mathrm{B}_f$, which by \sref[0]{0.1.3.2} is an extension of free modules.
\end{proof}

\begin{corollary}[6.9.3]
\label{IV.6.9.3}
Let $\mathrm{S}$ be a Noetherian prescheme, $u:\mathrm{X}\to\mathrm{S}$ a morphism of finite type, $\sh{F}$ a coherent $\sh{O}_\mathrm{X}$-Module. There then exists a partition $(\mathrm{S}_i)_{1\leqslant i\leqslant n}$ of $\mathrm{S}$ into a finite number of locally closed subsets of $\mathrm{S}$, such that, if one still denotes by $\mathrm{S}_i$ the reduced sub-prescheme having $\mathrm{S}_i$ for underlying space, and if one sets $\mathrm{X}_i=\mathrm{X}\times_\mathrm{S}\mathrm{S}_i$, the $\sh{O}_{\mathrm{X}_i}$-Module $\sh{F}_i=\sh{F}\otimes_{\sh{O}_\mathrm{S}}\sh{O}_{\mathrm{S}_i}$ is flat over $\mathrm{S}_i$.
\end{corollary}

\begin{proof}
We proceed by Noetherian induction \sref[0]{0.2.2.2} on the set of closed subsets $\mathrm{T}$ of $\mathrm{S}$ such that the reduced closed sub-prescheme with underlying space $\mathrm{T}$ satisfies the conclusion of \sref{IV.6.9.3}. One can therefore limit oneself to proving the corollary for $\mathrm{S}$ by supposing that it is true for every reduced closed sub-prescheme whose underlying space is a closed subset distinct from $\mathrm{S}$. Since the morphism $\mathrm{S}_\mathrm{red}\to\mathrm{S}$ is of finite type and surjective, one can replace $\mathrm{S}$ by $\mathrm{S}_\mathrm{red}$, hence suppose $\mathrm{S}$ reduced and non-empty. Since $\mathrm{S}$ is Noetherian, the interior $\mathrm{T}$ of an irreducible component of $\mathrm{S}$ is non-empty, and the induced prescheme on $\mathrm{T}$ is integral; there therefore exists by virtue of \sref{IV.6.9.1} a non-empty open $\mathrm{U}\subset\mathrm{T}$

\oldpage[IV-2]{155}

such that $\sh{F}|u^{-1}(\mathrm{U})$ is flat over $\mathrm{U}$. If $\mathrm{Y}$ is then the reduced closed sub-prescheme of $\mathrm{S}$ with underlying space $\mathrm{S}-\mathrm{U}$, there exists by hypothesis a partition $(\mathrm{Y}_i)$ of $\mathrm{Y}$ into locally closed subsets of $\mathrm{Y}$ (hence of $\mathrm{S}$) such that $\sh{F}_i=\sh{F}\otimes_{\sh{O}_\mathrm{S}}\sh{O}_{\mathrm{Y}_i}$ is flat over $\mathrm{Y}_i$ for every $i$; it is clear that the $\mathrm{Y}_i$ and $\mathrm{U}$ form a partition answering the question.
\end{proof}

\subsection{Dimension and depth of a module normally flat along a closed sub-prescheme}
\label{subsection:IV.6.10}

\begin{env}[6.10.1]
\label{IV.6.10.1}
Let $\mathrm{X}$ be a locally Noetherian prescheme, $\sh{J}$ a quasi-coherent Ideal of $\sh{O}_\mathrm{X}$, $\mathrm{Y}$ the closed sub-prescheme of $\mathrm{X}$ defined by $\sh{J}$, $j:\mathrm{Y}\to\mathrm{X}$ the canonical injection. For every integer $k\geqslant 0$, the $\sh{O}_\mathrm{X}$-Module $\sh{J}^k\sh{F}/\sh{J}^{k+1}\sh{F}$ is annihilated by $\sh{J}$, hence of the form $j_*(\sh{G}_k)$, where $\sh{G}_k=j^*(\sh{J}^k\sh{F}/\sh{J}^{k+1}\sh{F})=j^*(\sh{J}^k\sh{F})$ is a coherent $\sh{O}_\mathrm{Y}$-Module. By abuse of language, we shall denote by $\operatorname{gr}^{\bullet}_\sh{J}(\sh{F})$ the graded $\sh{O}_\mathrm{Y}$-Module equal to the direct sum
\[
\bigoplus_{k=0}^{\infty}\sh{G}_k = j^*\Big(\bigoplus_{k=0}^{\infty}\sh{J}^k\sh{F}/\sh{J}^{k+1}\sh{F}\Big);
\]
  in particular, we have $\operatorname{gr}^0_\sh{J}(\sh{F})=\sh{G}_0=\sh{F}\otimes_{\sh{O}_\mathrm{X}}\sh{O}_\mathrm{Y}=j^*(\sh{F})$. We shall say (following Hironaka) that $\sh{F}$ is \emph{normally flat along} $\mathrm{Y}$ if $\operatorname{gr}^{\bullet}_\sh{J}(\sh{F})$ is a flat $\sh{O}_\mathrm{Y}$-Module. This amounts to the same \sref{IV.2.1.12} and \sref[0]{0.6.1.2} as saying that each of the $\sh{O}_\mathrm{Y}$-Modules $\sh{G}_k=j^*(\sh{J}^k\sh{F})$ is \emph{locally free}.
\end{env}

\begin{proposition}[6.10.2]
\label{IV.6.10.2}
Let $\mathrm{X}$ be a locally Noetherian prescheme, $\mathrm{Y}$ a closed integral sub-prescheme of $\mathrm{X}$. For every coherent $\sh{O}_\mathrm{X}$-Module $\sh{F}$, there exists an open $\mathrm{U}$ of $\mathrm{X}$ such that $\mathrm{Y}\cap\mathrm{U}\neq\varnothing$ and that $\sh{F}|\mathrm{U}$ is normally flat along $\mathrm{Y}\cap\mathrm{U}$.
\end{proposition}

\begin{proof}
Indeed, let $\sh{J}$ be the coherent ideal of $\sh{O}_\mathrm{X}$ defining $\mathrm{Y}$. The $\sh{O}_\mathrm{Y}$-algebra $\sh{B}=\operatorname{gr}^{\bullet}_\sh{J}(\sh{O}_\mathrm{X})$ is quasi-coherent and \emph{of finite type}, since it is generated by $\operatorname{gr}^1_\sh{J}(\sh{O}_\mathrm{X})$, the inverse image of $\sh{J}/\sh{J}^2$. Since $\operatorname{gr}^{\bullet}_\sh{J}(\sh{F})$ is a quasi-coherent $\sh{B}$-module generated by $\operatorname{gr}^0_\sh{J}(\sh{F})$, it is a $\sh{B}$-Module \emph{of finite type}. If one sets $\mathrm{Z}=\operatorname{Spec}(\sh{B})$, the structural morphism $u:\mathrm{Z}\to\mathrm{Y}$ is of finite type, and if $\sh{H}$ is the coherent $\sh{O}_\mathrm{Z}$-Module such that $u_*(\sh{H})=\operatorname{gr}^{\bullet}_\sh{J}(\sh{F})$, it suffices to apply to $u$ and to $\sh{H}$ the theorem of generic flatness \sref{IV.6.9.1} to prove the proposition.
\end{proof}

\begin{proposition}[6.10.3]
\label{IV.6.10.3}
The notations being those of \sref{IV.6.10.1}, suppose that $\sh{F}$ is normally flat along $\mathrm{Y}$. Then:
\begin{enumerate}[label=(\roman*)]
  \item $\mathrm{Y}\cap\operatorname{Supp}(\sh{F})$ is a part both open and closed of $\mathrm{Y}$ (in other words, a union of connected components of $\mathrm{Y}$).
  \item If $\sh{J}$ is locally nilpotent, $\operatorname{Supp}(\sh{F})$ is an open and closed part of $\mathrm{X}$, and for every $x\in\operatorname{Supp}(\sh{F})$, we have
\end{enumerate}
\begin{align}
  \dim(\sh{F}_x) &= \dim(\sh{O}_{\mathrm{Y},x}) \tag{6.10.3.1}\label{IV.6.10.3.1}\\
  \operatorname{prof}(\sh{F}_x) &= \operatorname{prof}(\sh{O}_{\mathrm{Y},x}) \tag{6.10.3.2}\label{IV.6.10.3.2}\\
  \operatorname{coprof}(\sh{F}_x) &= \operatorname{coprof}(\sh{O}_{\mathrm{Y},x}). \tag{6.10.3.3}\label{IV.6.10.3.3}
\end{align}
\end{proposition}

\begin{proof}
(i) With the notations of \sref{IV.6.10.1}, we have $\mathrm{Y}\cap\operatorname{Supp}(\sh{F})=\operatorname{Supp}(\sh{G}_0)$ \sref[I]{I.9.1.13}, and as by hypothesis $\sh{G}_0$ is a locally free $\sh{O}_\mathrm{Y}$-Module of finite type, its support is at the same time open and closed in $\mathrm{Y}$.

\oldpage[IV-2]{156}

(ii) The hypothesis that $\sh{J}$ is locally nilpotent implies that the underlying spaces of $\mathrm{X}$ and $\mathrm{Y}$ are the same, whence the two first assertions of (ii), taking into account (i); it remains to prove \sref{IV.6.10.3.2} and \sref{IV.6.10.3.3} by difference. Let $(f_i)_{1\leqslant i\leqslant p}$ be a finite family of elements of the maximal ideal of $\sh{O}_{\mathrm{X},x}$ whose images in $\sh{O}_{\mathrm{Y},x}=\sh{O}_{\mathrm{X},x}/\sh{J}_x$ form a maximal $\sh{O}_{\mathrm{Y},x}$-regular sequence; the hypothesis on $\sh{F}$ and $\sh{J}$ implies that the $\sh{O}_{\mathrm{Y},x}$-module $\operatorname{gr}^{\bullet}_{\sh{J}_x}(\sh{F}_x)$ is free of finite type, hence the suite $(f_i)$ is $\operatorname{gr}^{\bullet}_{\sh{J}_x}(\sh{F}_x)$-regular; one deduces that this suite is also $\sh{F}_x$-regular \sref[0]{0.15.1.19}. On the other hand, let $n$ be the largest integer such that $\sh{F}^{(n)}_x=\sh{J}^n_x\sh{F}_x\neq 0$. The hypothesis implies also that $\operatorname{gr}^{\bullet}_{\sh{J}_x}(\sh{F}_x/\sh{F}^{(n)}_x)$ is free of finite type, hence the suite $(f_i)$ is also $(\sh{F}_x/\sh{F}^{(n)}_x)$-regular (\emph{loc.\ cit.}). Applying the lemma \sref{IV.3.4.1.4}, one concludes by induction on $i$ that there is an exact sequence
\[
0\to\sh{F}^{(n)}_x\Big/\Big(\sum_{i=1}^p f_i\sh{F}^{(n)}_x\Big) \to \sh{F}_x\Big/\Big(\sum_{i=1}^p f_i\sh{F}_x\Big).
\]
But the hypothesis implies that $\sh{F}^{(n)}_x$ is a free $\sh{O}_{\mathrm{Y},x}$-module of finite type and $\neq 0$, hence $\sh{F}^{(n)}_x/(\sum_{i=1}^p f_i\sh{F}^{(n)}_x)$ is isomorphic to a module of the form $(\sh{O}_{\mathrm{Y},x}/(\sum_{i=1}^p f_i\sh{O}_{\mathrm{Y},x}))^m$ with $m>0$; as $\operatorname{prof}(\sh{O}_{\mathrm{Y},x}/(\sum_{i=1}^p f_i\sh{O}_{\mathrm{Y},x}))=0$ \sref[0]{0.16.4.6} one also has, by the characterisation \sref[0]{0.16.4.6},~(i) of modules of zero depth, $\operatorname{prof}(\sh{F}^{(n)}_x/(\sum_{i=1}^p f_i\sh{F}^{(n)}_x))=0$, then $\operatorname{prof}(\sh{F}_x/(\sum_{i=1}^p f_i\sh{F}_x))=0$; the $(f_i)_x$ belonging to the maximal ideal of $\sh{O}_{\mathrm{Y},x}$ and forming a maximal $\sh{F}_x$-regular suite, this shows that $\operatorname{prof}(\sh{F}_x)=p$ \sref[0]{0.16.4.6},~(ii).
\end{proof}

\begin{corollary}[6.10.4]
\label{IV.6.10.4}
Let $\mathrm{U}_{\mathrm{S}_n}(\sh{F})$ (resp.\ $\mathrm{U}_{\mathrm{C}_n}(\sh{F})$) be the set of $x\in\mathrm{X}$ such that $\sh{F}$ satisfies $(\mathrm{S}_n)$ (resp.\ such that $\operatorname{coprof}(\sh{F}_x)\leqslant n$). If $\sh{F}$ is normally flat along $\mathrm{Y}$, if $\operatorname{Supp}(\sh{F})=\mathrm{X}$, and if $\sh{J}$ is locally nilpotent, then $\mathrm{U}_{\mathrm{S}_n}(\sh{F})=\mathrm{U}_{\mathrm{S}_n}(\sh{O}_\mathrm{Y})$ and $\mathrm{U}_{\mathrm{C}_n}(\sh{F})=\mathrm{U}_{\mathrm{C}_n}(\sh{O}_\mathrm{Y})$.
\end{corollary}

\begin{proposition}[6.10.5]
\label{IV.6.10.5}
The notations being those of \sref{IV.6.10.1}, suppose that $\mathrm{Y}$ is connected and non-empty, and that $\sh{F}$ is normally flat along $\mathrm{Y}$. For every integer $n\geqslant 0$, set
\begin{equation}
  r(n) = \operatorname{rg}_{\sh{O}_\mathrm{Y}}(\operatorname{gr}^n_\sh{J}(\sh{F})) \tag{6.10.5.1}\label{IV.6.10.5.1}
\end{equation}
(the locally free $\sh{O}_\mathrm{Y}$-Module $\operatorname{gr}^n_\sh{J}(\sh{F})$ being necessarily of constant rank). Then:
\begin{enumerate}[label=(\roman*)]
  \item There exists a polynomial $\mathrm{P}\in\bb{Q}[\mathrm{T}]$ such that $\mathrm{P}(n)=r(n)$ for all sufficiently large $n$.
  \item We have $\mathrm{Y}\cap\operatorname{Supp}(\sh{F})=\varnothing$ (in other words $\sh{F}/\sh{J}\sh{F}=0$), or $\mathrm{Y}\cap\operatorname{Supp}(\sh{F})=\mathrm{Y}$ (in other words $\mathrm{Y}\subset\operatorname{Supp}(\sh{F})$).
\end{enumerate}
In the second case, denote $d-1$ the degree of $\mathrm{P}$; for every maximal point $y$ of $\mathrm{Y}$, we have
\begin{equation}
  \dim(\sh{F}_y) = d \tag{6.10.5.2}\label{IV.6.10.5.2}
\end{equation}
and in particular
\begin{equation}
  \operatorname{codim}(\mathrm{Y},\operatorname{Supp}(\sh{F})) = d. \tag{6.10.5.3}\label{IV.6.10.5.3}
\end{equation}
\begin{enumerate}[label=(\roman*)]
\setcounter{enumi}{2}
  \item Suppose $\mathrm{Y}\subset\operatorname{Supp}(\sh{F})$. For every $x\in\mathrm{Y}$, we have
\end{enumerate}
\begin{equation}
  \dim(\sh{F}_x) = \dim(\sh{O}_{\mathrm{Y},x}) + d. \tag{6.10.5.4}\label{IV.6.10.5.4}
\end{equation}
More precisely, there exist in $\sh{J}_x$ elements $f_i$ ($1\leqslant i\leqslant d$) forming part of a system of parameters for $\sh{F}_x$ \sref[0]{0.16.3.6} and such that, in $\operatorname{Spec}(\sh{O}_{\mathrm{X},x})$, we have $\mathrm{V}(\sh{J}_x)=\mathrm{V}(\sum_{i=1}^d f_i\sh{O}_{\mathrm{X},x})\cap\operatorname{Supp}(\sh{F})$.
\end{proposition}

\begin{proof}
\oldpage[IV-2]{157}
As $\mathrm{Y}$ is assumed connected, the first assertion of (ii) follows from \sref{IV.6.10.3},~(i). If $\mathrm{Y}\cap\operatorname{Supp}(\sh{F})=\varnothing$, assertion (i) is trivial with $\mathrm{P}=0$. Suppose now that $\operatorname{Supp}(\sh{F})\supset\mathrm{Y}$, and let $y$ be a maximal point of $\mathrm{Y}$. Since then $\mathrm{Y}=\operatorname{Supp}(\sh{F}/\sh{J}\sh{F})$, the module $\sh{F}_y/\sh{J}_y\sh{F}_y$ has finite length \sref{IV.3.1.4}. Setting $s(n)=\operatorname{length}(\operatorname{gr}^n_{\sh{J}_y}(\sh{F}_y))$, we know that there exists a polynomial $\mathrm{R}\in\bb{Q}[\mathrm{T}]$ such that $s(n)=\mathrm{R}(n)$ for $n$ large enough, namely $\mathrm{R}(n)=\mathrm{H}(n+1)-\mathrm{H}(n)$, where $\mathrm{H}$ is the Hilbert--Samuel polynomial of $\sh{F}_y$ for the $\sh{J}_y$-adic filtration \sref[0]{0.16.2.1}. Since the $\sh{O}_{\mathrm{Y},y}$-modules $\operatorname{gr}^n_{\sh{J}_y}(\sh{F}_y)$ are free by hypothesis, we have $r(n)=s(n)/m$, denoting by $m$ the length of the $\sh{O}_{\mathrm{X},y}$-module $\sh{O}_{\mathrm{X},y}/\sh{J}_y$; thus (i) holds with $\mathrm{P}=\frac{1}{m}\mathrm{R}$. We know furthermore \sref[0]{0.16.2.3} that $\deg(\mathrm{H})=\dim(\sh{F}_y)$, whence the relation \sref{IV.6.10.5.2}; the relation \sref{IV.6.10.5.3} follows from \sref{IV.5.1.12.2}.

It remains to prove (iii) for any point $x\in\mathrm{Y}$. Set $\mathrm{A}=\sh{O}_{\mathrm{X},x}$, $\mathfrak{J}=\sh{J}_x$, and $\mathrm{M}=\sh{F}_x$, so that $\sh{O}_{\mathrm{Y},x}=\mathrm{A}/\mathfrak{J}$; let $\mathrm{S}=\operatorname{gr}^{\bullet}_\mathfrak{J}(\mathrm{A})$, which is a graded $(\mathrm{A}/\mathfrak{J})$-algebra of finite type, with positive degrees, such that $\mathrm{S}_0=\mathrm{A}/\mathfrak{J}$, and generated by its homogeneous elements of degree $1$; let finally $\mathrm{N}=\operatorname{gr}^{\bullet}_\mathfrak{J}(\mathrm{M})$, which is a graded $\mathrm{S}$-module of finite type, each homogeneous component $\mathrm{N}_n$ being by hypothesis a free $(\mathrm{A}/\mathfrak{J})$-module of rank $r(n)$. Let $\mathfrak{m}$ be the maximal ideal of $\mathrm{A}$, $k=\mathrm{A}/\mathfrak{m}$ its residue field; $\mathrm{B}=\mathrm{S}\otimes_\mathrm{A}k$ is a graded $k$-algebra of finite type, with positive degrees, generated by its homogeneous elements of degree $1$ and such that $\mathrm{B}_0=k$, so that $\mathfrak{q}=\mathrm{B}^+=\bigoplus_{n\geqslant 1}\mathrm{B}_n$ is a maximal ideal of $\mathrm{B}$; $\mathrm{E}=\mathrm{N}\otimes_\mathrm{A}k$ is a graded $\mathrm{B}$-module of finite type such that $\operatorname{rg}_k(\mathrm{E}_n)=r(n)$. Applying \sref[0]{0.16.2.7} to $\mathrm{B}$ and $\mathrm{E}$, choose homogeneous elements $t_i\in\mathrm{B}_{n_i}$ ($1\leqslant i\leqslant d$), and lift each $t_i$ to an element $f_i\in\mathfrak{J}^{n_i}$. For $n\geqslant\sup(n_i)$, consider the $\mathrm{A}$-submodule $\sum_{i=1}^d f_i\mathfrak{J}^{n-n_i}\mathrm{M}$ of $\mathrm{M}$; as the homogeneous component of degree $n$ of the submodule $\sum_i t_i\mathrm{E}$ of $\mathrm{E}$ is equal to $\mathrm{E}_n$ as soon as $n$ is large enough, one sees that, for $n$ large enough, we have
\[
\mathfrak{J}^n\mathrm{M} = \sum_{i=1}^d f_i\mathfrak{J}^{n-n_i}\mathrm{M}+\mathfrak{m}\mathfrak{J}^n\mathrm{M}
\]
and as $\mathfrak{J}^n\mathrm{M}$ is an $\mathrm{A}$-module of finite type, this implies, by Nakayama's lemma,
\[
\mathfrak{J}^n\mathrm{M} = \sum_{i=1}^d f_i\mathfrak{J}^{n-n_i}\mathrm{M} \subset \sum_{i=1}^d f_i\mathrm{M}.
\]
If $\mathfrak{a}$ is the annihilator of $\mathrm{M}$, we have therefore \sref[0]{0.1.7.5} in $\operatorname{Spec}(\mathrm{A})$
\[
\mathrm{V}\Big(\sum_{i=1}^d f_i\mathrm{A}\Big)\cap\mathrm{V}(\mathfrak{a})\subset\mathrm{V}(\mathfrak{J}^n)\cap\mathrm{V}(\mathfrak{a})=\mathrm{V}(\mathfrak{J})\cap\mathrm{V}(\mathfrak{a})=\mathrm{V}(\mathfrak{J})
\]
since by hypothesis $\mathrm{Y}\cap\operatorname{Spec}(\mathrm{A})=\mathrm{V}(\mathfrak{J})\cap\operatorname{Supp}(\mathrm{M})=\mathrm{V}(\mathfrak{J})$. Since on the other hand the $f_i$ belong to $\mathfrak{J}$, we have $\mathrm{V}(\mathfrak{J})\subset\mathrm{V}(\sum_{i=1}^d f_i\mathrm{A})$, which proves the last relation of (iii). It remains to show that the $f_i$ form part of a system of parameters for $\mathrm{M}$. Replacing $\mathrm{A}$ by $\mathrm{A}/\mathfrak{a}$ and the $f_i$ by their images in $\mathrm{A}/\mathfrak{a}$, one can limit oneself to the case where $\mathfrak{a}=0$; we have just seen that
\[
\dim(\mathrm{A}/\mathfrak{J})=\dim\Big(\mathrm{A}/\Big(\sum_{i=1}^d f_i\mathrm{A}\Big)\Big),
\]
and it remains therefore to prove \sref[0]{0.16.3.7} that we have
\[
\dim(\mathrm{A})\geqslant\dim(\mathrm{A}/\mathfrak{J})+d.
\]

Now, let $\mathfrak{p}$ be a prime ideal of $\mathrm{A}$ containing $\mathfrak{J}$ and such that $\dim(\mathrm{A}/\mathfrak{J})=\dim(\mathrm{A}/\mathfrak{p})$, so that $\mathfrak{p}$ is minimal among the prime ideals containing $\mathfrak{J}$. We have therefore $\mathfrak{p}=\mathfrak{i}_y$ where $y\in\operatorname{Spec}(\mathrm{A})$ is a maximal point of $\mathrm{Y}$. But by virtue of \sref{IV.6.10.5.2} and the hypothesis $\operatorname{Spec}(\mathrm{A})=\operatorname{Supp}(\mathrm{M})$, we have $\dim(\mathrm{A}_\mathfrak{p})=d$; the inequality $\dim(\mathrm{A}/\mathfrak{p})+\dim(\mathrm{A}_\mathfrak{p})\leqslant\dim(\mathrm{A})$ \sref[0]{0.16.1.4} completes the proof.
\end{proof}

\begin{proposition}[6.10.6]
\label{IV.6.10.6}
Let $\mathrm{X}$ be a locally Noetherian prescheme, $\sh{F}$ a coherent $\sh{O}_\mathrm{X}$-Module, $\mathrm{Y}$ a closed irreducible sub-prescheme of $\mathrm{X}$, of generic point $y\in\operatorname{Supp}(\sh{F})$. There then exists a neighbourhood $\mathrm{U}$ open and non-empty of $y$ in $\mathrm{X}$ such that, for every $x\in\mathrm{U}\cap\mathrm{Y}$, we have
\begin{align}
  \dim(\sh{F}_x) &= \dim(\sh{F}_y)+\dim(\sh{O}_{\mathrm{Y},x}) \tag{6.10.6.1}\label{IV.6.10.6.1}\\
  \operatorname{prof}(\sh{F}_x) &= \operatorname{prof}(\sh{F}_y)+\operatorname{prof}(\sh{O}_{\mathrm{Y},x}) \tag{6.10.6.2}\label{IV.6.10.6.2}\\
  \operatorname{coprof}(\sh{F}_x) &= \operatorname{coprof}(\sh{F}_y)+\operatorname{coprof}(\sh{O}_{\mathrm{Y},x}). \tag{6.10.6.3}\label{IV.6.10.6.3}
\end{align}
\end{proposition}

\begin{proof}
\oldpage[IV-2]{158}
Let $\mathrm{Y}'=\mathrm{Y}_\mathrm{red}$, which is a closed integral sub-prescheme of $\mathrm{Y}$ having the same underlying space, and defined by a locally nilpotent Ideal $\sh{J}$ of $\sh{O}_\mathrm{Y}$ \sref[I]{I.6.1.6}. It follows that $\operatorname{gr}^{\bullet}_{\sh{J}}(\sh{O}_\mathrm{Y})$ is a coherent $\sh{O}_{\mathrm{Y}'}$-Module, and as $\mathrm{Y}'$ is integral, there exists an open neighbourhood $\mathrm{V}$ of $y$ in $\mathrm{Y}'$ such that $\operatorname{gr}^{\bullet}_{\sh{J}}(\sh{O}_\mathrm{Y})|\mathrm{V}$ is locally free \sref[0]{0.5.2.7}; otherwise said, replacing $\mathrm{X}$ by a neighbourhood of $y$, one can suppose that $\sh{O}_\mathrm{Y}$ is \emph{normally flat} along $\mathrm{Y}'$; one then deduces from \sref{IV.6.10.3} that we have
\[
\dim(\sh{O}_{\mathrm{Y},x})=\dim(\sh{O}_{\mathrm{Y}',x}),\qquad \operatorname{prof}(\sh{O}_{\mathrm{Y},x})=\operatorname{prof}(\sh{O}_{\mathrm{Y}',x})
\]
for every $x\in\mathrm{Y}$; this allows us to limit ourselves to the case where the closed sub-prescheme $\mathrm{Y}$ is \emph{integral}.

This being so, by virtue of \sref{IV.6.10.2}, one can, replacing $\mathrm{X}$ by an open neighbourhood of $y$, suppose that $\sh{F}$ is \emph{normally flat} along $\mathrm{Y}$. Since $y\in\operatorname{Supp}(\sh{F})$, relation \sref{IV.6.10.6.1} follows from \sref{IV.6.10.5.2} and \sref{IV.6.10.5.4}. Suppose now $p=\operatorname{prof}(\sh{F}_y)$; replacing at need $\mathrm{X}$ by an open neighbourhood of $y$, one can further suppose that there exist $p$ sections $f_i$ ($1\leqslant i\leqslant p$) of $\sh{O}_\mathrm{X}$ above $\mathrm{X}$, forming an $\sh{F}$-regular sequence, such that the $(f_i)_y$ belong to the maximal ideal $\mathfrak{m}_y$ of $\sh{O}_{\mathrm{X}}$ at the point $y$ and that $\operatorname{prof}(\sh{F}_y/(\sum_{i=1}^p (f_i)_y\sh{F}_y))=0$ \sref[0]{0.16.4.6}. If $\sh{J}$ is the Ideal of $\sh{O}_\mathrm{X}$ defining $\mathrm{Y}$, we have $\sh{J}_y=\mathfrak{m}_y$ and, replacing $\mathrm{X}$ once more by a neighbourhood of $y$ if necessary, we may suppose that $f_i\in\Gamma(\mathrm{X},\sh{J})$ for $1\leqslant i\leqslant p$. Moreover, setting $\sh{G}=\sh{F}/(\sum_{i=1}^p f_i\sh{F})$, the hypothesis $\operatorname{prof}(\sh{G}_y)=0$ implies that $\sh{G}_y$ contains a sub-module isomorphic to $\sh{O}_{\mathrm{X},y}/\mathfrak{m}_y$ \sref[0]{0.16.4.6}; reasoning as above (taking into account \sref[0]{0.5.3.8} and \sref[0]{0.5.2.7}) shows that one can suppose that there exists a sub-$\sh{O}_\mathrm{X}$-Module $\sh{G}'$ of $\sh{G}$ \emph{isomorphic to} $\sh{O}_\mathrm{X}/\sh{J}$ so that $\sh{G}'_x=\sh{O}_{\mathrm{Y},x}$ for every $x\in\mathrm{Y}$. Setting $\sh{G}''=\sh{G}/\sh{G}'$; replacing $\mathrm{X}$ by an open neighbourhood of $y$, one can, by virtue of \sref{IV.6.10.2}, suppose that $\sh{G}'$ and $\sh{G}''$ are \emph{normally flat} along $\mathrm{Y}$. Let then $x$ be any point of $\mathrm{Y}$, and set $q=\operatorname{prof}(\sh{O}_{\mathrm{Y},x})=\operatorname{prof}(\sh{G}'_x)$; let $(g_j)_{1\leqslant j\leqslant q}$ be a maximal $\sh{G}'_x$-regular sequence of elements of the maximal ideal $\mathfrak{m}_x$ of $\sh{O}_{\mathrm{X},x}$. Each of the homogeneous components of $\operatorname{gr}^{\bullet}_{\sh{J}_x}(\sh{G}''_x)$ is a flat $\sh{O}_{\mathrm{Y},x}$-module of finite type by hypothesis, hence a free $\sh{O}_{\mathrm{Y},x}$-module, and consequently $\operatorname{gr}^{\bullet}_{\sh{J}_x}(\sh{G}''_x)$ is free as well; since the sequence $(g_j)$ is $\sh{O}_{\mathrm{Y},x}$-regular, it is therefore $\operatorname{gr}^{\bullet}_{\sh{J}_x}(\sh{G}''_x)$-regular, whence it is $\sh{G}''_x$-regular \sref[0]{0.15.1.19}. Applying the lemma \sref[0]{0.15.1.18} to the exact sequence
\[
0\to\sh{G}'_x\to\sh{G}_x\to\sh{G}''_x\to 0
\]
by induction on $j$, one concludes an exact sequence
\[
0\to\sh{G}'_x\Big/\Big(\sum_{j=1}^q g_j\sh{G}'_x\Big)\to\sh{G}_x\Big/\Big(\sum_{j=1}^q g_j\sh{G}_x\Big).
\]
But by hypothesis $\operatorname{prof}(\sh{G}'_x/(\sum_{j=1}^q g_j\sh{G}'_x))=0$ \sref[0]{0.16.4.6}; by the characterisation \sref[0]{0.16.4.6} of modules of zero depth, one concludes that $\operatorname{prof}(\sh{G}_x/(\sum_{j=1}^q g_j\sh{G}_x))=0=\operatorname{prof}(\sh{F}_x/(\sum_{i=1}^p f_i\sh{F}_x+\sum_{j=1}^q g_j\sh{F}_x))$; the sequence $(g_j)$, being $\sh{G}'_x$-regular and $\sh{G}''_x$-regular, is also $\sh{G}_x$-regular; finally, the sequence $((f_i)_x)$ is $\sh{F}_x$-regular by hypothesis and consists of elements of the maximal ideal $\mathfrak{m}_x$; one deduces \sref[0]{0.16.4.6} that $\operatorname{prof}(\sh{F}_x)=p+q$, which completes the proof.
\end{proof}


\subsection{Criteria for the sets \texorpdfstring{$\mathrm{U}_{\mathrm{S}_n}(\sh{F})$}{U\_Sn(F)} or \texorpdfstring{$\mathrm{U}_{\mathrm{C}_n}(\sh{F})$}{U\_Cn(F)} to be open}
\label{subsection:IV.6.11}
\label{IV.6.11}

\begin{lemma}[6.11.1]
\label{IV.6.11.1}
Let $\mathrm{X}$ be a locally Noetherian prescheme, $\sh{F}$ an $\sh{O}_\mathrm{X}$-Module
coherent; then the function $x\mapsto\operatorname{dim.proj}(\sh{F}_x)$ is upper semi-continuous in $\mathrm{X}$.
\end{lemma}

\begin{proof}
One can restrict to the case where $\mathrm{X}=\operatorname{Spec}(\mathrm{A})$ is the spectrum of a Noetherian ring
and $\sh{F}=\widetilde{\mathrm{M}}$, where $\mathrm{M}$ is an $\mathrm{A}$-module of finite type. Suppose that for an $x\in\mathrm{X}$, we have
$\operatorname{dim.proj}(\mathrm{M}_x)=n<+\infty$ (if $n=+\infty$ there is nothing to prove); there exists a resolution
\[
  \mathrm{L}_{n-1}\to\mathrm{L}_{n-2}\to\cdots\to\mathrm{L}_0\to\mathrm{M}\to 0
\]
\oldpage[IV-2]{159}
where the $\mathrm{L}_i$ are free $\mathrm{A}$-modules of \emph{finite type} ($\mathrm{A}$ being Noetherian), whence an exact sequence
\[
  0\to\mathrm{R}\to\mathrm{L}_{n-1}\to\mathrm{L}_{n-2}\to\cdots\to\mathrm{L}_0\to\mathrm{M}\to 0
\]
where $\mathrm{R}$ is an $\mathrm{A}$-module of finite type; one deduces an exact sequence
\[
  0\to\mathrm{R}_x\to(\mathrm{L}_{n-1})_x\to\cdots\to(\mathrm{L}_0)_x\to\mathrm{M}_x\to 0
\]
where the $(\mathrm{L}_i)_x$ are free $\mathrm{A}_x$-modules of finite type; as by hypothesis
$\operatorname{dim.proj}(\mathrm{M}_x)=n$, this implies that $\mathrm{R}_x$ is a projective $\mathrm{A}_x$-module of finite type
(\textbf{M},~VI,~2.1) and hence a free $\mathrm{A}_x$-module of finite type $(\mathbf{0}_{\mathrm{III}},~10.1.3)$. We conclude
that there exists an open neighbourhood $\mathrm{U}$ of $x$ in $\mathrm{X}$ such that, for every $x'\in\mathrm{U}$, $\mathrm{R}_{x'}$ is a
free $\mathrm{A}_{x'}$-module $(\mathbf{0}_\mathrm{I},~5.2.7)$, whence $\mathrm{M}_{x'}$ admits a projective resolution of length $n$,
in other words $\operatorname{dim.proj}(\mathrm{M}_{x'})\leqslant n$, which proves the lemma.
\end{proof}

\begin{proposition}[6.11.2]
\label{IV.6.11.2}
Let $\mathrm{X}$ be a locally Noetherian prescheme and locally immersible
in a regular scheme \sref{IV.5.11.1}; let $\sh{F}$ be an $\sh{O}_\mathrm{X}$-Module coherent. Then:
\begin{enumerate}[label=(\roman*)]
  \item \textnormal{(M.~Auslander)} The function $x\mapsto\operatorname{coprof}(\sh{F}_x)$ is upper semi-continuous in $\mathrm{X}$
(in other words, for every integer $n$, the set $\mathrm{U}_{\mathrm{C}_n}(\sh{F})$ of $x\in\mathrm{X}$ such that $\operatorname{coprof}(\sh{F}_x)\leqslant n$
is \emph{open}).
  \item For every integer $n$, the set $\mathrm{U}_{\mathrm{S}_n}(\sh{F})$ of $x\in\mathrm{X}$ such that $\sh{F}$ has the property $(\mathrm{S}_n)$
at the point $x$ is open.
\end{enumerate}
\end{proposition}

\begin{proof}
(i) The question being local on $\mathrm{X}$, one can, by virtue of the hypothesis, restrict to the
case where $\mathrm{X}$ is a closed sub-prescheme of an affine regular scheme $\mathrm{Y}$. If $j:\mathrm{X}\to\mathrm{Y}$ is
the canonical injection, and $\sh{G}=j_*(\sh{F})$, one knows that one has then, for every $x\in\mathrm{X}$,
$\dim(\sh{F}_x)=\dim(\sh{G}_x)$ and $\operatorname{prof}(\sh{F}_x)=\operatorname{prof}(\sh{G}_x)$, hence $\operatorname{coprof}(\sh{F}_x)=\operatorname{coprof}(\sh{G}_x)$, and as
$\operatorname{coprof}(\sh{G}_x)=0$ for $y\notin\mathrm{X}\text{---}\mathrm{Y}$, one is reduced to proving the property for $\sh{G}$;
in other words, one can restrict to the case where $\mathrm{X}$ is an affine \emph{regular} scheme. One knows then $(\mathbf{0},~17.3.4)$
that one has
\[
  \operatorname{prof}(\sh{F}_x)=\dim(\sh{O}_{\mathrm{X},x})-\operatorname{dim.proj}(\sh{F}_x).
\]

On the other hand, if $\mathrm{S}$ is the unique closed sub-prescheme of $\mathrm{X}$ having for underlying space
$\operatorname{Supp}(\sh{F})$, one has $\dim(\sh{F}_x)=\operatorname{codim}(\{x\},\mathrm{S})$ $(5.1.12.1)$, and, by virtue of $(5.1.9)$
\[
  \dim(\sh{F}_x)=\dim(\sh{O}_{\mathrm{X},x})-\operatorname{codim}_x(\mathrm{S},\mathrm{X})
\]
since $\sh{O}_{\mathrm{X},x}$ is a regular ring, and \emph{a fortiori} equidimensional $(\mathbf{0},~16.5.12)$. We
can therefore write
\begin{equation}
\tag{6.11.2.1}
\label{IV.6.11.2.1}
  \operatorname{coprof}(\sh{F}_x)=\operatorname{dim.proj}(\sh{F}_x)-\operatorname{codim}_x(\mathrm{S},\mathrm{X})
\end{equation}
and the proposition follows then from \sref{IV.6.11.1} and from $(\mathbf{0},~14.2.6)$.

\oldpage[IV-2]{160}

(ii) As the $\mathrm{U}_{\mathrm{C}_n}(\sh{F})$ are open for every $n$ by (i), the $\mathrm{Z}_n=\mathrm{X}-\mathrm{U}_{\mathrm{C}_n}(\sh{F})$
are closed in $\mathrm{X}$; furthermore it is clear that $\mathrm{Z}_{n+1}\subset\mathrm{Z}_n$, and as $\dim(\sh{F}_x)$ is finite
for every $x\in\mathrm{X}$ and $\operatorname{coprof}(\sh{F}_x)\leqslant\dim(\sh{F}_x)$, the intersection of the $\mathrm{Z}_n$ is \emph{empty}; as one can
restrict to the case where $\mathrm{X}$ is affine, hence quasi-compact, one can suppose that there exists an $m$
such that $\mathrm{Z}_m=\varnothing$. Now, it follows from $(5.7.4)$ that the relation $x\in\mathrm{U}_{\mathrm{S}_k}(\sh{F})$ is equivalent to the set of
relations
\begin{equation}
\tag{6.11.2.2}
\label{IV.6.11.2.2}
  \operatorname{codim}_x(\mathrm{Z}_n,\mathrm{S})>n+k
\end{equation}
for every $n\geqslant 0$; but for $n\geqslant m$ this relation is automatically satisfied, hence one
need only consider the relations \sref{IV.6.11.2.2} for $0\leqslant n<m$. Now $(\mathbf{0},~14.2.6)$
the set $\mathrm{V}_{n,k}$ of $x$ satisfying \sref{IV.6.11.2.2} is open, and $\mathrm{U}_{\mathrm{S}_k}(\sh{F})$, intersection of the $\mathrm{V}_{n,k}$
for $0\leqslant n<m$, is also open.

This completes the proof.
\end{proof}

Let us recall that the hypothesis that $\mathrm{X}$ is locally immersible in a regular scheme
is always satisfied when $\mathrm{X}$ is a prescheme \emph{locally of finite type over a field $k$}
$(5.8.3)$.

\begin{corollary}[6.11.3]
\label{IV.6.11.3}
Under the hypotheses of \sref{IV.6.11.2}, the set $\mathrm{CM}(\sh{F})$ of points
$x\in\mathrm{X}$ such that $\sh{F}_x$ is a Cohen-Macaulay module is open in $\mathrm{X}$.
\end{corollary}

\begin{proof}
Indeed, it is the set $\mathrm{U}_{\mathrm{C}_0}(\sh{F})$.
\end{proof}

\begin{remarks}[6.11.4]
\label{IV.6.11.4}
(i) The reasoning of \sref{IV.6.11.2},~(ii)) proves that (without any
hypothesis on $\mathrm{X}$) when $\mathrm{U}_{\mathrm{C}_n}(\sh{F})$ is open for \emph{every} integer $n$, then $\mathrm{U}_{\mathrm{S}_n}(\sh{F})$ is open
for every integer $n$.

(ii) We have $\mathrm{CM}(\sh{F})=\bigcap\mathrm{U}_{\mathrm{S}_n}(\sh{F})$. If every point $x\in\mathrm{X}$ admits an open neighbourhood $\mathrm{V}$
of \emph{finite} dimension, the sequence of intersections $\mathrm{V}\cap\mathrm{U}_{\mathrm{S}_n}(\sh{F})$ is \emph{stationary} since
there exists an $m$ such that $\dim(\sh{F}_x)\leqslant m$ for every $x\in\mathrm{V}$; consequently, if the $\mathrm{U}_{\mathrm{S}_n}(\sh{F})$ are open
in $\mathrm{X}$, then $\mathrm{CM}(\sh{F})$ is also open in $\mathrm{X}$.

(iii) We shall write $\mathrm{U}_{\mathrm{S}_n}(\mathrm{X})$, $\mathrm{U}_{\mathrm{C}_n}(\mathrm{X})$, $\mathrm{CM}(\mathrm{X})$ instead of $\mathrm{U}_{\mathrm{S}_n}(\sh{O}_\mathrm{X})$, $\mathrm{U}_{\mathrm{C}_n}(\sh{O}_\mathrm{X})$, $\mathrm{CM}(\sh{O}_\mathrm{X})$.
\end{remarks}

\begin{proposition}[6.11.5]
\label{IV.6.11.5}
Let $\mathrm{A}$ be a local Noetherian ring, $\mathrm{M}$ an $\mathrm{A}$-module of finite
type. For every prime ideal $\mathfrak{p}$ of $\mathrm{A}$, we have
\begin{equation}
\tag{6.11.5.1}
\label{IV.6.11.5.1}
  \operatorname{coprof}_{\mathrm{A}_\mathfrak{p}}(\mathrm{M}_\mathfrak{p})\leqslant\operatorname{coprof}_\mathrm{A}(\mathrm{M}).
\end{equation}
\end{proposition}

\begin{proof}
One can restrict to the case where $\mathrm{M}_\mathfrak{p}\neq 0$. As $\widehat{\mathrm{A}}$ is a faithfully flat $\mathrm{A}$-module
$(\mathbf{0}_\mathrm{I},~7.3.5)$, there exists an ideal $\mathfrak{q}$ of $\widehat{\mathrm{A}}$ lying over $\mathfrak{p}$ $(\mathbf{0}_\mathrm{I},~6.5.1)$; as
$\widehat{\mathrm{M}}_\mathfrak{q}=(\mathrm{M}\otimes_\mathrm{A}\widehat{\mathrm{A}})_\mathfrak{q}=\mathrm{M}\otimes_\mathrm{A}\widehat{\mathrm{A}}_\mathfrak{q}$ and $\widehat{\mathrm{A}}_\mathfrak{q}$ is a flat $\mathrm{A}$-module, hence also a flat
$\mathrm{A}_\mathfrak{p}$-module
$(\mathbf{0}_\mathrm{I},~6.2.1)$, it follows from $(6.3.2)$, applied to the local homomorphism $\mathrm{A}_\mathfrak{p}\to\widehat{\mathrm{A}}_\mathfrak{q}$, to $\mathrm{M}_\mathfrak{p}$
and $\widehat{\mathrm{M}}_\mathfrak{q}$, that $\operatorname{coprof}_{\mathrm{A}_\mathfrak{p}}(\mathrm{M}_\mathfrak{p})\leqslant\operatorname{coprof}_{\widehat{\mathrm{A}}_\mathfrak{q}}(\widehat{\mathrm{M}}_\mathfrak{q})$. On the other part, $\mathrm{X}=\operatorname{Spec}(\widehat{\mathrm{A}})$
is isomorphic to a sub-prescheme of a regular scheme by virtue of the Cohen structure theorem
$(\mathbf{0},~19.8.8,~\text{(i)})$. One deduces therefore from \sref{IV.6.11.2} that $\operatorname{coprof}_{\widehat{\mathrm{A}}_\mathfrak{q}}(\widehat{\mathrm{M}}_\mathfrak{q})\leqslant\operatorname{coprof}_{\widehat{\mathrm{A}}}(\widehat{\mathrm{M}})$; finally,
one knows that $\operatorname{coprof}_{\widehat{\mathrm{A}}}(\widehat{\mathrm{M}})=\operatorname{coprof}_\mathrm{A}(\mathrm{M})$ $(\mathbf{0},~16.4.10)$, which finishes the proof of \sref{IV.6.11.5}.
This proposition justifies the definition of the coprofundity of an $\mathrm{A}$-module when $\mathrm{A}$
is not a local ring, given in $(5.7.12)$.
\end{proof}

\begin{proposition}[6.11.6]
\label{IV.6.11.6}
Let $\mathrm{X}$ be a locally Noetherian prescheme, $n$ an integer $>0$, $\sh{F}$ a coherent $\sh{O}_\mathrm{X}$-Module.
Suppose that for every closed integral sub-prescheme $\mathrm{Y}$ of $\mathrm{X}$, there exists
an open non-empty part $\mathrm{W}$ of $\mathrm{Y}$ such that the sub-prescheme of $\mathrm{X}$ induced on $\mathrm{W}$ satisfies $(\mathrm{S}_n)$.
Then the set $\mathrm{U}_{\mathrm{S}_n}(\sh{F})$ is open in $\mathrm{X}$.
\end{proposition}

\begin{proof}
The question being local on $\mathrm{X}$, one can restrict to the case where $\mathrm{X}$ is Noetherian.
We shall reason by induction on $n$: for $n=1$, the set $\mathrm{U}_{\mathrm{S}_1}(\sh{F})$ is open
in $\mathrm{X}$, for the set of $x\in\mathrm{X}$ where $\sh{F}$ does not satisfy $(\mathrm{S}_1)$ is the set of $x$ such that $\sh{F}_x$
admits immersed associated prime ideals $(5.7.5)$; if $(\mathrm{Z}_j)$ is the finite family of the
prime cycles associated to $\sh{F}$ which are immersed, then $\mathrm{U}_{\mathrm{S}_1}(\sh{F})=\mathrm{X}-\bigcup_j\mathrm{Z}_j$,
whence our assertion since the $\mathrm{Z}_j$ are closed. We shall therefore suppose henceforth that
$n>1$. In the second place, one can restrict to the case where $\operatorname{Supp}(\sh{F})=\mathrm{X}$, for there is a
closed sub-prescheme $\mathrm{T}$ of $\mathrm{X}$ having $\operatorname{Supp}(\sh{F})$ as underlying space, and an $\sh{O}_\mathrm{T}$-Module
coherent $\sh{G}$ such that $\sh{F}=j_*(\sh{G})$, $j:\mathrm{T}\to\mathrm{X}$ being the canonical injection $(\textbf{Err}_\textbf{III},~30)$;
as it amounts to the same to prove that $\sh{F}$ or $\sh{G}$ satisfies $(\mathrm{S}_n)$ at a point $x\in\operatorname{Supp}(\sh{F})$,
one can restrict to considering the case where $\mathrm{T}=\mathrm{X}$. Let us note finally that by definition $(5.7.2)$,
$\mathrm{U}_{\mathrm{S}_n}(\sh{F})$ is stable by generization. We shall now use the following lemma:

\begin{lemma}[6.11.6.1]
\label{IV.6.11.6.1}
Let $\mathrm{X}$ be a Noetherian space in which every irreducible closed subset
admits a generic point. For $\mathrm{E}$ to be an open subset of $\mathrm{X}$, it is necessary and it
suffices that $\mathrm{E}$ be stable by generization, and that, for every open part $\mathrm{V}$ of $\mathrm{X}$ and every irreducible
closed subset $\mathrm{Y}$ of $\mathrm{V}$, such that $\mathrm{V}-\mathrm{Y}\subset\mathrm{E}$ and that the generic point of $\mathrm{Y}$ belongs to $\mathrm{E}$,
$\mathrm{E}\cap\mathrm{Y}$ contains an open non-empty part of $\mathrm{Y}$.
\end{lemma}

\begin{proof}
We observe that this criterion implies that of $(\mathbf{0}_{\mathrm{III}},~9.2.6)$ when every irreducible
closed subset admits a generic point. There is obviously nothing other than to prove
the sufficiency of the conditions stated.

Consider the interior $\mathrm{U}$ of $\mathrm{E}$; the closed set $\mathrm{X}-\mathrm{U}$ is a union of its
irreducible components, which are finite in number and closed in $\mathrm{X}$. If $\mathrm{E}\neq\mathrm{U}$,
the hypothesis that $\mathrm{E}$ is stable by generization would imply that the generic point $z$ of one
of the irreducible components of $\mathrm{X}-\mathrm{U}$ would belong to $\mathrm{E}$. Now, $z$ belongs to only
one of the irreducible components of $\mathrm{X}-\mathrm{U}$; if $\mathrm{T}$ is the union of the other
irreducible components of $\mathrm{X}-\mathrm{U}$, $\mathrm{V}=\mathrm{X}-\mathrm{T}$ is open in $\mathrm{X}$, union of $\mathrm{U}$ and of
the set $\mathrm{Y}=\overline{\{z\}}\cap\mathrm{V}$ closed in $\mathrm{V}$ and irreducible. By hypothesis $\mathrm{E}\cap\mathrm{Y}$ contains
an open part $\mathrm{W}$ of $\mathrm{Y}$; we conclude that $\mathrm{U}\cup\mathrm{W}$ is open \emph{in} $\mathrm{V}$, hence in $\mathrm{X}$,
which is absurd since $\mathrm{U}$ is supposed to be the interior of $\mathrm{E}$.

\oldpage[IV-2]{161}

By virtue of this lemma, one can suppose that there exists in $\mathrm{X}$ an integral sub-prescheme
$\mathrm{Y}$ of generic point $y$ such that $\sh{F}$ satisfies $(\mathrm{S}_n)$ at the point $y$ and at every point of
$\mathrm{X}-\mathrm{Y}$, and restrict to verifying that there exists in $\mathrm{X}$ an open neighbourhood of $y$ such
that $\sh{F}$ satisfies $(\mathrm{S}_n)$ in this neighbourhood. We distinguish two cases:

$1^\circ$ $y$ is a maximal point of $\mathrm{X}$; as there exists an open neighbourhood of $y$ meeting
no irreducible component of $\mathrm{X}$ other than $\overline{\{y\}}$, one can suppose
that $\mathrm{X}$ is irreducible, hence a \emph{fortiori} has the same underlying space as $\mathrm{Y}$, so that $\mathrm{Y}$ is defined
by the Nilradical of $\sh{O}_\mathrm{X}$, which is \emph{nilpotent}. On the other hand, on can, replacing $\mathrm{X}$ by
an open neighbourhood of $y$, suppose that $\sh{F}$ is \emph{normally flat} along $\mathrm{Y}$ $(6.9.1)$;
it follows then from $(6.10.4)$ that $\mathrm{U}_{\mathrm{S}_n}(\sh{F})=\mathrm{U}_{\mathrm{S}_n}(\sh{O}_\mathrm{Y})$, and as the latter is by hypothesis an open
neighbourhood of $y$ in $\mathrm{X}$, this proves the proposition in this case.

$2^\circ$ $y$ is not a maximal point of $\mathrm{X}$, in other words $\operatorname{Supp}(\sh{F})=\mathrm{X}$,
$\dim(\sh{F}_y)\geqslant 1$, hence also $\operatorname{prof}(\sh{F}_y)\geqslant 1$ since $\sh{F}$ satisfies by hypothesis $(\mathrm{S}_n)$ (and \emph{a fortiori} $(\mathrm{S}_1)$)
at the point $y$. Replacing $\mathrm{X}$ by an open neighbourhood of $y$ if necessary, one can therefore suppose
that there exists a section $f$ of $\sh{O}_\mathrm{X}$ over $\mathrm{X}$, $\sh{F}$-regular and such that $f_y\in\mathfrak{m}_y$, or
equivalently $f(y)=0$ $(\mathbf{0},~15.2.4)$; one has then also $f(x)=0$ for every $x\in\mathrm{Y}$ $(\mathbf{0}_\mathrm{I},~5.5.2)$.
One knows that $\sh{F}/f\sh{F}$ satisfies $(\mathrm{S}_{n-1})$ at the point $y$ $(5.7.6)$. Applying the induction hypothesis,
and replacing $\mathrm{X}$ by an open neighbourhood of $y$, one can therefore
suppose that $\sh{F}/f\sh{F}$ satisfies $(\mathrm{S}_{n-1})$ at every point of $\mathrm{X}$. But for every $x\in\mathrm{Y}$, the relation
$f(x)=0$ implies that $\operatorname{prof}(\sh{F}_x|f_x\sh{F}_x)=\operatorname{prof}(\sh{F}_x)-1$ and $\dim(\sh{F}_x|f_x\sh{F}_x)=\dim(\sh{F}_x)-1$
$(\mathbf{0},~16.3.4$ and $16.4.6)$; the relation
\[
  \operatorname{prof}(\sh{F}_x|f_x\sh{F}_x)\geqslant\inf(n-1,\dim(\sh{F}_x|f_x\sh{F}_x))
\]
is therefore equivalent to
\[
  \operatorname{prof}(\sh{F}_x)\geqslant\inf(n,\dim(\sh{F}_x)).
\]
As we supposed that $\sh{F}$ satisfies $(\mathrm{S}_{n-1})$ at every point of $\mathrm{X}-\mathrm{Y}$, this finishes the proof.
\end{proof}

\begin{corollary}[6.11.7]
\label{IV.6.11.7}
With the notations being those of \sref{IV.6.11.6}:
\begin{enumerate}[label=(\roman*)]
  \item The set $\mathrm{U}_{\mathrm{S}_1}(\sh{F})$ is open in $\mathrm{X}$.
  \item For the set $\mathrm{U}_{\mathrm{S}_n}(\sh{F})$ to be open, it suffices that every maximal point $x$ of $\operatorname{Supp}(\sh{F})$,
belonging to $\mathrm{U}_{\mathrm{S}_n}(\sh{F})$, is interior to $\mathrm{U}_{\mathrm{S}_n}(\sh{F})$.
\end{enumerate}
\end{corollary}

\begin{proof}
Assertion (i) was proved in the course of the proof of \sref{IV.6.11.6};
on the other hand, for $n=2$ the case $2^\circ$ of the proof of \sref{IV.6.11.6} is valid without any
hypothesis on $\mathrm{X}$, since (with the same notations) $\mathrm{U}_{\mathrm{S}_1}(\sh{F})$ and $\mathrm{U}_{\mathrm{S}_1}(\sh{F}/f\sh{F})$ are open
in $\mathrm{X}$. As for the case $1^\circ$ of this proof, the hypothesis ensures precisely that it
is unnecessary to consider it.
\end{proof}

\oldpage[IV-2]{162}

\begin{proposition}[6.11.8]
\label{IV.6.11.8}
Let $\mathrm{X}$ be a locally Noetherian prescheme satisfying the following property:
\begin{itemize}
  \item[\textnormal{(CMU)}] Every closed integral sub-prescheme $\mathrm{Y}$ of $\mathrm{X}$ contains an open non-empty $\mathrm{W}$ such that the prescheme induced by $\mathrm{X}$ on $\mathrm{W}$
is a Cohen-Macaulay prescheme.
\end{itemize}
Then, for every coherent $\sh{O}_\mathrm{X}$-Module $\sh{F}$, the function $x\mapsto\operatorname{coprof}(\sh{F}_x)$ is locally constructible and upper
semi-continuous; the sets $\mathrm{U}_{\mathrm{C}_n}(\sh{F})$ and $\mathrm{U}_{\mathrm{S}_n}(\sh{F})$ are open in $\mathrm{X}$.
\end{proposition}

\begin{proof}
Indeed, let $\mathrm{Y}$ be a closed integral sub-prescheme of $\mathrm{X}$, of generic point $y$; by virtue of $(6.10.6)$, there is a
neighbourhood $\mathrm{V}$ of $y$ in $\mathrm{Y}$ such that, for every $x\in\mathrm{V}\cap\mathrm{Y}$, we have
\begin{equation}
\tag{6.11.8.1}
\label{IV.6.11.8.1}
  \operatorname{coprof}(\sh{F}_x)=\operatorname{coprof}(\sh{F}_y)+\operatorname{coprof}(\sh{O}_{\mathrm{Y},x}).
\end{equation}

But by hypothesis there exists an open non-empty $\mathrm{W}$ of $\mathrm{Y}$ such that, for $x\in\mathrm{V}\cap\mathrm{W}$, $\operatorname{coprof}(\sh{O}_{\mathrm{Y},x})=0$,
hence $\operatorname{coprof}(\sh{F}_x)$ is constant in a neighbourhood of $y$ in $\mathrm{Y}$, which proves that the function
$x\mapsto\operatorname{coprof}(\sh{F}_x)$ is
locally constructible $(\mathbf{0}_{\mathrm{III}},~9.3.2)$; it follows furthermore from \sref{IV.6.11.5} and from $(\mathbf{0}_{\mathrm{III}},~9.3.4)$ that this function
is upper semi-continuous. The last assertion follows from \sref{IV.6.11.4},~(i), or from \sref{IV.6.11.6}.
\end{proof}

\begin{remarks}[6.11.9]
\label{IV.6.11.9}
(i) If $\mathrm{X}$ satisfies (CMU), then every prescheme $\mathrm{X}'$ locally of finite type over $\mathrm{X}$ also satisfies (CMU). Let $\mathrm{Y}'$ be a closed integral subprescheme of $\mathrm{X}'$, let $y'$ be its generic point, put $y=u(y')$, and let $\mathrm{Y}$ be the integral closed subprescheme of $\mathrm{X}$ with underlying space $\overline{\{y\}}$. Then $u|_{\mathrm{Y}'}$ factors as
\[
  \mathrm{Y}'\xrightarrow{v}\mathrm{Y}\xrightarrow{j}\mathrm{X},
\]
where $j$ is the canonical injection and $v$ is locally of finite type.

Restricting to affine open neighbourhoods of $y$ and $y'$, we may suppose that $\mathrm{X}=\operatorname{Spec}(\mathrm{A})$, where $\mathrm{A}$ is an integral Cohen--Macaulay ring, and $\mathrm{X}'=\operatorname{Spec}(\mathrm{A}')$, where $\mathrm{A}'$ is an integral finite-type $\mathrm{A}$-algebra containing $\mathrm{A}$.
\oldpage[IV-2]{163}
After replacing $\mathrm{A}$ by $\mathrm{A}_f$ for some $f\neq 0$, we may suppose that $\mathrm{A}'$ contains a polynomial ring $\mathrm{A}[\mathrm{T}_1,\ldots,\mathrm{T}_n]=\mathrm{A}''$ and is finite over $\mathrm{A}''$. The ring $\mathrm{A}''$ is Cohen--Macaulay \sref{IV.6.3.6}; hence we may reduce to the case where $\mathrm{A}'$ is finite over $\mathrm{A}$. After one further localisation, $\mathrm{A}'$ is a finite free $\mathrm{A}$-module. It is therefore a Cohen--Macaulay $\mathrm{A}$-module \sref[0]{0.16.5.1}; since it is finite over $\mathrm{A}$, it is also a Cohen--Macaulay $\mathrm{A}'$-module \sref[0]{0.16.5.3}, and hence a Cohen--Macaulay ring.

(ii) Suppose that there exists a coherent $\sh{O}_\mathrm{X}$-Module $\sh{F}$ such that $\operatorname{Supp}(\sh{F})=\mathrm{X}$ and that $\sh{F}$ is an $\sh{O}_\mathrm{X}$-Module
\emph{of Cohen-Macaulay}. Then $\mathrm{X}$ satisfies the condition (CMU): indeed, with the notations of the proof of \sref{IV.6.11.8},
the relation \sref{IV.6.11.8.1} shows that $\operatorname{coprof}(\sh{O}_{\mathrm{Y},x})=0$ in a neighbourhood (with respect to $\mathrm{Y}$) of the
generic point $y$ of $\mathrm{Y}$.

It is unknown whether there exist locally Noetherian preschemes of dimension $\geqslant 2$ that do not satisfy (CMU)
(if $\dim(\mathrm{X})=1$, it is immediate that every maximal point of $\mathrm{X}_{\mathrm{red}}$ admits an open neighbourhood of dimension $1$,
hence a Cohen-Macaulay ring).
\end{remarks}

\subsection{Nagata's criteria for \texorpdfstring{$\operatorname{Reg}(\mathrm{X})$}{Reg(X)} to be open}
\label{subsection:IV.6.12}
\label{IV.6.12}

\begin{env}[6.12.1]
\label{IV.6.12.1}
Given a locally Noetherian prescheme $\mathrm{X}$, we call the \emph{singular
locus} of $\mathrm{X}$ and denote by $\operatorname{Sing}(\mathrm{X})$ the set of points $x\in\mathrm{X}$ such that $\mathrm{X}$ is not
regular at the point $x$ (in other words, such that the local ring $\sh{O}_x$ is not regular); the
complement $\mathrm{X}-\operatorname{Sing}(\mathrm{X})$, set of $x\in\mathrm{X}$ where $\mathrm{X}$ is regular, is denoted $\operatorname{Reg}(\mathrm{X})$.
We propose in this subsection to give conditions ensuring that
$\operatorname{Sing}(\mathrm{X})$ is \emph{closed} (i.e.\ $\operatorname{Reg}(\mathrm{X})$ \emph{open}).
\end{env}

\begin{proposition}[6.12.2]
\label{IV.6.12.2}
Let $\mathrm{X}$ be a locally Noetherian prescheme. The following conditions
are equivalent:
\begin{enumerate}[label=\emph{\alph*)}]
  \item $\operatorname{Reg}(\mathrm{X})$ is open in $\mathrm{X}$.
  \item For every $x\in\operatorname{Reg}(\mathrm{X})$, there exists an open non-empty part of $\overline{\{x\}}$ contained in $\operatorname{Reg}(\mathrm{X})$.
\end{enumerate}
Moreover, these conditions are implied by the following:
\begin{enumerate}[label=\emph{\alph*)}]
\setcounter{enumi}{2}
  \item For every $x\in\operatorname{Reg}(\mathrm{X})$, if we denote by $\mathrm{Y}$ the reduced closed sub-prescheme of $\mathrm{X}$ having $\overline{\{x\}}$
as underlying space, $\operatorname{Reg}(\mathrm{Y})$ is a neighbourhood of $x$ in $\mathrm{Y}$.
\end{enumerate}
\end{proposition}

\begin{proof}
The equivalence of $a)$ and $b)$ follows from the fact that $\operatorname{Reg}(\mathrm{X})$ is stable by generization
$(\mathbf{0},~17.3.2)$ and from $(\mathbf{0}_{\mathrm{III}},~9.2.6)$. To prove that $c)$ implies $b)$, one can restrict to the
case where $\mathrm{X}=\operatorname{Spec}(\mathrm{A})$ is the spectrum of an affine ring of $x$; if $(t_i)_{1\leqslant i\leqslant n}$ is a regular
system of parameters $(\mathbf{0},~17.1.6)$ of the regular local ring $\mathrm{A}_x$, one can suppose
(replacing $\mathrm{X}$ by an open neighbourhood of $x$ if necessary) that $t_i=(s_i)_x$ for every $i$,
where $s_i\in\mathrm{A}$ and where the family $(s_i)_{1\leqslant i\leqslant n}$ is $\mathrm{A}$-regular $(\mathbf{0},~15.2.4)$. One has then $\mathrm{Y}=\operatorname{Spec}(\mathrm{A}/\mathfrak{p})$,
where $\mathfrak{p}=\mathfrak{j}_x$; as the $t_i$ generate the maximal ideal $\mathfrak{m}=\mathfrak{p}\mathrm{A}_x$ of $\mathrm{A}_x$, one can further
suppose (replacing $\mathrm{X}$ by a smaller neighbourhood of $x$) that the $s_i$ generate $\mathfrak{p}$.
For every $y\in\mathrm{Y}$, the $(s_i)_y$ generate therefore $\mathfrak{p}_y$; as they form a suite $\mathrm{A}_y$-regular,
they also form a suite $\operatorname{gr}^0_{\mathfrak{p}_y}(\sh{O}^{\mathrm{Y}}_y)$-regular, so they form a
$\operatorname{gr}_{\mathfrak{p}_y}(\sh{O}^{\mathrm{Y}}_y)$-regular sequence and it is thus also $\sh{O}^{\mathrm{Y}}_y$-regular,
whence the conclusion.
\end{proof}

\begin{corollary}[6.12.3]
\label{IV.6.12.3}
Let $\mathrm{X}$ be a locally Noetherian prescheme. The following conditions
are equivalent:
\begin{enumerate}[label=\emph{\alph*)}]
  \item For every sub-prescheme $\mathrm{Y}$ of $\mathrm{X}$, $\operatorname{Reg}(\mathrm{Y})$ is open in $\mathrm{Y}$.
\oldpage[IV-2]{164}
  \item For every closed integral sub-prescheme $\mathrm{Y}$ of $\mathrm{X}$, $\operatorname{Reg}(\mathrm{Y})$ contains an open non-empty
part of $\mathrm{Y}$.
\end{enumerate}
\end{corollary}

\begin{proof}
It is clear that $a)$ implies $b)$. Conversely, to see that $b)$ implies $a)$, consider
a closed integral sub-prescheme $\mathrm{Z}$ of $\mathrm{Y}$ with generic point $z$; if $\mathrm{Y}'$ is the
integral closed sub-prescheme of $\mathrm{X}$ whose underlying space is the closure $\mathrm{Y}'=\overline{\{z\}}$ of $z$
in $\mathrm{X}$, $\mathrm{Z}$ is open in $\mathrm{Y}'$ and the sub-prescheme $\mathrm{Z}$ of $\mathrm{Y}'$, being reduced, is induced by $\mathrm{Y}'$
on the open subset of $\mathrm{Z}$ in $\mathrm{Y}'$. Now the hypothesis implies that $\operatorname{Reg}(\mathrm{Y}')$ is
a neighbourhood of $z$ in $\mathrm{Y}'$, hence $\operatorname{Reg}(\mathrm{Z})$ is a neighbourhood of $z$ in $\mathrm{Z}$; it suffices then
to apply \sref{IV.6.12.2} by replacing $\mathrm{X}$ by $\mathrm{Y}$ and $\mathrm{Y}$ by $\mathrm{Z}$.
\end{proof}

\begin{theorem}[6.12.4]
\label{IV.6.12.4}
\textnormal{(Nagata).} --- Let $\mathrm{A}$ be a Noetherian ring, $\mathrm{X}=\operatorname{Spec}(\mathrm{A})$. The
following conditions are equivalent:
\begin{enumerate}[label=\emph{\alph*)}]
  \item For every prescheme $\mathrm{X}'$ locally of finite type over $\mathrm{X}$, $\operatorname{Reg}(\mathrm{X}')$ is open in $\mathrm{X}'$.
  \item For every $\mathrm{A}$-algebra $\mathrm{A}'$ finite and integral, there exists an open non-empty subset of
$\mathrm{X}'=\operatorname{Spec}(\mathrm{A}')$ contained in $\operatorname{Reg}(\mathrm{X}')$.
  \item For every prime ideal $\mathfrak{p}$ of $\mathrm{A}$ and every finite radicial extension $\mathrm{K}'$ of the field of fractions $\mathrm{K}$
of $\mathrm{A}/\mathfrak{p}$, there exists a sub-$\mathrm{A}$-algebra $\mathrm{A}'$ of $\mathrm{K}'$, finite over $\mathrm{A}$, having $\mathrm{K}'$ as field of fractions, and
such that there exists in $\mathrm{X}'=\operatorname{Spec}(\mathrm{A}')$ an open non-empty subset contained in $\operatorname{Reg}(\mathrm{X}')$.
\end{enumerate}
\end{theorem}

\begin{proof}
It is clear that $a)$ implies $b)$. To see that $b)$ implies $c)$, it suffices to note
that one can take as generators of the extension $\mathrm{K}'$ of $\mathrm{K}$ elements integral over $\mathrm{A}/\mathfrak{p}$
(and \emph{a fortiori} over $\mathrm{A}$), and as these elements are finite in number, they generate a finite
$\mathrm{A}$-algebra $\mathrm{A}'$ whose field of fractions is $\mathrm{K}'$; one can then apply $b)$ to $\mathrm{A}'$.
It remains therefore to prove that $c)$ implies $a)$. The question being local on $\mathrm{X}'$, one can
first suppose that $\mathrm{X}'=\operatorname{Spec}(\mathrm{A}')$, where $\mathrm{A}'$ is an $\mathrm{A}$-algebra of finite type; in view of
\sref{IV.6.12.2}, it suffices to prove that for every closed integral sub-prescheme $\mathrm{Y}'$
of $\mathrm{X}'$, $\operatorname{Reg}(\mathrm{Y}')$ contains an open non-empty part of $\mathrm{Y}'$. In other words, one can restrict
to proving that if $\mathrm{A}'$ is an $\mathrm{A}$-algebra integral \emph{of finite type} and $\mathrm{X}'=\operatorname{Spec}(\mathrm{A}')$,
$\operatorname{Reg}(\mathrm{X}')$ contains an open non-empty part. Let $\mathrm{K}'$ be the field of fractions of $\mathrm{A}'$;
if $\mathfrak{p}$ is the canonical image in $\mathrm{X}$ of the generic point of $\mathrm{X}'$, $\mathrm{K}'$ is an extension of the
field of fractions $\mathrm{K}$ of $\mathrm{A}/\mathfrak{p}$, and $\mathrm{A}/\mathfrak{p}$ is identified with a subring of $\mathrm{A}'$, $\mathrm{A}'$ being a finite
$(\mathrm{A}/\mathfrak{p})$-algebra of finite type. One can obviously restrict in what follows to the proof
when $\mathfrak{p}=0$. We distinguish two cases:

I) $\mathrm{K}'$ \emph{is a separable extension} of $\mathrm{K}$. --- One is then reduced to proving the

\begin{lemma}[6.12.4.1]
\label{IV.6.12.4.1}
Let $\mathrm{A}$ be a Noetherian integral ring, $\mathrm{A}'$ an $\mathrm{A}$-algebra integral
of finite type, such that the field of fractions $\mathrm{K}'$ of $\mathrm{A}'$ is a separable extension of the
field of fractions $\mathrm{K}$ of $\mathrm{A}$. If $\operatorname{Spec}(\mathrm{A})$ contains an open non-empty subset consisting of regular points, then
so does $\operatorname{Spec}(\mathrm{A}')$.
\end{lemma}

\begin{proof}
Replacing $\mathrm{A}$ by a ring of fractions $\mathrm{A}_f$, so that $\mathrm{D}(f)\subset\operatorname{Reg}(\operatorname{Spec}(\mathrm{A}))$,
one can already suppose the ring $\mathrm{A}$ \emph{regular}. By hypothesis there exists in $\mathrm{K}'$ a system
$(t_i)_{1\leqslant i\leqslant n}$ of elements algebraically independent over $\mathrm{K}$ and such that $\mathrm{K}'$ is a
finite separable algebraic extension of $\mathrm{K}(t_1,\ldots,t_n)$; considering a finite system
\oldpage[IV-2]{165}
of generators $t'_j$ $(1\leqslant j\leqslant m)$ of $\mathrm{K}'$ over $\mathrm{K}(t_1,\ldots,t_n)$, which one can suppose to be integral
over $\mathrm{A}_1=\mathrm{A}[t_1,\ldots,t_n]$, one sees that $\mathrm{A}'_1=\mathrm{A}_1[t'_1,\ldots,t'_m]$ is \emph{finite} over $\mathrm{A}_1$ and has $\mathrm{K}'$ as
field of fractions. If one sets $\mathrm{X}'=\operatorname{Spec}(\mathrm{A}')$, $\mathrm{X}'_1=\operatorname{Spec}(\mathrm{A}'_1)$, it follows from the fact that the fields
of rational functions $\mathrm{R}(\mathrm{X}')$ and $\mathrm{R}(\mathrm{X}'_1)$ are both isomorphic to $\mathrm{K}'$ and from the fact that $\mathrm{X}'$
and $\mathrm{X}'_1$ are $\mathrm{A}$-preschemes of finite type, that there exists an open subset $\mathrm{U}'\subset\mathrm{X}'$ and an open
$\mathrm{U}'_1\subset\mathrm{X}'_1$ which are $\mathrm{A}$-\emph{isomorphic} $(\mathbf{I},~6.5.5)$. One is therefore reduced to proving that $\operatorname{Reg}(\mathrm{X}'_1)$
contains a non-empty open subset, in other words one can suppose that $\mathrm{A}'$ is an $\mathrm{A}_1$-algebra
\emph{finite}. Now, one knows $(\mathbf{0},~17.3.7)$ that $\mathrm{A}_1$ is a regular ring, and one can therefore restrict to the
case where $\mathrm{A}'$ is an $\mathrm{A}$-algebra \emph{finite} and $\mathrm{K}'$ an extension \emph{separable} and \emph{finite} of $\mathrm{K}$. Let $\xi$ be the
generic point of $\mathrm{X}$, $\mathrm{A}'_\xi=\mathrm{K}'$ is then a free module over $\mathrm{A}_\xi=\mathrm{K}$, hence (Bourbaki, \emph{Alg.},
\emph{comm.}, chap.~II, \S~5, no~1, cor.\ of prop.~2) one can, replacing $\mathrm{A}$ by a ring of fractions $\mathrm{A}_f$,
suppose that $\mathrm{A}'$ is a free $\mathrm{A}$-module of finite type. Let then $(x_h)_{1\leqslant h\leqslant r}$ be a basis of
this $\mathrm{A}$-module, and set
\[
  d=\det(\operatorname{Tr}_{\mathrm{A}'/\mathrm{A}}(x_h x_k))=\det(\operatorname{Tr}_{\mathrm{K}'/\mathrm{K}}(x_h x_k))\in\mathrm{A}.
\]

As $\mathrm{K}'$ is separable over $\mathrm{K}$, one knows (Bourbaki, \emph{Alg.}, chap.~IX, \S~2, prop.~5)
that $d\neq 0$; replacing $\mathrm{A}$ by a ring of fractions $\mathrm{A}_z$, one can therefore suppose $d$
\emph{invertible} in $\mathrm{A}$. But then, for every $z\in\operatorname{Spec}(\mathrm{A})$, if one denotes by $\overline{x}_h$ $(1\leqslant h\leqslant r)$
the canonical image of $x_h$ in $\mathrm{A}'(z)=\mathrm{A}'\otimes_\mathrm{A}\boldsymbol{k}(z)$, one has $\det(\operatorname{Tr}_{\mathrm{A}'(z)/\boldsymbol{k}(z)}(\overline{x}_h\overline{x}_k))=\overline{d}$, canonical
image of $d$ in $\boldsymbol{k}(z)=\mathrm{A}_z/\mathfrak{m}_z$, and as $\overline{d}$ is invertible (hence $\neq 0$) in $\boldsymbol{k}(z)$, one
knows (\emph{loc.~cit.}) that $\mathrm{A}'(z)$ is a $\boldsymbol{k}(z)$-algebra \emph{separable}, hence a direct product of finite separable field
extensions of $\boldsymbol{k}(z)$. Since such an algebra is a \emph{regular} ring, one sees
that the morphism $g:\mathrm{X}'\to\mathrm{X}$ is \emph{flat} and that its fibres $g^{-1}(z)$ are \emph{regular} for every
$z\in\mathrm{X}$; one concludes then from $(6.5.2,~\text{(ii)})$ that $\mathrm{X}'$ is \emph{regular}, which finishes the proof in this case.
\end{proof}

II) \emph{General case.} --- As $\mathrm{A}'$ is an $\mathrm{A}$-module without torsion, $\mathrm{A}'\otimes_\mathrm{A}\mathrm{K}$ is identified with a
subring of $\mathrm{K}'$, hence $\mathrm{X}''=\operatorname{Spec}(\mathrm{A}'\otimes_\mathrm{A}\mathrm{K})$ is an integral $\mathrm{K}$-prescheme whose field of
rational functions is $\mathrm{K}'$. One knows $(4.6.6)$ that there exists a finite radicial extension
$\mathrm{K}_1$ of $\mathrm{K}$ such that $\mathrm{X}'_1=\operatorname{Spec}(\mathrm{A}'\otimes_\mathrm{A}\mathrm{K}_1)=\mathrm{X}''\otimes_\mathrm{K}\mathrm{K}_1$, $(\mathrm{X}'_1)_{\mathrm{red}}$ is a $\mathrm{K}_1$-prescheme \emph{geometrically
reduced} of finite type; furthermore, the morphism $\operatorname{Spec}(\mathrm{K}_1)\to\operatorname{Spec}(\mathrm{K})$ being radicial,
finite and surjective, it is a universal homeomorphism $(2.4.5)$, hence $\mathrm{X}'_1$ is homeomorphic
to $\mathrm{X}''$, and consequently $(\mathrm{X}'_1)_{\mathrm{red}}$ is integral; its field of rational functions $\mathrm{K}'_1$
is a finite radicial extension of $\mathrm{K}'$ and a \emph{separable} extension of finite type of $\mathrm{K}_1$ $(4.6.1)$.
By virtue of hypothesis $c)$ of the statement, there is a sub-$\mathrm{A}$-algebra \emph{finite} $\mathrm{A}_1$ of $\mathrm{K}_1$, having $\mathrm{K}_1$
as field of fractions and such that if one sets $\mathrm{X}_1=\operatorname{Spec}(\mathrm{A}_1)$, $\operatorname{Reg}(\mathrm{X}_1)$ contains an
open non-empty subset. Let $\mathrm{A}'_1$ be the image of the canonical homomorphism $\mathrm{A}'\otimes_\mathrm{A}\mathrm{A}_1\to\mathrm{K}'_1$
and let $\mathrm{X}'_1=\operatorname{Spec}(\mathrm{A}'_1)$; $\mathrm{A}'_1$ is an integral ring which is an $\mathrm{A}'$-algebra \emph{finite} and whose
field of fractions is $\mathrm{K}'_1$; furthermore, as the composite homomorphism
$\mathrm{A}'\to\mathrm{A}'\otimes_\mathrm{A}\mathrm{A}_1\to\mathrm{A}'_1\to\mathrm{K}'_1$ is identical to $\mathrm{A}'\to\mathrm{K}'\to\mathrm{K}'_1$, hence injective, the homomorphism
$\mathrm{A}'\to\mathrm{A}'_1$ is injective; the morphism $g:\mathrm{X}'_1\to\mathrm{X}'$ is therefore finite and surjective $(\mathbf{II},~6.1.10)$.
This being the hypothesis on $\mathrm{A}_1$ and part I) of the proof imply that $\operatorname{Reg}(\mathrm{X}'_1)$
\oldpage[IV-2]{166}
contains an open non-empty $\mathrm{V}'_1$; as $g$ is closed $(\mathbf{II},~6.1.10)$, one can suppose
that $\mathrm{V}'_1=g^{-1}(\mathrm{V}')$, where $\mathrm{V}'$ is an open affine of $\mathrm{X}'$; replacing $\mathrm{A}'$ by the ring of $\mathrm{V}'$
and $\mathrm{A}'_1$ by that of $\mathrm{V}'_1$, one can therefore suppose that $\mathrm{X}'_1$ is \emph{regular}. Furthermore, the same
reasoning as in I) applied to the generic point $\xi'$ of $\mathrm{X}'$ allows (replacing $\mathrm{X}'$
by a neighbourhood of $\xi'$) to suppose that $\mathrm{A}'_1$ is an $\mathrm{A}'$-module \emph{free}, and hence
that the morphism $g$ is \emph{flat}. But it follows then from $(6.5.2,~\text{(i)})$ that $\mathrm{X}'$ is \emph{regular}.
\end{proof}

\begin{corollary}[6.12.5]
\label{IV.6.12.5}
\textnormal{(Zariski).} --- Let $k$ be a field; for every $k$-prescheme $\mathrm{X}$ locally
of finite type over $k$, the set of $x\in\mathrm{X}$ where $\mathrm{X}$ is regular (resp.\ \emph{geometrically regular} over $k$)
is open in $\mathrm{X}$.
\end{corollary}

\begin{proof}
The assertion of the corollary relative to the property of being regular follows from
\sref{IV.6.12.4} by taking $\mathrm{A}=k$. The assertion relative to the property of being geometrically
regular follows already from $(6.8.7)$; one can also deduce it from \sref{IV.6.12.4} in the following
way: set $k'=k^{p^{-\infty}}$ ($p$ being the characteristic exponent of $k$); as the morphism
$\operatorname{Spec}(k')\to\operatorname{Spec}(k)$ is radicial, integral and surjective, it is a universal homeomorphism
$(2.4.5)$, hence the morphism projection $\mathrm{X}\otimes_k k'\to\mathrm{X}$ is a universal homeomorphism.
The projection in $\mathrm{X}$ of $\operatorname{Reg}(\mathrm{X}\otimes_k k')$ is the set of points of $\mathrm{X}$ where $\mathrm{X}$ is geometrically
regular over $k$, by virtue of $(6.7.7,~e))$; hence this set is open from what
precedes.
\end{proof}

\begin{corollary}[6.12.6]
\label{IV.6.12.6}
Let $\mathrm{A}$ be a ring having one of the following properties:
\begin{enumerate}[label=(\roman*)]
  \item $\mathrm{A}$ is a Dedekind ring and its field of fractions $\mathrm{K}$ is of characteristic $0$.
  \item $\mathrm{A}$ is a semi-local Noetherian ring of dimension $\leqslant 1$.
\end{enumerate}
Then, for every prescheme $\mathrm{X}$ locally of finite type over $\mathrm{A}$, $\operatorname{Reg}(\mathrm{X})$ is open in $\mathrm{X}$.
\end{corollary}

\begin{proof}
Let us verify in both cases the condition $c)$ of \sref{IV.6.12.4}. In both cases, a prime ideal
$\mathfrak{p}$ of $\mathrm{A}$ is maximal or minimal; if $\mathfrak{p}$ is maximal, an $(\mathrm{A}/\mathfrak{p})$-algebra finite and
integral is a field and the condition $c)$ of \sref{IV.6.12.4} is trivially satisfied. Suppose therefore
$\mathfrak{p}$ non-maximal, and distinguish the two cases of the statement.

(i) If $\mathrm{K}$ is of characteristic $0$, there is no radicial extension of $\mathrm{K}$ other than $\mathrm{K}$
itself; as a Dedekind ring is regular $(\mathbf{0},~17.1.4)$, the condition $c)$
of \sref{IV.6.12.4} is trivially satisfied.

(ii) One can then suppose $\mathrm{A}$ integral \sref{IV.6.12.2}, let $\mathrm{K}$ be its field of fractions; if $\mathrm{K}'$ is
a finite radicial extension of $\mathrm{K}$, $\mathrm{A}'$ a sub-$\mathrm{A}$-algebra of $\mathrm{K}'$ generated by a
finite system of generators of $\mathrm{K}'$ over $\mathrm{K}$, integral over $\mathrm{A}$, $\mathrm{A}'$ is a semi-local integral ring
of dimension $1$ $(\mathbf{0},~16.1.5)$, and consequently, in $\mathrm{X}'=\operatorname{Spec}(\mathrm{A}')$, the reduced set at the
generic point is open and evidently contained in $\operatorname{Reg}(\mathrm{X}')$, which proves the
condition $c)$ of \sref{IV.6.12.4} in this case.
\end{proof}

This corollary applies notably when $\mathrm{A}=\bb{Z}$.

\begin{theorem}[6.12.7]
\label{IV.6.12.7}
\textnormal{(Nagata).} --- Let $\mathrm{A}$ be a local Noetherian complete ring,
$\mathrm{X}=\operatorname{Spec}(\mathrm{A})$. Then $\operatorname{Reg}(\mathrm{X})$ is open in $\mathrm{X}$.
\end{theorem}

\begin{proof}
In view of \sref{IV.6.12.2}, one is reduced to the case where moreover $\mathrm{A}$ is \emph{integral}, and to
proving that $\operatorname{Reg}(\mathrm{X})$ contains an open non-empty part of $\mathrm{X}$. We distinguish
two cases:

\oldpage[IV-2]{167}

I) \emph{The field of fractions} $\mathrm{K}$ \emph{of} $\mathrm{A}$ \emph{is of characteristic} $0$. --- One knows then $(\mathbf{0},~19.8.8)$
that there exists a complete discrete valuation ring $\mathrm{C}$ and a subring $\mathrm{B}$ of $\mathrm{A}$ such
that $\mathrm{A}$ is a $\mathrm{B}$-algebra \emph{finite} and that $\mathrm{B}$ is isomorphic to a formal power series ring
$\mathrm{C}[[\mathrm{T}_1,\ldots,\mathrm{T}_r]]$. As $\mathrm{C}$ is regular $(\mathbf{II},~7.1.6)$, so is $\mathrm{B}$ $(\mathbf{0},~17.3.8)$;
furthermore, the field of fractions $\mathrm{L}$ of $\mathrm{B}$ being of characteristic $0$, $\mathrm{K}$ is a \emph{separable} extension
finite of $\mathrm{L}$, hence one can apply \sref{IV.6.12.4.1} to $\mathrm{B}$, and this proves the
proposition.

II) \emph{The field of fractions} $\mathrm{K}$ \emph{of} $\mathrm{A}$ \emph{is of characteristic} $p>0$. --- Then $\mathrm{A}$ contains the
prime field $\mathbf{F}_p$, and the theorem has been proved in this case $(\mathbf{0},~22.7.6)$.
\end{proof}

\begin{corollary}[6.12.8]
\label{IV.6.12.8}
Let $\mathrm{A}$ be a local Noetherian complete ring. Then, for every
prescheme $\mathrm{X}$ locally of finite type over $\mathrm{A}$, $\operatorname{Reg}(\mathrm{X})$ is open in $\mathrm{X}$.
\end{corollary}

\begin{proof}
Let us verify the condition $c)$ of \sref{IV.6.12.4}. If $\mathfrak{p}$ is prime in $\mathrm{A}$, $\mathrm{A}/\mathfrak{p}$ is again a
local Noetherian complete ring; if $\mathrm{K}'$ is a finite extension of the field of fractions $\mathrm{K}$
of $\mathrm{A}/\mathfrak{p}$, $\mathrm{K}'$ is the field of fractions of a sub-$\mathrm{A}$-algebra finite $\mathrm{A}'$ of $\mathrm{K}'$, generated by
a finite system of generators of $\mathrm{K}'$ over $\mathrm{K}$, integral over $\mathrm{A}$. One knows then that $\mathrm{A}'$ is a
semi-local complete ring (Bourbaki, \emph{Alg.~comm.}, chap.~IV, \S~2, no~5, cor.~3 of prop.~9), hence
a product of local complete rings, and as $\mathrm{A}'$ is integral, it is a local
complete ring; and by virtue of \sref{IV.6.12.7}, if $\mathrm{X}'=\operatorname{Spec}(\mathrm{A}')$, $\operatorname{Reg}(\mathrm{X}')$ is open and non-empty,
whence the conclusion.
\end{proof}

\begin{proposition}[6.12.9]
\label{IV.6.12.9}
Let $\mathrm{X}$ be a locally Noetherian prescheme such that $\operatorname{Reg}(\mathrm{X})$ is
open; then, for every $n\geqslant 0$, the set $\mathrm{U}_{\mathrm{R}_n}(\mathrm{X})$ of points $x\in\mathrm{X}$ where $\mathrm{X}$ satisfies $(\mathrm{R}_n)$ is open.
\end{proposition}

\begin{proof}
Indeed $(5.8.2)$, $\mathrm{U}_{\mathrm{R}_n}(\mathrm{X})$ is the union of $\operatorname{Reg}(\mathrm{X})$ and of the set of $x\in\operatorname{Sing}(\mathrm{X})$
such that the generic points $z_i$ of the irreducible components of $\operatorname{Sing}(\mathrm{X})$ containing $x$
satisfy the relation $\dim(\sh{O}_{\mathrm{X},z_i})\geqslant n+1$ $(5.1.2)$; in other words, these
points $x\in\operatorname{Sing}(\mathrm{X})$ are those for which $\operatorname{codim}_x(\operatorname{Sing}(\mathrm{X}),\mathrm{X})\geqslant n+1$;
as the function $x\mapsto\operatorname{codim}_x(\operatorname{Sing}(\mathrm{X}),\mathrm{X})$ is lower semi-continuous $(\mathbf{0},~14.2.6)$,
the set of these points is open in $\operatorname{Sing}(\mathrm{X})$, which finishes the proof.
\end{proof}

Let us also note the following elementary result:

\begin{proposition}[6.12.10]
\label{IV.6.12.10}
Let $\mathrm{X}$ be a locally Noetherian prescheme. The set $\mathrm{U}_{\mathrm{R}_0}(\mathrm{X})$ is
open in $\mathrm{X}$. For $\mathrm{U}_{\mathrm{R}_1}(\mathrm{X})$ to be open in $\mathrm{X}$, it is necessary and it suffices that every maximal $x\in\mathrm{X}$
such that $x\in\mathrm{U}_{\mathrm{R}_1}(\mathrm{X})$ (which means that $\mathrm{X}$ is reduced at $x$) be interior to $\mathrm{U}_{\mathrm{R}_1}(\mathrm{X})$.
\end{proposition}

\begin{proof}
By definition $(5.8.2)$, $\mathrm{U}_{\mathrm{R}_0}(\mathrm{X})$ is the set $\mathrm{X}-\bigcup_\alpha\mathrm{X}'_\alpha$, where $\mathrm{X}'_\alpha$ ranges over the set of
irreducible components of $\mathrm{X}$ such that $\mathrm{X}$ is not reduced at their generic point;
as the set of irreducible components of $\mathrm{X}$ is locally finite, $\mathrm{U}_{\mathrm{R}_0}(\mathrm{X})$ is
open. As for $\mathrm{U}_{\mathrm{R}_1}(\mathrm{X})$, the stated condition is trivially necessary;
let us prove that it is sufficient. Let $\mathrm{X}''_\beta$ be the irreducible components of $\mathrm{X}$ such that,
at their generic point $x''_\beta$, the prescheme $\mathrm{X}$ is reduced; by hypothesis there exists an open
$\mathrm{U}\subset\mathrm{U}_{\mathrm{R}_1}(\mathrm{X})$ containing all the $x''_\beta$. Denote then by $\mathrm{Z}_\lambda$ the irreducible components
of the closed set $\mathrm{Z}=\mathrm{X}-\mathrm{U}$ whose generic point $z_\lambda$ is such that $\dim(\sh{O}_{\mathrm{X},z_\lambda})\leqslant 1$
and that $\sh{O}_{\mathrm{X},z_\lambda}$ is non-regular. No point belonging to one of the $\mathrm{Z}_\lambda$ can belong
\oldpage[IV-2]{168}
to $\mathrm{U}_{\mathrm{R}_1}(\mathrm{X})$; but conversely, if $x\in\mathrm{Z}$ does not belong to any of the $\mathrm{Z}_\lambda$, then, for every
generization $x'$ of $x$, either $x'$ belongs to $\mathrm{U}$, or
$\dim(\sh{O}_{\mathrm{X},x'})\geqslant 2$, or
$\dim(\sh{O}_{\mathrm{X},x'})=1$, and as $x'$ does not belong to any $\mathrm{Z}_\lambda$, $\overline{\{x'\}}$ is necessarily
one of the irreducible components of $\mathrm{Z}$, of codimension $1$ in $\mathrm{X}$, distinct from the $\mathrm{Z}_\lambda$, hence $\sh{O}_{\mathrm{X},x'}$
is regular by definition. One concludes that $\mathrm{U}_{\mathrm{R}_1}(\mathrm{X})=\mathrm{X}-\bigcup_\lambda\mathrm{Z}_\lambda$; as the set of
$\mathrm{Z}_\lambda$ is locally finite in $\mathrm{Z}$, $\bigcup_\lambda\mathrm{Z}_\lambda$ is closed, which finishes the proof that $\mathrm{U}_{\mathrm{R}_1}(\mathrm{X})$
is open in $\mathrm{X}$.
\end{proof}

\subsection{Criteria for \texorpdfstring{$\operatorname{Nor}(\mathrm{X})$}{Nor(X)} to be open}
\label{subsection:IV.6.13}
\label{IV.6.13}

\begin{env}[6.13.1]
\label{IV.6.13.1}
Given a locally Noetherian prescheme $\mathrm{X}$, we shall denote by
$\operatorname{Nor}(\mathrm{X})$ the set of $x\in\mathrm{X}$ where $\mathrm{X}$ is \emph{normal}; this set contains $\operatorname{Reg}(\mathrm{X})$ and is
contained in the (open) set of points $x$ such that $\sh{O}_{X,x}$ is integral (i.e.\ the set
of points $x$ where $\mathrm{X}$ is reduced, and which belong to only one irreducible component
of $\mathrm{X}$).
\end{env}

\begin{proposition}[6.13.2]
\label{IV.6.13.2}
Let $\mathrm{X}$ be a locally Noetherian prescheme reduced, $\mathrm{X}'$ its
normalization $(\mathbf{II},~6.3.8)$. If the canonical morphism $f:\mathrm{X}'\to\mathrm{X}$ is finite, $\operatorname{Nor}(\mathrm{X})$ is open in $\mathrm{X}$.
\end{proposition}

\begin{proof}
Indeed, on $\mathrm{X}'=\operatorname{Spec}(f_*(\sh{O}_{\mathrm{X}'}))$ and $f_*(\sh{O}_{\mathrm{X}'})$ is a coherent $\sh{O}_\mathrm{X}$-Algebra by
hypothesis $(\mathbf{II},~6.1.3)$. To say that $\mathrm{X}$ is normal at a point $x$ means that the canonical
homomorphism $\sh{O}_x\to(f_*(\sh{O}_{\mathrm{X}'}))_x$ is bijective $(\mathbf{II},~6.3.4)$; but the set of these points is \emph{open}
since $\sh{O}_\mathrm{X}$ and $f_*(\sh{O}_{\mathrm{X}'})$ are coherent $(\mathbf{0}_\mathrm{I},~5.2.7)$.
\end{proof}

\begin{corollary}[6.13.3]
\label{IV.6.13.3}
If $\mathrm{A}$ is an integral Noetherian \emph{Japanese} ring, $\operatorname{Nor}(\operatorname{Spec}(\mathrm{A}))$
is open in $\operatorname{Spec}(\mathrm{A})$.
\end{corollary}

\begin{proposition}[6.13.4]
\label{IV.6.13.4}
Let $\mathrm{X}$ be a locally Noetherian prescheme. For $\operatorname{Nor}(\mathrm{X})$
to be open, it is necessary and it suffices that every maximal point $x$ of $\mathrm{X}$ such that $x\in\operatorname{Nor}(\mathrm{X})$ (which means
that $\mathrm{X}$ is reduced at the point $x$) be interior to $\operatorname{Nor}(\mathrm{X})$.
\end{proposition}

\begin{proof}
There is only to prove the sufficiency of the condition. By virtue of the remark made
in \sref{IV.6.13.1}, one can restrict (replacing $\mathrm{X}$ by an open of $\mathrm{X}$ in which $\mathrm{X}$
is reduced), i.e.\ to suppose that the maximal points of $\mathrm{X}$ belong
to $\mathrm{U}_{\mathrm{S}_1}(\mathrm{X})$ and to $\mathrm{U}_{\mathrm{R}_1}(\mathrm{X})$; as $\operatorname{Nor}(\mathrm{X})=\mathrm{U}_{\mathrm{S}_1}(\mathrm{X})\cap\mathrm{U}_{\mathrm{R}_1}(\mathrm{X})$ by virtue of Serre's
criterion \sref{IV.5.8.6}, one deduces from \sref{IV.6.11.7} and \sref{IV.6.12.10} that these two sets are
open; hence $\operatorname{Nor}(\mathrm{X})$ is open in $\mathrm{X}$.
\end{proof}

\begin{corollary}[6.13.5]
\label{IV.6.13.5}
If $\mathrm{X}$ is such that $\operatorname{Reg}(\mathrm{X})$ is open in $\mathrm{X}$, then $\operatorname{Nor}(\mathrm{X})$ is
open in $\mathrm{X}$.
\end{corollary}

\begin{proof}
Indeed, every maximal point $x$ where $\mathrm{X}$ is reduced (hence $\sh{O}_{\mathrm{X},x}$ is a field) belongs
to $\operatorname{Reg}(\mathrm{X})$, hence is interior to $\operatorname{Reg}(\mathrm{X})$ by hypothesis, and \emph{a fortiori} is interior to $\operatorname{Nor}(\mathrm{X})$.
\end{proof}

\begin{proposition}[6.13.6]
\label{IV.6.13.6}
\textnormal{(Nagata).} --- Let $\mathrm{A}$ be a Noetherian integral ring, $\mathrm{K}$ its field of
fractions, $\mathrm{K}'$ a finite extension of $\mathrm{K}$, $\mathrm{A}'$ the integral closure of $\mathrm{A}$ in $\mathrm{K}'$. For $\mathrm{A}'$
to be a finite $\mathrm{A}$-algebra, it is necessary and it suffices that the two following conditions be satisfied:
\begin{enumerate}[label=(\roman*)]
  \item There exists $f\neq 0$ in $\mathrm{A}$ such that the integral closure $\mathrm{A}'_f$ of the ring of fractions $\mathrm{A}_f$
in $\mathrm{K}'$ is a finite $\mathrm{A}_f$-algebra.
  \item For every prime ideal $\mathfrak{p}\in\operatorname{Spec}(\mathrm{A})$, the integral closure $\mathrm{A}'_\mathfrak{p}$ of the local ring $\mathrm{A}_\mathfrak{p}$
in $\mathrm{K}'$ is a finite $\mathrm{A}_\mathfrak{p}$-algebra.
\end{enumerate}
\end{proposition}

\begin{proof}
The conditions are necessary, for for every multiplicative subset $\mathrm{S}$ of $\mathrm{A}$,
$\mathrm{S}^{-1}\mathrm{A}'$ is the integral closure of $\mathrm{S}^{-1}\mathrm{A}$ in $\mathrm{K}'$ and is by hypothesis a $\mathrm{S}^{-1}\mathrm{A}$-module
of finite type. To see that the conditions are sufficient, consider the filtered increasing
family $(\mathrm{B}_\lambda)$ of sub-$\mathrm{A}$-algebras of $\mathrm{K}'$ which are $\mathrm{A}$-algebras finite and have $\mathrm{K}'$ as
field of fractions. Set $\mathrm{Y}=\operatorname{Spec}(\mathrm{A})$, $\mathrm{X}_\lambda=\operatorname{Spec}(\mathrm{B}_\lambda)$, denote by $u_\lambda$ the
morphism $\mathrm{X}_\lambda\to\mathrm{Y}$ and let $\mathrm{S}_\lambda=\mathrm{X}_\lambda-\operatorname{Nor}(\mathrm{X}_\lambda)$, $\mathrm{T}_\lambda=u_\lambda(\mathrm{S}_\lambda)$. One can restrict to the $\mathrm{B}_\lambda$
which contain a finite set whose image in $\mathrm{A}'_f$ is a system of generators of the
$\mathrm{A}_f$-module $\mathrm{A}'_f$; this implies, by virtue of hypothesis (i), that $(\mathrm{B}_\lambda)_f=\mathrm{A}'_f$ for every $\lambda$,
or equivalently that $u_\lambda^{-1}(\mathrm{D}(f))$ is contained in $\operatorname{Nor}(\mathrm{X}_\lambda)$. By virtue of \sref{IV.6.13.2}, $\mathrm{S}_\lambda$ is therefore
\emph{closed} in $\mathrm{Y}$. But for every
$\mathfrak{p}\in\operatorname{Spec}(\mathrm{A})$, it exists by virtue of (ii) a $\lambda$ such that $(\mathrm{B}_\lambda)_\mathfrak{p}=\mathrm{A}'_\mathfrak{p}$, and consequently all points
\oldpage[IV-2]{169}
of $\mathrm{X}_\lambda$ lying over $\mathfrak{p}$ belong to $\operatorname{Nor}(\mathrm{X}_\lambda)$; in other words $\mathfrak{p}\notin\mathrm{T}_\lambda$; as
$\mathfrak{p}\in\operatorname{Spec}(\mathrm{A})$, it exists by virtue of (ii) a $\lambda$ such that $(\mathrm{B}_\lambda)_\mathfrak{p}=\mathrm{A}'_\mathfrak{p}$, and consequently all points
of $\mathrm{X}_\lambda$ lying over $\mathfrak{p}$ belong to $\operatorname{Nor}(\mathrm{X}_\lambda)$; in other words $\mathfrak{p}\notin\mathrm{T}_\lambda$; as
$\mathrm{Y}$ is Noetherian and the $\mathrm{T}_\lambda$ are closed, there exists a $\lambda_0$ such that $\mathrm{T}_{\lambda_0}=\varnothing$; hence $\mathrm{B}_{\lambda_0}$
is integrally closed, and as its field of fractions is $\mathrm{K}'$, one has $\mathrm{B}_{\lambda_0}=\mathrm{A}'$.
\end{proof}

\begin{proposition}[6.13.7]
\label{IV.6.13.7}
Let $\mathrm{A}$ be a Noetherian ring, $\mathrm{X}=\operatorname{Spec}(\mathrm{A})$. The following conditions are equivalent:
\begin{enumerate}[label=\emph{\alph*)}]
  \item For every prescheme $\mathrm{X}'$ locally of finite type over $\mathrm{X}$, $\operatorname{Nor}(\mathrm{X}')$ is open in $\mathrm{X}'$.
  \item For every $\mathrm{A}$-algebra finite and integral $\mathrm{A}'$, $\operatorname{Nor}(\operatorname{Spec}(\mathrm{A}'))$ is open in $\operatorname{Spec}(\mathrm{A}')$.
  \item For every prime ideal $\mathfrak{p}$ of $\mathrm{A}$ and every finite radicial extension $\mathrm{K}'$ of the field of fractions $\mathrm{K}$ of $\mathrm{A}/\mathfrak{p}$, there exists a
sub-$\mathrm{A}$-algebra $\mathrm{A}'$ of $\mathrm{K}'$, finite over $\mathrm{A}$, having $\mathrm{K}'$ as field of fractions and such that $\operatorname{Nor}(\operatorname{Spec}(\mathrm{A}'))$ is open in $\operatorname{Spec}(\mathrm{A}')$.
\end{enumerate}
\end{proposition}

\begin{proof}
The fact that $a)$ implies $b)$ and that $b)$ implies $c)$ is proved as in \sref{IV.6.12.4}. To show that $c)$
implies $a)$, one reduces as in \sref{IV.6.12.4} to proving that if $\mathrm{A}'$ is an $\mathrm{A}$-algebra
integral and of finite type, the generic point of $\operatorname{Spec}(\mathrm{A}')$ is interior to $\operatorname{Nor}(\operatorname{Spec}(\mathrm{A}'))$. One distinguishes two cases
as in \sref{IV.6.12.4} by proving first the

\begin{lemma}[6.13.7.1]
\label{IV.6.13.7.1}
Let $\mathrm{A}$ be a Noetherian integral ring, $\mathrm{A}'$ an $\mathrm{A}$-algebra integral of finite type, containing $\mathrm{A}$ and such
that the field of fractions $\mathrm{K}'$ of $\mathrm{A}'$ is a separable extension of the field of fractions $\mathrm{K}$ of $\mathrm{A}$. If $\operatorname{Nor}(\operatorname{Spec}(\mathrm{A}))$ is open in
$\operatorname{Spec}(\mathrm{A})$, $\operatorname{Nor}(\operatorname{Spec}(\mathrm{A}'))$ is open in $\operatorname{Spec}(\mathrm{A}')$.
\end{lemma}

\begin{proof}
It suffices to prove (in view of \sref{IV.6.13.2}) that the generic point of $\operatorname{Spec}(\mathrm{A}')$ is interior to
$\operatorname{Nor}(\operatorname{Spec}(\mathrm{A}'))$. The proof follows the same approach as that of \sref{IV.6.12.4.1}, from which we conserve the notations.
We first note that we may assume that $\mathrm{A}$ is integrally closed; then $\mathrm{A}_1=\mathrm{A}[\mathrm{t}_1,\ldots,\mathrm{t}_n]$
is integrally closed (Bourbaki, \emph{Alg.~comm.}, chap.~V, \S~1, no~3, cor.~2 of prop.~13). We next reduce to the
case where $\mathrm{A}'$ is a free $\mathrm{A}$-module of \emph{finite type}; the reasoning of \sref{IV.6.12.4.1} then shows
(replacing $\mathrm{A}$, if necessary, by a ring $\mathrm{A}_f$ with $f\neq 0$) that the fibres $g^{-1}(z)$
of the morphism $g:\mathrm{X}'\to\mathrm{X}$ are regular and \emph{a fortiori} normal.
Furthermore $g$ is flat and $\mathrm{X}$ is normal, hence $(6.5.4,~\text{(ii)})$ $\mathrm{X}'$ is normal.
\end{proof}

This lemma being proved, one passes to the general case as in $(\text{6.12.4; II})$, from which we conserve the
notations; applying hypothesis $c)$, one sees this time that $\operatorname{Nor}(\mathrm{X}'_1)$ is open and one reduces thus to the case
where $\mathrm{X}'_1$ is normal and $g:\mathrm{X}'_1\to\mathrm{X}'$ is flat and surjective; one concludes by applying $(6.5.4,~\text{(i)})$.
\end{proof}

\subsection{Base change and integral closure}
\label{subsection:IV.6.14}

\begin{proposition}[6.14.1]
\label{IV.6.14.1}
Let $\mathrm{Y}$, $\mathrm{Y}'$ be two locally Noetherian preschemes, $g:\mathrm{Y}'\to\mathrm{Y}$
a normal morphism $(6.8.1)$. Then, for every normal $\mathrm{Y}$-prescheme $\mathrm{X}$, the prescheme $\mathrm{X}'=\mathrm{X}\times_\mathrm{Y}\mathrm{Y}'$
is normal.
\end{proposition}

\begin{proof}
We note that in this statement one \emph{does not suppose} $\mathrm{X}$ locally Noetherian.
\end{proof}

\begin{lemma}[6.14.1.1]
\label{IV.6.14.1.1}
Let $\mathrm{R}$ be a ring that is a direct product of finitely many fields.
\begin{enumerate}[label=(\roman*)]
  \item For a subring $\mathrm{A}$ of $\mathrm{R}$ having $\mathrm{R}$ as total ring of fractions to be normal,
it is necessary and it suffices that it be integrally closed in $\mathrm{R}$.
  \item Let $(\mathrm{A}_\lambda)$ be a family of normal subrings of $\mathrm{R}$; if $\mathrm{A}=\bigcap_\lambda\mathrm{A}_\lambda$ admits $\mathrm{R}$ as
total ring of fractions, $\mathrm{A}$ is normal.
\end{enumerate}
\end{lemma}

\begin{proof}
(i) As $\mathrm{A}$ is a subring of $\mathrm{R}$, $\mathrm{A}_\mathfrak{p}$ is a subring of $\mathrm{R}_\mathfrak{p}$ for every
prime ideal $\mathfrak{p}$ of $\mathrm{A}$; furthermore $\mathrm{R}_\mathfrak{p}$ is a ring of fractions of $\mathrm{A}_\mathfrak{p}$; furthermore $\mathrm{R}_\mathfrak{p}$ is necessarily
a direct product of finitely many fields, hence every non-zero divisor
of $0$ in $\mathrm{R}_\mathfrak{p}$ is invertible; this proves that $\mathrm{R}_\mathfrak{p}$ is the total ring of fractions of $\mathrm{A}_\mathfrak{p}$.
If $\mathrm{A}$ is integrally closed in $\mathrm{R}$, $\mathrm{A}_\mathfrak{p}$ is therefore integrally closed in $\mathrm{R}_\mathfrak{p}$; but
if $\mathrm{R}_\mathfrak{p}$ is a direct product of at least two rings, the integral closure of $\mathrm{A}_\mathfrak{p}$ in $\mathrm{R}_\mathfrak{p}$ is
a direct product of at least two nonzero rings, which is absurd since $\mathrm{A}_\mathfrak{p}$
is a local ring; hence $\mathrm{R}_\mathfrak{p}$ is necessarily a field and $\mathrm{A}_\mathfrak{p}$ is integral and integrally
closed, which by definition means that $\mathrm{A}$ is normal. Conversely, if $\mathrm{A}$ is normal,
$\mathrm{R}_\mathfrak{p}$, total ring of fractions of an integral ring $\mathrm{A}_\mathfrak{p}$, is a field and $\mathrm{A}_\mathfrak{p}$ is integrally
closed in $\mathrm{R}_\mathfrak{p}$; hence $\mathrm{R}$ is an integral element over $\mathrm{A}$, its image in each $\mathrm{R}_\mathfrak{p}$ is
integral over $\mathrm{A}_\mathfrak{p}$, hence belongs to $\mathrm{A}_\mathfrak{p}$; we conclude that $x\in\mathrm{A}$ (Bourbaki, \emph{Alg.~comm.},
chap.~II, \S~3, no~3, cor.~1 of th.~1), and hence $\mathrm{A}$ is integrally closed in $\mathrm{R}$.

\oldpage[IV-2]{170}

(ii) As $\mathrm{A}\subset\mathrm{A}_\lambda\subset\mathrm{R}$ for every $\lambda$, $\mathrm{R}$ is also the total ring of fractions of each $\mathrm{A}_\lambda$.
In view of the characterisation of normal rings having $\mathrm{R}$ as total ring of fractions,
the assertion (ii) follows from Bourbaki, \emph{Alg.~comm.}, chap.~V,
\S~1, no~3, prop.~12.
\end{proof}

This lemma being proved, the proof of \sref{IV.6.14.1} proceeds in several steps.

I) \emph{Reduction to the case where} $\mathrm{Y}=\operatorname{Spec}(\mathrm{A})$, $\mathrm{Y}'=\operatorname{Spec}(\mathrm{A}')$, $\mathrm{X}=\operatorname{Spec}(\mathrm{B})$, $\mathrm{A}$, $\mathrm{A}'$ \emph{being
local Noetherian rings,} $\mathrm{A}$ integral, $\mathrm{B}$ the integral closure of $\mathrm{A}$. --- It suffices to prove that
for every $x'\in\mathrm{X}'$, the local ring $\sh{O}_{\mathrm{X}',x'}$ is integral and integrally closed; let $x$, $y$, $y'$
be the canonical images of $x'$ in $\mathrm{X}$, $\mathrm{Y}$, $\mathrm{Y}'$ respectively; if one sets $\mathrm{A}=\sh{O}_{\mathrm{Y},y}$,
$\mathrm{A}'=\sh{O}_{\mathrm{Y}',y'}$, $\mathrm{B}=\sh{O}_{\mathrm{X},x}$, the rings $\sh{O}_{\mathrm{X},x}$, $\sh{O}_{\mathrm{Y},y}$, $\sh{O}_{\mathrm{Y}',y'}$ for $x$, $y$, $y'$ fixed, are the local rings at the
prime ideals of the ring $\mathrm{B}'=\mathrm{B}\otimes_\mathrm{A}\mathrm{A}'$ $(\mathbf{I},~3.6.5)$, and so it suffices to prove that $\mathrm{B}'$
is a \emph{normal} ring. Let us note on the other hand that $\operatorname{Spec}(\mathrm{A}')\to\operatorname{Spec}(\mathrm{A})$ is a normal
morphism by definition $(6.8.1$ and $\mathbf{I},~3.6.5)$; one is reduced to the case where $\mathrm{Y}$ and $\mathrm{Y}'$ are
local schemes, that is to say to the case where $\mathrm{A}$ is an integral and integrally
closed local ring $(\mathbf{0}_\mathrm{I},~6.8.1$ and $\mathbf{I},~3.6.5)$; one is reduced to the case where $\mathrm{Y}$ and $\mathrm{Y}'$
are local schemes. Let us designate by $(\mathrm{B}_\alpha)$ the family of \emph{integral closures} of the sub-$\mathrm{A}$-algebras \emph{of finite type} of $\mathrm{B}$;
it is clear that $\mathrm{B}$ is the union of the filtered increasing family of the $\mathrm{B}_\alpha$. As $\varinjlim$
commutes with the tensor product, $\mathrm{B}'$ is therefore also the union of the filtered increasing family
$\overline{\mathrm{B}'_\alpha}=\mathrm{B}_\alpha\otimes_\mathrm{A}\mathrm{A}'$. To prove that $\mathrm{B}'$ is normal, it suffices, by virtue of \sref{IV.5.13.6}, to
show that the rings $\mathrm{B}'_\alpha$ are \emph{normal} and that, for $\alpha\leqslant\beta$, every irreducible component
of $\operatorname{Spec}(\mathrm{B}'_\beta)$ \emph{dominates} one irreducible component of $\operatorname{Spec}(\mathrm{B}'_\alpha)$. But this
last property follows from the hypothesis that $\mathrm{A}'$ is an $\mathrm{A}$-module flat and from $(2.3.7,~\text{(ii)})$.

One can therefore restrict to proving that $\mathrm{B}'=\mathrm{B}\otimes_\mathrm{A}\mathrm{A}'$ is normal when $\mathrm{B}$
is the integral closure of a $\mathrm{A}$-algebra integral \emph{of finite type} $\mathrm{C}$; in this case, if one sets
$\mathrm{C}'=\mathrm{C}\otimes_\mathrm{A}\mathrm{A}'$, the morphism $\operatorname{Spec}(\mathrm{C}')\to\operatorname{Spec}(\mathrm{C})$ is normal $(6.8.2)$; as
$\mathrm{B}'=\mathrm{B}\otimes_\mathrm{C}\mathrm{C}'$,
one sees that one can replace $\mathrm{A}$ and $\mathrm{A}'$ by $\mathrm{C}$ and $\mathrm{C}'$ respectively, hence suppose that $\mathrm{A}$
is integral and that $\mathrm{B}$ is the \emph{integral closure} of $\mathrm{A}$. Finally, the procedure from the beginning allows to
restrict to the case where $\mathrm{A}$ is \emph{local} (taking into account the fact that, if $\mathrm{B}$ is the integral closure of $\mathrm{A}$,
$\mathrm{B}_\mathfrak{p}$ is the integral closure of $\mathrm{A}_\mathfrak{p}$ for every prime ideal $\mathfrak{p}$ of $\mathrm{A}$).

II) \emph{Reduction to the case where} $\mathrm{A}$ \emph{is a local integral ring of dimension $1$,} $\mathrm{B}$ \emph{a discrete
valuation ring, integral closure of} $\mathrm{A}$ \emph{and the morphism} $\operatorname{Spec}(\mathrm{B})\to\operatorname{Spec}(\mathrm{A})$ \emph{radicial.} ---
Let $\mathrm{K}$ be the field of fractions of $\mathrm{A}$. One knows $(\mathbf{0},~23.2.7)$ that $\mathrm{B}$ is intersection of a
family $(\mathrm{V}_\lambda)$ of discrete valuation rings such that, for every $x\in\mathrm{K}$, one has $x\in\mathrm{V}_\lambda$
except for a finite number of indices $\lambda$. One therefore has an exact sequence of $\mathrm{A}$-modules
\[
  0\to\mathrm{B}\to\mathrm{K}\to\bigoplus_\lambda(\mathrm{K}/\mathrm{V}_\lambda).
\]

Set $\mathrm{K}'=\mathrm{K}\otimes_\mathrm{A}\mathrm{A}'$, $\mathrm{V}'_\lambda=\mathrm{V}_\lambda\otimes_\mathrm{A}\mathrm{A}'$; by flatness, one deduces from the preceding
exact sequence a new exact sequence
\[
  0\to\mathrm{B}'\to\mathrm{K}'\to\bigoplus(\mathrm{K}'_\lambda/\mathrm{V}'_\lambda)
\]

\oldpage[IV-2]{171}

and hence $\mathrm{B}'=\bigcap_\lambda\mathrm{V}'_\lambda$. Furthermore $\operatorname{Spec}(\mathrm{K}')$ is the fibre of the morphism $\operatorname{Spec}(\mathrm{A}')\to\operatorname{Spec}(\mathrm{A})$
at the generic point of $\operatorname{Spec}(\mathrm{A})$, hence $\mathrm{K}'$ is a Noetherian ring, \emph{normal}, and hence
a direct product of finitely many integrally closed integral rings; the direct product
$\mathrm{L}'$ of the fields of fractions of these latter is the total ring of fractions of $\mathrm{A}'$,
hence also that of $\mathrm{B}'$. One can therefore apply the lemma \sref{IV.6.14.1.1}, and if one proves
that each of the $\mathrm{V}'_\lambda$ is a \emph{normal} ring, then so is $\mathrm{B}'$.

One knows on the other hand $(\mathbf{0},~23.2.7)$ that one can take the $\mathrm{V}_\lambda$ such that there is a
sub-$\mathrm{A}$-algebra finite $\mathrm{C}$ of $\mathrm{K}$ such that $\mathrm{V}_\lambda$ is the integral closure of $\mathrm{C}_{\mathfrak{p}_\lambda}$ where $\mathfrak{p}_\lambda$ is a
prime ideal of height $1$ in $\mathrm{C}$. If one sets $\mathrm{C}'_{\mathfrak{p}_\lambda}=\mathrm{C}_{\mathfrak{p}_\lambda}\otimes_\mathrm{A}\mathrm{A}'$, the morphism
$\operatorname{Spec}(\mathrm{C}'_{\mathfrak{p}_\lambda})\to\operatorname{Spec}(\mathrm{C}_{\mathfrak{p}_\lambda})$ is normal $(6.8.2)$ and one has $\mathrm{V}'_\lambda=\mathrm{V}_\lambda\otimes_{\mathrm{C}_{\mathfrak{p}_\lambda}}\mathrm{C}'_{\mathfrak{p}_\lambda}$; one can therefore
replace $\mathrm{B}$ by $\mathrm{V}_\lambda$ and $\mathrm{A}$ by $\mathrm{C}_{\mathfrak{p}_\lambda}$, hence suppose that $\mathrm{A}$ is local, integral and \emph{of dimension $1$},
$\mathrm{B}$ its integral closure and a discrete valuation ring. There is a sub-$\mathrm{A}$-algebra \emph{finite} $\mathrm{A}_1$
of $\mathrm{B}$ such that the morphism $\operatorname{Spec}(\mathrm{B})\to\operatorname{Spec}(\mathrm{A}_1)$ is radicial $(\mathbf{0},~23.2.5)$, which implies
in particular that $\mathrm{A}_1$ is also a local ring (obviously of dimension $1$); furthermore
one can suppose that $\mathrm{B}$ and $\mathrm{A}_1$ have the same residue field (\emph{loc.~cit.}). One can therefore by the
same method replace $\mathrm{A}$ by $\mathrm{A}_1$. Applying if necessary the procedure from the beginning of I),
one can suppose finally that $\mathrm{A}'$ is also a local ring and that the homomorphism
$\mathrm{A}\to\mathrm{A}'$ is \emph{local}.

III) \emph{End of the proof.} --- We shall first establish the following lemma:
\end{proof}

\begin{lemma}[6.14.1.2]
\label{IV.6.14.1.2}
Let $\mathrm{A}$ be a local Noetherian integral ring of dimension $1$, $\mathrm{A}'$ a local
Noetherian integral ring, $\mathrm{A}\to\mathrm{A}'$ a local homomorphism such that the corresponding morphism
$\operatorname{Spec}(\mathrm{A}')\to\operatorname{Spec}(\mathrm{A})$ is \emph{normal}. Let $\mathrm{K}$ be the field of fractions of $\mathrm{A}$, $\mathfrak{m}$ the maximal ideal of $\mathrm{A}$,
$\boldsymbol{k}=\mathrm{A}/\mathfrak{m}$ its residue field.
\end{lemma}

\begin{env}
\label{IV.6.14.1.2.env}
\begin{enumerate}[label=(\roman*)]
  \item If $\mathfrak{q}'_j$ $(1\leqslant j\leqslant r)$ are the minimal ideals of $\mathrm{A}'$, $\mathrm{K}'=\mathrm{K}\otimes_\mathrm{A}\mathrm{A}'$ is a direct product
of integral rings and integrally closed $\mathrm{K}'_j\supset\mathrm{A}'/\mathfrak{q}'_j$ $(1\leqslant j\leqslant r)$, $\mathrm{K}'_j$ having as field of fractions
the field of fractions $\mathrm{L}'_j$ of $\mathrm{A}'/\mathfrak{q}'_j$, so that the total ring of fractions $\mathrm{L}'$ of $\mathrm{A}'$
is identified with the direct product of the $\mathrm{L}'_j$ and $\mathrm{A}'$ is a subring of $\mathrm{K}'$.
  \item The ideal $\mathfrak{p}'=\mathfrak{m}\mathrm{A}'$ of $\mathrm{A}'$ is prime; the ring $\mathrm{A}'_{\mathfrak{p}'}$ is of dimension $1$ and is identified
with a sub-product of the fields $\mathrm{L}'_j$ for the indices $j$ such that $\mathfrak{q}'_j\subset\mathfrak{p}'$.
  \item If $\overline{\mathrm{A}'_{\mathfrak{p}'}}$ denotes the subring of $\mathrm{L}'$, product of $\mathrm{A}'_{\mathfrak{p}'}$ and of the $\mathrm{L}'_j$ such that $\mathfrak{q}'_j\not\subset\mathfrak{p}'$, one has
\begin{equation}
\tag{6.14.1.3}
\label{IV.6.14.1.3}
  \mathrm{A}'=\mathrm{K}'\cap\overline{\mathrm{A}'_{\mathfrak{p}'}}.
\end{equation}
\end{enumerate}
\end{env}

\begin{proof}
\oldpage[IV-2]{172}
Assertion (i) has already been seen in the course of the proof, independent of the hypothesis on the dimension of $\mathrm{A}$. The hypothesis that the morphism
$\operatorname{Spec}(\mathrm{A}')\to\operatorname{Spec}(\mathrm{A})$ is normal implies that $\mathrm{A}'\otimes_\mathrm{A}\boldsymbol{k}=\mathrm{A}'/\mathfrak{m}\mathrm{A}'$ is a normal local ring
and \emph{integral}, which shows already that $\mathfrak{p}'=\mathfrak{m}\mathrm{A}'$ is prime. One has,
by \sref{IV.6.1.2}, $\dim(\mathrm{A}'_{\mathfrak{p}'})=\dim(\mathrm{A})+\dim(\mathrm{A}'_{\mathfrak{p}'}/\mathfrak{m}\mathrm{A}'_{\mathfrak{p}'})$. But $\mathrm{A}'_{\mathfrak{p}'}/\mathfrak{m}\mathrm{A}'_{\mathfrak{p}'}=\mathrm{A}'_{\mathfrak{p}'}/\mathfrak{p}'\mathrm{A}'_{\mathfrak{p}'}$ is
the residue field of $\mathrm{A}'_{\mathfrak{p}'}$, hence $\dim(\mathrm{A}'_{\mathfrak{p}'})=1$. The fact that $\mathrm{A}'_{\mathfrak{p}'}$ is contained in the composed
direct of the $\mathrm{L}'_j$ for these indices $j$ such that $\mathfrak{q}'_j\subset\mathfrak{p}'$, for these indices $j$ are
the minimal prime ideals of $\mathrm{A}'_{\mathfrak{p}'}$.

It remains to prove \sref{IV.6.14.1.3}. Conversely, let $y'$ be an
element of this intersection; let $a$ be a ``parameter'' for $\mathrm{A}$, i.e.\ an element such that $\mathrm{A}a$ is $\mathfrak{m}$-primary,
and let $a'$ be the image of $a$ in $\mathrm{A}'$; every element of $\mathrm{K}$ can be written $x/a^n$ for $x\in\mathrm{A}$ and
$n>0$ since $\mathrm{A}a$ contains a power of $\mathfrak{m}$; hence one can write $y'=x'/a'^n$ with $x'\in\mathrm{A}'$.
Note now that $\mathfrak{p}'$ is the \emph{only} prime ideal associated to $\mathrm{A}'/a'^n\mathrm{A}'=(A/a^n\mathrm{A})\otimes_\mathrm{A}\mathrm{A}'$:
as $\mathrm{A}'$ is a flat $\mathrm{A}$-module, this follows from $(3.3.1)$, $\mathfrak{m}$ being the only prime ideal
associated to $\mathrm{A}/a^n\mathrm{A}$ since $\mathrm{A}$ is integral. Hence, $a'^n\mathrm{A}'$ is the inverse image in $\mathrm{A}'$
of $a'^n\mathrm{A}'_{\mathfrak{p}'}$; as by hypothesis $y'\in\widehat{\mathrm{A}'_{\mathfrak{p}'}}$, the image of $x'$ in $\mathrm{A}'_{\mathfrak{p}'}$ belongs to $a'^n\mathrm{A}'_{\mathfrak{p}'}$;
whence $x'\in a'^n\mathrm{A}'$ and $y'\in\mathrm{A}'$, which finishes the proof.

In the situation in which we have arrived at the end of II),
$\mathrm{B}'$ is \emph{radicial} over $\mathrm{A}'$ $(\mathbf{I},~3.5.7)$ and is hence also a local ring; furthermore $\mathrm{B}'$ is
integer over $\mathrm{A}'$, hence, if $\mathfrak{q}'$ is the unique prime ideal of $\mathrm{B}'$ lying over $\mathfrak{p}'$, $\mathfrak{q}'\mathrm{B}'_{\mathfrak{p}'}$ is the unique
maximal ideal of $\mathrm{B}'_{\mathfrak{p}'}$ (Bourbaki, \emph{Alg.~comm.}, chap.~V, \S~2, no~1, prop.~1), hence
$\mathrm{B}'_{\mathfrak{q}'}=\mathrm{B}'_{\mathfrak{p}'}=\mathrm{B}\otimes_{\mathrm{A}_\mathfrak{p}}\mathrm{A}'_{\mathfrak{p}'}$. We shall first prove that $\mathrm{B}'_{\mathfrak{p}'}$ is \emph{Noetherian}
and \emph{normal}. $\mathrm{B}'_{\mathfrak{p}'}$ contains $\mathrm{A}'_{\mathfrak{p}'}$ and is contained in $\mathrm{K}'_{\mathfrak{p}'}$, hence in the product $\mathrm{L}_1$ of the $\mathrm{L}'_j$ such that $\mathfrak{q}'_j\subset\mathfrak{p}'$. The composed
direct $\mathrm{L}''$ of the fields of fractions of these latter is also the total ring of fractions of $\mathrm{A}'$,
hence also that of $\mathrm{B}'_{\mathfrak{p}'}$. For every index $j$ such that $\mathfrak{q}'_j\subset\mathfrak{p}'$, let $\mathfrak{a}'_j$ be the product of the $\mathrm{L}'_h$ for these indices $h$ such that $h\neq j$, so that
$\mathfrak{a}'_j\cap\mathrm{A}'_{\mathfrak{p}'}=\mathfrak{q}'_j\mathrm{A}'_{\mathfrak{p}'}$; as every element of $\mathrm{A}'-\mathfrak{p}'$ is regular in $\mathrm{L}_1$, one has also
\[
  \mathfrak{a}'_j\cap\mathrm{A}'_{\mathfrak{p}'}=\mathfrak{q}'_j\mathrm{A}'_{\mathfrak{p}'}.
\]

Let $\mathfrak{r}'_j=\mathrm{B}'_{\mathfrak{p}'}\cap\mathfrak{a}'_j$, so that $\mathrm{B}'_{\mathfrak{p}'}/\mathfrak{r}'_j$ is isomorphic to the projection of $\mathrm{B}'_{\mathfrak{p}'}$ in $\mathrm{L}'_j$;
hence $\mathrm{B}'_{\mathfrak{p}'}/\mathfrak{r}'_j$ contains the local integral ring of \emph{dimension $1$}, $\mathrm{A}'_{\mathfrak{p}'}/\mathfrak{q}'_j\mathrm{A}'_{\mathfrak{p}'}$, and is contained
in its field of fractions $\mathrm{L}'_j$; it is hence a Noetherian ring by virtue of the theorem of
Krull-Akizuki (Bourbaki, \emph{Alg.~comm.}, chap.~VII, \S~2, no~5, prop.~5).

\oldpage[IV-2]{173}

As the intersection of the $\mathfrak{r}'_j$ is reduced to $0$, one deduces that $\mathrm{B}'_{\mathfrak{p}'}$ itself is
Noetherian, by virtue of the classical lemma:

\begin{lemma}[6.14.1.4]
\label{IV.6.14.1.4}
Let $\mathrm{R}$ be a ring, $\mathfrak{a}$ and $\mathfrak{b}$ two ideals of $\mathrm{R}$; if $\mathrm{R}/\mathfrak{a}$ and $\mathrm{R}/\mathfrak{b}$
are Noetherian, so is $\mathrm{R}/(\mathfrak{a}\cap\mathfrak{b})$.
\end{lemma}

\begin{proof}
Indeed, let $\mathfrak{c}$ be an ideal of $\mathrm{R}$ such that $\mathfrak{a}\cap\mathfrak{b}\subset\mathfrak{c}$; one has therefore $\mathfrak{a}\cap\mathfrak{b}\subset\mathfrak{a}\cap\mathfrak{c}\subset\mathfrak{c}$; now
$\mathfrak{c}/(\mathfrak{a}\cap\mathfrak{c})$ is an $\mathrm{R}$-module isomorphic to $(\mathfrak{a}+\mathfrak{c})/\mathfrak{a}$, hence an ideal of $\mathrm{R}/\mathfrak{a}$, and hence
is of finite type; on the other hand $(\mathfrak{a}\cap\mathfrak{c})/(\mathfrak{a}\cap\mathfrak{b})$ is a sub-$\mathrm{R}$-module of $\mathfrak{a}/(\mathfrak{a}\cap\mathfrak{b})$, and
this latter is isomorphic to $(\mathfrak{a}+\mathfrak{b})/\mathfrak{b}$, hence is of finite type as an ideal of $\mathrm{R}/\mathfrak{b}$;
hence $(\mathfrak{a}\cap\mathfrak{c})/(\mathfrak{a}\cap\mathfrak{b})$ is also an $\mathrm{R}$-module of finite type, and it follows finally that
$\mathfrak{c}/(\mathfrak{a}\cap\mathfrak{b})$ is also an $\mathrm{R}$-module of finite type.
\end{proof}

Let us note on the other hand that $\operatorname{Spec}(\mathrm{A})$ is made up of the closed point $\mathfrak{m}$ and of the generic
point (o), of residue fields $\boldsymbol{k}$ and $\mathrm{K}$ respectively; as in the case where we
are placed, the fields of fractions of $\mathrm{B}$ and its residue field are respectively
isomorphic to $\mathrm{K}$ and $\boldsymbol{k}$, the fibres of the morphism $u:\operatorname{Spec}(\mathrm{B}'_{\mathfrak{p}'})\to\operatorname{Spec}(\mathrm{B})$ at the closed point
and at the generic point of $\operatorname{Spec}(\mathrm{B})$ are respectively \emph{isomorphic} to the fibres of the morphism
$\operatorname{Spec}(\mathrm{A}'_{\mathfrak{p}'})\to\operatorname{Spec}(\mathrm{A})$ at the closed point and at the generic point of $\operatorname{Spec}(\mathrm{A})$ $(\mathbf{I},~3.6.4)$, hence
are \emph{geometrically normal}; as on the other hand the morphism $u$ is flat,
we conclude that it is \emph{normal} $(6.8.1)$. But $\mathrm{B}$ and $\mathrm{B}'_{\mathfrak{p}'}$ are Noetherian and $\mathrm{B}$ is normal, hence
$(6.5.4)$ $\mathrm{B}'_{\mathfrak{p}'}$ is \emph{normal}.

This being so, $\mathrm{B}$ is union of the filtered increasing family of its sub-$\mathrm{A}$-algebras of finite type
$\mathrm{A}_\alpha$; by flatness, $\mathrm{B}'$ is union of the filtered increasing family of the $\mathrm{A}'_\alpha=\mathrm{A}_\alpha\otimes_\mathrm{A}\mathrm{A}'$;
if one sets $\mathfrak{p}'_\alpha=\mathfrak{q}'\cap\mathrm{A}'_\alpha$, $\mathrm{B}'_{\mathfrak{q}'}$ is also union of the filtered increasing family of the $(\mathrm{A}'_\alpha)_{\mathfrak{p}'_\alpha}$
$(5.13.3)$. Denote by $\mathrm{L}''$ the direct product of the fields $\mathrm{L}'_j$ such that $\mathfrak{q}'_j\not\subset\mathfrak{p}'$; then for
every $\alpha$, $(\mathrm{A}'_\alpha)_{\mathfrak{p}'_\alpha}$ is contained in $\mathrm{K}'_{\mathfrak{p}'}$, and the ring $\overline{\mathrm{B}'_{\mathfrak{q}'}}=\mathrm{L}''\times\mathrm{B}'_{\mathfrak{q}'}$ is therefore union of the
rings $\mathrm{L}''\times(\mathrm{A}'_\alpha)_{\mathfrak{p}'_\alpha}=(\overline{\mathrm{A}'_\alpha})_{\mathfrak{p}'_\alpha}$. But each of the $\mathrm{A}_\alpha$ is local, Noetherian, integral and of
dimension $1$, and the morphism $\operatorname{Spec}(\mathrm{A}'_\alpha)\to\operatorname{Spec}(\mathrm{A}_\alpha)$ is normal $(6.8.2)$, hence one can
apply \sref{IV.6.14.1.2} and one has
\[
  \mathrm{A}'_\alpha=\mathrm{K}'\cap(\overline{\mathrm{A}'_\alpha})_{\mathfrak{p}'_\alpha}
\]
for every $\alpha$; taking the inductive limit of each of the two members, it follows that
\[
  \mathrm{B}'=\mathrm{K}'\cap\overline{\mathrm{B}'_{\mathfrak{q}'}}.
\]

But $\overline{\mathrm{B}'_{\mathfrak{q}'}}$, being the direct product of the normal rings $\mathrm{L}''$ and $\mathrm{B}'_{\mathfrak{q}'}$, is normal, and as
it has the same total ring of fractions as $\mathrm{K}'$, the lemma \sref{IV.6.14.1.1} shows that $\mathrm{B}'$ is normal.
\end{proof}

\begin{corollary}[6.14.2]
\label{IV.6.14.2}
Let $k$ be a field, $\mathrm{X}$ a $k$-prescheme normal (not necessarily
locally Noetherian). Then, for every separable extension $k'$ of $k$, $\mathrm{X}_{(k')}=\mathrm{X}\otimes_k k'$ is
\emph{normal}.
\end{corollary}

\begin{proof}
One knows indeed $(6.7.6)$ that the morphism $\operatorname{Spec}(k')\to\operatorname{Spec}(k)$ is then normal.
\end{proof}

\begin{corollary}[6.14.3]
\label{IV.6.14.3}
Let $k$ be a field, $\mathrm{X}$, $\mathrm{Y}$ two $k$-preschemes integral and normal,
whose fields of rational functions $\mathrm{K}=\mathrm{R}(\mathrm{X})$, $\mathrm{L}=\mathrm{R}(\mathrm{Y})$ are separable extensions of $k$.
Then the prescheme $\mathrm{X}\times_k\mathrm{Y}$ is \emph{normal}.
\end{corollary}

\begin{proof}
\oldpage[IV-2]{174}
The question being local on $\mathrm{X}$ and $\mathrm{Y}$, one can suppose that $\mathrm{X}=\operatorname{Spec}(\mathrm{A})$, $\mathrm{Y}=\operatorname{Spec}(\mathrm{B})$,
where $\mathrm{A}$ and $\mathrm{B}$ are two integral rings and integrally closed, of fields of fractions $\mathrm{K}$
and $\mathrm{L}$ respectively; suppose first that $\mathrm{A}$ is a $k$-algebra of finite type, hence $\mathrm{K}$
is a finite type extension of $k$. By flatness, $\mathrm{A}\otimes_k\mathrm{B}$ is a sub-$k$-algebra of $\mathrm{K}\otimes_k\mathrm{L}$;
now $\mathrm{K}\otimes_k\mathrm{L}$ is a Noetherian normal ring (by virtue of $(6.7.1)$ or \sref{IV.6.14.2}),
hence a direct product of finitely many integral rings $\mathrm{C}_i$, so that if $\mathrm{E}_i$ is the
field of fractions of $\mathrm{C}_i$, the direct product $\mathrm{E}$ of the $\mathrm{E}_i$ is the total ring of fractions
of $\mathrm{K}\otimes_k\mathrm{L}$; it is also the total ring of fractions of $\mathrm{A}\otimes_k\mathrm{B}$, since $\mathrm{K}\otimes_k\mathrm{L}$
is a ring of fractions of $\mathrm{A}\otimes_k\mathrm{B}$. This being so, one can write $\mathrm{A}\otimes_k\mathrm{B}=(\mathrm{A}\otimes_k\mathrm{L})\cap(\mathrm{K}\otimes_k\mathrm{B})$,
and it follows from \sref{IV.6.14.2} that $\mathrm{A}\otimes_k\mathrm{L}$ and $\mathrm{K}\otimes_k\mathrm{B}$ are normal, hence $\mathrm{A}\otimes_k\mathrm{B}$ is normal
by virtue of \sref{IV.6.14.1.1}.

Consider now the general case; $\mathrm{A}$ is then union of the filtered increasing
family of its sub-$k$-algebras of finite type $\mathrm{A}_\alpha$, hence $\mathrm{A}\otimes_k\mathrm{B}$ is the inductive limit of the
normal rings $\mathrm{A}_\alpha\otimes_k\mathrm{B}$. To prove that $\mathrm{A}\otimes_k\mathrm{B}$ is normal, it suffices hence $(5.13.6)$
to prove that every irreducible component of $\operatorname{Spec}(\mathrm{A}_\beta\otimes_k\mathrm{B})$ dominates one irreducible component
of $\operatorname{Spec}(\mathrm{A}_\alpha\otimes_k\mathrm{B})$ for $\mathrm{A}_\alpha\subset\mathrm{A}_\beta$, which follows immediately from the fact that $\mathrm{B}$ is
a $k$-module flat, from $(2.3.7,~\text{(ii)})$ and from the fact that $\mathrm{A}_\alpha$ and $\mathrm{A}_\beta$ are integral rings.
\end{proof}

\begin{proposition}[6.14.4]
\label{IV.6.14.4}
Let $\mathrm{A}$ be a Noetherian ring, $\mathrm{A}'$ a Noetherian $\mathrm{A}$-algebra
such that the morphism $\operatorname{Spec}(\mathrm{A}')\to\operatorname{Spec}(\mathrm{A})$ is normal. Let $\mathrm{B}$ be an $\mathrm{A}$-algebra, $\mathrm{C}$ the integral
closure of $\mathrm{A}$ in $\mathrm{B}$; set $\mathrm{B}'=\mathrm{B}\otimes_\mathrm{A}\mathrm{A}'$, $\mathrm{C}'=\mathrm{C}\otimes_\mathrm{A}\mathrm{A}'$, $\mathrm{C}'$ being identified with a subring
of $\mathrm{B}'$; then $\mathrm{C}'$ is the integral closure of $\mathrm{A}'$ in $\mathrm{B}'$.
\end{proposition}

\begin{proof}
We shall proceed in several steps.

I) \emph{Reduction to the case where} $\mathrm{B}$ \emph{is reduced.} --- Set $\mathrm{B}_0=\mathrm{B}_{\mathrm{red}}=\mathrm{B}/\mathfrak{N}$ where $\mathfrak{N}$ is the
nilradical of $\mathrm{B}$, and let $\mathrm{C}_0$ be the integral closure of $\mathrm{A}$ in $\mathrm{B}_0$. One has the following lemma:

\begin{lemma}[6.14.4.1]
\label{IV.6.14.4.1}
Let $\mathrm{A}$ be a ring, $\mathrm{B}$ an $\mathrm{A}$-algebra, $\mathrm{B}_0=\mathrm{B}/\mathfrak{n}$ the quotient of $\mathrm{B}$
by a nilideal. If $\mathrm{C}_0$ is the integral closure of $\mathrm{A}$ in $\mathrm{B}_0$, the inverse image $\mathrm{C}$ of $\mathrm{C}_0$ by
the canonical homomorphism $\varphi:\mathrm{B}\to\mathrm{B}_0$ is the integral closure of $\mathrm{A}$ in $\mathrm{B}$.
\end{lemma}

\begin{proof}
Indeed, if $x\in\mathrm{B}$ is such that $\varphi(x)$ satisfies an integral dependence equation
over $\mathrm{A}$, one deduces that $x$ satisfies a relation of the form
$x^n+a_1 x^{n-1}+\ldots+a_n\in\mathfrak{n}$ with $a_k\in\mathrm{A}$, whence, by raising to a sufficiently large power, an
integral dependence equation for $x$, with coefficients in $\mathrm{A}$.
\end{proof}

One has therefore the exact sequence of $\mathrm{A}$-modules
\[
  0\to\mathrm{C}\to\mathrm{B}\to\mathrm{B}_0/\mathrm{C}_0
\]
whence, by flatness, the exact sequence
\[
  0\to\mathrm{C}'\to\mathrm{B}'\to\mathrm{B}'_0/\mathrm{C}'_0
\]
where one has set $\mathrm{B}'_0=\mathrm{B}_0\otimes_\mathrm{A}\mathrm{A}'$, $\mathrm{C}'_0=\mathrm{C}_0\otimes_\mathrm{A}\mathrm{A}'$. If one proves that $\mathrm{C}'_0$ is the integral
closure of $\mathrm{A}'$ in $\mathrm{B}'_0$, the lemma \sref{IV.6.14.4.1} will prove that $\mathrm{C}'$ is the integral
closure of $\mathrm{A}'$ in $\mathrm{B}'$. One can therefore restrict to the case where $\mathrm{B}$ is an $\mathrm{A}$-algebra of
finite type; let $\mathfrak{q}_i$ $(1\leqslant i\leqslant n)$ be its minimal prime ideals; as $\mathrm{B}$
is supposed \emph{reduced}, it is identified with a \emph{subring} of the product $\mathrm{B}_0$ of the $\mathrm{B}/\mathfrak{q}_i$; if $\mathrm{C}_0$ is the integral
closure of $\mathrm{A}$ in $\mathrm{B}_0$, one has $\mathrm{C}=\mathrm{B}\cap\mathrm{C}_0$; if one sets $\mathrm{B}'_0=\mathrm{B}_0\otimes_\mathrm{A}\mathrm{A}'$, $\mathrm{C}'_0=\mathrm{C}_0\otimes_\mathrm{A}\mathrm{A}'$, one has,
by flatness, $\mathrm{C}'=\mathrm{B}'\cap\mathrm{C}'_0$ $(\mathbf{0}_\mathrm{I},~6.1.3)$; it suffices therefore to prove that $\mathrm{C}'_0$ is the integral
closure of $\mathrm{A}'$ in $\mathrm{B}'_0$. But $\mathrm{C}_0$ is the direct product of the $\mathrm{C}_i$, integral closures of $\mathrm{A}$
in the $\mathrm{B}_i=\mathrm{B}/\mathfrak{q}_i$; hence $\mathrm{C}'_0$ is the direct product of the $\mathrm{C}'_i=\mathrm{C}_i\otimes_\mathrm{A}\mathrm{A}'$ and it suffices to show
that $\mathrm{C}'_i$ is the integral closure of $\mathrm{A}'$ in $\mathrm{B}'_i=\mathrm{B}_i\otimes_\mathrm{A}\mathrm{A}'$. One is thus reduced to the case
where $\mathrm{B}$ is integral and an $\mathrm{A}$-algebra of finite type; if $\mathfrak{p}$ is the kernel of the homomorphism $\mathrm{A}\to\mathrm{B}$,
one has also $\mathrm{B}'=\mathrm{B}\otimes_{\mathrm{A}/\mathfrak{p}}(\mathrm{A}'/\mathfrak{p}\mathrm{A}')$; as the morphism $\operatorname{Spec}(\mathrm{A}'/\mathfrak{p}\mathrm{A}')\to\operatorname{Spec}(\mathrm{A}/\mathfrak{p})$ is
normal, one can replace $\mathrm{A}$ by $\mathrm{A}/\mathfrak{p}$ and $\mathrm{A}'$ by $\mathrm{A}'/\mathfrak{p}\mathrm{A}'$, and suppose hence
that $\mathrm{A}\subset\mathrm{B}$.

\oldpage[IV-2]{175}

II) \emph{Reduction to the case where} $\mathrm{B}$ \emph{is an integral} $\mathrm{A}$\emph{-algebra of finite type containing} $\mathrm{A}$. --- Let $(\mathrm{B}_\alpha)$
be the filtered increasing family of sub-$\mathrm{A}$-algebras of finite type of $\mathrm{B}$, and let $\mathrm{C}_\alpha$ be the integral
closure of $\mathrm{A}$ in $\mathrm{B}_\alpha$. It follows immediately from the definition of the integral closure that $\mathrm{C}$
is union of the $\mathrm{C}_\alpha$; if one sets $\mathrm{B}'_\alpha=\mathrm{B}_\alpha\otimes_\mathrm{A}\mathrm{A}'$, $\mathrm{C}'_\alpha=\mathrm{C}_\alpha\otimes_\mathrm{A}\mathrm{A}'$, $\mathrm{C}'_\alpha$ is contained in $\mathrm{B}'_\alpha$
by flatness, and for the same reason $\mathrm{B}'$ is union of the filtered increasing family of the
$\mathrm{B}'_\alpha$ and $\mathrm{C}'$ union of the filtered increasing family of the $\mathrm{C}'_\alpha$. If one proves that $\mathrm{C}'_\alpha$
is the integral closure of $\mathrm{A}'$ in $\mathrm{B}'_\alpha$, it will follow immediately that $\mathrm{C}'$ is the integral
closure of $\mathrm{A}'$ in $\mathrm{B}'$. One can therefore restrict to the case where $\mathrm{B}$ is an $\mathrm{A}$-algebra of
finite type; let $\mathfrak{q}_i$ $(1\leqslant i\leqslant n)$ be its minimal prime ideals; as $\mathrm{B}$
is supposed \emph{reduced}, it is identified with a \emph{subring} of the product $\mathrm{B}_0$ of the $\mathrm{B}/\mathfrak{q}_i$; if $\mathrm{C}_0$ is the integral closure of $\mathrm{A}$ in $\mathrm{B}_0$, one has $\mathrm{C}=\mathrm{B}\cap\mathrm{C}_0$; if one sets $\mathrm{B}'_0=\mathrm{B}_0\otimes_\mathrm{A}\mathrm{A}'$, $\mathrm{C}'_0=\mathrm{C}_0\otimes_\mathrm{A}\mathrm{A}'$, one has,
by flatness, $\mathrm{C}'=\mathrm{B}'\cap\mathrm{C}'_0$ $(\mathbf{0}_\mathrm{I},~6.1.3)$; it suffices therefore to prove that $\mathrm{C}'_0$ is the integral
closure of $\mathrm{A}'$ in $\mathrm{B}'_0$. But $\mathrm{C}_0$ is the direct product of the $\mathrm{C}_i$, integral closures of $\mathrm{A}$ in the $\mathrm{B}_i=\mathrm{B}/\mathfrak{q}_i$; hence $\mathrm{C}'_0$ is the direct product of the $\mathrm{C}'_i=\mathrm{C}_i\otimes_\mathrm{A}\mathrm{A}'$ and it suffices to show that $\mathrm{C}'_i$ is the integral closure of $\mathrm{A}'$ in $\mathrm{B}'_i=\mathrm{B}_i\otimes_\mathrm{A}\mathrm{A}'$. One is thus reduced to the case where $\mathrm{B}$ is integral and $\mathrm{A}$-algebra of finite type; if $\mathfrak{p}$ is the kernel of the homomorphism $\mathrm{A}\to\mathrm{B}$, one has also $\mathrm{B}'=\mathrm{B}\otimes_{\mathrm{A}/\mathfrak{p}}(\mathrm{A}'/\mathfrak{p}\mathrm{A}')$; as the morphism $\operatorname{Spec}(\mathrm{A}'/\mathfrak{p}\mathrm{A}')\to\operatorname{Spec}(\mathrm{A}/\mathfrak{p})$ is normal, one can replace $\mathrm{A}$ by $\mathrm{A}/\mathfrak{p}$ and $\mathrm{A}'$ by $\mathrm{A}'/\mathfrak{p}\mathrm{A}'$, and suppose hence that $\mathrm{A}\subset\mathrm{B}$.

III) \emph{Case where} $\mathrm{A}$ \emph{is a field,} $\mathrm{B}$ \emph{a field, finite extension of} $\mathrm{A}$ \emph{and such that} $\mathrm{A}$ \emph{is algebraically
closed in} $\mathrm{B}$. --- One has then $\mathrm{C}=\mathrm{A}$, and $\mathrm{A}'$ is a Noetherian geometrically
normal $\mathrm{A}$-algebra, hence a direct product of integrally closed integral rings $\mathrm{A}'_i$, the
fields of fractions $\mathrm{K}'_i$ of the $\mathrm{A}'_i$ being \emph{separable} extensions of $\mathrm{A}$ $(4.6.1)$. The $\mathrm{A}'$-algebra $\mathrm{B}'$
is therefore the direct product of the $\mathrm{B}\otimes_\mathrm{A}\mathrm{A}'_i=\mathrm{B}'_i$. One can therefore restrict to the case where $\mathrm{A}'$ is
integral, its field of fractions $\mathrm{K}'$ being a separable extension of $\mathrm{A}$. Now, as $\mathrm{K}'$
is a separable extension of $\mathrm{A}$, $\mathrm{B}\otimes_\mathrm{A}\mathrm{K}'$ is reduced $(4.3.7)$, and as furthermore $\mathrm{B}$ is a
primary extension of $\mathrm{A}$, $\mathrm{B}\otimes_\mathrm{A}\mathrm{K}'$ is integral $(4.3.2)$; it is hence of the same type as $\mathrm{B}'$.
Let $\mathrm{L}'$ be the field of fractions of $\mathrm{B}\otimes_\mathrm{A}\mathrm{K}'$ (which is also that of $\mathrm{B}'$); it is obviously a
composed field of $\mathrm{B}(\mathrm{K}')$ of its sub-fields $\mathrm{B}$ and $\mathrm{K}'$; as $\mathrm{A}$ is algebraically closed in $\mathrm{B}$,
that $\mathrm{K}'$ is a separable extension of $\mathrm{A}$, and that $\mathrm{B}$ and $\mathrm{K}'$ are linearly disjoint over $\mathrm{A}$,
$\mathrm{K}'$ is algebraically closed in $\mathrm{B}(\mathrm{K}')=\mathrm{L}'$ (Bourbaki, \emph{Alg.}, chap.~V, \S~9, exerc.~2);
\emph{a fortiori}, $\mathrm{A}'$, being integrally closed, is integrally closed in $\mathrm{L}'$, hence in $\mathrm{B}'$,
which proves the proposition in this case.

IV) \emph{Case where} $\mathrm{A}$ \emph{is integral,} $\mathrm{B}$ \emph{a field, finite extension of the field of fractions} $\mathrm{K}$ \emph{of} $\mathrm{A}$. ---
Then $\mathrm{B}'=(\mathrm{A}'\otimes_\mathrm{A}\mathrm{K})\otimes_\mathrm{K}\mathrm{B}$ is a $\mathrm{B}$-algebra Noetherian geometrically normal, and
hence $\mathrm{B}'$ is a direct product of integral rings; the total ring of fractions $\mathrm{L}'$ of $\mathrm{B}'$
is then a direct product of finitely many fields. This being so, as $\mathrm{B}$ is the field of
fractions of $\mathrm{C}$, $\mathrm{B}'$ is a ring of fractions of $\mathrm{C}'$ and $\mathrm{L}'$ is therefore also the total ring of
fractions of $\mathrm{C}'$. Now $(\mathbf{0},~6.14.1)$ as $\mathrm{C}$ is a normal ring and $\operatorname{Spec}(\mathrm{A}')\to\operatorname{Spec}(\mathrm{A})$
a normal morphism of Noetherian rings, $\mathrm{C}'$ is a normal ring; but it follows from
\sref{IV.6.14.1.1} that $\mathrm{C}'$ is then integrally closed in $\mathrm{L}'$, and \emph{a fortiori} in $\mathrm{B}'$, and hence
is the integral closure of $\mathrm{A}'$ in $\mathrm{B}'$.

\oldpage[IV-2]{176}

V) \emph{End of the proof.} --- By part II), one can suppose that $\mathrm{B}$ is an $\mathrm{A}$-algebra
integral of finite type containing $\mathrm{A}$. Let $\mathrm{K}$ be the field of fractions of $\mathrm{A}$, $\mathrm{L}$ that of $\mathrm{B}$, which
is a finite type extension of $\mathrm{K}$. Let $\mathrm{M}$ be the algebraic closure of $\mathrm{K}$ in $\mathrm{L}$, which
is a finite algebraic extension of $\mathrm{K}$; let $\mathrm{C}_0$ be the integral closure of $\mathrm{A}$ in $\mathrm{M}$,
which is also the integral closure of $\mathrm{A}$ in $\mathrm{L}$; if one sets
$\mathrm{C}'_0=\mathrm{C}\otimes_\mathrm{A}\mathrm{A}'$, one has hence $\mathrm{C}'=\mathrm{B}'\cap\mathrm{C}'_0$ by flatness $(\mathbf{0}_\mathrm{I},~6.1.3)$. Now, it follows from IV)
that $\mathrm{C}'_0$ is the integral closure of $\mathrm{A}'$ in $\mathrm{M}'=\mathrm{M}\otimes_\mathrm{A}\mathrm{A}'$; furthermore, $\mathrm{M}'$ is a Noetherian
ring and the morphism $\operatorname{Spec}(\mathrm{M}')\to\operatorname{Spec}(\mathrm{M})$ is normal (as we have seen in the
course of IV)); applying III) to $\mathrm{M}$ and $\mathrm{M}'$ in place of $\mathrm{A}$ and $\mathrm{A}'$ and to $\mathrm{L}$ in place of $\mathrm{B}$, one deduces
that $\mathrm{C}'_0$ is the integral closure of $\mathrm{A}'$ in $\mathrm{L}'$, and $\mathrm{C}'=\mathrm{C}'_0\cap\mathrm{B}'$ is the integral
closure of $\mathrm{A}'$ in $\mathrm{L}'$, and hence $\mathrm{C}'$ is the integral closure of $\mathrm{A}'$ in $\mathrm{B}'$.
\end{proof}

\begin{corollary}[6.14.5]
\label{IV.6.14.5}
Let $\mathrm{A}$ be a Noetherian ring, $\mathrm{A}'$ a Noetherian $\mathrm{A}$-algebra
such that the morphism $\operatorname{Spec}(\mathrm{A}')\to\operatorname{Spec}(\mathrm{A})$ is normal. Let $\mathrm{B}$, $\mathrm{C}$ be two $\mathrm{A}$-algebras, $\varphi:\mathrm{B}\to\mathrm{C}$
an $\mathrm{A}$-homomorphism, $\mathrm{D}$ the integral closure of $\mathrm{B}$ in $\mathrm{C}$. If one
sets $\mathrm{B}'=\mathrm{B}\otimes_\mathrm{A}\mathrm{A}'$, $\mathrm{C}'=\mathrm{C}\otimes_\mathrm{A}\mathrm{A}'$, $\mathrm{D}'=\mathrm{D}\otimes_\mathrm{A}\mathrm{A}'$, $\mathrm{D}'$ is the integral closure of $\mathrm{B}'$ in $\mathrm{C}'$.
\end{corollary}

\begin{proof}
Let $(\mathrm{B}_\lambda)$ be the filtered increasing family of sub-$\mathrm{A}$-algebras of finite type of $\mathrm{B}$, and
set $\mathrm{B}'_\lambda=\mathrm{B}_\lambda\otimes_\mathrm{A}\mathrm{A}'$, so that $\mathrm{B}'$ is union of the filtered increasing family of the
sub-$\mathrm{A}'$-algebras of finite type $\mathrm{B}'_\lambda$. Let $\mathrm{D}_\lambda$ be the integral closure of $\mathrm{B}_\lambda$ in $\mathrm{C}$, and set
$\mathrm{D}'_\lambda=\mathrm{D}_\lambda\otimes_\mathrm{A}\mathrm{A}'$, so that $\mathrm{D}$ is union of the $\mathrm{D}_\lambda$ and $\mathrm{D}'$ union of the $\mathrm{D}'_\lambda$; if we prove
that $\mathrm{D}'_\lambda$ is the integral closure of $\mathrm{B}'_\lambda$ in $\mathrm{C}'$, it will follow that $\mathrm{D}'$ is the integral
closure of $\mathrm{B}'$ in $\mathrm{C}'$. Now $\mathrm{B}_\lambda$ and $\mathrm{B}'_\lambda$ are Noetherian and the morphism $\operatorname{Spec}(\mathrm{B}'_\lambda)\to\operatorname{Spec}(\mathrm{B}_\lambda)$
is normal $(6.8.2)$; one can therefore apply \sref{IV.6.14.4} by replacing $\mathrm{A}$, $\mathrm{A}'$ and $\mathrm{B}$ by $\mathrm{B}_\lambda$,
$\mathrm{B}'_\lambda$ and $\mathrm{C}$ respectively, whence the corollary.
\end{proof}

One can for example apply \sref{IV.6.14.5} when $\mathrm{A}$ is its local ring \emph{excellent}
and $\mathrm{A}'$ its completion $\widehat{\mathrm{A}}$, since in this case $\operatorname{Spec}(\mathrm{A}')\to\operatorname{Spec}(\mathrm{A})$ is a regular
morphism $(7.8.2)$.

\subsection{Geometrically unibranch preschemes}
\label{subsection:IV.6.15}

\begin{env}[6.15.1]
\label{IV.6.15.1}
We shall say that a prescheme $\mathrm{X}$ is \emph{unibranch} (resp.\ \emph{geometrically
unibranch}) at a point $x$, or that the point $x$ is unibranch (resp.\ geometrically unibranch)
in $\mathrm{X}$ if the local ring $\sh{O}_{\mathrm{X},x}$ is unibranch (resp.\ geometrically unibranch) $(\mathbf{0},~23.2.1)$;
we say that $\mathrm{X}$ is \emph{unibranch} (resp.\ \emph{geometrically unibranch}) if it is so at every point. It
follows from this definition that, for $\mathrm{X}$ to be unibranch (resp.\ geometrically
unibranch) at a point, it is necessary and it suffices that $\mathrm{X}_{\mathrm{red}}$ be so at this point.
\end{env}

\begin{env}[6.15.2]
\label{IV.6.15.2}
One must take care that with the definitions of \sref{IV.6.15.1}, a \emph{local
ring} $\mathrm{A}$ can be unibranch (resp.\ geometrically unibranch) \emph{without} $\operatorname{Spec}(\mathrm{A})$
being so; in other terms, it can happen that there are prime ideals $\mathfrak{p}$ of $\mathrm{A}$ such
that $\mathrm{A}_\mathfrak{p}$ is not unibranch; it amounts to the same to say that on a prescheme,
the notion of geometrically unibranch point \emph{is not stable by generization}. We
shall see this on the following example.
\end{env}

\oldpage[IV-2]{177}

Let $\mathrm{K}$ be an algebraically closed field of characteristic $0$, $\mathrm{B}$ the integral ring
$\mathrm{K}[\mathrm{U},\mathrm{V},\mathrm{W}]/(\mathrm{U}^2(\mathrm{U}-\mathrm{W})-\mathrm{V}^2(\mathrm{U}+\mathrm{W}))$ ($\mathrm{U}$, $\mathrm{V}$, $\mathrm{W}$ indeterminates) so that
$\mathrm{Y}=\operatorname{Spec}(\mathrm{B})$ is a ``cone with vertex the origin, having a double generatrix''.
We shall designate by $u$, $v$, $w$ the images of $\mathrm{U}$, $\mathrm{V}$, $\mathrm{W}$ in $\mathrm{B}$. Let $\mathrm{R}$ be the field of
fractions of $\mathrm{B}$, and consider in $\mathrm{R}$ the element $t=v(u+w)/u$, which does not belong
to $\mathrm{B}$; let us show that $\mathrm{C}=\mathrm{B}[t]$ is the \emph{integral closure} of $\mathrm{B}$. One has indeed $t^2=u^2-w^2$,
hence $t$ is integral over $\mathrm{B}$, and $v=tu/(u+w)$; the ring $\mathrm{C}_1=\mathrm{K}[t,u,w]$ is integrally
closed, for it is isomorphic to $\mathrm{K}[\mathrm{T},\mathrm{U},\mathrm{W}]/(\mathrm{T}^2-\mathrm{U}^2+\mathrm{W}^2)$ and is therefore the integral
closure of the integrally closed ring $\mathrm{K}[\mathrm{U},\mathrm{W}]$ in the quadratic extension
$\mathrm{K}(\mathrm{U},\mathrm{W})(\sqrt{\mathrm{U}^2-\mathrm{W}^2})$ of its field of fractions (Bourbaki, \emph{Alg.~comm.}, chap.~V, \S~1,
no~6, prop.~18). The ring of fractions $\mathrm{K}[t,u,w,1/(u+w)]$ of $\mathrm{C}_1$ is hence
integrally closed. Set $\mathrm{C}_2=\mathrm{K}[t,v,w]$ and $a=t-v$. The preceding identities give
$t(t-v)^2-w^2(2v-t)=0$, or $t(a^2-w^2)=-2aw^2$.
In $(\mathrm{C}_2)_a$, the identity $(1-t/(2a))(a^2-w^2)=a^2$ shows that
$a^2-w^2$ is invertible. Thus
\[
(\mathrm{C}_2)_a=\mathrm{K}[a,w,1/a,1/(a^2-w^2)],
\]
since $t=-2aw^2/(a^2-w^2)$ and $v=t-a$. Moreover, $u=vw/a$, so
$\mathrm{K}(a,w)=\mathrm{R}$; hence $a$ and $w$ are algebraically independent over $\mathrm{K}$.
This localization of a polynomial ring is integrally closed. In particular,
$\mathrm{K}[t,v,w,1/(t-v)]=\mathrm{K}[t,\frac{tu}{u+w},w,\frac{u+w}{tw}]$ is integrally closed.\footnote{The printed text omits the square on $t-v$ in the equation and asserts that $\mathrm{C}_2$ is integrally closed. The corrected argument proves the property directly for the localization used here.}
Finally, taking into
account the fact that $u$ and $w$ are algebraically independent over $\mathrm{K}$, one proves easily
that $\mathrm{C}=\mathrm{K}[t,u,v,w]=\mathrm{K}[t,u,w,1/(u+w)]\cap\mathrm{K}[t,v,w,1/(t-v)]$, which finishes
showing that $\mathrm{C}$ is the integral closure of $\mathrm{B}$. It is immediate that if $\mathfrak{m}_0$ is the maximal ideal
of $\mathrm{B}$ generated by $u$, $v$, $w$ (``vertex of the cone''), there is only one maximal
ideal $\mathfrak{n}_0$ of $\mathrm{C}$ lying over $\mathfrak{m}_0$, namely the ideal generated by $t$, $u$, $v$, $w$. If one sets $\mathrm{A}=\mathrm{B}_{\mathfrak{m}_0}$, one deduces
easily that $\mathrm{A}'=\mathrm{C}_{\mathfrak{n}_0}$ is the integral closure of $\mathrm{A}$, which is hence unibranch, and hence
also geometrically unibranch since its residue field $\mathrm{K}$ is algebraically closed.
But in $\mathrm{A}$ the prime ideal $\mathfrak{p}$ generated by $u$ and $v$ is such that the integral
closure $\mathrm{A}_\mathfrak{p}[t]$ of $\mathrm{A}_\mathfrak{p}$ is not a local ring.

We shall see further below $(9.7.10)$ that when $\mathrm{X}$ is a locally Noetherian prescheme
such that, if $\mathrm{X}'$ is the normalization of $\mathrm{X}_{\mathrm{red}}$, the canonical morphism $\mathrm{X}'\to\mathrm{X}$
is \emph{finite} (which will be the case if $\mathrm{X}$ is such that the rings of its affine open subsets are
\emph{universally Japanese} $(\mathbf{0},~23.1.1)$)), then the set of points $x\in\mathrm{X}$ where $\mathrm{X}$ is geometrically
unibranch is \emph{locally constructible}.

\begin{env}[6.15.3]
\label{IV.6.15.3}
We shall say for short that a morphism $f:\mathrm{X}\to\mathrm{Y}$ is \emph{radicial at a
point} $y\in\mathrm{Y}$, if $f^{-1}(y)$ is empty or reduced to a single point $x$ and if $\boldsymbol{k}(x)$ is a radicial extension
of $\boldsymbol{k}(y)$. It amounts to the same to say that $f$ is radicial at every point of $\mathrm{Y}$
or that $f$ is radicial $(\mathbf{I},~3.5.8)$. If $g:f^{-1}(y)\to\operatorname{Spec}(\boldsymbol{k}(y))$ is the base change inverse
image of $f$ by the base change $\operatorname{Spec}(\boldsymbol{k}(y))\to\mathrm{Y}$, it amounts to the same to say that $f$
is radicial at the point $y\in\mathrm{Y}$, or that $g$ is \emph{radicial}.
\end{env}

\begin{lemma}[6.15.3.1]
\label{IV.6.15.3.1}
(i) Let $f:\mathrm{X}\to\mathrm{Y}$, $g:\mathrm{Y}\to\mathrm{Z}$ be two morphisms. If $f$ is radicial
at a point $y$ and $g$ radicial at the point $z=g(y)$, then $g\circ f$ is radicial at the point $g(y)$. The converse
is true if $f$ is surjective.

(ii) Let $f:\mathrm{X}\to\mathrm{Y}$, $h:\mathrm{Y}'\to\mathrm{Y}$ be two morphisms, and let $f'=f_{|\mathrm{Y}'}:\mathrm{X}\times_\mathrm{Y}\mathrm{Y}'\to\mathrm{Y}'$.
Let $y'\in\mathrm{Y}'$ and $y=h(y')$; then for $f$ to be radicial at the point $y$, it is necessary and it suffices that $f'$ be
radicial at the point $y'$.
\end{lemma}

\begin{proof}
\oldpage[IV-2]{178}
(i) If $f$ is radicial at the point $y$ and $g$ radicial at the point $z$, $g^{-1}(z)$ is reduced to $y$ and
$f^{-1}(g^{-1}(z))=f^{-1}(y)$ is empty or reduced to a single point $x$; furthermore $\boldsymbol{k}(x)$ is a radicial extension
of $\boldsymbol{k}(y)$ and $\boldsymbol{k}(y)$ a radicial extension of $\boldsymbol{k}(z)$, hence $\boldsymbol{k}(x)$ is a radicial extension of $\boldsymbol{k}(z)$.
Conversely, suppose $f$ surjective; if $g\circ f$ is radicial at the point $z=g(y)$, $g^{-1}(z)$ is reduced
to the point $y$, for otherwise $f^{-1}(g^{-1}(z))$ would have at least two distinct points; furthermore,
$f^{-1}(y)=f^{-1}(g^{-1}(z))$ is reduced to a single point $x$, and by hypothesis one has $\boldsymbol{k}(z)\subset\boldsymbol{k}(y)\subset\boldsymbol{k}(x)$,
and $\boldsymbol{k}(x)$ is a radicial extension of $\boldsymbol{k}(z)$, hence $\boldsymbol{k}(y)$ is a radicial extension of $\boldsymbol{k}(z)$ and $\boldsymbol{k}(x)$ radicial over $\boldsymbol{k}(y)$.

\begingroup\microtypesetup{protrusion=false,expansion=false}
(ii) Let $g:f^{-1}(y)\to\operatorname{Spec}(\boldsymbol{k}(y))$, $g':f'^{-1}(y')\to\operatorname{Spec}(\boldsymbol{k}(y'))$ be the morphisms
inverse images of \mbox{\ensuremath{f}} and $f'$ respectively; as $f'^{-1}(y')=f^{-1}(y)\otimes_{\boldsymbol{k}(y)}\boldsymbol{k}(y')$
$(\mathbf{I},~3.6.4)$, $g'$ is the inverse image of $g$ and the assertion reduces to $(2.6.1,~\text{(v)})$.
\par\endgroup
\end{proof}

Suppose then $\mathrm{X}$ a prescheme \emph{reduced} not having more than a \emph{finite} number of irreducible components,
and let $\mathrm{X}'$ be its \emph{normalization} $(\mathbf{II},~6.3.8)$; we say (\emph{loc.~cit.}) that the canonical morphism
$f:\mathrm{X}'\to\mathrm{X}$ is surjective. The definition \sref{IV.6.15.1} shows then that for $\mathrm{X}$ to be geometrically
unibranch at a point, it is necessary and it suffices that $f$ be \emph{radicial at this point}; for
$\mathrm{X}$ to be geometrically unibranch, it is necessary and it suffices that $f$ be \emph{radicial}.

\begin{env}[6.15.4]
\label{IV.6.15.4}
Generalizing the definition given in $(\mathbf{I},~2.2.9)$, we shall say,
for two reduced preschemes $\mathrm{X}$, $\mathrm{Y}$, a morphism $f:\mathrm{X}\to\mathrm{Y}$ is \emph{birational} if the restriction
of $f$ to the set of maximal points of $\mathrm{X}$ is a \emph{bijection} of this set onto
the set of maximal points of $\mathrm{Y}$, and if, for every maximal point $y$ of $\mathrm{Y}$, the
morphism $\mathrm{X}\times_\mathrm{Y}\operatorname{Spec}(\sh{O}_y)\to\operatorname{Spec}(\sh{O}_y)$ induced by $f$ is an isomorphism (in other words,
if for the fibre $f^{-1}(y)$ reduced to a single point $x$ (maximal in $\mathrm{X}$) and if the homomorphism
$\boldsymbol{k}(y)\to\boldsymbol{k}(x)$ induced by $f$ is bijective); this notion reduces to that of $(\mathbf{I},~2.2.9)$ when $\mathrm{X}$
and $\mathrm{Y}$ have only a \emph{finite} number of irreducible components.
\end{env}

\begin{lemma}[6.15.4.1]
\label{IV.6.15.4.1}
Let $f:\mathrm{X}\to\mathrm{Y}$ be a morphism, $g:\mathrm{Y}'\to\mathrm{Y}$ a flat morphism,
$\mathrm{X}'=\mathrm{X}\times_\mathrm{Y}\mathrm{Y}'$, $f'=f_{|\mathrm{Y}'}:\mathrm{X}'\to\mathrm{Y}'$. If $f_{\mathrm{red}}$ is birational, so is $f'_{\mathrm{red}}$.
\end{lemma}

\begin{proof}
Indeed, if $y'$ is a maximal point of $\mathrm{Y}'$, one knows $(2.3.4,~\text{(ii)})$ that $y=g(y')$ is
a maximal point of $\mathrm{Y}$; there is a unique point $x$ of $\mathrm{X}$ in $f^{-1}(y)$ by hypothesis; since moreover
$\boldsymbol{k}(x)=\boldsymbol{k}(y)$. It follows therefore from $(\mathbf{I},~3.4.9)$ that
$f'^{-1}(y')$ is reduced to a single point $x'$ and that $\boldsymbol{k}(x')=\boldsymbol{k}(y')$. We conclude from the reasoning
that, noting $(2.3.7,~\text{(ii)})$, every irreducible component of $\mathrm{X}'$ dominates
one irreducible component of $\mathrm{Y}'$.
\end{proof}

\begin{proposition}[6.15.5]
\label{IV.6.15.5}
Let $f:\mathrm{X}\to\mathrm{Y}$ be a morphism such that $f_{\mathrm{red}}$ is entire and birational
$(6.15.4)$, hence surjective.
\begin{enumerate}[label=(\roman*)]
  \item For $\mathrm{Y}$ to be geometrically unibranch at a point $y$, it is necessary and it suffices that $f$ be
radicial at the point $y$ and that $\mathrm{X}$ be geometrically unibranch at the unique point $x$ of $f^{-1}(y)$.
  \item For $\mathrm{Y}$ to be geometrically unibranch, it is necessary and it suffices that $\mathrm{X}$ be geometrically
unibranch and that $f$ be radicial.
\end{enumerate}
\end{proposition}

\begin{proof}
One can obviously restrict to the case where $\mathrm{X}$ and $\mathrm{Y}$ are \emph{reduced}; the fact that $f$ is
\emph{surjective} follows from the fact that $f$ is closed $(\mathbf{II},~6.1.10)$ and that $f(\mathrm{X})$ contains
all the maximal points of $\mathrm{Y}$. This also shows immediately that (i) implies (ii).
Using $(\mathbf{I},~3.6.5)$ and $(\mathbf{II},~6.1.5)$, one can suppose, to prove (i), that $\mathrm{Y}=\operatorname{Spec}(\sh{O}_{\mathrm{Y},y})$,
\oldpage[IV-2]{179}
in other words $\mathrm{Y}=\operatorname{Spec}(\mathrm{A})$ (with $\mathrm{A}=\sh{O}_{\mathrm{X},x}$) is a local scheme. Note that $\mathrm{A}'=\sh{O}_{\mathrm{X}',x'}$
is by hypothesis a faithfully flat $\mathrm{A}$-module, hence containing $\mathrm{A}$, and hence
the hypothesis that $\mathrm{X}$ is geometrically unibranch at the point $x$, or the hypothesis that $\mathrm{X}'$ is
geometrically unibranch at the point $x'$, both imply that $\mathrm{A}$ is a local ring
\emph{integral} (since $\mathrm{A}$, being reduced, is also isomorphic to a subring of $\mathrm{A}'_{\mathrm{red}}$).
Let $\mathrm{K}$ be the field of fractions of $\mathrm{A}$, $\mathrm{B}$ the integral closure of $\mathrm{A}$, and set $\mathrm{Y}=\operatorname{Spec}(\mathrm{B})$;
it amounts to the same to say that $\mathrm{X}$ is geometrically unibranch at the point $x$ or that
the morphism $g:\mathrm{Y}\to\mathrm{X}$ is radicial at the point $x$ $(6.15.3)$. Let $\mathrm{Y}'=\mathrm{Y}\otimes_k k'=\mathrm{Y}\times_\mathrm{X}\mathrm{X}'$
so that one has the commutative diagram
\[
\xymatrix{
  \mathrm{Y} \ar[r]^{h} & \mathrm{Y}' \\
  \mathrm{X} \ar[u]^{g} \ar[r]_{f} & \mathrm{X}' \ar[u]_{g'}
}
\]

Note that $f$ is a flat morphism; hence $(6.15.4.1)$ the entire morphism $g'_{\mathrm{red}}$ is
birational. As $\mathrm{Y}$ is normal, $\mathrm{Y}'$ is geometrically unibranch
$(6.15.6)$. For $\mathrm{X}$ to be geometrically unibranch at the point $x$, it is necessary and it suffices therefore
that $g'$ be radicial at the point $x'$ $(6.15.5)$. Now, $x$ is the projection of $x'$ in $\mathrm{X}$ and it is
equivalent to say that $g$ is radicial at the point $x$ or that $g'$ is radicial at the point $x'$
$(6.15.3.1)$; finally, for $\mathrm{X}$ to be geometrically unibranch at this point, it is necessary and it
suffices that $g$ be radicial at this point, whence the proposition.
\end{proof}

\begin{proposition}[6.15.6]
\label{IV.6.15.6}
Let $k$ be a field, $\mathrm{X}$ a $k$-prescheme. If $\mathrm{X}$ is normal, then,
for every extension $k'$ of $k$, $\mathrm{X}'=\mathrm{X}\otimes_k k'$ is \emph{geometrically unibranch}.
\end{proposition}

\begin{proof}
One knows that $k'$ is an algebraic extension of a pure extension $k_0$ of $k$, and if $k''$
is the largest separable extension of $k_0$ contained in $k'$, $k'$ is a radicial extension
of $k''$ and $k''$ a separable extension of $k$. One knows $(6.14.2)$ that $\mathrm{X}''=\mathrm{X}\otimes_k k''$ is normal;
as $\mathrm{X}\otimes_k k'=\mathrm{X}''\otimes_{k''}k'$, one sees that one can restrict to the case where the extension $k'$ of $k$
is \emph{radicial}. Furthermore $(\mathbf{I},~3.6.5)$, one can suppose that $\mathrm{X}=\operatorname{Spec}(\mathrm{A})$, where $\mathrm{A}$ is a
\emph{local integral and integrally closed} ring (since $\mathrm{X}$ is normal). The projection morphism $f:\mathrm{X}'\to\mathrm{X}$
is a universal homeomorphism, since $\operatorname{Spec}(k')\to\operatorname{Spec}(k)$ is a universal homeomorphism
$(2.4.5)$; as $\mathrm{X}'=\operatorname{Spec}(\mathrm{A}')$ where $\mathrm{A}'=\mathrm{A}\otimes_k k'$, one sees therefore that $\mathrm{A}'$ is a
local ring whose nilradical $\mathfrak{N}'$ is the unique minimal prime ideal, whence $\mathrm{X}'_{\mathrm{red}}=\operatorname{Spec}(\mathrm{A}'_0)$,
where $\mathrm{A}'_0=\mathrm{A}'/\mathfrak{N}'$ is a local integral ring; furthermore, if $\mathrm{K}$ is the field of fractions of $\mathrm{A}$,
the field of fractions $\mathrm{K}'_0$ of $\mathrm{A}'_0$ is radicial over $\mathrm{K}$, since the morphism $f$ is radicial.
But $\mathrm{A}'_0$ is integer over $\mathrm{A}$, its integral closure $\mathrm{B}$ is also the integral closure of $\mathrm{A}$,
and one knows (Bourbaki, \emph{Alg.~comm.}, chap.~V,
\S~2, no~3, lemma~3) that $\mathrm{B}$ is the set of $x\in\mathrm{K}'_0$ of which a power $p^m$-th (for $m$
large enough) belongs to $\mathrm{A}$ ($p$ being the characteristic exponent of $\mathrm{K}$); furthermore there exists
only a single prime ideal of $\mathrm{B}$ lying over every prime ideal of $\mathrm{A}$; in particular $\mathrm{B}$
is a local ring and its residue field is a radicial extension of that of $\mathrm{A}$, and
\emph{a fortiori} of the same type for $\mathrm{A}'_0$, which proves that $\mathrm{A}'_0$ is geometrically unibranch, and hence
so is $\mathrm{X}'$ as well.
\end{proof}

\begin{proposition}[6.15.7]
\label{IV.6.15.7}
Let $k$ be a field, $\mathrm{X}$ a $k$-prescheme, $k'$ an extension of $k$,
$\mathrm{X}'=\mathrm{X}\otimes_k k'$. Let $x'$ be a point of $\mathrm{X}'$, $x$ its projection in $\mathrm{X}$. For $\mathrm{X}$ to be geometrically
unibranch at the point $x$, it is necessary and it suffices that $\mathrm{X}'$ be geometrically unibranch at the point $x'$. For
$\mathrm{X}$ to be geometrically unibranch, it is necessary and it suffices that $\mathrm{X}'$ be so.
\end{proposition}

\begin{proof}
The second assertion follows from the first and from the fact that the projection $f:\mathrm{X}'\to\mathrm{X}$
is a surjective morphism. To prove the first assertion, one can, by virtue of
$(\mathbf{I},~5.1.8)$, restrict to the case where $\mathrm{X}$ is reduced, and by virtue of $(\mathbf{I},~3.6.5)$, to the case
\oldpage[IV-2]{180}
where $\mathrm{X}=\operatorname{Spec}(\mathrm{A})$ (with $\mathrm{A}=\sh{O}_{\mathrm{X},x}$) is a local scheme. Note that $\mathrm{A}'=\sh{O}_{\mathrm{X}',x'}$ is
by hypothesis a faithfully flat $\mathrm{A}$-module, hence containing $\mathrm{A}$, and hence
the hypothesis that $\mathrm{X}$ is geometrically unibranch at the point $x$, or the hypothesis that $\mathrm{X}'$ is
geometrically unibranch at the point $x'$, both imply that $\mathrm{A}$ is a local ring
\emph{integral} (since $\mathrm{A}$, being reduced, is also isomorphic to a subring of $\mathrm{A}'_{\mathrm{red}}$).
Let $\mathrm{K}$ be the field of fractions of $\mathrm{A}$, $\mathrm{B}$ the integral closure of $\mathrm{A}$, and set $\mathrm{Y}=\operatorname{Spec}(\mathrm{B})$;
it amounts to the same to say that $\mathrm{X}$ is geometrically unibranch at the point $x$ or that
the morphism $g:\mathrm{Y}\to\mathrm{X}$ is radicial at the point $x$ $(6.15.3)$. Let $\mathrm{Y}'=\mathrm{Y}\otimes_k k'=\mathrm{Y}\times_\mathrm{X}\mathrm{X}'$
so that one has the commutative diagram
\[
\xymatrix{
  \mathrm{Y} \ar[d]_{g} & \mathrm{Y}' \ar[l]^{h} \ar[d]_{g'} \\
  \mathrm{X} & \mathrm{X}' \ar[l]_{f}
}
\]

Note that $f$ is a flat morphism; hence $(6.15.4.1)$ the entire morphism $g'_{\mathrm{red}}$ is
birational. On the other hand, since $\mathrm{Y}$ is normal, $\mathrm{Y}'$ is geometrically unibranch $(6.15.6)$.
For $\mathrm{X}'$ to be geometrically unibranch at $x'$, it is therefore necessary and sufficient that $g'$
be radicial at the point $x'$ $(6.15.5)$. Now $x$ is the projection of $x'$ in $\mathrm{X}$, and it is equivalent
to say that $g$ is radicial at $x$ or that $g'$ is radicial at $x'$ $(6.15.3.1)$; finally, for $\mathrm{X}$
to be geometrically unibranch at $x$, it is necessary and sufficient that $g$ be radicial at this point,
whence the proposition.
\end{proof}

\begin{lemma}[6.15.8]
\label{IV.6.15.8}
Let $k$ be a separably closed field (in other words, such that the
algebraic closure of $k$ is \emph{radicial} over $k$), $\mathrm{X}$ a $k$-prescheme locally of finite type over $k$,
$x$ a closed point of $\mathrm{X}$. If $\mathrm{X}$ is unibranch at the point $x$, it is \emph{geometrically unibranch} at this point.
\end{lemma}

\begin{proof}
Indeed, one knows $(\mathbf{I},~6.4.2)$ that $\boldsymbol{k}(x)$ is an \emph{algebraic} extension of $k$; as the
residue field of the integral closure $(\sh{O}_x)_{\mathrm{red}}$ (integral closure which is by hypothesis
a local ring) is an algebraic extension of $\boldsymbol{k}(x)$, it is a radicial extension of $k$
by hypothesis, hence also of $\boldsymbol{k}(x)$.
\end{proof}

\begin{corollary}[6.15.9]
\label{IV.6.15.9}
Let $k$ be a field, $\mathrm{X}$ a $k$-prescheme locally of finite type over $k$,
$k'$ a separable extension close of $k$ (in other words, such that the algebraic closure of $k$
is radicial over $k'$); then, for $\mathrm{X}$ to be geometrically unibranch, it is necessary and it suffices
that $\mathrm{X}'=\mathrm{X}\otimes_k k'$ be unibranch.
\end{corollary}

\begin{proof}
In view of \sref{IV.6.15.7}, one is reduced to proving that if $\mathrm{X}$ is unibranch and $k$
separably closed, then $\mathrm{X}$ is geometrically unibranch. Now, $\mathrm{X}$ is geometrically
unibranch at its closed points \sref{IV.6.15.8}; we conclude from $(10.4.9)$ that $\mathrm{X}$
is geometrically unibranch at all its points, relying on the fact (pointed out
at the end of \sref{IV.6.15.2}) that the set of points where $\mathrm{X}$ is geometrically unibranch
is constructible (of course, the corollary \sref{IV.6.15.9} will not be used before that point).
\end{proof}

This result justifies, to a certain extent, the terminology of ``geometrically
unibranch''.

\oldpage[IV-2]{181}
\begin{proposition}[6.15.10]
\label{IV.6.15.10}
Let $\mathrm{Y}$, $\mathrm{Y}'$ be two locally Noetherian preschemes, $g:\mathrm{Y}'\to\mathrm{Y}$ a
normal morphism $(6.8.1)$, $f:\mathrm{X}\to\mathrm{Y}$ a morphism; set $\mathrm{X}'=\mathrm{X}\times_\mathrm{Y}\mathrm{Y}'$ and let $p:\mathrm{X}'\to\mathrm{X}$
be the canonical projection, $x'$ a point of $\mathrm{X}'$. Then, if $\mathrm{X}$ is reduced (resp.\ geometrically unibranch,
resp.\ integral and geometrically unibranch) at the point $x=p(x')$, $\mathrm{X}'$ is reduced (resp.\ geometrically
unibranch, resp.\ integral and geometrically unibranch) at the point $x'$.
\end{proposition}

\begin{proof}
By virtue of $(\mathbf{I},~3.6.5)$, one can restrict to the case where $\mathrm{Y}=\operatorname{Spec}(\mathrm{A})$, $\mathrm{Y}'=\operatorname{Spec}(\mathrm{A}')$,
$\mathrm{X}=\operatorname{Spec}(\mathrm{B})$, where $\mathrm{A}$, $\mathrm{A}'$, $\mathrm{B}$ are local rings, the homomorphisms $\mathrm{A}\to\mathrm{A}'$
and $\mathrm{A}\to\mathrm{B}$ being local, $\mathrm{A}$, $\mathrm{A}'$ Noetherian rings, and $x$ being the closed point of $\mathrm{X}$;
it suffices to prove that if $\mathrm{B}$ is reduced (resp.\ geometrically unibranch,
resp.\ integral and geometrically unibranch) then $\mathrm{B}'=\mathrm{B}\otimes_\mathrm{A}\mathrm{A}'$ is reduced (resp.\ $\operatorname{Spec}(\mathrm{B}')$ is geometrically
unibranch at the points of $p^{-1}(x)$, resp.\ $\operatorname{Spec}(\mathrm{B}')$ is integral and geometrically
unibranch at these points). Suppose first $\mathrm{B}$ reduced; $\mathrm{B}$ is
union of the filtered increasing family of its sub-$\mathrm{A}$-algebras of finite type $\mathrm{B}_\lambda$, and by flatness, $\mathrm{B}'$ is
union of the filtered increasing family of the sub-$\mathrm{A}'$-algebras $\mathrm{B}'_\lambda=\mathrm{B}_\lambda\otimes_\mathrm{A}\mathrm{A}'$. Now, the
morphism $\operatorname{Spec}(\mathrm{B}'_\lambda)\to\operatorname{Spec}(\mathrm{B}_\lambda)$ is normal $(6.8.2)$, and \emph{a fortiori} reduced; the fact that $\mathrm{B}'_\lambda$
is reduced is consequence of $(3.3.5)$; one concludes that $\mathrm{B}'$ itself is reduced
by $(5.13.2)$.

Suppose now $\mathrm{B}$ geometrically unibranch; in view of $(\mathbf{I},~5.1.8)$,
one can furthermore suppose $\mathrm{B}$ reduced, hence \emph{integral}. Let $\mathrm{C}$ be its integral closure, which is
by hypothesis a local ring; set $\mathrm{Z}=\operatorname{Spec}(\mathrm{C})$, $\mathrm{C}'=\mathrm{C}\otimes_\mathrm{A}\mathrm{A}'$, $\mathrm{Z}'=\operatorname{Spec}(\mathrm{C}')=\mathrm{Z}\times_\mathrm{X}\mathrm{X}'$,
so that one has the commutative diagram
\[
\xymatrix{
  \mathrm{Z} \ar[d]_{l} & \mathrm{Z}' \ar[l] \ar[d]_{l'} \\
  \mathrm{X} & \mathrm{X}' \ar[l]_{p}
}
\]

By virtue of the first part of the reasoning, $\mathrm{X}'$ is \emph{reduced}; on the other hand, as
$\mathrm{Z}'=\mathrm{Z}\times_\mathrm{Y}\mathrm{Y}'$, it follows from \sref{IV.6.14.1} that $\mathrm{Z}'$ is normal (and \emph{a fortiori} geometrically
unibranch). As $l$ is entire and birational, so is $l'$ by
\sref{IV.6.15.4.1}, since $p$ is flat. Finally, the hypothesis that $\mathrm{X}$ is geometrically
unibranch at the point $x$ implies that $l$ is radicial at this point $(6.15.5)$; we conclude
hence from $(6.15.3.1)$ that $l'$ is radicial at every point $x'\in p^{-1}(x)$, and it follows then from
$(6.15.5)$ that $\mathrm{X}'$ is geometrically unibranch (hence integral since it is reduced)
at each of these points.
\end{proof}

\begin{remarks}[6.15.11]
\label{IV.6.15.11}
(i) In the proof of \sref{IV.6.15.10}, one cannot dispense with
appealing to \sref{IV.6.14.1}, even when $\mathrm{X}=\mathrm{Y}$, for the argument then
uses the integral closure of the ring $\mathrm{B}$, which need not be Noetherian even
when $\mathrm{B}$ is.

(ii) The example $(6.5.5,~\text{(ii)})$ shows that in \sref{IV.6.15.10}, one cannot replace
``geometrically unibranch'' by ``integral'' in the statement, even if the residue field
$\boldsymbol{k}(x)$ is algebraically closed and the morphism $g$ is \'etale.

Nor can ``geometrically unibranch'' be replaced there by ``unibranch''.
Indeed, let
\[
 \mathrm{A}=\mathbb{R}[[\mathrm{U},\mathrm{V}]]/(\mathrm{U}^2+\mathrm{V}^2)
\]
and let $u,v$ be the images of $\mathrm{U},\mathrm{V}$ in $\mathrm{A}$; the maximal
ideal of $\mathrm{A}$ is $\mathrm{A}u+\mathrm{A}v$.
The integral closure of $\mathrm{A}$ is the ring $\mathrm{A}[t]$, where $t=u/v$
satisfies $t^2=-1$, so that $\mathrm{A}[t]$ is isomorphic to the local ring
$\mathbb{C}[[\mathrm{U}]]$.
Since the residue field of $\mathbb{C}[[\mathrm{U}]]$ is $\mathbb{C}$,
$\mathrm{A}$ is unibranch but not geometrically unibranch.
Yet $\mathrm{A}\otimes_{\mathbb{R}}\mathbb{C}$ is not integral: it is isomorphic to
\[
 \mathbb{C}[[\mathrm{U},\mathrm{V}]]/
 ((\mathrm{U}-i\mathrm{V})(\mathrm{U}+i\mathrm{V})).
\]
\end{remarks}
% END NATIVE INPUT source_aligned/ega4-6
% BEGIN NATIVE INPUT source_aligned/ega4-7
\section{Relations between a Noetherian local ring and its completion. Excellent rings}
\label{section:IV.7}

This section is mainly devoted to the exposition of certain properties of
Noetherian rings, generally stable under passage to an algebra of finite type
or under localization, which hold for the rings most often encountered
(algebras of finite type over $\mathbb Z$, or over a field, or over a complete
Noetherian local ring), without holding for all Noetherian rings. A first
series of properties, connected with dimension theory and notably with the
chain condition, is the subject of nos.~1 and 2; all the properties considered
hold for quotients of regular Noetherian rings, and their proofs in this case
are for the most part easy and well known [30]; consequently, for these rings
the developments of \S\S~1 and 2 are of no interest.

This is no longer the case for the properties studied in nos.~3, 4 and 5
(which do not depend on nos.~1 and 2); their proofs, even in the case of
localizations of algebras of finite type over a field or over $\mathbb Z$
(where they are due for the most part to Zariski and Nagata), are most often
delicate. Classical examples are the equivalences
$\mathrm{A}$ normal $\Longleftrightarrow\widehat{\mathrm{A}}$ normal, and
$\mathrm{A}$ reduced $\Longleftrightarrow\widehat{\mathrm{A}}$ reduced. We
develop a systematic method for formulating and proving such properties, with
the aid of the properties of the \emph{fibres} of the canonical morphism
$\Spec(\widehat{\mathrm{A}})\to\Spec(\mathrm{A})$, and using the notions and
results set out in \S~6; the success of this method ultimately rests on the
excellent properties of \emph{complete local rings} and of the rings deduced
from them by localization or passage to an algebra of finite type. In this
regard, Nagata's theorem \sref{IV.6.12.7}, saying that for such a ring the
singular locus in $\Spec(\mathrm{A})$ is closed, plays a crucial role; the same
is true (as concerns permanence properties) of the regularity criterion
\sref[0]{0.22.3.4}, which is more technical in nature.

Finally, in nos.~6 and 7, we apply the results obtained to the study of the
finiteness of the integral closure of Noetherian integral rings.

The reader will note that the most delicate results of \S~7 will be used only
sparingly in the remainder of Chapter~IV, and even in the later chapters.


\subsection{Formal equidimensionality and formally catenary rings}
\label{subsection:IV.7.1}

\oldpage[IV-2]{183}

\begin{definition}[7.1.1]
\label{IV.7.1.1}
We say that a local Noetherian ring $\mathrm{A}$ is \emph{formally equidimensional} if its completion $\widehat{\mathrm{A}}$ is equidimensional \sref[0]{0.16.1.4}.
\end{definition}

\begin{proposition}[7.1.2]
\label{IV.7.1.2}
Let $\mathrm{A}$ be a local Noetherian ring, $\mathfrak{p}_i$ ($1\leqslant i\leqslant n$) its minimal prime ideals.
For $\mathrm{A}$ to be formally equidimensional, it is necessary and sufficient that the $\mathrm{A}/\mathfrak{p}_i$ be formally equidimensional and that $\mathrm{A}$ be equidimensional (in other words, that the $\mathrm{A}/\mathfrak{p}_i$ have the same dimension).
\end{proposition}

\begin{proof}
Indeed, this follows from the fact that for every prime ideal $\mathfrak{p}'$ of $\widehat{\mathrm{A}}$, $\mathfrak{p}'\cap\mathrm{A}$ contains one of the $\mathfrak{p}_i$, hence $\mathfrak{p}'$ contains one of the $\mathfrak{p}_i\widehat{\mathrm{A}}$, and $\widehat{\mathrm{A}}/\mathfrak{p}_i\widehat{\mathrm{A}}=(\mathrm{A}/\mathfrak{p}_i)\otimes_\mathrm{A}\widehat{\mathrm{A}}$ is the completion of the local ring $\mathrm{A}/\mathfrak{p}_i$ \sref[0]{0.7.3.3}, hence has the same dimension.
Every maximal chain of prime ideals of $\widehat{\mathrm{A}}$ therefore identifies canonically with a maximal chain of prime ideals of one of the $\widehat{\mathrm{A}}/\mathfrak{p}_i\widehat{\mathrm{A}}$, and conversely, whence the conclusion.
\end{proof}

\begin{proposition}[7.1.3]
\label{IV.7.1.3}
Let $\mathrm{A}$, $\mathrm{A}'$ be two local Noetherian rings, $\mathfrak{m}$ the maximal ideal of $\mathrm{A}$, $\varphi:\mathrm{A}\to\mathrm{A}'$ a local homomorphism; suppose that $\mathrm{A}'$ is a flat $\mathrm{A}$-module, and that $\mathrm{A}'$ is equidimensional and catenary.
Then:
\begin{enumerate}[label=\emph{\roman*)}]
  \item $\mathrm{A}$ is equidimensional and catenary.
  \item Suppose in addition that $\mathfrak{m}\mathrm{A}'$ is an ideal of definition of $\mathrm{A}'$.
  Then, if $\mathfrak{a}$ is an ideal of $\mathrm{A}$, for $\mathrm{A}'/\mathfrak{a}\mathrm{A}'$ to be equidimensional, it is necessary and sufficient that $\mathrm{A}/\mathfrak{a}\mathrm{A}$ be so.
  In particular, for every prime ideal $\mathfrak{p}$ of $\mathrm{A}$, $\mathrm{A}'/\mathfrak{p}\mathrm{A}'$ is equidimensional.
\end{enumerate}
\end{proposition}

\begin{proof}
Let us first prove the second assertion of~(ii).
Set $\mathrm{X}=\Spec(\mathrm{A})$, $\mathrm{X}'=\Spec(\mathrm{A}')$, $\mathrm{Y}=\mathrm{V}(\mathfrak{p})=\Spec(\mathrm{A}/\mathfrak{p})$, $\mathrm{Y}'=\mathrm{V}(\mathfrak{p}\mathrm{A}')=\Spec(\mathrm{A}'/\mathfrak{p}\mathrm{A}')$; if $f:\mathrm{X}'\to\mathrm{X}$ is the morphism corresponding to $\varphi$, $\mathrm{Y}'$ is also the closed sub-prescheme $f^{-1}(\mathrm{Y})$ of $\mathrm{X}'$.
As $\mathfrak{m}\mathrm{A}'$ is an ideal of definition of $\mathrm{A}'$, we have
\[
  \dim(\mathrm{X})=\dim(\mathrm{X}')
  \tag{7.1.3.1}
\label{IV.7.1.3.1}
\]
by virtue of \sref{IV.6.1.3}; similarly $\mathfrak{m}\mathrm{A}'/\mathfrak{p}\mathrm{A}'$ is an ideal of definition of $\mathrm{A}'/\mathfrak{p}\mathrm{A}'$ and $\mathrm{Y}'$ is flat over $\mathrm{Y}$ \sref{IV.2.1.4}, hence \sref{IV.6.1.3}
\[
  \dim(\mathrm{Y})=\dim(\mathrm{Y}').
  \tag{7.1.3.2}
\label{IV.7.1.3.2}
\]
Let $y$ be the generic point of $\mathrm{Y}$, $\mathrm{Y}'_i$ ($1\leqslant i\leqslant r$) the irreducible components of $\mathrm{Y}'$, $y'_i$ the generic point of $\mathrm{Y}'_i$ ($1\leqslant i\leqslant r$); we have $f(y'_i)=y$ for all $i$ \sref{IV.2.3.4}.
On the other hand, $y'_i$ is the maximal point of the fibre $f^{-1}(y)$ \sref[0]{0.2.1.8}, hence $\sh{O}_{\mathrm{X}',y'_i}\otimes_{\sh{O}_{\mathrm{X},y}}\boldsymbol{k}(y)=\sh{O}_{\mathrm{X}',y'_i}/\mathfrak{m}_y\sh{O}_{\mathrm{X}',y'_i}$ is of dimension $0$, in other words $\mathfrak{m}_y\sh{O}_{\mathrm{X}',y'_i}$ is an ideal of definition of $\sh{O}_{\mathrm{X}',y'_i}$; we can again apply \sref{IV.6.1.3}, which gives
\[
  \dim(\sh{O}_{\mathrm{X}',y'_i})=\dim(\sh{O}_{\mathrm{X},y})
  \label{IV.7.1.3.3}
\tag{7.1.3.3}
\]
that is to say \sref{IV.5.1.2}
\[
  \operatorname{codim}(\mathrm{Y}'_i,\mathrm{X}')=\operatorname{codim}(\mathrm{Y},\mathrm{X})
  \tag{7.1.3.4}
\label{IV.7.1.3.4}
\]
\oldpage[IV-2]{184}
for all $i$.
As $\mathrm{A}'$ is equidimensional and catenary, we have, by virtue of \sref[0]{0.14.3.5}
\[
  \dim(\mathrm{X}')=\dim(\mathrm{Y}'_i)+\operatorname{codim}(\mathrm{Y}'_i,\mathrm{X}')
  \label{IV.7.1.3.5}
\tag{7.1.3.5}
\]
for all $i$; by virtue of {\hyperlink{ega4.c04.eq.7-1-3-1}{(7.1.3.1)}}, {\hyperlink{ega4.c04.eq.7-1-3-2}{(7.1.3.2)}} and {\hyperlink{ega4.c04.eq.7-1-3-4}{(7.1.3.4)}} and the inequality $\dim(\mathrm{Y}'_i)\leqslant\dim(\mathrm{Y}')$, we deduce
\[
  \dim(\mathrm{X})\leqslant\dim(\mathrm{Y})+\operatorname{codim}(\mathrm{Y},\mathrm{X})
  \tag{7.1.3.6}
\label{IV.7.1.3.6}
\]
hence the two sides of this inequality are equal by \sref[0]{0.14.2.2.2}; moreover this equality implies
\[
  \dim(\mathrm{Y}'_i)=\dim(\mathrm{Y}')=\dim(\mathrm{Y})
  \label{IV.7.1.3.7}
\tag{7.1.3.7}
\]
for all $i$, which proves the second assertion of~(ii).

Let us now prove the first assertion of~(ii).
Let $\mathfrak{p}_j$ ($1\leqslant j\leqslant n$) be the minimal prime ideals of $\mathrm{A}$ among those containing $\mathfrak{a}$, and let $\mathfrak{p}'_{jh}$ be the prime ideals of $\mathrm{A}'$ minimal among those containing $\mathfrak{p}_j\mathrm{A}'$ ($1\leqslant h\leqslant m_j$); we know then \sref{IV.2.3.4} that the $\mathfrak{p}'_{jh}$ ($1\leqslant j\leqslant n$, $1\leqslant h\leqslant m_j$ for all $j$) are also the minimal prime ideals among those containing $\mathfrak{a}\mathrm{A}'$.
From what we have just seen, for all $j$, the dimensions of the rings $\mathrm{A}'/\mathfrak{p}'_{jh}$ ($1\leqslant h\leqslant m_j$) are all equal and are equal to $\dim(\mathrm{A}'/\mathfrak{p}_j\mathrm{A}')$, itself equal to $\dim(\mathrm{A}/\mathfrak{p}_j)$ by {\hyperlink{ega4.c04.eq.7-1-3-2}{(7.1.3.2)}}.
To say that $\mathrm{A}'/\mathfrak{a}\mathrm{A}'$ is equidimensional therefore means that the dimensions of the $\mathrm{A}/\mathfrak{p}_j$ are all equal, that is to say that $\mathrm{A}/\mathfrak{a}$ is equidimensional.

Let us finally prove~(i).
Let $\mathfrak{r}'$ be a prime ideal of $\mathrm{A}'$ minimal among those containing $\mathfrak{m}\mathrm{A}'$, and set $\mathrm{A}''=\mathrm{A}'_{\mathfrak{r}'}$; as $\mathrm{A}''$ is a flat $\mathrm{A}'$-module, it is also a flat $\mathrm{A}$-module; moreover $\mathrm{A}''$ is equidimensional and catenary \sref[0]{0.16.1.4}, and $\mathfrak{m}\mathrm{A}''$ is an ideal of definition of $\mathrm{A}''$.
We can therefore restrict ourselves to the case where $\mathfrak{m}\mathrm{A}'$ is an ideal of definition of $\mathrm{A}'$.
If then $\mathfrak{p}$ and $\mathfrak{q}\subset\mathfrak{p}$ are two prime ideals of $\mathrm{A}$, we can apply~(ii) to $\mathrm{A}/\mathfrak{q}$ and $\mathrm{A}'/\mathfrak{q}\mathrm{A}'$, which is equidimensional and flat over $\mathrm{A}/\mathfrak{q}$; if $\mathrm{Z}=\Spec(\mathrm{A}/\mathfrak{q})$ and $\mathrm{Y}=\Spec(\mathrm{A}/\mathfrak{p})$, we have
$\dim(\mathrm{Z})=\dim(\mathrm{Y})+\operatorname{codim}(\mathrm{Y},\mathrm{Z})$;
moreover one can apply {\hyperlink{ega4.c04.eq.7-1-3-6}{(7.1.3.6)}}, which shows that $\operatorname{codim}(\mathrm{Y},\mathrm{X})=\dim\mathrm{X}-\dim\mathrm{Y}$; as $\mathrm{X}$ is equicodimensional, these relations show that it is biequidimensional by virtue of \sref[0]{0.14.3.3}.
\end{proof}

\begin{corollary}[7.1.4]
\label{IV.7.1.4}
Let $\mathrm{A}$ be a local Noetherian ring that is formally equidimensional.
Then:
\begin{enumerate}[label=\emph{\roman*)}]
  \item $\mathrm{A}$ is equidimensional and catenary (in other words, biequidimensional).
  \item Let $\mathfrak{a}$ be an ideal of $\mathrm{A}$; for $\mathrm{A}/\mathfrak{a}$ to be equidimensional, it is necessary and sufficient that $\mathrm{A}/\mathfrak{a}$ be formally equidimensional.
  In particular, for every prime ideal $\mathfrak{p}$ of $\mathrm{A}$, $\mathrm{A}/\mathfrak{p}$ is formally equidimensional.
\end{enumerate}
\end{corollary}

\begin{proof}
If $\mathfrak{m}$ is the maximal ideal of $\mathrm{A}$, $\mathfrak{m}\widehat{\mathrm{A}}$ is an ideal of definition of $\widehat{\mathrm{A}}$.
By hypothesis $\widehat{\mathrm{A}}$ is equidimensional, and on the other hand it is known to be catenary \sref{IV.5.6.4}; it suffices then to apply \sref{IV.7.1.3} with $\mathrm{A}'=\widehat{\mathrm{A}}$.
\end{proof}

\begin{corollary}[7.1.5]
\label{IV.7.1.5}
Let $\mathrm{A}$ be a local Noetherian ring such that there exists a finite type $\mathrm{A}$-module $\mathrm{M}$ that is a Cohen-Macaulay $\mathrm{A}$-module and for which $\operatorname{Supp}(\mathrm{M})=\Spec(\mathrm{A})$ (this holds for example if $\mathrm{A}$ is a Cohen-Macaulay ring).
Then $\mathrm{A}$ is formally equidimensional, hence \sref{IV.7.1.4} for a quotient ring $\mathrm{B}$ of $\mathrm{A}$ to be formally equidimensional, it is necessary and sufficient that it be equidimensional.
\end{corollary}

\begin{proof}
\oldpage[IV-2]{185}
Indeed, $\widehat{\mathrm{M}}=\mathrm{M}\otimes_\mathrm{A}\widehat{\mathrm{A}}$ is a Cohen-Macaulay $\widehat{\mathrm{A}}$-module \sref[0]{0.16.5.2}, with support equal to $\Spec(\widehat{\mathrm{A}})$; hence \sref[0]{0.16.5.4}, $\widehat{\mathrm{A}}$ is equidimensional.
\end{proof}

\begin{remark}[7.1.6]
\label{IV.7.1.6}
We have seen \sref{IV.6.3.8} that under the hypotheses of \sref{IV.7.1.5}, for every quotient ring $\mathrm{B}$ of $\mathrm{A}$, the fibres of the canonical morphism $\Spec(\widehat{\mathrm{B}})\to\Spec(\mathrm{B})$ are preschemes of Cohen-Macaulay; hence \sref{IV.6.4.3} and \sref{IV.5.7.5}, \emph{for $\mathrm{B}$ to have no embedded associated prime ideals, it is necessary and sufficient that the same hold for $\widehat{\mathrm{B}}$}.
\end{remark}

\begin{lemma}[7.1.7]
\label{IV.7.1.7}
Let $\mathrm{A}$ be a local Noetherian ring, $\mathfrak{p}_i$ ($1\leqslant i\leqslant r$) its minimal prime ideals.
Suppose that each of the rings $\mathrm{A}/\mathfrak{p}_i$ is formally equidimensional.
Then, for every prime ideal $\mathfrak{q}$ of $\mathrm{A}$ and every $i$ such that $\mathfrak{p}_i\subset\mathfrak{q}$, the ring $\mathrm{A}_\mathfrak{q}/\mathfrak{p}_i\mathrm{A}_\mathfrak{q}$ is formally equidimensional.
\end{lemma}

\begin{proof}
As $\mathrm{A}_\mathfrak{q}/\mathfrak{p}_i\mathrm{A}_\mathfrak{q}$ is the local ring of $\mathrm{A}/\mathfrak{p}_i$ at the prime ideal $\mathfrak{q}/\mathfrak{p}_i$, we can restrict ourselves to the case where $\mathrm{A}$ is integral and formally equidimensional.
Set $\mathrm{A}'=\widehat{\mathrm{A}}$, and let $\mathfrak{q}'$ be one of the prime ideals of $\mathrm{A}'$ minimal among those containing $\mathfrak{q}\mathrm{A}'$; if we set $\mathrm{B}=\mathrm{A}_\mathfrak{q}$, $\mathrm{B}'=\mathrm{A}'_{\mathfrak{q}'}$, then $\mathrm{B}'$ is a flat $\mathrm{B}$-module \sref[0]{0.6.3.2}.
Set $\mathrm{C}=\widehat{\mathrm{B}}$, $\mathrm{C}'=\widehat{\mathrm{B}'}$; as $\mathrm{B}'$ is a flat $\mathrm{B}$-module, we know that $\mathrm{C}'$ is a flat $\mathrm{C}$-module (Bourbaki, \emph{Alg.\ comm.}, chap.~III, \textsection\,5, n\textsuperscript{o}\,4, prop.~4).
As $\mathrm{A}'$ is catenary \sref{IV.5.6.4} and equidimensional by hypothesis, the same holds for $\mathrm{B}'=\mathrm{A}'_{\mathfrak{q}'}$ \sref[0]{0.16.1.4}; moreover, as $\mathrm{A}'$ is isomorphic to a quotient of a regular ring by the theorem of Cohen \sref[0]{0.19.8.8}, the same holds for $\mathrm{B}'$ \sref[0]{0.17.3.9}; we conclude therefore from \sref{IV.7.1.5} that $\mathrm{C}'$ is equidimensional.
On the other hand, $\mathrm{C}'$ is catenary \sref{IV.5.6.4}, hence $\mathrm{C}$ is equidimensional by virtue of \sref{IV.7.1.3},~(ii).
\end{proof}

\begin{theorem}[7.1.8]
\label{IV.7.1.8}
Let $\mathrm{A}$ be a local Noetherian ring.
The following conditions are equivalent:
\begin{enumerate}[label=\emph{\alph*)}]
  \item Every integral quotient ring of $\mathrm{A}$ is formally equidimensional.
  \item The quotient rings of $\mathrm{A}$ by its minimal prime ideals are formally equidimensional.
  \item Every local ring $\mathrm{B}$ that is an A-algebra essentially of finite type \sref{IV.1.3.8} and that is equidimensional, is formally equidimensional.
\end{enumerate}
\end{theorem}

\begin{proof}
It is trivial that $c)$ implies $a)$ and that $a)$ implies $b)$.
On the other hand, as every prime ideal of $\mathrm{A}$ contains a minimal prime ideal of $\mathrm{A}$, $b)$ implies $a)$ by virtue of \sref{IV.7.1.4},~(ii).
It remains to see that $a)$ implies $c)$.

It suffices to prove that the quotients of $\mathrm{B}$ by its minimal prime ideals are formally equidimensional \sref{IV.7.1.2}; as every quotient ring of $\mathrm{B}$ is also an A-algebra essentially of finite type \sref{IV.1.3.9}, we can first suppose $\mathrm{B}$ \emph{integral}.
If $\mathfrak{q}$ is the inverse image in $\mathrm{A}$ of the maximal ideal of $\mathrm{B}$, we know that $\mathrm{B}$ is also an $\mathrm{A}_\mathfrak{q}$-algebra essentially of finite type \sref{IV.1.3.10}, and it follows from $a)$ and from \sref{IV.7.1.7} that every integral quotient of $\mathrm{A}_\mathfrak{q}$ is formally equidimensional; we can therefore suppose that the homomorphism $\mathrm{A}\to\mathrm{B}$ is a \emph{local} homomorphism.
The kernel $\mathfrak{r}$ of this homomorphism is then a prime ideal of $\mathrm{A}$, and by virtue of $a)$, every integral quotient ring of $\mathrm{A}/\mathfrak{r}$ is formally equidimensional; as $\mathrm{B}$ is an $(\mathrm{A}/\mathfrak{r})$-algebra essentially of finite type, we see that we can also suppose
\oldpage[IV-2]{186}
that $\mathrm{A}$ is \emph{integral} and is a subring of $\mathrm{B}$.
We know \sref{IV.1.3.11} that $\mathrm{B}$ is a quotient of a local ring of the form $\mathrm{C}_\mathfrak{p}$, where $\mathrm{C}=\mathrm{A}[\mathrm{T}_1,\ldots,\mathrm{T}_n]$ is a polynomial algebra and $\mathfrak{p}$ is a prime ideal of $\mathrm{C}$ above the maximal ideal $\mathfrak{m}$ of $\mathrm{A}$.
By virtue of \sref{IV.7.1.7}, it suffices to prove that $\mathrm{C}_\mathfrak{p}$ is formally equidimensional, in other words one can restrict to the case where $\mathrm{B}=\mathrm{C}_\mathfrak{p}$.

Set $\mathrm{A}'=\widehat{\mathrm{A}}$, and $\mathrm{C}'=\mathrm{A}'[\mathrm{T}_1,\ldots,\mathrm{T}_n]=\mathrm{C}\otimes_\mathrm{A}\mathrm{A}'$; there is a unique prime ideal $\mathfrak{p}'$ of $\mathrm{C}'$ above $\mathfrak{p}$, hence above the maximal ideal $\mathfrak{m}\mathrm{A}'$ of $\mathrm{A}'$; set $\mathrm{B}'=\mathrm{C}'_{\mathfrak{p}'}$; the homomorphism $\mathrm{B}\to\mathrm{B}'$ is local and makes $\mathrm{B}'$ a flat B-module.
We know that $\widehat{\mathrm{B}'}$ is a flat $\widehat{\mathrm{B}}$-module (Bourbaki, \emph{Alg.\ comm.}, chap.~III, \textsection\,5, n\textsuperscript{o}\,4, prop.~4); as $\widehat{\mathrm{B}'}$ is catenary \sref{IV.5.6.4}, it will suffice to prove that $\widehat{\mathrm{B}'}$ is \emph{equidimensional}, in order to deduce by \sref{IV.7.1.3},~(i) that $\widehat{\mathrm{B}}$ is so as well, which will complete the proof.

Now, $\mathrm{A}'$ is a quotient of a regular ring by the theorem of Cohen \sref[0]{0.19.8.8}, so the same holds for $\mathrm{B}'$ \sref[0]{0.17.3.9}; by virtue of \sref{IV.7.1.5}, to show that $\mathrm{B}'$ is formally equidimensional, it suffices to prove that $\mathrm{B}'$ is equidimensional; for this it suffices to show that $\mathrm{C}'$ is equidimensional, since $\mathrm{C}'$ is a quotient of a regular ring, hence catenary \sref{IV.5.6.4}.
But the minimal prime ideals of $\mathrm{C}'$ are the ideals $\mathfrak{q}'\mathrm{C}'$, where $\mathfrak{q}'$ ranges over the minimal prime ideals of $\mathrm{A}'$ \sref{IV.5.5.3}, and we have $\mathrm{C}'/\mathfrak{q}'\mathrm{C}'=(\mathrm{A}'/\mathfrak{q}')[\mathrm{T}_1,\ldots,\mathrm{T}_n]$.
As $\mathrm{A}'$ is equidimensional by hypothesis (since $\mathrm{A}$ is integral), the same holds for $\mathrm{C}'$ by \sref{IV.5.5.4}.
\end{proof}

\begin{definition}[7.1.9]
\label{IV.7.1.9}
When the equivalent conditions of \sref{IV.7.1.8} are satisfied, we say that $\mathrm{A}$ is a \emph{formally catenary} ring.
\end{definition}

For a local Noetherian ring, it is therefore equivalent to say that it is formally equidimensional or formally catenary and equidimensional, by virtue of \sref{IV.7.1.2}.

\begin{proposition}[7.1.10]
\label{IV.7.1.10}
Every local quotient ring of a local Cohen-Macaulay ring is formally catenary.
\end{proposition}

\begin{proof}
This follows immediately from \sref{IV.7.1.8} and \sref{IV.7.1.5} applied to the integral quotient rings of $\mathrm{A}$.
\end{proof}

\begin{proposition}[7.1.11]
\label{IV.7.1.11}
\begin{enumerate}[label=\emph{\roman*)}]
  \item A local Noetherian ring that is formally catenary is universally catenary.
  \item If $\mathrm{A}$ is a local Noetherian ring that is formally catenary, every local ring $\mathrm{B}$ that is an A-algebra essentially of finite type is formally catenary.
\end{enumerate}
\end{proposition}

\begin{proof}
Assertion~(ii) follows immediately from condition~$c)$ of \sref{IV.7.1.8} and from the fact that if $\mathrm{C}$ is a B-algebra essentially of finite type, $\mathrm{C}$ is also an A-algebra essentially of finite type \sref{IV.1.3.9}.
For the proof of~(i), let us first note that by virtue of condition~$b)$ of \sref{IV.7.1.8}, a local Noetherian ring that is formally catenary $\mathrm{A}$ is catenary, since the quotients of $\mathrm{A}$ by its minimal prime ideals are so \sref{IV.7.1.4}.
If now $\mathrm{E}$ is an A-algebra of finite type, it follows from what precedes and from~(ii) that every local ring of $\mathrm{E}$ at a prime ideal is catenary, hence $\mathrm{E}$ is catenary.
Therefore $\mathrm{A}$ is universally catenary.
\end{proof}

\oldpage[IV-2]{187}

\begin{remarks}[7.1.12]
\label{IV.7.1.12}
\begin{enumerate}[label=\emph{\roman*)}]
  \item We do not know if, conversely, a local Noetherian ring that is universally catenary is always formally catenary.

  \item A local Noetherian ring $\mathrm{A}$ of dimension $1$ is necessarily equidimensional, since its maximal ideal cannot be minimal (for in that case $\mathrm{A}$ would have dimension~$0$).
  As $\dim(\widehat{\mathrm{A}})=1$, this shows that $\mathrm{A}$ is even \emph{formally equidimensional}, and \emph{a fortiori} formally catenary (hence \emph{universally catenary}).
  By contrast, the local ring of dimension $2$ defined in \sref{IV.5.6.11}, which is catenary but not universally catenary, is \emph{a fortiori} not formally catenary.
\end{enumerate}
\end{remarks}

\begin{corollary}[7.1.13]
\label{IV.7.1.13}
Let $\mathrm{Y}$ be an irreducible prescheme locally Noetherian of dimension~$1$, $\mathrm{X}$ an irreducible prescheme, $f:\mathrm{X}\to\mathrm{Y}$ a dominant morphism of finite type, $\xi$, $\eta$ the generic points of $\mathrm{X}$ and $\mathrm{Y}$ respectively.
Then, for all $y\in f(\mathrm{X})$, the dimensions of the irreducible components of $f^{-1}(y)$ are all equal to $\deg.\operatorname{tr}_{\boldsymbol{k}(\eta)}\boldsymbol{k}(\xi)$.
\end{corollary}

\begin{proof}
By hypothesis, $f^{-1}(\eta)$ is irreducible of generic point $\xi$ and of dimension $\deg.\operatorname{tr}_{\boldsymbol{k}(\eta)}\boldsymbol{k}(\xi)$ \sref{IV.5.2.1}.
As $\mathrm{Y}$ is irreducible and of dimension~$1$, we have $\dim(\sh{O}_y)=1$ for all $y\neq\eta$, and $\sh{O}_y$ is universally catenary by virtue of \sref{IV.7.1.12},~(ii).
If $y\in f(\mathrm{X})$ and if $z$ is a generic point of an irreducible component $\mathrm{Z}$ of $f^{-1}(y)$, we have therefore, by \sref{IV.5.6.5}
\[
  \dim(\mathrm{Z})=1+\dim(f^{-1}(\eta))-\dim(\sh{O}_{\mathrm{X},z}).
\]
But we have $\dim(\sh{O}_{\mathrm{X},z})\leqslant\dim(\sh{O}_y)$ \sref[0]{0.16.3.9}, and on the other hand, as $z$ is not the generic point of $\mathrm{X}$, $\dim(\sh{O}_{\mathrm{X},z})>0$; hence $\dim(\sh{O}_{\mathrm{X},z})=1$ and $\dim(\mathrm{Z})=\dim(f^{-1}(\eta))$.
\end{proof}


\newpage % Keep this contents destination with the following page-top heading.
\subsection{Strictly formally catenary rings}
\label{subsection:IV.7.2}

\begin{env}[7.2.1]
\label{IV.7.2.1}
\emph{Notations.} ---
For an integral Noetherian ring $\mathrm{A}$, we denote by $\mathrm{A}^{(1)}$ the intersection of the local rings $\mathrm{A}_\mathfrak{p}$, where $\mathfrak{p}$ ranges over the set of prime ideals of $\mathrm{A}$ of height~$1$ \sref{IV.5.10.17}.
If $\mathrm{A}$ is a \emph{local} Noetherian integral ring, we denote by $\mathrm{A}^{(\omega)}$ the intersection of the $\mathrm{A}_\mathfrak{p}$ where this time $\mathfrak{p}$ ranges over the set of all the prime ideals of $\mathrm{A}$ \emph{distinct from the maximal ideal}.

We will say that a local Noetherian ring $\mathrm{A}$ is \emph{strictly equidimensional} if, for every prime ideal $\mathfrak{p}\in\operatorname{Ass}(\mathrm{A})$, we have $\dim(\mathrm{A}/\mathfrak{p})=\dim(\mathrm{A})$; this amounts to saying that $\mathrm{A}$ is equidimensional and \emph{without embedded associated prime ideals}.
\end{env}

\begin{env}[7.2.1.1]
\label{IV.7.2.1.1}
\emph{Example.} ---
Let $\mathrm{A}$ be a local Noetherian ring \emph{of dimension}~$1$.
Then the following conditions are equivalent:
\begin{enumerate}[label=\emph{\alph*)}]
  \item $\mathrm{A}$ has no embedded associated prime ideals.
  \item[\emph{a$'$)}] $\widehat{\mathrm{A}}$ has no embedded associated prime ideals.
  \item $\mathrm{A}$ is strictly equidimensional.
  \item[\emph{b$'$)}] $\widehat{\mathrm{A}}$ is strictly equidimensional.
  \item $\mathrm{A}$ is a Cohen-Macaulay ring.
\end{enumerate}
\end{env}

\oldpage[IV-2]{188}

Indeed, $a)$ means that the maximal ideal $\mathfrak{m}$ of $\mathrm{A}$ is not associated to $\mathrm{A}$, hence the associated prime ideals to $\mathrm{A}$ are the minimal ideals of $\mathrm{A}$, all distinct from $\mathfrak{m}$, and for such an ideal $\mathfrak{p}$, we necessarily have $\dim(\mathrm{A}/\mathfrak{p})=1$; hence $a)$ implies $b)$.
Conversely, $b)$ implies that every associated prime ideal $\mathfrak{p}\in\operatorname{Ass}(\mathrm{A})$ is distinct from $\mathfrak{m}$, hence minimal, and therefore $b)$ implies $a)$.
We already know that $a)$ and $c)$ are equivalent for a local ring of dimension~$1$ \sref{IV.5.7.8}.
As it comes to the same thing to say that $\mathrm{A}$ is a Cohen-Macaulay ring or that $\widehat{\mathrm{A}}$ is a Cohen-Macaulay ring \sref[0]{0.16.5.2}, and since $\dim(\widehat{\mathrm{A}})=1$, we see finally that $a')$ and $b')$ are equivalent to $c)$.

\begin{proposition}[7.2.2]
\label{IV.7.2.2}
Let $\mathrm{A}$ be a local Noetherian integral ring; set $\mathrm{A}'=\widehat{\mathrm{A}}$, $\mathrm{X}=\Spec(\mathrm{A})$, $\mathrm{X}'=\Spec(\mathrm{A}')$ and denote by $f:\mathrm{X}'\to\mathrm{X}$ the canonical morphism.
Let $a$ be the closed point of $\mathrm{X}$, $j:\mathrm{X}-\{a\}\to\mathrm{X}$ the canonical injection.
Let $\sh{F}$ be a coherent $\sh{O}_\mathrm{X}$-Module, $\sh{F}'=f^*(\sh{F})=\sh{F}\otimes_{\sh{O}_\mathrm{X}}\sh{O}_{\mathrm{X}'}$; then the following conditions are equivalent:
\begin{enumerate}[label=\emph{\alph*)}]
  \item The $\sh{O}_\mathrm{X}$-Module $j_*(\sh{F}|\mathrm{X}-\{a\})$ is coherent.
  \item For all $x'\in\operatorname{Ass}(\sh{F}')$, we have $\dim(\overline{\{x'\}})\geqslant 2$.
\end{enumerate}
\end{proposition}

\begin{proof}
Let $a'$ be the closed point of $\mathrm{X}'$, which is the unique point of the fibre $f^{-1}(a)$, and let $j':\mathrm{X}'-\{a'\}\to\mathrm{X}'$ be the canonical injection.
As the morphism $f$ is faithfully flat and quasi-compact, it is equivalent to say that $j_*(\sh{F}|\mathrm{X}-\{a\})$ is coherent or that $j'_*(\sh{F}'|\mathrm{X}'-\{a'\})$ is coherent.
On the other hand, as $\mathrm{A}'$ is isomorphic to a quotient of a regular local ring by the theorem of Cohen \sref[0]{0.19.8.8}, it follows from \sref{IV.5.11.4} that condition~$b)$ of the statement is equivalent to the fact that $j'_*(\sh{F}'|\mathrm{X}'-\{a'\})$ is coherent.
Whence the proposition.

Instead of applying \sref{IV.5.11.4} and using the theorem of Cohen, which (by \sref{IV.5.11.2}) implicitly appeals to the cohomological results of chap.~III, one can also use the fact that $\mathrm{A}'$ is universally catenary \sref{IV.5.6.4} and universally Japanese by virtue of \sref{IV.7.6.5} (the proof of this last result not using \sref{IV.7.2.2}).
\end{proof}

\begin{proposition}[7.2.3]
\label{IV.7.2.3}
Let $\mathrm{A}$ be a local Noetherian integral ring.
The notations being as in \sref{IV.7.2.2}, the following conditions are equivalent:
\begin{enumerate}[label=\emph{\alph*)}]
  \item $\mathrm{A}^{(1)}$ is a finite A-algebra.
  \item For every closed subset $\mathrm{T}$ of $\mathrm{X}$ of codimension $\geqslant 2$, if we denote by $i:\mathrm{X}-\mathrm{T}\to\mathrm{X}$ the canonical injection, $i_*(\sh{O}_{\mathrm{X}-\mathrm{T}})$ is a coherent $\sh{O}_\mathrm{X}$-Module.
  \item For all $x'\in\operatorname{Ass}(\sh{O}_{\mathrm{X}'})$ and every closed subset $\mathrm{T}$ of $\mathrm{X}$ of codimension $\geqslant 2$, we have $\operatorname{codim}(f^{-1}(\mathrm{T})\cap\overline{\{x'\}},\overline{\{x'\}})\geqslant 2$.
\end{enumerate}
\end{proposition}

\begin{proof}
Set $\mathrm{Z}=\mathrm{Z}^{(2)}(\mathrm{X})$ \sref{IV.5.10.13} and $\mathrm{Z}'=f^{-1}(\mathrm{Z})$, which are stable parts under specialisation; the conditions~$a)$ and~$b)$ are equivalent respectively to the following properties:
$a_1$) $\sh{H}^0_{\mathrm{X}/\mathrm{Z}}(\sh{O}_\mathrm{X})$ is coherent;
$b_1$) $\sh{H}^0_{\mathrm{X}/\mathrm{T}}(\sh{O}_\mathrm{X})$ is coherent for every closed subset $\mathrm{T}$ of $\mathrm{X}$ of codimension $\geqslant 2$.
Taking into account \sref{IV.5.9.5}, these two properties are equivalent respectively to:
$a'_1$) $\sh{H}^0_{\mathrm{X}'/\mathrm{Z}'}(\sh{O}_{\mathrm{X}'})$ is coherent;
$b'_1$) setting $\mathrm{T}'=f^{-1}(\mathrm{T})$, $\sh{H}^0_{\mathrm{X}'/\mathrm{T}'}(\sh{O}_{\mathrm{X}'})$ is coherent for every closed subset $\mathrm{T}$ of $\mathrm{X}$ of codimension $\geqslant 2$.
Now, every point of $\operatorname{Ass}(\sh{O}_{\mathrm{X}'})$ projects to the generic point of $\mathrm{X}$ since $f$ is a flat morphism \sref{IV.3.3.2}; as by definition this point does not belong to $\mathrm{Z}$, we see that $\operatorname{Ass}(\sh{O}_{\mathrm{X}'})$ does not meet $\mathrm{Z}'$; the equivalence of $a'_1$ and $b'_1$ therefore follows from \sref{IV.5.11.5}, since $\mathrm{A}'$ is isomorphic to a quotient of a regular local ring by the theorem of Cohen \sref[0]{0.19.8.8} (or alternatively
\oldpage[IV-2]{189}
because $\mathrm{A}'$ is universally Japanese \sref{IV.7.6.5}).
For this reason, conditions~$b'_1$ and~$c)$ are also equivalent by virtue of \sref{IV.5.11.4}.
\end{proof}

\begin{proposition}[7.2.4]
\label{IV.7.2.4}
Let $\mathrm{A}$ be a local Noetherian ring and set $\mathrm{X}=\Spec(\mathrm{A})$.
The following conditions are equivalent:
\begin{enumerate}[label=\emph{\alph*)}]
  \item For every integral quotient ring $\mathrm{B}$ of $\mathrm{A}$, the ring $\mathrm{B}^{(1)}$ is a finite B-algebra.
  \item For every coherent $\sh{O}_\mathrm{X}$-Module $\sh{F}$ and every subset $\mathrm{Z}$ stable under specialisation, and such that for all $x\in\operatorname{Ass}(\sh{F})\cap(\mathrm{X}-\mathrm{Z})$, we have $\operatorname{codim}(\overline{\{x\}}\cap\mathrm{Z},\overline{\{x\}})\geqslant 2$, the $\sh{O}_\mathrm{X}$-Module $\sh{H}^0_{\mathrm{X}/\mathrm{Z}}(\sh{F})$ is coherent.
  \item For every closed subset $\mathrm{T}$ of $\mathrm{X}$ and every coherent $\sh{O}_\mathrm{U}$-Module $\sh{G}$ (where $\mathrm{U}=\mathrm{X}-\mathrm{T}$) such that, for all $x\in\operatorname{Ass}(\sh{G})$, we have $\operatorname{codim}(\overline{\{x\}}\cap\mathrm{T},\overline{\{x\}})\geqslant 2$, $i_*(\sh{G})$ (where $i:\mathrm{U}\to\mathrm{X}$ is the canonical injection) is a coherent $\sh{O}_\mathrm{X}$-Module.
  \item For every integral quotient ring $\mathrm{B}$ of $\mathrm{A}$, and every ideal $\mathfrak{J}$ of height $\geqslant 2$ in $\mathrm{B}$, the ring $\bigcap_{\mathfrak{p}\not\supset\mathfrak{J}}\mathrm{B}_\mathfrak{p}$ is a finite B-algebra.
  \item For every integral quotient ring $\mathrm{B}$ of $\mathrm{A}$ and every local ring $\mathrm{C}=\mathrm{B}_\mathfrak{q}$ at a prime ideal $\mathfrak{q}$ of $\mathrm{B}$, with $\dim(\mathrm{C})\geqslant 2$, the ring $\mathrm{C}^{(\omega)}$ is a finite C-algebra (or, what amounts to the same, if $\mathrm{Y}=\Spec(\mathrm{C})$ and $\mathrm{U}$ is the complement in $\mathrm{Y}$ of the closed point of $\mathrm{C}$ and $j:\mathrm{U}\to\mathrm{Y}$ the canonical injection, $j_*(\sh{O}_\mathrm{U})$ is a coherent $\sh{O}_\mathrm{Y}$-Module).
  \item For every integral quotient ring $\mathrm{B}$ of $\mathrm{A}$, every local ring $\mathrm{C}=\mathrm{B}_\mathfrak{q}$ at a prime ideal $\mathfrak{q}$ of $\mathrm{B}$, with $\dim(\mathrm{C})\geqslant 2$, and every ideal $\mathfrak{t}\in\operatorname{Ass}(\widehat{\mathrm{C}})$, we have $\dim(\widehat{\mathrm{C}}/\mathfrak{t})\geqslant 2$.
\end{enumerate}
\end{proposition}

\begin{proof}
We already know \sref{IV.5.11.6} that $a)$ and $b)$ are equivalent, as are $c)$ and $d)$, and that $a)$ implies $d)$.
The equivalence of $a)$ and $d)$ in the present case follows from \sref{IV.7.2.3} applied to an integral quotient ring $\mathrm{B}$ of $\mathrm{A}$, condition~$d)$ being an equivalent formulation of condition~\sref{IV.7.2.3},~$b)$ by virtue of \sref{IV.5.9.3.1}.
The equivalence of $e)$ and $f)$ is a consequence of \sref{IV.7.2.2} applied to the coherent $\sh{O}_\mathrm{Y}$-Module $\sh{O}_\mathrm{Y}$ itself.
We have already remarked \sref{IV.5.11.7},~(ii) that if $\mathrm{A}$ satisfies~$a)$, the same holds for every finite A-algebra and every ring of fractions of $\mathrm{A}$.
Condition~$a)$ for $\mathrm{A}$ therefore implies (with the notations of~$e)$) that the ring $\mathrm{C}$ satisfies~$a)$, hence also $c)$; but as $\mathrm{C}$ is integral and $\dim(\mathrm{C})\geqslant 2$, one can apply condition~$c)$ to $\mathrm{C}$ by taking $\mathrm{T}$ to be the closed point of $\mathrm{Y}=\Spec(\mathrm{C})$, which proves that $a)$ implies $e)$.
It remains to prove that $e)$ implies $a)$;
we can obviously restrict ourselves to the case where $\mathrm{A}$ is integral and $\mathrm{B}=\mathrm{A}$, and it amounts to showing that condition~$e)$ implies the condition \sref{IV.7.2.3},~$c)$).
Let us consider the notations of \sref{IV.7.2.3},~$c)$), and let $y'\in f^{-1}(\mathrm{T})\cap\overline{\{x'\}}$, and set $y=f(y')$ and $\mathrm{C}=\sh{O}_{\mathrm{X},y}$; as $y\in\mathrm{T}$, $\dim(\mathrm{C})\geqslant 2$, and hence (with the notations of $e)$), $j_*(\sh{O}_\mathrm{U})$ is a coherent $\sh{O}_\mathrm{Y}$-Module.
Set $\mathrm{Y}'=\mathrm{X}'\times_\mathrm{X}\mathrm{Y}$; since the morphism $f:\mathrm{X}'\to\mathrm{X}$ is flat, so is $g=f_{(\mathrm{Y})}:\mathrm{Y}'\to\mathrm{Y}$.
The underlying space is identified with $f^{-1}(\mathrm{Y})$ \sref[I]{I.3.6.5}, and as $\mathrm{A}$ is integral and $f$ flat, $\operatorname{Ass}(\sh{O}_{\mathrm{X}'})$ is contained in the fibre of the generic point of $\mathrm{X}$ \sref{IV.3.3.2}; this last being contained in $\mathrm{Y}$, we have $x'\in\mathrm{Y}'$.
Let $\mathrm{U}'=g^{-1}(\mathrm{U})$, and let $j'$ be the canonical injection $\mathrm{U}'\to\mathrm{Y}'$; as $j_*(\sh{O}_\mathrm{U})$ is a coherent $\sh{O}_\mathrm{Y}$-Module and $g$ is flat, it follows from \sref{IV.5.9.4} that $j'_*(\sh{O}_{\mathrm{U}'})$ is a coherent $\sh{O}_{\mathrm{Y}'}$-Module.
But we have $x'\in\operatorname{Ass}(\sh{O}_{\mathrm{U}'})$, hence we conclude from \sref{IV.5.10.10} that, in $\mathrm{Y}'$,
\oldpage[IV-2]{190}
$\operatorname{codim}(\overline{\{y'\}}\cap\overline{\{x'\}},\overline{\{x'\}})\geqslant 2$.
As $y'$ is arbitrary in $f^{-1}(\mathrm{T})\cap\overline{\{x'\}}$, we indeed have
$\operatorname{codim}(f^{-1}(\mathrm{T})\cap\overline{\{x'\}},\overline{\{x'\}})\geqslant 2$ in $\mathrm{X}'$.
\end{proof}

\begin{theorem}[7.2.5]
\label{IV.7.2.5}
Let $\mathrm{A}$ be a local Noetherian ring.
The following conditions are equivalent:
\begin{enumerate}[label=\emph{\alph*)}]
  \item For every integral quotient ring $\mathrm{B}$ of $\mathrm{A}$, the completion $\widehat{\mathrm{B}}$ is strictly equidimensional \sref{IV.7.2.1}.
  \item $\mathrm{A}$ is formally catenary \sref{IV.7.1.9} and the fibres of the canonical morphism $\Spec(\widehat{\mathrm{A}})\to\Spec(\mathrm{A})$ satisfy $(\mathrm{S}_1)$ (in other words, have no embedded associated prime cycle).
  \item $\mathrm{A}$ is universally catenary \sref{IV.5.6.2} and for every integral quotient ring $\mathrm{B}$ of $\mathrm{A}$, the ring $\mathrm{B}^{(1)}$ is a finite B-algebra (cf.\ \sref{IV.7.2.4}).
  \item $\mathrm{A}$ is universally catenary and for every integral quotient ring $\mathrm{B}$ of $\mathrm{A}$ and every local ring $\mathrm{C}=\mathrm{B}_\mathfrak{q}$ at a prime ideal $\mathfrak{q}$ of $\mathrm{B}$, with $\dim(\mathrm{C})\geqslant 2$, the ring $\mathrm{C}^{(\omega)}$ is a finite C-algebra.
  \item $\mathrm{A}$ is universally catenary and for every integral quotient ring $\mathrm{B}$ of $\mathrm{A}$ and every local ring $\mathrm{C}=\mathrm{B}_\mathfrak{q}$ at a prime ideal $\mathfrak{q}$ of $\mathrm{B}$, with $\dim(\mathrm{C})\geqslant 2$, the completed ring $\widehat{\mathrm{C}}$ is such that $\Spec(\widehat{\mathrm{C}})$ has no associated prime cycle of dimension~$1$.
\end{enumerate}
Moreover, when these conditions are satisfied, for every integral quotient ring $\mathrm{B}$ of $\mathrm{A}$ that is strictly equidimensional, the completion $\widehat{\mathrm{B}}$ is strictly equidimensional.
\end{theorem}

\begin{proof}
The equivalence of $c)$, $d)$ and $e)$ was proved in \sref{IV.7.2.4}.
To show that $a)$ and $b)$ are equivalent, recall \sref{IV.7.1.9} that to say that $\mathrm{A}$ is formally catenary means that for every integral quotient ring $\mathrm{B}$ of $\mathrm{A}$, $\widehat{\mathrm{B}}$ is equidimensional.
On the other hand, every fibre of the morphism $\Spec(\widehat{\mathrm{A}})\to\Spec(\mathrm{A})$ at a point $x\in\Spec(\mathrm{A})$ is the fibre of the morphism $\Spec(\widehat{\mathrm{B}})\to\Spec(\mathrm{B})$ at the generic point $x$ of $\Spec(\mathrm{B})$, where $\mathrm{B}=\mathrm{A}/\mathfrak{I}_x$; to say that this fibre has no embedded associated prime cycle is equivalent to saying that $\widehat{\mathrm{B}}$ has no embedded associated prime ideals, by virtue of \sref{IV.3.3.3}; this proves the equivalence of $a)$ and $b)$.
The same reasoning shows that if $a)$ is satisfied, then, for every integral quotient ring $\mathrm{B}$ of $\mathrm{A}$ that has no embedded associated prime ideals, the completion $\widehat{\mathrm{B}}$ is also without embedded associated prime ideals; on the other hand, the hypothesis that $\mathrm{A}$ is formally catenary implies that if $\mathrm{B}$ is equidimensional, it is also the case for $\widehat{\mathrm{B}}$ \sref{IV.7.1.8}; this establishes the last assertion of the theorem.

Let us show now that $a)$ implies $c)$.
Condition~$a)$ implies that $\mathrm{A}$ is universally catenary \sref{IV.7.1.11}; let us show on the other hand that for every integral quotient ring $\mathrm{B}$ of $\mathrm{A}$, $\mathrm{B}^{(1)}$ is then a finite B-algebra.
Taking into account \sref{IV.7.2.3} applied to $\mathrm{B}$, it amounts to showing that if $\mathrm{X}=\Spec(\mathrm{B})$, $\mathrm{T}$ is a closed subset of codimension $\geqslant 2$ in $\mathrm{X}$, $\mathrm{X}'=\Spec(\widehat{\mathrm{B}})$, and $g:\mathrm{X}'\to\mathrm{X}$ is the canonical morphism, then for all $x'\in\operatorname{Ass}(\sh{O}_{\mathrm{X}'})$, $\operatorname{codim}(g^{-1}(\mathrm{T})\cap\overline{\{x'\}},\overline{\{x'\}})\geqslant 2$.
But by hypothesis $\mathrm{X}'$ has no associated prime cycle of dimension $1$, hence $\inf(\operatorname{codim}(g^{-1}(\mathrm{T})\cap\overline{\{x'\}},\overline{\{x'\}}))$ when $x'$ ranges over $\operatorname{Ass}(\sh{O}_{\mathrm{X}'})$, is nothing other than $\operatorname{codim}(g^{-1}(\mathrm{T}),\mathrm{X}')$.
Now, as $g$ is a faithfully flat morphism, we have \sref{IV.6.1.4}
\[
  \operatorname{codim}(g^{-1}(\mathrm{T}),\mathrm{X}')=\operatorname{codim}(\mathrm{T},\mathrm{X})\geqslant 2.
\]

\oldpage[IV-2]{191}

It remains to prove that $c)$ implies $a)$.
Let us reason by induction on $n=\dim(\mathrm{A})$, the theorem being trivial for $n=0$; we can moreover restrict to the case where $\mathrm{A}=\mathrm{B}$ is integral.
We proceed in two steps, $n$ being $\geqslant 1$.

\medskip\noindent
I) \emph{Suppose first that $\mathrm{A}$ satisfies $(\mathrm{S}_2)$.} ---
Set $\mathrm{X}=\Spec(\mathrm{A})$, $\mathrm{A}'=\widehat{\mathrm{A}}$, $\mathrm{X}'=\Spec(\mathrm{A}')$ and let $u:\mathrm{X}'\to\mathrm{X}$ be the canonical morphism; it amounts to showing that for all $x'\in\operatorname{Ass}(\sh{O}_{\mathrm{X}'})$, we have $\dim(\overline{\{x'\}})=n$.
Let $f\neq 0$ be an element of the maximal ideal of $\mathrm{A}$, and set $\mathrm{C}=\mathrm{A}/f\mathrm{A}$; we know \sref{IV.5.7.6} that $\mathrm{A}/f\mathrm{A}$ satisfies $(\mathrm{S}_1)$; moreover, as the prime ideals of $\mathrm{A}$ minimal among those containing $f$ have height $1$ by virtue of the Hauptidealsatz, and as $\mathrm{C}$ is catenary, $\mathrm{C}=\mathrm{A}/f\mathrm{A}$ is \emph{strictly equidimensional} and $\dim(\mathrm{C})=n-1$.
We have $\widehat{\mathrm{C}}=\mathrm{C}\otimes_\mathrm{A}\mathrm{A}'=\mathrm{A}'/f\mathrm{A}'$, and $f$ is $\mathrm{A}'$-regular by flatness \sref[0]{0.6.3.4}; if we set $\mathrm{Y}'=\mathrm{V}(f\mathrm{A}')=\Spec(\widehat{\mathrm{C}})$, then, for every maximal point $y'$ of $\mathrm{Y}'\cap\overline{\{x'\}}$, we have $y'\in\operatorname{Ass}(\sh{O}_{\mathrm{Y}'})$ by virtue of \sref{IV.3.4.3}; on the other hand, we have $y'\neq x'$ since $u(x')$ is the generic point of $\mathrm{X}$ \sref{IV.3.3.2} and $u(y')\in\mathrm{V}(f\mathrm{A})$; we conclude from \sref{IV.5.1.8} that
\[
  \dim(\overline{\{y'\}})=\dim(\overline{\{x'\}})-1.
  \tag{7.2.5.1}
\label{IV.7.2.5.1}
\]

But the quotient ring $\mathrm{C}=\mathrm{A}/f\mathrm{A}$ also satisfies hypothesis $c)$ of the statement \sref{IV.5.6.1} and \sref{IV.5.11.7},~(ii)), hence by the induction hypothesis, we have $\dim(\overline{\{y'\}})=\dim(\mathrm{C})=n-1$; the conclusion therefore follows in this case from {\hyperlink{ega4.c04.eq.7-2-5-1}{(7.2.5.1)}}.

\medskip\noindent
II) \emph{General case.} ---
As by hypothesis $\mathrm{A}^{(1)}$ is a finite A-algebra, it is a semi-local ring and satisfies $(\mathrm{S}_2)$ by virtue of \sref{IV.5.10.17},~(i); the same is true of each of its local rings $({\mathrm{A}^{(1)}})_\mathfrak{n}$ at a maximal ideal $\mathfrak{n}$.
Moreover, as $\mathrm{A}$ is universally catenary, the rings $({\mathrm{A}^{(1)}})_\mathfrak{n}$ (in finite number) are all of \emph{dimension} $n$ \sref{IV.5.6.10}; these rings also satisfy hypothesis $c)$ of the statement \sref{IV.5.6.1} and \sref{IV.5.11.7},~(ii).
We know that the completion of $\mathrm{A}^{(1)}$, equal to $\widehat{\mathrm{A}}\otimes_\mathrm{A}\mathrm{A}^{(1)}$, is the direct product of the completions of the local rings $(\mathrm{A}^{(1)})_\mathfrak{n}$.
Set $\mathrm{X}_1=\Spec(\mathrm{A}^{(1)})$, $\mathrm{X}'_1=\Spec((\mathrm{A}^{(1)})\widehat{\phantom{x}})=\mathrm{X}'\times_\mathrm{X}\mathrm{X}_1$; for all $x'_1\in\operatorname{Ass}(\sh{O}_{\mathrm{X}'_1})$, it follows from what precedes and from case I) that $\dim(\overline{\{x'_1\}})=n$.
Let $u_1=u_{(\mathrm{X}_1)}:\mathrm{X}'_1\to\mathrm{X}_1$ be the canonical morphism; as $\mathrm{A}$ and $\mathrm{A}^{(1)}$ have the same field of fractions, the inverse image of the generic point $x$ of $\mathrm{X}$ by the projection $\mathrm{X}_1\to\mathrm{X}$ reduces to the generic point $x_1$ of $\mathrm{X}_1$; the inverse image by the projection $\mathrm{X}'_1\to\mathrm{X}'$ of the fibre $u^{-1}(x)$ is the fibre $u_1^{-1}(x_1)$, and this projection induces an isomorphism of the prescheme $u_1^{-1}(x_1)$ onto the prescheme $u^{-1}(x)$.
This being the case, the points of $\operatorname{Ass}(\sh{O}_{\mathrm{X}'})$ (resp.~$\operatorname{Ass}(\sh{O}_{\mathrm{X}'_1})$) are the generic points of the associated prime cycles to $u^{-1}(x)$ (resp.~$u_1^{-1}(x_1)$) \sref{IV.3.3.1}.
For every $x'\in\operatorname{Ass}(\sh{O}_{\mathrm{X}'})$, there is therefore an $x'_1\in\operatorname{Ass}(\sh{O}_{\mathrm{X}'_1})$ lying over $x'$; and if $\mathrm{Z}'$ (resp.~$\mathrm{Z}'_1$) is the reduced closed sub-prescheme of $\mathrm{X}'$ (resp.~$\mathrm{X}'_1$) having $\overline{\{x'\}}$ (resp.~$\overline{\{x'_1\}}$) as underlying space, the projection $\mathrm{Z}'_1\to\mathrm{Z}'$ is a finite and surjective morphism; we conclude from \sref{IV.5.4.2} that $\dim(\overline{\{x'\}})=\dim(\overline{\{x'_1\}})=n$.
\end{proof}

\begin{definition}[7.2.6]
\label{IV.7.2.6}
When a local Noetherian ring $\mathrm{A}$ satisfies the equivalent conditions of \sref{IV.7.2.5}, we say that it is \emph{strictly formally catenary}.
\end{definition}

\begin{corollary}[7.2.7]
\label{IV.7.2.7}
Let $\mathrm{A}$ be a local Noetherian ring such that there exists an $\mathrm{A}$-module $\mathrm{M}$ of finite type that is a Cohen-Macaulay $\mathrm{A}$-module and for which $\operatorname{Supp}(\mathrm{M})=\Spec(\mathrm{A})$; then $\mathrm{A}$ is strictly formally catenary.
\end{corollary}

\begin{proof}
\oldpage[IV-2]{192}
Indeed, $\mathrm{A}$ is formally catenary \sref{IV.7.1.5} and \sref{IV.7.1.9}, and on the other hand the fibres of the canonical morphism $\Spec(\widehat{\mathrm{A}})\to\Spec(\mathrm{A})$ are preschemes of Cohen-Macaulay \sref{IV.6.3.8}, hence \emph{a fortiori} satisfy $(\mathrm{S}_1)$.
\end{proof}

\begin{corollary}[7.2.8]
\label{IV.7.2.8}
If $\mathrm{A}$ is a local Noetherian ring that is strictly formally catenary, every local ring $\mathrm{B}$ that is an A-algebra essentially of finite type is strictly formally catenary.
\end{corollary}

\begin{proof}
Indeed, $\mathrm{B}$ is formally catenary \sref{IV.7.1.11} and it follows from \sref{IV.7.4.4}\footnote{The corollary \sref{IV.7.2.8} is not used to prove \sref{IV.7.4.4}.} that the fibres of $\Spec(\widehat{\mathrm{B}})\to\Spec(\mathrm{B})$ satisfy $(\mathrm{S}_1)$, whence the conclusion.
\end{proof}

\begin{corollary}[7.2.9]
\label{IV.7.2.9}
Every local Noetherian ring of dimension $1$ is strictly formally catenary.
\end{corollary}

\begin{proof}
Indeed, every \emph{integral} quotient ring $\mathrm{B}$ of such a ring $\mathrm{A}$ is of dimension $0$ or $1$, hence its completion is strictly equidimensional \sref{IV.7.2.1.1}.
\end{proof}

\begin{remarks}[7.2.10]
\label{IV.7.2.10}
\begin{enumerate}[label=\emph{\roman*)}]
  \item It follows from \sref{IV.7.2.7} and \sref{IV.7.2.8} that every quotient ring of a Cohen-Macaulay local ring is strictly formally catenary.
  A local Noetherian ring of dimension $2$ satisfying $(\mathrm{S}_2)$ is strictly formally catenary, since it is a Cohen-Macaulay ring.
  Recall however that there exist integral local Noetherian rings of dimension~$2$ that are not universally catenary, nor \emph{a fortiori} strictly formally catenary \sref{IV.5.6.11}.

  \item We do not know if a ring that is formally catenary \sref{IV.7.1.9} is strictly formally catenary; this is because we do not know if for a local Noetherian integral ring $\mathrm{B}$, $\widehat{\mathrm{B}}$ has no embedded associated prime ideals \sref{IV.6.4.3}.
  We also do not know if, in the equivalent conditions $c)$, $d)$, $e)$ of \sref{IV.7.2.5}, one can replace the hypothesis that $\mathrm{A}$ is universally catenary by the hypothesis that $\mathrm{A}$ is catenary; one can show that this is the case when $\mathrm{A}$ is Henselian \sref{IV.18.9.6}.
\end{enumerate}
\end{remarks}


\subsection{Formal fibres of Noetherian local rings}
\label{subsection:IV.7.3}

\begin{env}[7.3.1]
\label{IV.7.3.1}
We will consider in this number and the two following properties $\boldsymbol{P}(\mathrm{Z},k)$ of the following form:
\begin{quote}
``$\mathrm{Z}$ is a locally Noetherian prescheme over a field $k$,

and for all $z\in\mathrm{Z}$, we have $\boldsymbol{Q}(\sh{O}_z,k)$''
\end{quote}
where $\boldsymbol{Q}(\mathrm{A},k)$ is a property of a local Noetherian ring $\mathrm{A}$ that is a $k$-algebra; one assumes that if $k'$ is a field isomorphic to $k$, and if $\mathrm{A}'$ is a local Noetherian ring that is a $k'$-algebra di-isomorphic to $\mathrm{A}$, the properties $\boldsymbol{Q}(\mathrm{A},k)$ and $\boldsymbol{Q}(\mathrm{A}',k')$ are equivalent.
\end{env}

\oldpage[IV-2]{193}

If $\mathrm{X}$, $\mathrm{Y}$ are two locally Noetherian preschemes, we will say that a morphism $f:\mathrm{X}\to\mathrm{Y}$ is a \emph{$\boldsymbol{P}$-morphism} if:
\begin{enumerate}[label=\arabic{enumi}\textsuperscript{o}]
  \item $f$ is \emph{flat};
  \item for all $y\in\mathrm{Y}$, the property $\boldsymbol{P}(f^{-1}(y),\boldsymbol{k}(y))$ is true.
\end{enumerate}

\begin{lemma}[7.3.2]
\label{IV.7.3.2}
Let $\mathrm{X}$, $\mathrm{Y}$ be two locally Noetherian preschemes, $f:\mathrm{X}\to\mathrm{Y}$ a morphism.
The following properties are equivalent:
\begin{enumerate}[label=\emph{\alph*)}]
  \item $f$ is a $\boldsymbol{P}$-morphism.
  \item For all $x\in\mathrm{X}$, if we set $y=f(x)$, then $\sh{O}_x$ is a flat $\sh{O}_y$-module and the property $\boldsymbol{Q}(\sh{O}_x\otimes_{\sh{O}_y}\boldsymbol{k}(y),\boldsymbol{k}(y))$ is true.
  \item For all $x\in\mathrm{X}$, if we set $y=f(x)$, the morphism $\Spec(\sh{O}_x)\to\Spec(\sh{O}_y)$ corresponding to the homomorphism $\sh{O}_y\to\sh{O}_x$ is a $\boldsymbol{P}$-morphism.
  \item[\emph{c$'$)}] For every closed point $x\in\mathrm{X}$, the morphism $\Spec(\sh{O}_x)\to\Spec(\sh{O}_y)$ is a $\boldsymbol{P}$-morphism.
\end{enumerate}
\end{lemma}

\begin{proof}
The equivalence of $a)$ and $b)$ follows from the definitions; the same holds for $b)$ and $c)$, taking into account \sref[I]{I.2.4.2}; finally the equivalence of $b)$ and $c'$) follows from the fact that, by \sref{IV.5.1.11}, $\overline{\{x\}}$ contains a closed point.
\end{proof}

\begin{corollary}[7.3.3]
\label{IV.7.3.3}
Let $\mathrm{A}$, $\mathrm{B}$ be two Noetherian rings, $\varphi:\mathrm{A}\to\mathrm{B}$ a homomorphism such that the corresponding morphism $\Spec(\mathrm{B})\to\Spec(\mathrm{A})$ is a $\boldsymbol{P}$-morphism.
Then, for every multiplicative subset $\mathrm{S}$ of $\mathrm{A}$, the morphism $\Spec(\mathrm{S}^{-1}\mathrm{B})\to\Spec(\mathrm{S}^{-1}\mathrm{A})$ is a $\boldsymbol{P}$-morphism.
\end{corollary}

\begin{proof}
This follows immediately from \sref{IV.7.3.2} and \sref[I]{I.1.6.2}.
\end{proof}

\begin{env}[7.3.4]
\label{IV.7.3.4}
We will suppose throughout what follows that the property $\boldsymbol{Q}$ is such that the following three conditions are satisfied:
\begin{enumerate}
  \item[$(\mathrm{P_I})$] \emph{(transitivity).} ---
  If $f:\mathrm{X}\to\mathrm{Y}$ is a regular morphism \sref{IV.6.8.1} and $g:\mathrm{Y}\to\mathrm{Z}$ is a $\boldsymbol{P}$-morphism, then $g\circ f$ is a $\boldsymbol{P}$-morphism.
  \item[$(\mathrm{P_{II}})$] \emph{(descent).} ---
  If $f:\mathrm{X}\to\mathrm{Y}$ and $g:\mathrm{Y}\to\mathrm{Z}$ are morphisms of locally Noetherian preschemes such that $f$ is faithfully flat and that $g\circ f$ is a $\boldsymbol{P}$-morphism, then $g$ is a $\boldsymbol{P}$-morphism.
  \item[$(\mathrm{P_{III}})$] ---
  For every field $k$, the property $\boldsymbol{P}(\Spec(k),k)$ is true.
\end{enumerate}
\end{env}

\begin{remarks}[7.3.5]
\label{IV.7.3.5}
\begin{enumerate}[label=\emph{\roman*)}]
  \item Let us note that the conditions $(\mathrm{P_I})$ and $(\mathrm{P_{III}})$ imply that \emph{every regular morphism is a $\boldsymbol{P}$-morphism}.

  \item Let us note that the hypotheses of $(\mathrm{P_I})$ (resp.~$(\mathrm{P_{II}})$) imply that $h=g\circ f$ (resp.~$g$) is flat \sref{IV.2.2.13}; the hypotheses of $(\mathrm{P_I})$ or of $(\mathrm{P_{II}})$ imply on the other hand that for all $z\in\mathrm{Z}$, $f_z:h^{-1}(z)\to g^{-1}(z)$ is flat \sref{IV.2.1.4}; as, for all $y\in g^{-1}(z)$, $f_z^{-1}(y)$ is isomorphic to $f^{-1}(y)$ \sref[I]{I.3.6.4}, this shows that it suffices, to verify conditions $(\mathrm{P_I})$ and $(\mathrm{P_{II}})$, to do it only when $\mathrm{Z}$ is the \emph{spectrum of a field}.

  \item In certain cases, the property $\boldsymbol{Q}$ will be such that the following condition is also satisfied:
  \begin{enumerate}
    \item[$(\mathrm{P'_I})$] ---
    If $f:\mathrm{X}\to\mathrm{Y}$ and $g:\mathrm{Y}\to\mathrm{Z}$ are two $\boldsymbol{P}$-morphisms, then $g\circ f$ is a $\boldsymbol{P}$-morphism.
  \end{enumerate}
\end{enumerate}
\end{remarks}

\oldpage[IV-2]{194}

\begin{env}[7.3.6]
\label{IV.7.3.6}
We will say that the property $\boldsymbol{P}$ is \emph{geometric} if it satisfies (in addition to $(\mathrm{P_I})$, $(\mathrm{P_{II}})$ and $(\mathrm{P_{III}})$) the condition:
\begin{enumerate}
  \item[$(\mathrm{P_{IV}})$] \emph{(extension of finite type of the base field).} ---
  If $\boldsymbol{P}(\mathrm{Z},k)$ is true, then, for every extension $k'$ of finite type over $k$, $\boldsymbol{P}(\mathrm{Z}\otimes_k k',k')$ is true.
\end{enumerate}
\end{env}

\begin{lemma}[7.3.7]
\label{IV.7.3.7}
Let $f:\mathrm{X}\to\mathrm{Y}$ be a $\boldsymbol{P}$-morphism of locally Noetherian preschemes, $g:\mathrm{Y}'\to\mathrm{Y}$ a morphism locally of finite type.
Then, if $\boldsymbol{P}$ satisfies the condition $(\mathrm{P_{IV}})$, the morphism $f'=f_{(\mathrm{Y}')}:\mathrm{X}\times_\mathrm{Y}\mathrm{Y}'\to\mathrm{Y}'$ is a $\boldsymbol{P}$-morphism.
\end{lemma}

\begin{proof}
Indeed, for all $y'\in\mathrm{Y}'$, if we set $y=g(y')$, $\boldsymbol{k}(y')$ is an extension of finite type of $\boldsymbol{k}(y)$ \sref[I]{I.6.4.11} and $f'^{-1}(y')=f^{-1}(y)\otimes_{\boldsymbol{k}(y)}\boldsymbol{k}(y')$ \sref[I]{I.3.6.4}; it suffices therefore to apply $(\mathrm{P_{IV}})$.
\end{proof}

\begin{env}[7.3.8]
\label{IV.7.3.8}
\emph{Examples.} ---
The following properties $\boldsymbol{P}(\mathrm{Z},k)$ satisfy the conditions $(\mathrm{P_I})$, $(\mathrm{P_{II}})$ and $(\mathrm{P_{III}})$:
\begin{enumerate}[label=(\roman*)]
  \item (also denoted $(\mathrm{i}_n)$) $\mathrm{Z}$ is of coprofundity $\leqslant n$.
  \item $\mathrm{Z}$ is a prescheme of Cohen-Macaulay.
  \item (also denoted $(\mathrm{iii}_n)$) $\mathrm{Z}$ satisfies $(\mathrm{S}_n)$.
  \item $\mathrm{Z}$ is regular.
  \item (also denoted $(\mathrm{v}_n)$) $\mathrm{Z}$ satisfies $(\mathrm{R}_n)$.
  \item $\mathrm{Z}$ is reduced.
  \item $\mathrm{Z}$ is normal.
\end{enumerate}
For properties~(ii) to~(vii), this follows from \sref{IV.6.6.1}, which in fact proves the stronger condition $(\mathrm{P}'_\mathrm{I})$.
For~(i), the property $(\mathrm{P_{II}})$ follows from \sref{IV.6.6.2}; the property $(\mathrm{P_I})$ follows, by the same reasoning as in \sref{IV.6.6.1},~(i), from \sref{IV.6.3.2} and from the fact that a regular prescheme has codepth~$0$.

Moreover, it follows from \sref{IV.6.7.1} that properties~(i), (ii), and~(iii) are geometric; by virtue of \sref{IV.6.7.8}, the same is true of the following:
\begin{enumerate}[label=(\roman*')]
  \item[(iv')] $\mathrm{Z}$ is geometrically regular.
  \item[(v')] (also denoted $(\mathrm{v}'_n)$) $\mathrm{Z}$ possesses the property $(\mathrm{R}_n)$ geometric.
  \item[(vi')] $\mathrm{Z}$ is geometrically reduced.
  \item[(vii')] $\mathrm{Z}$ is geometrically normal.
\end{enumerate}
\end{env}

\begin{remarks}[7.3.9]
\label{IV.7.3.9}
\begin{enumerate}[label=\emph{\roman*)}]
  \item Each of the properties (iv'), $(\mathrm{v}'_n)$, (vi'), (vii') implies respectively the corresponding property (iv), $(\mathrm{v}_n)$, (vi), (vii).
  The property~(iv) implies all the properties (i) to (vii), and the property~(iv') implies all the properties enumerated in \sref{IV.7.3.8}.
  Recall also that $(\mathrm{i}_0)$ is equivalent to~(ii); the conjunction of $(\mathrm{iii}_1)$ and of $(\mathrm{v}_0)$ (resp.~of $(\mathrm{iii}_1)$ and of $(\mathrm{v}'_0)$) is equivalent to~(vi) (resp.~(vi')) \sref{IV.5.8.5}; the conjunction of $(\mathrm{iii}_2)$ and of $(\mathrm{v}_1)$ (resp.~of $(\mathrm{iii}_2)$ and of $(\mathrm{v}'_1)$) is equivalent to~(vii) (resp.~(vii')) \sref{IV.5.8.6}.

  \item In all the examples of \sref{IV.7.3.8}, the property $\boldsymbol{Q}(\sh{O}_z,k)$ which serves to define $\boldsymbol{P}$ is such that, for every \emph{generization} $z'$ of $z$ in $\mathrm{Z}$, $\boldsymbol{Q}(\sh{O}_z,k)$ \emph{implies} $\boldsymbol{Q}(\sh{O}_{z'},k)$.
  By virtue of \sref{IV.2.3.4}, it suffices to verify this for properties~(i) to~(vii) (and even~(i) to~(v), by \sref{IV.7.3.9},~(i)): it is part of the definition for~(iii) and~(v) \sref{IV.5.7.2} and \sref{IV.5.8.2}, it follows from \sref{IV.6.3.5},~(i), for~(i),\footnote{The printed text cites (6.3.9), but {\hyperlink{ega4.c04.subsec.6-3}{\S6.3}} ends at Proposition~{\hyperlink{ega4.c04.prop.6-3-8}{(6.3.8)}}; the codepth-under-generization statement used here is Corollary~{\hyperlink{ega4.c04.cor.6-3-5}{(6.3.5)}}(i).} and finally from \sref[0]{0.16.5.10} and \sref{IV.17.3.2} for~(ii) and~(iv).
\end{enumerate}
\end{remarks}

\oldpage[IV-2]{195}

\begin{env}[7.3.10]
\label{IV.7.3.10}
We will say that a property $\boldsymbol{P}(\mathrm{Z},k)$ is of the \emph{first type} if the property $\boldsymbol{Q}(\mathrm{A},k)$ which serves to define it defines a property $\boldsymbol{R}(\mathrm{A})$ \emph{not involving $k$} and is true when $\mathrm{A}$ is a \emph{field}; we will say that $\boldsymbol{P}(\mathrm{Z},k)$ is of the \emph{second type} if the property $\boldsymbol{Q}(\mathrm{A},k)$ is of the form
\begin{quote}
``For every extension $k'$ of finite type over $k$ and every point $z'$ of $\Spec(\mathrm{A}\otimes_k k')$ above the closed point of $\Spec(\mathrm{A})$, we have $\boldsymbol{R}(\sh{O}_{z'})$''
\end{quote}
$\boldsymbol{R}(\mathrm{A})$ being assumed true when $\mathrm{A}$ is a field.
It is clear that in the examples of \sref{IV.7.3.8}, the properties~(i) to~(vii) are of the first type, the properties~(iv') to~(vii') of the second type.

This being the case, if we take up the reasoning of \sref{IV.6.6.1} and \sref{IV.6.8.3}, we observe that when $\boldsymbol{P}$ is of the first or second type for a property $\boldsymbol{R}(\mathrm{A})$, \emph{the conditions $(\mathrm{P_I})$ and $(\mathrm{P_{II}})$ have the following consequences on $\boldsymbol{R}$}:

$(\mathrm{R_I})$ Let $\mathrm{X}$, $\mathrm{Y}$ be two locally Noetherian preschemes, $f:\mathrm{X}\to\mathrm{Y}$ a \emph{regular} morphism; then, for all $x\in\mathrm{X}$, the property $\boldsymbol{R}(\sh{O}_{f(x)})$ implies the property $\boldsymbol{R}(\sh{O}_x)$.

$(\mathrm{R_{II}})$ Let $\varphi:\mathrm{A}\to\mathrm{B}$ be a homomorphism of local Noetherian rings making $\mathrm{B}$ a flat $\mathrm{A}$-module; then $\boldsymbol{R}(\mathrm{B})$ implies $\boldsymbol{R}(\mathrm{A})$.

Moreover, the property $(\mathrm{P_{III}})$ results from the fact that $\boldsymbol{R}(k)$ is true for every field $k$; finally, when $\boldsymbol{P}$ is of the second type, the condition $(\mathrm{P_{IV}})$ is a consequence of the definition of $\boldsymbol{P}$ and of the transitivity of extensions of finite type.
\end{env}

\begin{remark}[7.3.11]
\label{IV.7.3.11}
We leave to the reader the task of formulating the property $\boldsymbol{R}$ corresponding to each of the examples of \sref{IV.7.3.8}, and of verifying the conditions $(\mathrm{R_I})$ and $(\mathrm{R_{II}})$ in each case, with the help of the results of \textsection\,6.
In fact, except for example~(i) of \sref{IV.7.3.8}, the property $\boldsymbol{R}$ satisfies in the other cases the following condition:
\begin{enumerate}
  \item[$(\mathrm{R}'_\mathrm{I})$] Let $\mathrm{X}$, $\mathrm{Y}$ be two locally Noetherian preschemes, $f:\mathrm{X}\to\mathrm{Y}$ a $\boldsymbol{P}$-morphism (for the property $\boldsymbol{P}$ of the first or second type defined by $\boldsymbol{R}$); then, for all $x\in\mathrm{X}$, the property $\boldsymbol{R}(\sh{O}_{f(x)})$ implies the property $\boldsymbol{R}(\sh{O}_x)$.
\end{enumerate}
The reasoning of \sref{IV.6.6.1} and \sref{IV.6.8.3} proves that if $\boldsymbol{R}$ satisfies the conditions $(\mathrm{R}'_\mathrm{I})$ and $(\mathrm{R_{II}})$, then $\boldsymbol{P}$ satisfies the conditions $(\mathrm{P'_I})$ \sref{IV.7.3.5},~(iii), and $(\mathrm{P_{II}})$.
\end{remark}

\begin{proposition}[7.3.12]
\label{IV.7.3.12}
Let $\boldsymbol{P}$ be a property of the first or the second type defined from a property $\boldsymbol{R}$ verifying the conditions $(\mathrm{R_I})$ and $(\mathrm{R_{II}})$ (resp.~$(\mathrm{R}'_\mathrm{I})$ and $(\mathrm{R_{II}})$).
If for every locally Noetherian prescheme $\mathrm{X}$, $\mathrm{U}_{\boldsymbol{R}}(\mathrm{X})$ denotes the set of $x\in\mathrm{X}$ such that $\boldsymbol{R}(\sh{O}_x)$ is true, then, for every regular morphism (resp.~every $\boldsymbol{P}$-morphism) $f:\mathrm{X}\to\mathrm{Y}$ of locally Noetherian preschemes, we have
\[
  \mathrm{U}_{\boldsymbol{R}}(\mathrm{X})=f^{-1}(\mathrm{U}_{\boldsymbol{R}}(\mathrm{Y})).
  \tag{7.3.12.1}
\label{IV.7.3.12.1}
\]
\end{proposition}

\begin{proof}
This is an immediate consequence of the definitions.
\end{proof}

\begin{env}[7.3.13]
\label{IV.7.3.13}
Given a semi-local Noetherian ring $\mathrm{A}$, we will call \emph{formal fibres} of $\mathrm{A}$ the fibres of the canonical morphism $f:\Spec(\widehat{\mathrm{A}})\to\Spec(\mathrm{A})$; for every prime ideal $\mathfrak{p}$ of $\mathrm{A}$, the formal fibre at $\mathfrak{p}$ is therefore the scheme $\Spec(\widehat{\mathrm{A}}\otimes_\mathrm{A}\boldsymbol{k}(\mathfrak{p}))$; as the completion of the semi-local ring $\mathrm{A}/\mathfrak{p}$ is $\widehat{\mathrm{A}}/\mathfrak{p}\widehat{\mathrm{A}}=(\mathrm{A}/\mathfrak{p})\otimes_\mathrm{A}\widehat{\mathrm{A}}$, we see that the formal fibre of $\mathrm{A}$ at $\mathfrak{p}$ is also the \emph{formal fibre of $\mathrm{A}/\mathfrak{p}$ at the generic point} (o) of $\Spec(\mathrm{A}/\mathfrak{p})$.
\end{env}

\oldpage[IV-2]{196}

The property $\boldsymbol{P}$ being defined as in \sref{IV.7.3.1}, we will say that \emph{the formal fibres of $\mathrm{A}$ have the property $\boldsymbol{P}$}, or that $\mathrm{A}$ is a \emph{$\boldsymbol{P}$-ring}, if, for all $x\in\Spec(\mathrm{A})$, $\boldsymbol{P}(f^{-1}(x),\boldsymbol{k}(x))$ is true.
As the morphism $f$ is \emph{flat}, this amounts to saying that $f$ is a $\boldsymbol{P}$-morphism.

\begin{proposition}[7.3.14]
\label{IV.7.3.14}
Let $\mathrm{A}$ be a semi-local Noetherian ring, $\mathfrak{m}_i$ ($1\leqslant i\leqslant r$) its maximal ideals, and set $\mathrm{A}_i=\mathrm{A}_{\mathfrak{m}_i}$; for $\mathrm{A}$ to be a $\boldsymbol{P}$-ring, it is necessary and sufficient that each of the $\mathrm{A}_i$ be so.
\end{proposition}

\begin{proof}
Indeed, $\widehat{\mathrm{A}}$ is the direct product of the $\widehat{\mathrm{A}}_i$, hence the formal fibre of $\mathrm{A}$ at a point $x\in\Spec(\mathrm{A})$ is the \emph{sum} of the formal fibres at $x$ of those $\mathrm{A}_i$ such that $x\in\Spec(\mathrm{A}_i)$.
\end{proof}

\begin{proposition}[7.3.15]
\label{IV.7.3.15}
Let $\mathrm{A}$ be a semi-local Noetherian ring.
\begin{enumerate}[label=\emph{\roman*)}]
  \item If $\mathrm{A}$ is a $\boldsymbol{P}$-ring, every quotient ring of $\mathrm{A}$ is a $\boldsymbol{P}$-ring.
  \item If in addition $\boldsymbol{P}$ satisfies the condition $(\mathrm{P_{IV}})$ \sref{IV.7.3.6}, then every finite A-algebra is a $\boldsymbol{P}$-ring.
\end{enumerate}
\end{proposition}

\begin{proof}
For every ideal $\mathfrak{a}$ of $\mathrm{A}$, $\widehat{\mathrm{A}}\otimes_\mathrm{A}(\mathrm{A}/\mathfrak{a})$ is the completion of $\mathrm{A}/\mathfrak{a}$, and the formal fibres of $\mathrm{A}/\mathfrak{a}$ are therefore the formal fibres of $\mathrm{A}$ at the points of $\mathrm{V}(\mathfrak{a})$, whence~(i).
On the other hand, if $\mathrm{B}$ is a finite $\mathrm{A}$-algebra, $\mathrm{B}$ is a semi-local ring, we have $\widehat{\mathrm{B}}=\mathrm{B}\otimes_\mathrm{A}\widehat{\mathrm{A}}$, and~(ii) follows from \sref{IV.7.3.7}.
\end{proof}

\begin{proposition}[7.3.16]
\label{IV.7.3.16}
Suppose that the property $\boldsymbol{P}$ is of the first (resp.~second) type, defined from a property $\boldsymbol{R}$ \sref{IV.7.3.10}.
When $\boldsymbol{P}$ is of the second type, suppose moreover that the relation
\begin{quote}
``for every finite extension $k'$ of $k$, and every $z'\in\mathrm{Z}\otimes_k k'$, $\boldsymbol{R}(\sh{O}_{z'})$ is true''
\end{quote}
implies $\boldsymbol{P}(\mathrm{Z},k)$ (this is verified for the examples~(iv') to~(vii') of \sref{IV.7.3.8}, by virtue of \sref{IV.6.7.7} and \sref{IV.4.6.1}).

Let $\mathrm{A}$ be a semi-local Noetherian ring.
If the property $\boldsymbol{R}$ satisfies the conditions $(\mathrm{R_I})$ and $(\mathrm{R_{II}})$, the following properties are equivalent:
\begin{enumerate}[label=\emph{\alph*)}]
  \item $\mathrm{A}$ is a $\boldsymbol{P}$-ring.
  \item For every integral quotient ring $\mathrm{B}$ of $\mathrm{A}$ (resp.~every finite integral A-algebra $\mathrm{B}$) and every prime ideal $\mathfrak{q}$ of $\widehat{\mathrm{B}}$ whose inverse image in $\mathrm{B}$ is $0$, $\boldsymbol{R}(\widehat{\mathrm{B}}_\mathfrak{q})$ is true.
\end{enumerate}
If in addition $\boldsymbol{R}$ satisfies the condition $(\mathrm{R}'_\mathrm{I})$, the properties $a)$ and $b)$ are also equivalent to:
\begin{enumerate}[label=\emph{\alph*)}]
  \setcounter{enumi}{2}
  \item For every quotient ring $\mathrm{B}$ of $\mathrm{A}$ (resp.~every finite A-algebra $\mathrm{B}$), if we set $\mathrm{Y}=\Spec(\mathrm{B})$, $\mathrm{Y}'=\Spec(\widehat{\mathrm{B}})$, and if $g:\mathrm{Y}'\to\mathrm{Y}$ is the canonical morphism, we have (with the notations of \sref{IV.7.3.12})
  \[
    \mathrm{U}_{\boldsymbol{R}}(\mathrm{Y}')=g^{-1}(\mathrm{U}_{\boldsymbol{R}}(\mathrm{Y})).
    \label{IV.7.3.16.1}
\tag{7.3.16.1}
  \]
\end{enumerate}
\end{proposition}

\begin{proof}
The equivalence of $a)$ and $b)$ is immediate when $\boldsymbol{P}$ is of the first type: indeed, for every prime ideal $\mathfrak{p}$ of $\mathrm{A}$, the formal fibre of $\mathrm{B}=\mathrm{A}/\mathfrak{p}$ at the generic point $\mathfrak{p}\in\Spec(\mathrm{A})$ is nothing other than the formal fibre of $\mathrm{A}$ at the point $\mathfrak{p}$ \sref{IV.7.3.13}.
When $\boldsymbol{P}$ is of the second type, the equivalence of $a)$ and $b)$ follows from the more general lemma below.
\end{proof}

\begin{lemma}[7.3.16.2]
\label{IV.7.3.16.2}
Let $\mathrm{A}$, $\mathrm{A}'$ be two rings, $\varphi:\mathrm{A}\to\mathrm{A}'$ a homomorphism, $f:\Spec(\mathrm{A}')\to\Spec(\mathrm{A})$ the corresponding morphism.
In order that for all $x\in\Spec(\mathrm{A})$, every

\oldpage[IV-2]{197}

finite extension $k$ of $\boldsymbol{k}(x)$ and every point $z$ of $f^{-1}(x)\otimes_{\boldsymbol{k}(x)} k=\mathrm{Z}$, the property $\boldsymbol{R}(\sh{O}_{\mathrm{Z},z})$ be true, it is necessary and sufficient that: for every integral finite A-algebra $\mathrm{B}$, if $g:\Spec(\mathrm{A}'\otimes_\mathrm{A}\mathrm{B})\to\Spec(\mathrm{B})$ is the morphism deduced from $f$ by base extension to $\mathrm{B}$, and if $\mathrm{T}=g^{-1}(\xi)$ is the fibre of the generic point $\xi$ of $\Spec(\mathrm{B})$, then, for all $t\in\mathrm{T}$, $\boldsymbol{R}(\sh{O}_{\mathrm{T},t})$ is true.
\end{lemma}

\begin{proof}
The condition is trivially necessary since $x$ is the point of $\Spec(\mathrm{A})$ above which $\xi$ lies, $\boldsymbol{k}(\xi)$ is a finite extension of $\boldsymbol{k}(x)$.
Conversely, let us consider a point $x\in\Spec(\mathrm{A})$ and let $k$ be a finite extension of $\boldsymbol{k}(x)$.
If we set $\mathfrak{p}=\mathfrak{j}_x$, then $\boldsymbol{k}(x)$ is the field of fractions of $\mathrm{A}/\mathfrak{p}$, and there is a basis of $k$ over $\boldsymbol{k}(x)$ consisting of elements integral over $\mathrm{A}/\mathfrak{p}$.
If $\mathrm{B}$ is the subring of $k$ generated by these elements, then $\mathrm{B}$ is a finite integral $\mathrm{A}$-algebra and $k$ is the field $\boldsymbol{k}(\xi)$ at the generic point $\xi$ of $\Spec(\mathrm{B})$.
Since $x$ is the image of $\xi$ in $\Spec(\mathrm{A})$, the fibre $g^{-1}(\xi)$ is none other than $f^{-1}(x)\otimes_{\boldsymbol{k}(x)}k$, which completes the proof of the lemma.
\end{proof}

\begin{proof}[End of the proof of \sref{IV.7.3.16}]
That $a)$ implies $c)$ is an immediate consequence of \sref{IV.7.3.15} and \sref{IV.7.3.12}.
On the other hand, specialise $c)$ to the case where $\mathrm{B}$ is integral, and let $y$ denote the generic point of $\mathrm{Y}=\Spec(\mathrm{B})$.
Then $\sh{O}_{\mathrm{Y},y}=\boldsymbol{k}(y)$, so $y\in\mathrm{U}_{\boldsymbol{R}}(\mathrm{Y})$, since $\boldsymbol{R}(k)$ is true for every field $k$.
Expressing that every point $y'\in g^{-1}(y)$ belongs to $\mathrm{U}_{\boldsymbol{R}}(\mathrm{Y}')$ gives statement $b)$; thus $c)$ implies $b)$.
\end{proof}

\begin{proposition}[7.3.17]
\label{IV.7.3.17}
Suppose that the property $\boldsymbol{R}$ satisfies the conditions $(\mathrm{R}'_\mathrm{I})$ and $(\mathrm{R_{II}})$, and that $\boldsymbol{P}$ is the property of the first (resp.~second) type defined from $\boldsymbol{R}$ \sref{IV.7.3.10}.
Then, if $\mathrm{A}$ is a $\boldsymbol{P}$-ring, the properties $\boldsymbol{R}(\mathrm{A})$ and $\boldsymbol{R}(\widehat{\mathrm{A}})$ are equivalent.
\end{proposition}

\begin{proof}
This follows from \sref{IV.7.3.12.1} applied to $\mathrm{Y}=\Spec(\mathrm{A})$ and $\mathrm{X}=\Spec(\widehat{\mathrm{A}})$.
\end{proof}

\begin{proposition}[7.3.18]
\label{IV.7.3.18}
Suppose that $\boldsymbol{P}$ is a property of the first (resp.~second) type defined from a property $\boldsymbol{R}$ satisfying the conditions $(\mathrm{R}'_\mathrm{I})$ and $(\mathrm{R_{II}})$, and also the condition following:

$(\mathrm{R_{III}})$ For every local Noetherian complete ring $\mathrm{C}$, if we set $\mathrm{Z}=\Spec(\mathrm{C})$, the set $\mathrm{U}_{\boldsymbol{R}}(\mathrm{Z})$ \sref{IV.7.3.12} is open in $\mathrm{Z}$.

Then, if $\mathrm{A}$ is a $\boldsymbol{P}$-ring and if we set $\mathrm{X}=\Spec(\mathrm{A})$, the set $\mathrm{U}_{\boldsymbol{R}}(\mathrm{X})$ is open in $\mathrm{X}$.
\end{proposition}

\begin{proof}
Indeed, if $\mathrm{X}'=\Spec(\widehat{\mathrm{A}})$ and if $f:\mathrm{X}'\to\mathrm{X}$ is the canonical morphism, $\mathrm{U}_{\boldsymbol{R}}(\mathrm{X}')=f^{-1}(\mathrm{U}_{\boldsymbol{R}}(\mathrm{X}))$ \sref{IV.7.3.12}.
As $f$ is faithfully flat and quasi-compact \sref{IV.2.3.12}, and as by hypothesis $\mathrm{U}_{\boldsymbol{R}}(\mathrm{X}')$ is open in $\mathrm{X}'$, the conclusion follows from \sref{IV.2.3.12}.
\end{proof}

\begin{remarks}[7.3.19]
\label{IV.7.3.19}
\begin{enumerate}[label=\emph{\roman*)}]
  \item The property $\boldsymbol{R}$ which serves to define $\boldsymbol{P}$ satisfies the condition $(\mathrm{R_{III}})$ in all the examples enumerated in \sref{IV.7.3.8}.
  For~(i), (ii) and~(iii), this follows from \sref{IV.6.11.2} and the theorem of Cohen \sref[0]{0.19.8.8}; when $\boldsymbol{P}$ is one of the properties~(iv), (iv'), (v), (v'), this follows from \sref{IV.6.12.7} and \sref{IV.6.12.9}; when $\boldsymbol{P}$ is one of the properties~(vii), (vii'), it follows from \sref{IV.6.12.7} and \sref{IV.6.13.4}; finally, for~(vi) and~(vi'), the assertion of \sref{IV.7.3.18} is trivial, being true for every locally Noetherian prescheme.

  \oldpage[IV-2]{198}

  \item We have already noted \sref{IV.6.4.3} that when $\boldsymbol{P}$ is one of the properties~(ii) or $(\mathrm{iii}_1)$ of \sref{IV.7.3.8}, we do not know whether every local Noetherian ring is a $\boldsymbol{P}$-ring.
  Recall, however, that when $\mathrm{A}$ is a quotient of a Cohen--Macaulay ring, the formal fibres of $\mathrm{A}$ are Cohen--Macaulay schemes \sref{IV.6.3.8}.

  \item The most interesting instance of the notion of $\boldsymbol{P}$-ring is that corresponding to the strongest property~(iv') of \sref{IV.7.3.8}: the rings whose formal fibres are geometrically regular.
  Fields, and more generally complete Noetherian local rings, trivially have this property.

  \item Let $\mathrm{A}$ be a local Noetherian ring of dimension~$1$.
  Then $\Spec(\mathrm{A})$ consists of the closed point $a$, corresponding to the maximal ideal $\mathfrak{m}$, and the maximal points $b_i$ ($1\leqslant i\leqslant r$), corresponding to the minimal prime ideals of $\mathrm{A}$.
  One has $\dim(\widehat{\mathrm{A}})=1$, and the maximal ideal of $\widehat{\mathrm{A}}$ is $\mathfrak{m}\widehat{\mathrm{A}}$.
  The formal fibre of $\mathrm{A}$ at $a$ is therefore $\Spec(k)$, where $k=\mathrm{A}/\mathfrak{m}$ is the residue field of $\mathrm{A}$; the formal fibre at each $b_i$ is the spectrum of an Artinian ring whose residue fields are the residue fields $L_{ij}$ of $\Spec(\widehat{\mathrm{A}})$ at the maximal points $b_{ij}$ lying over $b_i$.
  Since an Artinian ring is Cohen--Macaulay, $\mathrm{A}$ is a $\boldsymbol{P}$-ring when $\boldsymbol{P}$ is property~(ii) of \sref{IV.7.3.8}.
  Moreover, since a reduced Artinian ring is a direct sum of fields, the following properties are equivalent \sref{IV.6.7.7}:
  \begin{enumerate}[label=\emph{\alph*)}]
    \item the formal fibres of $\mathrm{A}$ are geometrically reduced;
    \item the formal fibres of $\mathrm{A}$ are geometrically normal;
    \item the formal fibres of $\mathrm{A}$ are geometrically regular.
  \end{enumerate}
  When $\mathrm{A}$ is reduced, they are also equivalent to:
  \begin{enumerate}[label=\emph{\alph*)}]
    \setcounter{enumii}{3}
    \item $\widehat{\mathrm{A}}$ is reduced and $L_{ij}$ is a separable extension of $K_i$ for every pair $(i,j)$ \sref{IV.4.6.1}.
  \end{enumerate}
  In particular, if $\mathrm{A}$ is a discrete valuation ring, $K$ its field of fractions, and $\widehat{K}$ the completion of $K$ for the valuation corresponding to $\mathrm{A}$ (the field of fractions of $\widehat{\mathrm{A}}$), then the formal fibres of $\mathrm{A}$ are geometrically regular if and only if $\widehat{K}$ is a separable extension of $K$.
  This is always so when $K$ has characteristic~$0$.
\end{enumerate}
\end{remarks}


\newpage % Keep this contents destination with the following page-top heading.
\subsection{Permanence of properties of formal fibres}
\label{subsection:IV.7.4}

\begin{env}[7.4.1]
\label{IV.7.4.1}
Throughout the rest of this number, we suppose that the property $\boldsymbol{P}$ is of the form defined in \sref{IV.7.3.1}, and \emph{satisfies the conditions $(\mathrm{P_I})$, $(\mathrm{P_{II}})$ and $(\mathrm{P_{III}})$} of \sref{IV.7.3.4}.
Moreover, we suppose that the property $\boldsymbol{Q}$ which serves to define $\boldsymbol{P}$ is such that, for every generization $z'$ of $z$ in $\mathrm{Z}$, $\boldsymbol{Q}(\sh{O}_z,k)$ implies $\boldsymbol{Q}(\sh{O}_{z'},k)$.
\end{env}

\begin{lemma}[7.4.2]
\label{IV.7.4.2}
Let $\mathrm{A}$, $\mathrm{A}'$ be two local Noetherian rings, $\varphi:\mathrm{A}\to\mathrm{A}'$ a local homomorphism such that $f={}^a\!\varphi:\Spec(\mathrm{A}')\to\Spec(\mathrm{A})$ is a $\boldsymbol{P}$-morphism.
Then, if the formal fibres of $\mathrm{A}'$ are geometrically regular, $\mathrm{A}$ is a $\boldsymbol{P}$-ring.
\end{lemma}

\oldpage[IV-2]{199}

\begin{proof}
Consider the completed homomorphism $\widehat{\varphi}:\widehat{\mathrm{A}}\to\widehat{\mathrm{A}'}$ and the corresponding morphism $\widehat{f}={}^a\!\widehat{\varphi}$; we have the commutative diagram
\[
\begin{tikzpicture}[>=latex]
\node (tl) at (0,1.5) {$\Spec(\widehat{\mathrm{A}})$};
\node (tr) at (4,1.5) {$\Spec(\widehat{\mathrm{A}'})$};
\node (bl) at (0,0) {$\Spec(\mathrm{A})$};
\node (br) at (4,0) {$\Spec(\mathrm{A}')$};
\draw[->] (tr) -- node[above] {$\widehat{f}$} (tl);
\draw[->] (tl) -- node[left] {$g$} (bl);
\draw[->] (tr) -- node[right] {$g'$} (br);
\draw[->] (br) -- node[below] {$f$} (bl);
\end{tikzpicture}
\]
where $g$ and $g'$ are the canonical morphisms.
As by hypothesis $f$ is a $\boldsymbol{P}$-morphism and $g'$ a regular morphism, it follows from $(\mathrm{P_I})$ that $f\circ g'=g\circ\widehat{f}$ is a $\boldsymbol{P}$-morphism.
On the other hand, the hypothesis that $f$ is a $\boldsymbol{P}$-morphism implies that $f$ is flat, hence it is the same for $\widehat{f}$ (Bourbaki, \emph{Alg.\ comm.}, chap.~III, \textsection\,5, n\textsuperscript{o}\,4, cor.~of prop.~3), which is moreover a local homomorphism, hence faithfully flat \sref[0]{0.6.6.2}; it follows then from $(\mathrm{P_{II}})$ that $g$ is a $\boldsymbol{P}$-morphism.
\end{proof}

\begin{corollary}[7.4.3]
\label{IV.7.4.3}
\begin{enumerate}[label=\emph{\roman*)}]
  \item Let $\mathrm{A}$ be a $\boldsymbol{P}$-ring local Noetherian, $\mathrm{A}'=\widehat{\mathrm{A}}$ its completion.
  Suppose that for every prime ideal $\mathfrak{p}'$ of $\mathrm{A}'$, the formal fibres of $\mathrm{A}'_{\mathfrak{p}'}$ are geometrically regular; then, for every prime ideal $\mathfrak{p}$ of $\mathrm{A}$, $\mathrm{A}_\mathfrak{p}$ is a $\boldsymbol{P}$-ring.

  \item Suppose that $\boldsymbol{P}$ satisfies the condition $(\mathrm{P_{IV}})$ \sref{IV.7.3.6}.
  Let $\mathrm{A}$ be a $\boldsymbol{P}$-ring local Noetherian, $\mathrm{B}$ a local A-algebra essentially of finite type, and such that the homomorphism $\mathrm{A}\to\mathrm{B}$ is local.
  Set $\mathrm{A}'=\widehat{\mathrm{A}}$, and let $\mathfrak{n}'$ be the unique prime ideal of $\mathrm{B}'=\mathrm{B}\otimes_\mathrm{A}\mathrm{A}'$ above the maximal ideal of $\mathrm{B}$ and the maximal ideal of $\mathrm{A}'$.
  If the formal fibres of $\mathrm{B}'_{\mathfrak{n}'}$ are geometrically regular, then $\mathrm{B}$ is a $\boldsymbol{P}$-ring.
\end{enumerate}
\end{corollary}

\begin{proof}
\begin{enumerate}[label=(\roman*)]
  \item As by hypothesis $\Spec(\mathrm{A}')\to\Spec(\mathrm{A})$ is a $\boldsymbol{P}$-morphism, the same holds for $\Spec(\mathrm{A}'_{\mathfrak{p}'})\to\Spec(\mathrm{A}_\mathfrak{p})$ by virtue of \sref{IV.7.3.2},~$b)$), for every prime ideal $\mathfrak{p}'$ of $\mathrm{A}'$ above $\mathfrak{p}$ and every prime ideal $\mathfrak{p}$ of $\mathrm{A}$.
  It suffices then to apply the lemma \sref{IV.7.4.2} to this morphism, noting that the morphism $\Spec(\mathrm{A}')\to\Spec(\mathrm{A})$ is surjective.

  \item By virtue of \sref{IV.7.4.2}, it suffices to prove that the morphism $\Spec(\mathrm{B}'_{\mathfrak{n}'})\to\Spec(\mathrm{B})$ is a $\boldsymbol{P}$-morphism.
  Now we have $\mathrm{B}=\mathrm{C}_\mathfrak{n}$, where $\mathrm{C}$ is a finite type A-algebra and $\mathfrak{n}$ is a prime ideal of $\mathrm{C}$ above the maximal ideal of $\mathrm{A}$.
  If we set $\mathrm{C}'=\mathrm{C}\otimes_\mathrm{A}\mathrm{A}'$, it follows from the hypotheses and from \sref{IV.7.3.7} that $\Spec(\mathrm{C}')\to\Spec(\mathrm{C})$ is a $\boldsymbol{P}$-morphism; as $\mathrm{B}'_{\mathfrak{n}'}$ is a local ring of $\mathrm{C}'$ at a prime ideal of $\mathrm{C}'$ above $\mathfrak{n}$, the conclusion follows from \sref{IV.7.3.2},~$b)$).
\end{enumerate}
\end{proof}

\begin{theorem}[7.4.4]
\label{IV.7.4.4}
The hypotheses on $\boldsymbol{P}$ being those of \sref{IV.7.4.1}, let $\mathrm{A}$ be a $\boldsymbol{P}$-ring local Noetherian.
\begin{enumerate}[label=\emph{\roman*)}]
  \item For every prime ideal $\mathfrak{p}$ of $\mathrm{A}$, $\mathrm{A}_\mathfrak{p}$ is a $\boldsymbol{P}$-ring.
  \item Suppose in addition that $\boldsymbol{P}$ satisfies the condition $(\mathrm{P_{IV}})$.
  Then every local ring that is an A-algebra essentially of finite type is a $\boldsymbol{P}$-ring.
\end{enumerate}
\end{theorem}

\begin{proof}
\begin{enumerate}[label=(\roman*)]
  \item Applying \sref{IV.7.4.3},~(i), everything reduces to seeing that for every prime ideal $\mathfrak{p}'$ of $\mathrm{A}'=\widehat{\mathrm{A}}$, $\mathrm{A}'_{\mathfrak{p}'}$ has its formal fibres geometrically regular.
  Now, this has been proved in \sref[0]{0.22.3.3} and \sref[0]{0.22.5.8}.

  \oldpage[IV-2]{200}
  \item Let $\mathrm{B}$ be an A-algebra essentially of finite type that is a local ring; if $\mathfrak{p}$ is the prime ideal of $\mathrm{A}$ above which the maximal ideal of $\mathrm{B}$ lies, $\mathrm{B}$ is also an $\mathrm{A}_\mathfrak{p}$-algebra essentially of finite type \sref{IV.1.3.8}; by virtue of~(i), we may therefore suppose that $\mathfrak{p}$ is the maximal ideal $\mathfrak{m}$ of $\mathrm{A}$.
  We have then $\mathrm{B}=\mathrm{C}_\mathfrak{q}$, where $\mathrm{C}$ is a finite type A-algebra and $\mathfrak{q}$ is a prime ideal of $\mathrm{C}$ above $\mathfrak{m}$. Choose a maximal ideal $\mathfrak{n}\supset\mathfrak{q}$ of $\mathrm{C}$, necessarily above $\mathfrak{m}$. If $k=\mathrm{A}/\mathfrak{m}$, then $\mathrm{C}/\mathfrak{m}\mathrm{C}$ is a finite type $k$-algebra, so the residue field $k'$ of $\mathrm{C}$ at $\mathfrak{n}$ is finite over $k$ \sref[I]{I.6.4.2}. By~(i), it is enough to prove that $\mathrm{C}_\mathfrak{n}$ is a $\boldsymbol{P}$-ring, since $\mathrm{C}_\mathfrak{q}$ is a local ring at a prime ideal of $\mathrm{C}_\mathfrak{n}$. We are thus reduced to the following lemma.
\end{enumerate}
\end{proof}

\begin{lemma}[7.4.4.1]
\label{IV.7.4.4.1}
Let $\mathrm{A}$ be a $\boldsymbol{P}$-ring local Noetherian, $k$ its residue field, $\mathrm{C}$ a finite type A-algebra, $\mathrm{B}$ a local ring in a prime ideal $\mathfrak{n}$ of $\mathrm{C}$, such that:
$1$\textsuperscript{o} the homomorphism $\mathrm{A}\to\mathrm{B}$ is local; $2$\textsuperscript{o} the residue field $k'$ of $\mathrm{B}$ is a finite extension of $k$.
If $\boldsymbol{P}$ satisfies $(\mathrm{P_{IV}})$, then $\mathrm{B}$ is a $\boldsymbol{P}$-ring.
\end{lemma}

\begin{proof}
Let $(x_i)_{1\leqslant i\leqslant m}$ be a system of generators of the $\mathrm{A}$-algebra $\mathrm{C}$; let us show first that we can argue by induction on $m$.
Let $\mathrm{C}'$ be the sub-$\mathrm{A}$-algebra of $\mathrm{C}$ generated by $x_1,\ldots,x_{m-1}$, and let $\mathfrak{n}'=\mathfrak{n}\cap\mathrm{C}'$.
The homomorphism $\mathrm{A}\to\mathrm{C}_\mathfrak{n}$ factors through the local homomorphisms $\mathrm{A}\to\mathrm{C}'_{\mathfrak{n}'}\to\mathrm{C}_\mathfrak{n}$; if $k''$ is the residue field of $\mathrm{C}'_{\mathfrak{n}'}$, the extension $k\to k'$ factors through $k''$, hence $k''$ is a finite extension of $k$ and $k'$ a finite extension of $k''$.
The induction hypothesis implies that $\mathrm{C}'_{\mathfrak{n}'}$ is a $\boldsymbol{P}$-ring. If $\mathrm{S}'=\mathrm{C}'-\mathfrak{n}'$, then $\mathrm{C}_\mathfrak{n}$ is a local ring of $\mathrm{S}'^{-1}\mathrm{C}=\mathrm{C}'_{\mathfrak{n}'}[x_m/1]$; the induction hypothesis applied once more shows that $\mathrm{C}_\mathfrak{n}$ is a $\boldsymbol{P}$-ring. Thus the problem reduces to the case where $\mathrm{C}$ is an $\mathrm{A}$-algebra generated by a \emph{single} element $t$.

Applying \sref{IV.7.4.3},~(ii), with $\mathrm{A}'=\widehat{\mathrm{A}}$, the ring $\mathrm{B}'=\mathrm{B}\otimes_\mathrm{A}\mathrm{A}'$ is a local ring of $\mathrm{C}\otimes_\mathrm{A}\mathrm{A}'$ at a prime ideal above $\mathfrak{n}$; as $\mathrm{A}$ and $\mathrm{A}'$ have the same residue field $k$, the residue field of $\mathrm{C}\otimes_\mathrm{A}\mathrm{A}'$ at this prime ideal is equal to $k'$.
Moreover $\mathrm{C}\otimes_\mathrm{A}\mathrm{A}'$ is generated by a single element over $\mathrm{A}'$.
It follows therefore from \sref{IV.7.4.3},~(ii) that one can restrict to proving \sref{IV.7.4.4.1} when $\boldsymbol{P}$ is the property~(iv') of \sref{IV.7.3.8}, that $\mathrm{A}$ is \emph{complete} and $\mathrm{C}$ generated by a single element $t$.

To show that the formal fibres of $\mathrm{B}=\mathrm{C}_\mathfrak{n}$ are then geometrically regular, let us apply the criterion \sref{IV.7.3.16},~$b)$).
Let $\mathrm{B}_1$ be a finite integral $\mathrm{B}$-algebra, hence generated by a finite number of elements integral over $\mathrm{B}$.
By multiplying these elements by an element of $\mathrm{S}=\mathrm{C}-\mathfrak{n}$, one can suppose that they are integral over $\mathrm{C}$, and we can therefore write $\mathrm{B}_1=\mathrm{S}^{-1}\mathrm{C}_1$, where $\mathrm{C}_1$ is a sub-C-algebra of $\mathrm{B}_1$ generated by a finite number of integral elements over $\mathrm{C}$, hence a finite integral C-algebra.
On the other hand, $\mathrm{B}_1$ is a semi-local ring, and every local ring $\mathrm{B}_2$ of $\mathrm{B}_1$ at a maximal ideal is a local ring of $\mathrm{C}_1$, such that $\mathrm{A}\to\mathrm{B}_2$ is a local homomorphism; moreover, the residue field of $\mathrm{B}_2$ is a finite extension of $k'$, hence of $k$.
We see therefore (taking into account \sref{IV.7.3.14} and~(i)) that we are reduced to the following question: let $\mathrm{C}$ be a Noetherian integral ring containing a subring $\mathrm{C}_0$ which is an $\mathrm{A}$-algebra generated by a single element $t$ and such that $\mathrm{C}$ is a finite $\mathrm{C}_0$-algebra; if $\mathfrak{n}$ is a maximal ideal of $\mathrm{C}$
\oldpage[IV-2]{201}
above the maximal ideal of $\mathrm{A}$, and $\mathrm{B}=\mathrm{C}_\mathfrak{n}$, the question is to show that for every prime ideal $\mathfrak{q}$ of $\widehat{\mathrm{B}}$ whose inverse image in $\mathrm{B}$ is $0$, the ring $\widehat{\mathrm{B}}_\mathfrak{q}$ is regular.
One may moreover replace $\mathrm{A}$ by its image in $\mathrm{C}$, which is a complete local ring (as a quotient of $\mathrm{A}$) and integral. The desired conclusion follows from the more general lemma below.
\end{proof}

\begin{lemma}[7.4.4.2]
\label{IV.7.4.4.2}
Let $\mathrm{A}$ be a complete Noetherian local integral ring, $k$ its residue field, $\mathrm{C}$ an integral ring containing $\mathrm{A}$, such that there exists $t\in\mathrm{C}$ for which $\mathrm{C}$ is a finite $\mathrm{A}[t]$-algebra.
Let $\mathfrak{n}$ be a maximal ideal of $\mathrm{C}$ above the maximal ideal $\mathfrak{m}$ of $\mathrm{A}$; set $\mathrm{B}=\mathrm{C}_\mathfrak{n}$, $\mathrm{X}=\Spec(\mathrm{B})$, $\mathrm{B}'=\widehat{\mathrm{B}}$, $\mathrm{X}'=\Spec(\widehat{\mathrm{B}})$; if $\mathrm{U}=\operatorname{Reg}(\mathrm{X})$, $\mathrm{U}'=\operatorname{Reg}(\mathrm{X}')$, and if $f:\mathrm{X}'\to\mathrm{X}$ is the canonical morphism, then $f^{-1}(\mathrm{U})\subset\mathrm{U}'$.
\end{lemma}

\begin{proof}
The assertion needed for \sref{IV.7.4.4.1} follows from this lemma because, since $\mathrm{B}$ is integral, the generic point of $\mathrm{X}$ belongs to $\mathrm{U}$.
Since $\mathrm{C}$ is an $\mathrm{A}$-algebra of finite type and $\mathfrak{n}$ is a maximal ideal of $\mathrm{C}$, the residue field $k'$ of $\mathrm{C}_\mathfrak{n}=\mathrm{B}$ (hence also of $\mathrm{B}'$) is a finite extension of $k$ \sref[I]{I.6.4.11} and \sref{IV.6.4.2}.

Set $\mathrm{Y}=\Spec(\mathrm{C})$. It follows from \sref{IV.6.12.8} that $\operatorname{Reg}(\mathrm{Y})$ is open in $\mathrm{Y}$; since the local rings of $\mathrm{X}$ are local rings of $\mathrm{Y}$ \sref[I]{I.2.4.2}, we have $\mathrm{U}=\mathrm{X}\cap\operatorname{Reg}(\mathrm{Y})$, hence $\mathrm{U}$ is open in $\mathrm{X}$. On the other hand, \sref{IV.6.12.7} shows that $\mathrm{U}'$ is open in $\mathrm{X}'$, hence $\mathrm{S}'=\mathrm{X}'-\mathrm{U}'$ is closed. Consequently $\mathrm{S}'\cap f^{-1}(\mathrm{U})$ is locally closed in $\mathrm{X}'$, and the whole question is to prove that this set is empty.
If it were non-empty, \sref{IV.5.1.10} would give a prime ideal $\mathfrak{p}'$ in $\mathrm{S}'\cap f^{-1}(\mathrm{U})$ such that $\dim(\mathrm{B}'/\mathfrak{p}')\leqslant 1$.
The prime ideal $\mathfrak{p}'$ cannot be the maximal ideal $\mathfrak{m}\mathrm{B}'$ of $\mathrm{B}'$, where $\mathfrak{m}$ is the maximal ideal of $\mathrm{B}$. Indeed, this would mean that $\mathrm{B}=\mathrm{B}_\mathfrak{m}$ is regular, hence so is $\mathrm{B}'=\widehat{\mathrm{B}}$ \sref[0]{0.17.1.5}, and we would have $\mathfrak{m}\mathrm{B}'\in\mathrm{U}'$, contrary to the hypothesis. We must therefore have $\dim(\mathrm{B}'/\mathfrak{p}')=1$.

Set $\mathfrak{p}=\mathrm{B}\cap\mathfrak{p}'$. By hypothesis $\mathrm{B}_\mathfrak{p}$ is regular, whereas $\mathrm{B}'_{\mathfrak{p}'}$ is not. Since $\mathrm{B}'_{\mathfrak{p}'}$ is a flat $\mathrm{B}_\mathfrak{p}$-module, \sref{IV.6.5.2} shows that the fibre $\mathrm{Z}$ of $f$ at $\mathfrak{p}$ is not regular at $\mathfrak{p}'$.
We may reduce to the case $\mathfrak{p}=0$. Indeed, in general, set $\mathfrak{q}=\mathfrak{p}\cap\mathrm{C}$ and $\mathfrak{r}=\mathfrak{p}\cap\mathrm{A}$. Then $\mathrm{C}/\mathfrak{q}$ is a finite $(\mathrm{A}/\mathfrak{r})[\bar{t}]$-algebra, where $\bar{t}$ is the class of $t$ modulo $\mathfrak{q}$; since $\mathfrak{p}=\mathfrak{q}\mathrm{B}$, the ring $\mathrm{B}/\mathfrak{p}$ is equal to $(\mathrm{C}/\mathfrak{q})_{\mathfrak{n}/\mathfrak{q}}$, and $\mathfrak{n}/\mathfrak{q}$ is a maximal ideal of $\mathrm{C}/\mathfrak{q}$ lying over the maximal ideal $\mathfrak{m}/\mathfrak{r}$ of $\mathrm{A}/\mathfrak{r}$. Thus the hypotheses of \sref{IV.7.4.4.2} are still satisfied by $\mathrm{A}/\mathfrak{r}$, $\mathrm{C}/\mathfrak{q}$, and $\mathrm{B}/\mathfrak{p}$; since the completion of $\mathrm{B}/\mathfrak{p}$ is $\mathrm{B}'/\mathfrak{p}\mathrm{B}'$, this proves the reduction.

Suppose therefore that $\mathfrak{p}=0$, so that $\mathrm{Z}$ is the generic fibre of $f$ and the homomorphism $\mathrm{B}\to\mathrm{B}'/\mathfrak{p}'$ is injective. Set $\mathrm{V}=\mathrm{B}'/\mathfrak{p}'$, and distinguish two cases.

\emph{I) $\mathrm{V}$ is a finite $\mathrm{A}$-algebra.} Since $\mathrm{B}\subset\mathrm{V}$, $\mathrm{B}$ is \emph{a fortiori} a finite $\mathrm{A}$-algebra; as $\mathrm{A}$ is complete, so is $\mathrm{B}$ (Bourbaki, \emph{Alg.\ comm.}, chap.~IV, \textsection\,2, n\textsuperscript{o}\,5, cor.~3 of prop.~9). Hence $\mathrm{B}'=\mathrm{B}$ and $\mathfrak{p}'=0$, so $\mathrm{B}'_{\mathfrak{p}'}$ is a field, contrary to the hypothesis.

\emph{II) $\mathrm{V}$ is not a finite $\mathrm{A}$-algebra.} Since the local ring $\mathrm{A}$ is complete, $\mathrm{V}$ is not a quasi-finite $\mathrm{A}$-algebra \sref[0]{0.7.4.1} and \sref[0]{0.7.4.2}; but the residue field $k'$ of $\mathrm{V}$ is a finite extension of the residue field $k$ of $\mathrm{A}$,

\oldpage[IV-2]{202}

hence \sref[0]{0.7.4.4} the ideal $\mathfrak{m}\mathrm{V}$ is not an ideal of definition of $\mathrm{V}$. Since $\mathrm{V}$ is a one-dimensional Noetherian local integral ring, $0$ is the only ideal of $\mathrm{V}$ that is not an ideal of definition, hence $\mathfrak{m}\mathrm{V}=0$. But $\mathrm{A}\subset\mathrm{V}$ and $\mathrm{V}$ is integral, so $\mathfrak{m}=0$ and $\mathrm{A}=k$ is a field.
It follows first that $\dim(\mathrm{C})\leqslant 1$ \sref[0]{0.16.1.5}; but since $\dim(\mathrm{V})=1$, the inequalities
\[
  \dim(\mathrm{V})\leqslant\dim(\mathrm{B}')=\dim(\mathrm{B})\leqslant\dim(\mathrm{C})
\]
show that $\dim(\mathrm{C})=\dim(\mathrm{B})=\dim(\mathrm{B}')=\dim(\mathrm{B}'/\mathfrak{p}')=1$, and consequently $\mathfrak{p}'$ is a minimal prime ideal of $\mathrm{B}'$. We therefore obtain a contradiction once we prove that $\mathrm{B}'_{\mathfrak{p}'}$ is a field, or equivalently that $\mathrm{B}'$ is reduced.
Since $\mathrm{C}$ is a finite-type $k$-algebra, its integral closure $\mathrm{C}_1$ is a finite $\mathrm{C}$-algebra (Bourbaki, \emph{Alg.\ comm.}, chap.~V, \textsection\,3, n\textsuperscript{o}\,2, th.~2). If $\mathrm{S}=\mathrm{C}-\mathfrak{n}$, then $\mathrm{B}_1=\mathrm{S}^{-1}\mathrm{C}_1$ is the integral closure of $\mathrm{B}$, hence a finite $\mathrm{B}$-algebra and therefore a one-dimensional semi-local Noetherian integral and integrally closed ring \sref[0]{0.16.1.5}. If $\mathfrak{m}_j$ ($1\leqslant j\leqslant h$) are its maximal ideals, the rings $(\mathrm{B}_1)_{\mathfrak{m}_j}$ are discrete valuation rings \sref[II]{II.7.1.6}, and the completion $\mathrm{B}'_1$ of $\mathrm{B}_1$ is the direct product of the completed discrete valuation rings $(\mathrm{B}_1)_{\mathfrak{m}_j}$ (Bourbaki, \emph{Alg.\ comm.}, chap.~III, \textsection\,2, n\textsuperscript{o}\,13, prop.~18). Thus $\mathrm{B}'_1$ is reduced; since the completion $\mathrm{B}'$ of $\mathrm{B}$ is a subring of $\mathrm{B}'_1$ (Bourbaki, \emph{Alg.\ comm.}, chap.~IV, \textsection\,2, n\textsuperscript{o}\,5, cor.~3 of prop.~9), $\mathrm{B}'$ is reduced as well.
\end{proof}

\begin{corollary}[7.4.5]
\label{IV.7.4.5}
Suppose that the property $\boldsymbol{P}$ satisfies the conditions $(\mathrm{P_I})$, $(\mathrm{P_{II}})$, $(\mathrm{P_{III}})$.
Let $\mathrm{A}$ be a Noetherian ring.
The following conditions are equivalent:
\begin{enumerate}[label=\emph{\alph*)}]
  \item For every prime ideal $\mathfrak{p}$ of $\mathrm{A}$, $\mathrm{A}_\mathfrak{p}$ is a $\boldsymbol{P}$-ring.
  \item For every maximal ideal $\mathfrak{m}$ of $\mathrm{A}$, $\mathrm{A}_\mathfrak{m}$ is a $\boldsymbol{P}$-ring.
\end{enumerate}
If in addition $\boldsymbol{P}$ satisfies $(\mathrm{P_{IV}})$, the properties $a)$ and $b)$ are also equivalent to:
\begin{enumerate}[label=\emph{\alph*)}]
  \setcounter{enumi}{2}
  \item For every finite type A-algebra $\mathrm{B}$ and every prime ideal $\mathfrak{q}$ of $\mathrm{B}$, $\mathrm{B}_\mathfrak{q}$ is a $\boldsymbol{P}$-ring.
\end{enumerate}
\end{corollary}

\begin{proof}
The equivalence of $a)$ and $b)$ follows from \sref{IV.7.4.4},~(i), and that of $a)$ and $c)$ from \sref{IV.7.4.4},~(ii).
When condition $a)$ of \sref{IV.7.4.5} is satisfied, we say that $\mathrm{A}$ is a $\boldsymbol{P}$-ring; for semi-local Noetherian rings, this definition coincides with that of \sref{IV.7.3.13}, when the conditions $(\mathrm{P_I})$, $(\mathrm{P_{II}})$ and $(\mathrm{P_{III}})$ are satisfied.
The ring $\bb{Z}$ is a $\boldsymbol{P}$-ring \sref{IV.7.3.19},~(iv); every local Noetherian complete ring is a $\boldsymbol{P}$-ring.
\end{proof}

\begin{proposition}[7.4.6]
\label{IV.7.4.6}
Suppose that the property $\boldsymbol{P}$ satisfies the conditions $(\mathrm{P_I})$, $(\mathrm{P_{II}})$ and $(\mathrm{P_{III}})$.
Let $\mathrm{A}$ be a Noetherian ring, $\mathfrak{J}$ an ideal of $\mathrm{A}$, $\widehat{\mathrm{A}}$ the separated completion of $\mathrm{A}$ for the $\mathfrak{J}$-adic topology.
Then, if $\mathrm{A}$ is a $\boldsymbol{P}$-ring \sref{IV.7.4.5}, the canonical morphism $\Spec(\widehat{\mathrm{A}})\to\Spec(\mathrm{A})$ is a $\boldsymbol{P}$-morphism.
\end{proposition}

\begin{proof}
Using \sref{IV.7.3.2},~$c')$), it suffices to prove that for every maximal ideal $\mathfrak{n}$ of $\widehat{\mathrm{A}}$ with inverse image $\mathfrak{m}$ in $\mathrm{A}$, the morphism $\Spec((\widehat{\mathrm{A}})_\mathfrak{n})\to\Spec(\mathrm{A}_\mathfrak{m})$ is a $\boldsymbol{P}$-morphism.
We know (Bourbaki, \emph{Alg.\ comm.}, chap.~III, \textsection\,3, n\textsuperscript{o}\,4, prop.~8) that the canonical homomorphism $\mathrm{A}_\mathfrak{m}\to(\widehat{\mathrm{A}})_\mathfrak{n}$ is injective, that the $\mathfrak{m}\mathrm{A}_\mathfrak{m}$-adic topology on $\mathrm{A}_\mathfrak{m}$ is induced by the $\mathfrak{n}(\widehat{\mathrm{A}})_\mathfrak{n}$-adic topology, and that $\mathrm{A}_\mathfrak{m}$ is dense in $(\widehat{\mathrm{A}})_\mathfrak{n}$ for this topology, so that the completion of $\mathrm{A}_\mathfrak{m}$ for the $\mathfrak{m}\mathrm{A}_\mathfrak{m}$-adic topology identifies with that of $(\widehat{\mathrm{A}})_\mathfrak{n}$ for the $\mathfrak{n}(\widehat{\mathrm{A}})_\mathfrak{n}$-adic topology.
We have therefore two morphisms
\[
  \Spec((\mathrm{A}_\mathfrak{m})\widehat{\phantom{x}})\xrightarrow{f}\Spec((\widehat{\mathrm{A}})_\mathfrak{n})\xrightarrow{g}\Spec(\mathrm{A}_\mathfrak{m})
\]

\oldpage[IV-2]{203}

such that $f$ is faithfully flat. Since, by hypothesis, $g\circ f$ is a $\boldsymbol{P}$-morphism, it follows from $(\mathrm{P_{II}})$ that $g$ is a $\boldsymbol{P}$-morphism as well.
\end{proof}

\begin{corollary}[7.4.7]
\label{IV.7.4.7}
Suppose that $\boldsymbol{P}$ satisfies the conditions $(\mathrm{P_I})$, $(\mathrm{P}'_\mathrm{I})$, $(\mathrm{P_{II}})$, $(\mathrm{P_{III}})$ and $(\mathrm{P_{IV}})$.
Then, if $\mathrm{A}$ is a $\boldsymbol{P}$-ring \sref{IV.7.4.5}, the canonical morphism $\Spec(\mathrm{A}[[\mathrm{T}_1,\ldots,\mathrm{T}_r]])\to\Spec(\mathrm{A})$ is a $\boldsymbol{P}$-morphism.
In particular, if in addition $\mathrm{A}$ is integral, $\mathrm{K}$ its field of fractions, and if $\mathfrak{p}$ is a prime ideal of $\mathrm{B}=\mathrm{A}[[\mathrm{T}_1,\ldots,\mathrm{T}_r]]$ such that $\mathfrak{p}\cap\mathrm{A}=0$, then the property $\boldsymbol{P}(\mathrm{B}_\mathfrak{p},\mathrm{K})$ is true.
\end{corollary}

\begin{proof}
Indeed, $\Spec(\mathrm{B})\to\Spec(\mathrm{A})$ factorises as
\[
  \Spec(\mathrm{A}[[\mathrm{T}_1,\ldots,\mathrm{T}_r]])\xrightarrow{f}\Spec(\mathrm{A}[\mathrm{T}_1,\ldots,\mathrm{T}_r])\xrightarrow{g}\Spec(\mathrm{A}).
\]
It is clear that the morphism $g$ is regular \sref[0]{0.17.3.7}; by virtue of \sref{IV.7.4.5}, $\mathrm{A}[\mathrm{T}_1,\ldots,\mathrm{T}_r]$ is a $\boldsymbol{P}$-ring, hence it follows from \sref{IV.7.4.6} that $f$ is a $\boldsymbol{P}$-morphism; as $g$ is regular, hence also a $\boldsymbol{P}$-morphism \sref{IV.7.3.5},~(i), it follows from $(\mathrm{P}'_\mathrm{I})$ that $g\circ f$ is a $\boldsymbol{P}$-morphism.

One will note that the conclusion is still valid if, instead of supposing that $\boldsymbol{P}$ satisfies $(\mathrm{P}'_\mathrm{I})$ and that every regular morphism is a $\boldsymbol{P}$-morphism, one supposes only that a composite $g\circ f$ is a $\boldsymbol{P}$-morphism whenever $g$ is regular and $f$ is a $\boldsymbol{P}$-morphism (a condition in some sense symmetric to $(\mathrm{P}'_\mathrm{I})$).
\end{proof}

\begin{remark}[7.4.8]
\label{IV.7.4.8}
The preceding results pose the following problems:
\begin{enumerate}[label=\rm\Alph*)]
  \item Let $\mathrm{A}$ be a Zariski \emph{complete} ring, and let $\mathfrak{J}$ be an ideal of definition of $\mathrm{A}$; if the ring $\mathrm{A}/\mathfrak{J}$ is a $\boldsymbol{P}$-ring, is $\mathrm{A}$ also a $\boldsymbol{P}$-ring?
  It would follow that for every Noetherian $\boldsymbol{P}$-ring $\mathrm{A}$ and every ideal $\mathfrak{J}$ of $\mathrm{A}$, the separated completion $\widehat{\mathrm{A}}$ for the $\mathfrak{J}$-adic topology would also be a $\boldsymbol{P}$-ring.
  \item Let $k$ be a valued field, complete and non-discrete; one calls \emph{ring of restricted formal power series} $k\{\mathrm{T}_1,\ldots,\mathrm{T}_n\}$ the subring of the formal power series ring $k[[\mathrm{T}_1,\ldots,\mathrm{T}_n]]$ formed of the series whose coefficients \emph{tend to} $0$.
  Is such a ring a $\boldsymbol{P}$-ring?
  \item If $\mathrm{A}$ is a linearly topologised $\boldsymbol{P}$-ring, $\mathrm{S}$ a multiplicative subset of $\mathrm{A}$, are the rings $\mathrm{A}\{S^{-1}\}$ \sref[0]{0.7.6.1} and $\mathrm{A}_{\{S\}}$ \sref[0]{0.7.6.15} $\boldsymbol{P}$-rings?
\end{enumerate}
\end{remark}


\subsection{A criterion for \texorpdfstring{$\boldsymbol{P}$}{P}-morphisms}
\label{subsection:IV.7.5}

\begin{env}[7.5.0]
\label{IV.7.5.0}
This number will not be used in the sequel of Chapter~IV and may therefore be omitted on a first reading. We shall moreover see later \sref{IV.7.9.8} that the results of the present number can be considerably improved when one has ``resolution of singularities'' at one's disposal.

In the sequel of this number, we shall consider a property $\boldsymbol{R}(\mathrm{A})$, and we shall denote by $\boldsymbol{P}(\mathrm{Z},k)$ the following property:
\begin{quote}
``$\mathrm{Z}$ is a locally Noetherian prescheme over a field $k$, and, for every finite extension $k'$ of $k$, if $\mathrm{Z}'=\mathrm{Z}\otimes_k k'$, the property $\boldsymbol{R}(\sh{O}_{\mathrm{Z}',z'})$ is true for every $z'\in\mathrm{Z}'.$'
\end{quote}
\end{env}

\begin{theorem}[7.5.1]
\label{IV.7.5.1}
Let $\mathrm{A}$ and $\mathrm{B}$ be two complete Noetherian local rings, let $\mathfrak{m}$ be the maximal ideal of $\mathrm{A}$, let $k=\mathrm{A}/\mathfrak{m}$ be its residue field, and let $\varphi:\mathrm{A}\to\mathrm{B}$ be a local homomorphism such that:
\begin{enumerate}[label=(\roman*)]
  \item the residue field of $\mathrm{B}$ is a finite extension of $k$;
  \item $\mathrm{B}$ is a flat $\mathrm{A}$-module.
\end{enumerate}
Let $\boldsymbol{R}(\mathrm{C})$ moreover be a property satisfying condition $(\mathrm{R}_{\mathrm{III}})$ \sref{IV.7.3.18} and the following condition:
\begin{quote}
$(\mathrm{R}_{\mathrm{IV}})$ For every local ring $\mathrm{C}$ at a prime ideal of a complete Noetherian local ring, and every $\mathrm{C}$-regular element $t$ in the maximal ideal of $\mathrm{C}$, the property $\boldsymbol{R}(\mathrm{C}/t\mathrm{C})$ implies $\boldsymbol{R}(\mathrm{C})$.
\end{quote}

\oldpage[IV-2]{204}

Let $f={}^a\varphi:\Spec(\mathrm{B})\to\Spec(\mathrm{A})$, and suppose that $\boldsymbol{P}(\Spec(\mathrm{B}\otimes_\mathrm{A}k),k)$ is true. Then, for every $x\in\Spec(\mathrm{A})$, the property $\boldsymbol{P}(f^{-1}(x),\boldsymbol{k}(x))$ is true.
\end{theorem}

\begin{proof}
In other words, property $\boldsymbol{P}$ for the fibre over the closed point of $\Spec(\mathrm{A})$ entails the same property for all fibres of $f$; that is, $f$ is a $\boldsymbol{P}$-morphism \sref{IV.7.3.1}. We proceed in several steps.

\emph{I) Reduction to the study of the local rings of the generic fibre.} Apply \sref{IV.7.3.16.2}. It suffices to show that, for every finite integral $\mathrm{A}$-algebra $\mathrm{A}'$ with field of fractions $\mathrm{K}'$, if $\mathrm{B}'=\mathrm{B}\otimes_\mathrm{A}\mathrm{A}'$, then every local ring of $\Spec(\mathrm{B}'\otimes_{\mathrm{A}'}\mathrm{K}')$ satisfies $\boldsymbol{R}$. The ring $\mathrm{A}'$ (resp.\ $\mathrm{B}'$) is complete and semi-local, the latter because it is a finite $\mathrm{B}$-algebra, and is therefore a direct product of complete local rings. Since $\mathrm{A}'$ is integral it is local. Every maximal ideal $\mathfrak{n}'$ of $\mathrm{B}'$ lies over the maximal ideal $\mathfrak{n}$ of $\mathrm{B}$, hence over $\mathfrak{m}$, and its inverse image in $\mathrm{A}'$ is consequently the maximal ideal $\mathfrak{m}'$ of $\mathrm{A}'$. Every local ring of $\Spec(\mathrm{B}'\otimes_{\mathrm{A}'}\mathrm{K}')$ is a local ring of $\Spec(\mathrm{B}')$, and hence of one of the $\Spec(\mathrm{B}'_{\mathfrak{n}'})$ over the generic point of $\mathrm{A}'$. Thus it is enough to prove that the local rings of
\[
  \Spec(\mathrm{B}'_{\mathfrak{n}'}\otimes_{\mathrm{A}'}\mathrm{K}')
\]
satisfy $\boldsymbol{R}$. Now $\mathrm{A}'$ and $\mathrm{B}'_{\mathfrak{n}'}$ are complete Noetherian local rings. The residue field of $\mathrm{B}'_{\mathfrak{n}'}$ is a finite extension of that of $\mathrm{B}$, hence also of $k$, and \emph{a fortiori} of the residue field $k'$ of $\mathrm{A}'$. This observation and \sref{IV.2.1.4} show that $\mathrm{A}'$ and $\mathrm{B}'_{\mathfrak{n}'}$ satisfy conditions~(i) and~(ii) of the statement. On the other hand, if $k''$ is a finite extension of $k'$, it is also a finite extension of $k$, and the local rings of
\[
  \Spec(\mathrm{B}'_{\mathfrak{n}'}\otimes_{\mathrm{A}'}k'')
\]
are also local rings of $\Spec(\mathrm{B}\otimes_\mathrm{A}k'')$. The hypothesis $\boldsymbol{P}(\Spec(\mathrm{B}\otimes_\mathrm{A}k),k)$ therefore implies
\[
  \boldsymbol{P}(\Spec(\mathrm{B}'_{\mathfrak{n}'}\otimes_{\mathrm{A}'}k'),k').
\]
We are thus reduced to the case where $\mathrm{A}$ is integral with field of fractions $\mathrm{K}$, and to proving that the local rings of $\Spec(\mathrm{B}\otimes_\mathrm{A}\mathrm{K})$ satisfy $\boldsymbol{R}$.

\emph{II) The case $\dim(\mathrm{A})=1$.} Let $\mathrm{A}'$ be the integral closure of $\mathrm{A}$. By \sref[0]{0.23.1.6}, $\mathrm{A}'$ is a finite $\mathrm{A}$-algebra and a complete local ring. If $\mathrm{B}'=\mathrm{B}\otimes_\mathrm{A}\mathrm{A}'$, then
\[
  \mathrm{B}\otimes_\mathrm{A}\mathrm{K}=\mathrm{B}'\otimes_{\mathrm{A}'}\mathrm{K}.
\]
The same argument as in~I) shows that it suffices, for every maximal ideal $\mathfrak{n}'$ of $\mathrm{B}'$, to prove that the local rings of $\mathrm{B}'_{\mathfrak{n}'}\otimes_{\mathrm{A}'}\mathrm{K}$ satisfy $\boldsymbol{R}$; it also shows that $\mathrm{A}'$ and $\mathrm{B}'_{\mathfrak{n}'}$ satisfy hypotheses~(i) and~(ii), and that
\[
  \boldsymbol{P}(\Spec(\mathrm{B}'_{\mathfrak{n}'}\otimes_{\mathrm{A}'}k'),k')
\]
is true, where $k'$ is the residue field of $\mathrm{A}'$. Since $\dim(\mathrm{A}')=1$ \sref[0]{0.16.1.5} and $\mathrm{A}'$ is integral and integrally closed, it is a complete discrete valuation ring. We may therefore suppose that $\mathrm{A}$ itself is a complete discrete valuation ring. If $u$ is a uniformizer of $\mathrm{A}$, then $u$ is $\mathrm{A}$-regular, and the flatness of $\mathrm{B}$ over $\mathrm{A}$ implies that $t=\varphi(u)$ is $\mathrm{B}$-regular and belongs to the maximal ideal of $\mathrm{B}$. The hypothesis on $\Spec(\mathrm{B}\otimes_\mathrm{A}k)$ implies in particular that $\boldsymbol{R}(\mathrm{B}/t\mathrm{B})$ is true; hence $\boldsymbol{R}(\mathrm{B})$ is true by $(\mathrm{R}_{\mathrm{IV}})$. Thus $\mathrm{U}_{\boldsymbol{R}}(\Spec(\mathrm{B}))$ contains the closed point. It is open by $(\mathrm{R}_{\mathrm{III}})$, and $\Spec(\mathrm{B})$ is the only open subset of itself containing its closed point. Consequently all local rings of $\Spec(\mathrm{B})$, and in particular those of $\Spec(\mathrm{B}\otimes_\mathrm{A}\mathrm{K})$, satisfy $\boldsymbol{R}$.

\emph{III) The general case.} The case $\dim(\mathrm{A})=0$ is trivial, since then $\mathrm{A}=k$; we may therefore suppose that $\dim(\mathrm{A})\geqslant1$. By \sref{IV.6.12.7}, $\Spec(\mathrm{A})$ contains a non-empty open subset $\mathrm{V}$ all of whose points are regular. The intersection $\mathrm{V}'$ of $\mathrm{V}$ with the complement of the closed point of $\Spec(\mathrm{A})$ is non-empty and therefore contains the generic point. It suffices to prove that
\[
  f^{-1}(\mathrm{V}')\subset\mathrm{U}_{\boldsymbol{R}}(\Spec(\mathrm{B})).
\]
Equivalently, the intersection $\mathrm{Z}$ of $f^{-1}(\mathrm{V}')$ with the complement of $\mathrm{U}_{\boldsymbol{R}}(\Spec(\mathrm{B}))$ must be empty. Suppose otherwise. Since $\mathrm{U}_{\boldsymbol{R}}(\Spec(\mathrm{B}))$ is open by $(\mathrm{R}_{\mathrm{III}})$, $\mathrm{Z}$ is locally closed in $\Spec(\mathrm{B})$; if non-empty, it contains a point $x$ such that $\dim(\overline{\{x\}})\leqslant1$ \sref{IV.5.1.10}. Since $f(x)\in\mathrm{V}'$ is distinct from the closed point of $\Spec(\mathrm{A})$, $x$ is distinct from the closed point of $\Spec(\mathrm{B})$, and therefore $\dim(\overline{\{x\}})=1$.

We show this is impossible. Let $\mathfrak{q}$ be the corresponding prime ideal of $\mathrm{B}$, so that $\dim(\mathrm{B}/\mathfrak{q})=1$, and put $\mathfrak{p}=\varphi^{-1}(\mathfrak{q})$. Then $\mathfrak{p}\neq\mathfrak{m}$ and $\mathrm{A}_\mathfrak{p}$ is regular. Since $\mathfrak{p}\neq\mathfrak{m}$, the ideal $\mathfrak{m}(\mathrm{B}/\mathfrak{q})$ is non-zero, hence is an ideal of definition of the one-dimensional Noetherian local integral ring $\mathrm{B}/\mathfrak{q}$. By~(i), the residue field of $\mathrm{B}/\mathfrak{q}$ is a finite extension of the residue field of $\mathrm{A}/\mathfrak{p}$, so $\mathrm{B}/\mathfrak{q}$ is a quasi-finite $(\mathrm{A}/\mathfrak{p})$-algebra \sref[0]{0.7.4.4}. But $\mathrm{A}/\mathfrak{p}$ is complete and $\mathrm{B}/\mathfrak{q}$ is separated for the $\mathfrak{m}$-preadic topology, which is its local-ring topology; hence $\mathrm{B}/\mathfrak{q}$ is finite over $\mathrm{A}/\mathfrak{p}$ by \sref[0]{0.7.4.1}. The injection $\mathrm{A}/\mathfrak{p}\to\mathrm{B}/\mathfrak{q}$ consequently gives
\[
  \dim(\mathrm{A}/\mathfrak{p})=\dim(\mathrm{B}/\mathfrak{q})=1
\]
by \sref[0]{0.16.1.5}. Apply~II) to $\mathrm{A}/\mathfrak{p}$ and $\mathrm{B}/\mathfrak{p}\mathrm{B}$. Their residue fields are respectively those of $\mathrm{A}$ and $\mathrm{B}$, $\mathrm{B}/\mathfrak{p}\mathrm{B}$ is flat over $\mathrm{A}/\mathfrak{p}$, and
\[
 (\mathrm{B}/\mathfrak{p}\mathrm{B})\otimes_{\mathrm{A}/\mathfrak{p}}k
 =\mathrm{B}\otimes_\mathrm{A}k=\mathrm{B}/\mathfrak{m}\mathrm{B}.
\]
It follows that the local rings of
\[
 \Spec\bigl((\mathrm{B}/\mathfrak{p}\mathrm{B})\otimes_{\mathrm{A}/\mathfrak{p}}\boldsymbol{k}(\mathfrak{p})\bigr)
\]
satisfy $\boldsymbol{R}$. In particular, $\mathrm{B}_\mathfrak{q}/\mathfrak{p}\mathrm{B}_\mathfrak{q}$ is one of these local rings. Moreover, $\mathrm{B}_\mathfrak{q}$ is flat over the regular local ring $\mathrm{A}_\mathfrak{p}$. We now use the following lemma.

\begin{lemma}[7.5.1.1]
\label{IV.7.5.1.1}
\oldpage[IV-2]{205}
Let $\mathcal{C}$ be a full subcategory of the category of Noetherian local rings such that every quotient ring of a ring in $\mathcal{C}$ still belongs to $\mathcal{C}$. Suppose $\boldsymbol{R}(\mathrm{C})$ has the property that, if $\mathrm{C}\in\mathcal{C}$ and $t$ is a regular element of the maximal ideal of $\mathrm{C}$ such that $\boldsymbol{R}(\mathrm{C}/t\mathrm{C})$ is true, then $\boldsymbol{R}(\mathrm{C})$ is true.

Let $\mathrm{C}$ be a regular local ring with residue field $k$, let $\mathrm{D}\in\mathcal{C}$ be a local ring, and let $\psi:\mathrm{C}\to\mathrm{D}$ be a local homomorphism making $\mathrm{D}$ a flat $\mathrm{C}$-module. If $\boldsymbol{R}(\mathrm{D}\otimes_\mathrm{C}k)$ is true, then $\boldsymbol{R}(\mathrm{D})$ is true.
\end{lemma}

\begin{proof}
Proceed by induction on $n=\dim(\mathrm{C})$. The assertion is true by hypothesis when $n=0$, for then $\mathrm{C}=k$. Let $t\in\mathfrak{m}-\mathfrak{m}^2$. Then $\mathrm{C}/t\mathrm{C}$ is regular and $\dim(\mathrm{C}/t\mathrm{C})=n-1$ \sref[0]{0.17.1.8} and \sref[0]{0.16.3.4}. Since
\[
  (\mathrm{D}/t\mathrm{D})\otimes_{\mathrm{C}/t\mathrm{C}}k
  =\mathrm{D}\otimes_\mathrm{C}k
\]
and $\mathrm{D}/t\mathrm{D}\in\mathcal{C}$, the induction hypothesis gives $\boldsymbol{R}(\mathrm{D}/t\mathrm{D})$. Moreover, $t$ is $\mathrm{C}$-regular, hence $\mathrm{D}$-regular by flatness \sref[0]{0.6.3.4}; the assumed property of $\boldsymbol{R}$ therefore gives $\boldsymbol{R}(\mathrm{D})$.
\end{proof}

To finish the proof of \sref{IV.7.5.1}, apply \sref{IV.7.5.1.1} with $\mathcal{C}$ the category of local rings of spectra of complete local rings, using $(\mathrm{R}_{\mathrm{IV}})$, and take $\mathrm{C}=\mathrm{A}_\mathfrak{p}$ and $\mathrm{D}=\mathrm{B}_\mathfrak{q}$.
\end{proof}

\begin{corollary}[7.5.2]
\label{IV.7.5.2}
Let $\mathrm{A}$ be a Noetherian local ring, let $\mathfrak{m}$ be its maximal ideal and $k=\mathrm{A}/\mathfrak{m}$ its residue field, let $\mathrm{B}$ be a Noetherian local ring, and let $\varphi:\mathrm{A}\to\mathrm{B}$ be a local homomorphism such that:
\begin{enumerate}[label=(\roman*)]
  \item the residue field of $\mathrm{B}$ is a finite extension of $k$;
  \item $\mathrm{B}$ is a flat $\mathrm{A}$-module.
\end{enumerate}
Let $\boldsymbol{R}(\mathrm{C})$ be a property satisfying $(\mathrm{R}_{\mathrm{II}})$ \sref{IV.7.3.10}, $(\mathrm{R}_{\mathrm{III}})$ \sref{IV.7.3.18}, and $(\mathrm{R}_{\mathrm{IV}})$ \sref{IV.7.5.1}. Let $\boldsymbol{R}'(\mathrm{C})$ be a property satisfying:
\begin{quote}
$(\mathrm{R}_{\mathrm{V}})$ If $\mathrm{C}$ and $\mathrm{D}$ are Noetherian local rings and $\psi:\mathrm{C}\to\mathrm{D}$ is a local homomorphism such that $g={}^a\psi:\Spec(\mathrm{D})\to\Spec(\mathrm{C})$ is a $\boldsymbol{P}$-morphism, then $\boldsymbol{R}'(\mathrm{C})$ implies $\boldsymbol{R}(\mathrm{D})$.
\end{quote}
Suppose that the canonical morphism $\Spec(\widehat{\mathrm{A}})\to\Spec(\mathrm{A})$ is a $\boldsymbol{P}'$-morphism, where $\boldsymbol{P}'$ is defined from $\boldsymbol{R}'$ as $\boldsymbol{P}$ is defined from $\boldsymbol{R}$ in \sref{IV.7.5.0}. Let $f={}^a\varphi:\Spec(\mathrm{B})\to\Spec(\mathrm{A})$, and suppose that
\[
  \boldsymbol{P}\bigl(\Spec(\widehat{\mathrm{B}}\otimes_{\widehat{\mathrm{A}}}k),k\bigr)
\]
is true. Then $\boldsymbol{P}(f^{-1}(x),\boldsymbol{k}(x))$ is true for every $x\in\Spec(\mathrm{A})$; in other words, $f$ is a $\boldsymbol{P}$-morphism.
\end{corollary}

\begin{proof}
We proceed in two steps.

\emph{I) Reduction to the study of the local rings of the generic fibre.} Apply \sref{IV.7.3.16.2} once more. The only difference from part~I) of the proof of \sref{IV.7.5.1} is that $\mathrm{A}'$ is now only a semi-local integral ring, not necessarily local. Every maximal ideal $\mathfrak{n}'$ of $\mathrm{B}'$ lies above a maximal ideal $\mathfrak{m}'$ of $\mathrm{A}'$. Every local ring of $\Spec(\mathrm{B}'\otimes_{\mathrm{A}'}\mathrm{K}')$ is a local ring of $\Spec(\mathrm{B}')$, hence of one of the $\Spec(\mathrm{B}'_{\mathfrak{n}'})$ over the generic point of $\Spec(\mathrm{A}'_{\mathfrak{m}'})$. Thus it is enough to prove that the local rings of
\[
  \Spec(\mathrm{B}'_{\mathfrak{n}'}\otimes_{\mathrm{A}'_{\mathfrak{m}'}}\mathrm{K}')
\]
satisfy $\boldsymbol{R}$. By definition of $\boldsymbol{P}'$, if $\boldsymbol{P}'(\mathrm{Z},k)$ is true then so is $\boldsymbol{P}'(\mathrm{Z}\otimes_k k',k')$ for every finite extension $k'$ of $k$; the same argument as in \sref{IV.7.3.7} shows that the formal fibres of $\mathrm{A}'_{\mathfrak{m}'}$ satisfy $\boldsymbol{P}'$. As in part~I) of the proof of \sref{IV.7.5.1}, the rings $\mathrm{A}'_{\mathfrak{m}'}$ and $\mathrm{B}'_{\mathfrak{n}'}$ satisfy conditions~(i) and~(ii) above. If $k'$ is the residue field of $\mathrm{A}'_{\mathfrak{m}'}$, then
\[
  \boldsymbol{P}\bigl(\Spec(\widehat{\mathrm{B}'_{\mathfrak{n}'}}\otimes_{\widehat{\mathrm{A}'_{\mathfrak{m}'}}}k'),k'\bigr)
\]
is true: indeed, the completion of $\mathrm{A}'_{\mathfrak{m}'}$ is one of the local components of $\widehat{\mathrm{A}}'=\widehat{\mathrm{A}}\otimes_\mathrm{A}\mathrm{A}'$, and the completion of $\mathrm{B}'_{\mathfrak{n}'}$ is one of the local components of
\[
  \widehat{\mathrm{B}}'=\widehat{\mathrm{B}}\otimes_\mathrm{B}\mathrm{B}'
  =\widehat{\mathrm{B}}\otimes_\mathrm{A}\mathrm{A}'
  =\widehat{\mathrm{B}}\otimes_{\widehat{\mathrm{A}}}\widehat{\mathrm{A}}'.
\]
Part~I) of the proof of \sref{IV.7.5.1} therefore proves the assertion. We may consequently suppose that $\mathrm{A}$ is integral and need only prove, for its field of fractions $\mathrm{K}$, that the local rings of $\Spec(\mathrm{B}\otimes_\mathrm{A}\mathrm{K})$ satisfy $\boldsymbol{R}$.

\emph{II) Reduction to the case where $\mathrm{A}$ and $\mathrm{B}$ are complete.} Let $\xi$ be the generic point of $\Spec(\mathrm{A})$ and let $y\in f^{-1}(\xi)$. We must prove $\boldsymbol{R}(\sh{O}_y)$. By $(\mathrm{R}_{\mathrm{II}})$, it is enough to find a point $z\in\Spec(\widehat{\mathrm{B}})$ over $y$ such that $\boldsymbol{R}(\sh{O}_z)$ is true. The hypothesis on $\Spec(\widehat{\mathrm{B}}\otimes_{\widehat{\mathrm{A}}}k)$ and \sref{IV.7.5.1} imply that the morphism
\[
  \widehat{f}:\Spec(\widehat{\mathrm{B}})\to\Spec(\widehat{\mathrm{A}})
\]
is a $\boldsymbol{P}$-morphism. For any point $z\in\Spec(\widehat{\mathrm{B}})$ above $y$---such a point exists because $\widehat{\mathrm{B}}$ is faithfully flat over $\mathrm{B}$---its image $x$ in $\Spec(\widehat{\mathrm{A}})$ lies in the fibre $h^{-1}(\xi)$ of the canonical morphism $h:\Spec(\widehat{\mathrm{A}})\to\Spec(\mathrm{A})$. By the hypothesis on $\mathrm{A}$, $\boldsymbol{R}'(\sh{O}_x)$ is true; hence $(\mathrm{R}_{\mathrm{V}})$ gives $\boldsymbol{R}(\sh{O}_z)$.
\end{proof}

\begin{examples}[7.5.3]
\label{IV.7.5.3}
Consider the following properties $\boldsymbol{R}(\mathrm{C})$, where $\mathrm{C}$ is a Noetherian local ring:
\begin{enumerate}[label=(\roman*)]
  \item also denoted $(\mathrm{i}_n)$, $\operatorname{coprof}(\mathrm{C})\leqslant n$;
  \oldpage[IV-2]{206}
  \item also denoted $(\mathrm{ii}_n)$, $\mathrm{C}$ satisfies $(\mathrm{S}_n)$;
  \item $\mathrm{C}$ is a Cohen--Macaulay ring;
  \item $\mathrm{C}$ is a regular ring;
  \item also denoted $(\mathrm{v}_n)$, $\mathrm{C}$ is integral, integrally closed, and satisfies $(\mathrm{R}_n)$;
  \item $\mathrm{C}$ is integral and integrally closed;
  \item $\mathrm{C}$ is integral;
  \item $\mathrm{C}$ is reduced.
\end{enumerate}
All these properties satisfy $(\mathrm{R}_{\mathrm{III}})$, as was seen in \sref{IV.7.3.19},~(i). They also satisfy $(\mathrm{R}_{\mathrm{IV}})$. For~(i), this follows from \sref[0]{0.16.4.10},~(ii); for~(ii) and~(iii), from \sref{IV.5.12.4} and the fact that a complete Noetherian local ring is catenary \sref{IV.5.6.4}; for~(iv), it is the special case \sref[0]{0.17.1.8}; for~(vii), it follows from \sref{IV.3.4.5}, and for~(viii), from \sref{IV.3.4.6}; for~(vi), from \sref{IV.5.12.7} and the catenarity just recalled; and finally, for~(v), from the preceding facts and \sref{IV.5.12.5}.

Theorem \sref{IV.7.5.1} is therefore applicable when $\boldsymbol{R}$ is any one of the above properties. We have also observed in \sref{IV.7.3.11} that properties~(i) to~(viii) satisfy $(\mathrm{R}_{\mathrm{II}})$, except that for~(vii) this follows from \sref{IV.2.1.14}. As for $(\mathrm{R}_{\mathrm{V}})$, if one takes $\boldsymbol{R}'=\boldsymbol{R}$, it reduces to condition $(\mathrm{R}'_{\mathrm{I}})$ of \sref{IV.7.3.11}; properties~(ii) to~(vi), and~(viii), satisfy this condition. For~(i), one may take for $\boldsymbol{R}'$ the property of being Cohen--Macaulay, by \sref{IV.6.3.2}. For~(vii), condition $(\mathrm{R}'_{\mathrm{I}})$ is not satisfied, as shown by \sref{IV.6.15.11},~(ii) or \sref{IV.6.5.5},~(ii). However, $(\mathrm{R}_{\mathrm{V}})$ is true if $\boldsymbol{R}'$ is the property of being regular: apply \sref{IV.7.5.1.1} with $\mathcal{C}$ the category of all Noetherian local rings and with $\boldsymbol{R}(\mathrm{C})$ the property of being integral; this is possible by \sref{IV.3.4.5}.

In particular, when $\boldsymbol{R}$ is any of properties~(i) to~(viii), the conclusion of \sref{IV.7.5.2} applies if the formal fibres of $\mathrm{A}$ are geometrically regular; see \sref{IV.7.8.2}.
\end{examples}

\begin{remarks}[7.5.4]
\label{IV.7.5.4}
\begin{enumerate}[label=(\roman*)]
  \item It would be interesting to know whether \sref{IV.7.5.1} remains true without the finiteness hypothesis~(i) on the residue field of $\mathrm{B}$. The answer is affirmative when the residue field $k$ of $\mathrm{A}$ has characteristic~$0$, by Hironaka's results on resolution of singularities; see \sref{IV.7.9.8}. A particular case of the question is the following. Let $\mathrm{A}$ and $\mathrm{B}$ be complete Noetherian rings, let $k$ be the residue field of $\mathrm{A}$, and let $\varphi:\mathrm{A}\to\mathrm{B}$ be a local homomorphism making $\mathrm{B}$ a formally smooth $\mathrm{A}$-algebra \sref[0]{0.19.3.1}; equivalently, $\mathrm{B}$ is flat over $\mathrm{A}$ and $\mathrm{B}\otimes_\mathrm{A}k$ is geometrically regular over $k$ \sref[0]{0.19.7.1}. Are the fibres of ${}^a\varphi:\Spec(\mathrm{B})\to\Spec(\mathrm{A})$ geometrically regular? This is so when the residue field $k'$ of $\mathrm{B}$ is a finite extension of $k$, by \sref{IV.7.5.1}; the answer is also affirmative when $k'$ is an extension of finite type of $k$. We do not know the answer when $\mathrm{B}\otimes_\mathrm{A}k$ is a field equal to a separable closure of $k$.

  \item One could state a result analogous to \sref{IV.7.5.1} for a finite-type $\mathrm{A}$-module $\mathrm{M}$ and a finite-type $\mathrm{B}$-module $\mathrm{N}$, concerning the properties of
  \[
    (\mathrm{M}\otimes_\mathrm{A}\mathrm{N})_y
    \otimes_{\sh{O}_{\mathrm{X},x}}\boldsymbol{k}(y),
  \]
  where $y$ runs through $\Spec(\mathrm{B})$ and $x=f(y)$, deduced from properties of the same type of $\mathrm{M}_x$ and $\mathrm{N}_y$.

  \item Suppose $\boldsymbol{R}$ satisfies $(\mathrm{R}'_{\mathrm{I}})$, $(\mathrm{R}_{\mathrm{II}})$, $(\mathrm{R}_{\mathrm{III}})$, and $(\mathrm{R}_{\mathrm{IV}})$, and hypotheses~(i) and~(ii) of \sref{IV.7.5.2} hold. Then $\boldsymbol{R}(\mathrm{A})$ and
  \[
    \boldsymbol{P}\bigl(\Spec(\widehat{\mathrm{B}}\otimes_{\widehat{\mathrm{A}}}k),k\bigr)
  \]
  imply $\boldsymbol{R}(\mathrm{B})$. By \sref{IV.7.5.3}, this holds for the properties listed there except~(i) and~(vii). For~(vii), it is nevertheless plausible that the following question has an affirmative answer. Let $\varphi:\mathrm{A}\to\mathrm{B}$ be a local homomorphism of Noetherian local rings making $\mathrm{B}$ a flat $\mathrm{A}$-module. Suppose $\mathrm{A}$ is complete, integral, and geometrically unibranch, and that $\Spec(\mathrm{B}\otimes_\mathrm{A}k)$, where $k$ is the residue field of $\mathrm{A}$, is geometrically locally integral. Is $\mathrm{B}$ then integral? This is known when $k$ has characteristic~$0$, by Hironaka's resolution of singularities \sref{IV.7.9.8}; it is also true when $\mathrm{B}$ is an $\mathrm{A}$-algebra essentially of finite type \sref{IV.1.3.8} (see \sref{IV.11.3.10} and \sref{IV.11.3.11}), or when $\Spec(\mathrm{B}\otimes_\mathrm{A}k)$ is geometrically normal, for then \sref{IV.7.5.3} shows that $\Spec(\mathrm{B})\to\Spec(\mathrm{A})$ is normal and one concludes by \sref{IV.6.15.10}. Even if the residue field of $\mathrm{B}$ equals that of $\mathrm{A}$ and $\mathrm{B}$ is complete, the general answer is unknown.

  \item Consider the property $\boldsymbol{R}(\mathrm{C})$: ``$\mathrm{C}$ is reduced, equidimensional, and satisfies $(\mathrm{R}_n)$''. It satisfies $(\mathrm{R}_{\mathrm{IV}})$ by \sref{IV.5.12.5}, but it is not known whether it satisfies $(\mathrm{R}'_{\mathrm{I}})$, $(\mathrm{R}_{\mathrm{II}})$, or $(\mathrm{R}_{\mathrm{III}})$; the difficulty lies in verifying equidimensionality.
\end{enumerate}
\end{remarks}

\oldpage[IV-2]{207}

In what follows of this section, we will apply \sref{IV.7.5.1} to the study of \emph{completed tensor products} $\mathrm{A}\widehat{\otimes}_k\mathrm{B}$ of local Noetherian rings which are algebras over a field $k$.

\begin{lemma}[7.5.5]
\label{IV.7.5.5}
Let $k$ be a field, $\mathrm{A}$, $\mathrm{B}$ two local Noetherian complete rings containing $k$, the residue field of $\mathrm{A}$ being a finite extension of $k$. Let $\mathrm{C}$ be the completed tensor product $\mathrm{A}\widehat{\otimes}_k\mathrm{B}$ \sref[0]{0.7.7.5}. Then:
\begin{enumerate}[label=(\roman*)]
  \item $\mathrm{C}$ is a semi-local Noetherian complete ring.
  \item $\mathrm{C}$ is a flat $\mathrm{A}$-module and a flat $\mathrm{B}$-module.
  \item If $\mathfrak{m}$ is the maximal ideal of $\mathrm{A}$, $\mathfrak{m}\mathrm{C}$ is contained in the radical of $\mathrm{C}$, and $\mathrm{C}/\mathfrak{m}\mathrm{C}$ is isomorphic to $(\mathrm{A}/\mathfrak{m})\otimes_k\mathrm{B}$.
  \item The residue fields of the components of $\mathrm{C}$ are finite extensions of the residue field of $\mathrm{B}$.
\end{enumerate}
The properties (i), (iii) and the first assertion of (ii) are particular cases of \sref[0]{0.19.7.1.2}. To prove that $\mathrm{C}$ is a flat $\mathrm{B}$-module, note that for every $h>0$,
\[
  \mathrm{C}/\mathfrak{m}^h\mathrm{C}
  =(\mathrm{A}/\mathfrak{m}^h)\otimes_k\mathrm{B}
\]
is a flat $\mathrm{B}$-module, since $k$ is a field; it therefore suffices to apply \sref[0]{0.10.2.6} to each local component of $\mathrm{C}$. Finally, if $\mathfrak{n}$ is the maximal ideal of $\mathrm{B}$, the residue fields of the local components of $\mathrm{C}$ are also those of the local components of the Artinian ring $(\mathrm{A}/\mathfrak{m})\otimes_k(\mathrm{B}/\mathfrak{n})$, and hence are finite extensions of $\mathrm{B}/\mathfrak{n}$.
\end{lemma}

\begin{proposition}[7.5.6]
\label{IV.7.5.6}
Let $k$ be a field, and let $\mathrm{A}$ and $\mathrm{B}$ be complete Noetherian local rings containing $k$, whose residue fields are finite extensions of $k$. Let $\boldsymbol{R}$ be a property satisfying $(\mathrm{R}'_{\mathrm{I}})$, $(\mathrm{R}_{\mathrm{III}})$ and $(\mathrm{R}_{\mathrm{IV}})$, with $\boldsymbol{P}$ defined from $\boldsymbol{R}$ as in \sref{IV.7.5.0}. Suppose that $\boldsymbol{R}(\mathrm{A})$ and $\boldsymbol{P}(\operatorname{Spec}(\mathrm{B}),k)$ hold; the latter means that, for every finite extension $k'$ of $k$ and each complete local ring $\mathrm{B}_h$ occurring as a component of $\mathrm{B}\otimes_k k'$, one has $\boldsymbol{R}(\mathrm{B}_h)$. Then $\boldsymbol{R}(\mathrm{C}_j)$ holds for every complete local component $\mathrm{C}_j$ of $\mathrm{A}\widehat{\otimes}_k\mathrm{B}$.
\end{proposition}

\begin{proof}
By \sref{IV.7.5.5},~(iii), each homomorphism $\mathrm{A}\to\mathrm{C}_j$ is local; by \sref{IV.7.5.5},~(ii), each $\mathrm{C}_j$ is flat over $\mathrm{A}$; and by \sref{IV.7.5.5},~(iv), its residue field is a finite extension of $\mathrm{A}/\mathfrak{m}$. Moreover,
\[
  \mathrm{C}/\mathfrak{m}\mathrm{C}
  =(\mathrm{A}/\mathfrak{m})\otimes_k\mathrm{B},
\]
so the definition of $\boldsymbol{P}$ and the fact that $\mathrm{A}/\mathfrak{m}$ is finite over $k$ imply
\[
  \boldsymbol{P}\bigl(\operatorname{Spec}(\mathrm{C}_j\otimes_\mathrm{A}\mathrm{A}/\mathfrak{m}),\mathrm{A}/\mathfrak{m}\bigr).
\]
Thus all the hypotheses of \sref{IV.7.5.1} hold for the local homomorphism $\mathrm{A}\to\mathrm{C}_j$, so $\operatorname{Spec}(\mathrm{C}_j)\to\operatorname{Spec}(\mathrm{A})$ is a $\boldsymbol{P}$-morphism. The conclusion follows from $(\mathrm{R}'_{\mathrm{I}})$ and $\boldsymbol{R}(\mathrm{A})$.
\end{proof}

\begin{corollary}[7.5.7]
\label{IV.7.5.7}
(Chevalley). --- Let $k$ be a perfect field (resp.\ algebraically closed), $\mathrm{A}$, $\mathrm{B}$ two local Noetherian complete rings containing $k$ and whose residue fields are finite extensions of $k$. Then, if $\mathrm{A}$ and $\mathrm{B}$ are reduced (resp.\ integral), the completed tensor product $\mathrm{A}\widehat{\otimes}_k\mathrm{B}$ is reduced (resp.\ integral).
\end{corollary}

\begin{proof}
I) Suppose $k$ is perfect and $\mathrm{A}$ and $\mathrm{B}$ are reduced. Let $\mathrm{A}_i$ (resp.\ $\mathrm{B}_j$) be the quotients of $\mathrm{A}$ (resp.\ $\mathrm{B}$) by its minimal prime ideals. Then $\mathrm{A}$ and $\mathrm{B}$ identify with subrings of $\prod_i\mathrm{A}_i$ and $\prod_j\mathrm{B}_j$. Since the tensor products are taken over a field, $\mathrm{A}\otimes_k\mathrm{B}$ identifies with a subring of $\prod_{i,j}(\mathrm{A}_i\otimes_k\mathrm{B}_j)$; the tensor-product topology is induced by the product of the corresponding topologies. Hence $\mathrm{A}\widehat{\otimes}_k\mathrm{B}$ embeds in $\prod_{i,j}(\mathrm{A}_i\widehat{\otimes}_k\mathrm{B}_j)$, and we are reduced to the case where $\mathrm{A}$ and $\mathrm{B}$ are integral.

Let $\mathrm{A}'$ and $\mathrm{B}'$ be their integral closures. By Nagata's finiteness theorem \sref[0]{0.23.1.6}, $\mathrm{A}'$ (resp.\ $\mathrm{B}'$) is a finite $\mathrm{A}$-module (resp.\ $\mathrm{B}$-module) and a complete local ring. It suffices to show that $\mathrm{A}\widehat{\otimes}_k\mathrm{B}$ embeds in $\mathrm{A}'\widehat{\otimes}_k\mathrm{B}$ and that the latter embeds in $\mathrm{A}'\widehat{\otimes}_k\mathrm{B}'$. This follows from the lemma below.

\begin{lemma}[7.5.7.1]
\label{IV.7.5.7.1}
Let $\mathrm{A}$, $\mathrm{B}$ be two local Noetherian complete rings containing a field $k$ and whose residue fields are finite extensions of $k$. Let $\mathfrak{m}$ be the maximal ideal of $\mathrm{A}$. For every $\mathrm{A}$-module $\mathrm{M}$ of finite type (equipped with the $\mathfrak{m}$-adic topology), the completed tensor product $\mathrm{M}\widehat{\otimes}_k\mathrm{B}$ \sref[0]{0.7.7.1} identifies canonically with $\mathrm{M}\otimes_\mathrm{A}(\mathrm{A}\widehat{\otimes}_k\mathrm{B})$.
\end{lemma}

\begin{proof}
Indeed, one has a canonical isomorphism $\mathrm{M}\otimes_k\mathrm{B}\xrightarrow{\sim}\mathrm{M}\otimes_\mathrm{A}(\mathrm{A}\otimes_k\mathrm{B})$, hence a canonical composite homomorphism
\[
  \varphi:\mathrm{M}\otimes_k\mathrm{B}
  \longrightarrow\mathrm{M}\otimes_\mathrm{A}(\mathrm{A}\otimes_k\mathrm{B})
  \longrightarrow\mathrm{M}\otimes_\mathrm{A}(\mathrm{A}\widehat{\otimes}_k\mathrm{B})
\]
and it is immediate that this homomorphism is continuous for the tensor product topologies; furthermore, $\mathrm{M}\otimes_\mathrm{A}(\mathrm{A}\widehat{\otimes}_k\mathrm{B})$ is separated and complete ($\sref{IV.7.5.5}$ and $\sref[0]{0.7.7.8}$), hence one also has by completion a continuous homomorphism
\[
  \widehat{\varphi}:\mathrm{M}\widehat{\otimes}_k\mathrm{B}\to\mathrm{M}\otimes_\mathrm{A}(\mathrm{A}\widehat{\otimes}_k\mathrm{B})
\]

\oldpage[IV-2]{208}

It is immediate that this homomorphism is bijective when $\mathrm{M}$ is free of finite type; as one has an exact sequence $\mathrm{L}'\to\mathrm{L}\to\mathrm{M}\to 0$, where $\mathrm{L}$ and $\mathrm{L}'$ are free $\mathrm{A}$-modules of finite type, one deduces a commutative diagram
\[
  \xymatrix{
    \mathrm{L}'\widehat{\otimes}_k\mathrm{B} \ar[r] \ar[d] & \mathrm{L}\widehat{\otimes}_k\mathrm{B} \ar[r] \ar[d] & \mathrm{M}\widehat{\otimes}_k\mathrm{B} \ar[r] \ar[d] & 0 \\
    \mathrm{L}'\otimes_\mathrm{A}(\mathrm{A}\widehat{\otimes}_k\mathrm{B}) \ar[r] & \mathrm{L}\otimes_\mathrm{A}(\mathrm{A}\widehat{\otimes}_k\mathrm{B}) \ar[r] & \mathrm{M}\otimes_\mathrm{A}(\mathrm{A}\widehat{\otimes}_k\mathrm{B}) \ar[r] & 0
  }
\]
where the rows are exact (the first by virtue of the definition of the completed tensor product and of $\sref[0]{0.13.2.2}$); the first two vertical arrows being bijections, it follows that so is the third.
\end{proof}

The fact that $\mathrm{A}\widehat{\otimes}_k\mathrm{B}$ identifies with a subring of $\mathrm{A}'\otimes_\mathrm{A}(\mathrm{A}\widehat{\otimes}_k\mathrm{B})$ follows from the inclusion $\mathrm{A}\subset\mathrm{A}'$ and from the fact that $\mathrm{A}\widehat{\otimes}_k\mathrm{B}$ is a flat $\mathrm{A}$-module \sref{IV.7.5.5}. One can thus further suppose that $\mathrm{A}$ and $\mathrm{B}$ are \emph{integrally closed}; as moreover $k$ is supposed \emph{perfect}, $\operatorname{Spec}(\mathrm{A})$ and $\operatorname{Spec}(\mathrm{B})$ are \emph{geometrically normal} over $k$, by virtue of \sref{IV.6.7.7},~$b)$); one can then apply \sref{IV.7.5.6} taking for $\boldsymbol{R}$ the property (vi) of \sref{IV.7.5.3}.

II) Suppose $k$ is algebraically closed and $\mathrm{A}$ and $\mathrm{B}$ are integral. The argument in I) again reduces us to the case where $\mathrm{A}$ and $\mathrm{B}$ are integrally closed. By \sref{IV.7.5.6}, $\operatorname{Spec}(\mathrm{A}\widehat{\otimes}_k\mathrm{B})$ is normal, so $\mathrm{A}\widehat{\otimes}_k\mathrm{B}$ is a finite product of integrally closed complete local rings. It remains only to show that it is local. If $\mathfrak{m}$ and $\mathfrak{n}$ are the maximal ideals of $\mathrm{A}$ and $\mathrm{B}$, it suffices to show that $(\mathrm{A}/\mathfrak{m})\otimes_k(\mathrm{B}/\mathfrak{n})$ is local. Since both residue fields are finite extensions of the algebraically closed field $k$, they are equal to $k$, and the conclusion follows.
\end{proof}

\begin{remark}[7.5.8]
\label{IV.7.5.8}
It would be interesting to determine whether, in \sref{IV.7.5.7}, one can replace the hypothesis that $k$ is perfect or algebraically closed by the hypothesis that $\operatorname{Spec}(\mathrm{A})$ or $\operatorname{Spec}(\mathrm{B})$ is \emph{geometrically integral} over $k$, or at least by the hypothesis that $\mathrm{A}$ (for example) is integral and contains a subring $\mathrm{A}_0$ isomorphic to a formal power series ring $k[\![\mathrm{T}_1,\ldots,\mathrm{T}_n]\!]$, such that the field of fractions $\mathrm{K}$ of $\mathrm{A}$ is \emph{separable} over the field of fractions $\mathrm{K}_0$ of $\mathrm{A}_0$ (one can show that this condition entails that $\mathrm{A}$ is geometrically integral over $k$, and that the two properties are equivalent when $[k:k^p]$ is finite ($p$ being the characteristic exponent of $k$)). Likewise, it would be desirable to develop variants of \sref{IV.7.5.6} and \sref{IV.7.5.7} where one would weaken the hypothesis of finiteness of the residue fields of $\mathrm{A}$ and $\mathrm{B}$, by supposing for example that only one of them is a finite extension of $k$, the other being arbitrary.
\end{remark}


\subsection{Applications: I. Japanese local rings}
\label{subsection:IV.7.6}

\begin{proposition}[7.6.1]
\label{IV.7.6.1}
Let $\mathrm{A}$ be a reduced local Noetherian ring whose formal fibres are geometrically normal. Then the completion $\widehat{\mathrm{A}}$ is reduced, the integral closure $\mathrm{A}'$ of $\mathrm{A}$ in its total ring of fractions is a finite $\mathrm{A}$-algebra (hence a semi-local Noetherian ring), and its completion $\widehat{\mathrm{A}'}$ is isomorphic to the integral closure of $\widehat{\mathrm{A}}$ in its total ring of fractions.
\end{proposition}

\begin{proof}
The formal fibres of $\mathrm{A}$ are \emph{a fortiori} geometrically reduced, so the reducedness of $\mathrm{A}$ implies that of $\widehat{\mathrm{A}}$ by \sref{IV.7.3.17}. Let $\mathfrak{p}_i$ be the minimal prime ideals of $\mathrm{A}$ and put $\mathrm{B}_i=\mathrm{A}/\mathfrak{p}_i$. The formal fibres of the local rings $\mathrm{B}_i$ are geometrically normal \sref{IV.7.3.15}; hence $\widehat{\mathrm{B}}_i$ is reduced, and \sref[0]{0.23.1.7},~(i), shows that the integral closure $\mathrm{B}'_i$ of $\mathrm{B}_i$ in its field of fractions is finite over $\mathrm{B}_i$, and therefore over $\mathrm{A}$. Since $\mathrm{A}'=\prod_i\mathrm{B}'_i$ (\textbf{II}, $6.3.8$), $\mathrm{A}'$ is a finite $\mathrm{A}$-module.

Let $\mathfrak{m}_j$ be the maximal ideals of the semi-local ring $\mathrm{A}'$. Its completion $\widehat{\mathrm{A}'}$ identifies with the product of the completions $\widehat{\mathrm{A}'_{\mathfrak{m}_j}}$ (Bourbaki, \emph{Alg.\ comm.}, chap.~III, \textsection\,2, n\textsuperscript{o}\,13, cor.~of prop.~19).

\oldpage[IV-2]{209}

By the hypothesis and \sref{IV.7.3.15}, the formal fibres of each $\mathrm{A}'_{\mathfrak{m}_j}$ are geometrically normal. Since $\mathrm{A}'$ is normal, so is each $\mathrm{A}'_{\mathfrak{m}_j}$, and \sref{IV.7.3.17} then shows that each $\widehat{\mathrm{A}'_{\mathfrak{m}_j}}$, and hence $\widehat{\mathrm{A}'}$, is normal.

Since $\mathrm{A}'$ is finite over $\mathrm{A}$, one has $\widehat{\mathrm{A}'}=\mathrm{A}'\otimes_\mathrm{A}\widehat{\mathrm{A}}$ (Bourbaki, \emph{Alg.\ comm.}, chap.~IV, \textsection\,2, n\textsuperscript{o}\,5, cor.~3 of prop.~9, and chap.~III, \textsection\,3, n\textsuperscript{o}\,4, th.~3). If $\mathrm{R}$ denotes the total ring of fractions of $\mathrm{A}$, then
\[
  \widehat{\mathrm{A}}\subset\widehat{\mathrm{A}'}
  \subset \mathrm{R}'=\mathrm{R}\otimes_\mathrm{A}\widehat{\mathrm{A}}.
\]
Flatness of $\widehat{\mathrm{A}}$ over $\mathrm{A}$ implies that every regular element of $\mathrm{A}$ remains regular in $\widehat{\mathrm{A}}$; hence $\mathrm{R}'$ embeds canonically in the total ring of fractions $\mathrm{R}''$ of $\widehat{\mathrm{A}}$ (Bourbaki, \emph{Alg.\ comm.}, chap.~II, \textsection\,2, n\textsuperscript{o}\,1, Remark~7). Since $\widehat{\mathrm{A}'}$ is normal and finite over $\widehat{\mathrm{A}}$, it is precisely the integral closure of $\widehat{\mathrm{A}}$ in $\mathrm{R}''$.
\end{proof}

\begin{corollary}[7.6.2]
\label{IV.7.6.2}
Under the hypotheses of \sref{IV.7.6.1}, there is a bijective correspondence between the maximal ideals $\mathfrak{m}_j$ of $\mathrm{A}'$ (in other words, the set of closed points of $\operatorname{Spec}(\mathrm{A}')$ above the closed point of $\mathrm{A}$) and the set of minimal prime ideals $\mathfrak{q}_j$ of $\widehat{\mathrm{A}}$ (in other words, the set of maximal points of $\operatorname{Spec}(\widehat{\mathrm{A}})$); in this correspondence, the completion $\widehat{\mathrm{A}'_{\mathfrak{m}_j}}$ of $\mathrm{A}'_{\mathfrak{m}_j}$ is isomorphic to the integral closure of $\widehat{\mathrm{A}}/\mathfrak{q}_j$.
\end{corollary}

\begin{proof}
  Indeed, one knows that the integral closure of $\widehat{\mathrm{A}}$ in its total ring of fractions is the direct product of the integral closures of the $\widehat{\mathrm{A}}/\mathfrak{q}_j$, which are complete local rings \sref[0]{0.23.1.6}.
\end{proof}

\begin{corollary}[7.6.3]
\label{IV.7.6.3}
Under the hypotheses of \sref{IV.7.6.1}, for $\widehat{\mathrm{A}}$ to be integral, it is necessary and sufficient that $\mathrm{A}$ be unibranch; for $\widehat{\mathrm{A}}$ to be geometrically unibranch, it is necessary and sufficient that $\mathrm{A}$ be geometrically unibranch.
\end{corollary}

\begin{proof}
This is a particular case of \sref{IV.7.6.2}.
\end{proof}

\begin{theorem}[7.6.4]
\label{IV.7.6.4}
(Zariski-Nagata). --- Let $\mathrm{A}$ be a semi-local Noetherian ring. The following conditions are equivalent:
\begin{enumerate}[label=\emph{\alph*)}]
  \item For every finite reduced $\mathrm{A}$-algebra $\mathrm{C}$, the completion $\widehat{\mathrm{C}}$ is a reduced ring.
  \item[\rm a$'$)] For every integral quotient ring $\mathrm{B}$ of $\mathrm{A}$, with field of fractions $\mathrm{K}$, the completion $\widehat{\mathrm{B}}$ is reduced, and the fields $\mathrm{L}_i$ of the total ring of fractions of $\widehat{\mathrm{B}}$ are separable extensions of $\mathrm{K}$.
  \item[\rm a$''$)] The formal fibres are geometrically reduced (in other words, for every integral quotient ring $\mathrm{B}$ of $\mathrm{A}$, with field of fractions $\mathrm{K}$, the $\mathrm{K}$-algebra $\widehat{\mathrm{B}}\otimes_\mathrm{B}\mathrm{K}$ is separable \sref{IV.4.6.2}).
  \item[\rm b)] Every integral quotient ring $\mathrm{B}$ of $\mathrm{A}$ is a Japanese ring.
\end{enumerate}
\end{theorem}

\begin{proof}
\renewcommand{\qedsymbol}{}
To show that $a)$ implies $a')$, it suffices to verify that the $\mathrm{L}_i$ are separable extensions of $\mathrm{K}$, or equivalently \sref{IV.4.6.1} that for every finite extension $\mathrm{K}'$ of $\mathrm{K}$, the ring $\widehat{\mathrm{B}}\otimes_\mathrm{B}\mathrm{K}'$ is reduced. The field $\mathrm{K}'$ is generated by finitely many elements integral over $\mathrm{B}$; these generate a finite $\mathrm{B}$-subalgebra $\mathrm{B}'$ of $\mathrm{K}'$ whose field of fractions is $\mathrm{K}'$. One has
\[
  \widehat{\mathrm{B}'}=\widehat{\mathrm{B}}\otimes_\mathrm{B}\mathrm{B}'
\]
\sref[0]{0.7.3.3} and Bourbaki, \emph{Alg.\ comm.}, chap.~IV, \textsection\,2, n\textsuperscript{o}\,5, cor.~3 of prop.~9.

\oldpage[IV-2]{210}

Therefore, by associativity of the tensor product,
\[
  \widehat{\mathrm{B}}\otimes_\mathrm{B}\mathrm{K}'
  =\widehat{\mathrm{B}'}\otimes_{\mathrm{B}'}\mathrm{K}'.
\]
Now $\mathrm{B}'$ is a finite integral $\mathrm{A}$-algebra, so $\widehat{\mathrm{B}'}$ is reduced by $a)$. Since $\widehat{\mathrm{B}'}$ is flat over $\mathrm{B}'$, the nonzero elements of $\mathrm{B}'$ are $\widehat{\mathrm{B}'}$-regular; hence $\widehat{\mathrm{B}'}\otimes_{\mathrm{B}'}\mathrm{K}'$ identifies with a subring of the total ring of fractions of $\widehat{\mathrm{B}'}$, and is therefore reduced.

To see that $a')$ implies $a'')$, consider any point $x$ of $\mathrm{X}=\operatorname{Spec}(\mathrm{A})$, and let $\mathrm{Y}$ be the integral closed subprescheme of $\mathrm{X}$ having $\overline{\{x\}}$ as its underlying space. Then $\mathrm{Y}=\operatorname{Spec}(\mathrm{B})$, where $\mathrm{B}$ is an integral quotient ring of $\mathrm{A}$, and $\operatorname{Spec}(\widehat{\mathrm{B}})=\mathrm{Y}\times_\mathrm{X}\operatorname{Spec}(\widehat{\mathrm{A}})$. Thus the formal fibre of $\mathrm{A}$ at $x$ is the same as that of $\mathrm{B}$ at its generic point, namely $\operatorname{Spec}(\widehat{\mathrm{B}}\otimes_\mathrm{B}\mathrm{K})$. Its local rings are those of $\operatorname{Spec}(\widehat{\mathrm{B}})$ at the points of this fibre, so the hypothesis implies that $\widehat{\mathrm{B}}\otimes_\mathrm{B}\mathrm{K}$ is reduced; the remaining condition in $a')$ then says that it is a separable $\mathrm{K}$-algebra \sref{IV.4.6.1}.

The condition $a'')$ implies $a)$, for it follows from \sref{IV.7.3.15} that the formal fibres of $\mathrm{C}$ are geometrically reduced, and if $\mathrm{C}$ is reduced, it follows hence that so is $\widehat{\mathrm{C}}$ (particular case of \sref{IV.7.3.17}).

The fact that $a')$ implies $b)$ is a particular case of \sref[0]{0.23.1.7}. It remains to prove that $b)$ implies $a)$. If $\mathrm{C}$ is a finite $\mathrm{A}$-algebra and $\mathfrak{q}$ is a prime ideal of $\mathrm{C}$, then the inverse image $\mathfrak{p}$ of $\mathfrak{q}$ in $\mathrm{A}$ is prime and $\mathrm{C}/\mathfrak{q}$ is a finite $(\mathrm{A}/\mathfrak{p})$-algebra. Thus $b)$ implies that every integral quotient ring of $\mathrm{C}$ is Japanese. We are therefore reduced to proving that, under $b)$, \emph{the completion of every integral quotient ring of $\mathrm{A}$ is reduced}.

We argue by induction on $n=\dim(\mathrm{A})$, the assertion being trivial for $n=0$. Replacing $\mathrm{A}$ by the quotients by its minimal prime ideals $\mathfrak{p}_i$, we may restrict to the case where $\mathrm{A}$ is integral. For every nonzero prime ideal $\mathfrak{p}$, the induction hypothesis already shows that the completion of $\mathrm{A}/\mathfrak{p}$ is reduced; it therefore suffices to prove that $\widehat{\mathrm{A}}$ is reduced.

Moreover, the integral closure $\mathrm{A}'$ of $\mathrm{A}$ is, by hypothesis, a finite $\mathrm{A}$-module, hence a semi-local Noetherian ring, and $\widehat{\mathrm{A}}$ identifies with a subring of $\widehat{\mathrm{A}'}$ \sref[0]{0.7.3.3}. It is therefore enough to prove that $\widehat{\mathrm{A}'}$ is reduced. As observed above, $b)$ also holds for $\mathrm{A}'$, and $\dim(\mathrm{A}')=n$ \sref[0]{0.16.1.5}; hence we may suppose that $\mathrm{A}$ is \emph{integrally closed}. Let $t\neq 0$ belong to the Jacobson radical of $\mathrm{A}$, and let $\mathfrak{q}_j$ $(1\leqslant j\leqslant r)$ be the prime ideals minimal among those containing $t\mathrm{A}$; then:
\begin{enumerate}[label=(\roman*)]
  \item $t$ is regular.
  \item $\mathrm{A}/t\mathrm{A}$ has no immersed associated prime ideals.
  \item The $\mathrm{A}_{\mathfrak{q}_j}$ are discrete valuation rings.
  \item The completions of the $\mathrm{A}/\mathfrak{q}_j$ are reduced.
\end{enumerate}
Indeed, (i) is trivial since $\mathrm{A}$ is integral and $t\neq 0$. As $\mathrm{A}$ is integrally closed, $\mathrm{A}/t\mathrm{A}$ satisfies $(\mathrm{S}_1)$, that is to say \sref{IV.5.7.7} is without immersed associated prime ideals (Bourbaki, \emph{Alg.\ comm.}, chap.~VII, \S~1, no~4, prop.~8). Since $\mathrm{A}$ is integrally

\oldpage[IV-2]{211}

closed, the $\mathrm{A}_{\mathfrak{q}_j}$ are, and one knows (\emph{loc.\ cit.}) that these rings are of dimension~1, hence they are discrete valuation rings, whence (iii). Finally, (iv) follows from the hypothesis of induction. The proof will be completed if we prove the

\begin{lemma}[7.6.4.1]
\label{IV.7.6.4.1}
(Zariski). --- Let $\mathrm{A}$ be a semi-local Noetherian ring, $t$ an element of its radical satisfying the conditions (i) to (iv) above; then $\widehat{\mathrm{A}}$ is reduced.
\end{lemma}

\begin{proof}
\renewcommand{\qedsymbol}{\openbox}
Set for simplicity $\mathrm{A}'=\widehat{\mathrm{A}}$, and consider $t$ as an element of $\mathrm{A}'$; one has $\mathrm{A}'/t\mathrm{A}'=(\mathrm{A}/t\mathrm{A})\otimes_\mathrm{A}\mathrm{A}'$, and as $\mathrm{A}'$ is a flat $\mathrm{A}$-module, it follows from \sref{IV.3.3.1} that the prime ideals of $\mathrm{A}'$ associated to the $\mathrm{A}'$-module $\mathrm{A}'/t\mathrm{A}'$ are the prime ideals $\mathfrak{q}'_{jh}$ such that $\mathfrak{q}'_{jh}$ lies above $\mathfrak{q}_j$ and is associated to the $k_j$-module $(\mathrm{A}'/t\mathrm{A}')\otimes_\mathrm{A} k_j$, where $k_j$ is the field of fractions of $\mathrm{A}/\mathfrak{q}_j$. Now, $\operatorname{Spec}((\mathrm{A}'/t\mathrm{A}')\otimes_\mathrm{A} k_j)$ is the formal fibre of $\mathrm{A}$ at the point $\mathfrak{q}_j$, or equivalently that of $\mathrm{A}/\mathfrak{q}_j$ at its generic point. By (iv), the completion $\mathrm{A}'/\mathfrak{q}_j\mathrm{A}'$ of $\mathrm{A}/\mathfrak{q}_j$ is reduced; hence so is its formal fibre at the generic point. This formal fibre therefore has no embedded associated prime, and its local rings at the generic points $\mathfrak{q}'_{jh}$ of its irreducible components are fields \sref{IV.3.2.1}. It follows from \sref{IV.3.3.3} and (ii) that $\mathfrak{q}_j\mathrm{A}'_{\mathfrak{q}'_{jh}}$ is the maximal ideal of $\mathrm{A}'_{\mathfrak{q}'_{jh}}$ for every $h$. Since $\mathrm{A}_{\mathfrak{q}_j}$ is a discrete valuation ring, its maximal ideal is principal; therefore the maximal ideal of the Noetherian local ring $\mathrm{A}'_{\mathfrak{q}'_{jh}}$ is principal, and this ring is a discrete valuation ring (Bourbaki, \emph{Alg.\ comm.}, chap.~VI, \textsection\,3, n\textsuperscript{o}\,6, prop.~9). We have thus verified the relevant hypotheses for the \emph{complete semi-local ring} $\mathrm{A}'$. It therefore suffices to show that if $\mathrm{A}$ is complete and satisfies (i) to (iii), then $\mathrm{A}$ is reduced.

The hypotheses (i) and (ii) imply that, if $\varphi:\mathrm{A}\to\prod_{j=1}^r\mathrm{A}_{\mathfrak{q}_j}$ is the canonical homomorphism, then for every $m\geqslant 0$ one has $t^m\mathrm{A}=\varphi^{-1}(\varphi(t^m\mathrm{A}))$ \sref{IV.3.4.9}. Since the image of $t\mathrm{A}$ in each discrete valuation ring $\mathrm{A}_{\mathfrak{q}_j}$ is a power of its maximal ideal, the $t\mathrm{A}$-adic topology on $\mathrm{A}$ is the inverse image under $\varphi$ of the product topology on $\prod_{j=1}^r\mathrm{A}_{\mathfrak{q}_j}$. As $t$ belongs to the Jacobson radical of $\mathrm{A}$, the $t\mathrm{A}$-adic topology is separated, hence $\varphi$ is injective. By (iii), the rings $\mathrm{A}_{\mathfrak{q}_j}$ are integral; their product is reduced, and therefore so is $\mathrm{A}$.
\end{proof}
\end{proof}

\begin{corollary}[7.6.5]
\label{IV.7.6.5}
Let $\mathrm{A}$ be a semi-local Noetherian ring satisfying the equivalent conditions of \sref{IV.7.6.4}; every semi-local $\mathrm{A}$-algebra essentially of finite type also satisfies the conditions of \sref{IV.7.6.4}.
\end{corollary}

\begin{proof}
The condition $a'')$ of \sref{IV.7.6.4} means indeed that $\mathrm{A}$ is a $\boldsymbol{P}$-ring, where $\boldsymbol{P}(\mathrm{Z},k)$ is the property ``$\mathrm{Z}$ is geometrically reduced over $k$''; the corollary is then a consequence of the general theorem \sref{IV.7.4.4}.
\end{proof}

\begin{corollary}[7.6.6]
\label{IV.7.6.6}
Let $\mathrm{A}$ be a semi-local Noetherian integral ring of dimension $1$, $\mathrm{K}$ its field of fractions. For $\mathrm{A}$ to be a Japanese ring, it is necessary and sufficient that the completion $\widehat{\mathrm{A}}$

\oldpage[IV-2]{212}

be reduced and that, if $\mathfrak{q}_j$ $(1\leqslant j\leqslant n)$ are the minimal prime ideals of $\widehat{\mathrm{A}}$, the fields of fractions of the integral rings $\widehat{\mathrm{A}}/\mathfrak{q}_j$ be separable extensions of $\mathrm{K}$.
\end{corollary}

\begin{proof}
The integral quotients of $\mathrm{A}$ are $\mathrm{A}$ itself and fields, hence the condition $b)$ of \sref{IV.7.6.4} is equivalent to the hypothesis that $\mathrm{A}$ is a Japanese ring, and the hypothesis $a')$ to the conditions on $\widehat{\mathrm{A}}$, whence the conclusion.
\end{proof}

\begin{remarks}[7.6.7]
\label{IV.7.6.7}
\begin{enumerate}[label=\emph{\roman*)}]
  \item The equivalent conditions of \sref{IV.7.6.6} mean again that the formal fibres of $\mathrm{A}$ are geometrically regular \sref{IV.7.3.19},~(iv)). One has already observed (\emph{loc.\ cit.}) that this condition is satisfied when $\mathrm{A}$ is a \emph{discrete valuation ring whose field of fractions is of characteristic} $0$.

  \item The same reasoning as in the first part of the proof of \sref{IV.7.6.4} shows that for a semi-local Noetherian ring $\mathrm{A}$, the two following conditions are equivalent:
\begin{enumerate}[label=\emph{\alph*)}]
  \item For every integral quotient ring $\mathrm{B}$ of $\mathrm{A}$, the completion $\widehat{\mathrm{B}}$ is reduced.
  \item[\rm a$'$)] The formal fibres of $\mathrm{A}$ are reduced.
\end{enumerate}
Taking into account \sref[0]{0.23.1.7},~(i), these two conditions imply the following:
\begin{enumerate}[label=\emph{\alph*)}]
  \item[b)] For every integral quotient ring $\mathrm{B}$ of $\mathrm{A}$, the integral closure of $\mathrm{B}$ is a finite $\mathrm{B}$-algebra.
\end{enumerate}
When $\mathrm{A}$ is a universally catenary ring, one can prove that the condition $b)$ conversely implies $a)$; we will not need to use this result.
\end{enumerate}
\end{remarks}


\subsection{Applications: II. Universally Japanese rings}
\label{subsection:IV.7.7}

\begin{env}[7.7.1]
\label{IV.7.7.1}
Recall \sref[0]{0.23.1.1} that one calls a ring \emph{universally Japanese} a ring $\mathrm{A}$ such that every integral $\mathrm{A}$-algebra \emph{of finite type} is a Japanese ring. This amounts to saying that for every integral $\mathrm{A}$-algebra $\mathrm{B}$ of finite type, the integral closure of $\mathrm{B}$ is a finite $\mathrm{B}$-algebra.
\end{env}

\begin{theorem}[7.7.2]
\label{IV.7.7.2}
(Nagata). --- Let $\mathrm{A}$ be a Noetherian ring. The following conditions are equivalent:
\begin{enumerate}[label=\emph{\alph*)}]
  \item $\mathrm{A}$ is a universally Japanese ring.
  \item Every integral quotient ring of $\mathrm{A}$ is a Japanese ring.
  \item For every maximal ideal $\mathfrak{m}$ of $\mathrm{A}$, every integral quotient ring of $\mathrm{A}_\mathfrak{m}$ is a Japanese ring, and for every integral quotient ring $\mathrm{B}$ of $\mathrm{A}$, with field of fractions $\mathrm{K}$, every finite radicial extension $\mathrm{K}'$ of $\mathrm{K}$ and every sub-$\mathrm{B}$-algebra finite $\mathrm{B}'$ of $\mathrm{K}'$, with field of fractions $\mathrm{K}'$, there exists $f'\neq 0$ in $\mathrm{B}'$ such that $\mathrm{B}'_{f'}$ is integrally closed (cf.\ \sref{IV.6.13.7}).
\end{enumerate}
Moreover, if $\mathrm{A}$ satisfies these conditions, so does every ring of fractions $\mathrm{S}^{-1}\mathrm{A}$ of $\mathrm{A}$ and every $\mathrm{A}$-algebra of finite type.
\end{theorem}

\begin{proof}
Let us first show that $b)$ implies $c)$; every integral quotient ring of $\mathrm{A}_\mathfrak{m}$ is a ring of fractions of an integral quotient ring of $\mathrm{A}$, hence a Japanese ring \sref[0]{0.23.1.1}. On the other hand, every quotient $\mathrm{B}$ of $\mathrm{A}$ by one of its prime ideals is a Japanese ring, and hence so is $\mathrm{B}'$ \sref[0]{0.23.1.1}. One deduces that the set $\operatorname{Nor}(\operatorname{Spec}(\mathrm{B}'))$ is open \sref{IV.6.13.3}, hence contains a non-empty open set $\mathrm{D}(f')=\operatorname{Spec}(\mathrm{B}'_{f'})$.

\oldpage[IV-2]{213}

Let us show secondly that, if $\mathrm{A}$ satisfies $c)$, so does every $\mathrm{A}$-algebra of finite type $\mathrm{B}$. In the first instance, for every prime ideal $\mathfrak{q}$ of $\mathrm{B}$, $\mathrm{B}_\mathfrak{q}$ is a local ring in $\mathrm{S}^{-1}\mathrm{B}$, where $\mathrm{S}=\mathrm{A}-\mathfrak{p}$, $\mathfrak{p}$ being the inverse image of $\mathfrak{q}$ in $\mathrm{A}$; as by hypothesis every integral quotient ring of $\mathrm{A}_\mathfrak{p}$ is a Japanese ring, so is every integral quotient ring of $\mathrm{B}_\mathfrak{q}$, since $\mathrm{S}^{-1}\mathrm{B}$ is an $\mathrm{A}_\mathfrak{p}$-algebra of finite type \sref{IV.7.6.5}. On the other hand, if $\mathrm{C}$ is an integral quotient ring of $\mathrm{B}$, with field of fractions $\mathrm{K}$, $\mathrm{K}'$ a finite radicial extension of $\mathrm{K}$ and $\mathrm{C}'$ a sub-$\mathrm{C}$-algebra finite of $\mathrm{K}'$, with field of fractions $\mathrm{K}'$, $\mathrm{C}'$ is an $\mathrm{A}$-algebra integral of finite type, and in virtue of \sref{IV.6.13.7}, the hypothesis $c)$ on $\mathrm{A}$ entails that the spectrum of $\mathrm{C}'$ contains a non-empty open part whose points are all normal; in other words, there is a $g'\neq 0$ in $\mathrm{C}'$ such that $\mathrm{C}'_{g'}$ is integrally closed, which proves our assertion.

This result shows that, to prove the equivalence of $a)$, $b)$ and $c)$, it suffices to show that $c)$ implies $b)$, and even that for every integral quotient ring $\mathrm{B}$ of $\mathrm{A}$, with field of fractions $\mathrm{K}$, the integral closure $\mathrm{B}'$ of $\mathrm{B}$ is a $\mathrm{B}$-module of finite type. Now, condition $c)$ implies that there exists $f\neq0$ in $\mathrm{B}$ such that $\mathrm{B}_f$ is integrally closed, and thus condition~(i) of \sref{IV.6.13.6} has been verified for $\mathrm{B}$. Condition~(ii) of \sref{IV.6.13.6} is also verified for $\mathrm{B}$: for every maximal ideal $\mathfrak{q}$ of $\mathrm{B}$, $\mathrm{B}_\mathfrak{q}$ is a quotient of a local ring $\mathrm{A}_\mathfrak{p}$, where $\mathfrak{p}$ is maximal in $\mathrm{A}$; since every integral quotient ring of $\mathrm{A}_\mathfrak{p}$ is a Japanese ring by hypothesis, so is $\mathrm{B}_\mathfrak{q}$. Moreover, for every prime ideal $\mathfrak{r}$ of $\mathrm{B}$, $\mathrm{B}_\mathfrak{r}$ is a ring of fractions of some $\mathrm{B}_\mathfrak{q}$, where $\mathfrak{q}$ is a maximal ideal of $\mathrm{B}$, and is therefore again a Japanese ring. This proves the assertion.

We have also seen above that if $\mathrm{A}$ satisfies the equivalent conditions $a)$, $b)$ and $c)$, the same holds for every $\mathrm{A}$-algebra of finite type. Moreover, if $\mathrm{A}$ satisfies $b)$, so does every ring of fractions $\mathrm{S}^{-1}\mathrm{A}$: every integral quotient ring of $\mathrm{S}^{-1}\mathrm{A}$ is of the form $\mathrm{S}^{-1}\mathrm{A}/\mathrm{S}^{-1}\mathfrak{p}$, where $\mathfrak{p}$ is a prime ideal of $\mathrm{A}$; since $\mathrm{A}/\mathfrak{p}$ is by hypothesis a Japanese ring, so is $\mathrm{S}^{-1}(\mathrm{A}/\mathfrak{p})$ \sref[0]{0.23.1.1}. This completes the proof of the last assertion of the statement.
\end{proof}

\begin{corollary}[7.7.3]
\label{IV.7.7.3}
Let $\mathrm{A}$ be a Noetherian ring satisfying the following conditions:
\begin{enumerate}[label=(\roman*)]
  \item For every maximal ideal $\mathfrak{m}$ of $\mathrm{A}$, the formal fibres of $\mathrm{A}_\mathfrak{m}$ are geometrically reduced.
  \item The equivalent conditions of \sref{IV.6.12.4} are satisfied.
\end{enumerate}
Then $\mathrm{A}$ is universally Japanese.
\end{corollary}

\begin{proof}
Indeed, it follows then from \sref{IV.7.6.4} that every integral quotient ring of $\mathrm{A}_\mathfrak{m}$ is a Japanese ring; whence the conclusion, by virtue of \sref{IV.7.7.2}.
\end{proof}

\begin{corollary}[7.7.4]
\label{IV.7.7.4}
A Dedekind ring whose field of fractions is of characteristic $0$ (in particular $\mathbf{Z}$) is a universally Japanese ring. A local Noetherian ring whose formal fibres are geometrically reduced (in particular a local Noetherian complete ring) is universally Japanese.
\end{corollary}

\begin{proof}
The first assertion follows from \sref{IV.7.7.3}, \sref{IV.7.6.7},~(i) and \sref{IV.6.12.6}. The second follows from \sref{IV.7.6.4} and \sref{IV.7.7.2}.
\end{proof}


\subsection{Excellent rings}
\label{subsection:IV.7.8}

\begin{env}[7.8.1]
\label{IV.7.8.1}
The previous paragraphs of this section (as well as \sref{IV.6.11}, \sref{IV.6.12} and \sref{IV.6.13}) have been devoted to the study of problems frequently encountered in the use of rings and \emph{Noetherian} preschemes, and which can be grouped into the following types:

\oldpage[IV-2]{214}

A) For a \emph{local} Noetherian ring $\mathrm{A}$, do the properties of $\mathrm{A}$ ``transfer'' to its \emph{completion} $\widehat{\mathrm{A}}$? For example, if $\mathrm{A}$ is reduced (resp.\ integral, resp.\ integral and integrally closed), is the same true of $\widehat{\mathrm{A}}$? Most of these questions are related to the properties of the local \emph{formal fibres} of $\mathrm{A}$ (recall that these are the fibres of the canonical morphism $\operatorname{Spec}(\widehat{\mathrm{A}})\to\operatorname{Spec}(\mathrm{A})$). An exception is formed by the properties related to the notion of dimension, in particular the notion of being equidimensional; it is then the condition of \emph{chains} and its various refinements that play the essential role.

B) For a locally Noetherian prescheme $\mathrm{X}$ (in particular for an affine scheme $\operatorname{Spec}(\mathrm{A})$), does the set of $x\in\mathrm{X}$ at which the local ring $\sh{O}_x$ possesses a certain property (for example being integrally closed, or being a Cohen-Macaulay ring, or being regular) form an \emph{open} set?

C) For an \emph{integral} ring $\mathrm{A}$, is the integral closure of $\mathrm{A}$ in a finite extension of its field of fractions a finite $\mathrm{A}$-module? Of course, this question can be reformulated for Noetherian preschemes (\textbf{II}, $6.3$).

The problems of type B) or C) may also be posed for \emph{local} rings, but it does not suffice in general that they be solved affirmatively for every local ring $\mathrm{A}_\mathfrak{p}$ of a ring $\mathrm{A}$ for them to be solved for $\mathrm{A}$ (cf.\ \sref{IV.6.13.6}).

Let us also emphasise that in the study of these problems, we have been systematically concerned with whether an affirmative answer to one of them is \emph{stable} under the two most important operations in Commutative Algebra: localisation and passage to an algebra of finite type.

The results obtained in the study of these problems lead to identifying a class of Noetherian rings whose behaviour in this regard is the best possible:
\end{env}

\begin{definition}[7.8.2]
\label{IV.7.8.2}
A ring $\mathrm{A}$ is called \emph{excellent} if it is Noetherian and satisfies the following conditions:
\begin{enumerate}[label=(\roman*)]
  \item $\mathrm{A}$ is universally catenary (or, equivalently by \sref{IV.5.6.3},~(i), for every prime ideal $\mathfrak{p}$ of $\mathrm{A}$, $\mathrm{A}_\mathfrak{p}$ is universally catenary).
  \item For every prime ideal $\mathfrak{p}$ of $\mathrm{A}$, the formal fibres of $\mathrm{A}_\mathfrak{p}$ are geometrically regular.
  \item For every integral quotient ring $\mathrm{B}$ of $\mathrm{A}$ and every finite radicial extension $\mathrm{K}'$ of the field of fractions $\mathrm{K}$ of $\mathrm{B}$, there exists a finite sub-$\mathrm{B}$-algebra $\mathrm{B}'$ of $\mathrm{K}'$, containing $\mathrm{B}$, having $\mathrm{K}'$ as field of fractions, and such that the set of regular points of $\operatorname{Spec}(\mathrm{B}')$ contains a non-empty open set.
\end{enumerate}
With this terminology, the essential content of the results of \S~7 and of part of \S~6 is summarised as follows:
\end{definition}

\begin{scholium}[7.8.3]
\label{IV.7.8.3}
(i) In the conditions (i) and (ii) of the definition \sref{IV.7.8.2} one may restrict to the maximal ideals $\mathfrak{p}$ in $\mathrm{A}$. For a local Noetherian ring to be excellent, it is necessary and

\oldpage[IV-2]{215}

sufficient that it be universally catenary and that its formal fibres be geometrically regular.

(ii) If $\mathrm{A}$ is an excellent ring, so is every ring of fractions $\mathrm{S}^{-1}\mathrm{A}$ and every $\mathrm{A}$-algebra of finite type.

(iii) A local complete ring (in particular a field) is excellent. A Dedekind ring whose field of fractions is of characteristic $0$ (in particular $\mathbf{Z}$) is excellent.

(iv) Let $\mathrm{A}$ be an excellent ring, $\mathrm{X}=\operatorname{Spec}(\mathrm{A})$; the set $\operatorname{Reg}(\mathrm{X})$ (resp.\ $\operatorname{Nor}(\mathrm{X})$, $\mathrm{U}_{\mathrm{R}_n}(\mathrm{X})$) of points where $\mathrm{X}$ is regular (resp.\ normal, resp.\ satisfies $(\mathrm{R}_n)$) is open in $\mathrm{X}$. For every coherent $\sh{O}_\mathrm{X}$-Module $\sh{F}$, the set $\mathrm{U}_{\mathrm{C}_n}(\sh{F})$ (resp.\ $\mathrm{U}_{\mathrm{S}_n}(\sh{F})$, $\mathrm{CM}(\sh{F})$) of $x\in\mathrm{X}$ where $\operatorname{coprof}(\sh{F}_x)\leqslant n$ (resp.\ the points where $\sh{F}$ satisfies $(\mathrm{S}_n)$, resp.\ the points where $\sh{F}_x$ is a Cohen-Macaulay $\sh{O}_x$-Module) is open in $\mathrm{X}$.

(v) Let $\mathrm{A}$ be an excellent ring, $\mathfrak{J}$ an ideal of $\mathrm{A}$, and $\widehat{\mathrm{A}}$ the separated completion of $\mathrm{A}$ for the $\mathfrak{J}$-adic topology. Then the canonical morphism $f:\operatorname{Spec}(\widehat{\mathrm{A}})\to\operatorname{Spec}(\mathrm{A})$ is regular (in other words, flat and with geometrically regular fibres). If we set $\mathrm{X}=\operatorname{Spec}(\mathrm{A})$, $\mathrm{X}'=\operatorname{Spec}(\widehat{\mathrm{A}})$, use the notation of (iv), and put $\sh{F}'=\sh{F}\otimes_{\sh{O}_\mathrm{X}}\sh{O}_{\mathrm{X}'}$, then

\begin{env}[7.8.3.1]
\label{IV.7.8.3.1}
\begin{align*}
  \operatorname{Reg}(\mathrm{X}')&=f^{-1}(\operatorname{Reg}(\mathrm{X})),\quad \operatorname{Nor}(\mathrm{X}')=f^{-1}(\operatorname{Nor}(\mathrm{X})),\quad \mathrm{U}_{\mathrm{R}_n}(\mathrm{X}')=f^{-1}(\mathrm{U}_{\mathrm{R}_n}(\mathrm{X}))\\
  \mathrm{U}_{\mathrm{C}_n}(\sh{F}')&=f^{-1}(\mathrm{U}_{\mathrm{C}_n}(\sh{F})),\quad \mathrm{U}_{\mathrm{S}_n}(\sh{F}')=f^{-1}(\mathrm{U}_{\mathrm{S}_n}(\sh{F})),\quad \mathrm{CM}(\sh{F}')=f^{-1}(\mathrm{CM}(\sh{F}))
\end{align*}
\end{env}

In particular, if $\mathfrak{J}$ is contained in the radical of $\mathrm{A}$ (for example if $\mathrm{A}$ is local and $\mathfrak{J}$ its maximal ideal), for $\mathrm{A}$ to be regular (resp.\ normal, resp.\ reduced, resp.\ to satisfy $(\mathrm{R}_n)$, resp.\ its codepth to be $\leqslant n$, resp.\ to satisfy $(\mathrm{S}_n)$, resp.\ to be a Cohen-Macaulay ring), it is necessary and sufficient that the same hold for $\widehat{\mathrm{A}}$; in particular, for $\mathrm{A}$ to have no immersed associated primes, it is necessary and sufficient that the same hold for $\widehat{\mathrm{A}}$.

(vi) An excellent ring is universally Japanese; in particular, if $\mathrm{A}$ is integral, its integral closure in every finite extension of its field of fractions is a finite $\mathrm{A}$-algebra.

(vii) Let $\mathrm{A}$ be a reduced excellent local ring. Then the completion $\widehat{\mathrm{A}}$ is reduced, the integral closure $\mathrm{A}'$ of $\mathrm{A}$ in its total ring of fractions is a finite $\mathrm{A}$-algebra (hence a semi-local Noetherian ring), and the integral closure of $\widehat{\mathrm{A}}$ in its total ring of fractions is isomorphic to the completion $\widehat{\mathrm{A}}'$ of $\mathrm{A}'$. Furthermore, there is a bijective canonical correspondence between the set of maximal points of $\operatorname{Spec}(\widehat{\mathrm{A}})$ (in other words, the set of minimal prime ideals of $\widehat{\mathrm{A}}$) and the set of closed points of $\operatorname{Spec}(\mathrm{A}')$ (in other words, the set of maximal ideals of $\mathrm{A}'$). For the local ring $\mathrm{A}$ to be unibranch, it is necessary and sufficient that $\widehat{\mathrm{A}}$ be integral; for $\mathrm{A}$ to be geometrically unibranch, it is necessary and sufficient that $\widehat{\mathrm{A}}$ be.

(viii) If $\mathrm{A}$ is an excellent integral ring, the integral closure of $\mathrm{A}$ is the intersection of the integral closures of the $\mathrm{A}_\mathfrak{p}$, where $\mathfrak{p}$ ranges over the set of prime ideals of $\mathrm{A}$ of height~$1$.

(ix) Let $\mathrm{A}$ be an excellent ring, $\mathrm{X}=\operatorname{Spec}(\mathrm{A})$, $\mathrm{Z}$ a closed subset of $\mathrm{X}$, $\mathrm{U}=\mathrm{X}-\mathrm{Z}$, $i:\mathrm{U}\to\mathrm{X}$ the canonical injection, $\sh{F}$ a coherent $\sh{O}_\mathrm{U}$-Module. For the $\sh{O}_\mathrm{X}$-Module $i_*(\sh{F})$ to be coherent, it is necessary and sufficient that, for every $x\in\operatorname{Ass}(\sh{F})$, one has $\operatorname{codim}(\overline{\{x\}}\cap\mathrm{Z},\overline{\{x\}})\geqslant 2$. In particular, if $\mathrm{A}$ is integral and $\sh{F}$ nonzero and torsion-free, for $i_*(\sh{F})$ to be coherent, it is necessary and sufficient that $\operatorname{codim}(\mathrm{Z},\mathrm{X})\geqslant 2$.

\oldpage[IV-2]{216}

(x) Let $\mathrm{A}$ be an excellent local ring. For every integral quotient ring $\mathrm{B}$ of $\mathrm{A}$, $\widehat{\mathrm{B}}$ is equidimensional. For $\mathrm{A}$ to be equidimensional, it is necessary and sufficient that $\widehat{\mathrm{A}}$ be.
\end{scholium}

\begin{proof}
Most of these results have already been proved.

(i) The first assertion follows from \sref{IV.5.6.3},~(i) and \sref{IV.7.4.4},~(ii); the second from \sref{IV.7.4.4}, \sref{IV.7.3.18}, \sref{IV.6.12.7} and \sref{IV.6.12.4}.

Assertion (ii) follows from \sref{IV.5.6.1}, \sref{IV.5.6.3},~(i), \sref{IV.7.4.4} and \sref{IV.6.12.4}.

(iii) The first assertion has been seen in \sref{IV.5.6.4} and \sref{IV.7.4.4}, taking into account (i). The second has been seen in \sref{IV.5.6.4} and \sref{IV.7.3.19},~(iv), taking into account (i).

Assertion (iv) follows from \sref{IV.6.12.4}, \sref{IV.6.13.5}, \sref{IV.6.12.9} and \sref{IV.6.11.8} (taking into account (ii)).

The first assertion of (v) is a consequence of \sref{IV.7.4.6}; the relations \sref{IV.7.8.3.1} follow, taking into account respectively \sref{IV.6.5.1}, \sref{IV.6.5.4}, \sref{IV.6.5.3}, \sref{IV.6.3.5} and \sref{IV.6.4.2}. The last assertion of (v) is a particular case corresponding to the property $(\mathrm{S}_n)$ for $n=1$ \sref{IV.5.7.5}.

Assertion (vi) follows from \sref{IV.7.7.3}, and assertion (vii) from \sref{IV.7.6.1}, \sref{IV.7.6.2} and \sref{IV.7.6.3}.

For assertion (viii), let us note that the ring $\mathrm{A}^{(1)}$, intersection of the $\mathrm{A}_\mathfrak{p}$ for all prime ideals $\mathfrak{p}$ of height~$1$ of $\mathrm{A}$, is a finite $\mathrm{A}$-algebra, as it follows from (vi), from \sref{IV.7.8.2},~(i) and \sref{IV.5.11.2}. As the integral closure $\mathrm{A}'$ of $\mathrm{A}$ is a finite $\mathrm{A}$-algebra, the prime ideals of height~$1$ of $\mathrm{A}'$ are exactly those which lie above prime ideals of height~$1$ of $\mathrm{A}$ \sref{IV.5.10.17},~(iv); the conclusion follows hence from \sref[0]{0.23.2.9}.

The first assertion of (ix) is a consequence of (vi) and \sref{IV.5.11.4}; the second is a particular case, since, when $\sh{F}$ is nonzero and torsion-free, $\operatorname{Ass}(\sh{F})$ is reduced to the generic point of $\operatorname{Spec}(\mathrm{A})$. Finally, to prove (x), it suffices to observe that $\mathrm{B}$ is an excellent local ring by (ii). Since $\mathrm{B}$ is universally catenary and universally Japanese, the ring $\mathrm{B}^{(1)}$ is a finite $\mathrm{B}$-algebra by (vi) and \sref{IV.5.11.2}; this shows that $\mathrm{B}$ is strictly formally catenary by \sref{IV.7.2.5},~$c)$), and completes the proof of (x).\footnote{The printed text has $\mathrm{A}$ in this last clause; the argument and the first assertion of (x) require $\mathrm{B}$.}
\end{proof}

\begin{remarks}[7.8.4]
\label{IV.7.8.4}
(i) If one considers only the Noetherian rings $\mathrm{A}$ that satisfy the conditions (ii) and (iii) of the definition \sref{IV.7.8.2}, one verifies by inspection of the preceding proofs that the assertions (ii), (iii), (iv), (v), (vi) and (vii) of \sref{IV.7.8.3} are still valid, replacing ``excellent'' by ``satisfying \sref{IV.7.8.2},~(ii) and (iii)''.

(ii) The catenary local ring $\mathrm{A}$ constructed in \sref{IV.5.6.11} satisfies the conditions (ii) and (iii) (but not condition (i)) of the definition \sref{IV.7.8.2}. Indeed, the formal fibre at the maximal ideal $\mathfrak{n}\mathrm{C}_\mathfrak{n}$ of $\mathrm{A}$ is trivially geometrically regular, and the formal fibre at every other prime ideal $\mathfrak{p}$ of $\mathrm{A}$ coincides with that of the ring $\mathrm{E}$ at $\mathfrak{p}$. Now $\mathrm{E}$ is a $\mathrm{V}$-algebra of finite type, and $\mathrm{V}$ is a discrete valuation ring, hence is excellent when its residue field $k_0$ has characteristic~$0$ by \sref{IV.7.8.3},~(iii). Consequently $\mathrm{E}$ is excellent by \sref{IV.7.8.3},~(ii), which proves that $\mathrm{A}$ satisfies condition \sref{IV.7.8.2},~(ii).

\oldpage[IV-2]{217}

On the other hand, the complement of the closed point in $\operatorname{Spec}(\mathrm{A})$ is isomorphic to an open subset of $\operatorname{Spec}(\mathrm{E})$, one of whose affine open subsets is the spectrum of an excellent ring and therefore satisfies condition \sref{IV.6.12.4},~$a)$) by \sref{IV.7.8.3},~(iv). It follows that $\mathrm{A}$ itself satisfies condition \sref{IV.6.12.4},~$a)$), hence also \sref{IV.6.12.4},~$c)$), which is precisely condition \sref{IV.7.8.2},~(iii). Nevertheless $\mathrm{A}$ does not satisfy condition \sref{IV.7.8.2},~(i), since we saw in \sref{IV.5.6.11} that it is not universally catenary.

(iii) One can replace the condition (i) of the definition \sref{IV.7.8.2} by the following (weaker in appearance, taking into account \sref{IV.5.6.10}):

(i$'$) \emph{For every minimal prime ideal} $\mathfrak{p}_i$ \emph{of} $\mathrm{A}$ $(1\leqslant i\leqslant r)$, \emph{let} $\mathrm{A}'_i$ \emph{be the integral closure of the integral ring} $\mathrm{A}_i=\mathrm{A}/\mathfrak{p}_i$; \emph{then, for every maximal ideal} $\mathfrak{m}'$ \emph{of} $\mathrm{A}'_i$, \emph{denoting by} $\mathfrak{m}$ \emph{the inverse image of} $\mathfrak{m}'$ \emph{in} $\mathrm{A}$,

\begin{env}[7.8.4.1]
\label{IV.7.8.4.1}
\[
  \dim(\mathrm{A}_\mathfrak{m})=\dim(\mathrm{A}'_{i,\mathfrak{m}'}).
\]
\end{env}

Indeed, Remark~(i) shows that, under only the hypotheses (ii) and (iii) of \sref{IV.7.8.2}, the analogues of the properties (ii), (v), (vi) and (vii) of \sref{IV.7.8.3} remain valid. It follows first from (vi) and (ii) that the $\mathrm{A}'_i$ satisfy \sref{IV.7.8.2},~(ii) and (iii); then (v) shows that the completions $((\mathrm{A}'_i)_{\mathfrak{m}'})^\wedge$ are integral (and integrally closed). Let now $\mathfrak{m}$ be an arbitrary maximal ideal of $\mathrm{A}$. For every $i$ such that $\mathfrak{p}_i\subset\mathfrak{m}$, the integral closure of $\mathrm{A}_\mathfrak{m}/\mathfrak{p}_i\mathrm{A}_\mathfrak{m}$ is a finite $(\mathrm{A}_\mathfrak{m}/\mathfrak{p}_i\mathrm{A}_\mathfrak{m})$-algebra by (vi), hence is semi-local, and its local components are of the form $(\mathrm{A}'_i)_{\mathfrak{m}'}$, where $\mathfrak{m}'$ is a prime ideal (necessarily maximal) of $\mathrm{A}'_i$ above $\mathfrak{m}/\mathfrak{p}_i$. It then follows from (vii), from what precedes, and from \sref{IV.7.8.4.1} that the quotients of the completion $(\mathrm{A}_\mathfrak{m}/\mathfrak{p}_i\mathrm{A}_\mathfrak{m})^\wedge$ by its minimal prime ideals all have the same dimension, equal to $\dim(\mathrm{A}_\mathfrak{m})$. Consequently, by \sref{IV.7.1.9} and \sref{IV.7.1.8},~$b)$), $\mathrm{A}_\mathfrak{m}$ is formally catenary, and \emph{a fortiori} universally catenary by \sref{IV.7.1.11}. The same is therefore true of $\mathrm{A}_\mathfrak{p}$ for every prime ideal $\mathfrak{p}$ of $\mathrm{A}$ by \sref{IV.5.6.3},~(i), hence also of $\mathrm{A}$.

In particular, if $\mathrm{A}$ is normal, or more generally if $\mathrm{A}_\mathfrak{m}$ is unibranch for every maximal ideal $\mathfrak{m}$ of $\mathrm{A}$, the condition (i$'$) is always satisfied, since $\mathfrak{m}$ can contain only a single prime ideal $\mathfrak{p}_i$ \sref[0]{0.16.1.5}; one can then omit (i) from the definition \sref{IV.7.8.2}. More particularly, if $\mathrm{A}$ is a Noetherian \emph{local unibranch ring}, it is excellent if and only if its formal fibres are geometrically regular.

(iv) Recall that a quotient ring of a local Cohen--Macaulay ring (and \emph{a fortiori} a quotient ring of a regular local ring) also satisfies the properties \sref{IV.7.8.3},~(ix) and (x) by virtue of \sref{IV.7.2.7}; but the interest of quotient rings of regular rings lies above all in their cohomological properties (chap.~III, 3${}^{\text{e}}$ Partie).

(v) We will see further on (\S~18) that if $k$ is a field with a complete non-discrete valuation, the ring $k\{\!\{T_1,\ldots,T_n\}\!\}$ of \emph{convergent power series} (in a neighbourhood of $0$ in $k^n$) is an \emph{excellent ring}.
\end{remarks}

\begin{env}[7.8.5]
\label{IV.7.8.5}
One says that a locally Noetherian prescheme $\mathrm{X}$ is \emph{excellent} if for some covering $(\mathrm{U}_\alpha)$ of $\mathrm{X}$ by affine opens, each of the rings $\mathrm{A}_\alpha$ of the $\mathrm{U}_\alpha$ is excellent; this property is then true for \emph{every} covering $(\mathrm{U}_\alpha)$ of $\mathrm{X}$ by affine opens.
\end{env}

\begin{proposition}[7.8.6]
\label{IV.7.8.6}
Let $\mathrm{X}$ be an excellent prescheme.
\begin{enumerate}[label=(\roman*)]
  \item If $f:\mathrm{X}'\to\mathrm{X}$ is a morphism locally of finite type, $\mathrm{X}'$ is excellent.

\oldpage[IV-2]{218}

  \item If $\mathrm{X}$ is reduced, its normalization $\mathrm{X}'$ (\textbf{II}, $6.3.8$) is finite over $\mathrm{X}$.
  \item The sets $\operatorname{Reg}(\mathrm{X})$, $\operatorname{Nor}(\mathrm{X})$, $\mathrm{U}_{\mathrm{R}_n}(\mathrm{X})$, as well as $\mathrm{U}_{\mathrm{C}_n}(\sh{F})$, $\mathrm{U}_{\mathrm{S}_n}(\sh{F})$ and $\mathrm{CM}(\sh{F})$ for every coherent $\sh{O}_\mathrm{X}$-Module $\sh{F}$, are open in $\mathrm{X}$.
\end{enumerate}
\end{proposition}

\begin{proof}
This follows immediately from \sref{IV.7.8.3},~(ii), (vi) and (iv). Let us also note that \sref{IV.7.8.3},~(ix) is also valid without modification of the statement when $\mathrm{X}$ is an excellent prescheme.
\end{proof}


\subsection{Excellent rings and resolution of singularities}
\label{subsection:IV.7.9}

\begin{env}[7.9.1]
\label{IV.7.9.1}
Given a \emph{locally Noetherian reduced} prescheme $\mathrm{X}$, one calls a \emph{resolving morphism} for $\mathrm{X}$ a morphism $f:\mathrm{X}'\to\mathrm{X}$ such that $\mathrm{X}'$ is \emph{regular} and that $f$ is \emph{proper} and \emph{birational}. When such a morphism exists, one says that one can \emph{resolve the singularities} of $\mathrm{X}$ (or more simply that one can \emph{resolve} $\mathrm{X}$). For example, if $\mathrm{A}$ is a \emph{Japanese} ring of dimension $1$, one can resolve $\mathrm{X}=\operatorname{Spec}(\mathrm{A})$ since the morphism $\mathrm{X}'\to\mathrm{X}$ (where $\mathrm{X}'$ is the \emph{normalization} of $\mathrm{X}$) is finite (hence proper) and birational, and the local rings of $\mathrm{X}'$ are integrally closed rings of dimension at most~$1$, hence fields or discrete valuation rings, hence regular.
\end{env}

\begin{env}[7.9.2]
\label{IV.7.9.2}
It is clear that if one can resolve $\mathrm{X}$, one can also resolve every prescheme induced on an \emph{open subset} of $\mathrm{X}$ and every local prescheme $\operatorname{Spec}(\sh{O}_x)$ of $\mathrm{X}$ (\textbf{II}, $5.4.2$). On the other hand, if one can resolve $\mathrm{X}$, it is clear that there exists in $\mathrm{X}$ an open subset \emph{everywhere} dense contained in $\operatorname{Reg}(\mathrm{X})$. These remarks show immediately that if, for every closed \emph{integral} sub-prescheme $\mathrm{Y}$ of $\mathrm{X}$ and every $\mathrm{Y}$-prescheme $\mathrm{Y}'$ \emph{integral, finite} and \emph{radicial} over $\mathrm{Y}$, one can resolve $\mathrm{Y}'$, then the affine opens of $\mathrm{X}$ \emph{satisfy condition} (iii) of \sref{IV.7.8.2}.

On the other hand, for local schemes, one has the following proposition:
\end{env}

\begin{proposition}[7.9.3]
\label{IV.7.9.3}
Let $\mathrm{A}$ be a local Noetherian reduced ring, and suppose that one can resolve $\operatorname{Spec}(\mathrm{A})$. Then the fibres of the canonical morphism $\operatorname{Spec}(\widehat{\mathrm{A}})\to\operatorname{Spec}(\mathrm{A})$ at the maximal points of $\operatorname{Spec}(\mathrm{A})$ are regular.
\end{proposition}

\begin{proof}
Set $\mathrm{X}=\operatorname{Spec}(\mathrm{A})$, $\mathrm{X}'=\operatorname{Spec}(\widehat{\mathrm{A}})$, and let $g:\mathrm{X}'\to\mathrm{X}$ be the canonical morphism. Let $f:\mathrm{Y}\to\mathrm{X}$ be a resolving morphism; set $\mathrm{Y}'=\mathrm{X}'\times_\mathrm{X}\mathrm{Y}$, and let $f':\mathrm{Y}'\to\mathrm{X}'$ and $g':\mathrm{Y}'\to\mathrm{Y}$ be the canonical projections, so that one has a commutative diagram
\[
  \xymatrix{
    \mathrm{X}' \ar[d]_g & \mathrm{Y}' \ar[l]_{f'} \ar[d]^{g'} \\
    \mathrm{X} & \mathrm{Y} \ar[l]^f
  }
\]

It will suffice to prove that the prescheme $\mathrm{Y}'$ is \emph{regular}: indeed, as $f$ is birational and of finite type, there is an open $\mathrm{U}$ dense in $\mathrm{X}$ such that the restriction $f^{-1}(\mathrm{U})\to\mathrm{U}$ of $f$ is an isomorphism and that $f^{-1}(\mathrm{U})$ is everywhere dense in $\mathrm{Y}$ (\textbf{I}, $6.5.5$); the preschemes induced respectively on the open $g^{-1}(\mathrm{U})$ of $\mathrm{X}'$ and the open $g'^{-1}(f^{-1}(\mathrm{U}))$ of $\mathrm{Y}'$ are therefore isomorphic; this proves that $g^{-1}(\mathrm{U})$ is regular, and \emph{a fortiori} so are the fibres of $g$ at the maximal points of $\mathrm{X}$, which are contained in $\mathrm{U}$.

\oldpage[IV-2]{219}

Let us note that the morphism $f'$ is proper (\textbf{II}, $5.4.2$), and in particular of finite type, hence $\mathrm{Y}'$ is Noetherian, since $\mathrm{X}'$ is Noetherian. Let $a'$ be the closed point of $\mathrm{X}'$; to see that $\mathrm{Y}'$ is regular, it will suffice to prove that \emph{for every} $y'\in f'^{-1}(a')$, \emph{the local ring} $\sh{O}_{\mathrm{Y}',y'}$ \emph{is regular}. Indeed, one will then have $f'^{-1}(a')\subset\operatorname{Reg}(\mathrm{Y}')$. But as $\mathrm{X}'$ is the spectrum of a complete local ring and $f'$ is of finite type, $\operatorname{Reg}(\mathrm{Y}')$ is \emph{open} in $\mathrm{Y}'$ \sref{IV.6.12.8}, and as the morphism $f'$ is \emph{closed} and $f'^{-1}(a')\subset\operatorname{Reg}(\mathrm{Y}')$, there exists a neighbourhood $\mathrm{V}'$ of $a'$ in $\mathrm{X}'$ such that $f'^{-1}(\mathrm{V}')\subset\operatorname{Reg}(\mathrm{Y}')$. As $\widehat{\mathrm{A}}$ is a local ring, one necessarily has $\mathrm{V}'=\mathrm{X}'$, hence $\operatorname{Reg}(\mathrm{Y}')=\mathrm{Y}'$ and the proposition will be proved.

Now, if $a$ is the closed point of $\mathrm{X}$, one has $g^{-1}(a)=\{a'\}$ and the residue fields $\boldsymbol{k}(a)$ and $\boldsymbol{k}(a')$ are isomorphic; hence $g'^{-1}(f^{-1}(a))=f'^{-1}(a')$ and the restriction $f'^{-1}(a')\to f^{-1}(a)$ of $g'$ is an \emph{isomorphism} of these two fibres. Let $y=g'(y')$ and set $\mathrm{B}=\sh{O}_{\mathrm{Y},y}$; if $\mathfrak{n}=\mathfrak{m}_y$ is the maximal ideal of $\mathrm{B}$, what precedes shows that there exists in the (Noetherian) ring $\mathrm{B}'=\mathrm{B}\otimes_\mathrm{A}\widehat{\mathrm{A}}$ a single maximal ideal $\mathfrak{n}'$ above $\mathfrak{n}$ and that $\sh{O}_{\mathrm{Y}',y'}=\mathrm{B}'_{\mathfrak{n}'}$. To show that the local ring $\mathrm{B}'_{\mathfrak{n}'}$ is regular, it suffices to establish that its \emph{completion} is regular \sref[0]{0.17.1.5}; but by hypothesis $\mathrm{B}$ is regular, hence so is its completion $\widehat{\mathrm{B}}$. The conclusion will follow from the following lemma:

\begin{lemma}[7.9.3.1]
\label{IV.7.9.3.1}
Let $\mathrm{A}$, $\mathrm{B}$ be two local Noetherian rings, $\rho:\mathrm{A}\to\mathrm{B}$ a local homomorphism such that the ring $\mathrm{B}'=\mathrm{B}\otimes_\mathrm{A}\widehat{\mathrm{A}}$ is Noetherian. Then, for every maximal ideal $\mathfrak{n}'$ of $\mathrm{B}'$ above the maximal ideal $\mathfrak{n}$ of $\mathrm{B}$ and the maximal ideal of $\widehat{\mathrm{A}}$, the completions of $\mathrm{B}$ and of $\mathrm{B}'_{\mathfrak{n}'}$ are isomorphic.
\end{lemma}

\begin{proof}
Let $\mathfrak{m}$ be the maximal ideal of $\mathrm{A}$ and set $\mathrm{B}''=\mathrm{B}'_{\mathfrak{n}'}$. For every integer $h>0$, one has $\mathrm{B}'/\mathfrak{m}^h\mathrm{B}'=\mathrm{B}\otimes_\mathrm{A}(\widehat{\mathrm{A}}/\mathfrak{m}^h\widehat{\mathrm{A}})$; but $\widehat{\mathrm{A}}/\mathfrak{m}^h\widehat{\mathrm{A}}$ is isomorphic to $\mathrm{A}/\mathfrak{m}^h$, hence $\mathrm{B}'/\mathfrak{m}^h\mathrm{B}'$ is isomorphic to $\mathrm{B}/\mathfrak{m}^h\mathrm{B}$, and in particular it is a local ring whose maximal ideal is $\mathfrak{n}'/\mathfrak{m}^h\mathrm{B}'$; consequently $\mathrm{B}''/\mathfrak{m}^h\mathrm{B}''$, isomorphic to $(\mathrm{B}'/\mathfrak{m}^h\mathrm{B}')_{\mathfrak{n}'}$, is B-isomorphic to $\mathrm{B}'/\mathfrak{m}^h\mathrm{B}'$, and finally isomorphic to $\mathrm{B}/\mathfrak{m}^h\mathrm{B}$; in particular, the maximal ideal of $\mathrm{B}''$ is equal to $\mathfrak{n}\mathrm{B}''$, and $\mathrm{B}''/\mathfrak{n}^h\mathrm{B}''$ is isomorphic to $\mathrm{B}/\mathfrak{n}^h$ from what precedes. As $\widehat{\mathrm{B}}''=\varprojlim_h(\mathrm{B}''/\mathfrak{n}^h\mathrm{B}'')$, the lemma is proved, as was \sref{IV.7.9.3}.
\end{proof}
\end{proof}

\begin{corollary}[7.9.4]
\label{IV.7.9.4}
Let $\mathrm{A}$ be a local Noetherian ring.
\begin{enumerate}[label=(\roman*)]
  \item If, for every integral quotient ring $\mathrm{B}$ of $\mathrm{A}$, one can resolve $\operatorname{Spec}(\mathrm{B})$, then the formal fibres of $\mathrm{A}$ are regular.
  \item Suppose that, for every integral quotient ring $\mathrm{B}$ of $\mathrm{A}$, and every integral ring $\mathrm{B}'$ containing $\mathrm{B}$, which is a finite $\mathrm{B}$-algebra and whose field of fractions is radicial over that of $\mathrm{B}$, one can resolve $\operatorname{Spec}(\mathrm{B}')$; then the formal fibres of $\mathrm{A}$ are geometrically regular.
\end{enumerate}
\end{corollary}

\begin{proof}
Every formal fibre of $\mathrm{A}$ being the formal fibre of an integral quotient ring of $\mathrm{A}$ at the generic point of its spectrum, it is clear that (i) is an immediate consequence of \sref{IV.7.9.3}, and (ii) follows from (i) and \sref{IV.6.7.7}.
\end{proof}

\begin{proposition}[7.9.5]
\label{IV.7.9.5}
Let $\mathrm{X}$ be a locally Noetherian prescheme, such that, for every integral prescheme $\mathrm{Y}$ finite over $\mathrm{X}$, one can resolve $\mathrm{Y}$. Then the rings of the affine opens of $\mathrm{X}$ satisfy the conditions (ii) and (iii) of \sref{IV.7.8.2}; if moreover $\mathrm{X}$ is universally catenary \sref{IV.5.6.3} (and

\oldpage[IV-2]{220}

in particular if $\mathrm{X}$ is \emph{locally immersible in a regular prescheme} \sref{IV.5.6.4}), $\mathrm{X}$ is an \emph{excellent prescheme}.
\end{proposition}

\begin{proof}
This follows immediately from \sref{IV.7.9.4} and \sref{IV.7.9.2}.
\end{proof}

\begin{remark}[7.9.6]
\label{IV.7.9.6}
It is possible that the converse of \sref{IV.7.9.5} is true, that is to say, that the hypothesis that $\mathrm{X}$ is reduced and that the rings of the affine opens of $\mathrm{X}$ satisfy the conditions (ii) and (iii) of \sref{IV.7.8.2} implies the possibility of resolving $\mathrm{X}$ (and consequently also the possibility of resolving every reduced prescheme locally of finite type over $\mathrm{X}$, by virtue of \sref{IV.7.4.4} and \sref{IV.6.12.4}). This is in any case true when one limits oneself to reduced Noetherian preschemes whose residue fields are of \emph{characteristic} $0$, as shown by the recent results of Hironaka~[35] (the latter states his results under hypotheses that are too restrictive, the reasoning being in fact valid under the only hypotheses (ii) and (iii) of \sref{IV.7.8.2} for the rings of the affine opens of $\mathrm{X}$, combined with the fact that the residue fields are of characteristic~$0$). On the contrary, in the general case, it is known so far only how to resolve algebraic schemes of dimension $2$ over a perfect field~[36]. Given the importance that the resolution of singularities is acquiring in the topological study of schemes (particularly concerning their homological and homotopical properties), this gives a particular interest to the category of excellent rings and excellent preschemes.

One may also note in this regard that the less delicate part of the reasoning of Hironaka (\emph{loc.\ cit.}) shows independently of any hypothesis on the characteristics of the residue fields, and using only the properties (ii) and (iii) of \sref{IV.7.8.2} for the local rings of $\mathrm{X}$, the problem of resolution of singularities of a prescheme $\mathrm{X}$ reduces to the case where $\mathrm{X}=\operatorname{Spec}(\mathrm{A})$, $\mathrm{A}$ being a local Noetherian integral ring and \emph{complete}. Consequently, if subsequent results were to disprove the more advanced conjecture above, and were to lead to formulating more restrictive conditions on the local rings of $\mathrm{X}$, these could only be through the intermediary of restrictive conditions imposed on \emph{complete local} rings (probably conditions on their residue fields, perhaps the condition $[k:k^p]<+\infty$ ($p$ being the characteristic exponent of $k$)).
\end{remark}

\begin{env}[7.9.7]
\label{IV.7.9.7}
Consider on the one hand a full subcategory $\mathcal{C}$ of the category of locally Noetherian preschemes, and on the other hand a property $\boldsymbol{R}(\mathrm{A})$, subject to the following conditions:

1${}^\circ$ For every $\mathrm{X}\in\mathcal{C}$, every prescheme locally of finite type over $\mathrm{X}$ belongs to $\mathcal{C}$. For every Noetherian ring $\mathrm{A}$ such that $\operatorname{Spec}(\mathrm{A})\in\mathcal{C}$ and every multiplicative subset $\mathrm{S}$ of $\mathrm{A}$, one has $\operatorname{Spec}(\mathrm{S}^{-1}\mathrm{A})\in\mathcal{C}$.

2${}^\circ$ For every $\mathrm{X}\in\mathcal{C}$, the set $\mathrm{U}_{\boldsymbol{R}}(\mathrm{X})$ of $x\in\mathrm{X}$ such that $\boldsymbol{R}(\sh{O}_{\mathrm{X},x})$ is true is open in $\mathrm{X}$.

3${}^\circ$ For every local Noetherian ring $\mathrm{A}$ such that $\operatorname{Spec}(\mathrm{A})\in\mathcal{C}$ and every regular element $t$ of the maximal ideal of $\mathrm{A}$, $\boldsymbol{R}(\mathrm{A}/t\mathrm{A})$ implies $\boldsymbol{R}(\mathrm{A})$.

Let us denote as in \sref{IV.7.5.0} by $\boldsymbol{P}(\mathrm{Z},k)$, for a field $k$ and a $k$-prescheme $\mathrm{Z}\in\mathcal{C}$, the following property:

``For every finite extension $k'$ of $k$, all the local rings of $\mathrm{Z}\otimes_k k'$ satisfy the property $\boldsymbol{R}$.''

\oldpage[IV-2]{221}

We will further suppose that $\boldsymbol{R}$ satisfies the following condition:

4${}^\circ$ If $\boldsymbol{P}(\mathrm{Z},k)$ is true, then, for every extension \emph{of finite type} $k''$ of $k$, all the local rings of $\mathrm{Z}\otimes_k k''$ satisfy the property $\boldsymbol{R}$.

Let us note that these conditions are satisfied when one takes for $\mathcal{C}$ the category of \emph{excellent} preschemes and for $\boldsymbol{R}$ one of the properties (i) to (viii) of \sref{IV.7.5.3}; for condition 1${}^\circ$, this follows from \sref{IV.7.8.3},~(ii), and for conditions 2${}^\circ$ and 3${}^\circ$, from the arguments of \sref{IV.7.5.3}, taking into account that every excellent ring is catenary. Finally, condition 4${}^\circ$ follows, for (i), (ii), and (iii), from \sref{IV.6.7.1}; for (iv), (v), and (vi), from \sref{IV.6.7.7}; and for (viii), from \sref{IV.4.6.1}. As concerns property (vii) of \sref{IV.7.5.3}, the corresponding property $\boldsymbol{P}(\mathrm{Z},k)$ implies that $\mathrm{Z}$ is locally integral (being locally Noetherian) and that each of the integral subpreschemes whose union is $\mathrm{Z}$ is geometrically integral, by \sref{IV.4.5.9} and \sref{IV.4.6.1}; condition 4${}^\circ$ follows in this case as well.

With these notations and hypotheses:
\end{env}

\begin{proposition}[7.9.8]
\label{IV.7.9.8}
Let $\mathrm{A}$ be a local Noetherian ring, $k$ its residue field, and $\mathrm{B}$ a local Noetherian ring such that $\operatorname{Spec}(\mathrm{B})\in\mathcal{C}$. Let $\varphi:\mathrm{A}\to\mathrm{B}$ be a local homomorphism making $\mathrm{B}$ a flat $\mathrm{A}$-module; set $\mathrm{Y}=\operatorname{Spec}(\mathrm{A})$, $\mathrm{X}=\operatorname{Spec}(\mathrm{B})$, and let $f:\mathrm{X}\to\mathrm{Y}$ be the morphism corresponding to $\varphi$. Suppose that:

1${}^\circ$ the property $\boldsymbol{P}(\mathrm{B}\otimes_\mathrm{A} k,k)$ is true;

2${}^\circ$ for every finite morphism $\mathrm{Y}_1\to\mathrm{Y}$, one can resolve $(\mathrm{Y}_1)_{\mathrm{red}}$.

Then, for every $y\in\mathrm{Y}$, the property $\boldsymbol{P}(f^{-1}(y),\boldsymbol{k}(y))$ is true.
\end{proposition}

\begin{proof}
Let us note that condition 2${}^\circ$ of \sref{IV.7.9.8} is still satisfied when one replaces $\mathrm{Y}$ by the spectrum of a local ring in a maximal ideal of a finite $\mathrm{A}$-algebra \sref{IV.7.9.2}; on the other hand, by virtue of condition 1${}^\circ$ of \sref{IV.7.9.7}, the spectrum of a ring of fractions of $\mathrm{B}\otimes_\mathrm{A}\mathrm{A}'$ also belongs to $\mathcal{C}$. The lemma \sref{IV.7.3.16.2} then shows (as in part I of the reasoning of \sref{IV.7.5.2}) that it suffices to prove that, if $\mathrm{Y}$ is integral and if $y$ is the generic point of $\mathrm{Y}$, the local rings of the fibre $f^{-1}(y)$ satisfy $\boldsymbol{R}$.

This being so, there exists by hypothesis a regular prescheme $\mathrm{Y}'$ and a proper birational morphism $g:\mathrm{Y}'\to\mathrm{Y}$. Since $\mathrm{Y}'$ is Noetherian and locally integral, the fact that $g$ is birational implies that $\mathrm{Y}'$ is integral. Put $\mathrm{X}'=\mathrm{X}\times_\mathrm{Y}\mathrm{Y}'$, so that one has a commutative diagram

\begin{env}[7.9.8.1]
\label{IV.7.9.8.1}
\[
  \xymatrix{
    \mathrm{X} \ar[d]_f & \mathrm{X}' \ar[l]_{g'} \ar[d]^{f'} \\
    \mathrm{Y} & \mathrm{Y}' \ar[l]^g
  }
\]
\end{env}

where $f'$ and $g'$ are the canonical projections. Taking into account condition 2${}^\circ$ of \sref{IV.7.9.7}, the same reasoning as at the beginning of \sref{IV.7.9.3} shows that it suffices to prove

\oldpage[IV-2]{222}

that $\mathrm{U}_{\boldsymbol{R}}(\mathrm{X}')=\mathrm{X}'$. If $a$ denotes the closed point of $\mathrm{X}$, it is enough to show that $g'^{-1}(a)\subset\mathrm{U}_{\boldsymbol{R}}(\mathrm{X}')$. Let $x'\in g'^{-1}(a)$, and set $z'=f'(x')$, so that $b=g(z')$ is the closed point of $\mathrm{Y}$. We have
\[
  f'^{-1}(z')=f^{-1}(b)\otimes_k\boldsymbol{k}(z').
\]
Since $\boldsymbol{k}(z')$ is an extension of finite type of $k$ (because $g$ is proper, hence of finite type), the hypothesis that $\boldsymbol{P}(f^{-1}(b),k)$ is true implies, by condition 4${}^\circ$ of \sref{IV.7.9.7}, that $\boldsymbol{P}(f'^{-1}(z'),\boldsymbol{k}(z'))$ is true. Moreover, $\sh{O}_{\mathrm{Y}',z'}$ is regular and $\operatorname{Spec}(\sh{O}_{\mathrm{X}',x'})\in\mathcal{C}$ by condition 1${}^\circ$ of \sref{IV.7.9.7}; Lemma \sref{IV.7.5.1.1} therefore proves that $\boldsymbol{R}(\sh{O}_{\mathrm{X}',x'})$ is true, completing the proof.
\end{proof}

\begin{corollary}[7.9.9]
\label{IV.7.9.9}
Let $\mathrm{A}$, $\mathrm{B}$ be two local Noetherian rings, $\varphi:\mathrm{A}\to\mathrm{B}$ a local homomorphism making $\mathrm{B}$ a flat $\mathrm{A}$-module. Suppose that $\mathrm{A}$ is integral and geometrically unibranch, and that for every finite integral $\mathrm{A}$-algebra $\mathrm{A}'$ one can resolve $\operatorname{Spec}(\mathrm{A}')$ (this will be the case if the formal fibres of $\mathrm{A}$ are geometrically regular and the residue field of $\mathrm{A}$ is of characteristic $0$, by virtue of the results of Hironaka~[35]). Let $k$ be the residue field of $\mathrm{A}$ and suppose that $\operatorname{Spec}(\mathrm{B}\otimes_\mathrm{A} k)$ is geometrically punctually integral \sref{IV.4.6.9}. Then $\mathrm{B}$ is integral.
\end{corollary}

\begin{proof}
We apply \sref{IV.7.9.8}, taking for $\mathcal{C}$ the category of all locally Noetherian preschemes and for $\boldsymbol{R}$ the property of being integral. The arguments of \sref{IV.7.5.3} and \sref{IV.7.9.7} show that conditions 1${}^\circ$, 2${}^\circ$, 3${}^\circ$ and 4${}^\circ$ of \sref{IV.7.9.7} are satisfied in this case. The hypothesis on $\mathrm{B}\otimes_\mathrm{A}k$ and Proposition \sref{IV.7.9.8} therefore show (with the notation of \sref{IV.7.9.8}) that the fibres $f^{-1}(y)$ are geometrically punctually integral for $y\in\mathrm{Y}$. In particular the $f^{-1}(y)$ are reduced; since $\mathrm{A}$ is integral, and hence reduced, it follows from \sref{IV.3.3.5} that $\mathrm{B}$ is reduced. It remains to prove that $\mathrm{X}=\operatorname{Spec}(\mathrm{B})$ is \emph{irreducible}. Now, the proof of \sref{IV.7.9.8} shows that $\mathrm{X}'$ is locally integral (being locally Noetherian). Since $g':\mathrm{X}'\to\mathrm{X}$ is surjective, it is therefore enough to show that $\mathrm{X}'$ is irreducible, or, since $\mathrm{X}'$ is locally integral, that $\mathrm{X}'$ is \emph{connected}. For this it is enough to prove that the fibre $g'^{-1}(a)$ is connected. Indeed, if $\mathrm{X}'$ were the sum of two non-empty open preschemes $\mathrm{X}'_1$ and $\mathrm{X}'_2$, then, for example, $g'^{-1}(a)\cap\mathrm{X}'_2=\varnothing$; but the restriction $\mathrm{X}'_2\to\mathrm{X}$ of $g'$ is proper (since $\mathrm{X}'_2$ is also closed in $\mathrm{X}'$), so $g'(\mathrm{X}'_2)$ would be a non-empty closed subset of $\mathrm{X}$ not containing the closed point $a$, a contradiction. Finally,
\[
  g'^{-1}(a)=g^{-1}(b)\otimes_{\boldsymbol{k}(b)}\boldsymbol{k}(a).
\]
The morphism $g:\mathrm{Y}'\to\mathrm{Y}$ is proper and birational; hence, in its Stein factorisation $\mathrm{Y}'\xrightarrow{v}\mathrm{Y}''\xrightarrow{u}\mathrm{Y}$ (\textbf{III}, $4.3.3$), the finite morphism $u$ is also birational, and consequently $\mathrm{Y}''$ is integral. Since $\mathrm{A}$ is geometrically unibranch, it follows from (\textbf{III}, $4.3.4$) that $g^{-1}(b)$ is geometrically connected, completing the proof.
\end{proof}

\begin{remarks}[7.9.10]
\label{IV.7.9.10}
(i) We will show further on \sref{IV.18.8.11} that for a local Noetherian reduced ring $\mathrm{B}$ to be geometrically unibranch, it is necessary and sufficient that for every étale morphism $h:\mathrm{X}'\to\mathrm{X}=\operatorname{Spec}(\mathrm{B})$ and every point $x'\in\mathrm{X}'$ above the closed point of $\mathrm{X}$, $\sh{O}_{\mathrm{X}',x'}$ is integral. One deduces that if, in \sref{IV.7.9.9} one supposes, not only that $\operatorname{Spec}(\mathrm{B}\otimes_\mathrm{A} k)$ is geometrically punctually integral, but also geometrically unibranch, then one can conclude that $\mathrm{B}$ is \emph{geometrically unibranch}.

\oldpage[IV-2]{223}

(ii) Let $\mathrm{X}$, $\mathrm{Y}$ be two excellent preschemes, $f:\mathrm{X}\to\mathrm{Y}$ a flat morphism, and suppose that for every finite morphism $\mathrm{Y}_1\to\mathrm{Y}$, one can resolve $(\mathrm{Y}_1)_{\mathrm{red}}$. It follows from \sref{IV.7.9.8}, taking for $\boldsymbol{R}$ the property of being \emph{regular}, that the set $\mathrm{U}$ of $x\in\mathrm{X}$ such that the fibre at the point $f(x)$ of the morphism
\[
  \operatorname{Spec}(\sh{O}_x\otimes_{\sh{O}_{f(x)}}\boldsymbol{k}(f(x)))\to\operatorname{Spec}(\boldsymbol{k}(f(x)))
\]
is \emph{geometrically regular}, is \emph{stable under generization}. We do not know if this set is \emph{open} (or, which amounts to the same $(0_{\mathrm{III}}, 9.2.5)$, if it is \emph{constructible}) even in the particular case where $\mathrm{Y}$ is the spectrum of $\mathbf{Z}$ or of a polynomial ring in one indeterminate $k[\mathrm{T}]$ over a field $k$, and where consequently the condition of ``resolvability'' imposed on $\mathrm{Y}$ is trivially satisfied.
\end{remarks}

\begin{center}
\emph{(To be continued.)}
\end{center}

\clearpage
\oldpage[IV-2]{224}
\phantomsection
\addcontentsline{toc}{section}{Bibliography (continued)}
\section*{Bibliography \textnormal{\emph{(continued)}}}

\begingroup
\small
\setlength{\parindent}{0pt}
\setlength{\parskip}{4pt}

\textbf{[33]} M.~\textsc{Nagata}, ``Note on a chain condition for prime
ideals,'' \emph{Mem. Coll. Sci. Kyoto}, vol.~32 (1959), pp.~85--90.

\textbf{[34]} H.~\textsc{Hironaka}, forthcoming article on \emph{Projective
morphisms}.

\textbf{[35]} H.~\textsc{Hironaka}, ``Resolution of singularities of an
algebraic variety over a field of characteristic zero,'' \emph{Ann. of Math.},
vol.~79 (1964), pp.~109--326.

\textbf{[36]} S.~\textsc{Abhyankar}, ``Local uniformization on algebraic
surfaces over ground fields of characteristic $p\ne0$,'' \emph{Ann. of
Math.}, vol.~63 (1956), pp.~491--526.

\textbf{[37]} H.~\textsc{Cartan}, \emph{Séminaire de l'École Normale
Supérieure}, 13th year (1960--61).

\textbf{[38]} J.-P.~\textsc{Serre}, \emph{Groupes algébriques et corps de
classes}, Paris (Hermann), 1959.

\endgroup

\clearpage
\oldpage[IV-2]{225}
\phantomsection
\addcontentsline{toc}{section}{Index of notation}
\section*{Index of notation}

\begingroup
\small
\setlength{\parindent}{0pt}
\setlength{\parskip}{2pt}

$\mathfrak S(X)$ ($X$ a prescheme): 2.2.15 and 4.8.1.

$\operatorname{Ass}(\mathscr F)$ ($\mathscr F$ an $\mathcal O_X$-Module): 3.1.1.

$\operatorname{deg.tr.}_K L$ ($L$ an extension of the field $K$): 4.1.

$\lambda_x(\mathscr F)$ ($\mathscr F$ an $\mathcal O_X$-Module, $x$ a maximal
point of $\operatorname{Supp}(\mathscr F)$): 4.7.5.

$\Phi(\mathscr F)$ ($\mathscr F$ an $\mathcal O_X$-Module): 4.8.1.

$\operatorname{codim}(Y,X)$ ($Y$ any subset of a prescheme $X$): 5.1.3.

$\dim(\mathscr F)$ ($\mathscr F$ an $\mathcal O_X$-Module): 5.1.12.

$\operatorname{coprof}(\mathscr F)$, $\operatorname{coprof}(X)$ ($X$ a
prescheme, $\mathscr F$ an $\mathcal O_X$-Module): 5.7.1.

$\operatorname{coprof}_A(M)$ ($A$ a Noetherian ring, $M$ a finite type
$A$-module): 5.7.12.

$\mathcal H^0_{X/Z}(\mathscr F)$ ($X$ a prescheme, $Z$ a subset of $X$ stable
under specialization, $\mathscr F$ an $\mathcal O_X$-Module): 5.9.1.

$\rho_{X/Z}$ ($X$ a prescheme, $Z$ a subset of $X$ stable under
specialization): 5.9.7.

$\operatorname{prof}_T(\mathscr F)$ ($T$ a subset of a prescheme $X$,
$\mathscr F$ an $\mathcal O_X$-Module): 5.10.1.

$Z^{(n)}(X)$ ($X$ a prescheme, $n$ an integer $\geq0$): 5.10.13.

$A^{(1)}$ ($A$ an integral Noetherian ring): 5.10.17 and 7.2.1.

$\operatorname{gr}_{\mathscr J}^{\bullet}(\mathscr F)$ ($\mathscr F$ an
$\mathcal O_X$-Module, $\mathscr J$ an Ideal of $\mathcal O_X$): 6.10.1.

$U_{S_n}(\mathscr F)$, $U_{C_n}(\mathscr F)$ ($\mathscr F$ an
$\mathcal O_X$-Module): 6.11.2.

$\operatorname{CM}(\mathscr F)$ ($\mathscr F$ an $\mathcal O_X$-Module): 6.11.3.

$U_{S_n}(X)$, $U_{C_n}(X)$, $\operatorname{CM}(X)$ ($X$ a prescheme): 6.11.4.

$\operatorname{Reg}(X)$, $\operatorname{Sing}(X)$ ($X$ a prescheme): 6.12.1.

$U_{R_n}(X)$ ($X$ a prescheme): 6.12.9.

$\operatorname{Nor}(X)$ ($X$ a prescheme): 6.13.1.

$k\{T_1,\ldots,T_n\}$: 7.4.8.

\endgroup

\clearpage
\oldpage[IV-2]{226}
\phantomsection
\addcontentsline{toc}{section}{Terminological index}
\section*{Terminological index}

\begingroup
\footnotesize
\setlength{\parindent}{0pt}
\setlength{\parskip}{1.1pt}
\newcommand{\EGAIVtwoTermEN}[2]{#1: #2.\par}

\EGAIVtwoTermEN{Separable algebra over a field}{4.6.2}
\EGAIVtwoTermEN{Ring of restricted formal power series over a complete valued field}{7.4.8}
\EGAIVtwoTermEN{Excellent ring}{7.8.2}
\EGAIVtwoTermEN{Formally catenary ring}{7.1.9}
\EGAIVtwoTermEN{Formally equidimensional ring}{7.1.1}
\EGAIVtwoTermEN{Geometrically normal, geometrically reduced, geometrically
regular ring, or ring having geometric property $(R_k)$}{6.7.6}
\EGAIVtwoTermEN{Normal ring}{5.13.5}
\EGAIVtwoTermEN{$P$-ring}{7.3.13}
\EGAIVtwoTermEN{Ring satisfying property $(R_k)$}{5.8.2}
\EGAIVtwoTermEN{Strictly equidimensional ring}{7.2.1}
\EGAIVtwoTermEN{Strictly formally catenary ring}{7.2.6}
\EGAIVtwoTermEN{Universally catenary ring}{5.6.2}
\EGAIVtwoTermEN{Antifilter of subsets of a set}{5.10.8}

\medskip
\EGAIVtwoTermEN{Codimension of an arbitrary subset of a prescheme}{5.1.3}
\EGAIVtwoTermEN{Codepth of a Module at a point, codepth of a Module,
codepth of a prescheme}{5.7.1}
\EGAIVtwoTermEN{Field of definition}{4.8.4}
\EGAIVtwoTermEN{Separably closed field}{4.3.3}
\EGAIVtwoTermEN{Associated prime cycle of a Module or of a prescheme}{3.1.1}
\EGAIVtwoTermEN{Embedded associated prime cycle}{3.1.1}

\medskip
\EGAIVtwoTermEN{Irredundant decomposition, reduced irredundant decomposition
of a Module}{3.2.5}
\EGAIVtwoTermEN{Primary decomposition of $0$ in a Module}{3.2.5}
\EGAIVtwoTermEN{Dimension of a Module}{5.1.12}
\EGAIVtwoTermEN{Dimension of a prescheme over a field}{4.1.1}

\medskip
\EGAIVtwoTermEN{Inseparability exponent of an extension of a field}{4.7.4}
\EGAIVtwoTermEN{Primary extension of a field}{4.3.1}

\medskip
\EGAIVtwoTermEN{Separable algebraic closure of a field in an extension}{4.3.4}
\EGAIVtwoTermEN{Formal fibre of a ring at a point of its spectrum}{7.3.13}

\medskip
\EGAIVtwoTermEN{Universal homeomorphism}{2.4.2}
\EGAIVtwoTermEN{Homomorphism of Modules defined over a field}{4.8.4}

\medskip
\EGAIVtwoTermEN{Hironaka lemma}{5.12.8}
\EGAIVtwoTermEN{Geometric length of a Module at a point}{4.7.5}

\medskip
\EGAIVtwoTermEN{Integral module over a ring}{3.2.4}
\EGAIVtwoTermEN{$A$-Module satisfying property $(S_k)$}{5.7.2}
\EGAIVtwoTermEN{Reduced module over a ring}{3.2.2}
\EGAIVtwoTermEN{Cohen-Macaulay $\mathcal O_X$-Module at a point,
Cohen-Macaulay $\mathcal O_X$-Module}{5.7.1}
\EGAIVtwoTermEN{Faithfully flat $\mathcal O_X$-Module (relative to a morphism)}{2.2.4}
\EGAIVtwoTermEN{Geometrically integral $\mathcal O_X$-Module over a field}{4.6.19}

\clearpage
\oldpage[IV-2]{227}

\EGAIVtwoTermEN{Geometrically punctually integral $\mathcal O_X$-Module over a field}{4.6.22}
\EGAIVtwoTermEN{Geometrically reduced (or separable) $\mathcal O_X$-Module over a field}{4.6.17}
\EGAIVtwoTermEN{Geometrically reduced (or separable) $\mathcal O_X$-Module over a field at a point}{4.6.22}
\EGAIVtwoTermEN{Equidimensional $\mathcal O_X$-Module at a point}{5.1.12}
\EGAIVtwoTermEN{Integral $\mathcal O_X$-Module, integral at a point}{3.2.4}
\EGAIVtwoTermEN{Irredundant $\mathcal O_X$-Module}{3.2.4}
\EGAIVtwoTermEN{$\mathcal O_X$-Module normally flat along a sub-prescheme}{6.10.1}
\EGAIVtwoTermEN{$\mathcal O_X$-Module satisfying property $(S_k)$,
satisfying property $(S_k)$ at a point}{5.7.2}
\EGAIVtwoTermEN{Reduced $\mathcal O_X$-Module, reduced at a point}{3.2.2}
\EGAIVtwoTermEN{$Z$-closed, $Z$-pure $\mathcal O_X$-Module}{5.9.9}
\EGAIVtwoTermEN{$Z$-closed, $Z$-pure $\mathcal O_X$-Module at a point}{5.9.9}
\EGAIVtwoTermEN{Birational morphism of reduced preschemes}{6.15.4}
\EGAIVtwoTermEN{Connected morphism}{4.5.5}
\EGAIVtwoTermEN{Cohen-Macaulay morphism at a point, Cohen-Macaulay morphism}{6.8.1}
\EGAIVtwoTermEN{Morphism of codepth $\leq n$ at a point, morphism of codepth $\leq n$}{6.8.1}
\EGAIVtwoTermEN{Morphism defined over a field}{4.8.4}
\EGAIVtwoTermEN{Irreducible morphism}{4.5.5}
\EGAIVtwoTermEN{Smooth morphism at a point, smooth morphism}{6.8.1}
\EGAIVtwoTermEN{Normal morphism at a point, normal morphism}{6.8.1}
\EGAIVtwoTermEN{$P$-morphism}{7.3.1}
\EGAIVtwoTermEN{Morphism satisfying property $(R_k)$ at a point,
morphism satisfying property $(R_k)$}{6.8.1}
\EGAIVtwoTermEN{Morphism satisfying property $(S_k)$ at a point,
morphism satisfying property $(S_k)$}{6.8.1}
\EGAIVtwoTermEN{Quasi-faithfully flat morphism, quasi-flat morphism}{2.3.3}
\EGAIVtwoTermEN{Radicial morphism at a point}{6.15.3}
\EGAIVtwoTermEN{Reduced morphism at a point, reduced morphism}{6.8.1}
\EGAIVtwoTermEN{Regular morphism at a point, regular morphism}{6.8.1}
\EGAIVtwoTermEN{Resolving morphism}{7.9.1}
\EGAIVtwoTermEN{Universally open morphism, universally bicontinuous morphism}{2.4.2}
\EGAIVtwoTermEN{Radicial multiplicity, separable multiplicity, total
multiplicity of an extension of a field}{4.7.4}
\EGAIVtwoTermEN{Radicial multiplicity, separable multiplicity, total
multiplicity of a prescheme over a field}{4.7.4}
\EGAIVtwoTermEN{Radicial multiplicity of a point for an $\mathcal O_X$-Module
($X$ a prescheme over a field)}{4.7.5}
\EGAIVtwoTermEN{Radicial multiplicity of a maximal associated prime cycle of
an $\mathcal O_X$-Module ($X$ a prescheme over a field)}{4.7.8}
\EGAIVtwoTermEN{Total multiplicity of a point (of a maximal associated prime
cycle) for an $\mathcal O_X$-Module ($X$ a prescheme over a field)}{4.7.12}

\medskip
\EGAIVtwoTermEN{Geometric number of connected components or irreducible
components of a prescheme over a field}{4.5.2}

\medskip
\EGAIVtwoTermEN{Associated point of an $\mathcal O_X$-Module or of a prescheme}{3.1.1}
\EGAIVtwoTermEN{Cohen-Macaulay point of a prescheme}{5.7.1}
\EGAIVtwoTermEN{Geometrically unibranch point, unibranch point of a prescheme}{6.15.1}
\EGAIVtwoTermEN{Prescheme connected in codimension $d-1$}{5.10.8}
\EGAIVtwoTermEN{Cohen-Macaulay prescheme}{5.7.1}
\EGAIVtwoTermEN{Excellent prescheme}{7.8.5}
\EGAIVtwoTermEN{Geometrically connected prescheme (over a field)}{4.5.2}
\EGAIVtwoTermEN{Geometrically integral prescheme (over a field)}{4.6.2}
\EGAIVtwoTermEN{Geometrically punctually integral prescheme (over a field)}{4.6.9}
\EGAIVtwoTermEN{Geometrically irreducible prescheme (over a field)}{4.5.2}
\EGAIVtwoTermEN{Geometrically normal prescheme (over a field) at a point,
geometrically normal prescheme}{6.7.6}
\EGAIVtwoTermEN{Geometrically reduced prescheme (over a field)}{4.6.2}
\EGAIVtwoTermEN{Geometrically reduced prescheme (over a field) at a point}{4.6.9}
\EGAIVtwoTermEN{Geometrically regular prescheme (over a field) at a point,
geometrically regular prescheme}{6.7.6}

\clearpage
\oldpage[IV-2]{228}

\EGAIVtwoTermEN{Geometrically unibranch prescheme at a point,
geometrically unibranch prescheme}{6.15.1}
\EGAIVtwoTermEN{Integral prescheme at a point}{3.2.4}
\EGAIVtwoTermEN{Irredundant prescheme}{3.2.4}
\EGAIVtwoTermEN{Prescheme locally embeddable in a regular scheme}{5.11.1}
\EGAIVtwoTermEN{Punctually integral prescheme}{4.6.9}
\EGAIVtwoTermEN{Prescheme satisfying property $(R_k)$, satisfying property
$(R_k)$ at a point}{5.8.2}
\EGAIVtwoTermEN{Prescheme (over a field) having geometric property $(R_k)$ at
a point, having geometric property $(R_k)$}{6.7.6}
\EGAIVtwoTermEN{Prescheme satisfying property $(S_k)$, satisfying property
$(S_k)$ at a point}{5.7.2}
\EGAIVtwoTermEN{Prescheme regular in codimension $k$}{5.8.2}
\EGAIVtwoTermEN{Separable prescheme (over a field)}{4.6.2}
\EGAIVtwoTermEN{Separable prescheme (over a field) at a point}{4.6.9}
\EGAIVtwoTermEN{Unibranch prescheme at a point, unibranch prescheme}{6.15.1}
\EGAIVtwoTermEN{Universally catenary prescheme}{5.6.3}
\EGAIVtwoTermEN{Universally reduced prescheme (over a field)}{4.6.2}
\EGAIVtwoTermEN{Depth of an $\mathcal O_X$-Module at a point}{5.7.1}
\EGAIVtwoTermEN{Property $P$ of the first type (of the second type) associated
with a property $R$ of Noetherian local rings}{7.3.10}
\EGAIVtwoTermEN{Geometric property $P$}{7.3.6}

\medskip
\EGAIVtwoTermEN{Resolution of singularities}{7.9.1}

\medskip
\EGAIVtwoTermEN{Subset (of a prescheme) defined over a field}{4.8.4}
\EGAIVtwoTermEN{Sub-$\mathcal O_X$-Module defined over a field}{4.8.4}
\EGAIVtwoTermEN{Primary sub-$\mathcal O_X$-Module of an $\mathcal O_X$-Module}{3.2.4}
\EGAIVtwoTermEN{Sub-prescheme defined over a field}{4.8.4}
\EGAIVtwoTermEN{Primary closed sub-prescheme of a prescheme}{3.2.4}

\medskip
\EGAIVtwoTermEN{$Z$-closure of an $\mathcal O_X$-Module}{5.9.11}

\endgroup

\clearpage
\oldpage[IV-2]{229}
\phantomsection
\addcontentsline{toc}{section}{Original table of contents}
\section*{Original table of contents}

\begingroup
\small
\setlength{\parindent}{0pt}
\setlength{\parskip}{1.5pt}
\newcommand{\EGAIVtwoContentsLineEN}[3]{\textbf{#1}\quad #2\dotfill #3\par}

\EGAIVtwoContentsLineEN{Chapter IV.}{\textbf{Local study of schemes and
morphisms of schemes} \emph{(continued)}}{5}

\EGAIVtwoContentsLineEN{\S~2.}{Base change and flatness}{5}
\EGAIVtwoContentsLineEN{2.1.}{Flat Modules on preschemes}{5}
\EGAIVtwoContentsLineEN{2.2.}{Faithfully flat Modules on preschemes}{9}
\EGAIVtwoContentsLineEN{2.3.}{Topological properties of flat morphisms}{14}
\EGAIVtwoContentsLineEN{2.4.}{Universally open morphisms and flat morphisms}{19}
\EGAIVtwoContentsLineEN{2.5.}{Permanence of properties of Modules under faithfully flat descent}{22}
\EGAIVtwoContentsLineEN{2.6.}{Permanence of set-theoretic and topological properties of morphisms under faithfully flat descent}{27}
\EGAIVtwoContentsLineEN{2.7.}{Permanence of various properties of morphisms under faithfully flat descent}{29}
\EGAIVtwoContentsLineEN{2.8.}{Preschemes over a regular base of dimension 1; closure of a closed sub-prescheme of the generic fibre}{33}

\EGAIVtwoContentsLineEN{\S~3.}{Associated prime cycles and primary decompositions}{36}
\EGAIVtwoContentsLineEN{3.1.}{Associated prime cycles of a Module}{36}
\EGAIVtwoContentsLineEN{3.2.}{Irredundant decompositions}{40}
\EGAIVtwoContentsLineEN{3.3.}{Relations with flatness}{43}
\EGAIVtwoContentsLineEN{3.4.}{Properties of the sheaves $\mathscr F/t\mathscr F$}{46}

\EGAIVtwoContentsLineEN{\S~4.}{Change of base field in algebraic preschemes}{52}
\EGAIVtwoContentsLineEN{4.1.}{Dimension of algebraic preschemes}{52}
\EGAIVtwoContentsLineEN{4.2.}{Associated prime cycles on algebraic preschemes}{54}
\EGAIVtwoContentsLineEN{4.3.}{Review of tensor products of fields}{58}
\EGAIVtwoContentsLineEN{4.4.}{Irreducible and connected preschemes over an algebraically closed field}{59}
\EGAIVtwoContentsLineEN{4.5.}{Geometrically irreducible and geometrically connected preschemes}{61}
\EGAIVtwoContentsLineEN{4.6.}{Geometrically reduced algebraic preschemes}{68}

\clearpage
\oldpage[IV-2]{230}

\EGAIVtwoContentsLineEN{4.7.}{Multiplicities in the primary decomposition on an algebraic prescheme}{75}
\EGAIVtwoContentsLineEN{4.8.}{Fields of definition}{80}
\EGAIVtwoContentsLineEN{4.9.}{Fields of definition of a subset of a prescheme}{84}

\EGAIVtwoContentsLineEN{\S~5.}{Dimension, depth, and regularity in locally Noetherian preschemes}{86}
\EGAIVtwoContentsLineEN{5.1.}{Dimension of preschemes}{86}
\EGAIVtwoContentsLineEN{5.2.}{Dimension of an algebraic prescheme}{90}
\EGAIVtwoContentsLineEN{5.3.}{Dimension of the support of a Module and the Hilbert polynomial}{92}
\EGAIVtwoContentsLineEN{5.4.}{Dimension of the image of a morphism}{93}
\EGAIVtwoContentsLineEN{5.5.}{Dimension formula for a finite type morphism}{94}
\EGAIVtwoContentsLineEN{5.6.}{Dimension formula and universally catenary rings}{97}
\EGAIVtwoContentsLineEN{5.7.}{Codepth and property $(S_k)$}{103}
\EGAIVtwoContentsLineEN{5.8.}{Regular preschemes and property $(R_k)$. Serre's criterion for normality}{107}
\EGAIVtwoContentsLineEN{5.9.}{$Z$-pure and $Z$-closed Modules}{109}
\EGAIVtwoContentsLineEN{5.10.}{Property $(S_2)$ and $Z$-closure}{114}
\EGAIVtwoContentsLineEN{5.11.}{Coherence criteria for the Modules $\mathcal H^0_{X/Z}(\mathscr F)$}{122}
\EGAIVtwoContentsLineEN{5.12.}{Relations between the properties of a Noetherian local ring $A$ and of a quotient ring $A/tA$}{126}
\EGAIVtwoContentsLineEN{5.13.}{Permanence properties under passage to the inductive limit}{131}

\EGAIVtwoContentsLineEN{\S~6.}{Flat morphisms of locally Noetherian preschemes}{134}
\EGAIVtwoContentsLineEN{6.1.}{Flatness and dimension}{135}
\EGAIVtwoContentsLineEN{6.2.}{Flatness and projective dimension}{137}
\EGAIVtwoContentsLineEN{6.3.}{Flatness and depth}{138}
\EGAIVtwoContentsLineEN{6.4.}{Flatness and property $(S_k)$}{141}
\EGAIVtwoContentsLineEN{6.5.}{Flatness and property $(R_k)$}{143}
\EGAIVtwoContentsLineEN{6.6.}{Transitivity properties}{145}
\EGAIVtwoContentsLineEN{6.7.}{Application to base changes in algebraic preschemes}{145}
\EGAIVtwoContentsLineEN{6.8.}{Regular, normal, reduced, and smooth morphisms}{150}
\EGAIVtwoContentsLineEN{6.9.}{The generic flatness theorem}{153}
\EGAIVtwoContentsLineEN{6.10.}{Dimension and depth of a Module normally flat along a closed sub-prescheme}{155}
\EGAIVtwoContentsLineEN{6.11.}{Criteria for the sets $U_{S_n}(\mathscr F)$ or $U_{C_n}(\mathscr F)$ to be open}{158}
\EGAIVtwoContentsLineEN{6.12.}{Nagata criteria for $\operatorname{Reg}(X)$ to be open}{163}
\EGAIVtwoContentsLineEN{6.13.}{Criteria for $\operatorname{Nor}(X)$ to be open}{168}
\EGAIVtwoContentsLineEN{6.14.}{Base change and integral closure}{169}
\EGAIVtwoContentsLineEN{6.15.}{Geometrically unibranch preschemes}{176}

\clearpage
\oldpage[IV-2]{231}

\EGAIVtwoContentsLineEN{\S~7.}{Relations between a Noetherian local ring and its completion. Excellent rings}{182}
\EGAIVtwoContentsLineEN{7.1.}{Formal equidimensionality and formally catenary rings}{183}
\EGAIVtwoContentsLineEN{7.2.}{Strictly formally catenary rings}{187}
\EGAIVtwoContentsLineEN{7.3.}{Formal fibres of Noetherian local rings}{192}
\EGAIVtwoContentsLineEN{7.4.}{Permanence of properties of formal fibres}{198}
\EGAIVtwoContentsLineEN{7.5.}{A criterion for $P$-morphisms}{203}
\EGAIVtwoContentsLineEN{7.6.}{Applications: I. Japanese local rings}{208}
\EGAIVtwoContentsLineEN{7.7.}{Applications: II. Universally Japanese rings}{212}
\EGAIVtwoContentsLineEN{7.8.}{Excellent rings}{214}
\EGAIVtwoContentsLineEN{7.9.}{Excellent rings and resolution of singularities}{218}
\EGAIVtwoContentsLineEN{}{Bibliography \emph{(continued)}}{224}
\EGAIVtwoContentsLineEN{}{Index of notation}{225}
\EGAIVtwoContentsLineEN{}{Terminological index}{226}

\vspace{2em}
\begin{center}
\emph{Manuscript received December 15, 1963.}
\end{center}

\endgroup
\clearpage
% END NATIVE INPUT source_aligned/ega4-7

\clearpage
\phantomsection
\addcontentsline{toc}{section}{Consolidated bibliography}
\section*{Consolidated bibliography}
\nocite{*}
\bibliography{the}
\bibliographystyle{amsalpha}

\end{document}
